\xchapter{Chapter I}

Short URL for this page:

bit.ly/BURLAT1

p1

The continuity of history, which means the control of the present and future by the past, has become a commonplace, and chronological limits, which used to be considered important, are now recognised to have little significance except as convenient landmarks in a historical survey. Yet there are what we may call culminating epochs, in which the accumulating tendencies of the past, reaching a certain point, suddenly effect a visible transformation which seems to turn the world in a new direction. Such a culminating epoch occurred in the history of the Roman Empire at the beginning of the fourth century. The reign of Constantine the Great inaugurated a new age in a much fuller sense than the reign of Augustus, the founder of the Empire. The anarchy of the third century, when it almost seemed that the days of the Roman Empire were numbered, had displayed the defects of the irregular and heterogeneous system of government which Augustus had established to administer his immense dominion. His successors had introduced modifications and improvements here and there, but events made it clearer and clearer that a new system, more centralised and more uniform, was required, if the Empire was to be held together. To Diocletian, who rescued the Roman world at the brink of the abyss, belongs the credit of having framed a new system of administrative machinery. Constantine developed and completed the work of Diocletian by measures which were more radical and more far-reaching . The foundation of Constantinople as a second Rome inaugurated a permanent division between the Eastern and Western, the Greek and the Latin, halves of the Empire \texthebrew{—} a division to which events had already pointed \texthebrew{—} and

p2
affected decisively the whole subsequent history of Europe. Still more evidently and notoriously did Constantine mould the future by accepting Christianity as the State religion.

In the present work the history of the Roman Empire is taken up at a point about sixty years after Constantine's death, when the fundamental changes which he introduced have been firmly established and their consequences have emerged into full evidence. The new system of government has been elaborated in detail, and the Christian Church has become so strong that no enemies could prevail against it. Constantinople, created in the likeness of Rome, has become her peer and will soon be fully equipped for the great rôle which she is to play in Europe and Hither Asia for more than a thousand years. She definitely assumes now her historical position. For after the death of Theodosius the Great, who had ruled alone for a short time over a dominion extending from Scotland to Mesopotamia, the division of the Empire into two geographical portions, an eastern and a western, under two Emperors, a division which had been common during the past century, was finally established. This dual system lasted for eighty-five years, and but for the dismemberment of the western provinces by the Germans might have lasted indefinitely. In the constitutional unity of the Empire this arrangement caused no breach.

Again, the death of Theodosius marks the point at which the German danger, long imminent over the Empire, begins to move rapidly towards its culmination. We are on the eve of the great dismemberment of Roman dominion which, within seventy years, converted the western provinces into Teutonic kingdoms. The fourth century had witnessed the settlement of German peoples, as foederati , bound to military service, on Roman lands in the Balkan peninsula and in Gaul. Through the policy of Constantine Germans had become a predominant element in the Roman army, and German officers had risen to the highest military posts and had exercised commanding political influence. Outside, German peoples were pressing on the frontiers, waiting for opportunities to grasp at a share of the coveted wealth of the Roman world. The Empire was exposed to the double danger of losing provinces to these unwelcome claimants who desired to be taken within its border, and of the growing ascendancy

p3
of the German element in the army.\texthebrew{⁠}\footnote{The Roman fear of barbarisation is marked by a law of A.D. 370 or 373, which forbade marriages between provincials and barbarians on pain of death.

\emph{C. Th.} III.14.1 .

} The East was menaced as well as the West, and the great outstanding fact in the history of the fifth century is that the East survived and the West succumbed. The success of the Eastern government in steering through these perils was partly due to the fact that during this critical time it was on good terms, only seldom and briefly interrupted, with Persia, its formidable neighbour.

The diminished Roman Empire, now centering entirely in Constantinople, lasted for a thousand years, surrounded by enemies and frequently engaged in a struggle for life or death, but for the greater part of that long period the most power ­ful State in Europe. Its history is marked by distinct ages of expansion, decline, and resuscitation, which are easily remembered and help to simplify the long series of the annals of Byzantium.\texthebrew{⁠}\footnote{Cp. Bury, Appendix 9 to Gibbon, vol. V.

} Having maintained itself in the fifth century and won its way through the German peril, it found itself strong enough in the sixth to take the offensive and to recover Africa and Italy. Overstrain led to a decline, of which Persia took advantage, and when this danger had been overcome, the Saracens appeared as a new and more formidable force and deprived the Empire of important provinces in Asia, while at the same time European territory was lost to the Bulgarians and the Slavs (seventh century). Then a period of resuscitation in the eighth and ninth centuries led to a new age of brilliance and expansion (ninth to eleventh centuries). When the Saracens had ceased to be formidable, the Seljuk Turks appeared, and the Empire found it difficult to hold its own against this foe as well as against the western powers of Europe, and the barbarians of the north. This period ends with the disaster of 1204, when Constantinople fell into the hands of the Crusaders, who treated the city with more barbarity than the barbarian Alaric had treated Rome eight hundred years before. After this the cycle begins anew; first, the period of revival at Nicaea, which became the temporary capital; then the recovery of Constantinople (1261), followed by a period in which the Empire could assert its power; finally, from the middle of the fourteenth century, the decline, and the

p4
last death-struggle with the Ottomans, ending in the capture of the city in 1453.

The State which maintained itself in unbroken continuity throughout the vicissitudes of more than a thousand years is proverbial for its conservative spirit. It was conservative in its constitution and institutions, in the principles and the fashions of its civilisation, in its religion, in its political and social machinery. It may be conjectured that this conservatism is partly to be attributed to the influence of the legal profession.\texthebrew{⁠}\footnote{This point may be illustrated by the interesting section on \emph{L'Esprit légiste} in de Tocqueville's \emph{De la démocratie en Amérique} (part II. chap. VIII).

} Lawyers are always conservative and suspicious of change, and it would be difficult to exaggerate their importance and the power of their opinion in the later Empire. It was natural and just that their influence should be great, for it has well been observed that it was to the existence of a ``judicial establishment, guided by a published code, and controlled by a body of lawyers educated in public schools, that the subjects of the Empire were chiefly indebted for the superiority in civilisation which they retained over the rest of the world.''\texthebrew{⁠}\footnote{Finlay, \emph{Hist. of Greece}, I.411.

} But the conservatism of Byzantium is often represented as more rigid than it actually was. The State could not have survived if it had not been constantly adapting its institutions to new circumstances. We have seen how its external history may be divided into periods. But its administrative organisation, its literature, its art display equally well-defined stages.

One more introductory remark. The civilisation of the later Empire, which we know under the name of Byzantine, had its roots deep in the past. It was simply the last phase of Hellenic culture. Alexandria, the chief city of the Hellenic world since the third century B.C. , yielded the first place to Byzantium in the course of the fifth century. There was no breach in continuity; there was only a change of centre. And while the gradual ascendancy of Christianity distinguished and stamped the last phase, we must remember that Christian theology had been elaborated by the Greek mind into a system of metaphysics which Paul, the founder of the theology, would not have recognised, and which no longer seemed an alien product.

The Roman Empire was founded by Augustus, but for three centuries after its foundation the State was constitutionally a republic. The government was shared between the Emperor and the Senate; the Emperor, whose constitutional position was expressed by the title Princeps was limited by the rights of the Senate. Hence it has been found convenient to distinguish this period as the Principate or the Dyarchy. From the very beginning the Princeps was the predominant partner, and the constitutional history of the Principate turns on his gradual and steady usurpation of nearly all the functions of government which Augustus had attributed to the Senate. The republican disguise fell away completely before the end of the third century. Aurelian adopted external fashions which marked a king, not a citizen; and Diocletian and Constantine definitely transformed the State from a republic to an autocracy. This change, accompanied by corresponding radical reforms, was, from a purely constitutional point of view, as great a break with the past as the change wrought by Augustus, and the transition was as smooth. Augustus preserved continuity with the past by maintaining republican forms; while Constantine and his predecessors simply established on a new footing the supreme Imperial power which already existed in fact, discarding the republican mask which had worn too thin.

The autocracy brought no change in the principle of succession to the throne. Down to its fall in the fifteenth century the Empire remained elective, and the election rested with the Senate and the army. Either the Senate or the army could proclaim an Emperor, and the act of proclamation constituted a legitimate title. As a rule, the choice of one body was acquiesced in by the other; if not, the question must be decided by a struggle. Any portion of the army was considered, for this purpose, as representing the whole army, and thus in elections in Constantinople it was the troops stationed there with whom the decision lay. But whether Senate or army took the initiative, the consent of the other body was required; and the inauguration\texthebrew{⁠}\footnote{The term \textgreek{\texthebrew{ἀ}ναγόρευσις}, proclamation, was used for all the proceedings of the inauguration. In the case of Carus the Senate played no part. Mommsen, \emph{Staatsrecht}, II.843.

} of the new Emperor was not complete till he had

p6
been acclaimed by the people. Senate, army, and people, each had its place in the inaugural ceremonies.

But while the principle of election was retained, it was in actual practice most often only a form. From the very beginning the principle of heredity was introduced indirectly. The reigning Emperor could designate his successor by appointing a co-regent . In this way Augustus designated his stepson Tiberius, Vespasian his son Titus. The Emperors naturally sought to secure the throne for their sons, and if they had no son, generally looked within their own family. From the end of the fourth century it became usual for an Emperor to confer the Imperial title on his eldest son, whether an adult or an infant. The usual forms of inauguration were always observed; but the right of the Emperor to appoint co-regents was never disputed. The consequence was that the succession of the Roman Emperors presents a series of dynasties, and that it was only at intervals, often considerable, that the Senate and army were called upon to exercise their right of election.

The co-regent was a sleeping partner. He enjoyed the Imperial honours, his name appeared in official documents; but he did not share in the actual government, except so far as he might be specially authorised by his older colleague. This, at least, was the rule. Under the Principate the senior Imperator distinguished his own position from that of his colleague by raising to himself the title of Pontifex Maximus. Marcus Aurelius tried a new experiment and shared the full sovranty with Lucius Verus. This division of the sovranty was an essential part of the system of Diocletian, corresponding to the geographical partition of the Empire which he introduced. From his time down to A.D. 480, the Empire is governed by two (or even more) sovran colleagues, who have all equal rights and competence, and differ only in seniority. Sometimes the junior Emperor is appointed by the senior, sometimes he is elected independently and is recognised by the senior. Along with these there may be co-regents , who exercise no sovran power, but are marked out as eventual successors. Thus the child Arcadius was for nine years co-regent with the Emperors Valentinian II and Theodosius the Great. No formal title, however, raised the sovran above the co-regent , though the latter, for the sake of distinction, was often called ``the second

p7
Emperor,'' or if he was a child, ``the little Emperor.''\texthebrew{⁠}\footnote{\textgreek{\texthebrew{Ὁ} δεύτερος βασιλεύς, \texthebrew{ὁ} μικρ\texthebrew{ὸ}ς}. In later times the actual sovran was sometimes distinguished as the \textgreek{α\texthebrew{ὐ}τοκράτωρ}, but the plural, \textgreek{ο\texthebrew{ἱ} α\texthebrew{ὐ}τοκράτορες}, was used to designate all the Augusti.

} When towards the end of the fifth century the territorial partition of the Empire came to an end, the system of joint sovranty ceased, and henceforward, whenever there is more than one Augustus, only one exercises the sovran power.\footnote{There are indeed one or two exceptional cases.

}

But the Emperor could also designate a successor, without elevating him to the position of co-regent , by conferring on him the title of Caesar. This practice, which since Hadrian was usual under the Principate,\texthebrew{⁠}\footnote{Mommsen, \emph{Staatsrecht}, II.1139 \emph{sqq.}

} and was adopted by Constantine, is not frequent in the later Empire.\texthebrew{⁠}\footnote{Bury, \emph{Imp. Adm. System}, 36.

} If the Emperor has sons, he almost invariably creates his eldest son Augustus. If not, he may signify his will as to the succession by bestowing the dignity of Caesar. The Emperor before his death might raise the Caesar to the co-regency .\texthebrew{⁠}\footnote{So Justin II created the Caesar Tiberius Augustus, shortly before his death in 578; similarly Tiberius created his son-in\texthebrew{‑}law Maurice Caesar, and on his death-bed caused him to be proclaimed Augustus.

} If he died without having done this, the Caesar had to be elected in the usual way by the Senate and the army. This method of provisional and revocable designation was often convenient. An Emperor who had no male issue might wish to secure the throne to a son-in\texthebrew{‑}law, for instance, in case of his own premature death. If he conferred the Caesarship and if a male child were afterwards born to him,\texthebrew{⁠}\footnote{This occurred in the ninth century in the case of Theophilus. His children were daughters; he bestowed the rank of Caesar on his son-in\texthebrew{‑}law Musele, and a son, who succeeded him as Michael III, was born later. The Caesarship conferred on Bardas by his nephew Michael III is also a case in point.

} that child would be created Augustus, and the Caesar's claim would fall into abeyance.

When the Emperor had more than one son, it was usual to confer the title of Caesar on the younger.\texthebrew{⁠}\footnote{The only cases which occur before 800 are the three younger sons of Heraclius, and the second and third sons of Constantine V.

} Constitutionally this may be considered a provision for the contingency of the death of the co-regent . Practically it meant a title of dignity reserved for the members of the Imperial family. Sometimes the co-regency was conferred on more than one son. Theodosius the Great raised Honorius to the rank of Augustus as well as

p8
his elder son Arcadius. But it is to be observed that this measure was not taken till after the death of the West Emperor Valentinian II, and that its object was to provide two sovrans, one for the East and one for the West. If the division of the Empire had not been contemplated, Honorius would not have been created Augustus in A.D. 393. To avoid a struggle between brothers, the obvious policy was to confer the supreme rank on only one. Before the reign of Basil I in the ninth century, there were few opportunities to depart from this rule of expediency, and it was only violated twice, in both cases with unfortunate consequences.\footnote{By Heraclius and by Constans II.

}

But the Caesarship was not the only method employed to signalise an eventual successor. In the third century it became usual to describe the Caesar, the Emperor's adopted son, as nobilissimus . In the fourth, this became an independent title, denoting a dignity lower than Caesar, but confined to the Imperial family. On two occasions we find nobilissimus used as a sort of preliminary designation.\texthebrew{⁠}\footnote{Jovian conferred it on his infant son Valerian, and Honorius on his child-nephew Valentinian III.

} But it fell out of use in the fifth century, and apparently was not revived till the eighth, when it was conferred on the youngest members of the large family of Constantine V.\texthebrew{⁠}\footnote{Cp. Bury, \emph{Imp. Adm. System}, 35.

} In the sixth century Justinian introduced a new title, Curopalates, which, inferior to Caesar and nobilissimus, might serve either to designate or simply to honour a member of the Imperial family. We find it used both ways.\texthebrew{⁠}\footnote{As designation, of Justin II by his uncle Justinian, of Domentziolus by his uncle Phocas. As an honour it was conferred by Maurice and Heraclius on their brothers; by Leo III and Nicephorus I on their sons-in\texthebrew{‑}law. It was not confined to the Imperial family after the tenth century. Cp. Bury, \emph{ib.} 34.

} It was a less decided designation than the Caesarship, and a cautious or suspicious sovran might prefer it.

The principle of heredity, which was thus conciliated with the principle of election, gradually gave rise to the view that not only was the Emperor's son his \emph{legitimate} successor, but that if he had no male issue, the question of succession would be most naturally and satisfactorily settled by the marriage of a near female relative \texthebrew{—} daughter, sister, or widow, \texthebrew{—} and the election of her husband, who would thus continue the dynasty.\texthebrew{⁠}\footnote{In the fifth century we have two cases: Pulcheria (450) and Ariadne (491).

} There

p9
was a general feeling of attachment to a dynasty, and the history of the Later Empire presents a series of dynasties, with few and brief intervals of unsettlement. During the four centuries between 395 and 802, we have five dynasties, which succeed one another, except in two cases,\texthebrew{⁠}\footnote{Phocas, between the dynasty of Justin and that of Heraclius; and the period of anarchy between the Heraclian and Isaurian dynasties.

} without a break.

Though there was no law excluding women from the succession, yet perhaps we may say that up to the seventh or eighth century it would have been considered not merely politically impossible, but actually illegal, for a woman to exercise the sovran power in her own name. The highest authority on the constitution of the early Empire affirms that her sex did not exclude a woman from the Principate.\texthebrew{⁠}\footnote{Mommsen, \emph{Staatsrecht}, II.788. The evidence he adduces is not convincing.

} But the title Augusta did not include the proconsular Imperium and the tribunician potestas, which constituted the power of the Princeps, and it is not clear that these could have been conferred legally on a woman or that she could have borne the title Imperator. It is said, and may possibly be true, that Caligula, when he was ill, designated his favourite sister Drusilla as his successor;\texthebrew{⁠}\footnote{Suetonius, \emph{C. Caligula}, 24 .

} but this does not prove that she could legally have acted as Princeps. Several Empresses virtually shared the exercise of the Imperial authority, bore themselves as co-regents , and enjoyed more power than male co-regents ; but their power was \emph{de facto}, not \emph{de jure.} Some were virtually sovrans, but they were acting as regents for minors.\texthebrew{⁠}\footnote{Pulcheria; Placidia; Martina.

} Not till the end of the eighth century do we find a woman, the Empress Irene, exercising sovranty alone and in her own name.\texthebrew{⁠}\footnote{If the eligibility of a woman had been recognised, the principle would probably have been applied in the case of the Augusta Pulcheria (who had considerable experience of government, and enjoyed the respect and confidence of the Empire) in A.D. 450.

} This was a constitutional innovation. The experiment was only once repeated,\texthebrew{⁠}\footnote{In the eleventh century, when Zoe and Theodora reigned together. There would have been another instance if Stauracius, in 811, had succeeded in procuring the succession for his wife Theophano (Bury, \emph{Eastern Roman Empire}, p18).

} and only in exceptional circumstances would it have been tolerated. There was a general feeling against a female reign, both as inexpedient and as a violation of tradition.\texthebrew{⁠}\footnote{This feeling was strongly expressed towards Martina in A.D. 641.

} Between the fourth and eighth centuries, however, two circumstances may have combined

p10
to make it appear no longer illegal. The Greek official term for Imperator was Autokrator, and in the course of time, when Latin was superseded by Greek, and Imperator fell out of use and memory, Autokrator ceased to have the military associations which were attached to its Latin equivalent, and the constitutional incompatibility of the office with the female sex is no longer apparent. In the second place, female regencies prepared the way for Irene's audacious step. When a new Emperor was a minor, the regency might be entrusted to his mother or an elder sister, whether acting alone or in conjunction with other regents. Irene was regent for her son before she grasped the sole power for herself.

The title of Augusta was always conferred\texthebrew{⁠}\footnote{In the East, from the time of Arcadius. The wives of Honorius were not Augustae.

} on the wife of the Emperor and the wife of the co-regent , and from the seventh century it was frequently conferred on some or all of the Emperor's daughters. The reigning Augusta might have great political power. In the sixth century, Justinian and Theodora, and Justin II and Sophia, exercised what was virtually a joint rule, but in neither case did the constitutional position of the Empress differ from that of any other consort.

The diadem was definitely introduced by Constantine,\texthebrew{⁠}\footnote{See W. Sickel, \emph{B.Z.} VII.513 \emph{sqq.} According to Victor, \emph{Epit.} 35.5, it was all worn by Aurelian; according to John Lydus, \emph{De mag.} I.4 , by Diocletian. The diadem was a white browband, set with pearls.

} and it may be considered the supreme symbol of the autocratic sovranty which replaced the magistracy of the earlier Empire. Hitherto the distinguishing mark of the Emperor's costume had been the purple cloak of the Imperator; and ``to assume the purple'' continued to be the common expression for elevation to the throne. The crown was an importation from Persia, and it invested the Roman ruler with the same external dignity as the Persian king. In Persia it was placed on the king's head by the High Priest of the Magian religion.\texthebrew{⁠}\footnote{The rest of this paragraph is borrowed from my \emph{Constitution of the L.R.E.}

} In theory the Imperial crown should be imposed by a representative of those who conferred the sovran authority that it symbolised. And in the fourth century we find the Prefect Sallustius Secundus crowning Valentinian I, in whose election he had taken the most prominent part. But the Emperor seems to have felt some hesitation in

p11
receiving the diadem from the hands of a subject, and the selection of one magnate for the office was likely to cause jealousy. Yet a formality was necessary. In the fifth century the difficulty was overcome in an ingenious and tactful way. The duty of coronation was assigned to the Patriarch of Constantinople. In discharging this office the Patriarch was not envied by the secular magnates because he could not be their rival, and his ecclesiastical position relieved the Emperor from all embarrassment in receiving the diadem from a subject. There is, as we shall see, some evidence that this plan was adopted in A.D. 450 at the coronation of Marcian, but it seems certain that his successor Leo was crowned by the Patriarch in A.D. 457. Henceforward this was the regular practice. But it was only the practice. It was the regular and desirable mode of coronation, but was never legally indispensable for the autocrat's inauguration. The last of the East Roman Emperors, Constantine Palaeologus, was not crowned by the Patriarch; he was crowned by a layman.\texthebrew{⁠}\footnote{Nicephorus Bryennius (eleventh century) crowned himself. Anna Commena, º \emph{Alexiad.}, 1.4.

} This fact that coronation by the Patriarch was not constitutionally necessary is important. It shows that the Patriarch in performing the ceremony was not representing the Church. It is possible that the idea of committing the office to him was suggested by the Persian coronations by the High Priest. But the significance was not the same. The chief of the Magians acted as representative of the Persian religion, the Patriarch acted as representative of the State. If he had specially represented the Church, his co-operation could never have been dispensed with. The consent of the Church was not formally necessary to the inauguration of a sovran.

This point is further illustrated by the fact that when the Emperor appointed a colleague, the junior Augustus was crowned not by the Patriarch but by the Emperor who created him.\footnote{See

Const. Porph. \emph{De cer.} I.38, p194 . Sometimes he might commit the office to the Patriarch, who then acted simply as his delegate.

}

When Augustus founded the Empire, he derived his Imperial authority from the sovranty of the people; and the essence of this principle was retained throughout the duration not only of the Principate but also of the Monarchy; for the Imperial office remained elective, and the electors had the right of deposing the Emperor. But though these rights were never abrogated,

p12
there was a tendency, as time went on, to regard the majesty and power of the monarch as resting on something higher than the will of the people. The suggestion of divinity has constantly been the device of autocrats to strengthen and enhance their power; and modern theories of Divine Right are merely a substitute for the old pagan practice of deifying kings. Augustus attempted to throw a sort of halo round his authority by designating himself officially Divi Filius . But the glow of this consecration faded, and disappeared entirely with the fall of the Julio-Claudian dynasty. With Aurelian, who foreshadows the new Monarchy, the suggestion of divinity again appears.\texthebrew{⁠}\footnote{On his coins, Eckhel, \emph{Doct. num.} 7, 482.

} Diocletian and his colleague Maximian are designated as gods and parents of gods.\texthebrew{⁠}\footnote{Diis genitis et deorum creatoribus, CIL III.710 .

} The official deification of the Emperor, which seemed in sight at the beginning of the fourth century, was precluded by Christianity; but the consecration of the ruler's person was maintained in the epithets \emph{sacred} and \emph{divine}; and the Emperors came to regard themselves rather as vicegerents of God than as rulers set up by their people. Justinian, in one of his laws, speaks of the Emperor as sent down by God to be a living law.\texthebrew{⁠}\footnote{\emph{Nov.} 81.4. At the Council of Chalcedon, Marcian was acclaimed as ``priest and Emperor,'' \textgreek{τ\texthebrew{ῷ ἱ}ερε\texthebrew{ῖ} κα\texthebrew{ὶ} βασιλε\texthebrew{ῖ}} (Mansi, VII. p177).

} In the ninth century Basil I tells his son, ``You received the Empire from God.''\footnote{\emph{Paraenesis ad Leonem}, in \emph{P.G.} 107, p. xxv, cp. p.xxxii.

}

Under the Monarchy, the Emperor appropriated the full right of direct legislation, which had not belonged to him under the Principate.\texthebrew{⁠}\footnote{In one particular class of cases, namely the bestowal of civil rights on individuals and municipal rights on corporations, the Princeps had the power leges dare without the co-operation of the comitia. Mommsen, \emph{Staatsrecht}, II.888 \emph{sqq.}

} The Princeps possessed the right of initiating laws to be passed by the comitia of the people, but from the time of Tiberius legislation was seldom effected in this way, and after the first century it was exclusively in the hands of the Senate. The Emperor, communicating his instructions in the form of an oratio to the Senate, could have his wishes embodied in senatorial decrees ( senatus consulta ). But indirectly he possessed virtual powers of legislation by means of edicts and constitutions, which, though technically they were not laws, were for practical purposes equivalent.\texthebrew{⁠}\footnote{See Mommsen, \emph{ib.} 905 \emph{sqq.}

} The edict, unlike a law, did not necessarily contain

p13
a command; it was properly a public communication made by a magistrate to the people. But the legislative activity of the early Emperors was chiefly exercised in the form of constitutions, a term which in the stricter sense applied to decisions which were only brought to the notice of the persons concerned.\texthebrew{⁠}\footnote{Constitutiones is sometimes used in a wider sense to include leges and senatus consulta.

} This term included the Imperial correspondence and especially the mandates, or instructions addressed to officials. These ``acts'' had full validity, and the magistrates every year swore to observe them.\texthebrew{⁠}\footnote{Pomponius

(\emph{Dig.} I.2.2.14) : est principalis constitutio ut quod ipse princeps constituit pro lege servetur.

} But when an act required a dispensation from an existing law, the Imperial constitution was valid only during the lifetime of its author.

The power of dispensing from a law properly belonged to the Senate, and the earlier Emperors sought from the Senate a dispensation when necessary. Domitian began to encroach on this privilege. But the principle remained that the Princeps, who was constitutionally a magistrate, was bound by the laws; and when lawyers of the third century speak of the Princeps as legibus solutus , they refer to laws from which Augustus had formally obtained dispensation by the Senate.\footnote{Mommsen, \emph{ib.} 751, \emph{n.} 3.

}

Under the Monarchy the Emperors assumed full powers of legislation, and their laws took the form occasionally of an oratio to the Senate, but almost always of an edict. The term edict covered all the decisions which were formerly called constitutions, mandates, or rescripts, provided they had a general application.\texthebrew{⁠}\footnote{\emph{C. J.} I.14.3

(A.D. 426).

} And the Emperor not only legislated; he was the sole legislator, and reserved to himself the sole right of interpreting the laws.\texthebrew{⁠}\footnote{\emph{C. J.} I.14.12

(A.D. 529): si enim in praesenti leges condere soli imperatori concessum est, et leges interpretari solum dignum imperio esse oportet. Cp. \emph{ib.} 1.

} He possessed the dispensing power. But he always considered himself bound by the laws. An edict of A.D. 429 expresses the spirit of reverence for law, as something superior to the throne itself, which always animated the Roman monarchs. ``To acknowledge himself bound by the laws ( alligatum legibus ) is, for the sovran, an utterance befitting the majesty of a ruler. For the truth is that our authority depends on the authority of law. To submit our sovranty to the laws is verily a greater

p14
thing than Imperial power.''\texthebrew{⁠}\footnote{\emph{Ib.} 4 . So the Lawbook of the ninth century lays down that general laws are valid against the Emperor, and forbids rescripts which contradict them. \emph{Basilica} II.6.9.

} Deep respect for the rules of law, and their systematic observance characterised the Roman autocracy down to the fall of the Empire in the fifteenth century, and was one of the conditions of its long duration. It was never an arbitrary despotism, and the masses looked up to the Emperor as the guardian of the laws which protected against the oppression of nobles and officials.\footnote{Finlay has frequently insisted on this. Compare his remarks, and his comparison with the Saracen empire, in \emph{Hist. of Greece}, II.23\texthebrew{‑}24.

}

The laws, then, were a limitation on the power of the autocrat; and soon another means of limiting his power was discovered. In the fifth century, the duty of crowning a new Emperor at Constantinople was, as we saw, assigned to the Patriarch. In A.D. 491 the Patriarch refused to crown Anastasius unless he signed a written oath that he would introduce no novelty into the Church. This precedent was at first followed perhaps only in cases where a new Emperor was suspected of heretical tendencies, but by the tenth century\texthebrew{⁠}\footnote{Constantine Porph. \emph{De adm. imp.} p84.

} an oath of this kind seems to have been a regular preliminary to coronation. The fact that such capitulations could be and were imposed at the time of elevation shows that the autocracy was limited.

The essence of an autocracy is that no co-ordinate body exists which is able constitutionally to act as a check upon the monarch's will. The authority of the Senate or the Imperial Council might constitute a strong practical check upon an Emperor's acts, but if he chose to disregard their views, he could not be accused of acting unconstitutionally. The ultimate check on any autocracy is the force of public opinion. There is always a point beyond which the most arbitrary despot cannot go in defying it. In the case of a Roman Emperor, public opinion could exert this control constitutionally, by an extreme measure. The Emperor could be deposed. The right of deposition corresponded to the right of election. The deposition was accomplished not by any formal process, but by the proclamation of a new Emperor. If any one so proclaimed obtained sufficient support from the army, Senate, and people, the old Emperor was compelled to vacate the throne by \emph{force p15
majeure}; while the new Emperor was regarded as the legitimate monarch from the day on which he was proclaimed; the proclamation was taken as the legal expression of the general will. If he had not a sufficiently power ­ful following to render the proclamation effective and was suppressed, he was treated as a rebel; but during the struggle and before the catastrophe, the fact that the Senate or a portion of the army had proclaimed him gave him a presumptive constitutional status which the event might either confirm or annul. The method of deposition was, in fact, revolution; and we are accustomed to regard revolution as something essentially unconstitutional, an appeal from law to force; but under the Imperial system it was not unconstitutional; the government was, as has been said,\texthebrew{⁠}\footnote{By Mommsen.

} ``an autocracy tempered by the legal right of revolution.''\footnote{I have borrowed the last few sentences from my \emph{Constitution of the L.R.E.} 8\texthebrew{‑}9.

}

The transformation of the Principate into the Autocracy was accomplished by changes in the titular style of the Emperors, in their dress, in the etiquette of the court, which showed how entirely the old tradition of the republic had been forgotten.

The oriental conception of divine royalty is now formally expressed in the diadem; and it affects all that appertains to the Emperor. His person is divine; all that belongs to him is ``sacred.'' Those who come into his presence perform the act of adoration;\texthebrew{⁠}\footnote{Cp. Victor, \emph{Caes.} 39 (of Diocletian). See Godefroy's Comm. on \emph{C. Th.} vol. II, p83.

} they kneel down and kiss the purple. It had long been the habit to address the Imperator as dominus , ``lord''; in the fourth century the sovrans began to use it of themselves and Dominus Noster appears on their coins.\footnote{Mommsen, \emph{Staatsrecht}, II.760 \emph{sqq.} He observes that the terminological transition from princeps to dominus is a measure of the constitutional development towards autocracy. D.N. appears on brick-stamps towards end of 2nd cent.: CIL XV pp204\texthebrew{‑}5. \texthebrew{—} Probus, the consul of 406, in

his consular ivory diptych preserved at Aosta

(CIL V.6836 ) could describe himself as the famulus of Honorius.

}

Since the first century we can trace the use of Basileus to designate the Princeps, and Basileia to describe the Imperial power, in the eastern provinces of the Empire.\texthebrew{⁠}\footnote{Bréhier, ``L'Origine des titres impériaux à Byzance,'' \emph{B.Z.} XV.161 \emph{sqq.}

} Dion Chrysostom wrote a discourse on the Basileia; Fronto calls Marcus Aurelius ``the great Basileus, ruler of land and sea.'' Basileus was the equivalent of Rex, a title odious to Roman ears; but by the fourth century the Greek name had long ceased to wound

p16
any susceptibilities; it became the term regularly employed by Greek writers and in Greek inscriptions, and the Emperors began to employ it themselves. Usage soon went further. Basileus was reserved for the Emperor and the Persian king,\texthebrew{⁠}\footnote{Bréhier (p170) omits to note this important exception. The Abyssinian king seems to have been another. Cp. Bury, \emph{op. cit.} p20.

} and rex was employed to designate other barbarian royalties.

The Imperial Chancery was conservative, and it was not till the seventh century that the Emperor designated himself as Basileus in his constitutions and rescripts.\texthebrew{⁠}\footnote{This change was introduced by Heraclius.

} The official Greek equivalent of Imperator was Autokrator , which was similarly used as a praenomen.\texthebrew{⁠}\footnote{Justinian's style was: Imperator Iustinianus (or Imp. Caesar Flavius Iust.) pius felix inclitus victor ac triumphator semper Augustus (A.D. 529,

\emph{De Iust. cod. conf.}, at beginning of \emph{C. J.} ). In A.D. 534, this is expanded by a number of honourable epithets glorifying victories (Alamanicus, Gothicus, etc.) inserted immediately after Iustinianus. The Greek equivalent of the above is: \textgreek{α\texthebrew{ὐ}τοκράτωρ} (\textgreek{Κα\texthebrew{ῖ}σαρ Φλ}.) \textgreek{\texthebrew{Ἰ}ουστινιανός, ε\texthebrew{ὐ}σεβής, ε\texthebrew{ὐ}τυχής, \texthebrew{ἔ}νδοξος, νικητής, τροπαιο\texthebrew{ῦ}χος, \texthebrew{ἀ}εισέβαστος Α\texthebrew{ὐ}γο\texthebrew{ῦ}στος} (CIG III.8636). Cp. Bréhier, p171.

} The mint of Constantinople continued to inscribe the Imperial coins with Latin legends till the eighth century.\texthebrew{⁠}\footnote{The style is, \emph{e.g.} D(ominus) N(oster) Arcadius P(ius) F(elix) Aug(ustus). In the reign of Leo I, PP (or Perp) = Perpetuus was substituted for PF, and this was normal till the beginning of the eighth century.

} The earliest coins with Greek inscriptions have Basileus and Despotes .

The general use of Despotes is one of the most characteristic oriental features of the new Empire. It denoted the relation of a master to his slaves, and it was regularly used in addressing the Emperor from the time of Constantine to the fall of the Empire. Justinian expected this form of address. The subject spoke of himself as ``your slave.'' But this orientalism was a superficial etiquette; the autocrat seldom forgot that his subjects were freemen, that if he was a dominus , he was a dominus liberorum .

A few words may be said here about the unity of the Empire. From the reign of Diocletian to the last quarter of the fifth century, the Empire is repeatedly divided into two or more geographical sections \texthebrew{—} most frequently two, an Eastern and a Western \texthebrew{—} each governed by its own ruler. From A.D. 395 to A.D. 476, or rather 480, the division into two realms is practically continuous; each realm goes its own way, and the relations between them are sometimes even hostile. It has, naturally

p17
enough, proved an irresistible temptation to many modern writers to speak of them as if they were different Empires. To men of the fourth and fifth centuries such a mode of speech would have been unintelligible, and it is better to avoid it. To them there was and could be only one Roman Empire; and we should emphasise and not obscure this point of view.

But it is not merely a question of constitutional theory. The unity was not only formally recognised; it was maintained in practical ways. In the first place, the Imperial colleagues issued their laws under their joint names, and general laws promulgated by either and transmitted for publication to the chancery of his associate were valid throughout the whole Empire.\texthebrew{⁠}\footnote{\emph{C. Th.} I.1.5 .

} In the second place, on the death of either Emperor, the Imperial authority of the surviving colleague was constitutionally extended to the whole Empire until a successor was elected. Strictly speaking, it devolved upon him to nominate a new colleague. After the fall of the Theodosian House, some of the Emperors who were elected in Italy were not recognised at Constantinople, but the principle remained in force.

The unity of the Empire was also expressed in the arrangement for the nomination of the annual Consuls. Each Emperor named one of the two consuls for the year. As a general rule the names were not published together. The name of the Western consul was not known in the East, nor that of the Eastern in the West, in time for simultaneous publication.\footnote{There are exceptions to this rule. Occasionally the two Emperors held the consul­ship together, and this was prearranged. It also sometimes happened that one of them resigned his right of nomination to the other, and in this case the two names were published together. \emph{E.g.} in A.D. 437 Valentinian III nominated Aetius and Sigisvultus. The whole subject of the consul­ship in the fifth century and in the Ostrogothic period has been elucidated by Mommsen in \emph{Ostgothische Studien}, in \emph{Hist. Schr.} III.

}

Many passages in our narrative will show that the Empire throughout the fifth century was the one and undivided Roman Empire in all men's minds. There were ``the parts of the East,'' and ``the parts of the West,''\texthebrew{⁠}\footnote{Partes orientis et occidentis.

} but the Empire was one.\texthebrew{⁠}\footnote{Coniunctissimum imperium,

\emph{C. Th.} I.1.5 .

} No one would speak of two or more Roman Empires in the days of the sons of Constantine; yet their political relation to one another was exactly the same as that of Arcadius to Honorius or of Leo I to Anthemius. However independent of each other

p18
or even unfriendly the rulers from time to time may have been, the unity of the Empire which they ruled was theoretically unaffected. And the theory made itself felt in practice.

Although the dyarchy, or double government of Emperor and Senate, had come to an end, and autocracy, as we have seen, was established without reserve or disguise, the Senate remained as an important constitutional body, with rights and duties, and, though it was remodelled, it maintained many of its ancient traditions. The foundation of Constantinople had led to the formation of a second Senate, modelled on that of Rome \texthebrew{—} a great constitutional innovation. Constantine himself had not ventured upon this novelty. He did found a new senate in Byzantium, but his foundation seems rather to have resembled the senates of important cities like Antioch than the august Senatus Romanus.\texthebrew{⁠}\footnote{Cp. the (contemporary)

Anon. Vales. (Part 1) 6.30

senatum constituit secundi ordinis, claros vocavit.

} His son Constantius raised it from the position of a municipal to that of an Imperial body.\footnote{The exaltation of the senate by Constantius is touched on in the Presbeutic speech of Themistius addressed to the Emperor at Rome (\emph{Or.} 3).

}

The principles that senatorial rank was hereditary and that the normal way of becoming a member of the Senate itself was by holding a magistracy still remained in full force. The offices of aedile and tribune had disappeared, and by the end of the fourth century the quaestor ­ship was on the point of disappearing. Hence the praetor ­ship remained as the portal through which the sons of senators could enter the Senate. They not only could, but they were obliged. The sole duty of the Praetor now was to spend money on the exhibition of games or on public works. There were eight praetors in the East; the expenses were divided among them; and the Senate, which had the duty of designating them, named them ten years in advance, in order to enable them to economise or otherwise collect the necessary funds, as the cost of holding the office was extremely heavy.\texthebrew{⁠}\footnote{\emph{C. Th.} VI.4.13, § 2 . Olympiodorus, \emph{fr.} 44, mentions some sums spent on Praetorian games at Rome (£184,000; £92,000; £55,200). These were evidently quite exceptional. The expenses of a consul on the spectacles which he exhibited during the first week of the new year might amount to over £92,000, but were largely defrayed by the Imperial treasury, at least in the sixth century.

Procopius, \emph{H. A.} 26, p159 .

} The burden of the consul ­ship

p19
was not so severe, but that supreme dignity was bestowed only on men who were already senators.

Men who were not born in the senatorial order could be admitted to the Senate in various ways, whether by a decree of the Senate itself or by the Emperor, who might confer either upon an individual or upon a whole class of persons an order of rank which carried with it a seat in the Senate. Persons thus co-opted by the Senate were liable to the burden of the praetor ­ship, and likewise those whom the Emperor ennobled, unless special exemption were granted.

Exemption was granted frequently, and it took the form of adlectio .\texthebrew{⁠}\footnote{Lécrivain, \emph{Le Sénat romain}, 15\texthebrew{‑}23 , gives a lucid account.

} This was the term used in the early Empire for the process by which the Emperor could introduce into the Senate a candidate of his own and make him a member of the aedilician, for instance, or of the praetorian class, though he had never filled the corresponding magistracy. In the fourth century these classes disappeared and were replaced by the three orders of illustres , spectabiles , and clarissimi , in each of which there were certain subdivisions. The Emperor could confer these orders of rank on any one,\texthebrew{⁠}\footnote{It was done by means of a brief or patent of rank (codicilli). The older rank of perfectissimus, which did not carry senatorial rank, still survived, soon to disappear.

} and a person to whom he granted the clarissimate became thereby a member of the lowest order of the Senate, and belonged to the adlecti who were exempt from the praetor ­ship. Further, under the new administrative system which will be described in the following chapter, all the important offices carried with them the title illustris , or spectabilis , or clarissimus , and thus secured to their occupants eventually, if not immediately,\texthebrew{⁠}\footnote{\emph{C. J.} III.24.3

(law of Zeno) seems to imply that the quaestor s. pal., the mag. off., the praepositus s. cub. did not belong to the Senate, although they were illustres, till after they had laid down their offices.

} seats in the Senate. And in some cases, though by no means in all, this admission by virtue of office carried with it exemption. Again, there were many classes of subordinate functionaries who received, when they retired from office, the clarissimate or perhaps one of the higher titles, thus becoming senators, and these as a rule enjoyed exemption.

To resume: the Senate was recruited from men of senatorial origin, that is, sons of senators, and from men who, born outside the senatorial class, were ennobled by elevation to office, or

p20
on retiring from office, or occasionally by a special act of the Emperor or of the Senate. The praetor ­ship was the front gate for entering the Senate, but there was also a back gate, adlection, of which the Emperor held the key, and a large and increasing number of the second section entered by this way.

One of Constantine's administrative reforms was the opening to senators of all the official posts, which hitherto had been confined to the equestrian order, so that the careers open to a young man of senatorial birth were far more numerous and varied. The equestrian order gradually disappeared altogether. On the other hand, men of the lowest origin might rise through the inferior grades of the public service to higher posts which carried with them the right of admission to the Senate. Thus an aristocracy was formed, which was recruited every year by men whose fathers had not belonged to it, and was divided into grades depending on office or special Imperial favour, not on birth.\texthebrew{⁠}\footnote{Within the ranks of the three grades illustres, spectabiles, and clarissimi precedence was determined by office. Thus a Praetorian Prefect was superior to a Master of Soldiers; both were illustrious. A man who was created a spectabilis might be assimilated to a proconsul, a vicarius, or a dux, all of whom were spectabiles, but in descending rank. All these were superior to the viri consulares, who were practically coincident with the class of clarissimi (cp. \emph{C. J.} XII.17.2) . These viri consulares must be carefully distinguished from men who had held the consul­ship or had received the honorary consul­ship, and who were in the highest class of the illustres.

} Ancient tradition was so far preserved that those who had discharged the functions of consul (including honorary consuls) had the most exalted rank.\texthebrew{⁠}\footnote{But among the consuls, a Praetorian Prefect was superior to one who had not held that office, etc.

(\emph{C. Th.} VI.6, 1) .

} Next to the consuls came Patricians, a new order instituted by Constantine, not connected with any office, and conferred \texthebrew{—} at first very sparingly \texthebrew{—} by the Emperor on men highly distinguished for their services to the State.\footnote{In the fifth and sixth centuries the patriciate was bestowed more freely. By a law of Zeno

(\emph{C. J.} XII.3.3)

it could be conferred only on a man who had been Consul, Praetorian Prefect, Prefect of the City, Master of Soldiers, or Master of Offices. In later times, most ministers who would formerly have had the illustrious rank were patricians.

}

A large number of senators preferred living on their estates in the country to residence in the capitals, and of those who actually attended the meetings of the Senate\texthebrew{⁠}\footnote{The quorum for a meeting of the Senate in A.D. 356 was fixed at 50. There was no limit to the number of Senators. Themistius speaks of 2000 in his time (\emph{Or.} 34, \emph{ed. Dindorf}, p456). At the beginning of our period there were no Senators who had not the right to sit in the Senate. But there were some persons who had the clarissimate and yet were not Senators

(\emph{C. Th.} XVI.5.52)

\texthebrew{—}(p21)
apparently those who received this dignity without adlectio and had not discharged the office of praetor. Cp. Lécrivain, \emph{op. cit.} 12 .

} it is probable

p21
that the greater number were men who held official posts and that simple senators were few. We may conjecture that the highest and smallest class, the Illustrious, came to form the majority of the active members of the Senate, and that this fact caused the Emperors before the middle of the fifth century to permit the two inferior classes, Spectabiles and the Clarissimi, to live wherever they pleased.\texthebrew{⁠}\footnote{\emph{C. J.} XII.1.15 .

} A few years later all members of these classes who lived in the provinces were relieved from the Praetorship, and were graciously recommended to stay at home and enjoy their dignities.\texthebrew{⁠}\footnote{\emph{Ib.} 2.1

(A.D. 450).

} This meant that while they belonged to the senatorial class and paid the senatorial taxes, they were expressly discouraged from sitting in the Senate. The next step was to exclude entirely the two lower classes and confine the right of deliberating in the Senate to Illustres, and by the end of the fifth century this seems to have been the rule.\footnote{Cp. Lécrivain, \emph{op. cit.} 66. Add to his references

\emph{Digest}, I.9.12 .

}

The functions of the Senates of Rome and Constantine were both municipal and Imperial. As the funds contributed by the praetors were exclusively applied for the benefit of the capital cities, the nomination of these magistrates and the control exercised over the distribution of the funds belonged to the municipal part of their duties. The Prefect of the City acted as chief of the Senate and as its executive officer, and conducted all its communications with the Emperor.\texthebrew{⁠}\footnote{Illustrated by the \emph{Relationes} of Symmachus, Praef. urb. Cp.
Cassiodorus, \emph{Var.} VI.4 . Under the Prefect was a staff of censuales, who kept the lists, investigated the incomes of the senators, and managed the financial business. Cp. \emph{C. Th.} VI.4.13 and 26 .

} He was the guardian of the rights of senators;\texthebrew{⁠}\footnote{Symmachus, \emph{Rel.} 48 praefecturae urbanae proprium negotium est senatorum iura tutari.

} and that body acted with him as an advisory council on such matters as the food supply of the capital, or the regulation of the public instruction given by professors and rhetors.

We have already seen the constitutional importance of the Senate when a vacancy on the throne occurred. It could pass resolutions ( senatus consulta ) which the Emperor might adopt and issue in the form of edicts.\texthebrew{⁠}\footnote{This is obviously the case with Valentinian III,

\emph{Nov.} 14 ; possibly with Theodosius II,

\emph{Nov.} 15

(as Lécrivain has suggested). The Senate of Rome retained in theory the right leges constituere; but this perhaps never belonged to the Senate of New Rome.

} It could thus suggest Imperial

p22
legislation, and it acted from time to time as a consultative body in co-operation with the Imperial Council. Some of the Imperial laws took the form (we do not know on what principle) of ``Orations to the Senate,'' and were read aloud before that body.\texthebrew{⁠}\footnote{\emph{E.g.} \emph{C. Th.} VIII. 18.9, 10 ;

19.1 . Cp. Symmachus, \emph{Ep.} X.2.

} Valentinian III, in A.D. 446, definitely formulated a legislative procedure which granted to the Senate the right of co-operation . When any new law was to be promulgated, it was first to be discussed at meetings of the Senate and the Council; if agreed to, it was to be drafted (by the Quaestor), and then submitted again to the same bodies, after which it was to be confirmed by the Emperor.\texthebrew{⁠}\footnote{\emph{C. J.} I.14.8 . We cannot be sure whether this procedure was adopted in the East, though it is included in Justinian's Code.

} This regulation points to the probability that it was already the habit frequently to consult the Senate.\footnote{For instance cp. \emph{C. Th.} VI.24.11 ; Marcian, \emph{Nov.} 5, \emph{ad in.}

}

The Senate might act as a judicial court, if the Emperor so pleased, and trials for high treason were sometimes entrusted to it.\texthebrew{⁠}\footnote{John Lydus, \emph{De mag.} III.10 \textgreek{τ\texthebrew{ῶ}ν βασιλέων \texthebrew{ἅ}μα τ\texthebrew{ῇ} βουλ\texthebrew{ῇ} δίκας \texthebrew{ἀ}κροωμένων}, referring apparently to the time of Arcadius. \emph{Ib.} 27 , a reference of an appeal case to the Senate for revision is mentioned.

} For ordinary crimes, Senators were judged by a court consisting of the Prefect of the City and five Senators chosen by lot.\footnote{Quinquevirale iudicium, \emph{C. Th.} IX.1.13 ,

II.1.12 .

}

There were two Senate-houses at Constantinople, one, built by Constantine, on the east side of the Augusteum, close to the Imperial Palace;\texthebrew{⁠}\footnote{\emph{Not. Urbis Cpl.} p231. Sozomen, II.3;

Zosimus, V.24 ;

Procopius, \emph{Aed.} I.10 . \textgreek{σενάτον} is the Greek for Senate-house.

} the other on the north side of the Forum of Constantine.\texthebrew{⁠}\footnote{First referred to in connexion with Theodosius II: \textgreek{Παραστάσεις}, ed. Preger, 50. It was burnt down in the reign of Leo I. Cedrenus, I. pp 610

and 565 .

} It is not clear why two houses were required.\texthebrew{⁠}\footnote{Cedrenus, I.610 (= Zonaras III.125 ) says that the Emperor, when he assumed the consul­ship, was invested with the consular robes in the Senate-house in the Forum. He also mentions in the same passage another house, used for senatorial deliberations, in the Forum of Taurus. Of this we do not hear elsewhere.

} But in the sixth century we are told that the Senate had ceased to meet in its own place and used to assemble in the Palace.\texthebrew{⁠}\footnote{John Lydus, \emph{De mag.} II.9 .

} This change was probably connected with its co-operation with the Imperial Council.

Important decisions as to legislation and public policy were not usually taken by the Emperor on the single advice of the

p23
minister specially concerned. He was assisted by the Consistorium or Imperial Council, which was constantly summoned to deliberate on questions of moment, and we must always remember that, while the Emperor was officially and legally sole author of all laws and responsible for acts of state, the deliberations of the Imperial Council had a large share in the conduct of public affairs. The Consistorium was derived from the legal Consilium of Hadrian, enlarged in its functions and altered in its constitution by Diocletian and Constantine.\texthebrew{⁠}\footnote{Diocletian changed the old name consilium to consistorium, because, under the new autocracy, the members stood (consistere) in the Emperor's presence. Hadrian's consilium had no permanent members; those who assisted at each meeting were summoned \emph{ad hoc.} Constantine instituted permanent members, with the title of comites consistoriani, and included military as well as legal members. Comites was an appropriate name, as the Council accompanied the Emperor as he moved about from camp to camp, or city to city. Constantine bestowed the title of comes (of first, second, or third class) as an honorary distinction, and it was attached to many offices. It corresponded in some ways to our Privy Councillor. Cp. Seeck, \emph{Untergang}, II.76 \emph{sqq.}

} It acted as a high Court, before which important cases, such as treason, might be tried. It was consulted generally on matters of legislation and policy. The Quaestor was its president. It included the two financial Ministers and the Master of Offices; and probably the Praetor Prefect and the Masters of Soldiers who were in residence at the capital generally attended. We have very little information about its size or its constitution; nor do we know how often it met. We have good reason to suppose that it met at stated times, and not merely when convened for a special purpose.\texthebrew{⁠}\footnote{\emph{Nov. Theod. II.} xxiv (A.D. 443). A report concerning the strength of the military forces on the frontiers is to be made quotannis mense Ianuario in sacro consistorio.

} That the transaction of a considerable amount of ordinary business devolved upon it may be inferred from the fact that it disposed of a large bureau of secretaries and officials known as Tribunes and Notaries. These clerks, who had their office in the Palace, drafted the proceedings and resolutions of the Consistorium, and were sometimes employed to execute missions in pursuance of its decisions.\footnote{A certain number of the tribuni et notarii were appointed to special duties as legal secretaries of the Emperor and were often employed on special missions. They were called referendarii. For their functions and appointment see Bury, \emph{Magistri scriniorum}, etc.

}

Among the ordinary duties of the Council was that of receiving deputations from the provinces.\texthebrew{⁠}\footnote{\emph{C. Th.} XII.12.6\texthebrew{‑}10 . In these constitutions the Consistorium is called comitatus noster and sacrarium nostrum.

} But the most important part

p24
of its regular work seems to have been judicial. In serious cases, senators who did not belong to the Council were frequently called to assist.\texthebrew{⁠}\footnote{See the law of Justinian \emph{De ordine senatus},

\emph{Nov.} 62.1

(A.D. 537).

} The technical term for a meeting of the Council was silentium ; a meeting in which the Senate took part was called silentium et conventus .\texthebrew{⁠}\footnote{\emph{Ib.} John Malalas, p438

\textgreek{γενομένου σιλεντίου κομβέντου} (to try a patrician for libelling the Emperor Justinian). Peter Patr. apud Const. Porph. \emph{Cer.} 1. 92, p422 ,

95, p433 , \textgreek{σιλέντιον κα\texthebrew{ὶ} κομέντον} (where we should read \textgreek{κοβέντον}).

} But the words et conventus were frequently dropped;\texthebrew{⁠}\footnote{Justinian, \emph{Nov. cit.}, Etsi non addatur conventus vocabulum. Thus in Theophanes, p246.14 (\textgreek{\texthebrew{ἐ}π\texthebrew{ὶ} σελεντίου}) a silentium et conventus is meant, as is shown by the words \textgreek{παρουχί\texthebrew{ᾳ} τ\texthebrew{ῆ}ς συγκλήτου} below (l. 24).

} and thus it becomes difficult to say in a given case whether a silentium means the Council only or the Council and Senate.\footnote{In connexion with the relations of council and Senate it is worth noti­cing that the words in amplissimo et venerabili ordine (\emph{sc.} the Senate) in a law of Theodosius II,

\emph{C. Th.} VI.23.1 , are replaced in

\emph{C. J.} XII.16.1

by in nostro consistorio.

next \texthebrew{▸}Images with borders lead to more information.

The thicker the border, the more information.

(Details here.)

UP TO:

J. B. Bury
Homepage

Latin \& Greek

Texts

LacusCurtius

Home}

It would seem that, while the Senate and Council continued to be formally distinct, the Senate came virtually to be a larger Council and met in the great hall of council, the Consistorium in the Palace. The Emperor, at his discretion, referred political questions either to this larger body or to a smaller body of functionaries which corresponded to the old Imperial Council. The chief occasions on which the Senate could exercise independent political action were when a vacancy to the throne occurred; but some cases are recorded in which it seems to have taken the initiative in recommending political measures.

\xchapter{Chapter II}

Short URL for this page:

bit.ly/BURLAT2

We pass from the constitution of the monarchy to the bureaucratic system of government which it created. This system, constructed with the most careful attention to details, was a solution of the formidable problem of holding together a huge heterogeneous empire, threatened with dissolution and bankruptcy, an empire which was far from being geographically compact and had four long, as well as several smaller, frontiers to defend. To govern a large state by two independent but perfectly similar machines, controlled not from one centre but from two foci, without sacrifi­cing its unity was an interesting and entirely new experiment. These bureaucratic machines worked moderately well, and their success might have been extraordinary if the monarchs who directed them had always been men of superior ability. Blots of course and defects there were, especially in the fields of economy and finance:

The political creation of the Illyrian Emperors was not unworthy of the genius of Rome.

The old provinces had been split up by Diocletian into small parts, and these new provinces placed under governors whose powers were purely civil. A number of adjacent provinces were grouped together in a circumscription which was called a Diocese (resembling in extent the old province), and the Diocese was under the control of an official whose powers were likewise purely civil. The Dioceses in turn were grouped in four vast

p26
circumscriptions,\texthebrew{⁠}\footnote{During the fourth century, the number of Prefectures was sometimes four, sometimes three; for at times, Italy and Illyricum were under one Prefect. The division of the Empire in 395 stereotyped the quadruple division. Cp. Mommsen, \emph{Hist. Schr.} III.284 \emph{sqq.} \texthebrew{—} For the administrative fabric of the fourth and fifth centuries a main source is the \emph{Notitia dignitatum}, which consists of two distinct documents, the \emph{Not. in partibus Orientis}, and the \emph{Not. in partibus Occidentis.} It was the function of a high official, the primicerius notariorum, to prepare and issue the codicilli or diplomas of their appointments to all the higher officials of the Empire from Praetorian Prefects to provincial governors. The insignia of the office were represented on the codicil, for instance in the case of a Master of Soldiers the shields of the regiments which were under his command. For this purpose the primicerius of the West, and the primicerius of the East, had each a list (laterculum maius) of all the officials in order of precedence, with information as to their staffs and subordinates. The text which we by a lucky chance possess is derived from the lists which were probably in the hands of the primicerius of the West in A.D. 427 or not much later. The

\emph{Not. Or.}

did not strictly concern him, but it was useful for reference, and a copy brought up to date had been sent to him from Constantinople. Compare

Bury, \emph{The Notitia dignitatum}, in \emph{J. R. S.} X .

} under Praetorian Prefects, who were at the head of the whole civil administration and controlled both the diocesan and the provincial governors. This system, it will be observed, differed from the previous system in three principal features: military and civil authority were separated; the provincial units were reduced in size; and two higher officials were interposed between the Emperor and the provincial governor. Perhaps we should add a fourth; for the Praetorian Prefect (whom Constantine had shorn of his military functions) possessed, so far as civil administration was concerned, an immensely wider range of power than any provincial governor had possessed under the system of Augustus.

At the end of the fourth century, then, the whole Empire, for purposes of civil government, was divided into four great sections, distinguished as the Gauls, Italy, Illyricum, and the East ( Oriens ). \emph{The Gauls}, which included Britain, Gaul, Spain, and the north-western corner of Africa, and \emph{Italy}, which included Africa, Italy, the provinces between the Alps and the Danube, and the north-western portion of the Illyrian peninsula, were subject to the Emperor who resided in Italy. \emph{Illyricum}, the smallest of the Prefectures, which comprised the provinces of Dacia, Macedonia, and Greece, and \emph{the East}, which embraced Thrace in the north and Egypt in the south, as well as all the Asiatic territory, were subject to the Emperor who resided at Constantinople. Thus each of the Praetorian Prefects had authority over a region which is now occupied by several modern

p27
States. The Prefecture of \emph{the Gauls} was composed of four Dioceses: Britain, Gaul, Viennensis (Southern Gaul), and Spain; \emph{Italy} of three: Africa, the Italies,\texthebrew{⁠}\footnote{The Italies were divided into two districts, under two Vicarii: the V. urbis Romae, whose district comprised all Italy south of Tuscany and Umbria (inclusive) with Sardinia, Corsica, and Sicily; the V. Italiae, who was over the rest of Italy and Raetia.

} and Illyricum; \emph{Illyricum} of two: Dacia and Macedonia; \emph{the East} of five: Thrace, Asiana, Pontus, Oriens, and Egypt. Each of the diocesan governors had the title of Vicarius ,\texthebrew{⁠}\footnote{There was no Vicarius of Dacia; it was directly subject to the Prefect.

} except in the cases of Oriens where he was designated Comes Orientis , and of Egypt where his title was Praefectus Augustalis .\texthebrew{⁠}\footnote{Egypt had been part of the Diocese of the East till about A.D. 380\texthebrew{‑}382, when it was made a distinct Diocese, and the praefectus Aegypti received the title of Augustalis. Cp. M. Gelzer, \emph{Studien zur byz. Verw. Ägyptens}, 7. The Augustalis seems to have acted at the same time as praeses of the province of Egypt.

} It is easy to distinguish the Prefecture of the Oriens from the Diocese of Oriens (Syria and Palestine); but more care is required not to confound the Diocese with the Prefecture of Illyricum.

The subordination of these officials to one another was not complete or strictly graded. A comparison of the system to a ladder of four steps, the Emperor at the top, the provincial governor at the foot, with the Prefect and the Vicarius between, would be misleading. For not only were the relations between the provincial governor and the Prefect direct, but the Emperor might communicate directly both with the governor of the diocese and with the governor of the province. Two provinces had a special privilege: the proconsuls of Africa and of Asia\texthebrew{⁠}\footnote{Under the proconsul of Asia were two provinces, Hellespontus and Insulae (the islands along the coast of Asia Minor):

\emph{Not. Or.} XX .

} were outside the jurisdiction either of Vicarius or of Prefect, and were controlled immediately by the Emperor.\footnote{The governor of one other province, Achaia, bore the old title of proconsul; the others were consulares or correctores or praesides. The governor had judicial as well as administrative powers. His court was the court of first instance in his province; but an appeal lay either to the court of the Vicarius or to that of the Prefect. He had also the duty of supervising the collection of taxes.

}

The Praetorian Prefect of the East, who resided at Constantinople, and the Praetorian Prefect of Italy were in rank the highest officials in the Empire; next to them came respectively the Prefect of Illyricum, who resided at Thessalonica, and the Prefect of the Gauls. The functions of the Prefect embraced a wide sphere; they were administrative, financial, judicial,

p28
and even legislative. The provincial governors were appointed at his recommendation, and with him rested their dismissal, subject to the Emperor's approval. He received regular reports of the administration throughout his prefecture from the Vicarii and from the governors of the provinces. He had treasuries of his own, and the payment and the food supplies of the army devolved upon him. He was also a supreme judge of appeal; in cases which were brought before his court from a lower tribunal there was no further appeal to the Emperor. He could issue, on his own authority, praetorian edicts, but they concerned only matters of detail. The most important Imperial enactments were usually addressed to the Prefects, because they were the heads of the provincial administration, and possessed the machinery for making the laws known throughout the Empire.

The exalted position of the Praetorian Prefect was marked by his purple robe, or mandyes , which differed from that of the sovran only in being shorter, reaching to the knees instead of to the feet. His large silver inkstand, his pen-case of gold weighing \texthebrew{•} 100 lbs., his lofty chariot, are mentioned as three official symbols of his office. On his entry all military officers were expected to bend the knee, a survival of the fact that his office was originally not civil but military.

Rome and Constantinople, with their immediate neighbourhoods, were exempt from the authority of the Praetorian Prefect and under the jurisdiction of the Prefect of the City.\texthebrew{⁠}\footnote{\textgreek{\texthebrew{Ὁ ἔ}παρχος τ\texthebrew{ῆ}ς πόλεως.}} The Prefect of Constantinople had the same general powers and duties as the Prefect of Rome, though in some respects the arrangements were different. He was the head of the Senate, and in rank was next to the Praetorian Prefects. While all the other great officials, even though their functions were purely civil, had a military character, in token of which they wore military dress and the military belt, the Prefect of the City retained his old civil character and wore the toga. He was the chief criminal judge in the capital. For the maintenance of further order the Roman Prefect had under his control a force of city cohorts, as well as police. We hear nothing of any institution at Constantinople corresponding to the city cohorts, but the police ( vigiles ) were organised as at Rome under a

p29
praefectus vigilum ,\texthebrew{⁠}\footnote{\textgreek{Ηυκτέπαρχος}. The page of the

\emph{Not. dig.}

appertaining to the Prefect of Constantinople is unfortunately lost.

} subject to the Prefect. For the care of the aqueducts and the supervision of the markets the Prefect was responsible. One of his most important duties was to superintend the arrangements for supplying the city with corn. º \texthebrew{⁠}\footnote{In Rome there was a subordinate official, praefectus annonae, who presided over this department; and there was a praefectus annonae in Africa, who was under the Praetorian Prefect. At Constantinople there was no pr. ann., but the pr. ann. at Alexandria, where the corn was shipped, seems to have been under the Prefect of the City.

} He had also control over the trade corporations ( collegia ) of the capital.

The supreme legal minister was the Quaestor of the Sacred Palace. His duty was to draft the laws, and the Imperial rescripts in answer to petitions. A thorough knowledge of jurisprudence and a mastery of legal style were essential qualifications for the post.\footnote{His functions in regard to petitions involved co-operation with the Magistri scriniorum, and the Scrinia supplied him with assistants; he had no special staff of his own.

}

The post of Master of Offices ( magister officiorum ) had grown from small beginnings and by steps which are obscure into one of the most important ministries.\texthebrew{⁠}\footnote{In the

\emph{Not. dig.}

he precedes in rank the Quaestor, but this was only a temporary arrangement. \textgreek{Μάγιστρος}, when unqualified, in Greek writers always means the Mag. Off.

} It comprised a group of miscellaneous departments, unrelated to each other, and including some of the functions which had belonged to the pre-Constantinian Praetorian Prefects. Officium was the word for the body of civil servants ( officiales ) who constituted the staff of a minister or governor, and the Master of Offices was so called from the authority which he exercised over the civil service, but especially over the secretarial departments in the Palace.

There were three principal secretarial bureaux ( scrinia ), which had survived from the early Empire, and retained their old names: memoriae , epistularum , and libellorum .\texthebrew{⁠}\footnote{See Karlowa, \emph{op. cit.} I.834 \emph{sqq.}; Schiller, \emph{op. cit.} II.102 \emph{sqq.}; Bury, \emph{Magistri Scriniorum}, etc.

} At Constantinople the second bureau had two departments, one for Latin and one for Greek official correspondence. The secretarial business was conducted by magistri scriniorum ,\texthebrew{⁠}\footnote{The Greek title was \textgreek{\texthebrew{ἀ}ντιγραφ\texthebrew{ῆ}ς}, Bury, \emph{ib.} 24 \emph{sqq.}

} who were in direct touch with the Emperor and were not subordinate to any higher official. They were not, however, heads of the

p30
bureaux, but the bureaux, which were under the control of the Master of Offices, supplied them with assistants and clerks.\footnote{The Magister memoriae drafted brief Imperial decisions (adnotationes, on the margin of documents), answered petitions, and probably threw into their final form many of the documents emanating from the offices of the other magistri. The Magistri epistularum and epistularum Graecarum dealt with answers to communications from foreign powers and to deputations from the provinces; examined the questions addressed to the Emperor by officials; and also dealt with petitions. The duties of the Magister libellorum were concerned chiefly with appeal cases (cognitiones) and petitions which involved specifically legal questions. We have not sufficient information to draw a sharp line between the functions of these three ministers, which seem at many points to have overlapped and involved constant co-operation. They must also have been in constant touch both with the Master of Offices and with the Quaestor.

}

With the three ancient and homogeneous scrinia was associated a fourth,\texthebrew{⁠}\footnote{They are sometimes grouped together as sacra scrinia nostra.

} of later origin and at first inferior rank, the scrinium dispositionum , of which the chief official was the comes dispositionum . His duty, under the control of the Master of Offices, was to draw up the programme of the Emperor's movements and to make corresponding arrangements.

The Master of Offices was responsible for the conduct of court ceremonies, and controlled the special department\texthebrew{⁠}\footnote{Officium admissionum under a magister.

} which dealt with ceremonial arrangements and Imperial audiences. The reception of foreign ambassadors thus came within his scope, and he was the head of the corps of interpreters of foreign languages. In the Roman Empire the administrations of foreign and internal affairs were not sharply separated as in modern states, but the Master of Offices is the minister who more than any other corresponds to a Minister of Foreign Affairs. As director of the State Post ( cursus publicus )\texthebrew{⁠}\footnote{It had been under the control of the Praet. Prefect, who still retained the right of issuing passes or orders for its use. The change was made in A.D. 396; see below,

p115 .

} he made arrangements for the journeys of foreign embassies to the capital.

One of his duties was the control of the agentes in rebus , a large body of officials who formed the secret service of the State and were employed as Imperial messengers and on all kinds of confidential missions. As secret agents they were ubiquitous in the provinces, spying upon the governors, reporting the misconduct of officials, and especially vigilant to secure that the state post was not misused. Naturally they were open to bribery and corruption. The body or schola of agentes was strictly organised in grades, and when they had risen by regular

p31
promotion, they were appointed to be heads ( principes ) of the official staffs of diocesan and provincial governors, and might rise to be governors themselves. Their number, in the East, was over 1200.\footnote{They are often called magistriani (as under the authority of the Mag. Off.). In 430 there were more than 1174

(\emph{C. Th.} VI.27.23) ; in the reign of Leo I the number was 1248

(\emph{C. J.} XII.20.3) .

}

The Scholarian bodyguards, organised by Constantine,\texthebrew{⁠}\footnote{See below,

p37 .

} were subject to the authority of the Master of Offices, so that in this respect he may be regarded as a successor of the old Praetorian Prefect. He also possessed a certain control over the military commanders in frontier provinces.\texthebrew{⁠}\footnote{See

\emph{C. J.} XII.59.8 ; \emph{Nov. Theodosii} 24. Perhaps he inherited this duty from the Praet. Pref. in A.D. 396.

} He became (in A.D. 396) the director of the state factories of arms. In the Eastern half of the Empire there were fifteen of these factories ( fabricae ), six in the Illyrian peninsula, and nine in the Asiatic provinces.

One of the most striking features of the administrative system was the organisation of the subordinate officials, who were systematically graded and extremely numerous.\footnote{A short survey of this complicated subject will be found in Karlowa, \emph{Röm. Recht}, I.875 \emph{sqq.}

}

Our use of the words ``office'' and ``official'' is derived from the technical meaning of officium , which, as was mentioned above, denoted the staff of a civil or military dignitary.\texthebrew{⁠}\footnote{In Greek, \textgreek{τάξις} was used as well as \textgreek{\texthebrew{ὀ}φφίκιον}, and, for the members, \textgreek{ταξε\texthebrew{ῶ}ται} as well as \textgreek{\texthebrew{ὀ}φφικιάλιοι}. Apparitores (used in the early Empire for officials) is sometimes applied to members of the more important, cohortalini to those of the least important, offices. In the military officia the posts were confined to soldiers.

} Most ministers, every governor, all higher military commanders, had an officium, and its members were called officiales . Theoretically, the civil as well as the military officials were supposed to be soldiers of the Emperor; their service was termed militia , its badge was the military belt, which was discarded when their term of service expired, and their retirement from service was called in military language ``honourable dismissal'' ( honesta missio ). But these usages were a mere survival, and the state service was really divided into military, civil, and palatine offices. The term palatine in this connexion meant particularly the staffs of the financial ministers, the Counts of the Sacred Largesses and the Private Estates.

p32
The number of subalterns in each office was fixed. To obtain a post an Imperial rescript was required, and advancement was governed by seniority. Those who had served their regular term in the higher offices became eligible for such a post as the governor ­ship of a province and might rise to the highest dignities in the Empire.

Offices, such as those of a Praetorian Prefect, a vicar, or a provincial governor, were divided into a number of departments or bureaux ( scrinia ), each under a head. On these permanent officials far more than on their superior, who might only hold his post for a year, the efficiency of the administration depended. The bureaux differed in nature and name according to the functions of the ministry. Those in the office of the Praetorian Prefecture differed entirely from those of the financial ministries or those of the Master of Offices. But the offices of all the governors who were under the Praetorian Prefect reproduced in their chief departments the office of the Prefect himself. Each of these had a princeps ,\texthebrew{⁠}\footnote{The princeps of the Prefect, the vicars, and the proconsuls was selected from the agentes in rebus. Strictly speaking he was outside the officium, though he is included in it in the

\emph{Not. dig.}

The officium consisted of the cornicularius, who assisted the chief in administering justice; a criminal department under a commentariensis, who brought the accused to trial, drew up the acts of the process, executed judgment, superintended prisons; a section of accountants (numerarii), who dealt with fiscal business; the adiutor (\textgreek{βοηθός}), and some others. Outside the officium there were attached a number of organised bodies (scholae) of clerks and assistants of various kinds, who were at the disposal of the officials, especially the school of exceptores or shorthand writers, the most expert of whom formed an inner college of augustales (cp. John Lydus, \emph{De mag.} III.9 ). Other schools were the chartularii; the singulares (employed as messengers to the provinces); the scriniarii. From these the chief officials selected their clerks, who then became members of the officium.\texthebrew{—} The military staffs had a princeps and a commentariensis, but as they had no jurisdiction in civil cases they did not require a cornicularius or adiutor.

} who was the right hand of the chief and had a general control over all departments of the office.

The State servants were paid originally (like the army) both in kind and coin, but as time went on the annona or food ration was commuted into money. They were so numerous that their salaries were a considerable item in the budget. We have no information as to the total number of State officials; but we have evidence which may lead us to conjecture that the civil servants in the Prefectures of the East and Illyricum, including the staffs of the diocesan and provincial governors, cannot

p33
have been much fewer than 10,000.\texthebrew{⁠}\footnote{The offices of the provincial governors in Illyricum consisted of about 100 persons

(\emph{C. J.} XII.57.9) ; the maximum number in the vicariates was fixed at 300 (\emph{ib.} I. 15.5 , cp. 12

A.D. 386), but that of the vicariate of Asia was 200, and that of the Count of the East 600 (\emph{ib.} I. 15.13 ;

I.13.1 ). A calculation based on these figures for the dioceses and provinces of the Orient and Illyricum, as enumerated in

\emph{Not. dig.}

would give about 8000, to which we must add probably more than 1000 for the offices of the Prefects. Justinian's ordinances

(\emph{C. J.} I.27.1) , creating the Pr. Prefecture of Africa in the sixth century, gives the numbers and salaries of the officials both of the Prefect and of the provincial governors. There were 396 in the bureaux of the Prefect's office (including the scholiae), and each of the seven civil governors had a staff of 50. Including the salaries of the Prefect and the governors, the total cost amounted to nearly £11,000. The salary of the Prefect was 7200 solidi (£4500), that of a governor, 448 (£280). The staffs of the five military governors (dukes) were paid at a higher rate than those of the civil and the total cost of their establishments was £7050. The incomes of the subordinate officials, who handled legal matters, were considerably increased by fees; the salaries of all the subalterns were miserable.

} To this have to be added the staffs of the military commanders, of the financial and other central ministries.

It was a mark of the new monarchy that the eunuchs and others who held posts about the Emperor's person and served in the palace should be regarded as standing on a level of equality with the State officials and have a recognised position in the public service. The Grand Chamberlain ( praepositus sacri cubiculi ), who was almost invariably a eunuch, was a dignitary of the highest class. In the case of weak sovrans his influence might be enormous and make him the most power ­ful man in the State; in the case of strong Emperors who were personally active he seldom played a prominent part in politics. It is probable that he exercised a general authority over all officers connected with the Court and the Imperial person, but this power may have depended rather on a right of co-operation than on formal authority.\texthebrew{⁠}\footnote{The pages relating to the praepositus in the

\emph{Not. dig.}

(both \emph{Or.} and \emph{Occ.}) are lost. The primicerius sacri cubiculi, chief of staff of the bedchamber, may have been nominally or partially independent of the praepositus; he was a spectabilis ranking immediately after the Counts of the Domestics. It is not clear what the relations of the praepositus were to the castrensis sacri palatii, who appears to have controlled many of the servants of the Great Palace at Byzantium, besides supervising stewards and caretakers in the various Imperial residences (curae palatiorum). Imperial rescripts were sometimes addressed to him. The Count of the Wardrobe (com. sacrae vestis) was probably under the praepositus, as were the decurions and the silentiarii, ushers who kept guard at the doors during meetings of the Imperial Council and Imperial audiences.\texthebrew{—} The Empress had a staff of

cubicularii

of her own; and there was a praepositus sacri cubiculi Augustae, at least in the reign of Anastasius I.

(\emph{C. J.} XII.5.3 and 5) .

} At Constantinople the Grand

p34
Chamberlain had a certain control over the Imperial estates in Cappadocia which supplied the Emperor's privy purse.\footnote{See below,

p52 .

}

We have already seen\texthebrew{⁠}\footnote{See above,

p19 .

}that all the higher officials in the Imperial service belonged to one or other of the three classes of rank, the illustres , spectabiles , and clarissimi ,\texthebrew{⁠}\footnote{\textgreek{\texthebrew{Ἰ}λλούστριοι, περίβλεπτοι, λαμπρότατοι} were the official equivalents in Greek. Between A.D. 460 and 550 all illustres seem to have also a right to be addressed by the title of magnificus, \textgreek{μεγαλοπρεπής}. See Koch, \emph{Die byzantinischen Beamtentitel}, 51 (a book which must be used with caution).

} and were consequently members of the senatorial order. The heads of the great central ministries, the commanders-in\texthebrew{‑}chief of the armies,\texthebrew{⁠}\footnote{Also the Comites domesticorum.

} the Grand Chamberlain, were all illustres . The second class included proconsuls, vicars, the military governors in the provinces, the magistri scriniorum , and many others. The title clarissimus , which was the qualification for the Senate, was attached \emph{ex officio} to the governor ­ship of a province, and to other lesser posts. It was possessed by a large number of subaltern civil servants and was bestowed on many after their retirement. The liberality of the Emperors in conferring the clarissimate gradually detracted from its value. In consequence of this it was found expedient to raise many officials, who would formerly have been clarissimi to the rank of spectabiles ; and this in turn led to a cheapening of the rank of illustres . The result was that before the middle of the sixth century a new rank of gloriosi \texthebrew{⁠}\footnote{Also gloriosissimi. In Greek, \textgreek{\texthebrew{ἐ}νδοξότατοι} (also \textgreek{\texthebrew{ἔ}νδοξοι}). The gloriosissimi senatores are clearly marked off from the illustres in Justinian, \emph{Nov.} 43.1 (A.D. 537).

} was instituted, superior to that of illustres, and the highest officials are henceforward described as gloriosi .

The principal features in which the military establishment of the fourth century\texthebrew{⁠}\footnote{Mommsen, \emph{Das röm. Militärwesen seit Diocletian} (\emph{Hist. Schr.} III.206 \emph{sqq.}) is the principal work on the subject. A summary of the reorganisation by Reid will be found in \emph{C. Med. H.} I.44 \emph{sqq.} It is treated very fully in Grosse, \emph{Röm. Militargeschichte.} \texthebrew{—} Recent investigation has shown that Gallienus initiated changes, especially in regard to the organisation of the cavalry, that prepared the way for the reforms of Diocletian. Cp. Homo, ``L'Empereur Gallien,'' in \emph{Rev. Hist.} 1 \emph{sqq.}, 225 \emph{sqq.}, 1913; Ritterling, ``Zum röm. Heerwesen,'' in \emph{Festschr. f. O. Hirschfeld}, 1903.

} differed from that of the Principate were the existence of a mobile field army, the organisation of the

p35
cavalry in bodies independent of the infantry, and the smaller size of the legionary units.

Diocletian had created, and Constantine had developed, a field army in which the Emperor could move to any part of his dominion that happened to be threatened, while at the same time all the frontiers were defended by troops permanently stationed in the frontier provinces. The military forces, therefore, consisted of two main classes: the mobile troops or comitatenses , which accompanied the Emperor in his movements and formed a ``sacred retinue'' ( comitatus ); and the frontier troops or limitanei .

The strength of the old Roman legion was 6000 men. The legion of this type was retained in the case of the limitanei ; but it is broken up into detachments of about 1000 (corresponding to the old cohort), which are stationed in different quarters, sometimes in different provinces. And these detachments are no longer associated with a number of foot-cohorts and squadrons of horse, as of old, when the legatus of a legion commanded a body of about 10,000 men. The cavalry and the cohorts are under separate commanders.\footnote{There were cohorts as a rule among the frontier troops, but on the Danube, where there were auxilia, cohorts are exceptional. The cavalry squadron, ala, is generally 600 strong. Other classes of the cavalry of the limitanei were known as cunei equitum and equites.

}

The field army consisted of two classes of troops, the simple comitatenses and the palatini .\texthebrew{⁠}\footnote{Constantine, who formed the Palatini, increased the field army and withdrew many troops from the frontier provinces for the purpose. These new bodies were called pseudo-comitatenses (18 legions in the west, 20 in the east).

} The palatini, who took the place of the old Praetorian guards, were a privileged section of the comitatenses and retained the special character of Imperial guards, in so far as most of them were stationed in the neighbourhood of Constantinople or in Italy.\texthebrew{⁠}\footnote{Of the 12 palatine legions in the west, 8 were in Italy, 3in Africa, 1 in Gaul. Of the 13 in the east, 12 were near Constantinople, 1 in Illyricum. Of the 65 auxilia in the west, 21 were in Italy; of 43 in the east, 35 were near the capital. So the

\emph{Not. dig.}

See Mommsen, \emph{op. cit.} 235.

} The infantry of the field army was composed of small legions of 1000, and bodies of light infantry known as auxilia which were now mainly recruited from Gauls, and from Franks and other Germans. The cavalry, under a separate command, consisted of squadrons, called vexillationes , 500 strong.

Each of these units, \texthebrew{—} the legion, the auxilium , the vexillatio

p36
of the comitatenses, the legionary detachment, the cohort of the limitanei, \texthebrew{—} was as a rule under the command of a tribune, in some cases of a praepositus.\texthebrew{⁠}\footnote{Cp. Grosse, on tribune, praepositus praefectus, \emph{op. cit.} pp143\texthebrew{‑}151.

} The tribune corresponded roughly to the modern colonel.

All these armies were under the supreme command of Masters of Soldiers, magistri militum . The organisation of this command in the east, as it was finally ordered by Theodosius I, differed fundamentally from that in the west. In the east there were five Masters of Horse and Foot. Two of these, distinguished as Masters in Presence ( in praesenti , in immediate attendance on the Emperor), resided at Constantinople, and each of them commanded half of the Palatine troops. The three others exercised independent authority over the armies stationed in three large districts, the West, Thrace, and Illyricum.\footnote{The magistri in praes. had precedence over the others, and seem to have exercised some slight control (cp. \emph{C. J.} XII.35.18 ), but not so as to violate the principle of co-ordination.

}

It was otherwise in the west. Here instead of five co-ordinate commanders we find two masters in praesenti , one of infantry and one of cavalry. The Master of Foot was the immediate commander of the infantry in Italy and had superior authority over all the infantry of the field army in all the dioceses, and also over the commanders of the limitanei. In the dioceses the commanders of the comitatenses had the title of military counts.\footnote{Comites rei militaris.\texthebrew{—} The comitatenses in Africa were under the immediate command of the duces of the limitanei. In regard to the titles comes and dux it is to be observed that every dux had the rank of comes, but usually of the second class. When he was a comes of the first class he was called comes et dux, and then simply comes.

}

According to this scheme the Master of Horse in praesenti was co-ordinate with the Master of Foot. But this arrangement was modified by investing the Master of Foot with authority over both cavalry and infantry; he was then called Master of Horse and Foot, or Master of Both Services, magister utriusque militiae , and had a superior authority over the Master of Horse. In the last years of Theodosius the command of the western armies was thus centralised in the hands of Stilicho, and throughout the fifth century this centralisation, giving enormous power and responsibility to one man, was, as we shall see, the rule.

The limitanei were under the command of dukes, the successors of the old legati pro praetore of the Augustan system. In the west the duke was subordinate to the Master of Foot;

p37
in the east to the Master of Soldiers in the military district to which his province belonged.\footnote{\emph{E.g.} the dux of Osroene to the mag. mil. per orientem. There were 12 dukes in the west; 13 in the east, where there were also two of superior rank, the count of the limes of Egypt and the count of Isauria. The province of Isauria was treated exceptionally like frontier provinces on account of the wild, insubordinate character of its uncivilised mountaineers. For the same reason the civil powers were invested in the military governor; the count was also the praeses. Other exceptions to the rule of separating civil from military functions were Arabia and Mauretania Caesariensis. The union of functions was sometimes temporarily introduced, \emph{e.g.} in Sardinia ( \emph{C. Th.} IX.27.3 , A.D. 382), Tripolitania ( \emph{ib.} XII.1.133 , A.D. 393). Before A.D. 450 the duke of the Thebaid, which had been divided into two provinces, was praeses of the upper province (cp. Gelzer, \emph{Byz. Verw. Ägyptens}, º p10); and on some occasions the Augustal Prefect of Egypt was invested with military powers.

}

The Palatine legions were the successors of the old Praetorian guards, but Constantine or one of his predecessors organised guard troops who were more closely attached to the Imperial person.\texthebrew{⁠}\footnote{Cp. Babut, \emph{La Garde impériale}, § XI. p262, who thinks that they replaced the Equites singulares Augusti.

} These were the Scholae, destined to have a long history. We associate the name of School with the ancient Greek philosophers, who gave leisurely instruction to their schools of disciples in Athenian porticoes. It was applied to Constantine's guards because a portico was assigned to them in the Palace\texthebrew{⁠}\footnote{Procop. \emph{H.A.} 14.

} where they could spend idle hours waiting for Imperial orders. The Scholarians were all picked men, and till the middle of the fifth century chiefly Germans; mounted, better equipped and better paid than the ordinary cavalry of the army. There were seven schools at Constantinople, each 500 strong\texthebrew{⁠}\footnote{\emph{Not. dig.} \emph{Or.} 11 . Five in the west

(\emph{Occ.} 9)

and this was perhaps the original number.

} and commanded by a tribune who was generally a count of the first rank.\texthebrew{⁠}\footnote{\emph{C. Th.} VI.13.1 ; \emph{Nov. Theod.} 21. The title tribune was dropped in the course of the fifth century; and these officers were known till late times as Counts of the Schools (\textgreek{κόμητες σχολ\texthebrew{ῶ}ν}).

} We have already seen that the whole guard was under the control of the Master of Offices. Closely associated with the Scholarians was a special body of guards, called candidati from the white uniforms which they wore.

While the Scholarians and Candidates were in a strict sense bodyguards of the Imperial person and never left the Court except to accompany the Emperor, there was another body of guards, the Domestici, consisting both of horse and foot, who as a rule were stationed at the Imperial Court, but

p38
might be sent elsewhere for special purposes.\texthebrew{⁠}\footnote{\emph{C. Th.} VI.24.3

where praesentales are distinguished from non in praes. The full title of the domestics was protectores et domestici.\texthebrew{—} The question of the protectores is difficult. We have to distinguish the Protectors who formed the Schola prima scutariorum in the Scholarian Guards from the Protectors who belonged to a sort of school for officers and were under the orders of the Masters of Soldiers. The discussion of Babut, \emph{op. cit.}, has not definitely cleared up the questions connected with the Protectors. See also Grosse, \emph{op. cit.} 138 \emph{sqq.}

} They were under the command of Counts ( comites domesticorum ) who were independent of the Master of Soldiers.\texthebrew{⁠}\footnote{In the

\emph{Not. dig.}

we find two comites, a comes equitum and a comes peditum, in both east and west, but it seems probable that the command was not always thus divided. For the evidence see Seeck, \emph{sub} ``Comites'' in P.\texthebrew{‑}W. col.548.

} It will be observed that most of the new military creations of the third and fourth centuries had names indicating their close relation to the autocrat, comitatenses, soldiers of the retinue; palatines, soldiers of the palace; domestics, soldiers of the household.

The army of this age had a large admixture of men of foreign birth, and for the historian this perhaps is its most important feature. In the early Empire the foreigner was excluded from military service; the legions were composed of Roman citizens, the auxilia of Roman subjects. Every able-bodied citizen and subject was liable to serve. Under the autocracy both these principles were reversed. The auxilia were largely recruited from the barbarians outside the Roman borders; new troops were formed, designated by foreign names; and the less civilised these soldiers were the more they were prized.\texthebrew{⁠}\footnote{Mommsen, \emph{ib.} 247.

} Some customs and words\texthebrew{⁠}\footnote{Drungus (\textgreek{δρο\texthebrew{ῦ}γγος}), a body of infantry in close formation (cp. Vegetius, \emph{Ep. r. mil.} III.16) is Germanic, and so is bandum (\textgreek{βάνδον}), which the Greeks used as the regular term for military standard (\textgreek{σημε\texthebrew{ῖ}ον}). It may be noted here that in the fourth and fifth centuries the standard of the legion and the legionary detachment seems to have been the dragon. Though the eagle, the standard of the old legion, is sometimes mentioned, it probably went out of use gradually. See Grosse, \emph{op. cit.} 230 \emph{sqq.}

} illustrate the influence which the Germans exercised in the military world. The old German battle-noise , the barritus , was adopted as the cry of the Imperial troops when they went into battle. The custom of elevating a newly-proclaimed Emperor on a shield was introduced by German troops in the fourth century. It would be interesting to know how many Germans there were in the army. The fact that most of the soldiers whom we know to have held the highest posts of command in the last quarter of the fourth century were of German origin speaks for itself.

p39
The legions continued to be formed from Roman citizens; but the distinction between citizens and subjects had disappeared since the citizen ­ship had been bestowed, early in the third century, upon all the provincials, and it was from the least civilised districts of the Empire, from the highlands of Illyricum, Thrace, and Isauria, from Galatia and Batavia, that the mass of the citizen soldiers were drawn. From a military point of view highly civilised provinces like Italy and Greece no longer counted. The legions and citizen cavalry ceased to have a privileged position. For instance, the auxilia on the Danube frontier, who were chiefly of barbarian race, were superior in rank to the legionary troops under the same command.

It was a natural consequence of this new policy, in which military considerations triumphed over the political principle of excluding foreigners, that the other political principle of universal liability to service should also be relinquished. It was allowed to drop. In the fifth century it had become a dead letter, and Valentinian III expressly enacted that ``no Roman citizen should be compelled to serve,'' except for the defence of his town in case of danger.\footnote{\emph{Nov.} 5.

}

A third ancient principle of the Roman State, that only freemen could serve in the army, was theoretically maintained,\texthebrew{⁠}\footnote{\emph{C. Th.} VII.13.8 ;

\emph{Digest}, XLIX.16.11 .

} and though it was often practically evaded and occasionally in a crisis suspended,\texthebrew{⁠}\footnote{Mommsen, 250-251. In the danger of Italy, invaded by barbarians, in A.D. 406 slaves were invited to serve for the reward of liberty,

\emph{C. Th.} VII.13.16 .

} it is probable that there were never many slaves enrolled.

If we examine the means by which the army was kept up, we find that the recruits may be divided into four classes. (1) There were the numerous poor adventurers, Roman or foreign, who voluntarily offered themselves to the recruiting officer and received from him the pulveraticum (`` dust-money ,'' or travelling expenses), the equivalent of the King's shilling. (2) There were the recruits supplied by landed proprietors from among their serf-tenants . This was a State burden, but it fell only on the estates in certain provinces.\texthebrew{⁠}\footnote{\emph{C. Th.} VII.13.2

per eas provincias a quibus corpora flagitantur. In other provinces the proprietors could make a money payment instead of furnishing the men.

} (3) The son of a soldier was bound to follow his father's profession. But this hereditary military

p40
service fell into abeyance before the time of Justinian. (4) The settlements of foreign barbarians within the Empire were another source of supply. These foreigners ( gentiles ), incorporated in the Empire but not enjoying the personal rights of a Roman,\texthebrew{⁠}\footnote{For instance, such a foreigner could not marry a Roman woman. See Mommsen, \emph{Hist. Sch.} III.168.

} were chiefly Germans and Sarmatians, and they were organised in communities under the control of Roman officers. They are found in Gaul, where they had the special name of laeti ,\texthebrew{⁠}\footnote{\emph{C. Th.} VII.20.12

laetus Alamannus Sarmata; in

\emph{Not. Occ.}

we meet laeti Franci, l. gentiles Suebi, Sarmatae et Taifali gentiles.

} and in the Alpine districts of Italy.

The Imperial army was democratic in the sense that the humblest soldier, whatever his birth might be, might attain to the highest commands by sheer talent and capacity. The first step was promotion to the posts of centenarius and ducenarius , who discharged the duties of the old centurions and our non-commissioned officers.\texthebrew{⁠}\footnote{Cp. Vegetius, \emph{op. cit.} II.8 .

} Having served in these ranks the soldier could look forward to becoming a tribune, with the command of a military unit,\texthebrew{⁠}\footnote{Before becoming a tribune, it was usually necessary perhaps to serve in the school of protectors. The three ranks protector, tribunus, comes (et tribunus) appear \emph{e.g.} in

Ammian. XXX.7.3 ;

Vegetius, III.10 , and can be illustrated by inscriptions. But I do not think that Babut (\emph{op. cit.}) is right in regarding the \emph{protectors} as equivalent to the centurions under a new name and organisation.

} and the efficient tribune would in due course receive the rank of comes .

In order to follow the history of the fifth century intelligently and understand the difficulties of the Imperial government in dealing with the barbarian invaders it would be of particular importance to know precisely the strength of the military forces at the death of Theodosius.

The strength of the Roman military establishment at the beginning of the third century seems to have been about 300,000. It was greatly increased under Diocletian; and considerable additions were made in the course of the fourth century. The data of the \emph{Notitia dignitatum} would lead to the conclusion that about A.D. 428 the total strength considerably exceeded 600,000.\footnote{Mommsen's estimate (\emph{op. cit.} 263) based on the \emph{Notitia} is: Limitanei (infantry 249,500, cavalry 110,500) 360,000; Comitatenses (infantry 148,000, cavalry 46,500) 194,500. Total, 554,500. But to this have to be added the limitanei of Italy, Africa, Gaul, and Britain, and they must have amounted to not much less than 100,000. If we estimate them at 90,000 we should get the figure 645,000, which according to Agathias

(V.13)

ought to represent the total force of the Empire. Agathias must have derived this figure from some document of the fourth century. John Lydus (\emph{De (p41)
mens.} I.27) states that under Diocletian the strength of the army was 389,074, and that Constantine doubled it (the latter part of the statement is certainly an exaggeration). We are told that it was further increased by Valentinian I

(Zosimus IV.12.1)

but declined under Theodosius (\textgreek{μεμείωτο,}\emph{ib.} 29.1 ).

}

p41
We have, however, to reckon with the probability that the legions and other military units enumerated in the \emph{Notitia} were not maintained at their normal strength and in some cases may have merely existed on paper. We may conjecture that if the army once actually reached the number of 650,000 it was not after the death of Theodosius, but before the rebellions of Maximus and Eugenius, in which the losses on both sides must have considerably reduced the strength of the legions. But if we confine ourselves to the consideration of the field army, there seems no reason to doubt that in A.D. 428 it was nearly 200,000 strong. It was unequally divided between east and west, the troops assigned to the west being more numerous. In Italy there were about 24,500 infantry and 3500 cavalry.\footnote{The distribution of the troops in the west \emph{c.} 428 is given in

\emph{Not. Occ.} VII ; there is no corresponding section in

\emph{Or.}

In Africa there were 11,500 infantry, 9500 cavalry; in Spain 10,500 infantry; in western Illyricum 14,000 infantry; in Gaul 39,000 infantry, 5500 cavalry. Cp. Bury, \emph{Not. Dig.}

}

The military organisation of Rome, as it existed at the end of the fourth century, was to be completely changed throughout the following hundred years. We have no material for tra­cing the steps in the transformation; of the battles which were fought in this period not a single description has come down to us. But we shall see, when we come to the sixth century, for which we have very full information, that the military forces of the Empire were then of a different character and organised on a different system from those which were led to victory by Theodosius the Great. These changes partly depended on a change in military theory. The conquests of Rome had always been due to her infantry, the cavalry had always been subsidiary, and, down to the second half of the fourth century and the success ­ful campaigns of Julian on the Rhine, experience had consistently confirmed the theory that battles were won by infantry and that squadrons of horse were only a useful accessory arm. The battle of Hadrianople, in which the East German horsemen rode down the legions, shook this view, and the same horsemen who had defeated Valens showed afterwards in the battles which they helped Theodosius to win, how effective might be large bodies of heavy cavalry, armed with lance and

p42
sword. The lesson was not lost on the Romans, who during the following generations had to defend their provinces against the inroads of East German horsemen, and the leading feature of the transformation of the Imperial army was the gradual degradation of the infantry until it became more or less subsidiary to the cavalry on which the generals depended more and more to win their victories. In the sixth century we shall see that the battles are often fought and won by cavalry only. It is obvious that this revolution in tactic must have reacted on the organisation and carried with it a gradual modification of the legionary system. Another tactical change was the increased importance of archery, brought about by the warfare on the eastern frontier.

Rome did not depend only on her own regular armies to protect her frontiers. She relied also on the aid of the small Federate States which lay beyond her provincial boundaries but within her sphere of influence and under her control. The system of client states goes back to the time of the Republic. The princes of these peoples were bound by a definite treaty of alliance \texthebrew{—} foedus , whence they were called foederati \texthebrew{—} to defend themselves and thereby the Empire against an external foe, and in return they received protection and were dispensed from paying tribute. In the later period with which we are concerned the treaty generally took a new form. The client prince received from the Emperor a fixed yearly sum,\texthebrew{⁠}\footnote{Annonae foederaticae (\textgreek{σιτήσεις}). Perhaps at first it was paid in kind. The subject of the frontier Federates has been clearly and briefly elucidated by Mommsen, \emph{Hist. Sch.} III.225 \emph{sqq.}

} supposed to be the pay of the soldiers whom he was prepared to bring into the field. We shall meet many of these federates, such as the Abasgians and Lazi of the Caucasus, the Saracens on the Euphrates, the Ethiopians on the frontier of Egypt. It was on the basis of a contract of this kind that the Visigoths were settled south of the Danube by Theodosius the Great, and it was by similar contracts that most of the German peoples who were to dismember the western provinces would establish, in the guise of Federates, a footing on Imperial soil.

It may be added that ``federation'' was extended so as to facilitate and regulate the practice of purchasing immunity

p43
from foreign foes, such as the Huns and Persians, a device to which the rulers of the Empire as its strength declined were often obliged to resort. The tribute which was paid for this purpose was designated by the same name ( annonae ) as the subsidies which were allowed to the client princes.

While the Federate system was continued and developed, a new class of troops began to be formed in the fifth century to whom the name Federates was also applied, and who must be carefully distinguished. These troops were drawn indifferently from foreign peoples; they were paid by the government, were commanded by Roman officers, and formed a distinct section of the military establishment. We shall see that, in the course of the sixth century, these mixed Federate troops had come to be the most important and probably the most efficient soldiers in the Imperial army.

The origin of another class of fighting men who were to play a considerable part in the wars of the sixth century goes back to much the same time as that of the Federates. These were the Bucellarians, or private retainers.\texthebrew{⁠}\footnote{Olympiodorus fr. 7. (It was also used as an official term, for in the

\emph{Not. dig.}, \emph{Or.} 7 , we find a second of comites catafractarii bucellarii iuniores.) The bucellarians were largely drawn from Goths, Isaurians, and Galatians. Cp. Mommsen, \emph{op. cit.} 241 \emph{sqq.}; Benjamin, \emph{De Iust. imp. aet.} 18 \emph{sqq.}

} It became the practice of power ­ful generals, and sometimes even civilians, to form an armed retinue or private bodyguard.\texthebrew{⁠}\footnote{We have the cases of Rufinus

(Claudian, \emph{In Ruf.} II.76) ; Stilicho

(Zosimus V.11 ; on the other hand, cp. Claudian,

\emph{In cons. Stil.} III.220 \emph{sqq.} ); º Aetius (Prosper, \emph{sub a.}, 455); Aspar (Malalas, frag. in \emph{Hermes}, VI. p369, where Patzig has shown that the words \textgreek{ο\texthebrew{ὓ}ς \texthebrew{ἐ}κάλεσε φοιδεράτους} are not genuine, see \emph{Unerkannt und unbekannt gebliebene Malalasfragmente}, 13). The bucellarii are recognised as a regular institution in Spain in the laws of Euric (\emph{Leges Visigotorum}, p13). It is generally supposed that this custom was adopted by the Romans from the Germans.

} These soldiers were called bucellarii , from bucella , the military biscuit. Such private armed forces were strictly illegal, but notwithstanding Imperial prohibitions\texthebrew{⁠}\footnote{\emph{C. J.} IX.12.10

(A.D. 468).

} the practice increased, the number of retainers was limited only by the wealth of their master, and officers of subordinate rank had their private armed followers. In the sixth century Belisarius had a retinue of 7000 horse, and these private troops formed a substantial fraction of the fighting strength of the Empire. When they entered the service of their master they took an oath of loyalty to the Emperor.

If the expense of maintaining the army formed a large item

p44
in the annual budget the navy cost little. It would be almost true to say that the Empire at the period had no naval armaments. There were indeed fleets at the old naval stations which Augustus had established at Misenum and Ravenna, and another squadron (classis Venetum) was maintained at Aquileia. But it is significant that the prefects of these fleets, which were probably very small, were under the control of the Master of Soldiers in Italy.\texthebrew{⁠}\footnote{Under him, too, were flotillas on Lake Como, Lake Neuchâtel, and on the Loire and Seine. Those on the Middle Danube, Lake Constance, and in the British channel were under the local military commanders. The Britannic fleet was important in the fourth century, but in the fifth we find instead a classis Sambrica, stationed apparently at Étaples (cp. Lot, \emph{Les Migrations sax.} p5), and under the duke of Belgica Secunda. The care of the government is no longer to protect the coasts of Britain but to protect the other side of the channel. On all these fleets and flotillas see Fiebiger, art. ``Classis'' in \emph{P.\texthebrew{‑}W.} John Lydus (\emph{loc. cit.}) says that in Diocletian's time the number of sailors employed in the fleets both on sea and rivers was 45,000 and that Constantine increased it.

} There was no independent naval command. In the east we find no mention of fleets or naval stations\texthebrew{⁠}\footnote{The classis Carpathia, the classis Alexandrina, and the classis Seleuciae (\emph{C. J.} XI. 2.4

and

13.1 ) were merely fleets of transports, \texthebrew{—} the former two being part of the service for conveying the grain supplies from Egypt to Constantinople.

} with the exception of the small flotillas which patrolled the Lower Danube under the direction of military commanders on that frontier. For centuries the Mediterranean had been a Roman lake, and it was natural that the navy should come to be held as an almost negligible instrument of war. In the third century it had been neglected so far as even to be inadequate to the duty of poli­cing the waters and protecting the coasts against piracy. An amazing episode in the reign of Probus illustrates its inefficiency.\texthebrew{⁠}\footnote{Zosimus I.71.

} A party of Franks, settled on the shores of the Black Sea, seized some vessels, sailed through the Propontis, plundered Carthage, Syracuse, and other cities, and then passing into the Atlantic safely reached the mouths of the Rhine. Yet in the contest between Constantine and Licinius navies played a decisive part, and the two adversaries seemed to have found many useful vessels in the ports of Greece, Syria, Egypt, and Asia Minor. The fleet of Licinius numbered 350 ships and that of Constantine 200, some of which he built for the occasion. It is not clear what the status of these ships was. In the fifth century the Empire was to feel the want of an efficient navy, when the Mediterranean ceased to be an entirely Roman sea and a new German power in Africa contested the

p45
supremacy of its waters. But the failures and defeats which marked the struggle with the vandals did not impress the government of Constantinople with the need of building up a strong navy. The sea forces continued to be regarded as subsidiary, and in overseas expeditions the fleets which convoyed the transports were never placed under an independent naval command. Not until the seventh century, when the Empire had to fight for its very existence with an enemy more formidable than the Vandals, was a naval establishment effectively organised and an independent Ministry of Marine created.

There are three things which it is important to know about the finances of the Empire. The first is, the sources of revenue, and how they were collected; the second is the total amount of the revenue; the third is the total amount of the normal expenditure. As to the first we are fairly well informed; we know a good deal, from first-hand sources, about the system of taxation and the financial machinery. As to the second and third we are in the dark. No official figures as to the annual budget at any period of the later Roman Empire have been preserved, and all attempts to calculate the total of either income or outgoings are guess work, and are based on assumptions which may or may not be true. The utmost that can be done is to fix a minimum.

The financial, like every other department of administration under the autocracy, differed in its leading features from that of the Principate. In raising the revenue the ideal aimed at was equalisation and uniformity; to treat the whole Empire alike, to abolish privileges and immunities. Italy, which had always been free from the burdens borne by the provinces, was largely deprived of this favoured position by the policy of Diocletian.\texthebrew{⁠}\footnote{Aurelius Victor, \emph{Caes.} 39.

} The ideal was not entirely attained; some anomalies and differences survived; but on the whole, uniformity in taxation is the striking characteristic of the new system in contrast with the old. Another capital difference had been gradually brought about. The device of committing the collection of the revenue to middlemen, the publicans, who

p46
realised profits altogether disproportionate to their services, was superseded partly by the direct collection of the taxes by Imperial officials, partly through the agency of the local magistracies of the towns. Moreover, when we survey the sources of revenue at the end of the fourth century, we find that many of the old imposts of the Principate have disappeared, that new taxes have taken their place, and that the modes of assessment have been changed.

The most important and productive source of revenue was the tax on land and agricultural labour. This tax consisted of two distinct parts, the ground tax proper, which represented the old

tributum

imposed on conquered territories, and the annona . The tribute was paid only by those communities and in those districts which had always been liable; it was not extended to those which had been exempted under the Principate. It was paid in coin. The annona which was paid in kind was universal, and was a much heavier burden; no land was exempt; the Imperial estates and the domains of ecclesiastical communities had to pay it as well as the lands of private persons.

Originally the annona\texthebrew{⁠}\footnote{Much light has been thrown on the history of the annona by Seeck in \emph{Die Schatzungsordnung Diocletians}, (see

Bibliography, II.2, C ) and \emph{Gesch. des Untergangs der antiken Welt}, II.

} was an exceptional tax imposed on certain provinces in emergencies, especially to supply Rome with corn in case of a famine, or to feed the army in case of a war. The amount of this extraordinary burden, and its distribution among the communities which were affected by it, were fixed by a special order of the Emperor, known as an indiction. During the civil wars of the third century indictions became frequent. The scarcity of the precious metals and the depreciation of the coinage led to a change in the method of paying the soldiers. They no longer received their wages in coin. Money donations were bestowed on them from time to time, but their regular salary consisted in allowances of food. This practice was systematically organised by Diocletian. The supply of provisions, \texthebrew{—} consisting of corn, oil, wine, salt, pork, mutton \texthebrew{—} necessary to feed a soldier for a year, was calculated, and was called an annona.\texthebrew{⁠}\footnote{This annona was a unit; the officers received, according to their rank, so many annonae. There was also an allowance for horses (capitum). For the distribution of the annona militaris (\textgreek{\texthebrew{ῥ}όγα}) in the sixth cent. cp. \emph{Pap. Cairo}, II.67145.

} In the course of the fourth century

p47
the principle was extended and civil officials received salaries in kind.

This new method of paying the army was the chief consideration which determined the special character of Diocletian's reform in taxation. He made the annona a regular instead of an extraordinary tax, and he imposed, as was perfectly fair, on all parts of the Empire. But he did not fix it at a permanent amount. It was still imposed by an indiction; only an indiction was declared every year. Thus it could be constantly modified and varied, according to the needs of the government or the circumstances of the provinces; and it was intended that it should be revised from time to time by a new land survey.\footnote{Seeck has made it probable that a survey or census of the Empire was made every five years, beginning with A.D. 297 (then 302, 307, 312, etc.). See his article ``Die Entstehung des Indictionencyclus,'' in \emph{Deutsche Zeitschrift f. Geschichtwissenschaft}, XII. In later times a cycle of 15 indictions came to be used officially as a method of chronological reckoning. This cycle is usually counted as starting with A.D. 312, but it comes to the same thing if it is supposed to begin with A.D. 297.

}

The valuation of the land was the basis of the new system. All the territory of the Empire was surveyed, and landed property was taxed not according to its mere acreage but with reference to its value in produ­cing corn or wine or oil. Thus there was a unit ( iugum ) of arable land, and the number of acres in the unit might vary in different places according to the fertility of the soil; there were units for vineyards and for olives; and the tax was calculated on these units.\texthebrew{⁠}\footnote{In Syria there were seven classes of land; the same tax was paid on 5 acres of vineyard as on 20 of the best kind of tilled land and as on 225 of the best kind of olive land. The tax on the seventh class, mountain and pasture, was fixed according to the actual profits. See Bruns and Sachau, \emph{Syrorömisches Rechtsbuch} (1880), pp37, 287. The unit of the iugum was not universal. In Italy there was a larger unit, the millena. In Africa the unit was the centuria =100 acres, and no distinction was made between different classes of land.

} The unit was supposed to represent the portion of land which one able-bodied peasant ( caput ) could cultivate and live on. Thus a property of a hundred iuga meant a property of a hundred labourers or capita , human heads.\footnote{That the iugatio and the capitatio were not two different taxes (as Savigny held and Seeck and others still hold) but the same land tax seems to me to have been proved by Piganiol in his \emph{L'Impôt de capitation.} In most cases the terms could be applied indifferently; but in the case, for instance, mentioned in the text, of a proprietor reserving a part of his estate the term capitation would be inappropriate, as there were no capita (colons).

}

Apart from Imperial estates, the greater part of the soil of the Empire belonged to large proprietors ( possessores ). In country

p48
districts they were generally of the senatorial class; in the neighbourhood of the towns they were probably more often simple curials, members of the local municipal senate. Their lands were parcelled out among tenants who paid a rent to the proprietor and defrayed the land tax. The tenants were known as coloni and, as we shall see later, were practically serfs. Their names and descriptions were entered in the public registers of the land tax, and hence they were called adscriptitii .\texthebrew{⁠}\footnote{That is, censibus adscripti. The Greek is, \textgreek{\texthebrew{ἐ}ναπόγραφοι}. Fragments of a tax roll for the island of Thera have been preserved in which the various denominations of land, the cattle, asses, sheep, slaves, and colons are all enumerated. CIG III.8656 = IG XII. fasc.3, 343\texthebrew{‑}349 (1898).

} As a rule, the proprietor would reserve some part of his estate as a domain for himself, to be cultivated by slaves, and for the tax on the iuga of this domain he would, of course, be directly liable.

Besides the large proprietors there were also small peasants who owned and cultivated their own land, and were distinguished from the serfs on the great estates by the name of plebeians. The tax which they paid was known as the capitatio plebeia . The meaning of this term has been much debated, but there seems little doubt that it is simply the land tax, assessed on the free peasant proprietors on the same principles as it was assessed on large estates.\footnote{This has been shown by Piganiol, \emph{op. cit.} 33 \emph{sqq.} The capitatio humana \texthebrew{—} another term which has caused much discussion \texthebrew{—} was probably (in the fifth century) a tax on slaves, paid by their owners, like the capitatio animalium which is usually associated with it in the laws. \emph{Ib.} 68 \emph{sqq.}

}

The Imperial domains and the private estates of the Emperors, let on leases whether perpetual or temporary, and their cultivators, were liable to the universal annona or capitation, and it was the same with lands held by monastic communities. As to the amount of the land taxes we have hardly any information.\footnote{When Julian went to Gaul, the tribute on each caput was 25 solidi. He reduced it to 7, including all the burdens (on the text of

Ammian. XVI.5.14

cp. Seeck, \emph{Rheinisches Museum}, XLIX.630). In Illyricum it appears that the amount required by provincial governors for their own supplies was at one time a solidus on 120 capita, and was increased, illegally, so that the same sum was paid on 60 capita, and finally on 13. This flagrant case drew a rescript from the Emperor in A.D. 412,

\emph{C. Th.} VII.4.32 .\texthebrew{—} It was the duty of the Praetorian Prefect to send to the provinces lists of the dues for which the taxpayers were liable every year, and on him principally rested the responsibility of deciding whether the ordinary taxation was sufficient to cover the expenses or an addition (superindictio) would be required which could only be imposed with the consent of the Emperor.

}

The ground-tax proper, or tribute, which was a trifle compared with the annona, seems to have been always paid in money,

p49
except in Africa and Egypt, which were the granaries of Rome and Constantinople. It was fixed on the basis of the same survey and was entered in the same book as the annona, but, as we have seen, it was not paid in the privileged territories which had always been exempt. As the currency gradually became established, after Constantine's reforms, the annona too was under certain conditions commuted into a money-payment , and this practice gradually became more frequent.\footnote{Adaeratio was the technical term for the commutation of species into pretia. Its extension in the fifth century can be traced in

\emph{C. Th.} VII.4

(cp. 28 ,

31 ,

32 ,

35 ,

36 ).

}

In the town territories the body of the decurions or magistrates of the town were responsible for the total sum of the taxes to which the estates and farms of the district were liable. The general control of the taxation in each province was entirely in the hands of the provincial governor, but the collection was carried out by officials appointed by the decurions of each town.\texthebrew{⁠}\footnote{The exactor, whose duty was to make known the financial ordinances of the provincial governor and to see that they were executed, in his community; the susceptores = procuratores = \textgreek{\texthebrew{ἐ}πιμεληταί}) who actually received the taxes.

} These collectors handed over their receipts to the compulsor , who represented the provincial governor, and he brought pressure to bear upon those who had not paid.\footnote{Cp. \emph{Nov. Majoriani}, VII.14. The procedure is briefly summed up in

\emph{C. Th.} I.14.1 , omnia tributa \emph{exigere suscipere} postremo \emph{conpellere} iubemus. Egyptian documents afford a good deal of illustration, see Gelzer, \emph{Studien zur byz. Verw. Ägyptens}, 42 \emph{sqq.}

}

Heavy taxes fell upon all classes of the population when a new Emperor came to the throne and on each fifth anniversary of his accession. On these occasions it was the custom to distribute a donation to the army, and a large sum of gold and silver was required.\texthebrew{⁠}\footnote{Each common soldier seems to have received more than £6. Seeck (\emph{Untergang}, II.281) calculates that the quinquennial donation, including presents to senators and others, must have cost the Emperor 3½ millions sterling at least. But before the sixth century the amount per soldier seems to have been reduced to 5 solidi (about 3 guineas);
Procopius, \emph{H. A.} 2A.

} The senators contributed an offertory ( aurum oblaticium ).\texthebrew{⁠}\footnote{The amount presented to Valentinian II in A.D. 385 was 1600 lbs. = about £73,000 (Symmachus, \emph{Rel.} 13).

} The decurions of every town had to scrape together gold which was presented originally in the form of crowns ( aurum coronarium ). Finally a tax was imposed on all profits arising from trade, whether on a large or a petty scale. This burden, which was known as Five-yearly Contribution ( lustralis collatio ) or Chrysargyron (``Gold and Silver'') fell upon prostitutes as well as upon merchants and shopkeepers, and was

p50
felt as particularly oppressive. It is said that parents sometimes sold their children into slavery or devoted their daughters to infamy to enable them to pay it.\footnote{Libanius, \emph{Or.} XLVI.22 (vol. III p389);
Zosimus, II.38 ;

\emph{C. Th.} XIII.1 ; and see below,

chap. XIII. § 3 . It was collected at the end of every four years, and yielded in the case of Edessa, a town of moderate size, about £450 a year.

}

The chief immunity which senators enjoyed was exemption from the urban rates. Besides the aurum oblaticium , and the obligation of the wealthier of their class to fill the office of consul or of praetor, they were liable to a special property tax paid in specie. It was commonly known as the follis \texthebrew{⁠}\footnote{The official name was collatio glebalis; it was also called gleba, and descriptio. See

\emph{C. Th.} VI.2 ;

Zosimus, II.38 ; Hesychius, \emph{fr.} 5; Seeck, \emph{Collatio glebalis} in \emph{P.\texthebrew{‑}W.} The follis was originally a bag of small coins. It was probably sealed at the mint and contained 3125 double denarii = 1 lb. gold, and was used in making large payments. The senatorial tax was known as follis because, as instituted by Constantine, the amount was fixed as so many bags. Popular usage transferred the name from the bag to the coin, and the double denarius itself was known as follis.

} and was scaled in three grades (1 lb., ½ lb., and ¼ lb. of gold according to the size of the property. Very poor senators paid seven solidi\texthebrew{⁠}\footnote{The magister census, who was subordinate to the Prefect of the City, decided (on the basis of the annona registers) at which rate each senator should be liable.

} (£4, 8s. 6d.).

The senators, however, were far from being overtaxed. Most of them were affluent, some of them were very rich, and proportionally to their means they paid far less than any other class. In Italy the income of the richest was sometimes as high as £180,000, in addition to the natural products of their estates which would fetch in the market £60,000. Such revenues were exceptional, but as a rule the senatorial landed proprietors, who had often estates in Africa and Spain as well as in Italy, varied from £60,000 to £40,000.\footnote{Olympiodorus, \emph{fr.} 44, gives these figures. Probus, \emph{c.} A.D. 424, spent £52,500 on his praetor­ship; Symmachus and Maximus £80,000 and £180,000 respectively on the praetor­ships of their sons. Symmachus had estates in Mauretania and in Italy (where he had 15 country houses); the Sallustii had estates in Spain; the domains of the Probi were in all parts of the Empire.\texthebrew{—} The reader may be reminded that the real value, or purchasing power, of gold was far greater than it is to\texthebrew{‑}day. It is generally reckoned that a gold coin in the sixteenth or seventeenth century was as useful as five of the same weight, say, in 1900. It is safe to assume that the proportion, 1:3, is not excessive for a practical comparison, in regard to the purchasing power of money in the nineteenth with the fourth and following centuries. In other words the purchasing power of a solidus approached that of £2 in 1900. This of course does not apply to every commodity, but to labour and commodities all round. Compare the useful remarks of Tenney Frank, \emph{Economic History of Rome} (1920), pp80\texthebrew{‑}83.

}

p51
Besides the yield of all these taxes, which ultimately fell on agricultural labour, the Emperor derived a large revenue from custom duties,\texthebrew{⁠}\footnote{In the early Empire custom duties ( vectigalia ) varied in different places, and were nowhere very high. In the east, at least, they were raised in the fourth century, and an apparently uniform tariff of 12\texthebrew{1/3} per cent (octavae) was imposed

(\emph{C. Th.} IV.61.7 and 8) . As no alteration was made in subsequent laws, this rate probably continued. For the whole volume of trade, we have no figures except Pliny's estimate of imports from the east in the first century A.D. (see below,

p54 ). These imports were undoubtedly the largest item.

} mines, state factories, and extensive Imperial estates. We have no figures for conjecturing the amount of their yield.

The central treasury, which represented the fisc of the early Empire, was presided over by the Count of the Sacred Largess.\texthebrew{⁠}\footnote{Comes sacrarum largitionum, so called because when the office was first instituted the chief duty of the comes was to arrange the largess to the soldiers. The Greek equivalent is \textgreek{κόμης λαργιτιώνων} or \textgreek{τ\texthebrew{ῶ}ν θείων θησαυρ\texthebrew{ῶ}ν.}} All the senatorial taxes, the aurum oblaticium , the collatio lustralis , the custom duties, the yield of the mines and of the public factories, that portion of the land-tax which represented the old tributum, the land-tax which was paid by the colons on the Imperial domains,\texthebrew{⁠}\footnote{Cp.

\emph{C. Th.} V.16.29 .

} all flowed into this treasury. The Count of the Largess administered the mint, the customs, and the mines.

Besides the central treasury, at the Imperial residence in each half of the Empire, there were the chests ( arcae ) of the Praetorian Prefects. These ministers, though they had lost their old military functions, were paymasters of the forces. They were responsible not only for regulating the amount but also for the distribution of the annona. As much of the annona collected in each province as was required for the soldiers stationed there was handed over immediately to the military authorities; the residue was sent to the chest of the Praetorian Prefect.\texthebrew{⁠}\footnote{\emph{C. Th.} VII.4 .

Zosimus, II.33 .

} These chests seem also to have paid the salaries of the provincial governors and their staffs.

The administration of the Imperial domains, which were extensive and were increased from time to time by the confiscation of the property of persons convicted of treason, demanded a separate department and a whole army of officials. At the head of this department was the Count of the Private Estates.\footnote{Comes rerum privatarum, \textgreek{κόμης τ\texthebrew{ῶ}ν πριβάτων.}}

p52
The Private Estate ( res privata ) had originally been organised by Septimius Severus, who determined not to incorporate the large confiscated estates of his defeated rivals in the Patrimony but to have them separately administered.\texthebrew{⁠}\footnote{Cp. Platnauer, \emph{Lucius Septimius Severus}, pp183\texthebrew{‑}184. Stein, \emph{Stud. z. Gesch. d. byz. Reiches}, p169.

} In the fourth century the Patrimony and the Private Estate were combined and placed under a minister of illustrious rank. His officials administered the domains and collected the rent from the colons. The greater part of the Imperial lands were treated as State property of which the income was used for public purposes. But certain domains were set aside to furnish the Emperor's privy purse. Thus the domains in Cappadocia were withdrawn from the control of the Count of Private Estates and placed under the control of the Grand Chamberlain.\texthebrew{⁠}\footnote{Stein (\emph{op. cit.} 171) is certainly right in pointing out that this transference meant the appropriation of the Cappadocian domains to the privy purse. In A.D. 379 these domains were under the comes r. p.

(\emph{C. Th.} VI.30.2) . I should conjecture that the change was made in the first years of Arcadius while the power­ful Eutropius was Chamberlain. The section on praepositus s. cub. in the west in the

\emph{Not. dig.}, \emph{Occ.}

is lost, but there is no evidence that a corresponding change was ever made in the west, or that the Imperial domains in Africa were ever under the praepositus. Stein's view that the change was common to both parts of the Empire, and that in the west the domus divina in Africa was restored to the comes r. p. before A.D. 409, seems to me to be unnecessary.

} And in the same way, in the west, certain estates in Africa ( fundi domus divinae per Africam ) were appropriated to the personal disposition of the Emperor, although they remained under the control of the Count.

What were the relations between the fisc or treasury of the Count of the Sacred Largess on one hand, and the chests of the Praetorian Prefects and the treasury of the Count of the Private Estates on the other? We may conjecture that the Prefects paid out of the treasuries directly the salaries of all the officials, both central and provincial, who were under their control; that in the same way the Count of the Private Estates paid out of the monies that came in from the domains all the officials who were employed in their administration; and that all that remained over, after the expenses of the departments had been defrayed, was handed over to the treasury of the Count of the Sacred Largess.\texthebrew{⁠}\footnote{Or, if not handed over, that the accounts were submitted to him so that he knew the surplus on which he could draw.

} This was the public treasury which had to supply the money required for all purposes with the four exceptions of the Emperor's privy purse, the upkeep of the administration

p53
of the Imperial domains, the maintenance of the civil service under the Praetorian Prefects, and the payment of the army.

It has already been observed that no figures are recorded either for the annual revenue or for the annual expenditure. We have no data to enable us to conjecture, however roughly, the yield of the mines or of the rents of the Imperial domains. There is some material for forming a minimum estimate of the money value of the land-tax in Egypt, but even here there is much uncertainty.\texthebrew{⁠}\footnote{For instance the figures as to the corn sent to Constantinople in Justinian, \emph{Edict} 13; and to Rome, in the time of Augustus, in Victor, \emph{Epitome}, c. 1 as to the amount of corn and of money taxes paid by Antaeopolis in the sixth century, in \emph{Pap. Cairo}, I. No. 67057; and other data furnished by papyri. A figure which has been over­looked is the incidental statement in a later document, but which may come from a sixth century source, that the annual money taxes of Egypt amounted to 36,500 pounds of gold = 2,508,000 solidi; see \emph{\textgreek{Διήγησις περ\texthebrew{ὶ} τ\texthebrew{ῆ}ς \texthebrew{ἁ}γ. Σοφίας}}, § 25 (cp. below,

Vol. II Chap. XV. § 6 ).

} Turning to expenditure, we find that the evidence points to 500,000 or thereabouts as the lowest figure we can assume for the strength of the army in the time of Theodosius the Great. The soldiers were paid from the annona. When this payment in kind was commuted into coin, it was valued at 25 or 30 solidi a year for each soldier.\texthebrew{⁠}\footnote{25 solidi

(\emph{C. Th.} VII.7.13) , 30 solidi (\emph{Nov. Valentin.} VI.3) were the sums which, in A.D. 397 and 443 respectively, persons liable to the furnishing of recruits might pay instead. In \emph{Pap. Brit. Mus.} III.985, we have a soldier's receipt for his pay, 30 solidi.

} The annual value of the annona must then have exceeded 12½ million solidi or nearly 8 million sterling. Of the salaries paid to the civil and military officials and their staffs we can only say that the total must have exceeded, and may have far exceeded, £400,000.\footnote{Based on various figures given in laws of Justinian (sixth century, but rates of pay were probably much the same). Cp. above,

p33, n1 . We have no material for conjecturing the cost of the numerous officials subordinate to the mag. off. and the financial ministers; and the chamberlains and staff of the Palaces are left entirely out of account. Bouchard (\emph{Étude sur l'admin. des finances de l'empire romain}, p49) calculated that the \emph{civil} service cost less than £250,000. Sundwall (\emph{Weström. Studien}, p156) has much higher figures which seem precarious. He thinks the cost of paying the civil officials in Gaul and Italy amounted to £2,000,000. He calculates the revenue from land-taxes under Honorius as about £13,200,000 (p155).

}

From the general consideration that the population of the Empire at the lowest estimate must have been 50 millions, we might assume as the minimum figure for the revenue 50 million solidi, on the ground that in a state which was severely taxed the taxation could not have been less than 1 solidus per head.\footnote{To illustrate this, in 1760 the population of England and Wales was over 6½ millions, and the revenue

(p54)
from taxes amounted in 1762 to £6,711,000. This was about £1 a head, and the country, which was still mainly agricultural, was not overburdened. The taxation would necessarily have been much higher but for the happy expedient of the Public Debt.

}

p54
That would be about £31,250,000. It is probably much under the mark.

Of the financial problems with which Diocletian and Constantine had to deal, one of the most difficult was the medium of exchange. In the third century the Empire suffered from scarcity of gold. The yield of the mines had decreased; and a considerable quantity of the precious metals was withdrawn from circulation by private people, who during that troubled period buried their treasures. But the chief cause of the scarcity was the drain of gold to the east in exchange for the Oriental wares which the Romans required. In the first century A.D. the annual export of gold to the east is said to have amounted (at the least) to a million pounds sterling.\texthebrew{⁠}\footnote{100,000,000 sesterces

(Pliny, \emph{N. H.} XII.18, § 84) , of which 55,000,000 went to India. The Emperor Gratian, about A.D. 374, legislated against the export of gold,

\emph{C. J.} IV.63.2 .

} The Emperors resorted to a depreciation of the coinage, and up to a certain point this perhaps was not particularly disadvantageous so far as internal trade was concerned, since the value of the metals had risen in consequence of the scarcity. When Diocletian came to the throne there was practically nothing in circulation but the double denarius, which ought to have been a silver coin equivalent to about 1s. 9d.), but was now made of copper, with only enough silver in it to give it a whitish appearance, and worth about a halfpenny. Both Aurelian and Diocletian made attempts to establish a stable monetary system, but the solution of the problem was reserved for Constantine. The Constantinian gold solidus or nomisma remained the standard gold coin and maintained its proper weight, with little variation, till the eleventh century. Seventy-two solidi went to the pound of gold, so that its value was about twelve shillings and sixpence.\texthebrew{⁠}\footnote{The legend CONOB, which appears on solidi minted at Constantinople (till the reign of Leo III) is an abbreviation of the name of the mint and of the word obryzum, refined gold.

} But the solidus was not treated as a coin in the proper sense; and it was not received as interchangeable into so many silver or copper pieces. The pound of gold was really the standard, and, when solidi were used in ordinary transactions, they were weighed. In the payment of taxes they were accepted at their nominal value, but for other

p55
purposes they were pieces of metal, of which the purity, not the weight, was guaranteed by the mint.\footnote{The siliqua was a silver coin = 1\texthebrew{⁄}24th of the solidus; but the silver coin most in use was the half-siliqua known as the nummus decargyrus. The silver miliarense (= 1\texthebrew{⁄}1000 lb. gold) was, according to Babelon, in the fifth and sixth centuries a \emph{monnaie de luxe} (cp. Justinian,

\emph{Nov.} 105.2 ); 12 (not 14) went to the solidus. The ratio between gold and silver in A.D. 397 is given in

\emph{C. Th.} XIII.2.1

as 1 lb. silver = 5 solidi = 5\texthebrew{⁄}72 lb. gold, and in A.D. 422, 1 lb. silver = 4 solidi = 1\texthebrew{⁄}18 lb. gold

(\emph{ib.} VIII.4.27) . Thus in these 25 years the ratio changed from 1:14\texthebrew{2/5} to 1:18, a considerable depreciation of silver. On the silver and copper coins of the fourth and fifth centuries see Babelon, \emph{Traité des monnaies grecques et romaines}, vol. I (1901) 566 \emph{sqq.}, and 612 \emph{sqq.} \texthebrew{—} It may be noted here that the ordinary rate of interest in the fourth and early fifth century was from 4 to 6 per cent. 12 per cent (the centesima) was the maximum allowed by law, but it would be an error to infer from the fulminations of Ambrose and Chrysostom against it that it was normal or typical in business transactions. It was only exacted in cases where there was no good security. See Billeter, \emph{Gesch. des Zinsfusses}, 236 \emph{sqq.} It was possibly due to clerical influence that senators were forbidden towards the end of the fourth century to lend on interest. The law was, of course, evaded and (after the fall of Chrysostom) they were allowed to receive interest up to 6 p. c. (See

\emph{C. Th.} II.33.3 and 4 , with the commentary of Gothofredus, vol. I p274\texthebrew{‑}275).

}

Diocletian and Constantine had to seek solutions not only of political but also of more difficult economic problems. The troubles of the third century, the wars both domestic and foreign, the general disorder of the State, had destroyed the prosperity of the Empire and had rapidly developed sinister tendencies, which were inherent in ancient civilisation, and legislators whose chief preoccupation was the needs of the public treasury applied methods which in some ways did more to aggravate than to mitigate the evils. We find the State threatened with the danger that many laborious but necessary occupations would be entirely abandoned, and the fields left untilled for lack of labourers. The only means which the Emperors discovered for averting such consequences was compulsion. They applied compulsion to the tillers of the soil, they applied compulsion to certain trades and professions, and they applied it to municipal service. The results were serfdom and hereditary status. The local autonomy of the municipal communities,\texthebrew{⁠}\footnote{J. S. Reid, \emph{Municipalities of the Roman Empire} (1913). The early Roman Empire may be regarded ``as an organisation based upon a federation of municipalities forming an aggregate of civic communities enjoying a greater or less measure of autonomy, and having certain characteristics derived from an age when state and city were convertible terms'' (p3).

} the cities and towns

p56
which were the true units in the structure of the Empire, had been undermined in some ways under the Principate, but before Diocletian no attempt had been made to impose uniformity, and each community lived according to its own rules and traditions. The policy of uniform taxation, which Diocletian introduced, led to the strict control of the local bodies by the Imperial Government. The senates and the magistrates became the agents of the fisc; the municipalities lost their liberties and gradually decayed.

(1) For some centuries there had been a general tendency to substitute free for servile labour on large estates. The estate was divided into farms which were leased to free tenants, coloni , on various conditions, and this system of cultivation was found more remunerative.\texthebrew{⁠}\footnote{For the origins and history of the colonatus, see M. Rostowzew, \emph{Studien sur Geschichte des römischen Kolonates} (1910).

} But towards the end of the third century the general conditions of the Empire seem to have brought about an agrarian crisis. Many colons found themselves insolvent. They could not pay the rent and defray the heavy taxes. They gave up their farms and sought other means of livelihood. Proprietors themselves some sold their lands, and the tenants declined to hold their farms under the new owners. Thus land fell out of cultivation and the fiscal revenue suffered. Constantine's legislation, to solve this agrarian problem, created a new caste. He made the colons compulsory tenants. They were attached to the soil, and their children after them. They continued to belong legally to the free, not to the servile, class; they had many of the rights of freemen, such as that of acquiring property. But virtually they were unfree and were regarded as chattels. Severe laws prevented them from leaving their farms, and treated those who ran away as fugitive slaves. The conception of a colon as the chattel of his lord comes out clearly in a law which describes his flight as an act of theft; ``he steals his own person.''\texthebrew{⁠}\footnote{\emph{C. J.} XI.48.3 , sese . . . furari intelligatur. The oppression of the colons is graphically described by John Chrysostom, \emph{Homilia in Matth.} 61, 31 (\emph{P.G.} 58, 591).

} But the Emperors, whose principal aim in their agrarian legislation was to guard the interests of the revenue, protected the colons against exorbitant demands of rent on the part of the proprietors. And if a proprietor sold any part of his estate, he was not allowed to retain the tenants.\footnote{\emph{C. Th.} XIII.10.3 .

}

p57
At the same time the condition of rustic slaves was improved. The government interfered here too, for the same reason, and forbade masters to sell slaves employed on the land except along with the land on which they worked.\texthebrew{⁠}\footnote{\emph{C. J.} XI.48.7 .

} This limitation of the masters' rights tended to raise the condition of the slave to that of the colon.

The proprietor's power over his tenants was augmented by the fact that the State entrusted him with the duties of collecting the taxes for which each farm was liable,\texthebrew{⁠}\footnote{\emph{C. Th.} XI.1.14 .

} and of carrying out the conscription of the soldiers whom his estate was called upon to furnish. He also administered justice in petty matters and policed his domains. Thus the large proprietors formed an influential landed aristocracy, with some of the powers which the feudal lords of western Europe exercised in later times. They were a convenient auxiliary to the Government, but they were also a danger. The custom grew up for poor freemen to place themselves under the protection of wealthy landowners, who did not scruple to use their influence to divert the course of justice in favour of these clients, and were able by threats or bribery to corrupt the Government officials. Such patronage was forbidden by Imperial laws, but it was difficult to abolish it.\footnote{The evils of patronage (\textgreek{προστασία}) are portrayed in the oration of Libanius \textgreek{Περ\texthebrew{ὶ} τ\texthebrew{ῶ}ν προστασι\texthebrew{ῶ}ν} addressed to Theodosius I in A.D. 391 or 392 (\emph{Or.} XLVII. ed. Förster). Cp. F. de Zulueta, \emph{De patrociniis vicorum} (\emph{Oxford Studies in Legal and Social History}, ed. by Vinogradoff, 1909).

}

It had long been the custom for public bodies to grant the land which they owned on a perpetual lease, subject to the payment of a ground-rent ( vectigal ). It was on this principle that Rome had dealt with conquered territory. The former proprietors continued to possess their land, but subject to the owner ­ship ( dominium ) of the Roman people and liable to a ground-rent. In the fifth century this form of land tenure coalesced with another form of perpetual lease,

emphyteusis , which had its roots not in Roman but in Greek history. Emphyteusis meant the cultivation of waste land by planting it with olives or vines or palms.\texthebrew{⁠}\footnote{Cp. Rostowzew, \emph{op. cit.} 105, 267.

} To encourage such cultivation a special kind of tenure had come into use. The emphyteutes bound himself by contract to make certain improvements on the land; he paid a small fixed rent; his tenure was perpetual

p58
and passed to his heirs, lapsing only if he failed to fulfil his contract. In the course of time, all kinds of land, not only plantation land, might be held by emphyteutic tenure. Legally this agreement did not answer fully to the Roman conception either of a lease or of a sale, and lawyers differed as to its nature. It was finally ruled that it was neither a sale nor a lease, but a contract \emph{sui generis.}\texthebrew{⁠}\footnote{By Zeno,

\emph{C. J.} IV.66.1 . See Justinian, \emph{Instit.} 3, 24. This law provided that if part of an emphyteutic property became unproductive, the loss fell on the tenant; but if the whole, the owner was responsible.

} This kind of tenancy was the rule on the Imperial domains. But it was also to be found on the estates of private persons.

(2) The trades to which the method of compulsion was first and most harshly applied were those on which the sustenance of the capital cities, Rome and Constantinople, depended : the skippers who conveyed the corn supplies from Africa and Egypt, and the bakers who made it into bread. These trades, like many others, had been organised in corporations or guilds ( collegia ), and as a general rule the son probably followed the father in his calling. It was the most profitable thing he could do, if his father's capital was invested in the ships or in the bakery.\texthebrew{⁠}\footnote{Seeck, \emph{Untergang}, II.311.

} But this changed when Diocletian required the skippers to transport the public food supplies, and made their property responsible for the safe arrival of the cargoes. They had to transport not only the supplies for the population of the capital, but the annonae for the soldiers. This was a burden which tempted the sons of a skipper to seek some other means of livelihood. Compulsion was therefore introduced, and the sons were bound to their father's calling.\texthebrew{⁠}\footnote{\emph{C. Th.} XIII.5 . For the regulations about the navicularii see E. Gebhardt, \emph{Das Verpflegungswesen von R. und C.} Their services in transporting corn were remunerated by 4 per cent of the cargo

(\emph{C. Th.} XIII.5.7) .

} The same principle was applied to the bakers, and other purveyors of food, on whom the State laid public burdens. In the course of the fourth century the members of all the trade guilds were bound to their occupations. It may be noticed that the workmen in the public factories ( fabricae ) were branded, so that if they fled from their labours they could be recognised and arrested.

(3) The decline of municipal life, and the decay of the well-to\texthebrew{‑}do provincial citizen of the middle class, is one of the important social facts of the fourth and fifth centuries. The

p59
beginnings of this process were due to general economic conditions, but it was aggravated and hastened by Imperial legislation, and but for the policy of the Government might perhaps have been arrested.

The well-to\texthebrew{‑}do members of a town community, whose means made them eligible for member ­ship of the curia or local senate and for magistracy, formed the class of curiales .\texthebrew{⁠}\footnote{For the history and organisation of the curial bodies, see Kübler's article \emph{Decurio} in \emph{P.\texthebrew{‑}W.}

} The members of the senate were called decuriones . But in the period of decline these terms were almost synonymous. As the numbers of the curials declined, there was not one of them who was not obliged at some time or other to discharge the unwelcome functions of a decurion. In former times it had been a coveted honour to fulfil the unpaid duties of local administration, but the legislation of the Emperors, from the end of the third century onward, rendered these duties an almost intolerable burden. The curials had now not only to perform their proper work of local government, the collection of the rates, and all the ordinary services which urban councils everywhere discharge. They had also to do the work of Imperial officials. They had to collect the land-taxes of the urban district. And they were made responsible for the full amount of taxation, so that if there were defaulters, they were collectively liable for the deficiency.\texthebrew{⁠}\footnote{This seems to have been the rule, though the Emperors sometimes legislated otherwise; cp. \emph{C. J.} XI.59.16 ,

\emph{C. Th.} VII.22.1 . The decuriones themselves seem, so far as they could, to have made those whom they appointed to collect the taxes, liable for deficiencies. The results were not only cruel to the individual, but calamitous for the community. One of the forms of patronage, described by Libanius (\emph{op. cit.}) illustrates the difficulties of the tax-collector. Villages in the district of an urban community would place themselves under the protection of soldiers quartered in the district, who, in return for gifts in kind or corn, would help them to defy the tax-gatherer and drive him out of the village. The unfortunate man might have to sell his property to make up the sum which he was required to produce. And thus the number of curials was reduced. \textgreek{Βουλευτ\texthebrew{ὴ}ς βουλ\texthebrew{ῆ}ς \texthebrew{ἐ}ξαλείφεται . . . τα\texthebrew{ῦ}τ\texthebrew{’ ἐ}λάττους ποιε\texthebrew{ῖ} τ\texthebrew{ὰ}ς βουλ\texthebrew{ὰ}ς \texthebrew{ἀ}ντ\texthebrew{ὶ} μειζόνων} (\emph{ib.} 10). See also Libanius, \emph{Or.} II.33\texthebrew{‑}36 for the decline of the senates.

} They had also to arrange for the supply of horses and mules for the Imperial post, the upkeep of which, though its use was exclusively confined to Government officials, was laid upon the provincials and was a most burdensome \emph{corvée.}

The burdens laid upon the curials became heavier as their numbers diminished. Diocletian's reorganisation of the State

p60
service, with innumerable officials, invited the sons of well-to\texthebrew{‑}do provincial families, who in old times would have been content with the prospect of local honours, to embrace an official career by which they might attain senatorial rank; and senatorial rank would deliver them from all curial obligations.

In course of time the plight of the middle-class provincials, who were generally owners of small farms in the neighbourhood of their town and suffered under the heavy taxation, became so undesirable that many of them left their homes, enlisted in the army, took orders in the Church, or even placed themselves under the patronage of rich proprietors in the country. The danger was imminent that the municipal organisation would entirely dissolve. Here again the Emperors resorted to compulsion. The condition of the curial was made a hereditary servitude.\texthebrew{⁠}\footnote{The principle is laid down in

\emph{C. Th.} XII.1.22

(A.D. 336). This long Title, \emph{de decurionibus}, is a monument of merciless despotism. The decay of the curials is very fully treated by Dill, \emph{Roman Society}, Book III. chap. II.

} He was forbidden to leave his birthplace; if he wanted to travel, he had to obtain leave from the provincial governor. His sons were bound to be curials like himself; from their birth they were, in the expressive words of an Imperial law, like victims bound with fillets.\texthebrew{⁠}\footnote{\emph{Ib.} 122 veluti dicati infulis mysterium perenne custodiant. Men born in the curial class, who entered the army or the civil service, were sternly ``restored'' to their municipal duties, \emph{ib.} 137, 139, 146.

} He could only escape from his lot by forfeiting the whole or a part of his property. Restrictions were placed on his ordinary rights, as a Roman citizen, of selling his land or leaving it by will at his own discretion. Nothing shows the unenviable condition of the curial class more vividly than the practice of pressing a man into the curia as a punishment for misdemeanours.\footnote{The practice is forbidden \emph{ib.} 66 and 108.

}

The power of the local magistrates had been diminished in the second century by Trajan's institution of the curator civitatis , whose business was to superintend the finances of the municipality. The curator was indeed a townsman, but as a State servant he had ceased to belong to the curial order and he was appointed by the provincial governor. By the middle of the fourth century his prestige had declined because the right of appointing him had been transferred to the curia itself. He was overshadowed by the new office of defensor instituted by Valentinian I to protect the interests of the poorer classes against

p61
the oppression of the power ­ful.\texthebrew{⁠}\footnote{\emph{C. Th.} I.29 . See Seeck's art., \emph{Defensor civitatis}, in \emph{P.\texthebrew{‑}W.} Constantius had instituted (A.D. 361) defensores senatus in the provinces to protect members of the senatorial order against official oppression.

\emph{C. Th.} I.28 .

} The defensor was to be appointed by the Praetorian Prefect, and he was to be a man who filled some not unimportant post in the State service. But the institution did not prove a success. It was difficult to get the right sort of people to undertake the office, and it was soon bestowed for corrupt reasons on unsuitable persons. Theodosius the Great sought to remedy this by transferring the appointment of the defensor to the curials.\texthebrew{⁠}\footnote{\emph{C. Th.} I.29.6 .

} The prestige of the office at once declined, and the defensor ­ship like the curator ­ship became one more burden imposed upon the sorely afflicted curial class, without any real power to compensate for the duties which it involved. The influence of all the urban magistracies, which had become anything rather than an honour, was soon to be overshadowed by that of the bishop. And this reminds us of another feature in the decline of municipal life which deserves to be noticed.

That much-abused expression ``age of transition'' has a real meaning when some fundamental change forces a society to adapt itself slowly and painfully to new conditions. The period of the industrial transformation, brought about by the invention of machinery, in modern states is an example of a true age of transition. The expansion and triumph of Christianity in the third and fourth centuries rendered that period a genuine age of transition in the same sense, and the transition was marked by distress and destruction. Roman and Greek municipal life was inextricably bound up with pagan institutions \texthebrew{—} temples, cults, games. The interests and habits of the town communities were associated with these institutions, and when Christianity suppressed them, municipal life was deprived of a vital element. For the Church did not succeed in bringing her own institutions and practices into the same intimate connexion with municipal organisation.\texthebrew{⁠}\footnote{Cp. the excellent remarks of Vinogradoff, in \emph{C. Med. H.} I.554\texthebrew{‑}555.

} With the passing of paganism something went out of the vitality of ancient town life which could never be restored.

(4) The principle of compulsion was extended to military service. The sons of veterans were obliged to follow the

p62
profession of their fathers, with the uninviting alternative of being enrolled in the class of decurions. They were definitely debarred from a career in the civil service. The sons of civil servants too were expected to follow the career of their fathers.\footnote{See

\emph{C. Th.} VII.22.3 .

}

We might better understand the economic conditions which the Emperors sought to regulate by tyrannical legislation if we possessed some trustworthy statistics of the population of the Empire and its various provinces. In the eighteenth century, even after Hume had exploded the old delusion that the ancient states in Europe were far more populous than the modern, Gibbon estimated the population of the Empire in the time of Claudius as 120,000,000. It is now generally agreed that this figure is far too high. Any estimate rests on a series of conjectures, but perhaps half this figure would be nearer the truth. According to a recent calculation, which is probably below rather than over the mark, the population at the death of Augustus amounted to 54,000,000, of which 26,000,000 are assigned to the western provinces including the Danubian lands, and 28,000,000 to the Greek and Oriental provinces.\texthebrew{⁠}\footnote{Beloch, \emph{Die Bevölkerung der griechisch-römischen Welt} (cp. the Table, p507). His numbers for the Danubian lands are 2,000,000; for Greece, 3,000,000; for Spain, 6,000,000; for Narbonensis, 1,500,000; for the other Gallic provinces, 3,400,000. E. A. Foord has attempted to prove that in A.D. 395 the population was 120,000,000 (\emph{Byzantine Empire}, p10).

} By the beginning of the fourth century there seems some reason to suppose that the population had increased. This would be the natural result of the development of city life in Spain and Gaul, and the gradual civilisation of the Illyrian and Danubian provinces. On this basis of calculation, which, it must be repeated, involves many possibilities of error, we might conclude that in the time of Constantine the population of the Empire may have approached 70,000,000.

We have indeed some definite evidence that in the fourth century the government was not alarmed by the symptoms of a decline in numbers which had confronted the Emperor Augustus. It may be remembered that among the measures which Augustus adopted to arrest the fall in the birth-rate of Roman citizens he penalised bachelors by rendering them incapable of inheriting, and married people who were childless by allowing them to take only half of an inheritance which if they had children would

p63
fall to them entirely. It is significant that Constantine removed this disability from bachelors,\texthebrew{⁠}\footnote{\emph{C. Th.} VIII.16.1

(A.D. 320). In A.D. 410 Theodosius II abrogated the law of Augustus with regard to the childless; this applied only to the eastern half of the Empire.

\emph{Ib.} VIII.17.2 .

} while Theodosius II abrogated the law of Augustus with regard to the childless. This repeal of a law which had been so long in force may fairly be taken as an indication that in the fourth century no fears of a decline in population troubled the Imperial Government.

While in all ancient monarchies religion and sacerdotalism were a political as well as a social power, the position of the Christian Church in the Roman Empire was a new thing in the world, presenting problems of a kind with which no ruler had hitherto been confronted and to which no past experience offered a key. The history of the Empire would have been profoundly different if the Church had remained as independent of the State as it had been before Constantine, and if that Emperor and his successors had been content to throw the moral weight of their own example into the scale of Christianity and to grant the Church the same freedom and privileges which were enjoyed by pagan cults and priesthoods. But heresies and schisms and religious intolerance on one side, and the despotic instinct to control all social forces on the other, brought about a close union between State and Church which altered the character and spirit of the State and constituted perhaps the most striking difference between the early and the later Empire. The disorders caused by violent divisions in the Church on questions of doctrine called for the intervention of the public authorities, and rival sects were only too eager to secure the aid of the government to suppress their opponents. Hence at the very beginning Constantine was able to establish the principle that it devolved upon the Emperor not indeed to settle questions of doctrine at his own discretion, but to summon general ecclesiastical Councils for that purpose and to preside at them. The Council of Arles ( A.D. 314) was convoked by Constantine, and the Ecumenical Council of Nicaea exhibited the full claim of the Emperor to be head of the Church. But in this capacity he stood outside the ecclesiastical hierarchy; he

p64
assumed no title or office corresponding to that of Pontifex Maximus. Historical circumstances decided that this league of Church and State should develop on very different lines in the east and in the west. In the west it was to result in the independence and ultimately in the supremacy of the Church; in the east the Church was kept in subordination to the head of the State, and finally ecclesiastical affairs seem little more than a department of the Imperial Government. Even in the fourth century the bishop of Rome has a more independent position than the bishop of Constantinople.

At the beginning of our period the general lines of ecclesiastical organisation had been completed. The clergy were graded in a hierarchical scale of seven orders \texthebrew{—} bishops, priests, deacons, subdeacons, acolytes, exorcists, and readers. In general, the ecclesiastical divisions closely correspond to the civil.\texthebrew{⁠}\footnote{In the east this seems to have strictly prevailed.

} Every city has its bishop. Every province has its metropolitan, who is the bishop of the metropolis of the province. And above the provincial metropolitans is the exarch, whose jurisdiction corresponds to the civil diocese. A synod of bishops is held annually in each province.

But among the more important sees, four stood out pre-eminent \texthebrew{—} Rome, Constantinople, Alexandria, and Antioch. Of these Rome was acknowledged to be the first, but there was rivalry for the second place. Besides these the See of Jerusalem had, by virtue of its association with the birth of Christianity, a claim to special recognition. By the middle of the fifth century the positions of these great sees were defined, and their jurisdiction fixed. Their bishops were distinguished as Patriarchs,\texthebrew{⁠}\footnote{In early times the name Patriarch was sometimes given to simple bishops; cp. \emph{J.H.S.} VI.346 (archbishop of Hierapolis in Phrygia). See also Cassiodorus, \emph{Var.} IX.15; and cp. the ``Patriarchate'' of Aquileia. Duchesne, \emph{Églises séparées}, 262.

} though the bishop of Rome did not assume this title. The ecclesiastical map shows five great jurisdictions or Patriarchates. The authority of Rome extended over the whole western or Latin half of the Empire, and included the Praetorian Prefecture of Illyricum.\texthebrew{⁠}\footnote{The bishop of Thessalonica acted as Vicar of the Pope in Illyricum. The Patriarchs of Constantinople sometimes contested the Papal rights in this prefecture; \emph{e.g.} Atticus, who doubtless prompted the law of Theodosius II, in

\emph{C. Th.} XVI.2.45

(A.D. 421), claiming the jurisdiction for the Patriarch. On the whole subject see Duchesne, \emph{op. cit.} 229 \emph{sqq.}

} The Patriarchate of Constantinople ultimately

p65
embraced the civil dioceses of Thrace, Pontus, and Asia.\texthebrew{⁠}\footnote{This was settled at the Council of Chalcedon, A.D. 451.

} The Patriarchate of Alexandria, third in precedence, corresponded to the Diocese of Egypt. The Patriarchate of Antioch comprised the greater part of the Diocese of the East; the small Patriarchate of Jerusalem the three Palestinian provinces. The autocephalous Church of Cyprus stood apart and independent.\footnote{Its independence of Antioch was decreed by the Council of Ephesus A.D. 431.

}

The development of a graded hierarchy among the bishops revolutionised the character of the Church. For three centuries the Christian organisation had been democratic. Its union with the monarchical state changed that. The centralised hierarchical system enabled the Emperors to control it in a way which would have been impossible if the old democratic forms had continued.

Constantine and his successors knew how to attach to themselves the power ­ful organisation of which they had undertaken the direction. Valuable privileges were conceded to the clergy and the churches. Above all, the clergy, like the pagan priests, were exempted from taxation,\texthebrew{⁠}\footnote{\emph{C. Th.} XVI.2.1.2 ;

XI.1.1 .

} a privilege which attracted many to their ranks. The churches had an unrestricted right of receiving bequests, and they inherited from the pagan temples the privilege of affording asylum.\texthebrew{⁠}\footnote{\emph{Ib.} XVI.2.4 ;

\emph{C. J.} I.12.2 .

} The bishops received the right of acting as judges in civil cases which the parties concerned agreed to bring before them, and their decisions were without appeal.\texthebrew{⁠}\footnote{\emph{C. Th.} I.27 , episcopalis audientia.} It was the Imperial policy to make use of the ecclesiastical authorities in local administration, and as the old life of the urban communities declined the influence of the bishops increased. The bishop shared with the defensor civitatis the duty of protecting the poor against the oppression of the power ­ful and the exactions of government officials, and he could bring cases of wrongdoing to the ears of the Emperor himself. Ultimately he was to become the most influential person in urban administration.

The first century of Christianity in its new rôle as a state religion was marked by the development of ecclesiastical law. The canons of the Council of Nicaea formed a nucleus which was enlarged at subsequent councils. The first attempt to codify canon law was made at the beginning of the fifth century.

p66
The legislation of councils was of course only binding on the Church as such, but as time went on it became more and more the habit of the Emperors to embody ecclesiastical canons in Imperial constitutions and thus make them part of the law of the state. It is, however, to be noticed that canon law exerted little or no effect upon the Roman civil law before the seventh century.

\xchapter{Chapter III}

Short URL for this page:

bit.ly/BURLAT3

The history of a thousand years approved the wisdom of Constantine in choosing Byzantium for his new capital. A situation was needed from which the Emperor could exercise imminent authority over south-eastern Europe and Asia, and could easily reach both the Danube and the Euphrates. The water passage where Asia and Europe confront each other was one of the obvious regions to be considered in seeking such a central site. Its unique commercial advantages might have been alone sufficient to decide in its favour. It was the natural meeting-place of roads of trade from the Euxine, the Aegean, and northern Europe. When he determined to found his city by this double-gated barrier between seas and continents, there were a few sites between which his choice might waver. But there was none which in strategical strength could compare with the promontory of Byzantium at the entrance of the Bosphorus. It had indeed some disadvantages. The prevailing winds are north-easterly, and the arrival of sea-borne merchandise was often seriously embarrassed, a fact which the enemies of Constantine did not fail to insist on.\texthebrew{⁠}\footnote{Eunapius, \emph{Vit. Aedes.} p23.

} The frequency of earthquakes\texthebrew{⁠}\footnote{Thirteen are recorded between 395 and 565. The most serious were those of 447, 480, and 558.

} was another feature which might be set against the wonder ­ful advantages of Byzantium as a place for a capital of the Empire.

While the whole trend of the passage through which the waters of the Euxine reach the Aegean is from east to west, the channel of the Bosphorus runs from north to south.\texthebrew{⁠}\footnote{Dethier (\emph{Der Bosphor und Cpel.} p65) gives the length of the Bosphorus as exactly 27 kils. and the narrowest breadth between Rumili and Anatoli Hissar as 550 metres.

} At

p68
the point where it widens into the Propontis, the European shore is broken by a deep narrow inlet which penetrates for \texthebrew{•} more than six miles and forms the northern boundary of a hilly promontory, on which Byzantium was built. This inlet or harbour was known as the Golden Horn, and it is the feature which made the fortune of Constantine's city.

The shape of Constantinople is a trapezium, but the eastern side is so short that the city may be described as a triangle with a blunted apex. On three sides, north, east, and south, it is washed by water. The area of the city ``is \texthebrew{•} about four miles long and from one to four miles wide, with a surface broken up into hills and plains. The higher ground, which reaches an elevation of \texthebrew{•} some 250 feet, is massed in two divisions \texthebrew{—} a large isolated hill at the south-western corner of the promontory, and a long ridge, divided, more or less completely, by five cross valleys into six distinct eminences, overhanging the Golden Horn.'' These two masses of hill ``are separated by a broad meadow through which the stream of the Lycus flows athwart the promontory into the Sea of Marmora.''\footnote{Van Millingen, \emph{Byzantine Cple.} p2.

}

Constantine found the town\texthebrew{⁠}\footnote{Besides the miscellaneous notices in histories and chronicles, the chief sources for the topography of the city are: (1) \emph{Notitia urbis Constantino­politanae} , an inventory of the principal buildings and monuments in each of the fourteen regions. The author says in his preface that he describes it in its perfect completion, as it has been transformed and adorned by the labours of Theodosius II (invicti principis); and we can fix the date of its composition to A.D. 447\texthebrew{‑}450, as the \emph{double} wall of Theodosius is mentioned

(p242 ed. Seeck) . See

Bury, \emph{Eng. Hist. Review}, XXXI p442 (1916) . (2) The \emph{\textgreek{Πάτρια Κωνσταντινοπόλεως}}, a work of the end of the tenth century, first published by Banduri, and known as the Anonymus Banduri, but recently edited critically by Preger. The \emph{Antiquities} of Codinus is only a corrupt copy of this work. (3) The treatise \emph{De cerimoniis} of Constantine Porphyrogennetos (10th century). (4) Petrus Gyllius, \emph{De topographia Constantinopoleos} (16th century). Much information is also derived from the descriptions of other foreign visitors, in the later Middle Ages, which need not be enumerated here. Of modern books the older are of little value now, except Ducange's \emph{Constantinopolis Christiana.} For the more recent see

Bibliography, II.2, E .

} as it had been left by the Emperor Septimius Severus, who had first destroyed and then restored it.\texthebrew{⁠} a The area enclosed by his wall occupied only a small portion of the later city, lying entirely to the east of a line drawn southward from the modern bridge.\texthebrew{⁠}\footnote{The wall of Severus appears not to have reached the southern coast of the promontory but to have turned eastward, south of the Hippodrome.

} The central place in old Byzantium was the Tetrastoon, north of the Great Hippodrome which Severus built but left incomplete. In the

p69
north-east corner rose the fortified Acropolis, on which stood the chief temples. Against the eastern side of the hill, close to the shore, were a theatre and amphitheatre ( Kynêgion ); on the north a Stadion , for foot-races ; on the north-west, the Stratêgion , an open space for military drill.

The area of Constantine's city was about four times as large. He built a wall across the promontory from the Propontis to the Golden Horn, about two miles to the west of the wall of Severus. Of this wall of Constantine nothing is left, and its course can only be traced approximately; for within a city the city was enlarged, a new land fortification was built, and the founder's wall was allowed to fall into decay and gradually disappeared.\footnote{One of the gates, the Porta Aurea (also called Old Gate), survived the Turkish conquest and was destroyed by an earthquake in 1508. The Turks knew it as Isa Kapussi. Van Millingen, \emph{ib.} 21, 30.

}

The New Rome, as Constantinople was called, dissimilar as it was from the Old in all its topographical features, was nevertheless forced to resemble it, or at least to recall it, in some superficial points. It was to be a city of seven hills and of fourteen regions. One of the hills, the Sixth, lay outside the wall of Constantine, on the Golden Horn, and had a fortification of its own. This was the Fourteenth Region. The Thirteenth Region lay on the northern side of the Horn (in Galata) and corresponded to the Region beyond the Tiber in Rome.\footnote{Other points of resemblance were the proximity of the Great Palace to the Hippodrome, recalling that of the Circus Maximus to the palaces of the Palatine; and the erection of a building called the Capitolium on the Second Hill. The Milion in the Augusteum corresponded to the Milliarium in the Roman Forum. As Rome had a hieratic name, Flora, so the personified city of Constantinople had a corresponding secret name Anthûsa (Flowering). See John Lydus, \emph{De mens.} IV. 25 ,

50 , 51 ; Stephanus Byz. \emph{s.v.} \textgreek{Συκαί}; Paulus Silent. \emph{Hagia Sophia}, v.156 \textgreek{χρυσοχίτ\texthebrew{ῶ}ν \texthebrew{Ἀ}νθο\texthebrew{ῦ}σα}. In \emph{Chron. Pasch.}, \emph{s.a.} 328 , it is said that the Tyche or personification of the city was named Anthusa.

}

Constantine was more success ­ful perhaps than he had hoped in attracting inhabitants to his eastern capital. Constantinople was dedicated in A.D. 330 (May 11),\texthebrew{⁠}\footnote{The Encaenia of the city were celebrated annually on this date. Cp. Hesychius, \emph{Patria}, p154.

} and in the lifetime of two generations the population had far outgrown the limits of the town as he had designed it. The need of greater space was met partly by the temporary expedient of filling up the sea, here and there, close to the shore, and a suburban town was growing

p70
up outside the Constantinian wall.\texthebrew{⁠}\footnote{Cp. Himerius, \emph{Or.} vii.7, p522 (reign of Julian). For the growth of the population in the fourth century compare

Zosimus 2.35 ; Eunapius, \emph{Vit. Aedes.}, p22; Sozomen, II.3.

} The desirability of enlarging the city was forced upon the government,\texthebrew{⁠}\footnote{Themistius said (\emph{Or.} 18, p223), in A.D. 384, that ``if the city goes on growing as it has recently, it will require next year a new circuit of wall.''

} and early in the reign of Theodosius II the matter was taken in hand. Anthemius, Praetorian Prefect of the East and pilot of the State during the Emperor's minority, may be called, in a sense, the second founder of Constantinople; the stones of his great wall still stand, an impressive monument of his fame.

The new line of circuit was drawn \texthebrew{•} about a mile to the west of the old. The Anthemian wall did not extend the whole way from sea to sea. It was planned so as to take advantage of the fortification round the Sixth Hill, within which the Palace of Blachernae stood, but this north-western quarter of the city has been so changed, partly by subsequent constructions and partly by demolition, that it is impossible, at least without systematic excavation, to determine how the line of defence ran in the fifth century.\footnote{See the interesting discussion in van Millingen, \emph{op. cit.} chap. viii.

}

The wall which was constructed under the auspices of Anthemius ( A.D. 413)\texthebrew{⁠}\footnote{Cp. \emph{C. Th.} XV.1.51 ; Socrates, \emph{H.E.} vii.1.

} sustained extensive damages from an earthquake in A.D. 447. It was then restored and strengthened by the exertions of the Praetorian Prefect Constantine, and a new outer wall was erected.\texthebrew{⁠}\footnote{The building of the wall in sixty days is recorded in inscriptions, of which two, one in Latin, the other in Greek hexameters, are still to be read on the Porta Rhegii. The Latin runs:

Theodosii iussis, gemino nec mense peracto,

Constantinus ouans haec moenia firma locauit.

tam cito tam stabilem Pallas uix conderet arcem.

See van Millingen, \emph{op. cit.} p47.

Thayer's Note: For a photo of most of the Greek inscription, and another of the gate itself, see

Roberto Piperno's page .

} At this time the city might have been exposed at any moment to an attack of the Huns, and the whole work was executed with incredible rapidity in the course of a few months.

The fortification, thus completed and enlarged, was never afterwards structurally altered. It consists of five parts. The inner wall, which was the main defence, had a mean thickness of about 14 feet, and was strengthened by ninety-six towers, 60 feet high, about 60 yards apart. Each tower had two chambers, of which the upper, entered from the parapet of the wall, contained munitions, and was always occupied by watchmen.

p71
Between the inner and the outer wall was a terrace ( peribolos ) from 50 to 64 feet broad. The outer wall was only 2 to 6½ feet thick, and it was built for the most part in arches; it too had ninety-six towers, varying from 30 to 35 feet in height. Outside the wall was an embankment,\texthebrew{⁠}\footnote{\textgreek{Τ\texthebrew{ὸ ἔ}ξω παρατείχιον.}} 61 feet broad; and outside the embankment a ditch, of varying depth,\texthebrew{⁠}\footnote{It is still 22 feet deep in front of the Golden Gate.

} also 61 feet broad, and divided by low dams.

The fortification was pierced by ten gates, of which five were exclusively for military purposes. The two sets, civil and military, were arranged alternately. The chief and most famous entrance, nearest to the Sea of Marmora, was the Golden Gate. It may have been erected by Theodosius the Great as a triumphal arch in memory of his victory over the rebel Maximus. This imposing structure was pierced by three archways and was built of huge square blocks of polished marble. Above the central archway, on either front, it bore the following inscription in metal:

This designation of the arch as a gate suggests that Theodosius may have already contemplated the enclosure of the city by a new wall.\footnote{Against the view (of Strzygowski) stated in the text, E. Weigand (\emph{Das Goldene Tor}, in \emph{Ath. Mitt.} xxxix.1 \emph{sqq.}) has argued that Theudosius is Theodosius II and the tyrant John (see below,

p222 ), that decorat, construit auro, mean gilding, not building, and that the structure was originally built as a gate of the Anthemian wall.

}

The other four public gates were those known by the names of Melantias, Rhegion, St. Romanus, and Charisius.\texthebrew{⁠}\footnote{The road issuing from the Porta Melantiados led to Melantias and Selymbria. In later times it was called the Gate of Selybria and is now known as Selivri Kapussi. Since the later fifth century it was also known as the Gate of the Pêgê, from a holy well close at hand. The Gate of Rhegion (named from the town on the Marmora at Kuchuk Chekmeje) was also known as Porta Rusia (a reference to the Red Faction of the circus). The Gate of Romanus is that known to the Turks as Top Kapussi (Cannon Gate). The most northerly gate, that of Charisius, was also called the Gate of Polyandrion from the cemetery which lay outside the city near this point. The traveller to Hadrianople would quit the city by this egress, and it is called by the Turks the Gate of Hadrianople (Edirne Kapussi).

} The stretch of wall descending from the Gate of St. Romanus into the valley

p72
of the Lycus, and then ascending to the Gate of Charisius, was known as the Mesoteichion or Middle Wall, and when the city was attacked the enemy usually selected it as the most vulnerable portion of the defences. The gates divided the wall into six sections, each of which had its own division of the garrison, distinguished as the First, the Second, and so on. In each section, except in the short one between the Golden Gate and the sea which was manned by the First division, there was a military gate giving access to the terrace, and these gates were distinguished by the number of the division. Thus the military gate between the Porta Aurea and the Porta Melantiados was known as the gate of the Second.\texthebrew{⁠}\footnote{\textgreek{\texthebrew{Ἡ} πύλη το\texthebrew{ῦ} δευτέρου}. It is convenient in modern languages to call these gates the Second, Third, etc., military gate; but the true nomenclature prevents us from asking the question, where was the First?

} The gate of the Sixth, north of the Porta Charisii, was called the gate of the Xylokerkos, from a wooden circus which was near it.

It was twenty-five years after the completion of the wall of Anthemius that the sea-walls of the Constantinian city were extended along the Golden Horn and the Marmora to join the new line of fortification. This work seems to have been carried out under the direction of Cyrus, Prefect of the city, in A.D. 439.\footnote{\emph{Chron. Pasch., sub a.} Cyrus was afterwards credited with the subsequent additions to the land wall which were due to the Prefect Constantine, and has even been identified with him. Cp. van Millingen, p48.

}

The Thirteenth Region, beyond the Golden Horn, known as Sycae, and subsequently as Galata,\texthebrew{⁠}\footnote{We do not meet this, the modern name of the region, before the eighth century. See Theophanes, A.M. 6209 (A.D. 717).

} was not fortified, and, though formally a part of the city, it was virtually a suburb. The regular communication with this region was by ferry,\texthebrew{⁠}\footnote{The ferry started close to the Arsenal, near the modern outer bridge. A gate in the sea wall at this point was called the Gate of the Ferry (\textgreek{το\texthebrew{ῦ} περάματος}).

} but the Golden Horn was also crossed by a wooden bridge of which the southern end was at Blachernae.\texthebrew{⁠}\footnote{Pontem sublicium siue ligneum, \emph{Not. urb. Cpl.} p241. The stone bridge was built by Justinian, \emph{Chron. Pasch., sub a.} 528 .

} In the sixth century this was replaced by a bridge of stone.

The Golden Horn itself was the great port of Constantinople. But there were also small harbours on the Propontis. At the end of the fourth century there were two: the Harbour of Eleutherius or of Theodosius,\texthebrew{⁠}\footnote{At Vlanga Bostan.

} and farther east the Harbour of Julian, also known as the New Harbour, and after the sixth

p73
century as the Harbour of Sophia.\texthebrew{⁠}\footnote{At Kadriga Lamini.

} At these wharves the corn º -ships from Egypt were probably unloaded, for between them were situated the Alexandrine grain magazines.\texthebrew{⁠}\footnote{Horrea Alexandrina, also the Horreum Theodosianum, in the Ninth Region.

} In the fifth century the harbour of Eleutherius, which Theodosius the Great had improved and honoured with his own name, was filled up and disused, but a small new harbour was built near it known as the Portus Caesarii.\texthebrew{⁠}\footnote{Perhaps in the reign of Leo I. Van Millingen would identify this harbour with that which in later times was called Heptaskalon (seven piers), \emph{op. cit.} 301 \emph{sqq.}

} It was probably not till a later period, but before the end of the sixth century, that the port of Hormisdas (afterwards known as that of Bucoleon) was constructed.\texthebrew{⁠}\footnote{Another harbour, the Kontoskalion (short-pier), is first mentioned in the eleventh century. Van Millingen locates it between the harbours of Caesarius and Julian.

} These small harbours on the Propontis were a great convenience, indeed a necessity. For the frequently prevailing north winds often rendered it very difficult for ships to round the promontory and enter the Golden Horn. In that gulf the chief landing-place was the Portus Prosphorianus, also called the Bosporion, under the Acropolis and close to the Arsenal.

In founding a new city, one of the first things which the practical Romans provided was an abundant supply of water.\texthebrew{⁠} b The construction of aqueducts was a branch of engineering which they had brought to perfection, and it was a task of little difficulty to bring in water from the northern hills. A ruined bit of the old aqueduct is still a striking object in the centre of the city.\texthebrew{⁠}\footnote{An extension built by Valens A.D. 368.

} Many reservoirs and cisterns, both open and covered, supplied the inhabitants with water;\texthebrew{⁠}\footnote{The remains of the cisterns have been studied in full detail by Strzygowski and Forchheimer, \emph{Die Wasserbehälter Cpels.} Strzygowski had identified the Cisterna Modestiaca (A.D. 369) with Sarrâdshchane, near the aqueduct of Valens. The Cist. Aetii (\emph{c.} A.D. 368) was on the Sixth Hill near the Tekfur Serai; the Cist. Theodosiana near the mosque of Valideh. The Cist. Asparis (A.D. 459) is probably Kara Gumrûk, in north-west of the city, outside the Constantinian wall. The Cist. S. Mocii is Exi Marmara

(see plan) .

All these were open reservoirs. Of the covered may be mentioned Cist. Maxima, in the Forum of Constantine, and Cist. Philoxeni, near this Forum, neither of which has been discovered; Cist. Basilica (built by Justinian), adjoining the Basilica, identified (with certainty) with the Yeri Batan Serai; and Cist. Illi (A.D. 528), identified with Bin Bir Derek (= 1001 pillars), W. of the Hippodrome.

} and, a hundred years after the

p74
foundation of the city, there were eight public baths ( thermae ), and 153 private baths in the fourteen Regions.\footnote{See \emph{Not. urb. Cpl.}

}

Constantine accorded to the citizens of his new capital the same demoralising privilege which Rome had so long enjoyed, a free supply of bread at the public expense. The granaries of Africa were still appropriated to the needs of Rome; the fruitful lands of the Nile supplied Constantinople. There were five corn-stores; there were twenty public bakeries, and 117 ``steps,'' from which the bread was distributed to the people, in different parts of the city.\footnote{80,000 loaves were distributed daily. Socrates, II.13.

}

A visitor to Constantinople soon after its foundation would have been struck by the fact that there was no public sign of pagan worship. The gods of Greece and Rome were conspicuously absent. If he were a pagan, he might walk to the Acropolis and gaze sadly on the temples of Apollo, Artemis, and Aphrodite, in which the men of old Byzantium had sacrificed, and which Constantine had dismantled but allowed to stand as relics of the past.\texthebrew{⁠}\footnote{John Malalas, XIII. p324 . Theodosius I turned the temple of Aphrodite into a coachhouse for the chariot of the Praet. Prefect, \emph{ib.} 345.

} From its very inauguration the New Rome was ostensibly and officially Christian.\texthebrew{⁠}\footnote{Augustine, \emph{De civ. Dei}, V.25 ; Eusebius, \emph{Vit. Const.} III.48. There is, however, no reason to reject the statement that Constantine consulted the advice of astronomers in laying out the city (John Lydus, \emph{De mens.} IV.25) .

} Nor did the statue of the founder, as a sun-god, compromise his Christian intention. In the centre of the oval Forum, which he laid out on the Second Hill just outside the wall of the old Byzantium, he erected a high column with porphyry drums, on the top of which he placed a statue of Apollo, the work of an old Greek master, but the head of the god was replaced by his own. It was crowned with a halo of seven rays, and looked towards the rising sun.\texthebrew{⁠}\footnote{Constantine Rhodios (in his poem on the Church of the Apostles, 71 \emph{sqq.}, in \emph{Revue des Études grecques}, ix.) quotes four verses, as an inscription on this column, dedicating the City to Christ. But they are certainly not of the Constantinian epoch.

} The column, blackened by time and fire, and injured by earthquakes, still stands,\texthebrew{⁠}\footnote{It is commonly known as the Burnt Column. The Turks call it Chemberli Tash, hooped pillar.

Thayer's Note: For a photo and further details see the

Çemberlitaş

page at Turkey Travel Planner.

} the one monument of the founder which has survived. Within the pedestal beneath Constantine is said to have placed the Palladium of Rome and several Christian relics.

Lofty columns, as Imperial monuments, were a feature of

p75
Constantinople as of Rome. Theodosius the Great, Arcadius, Marcian, Justinian, all had their memorial pillars like Trajan and Marcus Aurelius. That of Marcian, the least interesting, still towers in the centre of the city;\texthebrew{⁠}\footnote{South of the Mosque of Mohammad the Conqueror. Incisions on the pedestal have made it possible to recover the inscription:

principis hanc statuam Marciani cerne torumque

Tereius vovit quod Tatianus opus.

The column which stands near the N.E. shore of the promontory, under the Acropolis, probably commemorated the victory of Claudius Gothicus over the Goths. It bears the inscription Fortunae reduci ob devictos Gothos.

} and the site of the sculptured column of Arcadius, erected by his son, is marked by the ruins of its high pedestal.

The Tetrastoon (Place of the Four Porticoes), on the First Hill, was the centre of old Byzantium. Constantine laid it out anew, and renamed it the Augusteum in honour of his mother, the Augusta Helena, whose statue he set up here.\texthebrew{⁠}\footnote{\emph{Chron. Pasch., sub a.} 328 ; Hesychius, \emph{Patria}, 40, 2. The site of the Augusteum is the place which the Turks call Aya Sofia Atmeïdan.

} Around it were grouped the buildings which played a principal part in the political life and history of the city. On the north side was the Great Church dedicated to St. Sophia, the Holy Wisdom, which was perhaps founded by Constantine, and certainly completed by his son Constantius.\texthebrew{⁠}\footnote{Dedicated in 360, Socrates, \emph{H.E.} ii.43. For the later sources ascribing the foundation to Constantine, see Antoniades, \textgreek{\texthebrew{Ἔ}κφρασις τ\texthebrew{ῆ}ς \texthebrew{ἁ}γ. Σοφίας}, i.3. Close to St. Sophia was St. Irene, which was certainly built by Constantine, Socrates, I.16, ii.16.

} On the east was the Senate-house , a basilica with the customary apse at the eastern end. On the south was the principal entrance to the Imperial Palace, and near it the Baths of Zeuxippus.\texthebrew{⁠}\footnote{Built by Severus, improved and adorned with statues by Constantine. The Zeuxippus was between the Augusteum and the Hippodrome, but did not touch the Hippodrome, as we know that there was a house, and therefore probably a passage, between. See the epigram of Leontius,

\emph{Anth. Pal.} IX.630 . º It seems likely that this passage is meant by the Diabatika of Achilleus, through which the Hippodrome could be reached from the Palace gate. The Achilleus was probably a statue (Bieliaev, \emph{Byzantina}, I p132), not a bath as some have supposed. The Zeuxippus was in the Augusteum, for acc. to \emph{Chron. Pasch., sub} 197 , it was in the middle of the Tetrastoon. Ebersolt places it outside the Aug. on his plan; but p20 places it ``between the Chalkê and the Milion.''

} The Augusteum was entered from the west, and here was the Milion (Milestone), a vaulted monument, from which the mileage was measured over the great network of roads which connected the most distant parts of the European provinces with Constantinople.\footnote{Ebersolt supposes that the Augusteum was entered through gates (\emph{Le Grand Palais}, p15). But the evidence relates only to a very late period (Nicolaus Mesarites, ed. Heisenberg, p21; beginning of thirteenth century).

}

p76
Passing the Milion one entered the great central thoroughfare of the city, the Mesê or Middle Street,\texthebrew{⁠} c which led, through the chief Fora and public places, direct to the Golden Gate. Descending the First and ascending the Second Hill, it passed on the right the palace of the rich eunuch Lausus,\texthebrew{⁠}\footnote{He lived in the reign of Theodosius II.

} which was a museum of art, and on the left the Praetorium, where the Prefect of the city administered justice.\texthebrew{⁠}\footnote{The section of the street between the Augusteum and the Forum was called the Regia (Royal Street). The colonnades on either side had been built by Constantine and were adorned with statues and marbles. \emph{Chron. Pasch., sub a.} 328 . There seem to have been colonnades (\textgreek{\texthebrew{ἔ}μβολοι}) along the whole length of the Mesê.

} Then it reached the oval Forum of Constantine, generally known as ``the Forum,'' on the north side of which was the second Senate-house . Continuing our way westward we reach the Forum of Taurus, adorned with the column of Theodosius the Great, which could be ascended by an interior staircase. In close proximity to this space was the Capitolium, in which, when a university was established, lecture-rooms were assigned to the professors.\texthebrew{⁠}\footnote{See below,

p231 .

} Just beyond the Forum was a monument known as the Philadelphion,\texthebrew{⁠}\footnote{It is said to have been so called from a representation, apparently plastic, of the meeting of the three sons of Constantine after their father's death. See \emph{Patria}, p177.

} perhaps an archway, where an important main street branched off, leading to the Church of the Holy Apostles and to the Gate of Charisius. Following Middle Street one passed through a place called the Amastrianos, and then bearing south-westward reached the Forum of Bous, so named from an oven shaped like an ox, in which calumnious legend said that Julian the Apostate had burned Christians.\texthebrew{⁠}\footnote{\emph{Ib.} 180.

} The street soon ascended the Sixth Hill and, passing through the Forum of Arcadius,\texthebrew{⁠}\footnote{Now called the Evret Bazaar. The Sixth Hill was known as the Xerolophos.

} reached the old Golden Gate in the wall of Constantine. Just outside this gate was the Exakionion, perhaps a pillar with a statue of Constantine, which gave its name to the locality.\texthebrew{⁠}\footnote{``The Exakionion was a land wall built by the great Constantine. . . . Outside it stood a pillar with a statue of C.; hence the name,'' \emph{Patria}, p180.

} Farther on, before reaching the Golden Gate of Theodosius, a street diverged leading to the Gate of Pêgê.

Many streets must have diverged from this thoroughfare, both northwards and southwards, but only for three have we direct evidence: the two already mentioned leading one to the Pêgê

p77
Gate, the other to the Church of the Apostles, and a third close to the Augusteum, which conducted to the Basilica and the quarter of the Bronzesmiths (Chalkoprateia),\texthebrew{⁠}\footnote{Cp. Bury, \emph{The Nika Riot}, p111.

} where the Empress Pulcheria built a famous church to the Mother of God. The site of the Basilica or law-court can be determined precisely, for the Emperor Justinian constructed beside it an immense covered cistern, which is still preserved,\texthebrew{⁠}\footnote{Technical description in Forchheimer and Strzygowski, \emph{op. cit.} 212 \emph{sqq.}

} a regular under ­ground pillared palace, well described by its Turkish name Yeri Batan Sarai. Julian had endowed the Basilica with a library of 150,000 books, and it was the haunt of students of law.\texthebrew{⁠}\footnote{Cp. Agathias, III.1 .

} The proximity of the cistern seems to have inspired an anonymous writer to pen the following epigram:\footnote{\emph{Anth. Pal.} IX.660.

}

The Church of the Holy Apostles stood in the centre of the city, on the summit of the Fourth Hill.\texthebrew{⁠}\footnote{This hill was called \textgreek{Μεσόλοφον} (central hill), and hence popularly \textgreek{Μεσόμφαλον} (navel), \emph{Patria}, p219.

} It was built in the form of a basilica by Constantine, and completed and dedicated by his son Constantius.\texthebrew{⁠}\footnote{It is described by Eusebius, \emph{Vit. Const.} iv.58. See Heisenberg, \emph{Apostelkirche}, 99, 110.

} Contiguous to the east end Constantine erected a round mausoleum, to receive the bodies of himself and his descendants.\texthebrew{⁠}\footnote{The mausoleums of Diocletian at Salona, of Augustus and Hadrian at Rome, would have naturally suggested the idea. Cp. Schultze, \emph{Konstantinopel}, 13, 15. Heisenberg (\emph{op. cit.} 100, 116), however, thinks that Constantine only contemplated his own burial in the rotunda, that the other twelve sarcophagi were meant as cenotaphs of the Apostles, and that Constantius converted the building into an Imperial mausoleum. The question is difficult, and depends on the interpretation of some phrases in Eusebius, \emph{loc. cit.}

} He placed his own sarcophagus in the centre, and twelve others (the number was suggested by the number of the Apostles) to right and left. This mausoleum remained intact till the Turkish conquest, and many emperors were laid to rest in it; but the church itself was rebuilt in the sixth century. In its new form it was the most magnificent ecclesiastical building in Constantinople, next to St. Sophia, but it was less fortunate than its greater rival. After the Turkish conquest it was destroyed to make room for the mosque of Mohammad the

p78
Conqueror, and no vestige remains of it or of the imperial burying-place .

The Great Palace lay east of the Hippodrome. Ultimately it was to occupy almost the whole of the First Region, extending over the terraced slopes of the first hill down to the sea-shore.\texthebrew{⁠}\footnote{Its northern limit near the shore was marked by the Topoi, a place which has been identified by a tier of seats. See van Millingen, \emph{op. cit.} p256.

} Thus gradually enlarged from age to age it came to resemble the mediaeval palaces of Japan or the Kremlin at Moscow,\texthebrew{⁠}\footnote{Or the Turkish Seraglio which replaced it.

} and consisted of many isolated groups of buildings, throne rooms, reception halls, churches, and summer houses amid gardens and terraces. But the original palace which was designed for Constantine, and to which few or no additions were made till the sixth century, was of more modest dimensions. It was on the top and upper slopes of the hill, and was perhaps not much larger than the fortified residence which Diocletian built for himself at Salona.\texthebrew{⁠}\footnote{For the construction and plan of this palace see Hébrard and Zeiller, \emph{Spalato.}

} It is reasonable to suppose that the two palaces resembled each other in some of their architectural features; but the plan of the palace at Salona can hardly serve as a guide for attempting to reconstruct the palace at Constantinople;\texthebrew{⁠}\footnote{Ebersolt was influenced by the plan of Spalato in his conjectural plan of Constantine's Palace, but I have shown that his reconstruction does not conform to our actual data (see \emph{B.Z.} 21, 210 \emph{sqq.}). He has also sought analogies at the palace of Mschatta in Syria.

} for not only were the topographical conditions different, but the arrangements requisite in the residence of a reigning sovereign could not be the same as those which sufficed for a prince living in retirement. It is indeed not improbable that Constantine's palace, like Diocletian's, was rectangular in form. It was bounded on the west by the Hippodrome, on the north by the Augusteum, and on this side was the principal entrance.\texthebrew{⁠}\footnote{Over this entrance was a painting representing the triumph of Christianity. Constantine with a cross above his head was depicted with his sons, and at their feet a dragon pierced by a dart sank into the abyss.

} This gate was known as the Chalkê, called so probably from the bronze roof of the vestibule. Immediately inside the entrance were the quarters of the Scholarian guards, and here one may notice a resemblance to the palace of Diocletian, in which the quarters of the guards

p79
were close to the chief entrance, the Porta Aurea.\texthebrew{⁠}\footnote{On the right side of the entrance. At Constantinople the Scholarian quarters were in front of the entrance and were traversed in order to reach the interior of the Palace.

} On the western side of the enclosure, towards the Hippodrome, was a group of buildings specially designated as the Palace of Daphne, of which the two most important were the Augusteus, a throne room, on the ceiling of which was represented a large cross wrought in gold and precious stones,\texthebrew{⁠}\footnote{The Augusteus is referred to by Eusebius in \emph{Vit. Const.} iii.49, and iv.66.

} and the Hall of the Nineteen Akkubita, which was used for ceremonial banquets.\texthebrew{⁠}\footnote{But this hall consisted of two parts, probably separated by curtains, one on a higher level in which the banquets were held, and the other a reception hall (triklinos). The building is ascribed to Constantine in \emph{Patria}, p144.

} It is possible that the Tribunal, a large open terrace, lay in the centre of the precincts. On the eastern side were the Consistorium,\texthebrew{⁠}\footnote{Probably a rectangular building like the Consistorium at Mschatta. It was used not only for meetings of the Council but also for the reception of embassies and other functions. In later times there was also a small Consistorium for use in winter.

} or Council Chamber, the Chapel of the Lord,\texthebrew{⁠}\footnote{\textgreek{\texthebrew{Ὁ} Κύριος}. Ascribed to Constantine (\emph{Patria}, p141). It contained relics of the true cross.

} and the quarters of the Candidati and the Protectors.\footnote{These porticoes (\emph{Chron. Pasch.}, \emph{loc. cit.}) were probably replaced in the same area by the Halls of the Excubitors and the Candidati after A.D. 532.

}

If all these buildings, with other apartments and offices,\texthebrew{⁠}\footnote{Ebersolt has not made due allowance in his plan for the private apartments of the Emperor and of the Empress, or for the quarters of the Chamberlains and numerous palace officials. The Master of Offices must have had a bureau in the Palace; likewise the two ministries of finance and the treasuries were doubtless within the precincts. He tacit­ly assumes that the Palace of Constantine as a whole remained intact when later additions were made and the Imperial family ceased to reside in Daphne. This assumption seems to be unwarranted. It is probable that many of Constantine's constructions were removed in later times to make way for others.

} were, as seems not improbable, arranged symmetrically in a rectangular enclosure, there was outside this enclosure another edifice contiguous and in close communication, which might be regarded either as a separate palace or as part of the Great Palace. This was the Magnaura.\texthebrew{⁠}\footnote{See \emph{Patria}, p144. The great Hall of the Magnaura was a basilica with three naves. In the tenth century it was a very magnificent building, but we cannot be sure that the descriptions of it apply to earlier times.

} It was situated on the east side of the Augusteum, close to the Senate-house , and the passage which connected the Great Palace with the precincts of the Magnaura was near the Chapel of the Lord.

On the sea-shore to the south of the Palace was the House of

p80
Hormisdas, which Constantine the Great is said to have assigned as a dwelling to Hormisdas, a Persian prince who had fled to him for protection. In later times this house was enclosed within the grounds of the Great Palace.\texthebrew{⁠}\footnote{The façade of the House of Hormisdas on the sea-shore is still preserved (generally known as the House of Justinian, who resided there before his accession). About 100 yards from here there were till recently remains of another imperial edifice. Both buildings doubtless formed parts of the Palace of Bucoleon. See van Millingen, p275 \emph{sqq.}

} The sea-shore and the lower slopes of the hill, for a long time after the foundation of the city, were covered with the private houses of rich senators, which were destined gradually to disappear as the limits of the Imperial residence were extended.\footnote{The author of the \emph{Notitia} of Constantinople describes the First Region as regiis nobiliumque domiciliis clara, and enumerates 118 mansions.

}

There was another Imperial Palace at Blachernae, in the north-west of the city. We know little of it in early times, but in the thirteenth century it superseded the Great Palace as the home of the Emperors.\footnote{It is mentioned in the \emph{Notitia.} For the position of the palace see van Millingen, 128 \emph{sq.}

}

Much more important in the fourth and fifth centuries was the Palace of Hebdomon on the shore of the Propontis not far from the Golden Gate. The place has been identified with Makri Keui, which is distant exactly \texthebrew{•} seven Roman miles from the Augusteum.\texthebrew{⁠}\footnote{See van Millingen, chap. xix; Bieliaev, \emph{Byzantina}, iii. p57 \emph{sqq.}

} Here there was a plain suitable for a military encampment, and it was called, in reminiscence of Rome, the Campus Martius. The Emperor Valens built a Tribune\texthebrew{⁠}\footnote{Van Millingen takes it for granted (p326) that the harbour was the little bay east of Makri Keui, but Bieliaev thinks that it was at Makri Keui itself, houses and gardens now covering the place where were once the waters and quays of the port.

} for the use of the Emperor when he was reviewing troops, and to him we may probably attribute the foundation of the palace which was afterwards enlarged or rebuilt by Justinian. The place was sanctified by several churches, especially that of the Prophet Samuel containing his remains, and that of John the Baptist which Theodosius I built to receive the sacred relic of the saint's head.\texthebrew{⁠}\footnote{Sozomen, VII.24.

} All the emperors who were elevated at New Rome from Valens to Zeno and Basiliscus were crowned and acclaimed at the Hebdomon. The Campus Martius was to witness many historical scenes, and more than once when the city was visited by earthquakes the panick-stricken populace found it a convenient refuge.

The site of the Hippodrome corresponds to the modern Atmeïdan, which is the Turkish equivalent of the word, and its orientation (NNE to SSW) is exactly marked by three monuments which lay in its axis and still stand in their original positions. Of its general structure and arrangements we can form an idea from what we know of the Circus Maximus at Rome, which seems to have served as its model when it was designed and begun by Septimius Severus before the end of the second century.\texthebrew{⁠}\footnote{Descriptions of the building will be found in Labarte and Oberhummer, \emph{opp. citt.}; in Murray's \emph{Handbook to Constantinople} (the part written by van Millingen), pp39 \emph{sqq.}; in Grosvenor's \emph{Constantinople}, i.319 \emph{sqq.} (a minute reconstruction, of which many details cannot be substantiated); in Paspatês, \emph{Great Palace}, 38 \emph{sqq.}

} But it was of smaller dimensions,\texthebrew{⁠}\footnote{The dimensions of the Circus are given by Pollack (\emph{Circus Maximus}, in Pauly-Wiss.) as follows: length of course = 590 metres (2000 Roman feet); length of building including carceres and semicircle = 635 m; º breadth of arena = 80 m; breadth of building = 150 m. Van Millingen estimates the Hippodrome as ``between 1200 and 1300 feet in length and about half as wide.'' Grosvenor makes it longer and narrower (1382 feet long, 395 feet wide). Van Millingen had probably exaggerated the width, but it is not unlikely that the area occupied by the seats was larger in the Hippodrome than in the Circus Maximus.

} and, completed by Constantine, it had many peculiarities of its own. As there was not enough level ground on the hill, the southern portion, which terminated in a semicircle (the sphendone ), was suspended on massive vaults, which can still be seen. The nature of the site determined an important difference from the arrangement of the Circus Maximus. There the main entrances were at the semi-circular extremity; here this was impossible, and the main entrances (if there was more than one) were on the western side.

At the northern end, as at Rome, were the carceres , stalls for the horses and chariots, and storehouses for all the appurtenances of the races and spectacles. But above this structure, which was an indispensable part of all Roman racecourses, arose the Kathisma, the unique and characteristic feature of the Hippodrome of Constantinople. This edifice, apparently erected by Constantine, was a small ``palace'' with rooms for the accommodation of the Emperor, communicating with the Great Palace by a spiral staircase.\texthebrew{⁠}\footnote{The earliest mention of the staircase (\textgreek{κοχλίας}) is in \emph{Chron. Pasch., s. a.} 380 . It is not clear whether the door of Decimus, which is connected with it here, was at the bottom or the top. The Kathisma could also be reached from the hippodrome itself, as is clear from the story of the Nika riot in A.D. 532.

} In front of it was the Imperial ``box,''

p82
from which the Emperors watched the races \texthebrew{—} the Kathisma or seat which gave its name to the whole building. Immediately below the palace there was a place, probably raised above the level of the course and known as the Stama ,\texthebrew{⁠}\footnote{Also known as the Pi. See Constantine Porph. \emph{Cer.} i.69, pp310, 338; 92, p423.

} which was perhaps occupied during the spectacles by Imperial guards.

Down the middle of the racecourse ran the spina (backbone), a long low wall at either end of which were the goals round which the chariots had to turn. The length of a race was generally seven circuits, and it is probable that the same device was used at Constantinople as at Rome for helping the spectators to remember at any moment the number of circuits already accomplished. At one extremity of the spina seven dolphins were conspicuously suspended, at the other seven eggs \texthebrew{—} emblems respectively of Neptune and of Castor and Pollux, deities associated with horses . As the foremost chariot passed the turning-point , an attendant removed a dolphin or an egg. The spina was adorned by works of art, and three of these ornaments have survived the Turkish conquest. An ancient Egyptian obelisk of Thothmes III, which had been brought from Heliopolis, was placed at the central point of the spina by Theodosius the Great, on a pedestal with bas-reliefs representing the Emperor and his family witnessing races.\texthebrew{⁠}\footnote{The obelisk is 60 feet high. The bas-relief on the north side represents (1) below \texthebrew{—} the erection of the obelisk; (2) above \texthebrew{—} the Kathisma with upper and lower balconies; Theodosius with his two sons is seated in the upper, on either side are courtiers and guards. On the east: (1) above \texthebrew{—} Kathisma, as before; Theodosius holds crown for the victor in a race, and in the lower balcony are a number of persons, including musicians; (2) below \texthebrew{—} a Latin inscription recording the erection of the obelisk. On the south: (1) above \texthebrew{—} Kathisma, Imperial family in upper balcony, courtiers in the lower; in front on steps two mandatores, addressing the people for the Emperor; (2) below \texthebrew{—} a chariot race. On the west: (1) above \texthebrew{—} Kathisma, Imperial family in upper balcony, barbarians bringing tribute in lower; (2) below \texthebrew{—} a Greek inscription on the erection of the obelisk. These reliefs supply some material for a conjectural construction of the front of the Kathisma.

Thayer's Note: For several photos of the Obelisk of Thothmes III; plus other photos of the remains of the Hippodrome, the Columns of Constantine, of Marcian, and of the Goths, and the Aqueduct of Valens, see

Roberto Piperno's page on Roman monuments still visible in Istanbul .

} The choice of the position for this monument was doubtless suggested by the fact that Augustus had placed in the centre of the spina of the Roman Circus the obelisk which now stands in the Piazza del Popolo. South of the memorial of Theodosius is a more illustrious relic of history, the bronze pillar shaped of three serpents whose heads had once supported the gold tripod which the Greeks dedicated to Apollo at Delphi after the great deliverance of Plataea. Constantine had carried it off from Delphi when he despoiled Hellas to adorn

p83
his new capital. The third monument, which stands farther south, is a column of masonry, which originally rose to the height of 94 feet and was covered with plates of gleaming bronze. The bronze has gone, and the upper half of the pillar.\texthebrew{⁠}\footnote{As to its date we only know from the inscription which remains on the pedestal that by the reign of Constantine VII in the tenth century it had suffered from the injuries of time (\textgreek{χρόν\texthebrew{ῳ} φθαρέν}) and required restoration. Paspatês (\emph{op. cit.} p42) gives the distance from the Egyptian obelisk to the bronze pillar as 94 paces.

} There were many statues and works of art, not only along the spina , but in other parts of the Hippodrome, especially in the long promenade which went round the building above the tiers of seats. The façade of the Kathisma was decorated with the four Horses of Lysippus,\texthebrew{⁠}\footnote{It is said that they were brought from Chios by Theodosius II.

Thayer's Note: For two excellent photos of the horses and comprehensive details on them, including a reconstruction of the Kathisma, see

The Horses at St. Mark's in Venice

by Tom Wukitsch.

} in gilt bronze, which were carried off to Venice by the Doge Dandolo, after the capture of the city by the brigands of the Fourth Crusade, and now adorn the front of San Marco.

The accommodation for spectators may have been larger than in the original Circus Maximus, where, according to a recent calculation, there may have been room for 70,000 or 80,000.\texthebrew{⁠}\footnote{In the time of Augustus; in that of Constantine, perhaps it was more than double (Hülsen, in Jordan, \emph{Top. d. S. Rom.} I.III.137). Paspatês calculates that the Hippodrome accommodated 60,000, Grosvenor 80,000.

Thayer's Note: For further discussion of the seating capacity of the structure in Rome, see the article

Circus Maximus

in Platner and Ashby's \emph{Topographical Dictionary of Ancient Rome}, a more recent work than Jordan's \emph{Topographie.}

} The tiers of seats rose higher; it appears that there were over thirty rows. Special seats, probably on the lowest row, were reserved for senators,\texthebrew{⁠}\footnote{Marcellinus, \emph{Chron., s. a.} 528.

} and it was customary for members of the Blue Faction to sit on the west side of the building, to the right of the throne, and those of the Green on the east.

The spectators entered the Hippodrome from the west. We know that there was one main entrance close to the Kathisma, and it was probably known as the Great Gate.\texthebrew{⁠}\footnote{Const. Porph. \emph{Cer.} I.68, p307 . The existence of a principal gate here is generally admitted. The position of the entrances is discussed by Labarte, \emph{loc. cit.} His assumption, on grounds of symmetry, that there were gates on the E side exactly opposite to those on the W is arbitrary. The question of the gates is important in connexion with the Nika riot of A.D. 532. See below,

Vol. II, Chap. XV . º

} We may consider it likely that there was another ingress farther south, though its existence is not expressly recorded.\texthebrew{⁠}\footnote{It is assumed by Labarte, and is probable on grounds of convenience (to avoid congestion).

} The only other issue of which we hear in early times was the Dead Gate, which, from is name, is supposed to have been used for carrying out corpses. It seems to have been somewhere in the eastern wall

p84
of the building.\texthebrew{⁠}\footnote{Labarte placed it near the Sphendone, but there is no evidence. If conjecture is permissible, it may have been in the centre of the eastern wall, where the Skyla gate was afterwards constructed (probably by Justinian II).

} In later times there was a gate into the Palace near the Kathisma, but in the fifth and sixth centuries the only passage from the Hippodrome to the Daphne Palace was through the Kathisma itself and the winding stair which has been mentioned.\footnote{The absence of any entrance here may be inferred from the circumstances of the suppression of the Nika riot. I have shown that in the seventh and eighth centuries there was a covered hippodrome on the E. side of the great Hippodrome (and about half as long) between it and the Palace grounds; but there is no evidence that it existed in the fifth or sixth century. See Bury, \emph{Covered Hippodrome}, 113\texthebrew{‑}115.

}

Since the establishment of the Empire, chariot-races had been a necessity of life for the Roman populace. Inscriptions, as well as literary records, of the early Empire abundantly illustrate the absorbing interest which was found by all classes in the excitement of the circus, and this passion, which Christianity did nothing to mitigate, was inherited by Constantinople. Theologians might fulminate against it, but their censures produced no greater effect than the declamations of pagan satirists. In the fifth and sixth centuries, charioteers were as wealthy a class as ever; Porphyrius was as popular an idol in the days of Anastasius as Scorpus and Thallus had been in the days of Domitian, or Diocles in those of Hadrian and Antoninus. Emperors, indeed, did not follow the unseemly example of Nero, Commodus, and other dissolute princes, and practise themselves the art of the charioteer, but they shared undisguisedly in the ardours of partisan ­ship for one or other of the Circus Factions, which played a far more conspicuous part at Constantinople for a couple of centuries than they had ever played at Rome.

The origin of the four Factions, named after their colours, the Blues, Greens, Reds, and Whites, is obscure. They existed in the last age of the Republic,\texthebrew{⁠}\footnote{The Reds and Whites, at least; some think that the Blues and Greens (Veneti and Prasini) arose under the Empire.

} and they were perhaps definitely organised by contractors who supplied the horses and chariots when a magistrate or any one else provided a public festival. The number of the rival colours was determined by the fact that four chariots generally competed in a race, and there consequently arose four rival companies or Factions, requiring considerable staffs of grooms, mechanics, and messengers, and supported

p85
by what they received from the givers of the festivals, who paid them according to a regular tariff.\footnote{Friedländer, \emph{Roman Life and Manners}, ii.27.

}

In every class of the community, from the Emperor down, people attached their sympathies to one or other of the rival factions. It would be interesting to know whether this partisan ­ship was, like political views, frequently hereditary. In the fourth century a portion of the urban populations, in the greater cities of the east, was officially divided into partisans of the four colours, and used for purposes which had no connexion with the hippodrome. They were organised as quasi-military bodies, which could be used at need for the defence of the city or for the execution of public works.\texthebrew{⁠}\footnote{The part taken by the demes in restoring and extending the walls of Theodosius II at Cple. is recorded in \emph{\textgreek{Πάτρια}}, p150. See van Millingen, \emph{Byz. Cple.} pp44, 79. The name of the 3rd military gate, Rusion, may refer to its construction by the Red deme. In later times we have cases of the demes defending the walls. For the organisation at Alexandria, cp. M. Gelzer, \emph{Studien}, p18.

} In consequence of this official organisation, embra­cing the dêmos or people, the parties of the hippodrome came to be designated as the demes,\texthebrew{⁠}\footnote{\textgreek{Δημότης} was used to designate the member of a deme, and \textgreek{δημοτεύω} was used in two senses \texthebrew{—} (1) neuter, to be a \textgreek{δημότης}; (2) trans., to arm \textgreek{δημόται} for military service (Theophanes, A.M. 6051). \textgreek{μέρη} was the ordinary word for the circus parties.

} and they were placed under the general control of demarchs, who were responsible to the Prefect of the city. We do not know on what principle the members of the demes were selected from the rest of the citizens, most of whom were attached in sympathy to one or other of the colours; but we may assume it to be probable that enrolment in a deme was voluntary.\footnote{There is abundant evidence to show that the demes included only a portion of the urban population (see Rambaud, \emph{De Byz. Hippodromo}, pp87, 88;

Reiske, \emph{Comm. ad Const. Porph. de Cer.} pp28, 29 ).

}

Like the princes of the early Empire, the autocrats of the fifth and sixth centuries generally showed marked favour towards one of the parties. Theodosius II was indulgent to the Greens,\texthebrew{⁠}\footnote{He changed the seats of the Greens from the right to the left of the Kathisma (John Mal. xiv. p351).

} Marcian favoured the Blues, Leo and Zeno the Greens, while Justinian preferred the Blues. These two parties had risen into such importance and popularity that they completely overshadowed the Reds and Whites, which were gradually sinking into insignificance\texthebrew{⁠}\footnote{Thus in an important passage of Theophylactus Simocatta (who wrote early in the seventh century), \emph{Hist.} VIII.7.11, only two parties are recognised, \textgreek{ε\texthebrew{ἰ}ς δύο γάρ χρωμάτων \texthebrew{ἐ}φέσεις τ\texthebrew{ὰ} τ\texthebrew{ῶ}ν}(p86)
\textgreek{\texthebrew{Ῥ}ωμαίων καταρ\texthebrew{ἔ}κτωκε πλήθη}.\texthebrew{—} The history of the demes has been investigated in the important article of Uspenski, \emph{Partii tsirka i Dimy v Kplie} in \emph{Viz. Vrem.} i.1 \emph{sqq.}

} and were destined ultimately, though they

p86
retained their names, to be merged in the organisations of the Greens and Blues respectively.

While the younger Rome inherited from her elder sister the passion for chariot races,\texthebrew{⁠}\footnote{The popularity of the circus with the Romans of the sixth century is noted in Cassiodorus, \emph{Var.} III.51, 11: illic supra cetera spectacula fervor animorum inconsulta gravitate rapiatur. transit prasinus, pars populi maeret: praecedit venetus et ocius turba civitatis affligitur. Cp. Salvian, \emph{De gub. Dei}, vi.20\texthebrew{‑}26.

} the Byzantine hippodrome acquired a political significance which had never been attached to the Roman circus. It was here that on the accession of a new Emperor the people of the capital acclaimed him and showed their approval of his election. Here they criticised openly his acts and clamoured for the removal of unpopular ministers. The hippodrome was again and again throughout later Roman history the scene of political demonstrations and riots which shook or threatened the throne, and a modern writer has described the spina which divided the racecourse as the axis of the Byzantine world.\texthebrew{⁠}\footnote{Rambaud, \emph{op. cit.} p19 quidam axis fuit quo Byzantinus orbis universus nitebatur.

} It may be said that the hippodrome replaced, under autocratic government, the popular Assembly of the old Greek city-state .

The Romans whom Constantine induced to settle in his new city found in its immediate neighbourhood as favourable conditions as they could desire for the villeggiatura which for hundreds of years had been a feature of Roman life. From Rome they had to travel up to Tibur or Tusculum or Lanuvium, or drive to the seaside resorts of Antium and Terracina, if they did not fare further and seek the attractions of the bay of Naples. At Constantine their villas were in the suburbs near the seashore and could easily be reached by boat. We may divide the suburbs into three principal groups: the western, extending from the Theodosian Wall to Hebdomon; the banks of the Bosphorus; and the Asiatic coast from Chrysopolis (Skutari) south-eastward to Karta Limên (Kartal). The suburb and palace of Hebdomon have already been described.

On the European side of the Bosphorus, outside Galata, was the suburban quarter of St. Mamas, where the Emperors had a

p87
house, which in the eighth and ninth centuries they often frequented.\texthebrew{⁠}\footnote{That there was an imperial house here in the fifth century seems to follow from the fact that in 469, on the occasion of the great fire, Leo I stayed at St. Mamas for six months. He constructed a harbour and portico (\emph{Chron. Pasch., sub a.} ). The question as to the locality was cleared up by the late J. Pargoire, who has definitely identified many of the more important suburbs in his valuable articles (see

Bibliography, ii.2, E ).

} Farther north was one of the two places specially known as the Anaplûs \texthebrew{—} a confusing term, which was also used in the more general sense of the whole European bank of the straits. This, the southern Anaplûs, corresponds to the modern Kuru-Chesme ; the other is at Rumili Hissar. Between these places were the suburbs of Promotus and Hestiae (Arnaut Keui), where there was a famous church of St. Michael, founded by Constantine and rebuilt by Justinian. This must not be confused with another church of the Archangel at Sosthenion, of which the name is preserved in Stenia, \texthebrew{•} about two miles north of Rumili Hissar. On the Asiatic side, opposite Stenia and in the neighbourhood of Kanlija, were the suburbs of Boradion and Anthemius.

Opposite Constantinople itself were the towns of Chrysopolis, beautifully situated on the western slopes of a hill, and Chalcedon, now Kadi Keui. South of Chalcedon the coast turns and trends south-eastward, to form the bay of Nicomedia. Here were the suburbs of Hieria (Fanar Bagche), Drys, the ``Oak'' (Jadi Bostan), Satyros, Bryas ( Mal-tepe ), and Karta Limên. At Drys was Rufinianae, the estate of the Praetorian Prefect Rufinus, where he built a monastery and a mansion; confiscated after his death it became imperial property, and we find the palace sometimes occupied by members of the Imperial family. At Hieria, Justinian built a famous palace as a summer retreat, and in the ninth century Theophilus chose Bryas for the same purpose. These suburbs look across to the group of the Princes' Islands, so admirably situated by their climate for villa-life ; but in the days of the Empire they were not to Constantinople what Capri and Ischia are to Naples and what they were to become in modern times; they were covered with convents and were used as honourable and agreeable prisons for fallen princes.

All these suburban quarters in both continents formed a greater Constantinople connected by water-roads . If we suppose that the population of the city itself and all these suburbs approached a million, we shall probably not be much over the

p88
mark. There are no data for a precise calculation. A writer of the fifth century declares that it was generally admitted that the new city had outstripped Rome in numbers as well as in wealth.\texthebrew{⁠}\footnote{Sozomen, \emph{H.E.} II.3. Chrysostom (\emph{In Acta Ap. Hom.} xi.3) gives 100,000 as the number of Christians and 50,000 as the number of poor (\emph{sc.} Christians) who need public assistance. But we can base no conclusion on figures which are clearly Chrysostom's own guesses. How wildly he guessed is shown by his estimate of the wealth of Constantinople in the same passage. He reckons the value of all the real and personal property to be a million pounds of gold (\emph{i.e.} over £45,000,000), ``or rather \emph{twice or thrice as much.}''

} But unfortunately the population of Rome at this time, and indeed throughout the Imperial period, is highly uncertain; recent computations vary from 800,000 to 2,000,000.\texthebrew{⁠}\footnote{Beloch and Lanciani respectively. There have been many estimates, based on area, the corn distribution, the number of houses and insulae (apartment-houses), etc.

} They vary from 500,000 to 1,000,000 for Constantinople; the probability is that in the fifth century its population was little less than a million.\footnote{It would be too long to go into the evidence, which has been thoroughly sifted and criticised by A. Andreades in his articles \textgreek{Περ\texthebrew{ὶ} το\texthebrew{ῦ} πληθυσμο\texthebrew{ῦ}} and \emph{De la population de Cple} (see

Bibliography, ii.2, C ), in which he has refuted the arguments of E. A. Foord (\emph{The Byzantine Empire}, 1911) that the population was 500,000. His conclusion is that the population was between 800,000 and 1,000,000 at the end of the fifth century. I may observe that the number of domus given in the \emph{Not. Urb. Const.} is 4388; the domus are the palaces and houses of the rich. The number of the insulae or apartment-houses in which the poorer lived is not given. Now in Rome, in the time of Constantine, the number of domus was about 1790, and the number of insulae more than 4400. It is reasonable to suppose that the number of insulae in Constantinople, though not more than double (like that of the domus) the number of insulae in Rome, was at least considerably over 2000; and this would bear out Sozomen's statement (see

penultimate note ) that the new city was more populous than Rome. \texthebrew{—} As to the population of Alexandria the available evidence tends to show that from the early period of the Empire down to the seventh century it was not less than 600,000. For Antioch, Libanius (\emph{Epp.} 1137) gives 150,000, which is much too small. Acc. to John Malalas (Bk. xvii. p420) 250,000 perished in the earthquake of A.D. 526. He was an Antiochene and a contemporary.

Thayer's Note: For a brief discussion of the population of Alexandria, see

Bevan's \emph{House of Ptolemy}, p97 .}

\xchapter{Chapter IV}

Short URL for this page:

bit.ly/BURLAT4

It was the mature judgment of the founder of the Empire that Roman dominion had then reached the due limit of its expansion, and it was a corollary of this opinion of Augustus that all the future wars of Rome should be wars in which defence and not aggression was the motive. His discernment was confirmed by the history of nearly fifteen hundred years. Through the long period of its duration, there were not many decades in which the Roman Empire was not engaged in warfare, but with few exceptions all its wars were waged either to defend its frontiers or to recover provinces which had been taken from it. The only clear exception was the conquest of Britain.\texthebrew{⁠}\footnote{Yet in this case too the motive was that the complete Romanisation of the Celts of Gaul could not be accomplished so long as the Celts of Britain, with whom they were in constant communication, remained free.

} For the motive of Trajan's conquest of Dacia and of the lands beyond the Tigris (which were almost immediately abandoned) was not the spirit of aggression or territorial greed or Imperial vanity, so much as the need of strengthening the defences of the Illyrian and eastern provinces. After Trajan there were few cases even of this kind. Diocletian's acquisitions on the Tigris were mainly designed for security, and if any war can be described as a war of self-defence it was that which carried Heraclius into the heart of Persia. There were, indeed, wars of conquest, in which the Roman government took the first step, but they were all to recover lands which had formerly belonged to Rome for centuries. If we regard unprovoked aggression against neighbours as the most heinous crime of which a state can be guilty, few states

p90
have a cleaner record than the later Roman Empire. But it was a crime which there was neither the temptation nor the power to commit. There was little temptation, because there was no pressure of population demanding more territory for expansion; and the Empire was seldom in a position to plan conquests, for all its available forces were required for self-preservation . As in the days of Augustus, there were perpetually two enemies to be faced:

In the east, Parthian was succeeded by Persian, Persian by Saracen, Saracen by Turk. In the west, after the German invasions had reduced the Empire to half its size and the Teutonic kingdoms had been shaped, the Roman rulers had to confront the Frank after the Lombard, the Norman after the Frank, and then the Crusaders. But this was not all. New enemies appeared in the north in the shape of Asiatic nomads and Slavs.

In this chapter we will glance at the three enemies with whom the Empire had to reckon in the fifth century, the Persians, the Germans, and the Huns.

When the Parthian power was overthrown by the revolution of A.D. 226, the Iranian state was renewed and strengthened under a line of monarchs who revived the glories of the ancient Achaemenids, of whom they considered themselves the true successors. Persia under the Sassanid dynasty was recognised by the Roman Empire as a power of equal rank with itself, a consideration which it showed to no other foreign state and had never accorded to the Parthian. The rise of the new dynasty occurred when the Empire was about to enter on a period of internal trouble which shook it to its foundations, and nothing shows more impressively the efficacy of the reforms which were carried out at the end of the third century than the fact that for the following three hundred years the Romans (notwithstanding the perpetual struggles which claimed their energy in Europe) were able to maintain their eastern frontiers, without any serious losses, against this formidable and well-organised enemy.

The two most conspicuous features of the Persian state were

p91
the hereditary nobility and the Zoroastrian church. The first was a point of sharp contrast, the second of remarkable resemblance, to the Roman Empire. The highest nobility were known as ``the people of the Houses,''\texthebrew{⁠}\footnote{Or ``of the seven Houses.'' On the seven families, which included the royal, see Nöldeke, Excurs 3 to \emph{Tabari}, p437.

} and probably all of them possessed large domains in which they exercised princely rights. But the soundest part of the nation seems to have been the inferior nobility, also landed proprietors, who were known as the Dikhāns. Relations of a sort which may be called feudal are supposed to have existed between the two classes of nobility, and the organisation of the army seems to have been connected with the feudal obligations. Some of the high offices of state were restricted by law to certain families, and the power of the great nobles was frequently opposed to the authority of the kings.

To admirers of ancient Greece and Rome one of the most pleasing features of their condition, compared with that of the subjects of the great Iranian monarchy which threatened them in the east, was the absence of a jealous religion controlled by a priesthood possessing immense power in the state and exerting an extreme conservative influence incompatible with the liberty which the city-states of Europe enjoyed. The establishment of Christianity brought Rome into line with Persia. Henceforward both states were governed by jealous gods. Both realms presented the spectacle of a power ­ful priesthood organised as a hierarchy, intolerant and zealous for persecution. Each district in a Persian province seems to have been under the spiritual control of a Magian high priest (corresponding to a bishop), and at the head of the whole sacerdotal hierarchy was the supreme Archi-mage .\texthebrew{⁠}\footnote{The Magian high priest was called Mobedh; the supreme head, Mobedhan-mobedh.

} In some respects the Magian organisation formed a state within a state. The kings often chafed under the dictation of the priests and there were conflicts from time to time, but the priests generally had the moral support of the nobility behind them. They might be defied for a few years, but their power inevitably reasserted itself.

Although both governments discouraged private peaceable intercourse between their subjects, following a policy which reminds us of China or mediaeval Russia,\texthebrew{⁠} a and the commerce

p92
between the two countries was carried on entirely on the frontiers, the influence of Persia on Roman civilisation was considerable. We have seen how the character of the Roman army was affected by the methods of Persian warfare. We have also seen how the founders of the Imperial autocracy imitated, in however modified a form, the royal ceremonial of the court of Ctesiphon; and from this influence must ultimately be derived the ceremonial usages of the courts of modern Europe. In the diplomatic intercourse between the Imperial and Persian governments we may find the origin of the formalities of European diplomacy.

It is a convention for modern sovrans to address each other as ``brother,'' and this was the practice adopted by the Emperor and the King of kings.\texthebrew{⁠}\footnote{Amm. Marc. XVII.5.10

victor terra marique Constantius semper Augustus fratri meo Sapori regi salutem plurimam dico. Cp. Kavad's letter in John Mal. XVIII. p449; etc. The Empress Theodora addressed the Persian queen as her sister (\emph{ib.} p467). The Emperor never gave the title \textgreek{βασιλε\texthebrew{ὺ}ς βασιλέων} (al\texthebrew{‑}Mudik) to the king; always simply \textgreek{βασιλε\texthebrew{ὺ}ς}. The king called him Kaisir i Rūm.

} Whatever reserves each might make as to his own superiority, they treated each other as equals, and considered themselves as the two lights of the world \texthebrew{—} in oriental figurative language, the sun of the east and the moon of the west.\texthebrew{⁠}\footnote{Malalas, \emph{ib.} p449. Cp. Peter Patr. \emph{fr.} 12, \emph{De leg. gent.} \textgreek{\texthebrew{ὡ}σπερανει δύο λαμπτήρες}. Theophylactus Sim. IV.11. See Güterbock, \emph{Byzanz und Persien}, for a detailed study of the diplomatic forms.

} When a new sovran ascended to either throne it was the custom to send an embassy to the other court to announce the accession,\texthebrew{⁠}\footnote{\textgreek{Κατ\texthebrew{ὰ} τ\texthebrew{ὸ} ε\texthebrew{ἰ}ωθός}, Menander, \emph{fr.} 15 \emph{De leg. Rom.} p188. Several particular instances are recorded.

} and it was consisted a most unfriendly act to omit this formality. The ambassadors enjoyed special privileges; their baggage was exempt from customs duties; and when they reached the frontier, the government to which they were sent provided for their journey to the capital and defrayed their expenses. At Constantinople it was one of the duties of the Master of Offices to make all the arrangements for the arrival of an ambassador, for his reception and entertainment, and, it must be added, for supervising his movements.\texthebrew{⁠}\footnote{The arrangement for the journey from Daras to Constantinople and the reception ceremonies in the sixth century are described by Peter the Patrician (apud Const. Porph. \emph{De Cer.} I.89, 90 ). The journey was very leisurely, 103 days were allowed. Five horses and thirty mules were placed at the envoy's disposal, by an agreement concluded ``in the Praetorian Prefecture of Constantine'' (\emph{ib.} p400). Perhaps this refers to Constantine who was Pr. Pr. in A.D. 505.

} For all important negotiations men of high rank were chosen, and were

p93
distinguished as ``great ambassadors'' from the envoys of inferior position who were employed in matters of less importance.\footnote{Cp. \emph{ib.} p398; Menander, \emph{fr.} 13 \emph{De leg. Rom.} p200; etc.

}

Of the details of the procedure followed in concluding treaties between ancient states we have surprisingly little information. But a very full account of the negotiations which preceded the peace of A.D. 562 between Rome and Persia, and of the manner in which the treaty was drafted, has come down to us, and illustrates the development of diplomatic formalities.\footnote{See below,

Vol. II. Chap. XVI .

}

We may conclude with great probability that it was the intercourse with the Persian court that above all promoted the elaboration of a precise system of diplomatic forms and etiquette at Constantinople. Such forms were carefully adhered to in the relations of the Emperor with all the other kings and princes who came within his political horizon. They were treated not as equals, like the Persian king, but with gradations of respect and politeness, nicely regulated to correspond to the position which they held in the eyes of the Imperial sovran. This strict etiquette, imposed by Constantinople, was the diplomatic school of Europe.

In the fourth century the eastern frontier of the Empire had been regulated by two treaties, and may roughly be represented by a line running north and south from the borders of Colchis on the Black Sea to Circesium on the Euphrates.

Jovian had restored to Persia, in A.D. 363, most, but not all, of the territories beyond the Tigris which Diocletian had conquered;\texthebrew{⁠}\footnote{There has been a confusion in the identification of the provinces recorded to have been conquered (Peter Patr.) and those recorded to have been surrendered

(Amm. Marc. XXV.7, 9) . The question has been recently discussed by Adonts, \emph{Armeniia v epokhu Iustiniana}, pp43, 44. Diocletian conquered the five provinces Arzanene, Zabdicene, Corduene, Sophene, and Ingilene (or Angilene). The first three were restored, the last two retained, in 363. The two districts which Ammian enumerates as also restored, Moxoene and Rehimene, were portions respectively of Arzanene and Zabdicene. Sophene means Little Sophene (NW of Anzitene), and is to be distinguished from Great Sophene = Sophanene (SE of Anzitene), of which Ingilene was a portion.

} and the new boundary followed the course of the Nymphius, which flows from the north into the upper Tigris, then a straight line drawn southward between Nisibis and Daras to the river Aborras, and then the course of the Aborras, which joins the Euphrates at Circesium. Thus of the great strongholds

p94
beyond the Euphrates, Nisibis and Singara were Persian; Amida and Martyropolis, Edessa, Constantia, and Resaina were Roman.\footnote{It will be convenient to enumerate here the following identifications of places mentioned in the eastern campaigns: Amida = Diarbekr; Apamea = Kalaat al\texthebrew{‑}Mudik; Batnae = Seruj; Beroae = Aleppo; Carrhae = Harran; Chalcis = Kinnesrin; Constantia = Uerancher; Edessa = Urfa; Emesa = Hims; º Epiphania = Hama; Hierapolis = Kara Membij; Marde = Mardin; Martyropolis = Mayafarkin; Melitene = Malatia; Resaina (Theodosiopolis) = Ras al\texthebrew{‑}Aïn; Samosata = Samsat; Singara = Sinjar; Theodosiopolis (in Armenia) = Erzerum.

}

The treaty of A.D. 387\texthebrew{⁠}\footnote{Faustus, VI.1 . The correct date has been established by Güterbock (\emph{Römisch-Armenien}, in \emph{Festgabe} of the Juristic Faculty at Königsberg in honour of J. Th. Schirmer, 1900). It is accepted by Baynes and Hübschmann. For the circumstances and the history of Armenia between 363 and 387, see Baynes, ``Rome and Armenia in the Fourth Century,'' in \emph{E.H.R.} XXV, Oct. 1910 . This article proves the value and trustworthiness of the history of Faustus.

} between Theodosius and Sapor III, which was negotiated by Stilicho, partitioned Armenia into two client states, of which the smaller (about one-fifth of the whole) was under a prince dependent on the Empire, the larger under a vassal of Persia. The Roman client, Arsaces, died in A.D. 390, leaving the government in the hands of five satraps. The Emperor gave him no successor, but committed the supervision of the satrapies to an official entitled the Count of Armenia, and this arrangement continued till the sixth century.\footnote{Procopius, \emph{De aed.} III.1 ;

\emph{C. J.} I.29.3 . Cp. Chapot, \emph{La Frontière de l'Euphrate}, p169.

Thayer's Note: For a sharply different Armenian viewpoint, see Vahan Kurkjian's History of Armenia,

Ch. 19, p130 and note 4 .

}

The Roman system of frontier defence, familiar to us in Britain and Germany, was not adopted in the east, and would hardly have been suitable to the geographical conditions. In Mesopotamia, or in the desert confines of Syria, we find no vestiges of a continuous barrier of \texthebrew{•} vallum and foss, such as those which are visible in Northumberland and Scotland and in the Rhinelands. The defensive works consisted of the modern system of chains of forts. The Euphrates was bordered by castles, and there was a series of forts along the Aborras (Khabur), and northward from Daras to Amida.\footnote{The best account of the defences of the eastern frontier will be found in Chapot, \emph{op. cit.}

}

The eastern frontier of Asia Minor followed the Upper Euphrates (the Kara\texthebrew{‑}Su branch), and the two most important bases were Melitene in the south and Satala (Sadagh) in the north.\texthebrew{⁠}\footnote{Roads from Melitene led westward (1) to Arabissus (Yarpuz), (2) to Caesarea (Kaisariyeh), and (3) to Sebastea (Sivas). Roads from Satala: (1) westward to Sebastea and Amasea; (2) northward to the coast; (3) eastward to Erzerum; while Colchis was reached by the Lycus (Chorok) valley.

} Melitene was equally distant from Antioch and Trebizond,

p95
and it could be reached from Samosata either by a direct road or by a longer route following the right bank of the Euphrates. Beyond the Euphrates lay Roman Armenia (as far as a line drawn from Erzerum to the Nymphius), which in itself formed a mountain defence against Persia.

The great desert which stretches east of Syria and Palestine to the Euphrates, and the waste country of southern Mesopotamia, were the haunt of the Nabatean Arabs, who were known to the Romans as Saracens or Scenites (people of the tents). They had no fixed abode, they lived under the sky, and a Roman historian graphically describes their life as a continuous flight: vita est illis semper in fuga .\texthebrew{⁠}\footnote{Amm. Marc. XIV.4.1 . He describes them as nec amici nobis umquam nec hostes optandi.

} They occupied all the strips of land which could be cultivated, and otherwise lived by pillage. They could raid a Roman province with impunity, for it was useless to pursue them into the desert. Vespasian used their services against the Jews. In the third century some of their tribes began to immigrate into Roman territory, and these settlements, which may be compared to the German settlements on other frontiers, were countenanced by the government. Beyond the frontier they remained brigands, profiting by the hostilities between Rome and Persia, and offering their services now to one power and now to the other. In the south many were converted to Christianity in the fourth and fifth centuries, through the influence of the hermits who set up their abodes in the wilderness.\texthebrew{⁠}\footnote{Socrates IV.36, Sozomen. VI.38. Duchesne, \emph{Eglises séparées}, 336 \emph{sqq.}

} These converts belonged chiefly to the tribe of Ghassan, and we shall find the Ghassanids acting, when it suited them, as dependents of the Empire; while their bitter foes, the Saracens of Hira,\texthebrew{⁠}\footnote{Hira was close to the site of the later Arabic foundation of Kufa.

} who had formed a power ­ful state to the south of Babylon, are under the suzerainty of Persia. These barbarians, undesirable either as friends or foes, played somewhat the same part in the oriental wars as the Red Indian tribes played in the struggle between the French and English in North America.

The defence of Syria against the Saracens of the waste was a chain of fortresses from Sura on the Euphrates to Palmyra, along an excellent road which was probably constructed by

p96
Diocletian.\texthebrew{⁠}\footnote{Diocletian organised a systematic defence of the frontier from Egypt to the Euphrates. John Mal. XII. p308.

} Palmyra was a centre of routes leading southward to Bostra, south-westward to Damascus, westward to Emesa, and to Epiphania and Apamea.\footnote{From Apamea a north road followed the valley of the Orontes to Antioch; while the north road from Epiphania ran by Chalcis and Beroea to Batnae (Chapot, 332 \emph{sqq.}). From Batnae an east road reached the Euphrates at Caeciliana (Kalaat al\texthebrew{‑}Najim) via Hierapolis. In north Syria the principal east highway was that from Antioch to Zeugma. Another led via Cyrrhus (Herup-Pshimber, cp. Chapot, \emph{op. cit.} p340) to Samosata.

}

The long fierce wars of the third and fourth centuries, in the course of which two Roman Emperors, Valerian and Julian, had perished, were succeeded by a period of 140 years ( A.D. 363\texthebrew{‑}502) in which peace was only twice broken by short and trifling interludes of hostility. This relief from war on the eastern frontier was of capital importance for the Empire, because it permitted the government of Constantinople to preserve its European provinces, endangered by the Germans and the Huns. This protracted period of peace was partly at least due to the fact that on the Oxus frontier Persia was constantly occupied by savage and power ­ful foes.

The leading feature of the history of Europe in the fifth century was the occupation of the western half of the Roman Empire by German peoples. The Germans who accomplished this feat were not, with one or two exceptions, the tribes who were known to Rome in the days of Caesar and of Tacitus, and whose seats lay between the Rhine and the Elbe. These West Germans, as they may be called, had attained more or less settled modes of life, and, with the exception of those who lived near the sea-coast, they played no part in the great migrations which led to the dismemberment of the Empire. The Germans of the movement which is known as the Wandering of the Peoples were the East Germans, who, on the Baltic coast, in the lands between the Elbe and the Vistula, had lived outside the political horizon of the Romans in the times of Augustus and Domitian and were known to them only by rumour. The evidence of their own traditions, which other facts seem to confirm, makes it probable that these peoples \texthebrew{—} Goths, Vandals, Burgundians, Lombards,

p97
and others \texthebrew{—} had originally lived in Scandinavia and in the course of the first millennium B.C. migrated to the opposite mainland.

It was in the second century A.D. that the East German group began to affect indirectly Roman history. When the food question became acute for a German people, as a consequence of the increase of population, there were two alternatives. They might become an agricultural nation, converting their pasture-lands into tillage, and reclaiming more land by clearing the forests which girdled their settlements and which formed a barrier against their neighbours; or they might migrate and seek a new and more extensive habitation. The East German barbarians were still in the stage in which steady habits of work seem repulsive and dishonourable. They thought that laziness consisted not in shirking toil but in ``acquiring by the sweat of your brow that which you might procure by the shedding of blood.''\texthebrew{⁠}\footnote{Tacitus, \emph{Germ.} c14 .

} Though the process is withdrawn from our vision, we may divine, with some confidence, that the defensive wars in which Marcus Aurelius was engaged against the Germans north of the Danube frontier were occasioned by the pressure of tribes beyond the Elbe driven by the needs of a growing population to encroach upon their neighbours. Not long after these wars, early in the third century, the Goths migrated from the lower Vistula to the northern shores of the Black Sea. This was the first great recorded migration of an East German people. In their new homes they appear divided into two distinct groups, the Visigoths and the Ostrogoths, each of which was destined to have a separate and independent history. How the Visigoths severed themselves from their brethren, occupied Dacia, and were gradually converted to Arian Christianity is a story of which we have only a meagre outline. They do not come into the full light of history until they pour into the Roman provinces, fleeing in terror before the invasion of the Huns, and are allowed to settle there as Federates by the Roman government. The battle in the plains of Hadrianople, where a Roman army was defeated and a Roman Emperor fell, foretold the nature of the danger which was threatening the Empire. It was to be dismembered, not only or chiefly by the attacks of professed enemies from without, but by the self-assertion of the barbarians

p98
who were admitted within the gates as Federates and subjects. The tactful policy of Theodosius the Great restored peace for a while. We shall see how soon hostilities were resumed, and how the Visigoths, beginning their career as a small federate people in a province in the Balkan peninsula, founded a great independent kingdom in Spain and Gaul.

Of the other East German peoples who made homes and founded kingdoms on Imperial soil, nearly all at one time or another stood to Rome in the relation of Federates. This is a capital feature of the process of the dismemberment of the Empire. Another remarkable fact may also be noticed. Not a single one of the states which the East Germans constructed was permanent. Vandals, Visigoths, Ostrogoths, Gēpīds, all passed away and are clean forgotten; Burgundians and Lombards are remembered only by minor geographical names. The only Germans who created on Roman territory states which were destined to endure were the Franks and Saxons, and these belonged to the Western group.

It is probable that the dismemberment of the Empire would have been, in general, a far more violent process than it actually was, but for a gradual change which had been wrought out within the Empire itself in the course of the third and fourth centuries, through the infiltration of Germanic elements. It is to be remembered in the first place that the western fringe of Germany had been incorporated in the Germanic provinces of Gaul. Cöln, Trier, Mainz were German towns. In the second place, many Germans had been induced to settle within the Empire as farmers (colons), in desolated tracts of country, after the Marcomannic Wars of Marcus Aurelius. Then there were the settlements of the laeti , chiefly in the Belgic provinces, Germans who came from beyond the Rhine, and received lands in return for which they were bound to military service. Towards the end of the fourth century we find similar settlers both in Italy and Gaul, under the name of gentiles , but these were not exclusively Germans.\texthebrew{⁠}\footnote{See above,

p40 .

} Further there was a German population in many of the frontier districts. This was not the result of a deliberate policy; Germans were not settled there as such. Lands were assigned to the soldiers ( milites limitanei ) who protected the frontiers, and as the army became more and more

p99
German, being recruited extensively from German colons, the frontier population became in some regions largely German.

In the third century German influence was not visible. The army had been controlled by the Illyrian element. The change begins in the time of Constantine. Then the German element, which had been gradually filtering in, is rising to the top. Constantine owed his elevation as Imperator by the army in Britain to an Alamannic chief; he was supported by Germans in his contest with the Illyrian Licinius; and to Germans he always showed a marked favour and preference, for which Julian upbraids him. Thus within the Empire the German star is in the ascendant from the end of the first quarter of the fourth century. We notice the adoption of German customs in the army. Both Julian and Valentinian I were, on their elevation, raised on the shields of soldiers, in the fashion of German kings. Henceforward German officers rise to the highest military posts in the State, such as Merobaudes, Arbogastes, Bauto and Stilicho, and even intermarry with the Imperial family. An Emperor of the fifth century, Theodosius II, has German blood in his veins.

At the death of Theodosius the Great the geography of the German world, so far as it can roughly be determined, was as follows. On the Rhine frontier there were the Franks in the north, and the federated group of peoples known as Alamanni in the south. The Franks fell into two distinct groups: the Salians, the future conquerors of Gaul, who were at this time Federates of the Empire, and dwelled on the left bank of the Rhine in the east of modern Belgium; and the Ripuarians, whose abodes were beyond the middle Rhine, extending perhaps as far south as the Main, where the territory of the Alamanni began. Behind these were the Frisian coast dwellers, in Holland and Frisia; the Saxons, whose lands stretched from the North Sea into Westphalia; the Thuringians, in and around the forest region which still bears their name. Neighbours of the Alamanni on the Upper Main were the Burgundians.\texthebrew{⁠}\footnote{Somewhere in this neighbourhood too were a portion of the Silings.

} More remote were the Angles near the neck of the Danish peninsula, the Marcomanni in Bohemia, the Silings (who belonged to the Vandal nation) in Silesia, to which they seem to have given their name. The Asdings, the other great section of the Vandals, were still on the Upper Theiss, where they had been settled since the end of the

p100
second century, and not far from them were the Rugians. Another East German people, the Gēpīds (closely akin to the Goths), inhabited the hilly regions of northern Dacia. Galicia was occupied by the Scirians; and on the north coast of the Black Sea were the Ostrogoths, and beyond them the Heruls, who in the third century had left Sweden to follow in the track of the Goths.\texthebrew{⁠}\footnote{A portion of them migrated to the neighbourhood of the Lower Rhine, where they appear in A.D. 286. In the following century they furnished auxilia to the Roman armies on the Rhine. Schmidt, \emph{Deutsche Stämme}, I.344.

} The Pannonian provinces were entirely in the hands of barbarians, Huns, Alans, and a section of the Ostrogoths, which had moved westward in consequence of the Hunnic invasion. Dacia was in the power of the Huns, whose appearance on the scene introduced the Romans to enemies of a new type, from whom European civilisation was destined to suffer for many centuries.

It must not be thought that the inhabitants of central and northern Europe were so numerous that each of the principal peoples could send a host of hundreds of thousands of warriors to plunder the Empire. ``The irregular divisions and the restless motions of the people of Germany dazzle our imagination, and seem to multiply their numbers.''\texthebrew{⁠}\footnote{Gibbon, \emph{Decline}, c. ix. \emph{ad fin.} The evidence as to the numbers is discussed

in an appendix to this chapter .

} Fear and credulity magnified tenfold the hosts of Goths and Vandals and other peoples who invaded and laid waste the provinces. A critical analysis of the evidence suggests that of the more important nations the total number may have been about 100,000, and that the number of fighting men may have ranged from 20,000 to 30,000.

The period of the invasions of the Empire by the East German peoples, from the middle of the fourth century till the middle of the sixth, was the ``heroic age'' of the Teutons, the age in which minstrels, singing to the harp at the courts of German kings, created the legendary tales which were to become the material for epics in later times, and passing into the Norse Eddas, the Nibelungenlied, and many other poems, were to preserve in dim outline the memory of some of the great historical chieftains who played their parts in dismembering the Empire.\texthebrew{⁠}\footnote{The facts are collected, ordered, and illuminated in Chadwick's \emph{The Heroic Age}, 1912.

} It has been the fashion to regard with indulgence these German leaders, who remade the map of Europe, as noble and attractive

p101
figures; some of them have even been described as chivalrous. This was the ``propaganda'' of the nineteenth century. When we coldly examine their acts, we find that they were as barbarous, cruel, and rapacious as in the days of Caesar's foe, Ariovistus, and that the brief description

of Velleius

still applies to them, in summa feritate uersutissimi natumque mendacio genus .

The nomad hordes, known to history as the Huns, who in the reign of Valens appeared west of the Caspian, swept over southern Russia, subjugating the Alans and the Ostrogoths, and drove the Visigoths from Dacia, seem to have belonged to the Mongolian division of the great group of races which includes also the Turks, the Hungarians, and the Finns.\texthebrew{⁠}\footnote{Cp. the classification of the Ural-Altaic languages in Peisker's brilliant chapter, ``The Asiatic Background,'' in

\emph{C. Med. H.} I. p333 . The \emph{Uralic} group includes the three classes, (1) Finnish: Fins, º Mordvins, etc.; (2) Permish; (3) Ugrian: Hungarians, Voguls, Ostyaks; the \emph{Altaic} includes (1) Turkish, (2) Mongolian, (3) Manchu-Tungusic. Peisker's chapter, to which I would refer, supersedes all previous studies of the Altaic nomads.

} It is probable that for many generations the Huns had established their pastures near the Caspian and Aral lakes. It is almost certain that political events in northern and central Asia, occasioning new movements of nomadic peoples, drove them westward; and the rise of the Zhu\texthebrew{‑}zhu, who were soon to extend their dominion from Corea to the borders of Europe, about the middle of the fourth century, is probably the explanation. As rulers of Tartar Asia, the Zhu\texthebrew{‑}zhu succeeded the Sien\texthebrew{‑}pi , and the Sien\texthebrew{‑}pi were the successors of the Hiung\texthebrew{‑}nu . It is supposed that the name \emph{Huns} is simply a Greek corruption of Hiung\texthebrew{‑}nu ; and this may well be so. The designation (meaning ``common slaves'') was used by the Chinese for all the Asiatic nomads. But the immediate events which precipitated the Huns into Europe had nothing directly to do with the collapse of the Hiung\texthebrew{‑}nu power which had occurred in the distant past.\footnote{Our knowledge of these revolutions is derived from the Annals of China; De Guignes, \emph{Histoire des Huns}, 4 vols., 1756\texthebrew{‑}1758; E. H. Parker, \emph{A Thousand Years of the Tartars}, 1895 (cp. his article on the ``Origin of the Turks'' in \emph{E.H.R.} XI.1896) ; L. Cahun, \emph{Introduction à l'histoire de l'Asie}, 1896; F. Hirth's article in \emph{S.-B.} of Bavarian Academy, \emph{Phil.-Hist. Kl.} II.245 \emph{sqq.}, 1899; Drouin, \emph{notice sur les Huns et Hioung\texthebrew{‑}nu}, 1894 (he dates the destruction of the Hiung\texthebrew{‑}nu empire by the Sien\texthebrew{‑}pi to \emph{c.} 221 A.D.).

}

The nomad life of the Altaic peoples in central Asia was

p102
produced by the conditions of climate. The word nomad, which etymologically means a grazer, is often loosely used to denote tribes of unsettled wandering habits. But in the strict and proper sense nomads are pastoral peoples who have two fixed homes far apart and migrate regularly between them twice a year, like migratory birds, the nomads of the air. In central Asia, northern tracts which are green in the summer supply no pasturage in winter, while the southern steppes, in the summer through drought uninhabitable, afford food to the herds in winter. Hence arises the necessity for two homes. Thus nomads are not peoples who roam promiscuously all over a continent, but herdsmen with two fixed habitations, summer and winter pasture-lands , between which they might move for ever, if they were allowed to remain undisturbed and if the climatic conditions did not change.\texthebrew{⁠}\footnote{Peisker

(\emph{op. cit.} 327\texthebrew{‑}328)

shows that ``the main cause of the nomad invasions of Europe is not increasing aridity but political changes.''

} Migrations to new homes would in general only occur if they were driven from their pastures by stronger tribes.

The structure of Altaic society was based on kinship. Those who lived together in one tent formed the unit. Six to ten tents formed a camp, and several camps a clan. The tribe consisted of several clans, and the highest unit, the il or people, of several tribes. In connexion with nomads we are more familiar with the word `` horde `` . But the horde was no ordinary or regular institution. It was only an exceptional and transitory combination of a number of peoples, to meet some particular danger or achieve some special enterprise; and when the immediate purpose was accomplished, the horde usually dissolved again into its independent elements.

Milk products are the main food of most of these nomad tribes. They may eke out their sustenance by fishing and hunting, but they seldom eat the flesh of their herds. Their habits have always been predatory. Persia and Russia suffered for centuries from their raids, in which they lifted not only cattle but also men, whom they sent to the slave markets.

The successive immigrations of nomads into Europe, of the ancient Scythians, of the Huns, and of all those who came after them, were due, as has already been intimated, to the struggle for existence in the Asiatic steppes, and the expulsion of the

p103
weakest. Those who were forced to migrate ``with an energetic Khan at their head, who organised them on military lines, such a horde transformed itself into an incomparable army, compelled by the instinct of self-preservation to hold fast together in the midst of the hostile population which they subjugated; for however superfluous a central government may be in the steppe, it is of vital importance to a conquering nomad horde outside it.''\texthebrew{⁠}\footnote{Peisker, \emph{op. cit.} p350.

} These invading hordes were not numerous; they were esteemed by their terrified enemies far larger than they actually were. ``But what the Altaian armies lacked in numbers was made up for by their skill in surprises, their fury, their cunning, mobility, and elusiveness, and the panic which preceded them and froze the blood of all peoples. On their marvellously fleet horses they could traverse immense distances, and their scouts provided them with accurate local information as to the remotest lands and their distances. Add to this the enormous advantage that among them even the most insignificant news spread like wildfire from aul to aul by means of voluntary couriers surpassing any intelligence department, however well organised.''\texthebrew{⁠}\footnote{\emph{Ib.}

} The fate of the conquered populations was to be partly exterminated, partly enslaved, and sometimes transplanted from one territory to another, while the women became a prey to the lusts of the conquerors. The peasants were so systematically plundered that they were often forced to abandon the rearing of cattle and reduced to vegetarianism. This seems to have been the case with the Slavs.\footnote{\emph{Op. cit.} p348.

}

Such was the horde which swept into Europe in the fourth century, encamped in Dacia and in the land between the Theiss and Danube, and held sway over the peoples in the south Russian steppes, the Ostrogoths, Heruls, and Alans.\footnote{Also over the Scirians in Galicia; probably over the Slavs in the Pripet region; perhaps already over the Gepids; and presently over the Rugians, who soon after 400 occupied the regions on the Upper Theiss vacated by the Asdings (cp. Schmidt, \emph{op. cit.} p327).

}

For fifty years after their establishment north of the Danube, we hear little of the Huns. They made a few raids into the Roman provinces, and they were ready to furnish auxiliaries, from time to time, to the Empire. At the time of the death of Theodosius they were probably regarded as one more barbarian

p104
enemy, neither more nor less formidable than the Germans who threatened the Danubian barrier. We may conjecture that the organisation of the horde had fallen to pieces soon after their settlement in Europe.\texthebrew{⁠}\footnote{It is uncertain whether Uldin, the Hun king whom we meet in the reign of Arcadius, was king of all the Huns or only of a portion.} No one could foresee that after a generation had passed Rome would be confronted by a large and aggressive Hunnic empire.

\xsubsection{ON THE NUMBERS OF THE BARBARIANS}

The question of the numbers of the German invaders of the Empire is so important that it seems desirable to collect here some of the principal statements of our authorities, so as to indicate the character of the evidence. These statements fall into two classes.

(1) Large numbers, running into hundreds of thousands.

\textgreek{α}. Eunapius appears to say that the fighting forces of the Visigoths when they crossed the Danube in A.D. 376 numbered 200,000, \emph{fr.} 6, \emph{De leg. gent.} p595. The text of the passage, however, is corrupt.

\textgreek{β}. The mixed host of barbarians who invaded Italy in A.D. 405\texthebrew{‑}406 is variously stated to be 400,000 , 200,000 , or more than 100,000 strong. See below,

Chap. V § 7 . It is to be observed that the lowest of these figures is given (by Augustine) in an argument where a high figure is effective.

\textgreek{γ}. Two widely different figures are recorded for the number of those who fell (on both sides) in the battle of Troyes in A.D. 451, 300,000 and 162,000. See below,

Chap. IX § 4 .

\textgreek{δ}. 150,000 is given (by Procopius) as the number of the Ostrogoths who besieged Rome in A.D. 537. This can be shown, from the circumstances, to be incredible. See below,

Chap. XVIII § 5 .

\textgreek{ε}. The Franks are made to boast, in A.D. 539, that they could send an army of 500,000 across the Alps (Procopius, \emph{B. G.} II .28, 10) . Then they were a great power and had many subjects. A few months before, one of their kings had invaded Italy with 100,000 men (\emph{ib.} 25, 2); but the number is highly suspicious.

(2) Small numbers.

\textgreek{α}. It is difficult to forgive Ammian, who was a soldier and well versed in military affairs, for not stating the number of the forces engaged on either side in the battle of Hadrianople in A.D. 378. The one indication he gives is that the Roman scouts by some

p105
curious mistake reported that the Visigothic forces numbered only 10,000. It is difficult to believe that this mistake could have been made if the Goths, with their associates, had had anything like 50,000 to 100,000 men (Hodgkin's estimate for the army of Alaric), much less the 200,000 of Eunapius. So far as it goes, the indication points rather to a host of not more than 20,000.

\textgreek{β}. After Alaric's siege of Rome in 408, it is stated that his army, reinforced by a multitude of fugitive slaves from Rome, was about 40,000 strong. See below,

Chap. VI § 1 .

\textgreek{γ}. The total number of the Vandal people (evidently including the Alans who were associated with them), not merely of the fighting forces, is stated to have been 80,000 in A.D. 429 (see below,

Chap. VIII § 2 ). They were then embarking for Africa and it was necessary to count them in order to know how many transport ships would be needed. This figure has, therefore, particular claims on our attention.

\textgreek{δ}. The facts we know about the Vandalic and Ostrogothic wars in the sixth century, as related by Procopius, consistently point to the conclusion that the fighting forces of the Vandals and the Ostrogoths were to be counted by tens, not by hundreds, of thousands. Procopius does not give figures (with the exception of one, which is a deliberate exaggeration, see above,

(1) \textgreek{δ} ), but the details of his very full narrative and the small number of the Roman armies which were sent against them and defeated them make this quite clear.

\textgreek{ε}. The total number of the warriors of the Heruls, who were a small people, in the sixth century was 4500 (Procopius, \emph{B. G.} III .34, 42\texthebrew{‑}43).

Intermediate between these two groups, but distinctly inclining towards the first, is the statement of

Orosius, \emph{Hist.} VII .32.11 , that the armed forces of the Burgundians on the Rhine numbered more than 80,000. If the figure has any value it is more likely to represent the total number of the Burgundian people at the beginning of the fifth century.

Schmidt has observed (\emph{Gesch. der deutschen Stämme}, I .46 \emph{sqq.}) that certain numbers in the enumerations of German forces by Roman writers constantly recur (300,000, 100,000, 60,000, etc.) and are therefore to be suspected.

Delbrück (\emph{Gesch. der Kriegskunst}, II .34 \emph{sqq.}) discusses the density of population in ancient Germany and concludes that it was from four to five to the square kilometre.

\xchapter{Chapter V, Part A}

\xsubsection{(Part 1 of 3)}

The Emperor Theodosius the Great died at Milan on January 17, A.D. 395. His wishes were that his younger son, Honorius, then a boy of ten years, should reign in the west, where he had already installed him, and that his elder son, Arcadius, whom he had left as regent at Constantinople when he set out against the usurper Eugenius, should continue to reign in the east.\texthebrew{⁠}\footnote{Flavius Arcadius was born in 377\texthebrew{‑}378, created Augustus Jan. 19, 383, at Constantinople, and was consul in 385. Honorius, born Sept. 9, 384, was created Augustus Jan. 10, 393. As to the succession, we are told that before his death Theodosius had made all the necessary arrangements: Ambrose, \emph{De obitu Theod.} 5 .

} But Theodosius was not willing to leave his youthful heirs without a protector, and the most natural protector was one bound to them by family ties. Accordingly on his deathbed he commended them to the care of Stilicho,\texthebrew{⁠}\footnote{Ambrose in the funeral oration he pronounced in the presence of Honorius says: liberos praesenti commendabat parenti (\emph{ib.}). We must reject the statement of Olympiodorus, fr. 2, that Theodosius appointed Stilicho legal guardian (\textgreek{\texthebrew{ἐ}πίτροπος}) of his sons. The relation of guardian and ward had no existence in constitutional law. Cp. Mommsen, \emph{Hist. Schr.} I. p516.

} an officer of Vandal birth, whom he had raised for his military and other talents to the rank of Master of Both Services in Italy,\texthebrew{⁠}\footnote{Originally serving in the Protectors, he had been raised to the post of Count of Domestics. Then he married Serena and was appointed magister equitum praesentalis (c. 385). After the victory in 394 over Eugenius and Arbogastes, he succeeded the latter as mag. utriusque militiae, and held this supreme command till his death. We do not know who succeeded him as mag. equitum in Italy, but in 401\texthebrew{‑}402 the post was held by one Jacobus, whose name happens to be recorded because he did not admire Claudian's verses

(Claudian, \emph{Carm. min.} 2) . That Stilicho's mother was a Roman may be inferred from Jerome's description of him as semibarbarus

(\emph{Epp.} 123) . His son Eucherius was named after the uncle, his daughter Thermantia after the mother, of Theodosius.

} and, deeming him worthy of an alliance with his own house, had united to

p107
his favourite niece, Serena. It was in this capacity, as the husband of his niece and a trusted friend, that Stilicho received the last wishes of the Emperor; it was as an elder member of the same family that he could claim to exert an influence over Arcadius. Of Honorius he was the natural protector, for he seems to have been appointed regent of the western realm during his minority.

Arcadius was in his seventeenth or eighteenth year at the time of his father's death. He was of short stature, of dark complexion, thin and inactive, and the dulness of his wit was betrayed by his speech and by his sleepy, drooping eyes.\texthebrew{⁠}\footnote{He was educated first by his mother Aelia Flaccilla, then by Arsenius a deacon, and finally by the pagan sophist Themistius. His personal appearance and that of Rufinus are described by Philostorgius

(\emph{H. E.} XI.3) , who lived at Constantinople and must have known them both by sight. That Arcadius seldom appeared outside the Palace has been inferred from the mention in Socrates, VI.23, of the crowds which flocked to see him when on one occasion he did appear in the streets (Seeck, \emph{Gesch. d. Untergangs}, V.545).

} His mental deficiency and the weakness of his character made it inevitable that he should be governed by the strong personalities of his court. Such a commanding personality was the Praetorian Prefect of the East, Flavius Rufinus, a native of Aquitaine, who presented a marked contrast to his sovran. He was tall and manly, and the restless movements of his keen eyes and the readiness of his speech, though his knowledge of Greek was imperfect, were no deceptive signs of his intellectual powers. He was ambitious and unprincipled, and, like most ministers of the age, avaricious, and he was a zealous Christian. He had made many enemies by acts which were perhaps more than commonly unscrupulous, but we cannot assume that all the prominent officials\texthebrew{⁠}\footnote{Promotus, Tatian, and Proclus

(Zosimus, V.51, 52) . Rufinus had become Master of Offices in 388 (cp. Seeck, \emph{Die Briefe des Libanius}, 256\texthebrew{‑}257); he was consul in 392, and in the same year became Praet. Pref. He was on friendly terms with the pagan sophist Libanius (Lib. \emph{Epp.} 784, 1025). His sister Salvia is remembered as one of the early pilgrims to the Holy Land (Palladius, \emph{Hist. Laus.} 142).

Thayer's Note: My copy of the Historia Lausiaca (Palladio, La Storia Lausiaca, critical text and commentary by G. J. M. Bartelink, ed. Fondazione Lorenzo Valla/Mondadori, 1974, p230), has Palladius writing, 55.1, ``\textgreek{τ\texthebrew{ὴ}ν μακαρίαν Σιλβανίαν τ\texthebrew{ὴ}ν παρθένον γυναικαδέλφην \texthebrew{Ῥ}ουφίνου}'' \texthebrew{—} ``the blessed Silvania, the virgin female relative of Rufinus''. While the apparatus indicates a variant reading \textgreek{\texthebrew{ἀ}δελφ\texthebrew{ὴ}ν δ\texthebrew{ὲ}} (``sister''), the only manuscript variant of the woman's name given there is Silvina (\textgreek{Σιλβίναν}). Bartelink's critical apparatus is, however, only partial.

Palladius does not exactly say that Rufinus' sister was on a pilgrimage to the Holy Land, either, merely that he was accompanying her on a trip from Jerusalem (``Aelia'') to Egypt.

} for whose fall he was responsible were innocent victims of his malice. But it is almost certain that he had formed the scheme of ascending the throne as the Imperial colleague of Arcadius.

This ambition of Rufinus placed him at once in an attitude of opposition to Stilicho,\texthebrew{⁠}\footnote{The antagonism was of older date. Theodosius, at the instance of Rufinus, had forbidden Stilicho to punish the Bastarnae who had slain Promotus, whom Rufinus had caused to be exiled.

Claudian, \emph{De laud. Stil.} I. (p108) 94\texthebrew{‑}115 . It may be noted that Zosimus, at the beginning of Book V, represents Rufinus and Stilicho as ethically on a level; but when his source is no longer Eunapius, but Olympiodorus, his tone towards Stilicho changes. Cp. Eunapius, \emph{fr.} 62, 63 \textgreek{\texthebrew{ἄ}μφω τ\texthebrew{ὰ} πάντα συνήργαζον \texthebrew{ἐ}ν τ\texthebrew{ῷ} πλούτ\texthebrew{ῳ} τ\texthebrew{ὸ} κράτος τιθέμενος}. Eunapius was also the source of

John Ant., \emph{frr.} 188\texthebrew{‑}190 (\emph{F. H. G.} IV. p610) .

} who was himself suspected of entertaining

p108
similar schemes, not however in his own interest, but for his son Eucherius. He certainly cherished the design of wedding his son to the Emperor's stepsister, Galla Placidia.\texthebrew{⁠}\footnote{This is unmistakably conveyed in Claudian,

\emph{De cons. Stil.} II.352\texthebrew{‑}361 , and hinted at again in

\emph{De VI. cons. Honorii}, 552\texthebrew{‑}554 .

} The position of the Vandal, who was connected by marriage with the Imperial family, gave him an advantage over Rufinus, which was strengthened by the generally known fact that Theodosius had given him his last instructions. Stilicho, moreover, was popular with the army, and for the present the great bulk of the forces of the Empire was at his disposal; for the regiments united to suppress Eugenius had not yet been sent back to their various stations. Thus a struggle was imminent between the ambitious minister who had the ear of Arcadius, and the strong general who held the command and enjoyed the favour of the army. Before the end of the year this struggle began and ended in a curious way; but we must first see how a certain scheme of Rufinus had been foiled by an obscurer but wilier rival nearer at hand.

It was the cherished project of Rufinus to unite Arcadius with his only daughter; once the Emperor's father-in\texthebrew{‑}law he might hope to become an Emperor himself. But he was thwarted by a subtle adversary, Eutropius, the lord chamberlain ( praepositus sacri cubiculi ), a bald old eunuch, who with oriental craftiness had won his way up from the meanest services and employments. Determining that the future Empress should be bound to himself and not to Rufinus, he chose Eudoxia, a girl of singular beauty, who had been brought up in the house of the widow and sons of one of the victims of Rufinus.\texthebrew{⁠}\footnote{Promotus. His sons had been playmates of Arcadius.

Zosimus, V.3.

} Her father was Bauto, a Frank soldier who had risen to be Master of Soldiers, and for a year or two the most power ­ful man in Italy, in the early years of Valentinian II.\texthebrew{⁠}\footnote{Ambrose, \emph{Epp.} I.24

Bauto qui sibi regnum sub specie pueri vindicavit (words quoted from the tyrant Maximus). In 385 Bauto was consul, as colleague of Arcadius.

} Her mother had doubtless been a Roman, and she received a Roman education, but she inherited, as a contemporary writer observes, barbaric

p109
traits from her German father.\texthebrew{⁠}\footnote{Philostorgius, XI.6

\textgreek{\texthebrew{ἐ}ν\texthebrew{ῆ}ν α\texthebrew{ὐ}τ\texthebrew{ῇ} το\texthebrew{ῦ} βαρβαρικο\texthebrew{ῦ} θράσους ο\texthebrew{ὐ}κ \texthebrew{ὀ}λίγον.}} Eutropius showed a picture of the maiden to the Emperor, and so successfully enlarged upon her merits and her charms that Arcadius determined to marry her; the intrigue was carefully concealed from the Praetorian Prefect;\texthebrew{⁠}\footnote{It is difficult to understand how Rufinus could have been so completely hoodwinked, unless the machinations of Eutropius were carried out during the absence of the Prefect from the court, and he was confronted on his return by a \emph{fait accompli.} We are entitled to conclude from the account of Zosimos (source, Olympiodorus) that Rufinus was absent at Antioch just before the marriage, having gone thither in order to punish Lucian, the Count of the East, for an offence which he had offered to an uncle of the Emperor. Seeck has argued that this visit to Antioch is wrongly dated by Zosimus and belongs to A.D. 393 (\emph{op. cit.} p447), but his reasoning is not convincing. Rufinus did visit Antioch in 393 (as letters of Libanius show), and was in a hurry, but he may have gone there again in 395.

} and till the last moment the public supposed that the bride for whose Imperial wedding preparations were being made was the daughter of Rufinus. The nuptials were celebrated on April 27, A.D. 395. It was a blow to Rufinus, but he was still the most power ­ful man in the east.

The event which at length brought Rufinus into collision with Stilicho was the rising of the Visigoths. They had been settled by Theodosius in the province of Lower Moesia, between the Danube and the Balkan mountains, and were bound in return for their lands to do battle for the Empire when their services were needed. They had accompanied the Emperor in his campaign against Eugenius, and had returned to their homes earlier than the rest of the army. In that campaign they had suffered severe losses, and it was thought that Theodosius deliberately placed them in the most dangerous post for the purpose of redu­cing their strength.\texthebrew{⁠}\footnote{See Sievers, \emph{Studien}, 326; Schmidt, \emph{Deutsche Stämme}, 1.191.

} This was perhaps the principal cause of the discontent which led to their revolt, but there can be no doubt that their ill humour was stimulated by one of their leaders, Alaric (of the family of the Balthas or Bolds), who aspired to a high post of command in the Roman army and had been passed over. The Visigoths had hitherto had no king. It is uncertain whether it was at this crisis\texthebrew{⁠}\footnote{Jordanes, \emph{Get.} 146 ;

Isidore, \emph{Hist. Goth.} (\emph{Chron. min.} II) p272 . But contemporary writers do not use the word \emph{king}, and Schmidt (\emph{ib.} 192) thinks that Alaric was on this occasion only nominated commander-in\texthebrew{‑}chief.

} or at a later stage in Alaric's career that he was elected king by the assembly of his people. In any case he was chosen leader

p110
of the whole host of the Visigoths, and the movements which he led were in the fullest sense national.

Under the leader ­ship of Alaric, the Goths revolted and spread desolation in the fields and homesteads of Thrace and Macedonia. They advanced close to the walls of Constantinople. They carefully spared certain estates outside the city belonging to Rufinus, but their motive was probably different from that which caused the Spartan king Archidamus to spare the lands of Pericles in the Peloponnesian war. Alaric may have wished, not to draw suspicions on the Prefect, but to conciliate his friendship and obtain more favourable terms. Rufinus went to the Gothic camp, dressed as a Goth.\texthebrew{⁠}\footnote{Claudian, \emph{In Rufin.} II.78 \emph{sqq.} Alaric must have moved very early in the spring; for it was still early in the year when Stilicho marched from Italy,

\emph{ib.} 101 .\texthebrew{—} It has been suggested (Seeck, \emph{Gesch. des Untergangs}, V.274) that the zeal of Rufinus against heretics (especially the Eunomians), displayed in a series of four edicts (\emph{C. Th.} XXVI.5.25 ,

XXVI.28, 29 ), was dictated by a superstitious belief that the calamities of the time were due to the anger of Heaven at laxity in the suppression of heresy.

} The result of the negotiations seems to have been that Alaric left the neighbourhood of the capital and marched westward.

At the same time the Asiatic provinces were suffering, as we shall see, from the invasions of other barbarians, and there were no troops to take the field against them, as the eastern regiments which had taken part in the war against Eugenius were still in the west. Stilicho, however, was already preparing to lead them back in person.\texthebrew{⁠}\footnote{He had been occupied with the task of driving out bands of German marauders who had invaded Pannonia and Noricum.

} He deemed his own presence in the east necessary, for, besides the urgent need of dealing with the barbarians, there was a political question which deeply concerned him, touching the territorial division of the Empire between the two sovrans.

Before A.D. 379 the Prefecture of Illyricum, which included Greece and the central Balkan lands, had been subject to the ruler of the west. In that year Gratian resigned it to his new colleague Theodosius, so that the division between east and west was a line running from Singidunum (Belgrade) westward along the river Save and then turning southward along the course of the Drina and reaching the Hadriatic º coast at a point near the lake of Scutari. It was assumed at Constantinople that this arrangement would remain in force and that the Prefecture would continue to be controlled by the eastern

p111
government. But Stilicho declared that it was the will of Theodosius that his sons should revert to the older arrangement, and that the authority of Honorius should extend to the confines of Thrace, leaving to Arcadius only the Prefecture of the East.\texthebrew{⁠}\footnote{Olympiodorus, \emph{fr.} 3. Cp. Mommsen, \emph{op. cit.} 517.

} Whether this assertion was true or not, his policy meant that the realm in which he himself wielded the power would have a marked predominance, both in political importance and in military strength, over the other section of the Empire.

It would perhaps be a mistake to suppose that this political aim of Stilicho, of which he never lost sight, was dictated by mere territorial greed, or that his main object was to increase the revenues. The chief reason for the strife between the two Imperial governments may have lain rather in the fact that the Balkan peninsula was the best nursery in Empire for good fighting men.\texthebrew{⁠}\footnote{This aspect of the question has hitherto been over­looked.

} The stoutest and most useful native troops in the Roman army were, from the fourth to the sixth century, recruited from the highlands of Illyricum and Thrace. It might well seem, therefore, to those who were responsible for the defence of the western provinces that a partition which assigned almost the whole of this great recruiting ground to the east was unfair to the west; and as the legions which were at Stilicho's disposal were entirely inadequate, as the event proved, to the task of protecting the frontiers against the Germans, it was not unnatural that he should have aimed at acquiring control over Illyricum.

It was a question on which the government of New Rome, under the guidance of Rufinus, was not likely to yield without a struggle, and Stilicho took with him western legions belonging to his own command as well as the eastern troops whom he was to restore to Arcadius. He marched overland, doubtless by the Dalmatian coast road to Epirus, and confronted the Visigoths in Thessaly, whither they had traced a devastating path from the Propontis.\footnote{Alaric had experienced a repulse at the hands of garrison soldiers in Thessaly \texthebrew{—} perhaps in attempting the pass of Tempe. See Socrates, \emph{H.E.} VII.10, a confused passage of which little can be made.

}

Rufinus was alarmed lest his rival should win the glory of crushing the enemy, and he induced Arcadius to send to Stilicho

p112
a peremptory order to dispatch the troops to Constantinople and depart himself whence he had come. The Emperor was led, legitimately enough, to resent the presence of his relative, accompanied by western legions, as an officious and hostile interference. The order arrived just as Stilicho was making preparations to attack the Gothic host in the valley of the Peneius. His forces were so superior to those of Alaric that victory was assured; but he obeyed the Imperial command, though his obedience meant the delivery of Greece to the sword of the barbarians. We shall never know his motives, and we are so ill-informed of the circumstances that it is difficult to divine them. A stronger man would have smitten the Goths, and then, having the eastern government at his mercy, would have insisted on the rectification of the Illyrian frontier which it was his cherished object to effect. Never again would he have such a favourable opportunity to realise it. Perhaps he did not yet feel quite confident in his own position; perhaps he did not feel sure of his army. But his hesitation may have been due to the fact that his wife Serena and his children were at Constantinople and could be held as hostages for his good behaviour.\texthebrew{⁠}\footnote{Cp. Mommsen, \emph{ib.} 521 See Claudian,

\emph{In Rufin.} II.95

and

\emph{Laus Serenae}, 232

(Serena kept Stilicho informed by letters of what was going on in the East).\texthebrew{—} The chronology presents a difficulty. Stilicho had set out in the spring, yet Gaïnas and the army did not reach Constantinople till November (see below).

} In any case he consigned the eastern troops to the command of a Gothic captain, Gaïnas, and departed with his own legions to Salona, allowing Alaric to proceed on his wasting way into the lands of Hellas. But he did not break up his camp in Thessaly without coming to an understanding with Gaïnas which was to prove fatal to Rufinus.

Gaïnas marched by the Via Egnatia to Constantinople,\texthebrew{⁠}\footnote{Claudian, \emph{In Rufin.} II.291

\texthebrew{—}percurritur Hebrus,

deseritur Rhodope Thracumque per ardua tendunt,

donec ad Herculei perventum nominis urbem.

The city of Herculean name, Heraclea, is the ancient Perinthus.

} and it was arranged that, according to a usual custom,\texthebrew{⁠}\footnote{Zosimus, V.7.5

\textgreek{ταύτης γ\texthebrew{ὰ}ρ τ\texthebrew{ῆ}ς τιμ\texthebrew{ῆ}ς \texthebrew{ἠ}ξι\texthebrew{ῶ}σθαι το\texthebrew{ὺ}ς στρατιώτας \texthebrew{ἔ}λεγε σύνηθες ε\texthebrew{ἶ}ναι.}} the Emperor and his court should come forth from the city to meet the army in the Campus Martius at Hebdomon. We cannot trust the statement of a hostile writer that Rufinus actually expected to be created Augustus on this occasion, and appeared at the Emperor's side prouder and more sumptuously arrayed than

p113
ever; we only know that he accompanied Arcadius to meet the army. It is said that, when the Emperor had saluted the troops, Rufinus advanced and displayed a studied affability and solicitude to please even towards individual soldiers. They closed in round him as he smiled and talked, anxious to secure their goodwill for his elevation to the throne, but just as he felt himself very nigh to supreme success, the swords of the nearest were drawn, and his body, pierced with wounds, fell to the ground (November 27, A.D. 395).\texthebrew{⁠}\footnote{\textgreek{Πλήθους \texthebrew{ὁ}πλιτευόντος \texthebrew{ἀ}θρό\texthebrew{ᾲ} κινήσει περιπεσών}, Asterius of Amasea, in his \textgreek{Λόγος κατηγορικ\texthebrew{ὸ}ς τ\texthebrew{ῆ}ς \texthebrew{ἐ}ορτης τ\texthebrew{ῶ}ν καλανδ\texthebrew{ῶ}ν}, \emph{P. G.} XL.24.

} His head, carried through the streets, was mocked by the people, and his right hand, severed from the trunk, was presented at the doors of houses with the requirement, ``Give to the insatiable!''

There can be no reasonable doubt that the assassination of Rufinus was instigated by Stilicho, as some of our authorities expressly tell us.\texthebrew{⁠}\footnote{Zosimus (source certainly Eunapius),

\emph{ib.} 3 .

Philostorgius, XI.3 . It is remarkable that Claudian does not mention Gaïnas, whose part in the affair we find in Zosimus.\texthebrew{—} On the confiscation of the large property of Rufinus see

\emph{C. Th.} IX.42.14 ; Symmachus, \emph{Epp.} VI.14.

} The details may have been arranged between him and Gaïnas, and he appears not to have concerned himself to conceal his complicity. The scene of the murder is described by a gifted but rhetorical poet, Claudius Claudianus, who now began his career as a trumpeter of Stilicho's praises by his poem

\emph{Against Rufinus} .\texthebrew{⁠}\footnote{See

Claudian, \emph{Carm. min.} XLI. vv. 13\texthebrew{‑}16 , which seem to imply that he came to Italy in the consul­ship of Probinus (and Olybrius), A.D. 395. Cp. Prosper, \emph{Chron., sub a.}

} He paints Stilicho and Rufinus as two opposing forces, powers of darkness and light: the radiant Apollo, deliverer of mankind, and the terrible Pytho, the scourge of the world. What we should call the crime of Stilicho is to him a glorious deed, the destruction of a monster, and though he does not say in so many words that his hero planned it, he does not disguise his responsibility. Claudian was a master of violent invective, and his portrait of Rufinus, bad man though he unquestionably was, is no more than a caricature. The poem concludes with a picture of the Prefect in hell before the tribunal of Rhadamanthys, who declares that all the iniquities of the tortured criminals are but a fraction of the sins of the latest comer, who is too foul even for Tartarus, and consigns him to an empty pit outside the confines of Pluto's domain.

p114

It was not only the European parts of the dominion of Arcadius that were ravaged, in this year, by the fire and sword of barbarians. Hordes of trans-Caucasian Huns poured through the Caspian gates, and, rushing southwards through the Armenian highlands and the plains of Mesopotamia, carried desolation into Syria. St. Jerome was in Palestine at this time, and in two of his letters we have the account of an eye-witness . ``As I was searching for an abode worthy of such a lady (Fabiola, his friend), behold, suddenly messengers rush hither and thither, and the whole East trembles with the news, that from the far Maeotis, from the land of the ice-bound Don and the savage Massagetae, where the strong works of Alexander on the Caucasian cliffs keep back the wild nations, swarms of Huns had burst forth, and, flying hither and thither, were scattering slaughter and terror everywhere. The Roman army was at that time absent in consequence of the civil wars in Italy. . . . May Jesus protect the Roman world in future from such beasts! They were everywhere, when they were least expected, and their speed outstripped the rumour of their approach; they spared neither religion nor dignity nor age; they showed no pity to the cry of infancy. Babes, who had not yet begun to live, were forced to die; and, ignorant of the evil that was upon them, as they were held in the hands and threatened by the swords of the enemy, there was a smile upon their lips. There was a consistent and universal report that Jerusalem was the goal of the foes, and that on account of their insatiable lust for gold they were hastening to this city. The walls, neglected by the carelessness of peace, were repaired. Antioch was enduring a blockade. Tyre, fain to break off from the dry land, sought its ancient island. Then we too were constrained to provide ships, to stay on the seashore, to take precautions against the arrival of the enemy, and, though the winds were wild, to fear a shipwreck less than the barbarians \texthebrew{—} making provision not for our own

p115
safety so much as for the chastity of our virgins.''\texthebrew{⁠}\footnote{\emph{Epp.} LXXVII.8. These Huns were doubtless the Sabeiroi, whom we shall meet again. Their seats were between the Caspian and the Euxine. See below,

p434, note . º

} In another letter, speaking of these ``wolves of the north,'' he says: ``How many monasteries were captured? the waters of how many rivers were stained with human gore? Antioch was besieged and the other cities, past which the Halys, the Cydnus, the Orontes, the Euphrates flow. Herds of captives were dragged away; Arabia, Phoenicia, Palestine, Egypt were led captive by fear.''\footnote{\emph{Epp.} LX.16. Jerome is dwelling on the miseries of human society (temporum nostrorum ruinas), which he also illustrates by the ravages of Alaric in Europe, and by the fates of Rufinus, Abundantius, and Timasius. The letter was written in 396.

}

After the death of Rufinus, the weak Emperor Arcadius passed under the influence of the eunuch Eutropius, who in unscrupulous greed of money resembled Rufinus and many other officials before and after, and, like Rufinus, has been painted blacker than he really was. All the evil things that were said of Rufinus were said of Eutropius; but in reading of the enormities of the latter we must make great allowance for the general prejudice existing against a person with his physical disqualifications.

The ambitious eunuch naturally looked on the Praetorian Prefects of the East, the most power ­ful men in the administration next to the Emperor, with jealousy and suspicion. To his influence we are probably justified in ascribing an innovation which was made by Arcadius. The administration of the cursus publicus , or office of the postmaster-general , and the supervision of the factories of arms, were transferred from the Praetorian Prefect to the Master of Offices.\footnote{John Lydus, \emph{De mag.} III.40 .

}

It has been supposed that a more drastic arrangement was made for the purpose of curtailing the far-reaching authority of the Praetorian Prefect of the East. There is evidence which has been interpreted to mean that during the three and a half years which coincided with the régime of Eutropius there were two Prefects holding office at the same time and dividing the spheres of administration between them. If this was so, it would have been a unique experiment, never essayed before or

p116
since. But the evidence is not cogent, and it is very difficult to believe that some of the contemporary writers would not have left a definite record of such a revolutionary change.\footnote{The evidence consists in the circumstance that in the Theodosian Code we find laws addressed to Caesarius Pr. Pr. from Nov. 30, 395, throughout 396 and 397, and on July 26, 398, and at the same time four laws addressed to Eutychianus Pr. Pr. in 396, six laws addressed to him in 397, and six in the first half of 398. Hence Seeck has argued that Caesarius and Eutychianus were colleagues in the Prefecture of the East during these years. The natural explanation is that the dates of some of the constitutions are wrong (\emph{viz.} six Eutychianus dates in 396 and 397, and one Caesarius date in 398) and that while Caesarius succeeded Rufinus Dec. 395, Eutychianus succeeded Caesarius between July 13

(\emph{C. Th.} VIII.15.8)

and Sept. 4

(\emph{ib.} VI.3.4)

397. Eutychianus held office till the fall of Eutropius in August 399. Seeck thinks that this series of errors is improbable (\emph{Gesch. des Untergangs}, V.551), but errors of date are very common, and in these years alone we find Caesarius addressed as Pr. Pr. in June 395; Aurelian in Oct. 396 ( IV.2.1

and

V.1.5 ) and Jan. 399; and Eutychianus in Dec. 399 (when Aurelian was Prefect). On the general question see Mommsen, \emph{Hist. Schr.} III. p290; Seeck, in \emph{Philologus}, 52, p449. The list of the laws which are concerned will be found in Mommsen's ed. of \emph{C. Th.} I. pp. clxxv\texthebrew{‑}vi.

}

The Empire was now falling into a jeopardy, by which it had been threatened from the outset, and which it had ever been trying to avoid. There were indeed two dangers which had constantly impended from its inauguration by Augustus to its renovation by Diocletian. The one was a cabinet of imperial freedmen, the other was a military despotism. The former called forth, and was averted by, the creation of a civil service system, to which Hadrian perhaps made the most important contributions, and which was elaborated by Diocletian, who at the same time met the other danger by separating the military and civil administrations. But both dangers revived in a new form. The danger from the army became danger from the Germans, who preponderated in it; and the institution of court ceremonial tended to create a cabinet of chamberlains and imperial dependents. This oriental ceremonial, so notorious a feature of ``Byzantinism,'' meant difficulty of access to the Emperor, who, living in the retirement of his palace, was tempted to trust less to his eyes than his ears, and saw too little of public affairs. Diocletian himself appreciated this disadvantage, and remarked that the sovran, shut up in his palace, cannot know the truth, but must rely on what his attendants and officers tell him. Autocracy, by its very nature, tends in this direction; for it generally means a dynasty, and a dynasty implies that there must sooner or later come to the throne weak men, inexperienced in public affairs,

p117
reared up in an atmosphere of flattery and illusion, at the mercy of intriguing chamberlains and eunuchs. In such conditions aulic cabals and chamber cabinets are a natural growth.

The greatest blot on the ministry of Eutropius (for, as he was the most trusted adviser of the Emperor, we may use the word ministry), was the sale of offices, of which the poet Claudian gives a vivid and exaggerated account.\texthebrew{⁠}\footnote{\emph{In Eutrop.} I.198

institor imperii, caupo famosus honorum, etc.

} This was a blot, however, that stained other power ­ful men in those days as well as Eutropius, and we must view it rather as a feature of the times than as a peculiar enormity. Of course, the eunuch's spies were ubiquitous; of course, informers of all sorts were encouraged and rewarded. All the usual stratagems for grasping and plundering were put into practice. The strong measures that a determined minister was ready to take for the mere sake of vengeance, may be exemplified by the treatment which the whole Lycian province received at the hands of Rufinus. On account of a single individual, Tatian, who had offended that minister, all the provincials were excluded from the public offices.\texthebrew{⁠}\footnote{Probably in A.D. 394. Tatian had been Praetorian Prefect of the East 389\texthebrew{‑}393 and Consul in 391. His son Proclus was Prefect of the City 389\texthebrew{‑}392. Both incurred the jealousy of Rufinus, who procured their arrest and condemnation. Proclus was beheaded, Tatian exiled to Lycia. Cp. Asterius, \emph{op. cit.} \emph{ib.}

} After the death of Rufinus, the Lycians were relieved from these disabilities; but the fact that the edict of repeal expressly enjoins ``that no one henceforward venture to wound a Lycian citizen with a name of scorn'' shows what a serious misfortune their degradation was.\footnote{\emph{C. Th.} IX.38.9 ;

Claudian, \emph{In Ruf.} I.232 .

}

The eunuch won considerable odium in the first year of his power ( A.D. 396) by bringing about the fall of two soldiers of distinction, whose wealth he coveted \texthebrew{—} Abundantius, to whose patronage he owed his rise in the world, and Timasius, who had been the commander-general in the East. The arts by which Timasius was ruined may illustrate the character of the intrigues that were spun at the Byzantine court.\footnote{Zosimus, V.8.

}

Timasius had brought with him from Sardis a Syrian sausage-seller , named Bargus, who, with native address, had insinuated himself into his good graces, and obtained a subordinate command in the army. The prying omniscience of Eutropius discovered that, years before, this same Bargus had been forbidden

p118
to enter Constantinople for some misdemeanour, and by means of this knowledge he gained an ascendancy over the Syrian, and compelled him to accuse his benefactor Timasius of a treasonable conspiracy and to support the charge by forgeries. The accused was tried,\texthebrew{⁠}\footnote{The general feeling in favour of Timasius, a man of the highest character, was so great that the Emperor gave up his first intention of presiding at the trial. The letter of Jerome (LX. \texthebrew{—} quoted above,

p114 ), which was written in 396, proves that Abundantius and Timasius were exiled in that year. Abundantius had been consul and mag. utr. mil., Timasius consul and mag. mil. Their fates are referred to by Asterius, \emph{ib.} (cp. M. Bauer, \emph{Asterios Bischof von Amaseia}, 1911, p12 \emph{sqq.}). See also Sozomen, VIII.17.

} condemned, and banished to the Libyan oasis, a punishment equivalent to death; he was never heard of more. Eutropius, foreseeing that the continued existence of Bargus might at some time compromise himself, suborned his wife to lodge very serious charges against her husband, in consequence of which he was put to death.

It seems probable that a serious plot was formed in the year 397, aiming at the overthrow of Eutropius. Though this is not stated by any writer, it seems a legitimate inference from a law\texthebrew{⁠}\footnote{Sept. 4;

\emph{C. Th.} IX.14.3 .

} which was passed in the autumn of that year, assessing the penalty of death to any one who had conspired ``with soldiers or private persons, including barbarians,'' against the lives of illustres who belong to our consistory or assist at our counsels,'' or other senators, such a conspiracy being considered equivalent to treason. Intent was to be regarded as equivalent to crime, and not only did the person concerned incur capital punishment, but his descendants were visited with disfranchisement. It is generally recognised that this law was an express protection for chamberlains; but we must suppose it to have been suggested by some actual conspiracy, of which Eutropius had discovered the threads. The mention of \emph{soldiers} and \emph{barbarians} points to a particular danger, and we may suspect that Gaïnas, who afterwards brought about the fall of Eutropius, had some connexion with it.

During this year, Stilicho was engaged in establishing his power in Italy and probably in courting a popularity which he had so far done little to deserve. He found time to pay a hurried visit\texthebrew{⁠}\footnote{Cp. Claudian, \emph{De cons. Stil.} I.218 \emph{sqq.} Perhaps it was at this time that the military administration on the Rhine frontier was reorganised by the institution of two new high commands, that of the dux Mogontiacensis

(\emph{Not., Occ.}, XLI)

and that of the comes Argentoratensis (\emph{ib.} XXVII), who had their seats at Mainz and Strassburg respectively. Cp. Seeck, art. \emph{Comites}, in P.\texthebrew{‑}W.

} to the Rhine provinces, to conciliate or pacify the federate

p119
Franks and other German peoples on the frontier, and perhaps to collect recruits for the army. We may conjecture that he also made arrangements for the return of his own family to Italy. He had not abandoned his designs on Eastern Illyricum, but he was anxious to have it understood that he aimed at fraternal concord between the courts of Milan and Byzantium and that the interests of Arcadius were no less dear to him than those of Honorius. The poet Claudian, who filled the rôle of an unofficial poet-laureate to Honorius, was really retained by Stilicho who patronised and paid him. His political poems are extravagant eulogies of the power ­ful general, and in some cases we may be sure that his arguments were directly inspired by his patron. In the panegyric for the Third Consulate of Honorius ( A.D. 396) which, composed soon after the death of Rufinus, suggests a spirit of concord between East and West, the writer calls upon Stilicho to protect the two brethren:

Such lines as this were written to put a certain significance on Stilicho's policy.

For Stilicho was preparing to intervene again in the affairs of the East. We must return here to the movements of Alaric who, when the Imperial armies retreated from Thessaly without striking a blow, had Greece at his mercy. Gerontius, the commander of the garrison at Thermopylae, offered no resistance to his passage; Antiochus, the pro-consul of Achaia, was helpless, and the Goths entered Boeotia, where Thebes alone escaped their devastation.\texthebrew{⁠}\footnote{For the invasion, besides Zosimus, see Socrates VII.10. It is noticed also by Eunapius, \emph{Vita Maximi} (I. p52), and \emph{Vita Prisci} (I. p67).

} They occupied Piraeus but Athens itself was spared, and Alaric was entertained as a guest in the city of Athene.\texthebrew{⁠}\footnote{The walls of Athens had been restored in the reign of Valerian

(Zosimus, I.29) , and Alaric was amenable to terms. The legend was that he saw Athene Promachus standing on the walls, and Achilles in front of them (\emph{ib.} V.6). Philostorgius says that Alaric ``took Athens''

(XII.2)

but he meant Piraeus. The mischief wrought by the Goths in Greece has often been exaggerated (see Gregorovius, \emph{Gesch. der Stadt Athen}, I.35 \emph{sqq.}; Bury, App. 13 to Gibbon, vol. III).

} But the great temple of the mystic goddess, Demeter and Persephone, at Eleusis was plundered by the barbarians; Megara, the next place on their southward route, fell; then Corinth, Argos, and Sparta. It is possible that Alaric entertained

p120
the design of settling his people permanently in the Peloponnesus.\texthebrew{⁠}\footnote{So Schmidt, \emph{ib.} 197.

} However this may be, he remained there for more than a year, and the government of Arcadius took no steps to dislodge him or arrange a settlement.

Then in the spring of A.D. 397,\texthebrew{⁠}\footnote{Since Koch's article in \emph{Rh. Museum}, XLIV. (1889), it has generally been recognised that Stilicho's second expedition to Greece must be placed in 397 (not 396). See Birt, \emph{Praef.} to Claudian, p. XXXI; Gibbon, \emph{Decline}, III. editor's App. 12. The spring of the year must be inferred from

Claudian, \emph{De c. Stil.} I.174 \emph{sqq.} That Elis was the scene of operations is proved by Pholoe in

Zos. V.7.2 , and more than one reference to the Alpheus in Claudian. The second Book \emph{In Rufinum} was not published till after this campaign (see

\emph{Praef.} 9 \emph{sqq.} ).

} Stilicho sailed across from Italy, and landing at Corinth marched to Elis to give the general's poet a pretext for singing of the slaughter of skin-clad warriors ( metitur pellita iuventus ).\texthebrew{⁠}\footnote{\emph{De IV. cons. Hon.} 466 . Cp. \emph{De cons. Stil.} 186

Alpheus Geticis angustus acervis.

} But the outcome was that the Gothic enemy was spared in Elis much as he had been spared in Thessaly. The Eastern government seems to have again intervened with success.\texthebrew{⁠}\footnote{See

Claudian, \emph{B. Got.} 517

sub nomine legum proditio regnique favor texisset Eoi.

} But what happened is unknown, except that Stilicho made some agreement with Alaric,\texthebrew{⁠}\footnote{Claudian \emph{ib.} 496

seems to imply that Alaric undertook not to cross the frontiers of the territory of Honorius.

} and Alaric withdrew to Epirus, where he appears to have come to terms with Arcadius and perhaps to have received the title he coveted of Master of Soldiers in Illyricum.\footnote{Cp. \emph{ib.} 535\texthebrew{‑}539

and

\emph{In Eutrop.} II.216 . Was this a breach of the agreement with Stilicho (cp. foedera rumpit,

\emph{ib.} 213 )? \texthebrew{—} It may have been during this absence of Stilicho that Serena embellished with marble the tomb of St. Nazarius at Milan as a vow for his safe return, CIL V.6250 , unless it were rather in the Raetian campaign of 401\texthebrew{‑}402.

}

That Stilicho had set out with the purpose of settling the question of Illyricum cannot be seriously doubted. That he withdrew for the second time without accomplishing his purpose was probably due to the news of a dangerous revolt in Africa to which the government of Arcadius was accessory. We can easily understand the indignation felt at Constantinople when it was known that Stilicho had landed in Greece with an army. It was natural that the strongest protest should be made, and Eutropius persuaded the Emperor and the Senate to declare him a public enemy.\footnote{Zosimus, V.11.

}

Of this futile expedition, Claudian has given a highly misleading

p121
account in his panegyric in honour of the Fourth Consulate of Honorius ( A.D. 398), which no allowance for conventional exaggeration can excuse. He overwhelms the boy of fourteen with the most extravagant adulations, pretending that he is greater \texthebrew{—} vicariously indeed, through the deeds of his general \texthebrew{—} than his father and grandfather. We can hardly feel able to accord the poet much credit when he declares that the western provinces are not oppressed by heavy taxes nor the treasury replenished by extortion.\footnote{496 \emph{sqq.}

Claudian is at his finest in his eulogies of Theodosius \emph{avus}, the hero of Africa and Britain, and Theodosius \emph{pater}, the Great.

}

Eighteen years before an attempt had been made by the Moor Firmus to create a kingdom for himself in the African provinces ( A.D. 379), and had been quelled by the armies of Theodosius, who had received valuable aid from Gildo, the brother and enemy of Firmus. Gildo was duly rewarded. He was finally appointed Count of Africa with the exceptional title of Master of Soldiers, and his daughter Salvina was united in marriage to a nephew of the Empress Aelia Flaccilla.\texthebrew{⁠}\footnote{Nebridius. Salvina was afterwards a friend of John Chrysostom.

} But the faith of the Moors was as the faith of the Carthaginians. Gildo refused to send troops to Theodosius in his expedition against Eugenius, and after the Emperor's death he prepared to assume a more decided attitude of independence and engaged many African tribes to support him in a revolt. The strained relations between the two Imperial courts suggested to him that the rebellion might assume the form of a transference of Africa from the sovranty of Honorius to that of Arcadius; and he entered into communication with Constantinople, where his overtures were welcomed. A transference of the diocese of Africa to Arcadius seemed quite an appropriate answer to the proposal of transferring the Prefecture of Illyricum to Honorius. But the Eastern government rendered no active assistance to the rebel.\footnote{Zosimus, \emph{ib.} It appears that embassies on the subject passed between Italy and Constantinople (Symmachus, \emph{Epp.} IV.5; Claudian, \emph{B. Gild.} 236 \emph{sqq.} ,

279 ,

\emph{De cons. Stil.} I.295 ;

Orosius, VII.36 ), and that Arcadius went so far as to issue edicts mena­cing any one who should attack Gildo, see

Claudian, \emph{De cons. Stil.} I.275 \emph{sqq.} \texthebrew{—}hoc Africa saevis

cinxerat auxiliis, hoc coniuratus alebat

insidiis Oriens. illinc edicta meabant

corruptura duces.

}

p122
For Rome and the Italians a revolt in Africa was more serious than rebellions elsewhere, since the African provinces were their granary. In the summer of A.D. 397 Gildo did not allow corn º ships to sail to the Tiber; this was the declaration of war. The prompt and efficient action of Stilicho prevented a calamity; corn supplies were obtained from Gaul and Spain sufficient to feed Rome during the winter months. Preparations were made to suppress Gildo, and Stilicho sought to ingratiate himself with the Senate by reverting to the ancient usage of obtaining its formal authority.\texthebrew{⁠}\footnote{Claudian, \emph{De cons. Stil.} I.326 \emph{sqq.} \texthebrew{—}non ante fretis exercitus adstitit ultor

ordine quam prisco censeret bella senatus

neglectum Stilicho per tot iam saecula morem

rettulit, etc.

} The Senate declared Gildo a public enemy, and during the winter a fleet of transports was collected at Pisa. In the early spring an army of perhaps 10,000 embarked.\texthebrew{⁠}\footnote{So Seeck (\emph{Forsch. zur d. Geschichte}, 24, 175 \emph{sqq.}), who identifies the troops (chiefly auxilia palatina) named by Claudian,

\emph{B. Gild.} 418\texthebrew{‑}423 . Orosius,

VII.36 , says 5000 (ut aiunt).

} Stilicho remained in Italy, and the command was entrusted to Mascezel, a brother of Gildo who had come to the court of Honorius to betray Gildo as Gildo had betrayed Firmus. The war was decided, the rebel subdued, almost without bloodshed, in the Byzacene province on the little river Ardalio between Tebessa and Haïdra. The forces of Gildo are said to have been 70,000 strong, but they offered no resistance. We may suspect that some of his Moorish allies had been corrupted by Mascezel, but Gildo himself was probably an unpopular leader. He tried to escape by ship, but was driven ashore again at Thabraca and put to death.\footnote{Claudian, \emph{De cons. Stil.}

I.359 ,

II.258 ;

\emph{In Eutrop.} I.410 ;

\emph{De VI. cons. Hon.} 381 . The date was July 31, \emph{Fasti Vind. pr., sub a.} 398 (\emph{Chron. min.} I. p398). According to Zosimus

(V.11)

Gildo took his own life.

}

Returning to Italy, Mascezel was welcomed as a victor, and might reasonably hope for promotion to some high post. But his swift and complete success was not pleasing to Stilicho, who desired to appropriate the whole credit for the deliverance of Italy from a grave danger; perhaps he saw in Mascezel a possible rival. Whether by accident or design, the Moor was removed from his path. The only writer who distinctly records the event, states that while he was crossing a bridge he was thrown into a river by Stilicho's bodyguards and that Stilicho gave the sign for the act.\texthebrew{⁠}\footnote{Zosimus, V.11.

Orosius (\emph{ib.}), who represents the Moor's death as a punishment for profaning a church, does not tell how it occurred; but occisus est means a violent end.

} The evidence is not good enough to justify us in

p123
bringing in a verdict of murder against Stilicho; Mascezel may have been accidentally drowned and the story of foul play may have been circulated by Stilicho's enemies. But if the ruler of Italy was innocent, he assuredly did not regret the capable executor of his plans. The order seems to have gone out that the commander of the expedition against Gildo was to have no share in the glory,\texthebrew{⁠}\footnote{Cp. CIL VI.1730 (see below,

p125 ). The question is discussed by Crees, \emph{Claudian} 102 \emph{sqq.} The inscription found in the Roman Forum,

armipotens Libycum defendit Honorius orbem,

may refer to the Gildonic War, CIL VI.31256 .

} and

the incomplete poem of Claudian on the \emph{Gildonic War}

tells the same tale.

This poem, which will serve as an example of Claudian's art, begins with an announcement of the victory and was probably composed when the first news of the success arrived in Italy. Redditus imperiis Auster , ``the South has been restored to our Empire; the twin sphere, Europe and Libya, are reunited; and the concord of the brethren is again complete.'' Iam domitus Gildo , the tyrant as already been vanquished, and we can hardly believe that this has been accomplished so quickly.

Having announced the glad tidings, Claudian goes back to the autumn and imagines Rome, the goddess of the city, in fear of famine and disaster, presenting herself in pitiable guise before the throne of Jupiter and supplicating him to save her from hunger. Are the labours and triumphs of her glorious history to be all in vain? Is the amplitude of her Empire to be her doom? Ipsa nocet moles . ``I am excluded from my granaries, Libya and Egypt; I am abandoned in my old age.''

The supplications of Rome are reinforced by the sudden appearance of Africa, who burst into the divine assembly with torn raiment, and in wild words demands that Neptune should submerge her continent rather than it should have to submit to the pollution of Gildo's rule.

Jupiter dismisses the suppliants, assuring them that ``Honorius will lay low the common enemy,'' and he sends Theodosius the

p124
Great and his father, who are both deities in Olympus, to appear to the two reigning Emperors in the night. Arcadius is reproached by his father for the estrangement from his brother, for his suspicions of Stilicho, for entertaining the proposals of Gildo; and he promises to do nothing to aid Gildo. Honorius is stimulated by his grandfather to rise without delay and smite the rebel. He summons Stilicho and proposes to lead an expedition himself. Stilicho persuades him that it would be unsuitable to his dignity to take the field against such a foe, and suggests that the enterprise should be committed to Mascezel. This is the only passage in which Mascezel is mentioned, and Claudian does not bestow any praise on him further than the admission that he does not resemble his brother in character ( sed non et moribus isdem ), but dwells on the wrongs he had suffered, and argues that to be crushed by his injured brother, the suppliant of the Emperor, will be the heaviest blow that could be inflicted on the rebel.

The military preparations are then described, and an inspiriting address to the troops, about to embark, is put into the mouth of Honorius, who tells them that the fate of Rome depends on their valour:

The fleet sails and safely reaches the African ports, and the first canto of the poem ends.\footnote{In the MSS. it is described as Liber primus.

}

It is all we have: a second canto was never written. Claudian evidently intended to sing the whole story of the campaign as soon as the story was known. The overthrow of ``the third tyrant,'' whom he represents as the successor of Maximus and Eugenius, deserved an exhaustive song of triumph. But it would have surpassed even the skill of Claudian to have told the tale without giving a meed of praise to the commander who carried the enterprise through to its victorious end. We need have little hesitation in believing that the motive which hindered the poet from completing the \emph{Gildonic War} was the knowledge that to celebrate the achievements of Mascezel would be no service to his patron.\footnote{The complications which resulted in Africa from the despotism of Gildo, and the efforts to right wrongs and restore property, lasted for many years. The large property which Gildo had amassed required a special

(p125)
official to administrate it, entitled comes Gildoniaci patrimonii. See

\emph{C. Th.} VII.8.7 , and

\emph{Notit. Occ.} XI .

}

p125
While the issue of the war was still uncertain, in the spring of A.D. 398,\texthebrew{⁠}\footnote{Claudian, \emph{De cons. Stil.} I.1\texthebrew{‑}5 .

} Stilicho's position as master of the west was strengthened by the marriage of his daughter Maria with the youthful Emperor. Claudian wrote an epithalamium for the occasion, duly extolling anew the virtues of his incomparable patron. We may perhaps wonder that, secured by this new bond with the Imperial house, and his prestige enhanced by the suppression of Gildo,\texthebrew{⁠}\footnote{An inscription in honour of Stilicho on a marble base, found at Rome (CIL VI.1730 ), celebrates the ``deliverance'' of Africa:

Flavio Stilichoni inlustrissimo viro, magistro equitum peditumque, comiti domesticorum, tribuno praetoriano et ab ineunte aetate per gradus clarissimae militiae ad columen gloriae sempiternae et regiae adfinitatis evecto, progenero divi Theodosi, comiti divi Theodosi Augusti in omnibus bellis adque º victoriis et ab eo in adfinitatem regiam cooptato itemque socero d. n. Honori Augusti Africa consiliis et provisione et liberata.

There is also an inscription to the two Emperors, belonging to some memorial erected by the Senate and Roman people, vindicata rebellione et Africae restitutione laetus; CIL VI.31256 . This is the titulus perennis of

Claudian, \emph{De VI. cons. Hon.} 372 . Cp. also CIL IX.4051 .

} Stilicho did not now make some attempt to carry out his project of annexing the Prefecture of Illyricum. The truth is that he had not abandoned it, but he was waiting for a favourable opportunity of intervention in the affairs of the east. It seems safe to infer his attitude from the drift of Claudian's poems, for Claudian, if he did not receive express instructions, had sufficient penetration to divine the note which Stilicho would have wished him to strike. In the \emph{Gildonic War} he had announced the restoration of concord between east and west: concordia fratrum plena redit ; it was the right thing to say at the moment, but the strain in the relations between the two courts had only relaxed a little. The discord broke out again, with more fury than ever, in

the two poems

in which he overwhelmed Eutropius with rhetoric no less savage than his fulminations against Rufinus four years before. The first was written at the beginning of A.D. 399, protesting against the disgrace of the Empire by the elevation of Eutropius to the consulate, the second in the summer, after the eunuch's fall. The significant point is that in both poems the intervention of Stilicho in eastern affairs is proposed.\texthebrew{⁠}\footnote{\emph{In Eutrop.}

I.500 \emph{sqq.} ,

II.591 \emph{sqq.}

} Stilicho did not overtly intervene; but it seems probable that he had an understanding with Gaïnas, the German commander in the east, who had been his instrument

p126
in the assassination of Rufinus. It is a suggestive fact that in describing the drama which was enacted in the east Claudian brings the minor characters on the stage but does not even pronounce the name of Gaïnas, who was the principal actor, or betray that he was aware of his existence. We must now pass to the east and follow the events of that drama.

In these years, in which barbarians were actively harrying the provinces of the Illyrian peninsula and the eastern provinces of Asia Minor, concord and mutual assistance between east and west were urgently needed. Unfortunately, the reins of government were in the hands of men who for different reasons were unpopular and in all their political actions were influenced chiefly by the consideration of their own fortunes. The position of Eutropius was insecure, because he was a eunuch; that of Stilicho, because he was a German. So far as the relation between the two governments was concerned the situation had been eased for a time after the fall of Rufinus, and it was doubtless with the consent and perhaps at the invitation of Eutropius that Stilicho had sailed to Greece in A.D. 397. For the eastern armies were not strong enough to contend at the same time against Alaric and against the Huns who were devastating in Asia. The generals who were sent to expel the invaders from Cappadocia and the Pontic provinces seem to have been incompetent, and Eutropius decided to take over the supreme command himself. It was probably in A.D. 398 that he conducted a campaign which was attended with success. The barbarians were driven back to the Caucasus and the eunuch returned triumphant to Constantinople.\texthebrew{⁠}\footnote{Claudian, \emph{In Eutr.} I.234\texthebrew{‑}286 . We can read clearly through the jeers and sarcasms of the poet that the martial adventure of Eutropius was a distinct success. It is not proved that he assumed the office and title of a Master of Soldiers, as Birt thinks (\emph{Preface}, xxxiv); but, however this may be, Birt is certainly wrong in his view that Eutropius ever filled the office of Praetorian Prefect. The expressions of Claudian which he cites (\emph{ib.} xxx) are far from proving it.

} His victory secured him some popularity for the moment, and he was designated consul for the following year.

The brief understanding between the courts of Milan and

p127
Byzantium had been broken as we saw by the attitude of the eastern government during the revolt of Gildo. There was an open breach. When the news came that Eutropius was nominated consul for A.D. 399, the Roman feelings of the Italians were deeply scandalised. A eunuch for a consul \texthebrew{—} it was an unheard-of , an intolerable violation of the tradition of the Roman Fasti.

wrote Claudian in the poem in which, at the beginning of the year, he castigated the minister of Arcadius.\texthebrew{⁠}\footnote{\emph{In Eutropium liber I}

(cp. Birt, \emph{ib.} xl).

} The west refused to recognise this monstrous consul ­ship.\texthebrew{⁠}\footnote{After this year, the practice was introduced of publishing the eastern and western consuls successively, in each part of the Empire. Simultaneous publication only occurred when the consuls had been fixed before Jan. 1 by special arrangement, as when two Emperors assumed the office together. Mommsen, \emph{Hist. Schr.} III.367.

} It was perhaps hardly less unpopular in the east.

The Grand Chamberlain, confidently secure through his possession of the Emperor's ear, had overshot the mark. His position was now threatened from two quarters. Gaïnas, the German officer who under the direction of Stilicho had led the eastern army back to Constantinople, had risen to the office of a Master of Soldiers.\texthebrew{⁠}\footnote{Socrates, VI.6 \textgreek{στρατηλάτης \texthebrew{Ῥ}ωμαίων \texthebrew{ἱ}ππικ\texthebrew{ῆ}ς τε κα\texthebrew{ὶ} πεζικ\texthebrew{ῆ}ς}, \emph{i.e.} Mag. mil. in praesenti;

Philostorgius, XI.8

\textgreek{\texthebrew{ὁ} στρατηγός}, cp. Sozomen, VIII.4 \emph{ad init.} Cp. Tillemont, \emph{Histoire}, V. p783.

} It is probable that he maintained communications with Stilicho, and his first object was to compass the downfall of Eutropius.

Less dangerous but not less hostile was the Roman party, which was equally opposed to the bedchamber administration of Eutropius and to the growth of German power. It consisted of senators and ministers attached to Roman traditions, who were scandalised by the nomination of the eunuch to the consul ­ship in A.D. 399 and alarmed by the fact that some of the highest military commands in the Empire were held by Germans. The leader of the party was Aurelian, son of Taurus (formerly a Praetorian Prefect of Italy), who had himself filled the office of Prefect of the City.

Gaïnas had some supporters among the Romans. The most power ­ful of his friends was an enigmatical figure, whose real name is unknown but who seems to have been a brother of

p128
Aurelian. Of this dark person, who played a leading part in the events of these years, we derive all we know from a historical sketch which its author Synesius of Cyrene cast into the form of an allegory and entitled

\emph{Concerning Providence or the Egyptians} . This distinguished man of letters, who was at this time a Platonist \texthebrew{—} some years later he was to embrace Christianity and accept a bishopric \texthebrew{—} was on terms of intimacy with Aurelian and was at Constantinople at this time.\texthebrew{⁠}\footnote{He was there for three years (A.D. 399\texthebrew{‑}402):

\emph{Hymns}, III.430\texthebrew{‑}434 ; he went home during the great earthquake of 402.

\emph{Epp.} 61 , p1404. Cp. Seeck, in \emph{Philologus} 52, p458.

Thayer's Note: A good readable introduction to the Cyrenean bishop, with a selection of his writings, may be found at

Livius.Org , and another excellent synopsis of Synesius' life and works is given by the article

Synesius of Cyrene

in the Catholic Encyclopedia; most, maybe all, of his works are online in English translation at

Livius .

} The argument is the contest for the kingship of Egypt between the sons of Taurus, Osiris and Typhos. Osiris embodies all that is best in human nature. Typhos is a monster, perverse, gross, and ignorant. Osiris is Aurelian; Typhos cannot be identified,\texthebrew{⁠}\footnote{On the interpretation of the allegory see Sievers, \emph{Studien}, 387 \emph{sqq.}, Seeck, 442 \emph{sqq.}, and \emph{Untergang}, V.314 \emph{sqq.}, Mommsen, \emph{Hist. Schr.} III.292 \emph{sqq.} Thebes is Constantinople, and the high priest (p1268) is Arcadius. Seeck has endeavoured to prove that Typhos is Caesarius, who succeeded Rufinus as Pr. Prefect of the East in 395 and held that office till 397, in which year he was consul (see laws in \emph{C. Th.} ed. Mommsen, I. p. clxxv,

Philostorgius XI.5 ). Mommsen has given cogent reasons for rejecting this view. If Typhos is Caesarius, it ought to have been mentioned that he had already held the office of king, but Synesius says (p1217) that he was \textgreek{ταμίας χρημάτων}, which would naturally mean comes rei privatae (Seeck interprets it as Praet. Pr., but Synesius describes it as a \textgreek{διάθεσις \texthebrew{ἐ}λάττων}), and then apparently a governor of part of the Empire (perhaps a vicarius).

} and we must call him by his allegorical name; the kingship of Egypt means the Praetorian Prefecture of the east.

In the race for political power Typhos allied himself with the German party, who welcomed him as a Roman of good family and position. Synesius dwells much on his profligacy, and on the frivolous habits of his wife, an ambitious and fashionable lady. She was her own tirewoman, a reproach which seems to mean that she was inordinately attentive to the details of her toilet.\texthebrew{⁠}\footnote{P. 1240 \textgreek{\texthebrew{ἑ}αυτ\texthebrew{ῆ}ς κομμώτρια, θεάτρου κα\texthebrew{ὶ ἀ}γορ\texthebrew{ᾶ}ς \texthebrew{ἄ}πληστος κτλ.}} She liked public admiration and constantly showed herself at the theatre and in the streets. Her love of notoriety did not permit her to be fastidious in her choice of society, she liked to have her salon filled, and her doors were not closed to professional courtesans. Synesius contrasts her with the modest wife of Aurelian, who never left her house, and asserts that the chief virtue of a woman is that neither her body nor her name should ever cross the threshold. This is a mere rhetorical flourish; the writer's friend and teacher, Hypatia the philosopher,

p129
whom he venerated, certainly did not stay at home. He was probably thinking of the piece of advice to women which Thucydides placed in the mouth of Pericles.

The struggle against the German power in the east began in the spring of A.D. 399. It was brought on by a movement on the part of Ostrogoths in Phrygia, but we have no distinct evidence to show that it was instigated by Gaïnas.\texthebrew{⁠}\footnote{Tribigild had visited the capital at the beginning of 399 to pay his respects to Eutropius the new consul, who on this occasion slighted him. It is possible that he arranged the plan of campaign with Gaïnas before he returned to Phrygia. But their complicity may have begun only after the fall of Eutropius.

} These Ostrogoths had been established as colons\texthebrew{⁠}\footnote{Claudian, \emph{In Eutrop.} II.153 , Ostrogothis colitur mixtisque Gruthungis Phryx ager. Gruthungi is only another name for Ostrogoths.

} by Theodosius the Great in fertile regions of that province (in A.D. 386), and contributed a squadron of cavalry to the Roman army. The commander, Tribigild, bore Eutropius a personal grudge, and he excited his Ostrogoths to revolt. The rebellion broke out just as Arcadius and his court were preparing to start for Ancyra, whither he was fond of resorting in summer to enjoy its pleasant and salubrious climate.

The barbarians were recruited by run away slaves and spread destruction throughout Galatia, Pisidia, and Bithynia. Two generals, Gaïnas and Leo, a friend of Eutropius \texthebrew{—} a good-humoured , corpulent man who was nicknamed Ajax \texthebrew{—} were sent to quell the rising.

It was at this time that Synesius, the philosopher of Cyrene, who had come to the capital to present a gold crown to Arcadius on behalf of his native city, fulfilled his mission and used the occasion to deliver a remarkable speech ``On the office of King.''\texthebrew{⁠}\footnote{\textgreek{Περ\texthebrew{ὶ} βασιλείας}, \emph{Opera} p1053 \emph{sqq.}

Thayer's Note: Available in English

at Livius.Org .

} It may be regarded as the anti-German manifesto of the party of Aurelian\texthebrew{⁠}\footnote{cp. Sievers \emph{Studien}, p379, and Güldenpenning, \emph{Gesch. d. oström. Reiches}, p106.

} with which Synesius had enthusiastically identified himself. The orator urged the policy of imposing disabilities on the Germans in order to eradicate the German element in the State. The argument depends on the Hellenic but by no means Christian principle that Roman and barbarian are different in kind and therefore their union is unnatural. The soldiers of a state should be its watchdogs, in Plato's phrase, but our armies are full of wolves in the guise of dogs. Our homes are full of German servants. A state cannot wisely give arms to

p130
any who have not been born and reared under its laws; the shepherd cannot expect to tame the cubs of wolves. Our German troops are a stone of Tantalus suspended over our State, and the only salvation is to remove the alien element.\texthebrew{⁠}\footnote{\textgreek{\texthebrew{Ἐ}κκρ\texthebrew{ῖ}ναι δ\texthebrew{ὲ} δε\texthebrew{ῖ} τ\texthebrew{ἀ}λλότριον}, p1089.

} The policy of Theodosius the Great was a mistake. Let the barbarians be sent back to their wilds beyond the Danube, or if they remain be set to till the fields as serfs. It was a speech which if it came to the ears of Gaïnas was not calculated to stimulate his zeal against the Germans he went forth to reduce.

The rebels, seeking to avoid an engagement with Leo's army, turned their steps to Pisidia and thence to Pamphylia, where they met unexpected resistance.\texthebrew{⁠}\footnote{Zosimus, V.16.

} While Gaïnas was inactive and writing in his reports to Constantinople that Tribigild was extremely formidable, Valentine, a landowner of Selge, gathered an armed band of peasants and slaves and laid an ambush near a narrow winding pass in the mountains between Pisidia and Pamphylia. The advancing enemy were surprised by showers of stones from the heights above them, and it was difficult to escape as there was a treacherous marsh all around. The pass was held by a Roman officer, and Tribigild succeeded in bribing him to allow his forces to cross it. But they had no sooner escaped than, shut in between two rivers, the Melas and the Eurymedon, they were attacked by the warlike inhabitants of the district. Leo meanwhile was advancing, and the insurrection might have been crushed if Gaïnas had not secretly reinforced the rebels with detachments from his own army. Then the German troops under his own command attacked and over­powered their Roman fellow-soldiers , and Leo lost his life in attempting to escape.\texthebrew{⁠}\footnote{Claudian, writing to put him in a ridiculous light, pretends that he was killed by fright \texthebrew{—} ualuit pro uulnere terror

(\emph{In Eutr.} II.453) . Leo was doubtless one of the two Masters of Soldiers in praesenti.

} Gaïnas and Tribigild were masters of the situation, but they still pretended to be enemies.

Gaïnas, posing as a loyal general, foiled by the superior power of the Ostrogoths, despatched a message to the Emperor urging him to yield to Tribigild's demand and depose Eutropius from power. Arcadius might not have yielded if a weightier influence had not been brought to bear upon him. The Empress Eudoxia, who had owed her fortune to the eunuch, had become jealous of the boundless power he had secured over have husband's

p131
mind; there was unconcealed antagonism between them; and one day Eudoxia appeared in the Emperor's presence, with her two little daughters,\texthebrew{⁠}\footnote{Flaccilla, born June 17, 397, and Pulcheria, born Jan. 19, 399 (\emph{Chr. Pasch., sub annis}). We never hear of Flaccilla again, she probably died in girlhood. The third child, Arcadia, was born April 3, 400; the youngest daughter, Marina, Feb. 11, 403.

} and made bitter complaint of the Chamberlain's insulting behaviour.

Eutropius realised his extreme peril when he heard of the demand of Gaïnas and he fled for refuge to the sanctuary of St. Sophia.\texthebrew{⁠}\footnote{The fall of Eutropius is recounted by the ecclesiastical historians, and by Zosimus

(V.18) .

} There he might not only trust in the protection of the holy place, but might expect that the Patriarch would stand by him in his extremity when he was deserted by his noonday friends. For it was through him that John Chrysostom, a Syrian priest of Antioch, had been appointed to the see of Constantinople in the preceding year. And the Patriarch's personal interference was actually needed. Arcadius had determined to sacrifice him, and Chrysostom had to stand between the cowering eunuch and those who would have dragged him from the altar. This incident seems to have occurred on a Saturday, and on the morrow, Sunday, there must have been strange excitement in the congregation which assembled to hear the eloquence of the preacher. Hidden under the altar, overwhelmed with fear and shame, lay the old man whose will had been supreme a few days before, and in the pulpit the Patriarch delivered a sermon on the moral of his fall, beginning with the words, ``Vanity of vanities, all is vanity.''\texthebrew{⁠}\footnote{\textgreek{\texthebrew{Ὀ}μιλία ε\texthebrew{ἰ}ς Ε\texthebrew{ὐ}τρόπιον}, \emph{P.G.} 52.391 \emph{sqq.} Asterius refers to the eunuch's fall in his sermon on the Calends (\emph{P.G.} 40.225), delivered Jan. 1, 400 (Bauer, \emph{Asterios}, p21). He mentions the enormous landed property the eunuch had acquired, \textgreek{\texthebrew{ἐ}κτησατο γ\texthebrew{ῆ}ν \texthebrew{ὅ}σην ο\texthebrew{ὐ}δ\texthebrew{ὲ} ε\texthebrew{ἰ}πε\texthebrew{ῖ}ν ε\texthebrew{ὔ}κολον.}} While he mercilessly exposed the levity and irreligion of Eutropius and his circle, he sought at the same time to excite the sympathy of his hearers.

The church was again entered by soldiers, and again Chrysostom interposed. Then Eutropius allowed himself to be removed on condition that his life was spared. He was deprived of his patrician rank, banished to Cyprus, and his property was confiscated. The imperial edict which pronounced this sentence is profuse of the language of obloquy.\texthebrew{⁠}\footnote{\emph{C. Th.} IX.40.17 , addressed to Aurelian Pr. Pr., but wrongly dated.

} The consul ­ship ``befouled and defiled by a filthy monster'' has been ``delivered

p132
from the foul stain of his tenure and from the recollection of his name and the base filth thereof,'' by erasing his name from the Fasti. All statues in bronze or marble, all coloured pictures set up in his honour in public or private places, are to be abolished ``that they may not, as a brand of infamy on our age, pollute the gaze of beholders.''

The fall of Eutropius involved the fall of Eutychian, the Praetorian Prefect of the east, who was presumably one of his creatures. There was a contest between the two brothers, Aurelian and Typhos, for the vacant office, which Synesius in his allegory designated as the kingship of Egypt. But though Gaïnas had succeeded in overthrowing the eunuch, he failed to secure the appointment of Typhos. The post was given to Aurelian, and this was a triumph for the anti-German party.\texthebrew{⁠}\footnote{The last constitution addressed to Eutychian is dated July 25, 399

(\emph{C. Th.} IX.40.18) , the first to Aurelian, Aug. 27

(\emph{ib.} II.8.23) . This gives limits for the fall of Eutropius, which may be placed in August.\texthebrew{—} There are many errors in the dates of the laws in \emph{C. Th.} from 395 to 400. The solution certainly does not lie in Seeck's theory that Caesarius and Eutychian held the Pr. Prefecture conjointly in 396 and 397. The dates of the six laws addressed to Eutychian between Feb. 26, 396 and April 1, 397; as well as that to Caesarius on July 26, 398, are simply false. See above,

p128, n2 . The succession was Caesarius, Nov. 395 to July or August 397; Eutychian, to July or August 399; Aurelian, Aug. 399 to Oct. 6 at least ( \emph{C. Th.} IV.21 ;

V.1.5 ); Typhos (no laws); Aurelian again, 400, perhaps continuing to 402, if we accept Seeck's corrections in \emph{C. Th.} IV.2.1 and

V.1.5

of Arc. A. \emph{v.} (for IIII) et Honor. A. \emph{v.} (for III) conss., \emph{i.e.} 402 (for 396); Eutychian again 403\texthebrew{‑}405. Cp. Seeck and Mommsen, \emph{opp. citt.}

} Aurelian was a man of considerable intellectual attainments; he was surrounded by men of letters such as Synesius, Troilus the poet, and Polyaemon the rhetor. His success was a severe blow to Typhos and his friends, and especially to his wife, who had been eagerly looking forward to the Prefecture for the sake of the social advantage of it. Synesius gives a curious description of the efforts of the profligate to console himself for his disappointment. He constructed a large pond in which he made artificial islands provided with warm baths, and in these retreats he and his friends, male and female, used to indulge in licentious pleasures.\footnote{\emph{Egyptians}, p1245 .

}

But if Aurelian's elevation was a blow to Typhos it was no less a blow to Gaïnas, who now threw off the mask and, openly declaring his true colours, acted no longer as a mediator for Tribigild, but as an adversary bargaining for terms. Tribigild and he met at Thyatira and advanced to the shores of the

p133
Propontis, plundering as they went. Gaïnas demanded and obtained an interview with the Emperor himself at Chalcedon. An agreement was made that he should be confirmed in his post as Master of Soldiers in praesenti ,\texthebrew{⁠}\footnote{Synesius describes the intrigues carried on by the wife of Typhos and the wife of Gaïnas, p1245. The Gothic lady is described as a \textgreek{βάρβαρος γρα\texthebrew{ῦ}ς κα\texthebrew{ὶ ἀ}νόητος.}} that he and Tribigild might cross over into Europe, and that three hostages should be handed over to him, Aurelian, Saturninus, one of Aurelian's chief supporters, and John, the friend (report said the lover) of the Empress. This meant the deposition of Aurelian from the Prefecture and the succession of Typhos. For the moment Gaïnas was master of the government of the east (end of A.D. 399).

The demand for the surrender of Aurelian had been pre-arranged with Typhos,\texthebrew{⁠}\footnote{Sozomen, VIII.4; Tillemont, V.461.

} and the intention seems to have been to put him to death. The Patriarch went over to Chalcedon to intercede for the lives of the three hostages, and Gaïnas contented himself with inflicting the humiliation of a sham execution and banishing them. He then entered Constantinople with his army.\texthebrew{⁠}\footnote{Tribigild disappears entirely from the scene; he perished soon afterwards.

} The rule of Gaïnas seems to have lasted for about six months (to July A.D. 400). But he was evidently a man of no ability. He had not even a definite plan of action, and of his short period of power nothing is recorded except that he tried to secure for the Arians a church of their own within the city, and failed through the intolerant opposition of the Patriarch; and that his plans to seize the Imperial Palace, and to sack the banks of the money-changers , were frustrated.

This episode of German tyranny came to an abrupt end early in July. The Goth suddenly decided to quit the capital. We know not why he found his position untenable, or what his intentions were. Making an excuse of illness he went to perform his devotions in a church about seven miles distant, and ordered his Goths to follow him in relays. Their preparations for departure frightened the inhabitants, ignorant of their plans, and the city was so excited that any trifle might lead to serious consequences. It happened that a beggar-woman was standing at one of the western gates early in the morning asking for alms. At the unusual sight of a long line of Goths issuing from the

p134
gate she thought it was the last day for Constantinople and prayed aloud. Her prayer offended a passing Goth, and as he was about to cut her down a Roman intervened and slew him. The incident led to a general tumult, and the citizens succeeded in closing the gates, so that the Goths who had not yet passed through were cut off from their comrades without. There were some thousands of them\texthebrew{⁠}\footnote{Synesius says they numbered 7000, rather more than one-fifth of the whole army of Gaïnas; which has hence been reckoned by modern writers as 30,000 strong. The number is probably much too high. In any case the church could not have been large enough to hold 7000.

} but not enough to cope with the infuriated people. They sought refuge in a church (near the Palace) which had been appropriated to the use of such Goths as had embraced the Catholic faith. There they suffered a fate like that which had befallen the oligarchs of Corcyra during the Peloponnesian war. The roof was removed and the barbarians were done to death under showers of stones and burning brands (July 12, A.D. 400).\footnote{These events are related by Synesius, p1261 \emph{sqq.} ,

Zosimus, V.19 ,

Philostorgius, XI.8 , Socrates, VI.6, Sozomen, VIII.4. Socrates had the poems of Eusebius and Ammonius (see below) before him. For date see \emph{Chron. Pasch., sub a.}

}

The immediate consequence of this deliverance was the fall of Typhos\texthebrew{⁠}\footnote{From Synesius we know that his tenure of the office was less than a year, p1256 : \textgreek{ο\texthebrew{ὐ} γ\texthebrew{ὰ}ρ \texthebrew{ἐ}νιαυτο\texthebrew{ὺ}ς \texthebrew{ἀ}λλ\texthebrew{ὰ} μ\texthebrew{ῆ}νας \texthebrew{ἔ}φη το\texthebrew{ὺ}ς ε\texthebrew{ἰ}μαρτο\texthebrew{ὺ}ς ε\texthebrew{ἴ}ναι.}} and the return of Aurelian, who at once replaced him in the Prefecture. The conduct of Typhos was judicially investigated, his treasonable collusion with the Germans was abundantly exposed, and he was condemned provisionally to imprisonment. He was afterwards rescued from the vengeance of the mob by his brother. His subsequent fate is as unknown to us as his name. Aurelian, who had been designated for the consul ­ship of the year 400, but had been unable to enter upon it in January, seems now to have been invested with the insignia,\texthebrew{⁠}\footnote{This seems to be the meaning of Synesius, p12 \textgreek{μετ\texthebrew{ὰ} συνθήματος μείζονος.}Zosimus, V.18.8 , is inaccurate.

} and the name of whatever person had been chosen to fill it by Typhos and Gaïnas was struck from the Fasti.

Gaïnas, in the meantime, a declared enemy, like Alaric three years before, marched plundering through Thrace. But he won little booty, for the inhabitants had retreated into the strong places which he was unable to take. He marched to the Hellespont, intending to pass over into Asia. But when he reached

p135
the coast opposite Abydos he found the Asiatic shore occupied by troops, who were supported by warships. These forces were under the command of Fravitta, a loyal pagan Goth who in the last years of Theodosius had played a considerable part in the politics of his own nation as leader of the philo-Roman party. He had since served under Arcadius, had been promoted to be Master of Soldiers in the east, and had cleared the eastern Mediterranean of pirates from Cilicia to Syria and Palestine.\texthebrew{⁠}\footnote{Zosimus, V.20.2.

The article in Suidas, \emph{s.v.} \textgreek{Φράβιθος}, may come from Eunapius (see

Müller, \emph{F. H. G.} IV.49 ).

} The Goths encamped on the shore, but when their provisions were exhausted they resolved to attempt the crossing and constructed rude rafts which they committed to the current. Fravitta's ships easily sank them, and Gaïnas, who had remained on shore when he saw his troops perishing, hastened northwards, beyond Mount Haemus, even beyond the Danube, expecting to be pursued. Fravitta did not follow him, but he fell into the hands of Uldin, king of the Huns, who cut off his head and sent it as a grateful offering to Arcadius (December 23, A.D. 400). History has no regrets for the fate of this brutal and incompetent barbarian.

It was significant of the situation in the Empire that a Gothic enemy should be discomfited by a Goth. Fravitta enjoyed the honour of a triumph, and was designated consul for A.D. 401. Arcadius granted him the only favour he requested, to be allowed to worship after the fashion of his fathers.

Thus the German danger hanging over the Empire was warded off from the eastern provinces. Stilicho could no longer hope to interfere in eastern affairs through the Goths of the eastern army. The episode was a critical one in Roman history, and its importance was recognised at the time. It was celebrated in two epic poems\texthebrew{⁠}\footnote{The \emph{Gaïnia} of Eusebius (a pupil of Troilus, Aurelian's friend) and a poem of Ammonius (recited in 437), of which two lines are preserved in the \emph{Etymologicum genuinum}, 588.4 \texthebrew{—}\textgreek{\texthebrew{ἤ}δη δ\texthebrew{’ ὑ}ψιτενής τε Μίμας \texthebrew{ὑ}πελείπετ\texthebrew{’ ὀ}πίσσω, λείπετο δ\texthebrew{’ ὑ}ψικάρηνον \texthebrew{ἕ}δος Πιμπληϊδος \texthebrew{ἄ}κρης,}which seem to come in the description of a voyage along the coast of Asia Minor.

} as well as in the myth of Synesius. Scenes from the revolt were represented in sculpture on the pillar of Arcadius which was set up in A.D. 403 in the Forum named after him.\footnote{See Strzygowski, in \emph{Jahrb. des kais. arch. Instituts}, VIII.203 \emph{sqq.} (1893); C. Gurlitt, \emph{Antike Denkmalsäulen in Konstantinopel} (1909).

}

The year 400, which witnessed the failure of the German bid

p136
for ascendancy at Constantinople, was the year of Stilicho's first consul ­ship. Claudian celebrated it in a poem which was worthy of a greater subject:

The hero's services to the Empire in war and peace outshine the merits and glories of the most famous figures in old Roman history. The poet himself aspired to be to Stilicho what Ennius had been to Scipio Africanus. Noster Scipiades Stilicho \texthebrew{—} a strange conjunction of names; but we forgive the poet his hyperboles for his genuine sense of the greatness of Roman history. The consul ­ship of the Vandal general inspired him with the finest verses he ever wrote, a passage which deserves a place among the great passages of Latin literature \texthebrew{—} the praise of Rome, beginning \texthebrew{—}He has expressed with memorable eloquence the Imperial ideal of the Roman State:

The approaching disruption of the Empire was indeed hidden from Claudian and all others at the end of the fourth century. The Empire still reached from the Euphrates to the Clyde. Theodosius, who ruled a larger realm than Augustus, had steered it safely through dangers apparently greater than any which now menaced, and Stilicho was the military successor of Theodosius. The sway of Rome, if the Roman only looked at the external situation, might seem the assured and permanent order of the world:

Yet there was a very uneasy feeling in these years that the end of Rome might really be at hand. It was due to superstition.

p137
The twelve vultures that appeared to Romulus had in ages past been interpreted to mean that the life of Rome would endure for twelve centuries, and for some reason it was thought that this period was now drawing to a close:

The ancient auspice seemed to be confirmed by exceptional natural phenomena \texthebrew{—} the appearance of a huge comet in the spring of A.D. 400\texthebrew{⁠}\footnote{Claudian, \emph{ib.} 243 \emph{sqq.} The comet is also referred to by eastern writers (\emph{e.g.} Socrates, VI.6), and its appearance is recorded in Chinese annals. In the same passage, 233 \emph{sqq.}, are mentioned the eclipses which occurred in Dec. 17, 400, June 12 and Dec. 6, 401.

} and three successive eclipses of the moon.\texthebrew{⁠} a Before these signs appeared, Honorius and Stilicho had allowed the altar of Victory which had been removed from the Senate-house by Theodosius to be brought back, a momentary concession to the fears of the Roman pagans. And it is very probably due to superstitious fears that the work of restoring the walls of Rome was now taken in hand.\footnote{Seeck, \emph{Untergang}, V.329.

}

When Stilicho went to Rome to enter upon his consul ­ship,\texthebrew{⁠}\footnote{A fine consular diptych is preserved in the Cathedral of Monza, which is probably Stilicho's (whether to be associated with his first consul­ship in 400 or with his second in 405). The consul is represented on the left leaf, a bearded man standing with a lance in his left hand. On the right leaf is a lady (Serena) with pearl earrings and a necklace, and an oriental turban like a wig (we see similar coiffures on coins), holding a boy (Eucherius) by the hand. See Molinier, \emph{Cat. des ivoires.} The robe of state (trabea) which Stilicho wears is embroidered with pictures of his wife and son, according to the custom of the time, and it is interesting to find that Claudian in his \emph{De cons. Stil.} describes such a trabea, on which scenes of Stilicho's family life (including the birth of Maria, Eucherius practising horseman­ship) were represented

(III.340 \emph{sqq.}) . A good reproduction will be found in the Album (vol. I pl. 1) to the \emph{Histoire des art indust.} of Labarte, who thought that it was a diptych of Aetius, with Placidia and Valentinian III. It was the custom for the consul of the year to present to senators these ivory diptychs (two pieces of ivory joined by hinges), to commemorate his year of office. They were generally inscribed with the consul's name and titles, and many specimens of them have survived from the fifth and sixth centuries.

Thayer's Note: Several photographs of the famous diptych are online, among which these have informative pages that look fairly permanent:

Circolo Sardegna di Brienza . In view of the transience of webpages, however,

this link

may be useful as well.

Notice that Bury seems to have described a mirror image of the diptych; either that, or all the photographs of it online are reversed \texthebrew{—} which is not as improbable as it sounds since they likely derive from one or two originals. I have not seen the actual object and cannot vouch for it either way, but since it was normal to hold the lance with the right arm and the shield with the left, I suspect the mistake is Bury's. (Amusingly, for those of you who did what I did and looked at your Britannica, my copy \texthebrew{—} of the notoriously substandard Fifteenth Edition \texthebrew{—} agrees with the Web photos, but then goes on to caption it as being of ebony, a clear mistake, almost certainly springing from a mistranslation of the original Alinari caption: eburneo in Italian means of \emph{ivory}, not \emph{ebony}. Never, ever trust anything you read.)

} Claudian accompanied him, and his verses richly deserved the statue which was erected at the instance of the senate in the Forum of Trajan ``to the most glorious of poets,'' although (the inscription runs) ``his written poems suffice to keep his memory eternal.''\footnote{CIL VI.1710 , from which we learn that Claudian was a tribunus et notarius. A distich in Greek is appended to the inscription:

\textgreek{ε\texthebrew{ἰ}ν \texthebrew{ἑ}ν\texthebrew{ὶ} Βιργιλίοιο νόον κα\texthebrew{ὶ} μο\texthebrew{ῦ}σαν \texthebrew{Ὀ}μήρου, Κλαυδιαν\texthebrew{ὸ}ν \texthebrew{Ῥ}ώμη κα\texthebrew{ὶ} βασιλ\texthebrew{ῆ}ς \texthebrew{ἔ}θεσαν.}Thayer's Note: For a full transcription and translation of the inscription, and a link to a photograph, see Maurice Platnauer's introduction to the Loeb edition of Claudian

(note 5) .}

\xchapter{Chapter V, Part B}

Short URL for this page:

bit.ly/BURLAT5B

\xsubsection{(Part 2 of 3)}

It was during the interlude in which Gaïnas and Typhos were supreme that Eudoxia, who had borne Arcadius two daughters, was crowned Augusta (January 9, A.D. 400).\texthebrew{⁠}\footnote{Copper coins of Ael. Eudoxia Aug., with Gloria Romanorum on the reverse, are ascribed by De Salis (\emph{Coins of the Eudoxias}) to A.D. 400, 401, before the coronation of the child Theodosius (in Jan. 402), and her gold coins with Salus Reipublicae and a seated Victory holding the monogram of Christ (and copper coins with the same legend) to the period after that event.

} Notwithstanding her German descent, she had no sympathies with the German party, though she had independently helped them to compass the fall of Eutropius. It is significant that of the hostages whom Gaïnas had demanded, John was notoriously her favourite and Saturninus was the husband of her intimate friend Castricia. The Empress was a woman of forceful character and impulsive temper,\texthebrew{⁠}\footnote{Compare above,

p109, n1 , and

Zosimus, V.24

\textgreek{πέρα τ\texthebrew{ῆ}ς φύσεως α\texthebrew{ὐ}θαδιζομένη.}} and after the eunuch's fall she won unbounded influence over her weak and sluggish husband. Her historical importance centres in the conflict into which she was drawn with Chrysostom, a drama which was to settle the future relations between the Imperial and the Patriarchal authority. No critical collision had occurred before. With the exception of Valens no Emperor had resided constantly at Constantinople before Arcadius, who never left the capital except for a summer holiday at Ancyra. Moreover, the see had only recently attained to the first rank in the Eastern Empire ( A.D. 381), and its primacy was hotly disputed by Alexandria. That the collision between Emperor and Patriarch occurred at this time was due principally to the aggressive and uncompromising character of Chrysostom.

John, the `` golden-mouthed `` preacher, was in his forty-sixth or forty-seventh year when he became bishop of Constantinople (February 26, A.D. 398).\texthebrew{⁠}\footnote{Besides the monographs of Thierry, Stephens, Ludwig, Puech (see

Bibliography ), there is a good article by E. Venables in the \emph{Dict. of Christian Biography.} The chief sources for his life are his own letters and sermons, the \emph{Dialogue} of Palladius (a very partial work), and the \emph{Histories} of Socrates and Sozomen.

} He was an independent and austere man, who in his own habits carried asceticism to excess, and his ways were rough and uncourtly. At Constantinople he found himself confronted by a superb court under the sway of Eudoxia. There is no reason to suppose that it was particularly vicious,

p139
but it was at least frivolous and embodied for him the pride of life and the pomps and vanities of the world.

Chrysostom stands alone among the great ecclesiastics of the later Empire in that his supreme interest lay not in controversial theology but in practical ethics. His aim was the moral reformation of the world, and as his work lay in two rich cities, Antioch and Constantinople, he conceived it to be one of his chief duties to strive against the flaunting luxury of the rich classes, and denounce the lavish expenditure of wealth on personal gratification, wealth which in his eyes should have been devoted to alleviating the lot of the poor. Thus we learn from his sermons, whether at Constantinople or at Antioch, many details as to the luxurious life of the higher classes. Many rich nobles possessed ten or twenty mansions and as many private baths; a thousand, if not wellnigh two thousand, slaves called them lord, and their halls were thronged with eunuchs, parasites, and retainers.\texthebrew{⁠}\footnote{For the description of such houses see the \emph{Homily} on \emph{Ps.} 48.17, \emph{P.G.} 55.510\texthebrew{‑}511.

} In their gorgeous houses doors were of ivory, the ceilings lined with gold, the floors inlaid with mosaics or strewn with rich carpets; the walls of the halls and bedrooms were of marble, and wherever commoner stone was used the surface was beautified with gold plate. Nude statues, to the scandal of strict ecclesiastics, decorated the halls. Spacious verandahs and baths adjoined the houses, which were surrounded by gardens with fountains. The beds were made of ivory or solid silver, or, if on a less expensive scale, of wood plated with silver or gold. Chairs and stools were usually of ivory, and the most homely vessels were often of the most costly metal; the semicircular tables or sigmas, made of gold or silver, were so heavy that two youths could hardly lift one. Oriental cooks were employed; and at banquets the atmosphere was heavy with all the perfumes of the East, while flute girls, whose virtue was as easy as in the old days of Greece and Rome, entertained the feasters.

To Chrysostom the contrast between the life of the higher classes and the miseries of the toiling populace was such a painful spectacle, that he was almost a socialist. If he inveighs against the men for their banquets, he is no less severe on the women for their sumptuous mule-cars, their rich dresses, their jewellery,

p140
their coquettish toilettes.\texthebrew{⁠}\footnote{See also the verses of Gregory Naz. \textgreek{Κατ\texthebrew{ὰ} γυναικ\texthebrew{ῶ}ν καλλωπιζομένων} (\emph{Carmina}, I. sect. ii.29, \emph{P.G.} 37.884).

} Their extravagance often involved their husbands in expenses which they could not afford. He denounces the use of silk and brocade. All ``evils'' which Chrysostom describes are characteristic \texthebrew{—} allowance being made for difference of environment \texthebrew{—} of all wealthy societies, pagan or christian. His passionate denunciations of the rich have the same import and value as the denunciations of modern European plutocrats by socialists.

The problem of marriage interested him, and he preached the unpopular doctrine that the two partners in marriage are equal, the woman having the same rights against an unfaithful husband as the man against an unfaithful wife. We should hardly require the express evidence with Chrysostom supplies, to know that marriages for money were frequent. He complains that children were excessively indulged, and that their fathers too often gave their sons the worst possible moral education.\texthebrew{⁠}\footnote{\emph{Ep. I ad Tim., Hom.} 9, \emph{P.G.} 64.546.

} It is interesting to learn from his homilies that the treatment of slaves was still often marked by much of the old brutality. People passing in the street might often hear the furious outbreaks of an angry mistress beating her maid. Chrysostom describes vividly how a wife summoned her husband to aid her in punishing an offending servant.\texthebrew{⁠}\footnote{\emph{Ep. ad Ephes., Hom.} 15, \emph{P.G.} 64.109.

} The girl is stripped, tied to the foot of the bed, whipped by the master, while the mistress exhausts her vocabulary of abuse. The offence was probably quite trivial, perhaps an awkwardness in assisting at the mistress's toilette.\texthebrew{⁠}\footnote{Cp. Juvenal VI.490 \emph{sqq.} ;

Martial II.66 .

} The condition of domestic slaves had in some respects changed little more than human nature since the days of Juvenal. But harsh and brutal treatment was not more universal than in those days. There were many masters (as other passages of Chrysostom show) who took the deepest interest in the well-being of their slaves. And there was also another side to the question. The servants were often trying and maleficent, slandering and spying upon their owners. The troubles which were caused by the lying tongues of maidservants are actually urged by Chrysostom as an argument against marriage.

p141
Christianity had not yet succeeded in abolishing all the old pagan customs from the celebrations of funerals and marriages. In the reign of Arcadius the usage was still maintained of hiring female mourners to sing dirges over the dead. Chrysostom considered it idolatry, and even threatened to excommunicate those who practised it. He also stigmatised the pagan practice of ablutions after the funeral ceremony, which were intended to purify from contact with the dead. The expense and ostentation which marked the funerals of the rich also earned his censure. More scandalous in the eyes of austere Christians were the survivals of pagan manners on the occasion of weddings. The Church had introduced an ecclesiastical ceremony in the presence of the bishop, but as soon as this was completed, the wedding was celebrated in the old way. The bride was conducted in procession at nightfall from the house of her father to that of the bridegroom. The procession was followed by troops of actors and actresses and dancing-girls , who were admitted to the house, where they danced indecently and sang indelicate songs. The

epithalamia

and the

odes

which Claudian composed on the occasions of the marriages of Honorius may give some idea of the licence which was still fashionable.

Chrysostom fought not only against the extravagance of the rich but also against the sensuality, gluttony, and avarice of the clergy and the monks, to whom his austerity was, in the words of his biographer, ``as a lamp burning before sore eyes.'' Women were introduced into the monasteries or shared the houses of priests as ``spiritual sisters,'' a practice which if often innocent was always a snare.\texthebrew{⁠}\footnote{See the vivid picture of such a ménage in the sermon \emph{Contra eos qui subintroductas habent virgines}, c. 9.

} Deaconesses, unable to adopt the meretricious apparel that had become the mode, arranged their coarse dresses with an immodest coquetry which made them more piquant than professional courtesans. º

The Patriarch had his own devoted female admirers. The most distinguished was the deaconess Olympias, a rich lady, who in her early girlhood had been a favourite of Gregory Nazianzene. Her bounty to the poor won the heart of Chrysostom, to whom she proved a most unselfish and devoted friend. Another of his friends was Salvina, daughter of the Moor Gildo, whom Theodosius had given in marriage to Nebridius his

p142
wife's nephew. In ``A Letter to a Young Widow'' Chrysostom contrasts the peaceful happiness of her life at Constantinople with the unrest of her father's turbulent career. A deacon named Serapion was the Patriarch's trusted and devoted counsellor, but his influence was not always wisely exerted. He had no judgment, and instead of trying to restrain the impetuous temper of Chrysostom, encouraged or incited him to rash acts.

With the common people the Patriarch enjoyed great popularity. He was no respecter of persons, and he interpreted Christianity in a socialistic sense which has not generally been countenanced or encouraged by the Church. Though it was not political but social inequality that he deprecated, and nothing was further from his thoughts than to upset the established order of things, the spirit of his teaching certainly tended to set the poor against the rich. On the occasion of an earthquake he said publicly that ``the vices of the rich caused it, and the prayers of the poor averted the worst consequences.'' It was easy for his enemies to fasten upon utterances like this and accuse him of ``sedu­cing the people.'' His friendships with Olympias and other women whom he sometimes received alone supplied matter for another slander. Having ruined his digestive organs by excessive asceticism, he made a practice of not dining in company, and in consequence of this unsocial habit he was suspected of private gluttony.

For three years Chrysostom and Eudoxia were on the best of terms. Chrysostom owed his see, Eudoxia her throne, to Eutropius, and they both refused to be his creatures. But early in A.D. 401 she did something which evoked a stern rebuke from the Archbishop, and the consequence of his audacity was that he was not received at Court. We learn of this in connexion with an episode which reveals Eudoxia herself in an amiable light.

Porphyrius, the bishop of Gaza, with other clergy of that diocese, visited Constantinople in the spring of A.D. 401, to persuade the government to take strong measures for the suppression of pagan practices. For the citizens of Gaza still obstinately held to the worship of their old deities, Aphrodite, the Sun, Persephone, and above all Marnas, the Cretan Zeus. When the clergy reached the capital and secured lodgings, their first act was to visit Chrysostom. ``He received us with great honour and courtesy, and asked us why we undertook the fatigue of

p143
the journey, and we told him. And he bade us not to despond but to have hope in the mercies of God, and said, 'I cannot speak to the Emperor, for the Empress excited his indignation against me because I charged her with a thing which she coveted and robbed. And I am not concerned about his anger, for it is themselves they hurt and not me, and even if they hurt my body they do the more good to my soul. . . . To\texthebrew{‑}morrow I will send for the eunuch Amantius, the castrensis (chamberlain) of the Empress, who has great influence with her and is really a servant of God, and I shall commit the matter to him.' Having received these injunctions and a recommendation to God, we proceeded to our inn. And on the next day we went to the bishop and found in his house the chamberlain Amantius, for the bishop had attended to our affair and had sent for him and explained it to him. And when we came in, Amantius stood up and did obeisance to the most holy bishops, inclining his face to the ground, and they, when they were told who he was, embraced him and kissed him. And the archbishop John bade them explain orally their affair to the chamberlain. And Porphyrius explained to him all the concealment of the idolaters, how licentiously they perform the unlawful rites and oppress the Christians. And Amantius, when he heard this, wept and was filled with zeal for God, and said to them, 'Be not despondent, fathers, for Christ can shield His religion. Do ye therefore pray, and I will speak to the Augusta.'

``The next day the chamberlain Amantius sent two deacons to bid us come to the Palace, and we arose and proceeded with all expedition. And we found him awaiting us, and he took the two bishops and introduced them to the Empress Eudoxia. And when she saw them she saluted them first and said, 'Give me your blessing, fathers,' and they did obeisance to her. Now she was sitting on a golden sofa. And she said to them, 'Excuse me, priests of Christ, on account of my situation, for I was anxious to meet your sanctity in the antechamber. But pray God on my behalf that I may be delivered happily of the child which is in my womb.' And the bishops, wondering at her condescension, said, 'May He who blessed the wombs of Sarah and Rebecca and Elizabeth, bless and quicken the child in thine.' After further edifying conversation, she said to them, 'I know why ye came, as the castrensis Amantius explained it

p144
to me. But if you are fain to instruct me, fathers, I am at your service.' Thus bidden, they told her all about the idolaters, and the impious rites which they fearlessly practised, and their oppression of the Christians, whom they did not allow to hold a public office nor to till their lands 'from whose produce they pay the dues to your Imperial sovereignty.' And the Empress said, 'Do not despond; for I trust in the Lord Christ, the Son of God, that I shall persuade the Emperor to do those things that are due to your saintly faith and to dismiss you hence well treated. Depart, then, to your privacy, for you are fatigued, and pray God to co-operate with my request.' She then commanded money to be brought, and gave three handfuls of money to the bishops, saying, 'In the meantime take this for your expenses.' And the bishops took the money and blessed her abundantly and departed. And when they went out they gave the greater part of the money to the deacons who were standing at the door, reserving little for themselves.

``And when the Emperor came into the apartment of the Empress, she told him all touching the bishops, and requested him that the heathen temples of Gaza should be pulled down. But the Emperor was put out when he heard it, and said, 'I know that city is devoted to idols, but it is loyally disposed in the matter of taxation and pays a large sum to the revenue. If then we overwhelm them with terror of a sudden, they will betake themselves to flight and we shall lose so much of the revenue. But if it must be, let us afflict them partially, depriving idolaters of their dignities and other public offices, and bid their temples be shut up and be used no longer. For when they are afflicted and straitened on all sides they will recognise the truth; but an extreme measure coming suddenly is hard on subjects.' The Empress was much vexed at this reply, for she was ardent in matters of faith, but she merely said, 'The Lord can assist his servants the Christians, whether we consent or decline.'

``We learned these details from the chamberlain Amantius. On the morrow the Augusta sent for us, and having first saluted the bishops according to custom, she bade them sit down. And after a long spiritual talk, she said, 'I spoke to the Emperor, and he was somewhat displeased. But do not despond, for, God willing, I cannot cease until ye be satisfied and depart,

p145
having succeeded in your pious purpose.' And the bishops made obeisance. Then the sainted Porphyrius, moved by the spirit, and recollecting the word of the thrice blessed anchoret Procopius, said to the Empress: 'Exert yourself for the sake of Christ, and in recompense for your exertions He can bestow on you a son whose life and reign you will see and enjoy for many years.' At these words the Empress was filled with joy, and her face flushed, and new beauty beyond that which she already had passed into her face; for the outward appearance shows what passes within. And she said, 'Pray, fathers, that according to your word, with the will of God, I may bear a male child, and if it so befall, I promise you to do all that ye ask. And another thing, for which ye ask not, I intend to do with the consent of Christ; I will found a church at Gaza in the centre of the city. Depart then in peace, and rest quiet, praying constantly for my happy delivery; for the time of the birth is near.' The bishops commended her to God and left the Palace. And prayer was made that she should bear a male child; for we believed in the words of Saint Procopius the anchoret.

``And every day we used to visit John, the archbishop, and had the fruition of his pious discourse, sweeter than honey and the honey comb. And Amantius the chamberlain used to come to us, sometimes bearing messages from the Empress, at other times merely to pay a visit. And after a few days she brought forth a male child [April 10], and he was called Theodosius after his grandfather Theodosius, the Spaniard, who reigned along with Gratian. And the child Theodosius was born in the purple, wherefore he was proclaimed Emperor at his birth. And there was great joy in the city, and men were sent to the cities of the Empire, bearing the good news, with gifts and bounties.

``But the Empress, who had only just been delivered, sent Amantius to us with this message: 'I thank Christ that God bestowed on me a son, on account of your holy prayers. Pray, then, fathers, for his life and for my lowly self, in order that I may fulfil those things which I promised you, Christ himself again consenting through your holy prayers.' And when the seven days of her lying-in were fulfilled, she sent for us and met us at the door of the chamber, carrying in her arms the infant in the purple robe. And she inclined her head and said, 'Draw

p146
nigh, fathers, unto me and the child which the Lord granted to me through your holy prayers.' And she gave them the child that they might seal it (with God's signet). And the bishops sealed both her and the child with the seal of the cross, and, offering a prayer, sat down. And when they had spoken many words full of edification, the lady says to them, 'Do ye know, fathers, what I resolved to do in regard to your affair?' [Here Porphyrius related a dream which he had dreamed the night before; then Eudoxia resumed:] 'If Christ permit, the child will be privileged to receive baptism in a few days. Do ye then depart and compose a petition and insert in it all the requests ye wish to make. And when the child comes forth from the baptismal rite, give the petition to him who holds the child in his arms; and I shall instruct him what to do.' Having received these directions we blessed her and the infant and went out. Then we composed the petition, inserting many things in the document, not only as to the overthrow of the idols but also that privileges and revenue should be granted to the holy Church and the Christians; for the Church was poor.

``The days ran by, and the day on which the young Emperor Theodosius was to be baptized arrived. And all the city was crowned with garlands and decked out in garments made of silk and gold jewels and all kind of ornaments, so that no one could describe the adornment of the city. One might behold the inhabitants, multitudinous as the waves, arrayed in all manner of garments. But it is beyond my power to describe the brilliance of that pomp; it is a task for those who are practised writers, and I shall pursue my true history. When the young Theodosius was baptized and came forth from the church to the Palace, you might behold the magnificence of the multitude of the magnates and their dazzling raiment, for all were dressed in white, and you would have thought they were covered with snow. The patricians headed the procession, with the illustres and all the other ranks, and the military contingents, all carrying wax candles, so that the stars seemed to shine on earth. And close to the infant, which was carried in arms, was the Emperor Arcadius himself, his face cheerful and more radiant than the purple robe he was wearing, and one of the magnates carried the infant in brilliant apparel. And we marvelled, beholding such glory. . . .

p147
``And we stood at the portal of the church, with our petition, and when he came forth from the baptism we called aloud, saying, 'We petition your Piety,' and held out the paper. And he who carried the child seeing this, and knowing our business, for the Empress had instructed him, bade the paper be showed to him, and when he received it halted. And he commanded silence, and having unrolled a part he read it, and folding it up, placed his hand under the head of the child and cried out, 'His majesty has ordered the requests contained in the petition to be ratified.' And all having seen marvelled and did obeisance to the Emperor, congratulating him that he had the privilege of seeing his son an emperor in his lifetime; and he rejoiced thereat. And that which had happened for the sake of her son was announced to the Empress, and she rejoiced and thanked God on her knees. And when the child entered the Palace, she met it and received it and kissed it, and holding in her arms greeted the Emperor, saying, 'You are blessed, my lord, for the things which your eyes have beheld in your lifetime.' And the king rejoiced thereat. And the Empress, seeing him in good humour, said, 'Please let us learn what the petition contains that its contents may be fulfilled.' And the Emperor ordered the paper to be read, and when it was read, said, 'The request is hard, but to refuse is harder, since it is the first mandate of our son.' ``

The petition was granted, and Eudoxia arranged a meeting between the quaestor, the minister on whom it devolved to draft the Imperial rescripts, and the bishops, that all the wishes of the latter might be incorporated in the edict. The execution of it, which was invidious and required a strong hand and will, was entrusted to Cynegius, and the bishops returned to Palestine, having received considerable sums of money from the Empress and Emperor, as well as the funds which the Empress had promised for the erection of a church at Gaza.

This narrative gives us an idea of the kind of little dramas that probably lay behind many of the formal decrees and rescripts preserved in the Imperial Codes. The wonder of the provincial bishops at the splendid apparel of the great of the earth, their edifying spiritual conversations with the Empress, with the eunuch, and with the archbishop, the ruse of Eudoxia to compass the success of the petition, all such details help

p148
us to realise the life of the time; while the hesitation of the pious Arcadius to root out the heathen ``abominations'' because the heathen were respectable taxpayers shows that even he, when the ghostly and worldly policies of the Empire clashed, was more inclined to be Emperor than churchman.

To return to Chrysostom. When he performed the ceremony of baptizing the Emperor's son and heir, there must have been a reconciliation with the court, but Eudoxia could not forget the incident, and henceforward she would be at least disposed to lend a patient ear to his enemies. And his enemies were many, both in clerical and in secular circles. Among the fashionable ladies who were particularly offended by his castigations of female manners were three who were intimate friends of the Empress \texthebrew{—} Marsa, wife of Promotus, in whose house Eudoxia had been brought up; Castricia, the wife of Saturninus, whom Chrysostom had helped to rescue from the vengeance of Gaïnas; and Eugraphia, whose house was a centre for all those who detested him.\texthebrew{⁠}\footnote{See Palladius, c. 4, c. 8. The author describes Eugraphia as \textgreek{\texthebrew{ἀ}μφιμανής τις, τ\texthebrew{ὰ} δ\texthebrew{ὲ} λοιπ\texthebrew{ὰ} α\texthebrew{ἰ}δο\texthebrew{ῦ}μαι κα\texthebrew{ὶ} λέγειν.}} It is easy to imagine how easily they could continue to poison Eudoxia's mind against a priest who was exceptionally tactless by twisting his invectives against the foibles of women into personal attacks upon herself.

But the agitation of irresponsible enemies might not have shaken his position, if he had not committed indiscretions in the domain of ecclesiastical policy. Antoninus, the bishop of Ephesus, had been accused of simony and other offences, and Chrysostom was appealed to. He determined to investigate the matter on the spot, and set out in the winter of A.D. 401.\texthebrew{⁠}\footnote{Cp. Seeck, \emph{Untergang}, V. p577.

} The inquiry disclosed abuses in many of the churches of western Asia Minor, and Chrysostom acted with more zeal than wariness. He deposed and replaced at least thirteen bishops, exceeding the rights of his jurisdiction, and, it was said, not giving a fair hearing to the cases. Naturally he stirred up many new enemies.

He was absent five months from Constantinople. He had deputed an eloquent Syrian, Severian, bishop of Gabala, to act for him during his absence. Severian seems to have joined the league of his enemies, and there was an open rupture between him and Serapion the deacon. When the Patriarch returned he found his own See disorganised, and a local council was held

p149
to hear the charges which Serapion brought against Severian. When Severian, who felt sure of support in high quarters, resisted the efforts of the bishops to induce him to be reconciled with the deacon, Chrysostom told him that it would be well for him to return to the see of Gabala which he had so long neglected. Severian, who seems to have entertained the ambition of repla­cing Chrysostom on the Patriarchal throne, now saw that he had gone too far, and he left the city. At Chalcedon he was recalled. The Empress had herself implored the Patriarch to reconcile himself with Severian. Throughout the quarrel popular opinion had been on Chrysostom's side, but it may be questioned whether his conduct was altogether creditable.\texthebrew{⁠}\footnote{The silence of Palladius is significant (cp. Seeck, \emph{op. cit.} 577).

} He yielded to Eudoxia's prayers, but it was necessary to tranquillise popular feeling, for which purpose he preached a pacific sermon which ended with the words, ``Receive our brother Severian the bishop.''\texthebrew{⁠}\footnote{Text in Papadopolous-Kerameus, \textgreek{\texthebrew{Ἀ}ναλ. \texthebrew{ἱ}εροσολ. σταχυολογίας}, i.15 \emph{sqq.}; Latin translation in Migne, \emph{P.G.} 52, 423 \emph{sqq.}

} Severian responded by a sermon of which the note was likewise peace. But the peace was hollow.

A new storm from another quarter was soon to burst over Chrysostom. Theophilus, the archbishop of Alexandria, bore no goodwill to the eloquent preacher who occupied the great see which had now precedence over his own. Theophilus, whose principal claim to be remembered is the destruction of the Serapeum, the famous stronghold of paganism at Alexandria, seems, so far as we can judge from his acts, to have been a domineering and unscrupulous prelate. He had probably been spoiled by the enjoyment of power. He is described as ``naturally impulsive, bold and precipitous in action, extraordinarily quarrelsome, impatient and determined in grasping at any object he had set his mind on.''\texthebrew{⁠}\footnote{Palladius, c. 9.

} He had hoped to secure for a candidate of his own the archiepiscopal chair of Constantinople after the death of Nectarius, and had not forgiven Chrysostom his disappointment; which was rendered particularly humiliating by the fact that Eutropius had forced him to take part in Chrysostom's consecration. Theophilus had held the heretical opinion of Origen, who rejected the anthropomorphic conception of the Deity which is suggested by many passages in the Hebrew Scripture. The same opinion was held in a monastic settlement

p150
in the desert of Nitria in Upper Egypt, over which four monks presided who were known, from their remarkable stature, as the Tall Brothers.\texthebrew{⁠}\footnote{Ammonius, Dioscorus (whom Theophilus made bishop of Hermupolis), Eusebius, and Euthymius.

} Theophilus, however, changed his view on the theological point and ( A.D. 401) issued a Paschal letter condemning Origen and his disciples. He then convoked a synod, which anathematised Origen and condemned the Nitrian monks. He had other reasons for desiring the destruction of the Tall Brothers, and he obtained troops from the augustal Prefect of Egypt to arrest them. The habitations of the monks were sacked and pillaged, and the Tall Brothers with their followers, clad in sheepskins, made their way to Palestine, where the bishops, admonished by letters from Theophilus, refused them shelter. Unable to find rest for the soles of their feet, they took ship for Constantinople to place themselves under the protection of Chrysostom. He received them kindly, but would not communicate with them until their cause had been examined, and he lodged them in the church of St. Anastasia,\texthebrew{⁠}\footnote{On this church see Du Cange, \emph{Cplis Christiana}, bk. IV pp98\texthebrew{‑}99; Paspates, \textgreek{Βυζ. μελέται}, 368 \emph{sqq.}

} where their wants were ministered to by his deaconesses.

The piety and virtues of the Tall Brothers were well known by repute at Constantinople, and the Empress was eager to exert herself in their behalf. Meeting one of them as she was driving through the city, she stopped her carriage, asked him to pray for her, and promised to arrange that a synod should be convoked and Theophilus summoned to attend it. The monks then drew up a petition to the Emperor, setting forth their charges against their archbishop, and an Imperial messenger was sent to Alexandria to compel Theophilus to come to Constantinople and answer for his conduct at a synod to be held there.

Theophilus had already instigated Epiphanius, bishop of Constantia in Cyprus, who was an authority on heresies, to convene a synod of the Cypriot bishops to condemn the opinion of Origen, and to circulate its decisions to the sees of the Church. This had been done, and Theophilus, finding himself in an awkward position by the peremptory summons to appear as a defendant in the capital, urged Epiphanius to go in person to Constantinople and obtain Chrysostom's signature to the decree of the Cypriote council. Epiphanius, persuaded by the crafty

p151
flatteries of the Alexandrian prelate that a crisis in the Church depended on his intervention, sailed for Constantinople (early in A.D. 403). But he was not a strong ally; he was out of place and bewildered amid the intrigues of the capital. Finally he became acquainted with the Tall Brothers, and when they told him that they had read his books\texthebrew{⁠}\footnote{For his works on heresies, the \emph{Panarion} and the \emph{Ancoratus}, see the article on him by R. A. Lipsius in \emph{Dict. of Chr. Biography.}

} with admiration, and remonstrated with him for condemning their writings, which he was obliged to confess he only knew from hearsay, he came to the conclusion that he had made a mistake and allowed himself to be used as a tool by Theophilus. Disgusted and dejected he set sail for home, but the fatigue and excitement had overtaxed his failing strength and he died on the voyage (May 12).

About a month later (in June) Theophilus arrived with a large retinue of bishops who came to support him from Egypt, Syria, and Asia Minor. He had been summoned to appear as an accused man before an ecclesiastical tribunal over which Chrysostom would preside, but he was determined to invert the parts, and be himself the judge, with Chrysostom at the bar. That he succeeded in his plan was due entirely to Chrysostom's indiscretions. The Empress had interested herself in the affair of the Tall Brothers, and it was due to her influence that Theophilus had been forced to come to answer for his conduct. If Chrysostom, who in that affair had shown admirable caution, had now exercised ordinary tact and self-restraint , he could have had Eudoxia entirely on his side and might have defied all the arts and intrigues of his Alexandrian rival. Eudoxia had shown her veneration for the saintly bishop Epiphanius, by asking him to pray for her infant son who was ill, and Chrysostom, offended by her graciousness towards a bishop who had been openly hostile to himself, preached a violent sermon against women, in which the word Jezebel was pronounced. The congregation interpreted it as allusive to the Empress, and the matter was soon brought to her ears.\texthebrew{⁠}\footnote{Socrates, VI.15, combined with Palladius, c. 8. Cp. below.

} She was furious at the insult, and prepared to exert all her influence to support the party which was planning the ruin of the archbishop. Theophilus, rejecting the hospitality which Chrysostom offered him, established himself in the palace of Placidia, close to the Great Palace, and his

p152
bribes, banquets, and flatteries drew thither all the ecclesiastics and fashionable ladies whom Chrysostom had offended.

Chrysostom seems hardly to have realised the danger of his position. Instead of attempting to turn away the wrath of the Empress, he adopted a weak and conciliatory attitude towards the archbishop of Alexandria. The question of the Tall Brothers, though it was now a secondary consideration, had to be disposed of before Theophilus could take any open steps against Chrysostom, and Chrysostom was invited by the Emperor to preside over an investigation into the charges they had preferred against Theophilus. But he declined on the ground that such an inquiry into things which had occurred in another diocese would be illegal. This decision at once freed Theophilus from his position as an accused person, and the board was clear for him to organise his attack on Chrysostom. A list of charges was drawn up, sufficient to move the Emperor, under his wife's influence, to summon a council to inquire into them. Witnesses were procured to substantiate the accusations.

Popular feeling ran so high in favour of Chrysostom that the authorities were afraid to hold the synod within the precincts of the city, and it met across the water in the palace of the Oak, which had been built by the Praetorian Prefect Rufinus in the suburbs of Chalcedon. Chrysostom refused to appear before a body which was packed with his enemies. The majority of the bishops present were Egyptians, prepared to do whatever their archbishop told them. The chief accuser of Chrysostom was John, his archdeacon. Among the numerous charges that were formulated for the synod to investigate were these: that he had sold the marble which Nectarius had set aside for decorating the church of St. Anastasia; that he had reviled the clergy as corrupt; that he had called Epiphanius a fool and a demon; that he had intrigued against Severian; that he received visits from women by themselves after he had sent every one else out of the room; that a bath was heated for him alone, and that after he had bathed Serapion emptied the bath so that no one else might use it; that he ate gluttonously alone, living like a Cyclops.\texthebrew{⁠}\footnote{The proceedings of the synod are known from a summary of the Acts in

Photius, \emph{Bibliotheca}, 59 . It sat for about a fortnight (there were 13 sessions), and there were 45 members. Palladius, c. 3, gives the number as 36, of whom 29 were Egyptians. Bishops friendly to

(p153)
Chrysostom did not attend. The date of the synod may probably be fixed to the end of June, or July. As Epiphanius died on May 12 (ASS. III.36) on his voyage to Cyprus, he must have left the city early in May. A short time elapsed before the arrival of Theophilus (\textgreek{ο\texthebrew{ὐ} πολ\texthebrew{ὺ}ς \texthebrew{ἐ}ν μέσ\texthebrew{ῳ} χρόνος}, Socrates, VI.15), who spent three weeks in organising the opposition to Chrysostom (\textgreek{τρε\texthebrew{ῖ}ς \texthebrew{ἐ}βδομάδας \texthebrew{ἡ}μερ\texthebrew{ῶ}ν}, Pallad. c. 8). This gives the middle of June as the earliest date for the opening of the synod, and the beginning of July as the earliest date for Chrysostom's exile. On the other hand, the synod cannot be placed later than some time in July. This accords with the indication of Palladius, in c. 9, \textgreek{παρίππασαν μ\texthebrew{ῆ}νες \texthebrew{ἐ}ννέα \texthebrew{ἢ} δέκα}, sc. from the first exile of Chrysostom to Lent or Easter 404, as the context shows. Cp. Tillemont, \emph{Mém. ecc.} XI. p601.

} The accusations which really demanded an inquiry

p153
concerned his conduct in deposing bishops in Asia and ordaining others without due investigation of their characters.

As Chrysostom, repeatedly summoned, refused to appear and plead, he was condemned, not as guilty of the crimes which were alleged against him, but because he refused to appear, and he was formally deposed from his see. A report of the result was communicated to the Emperor, with the suggestion that it was for him and not for the Council to deal with the charge that the archbishop had spoken treasonably of the Empress.\texthebrew{⁠}\footnote{By calling her Jezebel. See Palladius, c. 8. This charge was not formally included in the list of charges presented to the synod, and did not appear in the Acta.

} Arcadius confirmed the decree in a rescript which pronounced the sentence of banishment. To the archbishop's enemies the penalty may have seemed too lenient, but it roused the indignation of the people, who would not have their idol removed by the act of a small packed assembly like the Synod of the Oak. Loud clamours were raised for the assembling of a general Council of the Church. Flocking round St. Sophia and the archiepiscopal palace, the populace made it impossible for the Imperial officers to seize Chrysostom and expel him from the city for three days. He delivered two discourses in the church, in which he referred to the Empress as a Jezebel or a Herodias. ``One day she called me the thirteenth apostle, and now her name for me is Judas.''\texthebrew{⁠}\footnote{\emph{Homilia ante exilium}, \emph{P.G.} 52, 437; the other text, 427 \emph{sqq.}, seems to be a second version of the same discourse. As Seeck has observed, it was doubtless taken down in shorthand (\textgreek{\texthebrew{ὑ}π\texthebrew{ὸ} τ\texthebrew{ῶ}ν \texthebrew{ὀ}ξυγράφων}, Socrates, VI.4).

} But he had no intention of defying the Emperor or causing a sedition. He stole out from his palace at night, surrendered himself, was taken across to the Asiatic coast, and withdrew to Praenetus near Nicomedia.

When it was discovered that he had departed, the fury of the people burst out. The city was in an uproar. The populace clamoured for the recall of their pastor, and an earthquake

p154
which at this crisis shook the city and the Great Palace was interpreted to mean that the voice of the people was the voice of God.\texthebrew{⁠}\footnote{The earthquake is recorded only by Theodoret, V.34.5 \textgreek{σεισμο\texthebrew{ῦ} δ\texthebrew{ὲ} μεγίστου νύκτωρ γεγενημένουκ} º \textgreek{α\texthebrew{ἰ} δείματος τ\texthebrew{ὴ}ν βασιλίδα κατεσχηκότος}, and has been generally supposed to explain the vague words of Palladius, c. 9 \emph{ad init.}, \textgreek{θρα\texthebrew{ῦ}σίν τινά γενέσθαι \texthebrew{ἐ}ν τ\texthebrew{ῶ} κοιτ\texthebrew{ῶ}νι}. Seeck (\emph{op. cit.} p582) questions the earthquake and conjectures that the \textgreek{θρα\texthebrew{ῦ}σις} was the death of one of the Imperial children (Flaccilla). I cannot think that such an occurrence would have been referred to in this way.

} The Empress herself, who was very superstitious, was panic-stricken , and she sent one of her chamberlains with a letter to Chrysostom imploring him to return. In this conciliatory letter she disclaimed all responsibility for his exile. ``Let not your Holiness suppose,'' she wrote, ``that I was privy to what has been done. I am innocent of thy blood. Wicked and corrupt men devised this plot; God to whom I sacrifice is witness of my tears. I remember that my children were baptized by thy hands. I touched the knees of the Emperor and besought him: 'We have lost the priest, let us bring him back. Unless we restore him there is no hope for the Empire.' `` Chrysostom accepted her overtures and returned. When he was back in his palace, Eudoxia sent him a verbal message. ``My prayer has been fulfilled. My success is a crown more precious than my Imperial diadem. I have received the priest, restored the head to the body, the pilot to the ship, the shepherd to the flock, the bridegroom to the bridal chamber.'' She was generous in her amends, and the archbishop, not to be outdone in generosity, paid an extravagant tribute to her in a triumphant sermon he preached the next day in St. Sophia\texthebrew{⁠}\footnote{See the Homily after his return, \emph{P.G.} 52, 443 \emph{sqq.}

} (July). His eulogy of the Empress, who seems to have been very popular, was loudly applauded.

Chrysostom desired to regularise his position by a general Council which should inquire into his case and the proceedings of the Synod of the Oak. Theophilus began to spin new intrigues, and there were bloody frays between the populace and his partisans. Not having the countenance of the court, he did not dare to remain any longer in the city, and sailed with his followers back to Egypt.\texthebrew{⁠}\footnote{Socrates, VI.17, whose story shows that Theophilus remained some time in the capital or its neighbourhood after Chrysostom's return. From Palladius, c. 9 \emph{ad init.}, and from Chrysostom's letter to Innocent (\emph{apud Pallad.} c. 2), one would suppose he had sailed almost immediately, but these narratives are very condensed. Sozomen, VIII.19, dates his departure

(p155)
at the beginning of winter, and Seeck has therefore dated the Synod of the Oak to September (\emph{op. cit.} p584), which contradicts the other evidence. I take Sozomen's statement to be an error.

} If Chrysostom had now been able to control

p155
his temper, his reconciliation with the court might have been permanent, and all might have gone smoothly. But a trivial incident occurred which betrayed him into gross impoliteness towards the Empress.

Some months after his return,\texthebrew{⁠}\footnote{Palladius, \emph{ib.}, says that \textgreek{μετ\texthebrew{ὰ} δύο μ\texthebrew{ῆ}νας πάλιν} his enemies began to recover and seek new means of deposing him. He says nothing of Eudoxia's statue, so that this incident which is related by Socrates (VI.18) may have occurred somewhat later than September.

} a silver image of Eudoxia on a tall porphyry column was erected by Simplicius, Prefect of the City, in the middle of the Augusteum, and thus close to the vestibule of St. Sophia.\texthebrew{⁠}\footnote{The stylobate was discovered in 1848, with the inscriptions, on one side, D N Aeliae Eudoxiae semper Augustae vir clarissimus Simplicius Prf. U. dedicavit; on the other the hexameters

\textgreek{κίονα πορφυρέην κα\texthebrew{ὶ ῤ}γυρέην βασίλειαν δέρκεο, \texthebrew{ἔ}νθα πόληι θεμιστεύουσιν \texthebrew{ἄ}νακτες. ο\texthebrew{ὐ}νομα δ\texthebrew{}’ ε\texthebrew{ἰ} ποθέεις, Ε\texthebrew{ὐ}δόξια· τίς δ\texthebrew{’ ἀ}νέθηκεν; Σιμπλίκιος, μεγάλων \texthebrew{ὑ}πάτων γένος, \texthebrew{ἐ}σθλ\texthebrew{ὸ}ς \texthebrew{ὑ}παρχος.}See CIL III.736 ; Paspates, \textgreek{Βυζ. \texthebrew{ἀ}νάκτορα}, p97. Cp. Marcellinus, \emph{Chron., sub a.}

} The inaugural ceremonies were of a pagan character, and accompanied by dancing and music, and the loud noise of the merriment interrupted the service in St. Sophia. Chrysostom complained to the Prefect in no measured terms, and his denunciation of the heathenish rites was taken by the Empress as a personal affront. She was an impulsive woman, and she was now ready to side with his enemies, Severian of Gabala and the rest, who were lurking for an opportunity of vengeance. Chrysostom poured fuel on the flame by a sermon which began: ``Again Herodias is furiously raging, again she is dancing, again demanding the head of John on a charger.''\footnote{Socrates, \emph{ib.} The homily which is preserved containing these words (\emph{P.G.} 59, 485) is generally considered a fabrication, but Seeck defends it (\emph{op. cit.} p583).

}

Chrysostom had demanded a general Council;\texthebrew{⁠}\footnote{His position had been provisionally regularised by a meeting or synod of about 60 bishops, who declared his condemnation illegal (Sozomen, VIII.19).

} the summonses had been sent out; but Eudoxia was now eager that the Council should be so packed with his opponents that its structure would be not to rescind but to confirm the decree of the Synod of the Oak. At Christmas she and the Emperor refused to communicate with the pastor whom she had so warmly welcomed on his return, until the approaching Council should have tried his case. Theophilus refused to attend; his experiences at Constantinople did not encourage a second visit. But many

p156
of his bishops went, and he instructed them to make use of the canon of the Council of Antioch of A.D. 341, which laid down that if a bishop who had been deposed by a synod should then appeal to the secular power his deposition should be final and irrevocable. The Council met early in A.D. 404, but many supporters of Chrysostom were present; and his enemies, who did not propose to investigate the charges against him but to condemn him by virtue of the canon of Antioch, found themselves in an awkward position. For the Council of Antioch was deeply tainted with Arianism, and the canon was aimed at Athanasius. When it was suggested to them in the Emperor's presence that if the canon was to be accepted as authoritative they must subscribe to the acts of the Council in question, they were taken aback, but for very shame they promised to subscribe. It was a promise they could not possibly fulfil, for the Council was notoriously heretical. And so the matter hung fire, while Chrysostom continued to perform his ordinary duties. But Easter (April 17) was now approaching, and representations were made to the Emperor that it was impossible to allow the ceremonies of that high festival to be celebrated by a man who had been deposed and excommunicated by a synod. He was ordered to remain in his palace and not to enter the church, but he refused to comply unless he were compelled by force.

Easter Eve was the great day for the baptism of converts, and in this year there were three thousand candidates. Large multitudes assembled in St. Sophia, many having come in from the neighbouring towns. At night the church was crowded, when a body of soldiers entered and scattered the congregation. Women and children fled shrieking through the streets, but the clergy succeeded in reassembling the congregation in the Baths of Constantine, and preparations were made to celebrate the services there. But the flock was again dispersed by soldiers. On Easter Day the devoted followers of Chrysostom would not attend the services in St. Sophia, and celebrated Easter in an open field beyond the walls.

For two months longer Chrysostom was allowed to remain in his palace, but was prevented from leaving it. Arcadius felt some compunction about proceeding to extremities. But at length he yielded to the pressure of Severian and the other bishops, who were urging him to tranquillise the city by removing the

p157
cause of scandal and disturbance, on June 20 an Imperial mandate was delivered to Chrysostom, ordering him to leave the city. He submitted, and allowed himself to be conducted stealthily to one of the harbours and conveyed in a boat to the Bithynian coast.

On the same night a fire broke out in St. Sophia. It began at the chair of the archbishop and, flaming upwards, caught the roof and turned round the building like a serpent. There was a high wind, and the flames, blown southward, caught the senate-house . Both buildings were destroyed, but the destruction of the senate-house was the greater misfortune, because it was a museum of precious works of classical art. The statues of the nine Muses were burned, but the Zeus of Dodona and the Athene of Lindus escaped.\footnote{Zosimus, V.24.6.

He (\emph{i.e.} Eunapius) considers the party of Chrysostom responsible for the fire.

}

The cause of the conflagration was made a matter of judicial inquiry. Some attributed it to Chrysostom himself, others to his friends. It was made a pretext for a bitter and cruel persecution of all his adherents.\texthebrew{⁠}\footnote{The investigation was conducted by Studius, Prefect of the City, and he acted with great severity. See Socrates, VI.18, 19; Sozomen, VIII.23, 24; Palladius, c. 20. Tigrius, a presbyter, was tortured till his bones were dislocated; Serapion was cruelly beaten, and his teeth knocked out. Olympias was harshly treated and withdrew to Cyzicus. Socrates seems to be wrong in ascribing the cruelties to Optatus, who was a pagan and succeeded Studius not before September; but Optatus may have afterwards sought to force the clergy to communicate with Arsacius. As the inquisition led to no results the imprisoned clergy were released at the end of August

(\emph{C. Th.} XVI.2.37) .

} The deaconess Olympias was treated with great harshness; she fell ill and withdrew to Cyzicus. Many persons were punished for refusing to communicate with Arsacius,\texthebrew{⁠}\footnote{His tenure of the see was brief. He died Nov. 11, 405, and was succeeded by Atticus in March 406.

} the new archbishop, who was installed a few day later (June 26). He was a brother of Chrysostom's predecessor Nectarius, and was a gentle old man, whom Chrysostom's admirers described as muter than a fish and more inert than a frog. Partaking of the communion with him was a sort of test for discovering Johannites, as the followers of Chrysostom were called.

Chrysostom lived in exile for three years, at first in Cucusus on the borders of Cappadocia and Armenia, then at Arabissus.\footnote{In common with the inhabitants of these regions he endured considerable distress and anxiety from the depredations of the Isaurians, who wasted the villages round Cucusus, and slaughtered the villagers. We learn about their doings in his letters.

}

p158
From these places he conducted an active correspondence with his friends and admirers in all parts of Christendom, and his influence was so great that his enemies thought it prudent to procure his removal to a more remote spot, Pityus on the Euxine coast. On the way thither he died from exhaustion (September 14, A.D. 407).

The treatment of Chrysostom caused fresh trouble between the courts of Constantinople and Ravenna. Theophilus had first apprised Pope Innocent I of his deposition: letters from Chrysostom himself and his clergy, delivered a few days afterwards, probably convinced him that the proceedings had been extremely irregular, and this conviction was confirmed when he received from Theophilus a memorandum of the acts of the Synod of the Oak. He decided that the matter should be brought before a general Council, and meanwhile declined to desist from communion with the Patriarch, to whom he sent a letter of consolation. An Italian Synod was summoned, and declared the condemnation of Chrysostom illegal and demanded a general Council at Thessalonica.

Honorius had already written twice to Arcadius,\texthebrew{⁠}\footnote{One of these letters is preserved, \emph{Coll. Avell., Ep.} 38 (or Mansi, III.1122), probably written in July 404. He refers in it to the criticisms which the Imperial honours conferred on Eudoxia had evoked: quamvis super imagine muliebri novo exemplo per provincias circumlata et diffusa per universum orbem obtrectantium fama litteris aliis commonuerim, etc.

} deploring the tumults and conflagrations which had disgraced Constantinople, and criticising the inconvenient haste with which the sentence against the condemned had been carried out before the decision of the head of the Church had been ascertained. He wrote under the influence of Innocent, and definitely asserted the doctrine that ``the interpretation of divine things concerns churchmen, the observation of religion concerns us (the Emperors).'' After the meeting of the Italian Synod he wrote a third letter,\texthebrew{⁠}\footnote{Quoted in Palladius, c. 3.

} to be carried by a deputation of bishops and priests, who were to inform his brother of the opinion of the Italian Church. The envoys had reason to repent of their expedition. Escorted by soldiers from Athens to Constantinople, they were not permitted to land in that city, but were thrown into a Thracian fortress, forcibly deprived of the letters they bore, and at last hardly allowed to return to Italy ( A.D. 406). As they had

p159
been specially recommended by Honorius himself to Arcadius, the outrageous treatment they received was a grievous affront to the western court. The Eastern Emperor took no notice whatever of the proposal to summon a general Council, and the Imperial brothers seem never again to have held any communications. Honorius and Innocent could do no more; they had to abandon Chrysostom to his fate.\footnote{Theophilus wrote an Apologia for his own conduct, with violent invectives against Chrysostom. He sent it to Jerome, who translated it into Latin (\emph{Epp.} 113, 114), and we have some extracts from it in the \emph{Pro defensione trium capit.} of Facundus. Chrysostom is denounced as a sacrilegious persecutor, not a Christian, but ``worse than Belshazzar,'' a blasphemer against Christ who delighted the Arians; in the next world he will suffer eternal punishment. Theophilus probably was genuinely convinced that his adversary was a very bad man.

}

The Empress Eudoxia did not live to see the later phase of the episode in which she had played a considerable part, though rather as the instrument of unscrupulous ecclesiastics than as the directress of a conspiracy against a man whose probity she certainly respected. She died on October 6, A.D. 404, of a miscarriage.\footnote{Eunapius, apud Photium, \emph{Bibl.} 77 , Date: \emph{Chron. Pasch., sub a.} A few days before her death, there was a terrible shower of hail at Constantinople (\emph{ib.}), which the populace said was a mark of divine displeasure at the persecution of Chrysostom (Socrates, VI.19); perhaps it also alarmed Eudoxia. This particular hailstorm may have been in the mind of Philostorgius when he wrote

XI.7 . It seems probable that it was just after these occurrences that an amnesty was granted to the Johannites. The date depends on a passage in

Synesius, \emph{Ep.} 66

\textgreek{τουτ\texthebrew{ὶ} μ\texthebrew{ὲ}ν \texthebrew{ἔ}τος τρίτος \texthebrew{ἐ}ξήκει μετ\texthebrew{ὰ} τ\texthebrew{ὴ}ν \texthebrew{ἀ}μνηστίαν}. Seeck has given reasons for dating the letter to the end of 407, and drawn the inference (\emph{op. cit.} 585\texthebrew{‑}586).

}

Arcadius slumbered on his throne for three and a half years after her death, and died on May 1, A.D. 408. During this time the reins of power seem to have been in the hands of Anthemius, the Praetorian Prefect of the East, who was afterwards to prove himself an able minister.\texthebrew{⁠}\footnote{He was already Prefect on July 10, 405

(\emph{C. Th.} VII.10.1) , and had been raised to the rank of Patrician before April 28, 406

(\emph{ib.} IX.34.10) . He was grandson of Philippus, Praet. Pr. in 346. He had been comes sacr. larg. in 400, then mag. off., and was consul in 405.

} One of the principal concerns of the government during these years was the condition of the southern and eastern provinces of Asia Minor, exposed to the savagery of the Isaurian brigands. Their devastations continued from A.D. 404 to 407.\texthebrew{⁠}\footnote{For 404 and the campaign of Arbazacius see

Zosimus, V.25 , who says that this incompetent commander escaped punishment by bribing Eudoxia; Eunapius, \emph{frr.} 84, 86; Marcellinus, \emph{Chron., sub} 405; Sozomen, VIII.25 (all the cities between Caria and Phoenicia devastated). Philostorgius,

XI.8 , says that they subdued Cyprus. During the following years

(p160)
their ravages can be traced in Chrysostom's letters; see Clinton, \emph{F.R. sub annis.} \texthebrew{—} For the Armenian Arbazacius cp. CIL VI.31978 . He may be the same as the \textgreek{\texthebrew{Ἀ}ρταβάζακος} in Synesius, \emph{Epp.} 134 .\texthebrew{—} Isaurians on vessels in the Cilician ports are vividly described by Ammianus,

XIV.2.1 .

\texthebrew{◂} previous

next \texthebrew{▸}Images with borders lead to more information.

The thicker the border, the more information.

(Details here.)

UP TO:

J. B. Bury
Homepage

Latin \& Greek

Texts

LacusCurtius

Home

OFF

SITE:

Catholic Encyclopedia: St. John Chrysostom} We hear of the failure of a general to suppress them

p160
at the beginning of the movement, but we are not told how this civil war was brought to an end. Anthemius had also to keep a watchful eye on Alaric and Stilicho. To them we must now return.

\xchapter{Chapter V, Part C}

\xsubsection{(Part 3 of 3)}

We saw how Alaric and his Visigoths had withdrawn from the Peloponnesus into the province of New Epirus in A.D. 397, and that Alaric had been appointed to some imperial post, probably that of Master of Soldiers in Illyricum. For four years we hear nothing of him except that he took advantage of his official position to equip his followers with modern arms from the Roman arsenals in the Dacian diocese.\texthebrew{⁠}\footnote{Claudian, \emph{B. Goth.} 537 \emph{sqq.}

Zosimus, who omits the Italian campaigns of 402 and 403, and passes from 397 to 405, as though Alaric had remained eight years quiet in Epirus, says

(V.26) : \textgreek{τ\texthebrew{ὸ} παρ\texthebrew{ὰ} Στελίχωνος \texthebrew{ἀ}νέμενε σύνθημα τοιόνδε πως \texthebrew{ὄ}ν}, namely \textgreek{τ\texthebrew{ῇ Ὁ}νωρίου βασιλεί\texthebrew{ᾳ} τ\texthebrew{ὰ ἐ}ν \texthebrew{Ἰ}λλυρίοις \texthebrew{ἔ}θνη πάντα προσθε\texthebrew{ῖ}ναι} with Alaric's help. From this point Zosimus follows Olympiodorus instead of Eunapius.

} Then suddenly he determined to invade Italy. Perhaps it was the defeat of the attempt of Gaïnas to establish a German ascendancy at Constantinople that averted his covetous eyes from the Balkan lands and moved him to seek a habitation for his people in the realm of Honorius. It can hardly have been his hope to establish a permanent kingdom in Italy itself.\texthebrew{⁠}\footnote{This, indeed, is Schmidt's view, \emph{Deutsche Stämme}, I.204.

} We may take it that his intention was rather to frighten Honorius into granting lands and concessions in the Danube provinces. An opportune moment came when, towards the end of A.D. 401, a host of Vandals and other barbarians under a savage leader named Radagaisus had broken into Noricum and Raetia.\texthebrew{⁠}\footnote{Radagaisus was a German, as his name shows; he is called a ``Scythian'' in the sources.

} Alaric passed the Italian Alps in November,\texthebrew{⁠}\footnote{\emph{Fast. Vind. pr., sub a.} (\emph{Chron. min.} I.299).

} and advanced to Aquileia, which he appears to have captured.\texthebrew{⁠}\footnote{Jerome, \emph{C. Rufin.} III.21. Claudian's words deploratumque Timavo vulnus may refer either to the capture of the city or to a battle in the neighbourhood.

} The Italians were in consternation, and not least Honorius himself, who thought of fleeing to Gaul, and was with difficulty persuaded that he was safe behind the walls of Milan.\footnote{Claudian, \emph{ib.} 296 . Prudentius has briefly described the invasion, \emph{C. Symm.} II.696 \emph{sqq.}

}

p161
During the next two months the cities of Venetia opened their gates to the Goths, and Alaric was ready to march on Milan, where he hoped to seize the Emperor's sacred person.

At the moment Italy was defenceless, because Stilicho had led his mobile troops across the Alps to drive back Radagaisus and the invaders of Raetia. This winter campaign was success ­ful. The barbarians were checked, and Stilicho induced them to furnish him with auxiliaries against the Goths.\texthebrew{⁠}\footnote{Claudian, \emph{ib.} 279 \emph{sqq.} ,

321 \emph{sqq.} ,

364 \emph{sqq.} ,

414 . Cp.
Hodgkin, \emph{Italy and her Invaders}, I.711 \emph{sqq.} , and Bury, App. 15 to Gibbon, vol. III.

} Reinforced by this accession and also by troops hastily summoned from the Rhine frontier and from Britain, he came down to relieve Milan and deliver Italy (about the end of February, A.D. 402).\texthebrew{⁠}\footnote{Symmachus, \emph{Ep.} VII.13 . Stilicho seems to have marched to Raetia by the Splügen pass

(Claudian, \emph{ib.} 320)

and returned by the Brenner

(\emph{ib.} 488) ; see Seeck, \emph{Ges. d. Untergangs}, V.573). For the forces summoned from Gaul and Britain see

Claudian, \emph{ib.} 416 \emph{sqq.} The Rhine was left defended solo terrore. The Britannic legion came from the north-west,

uenit et extremis legio praetenta Britannis

quae Scotto dat frena truci;

probably it was a detachment of the old IInd legion. Cp. Bury, \emph{The Not. Dig.}, in \emph{J. R. S.} X .

} Alaric abandoned the siege and marched westward to Hasta (Asti), which he failed to take, and then went on to Pollentia (Pollenzo) on the river Tanarus, where he decided to make a stand against the forces of Stilicho who marched in pursuit. According to the poet who celebrated this campaign, a council was held in the Gothic camp, and one of the veterans who feared the issue of a trial of strength with Stilicho besought the king to withdraw from Italy while there was yet time. Alaric indignantly refuses; he was confident that he was destined to capture Rome; and he assured the assembled warriors that a clear voice had come to him from a grove, saying penetrabis ad urbem , ``thou shalt penetrate to the City.''

The battle was fought on Easter-day (April 6). Neither side could claim a decisive victory,\texthebrew{⁠}\footnote{Prosper, \emph{sub a.}, vehementer utriusque partis clade pugnatum est. Cp. Prudentius, \emph{ib.} 717\texthebrew{‑}720 . Crees, \emph{Claudian}, 175\texthebrew{‑}180, argues for 403 as the year of the battle.

} but the Romans occupied the Gothic camp, and Alaric's family among other captives fell into their hands. The Goths descended to the Ligurian coast and marched along the coast road in the direction of Etruria.\footnote{Alaric's goal was Rome. Cp. Claudian, \emph{De VI. cons. Hon.} 483 . The verses of Prudentius, \emph{l.c.}, reflect the profound relief felt in Rome at the success of Stilicho. An illustration of this is possibly to be found in an early Christian missal, where deliverance from a foe at Eastertide is referred to. See Grisar, I.37.

}

p162
Stilicho did not attempt to overtake and crush them. He opened negotiations and Alaric agreed to leave Italy, but we do not know what conditions were made.\footnote{Schmidt (\emph{ib.} 206) is mistaken in thinking that Alaric was now created a western Mag. mil., citing Sozomen, VIII.25. This passage refers to a later date (406\texthebrew{‑}407?).

}

When he retired from Italian soil in accordance with this treaty, he remained near the borders of the peninsula, dissatisfied with a bargain which perhaps the captivity of his wife and children had chiefly moved him to accept. At the end of a year, during which Stilicho strengthened the military forces in Italy, probably at the expense of the defences of Gaul, he crossed the Italian frontier again in the early summer ( A.D. 403) and attacked Verona.\texthebrew{⁠}\footnote{This second invasion used to be placed in 402, but Birt (\emph{Praef.} to his ed. of Claudian, liii. \emph{sqq.}) determined the true date as 403. Cp. Mommsen, \emph{Hist. Sch.} I.525, n4; Bury, App. 18 to Gibbon, vol. III. The sole source is Claudian,

\emph{ib.} 201 \emph{sqq.}

} Here defeated by Stilicho, and almost captured himself, he took the northward road to the Brenner pass, pursued by the Romans. The army of the Goths suffered from hunger and disease, and seems to have been entirely at the mercy of the Roman general. But Stilicho acted once more as he had acted in Thessaly, in the Peloponnesus, and in Liguria.\texthebrew{⁠}\footnote{Orosius

(VII.37)

notes these repeated releases of Alaric: taceo de Alarico . . . saepe victo saepeque concluso semperque dimisso.

} He came to an understanding with Alaric and allowed him to take up his quarters in the border districts between Dalmatia and Pannonia, where he was to hold himself in readiness to help Stilicho to carry out the plan of annexing Eastern Illyricum.\texthebrew{⁠}\footnote{Zosimus, V.26

(Olympiodorus).

} Here he seems to have remained for some time and then to have moved again into Epirus.

The story of these two critical years in Italy can hardly be said to be known. The slight chronicle which we can construct of Alaric's invasions is drawn from rhetorical poets and the scrappy notices of chroniclers. They do not tell us the things that would enable us to judge the situation. They do not tell us the number of the Gothic warriors, or the number and composition of the Imperial forces which opposed them; they do not tell us anything of the actual course of the fighting or the tactic employed at Pollentia or at Verona; and they are silent as to the precise conditions on which Stilicho spared Alaric. We know enough, however, to see that if another than this German general had been at the head of affairs, if the

p163
defence of the provinces had been in the hands of a Roman commander possessing the ability and character of Theodosius or Valentinian I, the Visigoths and their king would have been utterly crushed, and many calamities would have been averted, which ensued from the indulgent policy of the Vandal to whom Theodosius had unwisely entrusted the destinies of Rome.

The Emperor Honorius celebrated the repulse of the invader by a triumphal entry into Rome.\texthebrew{⁠}\footnote{The walls, towers, and gates of Rome had been renovated and fortified, at the instance of Stilicho, in 402. This is recorded in the identical inscriptions over the Portae Portuensis, Praenestina, and Tiburtina, where statues of the Emperors were placed; these inscriptions (CIL VI. º 1188\texthebrew{‑}1190) are prior to Feb. or March 402, as the name of Theodosius, who was created Augustus on Jan. 10, does not appear. See Seeck, \emph{Symmachus}, p. clxxxviii.

} It was probably in the summer or autumn of A.D. 402 that, menaced by Alaric's proximity, he had moved his home and court from Milan to Ravenna,\texthebrew{⁠}\footnote{The earliest extant rescript issued at Ravenna is dated Dec. 6, 402

(\emph{C. Th.} VII.13.15) .

} and, as future events were to prove, he could not have chosen a safer retreat. But he could now venture to Rome, which he had never visited before, enjoy the celebration of a triumph,\texthebrew{⁠}\footnote{Claudian, \emph{ib.} 523 \emph{sqq.} Prudentius, \emph{ib.} 726 \emph{sqq.}

} reside in the palace of the Caesars on the Palatine Hill, and enter upon his sixth consul ­ship ( A.D. 404) in the presence of the Senate and the Roman people. For the Romans, the triumphal entry of the Emperor was an event. Rome, which had not witnessed a triumph for more than a hundred years, had in certain ways changed much since the days of Diocletian. In external appearance the transformation from ancient into medieval Rome had already begun. Most of the great churches that still exist, though rebuilt, enlarged, or restored, had been founded in the fourth century. St. John in the Lateran, the basilica of Liberius on the Esquiline which was soon to become Sta. Maria Maggiore, and outside the wall St. Peter beyond the Tiber, and St. Paul on the road to Ostia, were all probably visited by Honorius.\texthebrew{⁠}\footnote{St. Paul's and the baptistery of St. Peter's are described by Prudentius in the

\emph{Peristephanon}, XII .

} The temples of the gods stood still unharmed, but derelict; more than twenty years before the altar of Victory had been removed from the Senate-house . Some distinguished senatorial families had been converted from their errors, like the Anicii and the

p164
Bassi,\texthebrew{⁠}\footnote{Also the Paullini and the Gracchi. Prudentius eagerly enumerates them, and the shortness of his list shows that they were in a small minority.

} but the greater number of the senators were still devoted to paganism and would have welcomed a new Julian on the Imperial throne. Of these pagans the most distinguished was Symmachus, who had been their eloquent spokesman when they vainly pleaded with Theodosius and Valentinian II to permit the restoration of the altar of Victory. And now during the visit of Honorius to Rome the Christian poet Prudentius took occasion to compose a poem confuting the arguments of Symmachus and exulting over the discomfiture of his cause.\texthebrew{⁠}\footnote{\emph{Contra Symmachum}, in 2 Parts.

} He affected to believe that the senators had freely and joyfully proscribed the pagan idols, and that there were few pagans left \texthebrew{—} ingenia obtritos aegre retinentia cultus . ``The Fathers,'' he says, ``the luminaries of the world, the venerable assembly of Catos, were impatient to strip themselves of their pontifical garment, to cast the skin of the old serpent, to assume the snowy robes of baptismal innocence, and to humble the pride of the consular forces before the tombs of the martyrs.''\footnote{\emph{Ib.} I.546 \emph{sqq.} Gibbon's Paraphrase (vol. III chap. xxviii). If Claudian read the verses of Prudentius he must have smiled at the liberties which that writer takes with the quantities of Greek words (\emph{e.g.} idŏla, hĕresis, cātholicus).

}

Prudentius concluded his work with an appeal to the Emperor to suppress gladiatorial shows:\footnote{II.1114 \emph{sqq.}

}

This appeal probably expressed a considerable volume of public opinion, and if it was not in this year that exhibitions of gladiators were finally forbidden, it must have been soon afterwards. Possibly it is not a mere legend that the immediate occasion of the abolition of these spectacles was the act of an aged monk named Telemachus, who rushed into the arena of the Colosseum to separate two combatants and was killed by the indignant populace with showers of stones.\footnote{Theodoret, \emph{H.E.} V.26.

}

The occasion of the Imperial visit to Rome was celebrated by Claudian with his unflagging enthusiasm. He had already,

in a poem on the \emph{Gothic War} , sung the repulse of Alaric at Pollentia \texthebrew{—}p165
and united the name of Stilicho with that of Marius as the protectors of Italy, imagining the bones of Cimbrians and Goths laid under a common trophy with the inscription

The campaign of Verona was celebrated in the poem which he composed at the end of the year for the

\emph{Sixth Consulship of Honorius} , immediately after the triumph. This was his last work. Our records are silent as to his fate, but the most probable conjecture is that death cut short his career and that he did not live to see the second consul ­ship of his patron ( A.D. 405), a theme which he could not have neglected.\footnote{Seeck's view is that in 404 he was accused of pagan practices, and thrown into prison by the Praetorian Prefect of Italy, Rufus Synesius Hadrianus (a man of Egyptian birth, who had been com. s. larg. in 395, mag. off. 397\texthebrew{‑}399, and was Pr. Pr. 400\texthebrew{‑}405 and again 413\texthebrew{‑}416), who owed him a grudge for a biting epigram

(\emph{Carm. min.} XXI) ; his friends were tortured and banished; and in prison Claudian wrote an appeal to the Prefect for mercy, the

\emph{Deprecatio ad Hadrianum} (\emph{Carm. min.} XXII) . Stilicho basely withdrew his protection. This theory depends on the interpretation of the \emph{Deprecatio}, which Birt (Preface, xi, xii) has otherwise explained.

}

Great allowances as the historian has to make for Claudian's partiality and rhetoric, he owes him an appreciable debt and would give much to have his guidance for the last obscure and critical five years of Stilicho's career. But apart from the information which he gives us, his poetry is one of the most interesting facts of the age. He was born at Alexandria,\texthebrew{⁠}\footnote{Suidas, \emph{s.v.}; Claudian, \emph{Carmina minora},

XIX.3

(nostro Nilo),

XXII.58 . His identity with the author of the Greek \emph{Gigantomachy}, of which fragments remain, and seven Greek epigrams is generally admitted, and explains

\emph{Carm. min.} XLI (\emph{ad Probinum}) 14

et Latiae accessit Graia Thalia togae.

} and his earliest literary work was in Greek, but we may take it that he had learned Latin as a child. He saturated himself in the poetical literature of Rome from Ennius to Juvenal, and his verses abound in echoes and reminiscences. His Roman feeling for Roman traditions is not compromised or embarrassed by any allegiance to the new religion; and the statement of his contemporary Augustine that he was a stranger to the name of Christ\texthebrew{⁠}\footnote{\emph{De civ. Dei}, V.26 .

} is borne out by his poems, from which, if they were the sole monument of the time, we should not suspect the existence of Christianity.\texthebrew{⁠}\footnote{With the exception of the short, perfunctory production, \emph{De Salvatore}

(\emph{Carm. min.} XXXII) , which stands alone curiously out of place in the collection.

} In talent and technical skill he is incomparably

p166
superior to the Christian poets of the day, Prudentius and Paulinus, and through his genuine feeling for the dignity and majesty of the Empire he has succeeded in shedding a certain lustre over the age of Stilicho and Alaric.

The provinces of the Upper Danube, Raetia, Noricum, and Pannonia, were at this time still under the effective control of Roman governors, and the principal towns still flourishing centres of Roman civility. In Pannonia indeed considerable districts had been occupied by Ostrogoths, Huns, and Alans, whom Gratian and Theodosius had settled after their victories over the Gothic invaders of A.D. 380. Of these the Ostrogoths had perhaps been settled in the north-western of the four Pannonian provinces, Pannonia Prima,\texthebrew{⁠}\footnote{Cp. Schmidt, \emph{Gesch. der deutschen Stämme}, I.115. Alaric's wife belonged to one of the Ostrogothic families of this colony.

} and it is probable that the north-eastern, Valeria, was occupied by the Huns.\footnote{Pannonia at this time consisted of four provinces: (1) the north-western, Pannonia Prima, including the towns of Vindobona (Vienna) and Carnuntum (Petronell), Savaria (Stein-am\texthebrew{‑}Anger), Scarpantia (Odenburg); (2) the south-western, Savia: chief town, Siscia on the Save; (3) the north-eastern, Valeria, bounded on north and east by the Danube and including Sopianae, Aquincum (Alt-Ofen), Brigetio (O-szöny), and Intercisa (Dunapentele); and (4) the south-eastern, Pannonia Secunda, including the regions of the lower Drave, in which the chief towns were Sirmium, Mursa (Eszeg), Ciballae (Vinkovce). Noricum was divided into two provinces, (1) Noricum Ripense: chief towns, Lauriacum and Commagenae (near Tulln) on the Danube, Juvavum and Ovilava (Wels); (2) Noricum Mediterraneum: Teurnia (near Spittal), Virunum, Aguntum, Celeia, and Poetovio. Both Noric provinces and Pannonia I were governed by praesides; Pannonia II by a consular, Savia by a corrector.

}

The line of division between Pannonia and Noricum ran from the neighbourhood of Tulln on the Danube to Pettau, while the course of the Aenus (Inn) formed the western boundary of Noricum, separating it from Raetia.\texthebrew{⁠}\footnote{In

\emph{Not. dig.}

Valeria has no civil governor. The Huns remained in Pannonia for 45 years or more (380 to 426 or 427); see Marcellinus, \emph{Chron. sub} 427.

} The most northerly point in the course of the Danube, which was the northern border of Raetia, was marked by Batava Castra (Ratisbon), and the province extended westward to the source of that river.\texthebrew{⁠}\footnote{For the western boundary of Raetia, see Jung, \emph{Römer und Romanen}, p30. There were two Raetian provinces: (1) Prima, the southern: chief town, Curia (Chur); (2) Secunda, the northern: chief town, Aug. Vindelicorum. Each had both a civil and a military governor, a praeses and a dux. On all these Danubian lands see Jung, \emph{op. cit.}, and Zeiller, \emph{Les Origines chrét. dans les prov. dan.}

} The

p167
most important highway from Italy to Raetia was the Via Claudia Augusta, which led through the Tirol by Meran and Vintschgau to Augusta Vindelicorum (Augsburg); the Brenner road was less used. Aquileia was the great centre of roads leading from Italy into Noricum, Pannonia, and the Balkan lands. The traveller to Pannonia would proceed from Aquileia to Celeia (Cilly) and Poetovio (Pettau), whence the high road continued to Savaria ( Stein-am\texthebrew{‑}Anger ) where several roads met, one leading northward to Carnuntum (Petronell), a second north-eastward, and a third south-eastward to Sopianae (Fünfkirchen). Three roads led from Aquileia over the Julian Alps: (1) to Aguntum (near Lienz); (2) to Virunum (Maria Saal near Klagenfurt), whence roads led to Juvavum (Salzburg) and to Lauriacum (Lorsch) and other places on the Danube, and (3) to Emona (Laibach), which belonged administratively to Venetia and was itself connected by a road over the mountains to Virunum. Here at Emona the two roads met of which one led into northern Pannonia, as we saw, by Celeia, and other through southern Pannonia along the valley to the Save, by Siscia (Siszek) to Sirmium (Mitrovica) and Singidunum (Belgrade), and thence to Constantinople. It should be observed that Pannonia was bounded on the south by the province of Dalmatia, for Dalmatia then included not only the coastlands of the Hadriatic as far south as Alessio, but also the lands which were afterwards to be known as Bosnia and Herzegovina, and a part of Istria, west of the river Arsia.

During the early years of Honorius, the defence of the Pannonian frontier was almost abandoned, and the Pannonian provinces suffered both from the barbarians who were within,\texthebrew{⁠}\footnote{Jerome, \emph{Epp.} 123.15 hostes Pannonii.

} and from those who were without. Of all this devastation we have no regular story; we have only the vague complaints and hints of contemporary writers.\texthebrew{⁠}\footnote{\emph{Id.} \emph{Epp.} 60.16 (A.D. 396); Prudentius, \emph{C. Symm.} II.716 ; Claudian,

\emph{In Ruf.} II.45 ;

\emph{De cons. Stil.} II.191 \emph{sqq.}

} But the alarm, even in those much tried lands, must have been great when in the last months of A.D. 405 a vast host of Germans, principally Ostrogoths, descended upon Italy.\texthebrew{⁠}\footnote{The chief sources are Olympiodorus, fr. 9;

Zosimus, V.26 ,

Augustine, \emph{De civ. Dei}, V.23 ;

Orosius, VII.37 ; Paulinus, \emph{Vit. Ambrosii}, 50. We may conjecture that the number of the invaders did not exceed 50,000. The huge figures of Zosimus and Orosius, 400,000 and 200,000, are absurd; Augustine's ``more than 100,000'' must also be a gross exaggeration.\texthebrew{—} With Gothofredus and Seeck, I have followed \emph{Fast. Vind. pr.}

(p168)
p299 and Marcellinus \emph{sub} 406 in pla­cing the invasion in 405 and the defeat in 406, contrary to the general view, which, on the authority of Prosper, places them a year earlier. The later date supplies the motive of the two constitutions issued at Ravenna in April 406

(\emph{C. Th.} VII.13.16 and 17) , calling for volunteers to defend the provinces against invaders in an emergency (pro imminentibus necessitatibus). Gothofredus,

\emph{C. Th.} II.389 .

} They were led by the adventurer

p168
Radagaisus, who had been repulsed from Raetia by Stilicho a few years before. As the home of the Ostrogothic people was still in the neighbourhood of the river Dniester, they had a long march by whatever route they came, and it may be presumed that they crossed the Danube on the Pannonian frontier. We are told nothing of their doings in the Danubian provinces, or by what roads they reached Aquileia, and its seems probable that Radagaisus, wishing to surprise Italy, did not tarry on his way to plunder the cities of Pannonia and Noricum. But we are told that the inhabitants of the districts through which they passed fled before them, seeking the refuge of Italy.\texthebrew{⁠}\footnote{Cp. \emph{C. Th.} X.10.25 .

} Italy was entered without resistance, and the barbarian host overran the northern provinces. After some time it is said that they divided into three companies,\texthebrew{⁠}\footnote{We hear nothing of the other two.

} of which the chief under Radagaisus attacked Florence. Stilicho, who had collected his forces at

Ticinum , numbering perhaps less than 20,000 comitatenses ,\texthebrew{⁠}\footnote{30 \textgreek{\texthebrew{ἀ}ριθμοί} (Zos. \emph{ib.}) The numerus might vary between 300 and 900 men.

} reinforced by Alans and Huns from beyond the Danube,\texthebrew{⁠}\footnote{Huns under Uldin, whose seat was north of the lower Danube.

} compelled him to withdraw to Fiesole. The Romans were able to cut off the supplies of the barbarians and then massacre them at their pleasure.\texthebrew{⁠}\footnote{Augustine declares that not a single Roman was wounded. Probably the work was done by the Huns.

} Radagaisus was captured and executed (Aug. 23, A.D. 406), and the victory, which was fondly declared to have extinguished the Gothic nation for ever, was celebrated by a triumphal arch in Rome.\texthebrew{⁠}\footnote{CIL VI.1196 . The arch was adorned with the statues of the Emperors (toto orbe victoribus) and trophies ad perenne indicium triumphorum quod Getarum nationem in omne aevum docuere extingui. In another inscription the services of Stilicho were recorded, but his name was erased after his fall (\emph{ib.} 31987; cp. 31988). It was perhaps in the Jan. of this year, while the Goths were wasting Italy, that Paulinus of Nola wrote his \emph{Poema} XXVI. For the date of the victory see \emph{Consul. Ital.} (\emph{Chron. min.} I.299). According to Olympiodorus (\emph{fr.} 9) Stilicho enrolled 12,000 Goths in the Roman army. Probably they did not belong to the band which attacked Florence.

Thayer's Note: For further details and sources on the arch, see the article

Arcus Arcadii Honorii et Theodosii

in Platner and Ashby's \emph{Topographical Dictionary of Ancient Rome}.

} But Italy must have suffered terribly, for the barbarians had been six months in the land.

It is clear from the meagre records of this invasion that when Radagaisus surprised Italy, the field army at the disposal of

p169
Stilicho was so small that he could not venture on a battle with the superior forces of the enemy until he had obtained help from the Huns. It is possible that some of the troops which had come from Gaul and Britain to oppose Alaric had been sent back, but, if so, the Gallic legionaries of the Rhine frontier must have again been summoned to fight against Radagaisus, and must have been retained. For the Rhine was virtually undefended at the end of A.D. 406, when hosts of Germans crossed the river and began a progress of destruction through Gaul. This event was decisive for the future history of Western Europe, though the government of Ravenna had little idea what its consequences would be. But Stilicho was at least bound to hasten to the rescue of the Gallic provincials. Instead of doing this, he busied himself ( A.D. 407) with his designs on Illyricum which the invasion of Radagaisus had compelled him to postpone. The unfriendliness which had long existed between the eastern and western courts came to a crisis when the ecclesiastics whom Honorius had sent to remonstrate with his brother on the treatment of Chrysostom were flung into prison.\texthebrew{⁠}\footnote{See above,

p158 .

} It was a sufficient pretext for Stilicho to close the Italian ports to the ships of the subjects of Arcadius and break off all intercourse between the two realms.\texthebrew{⁠}\footnote{The measure is referred to in

\emph{C. Th.} VII.16.1 .

} Alaric was warned to hold Epirus for Honorius; and Jovius was appointed, in anticipation, Praetorian Prefect of Illyricum.\texthebrew{⁠}\footnote{Sozomen, VIII.25 = IX.4.

} Stilicho was at Ravenna, making ready to cross the Hadriatic, when a report reached him that Alaric was dead. It was false, but it caused delay; and then came the alarming news that a certain Constantine, a soldier in Britain, had been proclaimed Emperor and had crossed over to Gaul. Once again the design of Stilicho was thwarted. He might look with indifference on the presence of barbarian foes in the provinces beyond the Alps, but he could not neglect the duty of devising measures against a rebel.\footnote{The course of events in Gaul will be related in the following chapter.

}

Alaric cared not at all for the difficulties of his paymaster, and chafed under the intolerable delay. Early in A.D. 408, threatened perhaps by preparations which the eastern government was making to defend Illyricum,\texthebrew{⁠}\footnote{Cp. \emph{C. Th.} XI.17.4

(April 11, 408) = XV.1.49

(April 9, 412, false date). Cp. Seeck's conjectures, \emph{op. cit.} V.591.

} he marched northward,

p170
and followed the high road from Sirmium to Emona. He halted there and instead of marching across the Julian Alps to Aquileia and Italy, he turned northwards by the road which led across the Loibl Pass to Virunum.\texthebrew{⁠}\footnote{Cp. Jung, \emph{Römer und Romanen}, p120.

} Here in the province of Noricum he encamped, and sent an embassy to Rome demanding compensation for all the trouble he had taken in the interest of Honorius. 4000 pounds of gold (£180,000) was named. The Senate assembled, and Stilicho's influence induced it to agree to the monstrous demand; but many were dissatisfied with a policy which played into the hands of the barbarians, and one senator bolder than the rest exclaimed, ``That is not a peace; it is a compact of thraldom.'' Such, however, was the power of the Emperor's father-in\texthebrew{‑}law , and such the awe in which he was held, that the rash speaker after the dissolution of the assembly deemed it prudent to seek refuge in a church. The money was paid to Alaric, and he was retained in the service of Honorius. Perhaps he might be employed against the usurper in Gaul.

But Stilicho's position was not so secure as it seemed. His daughter, the Empress Maria, was dead, but Honorius had been induced to wed her sister Aemilia Materna Thermantia,\texthebrew{⁠}\footnote{Early in 408.

Zosimus, V.28 ; Olympiodorus, \emph{fr.} 2. For her full name see CIL XV.7152 . The marriage was arranged through the efforts of Serena. We do not know when Maria died. She seems to have been buried at Rome in a porphyry sarcophagus, in St. Peter's. A remarkable golden bulla was found in the sarcophagus with the inscription:

Honori, Maria, Stelicho, º Ser\textgreek{η}na vivatis!

Stelicho, Serena, Thermantia, Eucheri vivatis!

See Dessau, 800.

Thayer's Note: For comprehensive details on her tomb, see the article

Sepulcrum Mariae

in Platner and Ashby's \emph{Topographical Dictionary of Ancient Rome} and the passages of Lanciani, Armellini and Hülsen linked there.

} and Stilicho might think that his influence over the Emperor was impregnable and still hope for the union of his son with Placidia. But any popularity he had won by the victory over Gildo, by the expulsion of Alaric from Italy, by the defeat of Radagaisus was ebbing away. The misfortunes in Gaul, which had been occupied by a tyrant and was being plundered by barbarians, were attributed to his incapacity or treachery, and his ambiguous relations with Alaric had only resulted in a new danger for Italy. It was whispered that his design on Eastern Illyricum only covered the intention of a triple division of the Empire, in which his own son Eucherius should be the third Imperial colleague. Both he and his wife Serena were detested by the pagan families of Rome who still possessed predominant influence in the capital. Nor

p171
was his popularity with the army secure. While he and Honorius were at Rome in the spring of A.D. 408, a friend warned him that the spirit of the troops stationed at Ticinum was far from friendly to his government.

Honorius had reached Bononia, on his way back to Ravenna, when the news of his brother's death arrived (May). He entertained the idea of proceeding to Constantinople to protect the interests of his child nephew Theodosius, and he summoned Stilicho for consultation. Stilicho dissuaded him from this plan, urging that it would be fatal for the legitimate Emperor to leave Italy while a usurper was in possession of Gaul; and he undertook to travel himself to the eastern capital; during his absence there would be no danger from Alaric, if he were given a commission to march against Constantine. The death of Arcadius had presented to Stilicho too good an opportunity for prosecuting his design on Illyricum to be lost. Honorius agreed, and official letters were drafted and signed, to Alaric instructing him to restore the Emperor's authority in Gaul, and to Theodosius regarding Stilicho's mission to Constantinople.

The Emperor then proceeded to Ticinum, and there a plot was woven for the destruction of the power ­ful and unsuspicious minister. Olympius, a palace official, who had opportunities of access to Honorius on the journey, let fall calumnious suggestions that Stilicho was planning to do away with Theodosius and place his own son on the eastern throne. At Ticinum he sowed the same suspicions among the troops, who were discontented and mutinous. His efforts brought about a military revolution, in which nearly all the highest officials who were in attendance on the Emperor, including the Praetorian Prefects of Italy and Gaul, were slain (August 13).\footnote{The list is: Limenius, Pr. Pr. of Gaul (who had come to Italy to escape from Constantine); the Master of Offices, the Quaestor, the Comes s. larg., the Comes r. priv., the two comites domest., and Longinianus, Pr. Pr. of Italy

(Zos. V.32) . Date: \emph{Cons. Ital.} p300.

}

The first thought of Stilicho, when the confused story of these alarming occurrences reached him at Bononia and it was doubtful whether the Emperor himself had not been killed, was to march at the head of the barbarian troops who were with him and punish the mutineers. But when he was reassured that the Emperor was safe, reflexion made him hesitate to use the barbarians against Romans. His German followers, conspicuous

p172
among them Sarus the Goth, were eager to act and indignant at the change of his resolve. He went himself to Ravenna, probably to assure himself of the loyalty of the garrison; but Honorius, at the instigation of Olympius, wrote to the commander instructions to arrest the great Master of Soldiers. Stilicho under cover of night took refuge in a church, but the next day allowed himself to be taken forth and imprisoned on the assurance that the Imperial order was not to put him to death, but to detain him under guard. Then a second letter arrived, ordering his execution. The foreign retainers of his household, who had accompanied him to Ravenna, attempted to rescue him, but he peremptorily forbade them to interfere and was beheaded (August 22, A.D. 408). His executioner, Heraclian, was rewarded by the post of Count of Africa. His son Eucherius was put to death soon afterwards at Rome, and the Emperor hastened to repudiate Thermantia, who was restored a virgin to her mother. The estates of the fallen minister were confiscated as a matter of course. There had been no pretence of a trial, his treason was taken for granted, but after his execution there was an inquisition to discover which of his friends and supporters were implicated in his criminal designs. Nothing was discovered; it was quite clear that if Stilicho meditated treason he had taken no one into his confidence.\footnote{Stilicho's designs for the advantage of his son were not necessarily treasonable, but the suspicion of treason is not confuted by the fact that Eucherius only held insignificant posts. Cp. Zos. V.34.7 .

}

The fall of Stilicho caused little regret in Italy. For thirteen and a half years this half-Romanised German had been master of western Europe, and he had signally failed in the task of defending the inhabitants and the civilisation of the provinces against the greedy barbarians who infested its frontiers. He had succeeded in driving Alaric out of Italy, but he had not prevented him from invading it. He had annihilated the host of Radagaisus, but Radagaisus had first laid northern Italy waste. It was while the helm of state was in his hands that, as we have yet to see, Britain was nearly lost to the Empire, and Gaul devastated far and wide by barbarians who were presently to be lords in Spain and Africa. The difficulties of the situation were indeed enormous; but the minister who deliberately provoked and prosecuted a domestic dispute over the government of Eastern Illyricum, and allowed his

p173
policy to be influenced by jealousy of Constantinople, when all his energies and vigilance were needed for the defence of the frontiers, cannot be absolved from responsibility for the misfortunes which befell the Roman state in his own lifetime and for the dismemberment of the western realm which soon followed his death. Many evils would have been averted, and particularly the humiliation of Rome, if he had struck Alaric mercilessly \texthebrew{—} and Alaric deserved no mercy \texthebrew{—} as he might have done more than once, and as a patriotic Roman general would not have hesitated to do. The Roman provincials might well feel bitter\texthebrew{⁠}\footnote{Orosius, \emph{H.E.} VII.37 and 38 . Here (as in \emph{Chron. Gall.} 55, p652) Stilicho is accused of having stirred up the barbarians against Gaul. The charge must be rejected, but it illustrates the general feeling that his policy was to blame for many of the disasters of the time. The fiercest attack on Stilicho is that of Rutilius Namatianus,

\emph{De reditu suo}, II.41 \emph{sqq.} , who designates him as proditor arcani imperii, but sees his chief crime in the burning of the Sibylline Books, aeterni fatalia pignora regni. The poet consigns him to Nero's place in Tartarus:

omnia tartarei cessent tormenta Neronis,

consumat Stygias tristior umbra faces.

} over the acts and policy of this German, whom the unfortunate favour of Theodosius had raised to the supreme command. When an Imperial law designated him as a public brigand who had worked to enrich and excite the barbarian races, the harsh words probably expressed the general opinion.\footnote{\emph{C. Th.} IX.42.22

(Nov. 22, 408).

}

The death of the man who had been proclaimed a public enemy at Constantinople altered the relations between the two Imperial governments. Concord and friendly co-operation succeeded coldness and hostility. The edict which Stilicho had caused Honorius to issue, excluding eastern traders from western ports, was rescinded. The Empire was again really as well as nominally one.\texthebrew{⁠}\footnote{The fact that after 396 Arcadius and Honorius never assumed the consul­ship together is significant, and illustrates the bad relations between the two courts.

\texthebrew{◂} previous

next \texthebrew{▸}Images with borders lead to more information.

The thicker the border, the more information.

(Details here.)

UP TO:

J. B. Bury
Homepage

Latin \& Greek

Texts

LacusCurtius

Home} The Romans of the west, like the Romans of the east, had shown that they did not wish to be governed by men of German race, and the danger did not occur again for forty years.

\xchapter{Chapter VI}

The fall of Stilicho was the signal for the Roman troops to massacre with brutal perfidy the families of the barbarian auxiliaries who were serving in Italy. The foreign soldiers, 30,000 of them, straightway marched to Noricum, joined the standard of Alaric, and urged him to descend on Italy.\texthebrew{⁠}\footnote{The number 30,000 is open to some suspicion. For if this army joined Alaric's forces (say 15,000 or 20,000 in invading Italy, the invaders would have been at least 45,000 strong; and we are told that Alaric, when he was reinforced by fugitive slaves after the siege of Rome (see below,

p177 ), was 40,000 strong. Possibly 30,000 does represent the total of the barbarian troops, but only some of them joined Alaric. In any case these numbers are useful in illustrating the strength of the Visigothic host (see above,

p105 ).

} Among the few who remained faithful to Honorius were the Goth Sarus and his followers.

The general conduct of affairs was now in the hands of Olympius, who obtained the post of Master of Offices. He was faced by two problems. What measures were to be taken in regard to Constantine, the tyrant who was reigning in Gaul? And what policy was to be adopted towards Alaric, who was urgently demanding satisfaction of his claims, in Noricum? The Goth made a definite proposal, which it would have been wise to accept. He promised to withdraw into Pannonia if a sum of money was delivered to him and hostages were interchanged. The Emperor and Olympius declined, but took no measures for defending Italy against the menace of a Gothic invasion.\footnote{For the following events the chief sources are Olympiodorus, frags. 3, 4, 6, 8, 10, 13;

Zosimus, V.36 \emph{sqq.} ; Sozomen, IX. \emph{sqq.} (both these writers used Olympiodorus);

Philostorgius, XII.3 ;

Orosius, VII.38\texthebrew{‑}40 .

}

p175
Alaric acted promptly. In the early autumn of A.D. 408 he crossed the Julian Alps, and entered Italy for the third time. He marched rapidly and unopposed, by Cremona, Bononia, Ariminum, and the Flaminian Way, seldom tarrying to reduce cities,\texthebrew{⁠}\footnote{Narnia is the only case recorded (see below). As this town blocked the Flaminian Way, and Alaric failed to take it, we may guess that, having turned off from that road, he approached Rome by the Via Salaria.

} for this time his goal was Rome itself. The story was told that a monk appeared in his tent and warned him to abandon his design. Alaric replied that he was not acting of his own will, but was constrained by some power incessantly urging him to the occupation of Rome. Here we have, in another form, the same \emph{motif} of Alaric's belief in his destiny to capture the City \texthebrew{—} penetrabis ad Urbem \texthebrew{—} to which Claudian ascribed his resolve to risk battle at Pollentia.

At length he encamped before the walls of Rome\texthebrew{⁠}\footnote{Probably in October, as Seeck argues (\emph{op. cit.} V.593\texthebrew{‑}594). For Honorius was still at Milan on Sept. 24

(\emph{C. Th.} IX.42.10) , but at Ravenna during the siege

(Zosimus, V.37) .

} and hoped soon to reduce by blockade a city which had made no provision for a siege. His hopes were well founded. The Senate was helpless and stricken with fear. One of their first acts shows the extremity of their panic. Serena, the widow of Stilicho, lived in Rome, and, as Stilicho's collusive dealings with Alaric were well known, it was suspected that she had an understanding with the Goth and might betray the city. They decided to put her to death, calculating that Alaric, learning that he had no ally within to open the gates to him, would abandon the siege. The fact that she was the niece of the great Theodosius did not save her; she was strangled; and it is said that her cousin, the Emperor's sister, Galla Placidia, approved of the cruel act, which was based on the merest, and perhaps unfounded, suspicion.\texthebrew{⁠}\footnote{Should we assign to this year the bronze tablet with D. n. Gallae Placidiae n. p. (\emph{i.e.} nobilissimae puellae)? CIL XV.7153 .

} The pagan historian who records it acquits Serena of any thought of treachery, but regards her fate as a divine punishment for a sacrilege which she had committed many years before. The story is that when Theodosius closed the temples of Rome, Serena, moved by curiosity, visited the temple of the Great Mother,\texthebrew{⁠}\footnote{It was on the Palatine.

Thayer's Note: For detailed information and sources, see the article

Aedes Magnae Matris

in Platner and Ashby's \emph{Topographical Dictionary of Ancient Rome}.

} and seeing a necklace on the neck of the goddess took it off and hung it round her own. An aged Vestal virgin who had accompanied her cried shame on the impiety, and when

p176
Serena ordered her to be removed imprecated curses upon her, her husband, and children. To the pagans it seemed a fitting retribution that the neck which had worn the necklace of Rhea should feel the cord of the executioner.

The death of Serena did not change the plans of Alaric. He hindered provisions from coming up the Tiber from Portus, and the Romans were soon pressed by hunger and then by plague. The streets were full of corpses. Help had been expected from Ravenna, and as none came the Senate at length decided to negotiate. There was a curious suspicion abroad that the besieging army was led not by Alaric himself but by a follower of Stilicho who was masquerading as the Gothic king. In order to assure themselves on this point, the Senate chose as one of the envoys John, the chief of the Imperial notaries, who was personally acquainted with Alaric. The envoys were instructed to say that the Romans were prepared to make peace, but that they were ready to fight and were not afraid of the issue. Alaric laughed at the attempt to terrify him with the armed populace of Rome, and informed them that he would only desist from the siege on the delivery of all the gold, silver, and movable property in the city and all the barbarian slaves. ``What will be left to us?'' they asked. ``Your lives,'' was the reply.

The pagan senators of Rome attributed the cruel disaster which had come upon them to the wrath of the gods at the abandonment of the old religion. The blockade, continued a few days longer, would force them to accept Alaric's cruel terms; the only hope lay in reconciling the angry deities, if perchance they might save the city. Encouraging news arrived at this time that in the Umbrian town of

Narnia , to which Alaric had laid siege on his march, sacrifices had been performed and miraculous fire and thunder had frightened the Goths into abandoning the siege. The general opinion was that the same means should be tried at Rome. The Prefect of the City, Pompeianus, thought it well that the Christians should share in the responsibility for such a violation of the laws and he laid the matter before the bishop, Innocent I.\texthebrew{⁠}\footnote{A.D. 402\texthebrew{‑}417.

} The Pope is said to have ``considered the safety of the city more important than his own opinion, and to have consented to the \emph{secret} performance of the necessary rites. But the priests said that the rites would not avail unless they

p177
were celebrated publicly on the Capitol in the presence of the Senate, and in the Forum. Then the half-heartedness of the Roman pagans of that day was revealed. No one could be found with the courage to perform the ceremonies in public.\footnote{Zosimus V.40.

Sozomen does not refer to the alleged consent of Innocent. The statement in the \emph{Vit. Melaniae iun.}, published by Surius, I. p769, that a Prefect was slain by the people praetextu penuriae panis at a time when barbarians were devastating the neighbourhood is referred by Tillemont (\emph{Hist.} V.569) to Pompeianus. The incident is not mentioned in the older \emph{Life} (\emph{Anal. Boll.} VIII. p34), but the arrival of Alaric at Rome shortly after Melania departed for Africa is noticed.

}

After this futile interlude, nothing remained but, in a chastened and humble spirit, to send another embassy to Alaric and seek to move his compassion. After prolonged negotiations he granted tolerable terms. He would depart, without entering the city, on receiving 5000 pounds of gold (about £225,000), 30,000 of silver, 4000 silk tunics, 3000 scarlet-dyed skins, and 3000 pounds of pepper, and the Senate was to bring pressure to bear on the Emperor to conclude peace and alliance with the Goths. As the treasury was quite empty, and the contributions of the citizens fell short of the required amount of gold and silver, the ornaments were stripped from the images of the gods, and some gold and silver statues were melted down, to make up the ransom of the city. Before delivering the treasure to Alaric, messengers were despatched to Ravenna to obtain the Emperor's sanction of the terms and his promise to hand over to Alaric some noble hostages and conclude a peace. Honorius agreed, and Alaric duly received the treasures of Rome. He then withdrew his army to the southern borders of Etruria to await the fulfilment of the Emperor's promise (December A.D. 408). The number of his followers was soon increased by the flight from Rome of a multitude of the barbarian slaves, whose surrender he had formerly demanded. They flocked to his camp, and it is said that his host, thus reinforced, was 40,000 strong.

The year came to an end, Honorius entered upon his eighth consul ­ship,\texthebrew{⁠}\footnote{His colleague was his nephew, the Emperor Theodosius II.

} and through the influence of Olympius, who was engaged in tracking down the friends and adherents of Stilicho, nothing was done to carry out the engagements to Alaric. The Goth grew impatient, Rome feared another attack, and the Senate sent three distinguished men to Ravenna to urge the government to send the hostages demanded by Alaric and

p178
compose a peace. One of these envoys was Priscus Attalus,\texthebrew{⁠}\footnote{Attalus was a pagan and had been a friend of Symmachus; eleven short letters addressed to him are preserved in the correspondence of Symmachus (\emph{Epp.} VII.15\texthebrew{‑}25). Seeck (Symm. \emph{Opp.} p. clxxi) thinks that he was son of the Ampelius who was Prefect of Rome in 370\texthebrew{‑}372. The portrait of Attalus on his medallions confirms his Greek origin.

} who belonged to a family of Ionia. The embassy was unsuccessful, but Attalus was appointed to the position of count of the Sacred Largesses, and his colleague Caecilian to that of Praetorian Prefect of Italy (January 16\texthebrew{‑}20, A.D. 409).\texthebrew{⁠}\footnote{The date is from \emph{C. Th.} XVI.5.46 , and

IX.2.5 .

} It was recognised, however, that something must be done to protect Rome, and a force of six thousand men were brought over from Dalmatia and sent to serve as a garrison in the menaced city. On the march thither they were intercepted by Alaric and almost all killed or captured. Attalus, who accompanied them, escaped. The Senate then sent another embassy, including as the principal delegate the bishop of Rome himself.

Before the siege of Rome Alaric had sent a message to his wife's brother, Athaulf, who was then in Pannonia, to join him in Italy. Athaulf with a force of Goths and Huns now crossed the Alps and marched to Etruria. Olympius collected some troops and sent them to intercept the new-comers . There was an engagement near Pisa, in which 300 Huns were said to have slain 1100 Goths, losing themselves only 17 men. But the success was not followed up, and the failure to hinder Athaulf from joining Alaric gave the enemies of Olympius, among whom were the eunuchs of the Palace, an opportunity to compass his fall. He fled to Dalmatia, and Jovius, his most formidable opponent, was created a patrician and appointed to the office of Praetorian Prefect of Italy.\texthebrew{⁠}\footnote{Feb. or March 409. Jovius was Pr. Pr. before April 1

(\emph{C. Th.} II.8.25) , but not before Feb. 1

(\emph{C. J.} II.4.7) .

} The first thing to be done was to induce the Emperor to remove adherents of Olympius who were in command of the military forces, and Jovius brought this about by secretly organising a meeting of the soldiers at Classis. The mutineers clamoured for the heads of the Masters of Soldiers, and Honorius was terrified into superseding them.\footnote{The changes in the military commands between August 408 and April 409 seem to have been as follows. After the death of Stilicho mag. utr. mil., Varanes became mag. ped., and Turpilio mag. equit.; while Vincentius and Salvius comites domesticorum equit. et ped. (see Mommsen, \emph{Hist. Schr.} I.552, \emph{note} 1, on the interpretation of

Zos. V.32 ) were succeeded by Vigilantius and Valens. In the following months there was a rearrangement: Varanes is deposed and succeeded by Turpilio; whose place is taken by Vigilantius,

(p179)
and his by Hellebich. Finally in March 409 Valens replaces Turpilio, and Hellebich Vigilantius. See Mendelssohn on Zos. V.47, p288 . Shortly afterwards, apparently Hellebich is removed, and Valens becomes, like Stilicho, mag. utr. mil. (Olympiodorus, \emph{fr.} 13). Just after the fall of Stilicho it was an obvious measure of policy to restore the old system of two \emph{co-ordinate} magistri. Mommsen, however (\emph{ib.} 557), questions the accuracy of the statements of Zosimus.

}

p179
Jovius, who had been a guest friend of Alaric, was anxious to bring about peace, and for this purpose he arranged an interview at Ariminum. The Goth demanded that the provinces of Venetia, Istria, Noricum, and Dalmatia should be ceded to him and his people as foederati, and that a certain annual supply of corn º and a money stipend should be granted. In his report of these demands to Honorius, Jovius suggested that Alaric might relax their severity if the honorary rank of Master of Both Services were conferred on him. But Honorius would not entertain the idea of bestowing on the barbarian or any of his kin an Imperial dignity; and he refused to grant the lands in which the Goths desired to settle.

Jovius opened the Emperor's answer in the presence of the king and read it aloud. The German deeply resented the language in which it was couched, and rising up in anger he ordered his barbarian host to march to Rome to avenge the insult which was offered to himself and all his kin. But in the meantime the government had been engaged in military preparations, and a large body of Huns had come to their assistance. And the food of the Goths was running short. Considering all things, Alaric thought it worth while to offer more moderate terms. Innocent, the bishop of Rome, which the Goths again threatened, was sent as an envoy to Ravenna, to press the Emperor to pause ere he exposed the city which had ruled the world for more than four hundred years to the fury of a savage foe. All that Alaric asked now was the two Noric provinces; he did not ask for Venetia nor yet for Dalmatia. Give the Goths Noricum and grant them annual supplies of grain; in return, they will fight for the Empire, and Italy will be delivered of their presence. Hard as it would have been to have had these barbarians so close to the threshold of Italy, it might have been better to have accepted these conditions. But Jovius, instead of advising peace, which he had desired before, advised a firm refusal. It appears that Honorius had taken him to task for his disposition to yield to Alaric at Ariminum, and that, fearing

p180
for his personal safety, he had leaped to the other extreme, and swore, and made others swear, by the head of the Emperor \texthebrew{—} a most solemn oath\texthebrew{⁠}\footnote{More binding, Jovius asserted, than an oath by Heaven,

Zos. V.50 .

} \texthebrew{—} to war to the death with Alaric. Honorius himself swore to the same effect.

Having met with this new refusal, Alaric marched to Rome (towards the end of A.D. 409) and called upon the citizens to rally to him against the Emperor. When this invitation was declined, he occupied Portus and blockaded the city for the second time. The corn stores lay at Portus, and he threatened that if the Senate did not comply with his demands he would use them for his own army. The Romans had no desire to submit again to the tortures of famine and they decided to yield. Alaric's purpose was to proclaim a new Emperor, who should be more pliable to his will than Honorius. He selected Priscus Attalus, the Prefect of the City,\texthebrew{⁠}\footnote{Attalus was appointed to this post at the time of the fall of Olympius.

} who was ready to play the part, and the Senate consented to invest him with the purple and crown him with the diadem. Attalus permitted himself to be baptized into the Arian religion by a Gothic bishop, but he had no thought of playing the part of a puppet. He and Alaric hoped each to use the other as a tool.\footnote{He seems to have given hostages to Alaric, one of whom perhaps was Aetius. See Merobaudes, \emph{Panegyr.} II.127 \emph{sqq.} (pignusque superbi foederis et mundi pretium fuit) and \emph{Carm.} IV.42 \emph{sqq.}; Renatus Frigeridus, in Gregory of Tours, \emph{H.F.} II.8 tribus annis Alarico obsessus.

}

It was evidently a condition of the arrangement that Alaric should receive a military command. He was appointed Master of the Foot,\texthebrew{⁠}\footnote{But with the title Master of Both Services, Sozomen, IX.8. See

Zos. VI.7 .

} while the Mastership of the Horse was entrusted to a Roman. His brother-in\texthebrew{‑}law Athaulf was appointed Count of the Domestics.\texthebrew{⁠}\footnote{\emph{Sc.} of the cavalry. His colleague, too, was probably a Roman.

} Lampadius, the same senator who had in the days of Stilicho protested in the Senate-house against the ``compact of servitude'' with Alaric, now accepted the Praetorian Prefecture.\texthebrew{⁠}\footnote{He had been Prefect of Rome in 398.

} And it is significant that he and Marcian, who became Prefect of the City, and Attalus himself, had in old days all belonged to the circle of Symmachus, the great pagan senator.\texthebrew{⁠}\footnote{As observed by Seeck (Symmachus, \emph{Opp.} p. cci): Tertullus, a member of the same group, was nominated consul in 410.

} We are told that the inhabitants of Rome were in high spirits,

p181
because the new ministers were well versed in the art of government.

The first problem which presented itself to Attalus and Alaric was how they were to act in regard to Africa, which was held by the count Heraclian, who was loyal to Honorius. They were not safe so long as they did not possess the African provinces, on which Rome depended for her supplies of corn. Alaric advised that a Gothic force should be sent to seize Africa; but Attalus would not consent, confident that he could win Carthage without fighting a battle. He sent thither a small company of Roman soldiers under Constans, while he himself marched with Alaric against Ravenna.

Honorius was overwhelmed with terror at the tidings that a usurper had arisen in Italy, and that Rome had given him her adhesion. He made ready ships in Classis, which, if it came to the worst, might bear him to the shelter of New Rome, and he sent an embassy, including Jovius and other ministers, to Attalus, proposing a division of the Empire. But Attalus had such high hopes that he would not consent to a compromise; he agreed to allow the legitimate Augustus to retire to an island and end his days as a private individual. So probable did it seem that the tottering throne of Honorius would fall, and so bright the prospects of his rival, that Jovius, who had sworn eternal enmity to Alaric, went over to the camp of the usurper. The policy of Jovius was ever, when he adopted a new cause, to go to greater lengths than any one else. And now, when he joined the side of Attalus, he went further than Attalus in hostility to Honorius, and recommended that the Emperor, when he was dethroned, should be deformed by bodily mutilation.\texthebrew{⁠}\footnote{So Olympiodorus. Philostorgius

(XII.3)

attributes the proposal of \emph{acroteriasm} to Attalus himself.

} But Attalus is said to have chidden him for this proposal; he did not guess that it was to be his own fate hereafter.

It seemed probable that Honorius would flee. But at this juncture the Eastern came to the assistance of the Western government, and Anthemius, the Praetorian Prefect of the East, sent about four thousand soldiers to Ravenna (end of A.D. 409). With these Honorius was able to secure the city of the marshes against the hostile army, and await the result of the operations of Constans, the emissary of Attalus in Africa. If Heraclian

p182
maintained the province loyally against the usurper, the war might be prosecuted in Italy against Alaric and Attalus; if, on the other hand, Africa accepted a change of rule, Honorius determined to abandon Italy.

The news soon arrived that Constans had been slain. At this point, the opposition between the ideas of Attalus and the ideas of Alaric began to reveal itself openly. Alaric wished to send an army to Africa; and Jovius supported the policy in a speech to the Roman Senate. But neither the Senate nor Attalus were disposed to send barbarians against a Roman province; such a course seemed \emph{indecent}\texthebrew{⁠}\footnote{Zosimus, VI.9

\textgreek{\texthebrew{ἀ}φε\texthebrew{ὶ}ς πρ\texthebrew{ὸ}ς α\texthebrew{ὐ}τ\texthebrew{ὴ}ν} [the Senate] \textgreek{\texthebrew{ἀ}πρεπ\texthebrew{ῆ} τινα \texthebrew{ῥ}ήματα.}} \texthebrew{—} unworthy of Rome.

Jovius, the shifty Patrician, decided, on account of the failure in Africa, to desert his allegiance to Attalus, and return to his allegiance to Honorius; and he attempted to turn Alaric away from his league with the Emperor whom he had created. But Alaric would not yet repudiate Attalus. He had said that he was resolved to persist in the blockade of Ravenna, but the new strength which Honorius had obtained from Byzantium seems to have convinced him that it would be futile to continue the siege. He marched through the Aemilian province compelling the cities to acknowledge the authority of Attalus, and, failing to take Bononia, which held out for Honorius, passed on to Liguria, to force that province also to accept the tyrant.

Attalus meanwhile returned to Rome, which he found in a sad plight. Count Heraclian had stopped the transport of corn and oil from the granary of Italy, and Rome was reduced to such extremities of starvation, that some one cried in the circus, Pretium impone carni humanae , ``set a price on human flesh.'' The Senate was now desirous to carry out the plan which it had before rejected with Roman dignity, and to send an army of barbarians to Africa; but Attalus again refused to consent to such a step.

Accordingly Alaric determined to pull down the tyrant whom he had set up; he had found that in Attalus, as well as in Honorius, the Roman temper was firm, and that he too was keenly conscious that the Visigoths were only barbarians. An arrangement was made with Honorius, who consented to pardon the usurper and those who had supported him. Near Ariminum Attalus was discrowned and divested of the purple

p183
robe with ceremonious solemnity (summer, A.D. 410); but Alaric provided for his safety, and retained him in his camp.\footnote{Along with his son Ampelius (\emph{ib.} 12). For date see Schmidt, \emph{op. cit.} I.125.

}

Alaric could now approach Honorius with a good chance, as he thought, of concluding a satisfactory settlement. Leaving his main army at Ariminum he had a personal interview with the Emperor a few miles from Ravenna (July, A.D. 410).\texthebrew{⁠}\footnote{The chronology of the events between spring 409 and August 410 cannot be determined with any precision. Attalus can hardly have been elevated before the last months of 409. The hunger in Italy, due to the measures of Heraclian, was probably felt before the beginning of 410; and probably affected the loyalty of the followers of Attalus, who had begun to desert to Honorius before Feb. 14 (see

\emph{C. Th.} IX.38.11 , cp. Schmidt, \emph{op. cit.} I.214). The deposition of Attalus must have been later than the beginning of April (as it was not known at Constantinople on April 24;

\emph{C. Th.} VII.16.2 , where tyrannici furoris et barbaricae feritatis refer to Attalus and Alaric), perhaps in May or June (Schmidt, \emph{ib.} 215).

} At this juncture the Visigoth Sarus appeared upon the scene and changed the course of history. He had been a rival of Alaric and a friend of Stilicho, and had deserted his people to enter the Roman service. Hitherto he had taken no part in the struggle between the Romans and his own nation, but had maintained a watching attitude in Picenum, where he was stationed with three hundred followers. He now declared himself for Honorius, and he resolved to prevent the conclusion of peace. His motives are not clear, but he attacked Alaric's camp. Alaric suspected that he had acted not without the Emperor's knowledge, and enraged at such a flagrant violation of the truce, he broke off the negotiations and marched upon Rome for the third time.

Having surrounded the city and once more reduced the inhabitants to the verge of starvation, he effected an entry at night through the

Salarian gate , doubtless by assistance from within,\texthebrew{⁠}\footnote{Sozomen, IX.9 \textgreek{προδοσί\texthebrew{ᾲ}}. One of the stories told in Procopius, \emph{B. V.} I.2. is that Anicia Faltonia Proba was the culprit. Unable to endure the sight of the sufferings of the people, she admitted the foe. The story, generally rejected, is accepted by Seeck (\emph{op. cit.} 413). Proba was the cousin and wife of Sextus Petronius Probus, who had a long and distinguished career recorded in many inscriptions. She was mother of three consuls. Cp. CIL VI.1754 \texthebrew{‑}5, and the genealogical tree of the Anicii in Seeck's edition of Symmachus, p. xci.

} on August 24, A.D. 410.\texthebrew{⁠}\footnote{The day is recorded in one MS. of Prosper's chronicle (\emph{Chron. min.} I.466, cp. 491), in the \emph{Excerpta Sangallensia} (\emph{ib.} 300, where 9 should evidently be read for 19 Kal. Sept.), and Theophanes, \emph{Chron.} A.M. 5903.

} This time the king was in no humour to spare the capital of the world. The sack lasted for two or three days.\texthebrew{⁠}\footnote{Orosius,

II.19.13 ;

VII.39.15 .

} It was confessed that some respect was

p184
shown for churches, and stories were told to show that the violence of the rapacious Goths was mitigated by veneration for Christian institutions.\texthebrew{⁠}\footnote{Alaric issued special orders that the churches of St. Peter and St. Paul were not to be violated. We hear that the silver tabernacle over the altar of the Lateran Basilica was stolen (\emph{Lib. Pont.} I.233); cp. Grisar, I.85. For the sack see (besides Orosius, and Sozomen)

Augustine, \emph{De civ. Dei}, I.7

(and cp. the following chapters); \emph{De urbis excidio} (\emph{P. L.} 40); Jerome, \emph{Epp.} 127, 128, 130; \emph{Prolog.} to Bks. I and III of \emph{Comm. in Ezechielem.}

} There is no reason to suppose that all the building and antiquities of the city suffered extensive damage. The palace of Sallust, in the north of the city, was burnt down, and excavations on the Aventine, then a fashionable aristocratic quarter, have revealed many traces of the fires with which the barbarian destroyed the houses they had plundered.\texthebrew{⁠}\footnote{Marcellinus, \emph{Chron. sub} 410, says that Alaric burned part of the city. The palace of the Valerii on the Caelian hill was partly burned, \emph{Vit. Melan. iun.} c. 14. The devastation in Rome and Italy is referred to in

\emph{C. Th.} VII.13.20 , which is to be dated to Feb. 411 (not 410), as Seeck has shown (\emph{Regesten}, p73). See further Lanciani, \emph{Destruction of Ancient Rome}, and A. Merlin, \emph{L'Aventin dans l'antiquité}, pp. 430\texthebrew{‑}433.

} A rich booty and numerous captives, among whom was the Emperor's sister, Galla Placidia, were taken.

On the third day, Alaric led his triumphant host forth from the humiliated city, which it had been his fortune to devastate with fire and sword. He marched southward through Campania, took Nola and Capua, but failed to capture Naples. He did not tarry over the siege of this city, for his object was to cross over to Africa, probably for the purpose of establishing himself and his people in that rich country. Throughout their movements in Italy the food-supply had been a vital question for the Goths, and to seize Africa, the granary of Italy, whether for its own sake or as a step to seizing Italy itself, was an obvious course. The Gothic host reached Rhegium; ships were gathered to transport it to Messina, but a storm suddenly arose and wrecked them in the straits. Without ships, Alaric was forced to retire on his footsteps, perhaps hoping to collect a fleet at Naples. But his days were numbered. He died at Cosentia (Cosenza) before the end of the year ( A.D. 410); his followers buried him in the Basentus, and diverted its waters into another channel, that his body might never be desecrated.\texthebrew{⁠}\footnote{Cp. Olympiodorus, \emph{fr.} 10. The same writer (\emph{fr.} 15) relates the legend that Alaric was hindered from crossing the straits by the miraculous warning of a statue. The story was suggested by an actual statue at Catona (near Reggio), the place of embarkation for Sicily, which was known as ad fretum ad statuam, CIL X.6950 . See Pace, \emph{I Barb. e Biz.} p6.

} It is related that the men

p185
who were employed on the work were all massacred, that the secret might not be divulged.\footnote{Jordanes, \emph{Get.} 158 .

}

Alaric's Ostrogothic brother-in\texthebrew{‑}law Athaulf was elected by the Visigoths to succeed him as their king.\texthebrew{⁠}\footnote{Alaric had children in 402, and Theoderic I was his grandson (see below,

p205 ). They may have died since or perhaps were girls. Athaulf was marked out by his capacity, and may have been the nearest surviving and eligible relative of Alaric.

} They must have remained for some time in southern Italy, perhaps still contemplating an invasion of Africa, but they finally abandoned the idea and marched northward along the west coast, to seek their fortunes in Gaul. Of their doings in Italy during the thirteen or fourteen months which elapsed between Alaric's death and their entry into Gaul we hear almost nothing. It is hardly probable that they visited Rome and plundered it again,\texthebrew{⁠}\footnote{As alleged by

Jordanes, \emph{Get.} 159 .

} but they laid Etruria waste. Five years later a traveller from Rome to Gaul preferred a journey by sea to traversing Tuscany devastated by Gothic sword and fire.

Athaulf crossed the Alps early in A.D. 412, perhaps by the pass of Mont Genèvre,\texthebrew{⁠}\footnote{\emph{Chron. Gall.} 87, p654; Schmidt, \emph{op. cit.} I.223. If the Goths had taken the coast-road, they would have had to do with Constantius, who was at Arles.

} to play a leading part in the troubled politics of Gaul. But to explain the situation which confronted him we must go back to A.D. 406 and follow the course of events of six years which were of decisive importance for the future histories of Gaul, Spain, and Britain.

On the last day of December A.D. 406 vast companies of Vandals, Suevians, and Alans began to cross the Rhine near Moguntiacum and pour into Gaul.\footnote{The sources for the events related in this section are Olympiodorus, \emph{frs.} 12, 14, 16; Zosimus, V. 27 ,

31 ,

32 , 43 , VI. 1\texthebrew{‑}6 ,

9 , 13 , and Sozomen, IX.11\texthebrew{‑}15 (both dependent on Olympiodorus); Orosius, VII. 38

and

40\texthebrew{‑}42 ; Prosper; \emph{Consularia Italica}; and Hydatius; Jerome, \emph{Ep.} 123 (\emph{ad Ageruchiam}, A.D. 409); Renatus Profuturus Frigeridus apud Gregory of Tours, \emph{H.F.} II.9; Orientius,

(p186)
\emph{Common.} II.165 \emph{sqq.}; Paulinus (his identity is uncertain), \emph{Epigramma} 10 \emph{sqq.}; Prosper, \emph{De prov. Dei}, 15 \emph{sqq.}; Salvian, \emph{De gub. Dei}, VI.15, VII.12. The most useful modern studies are Freeman's essay on \emph{Tyrants of Britain, Gaul, and Spain} ( \emph{E.H.R.} I, Jan. 1886 ; reissued in \emph{Western Europe in the Fifth Century}), and Schmidt's \emph{Gesch. der Wandalen}, 17 \emph{sqq.}

}

p186
The Asding Vandals, who, as we saw, invaded Raetia in A.D. 401, were finding their lands on the Theiss insufficient to support their growing numbers,\texthebrew{⁠}\footnote{Procopius, \emph{B. V.} I.22.3 \textgreek{λιμ\texthebrew{ῷ} πιεζόμενοι} (perhaps the tradition of the Vandals themselves).

} and joining with the Alans, who were living in Pannonia, and with Suevians, who probably represent the ancient Quadi, they migrated northward to the Main. We may conjecture that this movement had some connexion with the unsettled conditions beyond the Middle Danube, which caused Radagaisus and his followers to invade Italy; and that the smaller German peoples who lived in those regions found themselves pressed and harried by their more power ­ful neighbours the Huns and the Ostrogoths. The idea of wandering into Gaul was naturally suggested by the fact that the Rhine frontier was no longer adequately defended. A large number of the Roman troops stationed there had been withdrawn recently by Stilicho, for the defence of Italy. On the Main, the host was joined by the Siling Vandals, who lived there with the Burgundians, to the east of the Alamanni.

The Alans were the first to reach the Rhine. They were led by two kings, Goar and Respendial, but here Goar separated himself from his fellows and offered his services to the Romans. The Asdings, under their king Godegisel, were some distance behind, when their march was interrupted by the appearance of an army of Franks,\texthebrew{⁠}\footnote{Obviously the Ripuarian Franks, whose seats were along the Rhine north of the Alamanni (whose territory extended from the Main southward to the Lake of Constance).

} who as federates had undertaken the duty of protecting the Rhine for Rome. Godegisel was slain, and the Vandals would have been utterly destroyed had not Respendial returned to their aid. His Alans changed the fortunes of the battle, the Franks were defeated, and the invaders crossed the Rhine. Their first exploit was to plunder Mainz and massacre many of the inhabitants, who had sought refuge in a church. Then advancing through Germania Prima they entered Belgica, and following the road to Trier they sacked and set fire to that Imperial city. Still continuing their westward path they crossed the Meuse and the Aisne and wrought their will on Reims. From here they seem to have turned northward. Amiens,

p187
Arras and Tournay were their prey; they reached Térouanne,\texthebrew{⁠}\footnote{Teruanna, the town of the Morini.

} not far from the sea, due east of Boulogne, but Boulogne itself they did not venture to attack. After this diversion to the north, they pursued their course of devastation southward, crossing the Seine and the Loire into Aquitaine, up to the foot of the Pyrenees. Few towns could resist them. Toulouse was one of the few, and its success ­ful defence is said to have been due to the energy of its bishop Exuperius.

Such, so far as we can conjecture from the evidence of our meagre sources, was the general course of this invasion, but we may be sure that the barbarians broke up into several hosts and followed a wide track, dividing among them the joys of plunder and destruction. Pious verse-writers of the time, who witnessed this visitation, painted the miseries of the helpless provinces vaguely and rhetorically, but perhaps truthfully enough, in order to point a moral.

The terror of fire and sword was followed by the horror of hunger in a wasted land.

In Eastern Gaul too some famous cities suffered grievously from German foes. But the calamities of Strassburg, Speier, and Worms were perhaps not the work of the Vandals and their associates. The Burgundians seem to have taken advantage of the crisis to push down the Main, and at the expense of the Alamanni to have occupied new territory astride the Rhine. And it is probably these two peoples, especially the Alamanni dislodged from their homes, who were responsible for the havoc wrought in the province of Upper Germany.\footnote{Cp. Schmidt, \emph{op. cit.} 2A.

}

It may have been in the early summer of A.D. 407 that the situation was changed by the arrival of Roman legions not from Italy but from Britain. That island had the reputation of being a fertile breeder of tyrants, and before the end of the previous year the Britannic soldiers had denounced the authority of Honorius and set up an Emperor for themselves in the person of a certain Marcus. We have no knowledge of their reason for this step, but we may conjecture that the revolt was due to discontent with the rule of the German Stilicho, just as the revolt of Maximus had been aimed at the German general

p188
Merobaudes. There was a certain Roman spirit alive among the legionaries, jealous of the growth of German influence. And we can well understand that they were impatient of the neglect of the defence of the Britannic provinces by the central government. One of the legions which guarded the island had been withdrawn in A.D. 401\texthebrew{⁠}\footnote{See above,

p161 .

} for the defence of Italy, but we are not informed whether it was sent back. In any case the troops in the island were probably not kept up to their nominal strength and were insufficient to contend against the constant inroads of the Picts and the expeditions of the Irish from beyond their channel, as well as the raids of Saxon freebooters from the continent. To subdue these enemies had been a task which had demanded all the energy of Theodosius himself. A victory over the Picts seems to have been gained in the early years of Honorius, but it was not of great account,\texthebrew{⁠}\footnote{Claudian, \emph{In Eutrop.} I.393

fracto secura Britannia Picto. Had the success been considerable, Claudian would have made more of it.

} and when events in the south forced Stilicho to denude the Rhine of its defenders, little thought can have been taken at Rome or Ravenna for the safety of remoter Britain. It was a favourable opportunity for such an expedition as that which Irish Annals record to have been led against the southern coasts of Britain by the High King of Ireland in A.D. 405.\texthebrew{⁠}\footnote{Cp. Bury, \emph{Life of Saint Patrick}, p331.

} In such circumstances we can easily conceive that the troops longed for a supreme responsible authority on the spot.

Marcus was not a success. Soon after his elevation he was pronounced unfit and slain, to make way for Gratian, who reigned for four months ( A.D. 407) and then met the fate of Marcus. The third tyrant was a private soldier who bore the auspicious name of Constantine, and was to play a considerable part for a few years on the stage of western Europe.

The first act of Constantine was to cross with an army into Gaul. It has been supposed that he feared an invasion of Britain by the German hordes, who had indeed approached the Channel, and that he went forth to meet the danger. It seems more probable that he was following the example of Magnus Maximus, who had in like manner crossed over to the continent to wrest Gaul and Spain from Gratian. He landed at Boulogne. It appears to be commonly supposed that he took with him all

p189
the forces in Britain, not only the field army, but also the garrisons of the frontiers. This is highly improbable. For we cannot imagine that he did not intend to retain his hold on the island, and it has been inferred from the evidence of a coin that he set up a colleague before he sailed.\texthebrew{⁠}\footnote{See A. J. Evans, \emph{Numismatic Chronicle}, 3rd series, VII.191 \emph{sqq.}, 1887; Bury, App. 19 to Gibbon, vol. III.

} But he must have been accompanied by the whole field army, which was not very large, or the greater part of it.

Gaul sorely needed a Roman defender at the head of Roman legions, and the Gallic legions went over to Constantine. He inflicted a severe defeat on the barbarians, we know not where, and he is said to have guarded the Rhine more efficiently than it had been guarded since the reign of Julian \texthebrew{—} a statement which comes from a pagan admirer of the Apostate. The representatives of Honorius fled to Italy when Constantine passed into the Rhone valley and the south-eastern districts, which had escaped the ravages of the Germans. He seems to have made agreements with some of the intruders,\texthebrew{⁠}\footnote{Probably with Alamanni and Burgundians. See

Orosius, VII.40 .

} which they perfidiously violated. But we know nothing definite as to his dealings with them. ``For two years,'' writes a modern historian,\texthebrew{⁠}\footnote{Freeman, \emph{op. cit.}

} ``they and he both carry on operations in Gaul, each, it would seem, without any interruption from the other. And when the scene of action is moved from Gaul to Spain, each party carries on its operations there also with as little of mutual let or hindrance. It was most likely only by winking at the presence of the invaders and at their doings that Constantine obtained possession, so far as Roman troops and Roman administration were concerned, of all Gaul from the Channel to the Alps. Certain it is that at no very long time after his landing, before the end of the year 407, he was possessed of it. But at that moment no Roman prince could be possessed of much authority in central or western Gaul, where Vandals, Suevians, and Alans were ravaging at pleasure. The dominion of Constantine must have consisted of a long and narrow strip of eastern Gaul, from the Channel to the Mediterranean, which could not have differed very widely from the earliest and most extended of the many uses of the word Lotharingia. He held the imperial city on the Mosel, the home of Valentinian and the earlier Constantine.''

p190
When Constantine obtained possession of Arelate (Arles), then the most prosperous city of Gaul, it was time for Honorius and his general to rouse themselves. We saw how Stilicho formed the design of assigning to Alaric the task of subduing the adventurer from Britain, who had conferred upon his two sons, Constans, a monk, and Julian, the titles of caesar and nobilissimus respectively. But this design was not carried out. A Goth indeed, and a brave Goth, but not Alaric, crossed the Alps to recover the usurped provinces; and Sarus defeated the army which was sent by Constantine to oppose him. But he failed to take Valentia, and returned to Italy without having accomplished his purpose ( A.D. 408).

The next movement of Constantine was to occupy Spain.\texthebrew{⁠}\footnote{Zosimus, VI.4.

Terentius was appointed mag. mil., Apollinaris (grandfather of Sidonius the poet) Praetorian Prefect (\emph{ib.}), and Decimius º Rusticus Master of Offices (Greg. of Tours, II.9, quoting from Renatus Frigeridus).

} We need not follow the difficult and obscure operations which were carried on between Spanish kinsmen of Honorius and the troops which the Caesar Constans and his lieutenant Gerontius led across the Pyrenees.\texthebrew{⁠}\footnote{Freeman has shown that we are not justified in accepting the version of the story which states that the representatives of the Theodosian house were engaged in defending the northern frontier of the peninsula against the Vandals and their fellow-plunderers before Constantine attempted to occupy it.

} The defenders of Spain were overcome, and Caesaraugusta (Zaragoza) became the seat of the Roman Caesar. Thus in the realm of Constantine almost all the lands composing the Gallic prefecture were included; he might claim to be the lord of Britain; the province of Tingitana, beyond the straits of Gades, was the only province that had obeyed Honorius and did not in theory obey Constantine.

Constans, however, was soon recalled to Gaul by his father, and elevated to the rank of Augustus. But Constantine himself meanwhile, possessing the power of an Emperor, was not wholly content; he desired also to be acknowledged as a colleague by the son of Theodosius, and become legitimised. He sent an embassy for this purpose to Ravenna (early in A.D. 409), and Honorius, hampered at the time by the presence of Alaric, was too weak to refuse the pacific proposals.\texthebrew{⁠}\footnote{Constantine assumed the consul­ship in 409 in his dominions, as colleague of Honorius. See Liebenam, \emph{Fasti consulares}, p41. Captives of the Theodosian house, who had been taken in the Spanish expedition, were in the hands of Constantine, and a hope of their release seems to have been one of the motives of Honorius in sending the purple robe to the

(p191)
usurper; but before the embassy was sent the captives had been put to death. For the coinage of Constantine and Constans see Cohen, VIII.198 \emph{sqq.}

} Thus

p191
Flavius Claudius Constantinus was recognised as an Augustus and an Imperial brother by the legitimate emperor; but the fact that the recognition was extorted and soon repudiated, combined with the fact that he was never acknowledged by the other Augustus at New Rome, might justify us in refusing to include the invader from Britain who ruled at Arelate in the numbered list of Imperial Constantines. Some time afterwards another embassy, of whose purpose we are not informed, arrived at Ravenna, and Constantine promised to assist his colleague Honorius against Alaric, who was threatening Rome. Perhaps what Honorius was to do in return for the proffered assistance was to permit the sovran of Gaul to assume the consul ­ship. In any case it was suspected that Constantine aspired to add Italy to his realm as he had added Spain, and that the subjugation of Alaric was only a pretext for entering Italy, as it might have been said that the subjugation of the Vandals and their fellow-invaders had been only a pretext for his entering Gaul. Hellebich, Master of Soldiers ( equitum ), was also suspected of favouring the designs of the usurper, and the suspicion, whether true or false, cost him his life; Honorius caused him to be assassinated. When this occurred Constantine was already in Italy, and the fact that when the news reached him he immediately recrossed the mountains, strongly suggests that the suspicion was true, and that he depended on this general's treason for the success of his Italian designs.

Constans had left his general, Gerontius, a Briton, in charge of Spain. Barbarian federates, known as Honorians, had been used for the conquest of Spain by Constans, and to these was entrusted the defence of the passes of the Pyrenees. It was an unfortunate measure. The Spanish regular troops, who now acknowledged the authority of Constantine, thought that the charge ought to have been entrusted as before to the national militia, and they revolted.\texthebrew{⁠}\footnote{For the troops stationed there in the fifth century see

\emph{Not. dig., Occ.} XLII.25\texthebrew{‑}32 . One legion (Septima Gemina) º and four cohorts in Gallicia, and one cohort at Veleia in Tarraconensis.

} The Honorians betrayed or neglected their trust. It was the autumn of A.D. 409, and on a Tuesday, either September 28 or October 5, the host of barbarians who had been oppressing western Gaul for more than two years \texthebrew{—} the

p192
Asdings under King Gunderic, the Silings, the Sueves, and the Alans \texthebrew{—} crossed the mountains and passed into Spain.\footnote{The alternative dates are given by the Spanish chronicler Hydatius. They may have followed (as Schmidt thinks, \emph{op. cit.} 26) the main road from Bordeaux to Pampluna.

}

Constans imputed the troubles in Spain to the incapacity of Gerontius, and he returned from Gaul to supersede him and restore order. But Gerontius was not of a spirit to submit tamely. He seems to have come to terms with the legions, and he made some sort of league with the barbarians, by which a large part of the land was abandoned to them.\texthebrew{⁠}\footnote{The sources give confused and contradictory accounts as to the order of events, and uncertainty may be felt whether the revolt of Gerontius preceded the entry of the Vandals into Spain, as there is a suggestion in some writers that they were invited by him.

} He renounced the authority of Constantine, and though he did not assume the purple himself, he raised up a new Emperor, a certain Maximus, who was perhaps his own son.

Thus at the beginning of A.D. 410 there were six Emperors, legitimate and illegitimate, acknowledged in various parts of the Empire. Besides Honorius and his nephew Theodosius, there was Attalus at Rome, there were Constantine and Constans at Arles, and there was Maximus at Tarragona.

Constans soon fled before Gerontius and his barbarian allies to Gaul, and after some time \texthebrew{—} the chronology is very obscure \texthebrew{—} Gerontius, leaving Maximus to reign in state at Tarragona, marched into Gaul against the father and son who had once been his masters. It was apparently in A.D. 411 that Constans was captured and put to death at Vienne, and then his father Constantine was besieged at Arles.

But Honorius, now that Alaric was dead, although the Goths were still in Italy, was able to bethink him of the lands he had lost beyond the Alps, and he sent an army under two generals Constantius and Ulfila, to do what Sarus had failed to do and win back Gaul. Constantius was an Illyrian, born at Naissus, the birthplace of Constantine the Great, and for the next ten years the fortunes of Honorius were to depend upon him as before they had depended upon Stilicho. We may consider it certain that when he led the troops of Italy to Gaul he had already been raised to the post of Master of Both Services.\texthebrew{⁠}\footnote{In succession to Valens. Prosper describes him as mag. mil., \emph{sub} 412, as patricius, \emph{sub} 415. What post Ulfila held and who was mag. equitum is unknown.

} We have a slight portrait of his appearance and manners. He had

p193
large eyes, a broad head, and a long neck; he leaned low over the neck of his horse, and as his eyes shot swift glances right and left he seemed to beholders a man who might one day aim at the throne. On public occasions his look was stern, but in private, at table and at wine-parties , he was genial and agreeable. He was superior to the temptations of money, though at a later stage of his career he was to fall into the vice of avarice. His ambition was associated with love. He was passionately attached to the Emperor's step-sister Galla Placidia, who was now a captive in the hands of the Goths.

When Constantius and his Gothic subordinate Ulfila advanced along the coast road of Provence against Arles, the blockading army of Gerontius fled before the representatives of legitimacy. Gerontius returned to Spain and there his own troops turned against him. The house in which he took refuge was besieged; he and his Alan squire fought long and bravely for their lives; then the house was set on fire, and at length in despair he slew his squire and his wife at their own request and then stabbed himself.\texthebrew{⁠}\footnote{The story is given in great detail by Sozomen (IX.4), who praises Nunechia (she was a Christian) for imploring her husband to kill her.

} Maximus fled to find safety among some of the barbarian invaders who had supported his throne.

Meanwhile Constantine, with his second son Julian, was being besieged in Arles by the army of Italy which had replaced the army of Spain. The siege wore on for three months, and the hopes of the legitimised usurper depended upon the arrival of his general Edobich, who had been sent beyond the Rhine to gain reinforcements from the Alamanni and Franks. Edobich at length returned with a formidable army, but a battle, fought near the city, resulted in a victory for the besiegers. Edobich was slain by the treachery of a friend in whose house he sought shelter, and Constantine, seeing that his crown was irrecoverably lost, thought only of saving his life. He stripped off the Imperial purple and ``fled to a sanctuary, where he was ordained priest, and the victors gave a sworn guarantee for his personal safety. Then the gates of the city were thrown open to the besiegers, and Constantine was sent with his son to Honorius. But that Emperor, cherishing resentment towards them for his cousins, whom Constantine had slain, violated the oaths and ordered

p194
them to be put to death, \texthebrew{•} thirty miles from Ravenna''\texthebrew{⁠}\footnote{Olympiodorus, \emph{fr.} 16.

} (September, A.D. 411).

It was not long after the fall of Constantine that a new tyrant was elevated in Gaul. Jovinus, a Gallo-Roman , was proclaimed at Moguntiacum. This city, which had been wrecked by the barbarians five years before, was now in the power of the Burgundians, and it was their king, Gundahar, and Goar, the Alan chief (who, it will be remembered, had been enlisted in the service of Honorius), to whom Jovinus owed the purple. Constantius and Ulfilas, having done their work in overthrowing the tyrant of Arles, had returned to Italy, and the subjugation of Jovinus was reserved for the Visigoths.

It has already been related that the Visigoths, under the leader ­ship of King Athaulf, crossed the Alps early in A.D. 412. They took with them their captive Galla Placidia and the deposed Emperor Attalus. They had come to no agreement with Ravenna; if any agreement had been made, the restoration of Placidia would have been a condition. Athaulf was probably more inclined to side with Jovinus against Honorius than with Honorius against Jovinus. Circumstances decided him to champion the cause of legitimacy.

Attalus, from some motive which is not clear, persuaded him to offer his services to Jovinus. But it appears that the arrival of this unexpected help was not welcome to the tyrant. Perhaps his Burgundian friends did not look with favour on the coming of a people into Gaul who might prove rivals to themselves. Perhaps the terms which Athaulf proposed seemed exorbitant. Then Sarus, the Visigoth who had been in the service of Honorius, and who was the mortal enemy of Athaulf just as he had been the mortal enemy of Alaric, appeared on the scene with above a score of followers to attach himself to the fortunes of Jovinus, because Honorius had refused to grant him justice for the murder of a faithful domestic. Athaulf was incensed when he heard of his approach, and advanced with ten thousand

p195
to crush twenty men. Sarus did not shirk fighting against such appalling odds, and having performed deeds of marvellous heroism he was taken and put to death. This incident did not tend to smooth the negotiations with Jovinus, and when the tyrant proclaimed his brother Sebastian Augustus, against Athaulf's wishes,\texthebrew{⁠}\footnote{The reason of his objection is not stated. Schmidt (\emph{op. cit.} I.224) says that Athaulf aspired himself to be the colleague of Jovinus. That sounds incredible. I suggest that Athaulf's scheme was the elevation of Attalus and the division of Gaul between him and Jovinus.

} the Visigoth entered into communication with Dardanus the Praetorian Prefect, the only important official in Gaul who had not deserted the cause of Honorius. Envoys were sent to Ravenna, and Honorius accepted the terms of Athaulf, who promised to send him the heads of the two tyrants. Sebastian was defeated and slain immediately, and Jovinus fled to Valence, which, so recently besieged by Gerontius, was now to undergo another siege. It seems to have been taken by storm; Jovinus was carried to Narbonne and executed by the order of Dardanus (autumn, A.D. 413).\texthebrew{⁠}\footnote{Olympiodorus says that the heads of the two tyrants were exposed \textgreek{Καρθαγένης \texthebrew{ἔ}ξωθεν}, as those of Constantine and Julian had been (two years before). \textgreek{Καρθαγένη} might mean either Carthage or New Carthage (Carthagena) in Spain. It is generally explained to mean Carthage. I am inclined to think that Olympiodorus confused the two cities, and that while the heads of the earlier tyrants were exhibited at Carthagena, those of the later pair were taken to Carthage (in view of the revolt of Heraclian). Coins of Sebastian (silver) were issued during what must have been a very brief reign at Arles and Trier. For these and those of Jovinus see Cohen, VIII.202\texthebrew{‑}203.

} For the moment the authority of Honorius was supreme in Gaul.

It may be wondered why Constantius having suppressed Constantine did not return to Gaul to deal with Jovinus. The explanation probably is that his presence in Italy was required to prepare measures for dealing with another tyrant who had arisen in Africa. The revolt of the count Heraclian, the slayer of Stilicho, was instigated, we are told, by the examples of tyranny which he had observed in Gaul.\texthebrew{⁠}\footnote{See

Philostorgius, XII.6 , where Heraclian's name has been rightly restored.

} So infectious was ``tyranny'' that the man who three years before resisted the proposals of Attalus and the menaces of Alaric, loyally standing by the throne of Honorius, and who had been rewarded by the consul ­ship,\texthebrew{⁠}\footnote{Heraclian's consul­ship in 413 shows that his revolt began in that year (not in 412 as Hydatius, 51, suggests).

} now threatened his sovran without provocation. He did not wait to be attacked in Africa. With a large fleet,

p196
of which the size was grossly exaggerated at the time,\texthebrew{⁠}\footnote{3700 ships acc. to

Orosius, VII.42

and one of the two best MSS. of Marcellinus (\emph{sub} 413; the other gives 700 ships and 3000 soldiers).

} he landed in Italy, intending to march on Rome, but was almost immediately defeated,\texthebrew{⁠}\footnote{The words of Orosius, \emph{ib.}, suggest that he landed at the mouth of the Tiber and was defeated near the coast on his way to Rome (so Gibbon). But our other Spanish authority, Hydatius, 56, states that the battle was fought at Otricoli and 50,000 were slain.

Otricoli

is the first place where the Via Flaminia crosses the Tiber, after the Pons Mulvius.

Thayer's Note: That piece of writing should not suggest to the gentle reader that Otricoli is anywhere near the Pons Mulvius. Otricoli (the ancient Ocriculum) is some 70 km N of Rome. It is, however, the first place after the Flaminia next crosses the Tiber: that is, the Flaminia traverses much of the northern Latium while staying to the W of the river.

} and fled back to Africa in a single ship to find that the African provinces would have none of him. He was beheaded in the Temple of Memory at Carthage (summer, A.D. 413).\texthebrew{⁠}\footnote{The edict annulling the acts of Heraclian and obliterating his name

(\emph{C. Th.} XV.14.13)

is dated Aug. 3.

} His consul ­ship was declared invalid, and his large fortune was made over to Constantius, who was designated consul for the following year.

This revolt affected the course of events in Gaul. Honorius, whose mind did not travel far beyond his family and his poultry-yard , was bent on recovering his sister Placidia from the hands of the Visigoth, and this desire was ardently shared by Constantius, who aspired to the hand of this princess. Athaulf had agreed to restore her when the bargain had been made that in return for his services in crushing Jovinus he and his people should be supplied with corn and receive a Gallic province as Federates of Empire. But Africa was the corn-chamber of Italy, and when Heraclian stopped the transport of supplies\texthebrew{⁠}\footnote{Orosius, \emph{ib.}

} it became impossible to fulfil the engagement with Athaulf. There was hunger in the Gothic camp. Athaulf therefore refused to carry out his part of the compact and surrender Placidia. He made an attempt to take Marseilles, which he hoped might fall by treachery, but it was defended by ``the most noble'' Boniface, an officer who with afterwards to play a more conspicuous and ambiguous part in Africa. Athaulf himself was severely wounded by a stroke which the Roman dealt him. But he was more fortunate at Narbonne. He captured this town and made it his headquarters, and he also seized the important cities of Bordeaux and Toulouse.\footnote{Capta Tolosa,

Rutil. Namat. I.496 ; nostra ex urbe [\emph{sc.} Burdigala] Gothi, fuerant qui in pace recepti,

Paulinus Pell. \emph{Eucharisticos}, 312 . The notice in \emph{Chron. Gall.} p654, Aquitania Gothis tradita, relates to A.D. 414, but seems to be a mistaken anticipation of the settlement of 418 (cp. Schmidt, \emph{op. cit.} I.226).

}

Having established himself in Narbonensis and Aquitaine,

p197
Athaulf determined to give himself a new status by allying himself in marriage to the Theodosian house. Negotiations with Ravenna were doubtless carried on during his military operations, but he now persuaded Placidia, against the will of her brother, to give him her hand. The nuptials were celebrated in Roman form (in January, A.D. 414)\texthebrew{⁠}\footnote{The description comes from Olympiodorus, \emph{fr.} 24. Philostorgius

(XII.4)

compares the marriage to the union of pottery with iron (the fourth empire symbolised by the iron legs of the image in Daniel ii was explained as the Roman, see Sulpicius Severus, \emph{Chron.} II.3). See the note of Bidez, \emph{ad. loc.} Hydatius, 57, saw in it the fulfilment of Daniel's prophecy (xi.6) that the queen of the south would marry a king of the north.

} at Narbonne, in the house of Ingenius, a leading citizen, and the pride of Constantius, who had just entered upon his first consul ­ship, was spoiled by the news that the lady whom he loved was the bride of a barbarian. We are told that, arrayed in Roman dress, Placidia sat in the place of honour, the Gothic king at her side, he too dressed as a Roman. With other nuptial gifts Athaulf gave his queen fifty comely youths, apparelled in silk, each bearing two large chargers in his hands, filled one with gold, the other with priceless gems \texthebrew{—} the spoils of Rome. They had an ex-Emperor , Attalus, to conduct an epithalamium. The marriage festivities were celebrated with common hilarity by barbarians and Romans alike.

A contemporary writer\texthebrew{⁠}\footnote{Orosius, VII.42 .

} has recorded words said to have been spoken by Athaulf, which show that, perhaps under the influence of Placidia, he had come to adopt a new attitude to the Empire. ``At first,'' he said, ``I ardently desired that the Roman name should be obliterated, and that all Roman soil should be converted into an empire of the Goths; I longed that Romania should become Gothia\texthebrew{⁠}\footnote{Romania, ut vulgariter loquar. This early use of Romania for the territory of the Roman Empire deserves notice.

} and Athaulf be what Caesar Augustus was. But I have been taught by much experience that the unbridled licence of the Goths will never admit of their obeying \emph{laws}, and without laws a republic is not a republic. I have therefore chosen the safer course of aspiring to the glory of restoring and increasing the Roman name by Gothic vigour; and I hope to be handed down to posterity as the initiator of a Roman restoration, as it is impossible for me to change the form of the Empire.''

We can hardly be wrong in ascribing this change in the spirit and policy of Athaulf to the influence of Placidia, and conjecturing

p198
that his conversion to Rome was the condition of her consent to the marriage. We know too little of the personality of this lady who was to play a considerable part in history for thirty years. She was now perhaps in her twenty-sixth year, and she may have been younger.\texthebrew{⁠}\footnote{Theodosius married her mother Galla in 387 ( Zos. IV.43 ; so Gibbon, Clinton, Güldenpenning; in 386 acc. to Marcellinus \emph{sub a.}, so Tillemont, Sievers) towards the end of the year; so that Placidia may have been born in 388. Theodosius went to the west in that year and did not return to Constantinople, where Galla had remained during his absence, till Nov. 391, where he remained till the day after Galla's death in May 394. Galla died in childbirth, and the child died. It follows from these dates that Placidia might have been born in 392\texthebrew{‑}393. Of the two alternatives 388 appears to me to be the more probable.

} Her personal attractiveness is shown by the passion she inspired in Constantius, and the strength of her character by the incidents of her life. She can have been barely twenty years of age when she approved of the execution of her cousin Serena at Rome, and in defiance of her brother's wishes in uniting herself to the Goth she displayed her independence. She was in later years to become the ruler of the West.

The friendly advances which were now made to Honorius by the barbarian, who had been forced upon him as a brother-in\texthebrew{‑}law , were rejected. Athaulf then resorted to the policy of Alaric. He caused the old tyrant Attalus to be again invested with the purple. Constantius, the Master of Soldiers, went forth for a second time to Arles to suppress the usurper and settle accounts with the Goths. He prevented all ships from reaching the coast of Septimania, as the territory of Narbonensis was now commonly called. The Goths were deprived of the provisions which reached Narbonne by sea, and their position became difficult. Athaulf led them southward to Barcelona, probably hoping to establish himself in the province of Tarraconensis (early in A.D. 415). But before they left Gaul, the Goths laid waste southern Aquitaine and set Bordeaux on fire.\texthebrew{⁠}\footnote{We learn of these events from the \emph{Eucharisticos}, the poem of Paulinus of Pella, already cited (308 \emph{sqq.}). He describes the siege of Vasatae (Bazas), of which he was a witness. It was attacked by Goths and Alans, and was saved by the success of Paulinus in indu­cing the Alans to go over to the side of the Romans, \emph{ib.} 329 \emph{sqq.} The king of the Alans, an old friend of his, was probably Goar, whom we have already met (so Tillemont, Freeman, Schmidt).

} Attalus was left behind and abandoned to his fate, as he was no longer of any use to the Goths. Indeed his elevation had been a mistake. He had no adherents in Gaul, no money, no army, no one to support him

p199
except the barbarians themselves.\texthebrew{⁠}\footnote{Paulinus, who was grandson of the poet Ausonius and son of Hesperius, Praet. Prefect of Gaul in 379, accepted from Attalus the post of keeper of the privy purse, comes privatae largitionis (the title of an official subordinate to the comes r. priv., see

\emph{Not. dig., Occ.} XII.4 ) \texthebrew{—} a post, says Paulinus (\emph{Euchar.} 296),

quam sciret nullo subsistere censu,

iamque suo ipse etiam desisset fidere regno,

solis quippe Gothis fretus male iam sibi notis

quos ad praesidium uitae praesentis habere,

non etiam imperii poterat, per se nihil ipse

aut opibus propriis aut ullo milite nixus.

(This is a specimen of the doggerel written by the grandson of Ausonius.) Coins show that Attalus had obtained some recognition at Trier.

} He escaped from Gaul in a ship, but was captured and delivered alive to Constantius.\texthebrew{⁠}\footnote{He was captured in 416 (\emph{Chron. Pasch., sub a.}) . Cp. Prosper, \emph{sub} 415, and

Orosius, VII.42 . Philostorgius

(XII.4)

says he was surrendered by the Goths, after Athaulf's death.

} In A.D. 417, the eleventh consul ­ship of Honorius and the second of Constantius, the Emperor entered Rome in triumph with Attalus at the wheels of his chariot. He punished the inveterate tyrant by maiming him of a finger and thumb, and condemning him to the fate which Attalus had once been advised to inflict upon himself. He had not forgotten how the friend of Alaric had demanded with an air of patronising clemency that the son of Theodosius should retire to some small island, and he banished his prisoner to Lipara.

At Barcelona a son was born to Athaulf and Placidia. They named him Theodosius after his grandfather, and the philo-Roman feelings of Athaulf were confirmed. The death of the child soon after birth was a heavy blow; the body was buried, in a silver coffin, near the city.\texthebrew{⁠}\footnote{Olympiodorus, \emph{fr.} 27.

} Athaulf did not long survive him. He had been so unwise as to take into his service a certain Dubius, one of the followers of Sarus, who avenged his first by slaying his second master. The king had gone to the stable, as was his custom, to look after his own horses , and the servant, who had long waited for a favourable opportunity, stabbed him (September, A.D. 415).\texthebrew{⁠}\footnote{The news of his death reached Byzantium on Sept. 24 (\emph{Chron. Pasch., sub a.}) and was the occasion of games and rejoicings.

} He did not die till he had time to recommend his brother, who he expected would succeed to the kingship, to send Placidia back to Italy. But his brother did not succeed him. Singeric, the brother of Sarus \texthebrew{—} who probably had been privy to the deed of Dubius \texthebrew{—} seized the royalty and put to death the children of the dead king by his first wife, tearing them from the arms of the bishop Sigesar to whose protection they had fled for refuge. Placidia he treated with indignity and cruelty, compelling her to walk on foot for \texthebrew{•} twelve

p200
miles in the company of captives. But the reign of the usurper (for he had seized the power by violence without any legal election) endured only for seven days; he was slain, and Wallia was elected king.

For the moment Gaul was free from the presence of German invaders, with the exception of one region. The Burgundians, who had crossed the Rhine and occupied the province of Germania Superior, had been confirmed in their possession by the tyrant Constantine. After the fall of Jovinus, whom they had supported, Honorius was in no position to turn them out. He accepted them as Federates of the Empire;\texthebrew{⁠}\footnote{Prosper, \emph{sub} 413.

} they were bound to guard the Rhine against hostile invaders. Thus in A.D. 413 was founded the first Burgundian kingdom in Gaul, the kingdom of Worms (Borbetomagus). It is the Burgundy of the Nibelungenlied, which also preserves the name of the king, Gundahar (Gunther), who had gained for his people a footing west of the Rhine.

The island of Britain, when many of the troops were withdrawn by Constantine in A.D. 407, was left to defend itself as best it could against Picts, Scots, and Saxons. For a while the Vicar of the Diocese and the two military commanders of the frontier forces, the Count of the Saxon Shore in the south-east, and the Duke of the Britains in the north, were doubtless in communication with Constantine and taking their orders from him. When a great Saxon invasion devastated the country in A.D. 408,\texthebrew{⁠}\footnote{\emph{Chron. Gall.} 62 (p654). Here there are two successive entries: 61, hac tempestate praevaletudine (praevalente hostium multitudine, Mommsen) Romanorum vires attenuatae; 62, Britannia Saxonum incursione devastatae. Freeman (\emph{op. cit.} p149) was misled by the bad text of Roncalli's edition.

} the Emperor in Gaul was in no position to send troops to the rescue, and the inhabitants of Britain renounced his authority, armed themselves, and defended their towns against the invaders.\texthebrew{⁠}\footnote{Zosimus, VI.5.2

\textgreek{\texthebrew{ὅ}πλα \texthebrew{ἐ}νδύντες} : this was a violation of a Lex Julia.

} The news reached Italy, and Honorius seized the opportunity of writing, apparently to the local magistrates, authorising them to take all necessary measures for self-defence .\texthebrew{⁠}\footnote{\emph{Ib.} 10.2 \textgreek{\texthebrew{Ὁ}νωρίου δ\texthebrew{ὲ} γράμμασι πρ\texthebrew{ὸ}ς τ\texthebrew{ὰ}ς \texthebrew{ἐ}ν Βρεταννί\texthebrew{ᾲ} χρησαμένου πόλεις φυλάττεσθαι παραγγέλλουσι}. It may be noted that in the reign of Honorius, Anderida (Pevensey) on the Saxon Shore was repaired and a new fort built at Peak on the Yorkshire coast (Haverfield, \emph{C. Med. H.} I.379).

} We have no information as to the attitude of the Imperial garrisons and their commanders to the revolution. It is possible

p201
that they sympathised with the provincials and shared in it; most of these troops had the tradition of association with Britain for centuries. In any case, when Constantine fell, and the tyrant Jovinus had been crushed and Honorius was again master in Gaul, there can be little doubt that he and Constantius took measures to re-establish his power in Britain.\texthebrew{⁠}\footnote{This is contrary to the ordinary view. Cp. Sagot, \emph{La Bretagne romaine}, 251 \emph{sqq.}; Lot, \emph{Les Migrations saxonnes}, 11\texthebrew{‑}13. For the condition of Britain in the last period of Roman rule see Haverfield, \emph{Romanisation of Roman Britain}, and

his article on \emph{Britain (Roman)} in \emph{Encyclopaedia Britannica} (Ed. 11) ; and \emph{C. Med. H.} I.

} In the first place, it is not probable that the provincials would have been able to hold out against the Saxon foe for fifteen or sixteen years without regular military forces, and we know that the Saxon did not begin to get any permanent foothold in the island before A.D. 428.\texthebrew{⁠}\footnote{This is the British tradition. See Nennius, \emph{Historia Brittonum},

31

and 66

(\emph{Chron. min.} III. pp171\texthebrew{‑}209). The Saxon tradition, recorded in the \emph{Saxon Chronicle}, places the coming of the Saxons as permanent settlers in 449. \emph{Chron. Gall.} 126, p660, has the following entry: Britanniae usque ad hoc tempus variis cladibus eventibusque latae (late vexatae, Mommsen) in dicionem Saxonum rediguntur. The date given is the 19th year of the joint rule of Theodosius and Valentinian = A.D. 442\texthebrew{‑}443. (The argument of Freeman, \emph{ib.} p158, is spoiled by his reckoning it as the 18th year of Theodosius after the death of Arcadius.) A little later we have the appeal of the Britons for help to Aetius in Gaul recorded by Gildas (\emph{De excidio Britanniae}, c. 20), Agitio ter consuli gemitus Britannorum. A.D. 446 was the third consul­ship of Aetius. These notices taken together look as if the Saxons, having gained some footholds about 428, during the following fourteen years extended their power, and then about 442 Roman rule definitely disappeared. See

Bury, \emph{The Not. dig.}, \emph{J. R. S.} X . Cp. also W. M. F. Petrie, \emph{Neglected British History}, 1917. It is to be noted that communications between Britain and the continent were not broken off during the fifth century. Germanus, bishop of Auxerre, who had been sent there by the Pope in 429 to contend with the Pelasgian heresy (Prosper, \emph{sub a.}), and is said to have gained a bloodless victory over the Saxons and Picts near St. Albans (Constantius, \emph{Vit. Germ.} c. 17), visited the island a second time probably about 440 (\emph{ib.} c. 25). See Levison, ``Bischof Germanus von Auxerre,'' in \emph{Neues Archiv}, XXIX (1903). We have evidence too of communications in 475 (Sidonius Apoll. \emph{Epp.} IX.9.6.).

} And, in the second place, we have definite evidence that in or not long after that year there was a field army there under the Count of the Britains.\texthebrew{⁠}\footnote{The fact that the Imperial officials in Britain are all recorded in the

\emph{Not. dig., Occ.}

(\emph{c.} A.D. 428) would not be decisive, as they might not have been erased unless Britain had been definitely handed over by treaty to another power. But there is one section, VII. (\emph{Distributio numerorum}), which has been brought up to date, and here we find, under the comes Britanniarum, three numeri of infantry and six vexillationes, of which at least four and probably more are not recorded in the lists of the field forces which are under the supreme commands of the mag. ped. and the mag. eq. praes. (in sections V and VI). This must mean that these forces had been sent to Britain comparatively recently and had been entered under VII but not under V and VI. See Bury, \emph{op. cit.}

} At this time the Empire

p202
was hard set to maintain its authority in Gaul and Spain and Africa, and it could not attempt to reinforce or keep up to strength the regiments in Britain. But there is no reason to suppose that during the last ten years of the reign of Honorius, and for some time after, Roman government in Britain was not carried on as usual. Its gradual collapse and final disappearance belong to the reign of Valentinian III.

In these years of agony many British provincials fled from the terror-stricken provinces and sought a refuge across the sea in the north-western peninsula of Gaul. Maritime Armorica received a new Celtic population and a new name, Brittany, the lesser Britain.\footnote{See Freeman, \emph{op. cit.} 162 \emph{sqq.} We do not know whether any of the German invaders who crossed the Rhine in 406 had penetrated to Armorica. The enemies from whom we are told that Armorica suffered in the days of Constantine III were probably the Saxon pirates who infested the Channel and the western coast of Gaul. The Armoricans like the Britons resorted to self-help.

Zosimus, VI.5.2.

}

The Visigoths were far from sharing in the philo-Roman proclivities of Athaulf. Their new king Wallia was animated by a national Gothic spirit and was not disposed at first to assume a pacific attitude towards Rome. A Spaniard two years later\texthebrew{⁠}\footnote{Orosius, VII.43 .

} informs us that ``he was elected by the Goths just for the purpose of breaking the peace, while God ordained him for the purpose of confirming it.'' Circumstances forced him into becoming a Federate of Rome, for he found his position in Spain untenable. The other barbarians had occupied most of the peninsula except Tarraconensis, and the Visigoths were unable to settle there because Roman ships blockaded the ports and hindered them from obtaining supplies. They were threatened by famine. To Wallia now, as to Alaric before, Africa seemed the solution of the difficulty, and he marched to the south of Spain (early in A.D. 416). But it was not destined that the Goths should set foot on African soil. As the fleet of Alaric had been wrecked in the straits of Sicily, even so some of the ships which Wallia had procured were shattered in the straits of Gades, and whether from want of troops or from

p203
superstitious fear he abandoned the idea. He decided that the best course was to make peace, and he entered into negotiations with Constantius.

Placidia, though still retained as a hostage, had been well treated, and her brother and lover were willing to treat with Wallia as they would not have treated with Athaulf. An agreement was concluded by which the Emperor undertook to supply the Goths with 600,000 measures of corn, and Wallia engaged to restore Placidia and to make war in the name of the Empire against the barbarians in Spain (before June, A.D. 416).

These engagements were carried out. After five years spent among the Goths, as captive and queen, Placidia returned to Italy,\texthebrew{⁠}\footnote{She was escorted by Euplutius, an agens in rebus who had conducted the negotiations. Olympiodorus, \emph{fr.} 31.

} and she was persuaded, against her own wishes, to give her hand to the Patrician Constantius. They were married on January 1, A.D. 417, the day on which he entered on his second consul ­ship.\footnote{He was consul again in 420, and in that year Symmachus the Prefect of Rome put up some monument in his honour, of which the dedicatory inscription is preserved (CIL VI.1719 ). He is there described as reparatori reipublicae et parenti invictissimorum principium \texthebrew{—} comiti et magistro utriusque militiae, patricio et tertio cons. ordinario. At Trier is preserved a memorial of his second consul­ship: an inscription copied on stone (in the twelfth century) probably from one of his consular ivory diptychs (CIL XIII.3674 ), Fl. Constantius v. c. comes et mag. utriusq. mil. atq. patricius et secundo consul ordinarius.

}

Wallia set about the congenial task of making war on the four barbarian peoples who had crossed the Pyrenees seven years before and entered the fair land of Spain, rich in corn and crops, rich in mines of gold and precious stones. For two years they seem to have devastated it far and wide. Then they settled down with the intention of occupying permanently the various provinces. The Siling Vandals, under their king Fredbal, took Baetica in the south; the Alans, under their king Addac, made their abode in Lusitania, which corresponds roughly to Portugal;\texthebrew{⁠}\footnote{Hydatius, our chief authority for Spain in these years, says Lusitaniam et Carthaginiensem; but we may question whether Carth. was occupied as a whole.

} the Suevians, and the Asding Vandals, whose king was Gunderic, occupied the north-western province of Gallaecia north of the Douro. The eastern provinces of Tarraconensis and Carthaginiensis, though the western districts may have been seized, and though they were doubtless constantly harried by raids, did not pass under the power of the invaders.

p204
Wallia began operations by attacking the Silings in Baetica. Before the end of the year he had captured their king by a ruse and sent him to the Emperor. The intruders in Spain were alarmed, and their one thought was to make peace with Honorius, and obtain by formal grant the lands which they had taken by violence. They all sent embassies to Ravenna. The obvious policy of the Imperial Government was to sow jealousy and hostility among them by receiving favourably the proposals of some and rejecting those of others.\texthebrew{⁠}\footnote{I infer this from what actually happened, combined with the naïve statement of Orosius

(VII.43)

that all the barbarian kings had made representations to Honorius that he should allow them to fight it out in Spain, as their mutual slaughter would be to the interest of the Empire: tu cum omnibus pacem habe omniumque obsides accipe; nos nobis confligimus, nobis perimus, tibi vincimus, immortali vero quaestu reipublicae tuae, si utrique pereamus. When Orosius was writing this last chapter of his work (for which see below,

Chap. IX, § 6 ), the war was still raging between the Visigoths and their foes, and the latest news was that Wallia was strenuously working for the establishment of peace, apparently early in 418. He was writing in Africa.

} The Asdings and the Suevians appear to have been success ­ful in obtaining the recognition of Honorius as Federates, while the Silings and Alans were told that their presence on Roman soil would not be tolerated. Their subjugation by Wallia was a task of about two years.\texthebrew{⁠}\footnote{Sidonius Apollinaris, celebrating Wallia's grandson Ricimer, writes (\emph{Carm.} II.363) :

Tartesiacis avus huius Vallia terris

Vandalicas turmas et iuncti martis Halanos

stravit et occiduam texere cadavera Calpen.

} The Silings would not yield, and they were virtually exterminated. The king of the Alans was slain, and the remnant of the people who escaped the sword of the Goths fled to Gallaecia and attached themselves to the fortunes of the Asding Vandals. Gunderic thus became ``King of the Vandals and Alans,'' and the title was always retained by his successors.

After these success ­ful campaigns, the Visigoths were recompensed by receiving a permanent home. The Imperial government decided that they should be settled in a Gallic not a Spanish province, and Constantius recalled Wallia from Spain to Gaul. A compact was made by which the whole rich province of Aquitania Secunda, extending from the Garonne to the Loire, with parts of the adjoining provinces (Narbonensis and Novempopulana), were granted to the Goths. The two great cities on the banks of the Garonne, Bordeaux and Toulouse, were handed over to Wallia. But Narbonne and the Mediterranean coast were reserved for the Empire. As Federates the Goths had no

p205
authority over the Roman provincials, who remained under the control of the Imperial administration. And the Roman proprietors retained one-third of their lands; two-thirds were resigned to the Goths. Thus, from the point of view of the Empire, south-western Gaul remained an integral part of the realm; part of the land had passed into the possession of Federates who acknowledged the authority of Honorius; the provincials obeyed, as before, the Emperor's laws and were governed by the Emperor's officials. From the Gothic point of view, a Gothic kingdom had been established in Aquitaine, for the moment confined by restraints which it would be the task of the Goths to break through, and limited territorially by boundaries which it would be their policy to overpass. Not that at this time, or for long after, they thought of renouncing their relation to the Empire as Federates, but they were soon to show that they would seize any favourable opportunity to increase their power and extend their borders.

This final settlement of the Visigoths, who had moved about for twenty years, in the three peninsulas of the Mediterranean, to find at last a home on the shores of the Atlantic, was a momentous stage in that process of compromise between the Roman Empire and the Germans which had been going on for many years and was ultimately to change the whole face of western Europe. Constantius was doing in Gaul what Theodosius the Great had done in the Balkans. There were now two orderly Teutonic kingdoms on Gallic soil under Roman lordship, the Burgundian on the Rhine, the Visigothic on the Atlantic.

Wallia did not live to see the arrangements which he had made for his people carried into effect. He died a few months after the conclusion of the compact, and a grandson of Alaric\texthebrew{⁠}\footnote{See Sidonius, \emph{Carm.} VII.505 . There seems no reason why avus should not be understood literally, if we assume that Alaric was born \emph{c.} 360 A.D.

} was elected to the throne, Theoderic I ( A.D. 418). Upon him it devolved to superintend the partition of the lands which the Roman proprietors were obliged to surrender to the Goths. It must have taken a considerable time to complete the transfer. The Visigoths received lion's share. Each landlord retained one-third of his property for himself and handed over the remaining portion to one of the German strangers.\texthebrew{⁠}\footnote{See the fragments of laws of Euric in \emph{Leges Visig. ant.} p3.

} This arrangement

p206
was more favourable to the Goths than arrangements of the same kind which were afterwards made in Gaul and Italy, as we shall see in due course, with other intruders. For in these other cases it was the Germans who received one-third, the Romans retaining the larger share. And this was the normal proportion. For the principle of these arrangements was directly derived from the old Roman system of quartering soldiers on the owners of land. On that system, which dated from the days of the Republic, and was known as hospitalitas , the owner was bound to give one-third of the produce of his property to the guests whom he reluctantly harboured. This principle was now applied to the land itself, and the same term was used; the proprietor and the barbarian with whom he was compelled to share his estate were designated as host and guest ( hospites ).

This fact illustrates the gradual nature of the process by which western Europe passed from the power of the Roman into that of the Teuton. Transactions which virtually meant the surrender of provinces to invaders were, in their immediate aspect, merely the application of an old Roman principle, adapted indeed to changed conditions. Thus the process of the dismemberment of the Empire was eased; the transition to an entirely new order of things was masked; a system of Federate States within the Empire prepared the way for the system of independent states which was to replace the Empire. The change was not accomplished without much violence and continuous warfare, but it was not cataclysmic.

The problem which faced the Imperial Government in Gaul was much larger than the settlement of the Gothic nation in Aquitaine. The whole country required reorganisation, if the Imperial authority was to be maintained effectively as of old in the provinces. The events of the last ten years, the ravages of the barbarians, and the wars with the tyrants had disorganised the administrative system. The lands north of the Loire, Armorica in the large sense of the name, had in the days of the tyrant Constantine been practically independent, and it was the work of Exuperantius to restore some semblance of law and order in these provinces.\texthebrew{⁠}\footnote{Rutilius Nam., writing in 417, says

(\emph{De red. suo}, I.213) :

cuius Aremoricas pater Exuperantius oras

nunc postliminium pacis amare docet.

leges restituit, libertatemque reducit

et servos famulis non sinit esse suis.

Freeman suggested that Exuperantius was Praet. Pref. Germanus, who

(p207)
became in 418 bishop of Auxerre, seems to have been in the preceding years Dux tractus Armoricani et Nervicani, a military command which extended over five provinces (the two Aquitaines, and 2nd, 3rd, and 4th Lugdunensis). This is the natural identification of his ducatus (Constantius, \emph{Vit. Germ.} I. c. 1 p202), since his authority ran in Sens and Auxerre which were in Lugd. Quarta.

} Most of the great cities in the south and

p207
east had been sacked or burned or besieged. We saw how Imperial Trier, the seat of the Praetorian Prefect, had been captured and plundered by the Vandals; since then it had been, twice at least, devastated by the Franks with sword and fire.\texthebrew{⁠}\footnote{Apparently about A.D. 410\texthebrew{‑}412. Renatus Frigeridus (in Greg. Tur. \emph{Hist. Fr.} II.9 ): Treverorum civitas a Francis direpta incensaque est secunda irruptione; Salvian, \emph{De gub. Dei}, VI. c. 15, ter excisa, but VII. c. 2, quadruplici eversione prostrata.

} The Prefect of the Gauls translated his residence from the Moselle to the Rhone, and Arles succeeded to the dignity of Trier.

What Constantius and his advisers did for the restoration of northern Gaul is unknown, but the direction of their policy is probably indicated by the measure which was adopted in the south, in the diocese of the Seven Provinces. On April 17, A.D. 418, Honorius issued an edict enacting that a representative assembly was to meet every autumn at Arles, to debate questions of public interest. It was to consist of the seven governors of the Seven Provinces,\texthebrew{⁠}\footnote{It is provided that the governors (iudices) of Aquitania Sec. and Novempopulana, on account of their distance from Arles, might send deputies.

} of the highest class of the decurions,\texthebrew{⁠}\footnote{Honorati, retired decurions.

} and of representatives of the landed proprietors. The council had no independent powers; its object was to make common suggestions for the removal of abuses or for improvements in administration, on which the Praetorian Prefect might act himself or make representations to the central government. Or it might concert measures for common action in such a matter as a petition to the Emperor or the prosecution of a corrupt official.\footnote{We shall meet an instance in the prosecution of Arvandus: below,

Chap. X § 4 . º

}

Such a council was not a new experiment. The old provincial assemblies of the early Empire had generally fallen into disuse in the third century, but in the fourth we find provincial assemblies in Africa, and diocesan assemblies in Africa and possibly in Spain.\texthebrew{⁠}\footnote{See

\emph{C. Th.} XII.12.1 and 9 ; Guiraud, \emph{Les Assemblées provinciales dans l'empire romain}, p228.

} Already in the reign of Honorius a Praetorian Prefect, Petronius, had made an attempt to create a diocesan assembly in Southern Gaul, probably in the hope that time and labour might be saved, if the affairs of the various provinces

p208
were all brought before him in the same month of the year. The Edict of A.D. 418 was a revival of this idea, but had a wider scope and intention. It is expressly urged that the object of the assembly is not merely to debate public questions, but also to promote social intercourse and trade. The advantages of Arles \texthebrew{—} a favourite city of Constantine the Great, on which he had bestowed his name, Constantina \texthebrew{—} and its busy commercial life are described. ``All the famous products of the rich Orient, of perfumed Arabia and delicate Assyria, of fertile Africa, fair Spain, and brave Gaul, abound here so profusely that one might think the various marvels of the world were indigenous in its soil. Built at the junction of the Rhone with the Tuscan sea, it unites all the enjoyments of life and all the facilities of trade.''\footnote{The edict is addressed to Agrippa, the Pr. Pr. of Gaul. It was not included in the Theodosian Code, but has been preserved as a separate document in several MSS. The text will be found in Sirmond's ed. of Sidonius Apollinaris (ed. 2, 1659, p241), in Hänel's \emph{Corpus legum} (p238), and other collections; and also in Carette, \emph{Les Assemblées prov. de la Gaule romaine}, p460 (in this book a very full discussion will be found).

}

It must also have been present to the mind of Constantius that the Assembly, attracting every year to Arles a considerable number of the richest and most notable people from Aquitania Secunda and Novempopulana, would enable the provincials, surrounded by Visigothic neighbours, to keep in touch with the rest of the Empire, and would help to counteract the influence which would inevitably be brought to bear upon them from the barbarian court of Toulouse.

The prospect of a return to peace and settled life in Spain seemed more distant than in Gaul. Soon after the Visigoths had departed, war broke out between Gunderic, king of the Vandals, and Hermeric, king of the Suevians. The latter were blockaded in the Nervasian mountains, but suddenly Asterius, Count of the Spains,\texthebrew{⁠}\footnote{The military command in Spain, with the title comes Hispaniarum, was new and must have been established after the invasion of the barbarians in 409. The first mentions of it are in

\emph{Not. Occ.} VII.118

and in Hydatius 74. Asterius was created a Patrician in reward for his success (Renatus, in Gregory of Tours, \emph{H.F.} II.9).

} appeared upon the scene, and his operations compelled the Vandals to abandon the blockade. At Bracara a large number were slain by the Roman forces. Then the Vandals and Alans, who now formed one nation, left Gallaecia and migrated to Baetica. On their way they met the Master of Soldiers,

p209
Castinus,\texthebrew{⁠}\footnote{Castinus is designated as mag. mil., not as mag. utr. mil., in the sources. This may mean that after the elevation of Constantius in 421 (see below) Castinus was appointed mag. ped. praes., along with a co-ordinate mag. equit. praes. We find in 423 Crispinus mag. equit. in

\emph{C. Th.} II.23.1

(where Seeck and Sundwall are surely wrong in reading Castino). In 419, or 420, Castinus was Count of the Domestics and led a campaign against the Franks (Renatus, \emph{ib.}).

} who had come from Italy to restore order in the peninsula. He had a large army, including a force of Visigothic Federates, but he suffered a severe defeat, partly through the perfidious conduct of his Gothic allies. The Vandals established themselves in Baetica, but it does not appear whether the recognition they had received in Gallaecia as a Federate people was renewed when they took up their abode in the southern province ( A.D. 422).\footnote{For these events see Hydatius 77 and Prosper \emph{sub} 422.

}

When the Patrician Constantius had been virtual ruler of the western provinces of the Empire for ten years and had been for four a member of the Imperial family as the Emperor's brother-in\texthebrew{‑}law , Honorius was persuaded, apparently against his own wishes, to co-opt him as a colleague. On February 8, A.D. 421, Flavius Constantius was crowned Augustus,\texthebrew{⁠}\footnote{The day of the month of his elevation, and that of his death, come from Theophanes, A.M. 5913.

} and immediately afterwards the two Emperors crowned Galla Placidia as Augusta. Two children had already been born to Constantius, the elder Justa Grata Honoria ( A.D. 417 or 418) and the younger Placidus Valentinianus (July 3, A.D. 419).\footnote{Marcellinus, \emph{sub a.} Honoria was called Justa Grata after her mother's maternal aunts, sisters of Galla.

}

But the achievement of the highest dignity in the world was attended by a bitter mortification. The announcement of his elevation and that of Placidia was sent in the usual way to Constantinople, but Theodosius and his sister Pulcheria refused to recognise the new Augustus and Augusta. Their reasons for this attitude are not clear. Perhaps they had never forgiven Placidia for her marriage with Athaulf, and perhaps they had some idea of reuniting the whole Empire under the sway of Theodosius when his uncle died, and saw in Placidia's son Valentinian, on

p210
whom the title of nobilissimus was bestowed,\texthebrew{⁠}\footnote{Honorius reluctantly yielded to the pressure of Placidia to confer the title, whether before or after the death of Constantius. For the conjecture as to the project of Theodosius see Güldenpenning, \emph{op. cit.} 240.

} an obstacle to the project. Constantius, writhing under this insult, thought of resorting to arms to force the eastern court to recognise him.\texthebrew{⁠}\footnote{Olympiodorus, \emph{frs.} 34, 38, 39, is our source for the last years of Constantius.

} In other ways too he found the throne a disappointment. The restraints surrounding the Imperial person were intolerably irksome to him; he was not free to go and come as he used when he was still in a private state. His popularity, too, had dwindled, for during the last few years he had grown grasping and covetous. His health failed, and after a reign of seven months he died (September 2).\footnote{Olympiodorus adds that after his death petitions came in from all sides complaining of unjust acts he had committed to extort money.

}

After his death, Honorius, who had always been fond of his step-sister , displayed his affection by kisses and endearments which were embarrassing for her and caused considerable scandal. The love, however, was presently turned into hatred through the machinations of Placidia's attendants;\texthebrew{⁠}\footnote{Her old nurse Elpidia, a maid Spadusa, and Leonteus her curator or intendant, are mentioned. Olymp. \emph{fr.} 40.

} and the estrangement between the Emperor and his sister led to frays in the streets of Ravenna between the parties who espoused their causes. Goths who had accompanied the widow of Athaulf from Spain and remained in her service, and retainers of her second husband, fought for her name and fame. Castinus, the Master of Soldiers, was her enemy; we may conjecture that he hoped to succeed to the power and authority of Stilicho and Constantius. The breach widened, and at length Placidia, with her two children, was banished from Ravenna, and sought refuge with her kindred at Constantinople ( A.D. 423).\texthebrew{⁠}\footnote{Prosper, \emph{sub a.}

} There was a rumour that Honorius suspected her of appealing to an enemy power to come to her assistance.\texthebrew{⁠}\footnote{Cassiodorus, \emph{Chron., sub a.}

} If there is any truth in this, we may guess that ``enemies'' to whom she appealed were the Visigoths.

The reign of Honorius came to an end a few months later. He died of dropsy\texthebrew{⁠}\footnote{Philostorgius, XII.13 ; \emph{Narr. de imp. dom. Val.} p630.

} on August 15, A.D. 423. His name would be forgotten among the obscurest occupants of the Imperial throne were it not that his reign coincided with the fatal period

p211
in which it was decided that western Europe was to pass from the Roman to the Teuton. A contemporary, who was probably writing at Constantinople,\texthebrew{⁠}\footnote{\emph{Ib.} The writer was an admirer of Theodosius II and probably wrote soon after the death of Honorius.

} observed that many grievous wounds were inflicted on the State during his reign. Rome was captured and sacked; Gaul and Spain were ravaged and ruined by barbarian hordes; Britain had been nearly lost. It was significant of the state of the times that a princess of the Imperial house should be taken into captivity and should deign to marry a barbarian chieftain.\texthebrew{⁠}\footnote{The curious expression used of Placidia's marriage, statum temporum decolorat, indicates the criticism which her act evoked in the east.

\texthebrew{◂} previous

next \texthebrew{▸}Images with borders lead to more information.

The thicker the border, the more information.

(Details here.)

UP TO:

J. B. Bury
Homepage

Latin \& Greek

Texts

LacusCurtius

Home} The Emperor himself did nothing of note against the enemies who infested his realm, but personally he was extraordinarily fortunate in occupying the throne till he died a natural death and witnessing the destruction of the multitude of tyrants who rose up against him.

\xchapter{Chapter VII}

When Arcadius died his son Theodosius was only seven years old.\texthebrew{⁠}\footnote{Born April 10, 401; crowned Augustus Jan. 10, 402. For the children of Arcadius see the genealogical table of the house of Theodosius. º On the will of Arcadius, under which the Persian king Yezdegerd is said to have been appointed guardian of Theodosius, see below,

Chap. XIV § 1 . º

} Anthemius, the Praetorian Prefect of the East, acted as regent,\texthebrew{⁠}\footnote{We do not know by what legal form this was arranged or whether others were associated in the regency. For Anthemius see above,

p159 . In 408 he was made a Patrician. Chrysostom wrote to congratulate him on the Praetorian Prefecture, saying that the office was more honoured by his tenure than he by the office (\emph{Ep.} 147).

} while Antiochus, a palace eunuch, was entrusted with the care of the young prince. The guidance of the State through the first critical years of the new reign showed the competence of the regent. The measures which were passed during the six years in which he held the power exhibit an intelligent and sincere solicitude for the general welfare. The name of Anthemius is chiefly remembered for its association with the great western land wall of Constantinople, which was built under his direction and has been described in an earlier chapter.\texthebrew{⁠}\footnote{See

Chap. III .

} But this was only one of many services that he performed for the Empire. Harmony was established between the courts of Constantinople and Ravenna and, while this was rendered possible by the death of Stilicho, it must be ascribed largely to the efforts and policy of Anthemius. A new treaty was made which secured peace on the Persian frontier.\texthebrew{⁠}\footnote{\emph{C. J.} V.63.4 .

} An invasion of Lower Moesia by Uldin, the king of the Huns, who had executed Gaïnas, seemed at first serious and mena­cing, but was successfully repelled.\texthebrew{⁠}\footnote{Sozomen, IX.5.

} An

p213
immense horde of Sciri were in the Hun's host, and so many were taken prisoners that the government had some trouble in disposing of them. They were given to large landowners in Asia Minor to be employed as serfs. In order to secure the frontier against future invasions of Hun or German barbarians, Anthemius provided for the improvement of the fleet stationed on the Danube; many new ships were built to protect the borders of Moesia and Scythia, and the old crafts were repaired.\footnote{\emph{C. Th.} VII.17.1

(Jan. 28, 412). The Danube boats were called lusoriae. The flotillas are enumerated in

\emph{Not. dig., Or.}

For the Sciri see

\emph{C. Th.} V.4.3 .

}

Constantinople depended on Egypt for its bread, and it sometimes happened that there was a lack of transport ships at Alexandria and the corn º supplies did not arrive at the due time.\texthebrew{⁠}\footnote{Sometimes a dishonest skipper sold his cargo at some remote place. See

\emph{C. Th.} XIII.5.33 .

} This occurred in A.D. 408, and there was famine in the city. The populace was infuriated, and burned the house of Monaxius, the Prefect of the City, whose duty it was to distribute the corn.\texthebrew{⁠}\footnote{Marcellinus, \emph{Chron., sub} 409. \emph{Chron. Pasch., sub} 407.

} Anthemius and the Senate did their utmost to relieve the distress by procuring corn elsewhere,\texthebrew{⁠}\footnote{\emph{C. Th.} XIV.16.1 .

} and then Anthemius made permanent provision for a more efficient organisation of the supplies from Egypt.\texthebrew{⁠}\footnote{A.D. 409. The responsibility was transferred from the navicularii or naval collegia, to the summates of the fleets, whose recompense for their trouble was increased by the addition of a small remuneration. The island of Carpathus was the half-way station between Alexandria and Byzantium, and thus the care of the corn supplies now devolved conjointly on the Prefect of the City, the Prefect of Egypt, and the praeses insularum (the governor of the Islands along the coast of Asia Minor; he was subordinate to the Proconsul of Asia).

\emph{C. Th.} XIII.5.32

(Jan. 19).

} He also took measures to revive the prostrate condition of the towns of the Illyrian provinces, which had suffered sorely through the protracted presence of Alaric and his Visigoths.\texthebrew{⁠}\footnote{\emph{C. Th.} XII.1.177

(A.D. 412).

} Towards the close of his tenure of office, all the fiscal arrears for forty years ( A.D. 368\texthebrew{‑}407) were remitted in the provinces of the eastern Prefecture.\texthebrew{⁠}\footnote{A.D. 414, April 9.

\emph{C. Th.} XI.28.9 .

} It is interesting to observe that the most intimate friend and adviser of Anthemius is said to have been Troilus, a pagan sophist of Side, who seems to have been the leader of a literary circle at Constantinople.\footnote{Socrates, VII.1. Anthemius was celebrated by Theotimus, a pagan poet

(Synesius, \emph{Epp.} 49) . Synesius calls Anthemius \textgreek{το\texthebrew{ῦ} μεγάλου} ( \emph{Epp.} 73 , addressed to Troilus) and cp. CIL III.737 magno Anthemio.

}

In her sixteenth year Pulcheria was created Augusta (July 4, A.D. 414),\texthebrew{⁠}\footnote{Coins of Ael. Pulcheria, with salus reipublicae on the reverse, belong to the years 414\texthebrew{‑}421, before her brother's marriage. They may have been struck in 415 when Theodosius celebrated his third quinquennalia and issued coins with Gloria reipublicae vot. XV. mult. XX. (Cp. de Salis, \emph{Coins of the Eudoxias.})

} and assumed the regency in the name of her brother, who was two years younger than herself. Anthemius soon disappeared from the scene; we may conjecture that death removed him; and he was succeeded in the Prefecture of the East by Aurelian, who in the preceding reign had been the leader of the Roman party in resisting the designs of Gaïnas.\texthebrew{⁠}\footnote{According to the date of

\emph{C. Th.} VIII.4.26 , Anthemius was still Prefect on Feb. 17, 415. But according to \emph{Chron. Pasch., sub a.} , he was succeeded by Aurelian before Dec. 30, 414.

} It seems probable that he was the chief adviser of Pulcheria.

One of her first acts was to remove from the court the eunuch Antiochus,\texthebrew{⁠}\footnote{Cp. Theophanes, A.M. 5905. Antiochus is said to have been sent to Constantinople by King Yezdegerd, in order to fulfil the duties of guardian which he had accepted under the will of Arcadius. See below,

Chap. XV § 1 . º

} who had been her brother's tutor. She superintended and assisted in the education of Theodosius. It is said that she gave him special instruction in deportment; and she sought to protect him from falling under the influence of intriguing courtiers to which his weak character might easily have rendered him a prey. The new mode of palatial life, established in the reign of Arcadius, enabled women to make their influence increasingly felt in public affairs. The example had been set by Eudoxia, and throughout the whole space of the fifth and sixth centuries we meet remarkable ladies of the imperial houses playing prominent parts. The daughters of Eudoxia were unlike their mother, and the court of Theodosius II was very different from that of Arcadius. The princesses Pulcheria, Arcadia, and Marina, and the young emperor inherited the religious temperament of their father, with which Pulcheria combined her grandfather's strength of character. The court, as a contemporary says, assumed the character of a cloister, and pious practices and charitable works were the order of the day. Pulcheria resolved to remain a virgin, and prevailed upon her sisters to take the same resolution, in which they were confirmed by their spiritual adviser, the Patriarch Atticus, who wrote for them a book in praise of virginity.

p215
Theodosius had studious tastes, and he formed a remarkable collection of theological books,\texthebrew{⁠}\footnote{See Socrates, VII.22, who devotes a chapter to his virtues.

} but he was also interested in natural science including astronomy. He was of a gentle and kindly nature, and it is recorded that he was reluctant to inflict capital punishment.\texthebrew{⁠}\footnote{John Ant. \emph{fr.} 71, in \emph{Exc. de Virt.} p204.

} He seems to have possessed none of the qualities of a capable ruler either in peace or war.\footnote{Tillemont has some just remarks on the defects in his character, \emph{Hist. des Empereurs}, VI.23 \emph{sqq.}

}

To an unprejudiced observer in the reign of Arcadius it might have seemed that the Empire in its eastern parts was doomed to a speedy decline. One possessed of the insight of Synesius might have thought it impossible that it could last for eight hundred years more when he considered the threatening masses of barbarians who encompassed it, the oppression of the subjects, and all the evils which Synesius actually pointed out. The beginning of the fifth century was a critical time for the whole Empire. At the end of the same period we find that while the western half had been found wanting in the day of its trial, the eastern half had weathered the storm; we find strong and prudent Emperors ruling at New Rome. The improvement began in the reign of Theodosius. The truth is that this Emperor, though weak like his father, was far more intelligent, and had profited more by his education. Throughout the greater part of his reign the guidance of affairs seems to have been in the hands of prudent ministers who maintained the traditions of Anthemius and Aurelian. In the chronicles we do not hear much about the Senate; everything is attributed to Pulcheria or Theodosius. But it seems probable that the Senate exercised considerable influence on the policy of the rulers. The State was not threatened in this reign by the danger of a military dictator ­ship, and it was only towards its close that an unworthy eunuch enjoyed undue political power.

Soon after her accession to the responsibilities of government the young Empress was called upon to deal with serious troubles which had arisen in Egypt. The old capitals, Alexandria and Antioch, although they had been overshadowed by the greatness of Byzantium, were far from degenerating into mere provincial towns. They retained much of their old importance and all their old characteristics. In Alexandria, in the fifth century,

p216
with its population of perhaps 600,000 citizens,\texthebrew{⁠}\footnote{300,000 is the number of the citizens given for the time of Augustus

(Diodorus, XVII.52) . It excludes slaves and foreigners. Güldenpenning (p225) thinks it must have been nearly twice as much in the fifth century. Cp. above,

Chap. III § 5, p88 .

} life was as busy, as various, and as interesting as ever. The Romans had found no city in the Empire so difficult to govern as that of the quick-witted and quick-tempered Alexandrians; the streets were continually the scene of tumults between citizens and soldiers, and revolts against the Augustal Prefects. ``While in Antioch, as a rule, the matter did not go beyond sarcasm, the Alexandrian rabble took on the slightest pretext to stones and cudgels. In street uproar, says an authority, himself Alexandrian, the Egyptians are before all others; the smallest spark suffices here to kindle a tumult. On account of neglected visits, on account of the confiscation of spoiled provisions, on account of exclusion from a bathing establishment, on account of a dispute between the slave of an Alexandrian of rank and the Roman foot-soldier as to the value or non-value of their respective slippers, the legions were under the necessity of charging among the citizens of Alexandria.''\footnote{Mommsen, \emph{Hist. of Rome}, V. (II.264 Eng. tr.).

}

Instead of healing the discords and calming the intractable temper of this turbulent metropolis by diffusing a spirit of amity and long-suffering , Christianity only gave the citizens new things to quarrel about, new causes for tumult, new formulae and catchwords which they could use as pretexts for violence and rioting.

The troubles which agitated Alexandria, when Pulcheria became regent, were principally due to the bigotry and ambition of the Patriarch. In this office, Theophilus, whom we met as the enemy of Chrysostom, had been succeeded ( A.D. 412) by his nephew Cyril, who was no less ambitious to elevate the prestige of his see and was even more unscrupulous in the arts of intrigue. In the first years of his pontificate his chief objects were to exalt his own authority above that of the civil governor of Egypt, the Augustal Prefect, and to make Alexandria an irreproachably Christian city by extirpating paganism which still flourished in its schools, and by persecuting the Jews who for centuries had formed a large minority of the population. He was an ecclesiastical tyrant of the most repulsive type,

p217
and the unfortunate Hypatia was the most illustrious of his victims.

Hypatia was the daughter of Theon, a distinguished mathematician,\texthebrew{⁠}\footnote{His most important studies were on Euclid, Aratus, and Ptolemy. Nearly all our MSS. of the geometry of Euclid are based on his critical recension, and the scholia on Aratus, whom he exalted as an astronomer above Eudoxus, are derived from him. The character of his work has been elucidated by Heiberg and Maass. Hypatia wrote three mathematical books, (1) a memoir on Diophantus (who wrote a standard work on arithmetic of which about half is extant); (2) a commentary on the Conic Sections of Apollonius; (3) a commentary on the astronomical Canon (\emph{\textgreek{καν\texthebrew{ὼ}ν βασιλει\texthebrew{ῶ}ν}}) of Ptolemy. See the article on \emph{\textgreek{\texthebrew{Ὑ}πατία}} in Suidas, which is largely based on the Life of Isidore by Damascius (for the reconstruction of which see the study of J. Asmus in \emph{B.Z.} XVIII.424 \emph{sqq.}, XIX.265 \emph{sqq.}). The statement of Suidas that Hypatia was the wife of Isidore was due to a misunderstanding of his source. Palladas, the contemporary Alexandrian poet, wrote the following very poor verses on Hypatia

(\emph{Anth. Pal.} IX.400) :

\textgreek{\texthebrew{ὅ}ταν βλέπω σε, προσκυν\texthebrew{ῶ}, κα\texthebrew{ὶ} το\texthebrew{ὺ}ς λόγους, τ\texthebrew{ῆ}ς παρθένου τ\texthebrew{ὸ}ν ο\texthebrew{ἶ}κον \texthebrew{ἀ}στρ\texthebrew{ῷ}ον βλέπων· ε\texthebrew{ἰ}ς ο\texthebrew{ὐ}ραν\texthebrew{ὸ}ν γ\texthebrew{ὰ}ρ \texthebrew{ἐ}στι σο\texthebrew{ῦ} τ\texthebrew{ὰ} πράγματα, \texthebrew{Ὑ}πατία σέμνη, τ\texthebrew{ῶ}ν λόγων ε\texthebrew{ὐ}μορφία, \texthebrew{ἄ}χραντον \texthebrew{ἀ}στρ\texthebrew{ὸ}ν τ\texthebrew{ῆ}ς σοφ\texthebrew{ῆ}ς παιδεύσεως.}} who was a professor at the Museum or university of Alexandria. Trained in mathematics by her father, she left that pure air for the deeper and more agitating study of metaphysics, and probably became acquainted with the older Neoplatonism of Plotinus\texthebrew{⁠}\footnote{Plotinus and his master Ammonius Sacas belonged to the university, while the later Neoplatonists were not connected with it. This point \texthebrew{—} Hypatia's affiliation to Plotinus \texthebrew{—} is due to W. A. Meyer, whose careful little tract, \emph{Hypatia von Alexandria} (1886), has thrown much light on the subject. Hoche (in his article in \emph{Philologus}, XV.439 \emph{sqq.}, 1860) showed that the supposed journey of Hypatia to Athens is based on a mistranslation of Suidas. The date of her birth was probably about 370.

} which, in the Alexandrian Museum, had been transmitted untainted by the later developments of Porphyrius and Iamblichus. When she had completed her education she was appointed to the chair of philosophy, and her extraordinary talents, combined with her beauty, made her a centre of interest in the cultivated circles at Alexandria, and drew to her lecture-room crowds of admirers. Her free and unembarrassed intercourse with educated men and the publicity of her life must have given rise to many scandals and backbitings, and her own sex doubtless looked upon her with suspicion, and called her masculine and immodest. She used to walk in the streets in her academical gown ( \textgreek{τρίβων} , the philosopher's cloak) and explain to all who wished to learn, difficulties in Plato or Aristotle.\texthebrew{⁠}\footnote{I follow Meyer's translation of a passage in Suidas. The most pleasing passage in Socrates is that in which he speaks with admiration of Hypatia (\emph{H.E.} VII.15).

} Of the influence of her personality on her pupils we have still a record in some letters of Synesius

p218
of Cyrene, who, although his studies under her auspices did not hinder him from adopting Christianity, always remained at heart a semi-pagan , and was devotedly attached to his instructress. That some of her pupils fell in love with her is not surprising,\texthebrew{⁠}\footnote{One of her pupils is said to have declared his passion for her, and the tale went that she exorcised his desire by disarranging her dress and displaying \textgreek{τ\texthebrew{ὸ} σύμβολον τ\texthebrew{ῆ}ς \texthebrew{ἀ}καθάρτου γεννήσεως}: ``This, young man,'' she said, ``is what you are in love with, and nothing beautiful.'' This story, recorded by Suidas, was without doubt a contemporary scandal, and indicates what exaggerated stories were circulated about the independence and perhaps the free-spokenness of Hypatia. Seven letters of Synesius to ``the philosopher Hypatia'' are preserved. He addresses her

(\emph{Ep.} 16)

as ``mother, sister, and teacher.''

} but Hypatia never married.

The cause of the tragic fate, which befell her in March A.D. 415, is veiled in obscurity. We know that she was an intimate friend of the pagan Orestes, the Prefect of Egypt; and she was an object of hatred to Cyril, both because she was an enthusiastic preacher of pagan doctrines and because she was the Prefect's friend.

The hatred of the Jews for the Patriarch brought the strained relations between Cyril and Orestes to a crisis. On one occasion, seeing a notorious creature of Cyril present in an assembly, they cried out that the spy should be arrested, and Orestes gratified them by inflicting public chastisement on him. The menaces which Cyril, enraged by this act, fulminated against the Jews led to a bloody vengeance on the Christian population. A report was spread at night that the great church was on fire, and when the Christians flocked to the spot the Jews surrounded and massacred them. Cyril replied to this horror by banishing all Hebrews from the city and allowing the Christians to plunder their property, a proceeding which was quite beyond the Patriarch's rights, and was a direct and insulting interference with the authority of Orestes, who immediately wrote a complaint to Constantinople. At this juncture 500 monks of Nitria, sniffing the savour of blood and bigotry from afar, hastened to the scene. These fanatics insulted Orestes publicly, one of them hitting him with a stone; in fact the governor ran a serious risk of his life.\texthebrew{⁠}\footnote{It is to be remembered that the Aug. Prefect did not possess military powers. Subsequently some Prefects united civil and military functions (Florus under Marcian, Alexander under Leo I), but these cases were exceptional. Cp. M. Gelzer, \emph{Byz. Verw. Ägyptens}, p19.

} The culprit who hurled the missile was executed, and Cyril treated his body as the remains of a martyr.

p219
It was then that Hypatia fell victim in the midst of these infuriated passions. One day as she was returning home she was seized by a band of parabalani \texthebrew{⁠}\footnote{We find the form \textgreek{παραβαλανε\texthebrew{ῖ}ς} in Mansi, VI. p828.

} or lay brethren, whose duty it was to tend the sick and who were under the supervision of the Patriarch. These fanatics, led by a certain Peter, dragged her to a church and, tearing off her garments, hewed her in pieces and burned the fragments of her body.\texthebrew{⁠}\footnote{\textgreek{\texthebrew{ὀ}στράκοις \texthebrew{ἀ}νε\texthebrew{ῖ}λον} (Socrates, VII.14), killed her with either sharp shards or mussel shells. Gibbon (V.117) misunderstood \textgreek{\texthebrew{ἀ}νε\texthebrew{ῖ}λον} when he interpreted, ``her flesh was scraped from her bones.'' Philostorgius

(VIII.9)

says that she was torn in pieces (\textgreek{διασπασθ\texthebrew{ῆ}ναι}) by the Homousians.

} The reason alleged in public for this atrocity was that she hindered a reconciliation between Orestes and Cyril; but the true motive, as Socrates tells us, was envy. This ecclesiastical historian does not conceal his opinion that Cyril was morally responsible.

There can be no doubt that public opinion was deeply shocked not only in Alexandria but also in Constantinople. Whatever Pulcheria and Atticus may have thought, the Praetorian Prefect Aurelian, who was the friend of her friend Synesius, must have been horrified by the fate of Hypatia. It would seem that the Empress found it impossible to act on the partial and opposite reports which were received from Orestes and Cyril, and a special commissioner, Aedesius, was sent to Alexandria to investigate the circumstances and assign the guilt. We have no direct information concerning his inquiry, but it would appear that it was long drawn out and it was publicly recognised that the parabalani were dangerous. The government consequently reduced the numbers of their corporation, forbade them to appear at games or public assemblies, and gave the Prefect authority over them.\texthebrew{⁠}\footnote{\emph{C. Th.} XVI.2.42 , A.D. 416, Sept. 29. It was suspected that Aedesius was bribed by Cyril and his party, Suidas, \emph{s.v.} \textgreek{\texthebrew{Ὑ}πατία.}} But within little more than a year the influence of Cyril at the pious court of Pulcheria elicited a new decree, which raised the number of the parabalani from 500 to 600 and restored them to the Patriarch's authority.\texthebrew{⁠}\footnote{\emph{C. Th.} \emph{ib.} 43 , A.D. 418, Feb. 3.

} If condign punishment had been inflicted on the guilty we should probably have heard of it. The obscure murderers may have escaped, but ``the murder of Hypatia has imprinted an indelible stain on the character and religion of Cyril of Alexandria.''\texthebrew{⁠}\footnote{Gibbon, \emph{ib.}

Thayer's Note: For an understandably very different view of Cyril, and reasons that have been put forth for absolving him of the murder, see the Catholic Encyclopedia article

St. Cyril of Alexandria .

} He was an

p220
able theologian and we shall next meet him in the stormy scene of an ecumenical Council.

We are not told at what time the regency of Pulcheria formally came to an end. Perhaps we may suppose that on reaching the age of fifteen Theodosius was declared to have attained his majority. But for several years after his assumption of the supreme authority his sister continued to be the presiding spirit in affairs of state. The most influential minister during these years was probably Monaxius, who succeeded Aurelian as Praetorian Prefect of the East.\footnote{Before August 416; he held the post till 420.

}

Pulcheria chose a wife for her brother when he was twenty years of age. She seems to have been confident that her own influence would not be endangered. The story of the Athenian girl who was selected to share the throne of Theodosius was romantic.\texthebrew{⁠}\footnote{Gregorovius made \emph{Athenais} the subject of an interesting monograph (1882).

} Athenais was the daughter of Leontius, a pagan philosopher, and had been highly educated by her father in the pagan atmosphere of Athens. When he died, she had a dispute with her brothers about the inheritance of her father's property and she came to Constantinople to obtain legal redress. Her beauty and accomplishments won the notice and patronage of the Empress, who chose her as a suitable bride for the Emperor. She took the name of Eudocia and embraced Christianity. The marriage was celebrated on June 7, A.D. 421, and was followed by the birth of daughter, who was named Eudoxia after her grandmother.\texthebrew{⁠}\footnote{A.D. 422. Her full name was Licinia Eudoxia. It appears on those of her coins which were minted in Italy, after her marriage. She was created Augusta in her infancy, for she is so designated in Placidia's dedicatory inscription (see below,

p262 ), which belongs probably to \emph{c.} 426\texthebrew{‑}428. From the same inscription we learn that Eudocia had a son named Arcadius (born 423\texthebrew{‑}425?), who must have died very young; and Dessau is doubtless right (\emph{Insc. Lat.} 818) in holding that this child is the minor Arcadius mentioned in the Preface (l. 13) to the \emph{Cento} of Proba, a copy of which the writer of the Preface seems to have presented to Theodosius II. A second daughter was born later, Flaccilla, who died in 431 (Marcellinus, \emph{Chron., sub a.}; Nestorius, \textgreek{Πραγμ. \texthebrew{Ἡ}ρακλ}., tr. Nau, p331).\texthebrew{—} Coins of Ael. Eudocia Aug. are preserved which must have been issued soon after her coronation in Jan. 423, as the reverse legend is vot. XX mult. XXX. They correspond closely to coins of Theodosius, Pulcheria, and Honorius. As Theodosius kept his third quinquennalia in 415 (\emph{Chron. Pasch., sub a.}) , the presumption is that he celebrated his vicennalia in 420, and that in that year were issued these coins of himself, Pulcheria, and Honorius at Constantinople. The design of the reverse (a standing winged Victory holding a cross) on the coins of Eudocia differs from the others from having a star. We have

(p221)
also similar coins of Ael. Placidia Aug., with the star, evidently minted in 423 or 424, soon after her arrival at Constantinople (see below). Cp. de Salis, \emph{Coins of the Eudoxias.}

} In A.D. 423 (January 2) she was created Augusta.

p221
Though she was sincerely loyal to her new faith, wrote religious poems, and learned to interest herself in theology, she always retained some pagan leanings, and we may be sure that, when her influence began to assert itself, the strict monastic character of the court was considerably alleviated.

It was about this time that the Empress Placidia with her two children, driven from Ravenna by Honorius, came to Constantinople and sought the protection of their kinsfolk.\texthebrew{⁠}\footnote{See above,

p210 . º

} Then the news arrived that Honorius was dead, and the first care of the government was to occupy the port of Salona in Dalmatia.\texthebrew{⁠}\footnote{Socrates, VII.23. Epigraphic evidence indeed suggests that Salona was under Constantinople in 414\texthebrew{‑}415, see Jung, \emph{Römer und Romanen}, 186, n2.

} The event was then made public, and for seven days the Hippodrome was closed and Constantinople formally mourned for the deceased Emperor. The intervention of Theodosius at this crisis in the destinies of the west was indispensable, and two courses were open to him. He might overlook the claims of his cousin, the child Valentinian, son of the Augustus whom he had refused to recognise as a colleague, and might attempt to rule the whole Empire himself as his grandfather had ruled it without dividing the power. Or he might recognise those claims, and act as his cousin's protector. In either case there was fighting to be done, for a usurper, whose name was John, had been proclaimed Emperor at Ravenna. Theodosius and Pulcheria decided to take the second course and support the cause of Placidia and her son. It was an important decision. The eastern government was not blind to its own interests, and a bargain seems to have been made with Placidia that the boundary between the two halves of the Empire should be rectified by the inclusion of Dalmatia and part of Pannonia in the realm of Theodosius.\texthebrew{⁠}\footnote{The words patrui mei in

\emph{C. Th.} XI.20.5

need not point to the definite transference of the administration of Dalmatia in A.D. 424, for in that year Theodosius was sole Emperor. But the change was not regarded as definitely settled till the marriage of Valentinian and Eudoxia in 437. See below,

p226 .

} The measure of occupying Salona had been taken with a view

p222
to this change. It is probable that at the same time it was arranged that the future Emperor of the west should marry the infant daughter of the Emperor of the east. In any case Theodosius could contemplate a closer union between his own court and that of Ravenna, a union in which he would have the preponderating influence for about a dozen years to come during the minority of his cousin and the regency of his aunt; while he would have no direct responsibility for any further misfortunes which the western provinces might sustain from the rapacity of the German guests whom they harboured.

John, who had assumed the purple at Rome, was an obscure civil servant who had risen to the rank of primicerius notariorum .\texthebrew{⁠}\footnote{Renatus Frigeridus, in Gregory of Tours, \emph{H.F.} II.8. Was he the same John who was sent to negotiate with Alaric in 408? (above,

p176 ).

} It is evident that he owed his elevation to the party which was adverse to Placidia, and certain that he had behind him the Master of Soldiers Castinus, who had failed to win laurels in Spain,\texthebrew{⁠}\footnote{See above,

p209 .

} and was probably partly responsible for her exile. His envoys soon arrived at Constantinople to demand his recognition from the legitimate Emperor, and the answer of Theodosius was to banish them to places on the Propontis.\texthebrew{⁠}\footnote{Philostorgius, XII.13 . º

} Placidia was now recognised as Augusta, her son as nobilissimus\texthebrew{⁠}\footnote{Olympiodorus, \emph{fr.} 46; Marcellinus, \emph{sub} 424.

} \texthebrew{—} titles which Constantinople had refused to acknowledge when they had been conferred by Honorius; and the dead Constantius was posthumously accepted as a legitimate Augustus.\texthebrew{⁠}\footnote{This is shown by the fact that some laws issued in his name with Honorius and Theodosius were published in \emph{C. Th.} (\emph{e.g.} III.16.2 ); cp. Mommsen, \emph{C. Th.} p. ccxcvii.

} A large army was prepared against the usurper and placed under the command of Ardaburius, an officer of Alan descent, and his son Aspar. Placidia and her children accompanied the army, and at Thessalonica Valentinian was raised to the rank of Caesar ( A.D. 424).\texthebrew{⁠}\footnote{Probably towards the end of the year. Valentinian was designated consul (Flavius Placidus Valentinianus Caesar) as colleague of Theodosius for 425. John assumed the consul­ship in the west. See \emph{Fast. Cons., sub a.}

} When they reached Salona, the infantry under Ardaburius embarked and sailed across to the coast of Italy, and Aspar with the cavalry proceeded by land to Sirmium and thence over the Julian Alps to the great city of the Venetian march, Aquileia, of which they made themselves masters.\texthebrew{⁠}\footnote{Philostorgius, \emph{ib.} , Olympiodorus, \emph{ib.}, and Socrates, VII.23 are the chief sources.

} Here Placidia remained to await the issue of the struggle.

p223
Of the situation in Italy and the attitude of the Italians to the Emperor who had established himself at Ravenna we know nothing, except the fact that he was not acknowledged at Rome,\texthebrew{⁠}\footnote{This may be inferred from the issue of gold coins of Theodosius II at Rome, which may probably be assigned (so de Salis) to 424\texthebrew{‑}425. The Roman mint did not issue coins of John (for whose Ravenna coins see Cohen, VIII.207). The loyalty of Rome is also shown by an inscription of Faustus, Prefect of the City in 425, acknowledging the Caesarship of Valentinian (CIL VI.1677 ).

} although it was at Rome that he had assumed the purple. Castinus, whom one might have expected to play the leader's part, remained in the background; we are only told that he was thought to have connived at John's elevation.\texthebrew{⁠}\footnote{Prosper, \emph{sub} 423. He was consul in 424, and was not acknowledged in the east.

} But two younger men, whose names were to become more famous than that of the Master of Soldiers, were concerned in the conflict of parties. Boniface, an able soldier, who was perhaps already Count of Africa in A.D. 422, had been ordered to co-operate with Castinus in the ill-fated expedition against the Vandals in Spain, but he had quarrelled with the commander and returned to Africa.\texthebrew{⁠}\footnote{Cp. Prosper, \emph{sub} 422, and Hydatius. It is not quite clear whether Boniface seized the government of Africa without Imperial warrant, or, as seems more likely, he had received the appointment before his disobedience in refusing to go to Spain. The presence of an able military commander in Africa was urgently demanded by the hostilities of the Moors. See the discussion in Freeman, \emph{Western Europe}, 305 \emph{sqq.}

} We next find him espousing the cause of Placidia when she was banished by Honorius and helping her with money. He is not recorded to have taken any direct part in the conflict with John, but he could maintain the loyalty of Africa to the Theodosian house and could exercise influence by his control of the corn supplies. The other rising soldier who played a part in these events was Aetius, of whom we shall hear much more. He accepted the new Emperor and was appointed to the post as Steward of the Palace ( cura palatii ). When the news arrived that an eastern army was on its way to Italy, he was sent to Pannonia to obtain help for his master from the Huns. For this mission he was well qualified, as he had formerly lived among them as a hostage and was on friendly terms with their king.

Ardaburius had embarked at Salona, but his fleet was unfortunate, it was caught in a storm and scattered. The general himself, driven ashore near Ravenna, was captured by the soldiers of John. If the usurper had proceeded immediately against Aspar, he might have thwarted his enemies. But he

p224
did not take prompt advantage of his luck. He decided to wait for the arrival of the Hun auxiliaries whom Aetius had gone to summon to his aid.

Meanwhile Ardaburius employed the time of his captivity at Ravenna in forming connexions with the officers and ministers of the usurper and undermining their fidelity. He then succeeded in sending a message to his son, who waited uneasily and expectantly at Aquileia, bidding him advance against Ravenna without delay. Guided by a shepherd through the morasses which encompassed that city, the soldiers of Aspar entered it without opposition; some thought that the shepherd was an angel of God in disguise. John was captured and conducted to Aquileia, where Placidia doomed him to death. His right hand was cut off, and mounted on an ass he was exposed in the circus before his execution. Castinus, the Master of Soldiers, was banished.\footnote{The victory of Placidia must be placed in May or June. For on July 9 she issued a law at Aquileia restoring some ecclesiastical privileges which had been abolished by John. \emph{Sirmondianae}, 6; also

\emph{C. Th.} XVI.2.46 and 47 ;

XVI.5.62 and 63 . Cp. Seeck, \emph{Regesten}, p5, on these laws. Placidia and her son did not leave Aquileia before Aug. 6 (\emph{C. Th.} XVI.2.47

and

V.64 ). Philostorgius

(XII.13)

says that John reigned for a year and a half, a rough figure but, if he was elevated in Sept. 423, pointing to May as the date of his fall.

}

When all was over, Aetius arrived in Italy with 60,000 Huns; if he had come a few days sooner, the conflict would probably have had a different issue and the course of history would have been changed. At the head of this large army, Aetius was able to make terms for himself with the triumphant Empress. She was forced to pardon him and accept his services. The Huns were induced by a large donation of money to return to their homes.

Placidia then proceeded with her children to Rome, where Valentinian III was created Augustus on October 23, A.D. 425.\texthebrew{⁠}\footnote{Socrates, VII.24, \emph{Chron. Pasch., sub a.} On the date compare Tillemont, \emph{Hist. des Emp.} VI.621, Clinton, \emph{F.R., sub a.} Gold coins of Valentinian were issued in Constantinople, conjecturally in 426: on the reverse two Emperors, both nimbate, one large, the other small, with the legend Salus Reipublicae.

} Theodosius had himself started for Italy to crown his cousin with his own hand, but fell ill at Thessalonica, and empowered the Patrician Helion, the Master of Offices, to take his place. It seems certain that Valentinian's sister Honoria was crowned Augusta, if not on the same occasion, soon afterwards.\footnote{See below,

p288 . Helion had acted for the Emperor in conferring

(p225)
the Caesarship at Thessalonica and had doubtless accompanied Placidia to Italy. A mutilated metrical inscription at Sitifis in Mauretania would refer to the elevation of Valentinian if de Rossi's restoration were near the truth (CIL VIII.8481 ). It runs:

zzz Terra [about 16 letters]ni sidera regni

iam de . . . ans armorum fulmina condit

gra . . . tutela Valentinianus

. . . . . . . et Theodosius artem.

De Rossi proposed

fulgida conscendens terraeni s. r.

in 1,

Placidiae grandis tutela

in 3, and

pace fruens doctam exercet

in 4 (very improbable). Cp. Bücheler's note in \emph{Anth. Lat.} II.288.

}

p225
Ardaburius was rewarded for his success ­ful conduct of the war by the honour of the consul ­ship in A.D. 427. He and his son Aspar were the ablest generals Theodosius had, and their devotion to the Arian creed did not stand in the way of their promotion. Aspar received the consul ­ship in A.D. 434, when he was again commanding an army in the interests of Placidia, this time against a foreign foe, not against a rebel;\texthebrew{⁠}\footnote{In Africa. See below,

p248 .

} and we have an interesting memorial of the event in a silver disc, on which he is represented, a bearded man, with a sceptre in his left hand and a handkerchief in his raised right, presiding at the consular games.\texthebrew{⁠}\footnote{It was found near Florence and is preserved there. The inscription round the disc is: Fl. Ardabur Aspar vir inlustris com. et mag. militum et consul ordinarius. For a full description see W. Meyer, \emph{Zwei ant. Elf.} pp6\texthebrew{‑}7.

Thayer's Note: The plate was found in the Fosso Castione, a creek near Marsiliana in the comune of Manciano (GR) in Tuscany. A good photograph of it is found online at

the Law School of the University of Palermo .

} It was a more than ordinary honour that was paid to Aspar, for he was consul for the West, not for the East,\texthebrew{⁠}\footnote{The eastern consul of the year was Areobindus.

} and the designation may have been suggested by Placidia herself, who owed him much for his services in securing the diadem for her son.

Twelve years passed, and the marriage arranged between the cousins, Valentinian and Licinia Eudoxia, was, as we saw, celebrated at Constantinople, whither the bridegroom went for the occasion (October 29, A.D. 437).\texthebrew{⁠}\footnote{\emph{Chr. Pasch.}, and Prosper, \emph{sub a.} Coins were issued in honour of the occasion: on the face a full-faced bust of Theodosius, on the reverse three figures, Theodosius in the centre joining the hands of his daughter and Valentinian, with legend Feliciter Nubtiis. º

} Now, if not before, a considerable part of the Diocese of Illyricum \texthebrew{—} Dalmatia and Eastern Pannonia certainly \texthebrew{—} were transferred from the sway of Valentinian to the sway of Theodosius.\texthebrew{⁠}\footnote{Cassiodorus, \emph{Var.} XI.1.9 (Placidia) remisse administrat imperium . . . nurum denique sibi amissione Illyrici comparavit factaque est coniunctio regnantis divisio dolenda provinciis; Jordanes, \emph{Rom.} 329 datamque pro munere soceri sui totam Illyricum (sic). The totam of Jordanes does not authorise us to suppose with Tillemont (\emph{Hist. des Emp.} VI.75) that the

(p226)
cession included the provinces of Noricum or even all Pannonia. Dalmatia, Pannonia Secunda, and Valeria were probably ceded, and no more. Cp. Zeiller, \emph{Les Origines chrét. dans les prov. Dan.} pp6, 7.

} This political transaction

p226
was part of the matrimonial arrangement, and was looked upon as the price which Placidia paid for her daughter-in\texthebrew{‑}law . The new provinces were now controlled by the Praetorian Prefect of Illyricum, and his seat was transferred for some years from Thessalonica to Sirmium.\footnote{We learn this from a law of Justinian

(\emph{Nov.} xi) : cum enim in antiquis temporibus Sirmii praefectura fuerat constituta ibique omne fuerat Illyrici fastigium tam in civilibus quam in episcopalibus causis, postea autem Attilanis temporibus eiusdem locis devastatis Apraeemius º praefectus praetorio de Sirmitana civitate in Thessalonicam profugus venerat [c. A.D. 447, see below,

p275 ]. This prefect is otherwise unknown.

}

After the departure of her daughter the Empress probably felt lonely, and she undertook, in accordance with her husband's wishes, a pilgrimage to Jerusalem to return thanks to the Deity for the marriage of their daughter.\texthebrew{⁠}\footnote{See Socrates, VII.47. The following inscription, recorded as existing in the church of St. Peter ad vincula at Rome, seems also to refer to the fulfilment of a vow for Eudoxia's marriage:

Theodosius pater Eudocia cum coniuge votum,

Cumque suo supplex Eudoxia nomine solvit

(where cum suo nomine = suo nomine). De Rossi, II.1, p110.

} In this decision they seem to have been confirmed by a saintly lady of high reputation, Melania by name, a Roman of noble family, who had been forced into a repugnant marriage, and had afterwards, along with her husband, whom she converted to Christianity, taken up her abode at first in the land of Egypt, where she founded monastic houses, and then at Jerusalem. She had visited Constantinople to see her uncle Volusian, whom she converted before his death, and she exercised considerable influence with the Emperor and his household. The journey of Eudocia to Jerusalem (in spring, A.D. 438) was marked by her visit to Antioch, where she created a sensation by the elegant oration which she delivered, posing rather as one trained in Greek rhetoric and devoted to Hellenic traditions and proud of her Athenian descent, than as a pilgrim on her way to the great Christian shrine. Although there was a large element of theological bigotry both in Antioch and in Alexandria, yet in both these cities there was probably more appreciation of Hellenic style and polish than in Constantinople. The last words of Eudocia's oration brought down the house \texthebrew{—} a quotation from Homer,

p227
``I boast that I am of your race and blood.''\texthebrew{⁠}\footnote{Evagrius, \emph{H. E.} I.20 . The verse is an adaptation of \emph{Iliad}, VI.211. It has been suggested that Eudocia's oration consisted of a poem in hexameters (Ludwich, \emph{Eudociae fragmenta}, p12).

} The city that hated and mocked the Emperor Julian and his pagan Hellenism loved and fêted the Empress Eudocia with her Christian Hellenism; a golden statue was erected to her in the curia and one of bronze in the museum. Her interest in Antioch took a practical form, for she induced Theodosius to build a new basilica, restore the thermae, extend the walls, and bestow other marks of favour on the city.

Eudocia's visit to Aelia Capitolina, as Jerusalem was called, brings to the recollection the visit of Constantine's mother Helena, one hundred years before, and, although Christianity had lost some of its freshness in the intervening period, it must have been a strange and impressive experience for one whose youth was spent amid pagan memories in the gardens of the philosophers at Athens, who in New Rome, with its museums of ancient art and its men of many creeds, had not been entirely weaned from the ways and affections of her youth, to visit, with all the solemnity of an exalted Christian pilgrim, a city whose memories were typically opposed to Hellenism, and whose monuments were the bones and relics of saints.\texthebrew{⁠}\footnote{Of the relics which she received (the bishop of Jerusalem plied a trade in relics), especially remarkable were the chains with which Herod bound Peter. One of these she gave to her daughter Eudoxia, who founded a church in Rome (called originally after herself, and in later times St. Peter ad vincula), where it is still preserved. Cp. above,

p226, n2 . An account of Eudocia's visit to Jerusalem will be found in the \emph{Vita Melaniae iunioris.} Melania met the Empress at Sidon and acted as her companion and cicerone.

Thayer's Note: For comprehensive details on the church, see

S. Petri ad Vincula

in Hülsen's \emph{Chiese di Roma nel Medio Evo}, and the further references there, which include links to the full texts of Armellini and Titi.

} It was probably only this religious side that came under Eudocia's notice; for Jerusalem at this period was a strange mixture of piety with gross licence. We are told by an ecclesiastical writer of the age that it was more depraved than Gomorrah; and the fact that it was a garrison town had something to do with this depravity. But it drew pilgrims from all quarters of the world.

On her return from Palestine ( A.D. 439) Eudocia's influence at Court was still power ­ful.\texthebrew{⁠}\footnote{In this year, the 32nd º of his reign, Theodosius was consul for the 17th time, and the mint of Constantinople issued gold coins (1) of the Emperor with a helmeted Rome on the reverse and legend IMP xxxii. cos xvii. PP, (2) of the Empress, with Constantinople seated on the prow of a vessel and the same legend.

} She seems to have been on terms of intimate friendship with Cyrus of Panopolis, who held a very exceptional position. He filled at the same time the two high

p228
offices of Praetorian Prefect of the East and Prefect of the city.\texthebrew{⁠}\footnote{That Cyrus held these offices simultaneously is expressly stated by John Lydus, \emph{De mag.} II.12 , and by John Malalas, XIV.361. Malalas says that he held them for four years. It is probable that the source of this record was Priscus, see \emph{Chron. Pasch., sub} 439 . We know from Theodosius, \emph{Nov.} 18, that he was Pr. Pr. Or. in Nov. 439; and from

\emph{C. J.} VIII.11.21 , that he was Pr. Urb. in Jan. 440.

} He was a poet like his fellow-townsman Nonnus though of minor rank;\texthebrew{⁠}\footnote{John Lydus, \emph{ib.}, says contemptuously that he knew nothing except poetry. Some epigrams and short poems are extant. The most interesting of these is

\emph{Anth. Pal.} IX.136 , written before leaving the city in exile:

Would that my father had taught me to tend his flock in the pastures,

Where sitting under the shade of elm-trees or rocks overhanging

Sweetly piping and reeds I would charm dull care with my music.

O Pierian maids, let us flee from the fair-built city

Forth to another land. And there will I tell of the mischief

Wrought by the baleful drones to the bees who toil for the honey.

The first verse is imitated by Nonnus,

\emph{Dionys.} XX.372 .

} he was a student of art and architecture; and he was a ``Hellene'' in faith. It has been remarked that Imperial officialdom was beginning to assume in the East a more distinctly Greek complexion in the reign of Theodosius II, and Cyrus was a representative figure in this transition. He used to issue decrees in Greek, an innovation for which a writer of the following century expressly blames him.\texthebrew{⁠}\footnote{John Lydus, \emph{ib.}

} His prefecture was popular and long remembered at Constantinople, for he built and restored many buildings and improved the illumination of the town, so that the people enthusiastically cried on some occasions in the Hippodrome, ``Constantine built the city but Cyrus renewed it.''\texthebrew{⁠}\footnote{For the building of the sea walls see above,

Chap. III .

} He still held his offices in the autumn of A.D. 441,\texthebrew{⁠}\footnote{\emph{C. J.} I.55.10 .

} but it could not be long after this that he fell into disgrace. Perhaps his popularity made him an object of suspicion; his paganism furnished a convenient ground for accusation. He was compelled to take ecclesiastical orders and was made bishop of Cotyaeum in Phrygia. His first sermon, which his malicious congregation forced him to preach against his will, astonished and was applauded by those who heard it:

``Brethren, let the birth of God, our Saviour, Jesus Christ be honoured by silence, because the Word of God was conceived in the holy Virgin through hearing only. To him be glory for ever and ever. Amen.''\footnote{The anecdote is told by John Malalas, \emph{ib.} The right reading \textgreek{\texthebrew{ὁ} το\texthebrew{ῦ} θεο\texthebrew{ῦ} λόγος} (for \textgreek{λόγ\texthebrew{ῳ}}) is preserved in the corresponding passage of Theophanes, A.M. 5937. For the opening words cp. below,

Chap. XI p349, n3 .

}

The friendship between Cyrus and the Empress Eudocia,

p229
who was naturally sympathetic with a highly educated pagan, suggests the conjecture that his disgrace was not unconnected with the circumstances which led soon afterwards to her own fall. We may conjecture that harmony had not always existed between herself and her sister-in\texthebrew{‑}law , and differences seem to have arisen soon after her return from Palestine.\texthebrew{⁠}\footnote{They differed on the Eutychian controversy, but there were doubtless other causes of jealousy.

} Discord was fomented by the arts of a eunuch, Chrysaphius Zstommas, who was at this time beginning to establish his ascendancy over the Emperor.\texthebrew{⁠}\footnote{These intrigues are related by Theophanes, A.M. 5940 = A.D. 447\texthebrew{‑}448. But the chronology of Theophanes during these years is full of errors. We know from Marcellinus and other sources that Eudocia had retired to Jerusalem in 444. John Malalas tells the story of Eudocia's life consecutively without chronological indications.

} Pulcheria had enjoyed the privilege of having in her household the Chamberlain ( praepositus Augustae ) who was officially attached to the service of the reigning Empress. It would not have been unnatural if this arrangement had caused jealousy in the heart of Eudocia, and we are told that Chrysaphius urged her to demand from the Emperor that a High Chamberlain should also be assigned to her. When Theodosius decidedly refused, she urged, again at the suggestion of Chrysaphius, that Pulcheria should be ordained a deaconess, inasmuch as she had taken a vow of virginity. Pulcheria refused to be drawn into a contest for power. She sent her Chamberlain to Eudocia and retired to the Palace of Hebdomon.\texthebrew{⁠}\footnote{This story appears in a curious form in John of Nikiu

(\emph{Chron.} LXXXVII.29\texthebrew{‑}33) , who thoroughly disliked Pulcheria.

} When Chrysaphius had succeeded in removing one Empress from the scene, his next object was to remove the other, so that his own influence over the weak spirit of Theodosius might be exclusive and undivided. In accomplishing this end he was probably assisted by the orthodox party at court, who were devoted to Pulcheria and looked with suspicion on the Hellenic proclivities of her sister-in\texthebrew{‑}law . The Emperor's mind was poisoned against his wife by the suggestion that she had been unduly intimate with Paulinus,\texthebrew{⁠}\footnote{We have no means of knowing whether there was any truth in this charge, but it should be observed that in Marcellinus, \emph{Chron., sub} 421, the true reading is Eudociam Achivam, not moecham (found in one MS.), so that this writer does not, as Güldenpenning thinks (\emph{op. cit.} p325), stigmatize her as unfaithful. Contemporary evidence for the charge of adultery has recently come to light in the \emph{Book of Heraclides} of Nestorius (tr. Nau, p331). The ex-Patriarch writes, ``the demon-prince of adultery, who had thrown the Empress into shame and disgrace, has just died.'' Cp. E. W. Brooks, \emph{B.Z.}, 21, 94\texthebrew{‑}95.

} a

p230
handsome man who had been a comrade of the Emperor in his boyhood.

This is probably the kernel of truth in the legend of Eudocia's apple which is thus told by a chronicler.\footnote{John Malalas, XIV. p356.

}

Whatever may have been the circumstances it seems that Paulinus, Master of Offices, was sent to Cappadocia and put to death by the Emperor's command in A.D. 444.\texthebrew{⁠}\footnote{This is the year to which the context in the passage of Nestorius points, and is confirmed by \emph{Chron. Pasch.} Marcellinus places the death of Paulinus in 440.

} It is credible that her former intimacy with Paulinus was used to alienate Theodosius from his wife, and she found her position so intolerable that at last she sought and obtained the Emperor's permission to withdraw from the Court and betake herself to Jerusalem ( A.D. 443).\texthebrew{⁠}\footnote{Cedrenus and Zonaras place Eudocia's visit to Jerusalem in the 42nd year of Theodosius, ``also 450 was ganz irrig ist,'' says Gregorovius (\emph{Athenais}, p187). But the 42nd year is reckoned from 402 (not from 408) and = Jan. 10, 443 to Jan. 10, 444. This was the official reckoning of

(p231)
his regnal years as appears from the coins which were issued in this very year: reverse: a seated Victory holding a cruciger globe, star underneath, and buckler on the ground behind, with legend Imp. xxxxii cos xvii PP. This shows that the 42nd year fell between the 17th consul­ship 439 and the 18th, 444, and therefore fell in 443. At the same time were minted coins of Eudocia, Pulcheria, Valentinian and Eudoxia with the same reverse. See De Salis, \emph{Coins of the Eudoxias}, and Sabatier, \emph{Monn. byz.} Pl. V.1, VI.1 and 11.

} She was not deprived of Imperial honours and an

p231
ample revenue was placed at her disposal. In Jerusalem she kept such state and was so energetic in public works that the jealousy of Theodosius was aroused and he sent Saturninus, the commander of his guards, to inquire into her activities. Saturninus slew the priest Severus and the deacon John who were confidants of the Empress.\texthebrew{⁠}\footnote{Marcellinus, \emph{Chron., sub} 444.

} She avenged this act by permitting the death of Saturninus; the words of one of our authorities might lead us to suppose that she caused him to be assassinated,\texthebrew{⁠}\footnote{Besides Marcellinus, Priscus, speaking of the heiress of Saturninus, says: \textgreek{τ\texthebrew{ὸ}ν δ\texthebrew{ὲ} Σατορνίλον} º \textgreek{\texthebrew{ἀ}ν\texthebrew{ῃ}ρήκει \texthebrew{Ἀ}θηνα\texthebrew{ὶ}ς} (\emph{fr.} 3, \emph{De leg. Rom.} p146). See the discussion of Gregorovius, \emph{op. cit.} cap. XXIII.

} but it has been suggested that officious servants or an indignant mob may have too hastily anticipated her supposed wishes. Then by the Emperor's command she was compelled to reduce her retinue.

The last sixteen years\texthebrew{⁠}\footnote{She died Oct. 20, 460. Cyrillus, \emph{Vita Euthymii}, p74.

} of the life of this amiable lady were spent at Jerusalem where she devoted herself to charitable work, built churches, monasteries and hospices, and restored the walls of the city.\texthebrew{⁠}\footnote{Evagrius, I.22 ;

John of Nikiu, LXXXVII.22, 23 .

} She was drawn into the theological storm which swept over the East in the last years of Theodosius, an episode which will claim our notice in another place. It is said that before her death she repeated her denial of the slander that she had been unfaithful to her husband.\footnote{\emph{Chron. Pasch., sub} 444.

}

The three most important acts of the reign of Theodosius II were the fortification of the city by land and sea, which has already been described, the foundation of a university, and the compilation of the legal code called after his name. It would be interesting to know whether the establishment of a school for higher education in the capital was due to the influence of the young Empress, who had been brought up in the schools of

p232
Athens. The new university (founded February 27, A.D. 425) was intended to compete with the schools of Alexandria and the university of Athens, the headquarters of paganism \texthebrew{—} with which, however, the government preferred not to interfere directly \texthebrew{—} and thereby to promote the cause of Christianity. Lecture-rooms were provided in the Capitol. The Latin language was represented by ten grammarians or philologists and three rhetors, the Greek likewise by ten grammarians, but by five rhetors; one chair of philosophy was endowed and two chairs of jurisprudence. Thus the Greek language had two more chairs than the Latin, and this fact may be cited as marking a stage in the official Graecisation of the eastern half of the Roman Empire.\footnote{\emph{C. Th.} XIV.9.3 , and

VI.21.1 . For the lecture-rooms in a portico in the Capitol see

\emph{C. Th.} XV.1.53 .

}

In the year 429 Theodosius determined to form a collection of all the constitutions issued by the ``renowned Constantine, the divine Emperors who succeeded him, and ourselves.'' The new code was to be drawn up on the mode of the Gregorian and Hermogenian codes,\texthebrew{⁠}\footnote{The Gregorian Code (c. A.D. 300) contained constitutions from Hadrian to A.D. 294; the Hermogenian those from 296 to 324.

Thayer's Note: For more details and references, see

the article

in Smith's \emph{Dictionary of Greek and Roman Antiquities}.

} and the execution of the work was entrusted to a commission of nine persons, among whom was Apelles, professor of law at the new university. Nine years later the work was completed and published, but during the intervening years the members of the commission had changed; of the eight who are mentioned in the edict which accompanied the final publication only two, Antiochus and Theodorus, were among the original workers, and a constitution of A.D. 435, which conferred full powers on the committee for the completion of the work, mentions sixteen compilers.\footnote{See

\emph{C. Th.} I.1.5 , March 26, 429, I.1.6, Dec. 20, 435.

}

The code was issued conjointly by Theodosius and Valentinian, and thus expressed the unity of the Empire (February 15, A.D. 438). The visit of the younger Emperor to Constantine on the occasion of his marriage with his cousin Eudoxia facilitated this co-operation . On December 23 of the same year, at a meeting of the Senate of Old Rome, the code which had been drawn up by the lawyers of New Rome was publicly recognised, and an official account of the proceedings on that occasion \texthebrew{—} gesta in senatu Urbis Romae de recipiendo Codice Theodosiano \texthebrew{—}p233
may still be read. The Praetorian Prefect and consul of the year, Anicius Acilius Glabrio Faustus, spoke as follows:

The felicity of the eternal Emperors proceeds so far as to adorn with the ornaments of peace those whom it defends by warfare. Last year when we loyally attended the celebration of the most fortunate of all ceremonies, and when the marriage had been happily concluded, the most sacred Prince, our Lord Theodosius, was fain to add this dignity also to his world, and ordered the precepts of the laws to be collected and drawn up in a compendious form of sixteen books, which he wished to be consecrated by his most sacred name. Which thing the eternal Prince, our Lord Valentinian, approved with the loyalty of a colleague and the affection of a son.

And all the senators cried out in the usual form, ``Well spoken!'' ( nove diserte, vere diserte ). But instead of following the course of the \emph{gesta} in the Roman senate-house , it will be more instructive to read the Imperial constitution which introduced the great code to the Roman world.

The code of Theodosius was superseded at the end of a hundred years by the Code of Justinian, and to the jurist it is less indispensable than to the historian. The historian must always remember with gratitude the name of Theodosius and that of Antiochus, if we may credit this minister with having originated the idea of the work. For the full record of legislation which it preserves furnishes clear and authentic information on the social

p235
conditions of the Empire, without which our other historical sources would present many insoluble problems.\footnote{The object of the compilers of the Code was to include all the laws, whether edicts or rescripts, which they could find, not to make a selection of those which were still valid. One might have thought that a record of all imperial laws would have been carefully preserved in the eastern and western chanceries, but it was not so. Seeck's valuable investigation of the sources of the Code (\emph{Regesten der Kaiser und Päpste}) shows that in many cases there were no copies at Constantinople, and the texts had to be sought at provincial centres, \emph{e.g.} at Berytus. Of much legislation there was probably no trace to be found anywhere. But laws issued in the west were more abundantly preserved than those in the east. It is remarkable that though the Code includes laws of Theodosius up to 437, it does not include laws of Valentinian after 432.

}

The last ten years of the reign were unfortunate. The Illyrian provinces suffered terribly from the depredations of the Huns, and the payments which a weak government made to buy off the invaders depleted the treasury.\texthebrew{⁠}\footnote{The gold paid to the Huns during the eight years A.D. 443\texthebrew{‑}450 exceeded in value £1,000,000.

} The eunuch Chrysaphius, having succeeded in removing from the Palace the rival influences of the Emperor's wife and sister, completely swayed the mind of his sovran and seems to have controlled the policy of the government. It is said, and we can easily believe it, that Theodosius at this time was in the habit of signing state papers without reading them.\footnote{Theophanes, A.M. 5942.

}

The power of Chrysaphius remained unshaken\texthebrew{⁠}\footnote{He had an enemy in the Isaurian Zeno, Master of Soldiers, who seems to have threatened a revolt in A.D. 449. See John Ant. \emph{fr.} 84 (\emph{De ins.}), and Priscus, \emph{fr.} 5 (\emph{De leg. Rom.}).

} until a few months before the Emperor's death, when he fell out of favour and the influence of Pulcheria again re-asserted itself.\texthebrew{⁠}\footnote{Theophanes, \emph{ib.}

} Theodosius died on July 28, A.D. 450, of a spinal injury caused by a fall from his horse.\footnote{The accident happened near the River Lycus not far from the city. See John Mal. XIV.366.

}

As Theodosius had no male issue and had not co-opted a colleague, the government of the eastern half of the Empire ought automatically to have devolved upon his cousin and western colleague Valentinian III. But this devolution would not have pleased Theodosius himself, and would not have been tolerated by his subjects. And we are told that on his death-bed

p236
Theodosius indicated a successor. Among the senators who were present on that occasion were Aspar, Master of Soldiers, and Marcian, a distinguished officer who had served as Aspar's aide-de\texthebrew{‑}camp in more than one campaign. The Emperor said to Marcian, ``It has been revealed to me that you will reign after me.''\texthebrew{⁠}\footnote{Marcellinus and \emph{Chron. Pasch., sub a.}

} We may conjecture that this choice had been arranged beforehand by Pulcheria and her brother. For Pulcheria agreed to become the nominal wife of Marcian, and thus the Theodosian dynasty was formally preserved.\footnote{At the beginning of his reign Marcian issued gold coins both of himself and of Pulcheria with a side-faced Victory holding a cross on the reverse and the legend Victoria Auggg. See Sabatier, \emph{Monn. byz.} Pl. VI.6 and 13. An inscription found in Eastern Thrace (CIL, III.14207) describes Marcian as serius in regnum missus (he was nearly 60 years old) and applying prompt remedies (celeri medicina) to restore a falling world.

}

Marcian was crowned in the Hebdomon by the Empress (August 25),\texthebrew{⁠}\footnote{\emph{Chron. Pasch., sub a.}

} and it is possible that on this occasion the Patriarch Anatolius took part in the coronation ceremony.\texthebrew{⁠}\footnote{Theophanes, A.M. 5942, \emph{ad fin.} \textgreek{μεταστέλλεται} (sc. Pulcheria) \textgreek{τ\texthebrew{ὸ}ν πατριάρχη κα\texthebrew{ὶ} τ\texthebrew{ὴ}ν σύγκλητον κα\texthebrew{ὶ ἀ}ναγορεύει α\texthebrew{ὐ}τ\texthebrew{ὸ}ν βασιλέα \texthebrew{Ῥ}ωμαίων}. We are ignorant what was the authority of Theophanes for introdu­cing the Patriarch. See below,

p317 . According to John Malalas, XIV.367, Marcian was crowned ``by the Senate''; according to John of Nikiu, ed. Zotenberg, p472, and Zonaras, XIII.24 , by Pulcheria; according to Simeon, the Logothete, vers. Slav. ed. Sreznevski, p50 (= Theodosius Mel. p78 = Leo Gramm. p111) by Anatolius. This last tradition is accepted by W. Sickel, \emph{BZ.} VII.517, 539.

} The first act of the new reign was the execution of Chrysaphius,\texthebrew{⁠}\footnote{John Mal. XIV.368. Marcellinus, \emph{ib.} Pulcheriae nutu interemptus est.

} and it is worthy of notice that Chrysaphius had favoured the Green faction of the Circus, and that Marcian patronised the Blues. His reign was a period of calm, all the more striking when it is contrasted with the storms which accompanied the dismemberment of the Empire in the west. In later times it was looked back to as a golden age.\texthebrew{⁠}\footnote{Theoph. A.M. 5946. \textgreek{κα\texthebrew{ὶ ἦ}ν \texthebrew{ἐ}κε\texthebrew{ῖ}να τ\texthebrew{ὰ ἐ}τη κυρίως χρυσ\texthebrew{ᾶ} τ\texthebrew{ῇ} το\texthebrew{ῦ} βασιλέως χρηστότητι}. Cp. John Lydus, \emph{De mag.} III.42 (p132) \textgreek{Μαρκιαν\texthebrew{ὸ}ν τ\texthebrew{ὸ}ν μέτριον.}} The domestic policy of Marcian was marked by financial economy, which was the more necessary, as during the last years of his predecessor the treasury was emptied by the large sums which were paid to the Huns.

Marcian refused to pay this tribute any longer, and at his death he left a well-filled treasury.\texthebrew{⁠}\footnote{More than £4,500,000. John Lydus, \emph{ib.}

} He accomplished this, not by imposing new burdens on the people, but by wisely regulating

p237
his expenditure. He alleviated the pressure of taxes so far as Roman fiscal principles would permit. He assisted his subjects from the exchequer when any unwonted calamity befell them. One of his first acts was a remission of arrears of taxation.\texthebrew{⁠}\footnote{Marcian, \emph{Nov.} 2 (A.D. 450).

} He confined the burdensome office of the praetor ­ship to senators resident in the capital.\texthebrew{⁠}\footnote{\emph{C. J.} XII.2.1 , A.D. 450.

} He decreed that the consuls instead of distributing money to the populace should contribute to keeping the city aqueduct in repair.\texthebrew{⁠}\footnote{Marcellinus, \emph{sub} 452.

} He attempted to put an end to the system of selling administrative offices.\texthebrew{⁠}\footnote{Theodore Lector, I.2.

} Perhaps the act which gave most satisfaction to the higher classes was the abolition of the follis , the tax of seven pounds on the property of senators.\footnote{\emph{C. J.} XII.2.2 . Cp. above,

p50 .

}

One of his enactments may perhaps be regarded as characteristic. Constantine the Great, in order to preserve the purity of the senatorial class, had declared illegal the marriage of a senator with a slave, a freed woman, an actress, or a woman of no social status ( humilis ). Marcian ruled that this law should not bar marriage with a respectable free woman, however poor, or however lowly her birth might be, and professed to believe that Constantine himself would have approved of this interpretation.\texthebrew{⁠}\footnote{Marcian, \emph{Nov.} 4 (A.D. 454).

} The Emperor's most confidential minister was Euphemius, the Master of Offices, whose advice he constantly followed.\texthebrew{⁠}\footnote{Priscus, \emph{fr.} 12, \emph{De leg. gent.} Palladius was Praet. Pref. of the East during the greater part of the reign (see \emph{Novels}, and other laws in \emph{C. J.}).

} While Marcian was not engaged in hostilities with any great power, there were slight troubles in Syria with the Saracens of the desert, and there was warfare on the southern frontier of Egypt. Since the reign of Diocletian Upper Egypt had been exposed to incursions of the Blemyes and the Nobadae. For the purposes of strengthening the defences of the frontier Theodosius II divided the province of Thebais into two (upper and lower), and united the civil and the military administration of the upper province in the same hands.\texthebrew{⁠}\footnote{This arrangement was probably made in the latter half of the reign. The title of the governor was, as elsewhere, dux; cp. a Leyden papyrus in \emph{Archiv f. Papyrusforschung}, I.399, \textgreek{κόμιτα κα\texthebrew{ὶ} δο\texthebrew{ῦ}κα το\texthebrew{ῦ} Θεβαικο\texthebrew{ῦ} λιμίτου}. In this passage the barbarians are mentioned \textgreek{τ\texthebrew{ῶ}ν \texthebrew{ἀ}λιτηρίων βαρβάρων . . . τ\texthebrew{ῶ}ν τε Βλεμύων κα\texthebrew{ὶ} τ\texthebrew{ῶ}ν Νουβάδων}. Cp. M. Gelzer. \emph{Studien zur byz. Ver. Ägyptens}, p10.

} At the beginning of Marcian's reign Florus held this post and distinguished

p238
himself by driving the barbarians who were again annoying the province back into the desert.\texthebrew{⁠}\footnote{Jordanes, \emph{Rom.} p43; Priscus, \emph{fr.} 11, \emph{De leg. gent.};

Evagrius, II.5 . It was in these raids probably that the exiled Patriarch Nestorius was captured by the barbarians at Oasis, see Evagrius

(I.7)

who quotes his letters. Fragments of a heroic poem on a war with the Blemyes, preserved on papyrus, are supposed by some to refer to the campaign of Florus. They have been edited most recently by A. Ludwich under the title of \emph{Blemyomachia}, but it is very doubtful to what historical events they refer. All the names of persons are fictitious (Persinoos, etc.) with the possible exception of Germanus. Ludwich thinks that the hostilities described are imaginary and, on metrical grounds, he regards the poem as considerably prior to A.D. 450.

} The Blemyes expressed a desire to conclude a definite treaty with the Empire and for this purpose they sent ambassadors to Maximin, who seems to have been Master of Soldiers in the East. Terms were arranged, and it was conceded to the Blemyes that they might at stated times visit Philae in order to worship in the temple of Isis, in which the policy of the Emperors still suffered the celebration of old pagan rites. But we are told that when Maximin soon afterwards died the predatory tribes renewed their raids.

The act for which the reign of Marcian is best remembered by posterity is the assembling of the Fourth Ecumenical Council at Chalcedon. The decisions of this council gave deep satisfaction to the Emperor and Empress; they could not foresee the political troubles to which it was to lead. Pulcheria died in A.D. 453.\texthebrew{⁠}\footnote{Marcellinus, \emph{sub a.}

} By a life spent in pious and charitable works she had earned the eulogies of the Church, and she left all her possessions to the poor. Among the churches which claimed her as foundress may be mentioned three dedicated to the Mother of God. One was known as the church of Theotokos in Chalkoprateia,\texthebrew{⁠}\footnote{Theodore Lector, I.5. See Bieliaev, \emph{Khram Bog. Khalkopr.}, p87, n2.

} so called from its situation in the quarter of the bronze merchants, not far from St. Sophia. The church of Theotokos Hodegetria,\texthebrew{⁠}\footnote{Nicephorus Callistus, XIV. cap. 11. The picture was said to be the work of St. Luke.

} Our Lady who leads to victory, which she built on the eastern shore of the city under the first hill, was sanctified by an icon of the Virgin which her sister-in\texthebrew{‑}law sent her from Jerusalem. More famous than either of these was the church which she founded shortly before her death at Blachernae. This sanctuary was deemed worthy to possess a robe of the Virgin, brought from Jerusalem in the reign of Marcian's successor, who built a special chapel to receive it.\footnote{\textgreek{\texthebrew{ἡ ἁ}γία σορός.}Cedrenus, I.614 .

}

p239
In later days the people of Constantinople put their trust in this precious relic as a sort of palladium to protect their city.

Marcian died in the first month of A.D. 457,\texthebrew{⁠}\footnote{Sometime between 26th January and 7th February (Clinton, \emph{F.R. sub a.}); possibly on 26th January; there is a lacuna in Theodore Lector, I.12, where the date is mentioned.

\texthebrew{◂} previous

next \texthebrew{▸}Images with borders lead to more information.

The thicker the border, the more information.

(Details here.)

UP TO:

J. B. Bury
Homepage

Latin \& Greek

Texts

LacusCurtius

Home} and with him the Theodosian dynasty, to which through his marriage he belonged, ceased to reign at New Rome.

\xchapter{Chapter VIII}

During the first twelve years of the reign of Valentinian, the Empress Placidia ruled the West, and her authority was not threatened or contested. Unbroken concord with her nephew Theodosius, who considered himself responsible for the throne of his young relative, was a decisive fact in the political situation and undoubtedly contributed to her security. The internal difficulties of her administration were caused by the rivalries of candidates not for the purple but for the Mastership of Both Services, the post which gave its holder, if he knew how to take advantage of it, the real political power.

The man whom Placidia chose to fill the supreme military command was Felix, of whose character and capacities we know nothing. He remained in power for about four years ( A.D. 425\texthebrew{‑}429),\texthebrew{⁠}\footnote{Flavius Constantius Felix was consul in 428, and we have portraits of him on the two leaves of his consular diptych. See Gori, \emph{Thes.} I p129. For a dedicatory inscription, in fulfilment of a vow, by him and his wife Padusia, see de Rossi, II.1, p149: Dessau, 1293. To Felix we must attribute the reorganisation of the defences of the Danubian provinces in A.D. 427\texthebrew{‑}428 (for which we find evidence in the

\emph{Not. dig.} ; see Seeck, \emph{Hermes}, XI.75 \emph{sqq.}), after the Huns restored Valeria, see below,

Chap. IX § 2 \emph{ad init.}

} and, so far as we know, did not leave Italy. He did not attempt to play the active and prominent part which had been played by Constantius and by Stilicho. The Germans, who had penetrated into the Empire, were the great pressing problem, and in the dealings with them during these four years it is not the name of Felix that history records, but those of the two

p241
subordinate officers whom we have seen taking opposite sides in the struggle for the throne of Honorius \texthebrew{—} Boniface and Aetius.

Flavius Aetius was the son of Gaudentius, a native of Lower Moesia,\texthebrew{⁠}\footnote{Aetius was born at Durostorum (Silistria).

} and an Italian mother. The career of his father, who fought with Theodosius the Great against the tyrant Eugenius, had been in the west, and Aetius had been given, in his childhood, as a hostage to Alaric,\texthebrew{⁠}\footnote{This fact is known from Merobaudes, \emph{Carm.} IV.46 \emph{sqq.}:

uix puberibus pater sub annis

obiectus Geticis puer cateruis,

bellorum mora, foederis sequester,

intentas Latio faces remouit

ac mundi pretium fuit pauentis;

and \emph{Pan.} II.129. The occasion may have been in 405\texthebrew{‑}406, or perhaps after the first siege of Rome in 408. Cp. above,

p180, n3 .

} and some years later had been sent, again as a hostage, to the Huns, among whom he seems to have remained for a considerable time, and formed abiding bonds of friendship with King Rugila. This episode in his life had a considerable effect upon his career.

A panegyrical description of this soldier and statesman, on whom the fortunes of the Empire were to lean for a quarter of a century, has come to us from the pen of a contemporary.\texthebrew{⁠}\footnote{Renatus Profuturus Frigeridus, in Gregory of Tours,

\emph{H. Fr.} II.8 . Mommsen, in his brief sketch of the career of Aetius, regards him as overrated (\emph{Hist. Schr.} I.531 \emph{sqq.}).

} He was ``of middle height, of manly condition, well shaped so that his body was neither too weak nor too weighty, active in mind, vigorous in limb, a most dexterous horseman, skilled in shooting the arrow, and strong in using the spear. He was an excellent warrior and famous in the arts of peace; free from avarice and greed, endowed with mental virtues, one who never deviated at the instance of evil instigation from his own purpose, most patient of wrongs, a lover of work, dauntless in perils, able to endure the hardships of hunger, thirst, and sleeplessness.''

That Aetius should take a German to wife was characteristic of the age in which an Imperial princess wedded a Goth and an Emperor was on the throne who had Frank blood in his veins. The lady was of royal Gothic family, ``a descendant of heroes,''\texthebrew{⁠}\footnote{Heroum suboles, Merobaudes, \emph{op. cit.}; Sidonius Apoll. \emph{Paneg. in Maior.} 126 \emph{sq.} Her father's name was Carpilio.

} and they had a son, Carpilio, who was old enough in A.D. 425 to be delivered as a hostage to the Huns.\footnote{Priscus, \emph{fr.} 3 (\emph{De leg. Rom.} p128);

Cassiodorus, \emph{Var.} I.4.11 .

}

It was to Aetius that the defence of Gaul was now entrusted; he commanded the field army and soon received the title of

p242
Magister Equitum.\texthebrew{⁠}\footnote{We are only told that Placidia conferred on him the title of count

(Philostorgius, XII.14) . For the new post of mag. eq. per Gallias see

Bury, \emph{The Not. dig.} (\emph{J. R. S.} X) .

} He had to defend the southern provinces against the covetous desires of the Goths, and the north-eastern against the aggressions of the Franks. King Theoderic was bent upon winning the Mediterranean coast adjacent to his dominion, and Aetius established his military reputation by the relief of Arles, to which the Goths had laid siege in A.D. 427.\texthebrew{⁠}\footnote{Prosper, \emph{sub} 425; \emph{Chron. Gall.} p658. A success won in 430 over Gothic forces near Arles, mentioned by Hydatius, 92, may be the battle of Mons Colubrarius recorded by Merobaudes, \emph{Pan.} 1.10 (Vollmer, \emph{ad loc.}).

} Hostilities continued, but a peace was made in A.D. 430 confining the Goths to the territories which had been granted to Wallia. On this occasion the Roman government gave hostages to Theoderic, and it has been suggested that at the same time the Goths were recognised as an independent power, the Roman governors were withdrawn from Aquitania Secunda and Novempopulana, and the Gallo-Roman inhabitants of those provinces passed under the direct rule of Theoderic.\texthebrew{⁠}\footnote{Schmidt (\emph{op. cit.} I.235) has adduced arguments for this view \texthebrew{—} among others the fact that Theoderic made laws affecting relations between Goths and provincials (referred to in Euric's Code: \emph{Leg. Vis. ant.} 277, cp. Sidonius Apoll. \emph{Epp.} II.1). He holds that in 453, after the accession of Theoderic II, the Goths again became foederati of the Empire (\emph{ib.} 252). But what exactly happened, how the legal position was changed he leaves very vague \texthebrew{—} inevitably, as there is no clear evidence.

} It may be doubted whether this change came about so early, but in any case the attitude of the Visigoths towards the Imperial government for the ensuing twenty years was that of an independent and hostile nation.

The Salian Franks had been living for nearly seventy years in the north-eastern corner of Lower Belgica, in the district known as Thoringia, where they had been settled as Federates by the Emperor Constantius II and Julian. In these lands of the Meuse and Scheldt they seem to have lived peacefully enough within the borders assigned to them by Rome. They were ruled by more than one king, but the principal royal family, which was ultimately to extinguish all the others, was the Merovingian. They seemed to be the least formidable of all the German peoples settled within the Empire, though they were destined to become the lords of all Gaul. The first step on the path of expansion seems to have been taken by Chlodio, the first of the long-haired Merovingian kings whose name is

p243
recorded. Taking advantage of the weakening of the Roman power, which was manifest to all, he invaded Artois. Aetius led an army against him and defeated him at Vicus Helenae, about A.D. 428.\texthebrew{⁠}\footnote{Hélesmes (Nord). º The source is
Sidonius Apoll. \emph{Carm.} V.212 \emph{sqq.} Cp. Prosper, \emph{sub a.} It is hardly to this campaign of Aetius that Merobaudes refers when he says (\emph{Pan.} II.6 \emph{sq.}) that the Rhine \texthebrew{—}Hesperiis flecti contentus habenis

gaudet ab alterna Thybrim sibi crescere ripa,

words which point to some pacification of the Middle Rhine, apparently an arrangement with the Ripuarian Franks between Cologne and Mayence, the two rivers alluded to being the Moselle and the Main, and probably made at a later date.

} But before his death Chlodio seems to have succeeded in extending his power as far as the Somme, crossing the Carbonarian Forest (the Ardennes) and capturing Cambrai.\texthebrew{⁠}\footnote{See Greg. Tur. \emph{H. Fr.} II.9 . His source was no doubt a Frank legend, and its historical value might be doubted, were it not borne out, so far as Chlodio's aggressive policy is concerned, by the incident related by Sidonius (last note).

} This annexation was probably recognised by the Imperial government; for the Salians remained federates of the Empire and were to fight repeatedly in the cause of Rome.

If the units of the field army with which Aetius conducted the defence of Gaul were up to their nominal strength, he had somewhat less than 45,000 men under his command. We do not know whether he had the help of the federate Burgundians in his operations against Visigoths and Franks. But it is certain that the most useful and effective troops, on whom he relied throughout his whole career in withstanding German encroachments in Gaul, were the Huns, and without them he would hardly have been able to achieve his moderate successes. Here his knowledge of the Huns, his friendship with the ruling family, and the trust they placed in him stood the Empire in good stead.

The prestige which Aetius gained in Gaul was far from welcome to the Empress Placidia, who never forgave him for his espousal of the cause of John. But now he was able to impose his own terms, and extort from her the deposition of Felix and his own elevation to the post which Felix had occupied. He was appointed Master of Both Services in A.D. 429, and it is said that he then caused Felix to be killed on suspicion of treachery.\footnote{The brief notices we have of these events only excite our curiosity. Prosper says that Felix was created a Patrician and succeeded by Aetius in 429, and was slain by Aetius on suspicion of treachery in 430 along with his wife Padusia. John Ant. (\emph{fr.} 85, \emph{De ins.} p126) says that Felix was suborned by Placidia to kill Aetius. Hydatius (94) says that Felix was killed in a military riot.

}

p244
It was, no doubt, the power of the Hunnic forces, which he could summon at his will, that enabled him to force the hand of the Empress. The one man whom she would have liked to oppose to him was Boniface, formerly her loyal supporter. Boniface had been for some time enacting the part of an enemy of the ``Republic.'' We must now go back to follow the fatal course of events in Africa.\footnote{In 430 and 431 Aetius was occupied in pacifying the Danubian provinces, Vindelicia
(Sidon. Apoll. \emph{Carm.} VII.234) , Raetia (\emph{ib.} 233 , Hydatius 93, \emph{Chron. Gall.} p658), Noricum (Hydatius 93, 95, Sidon. \emph{ib.}).

}

Africa, far from the Rhine and Danube, across which the great East-German nations had been pouring into the Roman Empire, had not yet been violated by the feet of Teutonic foes. But the frustrated plans of Alaric and Wallia were intimations that the day might be at hand when this province too would have to meet the crisis of a German invasion. The third attempt was not to fail, but the granaries of Africa were not to fall to the Goths. The Vandal people, perhaps the first of the East-German peoples to cross the Baltic, was destined to find its last home and its grave in this land so distant from its cradle.\footnote{For the difficult questions connected with the Vandal invasion of Africa and the part played by Boniface the most important modern discussions are those of Freeman, \emph{op. cit.}; Schmidt, in \emph{Geschichte der Wandalen}; Martroye, in \emph{Genséric.}

}

We saw how the Vandals settled in Baetica, and how King Gunderic assumed the title of ``King of the Vandals and the Alans.''\texthebrew{⁠}\footnote{This remained the official style of all the Vandal kings in Africa.

} He conquered New Carthage and Hispalis (Seville), and made raids on the Balearic Islands and possibly on Mauretania Tingitana.\texthebrew{⁠}\footnote{Hydatius, 86, 89.

} He died in A.D. 428 and was succeeded by his brother Gaiseric, who had perhaps already shared the kingship with him.\texthebrew{⁠}\footnote{Cp. Martroye, \emph{Genséric}, p103.

} About the same time events in Africa opened a new and attractive prospect to the Vandals.

After the restoration of the legitimate dynasty and the coronation of Valentinian,\texthebrew{⁠}\footnote{For his adhesion to Placidia see above,

p223 . He seems to have gone to Ravenna immediately after the restoration and to have received the additional dignity of comes domesticorum. Cp. Augustine, \emph{Ep.} 220 § 4 nauigasti. See Seeck, \emph{Bonifacius}, in \emph{P.\texthebrew{‑}W.}

} the conduct of Count Boniface laid him open to the suspicion that he was aiming at a tyranny himself.

p245
It had been a notable part of his policy, since he assumed the military command in Africa, to exhibit deep devotion to the Church and co-operate cordially with the bishops. He ingratiated himself with Augustine, the bishop of Hippo, and a letter of Augustine casts some welcome though dim light on the highly ambiguous behaviour of the count in these fateful years. Notwithstanding his professions of orthodox zeal, and hypocritical pretences that he longed to retire into monastic life, Boniface took as his second wife\texthebrew{⁠}\footnote{\emph{Ib.} Her name was Pelagia, Marcellinus, \emph{Chron., sub} 432. There is no positive evidence for the opinion of Baronius that she was a relative of the Vandal king (\emph{Ann. ecc., sub} 427).

} an Arian lady, and allowed his daughter to be baptized into the Arian communion. This degeneracy shocked and grieved Augustine, but it was a more serious matter that instead of devoting all his energies to repelling the incursions of the Moors, he was working to make his own authority absolute in Africa.\texthebrew{⁠}\footnote{Augustine, \emph{ib.} § 7. Prosper, \emph{sub} 427. Bonifatio cuius intra Africam potentia gloriaque augebatur. He seems to have enjoyed a high military reputation (Olympiodorus, \emph{fr.} 42), but the only exploit recorded, before his campaigns in Africa against the Moors (of which we know no details), is his defence of Marseilles against Visigoths in A.D. 413.

} So at least it seemed to the court of Ravenna, and Placidia \texthebrew{—} doubtless by the advice of Felix\texthebrew{⁠}\footnote{Prosper, \emph{loc. cit.} Procopius (\emph{B.V.} I.3) makes Aetius act the perfidious part of instigating the Empress against Boniface, and at the same time secretly advising Boniface to defy her. But Aetius was at this time almost certainly in Gaul. Cp. Freeman, \emph{op. cit.} p337.

} \texthebrew{—} recalled him to account for his conduct. Boniface refused to come and placed himself in the position of an ``enemy of the Republic.'' An army was immediately sent against him under three commanders, all of whom were slain ( A.D. 427). Then at the beginning of A.D. 428 another army was sent under the command of Sigisvult the Goth, who seems to have been named Count of Africa, to replace the rebel.\texthebrew{⁠}\footnote{Prosper, \emph{loc. cit.}

} Sigisvult appears to have succeeded in seizing Hippo and Carthage,\texthebrew{⁠}\footnote{Cp. Augustine, \emph{Collatio cum Maximino}, \emph{P. L.} 42.709; Possidius, \emph{Vit. Aug.} c17. Maximin was an Arian bishop who had come with Sigisvult.

} and Boniface, despairing of overcoming him by his own forces, resorted to the plan of inviting the Vandals to come to his aid.\footnote{Prosper, \emph{loc. cit.}, places the calling of the Vandals in 427 before the arrival of Sigisvult: exinde gentibus quae uti navibus nesciebant, dum a concertantibus in auxilium vocantur, mare pervium factum est, bellique contra Bonifatium coepti in Segisvultum comitem cura translata est. The story in Procopius, \emph{B. V.} I.3 enables us to interpret the vague plural a concertantibus as referring to Boniface. I follow Martroye (\emph{op. cit.} p87) in supposing that Boniface turned to the Vandals after the coming of Sigisvult. But it is of course possible that he took this step when the news

(p246)
reached him that the expedition was being prepared. The invitation of Boniface is also recorded by

Jordanes, \emph{Get.} 167, 169

(following Cassiodorus).

}

p246
The proposal of Boniface was to divide Africa between himself and the Vandals, for whom he doubtless destined the three Mauretanian provinces, and he undertook to furnish the means of transport.\texthebrew{⁠}\footnote{Procopius, \emph{loc. cit.}, where it is said that a tripartite division was contemplated between Boniface, Gunderic, and Gaiseric. This would imply that Gunderic died after Boniface's negotiations began.

} Gaiseric accepted the invitation. He fully realised the value of the possession of Africa, which had attracted the ambition of two Gothic kings. The whole nation of the Vandals and Alans embarked in May A.D. 429 and crossed over to Africa.\texthebrew{⁠}\footnote{Hydatius, 90. It is stated here that before he crossed to Africa, Gaiseric led an expedition against the Suevians who were plundering in Baetica or neighbouring regions. Martroye (p106) argues from this that Gaiseric intended to leave the non-combatant population in their Spanish home until his success in Africa was assured; that he was ready to start in 428, and that the Suevian invasion forced him to postpone his departure till 429; and that as a matter of fact the mass of the Vandals remained in Spain till after the capture of Carthage, when the Visigoths conquered the country (Cassiodorus, \emph{Chron., sub} 427, a Gothis exclusa de Hispaniis). This hypothesis runs counter to the evidence. Hydatius, \emph{ib.}, says that the Vandals embarked ``with their families,'' and so Victor Vitensis, \emph{Hist. Vand.} I.1. The notice in Cassiodorus might have some importance if it were under a later year. As we know nothing of the circumstances, we have no means of conjecturing why Gaiseric found it imperative to attack the Suevians at this juncture. And in any case the natural inference from the notice of Hydatius is that the defeat of the Suevians belongs to 429.

} If the population numbered, as is said, 80,000, the fighting force might have been about 15,000.\footnote{Victor Vit. \emph{op. cit.} I.2. This is to be preferred to the statement of Procopius, \emph{B. V.} I.5. See Schmidt, \emph{Gesch. der Wandalen}, p37, and \emph{B.Z.} XV.620\texthebrew{‑}621. Cp. also Martroye, \emph{op. cit.} p104. The reason for numbering the people before the migration from Spain to Africa was obviously to find out how many vessels would be needed, and non-combatants as well as combatants had to be transported.

}

Their king Gaiseric stands out among the German leaders of his time as unquestionably the ablest. He had not only the military qualities which most of them possessed, but he was also master of a political craft which was rare among the German leaders of the migrations. His ability was so exceptional that his irregular birth \texthebrew{—} his mother was a slave\texthebrew{⁠}\footnote{Sidonius Apoll. \emph{Carm.} 57 famula satus. Cp. Procopius, \emph{B. V.} I.3.

} \texthebrew{—} did not diminish his influence and prestige. We have a description of him, which seems to come from a good source. ``Of medium height, lame from a fall of º his horse, he had a deep mind and was sparing of speech. Luxury he despised, but his anger was uncontrollable and he was covetous. He was far-sighted in indu­cing foreign peoples to act in his interests, and resource ­ful in sowing seeds of discord and

p247
stirring up hatred.''\texthebrew{⁠}\footnote{Jordanes, \emph{Get.} 168

(after Cassiodorus).

} All that we know of his long career bears out this suggestion of astute and perfidious diplomacy.

The unhappy population of the Mauretanian regions were left unprotected to the mercies of the invaders, and if we can trust the accounts which have come down to us,\texthebrew{⁠}\footnote{Possidius, \emph{Vit. August.} 28, Victor Vit. \emph{op. cit.} I.1\texthebrew{‑}3.

} they seem to have endured horrors such as the German conquerors of this age seldom inflicted upon defenceless provinces. The Visigoths were lambs compared with the Vandal wolves. Neither age nor sex was spared and cruel tortures were applied to force the victims to reveal suspected treasures. The bishops and clergy, the churches and sacred vessels were not spared. We get a glimpse of the situation in the correspondence of St. Augustine. Bishops write to him to ask whether it is right to allow their flocks to flee from the approaching danger and for themselves to abandon their sees.\texthebrew{⁠}\footnote{Augustine, \emph{Ep.} 228. Augustine said that the bishop should let the people flee, but not abandon his post, so long as his presence was needed.

} The invasion was a signal to other enemies whether of Rome or of the Roman government to join in the fray. The Moors were encouraged in their depredations, and religious heretics and sectaries, especially the Donatists, seized the opportunity to wreak vengeance on the society which oppressed them.\footnote{See Martroye, \emph{op. cit.} 113. The devastation is described in general terms in a letter addressed by the bishop of Carthage to the Council of Ephesus in summer of 431 (Mansi, IV.1207).

}

If Africa was to be saved, it was necessary that the Roman armies should be united, and Placidia immediately took steps to regain the allegiance of Boniface. A reconciliation was effected by the good offices of a certain Darius, of illustrious rank, whom she sent to Africa,\texthebrew{⁠}\footnote{See Augustine's letter to Darius, congratulating him on his success, \emph{Ep.} 229; the reply of Darius, \emph{Ep.} 230; and Augustine's answer, \emph{Ep.} 231. Boniface seems to have given Darius a hostage, pignus pacis, \emph{Ep.} 229.1 and 231.7, who was probably the Verimodus of 230.6.

} and he seems also to have concluded a truce with Gaiseric,\texthebrew{⁠}\footnote{\emph{Ep.} 229.2 ipsa bella verbo occidere, to which Darius replies (230.3) si non extinximus bella, certe distulimus.

} which was, however, of but brief duration, for Boniface's proposals were not accepted. Gaiseric was determined to pillage, if could not conquer, the rich eastern provinces of Africa. He entered Numidia, defeated Boniface, and besieged him in Hippo (May-June A.D. 430). The city held out for more than a year.\texthebrew{⁠}\footnote{For the siege see Possidius, \emph{Vit. August.} Augustine died at the beginning of the siege (August 28, A.D. 430).

} Then Gaiseric raised the siege (July A.D. 431). New

p248
forces were sent from Italy and Constantinople under the command of Aspar, the general of Theodosius; a battle was fought, and Aspar and Boniface were so utterly defeated that they could make no further effort to resist the invader. Hippo was taken soon afterwards,\texthebrew{⁠}\footnote{Possidius, c28. Aspar seems to have remained in Africa for some time. He was the western consul in 434 and was at that time in Carthage (\emph{Lib. de permissionibus}, \emph{P. L.} 51, 841). See above,

p225 .

} and the only important towns which held out were Carthage and Cirta.

Boniface returned to Italy, where Placidia received him with favour, and soon afterwards she deposed Aetius, who was consul of the year ( A.D. 432), and gave his military command to the repentant rebel, on whom at the same time she conferred the dignity of Patrician.\texthebrew{⁠}\footnote{See \emph{Consularia Italica}, p301.

} Aetius refused to submit. There was civil war in Italy. The rivals fought a battle near Ariminum, in which Boniface was victorious, but he died shortly afterwards from a malady, perhaps caused by a wound.\texthebrew{⁠}\footnote{Cp. Hyd. 99 with Prosper, \emph{sub} 432. If John Ant., \emph{ib.}, says that Boniface was out-generalled by Aetius, this may be taken to mean that Aetius succeeded in the end. The common source of John Ant. and Procopius \emph{may} have been Priscus.

} His son-in\texthebrew{‑}law Sebastian was appointed to the vacant post of Master of Both Services,\texthebrew{⁠}\footnote{Hyd. \emph{ib.}

} but did not hold it long. Aetius escaped to Dalmatia and journeyed to the court of his friend Rugila the king of the Huns. By his help, we know not how, he was able to reappear in Italy, to dictate terms to the court of Ravenna, to secure the banishment of Sebastian, and to obtain for himself reinstatement in his old office and the rank of Patrician ( A.D. 434).\footnote{See Prosper, \emph{ib.}, \emph{Chron. Gall.}, \emph{sub} 433, from which source we learn that ``Goths were summoned by the Romans to bring aid'' against the Huns. Hydatius, 103. Sebastian found a refuge at Constantinople, where he remained for ten years (Hyd. 104, 129), and he is said to have been the commander of a pirate squadron which served Theodosius II (see Suidas, \emph{sub} \textgreek{Θεοδόσιος}, a fragment ascribed by Niebuhr to Priscus,

by Müller to John Ant., \emph{F. H. G.} IV. \emph{fr.} 194 ). Falling out of favour in A.D. 444 he went to the court of Theoderic the Visigoth, who would not receive him. Then he managed to seize Barcelona. Driven from there he went to the Vandals, and was put to death by Gaiseric (A.D. 450, Hyd. 144), and has come down to fame as a Catholic martyr.

}

In the meantime, during this obscure struggle for power, the Vandals were extending their conquests in Numidia. In spite of his wonderfully rapid career of success Gaiseric was ready to come to terms with the Empire. Aetius, who was fully occupied in Gaul, where the Visigoths and Burgundians were actively aggressive, saw that the forces at his disposal were unequal to

p249
the expulsion of the Vandals, and it was better to share Africa with the intruders than to lose it entirely. Gaiseric probably wished to consolidate his power in the provinces which he had occupied, and knew that any compact he made would not be an obstacle to further conquests. Hippo, from which the inhabitants had fled, seems to have been reoccupied by the Romans,\texthebrew{⁠}\footnote{Cp. Martroye, p128.

} and here (February 11, A.D. 435) Trygetius, the ambassador of Valentinian, concluded a treaty with Gaiseric, on the basis of the \emph{status quo.} The Vandals were to retain the provinces which they had occupied, the Mauretanias and a part of Numidia, but were to pay an annual tribute, thus acknowledging the overlord ­ship of Rome.\footnote{Prosper, \emph{sub a.} 435. Cp. Isidore, \emph{Hist. Vandalorum}, in \emph{Chron min.} II. p297 . Procopius, \emph{B. V.} I.4. The king's son Huneric was sent as a hostage to Rome, but was soon released (apparently before 439).

}

Aetius had now firmly established his power and Placidia had to resign herself to his guidance. Valentinian was fifteen years of age, and the regency could not last much longer. The presence of the Master of Soldiers was soon demanded in Gaul, where the Visigoths were again bent on new conquests and the Burgundians invaded the province of Upper Belgica ( A.D. 435). Against the Burgundians he does not appear to have sent a Roman army; he asked his friends the Huns to chastise them. The Huns knew how to strike. It is said that 20,000 Burgundians were slain, and King Gundahar was one of those who fell ( A.D. 436). Thus came to an end the first Burgundian kingdom in Gaul, with its royal residence at Worms. It was the background of the heroic legends which passed into the German epic \texthebrew{—} the Nibelungenlied. The Burgundians were not exterminated, and a few years later the Roman government assigned territory to the remnant of the nation in Sapaudia (Savoy), south of Lake Geneva ( A.D. 443).\footnote{\emph{Chron. Gall.} p660; Prosper, \emph{sub} 435;
Sidonius Apoll. \emph{Carm.} VII.234 . The number of 20,000 is of course an exaggeration.

}

Narbonne was besieged by Theoderic in A.D. 436, but was relieved by Litorius,\texthebrew{⁠}\footnote{Sidonius, \emph{ib.} 244 \emph{sqq.} Cp. Merobaudes, \emph{Pan.} I.9, l. 23; \emph{Pan.} II., l. 16.

} who was probably the Master of Soldiers in Gaul. Three years later the same commander drove the

p250
Goths back to the walls of their capital Toulouse, and it is interesting to find him gratifying his Hun soldiers by the performance of pagan rites and the consultation of auspices. These ceremonies did not help him. He was defeated and taken prisoner in a battle outside the city.\texthebrew{⁠}\footnote{Merobaudes, \emph{Pan.} II.153 \emph{sqq.}; Hydatius, 116, 117;
Sidonius, \emph{ib.} 299 \emph{sqq.}

} Avitus, the Praetorian Prefect of Gaul, who had great influence with Theoderic, then brought about the conclusion of peace. In these years there were also troubles in the provinces north of the Loire,\texthebrew{⁠}\footnote{John Ant. \emph{fr.} 201.2 (source probably Priscus); Merobaudes, \emph{Pan.} II.8:

lustrat Aremoricos iam mitior incola saltus;

Sidonius, \emph{Carm.} V.210 , mentions the defence of alarmed Tours, and VII.246 relates that Litorius having subdued the Armoricans hurried his troops against the Goths. This suggests 438\texthebrew{‑}439 as the date. There was another Armorican rebellion in 442; Aetius sent Goar, the veteran chief of the Alans now settled near Orleans, to punish the rebels, and Germanus bishop of Auxerre went to Ravenna to plead the Armorican cause. See Constantius, \emph{Vita Germani}, II. c8, and \emph{Chron. Gall.}, \emph{ib.}

} where the Armoricans rebelled, and Aetius or his lieutenant Litorius was compelled to reimpose upon them the ``liberty'' of Imperial rule.

In A.D. 437 Aetius was consul for the second time, and in that year Valentinian went to Constantinople to wed his affianced bride, Licinia Eudoxia. Now assuredly, if not before, the regency was at an end, and henceforward Aetius had to do in all high affairs not with the Empress who distrusted and disliked him but with an inexperienced youth. Valentinian was weak and worthless. He had been spoiled by his mother, and grown up to be a man of pleasure who took no serious interest in his Imperial duties. He associated, we are told, with astrologers and sorcerers, and was constantly engaged in amours with other men's wives, though his own wife was exceptionally beautiful.\texthebrew{⁠}\footnote{Procopius, \emph{B. V.} I.3.10. Perhaps the source was Priscus.

} He had some skill in riding and in archery and was a good runner, if we may believe Flavius Vegetius Renatus, who dedicated to him a treatise on the art of war.\texthebrew{⁠}\footnote{\emph{Epitome rei militaris}, III.26 . This treatise throws little light on the warfare of the writer's own time. It is mainly antiquarian, and there are few references to contemporary conditions. The fleet of lusoriae guarding the Danube is mentioned
(IV.46) . The disuse of coats of mail and helmets is noted, and frequent defeats of Imperial forces by Gothic archers are attributed to this
(I.20) . Vegetius also says that the art of naval warfare is now less important than formerly quia iamdudum pacato mari cum barbaris nationibus agitur terrestre certamen, a remark which points to the conclusion that the book was composed before 440 when the Vandal navy began to show what it could do. That the Emperor to whom the book is dedicated was Valentinian was conjectured by Gibbon and virtually proved by Seeck (in \emph{Hermes}, XI.61 \emph{sqq.}).

} From the end of the regency

p251
to his own death, Aetius was master of the Empire in the west, and it must be imputed to his policy and arms that Imperial rule did not break down in all the provinces by the middle of the fifth century.

Of his work during these critical years we have no history. We know little more than what we can infer from some bald notices in chronicles written by men who selected their facts without much discrimination. If we possessed the works of the court poet of the time we might know more, for even from the few fragments which have survived we learn facts unrecorded elsewhere. The Spaniard, Flavius Merobaudes, did for Valentinian and Aetius what Claudian had done for Honorius and Stilicho, though with vastly inferior talent. Like Claudian, he enjoyed the honour having a bronze statue erected to him at Rome, in the Forum of Trajan.\texthebrew{⁠}\footnote{CIL VI.1724 , Sidonius \emph{Carm.} IX.296. It was set up in A.D. 435, and he refers to it in his prose work on the second consul­ship of Aetius, written in 437. \emph{Pan.} I.8. From this inscription we learn that he had seen some service as a soldier: inter arma litteris militabat et in Alpibus acuebat eloquium. In another, but fragmentary, inscription (CIL VI.31983 ), his name appears as Flavius Merobaudes orator.

} His name was known and appreciated at the court of Constantinople, for Theodosius conferred upon him the rank of patrician.\footnote{\emph{Ib.} p9 pro his denique nuper ad honorem maximi nomen ille nascenti soli proximus imperator euexit (rightly explained by Vollmer).

}

He celebrated the three consul ­ships of Aetius,\texthebrew{⁠}\footnote{That he celebrated the first is a probable inference, see Vollmer, p20. We have parts of his oration on the second, and his poem on the third (\emph{Pan.} II).

} and we have part of a poem which he wrote for the second birthday of the general's younger son Gaudentius.\texthebrew{⁠}\footnote{\emph{Carm.} IV. Gaudentius was probably born about 440.

} We may be as certain as of anything that has not been explicitly recorded, that he wrote an ode for the nuptials of Valentinian and Eudoxia, and it is little less probable that he celebrated the birth of their elder child Eudocia, who was born in A.D. 438. But of all the poems he composed for the court only two have partly been preserved, both composed soon after the birth of the Emperor's younger daughter Placidia.\texthebrew{⁠}\footnote{\emph{Carm.} I. II. As Placidia was already married to Olybrius when she was taken to Carthage in 455 (see below,

p325 ) she can hardly have been born later than in 440.

} One of these is a description of mosaic pictures in a room in the Palace of Ravenna, representing scenes from the Emperor's life. He and Eudoxia shone in the centre of the ceiling like bright stars, and all around were scenes in

p252
which he appeared with his mother, his sister, his children, and his cousin Theodosius.\footnote{See

Bury, \emph{Justa Grata Honoria}

(see Bibl.). º

}

Like another more famous man of letters, his younger contemporary Sidonius, Merobaudes was called upon to fill a high office and to assist Aetius in the work of maintaining order in the provinces. We are told that he was appointed Master of Both Services and went to his native province of Baetica to suppress a rebellion of turbulent peasants ( bacaudae ), that he successfully accomplished this task but was recalled to Rome through the machination of his enemies ( A.D. 443). His immediate predecessor in the command had been his father-in\texthebrew{‑}law , Asturius.\footnote{Our informant is the contemporary Spanish writer Hydatius (128), and his statement as to the office held by Asturius is confirmed by the consular diptych of that personage (A.D. 440), on which he is described as ex mag. utriusq. mil. See Meyer, \emph{Zwei ant. Elfenb.} p56.

}

It must not be thought that Asturius and Merobaudes, in bearing the title ``Master of Both Services,'' had succeeded to the post of Aetius and were supreme commanders of the army. Aetius had not resigned the supreme command; he was still Master of Both Services. The command which Asturius and Merobaudes held, and which Sigisvult had held two years before,\texthebrew{⁠}\footnote{Valentinian III, \emph{Nov.} 6.1 and \emph{Nov.} 9 (March and June A.D. 440).

} was simply that of the magister equitum praesentalis under a new name. Under Stilicho, Constantius, and Felix the magister equitum had been subordinate to the magister utriusque militiae , and this arrangement undoubtedly continued still, but some time before A.D. 440 he received the same title as his superior, doubtless because it was found convenient to place legions as well as cavalry under his command. The superior Master of Both Services, the Emperor's principal statesman and director of affairs, is from this time forward generally designated as ``the Patrician'' \texthebrew{—} the Emperor's Patrician, the Patrician in a superlative sense.\footnote{\emph{Id.} Nov. 9, patricium nostrum Aetium. Cp. John Ant. \emph{fr.} 84 (\emph{De ins.} p126) \textgreek{τ\texthebrew{ῆ}ς πατρικιότητος}. So in \emph{Cons. Ital.} p305 Messianus, patricius Aviti. That Aetius continued to hold the Mastership is shown by Val. III

\emph{Nov.} 17

(A.D. 445). In A.D. 446 the subordinate master was Vitus, who was sent to Spain against the Suevians who were ravaging the southern and eastern provinces.\texthebrew{—} A different view is taken by Sundwall (in \emph{Weströmische Studien}), who thinks that Asturius, Merobaudes, and Vitus were Masters of Soldiers in Gaul. They did not operate in Gaul but in Spain, and were surely sent direct from Italy with Italian troops, so that it seems perverse not to regard them as the successors

(p253)
of Sigisvult. The magistri equitum in Gaul had indeed a mixed command, but the first of them who certainly bore the title mag. ped. et equit. or mag. utr. mil. was Avitus in 455
(Sidonius, \emph{Carm.} VII.377) . Cassius is described as mag. militum Galliarum in \emph{Vita Hilarii}, 6, 9, \emph{P. L.} 50, but this may mean no more than mag. equit.\texthebrew{—} At a later date, we find beside Ricimer a second mag. utr. mil. in Italy, namely Flavius Theodobius Valila, in A.D. 471 (see CIL VI.32169 , 32221).

}

p253
The position of Aetius in these years as the supreme minister was confirmed by the betrothal of his son to the Emperor's daughter Placidia,\texthebrew{⁠}\footnote{Prosper, \emph{sub} 454. The son of Aetius was doubtless Gaudentius, and the princess must have been Placidia, as Eudoxia was betrothed to Huneric (see below,

p256 ).

} an arrangement which can hardly have been welcome to Galla Placidia, the Augusta. With Valentinian himself he can hardly have been on intimate terms. The fact that he had supported the tyrant John was probably never forgiven. And it cannot have been agreeable to the young Emperor that it was found necessary to curtail his income and rob his privy purse in order to help the State in its financial straits.\texthebrew{⁠}\footnote{\emph{C. Th.} XI.1.36 , A.D. 431. In later years the necessity was more imperative. For the condition of Africa see \emph{Nov. Valent.} XII, XIII, and I.1.

} Little revenue could come from Africa, suffering from the ravages of the Vandals, and in A.D. 439, as we shall see, the richest provinces of that country passed into the hands of the barbarians.

The income derived from Gaul must have been very considerably reduced, and we are not surprised to find the government openly acknowledging in A.D. 444 that ``the strength of our treasury is unable to meet the necessary expenses.'' In that year two new taxes were imposed, one on the senatorial class, and one on sales, expressly for the purpose of maintaining the army. New recruits were urgently wanted, and there was not enough money in the treasury to feed and clothe the existing regiments. Senators of illustrious rank were required to furnish the money for maintaining three soldiers, senators of the second class one, senators of the third one-third; that meant 90, 30, and 10 solidi respectively, as the annual cost of a soldier was estimated at 30.\texthebrew{⁠}\footnote{\emph{Nov. Valent.} VI.3.

} A duty of ¹\texthebrew{⁄₂₄}th was imposed on sales \texthebrew{—} a siliqua in a solidus \texthebrew{—} of which the seller and the buyer each paid half.\texthebrew{⁠}\footnote{\emph{Nov. Valent.} XV.

} The government would have done better if it had forced the rich senators of Italy to contribute substantial sums, as they could well have afforded to do, to the needs of the State.\footnote{Sundwall, \emph{Weström. Studien}, 158. He calculates that the state revenue from the land tax c. A.D. 450 was at most £4,800,000, as compared with 13 millions fifty years before. Whatever

(p254)
may be thought about his figures, the proportion of the decline is hardly overstated. In this work Sundwall also illustrates the growing distinction between the highest senatorial class (illustres), and the two lower classes, and argues that while the members of the Roman senate in 400 were about 3000, in 450 they were about 2000. He pertinently points out that out of the not very large amounts which the senators paid in taxes, many of them got much back in the salaries of the high posts (Prefectures, etc.) to which they were appointed.

}

The treaty of A.D. 435 was soon violated by Gaiseric. He did not intend to stop short of the complete conquest of Roman Africa. In less than five years Carthage was taken (October 19, A.D. 439).\texthebrew{⁠}\footnote{Prosper, \emph{sub a.}

} If there was any news that could shock or terrify men who remembered that twenty years before Rome herself had been in the hands of the Goths, it was the news that an enemy was in possession of the city which in long past ages had been her most formidable rival. Italy trembled, for with a foe master of Carthage she felt that her own shores and cities were not safe. And, in fact, not many months passed before it was known that Gaiseric had a large fleet prepared to sail, but its destination was unknown.\texthebrew{⁠}\footnote{Satis incertum est ad quam oram terrae possint naves hostium pervenire. Valentinian, \emph{Nov.} 9 (June 24, A.D. 440).

} Rome and Naples were put into a state of defence;\texthebrew{⁠}\footnote{Naples: for I think that we should refer to this year the following inscription found there (CIL X.1485 ): d.n. Placidus Valentinianus providentissimus omnium retro principum, salvo adque concordi d.n. Fl. Theodosio invictissimo Aug. ad decus nominis sui Neapolitanam civitatem, ad omnes terra marique incursus expositam et nulla securitate gaudentem, ingenti labore adque sumptu muris turribusq. munivit.

} Sigisvult, Master of Soldiers, took steps to guard the coasts; Aetius and his army were summoned from Gaul; and the Emperor Theodosius prepared to send help.\texthebrew{⁠}\footnote{\emph{Ib.} All these preparations are announced in this constitution, addressed to the Roman people, and intended to calm their fears. The Emperor had come to Rome before March 3, where steps were being taken to repair the walls (\emph{Nov.} 3).

} There was indeed some reason for alarm at Constantinople. The Vandal pirates could afflict the eastern as well as the western coasts of the Mediterranean; the security of commerce was threatened. It was even thought advisable to fortify the shore and harbours of Constantinople.

Gaiseric, aware that Italy was prepared, directed his attack upon Sicily, where he laid siege to Panormus.\texthebrew{⁠}\footnote{He is said to have been invited by Maximin, bishop of the Arian communities in Sicily (Cassiodorus, \emph{Chron., sub} 440), and he persecuted the Catholics while he was in the island (Hydatius, 120).

} This city defied

p255
him, but it is possible, though not certain, that he occupied Lilybaeum.\texthebrew{⁠}\footnote{Cp. Pope Leo I, \emph{Ep.} 3 (\emph{P. L.} 54.606); Martroye, p132.

} His fleet, however, returned to Africa, perhaps on account of the considerable preparations which were on foot at Constantinople.\texthebrew{⁠}\footnote{Prosper, \emph{sub} 440, ascribes his return to danger from the threat of an attack on Carthage by Sebastian (the son-in\texthebrew{‑}law of Boniface), invading Africa ab Hispania. Sebastian seems to have been in the service of Theodosius. Cp. above,

p248 .

} The government of Theodosius had made ready a large naval squadron which sailed in the following year ( A.D. 441), with the purpose of delivering Carthage from the Vandals.\texthebrew{⁠}\footnote{The Imperial fleet was under three generals, Areobindus, Ansila, and Germanus (Prosper, \emph{sub} 441). Theophanes is evidently referring to the same expedition \emph{sub} A.M. 5942 = A.D. 448\texthebrew{‑}449. He says that the fleet of transports numbered 1100 (which has a suspicious resemblance to the number of Leo's armada in A.D. 468), and adds the names of two other generals.

} The expedition arrived in Sicily, and Gaiseric was alarmed. He opened negotiations, pending which the Imperial fleet remained in Sicilian waters. These diplomatic conversations were protracted by the craft of Gaiseric, and in the meantime an invasion of the Huns compelled Theodosius to recall his forces. The Emperors were thus constrained to make a disadvantageous peace.

By the treaty of A.D. 442 Africa was divided anew between the two powers. This division nearly reversed that of A.D. 435, and was far more advantageous to the Vandals. The Empire retained the provinces of Tripolitania, Mauretania Sitifensis, Mauretania Caesariensis, and part of Numidia; while the Vandals were acknowledged masters of the rest of that province, of Byzacena, and of the Proconsular province or Zeugitana.\texthebrew{⁠}\footnote{The sources for this division are Valentinian III, \emph{Nov.} 33 and \emph{Nov.} 18; Victor Vit. I.4.

} Mauretania Tingitana was probably not mentioned in the treaty.\texthebrew{⁠}\footnote{The words of Victor, Byzacenam Abaritanam atque Getuliam, are obscure. Getulia seems to be the southern districts of Byzacena. The exact meaning of Abaritana (cp.

Plin. \emph{N. H.} XVI.36, 172 ) is unknown. It seems to be a district of the Proconsular province, as we find among bishops of that province in the reign of Huneric Felix abaritanus (\emph{Notitia prov. et civ. Afr.} p63). Schmidt (\emph{Gesch. der W.} 72) thinks Tingitana is meant, but this has not the least plausibility. Cp. Martroye, 135\texthebrew{‑}136.

} It was part of the diocese of Spain, not of the diocese of Africa, and it is probable that the Vandals never occupied it effectively. In any case it now belonged to the Empire, which, since the departure of the Vandals, had been in possession of all Spain, except the Suevian kingdom in the north-western corner.

This settlement was an even greater blow to the Empire

p256
than that which necessity had imposed upon Constantius of settling the Visigoths in Aquitaine. The fairest provinces of Africa were resigned to barbarians who had an even worse reputation than the Goths. But it was worth while to attempt to secure that the settlement, such as it was, should be permanent. Aetius saw that the best policy was to cultivate good relations with Gaiseric and to give that ambitious and unscrupulous monarch no pretext for attacking Sicily, or Sardinia, or Italy itself. And so he prevailed upon Valentinian to consent to a betrothal between his elder daughter, Eudocia, and Gaiseric's son, Huneric. It is probable that this arrangement was considered at the time of the treaty, though it may not have been definitely decided.\texthebrew{⁠}\footnote{It was prior to 446, the year of the third consul­ship of Aetius, for Merobaudes refers to it in his poem on that occasion, \emph{Pan.} II.27 \emph{sqq.}:

nunc hostem exutus pactis proprioribus

Romanam uincire fidem Latiosque parentes

adnumerare sibi sociamque intexere prolem.

In an earlier poem of Merobaudes (\emph{Carm.} I.17) the future marriage of the child princess is touched on.

} But Huneric was already married. The Visigothic king Theoderic had bestowed upon him his daughter's hand. Such an alliance between Vandals and Goths could not have been welcome to Aetius; it was far more in the interest of his policy to keep alive the hostility between these two peoples which seems to have dated from the campaigns of Wallia in Spain. The existence of the Gothic wife was no hindrance to Gaiseric, and a pretext for repudiating her was easily found. She was accused of having plotted to poison him.\texthebrew{⁠}\footnote{Jordanes, \emph{Get.} 184 .

} She was punished by the mutilation of her ears and nose, and in this plight she was sent back to her father. The incident meant undying enmity between Visigoth and Vandal. Theoderic soon sought a new ally by marrying another daughter to Rechiar, king of the Suevians ( A.D. 449).\texthebrew{⁠}\footnote{Hydatius, who mentions that in the same year he helped his son-in\texthebrew{‑}law to capture Ilerda (140, 142).

} Huneric was free to contract a more dazzling matrimonial alliance with an Imperial princess.

We are not informed whether in the treaty of A.D. 442 any provision was made for supplying Italy with the corn º of Africa on which the Romans had subsisted for centuries. In the absence of evidence to the contrary, we may safely assume that, throughout the duration of the Vandal kingdom, the surplus of the corn production of Africa was consumed as of old in

p257
Italy (except, perhaps, in the few years in which there were open hostilities); only now instead of being a tribute it was an export.\texthebrew{⁠}\footnote{A certain amount could be re­quisitioned in the old way from the Mauretanias so long as they remained in Roman hands. We may wonder how the African shipping corporations, whose offices are to be seen in the great square north of the theatre at Ostia (cp. Ashby, ``Recent Excavations at Ostia,'' \emph{J. R. S.} II. (1912) 180 ), were affected by the changed circumstances.

} It was obviously to the interest of the Vandal proprietors to send the grain they did not want to Italian markets.

The Vandals themselves settled in Zeugitana, and made Carthage their capital. They appropriated the lands of the proprietors in this province, who, unless they migrated elsewhere, were probably degraded to the position of serfs. The Vandals, as Arians, had from the very beginning assumed a definitely hostile attitude to the Catholic creed. When Carthage was taken the Catholic clergy were banished, and all the churches of the city were given up to Arian worship. The independent attitude of the Vandals towards the Empire is reflected in their adopting a chronological era of their own, beginning on October 19, A.D. 439, the date of the capture of Carthage.

It is to be observed that the Vandals now held a position of vantage in regard to the Empire that none of the other Teutonic nations ever occupied. In relation to the foreign peoples of northern Europe, the front of the Roman Empire was the Rhine and the Danube. And so we may say that the Vandals had come round to the back of the Empire and were able to attack it from behind. Another exceptional feature in their position was that, in the language of a chronicler, the sea was made pervious to them: they created a fleet of small light cruisers and attacked the Empire by sea, as no other Teutonic people had done or was to do in the Mediterranean, though the Saxons and other men of the north used ships to harry it in the northern and western oceans. Thus they were able to follow in the track of the Carthaginians of old, and extend their dominion over the western islands.

Till after the death of Valentinian ( A.D. 455) the naval expeditions of the Vandals seem to have been simply piratical,\footnote{These depredations, which extended to the Aegean, are mentioned, A.D. 457, by Nestorius, \emph{Book of Heraclides}, p331: ``Sicily, Rhodes, and many other great islands with Rome itself.'' Rome, however, was

(p258)
not attacked till 455. Other sources mention raids in Greece and southern Italy after 457 (Procopius, \emph{B. V.} I.5; Victory Vit. \emph{Hist. pers.} I.51).\texthebrew{—} There is no definite record that he troubled Sicily between 442 and 455. Pace (\emph{I Barb. e Biz.} p12) thinks he did, and that the services which Cassiodorus performed in defending the coasts of Bruttii and Sicily (Cassiod. \emph{Var.} 1.4) are to be referred to that period. But they may be connected with the events of 440.

}

p258
though Gaiseric may have definitely formed the design of conquering Sicily. But soon after that year he seems to have occupied without resistance the two Mauretanian provinces which the Empire had retained under the treaty of A.D. 442, and to have annexed Sardinia, Corsica, and the Balearic Islands.\texthebrew{⁠}\footnote{Victor Vit. \emph{op. cit.} I.13.

} Sicily itself was to pass somewhat later under his dominion.

The military and diplomatic successes of Gaiseric encouraged and enabled him to encroach on the liberties of his people. Among all the ancient Germanic peoples, the sovran power resided in the assembly of the folk, and in the case of those which formed permanent states on Imperial soil, like the Franks and the Visigoths, it was only by degrees that the kings acquired great but not absolute power. In the Vandal state alone the free constitution was succeeded by an autocracy, without any intermediate stages. The usurpation by the king of unconstitutional powers occasioned a conspiracy of the nobles, and it was bloodily suppressed.\texthebrew{⁠}\footnote{Prosper, \emph{sub} 442. Cp. Schmidt, \emph{op. cit.} 161.

} The old aristocracy seems to have been superseded by a new nobility who owed their position, not to birth, but to appointments in the royal service. It is probable that the assembly of the folk ceased to meet. Before his death Gaiseric issued a law regulating the succession to the throne,\texthebrew{⁠}\footnote{Procopius, \emph{B. V.} I.7.29. Here it is not called a testament, but a law, \emph{ib.} 9.12 (constitutio, Victor Vit. \emph{op. cit.} II.13). On this act cp. Schmidt, \emph{ib.} 165. Gaiseric had already done away with the sons and wife of his brother Gunderic (Victor, \emph{ib.} 14). Neither Gaiseric nor his successor Huneric struck coins with their own names or busts (it is just possible that Huneric issued a bronze coin with his bust, but the attribution is uncertain). Gunthamund (484\texthebrew{‑}496) struck silver, and all his successors silver and bronze, with names and busts. The Vandals seem to have made little use of a gold currency, and their gold coins are all of Imperial type. The large bronze coins, probably attributed to Gaiseric (issued perhaps, as Wroth suggests, about 435 when he captured Carthage), marked with XLII and XXI Nummi, are remarkable as an anticipation of the folles of Anastasius (see below

p446 ). º See Wroth, \emph{Coins of the Vandals, etc.} xvi \emph{sqq.}

} thus depriving the people of the right of election, and the royal authority was so firmly established that his will was apparently accepted without demur. By this law the kingship was treated as a personal inheritance and was confined to Gaiseric's male descendants, of whom the eldest was always to succeed.

p259
The policy of Gaiseric differed entirely from that of the Goths in Gaul. He aimed at establishing a kingdom which should be free, so far as possible, from Roman influence, and he saw that, for this purpose, it was necessary above all to guard jealously the Arian faith of his people, and not expose them to the danger of being led away by the propaganda of the Catholics. He was therefore aggressively Arian, and persecuted the Catholic clergy.\texthebrew{⁠}\footnote{The chief sources are Victor Vit. \emph{Hist. pers.} bk. I; Prosper, \emph{sub} 437, 439; Theodoret, \emph{Epp.} 52, 53 (\emph{P.G.} 83). The details are recounted in Martroye, \emph{Genséric}, 328 \emph{sqq.}

} He imposed the Arian creed on all persons who were in his own immediate environment. After the capture of Carthage he seized the Donatist bishop Quodvultdeus and other clergy, set them on board old and untrustworthy ships, and committed them to the mercy of the sea. They reached Italy safely. Throughout the proconsular province the bishops were expelled from their sees and stripped of their property. It was not till A.D. 454 that a new bishop was allowed to be ordained at Carthage, and some churches were reopened for Catholic worship. But after the death of Deogratias, at the end of three years, the old rigorous suppression was renewed; the sees were left vacant throughout the province, and the priests were forced to surrender their books and sacred vessels. The monasteries, however, were not suppressed. And the persecution was not general or ubiquitous. Particular persons were singled out and dealt with by the express order of the king. He did not give a free hand to his officers, and there were probably few cases of death or personal violence.

It was no less important for the ends of Gaiseric's policy to eliminate the power of the senatorial aristocracy. He did this by such drastic measures that a contemporary chronicler observed, ``It is impossible to say whether his hostility to men or to God was the more bitter.'' He deprived of their domains the nobles of the proconsular province, and told them to betake themselves elsewhere. They were not to be suffered to remain lords of the soil to organise an opposition to the king, and gradually to recover political influence under his successors. If they remained in the land they were threatened with perpetual slavery. After the capture of Carthage most of the senators had been compelled to leave the shores of Africa, some sailing to Italy,

p260
others to the East.\texthebrew{⁠}\footnote{See \emph{Vita Fulgentii}, c1; Theodoret, \emph{Epp.} 29\texthebrew{‑}36; Prosper, \emph{sub} 539.

} In the other parts of his realm Gaiseric does not appear to have adopted such extreme measures. He deemed it sufficient to make the royal capital and the central province safe.

The Empress Galla Placidia, who had been supreme ruler in the west for about ten years, and for fifteen more had probably exercised some influence on the direction of affairs, died at Rome in A.D. 450.\texthebrew{⁠}\footnote{Nov. 27, Prosper, \emph{sub a.}

} But her memory will always be associated with Ravenna, where the Imperial court generally resided\texthebrew{⁠}\footnote{All the laws during her regency \texthebrew{—} they are not numerous \texthebrew{—} were issued from Ravenna. Valentinian lived both at Rome and at Ravenna; during the last years of his reign, after his mother's death, almost entirely at Rome.

} and where she was buried in the mausoleum which she had built to receive her ashes.

Honorius had done one memorable thing which altered the course of history. He made the fortune of Ravenna. To escape the dangers of the German invasions he had moved his government and court from Milan to the retired city of the marshes, which amid its lagoons and islands could defy an enemy more confidently than any other city in the peninsula, and, as events proved, could hardly be captured except by a maritime blockade. Before Augustus it had been an obscure provincial town, noted chiefly for its want of fresh water, but had served as a useful refuge to Caesar before he crossed the Rubicon. Augustus had chosen it to be a naval station, and had supplied it with a good harbour, Classis, \texthebrew{•} three miles from the town, with which he connected it by a solid causeway across the lagoons. But nothing seemed more unlikely than that it should overshadow Milan and vie with Rome as the leading city of Italy. Through the act of Honorius, which though conceived in fear turned out to be an act of good policy, Ravenna became the home of emperors, kings, and viceroys, and throughout the vicissitudes of four centuries of crowded history was a name almost as familiar as Rome itself in the European world.

Ravenna has no natural amenities. Here are the impressions the place produced on a visitor from Gaul not many years after

p261
Placidia's death.\texthebrew{⁠}\footnote{Sidonius Apollinaris, \emph{Epp.} I.5 (A.D. 467) and 8 (A.D. 468), Dalton's translation.

} ``The Po divides the city, part flowing through, part round the place. It is diverted from its main bed by the State dykes, and is thence led in diminished volume through derivative channels, the two halves so disposed that one encompasses and moats the walls, the other penetrates and brings them trade \texthebrew{—} an admirable arrangement for commerce in general, and that of provisions in particular. But the drawback is that, with water all about us, we could not quench our thirst; there was neither pure-flowing aqueduct, nor filterable cistern, nor trickling source, nor unclouded well. On the one side the salt tides assail the gates; on the other, the movement of vessels stirs the filthy sediment in the canals, or the sluggish flow is fouled by the bargemen's poles, piercing the bottom slime.'' ``In that marsh the laws of everything are always the wrong way about; the waters stand and the walls fall, the towers float and the ships stick fast, the sick man walks and the doctor lies abed, the baths are chill and the houses blaze, the dead swim and the quick are dry, the powers are asleep and the thieves wide awake, the clergy live by usury and the Syrian chants the psalms, business-men turn soldiers and soldiers business-men , old fellows play ball and young fellows hazard, eunuchs take to arms and rough allies to letters.''

In this description the writer remarks the presence of the Syrian, a familiar figure to him in the cities of southern Gaul. But it was not only oriental traders whom the new Imperial residence attracted. It is probable that artistic craftsmen from Syria and Anatolia came to embellish the city of Honorius and Placidia, and to teach their craft to native artists. For it is difficult otherwise to explain the oriental inspiration which so conspicuously distinguishes the Ravennate school of art that it has been described as `` half-Syrian .''\footnote{Dalton, \emph{Byz. Art}, p8.

}

It was indeed in the artistic works with which its successive rulers enriched it that the great attraction of Ravenna lay and still lies. Many of these monuments have perished, but many have been preserved, and they show vividly the development of Christian art in Italy in the fifth and sixth centuries, under the auspices of Placidia, Theoderic, and Justinian, under the influence of the East. Brick was generally the material of

p262
these buildings, but their unimpressive exterior appearance was compensated by the rich decoration inside and the brilliant mosaics which shone on the walls. Ravenna is the city of mosaics. At Rome we have from the fourth and early fifth centuries fine examples of this form of pictorial art in the churches of S. Costanza and S. Pudenziana and S. Maria Maggiore,\texthebrew{⁠}\footnote{The Basilica Liberiana, built by Pope Liberius on the Esquiline in the fourth century, was dedicated by Pope Sixtus III to St. Mary \emph{c.} A.D. 432 and perhaps partly rebuilt. The dating of the mosaics has been much debated. Richter and Taylor contend that the mosaics of the nave are pre-Constantinian, in their \emph{Golden Age of Classic Christian Art.} All previous studies of the Church have been superseded by Wilpert's magnificent work \emph{Die röm. Mosaiken . . . vom IV. bis XIII. Jahrhundert} (see

Bibliography ), where the pictures can be studied in coloured reproductions. His conclusion is (vol. I.412 \emph{sqq.}) that the mosaics of the nave belong to the time of Liberius, those of the triumphal arch to that of Sixtus.

} but at Ravenna, in the days of Placidia, the art of painting with coloured cubes seems to enter upon a new phase and achieve more brilliant effects.\footnote{The blue and gold backgrounds strongly contrast with the pale effects at Rome.

}

No trace remains of the Imperial palace of the Laurelwood, but the churches of St. John the Evangelist and St. Agatha, the Oratory of St. Peter Chrysologus,\texthebrew{⁠}\footnote{He became archbishop of Ravenna in 433 and was succeeded by Neon either in 449 or 458. His monogram in mosaic survives in one of the arches in his chapel. The source for the early ecclesiastical history of Ravenna is

Agnellus, \emph{Lib. Pont.}

(ninth century).

} the Baptistery, and the little chapel dedicated to SS. Nazarius and Celsus which was built to receive the sarcophagi of the Imperial family, are all monuments of the epoch of Placidia.\texthebrew{⁠}\footnote{For the architecture of the churches of this period Rivoira (\emph{Lombardic Arch.} I.21\texthebrew{‑}39) supersedes previous studies. Structurally the Ravennate architects represent the Roman traditions. It is in the decoration that the oriental influence reveals itself. For the mosaics and sculptures see Diehl, \emph{Ravenne}; Dalton, \emph{op. cit.}

} The basilica of St. John was the accomplishment of a vow which the Empress had made to the saint when she and her two children were in peril of shipwreck on the Hadriatic.\texthebrew{⁠}\footnote{The dedicatory inscription is preserved (Galla Placidia cum filio suo Placido Valentiniano Augusto et filia sua Iusta Grata Honoria liberationis periculum maris votum solverunt; Agnellus, \emph{ib.} p68.

CIL XI.276 ). The incident may have occurred on the voyage from Italy to an Illyrian port in 423. But I conjecture that the same storm which dispersed the ships of Ardaburius, drove the Empress and her children back to the Dalmatian coast, and they then proceeded by land to Aquileia (see above,

p223 ). As it is not likely that Placidia delayed the fulfilment of her vow, we may place the building and inscription in 426\texthebrew{‑}427. Another inscription is recorded (De Rossi, II.1.435), in which Honoria is associated with her mother and brother (and which must therefore be prior to 437) dedicating the church of Santa Crux in Hierusalem at Rome, and probably in fulfilment of the same vow (Sanctae ecclesiae Hierusalem

(p263)
Valentinianus Placidia et Honoria Augusti votum solverunt).

} The story of their experiences was

p263
depicted on the pavement and the walls, but all the original decorations of the church have perished.\texthebrew{⁠}\footnote{The Ravennate school of builders were fond of the motive of arcading. The walls of both St. John and St. Agatha are externally decorated with blank arcades resting on a plinth (Rivoira, \emph{ib.} 21\texthebrew{‑}22). Again we find small arcades, springing from corbels between pilasters, on the Baptistery, the chapel of Chrysologus, and the church of St. Francesco (begun in A.D. 450), used to form a sort of fringe below the cornice (\emph{ib.} 36\texthebrew{‑}37).

} The Baptistery may have been begun in the lifetime of Placidia, but appears not to have been completed till after her death by the archbishop Neon. It is an octagonal building, with two tiers of round arches springing from columns, inside, crowned by a hemispherical dome, of which it has been observed that ``the ancient world affords no instance of so wide a vault constructed of tapering tubes.''\texthebrew{⁠}\footnote{\emph{Ib.} 39.

} The mosaics of the Baptistery and of Placidia's mausoleum have been wonderfully well preserved. The mausoleum, constructed about A.D. 440, is in the form of a small Latin cross, of which the centre is surmounted by a square tower closed by a conical dome.\texthebrew{⁠}\footnote{Rivoira writes (\emph{ib.} 28): ``So far as I am aware there is no record of churches or tombs older than this mausoleum having the form of a Latin cross, with rectangular extended arms and not mere apses opposite to one another, and starting directly from the central space.'' The portico in front of the mausoleum connected it with the basilica of the Holy Cross, which was built about 449.

} Here the artist in mosaics has achieved a signal triumph in the harmonious effects of his colours. The cupola is a heaven of exquisite blue, dotted with golden stars and arabesques, and in the midst a great cross of gold. Above the door and fa­cing it are two pictures, one perhaps of St. Laurence, the other of the Good Shepherd, but not the simple Shepherd of the Catacombs, bearing a sheep on his shoulder.\texthebrew{⁠}\footnote{Cp. Diehl, \emph{op. cit.} p50.

} Here he is seated on a rock in a meadow where six sheep are feeding, his tunic is golden, his cloak purple, his head, which suggests that of a Greek god, is surrounded by a golden halo.

Into this charming chapel Placidia removed the remains of her brother Honorius and her husband Constantius, and it was her own resting-place . The marble sarcophagus of Honorius is on the right, that of Constantius, in which the body of Valentinian III was afterwards laid, on the left. Her own sarcophagus of alabaster stands behind the altar, and her embalmed body in Imperial robes seated on a chair of cypress wood could be seen through a hole in the back till A.D. 1577, when all the contents

p264
of the tomb were accidentally burned through the carelessness of children.\footnote{See Hodgkin, \emph{Italy and her Invaders}, I.888 .\texthebrew{—} The surviving mosaics of the Placidian period, in the tomb and the baptistery, are only a small portion of the artistic work which then adorned the churches of Ravenna. Besides the mosaics of St. John, those of the cathedral (built in the early years of Honorius, before 410), St. Agatha, St. Laurence, the Holy Cross have disappeared. Cp. the list in Dalton, \emph{op. cit.} 365. The mosaics of the palace may have been carried off to Aachen to adorn the palace of Charles the Great.

}

The coins of the Empress show a conventional face, like those of her daughter and of the other Imperial ladies of the age. They do not portray her actual features, nor can we form any very distinct impression of her appearance from a gold medallion of which two specimens are preserved.\footnote{Cp. Delbrück, \emph{Porträts byz. Kais.} 375. The legend salus reipublicae suggests a date in the last years of Honorius.\texthebrew{—} On a stamped silver ingot, found north of Minden and now in the Hanover Museum, there is an impression of three Imperial heads, which have been supposed to be Valentinian III, Theodosius II, and Placidia. Babelon, \emph{Traité des monnaies gr. et rom.} I. p887.

\texthebrew{◂} previous

next \texthebrew{▸}Images with borders lead to more information.

The thicker the border, the more information.

(Details here.)

UP TO:

J. B. Bury
Homepage

Latin \& Greek

Texts

LacusCurtius

Home}

\xchapter{Chapter IX}

Short URL for this page:

bit.ly/BURLAT9

The misfortunes of the Balkan Peninsula have been almost uninterrupted from the fourth century to the present day. In the fifth and sixth centuries their plight was almost unendurable. They suffered not only from the terrible raids of nomad savages who had come from beyond the Volga, but also from the rapacious cruelty of the Germans. From the reign of Valens to that of Heraclius the unhappy inhabitants might any morning wake up to find a body of barbarians at their gates. As we shall be concerned in these volumes with the successive invasions of Huns, Ostrogoths, Slavs, and Bulgars, it will be well for the reader to have a general idea of the conformation and geography of the peninsula.\footnote{The following works have been useful: Jireček, \emph{Die Gesch. der Bulgaren} and \emph{Die Heerstrasse v. Belgrad nach Constantinopel}; Evans, \emph{Antiquarian Researches in Illyricum}, with good sketch maps; W. Tomaschek, \emph{Haemus-halbinsel} (in \emph{S.B.} of Vienna Acad. 1881); F. Kanitz, \emph{Römische Studien in Serbien} (\emph{Denksch.} of Vienna Acad., ph.-hist. Kl. XLI., 1892); Kiepert, \emph{Formae orbis antiqui}, Map xvii \emph{Illyricum et Thracia}; the maps in CIL vol. III. There is a good military map of Serbia, Montenegro, and Albania, attached to an article of O. Kreutzbruch v. Lilienfels, in Petermann's \emph{Mitteilungen}, Nov. 1912.

}

We may consider Mount Vitoš, and the town of Sardica, now Sofia, which lies at its base as the central point. Rising in the shape of an immense cone to a height of \texthebrew{•} 7500 feet, Vitoš affords to the climber who ascends it a splendid view of the various intricate mountain chains which diversify the surrounding lands \texthebrew{—} a view which has been pronounced finer than that at Tempe or that at Vodena. In the group of which this mountain and another named Ryl, to southward, are the highest peaks, two

p266
rivers of the lower Danube system, the Oescus (Isker) and the Nišava have their sources, as well as the two chief rivers of the Aegean system, the Hebrus (Maritsa) and the Strymon (Struma).

From this central region stretches in a south-easterly direction the double chain of Rhodope, cleft in twain by the valley of the Nestos (Mesta). The easterly range, Rhodope proper, forms the western boundary of the great plain of Thrace, while the range of Orbelos separates the Nestos valley from the Strymon valley.

The Haemus or Balkan chain which runs from west to east is also double, like Rhodope, but is not divided by a large river. The Haemus mountains begin near the sources of the Timacus (Timok) and the Margus (Morava), from which they stretch to the shores of the Euxine. To a traveller approaching them from the northern or Danubian side they do not present an impressive appearance, for the ascent is very gradual; plateau rises above plateau, or the transition is accomplished by gentle slopes, and the height of the highest parts is lost through the number of intervening degrees. But on the southern side the descent is precipitous, and the aspect is imposing and sublime. This contrast between the two sides of the Haemus range is closely connected with the existence of the second and lower parallel range, called the Srêdna Gora, which runs through Roumelia from Sofia to Sliven. It seems as if a convulsion of the earth had cloven asunder an original and large chain by a sudden rent, which gave its abrupt and sheer character to the southern side of the Haemus mountains, and interrupted the gradual upward incline from the low plain of Thrace.

The chain of Srêdna Gora, which is not to be confused with the northern chain of Haemus, is divided into three parts, which may be distinguished as the Karadža Dagh, the Srêdna Gora, and the Ichtimaner. The Karadža Dagh mountains are the most easterly, and are separated from Srêdna Gora by the river Strêma (a tributary of the Maritsa), while the valley of the Tundža (Taenarus), with its fields of roses and pleasantly situated towns, divides it from Mount Haemus. Srêdna Gora reaches a greater height than the mountains to east or to west, and is divided by the river Topolnitsa from the most westerly portion, the Ichtimaner mountains, which connect the Balkan system

p267
with the Rhodope system, whilst at the same time they are the watershed between the tributaries of the Hebrus and those of the Danube.

There are eight chief passes across the Haemus range from Lower Moesia to southern Thrace. If we begin from the eastern extremity, there is the coast pass which a traveller would take who, starting from Odessus (Varna), wished to reach Anchialus. The next pass was one of the most important. It crossed the Kamcija at Pannysus, and through it ran the road from Trajan's Marcianopolis (near Provad, between Šumla and Varna) southward. Farther west were the two adjacent passes of Veregava and Verbits (together known as the Gylorski pass).\texthebrew{⁠}\footnote{Veregava is now called the Rish pass, and is to be identified with the Iron Gate of Greek historians. These routes became important in the eighth century when the Bulgarians had built their royal capital at Aboba, near Šumla; for they connected it directly with the towns of Marcellae (Karnobad) and Diampolis (Jambol on the Tundža).

} Passing over the Kotel and Vratniti passes, which seem to have been little used for military purposes in the period which concerns us, we come to the celebrated pass of Šipka which connects the valley of the Jatrus (Jantra) with that of the Tundža. Through it ran the direct road from Novae (Šistova) on the Danube to Beroe (Stara Zagora), Philippopolis, and Hadrianople.

From this pass eastward extend the wildest regions of the Balkans, which have always been the favourite home of outlaws \texthebrew{—} scamars, as they were called, or klephts \texthebrew{—} who could defy law in thick forests and inaccessible ravines, regions echoing with the songs and romances of outlaw life.

The traveller from Novae or Oescus (at Gigen, at the mouth of the river Isker) could also reach Philippopolis by the pass of Trojan, close to the sources of the river Asemus (Osma). Finally the long pass of Succi lay on the road from Sardica to Constantinople.

The journey from Singidunum to Constantinople along the main road was reckoned as \texthebrew{•} 670 Roman miles. Singidunum (Belgrade), situated at the junction of the Save with the Danube, was the principal city of the province of Upper Moesia, and was close to the frontier between the eastern and western divisions of the Empire. The road ran at first along the right shore of the Danube, passing Margus (near the village of Dubravica, where the Margus or Morava joins the greater river), till it reached,

p268
ten miles from the Viminacium (close to Kostolats), an important station of the Danube flotilla. Here the traveller, instead of pursuing the eastward road to Durostorum (Silistria), turned southward and again reached the Morava at the town of Horreum Margi, one of the chief factories of arms in the peninsula. The next important town was Naissus (Niš), on the north bank of the Nišava, so strongly fortified that hitherto no enemy had ever captured it. To\texthebrew{‑}day it is the junction of railways, in old days it was the junction of many roads. The Byzantium route continued south-eastward, passing Remesiana (Ak Palanka) to Sardica, the chief town of the province of Dacia Mediterranea, beautifully situated in the large oval plain, under the great mountains, Vitoš on the west and Ryl to the south. From here south-westward ran a road to Ulpia Pautalia (Küstendil) and Dyrrhachium. The traveller pressing to Constantinople, when he left the plain of Sardica, ascended to the pass of Succi in the Ichtimaner mountains. This pass was considered the key of Thrace and was strongly fortified. Descending from this defile the road followed the left bank of the Heberus to Philippopolis (the chief city of the province of Thracia), standing on its three great syenite rocks, with a magnificent view of Mount Rhodope to the south-west. From Philippopolis to Hadrianople (the capital of the province of Haemimontus) was a journey of six days. On the way one passed the fort of Arzus, on a river of the same name (probably the Uzundža). Hadrianople lies at the junction of three rivers; here the Tonzus (Tundža) from the north, and the Artiscus (Arda) from the south, flow into the Hebrus. Another journey of six days brought the traveller to the shore of the Propontis. He passed Arcadiopolis (Lüle Burgas) the ancient Bergule, which the Emperor Arcadius had renamed, on a tributary of the river Erginus.\texthebrew{⁠}\footnote{Not far away was the port of Vrysis, now Bunar Hissar.

} He passed Drusipara (near Karištaran), from which a road led northward to Anchialus on the Black Sea. Then he came to Tzurulon (Corlu), and at last to Heraclea (the old Samian colony of Perinthus) on the sea, now a miserable village. Here the road joined the road from Dyrrhachium and Thessalonica, and the rest of the way ran close to the seashore, past Selymbria\texthebrew{⁠}\footnote{Arcadius renamed it Eudoxiopolis in honour of his wife, but the new name, like so many other names of the kind, soon fell out of use, though it appears in the \emph{Synecdemus} of Hierocles.

} and the fort of Athyras (near

p269
Boyuk-Chekmedže) and Rhegium (at Kuchuk-Chekmedže), to the Golden Gate, which the traveller who tarried not on his way would reach on the thirty-first day after he had left Singidunum.\footnote{From Aquileia the distance was calculated as 47 days. For a pilgrim to Jerusalem walking from Burdigala (Bordeaux) by Aquileia and Singidunum, the distance to Byzantium was 112 days. See the full description of the route in Jireček, \emph{Die Heerstrasse.}

}

When we turn to the western half of the Peninsula, the lands of Illyria and Macedonia, we find an irregular network of mountains, compared with which the configuration of Thrace is simple. In these highlands there are no great plains, and perhaps the first thing to be grasped is that the rivers which water them belong to the systems of the Black Sea and the Aegean, except in the south-west where the Drin and other smaller streams fall into the Hadriatic. Thus the line of watershed between the western and eastern seas runs near the Hadriatic eastward to the range of Scardus (Šar Dagh), which divides the streams that feed the Drilo (Drin) from the western tributaries of the Vardar. The Alpine lands of Dalmatia, using this name in its ancient and wider meaning, are watered by the river Drinus (Drina) and other tributaries of the Save. They are inhospitable and were thinly inhabited and their chief value lay in their mineral wealth.\texthebrew{⁠}\footnote{Especially iron and gold. Statius uses Dalmaticum metallum as a name for gold
(\emph{Silvae}, 1.2.154) . For the whole subject of the Illyrian mine-fields see Evans, \emph{op. cit.} III.6, \emph{sqq.}

} The principal roads connecting these highlands with the Hadriatic were those from Jader (Zara) to Siscia on the Save, and from Salona to Ad Matricem, which corresponds to the modern Sarajevo though it is not on the same site.

The Drina is the western boundary of modern Serbia which answers roughly to the ancient provinces of Moesia prima, Dacia mediterranea, and Dardania. In the centre of this country is the high range known as Kopaonik (mountain of Mines), which with the Yastrebac Planina and the Petrova Gora forms a huge triangle round which the two great branches of the river Morava flow in many curves and windings. The western branch is now known as the Ibar in its upper course and the eastern is sometimes called the Bulgarian Morava.\footnote{Also known as the Binačka.

}

The three places marked out to be the most important inland centres in Illyricum were Naissus, Scupi (Uskub), and Ulpiana. We have seen that the great road from Constantinople to

p270
Singidunum and the west passed Naissus, which lay near the right bank of the western branch of the Margus. Another road connected Naissus directly with Ratiaria (Widin) on the Danube, while south-westward it was linked by a route passing over the Prepolać saddle with Ulpiana,\texthebrew{⁠}\footnote{Afterwards Justiniana secunda.

} which was on the site of the modern village of Lipljan but corresponded in importance to Priština. This town was situated at the southern end of the Kossovo Polje, a plain \texthebrew{•} about twenty miles long, famous as a battlefield in the later Middle Ages. Through this plain ran a road to Ad Matricem which passed Arsa, close to the modern Novipazar, and then turning westward continued its course by Plevlje and Goradža. Two other roads converged at Ulpiana, one from Scupi, which followed the course of the Lepenac, a tributary of the Vardar, and crossed the Kačanik Pass. The other road led to the Hadriatic: crossing the hills it emerged in the open country watered by the upper streams of the Drilo, and known as Metochia, from which it descended to Scodra (Scutari), whence the coast was reached either at Ulcinium (Dulcigno) or at Lissus (Alessio).

Scupi lay on the great road through the valley of the Vardar which brought Thessalonica into communication with the central districts of Illyricum and the Danube. From this centre Naissus could be reached not only by the Kačanik Pass and Ulpiana, but also by another road which skirted the mountains of Kara Dagh and followed the course of the western Margus. The most important station between Thessalonica and Scupi was Stobi, where a north-eastward road diverged to Pautalia and Sardica, while a cross-road connected Stobi with Heraclea (Monastir).

The land communication of Constantinople and Thessalonica with the ports on the Hadriatic was by the great Via Egnatia.\texthebrew{⁠}\footnote{See Tafel, \emph{De via mil. Rom. Egnatia.}

} Westward of Thessalonica, this road ran through western Macedonia and Epirus by Pella, Edessa (Vodena), Heraclea, Lychnidus (Ochrida), Scampae (El Basan), and Clodiana, where it diverged in a northerly direction to Dyrrhachium and in a southerly to Apollonia and Aulon (Valona).\footnote{The provincial divisions of the Dioceses of Thrace and Dacia may here be enumerated.

The \emph{D. of Thrace} (which belonged to the Prefecture of the East) contained six provinces, two north and four south of the Haemus range.

The northern were:

(1) \emph{Lower Moesia \texthebrew{—}} towns: Marcianopolis, Odessus, Durostorum, Novae, Nicopolis (Nicup);

(2) \emph{Scythia} (corresponding to the Dobrudža) \texthebrew{—} towns: Tomi (near Constanza), Callatis

(p271)
(Mangalia), Tropaeum (Adamclissi).

The southern were:

(3) south-eastern, \emph{Europa \texthebrew{—}} towns: Selymbria, Heraclea, Arcadiopolis, Bizye;

(4) south-western, \emph{Rhodope \texthebrew{—}} towns: Aenus, Traianopolis, Maroneia, Rusion;

(5) north-western, \emph{Thrace \texthebrew{—}} towns: Philippopolis, Beroe;

(6) north-eastern, \emph{Haemimontus \texthebrew{—}} towns: Hadrianople, Anchialus.

The \emph{D. of Dacia} contained five provinces:

(1) \emph{Upper Moesia \texthebrew{—}} towns: Singidunum, Viminacium, Margum;

(2) \emph{Dacia ripensis \texthebrew{—}} towns: Bononia, Ratiaria, Castra Martis, Oescus (Gigen);

(3) \emph{Dacia mediterranea \texthebrew{—}} towns: Sardica, Naissus, Pautlia (Küstendil), Remesiana (Ak-Palanka);

(4) \emph{Dardania \texthebrew{—}} towns: Scupi, Ulpiana;

(5) \emph{Praevalitana \texthebrew{—}} towns: Scodra, Lissus.

The \emph{D. of Macedonia} contained besides (1) \emph{Thessaly}, (2) \emph{Achaea}, (3) \emph{Crete}, the provinces of

(4) \emph{Macedonia Prima \texthebrew{—}} towns: Thessalonica, Pella, Beroea, Edessa;

(5) \emph{Macedonia Secunda (Salutaris) \texthebrew{—}} towns: Stobi, Heraclea;

(6) \emph{Old Epirus \texthebrew{—}} towns: Nicopolis, Dodona;

(7) \emph{New Epirus \texthebrew{—}} towns: Dyrrhachium, Scampae, Apollonia, Aulon.

}

p271
Throughout the greater part of the peninsula, north of the Egnatian Way, Latin had become the general language when the Roman conquest was consolidated,\texthebrew{⁠}\footnote{That Latin prevailed in the central and northern provinces there are many indications. For instance, the bishop of Marcianopolis used it in his correspondence with the Council of Chalcedon. Priscus in describing his journey to the court of Attila (see below,

p283 ) says that Latin was the language everywhere. Nicetas, bishop of Remesiana (in the fourth century), who converted the Thracian Bessi, was a Latin writer. The Emperor Justinian, a native of Dardania, speaks of Latin as his own language. The first traces of the development of Latin into Roumanian are found in the sixth century.

} except in Thrace south of Mount Haemus and the southern towns of Macedonia near the coast-line, where the Greek tongue continued to be spoken.

At the beginning of the reign of Theodosius an invasion of the peninsula by a host of Huns was a prelude and a warning. They were led by Uldin, who boasted that he could subdue the whole earth or even the sun. He captured Castra Martis,\texthebrew{⁠}\footnote{On the Danube, near Oescus.

} but as he advanced against Thrace he was deserted by a large multitude of his followers, who joined the Romans in driving their king beyond the Danube. The Romans followed up their victory by defensive precautions. The strong cities in Illyricum were fortified, and new walls were built to protect Byzantium; the fleet on the Danube was increased and improved. But a payment of money was a more effectual barrier against the barbarians than walls, and about A.D. 424 Theodosius consented to pay 350 lbs. of gold to King Rugila.

The tribes of the Huns were ruled each by its own chieftain, but Rugila seems to have brought together all the tribes into a

p272
sort of political unity.\texthebrew{⁠}\footnote{About 430 there seems to have been at least three Hun kings \texthebrew{—} Rugila, his brother Mundiuch, and Octar (probably another brother). Socrates, VII.30;

Jordanes, \emph{Get.} 105 .

} He had established himself between the Theiss and the Danube. The treaty which the government of Ravenna made with Rugila, when the Huns withdrew from Italy in A.D. 425 after the subjugation of the tyrant John, seems to have included the provision that the Huns should evacuate the Pannonian province of Valeria which they had occupied for forty-five years.\texthebrew{⁠}\footnote{Marcellinus, \emph{Chron., sub} 427.

} But soon afterwards a new arrangement was made by which another part of Pannonia was surrendered to them, apparently districts on the lower Save,\texthebrew{⁠}\footnote{Priscus, \emph{fr.} 5, \emph{De leg. gent.} p579 \textgreek{τ\texthebrew{ὴ}ν πρ\texthebrew{ὸ}ς τ\texthebrew{ῷ} Σά\texthebrew{ῳ} ποταμ\texthebrew{ῷ} Παιόνων χώραν}. I am sure that Mommsen and others are wrong in assuming that the province of Savia is meant. The words can equally apply to the parts of Pannonia Secunda west and north of Sirmium, between the Save and the Drave, districts which (like Valeria) were only separated by the Danube from Hunland. I am inclined to suspicion that Valeria was again handed over to the Hun at the same time.

} but not including Sirmium. We may conjecture that this concession was made by Aetius in return for Rugila's help in A.D. 433.\footnote{See above,

p248 .

}

Rugila died soon after this,\texthebrew{⁠}\footnote{According to Socrates, VII.33 (cp. Theodoret, \emph{H.E.} V.37), he was killed by lightning in an invasion of Thrace.

} and he was succeeded by his nephews Bleda and Attila,\texthebrew{⁠}\footnote{The indications are that Bleda was older than Attila, cp. \emph{Chron. Gall.} (A.D. 434), p660 Rugila rex Chunorum, cum quo pax firmata, moritur cui Bleda successit; Marcellinus, \emph{sub} 442, Bleda et Attila. Bleda is the historical proto­type of Blödel, as Attila is of Etzel, in the Nibelungen Lied.

} the sons of Mundiuch, as joint rulers. Bleda played no part on the stage of history. Attila was a leading actor for twenty years, and his name is still almost a household word. He was not well favoured. His features, according to a Gothic historian, ``bore the stamp of his origin; and the portrait of Attila exhibited the genuine deformity of a modern Kalmuck: a large head, a swarthy complexion, small, deep-seated eyes, a flat nose, a few hairs in the place of a beard, broad shoulders, and a short square body of nervous strength though of a disproportioned form. The haughty step and demeanour of the king of the Huns expressed the consciousness of his superiority above the rest of mankind, and he had the custom of fiercely rolling his eyes as if he wished to enjoy the terror which he inspired.''\texthebrew{⁠}\footnote{Gibbon, III. p443, after

Jordanes, \emph{Get.} 182 . For Attila and his relations and wars with the Empire the main source was the history of Priscus. Of this we have one long and a good many small fragments; but we have a great deal of important matter derived from Priscus, through Cassiodorus, in Jordanes.

} He was versed in all the arts of diplomacy, but

p273
the chief aim of his policy was plunder. He was far less cruel than the great Mongolian conqueror of the thirteenth century, Chingiz Khan, with whom he has sometimes been compared; he was capable of pity and could sometimes pardon his enemies.

Attila had some reason for his haughty disdain if he could trace his line of ancestry back for a thousand years and was directly descended from the great chieftains of the Hiung\texthebrew{‑}nu ,\texthebrew{⁠}\footnote{See above,

Chap. IV, p101 .

} whose names have been recorded by early Chinese writers. And if we accept this descent as a genuine tradition, we can infer that he was not of pure Turkish blood. Some of his forefathers had married Chinese princesses, and there may also have been an admixture of the blood of Indo-Scythians .\footnote{The pedigree is preserved in John of Thurócz, \emph{Chronica Hungarorum}, in Schwandtner's \emph{Script. rer. Hung.} I. p81, and has been discussed and compared with Chinese records in the interesting inquiry of Hirth, \emph{Die Ahnentafel Attilas}. In this list Bendegus or Bendeguck appears as Attila's father, and Hirth gives reasons for believing that it is the same name as Mundiuch of the Greek sources. He also seems to succeed in identifying the names of some of the remoter ancestors of the list with the names of Hiung\texthebrew{‑}nu chiefs (between 200 and 60 B.C.) mentioned in Chinese documents.

}

At the beginning of the new reign several points of dispute which had arisen between Rugila and Theodosius were settled. The settlement was entirely to the advantage of the Huns. The Imperial government undertook to double the annual payment, which was thus raised to 700 lbs. of gold; not to receive Hun deserters; to surrender all those who had already deserted; to restore or pay a ransom for Roman prisoners who had escaped; not to form an alliance with any barbarian people at war with the Huns; and to place no restrictions on the trade between the two peoples. The prohibition of receiving fugitives from Attila's empire was particularly important, because the Roman army was largely recruited from barbarians beyond the Danube.

During the early years of his reign, from A.D. 434 to 441, he seems to have been engaged in extending his power in the east towards the Caucasian Mountains. But in A.D. 441 an irresistible opportunity offered itself for attacking the provinces of Theodosius, for in that year the Imperial armies were engaged in operations against both the Vandals and the Persians.

He condescended to allege reasons for his aggression. He complained that the tribute had not been regularly paid, and

p274
that deserters had not been restored. When the Imperial government disregarded his complaints,\texthebrew{⁠}\footnote{Cp. Güldenpenning, \emph{op. cit.} p341. The sources for these invasions of Attila in 441\texthebrew{‑}448 are fragments of Priscus (\emph{De leg. gent.} \emph{fr.} 1\texthebrew{‑}6; \emph{De leg. Rom.}, \emph{fr.} 2\texthebrew{‑}5; and \emph{fr.} 2 in \emph{F. H. G.}, V. p25); Marcellinus; \emph{Chron. Pasch.} (ultimate source: Priscus).

} he appeared on the Danube and laid siege to Ratiaria. Here Roman ambassadors arrived to remonstrate with him for breaking the peace. He replied by alleging that the bishop of Margus had entered the land of the Huns and robbed treasures from the tombs of their kings, and he demanded the surrender of these treasures as well as of deserters. The negotiations broke down, and, having captured and plundered Ratiaria, the Hunnic horsemen rode up the course of the Danube to take the great towns on the banks. Viminacium and Singidunum itself were overwhelmed in the onslaught. Margus, which faces Constantia on the opposite side of the river, fell by treachery; the same bishop whom Attila accused as a grave-robber betrayed a Roman town and its Christian inhabitants to the cruelty of the heathen destroyer. Advancing up the valley of the Margus, the invaders halted before the walls of Naissus, and though the inhabitants made a brave defence, the place yielded to the machines of Attila and the missiles of a countless host. Then the marauders rode south-eastward and approached Constantinople. He did not venture to attack the capital, but he took Philippopolis and Arcadiopolis and the fort of Athyras.\footnote{Theophanes, A.M. 5942. Güldenpenning, \emph{ib.} p344.

}

The strong fortress of Asemus on the Danube, in Lower Moesia,\texthebrew{⁠}\footnote{Asamum near Nicopoli. The name is preserved in that of the river Osma, which flows into the Danube near the place. See CIL III p141.

} won high praise for its valiant resistance to Hunnic squadrons, which separating from the main body had invaded Lower Moesia. They besieged Asemus, and the garrison so effectually harassed them by sallies that they were forced to retreat. A success ­ful defence was not enough for the men of Asemus. Their scouts discovered the times when plundering bands were returning to the camp with spoils, and these moments were seized by the garrison, who unexpectedly assailed these small bodies of Huns and rescued many Roman prisoners.

The Imperial troops, which had been operating against the Persians and the Vandals, must have been available for operations against the Huns in A.D. 442 or 443, but it is not recorded

p275
that Aspar or Areobindus took the field when they returned from Persia and Sicily. We hear that a battle was fought in the Thracian Chersonese and that Attila was victorious, and after this a peace was negotiated by Anatolius ( A.D. 443). The terms were humiliating for the Emperor. Henceforward the annual Hun-tribute of 700 lbs. of gold was to be trebled, and an additional payment of 6000 lbs. was to be made at once. All Hun deserters were to be surrendered to Attila, while Roman deserters were to be handed over to the Emperor for a payment of ten solidi a head.

Hitherto the realm of the Huns had been divided between the two brothers, Bleda and Attila. Of Bleda's government and deeds we hear nothing. We may conjecture that he ruled in the east, from the Lower Danube to the Volga, and Attila in the west. Soon after the Peace of Anatolius, Attila found means to put Bleda to death and unite all the Huns and vassal peoples under his own sway. For the next nine years ( A.D. 444\texthebrew{‑}453) he was the most power ­ful man in Europe.

The Illyrian and Thracian provinces enjoyed a respite from invasion for three years. But in A.D. 447 the Huns appeared again south of the Danube. The provinces of Lower Moesia and Scythia, which had suffered less in the previous incursions, were now devastated. Marcianopolis was taken, and the Roman general Arnegisclus fell in a battle on the banks of the river Utus (Wid). At the same time, another host of the enemy descended the valley of the Vardar and advanced, it is said, to Thermopylae.\texthebrew{⁠}\footnote{It was perhaps in this invasion that Sardica was destroyed. See Priscus, \emph{fr.} 3, \emph{De leg. Rom.} p123.

} Others approached Constantinople, and many of its inhabitants fled from it in terror. So we are told by a contemporary, who says that more than a hundred towns were taken, and that the monks and nuns in the monasteries near the capital were slain, if they had not already fled.\footnote{Callinicus, \emph{Vit. Hypatii}, p108. Callinicus wrote this life in 447\texthebrew{‑}450.

}

Attila was now in a position to enlarge his demands. A new peace was concluded ( A.D. 448) by which a district, along the right bank of the Danube, extending from Singidunum eastward to Novae, and of a breadth of five days' journey, should be left waste and uninhabited, as a march region between the two realms, and Naissus, which was now desolate, should mark the

p276
frontier.\texthebrew{⁠}\footnote{Priscus, \emph{fr.} 5, \emph{De leg. gent.} It seems to have been in 447\texthebrew{‑}448 that the Huns got possession of Sirmium (\emph{id. fr.} 3, \emph{De leg. Rom.} p133). The seat of the Praet. Prefect of Illyricum, which had been moved there in A.D. 437 (see above,

p221 ), was now moved back to Thessalonica.

} But Attila continued to vex the government at Constantinople with embassies, complaints, and demands, and as the drain on the treasury was becoming enormous, the eunuch Chrysaphius conceived the base idea of bribing an envoy of Attila to murder his master. Edecon, the principal minister of Attila, accepted the money and returned to his master's residence, which was somewhere between the rivers Theiss and Körös, in company of a Roman embassy the head of which was Maximin. But the plot was revealed to Attila. He respected the person of the ambassador, but he sent to Constantinople Orestes (a Roman provincial of Pannonia who served him as secretary) with the bag which had held the bribe tied round his neck, and ordered him to ask Chrysaphius in the Emperor's presence whether he recognised it. The punishment of the eunuch was to be demanded. The Emperor then sent two men of patrician rank, Anatolius (Master of Soldiers in praesenti ) and Nomus (formerly Master of Offices), to pacify the anger of the Hun. Attila treated them haughtily at first, but then showed surprising magnanimity and no longer insisted on the punishment of Chrysaphius. He promised to observe the treaty and not to cross the Danube ( A.D. 449\texthebrew{‑}450).

Until the end of the reign of Theodosius the oppressive Hun-money was paid to Attila, but, as we saw, Marcian refused to pay it any longer. It seemed that the Illyrian provinces would again be trampled under the horse-hoofs of the Hun cavalry, though little spoil can have been left to take. But Attila turned his eyes westward, where there was hope of richer plunder, and the realm of Valentinian, not that of Marcian, was now to be exposed to the fury of the destroyer.

Under the rule of Rugila and Attila the Hunnic empire had assumed an imposing size and seemed a formidable power. The extent of Attila's dominion has doubtless been exaggerated, but his sway was effective in the lands (to use modern names) of Austria, Hungary, Roumania, and Southern Russia. How

p277
far northward it may have reached cannot be decided. The most important of the German peoples who were subject to Attila were the Gepids (apparently in the mountainous regions of northern Dacia),\texthebrew{⁠}\footnote{Cp. Schmidt, \emph{Deutsche Stämme}, I.306\texthebrew{‑}307.

} the Ostrogoths (who had migrated westward from their old homes on the Euxine),\texthebrew{⁠}\footnote{\emph{Ib.} 124.

} and the Rugians (somewhere near the Theiss)\texthebrew{⁠}\footnote{\emph{Ib.} 327. The Scirians were also, no doubt, under Hun rule.

} \texthebrew{—} all in the neighbourhood of the lands where the Huns themselves had settled. The Gepid king, Ardaric, was Attila's most trusted counsellor, and next to him, Walamir, one of the Ostrogothic kings. On these peoples he could rely in his military enterprises. Before A.D. 440 the Huns had made an incursion into the Persian empire, and such was the prestige of their arms and Attila's power eight years later that the Roman officers talked of the chances of the overthrow of Persia and the possible consequences of such an event for the Roman world.

Attila indeed looked upon himself as over ­lord of all Europe, including the Roman Empire. Theodosius paid him a huge sum yearly, Valentinian paid him gold too; were they not then his tributaries and slaves? He dreamed of an empire reaching to the islands of the Ocean,\texthebrew{⁠}\footnote{Priscus, \emph{fr.} 3, \emph{De leg. Rom.} p141. Britain, as Mommsen suggests, is probably meant.

} and he was soon to make an attempt to extend it actually to the shores of the Atlantic.\texthebrew{⁠}\footnote{Priscus says that Attila thought himself destined to be lord of the whole world by virtue of the accidental discovery of ``the Sword of Mars.'' A cowherd one day had seen a heifer limping, and following the tracks of blood that had dripped from her wounded foot found a sword on which she had trodden, and brought it to Attila. It was declared to be the Sword of Mars \texthebrew{—} that

\textgreek{\texthebrew{ἀ}κινάκης σιδήρεος} to which the Scythians used to sacrifice animals and men

(Herodotus, IV.62) . So the Alans used to fix a naked sword in the ground and worship it as the god of war

(Amm. Marc. XXXI.2.23) . See

Jordanes, \emph{Get.} 183 , and Priscus, \emph{fr.} 3, \emph{De leg. Rom.}, p142.

} In his dealings with the Empire he had one great military advantage. We have already seen how the Imperial government depended on the Huns and on the Germans beyond the frontier for the recruiting of its armies. Without his Hunnic auxiliaries Aetius would hardly have been able to save as much of Gaul as he succeeded in saving from the rapacity of the German settlers. Attila was in a position to stop these sources of supply. He could refuse to send Hunnic contingents to help the Romans again their enemies; he could forbid individual Huns to leave

p278
their country and enter Roman service; and he could bring pressure to bear on his vassal German kings to issue a similar prohibition to their subjects. That he was fully conscious of this power and made it a feature of his policy, is shown by his stern insistence, in negotiating with Theodosius, that all Hun deserters should be surrendered; perhaps by the device of keeping a strip of neutral territory south of the Danube in order to make it more difficult for his own subjects to pass into the Roman provinces; and particularly by the fact that when his empire was broken up after his death, the empire was inundated by Germans seeking to make their fortunes in Roman service.

Since their entry into Europe the Huns had changed in some important ways their life and institutions. They were still a pastoral people, they did not learn to practise tillage, but on the Danube and the Theiss the nomadic habits of the Asiatic steppes were no longer appropriate or necessary. And when they became a political power and had dealings with the Roman Empire \texthebrew{—} dealings in which diplomacy was required as well as the sword \texthebrew{—} they found themselves compelled to adapt themselves, however crudely, to the habits of more civilised communities. Attila found that a private secretary who knew Latin was indispensable, and Roman subjects were hired to fill the post. But the most notable fact in the history of the Huns at this period is the ascendancy which their German subjects appear to have gained over them. The most telling sign of this influence is the curious circumstance that some of their kings were called by German names. The names of Rugila ,\texthebrew{⁠}\footnote{Priscus calls him Ruas (= Roas in Jordanes); Rugas in Socrates and Roilas in Theodoret (\emph{H.E.} V.37) point to the form Rugila, which is independently preserved in \emph{Chr. Gall., sub} 433. Ruga and Rugila are probably both right, the termination -ila being hypocoristic. Attila (as Mr. H. M. Chadwick informs me) could mean ``little father.'' The derivation of Marquart (\emph{Chron. der alttürk. Insch.} p77) from the Hunnic name of the Volga, Atil or Itil (\textgreek{\texthebrew{Ἀ}ττίλας} in Menander, p8, \emph{De leg. gent.}), should be rejected.

} Mundiuch (Attila's father), and Attila are German or Germanised. This fact clearly points to intermarriages, but it is also an unconscious acknowledgment of the Huns that their vassals were higher in the scale of civilisation. If the political situation had remained unchanged for another fifty years the Asiatic invader would probably have been as thoroughly

p279
Teutonised as the Alans, whom the Romans had now come to class among the Germanic peoples.\footnote{Cp. Jung, \emph{op. cit.} 210, 221. Gothic was the \emph{lingua franca} in Central Europe. Cp. below,

p283 .

}

Of Attila himself we have a clearer impression than of any of the German kings who played leading parts in the period of the Wandering of the Nations. The historian Priscus, who accompanied his friend Maximin, the ambassador to Attila, in A.D. 448, and wrote a full account of the embassy, drew a vivid portrait of the monarch and described his court. The story is so interesting that it will be best to reproduce it in a free translation of the original.\footnote{Priscus in \emph{Exc. de leg.} p123 \emph{sqq.}

}

If the western provinces of the Empire had hitherto escaped the depredations of the Huns, this was mainly due to the personality and policy of Aetius, who had always kept on friendly terms with the rulers. But a curious incident happened, when Attila was at the height of his power, which diverted his rapacity from the east to the west, and filled his imagination with a new vision of power.

Of the court of Valentinian, of his private life, of his relations to his wife and to his mother we know no details. We have seen that he was intellectually and morally feeble, as unfitted for the duties of the throne as had been his uncles Honorius and Arcadius. But his sister Justa Grata Honoria had inherited from her mother some of the qualities we should expect to find in a granddaughter of Theodosius and a great-granddaughter of the first Valentinian. Like Placidia, she was a woman of ambition and selfwill, and she had inherited the temperament of her father which chafed against conventionality. We saw that she had been elevated to the rank of an Augusta probably about the same time that the Imperial title had been conferred

p289
on her brother.\texthebrew{⁠}\footnote{See above

p224 . For her coins (solidi with dn Ivst Grat Honoria pf Avg) see Cohen, VIII.219; De Salis, \emph{Numismatic Chronicle}, VII.203;

Bury, \emph{Justa Grata Honoria}, p4 . The early date of her coronation can be inferred from her coins as well as from the inscription,

CIL XI.276

(above,

p262 ).

} During her girlhood and until Valentinian's marriage her position in the court was important, but when her nieces were born she had the chagrin of realising that henceforward from a political and dynastic point of view she would have to play an obscure part. She would not be allowed to marry except a thoroughly safe man who could be relied upon to entertain no designs upon the throne. We can understand that it must have irked a woman of her character to see the power in the hands of her brother, immeasurably inferior to herself in brain and energy; she probably felt herself quite as capable of conducting affairs of state as her mother had proved herself to be. We can divine that she was a thorn in the side of Valentinian, but we are given no glimpse into the domestic drama played in the Palaces of Ravenna and Rome.

She had passed the age of thirty when her discontent issued in action. She had a separate establishment of her own, within the precincts of the Palace, and a comptroller or steward to manage it. His name was Eugenius, and with him she had an amorous intrigue in A.D. 449.\texthebrew{⁠}\footnote{See Bury, \emph{op. cit.}, where it is shown that the date commonly accepted for the affair with Eugenius, 434, is due to an error of the chronicler Marcellinus and is inconsistent with the story of Priscus and the evidence of Merobaudes. The error arose from the indictional dating: 449 was a 2nd indiction, and Marcellinus made the entry inadvertently under 434, also a 2nd indiction. The sources for Honoria's life are Merobaudes, \emph{Carm.} I; Priscus, \emph{fr.} 2; 7, 8, \emph{De leg. gent.}; John of Antioch, \emph{fr.} 84 \emph{De insidiis} (based on, or taken from, Priscus); Jordanes,

\emph{Get.} 223\texthebrew{‑}224 , \emph{Rom.} 328 (sources: Cassiodorus, \emph{Gothic Hist.}, of which the source here was Priscus, and Marcellinus).

} She may have been in love with him, but love was subsidiary to the motive of ambition. She designed him to be her instrument in a plot to overthrow her detested brother. The intrigue was discovered,\texthebrew{⁠}\footnote{Marcellinus says she was pregnant (concepit), which may or may not be true. He also says that she was sent to Constantinople, but this is inconsistent with the story of Priscus.

} and her paramour was put to death. She was herself driven from the Palace, and betrothed compulsorily to a certain Flavius Bassus Herculanus, a rich senator of excellent character, whose sobriety assured the Emperor that a dangerous wife would be unable to draw him into revolutionary schemes.\footnote{\emph{Fasti cons., sub} 452; CIL IX.1371 .

}

The idea of this union was hateful to Honoria and she bitterly

p290
resented the compulsion. She must often have heard \texthebrew{—} she had perhaps been old enough to have some recollection herself \texthebrew{—} of the breach between her mother and her uncle after her father's death. In that crisis of her life Placidia had turned for help to a barbarian power. Her daughter now decided to do likewise. She despatched by the hands of a trustworthy eunuch, Hyacinthus, her ring and a sum of money to Attila, asking him to come to her assistance and prevent the hateful marriage. Attila was the most power ­ful monarch in Europe and she boldly chose him to be her champion.

The proposal of the Augusta Honoria was welcome to Attila, and was to determine his policy for the next three years. The message probably reached him in the spring of A.D. 450. º She had sent her ring to show that the message was genuine, but he interpreted, or chose to interpret, it as a proposal of marriage. He claimed her as his bride, and demanded that half the territory over which Valentinian ruled should be surrendered to her.\texthebrew{⁠}\footnote{His theory was that the subject territory of the Empire was the private property of the Emperors, in this case of Constantius I and Honorius, and that the children male or female had a claim to equal portions. Attila's Latin secretaries could have informed him that Roman constitutional custom did not recognise such a principle.

} At the same time he made preparations to invade the western provinces. He addressed his demand to the senior Emperor, Theodosius, and Theodosius immediately wrote to Valentinian advising him to hand over Honoria to the Hun. Valentinian was furious. Hyacinthus was tortured, to reveal all the details of his mistress's treason, and then beheaded. Placidia had much to do to prevail upon her son to spare his sister's life. When Attila heard how she had been treated, he sent an embassy to Ravenna to protest; the lady, he said, had done no wrong, she was affianced to him, and he would come to enforce her right to a share in the Empire. Attila longed to extend his sway to the shores of the Atlantic, and he would now be able to pretend that Gaul was the portion of Honoria.

Meanwhile Theodosius had died and we saw how Marcian refused to pay the annual tribute to the Huns. This determined attitude may have helped to decide Attila to turn his arms against the weak realm of Valentinian instead of renewing his attacks upon exhausted Illyrian lands which he had so often wasted. There was another consideration which urged

p291
him to a Gallic campaign. The King of the Vandals had sent many gifts to the King of the Huns and used all his craft to stir him up against the Visigoths. Gaiseric feared the vengeance of Theoderic for the shameful treatment of his daughter,\texthebrew{⁠}\footnote{See above,

p256 .

} and longed to destroy or weaken the Visigothic nation. We are told by a contemporary writer, who was well informed concerning the diplomatic intrigues at the Hun court, that Attila invaded Gaul ``to oblige Gaiseric.''\texthebrew{⁠}\footnote{Priscus, \emph{loc. cit.},

Jordanes, \emph{Get.} 184 .

} But that was only one of his motives. Attila was too wary to unveil his intentions. It was his object to guard against the possibility of the co-operation of the Goths and Romans and he pretended to be friendly to both. He wrote to Tolosa that his expedition was aimed against the enemies of the Goths, and to Ravenna that he proposed to smite the foes of Rome.\footnote{Prosper, \emph{sub} 451,

Jord. \emph{Get.} 185, 186 . A minor matter on the Gallic frontier also engaged Attila's attention. There had been a struggle for kingship among the Ripuarian Franks; they had appealed to him, and the claimant against whom he decided appealed to Aetius. The route which he chose for the invasion of Gaul was perhaps determined by this affair. When he was already on the march he sent another embassy to Ravenna, renewing his demand for the surrender of Honoria and transmitting her ring as a proof of the betrothal (Priscus, \emph{fr.} 8).

}

Early in A.D. 451\texthebrew{⁠}\footnote{For his account of the Gallic campaign Jordanes used Cassiodorus, and the account of Cassiodorus was derived from Priscus. This narrative, doubtless abbreviated and distorted in a reproduction at third hand, is supplemented by Sidonius Apoll. \emph{Carm.} VII and the Latin Chroniclers (Prosper, Marcellinus, etc.). Sidonius intended to write a history of the war, but only began it. Cp. \emph{Epp.} VIII.15.

} he set forth with a large army composed not only of his own Huns, but of the forces of all his German subjects. Prominent among these were the Gepids, from the mountains of Dacia, under their king Ardaric, and the Ostrogoths under their three chieftains, Walamir, Theodemir, and Widimir;\texthebrew{⁠}\footnote{We are not told where precisely the Ostrogoths were settled at this period. Schmidt (\emph{op. cit.} I.124) conjectures with probability that, after they came under the empire of the Huns, they moved westward from their old territory on the Black Sea.

} the Rugians from the regions of the Upper Theiss; the Scirians from Galicia; the Heruls from the shores of the Euxine; the Thuringians;\texthebrew{⁠}\footnote{Sid. Apoll. \emph{Carm.} VII.323 .

} Alans, and others. When they reached the Rhine they were joined by the division of the Burgundians who dwelled to the east of that river and by a portion of the Ripuarian Franks. The army poured into the Belgic provinces, took Metz (April 7),\texthebrew{⁠}\footnote{Hydatius, 150.

} captured many other cities, and laid waste the

p292
land. It is not clear whether Aetius had really been lulled into security by the letter of Attila disclaiming any intention of attacking Roman territory. Certainly his preparations seem to have been hurried and made at the last moment. The troops which he was able to muster were inadequate to meet the huge army of the invader. The federate Salian Franks, some of the Ripuarians, the federate Burgundians of Savoy, and the Celts of Armorica obeyed his summons.\texthebrew{⁠}\footnote{Also Saxons, which shows there were already some Saxon settlements north of the Loire, recognised by the government.

} But the chance of safety and victory depended on securing the co-operation of the Visigoths, who had decided to remain neutral. Avitus, whom we have already met as a \emph{persona grata} at the court of Tolosa, was chosen by Aetius to undertake the mission of persuading Theoderic. He was success ­ful; but it has been questioned whether his success was due so much to his diplomatic arts as to the fact that Attila was already turning his face towards the Loire.\texthebrew{⁠}\footnote{Schmidt, \emph{ib.} 246.

} There was a settlement of Alans\texthebrew{⁠}\footnote{Settled there by Aetius in A.D. 440 (Prosper).

} in the neighbourhood of Valence, and their king had secretly agreed to help Attila to the possession of that city. The objective then of Attila was Orleans, and the first strategic aim of the hastily cemented arrangement between the Romans and Goths was to prevent him from reaching it. The accounts of what happened are contradictory.\texthebrew{⁠}\footnote{The narrative of Jordanes-Priscus implies that Orleans was not besieged by the Huns. Nor do I think that the words of Sidonius Apoll. \emph{Ep.} 8, 15 oppugnatio, inruptio nec direptio (we must allow for the rhetoric) imply that the enemy entered the city. But in this passage we may see the beginning of the ecclesiastical legend, which is expanded in the late \emph{Vita Aniani.} Bishop Anianus probably did much to keep up the hopes and spirits of the alarmed inhabitants. The arrival of the allies in time to save the city would be interpreted as an answer to his prayers. It was natural to magnify the danger and augment the services of the Church by representing the enemy as already within the gates.

} The truth seems to be that the forces of the allies \texthebrew{—} the mixed army of Aetius, and the Visigothic host under Theoderic, who was accompanied by his son Thorismud \texthebrew{—} reached the city before the Huns arrived, and Attila saw that he would only court disaster if he attempted to assault their strongly fortified camp. No course was open but retreat. Aetius had won a bloodless strategic victory (summer A.D. 451).\footnote{The date implied in the \emph{Vita Aniani}, c. VII p113 octavodecimo kal. Iulias = June 14, for the relief of Orleans probably preserves a true tradition. Clinton (\emph{F.R.} I.642), combining Isidore

(\emph{Hist. Goth.} p278)

with Hydatius, puts the battle of

(p293)
Troyes after Sept. 27. Hodgkin conjectures that it was fought early in July

(II.124) . If the \emph{Vita} is right, the battle may be placed about June 20.

}

p293

The Huns took the road to Troyes (Tricasses), and not very far from this town, in a district known as the Mauriac place,\texthebrew{⁠}\footnote{The battle has been vulgarly known as the battle of Châlons, because some of the sources (Jordanes, Hydatius) vaguely describe it as fought in the Catalaunian Plains, an expression which probably denoted nearly the whole of Champagne. That the scene was near Troyes and not near Châlons (Durocatalaunum) is shown by the more precise notices in \emph{Chron. Gall.} p663 Tricassis pugnat loco Mauriacos and \emph{Consul. Ital.} (Prosper Havn.) p302 in quinto miliario de Trecas loco nuncupato Maurica in eo Campania. (Cp. Greg. Tur. II.7 Mauriacum campum; and in the \emph{Lex Burgundionum}, XVII.1, the battle is called pugna Mauriacensis.) It has been thought that the Mauriac name may be preserved in Méry-sur\texthebrew{‑}Seine, which is \texthebrew{•} about 20 miles N of Troyes, and it may be that the battle was fought between Méry and Troyes. But Méry cannot be identified with Maurica of \emph{Cons. Ital.} if the numeral quinto is right. Hodgkin

(\emph{ib.} 139\texthebrew{‑}142)

thinks that the claims of this locality (as against the neighbourhood of Châlons) are made more probable by the discovery in 1842 at Pouan, \texthebrew{•} ten miles from Méry, of bones, weapons, and gold ornaments (including a ring inscribed heva). Peigné-Delacourt was confident that here was the grave of Theoderic. There is, however, no evidence for connecting the bones and ornaments with the battle of 451.

} they halted, and prepared to oppose the confederate army which was marching close upon their heels.\texthebrew{⁠}\footnote{Aetius was posted on the right wing, Theoderic on the left; the Alans, whose treacherous designs were suspected, in the centre. On the other side, Attila was in the centre. The action began with a struggle to gain possession of a hill, in which the Romans and Goths were success­ful.

} The battle, which began in the afternoon and lasted into the night, was drawn; there was immense slaughter,\texthebrew{⁠}\footnote{Jordanes gives the absurd figure of 165,000 for the fallen.

} and king Theoderic was among the slain. Next day, the Romans found that Attila was strongly entrenched behind his wagons, and it was said that he had prepared a funeral pyre in which he might perish rather than fall into the hands of his foes. Thorismud, burning to avenge his father's death, was eager to storm the entrenchment. But this did not recommend itself to the policy of Aetius. It was not part of his design to destroy the Hunnic power, of which throughout his career he had made constant use in the interests of the Empire; nor did he desire to increase the prestige of his Visigothic allies. He persuaded Thorismud to return with all haste to Tolosa, lest his brothers should avail themselves of his absence to contest his succession to the kingship. He also persuaded the Franks to return immediately to their own land. Disembarrassed of these auxiliaries, he was able to pursue his own policy and permit Attila to escape with the remnant of his host.

The battle of Maurica was a battle of nations, but its significance

p294
has been enormously exaggerated in conventional history. It cannot in any reasonable sense be designated as one of the critical battles of the world. The Gallic campaign had really been decided by the strategic success of the allies in cutting off Attila from Orleans. The battle was fought when he was in full retreat, and its value lay in damaging his prestige as an invincible conqueror, in weakening his forces, and in hindering him from extending the range of his ravages. But can the invasion and the campaign regarded as a whole be said to assume the proportions of an ecumenical crisis? The danger did not mean so much as has been commonly assumed. If Attila had been victorious, if he had defeated the Romans and the Goths at Orleans, if he had held Gaul at his mercy and had translated \texthebrew{—} and we have no evidence that this was his design \texthebrew{—} the seat of his government and the abode of his people from the Theiss to the Seine or the Loire, there is no reason to suppose that the course of history would have been seriously altered. For the rule of the Huns in Gaul could only have been a matter of a year or two; it could not have survived here, any more than it survived in Hungary, the death of the great king, on whose brains and personal character it depended. Without depreciating the achievement of Aetius and Theoderic we must recognise that at worst the danger they averted was of a totally different order from the issues which were at stake on the fields of Plataea and the Metaurus. If Attila had succeeded in his campaign, he would probably have been able to compel the surrender of Honoria, and if a son had been born of their marriage and proclaimed Augustus in Gaul, the Hun might have been able to exercise considerable influence on the fortunes of that country; but that influence would probably not have been anti-Roman .

Attila lost little time in seeking to take revenge for the unexpected blow which had been dealt him. He again came forward as the champion of the Augusta Honoria, claiming her as his affianced bride,\texthebrew{⁠}\footnote{After 452 we hear nothing more of Honoria. We are left to wonder whether she was compelled to marry Herculanus, who was consul in that year. A word in John of Antioch seems to hint that some grave punishment befell her. Recording her escape from death in 450 he says that ``Honoria on that occasion (\textgreek{τότε}) escaped from chastisement,'' suggesting that afterwards she was less lucky.

} and invaded Italy in the following year ( A.D. 452). Aquileia, the city of the Venetian march, now fell before the Huns, and was razed to the ground, never to rise again;

p295
in the next century hardly a trace of it could be seen. Verona and Vicentia did not share this fate, but they were exposed to the violence of the invader, while

Ticinum

and Mediolanum were compelled to purchase exemption from fire and sword.

The path of Attila was now open to Rome. Aetius, with whatever forces he could muster, might hang upon his line of march, but was not strong enough to risk a battle. But the land south of the Po, and Rome herself, were spared the presence of the Huns. According to tradition, the thanks of Italy were on this occasion due not to Aetius but to Leo, the bishop of Rome. The Emperor, who was at Rome, sent Leo and two leading senators, Avienus\texthebrew{⁠}\footnote{For Gennadius Avienus, consul in 450, see Sidonius Apoll. \emph{Epp.} I.9.

} and Trygetius, to negotiate with the invader. Trygetius had diplomatic experience; he had negotiated the treaty with Gaiseric in A.D. 435. Leo was an imposing figure, and the story gives him the credit for having persuaded Attila to retreat. He was supported by celestial beings; the apostles Peter and Paul are said to have appeared to Attila and by their threats terrified him into leaving the soil of Italy.\footnote{Pope Leo in his \emph{Sermo in octavis Apost. Petri et Pauli}, LXXXIV (\emph{P. L.} 54), probably refers to this invasion, not to that of the Vandals in 455 (see Gregorovius, \emph{Rome in the Middle Ages}, I.200). Was it, he asks, the circus games or the protection of the saints that delivered Rome from death?

}

The fact of the embassy cannot be doubted. The distinguished ambassadors visited the Hun's camp near the south shore of Lake Garda. It is also certain that Attila suddenly retreated. But we are at a loss to know what considerations were offered him to induce him to depart.\texthebrew{⁠}\footnote{Hydatius, 154. According to Priscus

(Jordanes, \emph{Get.} 222)

it was Attila's own counsellors who decided him to abandon the idea of marching to Rome by reminding him that Alaric had died a few weeks after its capture. There may be something in this. Attila's secretaries were doubtless open to bribes.

} It is unreasonable to suppose that this heathen king would have cared for the thunders or persuasions of the Church. The Emperor refused to surrender Honoria, and it is not recorded that money was paid. A trustworthy chronicle hands down another account which does not conflict with the fact that an embassy was sent, but evidently furnishes the true reasons which moved Attila to receive it favourably. Plague broke out in the barbarian host and their food ran short,\texthebrew{⁠}\footnote{It may be noted that in the winter of 450\texthebrew{‑}451 Italy suffered from a severe famine. See

\emph{Novel} 32

of Valentinian (Jan. 31, 451) obscaenissimam famem per totam Italiam desaevisse.

} and at the same time troops arrived from the east, sent by Marcian to the aid of Italy.

p296
If his host was suffering from pestilence, and if troops arrived from the east, we can understand that Attila was forced to withdraw. But whatever terms were arranged, he did not pretend that they meant a permanent peace. The question of Honoria was left unsettled, and he threatened that he would come again and do worse things in Italy unless she were given up with the due portion of the Imperial possessions.\footnote{Jordanes, \emph{Get.} 223 .

}

Attila survived his Italian expedition by only one year. His attendants found him dead one morning, and the bride whom he had married the night before sitting beside his bed in tears.\texthebrew{⁠}\footnote{\emph{Ib.} 254 . Priscus doubtless is the source and there is no hint at foul play. Ildico was the name of the woman.

} His death was ascribed to the bursting of an artery, but it was also rumoured that he had been slain by the woman in his sleep.\footnote{Marcellinus, \emph{sub} 454. In Teutonic legend the tradition that he was murdered by a woman was preserved, but the lady was Gudrun, the sister of the Burgundian king. Cp. Chadwick, \emph{The Heroic Age}, 37, 156.

}

With the death of Attila, the Empire of the Huns, which had no natural cohesion, was soon scattered to the winds. Among his numerous children there was none of commanding ability, none who had the strength to remove his brothers and step into his father's place, and they proposed to divide the inheritance into portions. This was the opportunity of their German vassals, who did not choose to allow themselves to be allotted to various masters like herds of cattle. The rebellion was led by Ardaric, the Gepid, Attila's chief adviser. In Pannonia near the river Nedao another battle of the nations was fought, and the coalition of German vassals, Gepids, Ostrogoths, Rugians, Heruls and the rest, utterly defeated the host of their Hun lords ( A.D. 454). It is not improbable that the Germans received encouragement and support from the Emperor Marcian.\footnote{The source for the battle is

Jordanes, \emph{Get.} 260 . The Nedao cannot be identified.

}

This event led to considerable changes in the geographical distribution of the barbarian peoples. The Huns themselves were scattered to the winds. Some remained in the west, but the greater part of them fled to the regions north of the Lower Danube, where we shall presently find them, under two of Attila's sons, playing a part in the troubled history of the Thracian provinces. The Gepids extended their power over the whole of Dacia (Siebenbürgen), along with the plains between the Theiss

p297
and the Danube which had been the habitation of the Huns.\texthebrew{⁠}\footnote{Jordanes, \emph{Get.} 264 . It is probable that they also had part of Walachia, see Schmidt, \emph{op. cit.} I.308.

} The Emperor Marcian was deeply interested in the new disposition of the German nations, and his diplomacy aimed at arranging them in such a way that they would mutually check each other. He seems to have made an alliance with the Gepids which proved exceptionally permanent.\texthebrew{⁠}\footnote{Schmidt, 309. Jordanes, \emph{ib.} They received the yearly payments (annua sollemnia) granted to Federates, and this arrangement lasted till the fall of the Gepid power.

} He assigned to the Ostrogoths settlements in northern Pannonia, as federates of the Empire. The Rugians found new abodes on the north banks of the Danube, opposite to Noricum, where they also were for some years federates of Rome. The Scirians settled farther east, and were the northern neighbours and foes of the Ostrogoths in Pannonia; and the Heruls found territory in the same vicinity \texthebrew{—} perhaps between the Scirians and Rugians.\texthebrew{⁠}\footnote{Schmidt, 335.

} But from all these peoples there was a continual flow into the Roman Empire, men seeking military service. In the depopulated provinces of Illyricum and Thrace there was room and demand for new settlers. Rugians were settled in Bizye and Arcadiopolis;\texthebrew{⁠}\footnote{Jordanes, \emph{ib.} 261 . Cp. John Ant. \emph{fr.} 214 (\emph{De ins.} p137).

} Scirians in Lower Moesia.\footnote{Jordanes, \emph{ib.} 265 .

}

The battle of the Nedao was an arbitrament far more momentous than the battle of Maurica. The catastrophe of the Hun power was indeed inevitable, for the social fabric of the Huns and all their social instincts were opposed to the concentration and organisation which could alone maintain the permanence of their empire. But it was not the less important that the catastrophe arrived at this particular moment \texthebrew{—} important both for the German peoples and for the Empire. Although their power disappeared, at one stroke, into the void from which it had so suddenly arisen, we shall see, if we reflect for a moment, that it affected profoundly the course of history. The invasion of the nomads in the fourth century had precipitated the Visigoths from Dacia into the Balkan peninsula and led to the disaster of Hadrianople, and may be said to have determined the whole chain of Visigothic history. But apart from this special consequence of the Hun invasion, the Hun empire performed a function of much greater significance in European history. It

p298
helped to retard the whole process of the German dismemberment of the Empire. It did this in two ways: in the first place, by controlling many of the East German peoples beyond the Danube, from whom the Empire had most to fear; and in the second place, by constantly supplying Roman generals with auxiliaries who proved an invaluable resource in the struggle with the German enemies. The devastations which some of the Roman provinces suffered from the Huns in the last years of Theodosius II and Valentinian III must be esteemed a loss which was more than set off by the support which Hunnic arms had for many years lent to the Empire; especially if we consider that, as subsequent events showed, the Germans would have committed the same depredations if the Huns had not been there. This retardation of the process of dismemberment, enabling the Imperial government to maintain itself, for a longer period, in those lands which were destined ultimately to become Teutonic kingdoms, was all in the interest of civilisation; for the Germans, who in almost all cases were forced to establish their footing on Imperial territory as foederati , then by degrees converted this dependent relation into independent sovranty, were more likely to gain some faint apprehension of Roman order, some slight taste for Roman civilisation, than if their careers of conquest had been less gradual and impeded.

The reward of Aetius for supporting Valentinian's throne for nearly thirty years was that he should fall by Valentinian's hand. One of the most prominent senators and ministers since the later years of Honorius was Petronius Maximus.\texthebrew{⁠}\footnote{He was comes sacr. larg. in 415\texthebrew{‑}417, and Prefect of the City in 420\texthebrew{‑}421. On laying down this office a statue was erected to him by the Emperors, on the petition of the senate and people, in the Forum of Trajan, and there his distinguished pedigree is referred to, a proavis atavisque nobilitas (CIL VI.1749 ). His first consul­ship was in 433, the second in 443.

} He had been twice Prefect of Rome, twice Praetorian Prefect of Italy; he had twice held the consul ­ship; and in A.D. 445 we find him a Patrician. He had a distinguished pedigree, though we do not know it; perhaps he was connected with the great Anician gens. But he probably owed his prestige and influence more to his immense wealth than to his family or to his official career.

p299
He was a notable figure at Rome, ``with his conspicuous way of life, his banquets, his lavish expense, his retinues, his literary pursuits, his estates, his extensive patronage,''\texthebrew{⁠}\footnote{Sidonius, \emph{Epp.} II.13, tr. Dalton.

} In A.D. 454 he was approaching his sixtieth year. He bore personal enmity against Aetius and determined to oust him from power.

He discovered that the sentiments of Heraclius, a eunuch who had the Emperor's ear, were similar to his own. The two conspired together, and persuaded Valentinian that he would perish at the hands of Aetius unless he hastened to slay him first.\footnote{I follow John of Antioch (\emph{fr.} 85, \emph{ib.} p125), because I hold that he followed Priscus. That Maximus played a part in the fall of Actius is confirmed by Marcellinus; Valentinianus dolo Maximi patricii cuius etiam fraude Aetius perierat.

}

Valentinian listened to this counsel and devised death against his power ­ful general. One day, when Aetius was in the Palace, laying some financial statement before the Emperor, Valentinian suddenly leaping from his throne accused him of treason, and not allowing him time to defend himself, drew his sword and rushed upon the defenceless minister, who was at the same moment attacked by the chamberlain Heraclius. Thus perished the Patrician Aetius (September 21, A.D. 454). A poet wrote his epitaph:\footnote{Sidonius Apoll. \emph{Carm.} VII.359 .

}

and it is said that some one afterwards boldly told the truth to Valentinian, ``You have cut off your right hand with your left.'' Who was now to save Italy from the Vandals?

Petronius Maximus assuredly was not the man for the task. It was his ambition to be ``the Patrician'' of the Emperor, but he reckoned without Heraclius. The eunuch persuaded Valentinian that, being well rid of the oppressive influence of Aetius, he would act foolishly if he transferred the power to Maximus. Bitterly disappointed, Maximus wove another murderous plot. He sought out two barbarians, Optila and Thraustila, who had been personal retainers of Aetius, had fought in his campaigns, and enjoyed the favour of the Emperor.\texthebrew{⁠}\footnote{Cp. Prosper, \emph{sub a.} Gregory of Tours calls Optila Occila bucellarius Aetii (\emph{H.F.} II.8).

} He urged these men to avenge their master, and the issue may be told in a chronicler's words:

``It seemed good to Valentinian to ride in the Campus Martius with a few guards accompanied by Optila and Thraustila and

p300
their attendants. And when he dismounted and proceeded to practise archery, Optila and those with him attacked him.\texthebrew{⁠}\footnote{The scene of the attack was called the Two Laurels (Prosper, \emph{Chron. Pasch.} ). There was a place of the same name on the Via Labicana. Cp. Holder-Egger in \emph{Neues Archiv}, I.270, and

Hodgkin, II.198 .

} Optila struck Valentinian on the temple, and when the prince turned to see who struck him dealt him a second blow on the face and felled him. Thraustila slew Heraclius. And the two assassins taking the Imperial diadem and the horse hastened to Maximus. They escaped all punishment for their deed.''\texthebrew{⁠}\footnote{John Ant. \emph{ib.} p126. The story of Valentinian's adultery with the wife of Maximus may or may not be true. The Salmasian fragment, attributed by Müller to John Ant. (\emph{fr.} 200), belongs to some other writer. The story is also found in Procopius, \emph{B. V.} I.4.

} The day of the murder was March 16, A.D. 455.

These two bloody deeds mark the beginning of a new disastrous period in the history of the western provinces. The strong man who might have averted the imminent danger from the Vandals, and the weak man whose mere existence held Italy, Gaul, and Spain together, were removed; there was no general to take the place of Aetius, ``the last of the Romans,''\texthebrew{⁠}\footnote{Procopius, \emph{ib.}, where Aetius and Boniface are so described. The compliment to Aetius is weakened by the inclusion of Boniface.

} as there was no male member of the Theodosian house to succeed Valentinian. A chronicler speaks\texthebrew{⁠}\footnote{Marcellinus.

} of the Patrician Aetius as ``the great safety of the western republic'' ( magna occidentalis reipublicae salus ), the terror of king Attila; ``and with him the Hesperian realm fell, and up to the present day has not been able to raise its head.'' We can comprehend this judgment; the death of Aetius was a grave event. He was the greatest of the three Romans who had been responsible for the defence of Italy and the western provinces since the fall of Stilicho, and he was to have no Roman successor. Two years after his death the supreme command of the Imperial forces would again pass into the hands of a Romanised German. But we must not leave out of sight the importance of the death of his master Valentinian without male offspring. A legitimate heir of the Theodosian house would have prevented some of the troubles which befell Italy in the following years.

An amazing sequence of events had surprised the Empire after the death of Theodosius the Great. Provinces had been seized by barbarous invaders, and the very soil of Italy desecrated by German violence. The sight of Rome herself stricken and insulted, no longer able to speak the language of a mistress but compelled to bargain with the intruders on her own territory, could not fail to make men ask, ``What is the cause of these disasters? Civil wars there have been in the past, our frontiers have been crossed, our provinces invaded, but since the Gauls bore down on Rome nearly eight hundred years ago, the queen of the world has never been violated and plundered by a foreign enemy till now, and it hardly entered any man's dream that such a horror might some day come to pass.'' In that age there was probably no one who held the view that political and social changes depend on the series of antecedent events and that sudden catastrophes are no exception. It was in the will of heaven, the anger of divine tyrants, or the inscrutable operations of the stars, that men were prone to seek explanations of shocking or unexpected public calamities.

Pagan patriots had no difficulty in solving the problem. ``So long,'' they said, ``as the gods under whose favour Rome won her Empire were supreme, so long as the traditions of the ancient religion were preserved, our empire flourished and was impregnable. But now their temples are destroyed, impious hands have been laid on the altars, the worship of our divinities has been proclaimed a crime. And what is the result? Has the alien deity, who has usurped the time-honoured prerogatives, conducted the state to new glory or even to its old prosperity ? On the contrary, the result of his supremacy is rapine and ruin. The Empire is inundated by a wild tide of rapacious savages, the dominions of Rome are at their mercy, her sword is broken, and her lofty walls have been scaled. These are the gifts that Constantine and the religion of Galilee, which he embraced in a disastrous hour, have bestowed upon the world.''\footnote{This point of view appears in the writings of pagan historians like Eunapius and Zosimus.

}

p302
Similar arguments indeed had been urged long before. In the third century pagans had made Christianity answerable for plagues, droughts, and wars; nature herself, they cried, had changed, since the advent of this abominable religion. Two African divines had replied to the charge. Cyprian the bishop of Carthage declared\texthebrew{⁠}\footnote{See Cyprian, \emph{Ad Demetrianum.}

} that the disasters of his day were signs of the approaching end of the world, and the inference might be drawn that they did not much matter in view of the vast event so soon to happen. Arnobius of Sicca, half a century later, in his \emph{Seven Books against the Nations}, met the arguments of the heathen by pointing out that before the appearance of Christianity the world had been the scene of as great or rather of greater calamities.

But in the early fifth century there was stuff for a more telling indictment, and one to which the average Christian of that age might find it hard to produce a convincing answer. And the Christian himself might have his own difficulties. How, he might wonder, is it compatible with a wise and just government of the universe that the godly who hold the right opinion concerning the nature of the Trinity should suffer all these horrors at the hands of barbarians, and that those barbarians who believe in a blasphemous heresy, which places them as much as the heathen outside the Christian pale, should triumph over us and wrest our provinces from us.\footnote{The Eunomian historian Philostorgius, who graphically described the miseries which the provinces suffered in the reigns of Arcadius and Theodosius II, attributed the calamities to the persecution of the Eunomians. It seems likely that he shared the view of those who saw in these contemporary events the signs of the end of the world predicted by Jesus in the Gospels. This view is reflected in a clearly contemporary Apocalypse, preserved in Syriac and translated by Arendzen, in \emph{Journal of Theological Studies}, 1901, 401 \emph{sqq.} (Cp. Bidez, \emph{Proleg.} to his ed. of Philostorgius, CXV \emph{sqq.}). How soon the impression of great events may grow faint is illustrated by the fact that Socrates and Sozomen, who wrote in the middle of the fifth century, notice Alaric's capture of Rome as if it were no extraordinary event; Theodoret does not even mention it. We can see from the abridged notice which has been preserved how differently it struck Philostorgius

(XII.3) .

}

Such questionings evoked three books. Africa, Spain, and Gaul each contributed an answer, one a work of genius, the other two dull but remarkable each in its way.

The first, as it was the greatest, was

Augustine's \emph{City of God} . Augustine had been deeply impressed by the capture of Rome by Alaric, and he recognised that the situation of the world called

p303
for a Christian explanation in reply to the criticisms of the pagans who made the new religion responsible for Rome's misfortunes. The motive and occasion of the work, which seems to have outgrown its original scope, may account for some of its defects.\texthebrew{⁠}\footnote{It was composed in the years 413 to 426 and parts of it were published before it was completed. In reading Augustine it is always well to bear in mind that he was a Neoplatonist before he was a Catholic, and a Manichean before he was a Neoplatonist. The stages of his thought have recently been studied by P. Alfaric, \emph{L'Évolution intellectuelle de Saint Augustin}, 1918. It is also well to remember that he was a rhetorician. His most interesting work \emph{The Confessions} is marred for modern taste by its rhetoric.

} It is one of the greatest efforts of Christian speculation, but the execution is not equal to the conception, and the fundamental conception itself was not original. The work consists of two distinct sections which might just as well have formed two independent treatises. The first section (Bks. I\texthebrew{‑}X ) is a polemic against pagan religion and pagan philosophies, in which it is shown that polytheism is not necessary to secure happiness either in this world or in the next. The most effective argument is that which had been already used by Arnobius: the miseries which we suffer to\texthebrew{‑}day are no exception to the general course of experience, for we have only to read the history of Rome to find them paralleled or exceeded. The writer insists that earthly glory and prosperity are unnecessary for true happiness. These things were bestowed on Constantine the Great, but that was in order to prove that they are not incompatible with the life of a Christian. On the other hand, if the reign of Christian Jovian was shorter than that of the apostate Julian, and if Gratian was assassinated, these were divine intimations that glory and long life are not the true reward of the Christian faith.\texthebrew{⁠}\footnote{\emph{De civ. Dei}, V.25 .

} Such an argument was not likely to make much impression upon pagans.

But the answer of Augustine to the questions which were perplexing the world is not to be found in the first part of his work. He realised that any satisfactory solution of the problem must lie in discovering a harmony between the actual events of history and the general plan of the universe. The synthesis which he framed for the interpretation of history as part of a general scheme of things is an essay in that field of speculation which is known nowadays as the philosophy of history. It can hardly, however, be described as philosophical, for the premisses

p304
on which it is based are not derived from reason but from revelation.\footnote{On its influence in later times see J. N. Figgis, \emph{The Political Aspects of St. Augustine's City of God.}

}

Augustine's conception is that the key to the history of the human race is to be found in the coexistence side by side of two cities or states which are radically opposed to each other in their natures, principles, and ends, the Civitas Dei and the Civitas Terrena. It may be observed that this conception was not original; Augustine derived it from his Donatist friend Tychonius. The origins of both these states go back to a time when man did not yet exist; the City of God was founded by the creation of the angels, the other city by the rebellion of the angels who fell. Since the sin of Adam the history of each of these cities, ``intertwined and mutually mixed'' ( perplexas quodam modo invicemque permixtas ), has been running its course. The vast majority of the human race have been and are citizens of the earthly city, of which the end is death. The minority who belong to the heavenly city are during their sojourn on earth merely foreigners or pilgrims ( peregrini ) in the earthly city. Till the conversion of the first Gentile to Christianity the members of the City of God belonged exclusively to the Hebrew race and its patriarchal ancestors; and Augustine determines the chief divisions of universal history by the great epochs of the biblical record: the Flood, Abraham, David, the Captivity, and the birth of Christ.\texthebrew{⁠}\footnote{Augustine here depended on the chronological scheme worked out by Hippolytus, Sextus Julius Africanus, and Eusebius.

} This last event is the beginning of the sixth period, in which we are living at present; and the sixth period is the last. For the periods of history correspond to the days of Creation, and as God rested on the seventh day, so the seventh period will witness the triumph of the heavenly City and the eternal rest of its citizens. To the question how long will the sixth period last, Augustine replies that he does not know.\texthebrew{⁠}\footnote{\emph{De civ. Dei},

XVIII.53 ;

XXII.30 .

} In this connexion he tells us an interesting fact. An oracle was current among the pagans, and seems to have given them much consolation, that the Christian religion would disappear from the world at the end of 365 years. It was said that the disciple Peter had been able by his sorceries to impose upon the world the worship of Christ for this period, but at its termination the work of the wizard would dissolve like a dream.

p305
Augustine observes triumphantly, and perhaps with a certain relief, that more years than 365 had already elapsed since the Crucifixion, and that there was no sign of the fulfilment of the oracle.\footnote{\emph{De civ. Dei}, XVIII.53, 54 .

}

To a modern, and possibly also to an ancient, inquirer, Augustine's work would have been more interesting if he had seriously addressed himself to an historical study of the Babylonian and Roman Empires, which according to him were the two principal embodiments of the earthly City. But he entrenches himself and remains almost immovably fixed in his headquarters in Judaea, and the excursions which he makes into other regions are few and slight. Many of his notices of events in secular history are simply trivial.

Having completed his historical survey he devotes the last portion of his work to an exposition of the ultimate goal to which the world and the human race are travelling. He examines the question of the Last Judgment, expatiates on the fiery death which is the destiny of the earthly City, and ends with a discussion on the bliss which awaits the citizens of the City of God.

Among the thinkers of the Middle Ages the influence of Augustine's work went far and deep. But his fruitful conception was lodged in a somewhat dreary mansion. If the polemical section which he intends to be a preliminary defeat of the enemies of the City of God\texthebrew{⁠}\footnote{Cp. \emph{ib.} XI.1 .

} had been omitted, the work would have gained in simplicity. But the main argument itself, although it has a definite architectural scheme,\texthebrew{⁠}\footnote{He indicates the design repeatedly in the work itself and in his \emph{Retractationes.} Bks. XI\texthebrew{‑}XIV deal with the exortus of the two cities and come down to Cain and Abel; Bks. XV\texthebrew{‑}XVIII describe the procursus, and survey the historical growth of the Civ. Dei; and the last four Books deal with the future and the debiti fines of the two cities.

} is marred by diffuseness and digressions. Augustine did not possess the literary art or command the method of lucid exposition whereby the prince of Greek philosophers compels his readers to assist in the building of the City, ``of which a model perchance is in heaven,'' with breathless interest from page to page and from section to section. There is at least one part which may hold the attention of the reader, fascinated by the very horror, the Book in which this arch-advocate of theological materialism and vindictive punishment expends all his ingenuity in proving that the fire

p306
of hell is literal fire and spares no effort to cut off the slenderest chance that the vast majority of his fellow-beings will not be tormented throughout eternity.

Augustine had produced a book which transcended in importance its original motive. But it is this motive which concerns us here. It was to teach the world to take a right view of the misfortunes which were befalling the Empire, and to place them in their true perspective. He says in effect to the pagans, ``These misfortunes are nothing exceptional, they are simply part of the heritage of your City of sin and death.''\texthebrew{⁠}\footnote{It is to be observed that Augustine regarded the virtues of the pagans as vices, because vera religio was absent

(XIX c25) . Chrysostom seems to take a more lenient and human view, \emph{Hom.} 5, on \emph{Ep. ad Rom.}, \emph{P.G.} 60, 426 \emph{sqq.}

} To the Christians he said, ``These things do not really concern you. Your interests are not affected by the calamities of a country in which you are merely foreigners.'' This theory might be consolatory, but if it were pressed to its logical conclusion it would assuredly be destructive of the spirit of patriotism; and, though the author would doubtless have deprecated this criticism, he does not consider the secular duties of Christians towards the state of which they are citizens in the earthly sense.

He was conscious that his treatment of the history of Rome was casual and superficial, and he thought that a fuller development of his historical argument in reply to the pagans was desirable. He requested his friend Orosius, a Spanish priest, to supply this need. He said to Orosius, ``Search the annals of the past, collect all the calamities which they record, wars, plagues, famines, earthquakes, fires, and crimes, and write a history of the world. Thus my general refutation of the charges of the unbelievers who impute to our religion the present misfortunes, which they allege to be unusual, will be proved abundantly by a long array of facts.''\texthebrew{⁠}\footnote{See

Orosius, \emph{Hist., Prol.}

} A work entitled \emph{Histories to confute the Pagans} was the outcome of this request, and it may thus be regarded as a sort of supplement to the \emph{City of God.}\texthebrew{⁠}\footnote{The date of the work is A.D. 418.

} Perhaps it deserves more than any other book to be described as the first attempt at a universal history, and it was probably the worst. But it had considerable vogue in the Middle Ages, and gave currency to the idea of four great monarchies, the

p307
Babylonian, Carthaginian, Macedonian, and Roman, corresponding to the four points of the compass.\footnote{\emph{Hist., Prol.} II.1. The number was based on Daniel, chap. II. Sulpicius Severus (\emph{Chron.} II.3) makes the four kingdoms, the Chaldaean, Persian, Macedonian, and Roman.

}

Fifteen or twenty years after the completion of Augustine's work Salvian, a priest of Marseilles, wrote his treatise \emph{On the Government of God},\texthebrew{⁠}\footnote{\emph{De gubernatione Dei.} It was written not earlier than 439 and before Attila's invasion in 451.

} dealing from a different point of view with the same problem which had suggested the books of Augustine and Orosius. Salvian addresses his discourse expressly to Christians, for he has no hope that his arguments would have any effect upon pagans.\texthebrew{⁠}\footnote{III. V.

} He propounds the question: How comes it that we Christians who believe in the true God are more miserable than all men? Is God indifferent to us? Has he renounced the business of governing the world? If he regards human affairs, why are we weaker and more unfortunate than all other peoples? Why are we conquered by the barbarians? Salvian's answer is, We suffer these evils because we deserve them. If, living in such vice and wickedness as we do, we flourished and were happy, then indeed God might be accused of not governing. In support of his argument the author paints an appalling picture of the condition of the Empire. His descriptions of the corruptness of the administration and of the oppression of the poor by the rich furnish the modern historian with an instructive commentary on those Imperial laws which attempt to restrain the rapacity of public officials. Salvian does not forget to dwell, with the zeal of a churchman, on the general love of unedifying pleasures, the games of the circus and licentious plays in the theatre, amusements of which the average Christian was not less avid than the average pagan.

But, it might be objected, we, whatever our faults, have at least right theological beliefs, whereas the barbarians who are permitted to overcome us are heathen or heretics. That is true, replies Salvian; in just one point we are better than they; but otherwise they are better than we. He then proceeds to enlarge on the virtues of the barbarians, which he uses, somewhat as Tacitus did in the

\emph{Germania} , as a foil to Roman civilisation. Among the Germans, or even among the Huns, we do not see the poor oppressed by the rich. If the Alamanni are

p308
given to drunkenness, if the Franks and Huns are perjured and perfidious, if the Alans are rapacious, are not all these vices found among us? On the other hand, the Vandals have put the provincials to shame by their high standard of sexual morality, and if the Saxons are ferocious and the Goths perfidious, both these peoples are wonderfully chaste.

There is no relief in Salvian's gloomy picture. It must be accepted with the reserves with which we must always qualify the rhetoric of preachers or satirists when they denounce the vices of their age. But the tone of despondency is genuine. He says that ``the Roman Republic is either dead, or at least is drawing her last breath in those parts in which she still seems to be alive.''\texthebrew{⁠}\footnote{\emph{De gubernatione Dei}, IV.30.

} He speaks as if this were a fact which was beyond dispute and to which men had already become accustomed. More than thirty years had elapsed since the news of the Goths at Rome had surprised Jerome in his retreat at Bethlehem and extorted the cry, Quid salvum est si Roma perit? Meanwhile the Romans had quickly recovered from the shock and had almost forgotten it. The calamity of the provinces did not move them to alter their way of life or renounce their usual amusements. And the one phrase that is worth remembering in Salvian's gloomy, declamatory book is the epigram on Rome, Moritur et ridet .

The explanations of the calamities of the Empire which have been hazarded by modern writers are of a different order from those which occurred to witnesses of the events, but they are not much more satisfying. The illustrious historian whose name will always be associated with the ``Decline'' of the Roman Empire invoked the ``principle of decay,'' a principle which has itself to be explained. Depopulation, the Christian religion, the fiscal system have all been assigned as causes of the Empire's decline in strength.\texthebrew{⁠}\footnote{Gibbon, IV. chap. XXXVIII.173 \emph{sqq.}

Hodgkin, \emph{Italy and her Invaders}, II.532 \emph{sqq.} , enumerates as contributory causes Christianity, the destruction of the middle classes, and ``barbarous finance.'' Seeley, ``Roman Imperialism,'' in \emph{Macmillan's Magazine}, August 1869, makes depopulation mainly responsible. Over-taxation of the rich has sometimes been assigned as one of the causes of the ``fall'' of the Empire. It has been pointed out above

(p253) º that under-taxation of the rich was rather the trouble.

} If these or any of them were responsible

p309
for its dismemberment by the barbarians in the West, it may be asked how it was that in the East, where the same causes operated, the Empire survived much longer intact and united.

Consider depopulation. The depopulation of Italy was an important fact and it had far-reaching consequences.\texthebrew{⁠}\footnote{Seeck, \emph{Untergang}, vol. I.

} But it was a process which had probably reached its limit in the time of Augustus. There is no evidence that the Empire was less populous in the fourth and fifth centuries than in the first.\texthebrew{⁠}\footnote{Cp. above,

Chap. III § 4 .

} The ``sterility of the human harvest'' in Italy and Greece affected the history of the Empire from its very beginning, but does not explain the collapse in the fifth century. The truth is that there are two distinct questions which have been confused. It is one thing to seek the causes which changed the Roman State from what it was in the best days of the Republic to what it had become in the age of Theodosius the Great \texthebrew{—} a change which from certain points of view may be called a ``decline.'' It is quite another thing to ask why the State which could resist its enemies on many frontiers in the days of Diocletian and Constantine and Julian suddenly gave way in the days of Honorius. ``Depopulation'' may partly supply the answer to the first question, but it is not an answer to the second. Nor can the events which transferred the greater part of western Europe to German masters be accounted for by the numbers of the peoples who invaded it. The notion of vast hosts of warriors, numbered by the hundreds of thousands, pouring over the frontiers, is, as we saw, perfectly untrue.\texthebrew{⁠}\footnote{Cp. above,

Chap. IV § 3 .

} The total number of one of the large East German nations probably seldom exceeded 100,000, and its army of fighting men can rarely have been more than from 20 to 30,000. They were not a deluge, overwhelming and irresistible, and the Empire had a well-organised military establishment at the end of the fourth century, fully sufficient in capable hands to beat them back. As a matter of fact, since the defeat at Hadrianople which was due to the blunders of Valens, no very important battle was won by German over Imperial forces during the whole course of the invasions.

It has often been alleged that Christianity in its political effects was a disintegrating force and tended to weaken the power of Rome to resist her enemies. It is difficult to see that

p310
it had any such tendency, so long as the Church itself was united. Theological heresies were indeed to prove a disintegrating force in the East in the seventh century, when differences in doctrine which had alienated the Christians in Egypt and Syria from the government of Constantine facilitated the conquests of the Saracens. But, after the defeat of Arianism, there was no such vital or deep-reaching division in the West, and the effect of Christianity was to unite, not to sever, to check, rather than to emphasise, national or sectional feeling. In the political calculations of Constantine it was probably this ideal of unity, as a counterpoise to the centrifugal tendencies which had been clearly revealed in the third century, that was the great recommendation of the religion which he raised to power.\texthebrew{⁠}\footnote{Cp. below,

Chap. XI § 2 , on the political bearing of the law of 445 in favour of the Roman See.

} Nor is there the least reason to suppose that Christian teaching had the practical effect of making men less loyal to the Empire or less ready to defend it. The Christians were as pugnacious as the pagans. Some might read Augustine's \emph{City of God} with edification, but probably very few interpreted its theory with such strict practical logic as to be indifferent to the safety of the Empire. Hardly the author himself, though this has been disputed.

It was not long after Alaric's capture of Rome that Volusian, a pagan senator of a distinguished family,\texthebrew{⁠}\footnote{On the family of the Albini and Volusiani see Seeck, \emph{Praef.} to Symmachus, p. clxxiv \emph{sqq.} Augustine's correspondent was comes rei priv. in 408.

} whose mother was a Christian and a friend of Augustine, proposed the question whether the teaching of Christianity is not fatal to the welfare of a State, because a Christian smitten on one cheek would if he followed the precepts of the Gospel turn the other to the smiter. We have the letter\texthebrew{⁠}\footnote{\emph{Epp.} 138, addressed to their common friend Marcellinus.

} in which Augustine answers the question and skilfully explains the texts so as to render it consistent with common sense. And to show that warfare is not forbidden another text is quoted in which soldiers who ask ``What shall we do?'' are bidden to ``Do violence to no man, neither accuse any falsely, and be content with your wages.'' They are not told not to serve or fight. The bishop goes on to suggest that those who wage a just war are really acting misericorditer , in a spirit of mercy and kindness to their enemies, as it is to the true

p311
interests of their enemies that their vices should be corrected. Augustine's misericorditer laid down unintentionally a dangerous and hypocritical doctrine for the justification of war, the same principle which was used for justifying the Inquisition. But his definite statement that the Christian discipline does not condemn all wars was equivalent to saying that Christians were bound as much as pagans to defend Rome against the barbarians. And this was the general view. All the leading Churchmen of the fifth century were devoted to the Imperial idea, and when they worked for peace or compromise, as they often did, it was always when the cause of the barbarians was in the ascendant and resistance seemed hopeless.\footnote{The monastic movement was anti-social, but in the period in question it was young, and cannot have withdrawn so many young men from public service as to affect appreciably the strength of the State.

}

The truth is that the success of the barbarians in penetrating and founding states in the western provinces cannot be explained by any general consideration . It is accounted for by the actual events and would be clearer if the story were known more fully. The gradual collapse of the Roman power in this section of the Empire was the consequence of \emph{a series of contingent events.} No general causes can be assigned that made it inevitable.

The first contingency was the irruption of the Huns into Europe, an event resulting from causes which were quite independent of the weakness or strength of the Roman Empire. It drove the Visigoths into the Illyrian provinces, and the difficult situation was unhappily mismanaged. One Emperor was defeated and lost his life; it was his own fault. That disaster, which need not have occurred, was a second contingency.\texthebrew{⁠}\footnote{It may be remembered that Valens would not wait for Gratian, who was hastening to his help.

\texthebrew{◂} previous

next \texthebrew{▸}Images with borders lead to more information.

The thicker the border, the more information.

(Details here.)

UP TO:

J. B. Bury
Homepage

Latin \& Greek

Texts

LacusCurtius

Home} His successor allowed a whole federate nation to settle on provincial soil; he took the line of least resistance and established an unfortunate precedent. He did not foresee consequences which, if he had lived ten or twenty years longer, might not have ensued. His death was a third contingency. But the situation need have given no reason for grave alarm if the succession had passed to an Emperor like himself, or Valentinian I, or even Gratian. Such a man was not procreated by Theodosius and the government of the West was inherited by a feeble-minded boy. That

p312
was a fourth event, dependent on causes which had nothing to do with the condition of the Empire.

In themselves these events need not have led to disaster. If the guardian of Honorius and director of his government had been a man of Roman birth and tradition, who commanded the public confidence, a man such as Honorius himself was afterwards to find in Constantius and his successor in Aetius, all might have been tolerably well. But there was a point of weakness in the Imperial system, the practice of elevating Germans to the highest posts of command in the army. It had grown up under Valentinian I, Gratian, and Theodosius; it had led to the rebellion of Maximus, and had cost Valentinian II his life. The German in whom Theodosius reposed his confidence and who assumed the control of affairs on his death probably believed that he was serving Rome faithfully, but it was a singular misfortune that at a critical moment when the Empire had to be defended not only against Germans without but against a German nation which had penetrated inside, the responsibility should have devolved upon a German. Stilicho did not intend to be a traitor, but his policy was as calamitous as if he had planned deliberate treachery. For it meant civil war. The dissatisfaction of the Romans in the West was expressed in the rebellion of Constantine, the successor of Maximus, and if Stilicho had had his way the soldiers of Honorius and of Arcadius would have been killing one another for the possession of Illyricum. When he died the mischief was done; Goths had Italy at their mercy, Gaul and Spain were overrun by other peoples. His Roman successors could not undo the results of events which need never have happened.

The supremacy of a Stilicho was due to the fact that the defence of the Empire had come to depend on the enrolment of barbarians, in large numbers, in the army, and that it was necessary to render the service attractive to them by the prospect of power and wealth. This was, of course, a consequence of the decline in military spirit, and of depopulation, in the old civilised Mediterranean countries. The Germans in high command had been useful, but the dangers involved in the policy had been shown in the cases of Merobaudes and Arbogastes. Yet this policy need not have led to the dismemberment of the Empire, and but for that series of chances its western provinces would not

p313
have been converted, as and when they were, into German kingdoms. It may be said that a German penetration of western Europe must ultimately have come about. But even if that were certain, it might have happened in another way, at a later time, more gradually, and with less violence. The point of the present contention is that Rome's loss of her provinces in the fifth century was not an ``inevitable effect of any of those features which have been rightly or wrongly described as causes or consequences of her general 'decline.' `` The central fact that Rome could not dispense with the help of barbarians for her wars ( gentium barbararum auxilio indigemus ) may be held to be the cause of her calamities, but it was a weakness which might have continued to be far short of fatal but for the sequence of contingencies pointed out above.

\xchapter{Chapter X}

It was always a critical moment when an Emperor died without a designated successor or a member of his family marked out to claim the diadem. Theodosius I had created his sons Augusti; Arcadius had co-opted his infant son; Theodosius II had designated Marcian as his successor just before his death, and Marcian's title was sealed by his marriage with the Augusta Pulcheria. On Marcian's death the Theodosian dynasty had come to an end, and the choice of a new Emperor rested with the army and the Senate. There was one obvious candidate, Anthemius, who was the grandson of the great Praetorian Prefect and had married Marcian's daughter Euphemia. He had held the office of Master of Soldiers in Illyricum, and had been consul in A.D. 455. But Marcian had not designated him as his successor, and though the Senate perhaps would have liked to elect him,\texthebrew{⁠}\footnote{Sidonius, \emph{Carm.} II.214 quamquam te posceret ordo. The poet asserts that he did not covet the throne, 210. From this poem we learn that he distinguished himself in defending Illyricum against the Ostrogoths under Walamir.

} he was not favoured by the man of most authority in the army, the patrician Aspar, who with his father Ardaburius had distinguished himself thirty-five years before in the suppression of the usurper John. Being an Arian, as well as a barbarian, he could not hope to wear the Imperial diadem; the only course open to his ambition was to secure the elevation of one on whose pliancy he might count. He chose Leo, a native of Dacia and an orthodox Christian, who was tribune of the Mattiarii,\texthebrew{⁠}\footnote{Constantine Porph. \emph{De cer.} I p411. The Mattiarii seniores were under one mag. mil. in praes., the Mattiarii iuniores under the other. In the former case they are associated with Dacians.

(p315)
See \emph{Not. Dig., Or.} VI.42 ,

V.47 . Leo's Dacian origin is mentioned by Candidus, \emph{F. H. G.} IV p135; John Mal. XIV p369, says he was a Bessian. He had the rank of count.

} a legion

p315
belonging to the troops which were under the control of a Master of Soldiers in praesenti . Aspar doubtless held this post, as Leo was his domesticus . The Senate was unable to reject the general's nominee and (on February 7) Leo was crowned at the Palace of Hebdomon. As there was no Augustus or Augusta to perform the ceremony of coronation, this duty was assigned to the Patriarch Anatolius, who had perhaps taken some part in the coronation of Marcian.\texthebrew{⁠}\footnote{See above,

p236 .

} We have a contemporary description of the ceremonies connected with Leo's elevation, though the act of crowning is passed over.

The senators and officials, the Scholarian guards, the troops which were present in the capital, and the Patriarch gathered at the Campus in the Hebdomon. The military insignia, the labara and the standards, lay on the ground. All began to cry, ``Hear, O God, we call upon thee. Leo will be Emperor. The public weal demands Leo. The army demands Leo. The palace expects Leo. This is the wish of the palace, the army, and the Senate.'' Then Leo ascended the tribunal or raised platform, and a chain was placed on his head, and another in his right hand, by officers.\texthebrew{⁠}\footnote{Campiductores, army-guides. Perhaps they were attached to the legion of the Lanciarii, for a \textgreek{καμπιδούκτωρ τ\texthebrew{ῶ}ν λαγκιαρίων} performed the same office at the elevation of Anastasius

(Const. Porph. \emph{op. cit.} p423) . The Greek word for the chain or torc is \textgreek{μανιάκιν.}} Immediately the labara were collected, and all cried: ``Leo Augustus, thou conquerest!\texthebrew{⁠}\footnote{\textgreek{Συ νικ\texthebrew{ᾷ}ς}. But the Latin \textgreek{το\texthebrew{ὺ} βίγκας} (tu vincas) remained long a regular acclamation in the Byzantine Hippodrome. We also meet the hybrid \textgreek{σ\texthebrew{ὺ} βίγκας.}} God gave thee, God will keep thee. A long reign! God will protect the Christian Empire.'' Then the Candidati closed round him and held their locked shields over his head. At this stage he must have retired into the palace where he put on the Imperial robes and the actual coronation was performed.\texthebrew{⁠}\footnote{This may be inferred from the order of proceedings in the case of the coronation of Anastasius.

} He came forth again bearing the diadem, and was adored by all the officials, in order of precedence. Then he took a shield and spear and was acclaimed anew. When the cries ceased, he replied, through the mouth of the magister a libellis ,\texthebrew{⁠}\footnote{\textgreek{\texthebrew{Ὁ} λιβελλήσιος.}} in the following words:

``Imperator\texthebrew{⁠}\footnote{\textgreek{Α\texthebrew{ὐ}τοκράτωρ.}} Caesar Leo, Victorious, Ever August (saith): Almighty God and your choice, most valiant fellow-soldiers ,

p316
elected me Emperor of the Roman State.'' All: ``Leo Augustus, thou conquerest. He who chose thee will keep thee. God will protect his choice.'' Leo: ``Ye shall have me as your master and ruler, who shared the toils which as your fellow-soldier I learned to bear with you.'' All: ``Our good fortune! The army accepts thee as Emperor, O conqueror. We all desire thee.'' Leo: ``I have decided what donatives I shall give to the troops.'' All: ``Pious and power ­ful and wise!'' Leo: ``To inaugurate my sacred and fortunate reign, I will give five nomismata [about £3] and a pound of silver to each shield.''\texthebrew{⁠}\footnote{\textgreek{Καταβουκο\texthebrew{ῦ}λον}, which should obviously be \textgreek{κατ\texthebrew{ὰ} βούκολον}. The \textgreek{βούκολον} was the centre of the clipeus. The Latin version mistranslates pro singulis buccis ``to each mouth.''

} All: ``Pious, lavish! Author of honour, author of riches! May thy reign be fortunate, a golden age!'' Leo: ``God be with us!'' Then a procession was formed, and the Emperor returned to the city where more ceremonies awaited him.\footnote{The description is taken from an evidently contemporary document preserved in Constantine Porph. \emph{De cerimoniis}, I. c91. There can, I think, be little doubt that Constantine found it in the ceremonial book (\textgreek{Κατάστασις}) compiled by Peter the Patrician in the sixth century, from which we know that he derived other accounts of early ceremonies (see \emph{ib.} cc84, 85). It is to be noted that the description of the actual ceremonies of A.D. 457 comes down only as far as the words \textgreek{κατ\texthebrew{ὰ} τάξιν}, p412, l. 18; the rest of the piece is generalised (in the present tense) so as to apply to any Emperor who is crowned in the Hebdomon Palace. It describes the return to the city, the halt at Hellenianae (near the Forum of Arcadius) and ceremonies there, a second halt at the Forum of Constantine, a third at St. Sophia, before the Great Palace is reached.

}

The danger which had threatened the Empire in the reign of Arcadius through the power of Gaïnas and his German faction was now repeated, though perhaps in a less openly mena­cing shape, and the interest and importance of Leo's reign lie in the struggle for ascendancy between the foreign and native powers in the State. To have averted this peril was Leo's one achievement. The position of Aspar, who, though an Alan and not a German, represented the German interest,\texthebrew{⁠}\footnote{His wife may have been an Ostrogoth, for Theoderic, son of Triarius, was her nephew (Theophanes, A.M. 5970).

} was extremely strong. He was Master of Soldiers in praesenti , and his son Ardaburius was, if not already, at least soon after Leo's accession, Master of Soldiers in the East.\texthebrew{⁠}\footnote{John Malalas, XIV p369. For the character of Ardaburius, who in time of peace devoted himself to frivolous amusements \texthebrew{—} actors, jugglers, and stage entertainments, \texthebrew{—} see Suidas, \emph{sub \textgreek{\texthebrew{Ἀ}ρδαβούριος}}, where Priscus may be the source (cp. \emph{F. H. G.} IV. p100).

} The Emperor, however, whom

p317
Aspar hoped to use as a puppet, soon showed that he had a will of his own and would not be as amenable to his general's dictation as he had led the general to expect. But, though differences arose\texthebrew{⁠}\footnote{Cp. Candidus, p135. Brooks (\emph{Zenon and the Isaurians}, 211\texthebrew{‑}212) gives reasons for dating the incident, referred to here, to 459.

} and Aspar was unable always to have his own way, yet for at least six or seven years his influence was predominant. Leo had made two promises, to raise Aspar's son Patricius to the rank of Caesar,\texthebrew{⁠}\footnote{The eastern consul in 459 was Patricius, but it is improbable that this was Aspar's son. We must rather identify him with Patricius, magister officiorum, to whom several undated laws of Leo are addressed (\emph{C. J.} XII. 19.9 ;

20.3\texthebrew{‑}5 ;

50.22 ) and who played a public part after Leo's death. Ardaburius was raised to the rank of patrician (Marcellinus, \emph{sub a.} 471), but the date is unknown. A third brother, Ermanaric, was perhaps consul in 465, as colleague of Leo's brother-in\texthebrew{‑}law Basiliscus. At that time Severus was Emperor in the West, and, as Leo did not recognise him, both consuls belonged to the eastern realm.

} thereby designating him as successor to the throne, and to give the Caesar one of his daughters in marriage.\texthebrew{⁠}\footnote{Of Leo's two daughters, Ariadne was born before, Leontia after, his accession. Brooks (\emph{ib.}) thinks that Ariadne must have been betrothed to Patricius, because Leontia was too young, and because a marriage with the younger daughter would not have had the same significance. But Leo might have preferred to promise the infant \texthebrew{—} many things might occur before she was ripe for marriage; and against the second objection might be set the fact that Leontia was born in the purple. We must also take into account that when Zeno married Ariadne we do not hear that Aspar complained that Leo had broken his promise. Leontia married Marcian (son of the western Emperor, Anthemius) whom we shall meet again;

Eustathius, apud Evagr. III.26 .

} The second arrangement could probably not be carried out immediately because the girl was too young, and Leo managed to postpone the fulfilment of the first. In the meantime he discovered a means of establishing a counterpoise to the excessive influence of the Germans.

In order to neutralise the fact on which Aspar's power rested, namely that the bulk and the flower of the army consisted of Germans and foreigners \texthebrew{—} who since the fall of the Hun Empire had begun again to offer themselves as recruits \texthebrew{—} he formed the plan of recruiting regiments from native subjects no less valiant and robust. He chose the hardy race of Isaurian mountaineers who lived almost like an independent people in the wild regions of Mount Taurus and were little touched by Hellenism. The execution of this policy, begun by himself and carried out by his successor, counteracted the danger that the Germans would prevail in the East as they were prevailing in the West.

p318
Leo had recourse to Tarasicodissa,\texthebrew{⁠}\footnote{Zeno's name is variously given as Tarasikodissa (Candidus, p135, who as an Isaurian should have known; cp. \textgreek{Στρακωδισσεων} in the MS. of John Malalas), Arikmesios (Eustathius of Epiphania, º

ap. Evagr. II.15 ), Traskalissaios (Theoph. A.M. 5974, perhaps an error for \textgreek{Τρασκωδισσα\texthebrew{ῖ}ον}, cp. Agathias, IV.29 , \textgreek{Ταρασικωδίσαιος}. He was a native of Rousoumblada in Isauria (Candidus, \emph{ib.}; Ramsay, \emph{Hist., Geography of Asia Minor}, p370). His mother's name was Lallis. Of his brother Longinus we shall hear much.

} an Isaurian chieftain, who came to Constantinople, and presently married his daughter Ariadne ( A.D. 466 or 467),\texthebrew{⁠}\footnote{Cp. Brooks, \emph{op. cit.} 212 , and Kulakovski, \emph{Ist. Viz.} I.352. Theophanes records the marriage under A.M. 5956 = A.D. 459, which is certainly wrong. 467 is the latest possible date, as Leo, son of Zeno and Ariadne, was six years old at the end of 474 (Michael Syr. IX. c5, ed. Chabot, vol. II p143.)

} having changed his uncouth name to Zeno. For about four years there was a struggle for ascendancy between the two factions. A new corps of Palace guards was formed, and we may conjecture that it was recruited from stalwart Isaurians, with the title of Excubitors.\texthebrew{⁠}\footnote{John Lydus, \emph{De mag.} I.16 . The number was 300.

} The Excubitors are for many centuries to be an important section of the residential troops, and, when we meet them for the first time in the reign of Leo, they were, as we shall see, called upon to oppose the Germans.

When a great expedition sailed to Africa against the Vandals in A.D. 468,\texthebrew{⁠}\footnote{See below,

p335 .

} Leo entrusted the command, not to Aspar or his son, but to Basiliscus, the brother of the Empress Verina. The commander's incompetence led to the failure of the enterprise. It was alleged, but the charge was probably false, that Aspar, sympathising with the Vandals, bribed Basiliscus to betray the fleet with the promise of making him Emperor.\texthebrew{⁠}\footnote{Hydatius, \emph{Chron.} 247 . According to this chronicler Aspar was consequently degraded from office and one of his sons put to death.

} In the following year Zeno was consul. It is possible that he had already been appointed Master of Soldiers in praesenti ,\texthebrew{⁠}\footnote{John Mal. XIV p375. This statement seems probably correct, for if Zeno was mag. mil. of the East he would have had no business to defend Thrace. The danger he ran with the Thracian army determined his transference to the eastern command. The statement of Theophanes (A.M. 5962) is certainly not decisive, but the date = A.D. 469, is probably right, and it seems probable that Zeno had been appointed to the East before the end of the same year. He continued to hold this post till the summer of 471 at least (see

\emph{C. J.} X.3.29 , and Brooks, \emph{op. cit.} 212, n17 ).

} and in this capacity he took the field in Thrace apparently against an incursion of Huns.\texthebrew{⁠}\footnote{The invasion of Huns under Attila's son Denzic is recorded in this year by Marcellinus. He was opposed by Anagast, mag. mil. in Thrace, and slain. \emph{Chron. Pasch.} records this under 468.

} Some of his soldiers, at the instigation of Aspar, conspired

p319
to assassinate him, but forewarned of the plot he escaped to Sardica. After this he was nominated Master of Soldiers in the East, and left Constantinople for Isauria, where he suppressed the brigand Indacus, one of the most dangerous and daring of the Isaurian bandits.\footnote{John Ant. \emph{fr.} 90 (\emph{Exc. de Ins.} p130), and Suidas, \emph{sub}\textgreek{\texthebrew{Ἰ}νδακός} (source Priscus?). The fortress of Indacus was Cherris.

}

It was probably during the absence of his son-in\texthebrew{‑}law in the East that Leo was at length induced by Aspar to perform his old promise of conferring the rank of Caesar upon his son Patricius ( A.D. 469\texthebrew{‑}470).\texthebrew{⁠}\footnote{Theophanes places this event in or before 468 (A.M. 5961), Victor Tonn. in 470, to which Brooks inclines. I agree, for Aspar would have been able to press Leo more effectively in Zeno's absence.

} Aspar is said (whether on this or some previous occasion) to have seized the sovran by his purple robe and said, ``Emperor, it is not fitting that he who wears this robe should speak falsely,'' and Leo to have replied, ``Nor yet is it fitting that he should be constrained and driven like a slave.''\texthebrew{⁠}\footnote{Zonaras, XIV.1 (p122 ed. B.-W.).

} There was great displeasure in Byzantium at the elevation of an Arian to a rank which was a recognised step to the Imperial throne. It appears that a deputation of clergy and laymen waited on the Emperor, imploring him to choose a Caesar who was orthodox, and the public dissatisfaction was expressed in the Hippodrome by a riotous protest, in which monks played a prominent part. Leo pacified the excited crowd by declaring that Patricius was about to turn from his Arianism and profess the true faith.\texthebrew{⁠}\footnote{See \emph{Vita Marcelli} in Simeon Metaphrastes, \emph{P.G.} 116.74. Marcellus, archimandrite of the Sleepless monks, led the protest. Theophanes says that Leo created Patricius Caesar \textgreek{δι\texthebrew{ὰ} τ\texthebrew{ὸ ἐ}λκύσαι τ\texthebrew{ὸ}ν \texthebrew{Ἄ}σπαρα \texthebrew{ἐ}κ τ\texthebrew{ῆ}ς \texthebrew{Ἀ}ρειανικ\texthebrew{ῆ}ς δόξης}. Possibly Aspar was converted.

} The new Caesar was soon afterwards betrothed to Leontia, the Emperor's younger daughter.

Meanwhile Anagast, a German soldier who had been appointed Master of Soldiers in Thrace, threatened to rebel. Messengers from the court persuaded him to desist from his enterprise, and he alleged that he had been instigated by Ardaburius, whose letters he sent to the Emperor as evidence.\texthebrew{⁠}\footnote{John Ant. \emph{ib.} (date, consul­ship of Jordanes = 470).

} Having failed in this attempt, Ardaburius endeavoured to gain over the Isaurian troops in Constantinople\texthebrew{⁠}\footnote{These Isaurians were reinforced by a body of their fellow-countrymen who had descended on the island of Rhodes. Many of these brigands had been cut down there, but the remnant escaped to Constantinople and were received by Zeno. Brooks dates this incident to 469 (\emph{op. cit.} 213 , and \emph{C. Med. H.} I.470).

} to his father's faction. These intrigues

p320
were betrayed to Zeno,\texthebrew{⁠}\footnote{Candidus, \emph{ib.} Zeno did not enter the city but remained at Chalcedon till after the murder, Theoph. A.M. 5964.

} who, if he was still in the East, must have hastened back to the capital ( A.D. 471). The destruction of Aspar and his family was now resolved upon. There was only too good reason to regard them as public enemies, but foul means were employed for their removal. Aspar and Ardaburius were slain in the palace by eunuchs;\texthebrew{⁠}\footnote{Marcellinus, \emph{sub} 471. The murder is branded by Damascius as treacherous (\textgreek{\texthebrew{ἐ}δολοφόνησεν}, \emph{Vita Isidori} in Photius, \emph{Bibliotheca}, 242, p340 ed. Bekker).

} the Caesar Patricius was wounded, but unexpectedly recovered; the third son Ermanaric happened to be absent and escaped.\texthebrew{⁠}\footnote{Candidus, \emph{ib.} Zeno is said to have assisted Ermanaric's escape to Isauria, where he married a daughter of an illegitimate son of Zeno. After Leo's death he returned to Constantinople (Theoph. \emph{ib.}).

} From this act the Emperor received the name of Butcher (Makelles). It was an important act in the long struggle against the German danger in the East. But it inaugurated a period of Isaurian domination which was to involve the Empire in a weary civil war. This was the price which had to be paid for the defeat of the German generals who sought to appropriate the Empire.

But the German danger was not yet quite stamped out. The Gothic friends of Aspar were dismayed, and they determined to avenge him. Count Ostrys,\texthebrew{⁠}\footnote{Doubtless the same as the \textgreek{στρατηγός} Ostryas mentioned by Priscus, \emph{fr.} 2 (\emph{De leg. gent.} p589), in connexion with the Hun invasion of 469. As \textgreek{στρατηγός} means mag. mil., it may be conjectured that Ostrys succeeded Zeno as mag. mil. in praesenti in that year (cp. above,

p318, n6 ).

} an officer of high rank who belonged to Aspar's faction, burst into the palace with an armed troop, but in an encounter with the new guards, the Excubitors, they were worsted. Ostrys fled to Thrace, taking with him Aspar's Gothic concubine. The Byzantine populace, with whom the power ­ful general, Arian as he was, probably had not been unpopular, cried, ``A dead man has no friend save Ostrys.''\texthebrew{⁠}\footnote{John Mal. XIV p371.

} The fugitive found a refuge in the camp of the Ostrogothic chief of German federate troops, Theoderic Strabo, Aspar's relative, who, as soon as he heard tidings of the murder, replied by ravaging Thrace. Whether he was deeply incensed or not, he saw an opportunity of stepping into Aspar's place, and when he made his peace with Leo in A.D. 473, he was appointed to the post of Master of Soldiers in praesenti , which Aspar had held. The career of Strabo will claim our attention later.\footnote{Below,

Chap. XII § 5 .

}

p321
At this time it was a common practice for rich people to maintain in their service not only armed slaves but bands of free retainers, often barbarians. It was natural enough that this practice should grow up in provinces which were exposed to hostile depredations, as in Illyricum and in those parts of Asia Minor which were constantly threatened by the Isaurian freebooters. But it is noteworthy, in view of Leo's Isaurian policy, that in his reign Isaurians were themselves hired or retained by private persons and that the Emperor found it necessary to forbid this dangerous usage.\footnote{\emph{C. J.} IX.12.10

omnibus per civitates et agros habendi bucellarios vel Isauros armatosque servos licentiam volumus esse praeclusam (A.D. 468). See above,

p43 .

}

Leo was a man of no education, but he seems to have possessed a good deal of natural good sense. The historian Malchus, who hated him for his religious bigotry, describes him as a sewer of wickedness and condemns his administration as ruinously rapacious.\texthebrew{⁠}\footnote{Malchus, \emph{fr.} 2\emph{a} (\emph{F. H. G.} IV. p114).

} This accusation is probably untrue and malicious. The financial methods of the Empire were so oppressive that the charge of rapacity might be brought against any Emperor, but Leo seems to have done nothing to make the system more rigorous, and to have followed in the steps of Marcian in adopting particular measures of relief and clemency as occasion offered.\texthebrew{⁠}\footnote{Antioch was laid in ruins by an earthquake in Sept. 458 (cp. Clinton, \emph{F.R., sub a.}). Leo rebuilt the public edifices (John Mal. XIV p369,

Evagrius, II.12 ).

} He is reported to have said that a king should distribute pity to those on whom he looks, as the sun distributes heat to those on whom he shines, and he may at least in some degree have practised what he preached. An anecdote suggests that he encouraged petitions. His unmarried sister, Euphemia, resided in a house in the south-eastern corner of the Augusteum, close to the Hippodrome. The Emperor used to pay her a visit with affectionate regularity every week. She erected a statue to him beside her house, and on its base petitioners used to place their memorials, which were collected every morning by one of the palace servants.\footnote{The statue was hence called the Pittakes (from \textgreek{πιττάκια}, letters). See \emph{Patria}, p167 (cp. 65).

}

One of the destructive conflagrations which have so often ravaged Constantinople occurred in A.D. 465 (September 2). The fire broke out close to the arsenal,\texthebrew{⁠}\footnote{Near the Gate of the Neorion, now Baghtsche Kapu.

} and it was said that it was

p322
caused by an old woman who was careless with her candle. Superstitious people believed that a malignant demon had assumed the shape of the old woman.\texthebrew{⁠}\footnote{See

Evagrius, II.13 , the chief description of the fire (probably derived from Priscus, through Eustathius). Also Theodore Lector, I.22; John Mal. XIV.372; Zonaras, XIV.1, 16.

} The fire spread eastward to the Acropolis, as far as the old temple of Apollo, and southward to the Forum of Constantine, whence it devastated the porticoes and buildings of Middle Street westward as far as the Forum of Taurus, and also pursued a southward course to the House of Amantius or Church of St. Thomas and to the Harbour of Julian.\texthebrew{⁠}\footnote{From here it spread to the Church of Homonoia of which the position is unknown.

} It lasted three days. The senate-house on the north side of the Forum of Constantine was destroyed,\texthebrew{⁠}\footnote{Cedrenus, I. p610

(source unknown), mentions the \textgreek{Σενάτον} and \textgreek{Νυμφα\texthebrew{ῖ}ον}. The building in the forum Tauri to which he refers may be the Basilica Theodosiana. The fire was described in the work of Candidus (cp. \emph{F. H. G.} IV. p135), and Shestakov has tried to show that the accounts in Cedrenus and Zonaras are derived from him (\emph{Kandid Isauriski}, cp. Bibl. II.2, 13).

} and the Nymphaeum directly opposite to it, a building in which those who had not large enough houses of their own used to celebrate their weddings. Many magnificent private residences were burned down. It is said that Aspar ran about the streets with a pail of water on his shoulders, urging all to follow his example and offering silver coins to encourage them. There is no hint of the existence of a fire-brigade .\texthebrew{⁠}\footnote{In later times we hear of regular arrangements for extinguishing fires. See Michael, \emph{Vita Theodori Studitae} in \emph{P.G.} 99, p312 \textgreek{τ\texthebrew{ὴ}ν τ\texthebrew{ῶ}ν σιφώνων κατ\texthebrew{ὰ} τόπους παρασκευήν.}} The Emperor, alarmed by the disaster, withdrew across the Golden Horn to the palace of St. Mamas and remained there for six months.

In his ecclesiastical policy Leo followed Marcian and faithfully maintained orthodoxy as established by the Council of Chalcedon. No memorable feat of arms distinguished his reign\texthebrew{⁠}\footnote{Early in his reign a barbarian people, which invaded Pontus, was repelled and subjugated. We hear of this in a letter of bishops of Pontus to the Emperor (Mansi, VII. p600), and the same event seems to be referred to in other letters (\emph{ib.} 581, 583). Tillemont thought these barbarians must be Huns (\emph{Hist. des Empereurs}, VI.367), but it seems to me more probable that they were the Tzani (for whom see below,

p434 ).

} to counterbalance the disastrous issue of his ambitious expedition against the Vandals, which will be recounted in another place. The Illyrian peninsula was troubled by the restlessness of the Ostrogoths, but the brunt of their hostilities was to be borne by Leo's successor. He died on February 3, A.D. 474, having co-opted as

p323
Augustus (in October) his grandson Leo, an infant aged about six years.\footnote{For dates see John Mal. XIV p375. The details of the coronation of Leo II are preserved in a contemporary account which was probably included in a work of Peter the Patrician

(Const. Porph. \emph{De cer.} I. c94) . After the coronation of the child the two Leos would be distinguished as \textgreek{Λέων \texthebrew{ὁ} μέγας} and \textgreek{Λέων \texthebrew{ὁ} μικρός}, and this, I believe, must be the origin of the designation of Leo as ``the Great''; just as reversely Theodosius II was called ``the Small,'' because in his infancy he had been known as \textgreek{\texthebrew{ὁ} μικρ\texthebrew{ὸ}ς βασιλεύς} to distinguish him from Arcadius (see above,

p7 ). Leo never did anything which could conceivably earn him the title of Great in the sense in which it was bestowed by posterity on Alexander or Constantine. \texthebrew{—} Coins issued at the beginning of Leo's reign show Marcian's head, the legend being merely altered to Dn Leo Perpet Aug; and Pulcheria's coin stamp with Victoria Augg (see above,

p236 ) was used for Verina (Ael Verina Aug). Later coins of Leo have his portrait, a bearded man; but his face is seen better on a medallion (Sabatier, \emph{Pl.} VII.1), for which an old stamp of the seventh quinquennalia of Theodosius II was used (as the unaltered reverse shows).

}

If it was a critical moment at Constantinople at the death of Marcian, it had been a still more critical moment in Italy on the death of Valentinian III two years before ( A.D. 455). For not only was there no male heir of the house of Theodosius, but there was no minister or general of commanding influence, no Aetius or Aspar, to force a decision. Military riots were inevitable, a civil war was possible; and we read that ``Rome was in a state of disturbance and confusion, and the military forces were divided into two factions, one wishing to elevate Maximus, the other supporting Maximian (son of an Egyptian merchant) who had been the steward of Aetius.''\texthebrew{⁠}\footnote{John Ant. \emph{fr.} 85 (\emph{De ins.} p127). His account deserves credit because he drew his information from the contemporary historian Priscus.

} A third possible candidate was Majorian, brother-in\texthebrew{‑}arms of Aetius, with whom he had fought against the Franks,\texthebrew{⁠}\footnote{Aetius, however, dismissed him from whatever post he held. Sidonius attributes this injustice to the influence of the wife of Aetius who was jealous of Majorian's growing fame
(\emph{Carm.} V.126\texthebrew{‑}294) . Majorian retired to his country estate, but was recalled to military service after the Patrician's death, by Valentinian (\emph{ib.} 305\texthebrew{‑}308).

} and he had the good wishes of Eudoxia, the widowed Empress. If there had been time to consult the Emperor Marcian, we may conjecture that his influence would have been thrown into the scale for Majorian. But the money of Petronius Maximus\texthebrew{⁠}\footnote{The wealth of Maximus is noted by Sidonius Apollinaris, \emph{Epp.} II.13.

} decided the event in his favour, just as Pertinax had won the Empire after the death of Commodus by

p324
bribing the Praetorian guards. He was elevated to the throne on March 17, A.D. 455.

Maximus endeavoured to strengthen himself on the throne by forcing Eudoxia to marry him, and if she had yielded willingly, it is possible that the Italians might have rallied round him and he might have reigned securely. But though he was a member of the noble Anician house, he was not like Marcian; he was not one whom the Augusta could bring herself to tolerate even for cogent political reasons. If he was really related to the British tyrant Maximus, who had been subdued by Theodosius, the great-granddaughter of Theodosius had perhaps not forgotten the connexion; but the widow of Valentinian must have known or suspected the instigator of her lord's murder.\texthebrew{⁠}\footnote{The favour he showed to the assassins is recorded by Prosper, \emph{sub a.}

} In any case, the new Augustus was so hated and despised by Eudoxia that she was said to have taken the bold and fatal step of summoning Gaiseric the Vandal to overthrow the tyrant. There was indeed a particular reason for asking aid from Carthage, instead of appealing, as one might have expected her to do, to Constantinople. Maximus had not only forced her to wed him, but he also forced her daughter Eudocia to give her hand to his son Palladius whom he created Caesar. And Eudocia was the affianced wife of Huneric, the heir to the Vandal throne. The act of Maximus touched the honour of Gaiseric, and he would be likely to come to the rescue more promptly than Marcian. The story, therefore, of the appeal of the Empress to the Vandal is credible, though it is not certainly true.\footnote{Little is said of it by western writers ( Hydatius, 167 , refers to it as an evil rumour). The sources are John Ant. \emph{loc. cit.}; Marcellinus, \emph{Chron., sub a.}, Procopius, \emph{B. V.} I.4;

Evagrius, \emph{H. E.} II.7 . All these accounts are probably derived from Priscus, but it is evident from John Ant. (cp. \textgreek{ο\texthebrew{ἱ} δέ φασι}) that Priscus did not tell the story as definitely true, but admitted the possibility that Gaiseric might have come of his own accord. The part played by the mysterious Burgundio in
Sidonius, \emph{Carm.} VII.441 \emph{sqq.} , is not clear. For Palladius see Prosper, \emph{ib.}

}

Petronius Maximus enjoyed the sweets of power for two months and a half, but he found them far from sweet. The man who as a private individual was so great a figure, ``once made emperor and prisoned in the palace walls, was rueing his own success before the first evening fell.'' Formerly he used to live by the clock, but now he had to renounce his old regular life and his ``senatorial ease.'' His rule was ``from

p325
the first tempestuous, with popular tumults, tumults of soldiery, tumults of allies.'' An influential nobleman, who was often with him, used to hear him exclaim, ``Happy thou, O Damocles, whose royal duresse did not outlast a single banquet!''\footnote{Sidonius Apoll. \emph{Epp.} II.13 (Dalton).

}

In May it was known in Italy that Gaiseric had set sail. There was consternation at Rome, and a considerable exodus both of the higher and the lower classes. Maximus, when he heard that the Vandals had landed, thought only of flight. He was deserted by his bodyguard and all his friends, and as he was riding out of the city, some one cast a stone and hit him on the temple. The stroke killed him on the spot and the crowd tore his body limb from limb (May 31).\footnote{His end is described by Prosper, \emph{ib.}, and with more detail by John Ant. \emph{ib.};

Jordanes, \emph{Get.} 235 , names Ursus, a Roman soldier, as the assassin.

}

Three days later\texthebrew{⁠}\footnote{Victor Tonn. \emph{sub a.}

} Gaiseric and his Vandals entered Rome. Whether they came entirely of their own accord or in answer to a summons from the Empress, they were now bent only on rapine. The bishop of Rome, Leo I, met them at the gates. Although he did not succeed in protecting the city against pillage, violence, and ``vandalism,'' he preserved it by his intervention from the evils of massacre and conflagration. For fourteen days the enemy abode in the city, and plundered it coolly and methodically.\texthebrew{⁠}\footnote{Secura et libera scrutatione (Prosper).

} The palace on the Palatine was ransacked thoroughly. Precious works of art were carried off, and many of the gilt bronze tiles which roofed the temple of Jupiter Capitolinus were removed. The robbers added to their booty the golden treasures which Titus had taken from the temple of Jerusalem. When they had rifled the public and private wealth of Rome, and loaded their ships, they returned to Africa with many thousand captives, including the Empress Eudoxia and her two daughters, Eudocia and Placidia.\texthebrew{⁠}\footnote{For the sack of Rome see, besides Prosper, Procopius, \emph{B. V.} I.5 (he mentions that a ship laden with statues was lost on the way to Carthage). Gaudentius, son of Aetius, was one of the captives. Cp. Grisar, \emph{Hist. of Rome and the Popes} (Eng. tr.), I.96\texthebrew{‑}99; Martroye, \emph{Genséric}, 158 \emph{sqq.}

} It will be remembered that the idea of an alliance between Gaiseric's heir and a daughter of Valentinian had been suggested by Aetius. This plan was now carried out. Huneric married

p326
Eudocia. Her sister Placidia was already the wife of a distinguished Roman, Olybrius.\footnote{Priscus, \emph{fr.} 10, \emph{De leg. Rom.}; Procopius, \emph{B. V.} I.4.

Evagrius, \emph{H. E.} II.7

blunders. See Clinton, \emph{F. R.} II p127.

}

But the question was, who was to be Emperor? Rome was paralysed by the shock of the Vandal visitation, but Gaul intervened. Marcus Maecilius Flavius Eparchius Avitus, the man who had fought by the side of Aetius and at a great crisis had decided Theoderic the Visigothic king to march against the Huns, had been appointed by Maximus Master of Both Services in Gaul. It was important for the new Emperor to establish a friendly understanding with the Visigothic ruler, and no one was more fitted to bring this about than Avitus, the intimate friend of Theoderic I, and no less a \emph{persona grata} to Theoderic II. He was, in fact, at Tolosa when the news of the death of Maximus arrived, and Theoderic persuaded him that he was the necessary man.\texthebrew{⁠}\footnote{Sidonius, \emph{Carm.} VII.517 tibi pareat orbis, ni pereat.

} He was proclaimed Emperor by the Goths at Tolosa (July 9, or 10); five weeks later his assumption of the Imperial power was confirmed at a meeting of representative Gallo-Romans at Ugernum (Beaucaire), and he was formally invested at Arles with Imperial insignia.\footnote{\emph{Ib.} 522\texthebrew{‑}600. For dates see \emph{Fast. Vind. pr.} p304, and Victor Tonn. \emph{sub a.} We have a portrait of Avitus on gold coins which show his side face, bearded; on the reverse he is trampling on a captive.

}

Towards the end of the year Avitus crossed the Alps to assert his authority in Italy and assume the consul ­ship for A.D. 456. He was accompanied by a famous man of letters who was his son-in\texthebrew{‑}law, Caius Sollius Apollinaris Sidonius, son and grandson of Praetorian Prefects of Gaul.\texthebrew{⁠}\footnote{The father was Prefect in 448; the grandfather in 408.

} Sidonius had been born and educated at Lyons, and was now about twenty-five years of age. For a quarter of a century he was to play a considerable part in the relations between Gaul and Italy as well as in the internal affairs of Gaul. The poetical panegyric which he recited at Rome in honour of his father-in\texthebrew{‑}law 's consul ­ship\texthebrew{⁠}\footnote{\emph{Carm.} VII .

} marks the beginning of his public career; his statue was set up in the Forum of Trajan. But the Emperor Avitus, who was so much at home at Tolosa, was not welcome at Rome, though he was acknowledged by Marcian. He was acceptable neither to the soldiers nor to the Senate, and his

p327
behaviour did not tend to make him popular, although his reign was distinguished by military successes by land and sea.

Both the Vandals and the Suevians had been alert to take advantage of the difficulties which followed Valentinian's death. Gaiseric had been extending his authority over those African provinces which had been left to Rome by the treaty of A.D. 442. The Emperor Marcian had sent an embassy to remonstrate with him on the sack of Rome and the captivity of the Imperial ladies; Avitus sent him an a message warning him to observe the treaty. But Gaiseric was inflexibly hostile; he defied both Marcian and Avitus; and he sent a fleet of sixty ships to descend on Italy or Gaul. The general Ricimer, destined to be the leading figure in the West for about sixteen years, now makes his appearance on the scene. His mother was a daughter of the Visigothic king Wallia, and his father was a Sueve; he had risen in Roman service, and Avitus appointed him Master of Soldiers.\texthebrew{⁠}\footnote{Sidonius, \emph{Carm.} II.361 \emph{sqq.}

} He now went to Sicily with an army and a fleet; a Vandal descent on that island was evidently expected, and was apparently attempted in the neighbourhood of Agrigentum. The enemy was forced to retreat, but Ricimer followed them and gained a naval victory in Corsican waters ( A.D. 456).\footnote{We have to combine

Hydatius, 177 , with
Sidonius, \emph{ib.} 367 , and Priscus, \emph{fr.} 7, \emph{De leg. Rom.}

}

Theoderic II, who seems to have been chiefly responsible for the elevation of Avitus, had won the Gothic throne by murdering his brother Thorismund ( A.D. 453).\texthebrew{⁠}\footnote{There is an interesting account of Theoderic and the daily routine of his life in Sidonius, \emph{Epp.} I.2.

} He now showed his good will to the new Emperor by marching into Spain and making war upon the Suevians, who were perpetually harrying Roman provinces. But, though he went in the name of Avitus and the Roman Republic, we cannot doubt that he was deliberately preparing for the eventual fulfilment of the ambition of the Goths to possess Spain themselves, by weakening the Suevic power. The king of the Suevians, Rechiar, was his brother-in\texthebrew{‑}law , and to him Theoderic sent ambassadors calling upon him to desist from his raids into Roman territory. Rechiar defied him and invaded Tarraconensis, whereupon Theoderic led a host of Goths, reinforced by Burgundians, into Gallaecia, and defeated the Suevians in a battle on the river Urbicus, near Astorga (October 5, A.D. 456). The victor pushed on

p328
to Bracara, which he captured three weeks later, and his barbarous army committed all the acts of violence and rapine usual in sacks, short of massacre and rape. Sometime later Rechiar, who had fled, was captured at Portuscale (Oporto) and paid with his life for defying his brother-in\texthebrew{‑}law .\texthebrew{⁠}\footnote{Hydatius, 170\texthebrew{‑}175 .

} The battle of the Urbicus was an important event, for it shattered the power of the Suevians. Their kingdom indeed survived for 120 years, but it never recovered its old strength.

The crushing victory won by his German allies in Spain did not avail Avitus. Before the great battle was fought he had left Rome, virtually as a fugitive, on his way to Gaul, and was probably already a prisoner. The circumstances which led to his fall are thus related:\texthebrew{⁠}\footnote{By John Ant. \emph{fr.} 86, \emph{De ins.}

} \texthebrew{—}He was captured at Placentia by Ricimer and Majorian. He was deposed from the throne and elected bishop of the city which witnessed his discomfiture (October 17 or 18, A.D. 456), but died soon afterwards.\footnote{John Ant. \emph{ib.} says he was starved or strangled. There is a different story in

Gregory of Tours, \emph{Hist. Franc.} II.11 . For the date cp. \emph{Cons. Ital.} p304. Avitus had armed men with him, and there was a battle at Placentia in which ``his patrician'' Messianus was slain, \emph{ib.}

}

A new Emperor was not immediately elected. A temporary cessation of a separate Imperial rule in the West occurred on several occasions during the twenty years which followed the death of Valentinian. One of these intervals occurred now. They are often called interregnums; it is natural to say that from October A.D. 456 to April A.D. 457 there was an interregnum

p329
in the West, and the expression represents the actual situation. But we must not forget that in theory the phrase is incorrect. Legally, Marcian was the sole head of the Empire from the fall of Avitus to his own death at the end of January, and Leo was the sole head of the Empire for three months after the death of Marcian.\footnote{Coins of Marcian minted in Italy belong to this interval, and those of Leo to the longer period between the death of Majorian and the accession of Anthemius (461\texthebrew{‑}467). Cp. de Salis, \emph{Coins of the Eudoxias}, p215, who holds that the custom of striking coins at Italian mints in honour of the eastern colleague ceased at the beginning of the reign of Valentinian III.

}

The Master of Soldiers, Ricimer, whose prestige had been established by his naval victory, now held the destinies of Italy in his hands. He had succeeded to the post and the responsibilities of Stilicho, Constantius, and Aetius, but his task was vastly more difficult. For while those defenders of the Empire against the German enemies were supported by the secure existence of an established dynasty, Ricimer had to set up Emperors in whose name he could act. At the beginning of A.D. 457 the situations in Italy and at Constantinople were similar. In both cases the solution of the difficulty depended on the action of a military leader of barbarian birth; Aspar's position was as that of Ricimer. Both were the makers of Emperors, neither could aspire to be an Emperor himself. They were Arians as well as barbarians.\footnote{We have an inscription of Ricimer recording that he decorated with mosaics the Arian church of S. Agatha in Rome in accordance with a vow. Its date is later than 459, the year of his consul­ship. De Rossi, II.1, p438; Dessau, 1294.

}

The legitimacy of any Emperor set up in Italy depended on his being recognised as a colleague by the Emperor reigning at Constantinople. Avitus had been recognised by Marcian, and if the seat of his successor was to be firmly established it was indispensable that he should obtain similar recognition. The political importance of conforming to this constitutional necessity was realised by Ricimer, and we may confidently assume that after the fall of Avitus, he, acting probably through the Roman Senate, communicated with the Emperor of the East. Marcian's death postponed a settlement, but one of the early acts of Leo I was to nominate a colleague. That the suggestion of Majorian's name came from Rome we can hardly doubt. Julius Valerianus Majorianus was a thorough Roman and on that account most acceptable to the Senators. He had been, we saw, the candidate

p330
of Eudoxia after her husband's death. He was elevated to the throne on April 1, A.D. 457.\texthebrew{⁠}\footnote{He had been created magister militum in February (\emph{Fast. Vind. Pr.} p305). His address to the Senate (\emph{Nov.} 1 de ortu imperii divi Maioriani) announces the inauguration of a new era. Ricimer is thus mentioned: erit apud nos cum parente patricioque nostro Ricimere rei militaris pervigil cura.

} At the same time Leo conferred upon Ricimer the title of Patrician.

There were two tasks for the new Augustus to accomplish if he was to make his seat on the throne secure and exercise effective rule in the west. He had, in the first place, to quell the opposition in Gaul. The fall of Avitus had aroused the wrath both of his barbarian friends, Visigoths and Burgundians, and of the provincials. Gallic Avitus had failed to conciliate Italian goodwill; it was now to be seen whether Italian Majorian would succeed in solving the reverse problem. There was little love lost between the Romans and the trans-Alpine provincials, and there was now a serious danger, such as had often occurred before, that Gaul would attempt to dissociate itself politically from Italy, and have an Emperor to itself.

There are indeed signs of a gradually widening rift between Gaul and the rest of the Empire ever since the time of the tyrants in the reign of Honorius. It has been observed\texthebrew{⁠}\footnote{By Sundwall (\emph{Weströmische Studien}, p8), who has insisted rightly on the importance of the struggle for power between Gaul and Italy.

} that of the twenty-eight Praetorian Prefects of Gaul in the fifth century whose names are recorded, we know that eighteen were Gauls, and of the other ten none is known to be of Italian birth. This points to the conclusion that the feeling in Gaul was such that the central government considered it impolitic to appoint any one to that post outside the circle of Gallic senators. The loss of Africa probably accentuated the sectional feeling in both Italy and Gaul, and from this point of view the elevation of Avitus was a momentarily success ­ful attempt of the Gallic nobility to wrest from the Italians the political predominance which had hitherto been theirs. It was the business of Majorian to preserve for Italy her leading position and at the same time to conciliate the Gallic nobility.

Majorian entered Gaul with an army composed mainly of German mercenaries, and found the Burgundians in league with the inhabitants of Lugdunensis Prima against himself.\texthebrew{⁠}\footnote{\emph{Cons. Ital.} (Auct. Prosper, Havn.), \emph{sub} 457; Marius Avent. \emph{sub} 456.

} Lyons,

p331
which had received a Burgundian garrison, was compelled to surrender and was punished for its rebellion by the imposition of heavier taxation. This burden, however, was soon remitted, through the efforts of Sidonius Apollinaris, who delivered an enthusiastic Panegyric at Lyons on the man who had helped to dethrone his father-in\texthebrew{‑}law .\texthebrew{⁠}\footnote{\emph{Carm.} IV , V .

} The Visigoths were besieging Arelate, but Majorian's general, Aegidius, drove back Theoderic from its walls and firm compacts were made between the two potentates.\texthebrew{⁠}\footnote{Hydatius, 197 , A.D. 459. Majorian was in Gaul in 458\texthebrew{‑}459.

} The Burgundians were allowed peacefully to possess the province of Lugdunensis Prima.\texthebrew{⁠}\footnote{Not including Lyons.

} Honours were freely distributed to the Gallic nobility.

Majorian had accomplished one task; the other was more difficult. It was indispensable for an Emperor, who had not the prestige of belonging to a dynasty, to win general confidence by proving himself equal to the great emergency of the time; he must ``preserve the state of the Roman world.''\texthebrew{⁠}\footnote{Majorian, \emph{Nov.} 1 Romani orbis statum . . . propitia divinitate servemus.

} The deliverance of Arelate was a good beginning. But the great emergency was the hostility of the Vandals who in their ships harried the Roman provinces and infested the Mediterranean waters. The defeats which Ricimer had inflicted on their fleet at Corsica did not paralyse their hostilities. The words of an historian indicate that Avitus in fa­cing this danger had felt his inability to grapple with it: ``He was afraid of the wars with the Vandals.''\footnote{See Priscus, \emph{fr.} 13 (\emph{De leg. gent.} p585), who is almost verbally followed by John Ant. \emph{fr.} 87, p203.

}

Majorian prepared an expedition against Africa on a grand scale; his fleet numbered 300 ships and was collected off the coast of Spain. The hopes of all his subjects were awakened and their eyes fixed on his preparations. But a curious fatality attended all expeditions undertaken against the Vandals, whether they proceeded from Old Rome or from New Rome, or from both together. The expedition of Castinus had collapsed in A.D. 422, that of Aspar had failed in A.D. 431, the armament of Ardaburius did not even reach its destination in A.D. 441, and the expedition of Majorian came to naught in A.D. 460. Gaiseric ravaged the coasts of Spain and many of the Roman warships

p332
were surprised and captured in the bay of Alicante.\texthebrew{⁠}\footnote{Marius Avent. \emph{sub a.} 460 (where Elice = Alicante is misleadingly described as near New Carthage). Cp. Hydatius, 200 , and see Martroye, \emph{Genséric}, p192. Majorian made a ``disgraceful treaty'' with Gaiseric, John Ant. \emph{ib.} Probably he ceded the Roman provinces in Africa (the Mauretanias and Tripolitana) which Gaiseric had recently seized.

} Yet another expedition, and one on a grand scale, was soon to be fitted out and also to meet with discomfiture; and more than seventy years were to elapse until the numerous failures were to be retrieved by the victories of Justinian and Belisarius.

This misfortune led to the fall of Majorian. He returned from Spain to Gaul, and after a sojourn at Arles\texthebrew{⁠}\footnote{He celebrated games at which Sidonius Apollinaris was present (\emph{Epp.} I.11).

} passed into Italy without an army. In Italy, and at Rome, he was probably popular;\texthebrew{⁠}\footnote{During A.D. 458 Majorian had attempted much remedial legislation. He alleviated the public burdens by a remission of arrears (\emph{Nov.} 2) and resuscitated the office of defensor civitatis (\emph{Nov.} 3). He enacted a much-needed law for preserving the public buildings of Rome, to check the ``disfigurement of the face of the venerable city'' (\emph{Nov.} 4). He also endeavoured to deal with the social evil of celibacy (\emph{Nov.} 6).

} but now that he had proved himself unable to ``preserve the state of the Roman world,'' Ricimer, who was thoroughly dissatisfied with him, could venture to take action against him. At Tortona Majorian was seized by Ricimer's officers, stripped of the purple, and beheaded (August 2, A.D. 461).\texthebrew{⁠}\footnote{\emph{Fasti Vind. pr., sub a.} 464.

} He had done at once too little and too much. An Emperor who was just strong enough to act with independent authority, but not strong enough to contend with the enemies of the State, was useless to Ricimer, who himself seemed resolved not to leave Italy, probably judging that the constant presence of a capable general with considerable forces was necessary against descents of the Vandals. There were other enemies too against whom he had to defend it. He had to fight against the Ostrogoths of Pannonia, and to repel an invasion of Alans. But the great foe was Gaiseric, who hated him as the grandson of King Wallia.

Nearly three and a half months passed before Majorian was succeeded by Libius Severus, a Lucanian, who was elected by the Senate at the instance of Ricimer and proclaimed at Ravenna (November 19, A.D. 461). He was not recognised at Constantinople.

p333
He reigned as a figurehead; Ricimer was the actual ruler.\footnote{His monogram appears on the reverse of coins of Severus. There were a good many issues of coins during this reign. Perhaps one of the earliest of the solidi of Severus was that with the same reverse type which appears on solidi of Petronius Maximus and Majorian \texthebrew{—} an Emperor holding a cross and a globe surmounted by a Victory, with his right foot on a dragon's head.\texthebrew{—} A bronze weight with an inscription of Plotinus Eustathius, Prefect of Rome, may belong to the reign of Severus (CIL X.8072 ). It illustrates the position of Ricimer, whose name is associated with the Emperors: salvis dd. nn. et patricio Ricimere. On another tablet, of the Praet. Prefect Probianus, his name does not appear with those of Leo and Severus (Dessau, 811).

}

It might seem that at this juncture Italy might have received another Augustus from Gaul, and that Aegidius, Master of Both Services in Gaul\texthebrew{⁠}\footnote{Hydatius, 218 . It was doubtless Maximus who first conferred this higher rank and title on a Gallic commander (Avitus), hitherto a magister equitum. The change illustrates the political importance of Gaul at this time. See above,

p326 .

} and friend of Majorian, might have crossed the Alps to avenge his death. Aegidius acknowledged no allegiance to Ricimer's Emperor,\texthebrew{⁠}\footnote{Priscus, \emph{fr.} 14, \emph{De leg. gent.}

} but he was fully occupied with the defence of the Gallic provinces against the Visigoths, who were attempting to extend their power northward and eastward. We find him winning a battle at Orleans in A.D. 463,\texthebrew{⁠}\footnote{Cp. Hydatius, \emph{ib.}; Marius Avent. \emph{sub a.} Aegidius defeated Frederic, brother of King Theoderic, near Orleans. Before his death he was negotiating with Gaiseric, the plan being that the Vandals should attack Ricimer in Italy while Aegidius was making war on the Visigoths.

} and in the following year he died.

Ricimer had an opponent in another quarter, the count Marcellinus. In A.D. 461 this general was in Sicily, in command of an army chiefly consisting of Hun auxiliaries; he had probably been posted there by Majorian to protect the island against the Vandals. But the bribes of Ricimer prevailed upon the cupidity of the Huns and induced them to leave the service of Marcellinus and enter his own. Then Marcellinus, conscious that he could not vie with Ricimer in riches, went to Dalmatia, where he ruled under the authority of Leo, and perhaps with the title of Master of Soldiers in Dalmatia.\texthebrew{⁠}\footnote{Marcellinus had been a friend of Aetius and after his murder had withdrawn to Dalmatia. The Gallo-Romans offered him the Imperial crown in 458 before they accepted Majorian. (Procopius, \emph{B. V.} I.6; Sidonius, \emph{Epp.} I.11.6.) Damascius in his \emph{Vita Isidori} (Photius, \emph{Bibl.} 242 p342) describes him as \textgreek{α\texthebrew{ὐ}τοδέσποτος \texthebrew{ἡ}γεμών} of Dalmatia. I conjecture that his title was magister militum Dalmatiae, because after his death his nephew Julius Nepos held this exceptional title: see

\emph{C. J.} VI.61.5

(A.D. 473).

} On his departure Sicily was ravaged by the Vandals and Moors, and a pacific embassy from

p334
Ricimer had no effect. But another embassy sent at the same time by the Emperor Leo induced Gaiseric to come to terms at last in regard to the ladies of the Theodosian house, whose deliverance from their captivity in Carthage Marcian had vainly endeavoured to secure. Eudocia, the bride of Huneric, was retained, but her mother Eudoxia and her sister Placidia were sent to Constantinople. In return, Gaiseric bargained for a certain share of the property of Valentinian III as the dowry of Eudocia.\texthebrew{⁠}\footnote{A.D. 462.

Hydatius, 216 . Priscus, \emph{fr.} 10 (in \emph{De leg. Rom.}).

} He had already occupied and annexed the Mauretanian provinces, as well as Sardinia, Corsica, and the Balearic islands.

This concession had its definite political purpose which was soon revealed. The Vandal monarch now came forward as the champion of the Theodosian house against Ricimer and his upstart Emperor. Placidia had married Olybrius, a member of the noble Anician gens, and Gaiseric demanded that Olybrius should succeed to the throne in Italy. Threatened on one hand by the Vandals, on the other by the ruler of Dalmatia, Ricimer and the obedient Senate solicited the good offices of Leo. He was asked to bring about a reconciliation with Gaiseric and with Marcellinus. Leo consented. One envoy prevailed on Marcellinus not to wage war against the Romans, the other returned from Carthage without result. Gaiseric claimed in his daughter-in\texthebrew{‑}law 's name all the private property possessed by her father in Italy, and also the inheritance of Aetius, whose son Gaudentius he retained a prisoner. In pursuance of these claims he led a great expedition against Italy and Sicily, ravaged the country districts and undefended towns. There was no efficient navy to oppose him at sea.

The elevation of Olybrius, which would have been a restitution of the Theodosian dynasty, might have seemed a hopeful solution of some of the difficulties of the situation, but the fact that he was Gaiseric's candidate and relative was a reason against accepting him. For a year and eight months after the death of Severus (August 15, A.D. 465),\texthebrew{⁠}\footnote{\emph{Fast. Vind. pr., sub} 464.

According to the Chronicle of Cassiodorus , ut dicitur, Ricimeris fraude Severus Romae in palatio veneno peremptus est. If this is true, Ricimer had a hand in the death of no fewer than three, if not four, Emperors.

} no successor was appointed. Then Gaiseric made a raid on the Peloponnesus ( A.D. 467) and Leo determined to take decisive steps and act in close conjunction

p335
with the Italian government. Now that not only Italy and Sicily were threatened, but the entire commerce of the Mediterranean, the forces of the east were to be united with those of Italy and Dalmatia against the African foe. The first step was to find a suitable man to invest with Imperial authority in the west. The choice of Leo fell on the patrician Anthemius, who, as the son-in\texthebrew{‑}law of the Emperor Marcian, might be considered in some sort a representative of the house of Theodosius, and his pretensions might be set against those of Gaiseric's candidate, the husband of Placidia. The support of Ricimer was secured by an arrangement that he should marry the daughter of Anthemius. The elder Placidia had married Athaulf, her granddaughter Eudocia had married Huneric, both indeed under a certain compulsion; yet Anthemius afterwards professed to regard it as a great condescension to have given his daughter to the barbarian general. He arrived in Italy and was proclaimed Emperor near Rome on April 12, A.D. 467.\footnote{\emph{Cons. Ital., sub a} p305, Cassiodorus, \emph{Chron., sub a.} He was not created Emperor or crowned until he arrived in Italy, for he sent Heliocrates to Constantinople to announce his elevation and obtain formal recognition. Leo sent his image bound with bay leaves (\textgreek{τ\texthebrew{ὰ} λαυρεάτα}) to the cities of the east with a command that it should be honoured like his own, ``that all the cities may learn with joy that the powers of both sections of the Empire are united'' (Peter Patricius, in Constantine Porph. \emph{De cer.} I.87, where the ceremony of the reception of the ambassador of Anthemius is described).

}

The expedition which was organised to overthrow the kingdom of the Vandals was on a grand and impressive scale, but in ended in miserable failure, due to lukewarmness and even treachery both in the east and in the west.

The number of vessels that set sail from Constantinople ( A.D. 468) is said to have been 1113, and the total number of men who embarked was calculated as exceeding 100,000. But unfortunately Leo, under the influence of his wife Verina and his friend Aspar, appointed as general a man who was both incompetent and untrustworthy, his wife's brother Basiliscus. Aspar, it is said, was not over-anxious that Leo's position should be strengthened by such an exploit as the subversion of the Vandal kingdom; he schemed therefore to procure the election of a general whose success was extremely improbable.\texthebrew{⁠}\footnote{Compare

Hydatius, 247

Asparem degradatum ad privatam vitam filiumque eius occisum adversus Romanum imperium, sicut detectique sunt, Vandalis consulentes.

} The western armament obeyed a more competent commander. Marcellinus

p336
assumed the direction of the Italian fleet.\texthebrew{⁠}\footnote{Marcellinus, \emph{Chron., sub a.} 468, where it is mentioned that Marcellinus was a pagan. He had received from Leo (we may assume) the title of Patrician.

} But his participation in the enterprise alienated Ricimer, who was his personal enemy, and who seems to have been jealous of Anthemius already.

The plan of operations was that the eastern forces should be divided into two parts, and that the Vandals should be attacked at three points at the same time. Basiliscus himself was to sail directly against Carthage. Heraclius, another general, having taken up the forces of Egypt on his way, was to disembark in Tripolitana, and to march to Carthage by land. Marcellinus, with the Italian forces, was to surprise the Vandals in Sardinia, and sail thence to join the eastern armies at Carthage.

If the commander-in\texthebrew{‑}chief had not been Basiliscus, and if the opponent had not been Gaiseric, the expedition might easily have succeeded. But Gaiseric, though physically the least, was mentally the greatest of the barbarians of his time. Even as it was, though Basiliscus had such a foe to cope with, success was within the grasp of his hand. The invaders were welcome to the Catholics of Africa, who were persecuted by their Arian lords. Marcellinus accomplished his work in Sardinia without difficulty; Heraclius met no obstacle in executing his part of the scheme; and the galleys of Basiliscus scattered the fleet of the Vandals in the neighbourhood of Sicily. On hearing of this disaster, Gaiseric is said to have given up all for lost; the Roman general had only to strike a decisive blow and Carthage would have fallen into his hands. But he let the opportunity slip, and, taking up his station in a haven at some distance from Carthage, he granted to the humble prayers of his wily opponent a respite of five days, of which Gaiseric made good use. He prepared a new fleet and a number of fireships. The winds favoured his designs, and he suddenly bore down on the Roman armament, which, under the combined stress of surprise, adverse wind, and the destructive ships of fire, was routed and at least half destroyed. Basiliscus fled with the remnant to Sicily, to join Marcellinus, whose energy and resources might possibly have retrieved the disaster; but the hand of an assassin, inspired perhaps by Ricimer, rendered this hope futile.\texthebrew{⁠}\footnote{Marcellinus, \emph{Chron.}

} Heraclius, who had not reached Carthage when he heard of defeat of the

p337
fleet, retraced his steps, and Basiliscus returned to Constantinople, where amid popular odium\texthebrew{⁠}\footnote{He was obliged to seek refuge in the sanctuary of St. Sophia.

} he led a life of retirement at Heraclea on the Propontis, until he appeared on the scene of public life again after Leo's death.

The ill-success of this expedition, organised on such a grand scale that it might have seemed irresistible, must have produced a great moral effect. The Roman Empire had put forth all its strength and had signally failed against one barbarian nation. This event must have not only raised the pretensions and arrogance of the Vandals themselves, but increased the contempt of other German nations for the Roman power; it was felt to be a humiliating disaster by the government at Constantinople, while the government of Italy was too habituated to defeat to be gravely affected.

The cost of the armament was immense. Leo had found in the treasury a reserve of 100,000 lbs. of gold (over £4,500,000).\texthebrew{⁠}\footnote{John Lydus, \emph{De mag.} III.43 .

Thayer's Note: The difficulties of converting financial figures! Today (2007), a hundred thousand pounds of gold would be just over a billion dollars; £4.5 million would be \$260 million.

} This was exceeded by the expenses of equipping the ill-omened expedition,\texthebrew{⁠}\footnote{According to Procopius, \emph{B. V.} I.6, the total cost was 130,000 lbs. of gold; according to Lydus, \emph{ib.}, 65,000 lbs. of gold and 700,000 lbs. of silver, which (calculating the ratio of gold to silver as 1:18) would together amount to about 104,000 lbs. of gold. The statement of Lydus evidently rests on the same data as the interesting notice of the historian Candidus

(\emph{fr.} 4, in \emph{F. H. G.} IV p137) . The chests of the Praetorian Prefects (East and Illyricum) contributed 47,000 lbs. gold, the treasury of the Sacred Largess, 17,000 lbs. gold; in all 64,000 lbs.; while the 700,000 lbs. of silver were supplied partly \textgreek{\texthebrew{ἐ}κ δημευσίμων} (\emph{i.e.} from confiscated property, and therefore from the treasury of the Private Estate) and partly by the treasury of Anthemius. It is unfortunate that we have not the story of the expedition given by the contemporary Priscus, whose work was the source of Theophanes, A.M. 5961, and indirectly (through Eustathius of Epiphania) of Procopius. Cp. Haury, \emph{Proleg.} to his ed. of Procopius, IX. \emph{sqq.}

} and the consequence is said to have been that the treasury hovered on the brink of bankruptcy for more than thirty years.

The conciliation of Gaul was a problem which was no less important for Anthemius than it had been for Majorian. The situation there had changed for the worse. The Visigothic crown had passed to Euric, who had murdered his brother Theoderic in A.D. 466. Euric was perhaps the ablest of all the Visigothic kings, and he aimed at extending his rule over all Gaul. The

p338
Gallo-Romans felt themselves now in greater danger, and they looked to Anthemius for protection with an eagerness which they had not shown in the case of Majorian. They sent a deputation to the new Emperor at Rome, both to petition him to remedy some administrative abuses and to stimulate him to take adequate measures for the defence of the Gallic provinces. The most distinguished member of the deputation was Sidonius Apollinaris.\texthebrew{⁠}\footnote{In an interesting letter (\emph{Epp.} I.5) Sidonius describes his journey to Rome (in 467), where he arrived when the nuptials of Ricimer with Alypia were being celebrated, and the city was given over to rejoi­cing.

} The panegyrist of Avitus and Majorian was now called upon to compose a panegyric of a third Emperor, on the occasion of his consul ­ship.\texthebrew{⁠}\footnote{\emph{Carm.} II .

} It was publicly recited on the kalends of January A.D. 468. The poet emphasised the fact that the elevation of Anthemius was a restoration of the unity of the Empire. He hailed Constantinople in these words:

and praised the Byzantine education of the new Augustus of the West. He was rewarded by the Prefecture of Rome. The appointment was much more than a recognition of his personal merit; it was intended to conciliate Gallo-Roman sentiment.\footnote{Cp. \emph{Epp.} I.9. Dalton has rightly pointed out that the whole affair was prearranged; the panegyric was a pretext, not the motive, of the appointment.

}

The pleasure of Sidonius in holding this high office was somewhat marred by the sensational trial of Arvandus, the Praetorian Prefect of Gaul, with whom he was on terms of friendship. Arvandus had sunk deeply into debt and had peculated public funds. His prosecution was decided by the Council of the Seven Provinces, and he was brought to trial before the Roman Senate. If malversation had been the only charge, he might have escaped through the influence of his friends, but he had been guilty of treasonable communications with the enemy, and there was clear proof of this in a letter in his own handwriting to King Euric, on which his accusers had managed to lay hands. Sidonius did all he could to help him, but the confidence of Arvandus himself, who was unable till the last moment to believe that he could be condemned, refused the advice of his friends and frustrated their efforts to save him. His confidence indeed was so strange that it has been conjectured that his communications with Euric

p339
had been secretly prompted by Ricimer, and that he was trusting in the protection of the Emperor's son-in\texthebrew{‑}law.\texthebrew{⁠}\footnote{Martroye, \emph{Genséric}, p234.

} He was condemned to death ``and flung into the island of the Serpent of Epidaurus (Island of the Tiber). There,'' writes Sidonius, ``an object of compassion even to his enemies, his elegance gone, spewed as it were by Fortune out of the land, he now drags out by benefit of Tiberius' law his respite of thirty days after sentence, shuddering through the long hours at the thought of hook and \texthebrew{•} Gemonian stairs, and the noose of the brutal executioner.''\footnote{Sidonius, \emph{Epp.} I.7, Dalton's translation. The sentence was not carried out. Seronatus, governor of Aquitanica Prima, was less lucky. Accused of oppression and treacherous relations with the Goths by the people of Auvergne, he was executed. Sidonius calls him a Catiline, \emph{Epp.} II.1.

}

Anthemius made large concessions to the Burgundians in Gaul to ensure their aid against the Goths, but he was not success ­ful in resisting the aggression of Euric.\texthebrew{⁠}\footnote{His son Anthemiolus was a commander in operations against the Goths. \emph{Chron. Gall.} 649 (p664).

} In Italy he was not popular. He was a Greek; he was too fond of philosophy or thaumaturgy; he was inclined to paganism.\texthebrew{⁠}\footnote{Damascius (\emph{Vita Isidori}, p208) says that he cherished the hope of restoring pagan idolatry.

} His high standard of justice and honest attempts to administer the laws impartially did not overcome the prejudices of the Italians, and the failure of the Vandal expedition did not heighten his prestige. His relations to Ricimer gradually changed from mutual tolerance to distrust and hostility; the father-in\texthebrew{‑}law regretted that he had given his daughter to a barbarian; the son-in\texthebrew{‑}law retorted with the epithets Galatian and Greekling (Graeculus). In this contest the Senate and people of Rome preferred the Greek Emperor to the Suevian patrician.\texthebrew{⁠}\footnote{John Ant. \emph{fr.} 93 (\emph{loc. cit.} p131).

} The question of ``Roman'' or German ascendancy, which had underlain the situation for fifteen years, was now clearly defined.

As a result of these dissensions, Italy in A.D. 472 was practically divided into two kingdoms, the Emperor reigning at Rome, the Patrician at Milan. The venerated Epiphanius, bishop of

Ticinum , attempted in vain to bring about a reconciliation. It will be remembered that Gaiseric had wished to elevate to the Imperial throne Olybrius, the husband of the younger Placidia. At this time Olybrius was at Constantinople, and his Vandal connexion made him a suspicious person in the eyes of Leo, who is said to have planned a treacherous device to remove him.

p340
He sent Olybrius to Rome for the ostensible purpose of reconciling Anthemius and Ricimer. But he also sent a messenger to Anthemius with a letter instructing him to put Olybrius to death. Ricimer intercepted the letter, and Leo's stratagem led to the result which he least wished.\texthebrew{⁠}\footnote{This transaction is related by John Malalas, XIV p374, and is quite credible. Cp. my note on \emph{The Emperor Olybrius} in \emph{E. H. R.}, July 1886 . Olybrius had been consul in 464. He was descended from Sextus Petronius Probus, consul in 371. His grandson (by his daughter Juliana) was consul in 491, and married Irene, a niece of the Emperor Anastasius. That he was never recognised as Augustus in the East seems clear from the circumstances, and Stein (\emph{Stud. z. Gesch. des byz. Reiches}, p176) had adduced confirmatory evidence.

} Ricimer invested Olybrius with the purple (April).

The army of Ricimer soon besieged Rome. Leo had overcome the power of Aspar; was Anthemius to overcome the power of Ricimer? In the camp of the besiegers was the Scirian soldier Odovacar, son of Edecon, destined soon to play a more memorable rôle in Italian history than Ricimer himself. The Tiber was guarded and supplies were cut off; and the Romans pressed by hunger resolved to fight. An army under Bilimer, who was perhaps Master of Soldiers in Gaul, had come to assist them. The Imperial forces lost heavily in the battle, and Ricimer completed his victory by treachery.\texthebrew{⁠}\footnote{Ennodius, \emph{Vit. Epiph.} p344 \emph{sqq.}; John Ant. \emph{fr.} 93 (\emph{Exc. de ins.} p131); Paul. Diac. \emph{Hist. misc.} 15.4; Schmidt, \emph{op. cit.} 262.

} Anthemius, when his adherents had surrendered to the barbarians, disguised himself and mingled with the mendicants who begged in the church of St. Chrysogonus.\texthebrew{⁠}\footnote{This church, restored more than once, still stands, near Sta. Maria in Trastevere.

Thayer's Note: For details and sources, and several further links, including to the detailed section in Armellini, see the entry

[Ecclesia] S. Chrysogoni

in Christian Hülsen's \emph{Chiese di Roma nel Medio Evo}.

} There he was found by Gundobad, Ricimer's nephew, and beheaded (July 11, A.D. 472).\footnote{John Ant. \emph{fr.} 209 (\emph{F. H. G.} IV) , says ``Gundobad, Ricimer's brother,'' and afterwards speaks of Gundibalos as his nephew. The fact is that Ricimer's sister married Gundioc, the Burgundian king, and their son was Gundobad, now in Roman service, but soon to succeed to a Burgundian throne.

}

But the days of Ricimer were numbered. He survived his father-in\texthebrew{‑}law by six weeks,\texthebrew{⁠}\footnote{He died August 18, 472, from vomiting blood, John Ant. \emph{ib.}

} and the last Emperor he created died two months later.\texthebrew{⁠}\footnote{Nov. 2, of dropsy, \emph{ib.} The date in the \emph{Paschale Campanum} (\emph{Chron. min.} I p306), borne out by an \emph{Auctarium} to Prosper (\emph{ib.} p492), is to be preferred to Oct. 23 of the \emph{Fasti Vind. pr.} (\emph{ib.} p306).

} He is not an attractive figure, and it would be easy to do him injustice. Barred by his Arian faith as well as by his German birth from ascending the throne, Ricimer had the choice of two alternative policies \texthebrew{—} to maintain an

p341
Imperial succession in Italy or to recognise the sole authority of the Emperor at Constantinople. It would probably have been repugnant to the ideas and traditions of his training to have cast off all allegiance to the Empire and created in Italy a government on German foundations, formally as well as practically independent. His choice of the first of the two policies was doubtless decided by public opinion and the influence of the Roman Senate, perhaps also by his own attachment to the system under which he was the successor of the great Masters of Soldiers, Stilicho and Aetius. But Italy had a taste of the other alternative in those sometimes long intervals between the puppet Emperors, when Leo was its only legitimate ruler. The success of Ricimer in maintaining this system for so many years was partly due to his diplomatic skill in dealing with Leo. But it worked badly. For it was based on the assumption that the Emperor was to be a nonentity like Honorius and Valentinian, and except in the case of Severus (whom Leo never acknowledged) circumstances hindered Ricimer from choosing a man who was suited to the rôle. In the matter of the expedition against the Vandals he had shown but lukewarm loyalty to the interests of the Empire, but Italy owed much to him for having defended her shores, and for having kept in strict control the German mercenaries on whom her defence depended. The events which followed his death will be the best commentary on the significance of his rule and enable us to appreciate his work.

The accession of Euric to the Visigothic throne, which he won by murder, meant the breaking of the last weak federal links which attached the Visigoths to the Empire.\texthebrew{⁠}\footnote{He sent an embassy to Constantinople

(Hydatius 238) . It has been conjectured that the purpose was to denounce the status of foederati and claim full sovranty (Schmidt, I.260).

} Euric was probably the ablest of their kings. He aimed at extending his power over all Gaul and Spain, and he accomplished in the eighteen years of his reign a large part of his programme. He was a fanatical Arian. ``They say that the mere mention of the name of Catholic so embitters his countenance and heart that one might take him for the chief priest of his Arian sect rather than for the monarch of his nation.''\texthebrew{⁠}\footnote{Sidonius, \emph{Epp.} VII.6.6.

} The principal

p342
hope of those Gallo-Romans of the south, who clung passionately to the Roman connexion, lay in the Burgundian power, which had itself in recent years made large encroachments on the Imperial provinces. King Chilperic ruled in Lyons and Vienne in the west, and at Geneva in the east; the provinces of Lugdunensis Prima and Maxima Sequanorum were almost entirely under his sway. His Arianism was not like that of Euric; he was tolerant and on friendly terms with Catholic bishops; he was glad to enjoy the breakfasts of Patiens, the rich and hospitable archbishop of Lyons.\texthebrew{⁠}\footnote{Sidonius, \emph{Epp.} VI.12.3.

} The higher clergy, who were mostly men of means and good family, played prominent parts in the politics of the time, and did a great deal to preserve the Roman tradition.\texthebrew{⁠}\footnote{Among the prominent bishops were, besides Patiens, Lupus of Troyes, Fonteius of Vaison, Perpetuus of Tours, Graecus of Marseilles, Leontius of Arles, Faustus of Riez, Basilius of Aix, and Sidonius himself.

} In the north the Imperial cause depended much on the attitude of the Salian Franks, who, under their king Childeric, seem to have been consistently loyal to their federal obligations. But in the Belgic provinces Roman civilisation was gradually declining.\texthebrew{⁠}\footnote{\emph{Ib.} IV.17.2.

} The lands of the Moselle and the Somme had never recovered from the shocks they had experienced in the days of Honorius. As for north-western Gaul, the province of the Third Lugdunensis, which was at this time generally called Armorica, it seems since some years before Valentinian's death to have been virtually independent.

The first important success that Euric won was a victory over the Bretons on the Indre. This enabled him to seize Bourges and the northern part of Aquitanica Prima, which, under their king Riothamus, they had come to defend at the request of the Emperor Anthemius. But he was unable to advance beyond the Loire, which was bravely defended by a count Paulus. Soon afterwards he laid siege to Arles, and defeated an Imperial army which had advanced to relieve it under Anthemiolus, the Emperor's son. Arles he appears to have occupied and then to have marched up the valley of the Rhone, burning the crops, and taking the towns of Riez, Orange, Avignon, Viviers and Valence.\texthebrew{⁠}\footnote{\emph{Ib.} VI.12.

} He did not hold these places, for he was not prepared to go to war with the Burgundians, but he left the land ruined, and the people would have starved if the archbishop Patiens

p343
had not collected supplies of corn º at his own expense, and sent grain carts through the ravaged districts.

Euric was determined to annex the rich country of Auvergne, and here he met a stout and protracted resistance, of which Ecdicius,\texthebrew{⁠}\footnote{He seems to have held the post of Master of Soldiers in Gaul, see

Jordanes, \emph{Get.} 45 . Nepos created him patrician in 474 (Sidonius, \emph{Epp.} V.16).

} son of the Emperor Avitus, was the soul. He was supported by his brother-in\texthebrew{‑}law , Sidonius Apollinaris, now bishop of Clermont, which held out for nearly four years against repeated sieges. But no help came either from Italy or from Burgundy, and finally the Emperor Julius Nepos arranged a peace with Euric, which surrendered Auvergne and recognised the conquests which the Goths had already made in Spain as well as in Gaul ( A.D. 475).\texthebrew{⁠}\footnote{The negotiations were conducted by four south Gallic bishops (\emph{ib.} VII.6, 10), and also by Epiphanius, bishop of Ticinum (Ennodius, \emph{Vit. Epiph.} 81); but it is not clear whether or not two separate missions were sent to Euric.

} The Gallic portion of the Gothic kingdom was now bounded by the Loire, the Rhone, and the Pyrenees, and seems to have included Tours.

Sidonius was taken prisoner and confined in fort Livia, near Carcassonne.\texthebrew{⁠}\footnote{Sidonius was very bitter over the surrender of Auvergne, \emph{ib.} VII.7.

} Here he employed his time in editing or translating the life of Apollonius of Tyana, by Philostratus, and was so well treated that the worst he had to complain of was that when he lay down to sleep ``there were two old Gothic women established quite close to the window of my chamber who at once began their chatter \texthebrew{—} quarrelsome, drunken, and disgusting creatures.''\texthebrew{⁠}\footnote{\emph{Ib.} VIII.3.3.

} He was finally released through the influence of Leo, the principal minister of Euric and his own good friend.

The peace lasted for little more than a year. Then Euric found a pretext for denouncing it, invaded Provence, and seized Arles and Marseilles. Then a new arrangement was made, and southern Provence, with the consent of the Emperor Zeno, was conceded to the Goths.\footnote{Procopius, \emph{B. G.} I.12.20. Cp. Candidus, in \emph{F. H. G.} IV p136 . Schmidt, \emph{op. cit.} I.267.

}

Euric was now the most power ­ful of the German kings. The Burgundians hastened to make peace with him. Ostrogoths, Heruls, Saxons, Franks were to be seen at Toulouse or Bordeaux paying court to him. Even the Persian king thought it worth while to send envoys

p344
to his court.\texthebrew{⁠}\footnote{Sidonius, VIII.9.5.

} When he died in A.D. 484 the Spanish peninsula, except the Suevian kingdom in the north-west, was entirely under his dominion.\footnote{The capture of Caesaraugusta, the siege of Tarraco, and the capture of coast cities in A.D. 473 are recorded in \emph{Chron. Gall.} pp664\texthebrew{‑}665. Cp. Isidore, \emph{Hist. Goth.} 34 .

}

For the Gallic provincials the change of masters probably made very little difference. They and the Goths lived side by side, each according to their own law. The Roman magnate had to surrender a part of his estates, but he could live with as much freedom and ease, and in just the same way, under the Goth as under the Emperor. Some of these men were enlisted in the royal service, such as Leo of Narbonne; Namatius, who commanded the Gothic fleet in the Atlantic to guard the coasts against Saxon pirates;\texthebrew{⁠}\footnote{Sidonius, VIII.6.

} Victorius, who was made governor of Auvergne. Latin was the language of intercourse. It is probable that very few provincials learned any of the German tongues which were spoken by their masters. Syagrius, a man of letters, who lived much at the Burgundian court, mastered the Burgundian language, to the amazement of his friends. Sidonius bantered him on his feat. ``You can hardly conceive how amused we all are to hear that, when you are by, not a barbarian but fears to perpetuate a barbarism in his own language. Old Germans bowed with age are said to stand astounded when they see you interpreting their German letters; they actually choose you for arbiter and mediator in their disputes. You are a new Solon in the elucidation of Burgundian law. In body and mind these people are as stiff as stocks and very hard to form; yet they delight to find in you, and equally delight to learn, a Burgundian eloquence and a Roman spirit.''\texthebrew{⁠}\footnote{\emph{Ib.} V.5 (this and the other quotations are taken from Dalton's translation). Syagrius was a great-grandson of Flavius Afranius Syagrius, who was Pr. Pr. of Gaul and consul in A.D. 382.

} In this connexion it is significant that the early German codes of law were composed in Latin. The earliest that we know of was the code of Euric, of which some fragments are preserved;\texthebrew{⁠}\footnote{Edited by Zeumer in the \emph{Leges Visigothorum.}

} a little later come the Burgundian laws of Gundobad. It is legitimate to guess that the Visigothic law-book was drawn up under the supervision of Euric's minister Leo, who was a notable jurist.

Sidonius gives us occasional glimpses of the life and habits

p345
of the Germans, who were then moulding the destinies of Gaul. Writing to a friend, for instance, he describes the wedding of a Burgundian princess: the bridegroom,\texthebrew{⁠}\footnote{Sigismer, otherwise unknown. Sidonius, IV.20.

} walking amid his guards ``in flame-red mantle, with much glint of ruddy gold, and gleam of snowy silken tunic, his fair hair, red cheeks and white skin according with the three hues of his equipment.'' The chiefs who accompanied him were in martial accoutrement. ``Their feet were laced in boots of bristly hide reaching to the heels; ankles and legs were exposed. They wore high tight tunics of varied colour, hardly descending to the bare knees, the sleeves covering only the upper arm. Green mantles they had with crimson borders; baldrics supported swords hung from their shoulders, and pressed on sides covered with cloaks of skin secured by brooches. No small part of their adornment consisted of their arms; in their hands they grasped barbed spears and missile axes; their left sides were guarded by shields which flashed with tawny golden bosses and snowy silver borders, betraying at once their wealth and their good taste.''

Sidonius confesses that he did not like Germans,\texthebrew{⁠}\footnote{\emph{Ib.} VII.14.10.

} and it is the society of his own fellows, the country gentlemen of southern Gaul, among whom he had a wide acquaintance, that is mainly depicted in his correspondence. The life of these rich members of the senatorial class went on its even and tranquil way, little affected by the process which was gradually substituting Teuton for Roman power.\texthebrew{⁠}\footnote{In Dalton's Introduction to his translation of the Letters there is an admirable account of this society.

} They had generally town mansions, as well as country estates on which they lived, well provided with slaves, and amusing themselves by hunting, hawking, and fishing, ball-games, and dice. But the remarkable feature of the life of these Gallo-Roman magnates was that they did not confine themselves to the business of looking after their domains and the outdoor pursuits of country gentlemen, but were almost all men of literary tastes and culture. There were many poets and trained rhetoricians among them; they circulated their verses; and mutually admired one another's accomplishments. It is probable that in literary achievement Sidonius was considerably superior to his friends, but in any case his works show

p346
us the sad decadence in style to which the tendencies of the rhetorical schools of the Empire, in Gaul as elsewhere, had brought literary prose. Of his epistolary style it is enough to say that it gains in a good modern translation. He could write good verses, occasionally approaching Claudian, and bad verses, which remind us of Merobaudes.

Of the last thirty years of Imperial rule in northern Gaul we know virtually nothing. Childeric, the principal king among the Salian Franks, seems to have loyally maintained the federal bond with the Empire.\texthebrew{⁠}\footnote{He had fought with Aegidius against the Goths at Orleans (see above,

p333 ).

} The blue-eyed Saxons, who were at this time the scourge of the coasts of Gaul, in the west as well as in the north, had sailed up the Loire and seized Angers. We find Childeric aiding the Imperial commander Paul in his operations against this foe.\texthebrew{⁠}\footnote{For the dealings with the Saxons see

Gregory of Tours, \emph{H. F.} II.18, 19 . On their invasion and the \emph{Litus Saxonicum} in Gaul see Lot, \emph{Les Migrations saxonnes}, 6 and 13 \emph{sqq.}

} We have already seen Paul holding the line of the Loire against the Visigoths. We are not told his official rank or functions; he is designated by the title of Count, but we may fairly assume that he had succeeded Aegidius as Master of Soldiers. His name and that of Syagrius are the only two recorded names of Roman functionaries who maintained Imperial authority in northern Gaul after the death of Aegidius. Syagrius was the son of Aegidius, and on him devolved the defence of Belgic Gaul in the last years of Childeric.\footnote{We may conjecture that he owed his appointment either to Anthemius or to Julius Nepos, and that he succeeded Paul as Master of Soldiers. He is mentioned only by Gregory of Tours, unless, as some think, he is identical with the correspondent of Sidonius referred to above.

}

Childeric died in A.D. 481 and was succeeded by his son Clovis (Chlodwig), who entered upon new paths of policy. He saw clearly that the Imperial power in Gaul was now negligible. The few provinces that were still administered in the name of the Augustus at Constantinople were cut off from the rest of the Empire by the kingdoms of the Visigoths and the Burgundians. It was evidently the destiny of Gaul to be possessed entirely by German rulers, and Clovis determined that the Franks should have their share. He took the field against Syagrius soon after his accession and defeated him near Soissons ( A.D. 486).\texthebrew{⁠}\footnote{Syagrius fled to Toulouse, but King Alaric gave him up to Clovis, who put him to death.

} The province of Belgica Secunda, with the important cities of Soissons

p347
and Reims, immediately passed under his sway.\texthebrew{⁠}\footnote{For some years Clovis allowed the Imperial administration to continue unchanged in this province. See the letter which Remigius, archbishop of Reims, addressed to him (\emph{Epp. Austras.} 2, in \emph{Epp. Mer. et Kar. aevi}, vol. I), which must be dated after 496.

\texthebrew{◂} previous

next \texthebrew{▸}Images with borders lead to more information.

The thicker the border, the more information.

(Details here.)

UP TO:

J. B. Bury
Homepage

Latin \& Greek

Texts

LacusCurtius

Home} Of his subsequent advance westward to the Loire and the borders of peninsular Brittany we know nothing, probably because it was gradual and easy.

The victory of Soissons completely changed the political situation and prospects of Gaul. Two years before, when Euric died, the destinies of the land seemed to depend on the Goths and the Burgundians, and if any one had prophesied that the whole land would ultimately be ruled by Gothic kings, few outside Burgundy would have questioned the probability of the prediction. Yet twenty years later the formidable power which Euric had created was to go down before the Franks; afterwards it would be the turn of the Burgundians. The failure of the Goths to fulfil their early promise was due above all to their Arian faith, which deprived them of the support of the Church. When Clovis embraced Christianity in its Catholic form, ten years after the battle of Soissons, he made the fortune of the Franks.

The part which the Church was able to play throughout the critical age in which the country was passing from Roman to Teuton lords depended on the fact that the Gallic episcopate was recruited from the highly educated and propertied class. The most public-spirited members of the senatorial families found in the duties of a bishop an outlet for their energies. It was these bishops who mediated between the German kings and the Roman government, and after the Imperial power had disappeared, helped to guide and moderate the policy of the barbarian rulers towards the provincials, and to preserve in some measure Gallo-Roman traditions. The study of the society mirrored in the pages of Sidonius, himself a case in point, is an indispensable preparation for the study of the France created by Clovis, of which the early history is recorded by Gregory, the bishop of Tours.

\xchapter{Chapter XI}

Short URL for this page:

bit.ly/BURLAT11

The existence of the State Church made a profound difference in the political and social development of the Empire. The old State religion of Rome was often used as an instrument of policy, but perhaps its main political value was symbolic. It involved no theory of the universe, no body of dogma to divide the minds of men and engender disputes. The gods were not jealous, and it was compatible with the utmost variety of other cults and faiths. For the Christian Church, on the contrary, a right belief in theological dogmas was the breath of its life, and, as such questions are abstruse and metaphysical, it was impossible to define a uniform doctrine which all minds would accept. As the necessity of ecclesiastical unity was an axiom, the government had to deal with a new problem, and a very arduous and embarrassing one, such as had not confronted it in the days before Constantine. Doctrine had to be defined, and heretics suppressed. Again, the Church, which once had claimed freedom for itself, denied freedom to others when it was victorious, and would not suffer rival cults. Hence a systematic policy of religious intolerance, such as the Greek and Roman world had never known, was introduced. Another consequence of the Christianising of the State was the rise to power and importance of the institution of monasticism, which was not only influential economically and socially, but was also, as we shall see, a political force. The theological controversies, the religious persecution, and the growth of monasticism, in the fifth century, will be reviewed briefly in this chapter.

The great theological controversy which rent Christendom in twain in the fourth century had been finally closed through the energy and determination of Theodosius the Great, and unity was for a short time restored to the Church. Theodosius had been baptized in Thessalonica in A.D. 380, and immediately afterwards he issued an edict, commanding his subjects to accept the orthodox faith of the Council of Nicaea.\texthebrew{⁠}\footnote{\emph{C. Th.} XVI.1.2 . On the doctrine of Nicaea Harnack (\emph{History of Dogma}, IV.49) observes: ``One of its most serious consequences was that from this time forward Dogmatics were for ever separated from clear thinking and defensible conceptions, and got accustomed to what was anti-rational. The anti-rational \texthebrew{—} not immediately, but soon enough \texthebrew{—} came to be considered as the characteristic of the sacred.''

} He described it as the doctrine professed by the bishop of Rome and the bishop of Alexandria. Then he proceeded to hand over to the orthodox all the Arian churches in Constantinople, and to prohibit heretics from holding public worship in the city. In the meantime he had come to see that the best prospect of terminating discussion in the East would be by a Council which was not controlled either from Alexandria or from Rome. The Council which met at his summons in A.D. 381 at Constantinople was entirely eastern, and Meletius, the bishop of Antioch, presided. Seventy years later it came to be called an Ecumenical Council; in the West it was not recognised as such till the end of the fifth century. This assembly of eastern bishops ratified the doctrine of the Council of Nicaea, and declared that the Son is of the same substance with the Father. Theodosius, after a vain attempt to win over the Arians by a Council which he summoned two years later, proceeded to measures of suppression,\texthebrew{⁠}\footnote{\emph{C. Th.} XVI.1.4 .

} and Arianism gradually declined.

But, while the Arian heresy in itself led to no permanent schism in the Church,\texthebrew{⁠}\footnote{Harnack, \emph{ib.} 106: ``The educated laity in the East regarded the orthodox formula rather as a necessary evil and as an unexplainable mystery than as an expression of their Faith. The victory of the Nicene Creed was a victory of the priests over the faith of the Christian people. The Logos-doctrine has already become unintelligible to those who were not theologians. . . . The thought that Christianity is the revelation of something incomprehensible became more and more a familiar one to men's minds.'' He refers to a quotation in Socrates, III.7, from Evagrius the Anchorite, who will have nothing to do with theological categories, and says \textgreek{σιωπ\texthebrew{ῇ} προσκυνείσθω τ\texthebrew{ὸ ἄ}ρρητον}, ``Let the mystery be adored in silence.''

} new and closely related controversies soon agitated

p350
the eastern world and were destined to issue in lasting divisions. Once the divinity of Christ in the fullest sense was universally admitted, the question ensued how the union of his divine substance with his human nature is to be conceived. Was the Godhead mixed with humanity, or only conjoined? Did Mary bear the flesh only or the Logos along with the flesh? Did Christ's human nature survive the Resurrection? In the fourth century, there was no definite doctrine, but the problem was disturbing the minds of some metaphysical theologians.

Apollinaris of Laodicea argued that the union of a perfect God into a perfect man was out of the question. For the result of such a union would be a monster, not a uniform being. He concluded that Christ was not a perfect man, and that he adopted human nature, determining it in such a way that it did not involve free will, which would be inconsistent with his Godhead. His flesh was taken up into the nature of the Logos and was thus divine, and the Logos shared in the sufferings of the flesh. Further, Christ's mind was not human; for, if he had had a human mind, he would have had a duplicate personality.

It has been said that this theory of Apollinaris expressed the belief entertained at heart by all pious Greeks.\texthebrew{⁠}\footnote{Harnack, p155.

} But it was clear that it did not do justice to the humanity of Christ as depicted in the Gospels, and other theologians, who like Apollinaris himself belonged to the school of Antioch, sought to render intelligible the union of a perfect God with a perfect man. According to Theodore of Mopsuestia, the union of the two natures was a contact which became more intimate at each stage of human growth, and the indwelling of the Logos in the man was not substantial, but of the same order as the indwelling of God, by grace, in any human being. Each nature was itself a person, and the Logos did not become man. It was the man only who suffered. And Mary was not, in the strict sense, the mother of God.

In the reign of Theodosius II this insoluble problem raised a bitter controversy, which agitated the eastern world. When Sisinnius, Patriarch of Constantinople, died at the end of A.D. 427, the bishops, the clergy, and the monks could not agree on the appointment of a successor, and the nomination was committed

p351
to the Emperor; who, seeing that no possible candidate among the ecclesiastics of Constantinople would be generally acceptable, chose Nestorius,\texthebrew{⁠}\footnote{The Emperor's difficulties are shown in his remarkable conversation with the abbot Dalmatius, recorded by Nestorius in the \emph{Bazaar of Heraclides}; see the extract in Bethune-Baker, \emph{Nestorius and his Teaching}, p6, n3. He desired Dalmatius to choose, but Dalmatius refused.

} a monk of a convent at Antioch, who had a high reputation as a preacher. The eloquence of Nestorius was matched by his intolerance, and no sooner was he seated on the Patriarchal throne\texthebrew{⁠}\footnote{He was consecrated on April 10, 428.

} than he began an energetic campaign against heresies. But his forcible language in condemning Apollinarian views, which he discovered to be rife among the local clergy, soon gave the Patriarch of Alexandria, who was the natural enemy of any Patriarch of Constantinople, a welcome opportunity of accusing him of heresy himself. The rivalry between these great sees, bitter since the Council of A.D. 381, when precedence over all sees except Rome had been granted to New Rome,\texthebrew{⁠}\footnote{By the 3rd Canon of the Council: \textgreek{τ\texthebrew{ὰ} πρεσβε\texthebrew{ῖ}α τ\texthebrew{ῆ}ς τιμ\texthebrew{ῆ}ς μετ\texthebrew{ὰ} τ\texthebrew{ὸ}ν τ\texthebrew{ῆ}ς \texthebrew{Ῥ}ώμης δι\texthebrew{ὰ} τ\texthebrew{ὸ} ε\texthebrew{ἶ}ναι α\texthebrew{ὐ}τ\texthebrew{ὴ}ν νέαν \texthebrew{Ῥ}ώμην.}} had been aggravated by the struggle between Theophilus and Chrysostom.

The Patriarch Cyril and the Alexandrines held that the two natures of Christ were joined in an indissoluble, ``hypostatic'' or personal union, yet remained distinct, but that the human nature had no substance independently of the divine; that the Logos suffered without suffering, and that Mary is the mother of God inasmuch as she bare flesh which was united indissolubly with the Logos. Cyril's doctrine approached that of Apollinaris in so far as it denied the existence of an individual man in Christ, but was sharply opposed to it by its maintenance of the distinction of the two natures.

Nestorius leaned to the doctrine of Theodore of Mopsuestia, which was popular in Syria. He characterised as fables the statements that a God was wrapped in swaddling clothes and was nailed upon the cross, and he protested against the use of the designation ``Mother of God'' ( Theotokos ).

It is to be observed that in this controversy both parties agreed in condemning the theory of Apollinaris and in holding that there were two natures in Christ. The main difference between them concerned the formula by which the union of the two natures was to be expressed \texthebrew{—} Cyril maintaining a ``natural

p352
union''\texthebrew{⁠}\footnote{\textgreek{\texthebrew{Ἕ}νωσις φυσική} or \textgreek{\texthebrew{ὑ}ποστατική}. On the ambiguity of this phrase see Bethune-Baker, \emph{op. cit.} 171 \emph{sqq.} The term hypostasis, subsistence, is not quite synonymous with \textgreek{ο\texthebrew{ὐ}σία}, substance; the difference is thus explained by a Nestorian theologian: ``We apply the term hypostasis to the particular substance, which subsists in its own single being, numerically one and separate from the rest'' (Labourt, \emph{Le Christianisme}, pp283\texthebrew{‑}284). Cp. Bethune-Baker, \emph{op. cit.} 220 \emph{sqq.} Nestorius maintained there were two natures, two substances, and two hypostaseis in Christ.

} and Nestorius a less intimate ``contact.''\texthebrew{⁠}\footnote{\textgreek{Κατ\texthebrew{ὰ} συνάφειαν.}} The truth may be that the view of Nestorius was not so very different from that of Cyril as Cyril thought. It seems probable that the doctrine of two Persons, somehow joined together, which is commonly imputed to Nestorius, would have been repudiated by him.\texthebrew{⁠}\footnote{See Bethune-Baker, \emph{op. cit.} ch. VI. The main object of this book was to prove that Nestorius was orthodox and was not a ``Nestorian.'' The dialogue of Nestorius, the \emph{Bazaar of Heraclides}, or \emph{\textgreek{Πραγματεία \texthebrew{Ἡ}ρακλείδου}}, recently discovered in a Syriac version, supplies the important evidence that Nestorius survived till the eve of the Council of Chalcedon and agreed with the \emph{Dogmatic Epistle} of Pope Leo. Cp. Loofs' \emph{Nestorius} (pp21, 22), which gives a clear and interesting account of the tragedy of Nestorius. This theologian agrees with Bethune-Baker partially; he concludes that Nestorius can be considered orthodox \emph{according to the western interpretation} of the definition of Chalcedon (p100). On the meaning of the term \textgreek{πρόσωπον} (person) see pp76 \emph{sqq.}

} Cyril wrote to Theodosius, to Eudocia, to Pulcheria and her sisters, censuring the heretical opinion of Nestorius,\texthebrew{⁠}\footnote{See Mansi, IV.617, 680. Cyril counted on theological differences in the Imperial family. Theodosius and Eudocia were under the influence of Nestorius. Theodosius saw through Cyril's tactics, and wrote him a sharp letter (\emph{ib.} V.1109). Pulcheria took the other side. She had quarrelled with Nestorius, who is said to have repeatedly rebuffed and insulted her. Nestorius describes her as ``a bellicose woman, a queen, a young virgin, who quarrelled with me because I would not agree to her demand of comparing a person corrupted by men to the spouse of Christ'' (\emph{Book of Heraclides}, p89). It appears that he questioned her virtue. The grievances of Pulcheria against Nestorius are enumerated in a letter written after the death of Nestorius to a certain Cosmas of Antioch (\emph{ib.} App. I pp363\texthebrew{‑}364).

} and stirred up the Egyptian monks, who were ever ready for a theological fray. A heated correspondence ensued between the two Patriarchs, and both invoked the support of Celestine, the bishop of Rome. Pope Celestine was no theologian. He was guided by the political expediency of supporting Alexandria against Constantinople, and he evaded the real issue by bringing into the forefront of the controversy a minor point, namely the question whether Mary might properly be called the Mother of God. On this particular point Nestorius was ready to yield, but he would not recant his doctrine at the bidding of a Roman synod.\texthebrew{⁠}\footnote{Held early in August 430. It condemned the views of Nestorius.

} Anathemas and counter-anathemas flew between Alexandria and Constantinople, and then the Emperor, by the advice of Nestorius,

p353
summoned a Council on the neutral ground of Ephesus for Whitsuntide A.D. 431. The two antagonists arrived in good time, but John the Patriarch of Antioch was three weeks late. Cyril, who was accompanied by fifty bishops, would not wait for him; and the supporters of the Alexandrian party met and decreed the deposition of Nestorius, who refused to attend the assembly. When John and the Syrian contingent arrived, a rival but far less numerous Council was opened; the commissioner Candidian, Count of the Domestics, who represented the Emperor, presided; and Cyril was condemned and deposed. Then the Roman legates appeared upon the scene, attended the assembly of Cyril, and signed the decree against Nestorius.

The shameless proceedings of the satellites of Cyril and the rabble whom they are collected are graphically described by Nestorius, whose house was guarded by soldiers to protect him from violence. ``They acted in everything as if it was a war they were conducting, and the followers of the Egyptian and of Memnon (the bishop of Ephesus), who were abetting them, went about in the city girt and armed with clubs, men with high necks, performing strange antics with the yells of barbarians, snorting fiercely with horrible and unwonted noises, raging with extravagant doings, carrying bells about the city, and lighting fires in many places and casting into them all kinds of writings. Everything they did was a cause of amazement and fear; they blocked up the streets so that every one was obliged to flee and hide while they acted as masters of the situation, lying about drunk and besotted and shouting obscenities.''\texthebrew{⁠}\footnote{\emph{Book of Heraclides}, in Bethune-Baker, p39; in Nau's version, p236. The antagonists of Nestorius complained similarly of the violence of the other faction and also of the partiality of Candidian. The best and fullest account of the whole proceedings is probably that of Tillemont, \emph{Mémoires}, XIV.307 \emph{sqq.}, allowing for his prejudice against Nestorius.

} Such were the circumstances of the Third Ecumenical Council, which had gathered to pronounce on the true doctrine of the natures of Christ.

The Emperor had at first resolved to reject the decree against Nestorius, but afterwards he decided to carry out the rulings of both assemblies. The two Patriarchs were deposed; Nestorius retreated to his old convent at Antioch. But at Constantinople there was a strong ecclesiastical opposition to Nestorius; the clergy addressed a petition to the Emperor demanding justice

p354
for Cyril, and the monks, under the leader ­ship of Dalmatius, excited the people.\texthebrew{⁠}\footnote{Mansi, IV.1453; 1428.

} The popular demonstrations were aided by Cyril's intrigues and a lavish distribution of bribes;\texthebrew{⁠}\footnote{A list of Cyril's presents to the chamberlains and some of the ministers at the court is preserved (printed in Nau, \emph{op. cit.} p368). The most important persons to gain over were the Empress Eudocia and the Grand Chamberlain Chrysoretus, both of whom had been on the side of Nestorius. Two ladies-in\texthebrew{‑}waiting, Marcella and Droseria, received each £2250 for ``persuading'' Eudocia. Paul, who was probably the praepositus of Pulcheria, received the same amount and a number of valuable household things (carpets, ivory chairs, etc.). Similar but more numerous presents were made to the Grand Chamberlain to buy off his opposition, and he was promised £9000 for his help. Helleniana, the wife of the Praetorian Prefect of the East, was presented with gifts of the same kind and number, and was to receive £4500 if she enlisted her husband's help. And so on.

Thayer's Note: A reminder that Bury wrote in 1923. That these are very substantial bribes is made clear by converting them into today's money (2007): as a very rough estimate,

£2250 would be £131,000, or about \$262,000.

£4500: £262,000, or about \$524,000.

£9000: £524,000 \texthebrew{—} over a million dollars.

} Pulcheria doubtless threw her influence into the scale; and the Emperor was compelled to yield and to permit Cyril to resume his Patriarchal seat. Cyril then sought to come to terms with Antioch, and a new formula was invented \texthebrew{—} ``the unconfused union of two natures'' \texthebrew{—} which could be accepted both by the Alexandrines and by moderate men of the Antiochian school. Cyril subscribed to this creed in A.D. 433. Good Nestorians retreated to Edessa, and here their theology was in the ascendant until the Emperor Zeno ( A.D. 489) took measures to extirpate Nestorianism and succeeded in driving it beyond the frontier. The subsequent fortunes of the sect are connected with Persian and Saracen history.

It is clear that throughout the whole controversy personal dislike of Nestorius, who was not an amiable or courteous man, played a considerable part. He was permitted to remain peacefully in his monastery for a few years, notwithstanding the urgent request of Pope Celestine that such a firebrand should be removed from all contact with men. But at length the Emperor adopted harsh measures against him ( A.D. 435).\texthebrew{⁠}\footnote{It seems that John the Patriarch of Antioch, who had supported him at Ephesus, found his presence embarrassing and made representation at court. Pulcheria was believed to be responsible for the exile to Oasis. See Brière, ``La Légende syriaque de Nestorius,'' in \emph{Revue de l'orient chrétien} (1910), pp1\texthebrew{‑}25.

} He was denounced in an edict as sacrilegious, his books were condemned to the flames,\texthebrew{⁠}\footnote{\emph{C. Th.} XVI.5.66 .

} and he was banished at first to Petra and then to Oasis in Upper Egypt ( A.D. 435). He seems to have died in A.D. 451.\footnote{He was for some time a prisoner among the Blemmyes. For the hardships he endured see

Evagrius, I.7 . In the \emph{Book of Heraclides}, which he wrote shortly before his death, he describes the proceedings of the Second Council of Ephesus (449), and he implies that the faith which he regarded as true had triumphed (at

(p355)
Chalcedon), and Dioscorus had been defeated. See Bethune-Baker, \emph{op. cit.} 34\texthebrew{‑}35, and
\emph{Journal of Theol. Studies}, IX.601 .

}

p355
The compromise of A.D. 433 was not final. The question was opened again by Dioscorus, who had succeeded Cyril ( A.D. 444) in the see of Alexandria, and was jealous of the prestige of the theologians of Antioch. He set himself the task of destroying the Antiochian formula of ``two natures or hypostaseis and one Christ.'' His views found a warm supporter at Constantinople in a certain Eutyches, the archimandrite of a monastery, who had been prominent in the agitation against Nestorius, and enjoyed the favour of the eunuch Chrysaphius.\texthebrew{⁠}\footnote{Cp. Victor Tonn. \emph{sub a.} 450.

} Eutyches was charged with heresy; the Patriarch Flavian\texthebrew{⁠}\footnote{Nestorius had been succeeded by Maximian, 431, Proculus was elected in 434, Flavian in 446. There is a good article on Eutyches in the \emph{Dict. of Chr. Biogr.}

} took up the matter and procured his condemnation at a local synod ( A.D. 448). Eutyches appealed to Leo, the bishop of Rome; and Dioscorus urged the Emperor to summon a general Council. Theodosius, guided by the counsels of Chrysaphius who hated Flavian, yielded to the wishes of Alexandria, and the Council met at Ephesus in August A.D. 449.

In the meantime Leo had come to the conclusion that the views of Eutyches were heretical, and he wrote in this sense to the Emperor and the Patriarch. He claimed that he was himself the person who should decide and define the dogma by virtue of the authority residing in the see of St. Peter; there was no necessity for a General Council.\texthebrew{⁠}\footnote{The collection of Pope Leo's letters (\emph{P. L.} LIV; Mansi, VI) includes not only his own letters on the controversy to Theodosius, Pulcheria, Marcian, eastern ecclesiastics, etc., but also the correspondence of Galla Placidia and others with the Imperial family at Constantinople.

} But the Council was called, and Leo sent three delegates, committing to them a Dogmatic Epistle or Tome addressed to Flavian in which he formulated the true doctrine: the unity of two hypostatic natures in one person, wherein the properties of both natures were preserved.\texthebrew{⁠}\footnote{Leo, \emph{Ep.} 28 (\emph{P. L.} LIV.755 \emph{sqq.}).

} It was not explained how this union was possible, and a distinguished historian of dogma observes\texthebrew{⁠}\footnote{Harnack, \emph{ib.} 207.

} that Leo left off at the point where the speculation of Cyril began.

Dioscorus presided at the Council. The letter of Leo was not read, and the Roman representative did not vote. Eutyches was declared orthodox, and Flavian was deposed as having gone

p356
beyond the doctrine of the creed of Nicaea.\texthebrew{⁠}\footnote{He was banished to Hypaepa in Lydia, and died on his way thither in consequence of ill-treatment.

} Other more distinguished adherents of the Antiochian doctrine, including Theodoret, bishop of Cyrrhus, a notable theologian, were also deposed. The result of the proceedings was to annul the compromise of A.D. 433 and to reinstate the Cyrillian doctrine of the one incarnate nature of the God-Logos . The voting of many of the 115 bishops who signed the Acts was not free; they were overawed by the Imperial authorities and by the violence of a noisy crowd of monks from Syria. Yet it has been said, perhaps with truth, that this Council more than any other expressed the general religious feeling of the time, and would have permanently settled the controversy in the East if extraneous interests had not been involved.

The bishop of Rome denounced the ``Robber Council,'' as he called it, and prompted Valentinian III to propose to his cousin Theodosius the convention of a new Council in Italy. Theodosius replied that the recent Council had simply defended the rulings of Nicaea and Ephesus against the innovations of Flavian; no further action was called for; the Church was at peace. If the question had been simply doctrinal and no political considerations had intervened, the decision of the ``Robber Council'' might have been the last word in Eastern Christendom. But that Council had been a triumph for Alexandria, and the prestige which Dioscorus acquired was a menace not only to Old Rome \texthebrew{—} he promptly excommunicated Leo \texthebrew{—} but also to New Rome. This danger could not long be ignored, and the death of Theodosius was followed by a change of policy at Constantinople.

Marcian resolved to terminate the ecclesiastical despotism which the Alexandrian bishops sought to impose upon the East, and Anatolius, who through the influence of Dioscorus had succeeded Flavian as Patriarch, did not scruple to lend himself to a new policy and to subscribe the Dogmatic Epistle of Leo. Marcian wrote to Leo agreeing to his request for a new Council, but insisting that it should meet in the East. Then the Pope changed his tactic. He claimed, as before, that his own Epistle was sufficient to settle the whole matter, and did all he could to prevent the meeting of a Council.\texthebrew{⁠}\footnote{Leo, \emph{Epp.} 82\texthebrew{‑}86.

} But Marcian knew that,

p357
however wonder ­ful Leo's Epistle might be, a Council would be indispensable to satisfy public opinion in the Eastern Churches, and he summoned a Council for the autumn ( A.D. 451). Leo rather sulkily yielded.\texthebrew{⁠}\footnote{\emph{Ep.} 89.

} In October an unusually large assembly of ecclesiastics\texthebrew{⁠}\footnote{About 600.

} met at Chalcedon, and the presidency, which meant the right of first recording his vote, was given to the legate of the Pope.

It was the common object of Leo and of Marcian to procure the deposition of Dioscorus, and in this they succeeded, but not without exercising moral violence. Most of the bishops, including Anatolius who really agreed with Dioscorus, voted against their consciences and relinquished the formula in which they believed. But, while Leo desired that his epistle should be accepted as it stood, Marcian saw that a new formula, which should indeed take account of the Pope's statement, would be less unacceptable in the East. Accordingly the Council decreed that the true doctrine was contained in certain writings of Cyril\texthebrew{⁠}\footnote{Namely, his Synodal letters to Nestorius and the Orientals (\emph{Epp.} 4, 17; 39).

} as well as in Leo's epistle; and described Jesus Christ as complete in his humanity as well as in his divinity; one and the same Christ in two natures, without confusion or change, division or separation;\texthebrew{⁠}\footnote{\textgreek{\texthebrew{Ἐ}ν δύο φύσεσιν \texthebrew{ἀ}συγχύτως κτλ}. Cp. Loofs, \emph{op. cit.} 97.

} each nature concurring\texthebrew{⁠}\footnote{\textgreek{Συντρεχούσης.}} into one person and one hypostasis.

The doctrine of the Fourth Ecumenical Council is still accepted as authoritative in the Churches of Christendom. It is interesting to learn the judgment of one of the most learned living theologians. The Council of Chalcedon, ``which we might call the Robber and the Traitor Council, betrayed the secret of the Greek faith.'' ``The disgrace attaching to this Council consists in the fact that the great majority of the bishops who held the same views as Cyril and Dioscorus finally allowed a formula to be forced upon them, which was that of strangers, of the Emperor and the Pope, and which did not correspond to their belief.''\texthebrew{⁠}\footnote{Harnack, \emph{ib.} 196, 214. It is worth observing that the majority of the bishops of the Asiatic provinces absented themselves.

} But the truth is that the definition of Chalcedon might be interpreted in different ways. To Leo and the Western Church it meant one thing; to the followers of Cyril another; to

p358
Antiochians and Theodoret, something different which Nestorius himself could have accepted.\footnote{Loofs, \emph{op. cit.} p99. Duchesne (\emph{Hist. anc. de l'Église}, III.449) remarks: ``Vive la doctrine de Flavien et de Léon ! Anathème à Nestorius ! C'est tout le concile de Chalcédoine.''

}

Politically, the Council was a decisive triumph for Constantinople and a final blow to the pretensions of the see of Alexandria. Marcian completed what Theodosius the Great had begun. Three successive Patriarchs, Theophilus, Cyril, and Dioscorus, had aimed at attaining to the supreme position in Eastern Christendom and at ruling Egypt like kings. Alexandria could never again claim to lead the Church in theology. But the defeat of Alexandria was accompanied by an exaltation of Byzantium which was far from acceptable to Rome. By the twenty-eighth Canon equal privileges with Rome were granted to the see of Constantinople, all the episcopal sees of the Dioceses of Thrace, Asia, and Pontus were assigned to the jurisdiction of the Patriarch. The Roman legates protested against this Canon, and Leo refused to confirm it.\footnote{The change was bitterly felt by Ephesus, the premier see of Asia Minor, associated as it was with early apostolic history, the memories of Paul and Timothy, and of John the Evangelist, who was said to have died and been buried there. The archbishop lost not only his independence but even his rank, for he was placed second to the metropolitan of Cappadocian Caesarea. It was hardly much consolation that he was allowed the title of ``Exarch of the Diocese of Asia.''

}

Dioscorus was deposed by the Council, and was banished to Gangra. Feeling ran so high at Alexandria that the aid of soldiers was required to establish his successor Proterius.

In Egypt and Syria there was a solid mass of opinion loyal to the doctrine of one nature, and firmly opposed to the formula of Chalcedon. These Monophysites, as they were called, were far too numerous and earnest to be stamped out; they ultimately created the national Coptic Church of Egypt. Throughout the reign of Leo I the dispute over the meaning of the Incarnation led to scenes of the utmost violence in Alexandria and to occurrences hardly less scandalous in Antioch.

At Jerusalem the Monophysites obtained the upper hand after the Council of Chalcedon, and a reign of terror prevailed for some time. The episode derives interest from the association of the Empress Eudocia, who was living there in retirement, with the Monophysitic cause.\texthebrew{⁠}\footnote{See Cyrillus, \emph{Vita Euthymii}, p64 \emph{sqq.} Génier, \emph{Vie de Saint Euthyme}, 209 \emph{sqq.}

} A monk named Theodosius,

p359
who was a zealous supporter of Dioscorus, gained the ear of the people, and the bishop of Jerusalem, Juvenal, when he returned from the Council, was forced to flee for his life, because he refused to renounce the doctrine which he had subscribed. Theodosius was ordained bishop, and methods of the utmost violence were adopted to coerce those who refused to communicate with him. He was supported by Eudocia, who had been a devoted admirer of Cyril and was led to believe that Cyril's doctrine was identical with that of Dioscorus and had been condemned at Chalcedon. The Emperor Marcian at length took strong measures; Theodosius fled to Mount Sinai, and Juvenal was restored to his see.\texthebrew{⁠}\footnote{The usurpation of Theodosius lasted for 20 months, A.D. 452\texthebrew{‑}453.

Evagrius, II.5 . Theophanes, A.M. 5945.

} Eudocia after some years began to feel doubts about her theology and she consulted the pillar saint, Simeon, who recommended her to seek the advice of Euthymius, abbot of the convent of Sahel, a few miles east of Jerusalem. An interview with the monk showed the Empress the error of her ways, and she died in the faith of Chalcedon.

The Christian religion, with its theology which opened such a wide field for differences of opinion, had introduced into the Empire dangerous discords which were a sore perplexity to the government. In some ways it augmented, in others it weakened, the power of the State to resist its external enemies. It cannot be maintained \texthebrew{—} as we have already seen \texthebrew{—} that it was one of the causes which contributed to the dismemberment of the Empire in the West by the Teutonic peoples; and subsequently, the religious communion, which was preserved throughout political separation, helped the Empire to recover some of the territory it had lost. In the East, bitter theological divisions, consequent on the Council of Chalcedon,\texthebrew{⁠}\footnote{Gelzer (in Krumbacher, \emph{G. B. L.}, p919) designates the decisions of Chalcedon, regarded from a political aspect, as a most grievous misfortune for the east-Roman Empire.

} facilitated the Saracen conquest of the provinces of Syria and Egypt, but afterwards, in the diminished Empire, the State religion formed a strong bond and fostered the growth of a national spirit which enabled the Imperial power to hold out for centuries against surrounding foes.

The subtle questions on the nature of the Incarnation, which were so hotly disputed by the Greeks and Orientals, created little or no disturbance in western Europe. But in the early years of the fifth century the western provinces were agitated by a heresy of their own, on a subject which had more obviously practical bearings, but involved no less difficult theological metaphysics. The Pelagian controversy concerned free will and original sin. Pelagius, probably a Briton of Irish extraction,\texthebrew{⁠}\footnote{Jerome, \emph{Comm. in Jerem.}, \emph{P. L.} XXIV.680\texthebrew{‑}682, 757\texthebrew{‑}758. Cp. Bury, ``The Origin of Pelagius,'' in \emph{Hermathena}, XXX.26 \emph{sqq.} On the controversy, see Pelagius, \emph{Letter to Demetrias}, \emph{P. L.} XXXII.1100; the numerous writings of Augustine on the subject, \emph{P. L.} XLVII; Marius Mercator, \emph{Commonitorium super nomine Caelestii}, \emph{ib.} XLVIII; Jerome's Three Books \emph{Adversus Pelagium}, and his \emph{Letter to Ctesiphon} (\emph{P. L.} XXII.1152). A Commentary by Pelagius on the Pauline epistles existed in Ireland in the Middle Ages (Zimmer, \emph{Pelagius in Ireland}, 1901). The Patriarch Nestorius wrote treatises against Pelagianism of which Latin translations by Marius Mercator are preserved (\emph{P. L.} XLVIII).

} propagated the views that man possesses the power of choosing between good and evil, and that there is no sin where there is not a voluntary choice of evil; that sin is not inherited; that man can live, and some men actually have lived, sinless; and that unbaptized infants attain to eternal life.\texthebrew{⁠}\footnote{Pelagius distinguished eternal life from the bliss of Paradise.

} The controversy is memorable because these doctrines found their chief antagonist in Augustine and led him gradually to develop the predestinarian theories which had such a power ­ful influence on subsequent theology. He maintained that sin was transmitted to all men from Adam; that man, by the mere gift of free will, cannot choose aright without the constant operation of grace; that no man has ever lived a sinless life; that infants dying unbaptized are condemned, as a just punishment for the sin which they inherited. As time went on, Augustine developed his theory, which raised the whole question of the origin of evil into a system which, while it professed to admit the freedom of the will, really annulled it. God, he said, decided from eternity to save some members of the human race from the consequence of sin; he fixed the number of the saved, which can be neither increased nor diminished, and on these favoured few he bestows the gifts of grace which are necessary for their salvation. The rest perish eternally, if not

p361
through their own transgressions, through the effects of original sin. This is not unjust, because there is no reason why God should give grace to any man; by refusing to bestow it, he affirms the truth that none deserve it. Augustine allowed that in the eternal punishment which awaits all but the few there may be different degrees of pain.

Pelagius, along with his friend Caelestius whom he had converted to his views, went from Rome to Africa ( A.D. 409). Leaving Caelestius there, he proceeded himself to Palestine. Caelestius stated his views before a council of African bishops at Carthage and was excommunicated ( A.D. 412). Three years later a synod was held at Jerusalem, at which Pelagius was present, the question was discussed, and it was decided that it should be referred to Pope Innocent I ( A.D. 415), but some months later another synod at Diospolis acquitted Pelagius of heterodoxy. In the meantime Augustine was writing on the subject,\texthebrew{⁠}\footnote{\emph{De peccatorum meritis}; \emph{de natura et gratia}; and \emph{de perfectione iustitiae hominis} (A.D. 415).

} and the African bishops condemned the Pelagian doctrine and asked Innocent to express his approval.\texthebrew{⁠}\footnote{At the Synods of Carthage and Milevis (A.D. 416). Innocent replied (Jan. 417), condemning the heresy in strong language. The correspondence will be found in Innocent, \emph{Epp.} 26\texthebrew{‑}31, \emph{P. L.} XX.

} A decision on the matter devolved upon Innocent's successor Zosimus, who was elected on March 17, A.D. 417, and the ear of this Pope was gained by Caelestius, who had come to Rome. Zosimus censured the African bishops for condemning Caelestius, and intimated that he would decide, if the accusers came and appeared before him. Then he received a letter from Pelagius, which convinced him that Pelagius was a perfectly orthodox Catholic.\texthebrew{⁠}\footnote{Zosimus, \emph{Epp.} 2 and 3 (\emph{P. L.} XX.649, 654).

} But the African bishops were not convinced, and in defiance of the Pope's opinion, they condemned Pelagius and his teaching in a synod at Carthage (May 1, A.D. 418). Zosimus at last became aware that the doctrines of Pelagius were really heretical; he was obliged to execute a retreat,\texthebrew{⁠}\footnote{\emph{Id.}, \emph{Ep.} 12 (\emph{ib.} 675). In this letter he says quamvis patrium traditio apostolica sedi auctoritatem tantam tribuerit ut de eius iudicio disceptare nullus auderet.

} and he confirmed the findings of the African synod. Honorius issued a decree banishing Pelagius and Caelestius from Rome and inflicting the penalty of confiscation on their followers.\texthebrew{⁠}\footnote{Cp. Maassen, \emph{Gesch. der Quellen und der Lit. des kanonischen Rechts}, I p316. Caelestius was condemned at the Council of Ephesus in 431.

} Although the views of the British heretic were crushed by the

p362
arguments and authority of Augustine, they led to the formation of an influential school of opinion in Gaul\texthebrew{⁠}\footnote{Known as Semipelagians.

Thayer's Note: For an exhaustive treatment, see the article

Semipelagianism

in the Catholic Encyclopedia.

} which, though condemning Pelagianism, did not accept the extreme predestinarian doctrines of the great African divine.

In the list of Roman pontiffs the name of Zosimus is not one which the Catholic Church holds in high esteem. His brief pontificate fell at a critical period, when the Roman see was laying the foundations of the supremacy which it was destined to gain by astute policy, and propitious circumstances, over the churches of western Europe. Zosimus, through his rashness and indiscretion, did as much as could be done in two years to thwart the purposes which he was himself anxious to promote. In the matter of Pelagius he committed himself to a judgment which shows that he was either unpardonably ignorant of the doctrine which had been challenged, or that he considered orthodox in A.D. 417 what he condemned as heterodox in A.D. 418; and he exposed himself to a smart rebuff from the bishops of Africa.\texthebrew{⁠}\footnote{In another African affair he also compromised himself. A priest of Sicca, deposed by his bishop, appealed to Rome. Zosimus, in demanding his reinstatement, based his action on a canon which he alleged to be Nicene. The African bishops were unable to discover it among the canons of the Council of Nicaea. It was really a canon of the Council of Sardica.

} But his indiscretion in this affair was of less importance than the ill-considered policy on which he embarked on a question of administration in the Gallic Church, and which proved highly embarrassing to his successors.

The authority which the Roman see exercised in western Europe at this time, beyond its prestige and acknowledged primacy in Christendom, was twofold. Decrees of Valentinian I and Gratian had recognised it as a court to which clergy condemned by provincial synods might appeal.\texthebrew{⁠}\footnote{\emph{Epp. Impp. Pontt.}, ed. Günther, I p58 = \emph{P. L.} XIII.587; it was provided that the appeal might also be addressed to a council of fifteen neighbouring bishops.

} In the second place it was looked up to as a model, and when doubtful questions arose about discipline it was consulted by provincial bishops. The answers of the Popes to such questions were known as Decretals. They did not bind the bishops; they were responses, not ordinances. Appellate jurisdiction and the moral weight of the Decretals were the principal bases on which the power of the Roman see was gradually to be built up.\footnote{It may be added that Roman excommunication was recognised as exclusion from Catholic communion. Boniface, \emph{Ep.} 14, \emph{P. L.} XX.777. Cp. Babut, \emph{Le Concile de Turin}, p75.

}

p363
Zosimus entertained an idea of his authority which transcended these rights and anticipated the claims of his successors. Immediately after his election his ear was gained by Patroclus, the bishop of Arles, who desired to make his see an ecclesiastical metropolis of the first rank. In the three provinces of Viennensis, Narbonensis Prima, and Narbonensis Secunda, the bishops of Vienne, Narbonne, and Marseilles\texthebrew{⁠}\footnote{Marseilles was an exception to the rule that the civil was also the ecclesiastical metropolis. The civil metropolis of Narbonensis II was Aquae (Aix). Marseilles was in Narbonensis I.

} were the metropolitans; Arles was merely a bishopric in Narbonensis Prima. The idea of Patroclus was naturally enough suggested by the translation of the residence of the Praetorian Prefect of Gaul from Trier to Arles.\texthebrew{⁠}\footnote{\emph{C.} A.D. 413, cp. Mommsen, \emph{Chron. Min.} I p553 . Patroclus was a friend of the Patrician, afterwards Emperor, Constantius, and doubtless had the support of his influence. Prosper, \emph{sub} 412.

} Zosimus determined to deprive the bishops of Vienne, Narbonne, and Marseilles of their metropolitan rights, and to invest the bishop of Arles with jurisdiction over the three provinces. He also proposed to establish a new Metropolitan of Arles as a sort of Roman vicar, apparently over the whole of Gaul.\footnote{See Zosimus, \emph{Ep.} 1, \emph{P. L.} XX.642, addressed universis episcopis per Gallias et septem provincias constitutis.

}

The bishop of Narbonne yielded with a protest to this revolutionary assumption of sovranty. But the bishops of Marseilles and Vienne defied Zosimus and brought the question before a council of the Milanese diocese which met at Turin (Sept. 22, A.D. 417).\texthebrew{⁠}\footnote{The difficulties about this Council, its date, and its importance, were first elucidated by Babut, \emph{op. cit.} For the history of the struggle over the º Arles see also Gundlach, \emph{Der Streit der Bistümer Arles und Vienne}, 1890.

} The council at first decided against the pretensions of Arles, but finally compromised by dividing the Viennese province into two parts, of which the southern was to depend on Arles. Zosimus was not pleased, but deemed it prudent to concur. The bishop of Marseilles, who declined to yield, was excommunicated by a Roman synod, but remained quietly in his see. Thus a part of the Pope's plan was actually carried out, but the facts remained that the council of Turin had refused to recognise the supreme authority of Rome, and that Marseilles had resisted with impunity.

The indiscretions of Zosimus were a lesson for his successors.\footnote{After his death the bishops of Africa, in a letter to Pope Boniface, expressed a hope that they would not again be exposed to such arrogance, non sumus iam istum typhum passuri, \emph{P. L.} XX.752.

}

p364
Moreover, they recognised that the establishment of such a large and power ­ful see as that which Zosimus called into being was likely to be a rival rather than a vassal of Rome. Their aim was to undo what Zosimus had done, and in accomplishing this they acted with greater circumspection and increased the authority of their see. Both Boniface and Celestine\texthebrew{⁠}\footnote{Boniface, 418\texthebrew{‑}422; Celestine, 422\texthebrew{‑}432.

} did what they could to restrict the powers of the bishop of Arles. The first Narbonensis was withdrawn from his jurisdiction and restored to Narbonne.\texthebrew{⁠}\footnote{By Boniface. For Celestine's attitude see his letter, \emph{Ep.} 4, \emph{P. L.} L.429.

} But the situation was more difficult for Rome, because the monks of Lérins, whose influence was strong in southern Gaul, threw the weight of their interest into the scale of Arles. Their founder, Honoratus, had been elected to succeed Patroclus, and he was followed by his disciple Hilary, whose authority threatened to usurp that of Rome in the Gallic Church.\texthebrew{⁠}\footnote{Babut, \emph{op. cit.} 147 \emph{sqq.}

} The conflict between Hilary and Leo I, who was elected in A.D. 440, is not edifying. An appeal to Rome ( A.D. 444) gave the Pope a welcome opportunity of striking his opponent. He did not venture to excommunicate him, but he deprived him of the remnant of the province which Zosimus had created. This sentence could not be executed without the aid of the secular power. He had much influence with the Emperor and Galla Placidia, and he procured an edict, which was issued (July 8, A.D. 445) at the same time as his own decree.\texthebrew{⁠}\footnote{Valentinian III,

\emph{Nov.} 17 . Leo, \emph{Ep.} 10, \emph{P. L.} LIV.628. Cp. Tillemont, \emph{Mém.} 15, 82; Babut, p172.

} Arles was deprived of its metropolitan dignity.\footnote{Five years later, however, Leo restored this rank to Arles, giving it a part of the Viennese diocese. This was after Hilary's death.

}

But that edict of Valentinian III did much more than settle in Rome's favour this particular question. It assigned to the Roman see that supremacy over the provincial churches which the Popes had been endeavouring to establish, but which the African synods and the council of Turin had refused to acknowledge.\texthebrew{⁠}\footnote{The unwillingness of leading churchmen at the beginning of the fifth century to admit the exorbitant claims of Rome is illustrated by Jerome's letter to Evangelus, \emph{Ep.} 146, \emph{P. L.} XXII p1194; he observes, orbis maior est urbe.

} It ordained that ``the bishops of Gaul or any other province should take no decision contrary to the ancient rules of discipline without the consent and authority of the venerable Pope of the eternal city. They must conform to all the decrees

p365
of the Apostolic see. Bishops summoned before the tribunal of Rome must be compelled to appear by the civil authorities.''

It is the political bearing of this law that interests us here. When many of the western provinces had wholly or partly passed out of the Emperor's control, it was a matter of importance to strive to keep alive the idea of the Empire and the old attachment to Rome in the minds of the provincials who were now subject to German masters. The day might come when it would be possible to recover some of these lost lands, which the Imperial government never acknowledged to be really lost, and in the meantime a close ecclesiastical unity presented itself as a power ­ful means for preserving the bonds of sentiment, which would then prove an indispensable help. To accustom the churches in Gaul and Britain, Spain and Africa to look up to Rome and refer their disputes and difficulties to the Roman bishop was a wise policy from the secular point of view, and it was doubtless principally by urging considerations of this nature that Leo was able to induce the government to establish the supremacy of his see.

It is important to bear in mind that the administrative authority of the Pope, at this time, extended into the dominions of the eastern Emperors. The lands included in the Prefecture of Illyricum belonged to the Patriarchate of Rome, and constituted the Vicariate of Thessalonica, where the Pope's vicar, who was entrusted with the administration, resided. Theodosius II wished to place this ecclesiastical province under Constantinople and published an edict with this intent, but the remonstrances of Honorius induced him to retract it;\texthebrew{⁠}\footnote{See above

Chap. II, p64 . A.D. 421,

\emph{C. Th.} XVI.11.45 . Cp. Innocent, \emph{Ep.} 13. Gieseler, \emph{Lehrbuch}, II.217.

} and Greece, Macedonia, and Dacia remained under the see of St. Peter till the eighth century.

Persecution was an unavoidable consequence of Constantine's act in adopting Christianity. Two of the chief points in which this faith differed from the Roman State religion were its exclusiveness and the vital importance which it assigned to dogma. The first logically led to intolerance of pagan religions, the second to intolerance of heresies, and these consequences could not be

p366 averted when Christianity became the religion of the State. It might be suggested that Constantine would have done better if, when he decided to embrace it and favour its propagation, he had been content to deprive pagan cults of their official status and to allow Christianity to compete in a free field with its rivals, aided by the prestige which it would derive from the Emperor's personal adhesion and favour. But such a policy would have been an anachronism. A state, at that time, was unthinkable without a State cult, and if an Emperor became a Christian a logical result was that Christianity should be adopted as the official religion of the Empire, and a second that the old Roman policy of toleration should be thrown overboard. In an age of superstition this was demanded not only in the interest of the Church but in the interest of the State itself. The purpose of the official cults in the pagan State was to secure the protection of the deities; these were liberal and tolerant lords who raised no objection to other forms of worship; and toleration was therefore a principle of the State. But the god of the new official religion was a jealous master; he had said, ``thou shalt have none other gods before me,'' and idolatry was an offence to him; how could his protection and favour be expected in a state in which idolatry was permitted? Intolerance was a duty, and the first business of a patriotic ruler was to take measures to extirpate the errors of paganism.

But these consequences were not drawn immediately. It must never be forgotten that Constantine's revolution was perhaps the most audacious act ever committed by an autocrat in disregard and defiance of the vast majority of his subjects. For at least four-fifths of the population of the Empire were still outside the Christian Church.\texthebrew{⁠}\footnote{Estimates, based on highly conjectural data, of the number of Christians vary from one-twentieth to one-sixth of the total population. See V. Schultze, \emph{Der Untergang des gr.-röm. Heidentums}, I p22 \emph{sqq.}

} The army and all the leading men in the administration were devoted to paganism. It is not, therefore, surprising that Constantine, who was a statesman as well as a convert, made no attempt to force the pace. His policy did little more than indicate and prepare the way for the gradual conversion of the Empire, and was so mild and cautious that it has been maintained by some that his aim was to establish a parity between the two religions.

p367
He retained the title of Pontifex Maximus, and thereby the constitutional right of the Emperor to supervise the religious institutions. He withdrew the support of state funds from pagan rites, but made an exception in favour of the official cults at Rome. His most important repressive measure was the prohibition of the sacrifice of victims in the temples.\texthebrew{⁠}\footnote{The law is not preserved, but is recorded by Eusebius, \emph{Vita Const.} II.45, and referred to by Constantius,

\emph{C. Th.} XVI.10.2 .

} One reason for this measure was the dangerous practice of divination by entrails, often employed by persons who contemplated a rebellion and desired to learn from the higher powers their chances of success.

In some particular places cults were suppressed, but a pagan could still worship freely in the temples, could offer incense and make libations of wine, and might even perform sacrificial rites in a private house. The sons of Constantine\texthebrew{⁠}\footnote{Firmicus Maternus, in his \emph{De errore profanarum religionum}, urged them to drastic measures.

} were indeed inclined to adopt a stringent policy, and their laws might lead us to suppose that there was something like a severe persecution. Constantius, in reaffirming the prohibition of sacrifices, menaced transgressors with the avenging sword.\texthebrew{⁠}\footnote{\emph{C. Th.} XVI.10.4

(A.D. 342) gladio ultore sternatur.

} But the death penalty was never inflicted, and there was a vast difference between the letter of the law and the practice. In the same edict was ordained the closing of temples ``in all places and cities,'' but this order can only have been carried out here and there. Its execution depended on local circumstances, and on the sentiments of the provincial governors. In some places Christian fanatics took advantage of the Imperial decree to demolish heathen shrines, and the pagans were naturally very apprehensive. When Julian visited Ilion, he inspected the antiquities under the guidance of Pegasius, who was ``nominally a bishop of the Galilaeans,'' but really worshipped the Sun god.\texthebrew{⁠}\footnote{Julian, \emph{Ep.} 78, ed. Hertlein.

} He had taken orders and succeeded in becoming a bishop in order that he might have the means of protecting the heathen sanctuaries from Christian destruction.

When paganism was restored by Julian, it is probable that any temples which had been closed under the edict of Constantius were again reopened, and after his fall it would seem that they were allowed to remain open for worship, though sacrifices were regarded as unlawful.

p368
The Emperors Valentinian I\texthebrew{⁠}\footnote{Ammianus Marc. XXX.9

inter religionum diversitates medius stetit.

} and Valens were consistently tolerant. The mysteries of Eleusis were expressly permitted, for the proconsul of Achaia told Valentinian that if they were suppressed the Greeks would find life not worth living.\texthebrew{⁠}\footnote{Zosimus, IV.3.

} But a new religious policy was inaugurated by Gratian and Theodosius the Great. Gratian abandoned the title of Pontifex Maximus; he withdrew the public money which was devoted to the cults of Rome, and he ordered the altar of Victory to be removed from the Senate-house , to the deep chagrin of the senators. The fathers appealed to Valentinian II to revoke this order, and to restore the public maintenance of the religious institutions of the capital; but the moving petition of Symmachus, who was their spokesman, was overruled by the influence of Ambrose, the archbishop of Milan, who possessed the ear of Valentinian and of Theodosius.\footnote{Ambrose, \emph{Epp.} I.17 and 18 (\emph{P. L.} XVI.961 and 971). Symmachus, Relatio 3. Prudentius, \emph{Contra Symmachum} . Gracchus, Prefect of Rome in 376, demolished a cave-temple of Mithras at Rome (Jerome, \emph{Ep.} 107 \emph{ad Laetam}, \emph{P. L.} XXII.868; Prudentius, \emph{ib.} I.561 \emph{sqq.} ).

}

It remained for Theodosius to inflict a far heavier blow on the ancient cults of Greece and Rome. In the earlier years of his reign the extirpation of pagan worship does not seem to have been an aim of his policy. He was only concerned to enforce obedience to the laws prohibiting sacrifices, which had evidently been widely evaded. He decided on the closing of all sanctuaries in which the law had been broken. He entrusted to Cynegius, Praetorian Prefect of the East, a pious Christian, the congenial task of executing this order in Asia and Egypt. But otherwise temples were still legally open to worshippers.\texthebrew{⁠}\footnote{See Libanius, \emph{Or.} XXX § 8 (ed. Förster). This appeal which Libanius addressed to the Emperor on behalf of the temples was written in summer A.D. 388, as R. van Loy has satisfactorily shown (\emph{B. Z.} XXII.313 \emph{sqq.}). The orator refers to the campaign conducted by Cynegius, who had recently died.

} It is to be particularly noted that the Emperor did not desire to destroy but only to secularise such buildings as were condemned, and the cases of barbarous demolition of splendid buildings which occurred in these years were due to the fanatical zeal of monks and ecclesiastics. Monks wrought the destruction of the great temple of Edessa, and the Serapeum at Alexandria, which gave that city ``the semblance of a sacred world,''\texthebrew{⁠}\footnote{Eunapius, \emph{Vita Aedesii}, p43.

} was demolished

p369
under the direction of the archbishop Theophilus ( A.D. 389),\texthebrew{⁠}\footnote{The account of Sozomen, VII.15, is better than that of Socrates, V.16, 17. See also Eunapius, \emph{ib.} The pagans were not guiltless in this affair. They had attacked the Christians and fortified themselves in the buildings of the Serapeum; but they had been provoked to this outbreak by Theophilus, who had paraded religious symbols, taken from a temple of Dionysus (which the Emperor had permitted him to convert into a church), through the streets in derision of the pagan cults. The most unfortunate occurrence was the destruction of the library of the Serapeum

(Orosius, VI.15) .

} who thereby dealt an effective blow to the paganism of Alexandria.

But Theodosius and his ecclesiastical advisers thought that the time was now ripe to make a clean sweep of idolatry, and in A.D. 391 and 392 laws were issued which carried to its logical conclusion the act of Constantine. We may conjecture that this drastic legislation was principally due to the influence of the archbishop of Milan. To sacrifice, whether in public or in private, was henceforward to be punished as an act of treason. Fines were imposed on any who should frequent temples or shrines; and for worshipping images with incense, for hanging sacred fillets on trees, for building altars of turf, the penalty was confiscation of the house or property where such acts were performed.\footnote{\emph{C. Th.} XVI.10.10 and 11

(391);

12

(392).

}

In the insurrection of A.D. 392 the restoration of paganism was a capital feature in the programme of the general Arbogastes and Eugenius the creature whom he crowned, and the lure attracted some distinguished adherents. For a short time the altar of Victory was set up in the Roman Senate-house . After the suppression of the revolt Theodosius visited Rome, attended a meeting of the Senate, and though his tone was conciliatory, his firmness compelled that body to decree the abolition of the ancient religious institutions of Rome.\texthebrew{⁠}\footnote{A.D. 394. On the debate in the Senate see

Zosimus IV.59 ; Prudentius, \emph{Contra Symm.} I.415 \emph{sqq.} ; Hodgkin, \emph{Italy}, I.580 \emph{sqq.}

} Some of the pagan senators had Christian families,\texthebrew{⁠}\footnote{\emph{E.g.} Albinus, a pontifex. Jerome, \emph{Ep.} 107 \emph{ad Laetam} (\emph{P. L.} XXII.868). As to the small number of Christian senators cp. above,

p164 .

} and domestic influence may have reinforced the imperial will.

The last years of the fourth century mark an epoch in the decay of paganism. While the gods were irrevocably driven from Rome itself, time-honoured institutions of Greece also came to an end. The old oracles seem to have been silenced at

p370
a much earlier date. The ``last oracle'' of the Delphic god, said to have been delivered to Julian, is a sad and moving expression of the passing away of the old order of things.

The Olympian games were celebrated for the last time in A.D. 393, and the chryselephantine statue of Zeus, the greatest monument of the genius of Pheidias, was removed soon afterwards from Olympia to Constantinople.\texthebrew{⁠}\footnote{Cedrenus I.364
(cp. Moses of Chorene, III.40 ); Clinton, \emph{Fasti Hellenici}, III. p. xv. A passage of Julian (\emph{Ep.} 35) seems to imply that the Pythian, Nemean, and Isthmian games were celebrated in his day. Cp. \emph{C. Th.} XV.5.4

for games at Delphi (A.D. 424), and there is a record that the Olympian games came to an end in the reign of Theodosius II (\emph{Scholia in Lucianum, Praec. Rhet.} ed. Rabe, p174).

} The Eleusinian mysteries ceased three years later in consequence of the injuries wrought to the sanctuaries by the invasion of Alaric.\texthebrew{⁠}\footnote{Eunapius (\emph{Vita Maximi}) suggests that the destruction was wrought by a band of fanatical monks who accompanied the Gothic army.

} The legend that Athens was saved from the rapacity of the Goths by the appearance of Athene Promachos and the hero Achilles illustrates the vitality of pagan superstition. Athens had fared better than many other towns at the hands of the Emperors.\texthebrew{⁠}\footnote{Gregorovius, \emph{Gesch. d. Stadt Athen}, I.26.33.

} Constantine, who ransacked Hellenic shrines for works of art in order to adorn his new capital, spared Athens; and in the reign of Theodosius, when the Samian Hera of Lysippus, the Cnidian Aphrodite of Praxiteles, the Athene of Lindos were carried off, the Parthenon was not compelled to surrender the ivory and gold Athene of Pheidias. Soon after A.D. 429 this precious work was ravished from the Acropolis,\texthebrew{⁠}\footnote{Marinus, \emph{Vita Procli}, c30.

} but we do not know its fate. Nor do we know at what date the Parthenon was converted into a church of the Virgin.\footnote{Gregorovius, \emph{ib.} 64, conjectures in the reign of Justinian.

}

The ordinances of Theodosius did not, of course, avail immediately to stamp out everywhere the forbidden cults. Pagan practices still went on secretly, and in some places openly,

p371
and the government, generally perhaps yielding to ecclesiastical pressure, issued from time to time new laws to enforce the execution of the old or to supplement them.\texthebrew{⁠}\footnote{\emph{C. Th.} XVI.10.13

(395);

XVI.10.14

(396) abolishing some immunities still enjoyed by old priesthoods.

} Arcadius, under the influence of Chrysostom, issued an edict to destroy, not merely to close, temples in the country and to use the material for public buildings.\texthebrew{⁠}\footnote{\emph{Ib.} XVI.10.16.

} Chrysostom sent monks to Phoenicia to carry out the work of destruction there, but the money required was provided not by the state but by pious Christians, especially women.\texthebrew{⁠}\footnote{Theodoret V.9.

} We have seen how bishop Porphyrius of Gaza secured with the help of the Empress Eudoxia the demolition of the temple of Marnas. As a rule the Emperors desired that the ancient sanctuaries should be preserved and turned to other uses, and we find them interfering to prevent destruction.\texthebrew{⁠}\footnote{In southern Gaul,

\emph{C. Th.} XVI.10.15 ; in Africa, \emph{ib.} XVIII (399).

} In many country districts Christianity was only beginning to penetrate, and for the eradication of heathenism there was much missionary teaching to be done, such as was carried on by Martin in western Gaul, by Victricius, archbishop of Rouen, in the Belgic provinces, and by Nicetas of Remesiana in the Balkan highlands.\footnote{Sulpicius Severus, \emph{Dialogus}, III.2. Vacandard, \emph{Saint Victrice}, 1903. Burns, \emph{Life and Works of Nicetas of R.}, 1905.

}

Theodosius II at one time professed to believe that no pagans survived in his dominions,\texthebrew{⁠}\footnote{A.D. 423,

\emph{C. Th.} XVI.10.22 .

} but this sanguine view, if it was seriously held, was premature, for in a later year he repeated the prohibition of sacrifices and ordered anew the conversion of temples into churches;\texthebrew{⁠}\footnote{A.D. 435, \emph{ib.} 25.

} and Leo I legislated severely against heathen practices.\texthebrew{⁠}\footnote{\emph{C. Th.} I.11.8 ; subsequent laws against Hellenism by Leo, Zeno, or Anastasius (?), \emph{ib.} 9.10.

} It is to be observed that this persecution differed in one important respect from the ecclesiastical persecutions of later ages in western Europe. Only pagan acts were forbidden; \emph{opinion as such was tolerated}, and no restrictions were placed on the diffusion of pagan literature. Perhaps the only exception was the edict of Theodosius II shortly before his death,\texthebrew{⁠}\footnote{A.D. 448, \emph{ib.} I.1.3. The law says, the books of Porphyry ``or any one else.'' The anti-Christian work of Porphyry has perished, like those of Celsus and Julian. There is a new edition of the fragments by Harnack, \emph{Porphyrius ``Gegen die Christen''} 1916 (\emph{Abh.} of the Prussian Academy).

} ordering the books of Porphyry, whose dangerous

p372
treatise \emph{Against the Christians} had apparently shocked the Emperor or some of his advisers, to be burned. The same monarch had enacted that no Christian shall disturb or provoke Jews or pagans ``living peaceably.''\texthebrew{⁠}\footnote{In quiete degentibus,

\emph{C. Th.} XVI.10.24 .

} Indeed pagans could not be dispensed with in the civil service, and in the sixth century we still find them in prominent positions.\texthebrew{⁠}\footnote{A law issued at Ravenna in 408 excluded enemies of the Catholic faith from serving in the Palace, but was probably applied only temporarily.

\emph{C. Th.} XVI.5.42 . In 416 persons polluted with errors of pagan rites were excluded from state service,

\emph{ib.} XVI.10.21 , but this would not affect those who had not been found guilty of sacrifi­cing.

} Hellenism largely prevailed in the law schools, and was no bar to promotion, though it might be made a pretext for removing an official who had fallen out of favour. An able pagan, Tatian, enjoyed the confidence of the fanatical Theodosius the Great, and was appointed Praetorian Prefect of the East; and the same Emperor showed friendly regard towards spokesmen of the old religion like Libanius and Symmachus. The headquarters of unchristian doctrine, the university of Athens, was held in high esteem by Constantine and Constans,\texthebrew{⁠}\footnote{Gregorovius, \emph{ib.} 28, 29.

} and it continued throughout the fifth century unmolested as the home of a philosophy which was the most dangerous rival of Christian theology. Pagans also received appointments in the university of Constantinople.

In a hundred years the Empire had been transformed from a state in which the immense majority of the inhabitants were devoted to pagan religions, into one in which an Emperor could say, with gross exaggeration, but without manifest absurdity, that not a pagan survived. Such a change was not brought to pass by mere prohibition and suppression. It is not too much to say that the success of the Church in converting the gentile world in the fourth and fifth centuries was due to a process which may be described as a pagan transmutation of Christianity itself. If Christian beliefs and worship had been retained unaltered in the early simplicity of their spirit and form, it may well be doubted whether a much longer period would have sufficed to christianize the Roman Empire. But the Church permitted a compromise. All the religions of the age had common ground in crude superstition, and the Church found no difficulty in proffering to converts beliefs and cults similar to those to which they had been accustomed. It was a comparatively small

p373
matter that incense, lights, and flowers, the accessories of various pagan ceremonials, had been introduced into Christian worship. It was a momentous and happy stroke to encourage the introduction of a disguised polytheism. A legion of saints and martyrs replaced the old legion of gods and heroes, and the hesitating pagan could gradually reconcile himself to a religion, which, if it robbed him of his tutelary deity, whom it stigmatized as a demon, allowed him in compensation the cult of a tutelary saint. A new and banal mythology was created, of saints and martyrs, many of them fictitious; their bodies and relics, capable of working miracles like those which used to be wrought at the tombs of heroes, were constantly being discovered. The devotee of Athene or Isis could transfer his homage to the Virgin Mother. The Greek sailor or fisherman, who used to pray to Poseidon, could call upon St. Nicolas. Those who worshipped at stone altars of Apollo on hill-tops could pay the same allegiance to St. Elias. The calendar of Christian anniversaries corresponded at many points to the calendars of Greek and Roman festivals. Men could more easily acquiesce in the loss of the heathen celebrations connected with the winter solstice and the vernal equinox, when they found the joyous celebrations of the Nativity and the resurrection associated with those seasons, and they could transfer some of their old customs to the new feasts. The date of the Nativity was fixed to coincide with the birthday of Mithras ( natalis Invicti , December 25), whose religion had many affinities with the Christian. This process was not the result, in the first instance, of a deliberate policy. It was a natural development, for Christianity could not escape the influence of the ideas which were current in its environment. But it was promoted by the men of light and leading º in the Church.\footnote{On the origins of the cult of saints and martyrs see E. Lucius, \emph{Die Anfänge des Heiligenkults in der christlichen Kirche}, 1904; P. Saintyves, \emph{Les Saints successeurs des dieux}, 1907; J. Rendel Harris, \emph{The Dioscuri in the Christian Legends}, 1903, \emph{Cult of the Heavenly Twins}, 1906.

}

A particular form of miraculous healing illustrates the way in which Christianity appropriated pagan superstitions. The same dream-cures which used to be performed by Aesculapius or the Dioscuri for those who slept a night in the temple courts were still available; only the patient must resort to a sanctuary of Saints Cosmas and Damian,\texthebrew{⁠}\footnote{Known as the anargyroi, physicians who take no fee. For their miracles see Zeumer, \emph{De incubatione}, 69 \emph{sqq.}, where the whole subject is treated.

} the new Castor and Pollux, or of the archangel

p374
Michael\texthebrew{⁠}\footnote{In his church at Sosthenion on the Bosphorus, Sozomen, II.3.

} or some other Christian substitute. We have an interesting example of the method employed by ecclesiastical magnates in an incident which occurred in Egypt. Near Canopus there was a temple of Isis where such nocturnal cures were dispensed, and professing Christians continued to have recourse to this unhallowed aid. The Patriarch Cyril found a remedy. He discovered the bodies of two martyrs, Cyrus and John, in the church of St. Mark at Alexandria, and dislodging Isis he interred them, and dedicated a church to them, in the same place, where they freely exhibited the same mysterious medical powers which had been displayed by the great goddess.\footnote{See their Acts in \emph{P. G.} LXXXVII.3.3424 \emph{sqq.}

}

The more highly educated pagans offered a longer and more obdurate resistance to the appeals of Christianity than the vulgar crowd. Throughout the fourth and fifth centuries they retained higher education in their hands. The schools of rhetoric, philosophy, law, and science maintained the ancient traditions and the pagan atmosphere. In their writings, some pagans frankly showed their hostility to Christianity, others affected to ignore it. We saw how they threw upon this religion the responsibility for the invasion of the barbarians. But in general their attitude was one of resignation, and they found no difficulty in serving Christian Emperors and working with Christian colleagues.\texthebrew{⁠}\footnote{There seems to have been much mutual tolerance between Christians and pagans in private life. Chrysostom exhorts to goodwill and friendliness toward Hellenes. ``They are all children,'' he says, ``and, like children, when we talk about necessary things they do not attend but laugh.'' \emph{Hom.} 4 on \emph{Ep.} I \emph{ad Corinth.}, \emph{P. G.} LXI.38. In another sermon he describes a dispute between a Hellene and a Christian on the merits of Plato and Paul, the one asserting that Paul was rude and unlearned, the other that he was more learned and eloquent than Plato. Chrysostom's comment is that the Christian took a wrong line, and that the glory of the apostles lay in their rudeness and ignorance. (\emph{Hom.} iii \emph{ib.} 27). Elsewhere he disparages Socrates (\emph{Hom.} 4, \emph{ib.} p35), and Plato (\emph{Hom.} 4 on \emph{Acts}, \emph{P. G.} 60, 50).

} This spirit of resignation is expressed in the most interesting piece we have of the poet Palladas of Alexandria, occasioned by the sight of a Hermes lying in the roadway.

Throughout the fifth century Athens was the headquarters of what may be called higher paganism. The Stoic and Epicurean schools had died out in the third century, and in the fourth the most distinguished savants of the university like Proaeresius and Himerius were sophists, not philosophers. But the Platonic Academy continued to exist, independent of State grants, for it had its own private property produ­cing a revenue of more than £600 a year.\texthebrew{⁠}\footnote{1000 solidi or something more, Damascius, \emph{Vita Isidori}, § 158, referring to the time of Proclus. That the other professors were well paid from the public \textgreek{σιτήσεις} may be inferred from Libanius, \emph{Epp.} 1449.

Thayer's Note: A reminder again that Bury wrote in 1923. Today (2007), as a very rough estimate, £600 would be £35,000, or about \$70,000: the same kind of salary as a modern university professor.

} Its scholarchs, however, were not men of much talent or distinction, until the office was filled by Priscus,\texthebrew{⁠}\footnote{He was the pupil of Aedesius, the most distinguished pupil of Iamblichus who was himself a pupil of Porphyry.

} a Neoplatonist and a friend of Julian, after that Emperor's death. Priscus inaugurated the reign of Neoplatonism at Athens; with him the revival of the university, as a centre of philosophic study, began, and vastly increased under his successor Plutarch. Towards the end of the fourth century, Synesius had spoken in disparaging words of Athens and her teachers: her fame, he said, rests with her bee-keepers .\texthebrew{⁠} a He was jealous for the reputation of Alexandria, and with good reason, for under Plutarch and his successors Syrianus and Proclus Athens was to eclipse the Egyptian city. These Platonists attracted students from all parts of the East, and some who had begun their studies, like Proclus himself, at Alexandria, completed them at Athens.\footnote{We know a good deal about university life at Athens in the fourth and fifth centuries. (See Hertzberg, \emph{Gesch. Griechenlands}, III \emph{passim}; Sievers, \emph{Das Leben des Libanius}, 42 \emph{sqq.}). The principal sources for the fifth are Olympiodorus, \emph{fr.} 51; Marinus, \emph{Vita Procli}; Damascius, \emph{Vita Isidori}; various articles in Suidas.

}

The Athenian professors had always regarded themselves as the champions of Hellenism, but when the Neoplatonic philosophy became ascendant, the Hellenism of Athens was a more serious danger. At this time Neoplatonism was the most formidable rival of Christian theology among educated men of a speculative turn of mind. Augustine recognised this; we know how it attracted him.\texthebrew{⁠}\footnote{He had studied Plotinus and Porphyry in Latin translations. See Angus, \emph{Sources of De civ. Dei}, p268.

} The Neoplatonists taught a system fundamentally differing from the current Christian theology as to the position which was assigned to the creator of the world. According

p376
to Plotinus, Nous or Reason, the creator, emanated from and was subordinate to the absolute One, and Soul again emanated from Nous. His successors developed his principles by multiplying and dividing the emanations, and the growth of the philosophy culminated in the system which Proclus constructed by means of a dialectic which Hegel himself has described as ``extremely tiring.''\texthebrew{⁠}\footnote{\emph{Gesch. der philosophie}, 73\texthebrew{‑}74, in \emph{Werke}, XV. The most recent treatment of the metaphysics of Proclus will be found in Whittaker's \emph{The Neoplatonists.} Procopius, the famous sophist of Gaza, wrote a refutation of the theology of Proclus, which has been preserved under the name of Nicolaus of Methone (\textgreek{\texthebrew{Ἀ}νάπτυξις τ\texthebrew{ὴ}ς θεολογικ\texthebrew{ῆ}ς στοιχειώσεως Πρόκλου}) who simply transcribed Procopius, as has been shown by J. Dräseke, \emph{B. Z.} VI.55 \emph{sqq.}

} In all these phases, the Demiurge or Creator is subordinated to the One of which no divine attributes could be predicted, and thus an apparently impassable gulf was fixed between the later Platonic philosophers and Christian theologians. There was, indeed, at Alexander another school of Platonism, which held closer here and there to the teaching of Plato himself, and men who were trained in this school found the transition to Christian doctrine comparatively easy. We know something of the system of Hierocles, a leading Platonist at Alexandria in the fifth century.\texthebrew{⁠}\footnote{See the instructive article of K. Prächter, in \emph{B. Z.} XXI.1 \emph{sqq.} Of the works of Hierocles are preserved his \emph{Commentary on the Golden Words of Pythagoras}, in Mullach, \emph{Fr. phil. graec.} I.416 \emph{sqq.}, and fragments \emph{\textgreek{Περ\texthebrew{ὶ} προνοίας κα\texthebrew{ὶ} ε\texthebrew{ἱ}μαρμένης}} in Photius, \emph{Bibl.} 214, 251.

} In his system there was no One or any other higher principle above God the creator and legislator, who was above, and in no sense co-ordinate with, the company of sidereal gods; and he, like the Christian Deity, created the world out of nothing. Some of the pupils of Hierocles became Christians. It is a curious circumstance that Hierocles should have been condemned to exile at Constantinople on grounds which are unknown to us.\texthebrew{⁠}\footnote{Suidas, \emph{sub} \textgreek{\texthebrew{Ἱ}εροκλ\texthebrew{ῆ}ς: προσέκρουσε το\texthebrew{ῖ}ς κρατο\texthebrew{ῦ}σι κα\texthebrew{ὶ} ε\texthebrew{ἰ}ς δικαστήριον \texthebrew{ἀ}χθε\texthebrew{ὶ}ς \texthebrew{ἐ}τύπτετο τ\texthebrew{ὰ}ς \texthebrew{ἐ}ξ \texthebrew{ἀ}νθρώπων πληγάς . . . φυγ\texthebrew{ὴ}ν δ\texthebrew{ὲ} κατακριθε\texthebrew{ῖ}ς κα\texthebrew{ὶ ἐ}πανελθ\texthebrew{ὼ}ν χρόν\texthebrew{ῳ ὕ}στερον ε\texthebrew{ἰ}ς \texthebrew{Ἀ}λεξάνδρειαν κτλ}. The source of Suidas was Damascius, \emph{Vita Isidori}, see Photius, \emph{Bibl.} 242 (p338). It may be noted that in political philosophy the Neoplatonists held to Plato's theories. The only attempt at original speculation in the field of political science in this age is to be found in a tract \emph{\textgreek{Περ\texthebrew{ὶ} πολιτικ\texthebrew{ῆ}ς \texthebrew{ἐ}πιστήμης}}, much mutilated, published in Mai (\emph{Scriptores vet. nov. coll.} II.571 \emph{sqq.}) which has been elucidated by Prächter in \emph{B. Z.} IX.621 \emph{sqq.} It seems to have been written by a Christian, \emph{c.} A.D. 500, who was influenced by Neoplatonism, but did not swear by Plato, and made much use of Cicero's \emph{De republica.}

} It can hardly have been for his teaching, seeing that far more anti-Christian Platonists, who had their stronghold at Athens, were tolerated.

p377
But the danger and offence of the later Neoplatonists did not lie in their mystical metaphysics, but in the theurgy and pagan practices to which they were almost always addicted. Proclus in his public lectures as scholarch confined himself, doubtless, to the interpretation of Plato in the Neoplatonic sense, and to problems of dialectic, but he reserved for his chosen disciples esoteric teaching in theurgy, and venerated the gods as beneficent beings worthy of worship, though occupying a subordinate place in the hierarchy of existences. He believed that by fasting and purifications on certain days it was possible to get into communication with supernatural beings, and he recognised the gods of other nations as well as those of Greece. He said that the philosopher should not confine himself to the religious rites of one city or people, but should be ``a hierophant of the whole world.'' He was more scrupulous in observing the fasts of the Egyptians than the Egyptians themselves.\texthebrew{⁠}\footnote{Marinus, \emph{Vita Procli}, c19. Six hymns of Proclus addressed to Greek gods are extant; others celebrated Isis, Marnas the god of Gaza, Thyandrites an Arabian deity (\emph{ib.}).

} He had been initiated in the Eleusinian secrets by his friend Asclepigenia, the daughter of Plutarch,\texthebrew{⁠}\footnote{\emph{Ib.} c28. The learned Asclepigenia married a rich landowner Archiadas, who was very generous to the university. Their daughter, the younger Asclepigenia, married Theagenes, the richest Greek of the day and notably public-spirited in the use of his wealth. He became Archon of Athens.

} who had learned them from the last priest of Eleusis, and in one of his writings he told how he had seen Hecate herself. Athens believed in his magical powers; he was said to have constructed an instrument by which he could bring down rain.

The Hellenists, even in the days of Proclus, had not abandoned all hope of winning toleration for pagan worship. At any time some one might ascend the throne with Hellenic sympathies. The elevation of Anthemius in the West was a proof that this was not impossible, though Anthemius was able to do little to help the pagan interest. Proclus died in A.D. 485, and at that very time a former pupil of his was prominently associated with a rebellion\texthebrew{⁠}\footnote{See

next Chapter, § 2 .

} which, if it had been success ­ful, might have been followed by some temporary relaxation of the severe laws against polytheism and pagan worship. This was to be the last flutter of a dying cause.

The persecution of heretics was more resolute and severe than the persecution of pagans. Those who stood outside of the Church altogether were less dangerous than those members of it who threatened to corrupt it by false doctrine, and the unity of the Catholic faith in matters of dogma was considered of supreme importance. ``Truth, which is simple and one,'' wrote Pope Leo I, ``does not admit of variety.''\texthebrew{⁠}\footnote{\emph{Ep.} 172 varietatem veritas, quae est simplex atque una, non recipit (\emph{P. L.} XX p1216).

} A modern inquirer is accustomed to regard the growth of heresies as a note of vitality, but in old times it was a sign of the active operation of the enemy of mankind.

The heresy which was looked upon as the most dangerous and abominable of all was that of the Manichees, which it would be truer to regard as a rival religion than as a form of Christianity.\texthebrew{⁠}\footnote{The chief sources for Manichaeanism are: the \emph{acta Archelai}, and Alexander of Lycopolis (a Platonist, not a Christian), \textgreek{πρ\texthebrew{ὸ}ς τ\texthebrew{ὰ}ς Μανιχαίου δόξας} (for the early stage of its development); Epiphanius, \emph{De haer.}; Augustine, \emph{Contra Faustum} (and other treatises).

} It was based on a mixture of Zoroastrian and Christian ideas, along with elements derived from Buddhism, but the Zoroastrian principles were preponderant. This religion was founded by Manes in Persia in the third century, and in the course of the fourth it spread throughout the Empire, in the West as well as in the East. Augustine in his youth came under its influence. The fundamental doctrine was that of Zoroaster, the existence of a good and an evil principle, God and Matter, independent of each other. The Old Testament was the work of the Evil Being. Matter being thoroughly evil, Jesus Christ could not have invested himself with it, and therefore his human body was a mere appearance. The story of his life in the Gospels was interpreted mystically. The Manichees had no churches, no altars, no incense; their worship consisted in prayers and hymns; they did not celebrate Christmas, and their chief festival was the Bêma, in March, kept in memory of the death of their founder, who was said to have been flayed alive or crucified by Varahran I. They condemned marriage, and practised rigorous austerities.\footnote{They are accused of disgusting practices, Augustine, \emph{De haer.} 46. Cp. Salmon's art. ``Manicheans'' in \emph{Dict. Chr. Biog.} III.798.

}

The laws against the Manichees, which were frequent and

p379
drastic, began in the reign of Theodosius I. The heresy was insidious, because the heretics were difficult to discover; they often took part in Christian ceremonies and passed for orthodox, and they disguised their views under other names. Theodosius deprived them of civil rights and banished them from towns. Those who sheltered themselves under harmless names were liable to the penalty of death; and he ordered the Praetorian Prefect of the East to institute ``inquisitors'' for the purpose of discovering them.\texthebrew{⁠}\footnote{\emph{C. Th.} XVI.5.7

(A.D. 381);

9

(A.D. 382). Further legislation under Arcadius and Theodosius II will be found in the same title.

} This is a very early instance of the application of this word, which in later ages was to become so offensive, to the uses of religious persecution. When the government of Theodosius II, under the influence of Nestorius, made a vigorous effort to sweep heresy from the world, the Manichaeans were stigmatised as men who had ``descended to the lowest depths of wickedness,'' and were condemned anew to be expelled from towns, and perhaps to be put to death\texthebrew{⁠}\footnote{The words et ultimo supplicio tradendis, in

\emph{C. J.} I.5.5

are omitted in

\emph{C. Th.} XVI.5.65 .

} ( A.D. 428). Later legislation inflicted death unreservedly; they were the only heretics whose opinions exposed them to the supreme penalty.

Arcadius, at the beginning of his reign, reaffirmed all the pains and prohibitions which his predecessors had enacted against heretics.\texthebrew{⁠}\footnote{\emph{Ib.} 25 .

} In most cases, this meant the suppression of their services and assemblies and ordinations. The Eunomians, an extreme branch of the Arians, who held that the Son was unlike the Father, were singled out for more severe treatment and deprived of the right of executing testaments. This disability, however, was afterwards withdrawn, and it was finally enacted that a Eunomian could not bequeath property to a fellow-heretic .\texthebrew{⁠}\footnote{\emph{Ib.} 27 (395); 58 (415) .

} Thus there was a certain vacillation in the policy of the government, caused by circumstances and influences which we cannot trace.

The combined efforts of Church and State were success ­ful in virtually stamping out Arianism, which after the end of the fourth century ceased to be a danger to ecclesiastical unity. They were also success ­ful ultimately in driving Nestorianism out of the Empire. The same policy, applied to the Monophysitic heresy,

p380
failed. Marcian's law of A.D. 455 against the Eutychians was severe enough.\texthebrew{⁠}\footnote{\emph{C. J.} I.5.8 .

} They were excluded from the service of the State; they were forbidden to publish books criticising the Council of Chalcedon; and their literature, like that of the Nestorians, was condemned to be burned. But in Syria, where anti-Greek feelings were strong, and in Egypt, where national sentiment was beginning to associate itself with a religious symbol, all attempts to impose uniformity were to break down.

The severe measures taken by the State against the Donatists in Africa were chiefly due to their own fanaticism. Donatism was not properly a heresy, it was a schism, which had grown out of a double election to the see of Carthage in A.D. 311, and the question at issue between the Catholics and the Donatists was one of church discipline. We need not follow the attempts of Constantine and Constans to restore unity to the African church by military force. The cause of the Donatists was not recommended by their association with the violent madmen known as Circumcellions, who disdained death themselves, and inflicted the most cruel deaths on their opponents. The schismatics survived the persecution. At the death of Theodosius I the greater number of the African churches seem to have been in their hands, and during the usurpation of Gildo they persecuted the Catholics. When Augustine became bishop of Hippo, where the Donatists were in a great majority, he set himself the task of restoring ecclesiastical unity in Africa by conciliation.\texthebrew{⁠}\footnote{The numerous writings against the Donatists will be found in \emph{P. L.} 43.

} He and the Catholic clergy had some success in making converts, but the fanatics were so infuriated by these desertions that with their old allies the Circumcellions they committed barbarous outrages upon the Catholic clergy and churches; Augustine himself barely escaped from being waylaid. Such disorders demanded the intervention of the secular power. Some injured bishops presented themselves at Ravenna, and in A.D. 405 Honorius condemned the Donatists to severe penalties by several laws intended ``to extirpate the adversaries of the Catholic faith.''\footnote{\emph{C. Th.} XVI. 6.4 and 6 ;

5.38 and 39 .

}

The Donatists rejoiced at the death of Stilicho whom they regarded as the author of these laws, and disorders broke out afresh.\texthebrew{⁠}\footnote{Augustine, \emph{Ep.} CXI.

} When Alaric was in south Italy threatening Rome,

p381
the Emperor revoked his decrees and soon afterwards, at the request of the Catholics, he convoked a conference of the bishops of the two parties which met at Carthage ( A.D. 411) under the presidency of Marcellinus, one of the ``tribunes and notaries'' whom the Emperors employed for special services. Marcellinus was empowered not only to act as chairman but to judge between the rival claims. The appointment of a secular official to adjudicate did not mean that the civil power claimed to settle questions of doctrine. The controversy, which originally turned on a dispute about facts, had throughout concerned the government not in its ecclesiastical aspect but as a cause of grave disorders and disturbances. But the commission entrusted to Marcellinus shows that the bishop of Rome was not yet recognised as possessing the jurisdiction which in later times resided in his see. At the end of the discussions, Marcellinus decided against the Donatists; they were allowed a certain time to come into the Church.\texthebrew{⁠}\footnote{For the proceedings of the conference see Mansi, IV.51 \emph{sqq.}

} Some were convinced, but others appealed to the Emperor, who confirmed the decision of his deputy and enacted a new law against the schismatics, imposing heavy fines on the recalcitrants, and banishing the clergy.\texthebrew{⁠}\footnote{A.D. 412.

\emph{C. Th.} XVI.5.52 . Slaves and colons were to be beaten out of their false religion (a prava religione).

} Two years later they were deprived of civil rights.\texthebrew{⁠}\footnote{A.D. 414.

\emph{Ib.} 53

perpetua inustos infamia.

} These strong measures, which Augustine defended, alleging the text ``Compel them to come in,''\texthebrew{⁠}\footnote{Augustine, \emph{Ep.} CXIII.

} broke the strength of the schismatics, and though the Donatist sect continued to exist and was tolerated under the Vandals, it ceased to be of importance.

It must be allowed that if the government had been perfectly indifferent and impartial in matters of religion, it would have had ample excuse for adopting severe measures of repression against the fanatical sect who disturbed the peace of the African provinces and persecuted their opponents. The penalties were severe but they stopped short of death. It should be remembered to the credit of the Emperors that, in contrast with the Christian princes of later ages, they never proposed, in pursuing their policy of the suppression of heresy, to inflict the capital penalty, except in the case of the Manichaeans, who were regarded as almost outside the pale of humanity.\texthebrew{⁠}\footnote{Diocletian had legislated against Manichaeanism (A.D. 287) as destructive of morality.

} The same may be said

p382
for the leading and representative ecclesiastics, all of whom would have recoiled with horror if they could have foreseen the system of judicial murder which was one day to be established under the auspices of the Roman see.\texthebrew{⁠}\footnote{Chrysostom expresses his views on the repression of heretics in \emph{Hom.} 46 \emph{in Matth.} (\emph{P. G.} XLVIII p477), where he comments on the parable of the tares. They should be silenced but not put to death.

} Martin of Tours did all he could to stay the persecution of the Spanish bishop Priscillian, who, rightly or wrongly, was accused of heresies akin to Manichaeanism. Priscillian was put to death by the Emperor Maximus ( A.D. 385), but he was tried before a civil tribunal for a secular offence.\texthebrew{⁠}\footnote{Maleficium. Sulpicius Severus, \emph{Chron.} II.50 nec diffitentem obscenis se studuisse doctrinis, nocturnos etiam turpium feminarum egisse conuentus nudumque orare solitum.

Babut, \emph{Priscillien et le Priscilliénisme} (1909); Holmes, \emph{The Christian Church in Gaul} , chapters VIII, IX.

} It may well have been a miscarriage of justice, but, formally at least, he was not executed as a heretic.

Under the Christian Empire the Jews remained for the most part in possession of the privileges which they had before enjoyed.\texthebrew{⁠}\footnote{A.D. 404.

\emph{C. Th.} XVI.8.15

confirms the privileges of the Jewish hierarchy. For the pressure put on Emperors by churchmen not to afford proper protection to the Jews against Christian fanatics cp. Ambrose, \emph{Epp.} 40, 41. For the anti-Semitism of Chrysostom see his eight homilies \emph{Against the Jews} delivered at Antioch, \emph{P. G.} XLVIII.843 \emph{sqq.} He says that ``demons inhabit their souls,'' p852. There is a virulent attack on the Jews in

Rutilius Nam. \emph{De reditu suo}, I.382 \emph{sqq.}

} The Church was unable to persuade the State to introduce measures to suppress their worship or banish them from the Empire. They were forbidden to possess Christian slaves,\texthebrew{⁠}\footnote{\emph{Ib.} IX.4

(416);

IX.5

(423). Justinian extended this to pagans, Samaritans, and all heterodox persons

(\emph{C. J.} I.10.2) .

} and a law of Theodosius II excluded them from civil offices and dignities.\texthebrew{⁠}\footnote{\emph{C. J.} I.9.18

(439).

} But the legislator was perhaps more often concerned to protect them than to impinge upon their freedom.\footnote{Thus in 412, Christians were forbidden to disturb Jewish worship; and in 423 to burn or take away synagogues

(\emph{C. Th.} XVI.8.20 and 25) .

}

The same period, in which the Christian religion gradually won the upper hand in the Empire, witnessed a movement which was at first independent of the Church but was destined soon to become an important part of the ecclesiastical system.

p383
The germs of asceticism had been implanted in the Christianity from the very beginning, and the tendencies to a rigorous life of self-abnegation may have been stimulated by the example of the austerities of the Essenes, the Therapeutae, the monks of Serapis, and later by the influence of the semi-Christian Zoroastrian religion of the Manichees. Ascetic practices seem to have been a strong temptation to all men of an ardently religious temperament in these ages, whatever doctrines they might hold concerning the universe; Julian the Apostate is an eminent example. For the Christian Church and State the consequences were far-reaching and could not have been anticipated. In the course of the fourth and fifth centuries a large and ever-growing number of men and women withdrew themselves from society, severed themselves from family ties, and embraced, whether in cells in the desert or in recluse communities in town or country, a life of celibacy, prayer, and fasting. Gradually regularised and organised by disciplines of varying degrees of rigour, monasticism established itself firmly as one of the most influential institutions of the Christian world, thoroughly consonant with the spirit of the time and richly endowed by the liberality of the pious.

We have not to follow the history of its growth, but the reader may be reminded that Christian monasticism originated \emph{circa} A.D. 300 under the auspices of St. Anthony in Lower Egypt. At first it took the form of a solitary life in the desert, where ascetics lived independently of one another in neighbouring cells and devoted themselves to an otherwise idle existence of religious contemplation.\texthebrew{⁠}\footnote{The chief settlements were in the desert south of Alexandria, at Nitria (Wadi Natron) and Scete. At Nitria there were 5000 monks towards the end of the fourth century. The chief sources for Egyptian monasticism are Palladius, \emph{Historia Lausiaca}, and Rufinus, \emph{Historia monachorum.}

} Another variety of monasticism was soon afterwards founded in Upper Egypt by Pachomius. In his monasteries near Tentyra (Denderah) and Panopolis (Akhmim) the brethren lived in common and performed all kinds of work. The Antonian ideal was approved by Athanasius, and his influence went far to spread it in the West. It was introduced into Palestine by Hilarion, and into Syria, where the rigours of the hermit assumed their most extreme and repulsive shape. There was originated the grotesque idea of living for years on the top of a high pillar. Simeon, the first of these

p384
pillar-saints ( stylitae ),\texthebrew{⁠}\footnote{A.D. 388\texthebrew{‑}459. Having lived at first in an enclosed cell at Antioch, he built a low pillar in 423, and gradually raised it till in 430 it was \texthebrew{•} forty cubits high; at the top it was \texthebrew{•} three feet in circumference, according to

Evagrius, I.13 . It was situated at the ruins known as Kalat Semian, house of Simeon, described by

De Vogüé, \emph{Syrie Centrale}, I.141 \emph{sqq.}

Theodosius II wrote a letter to him, asking him to descend from his column. On his death, his body was taken to Antioch with the honours of a state funeral, and Leo I wished to have it transported to Constantinople (see
Evagrius I.13, 14 ; Theodore Lector II.41; \emph{Vita Sim. Styl.} ed. Lietzmann). Daniel, an imitator of Simeon, set up a pillar four miles north of Constantinople, and lived on it for thirty-three years, in the reigns of Leo I and Zeno. He was constantly frozen over with snow and ice, and his feet were covered with sores. Leo insisted on putting a shed over the top of the pillar. The saint descended from his perch in order to denounce the ecclesiastical policy of Belisarius. The Patriarch Euphemius attended him in his last moments. See the \emph{Vita Danielis.}

} had many followers, and such was the temper of the times that these abnormal self-tormentors , who could not have been more healthy in mind than in body, were universally revered and consulted as oracles.

The monastic movement engaged the attention of St. Basil, and awoke his enthusiasm. He came to the conclusion that monastic institutions, framed on right lines, would be useful to the Church, and he established a coenobitic community at Neocaesarea (about A.D. 360), and drew up minute regulations. The brethren were not required to take vows; the asceticism of their life was not immoderate; and they were expected to perform work in the fields. St. Basil's idea had an immediate success and he became the founder of Greek monasticism. Cloisters adopting his Rule\texthebrew{⁠}\footnote{The Rules will be found in his works, \emph{P. G.} XXXI.889 \emph{sqq.}

} sprang up throughout Asia Minor, and in the following century in Palestine. But here there flourished also the lauras, or enclosures in which the monks lived an almost eremitical life in separate cells, and these institutions were numerous in the plain of the Jordan.\texthebrew{⁠}\footnote{They are enumerated and located in Génier, \emph{Vie de Saint Euthyme le Grand}, chap. I.

} The most famous of the ascetics of Palestine were Euthymius, Sabas, and Theodosius.\texthebrew{⁠}\footnote{The lives of Euthymius and Sabas were written by Cyril of Scythopolis in the sixth century. On Theodosius we have a brief sketch by the same writer and a panegyric by Theodore, bishop of Petrae, probably delivered in 530. (See

Bibliography, I.2, A .)

} Euthymius founded the laura of Sahel, to the east of Jerusalem, in A.D. 428;\texthebrew{⁠}\footnote{Euthymius did much for the conversion of the Saracens, and founded the Parembole (to the east of his own laura), a large enclosure in which baptized Saracens were settled. They had their own bishop. The Parembole was ruined by the invasion of Al\texthebrew{‑}Mundhar in the sixth century.

} Sabas founded in A.D. 483 the Great Laura on the Cedron, with a grotto which nature had moulded

p385
into the form of a church, and many others; and Theodosius his coenobitic monastery at the grotto of the Magi near Bethlehem in A.D. 476. Sabas was appointed archimandrite of all the lauras, and Theodosius of all the coenobia, in the diocese of Jerusalem by the Patriarch Sallust ( A.D. 494). It would seem that the monks of the lauras were considered to have attained to a higher grade of spiritual life than those who lived in convents, which were regarded as a preparation in ascetic discipline.\texthebrew{⁠}\footnote{Cp. Génier, \emph{op. cit.} p11.

} As Sabas and one of his disciples walked one day from Jericho to the Jordan, they met a young and comely girl. ``Did you remark that girl?'' said the saint, ``she is one-eyed.'' ``No, Father,'' said the disciple, ``she had both her eyes.'' ``How do you know?'' ``I looked at her intently.'' ``What about the commandment, 'Fix not your eyes on her, neither let her take thee with her eyelids'?''\texthebrew{⁠}\footnote{Cyril, \emph{Vita Sabae}, XLVIII p251.

} And the saint sent the youth to a convent till he had learned better to control his eyes and his thoughts.

The history of monasticism at Constantinople begins with the abbot\texthebrew{⁠}\footnote{\emph{Archimandrite}, \emph{i.e.} head of the mandra or sheep-pen, often used instead of \textgreek{\texthebrew{ἡ}γούμενος}, the usual term, or \textgreek{\texthebrew{ἀ}ββ\texthebrew{ᾶ}ς}. In later times \emph{archimandrite} was confined to designate authority over several monasteries; \emph{exarch} was also used.

} Isaac, a Syrian, who in the reign of Theodosius I founded a convent in the quarter of Psamathia outside the Constantinian Wall. He was a typical fanatical ascetic and was buried with great pomp when he died.\texthebrew{⁠}\footnote{In 407\texthebrew{‑}408. He was an opponent of Chrysostom. For the beginnings of monasticism at Constantinople see Pargoire's article in \emph{Revue des questions historiques}, LXV, 1899, where many false traditions are exposed.

} He was succeeded by Dalmatius, an active organiser, who founded new houses under his own authority. The community of the Akoimetoi or Sleepless was established at Gomon, near the northern entrance to the Bosphorus, by one Alexander in the reign of Theodosius II, but his successor John transported the monks to a new cloister at Chibukli, on the Asiatic side of the straits opposite to Sosthenion,\texthebrew{⁠}\footnote{\emph{Vita Marcellis}, in Simeon Metaphrastes, º \emph{P. G.} CXVI p712. Cp. Pargoire, \emph{Anaple et Sosthène}, p64. The monks were called sleepless because they maintained choral service continuously by relays.

} where it became famous under the next abbot Marcellus, who presided for about forty years. Two other foundations deserve notice. The monastery of Drys, a suburb of Chalcedon, was established

p386
by Hypatius, who enforced a very strict discipline, about A.D. 400. Hypatius enjoyed considerable influence. Theodosius II used to visit him, and he was constantly consulted by the nobles and ladies of the capital.\texthebrew{⁠}\footnote{See Callinicus, \emph{Vita Hypatii.}

} The most famous of the monastic communities of Constantinople was founded by Studius, an ex-consul who had come from Rome,\texthebrew{⁠}\footnote{Consul in 454. But an epigram on the church represents him as rewarded by the consul­ship for building it (\emph{Anth. Graeca}, I.4).

} in the reign of Leo I. He dedicated a small basilica to St. John the Baptist, which is still preserved as a mosque,\texthebrew{⁠}\footnote{Emir Ahor Jamissi. Described by van Millingen, \emph{Byzantine Churches}, chap. II.

} not far from the Golden Gate, and subsequently attached to it a monastery, in which he established some of the Sleepless brethren, who had belonged to the convent of Marcellus.\texthebrew{⁠}\footnote{A.D. 462\texthebrew{‑}463. Theodore Lector, I.7 = Theophanes A.M. 5955. But the sleepless tradition was not maintained at Studion, and when we read of the Akoimetoi, the monks of the cloister on the Bosphorus are meant, not the Studites.

} The Studite community was to become the largest and most influential in Constantinople.

Of the countries of western Europe, early monasticism spread most widely in Gaul. Martin of Tours was the pioneer; he founded a monastery at Poictiers about A.D. 362. Some forty years later Cassian inaugurated monastic life at Marseilles, and Honoratus in the islands of Lérins off the coast of Provence. Both Cassian and Honoratus were under the direct influence of the theories of ascetic life which were practised by the Antonian monks of northern Egypt.\texthebrew{⁠}\footnote{Cassian's \emph{Collationes} profess to reproduce the teaching of the monastic leaders of Nitria and Scete.

} In the same period, monasteries both for men and for women \texthebrew{—} women already took their full share in the ascetic movement \texthebrew{—} were established at Rome and in Italian towns, and Augustine introduced monastic life in Africa. Spain, so far as our evidence goes, seems to have been little affected by the fashion before the sixth century.

We have no information that would enable us to conjecture the total number of the voluntary exiles from social life, who in the fifth century, whether in communities or lonely cells, mortified their bodies and their natural affections in order to assure themselves of eternal happiness. Ascetic enthusiasm was infectious, and the leading authorities of the church, such as Jerome, Ambrose, Augustine, Chrysostom, all held up the monastic life as the highest spiritual ideal, and outdid each other in their

p387
praises of celibacy and virginity. But the Church and the State soon found it necessary, in the interests of public order, to exercise control over the ascetics, who in the early period of the movement were each his own master and acknowledged no superior. The towns were often troubled by the invasion of vagrant monks, genuine or spurious, who formed a highly undesirable addition to the idle and mendicant portion of the populace.\texthebrew{⁠}\footnote{An edict of 390 (Verona) commands monks to remain in ``desert places and vast solitudes,''

\emph{C. Th.} XVI.3.1 , but in 392 Theodosius withdrew the prohibition of free entry into towns (\emph{ib.} 2). Chrysostom excited much resentment in monastic circles by his endeavour to suppress the monks who lounged about Constantinople.

} We have seen again and again the turbulence of the monks, who, in their religious zeal, were ready to commit any excess of violence and transgression of decency. Their fanaticism was responsible for the useless destruction of pagan temples. They played a leading part in the disturbances at Alexandria which ended in the murder of Hypatia. They were the chief offenders in the scandalous disorders which disgraced the Council of Ephesus. During the first half of the fifth century, the bishops seem to have been gradually acquiring some control over the cloisters, but the prevailing anarchy was definitely ended by the Council of Chalcedon.\texthebrew{⁠}\footnote{Canons 3, 4, 7, 16 (Mansi, VII.371 \emph{sqq.}). Cp. \emph{C. J.} I.3.53 .

} This assembly deplored the turbulence of the monks, and forbade them to abandon their holy life. It ordained that no one could found a monastery without a licence from the bishop of the diocese, and that no monk could leave his convent without the bishop's permission. Monastic communities were thus brought under ecclesiastical control.

The estates of the monasteries gradually increased through the donations of the rich and pious, and at the beginning of the sixth century a pagan historian writes thus of the ``so\texthebrew{‑}called monks'':\texthebrew{⁠}\footnote{Zosimus, V.23.

} ``They renounce legal marriages and fill their populous institutions in cities and villages with celibate people, useless either for war or for any service to the State; but gradually growing from the time of Arcadius to the present day they have appropriated the greater part of the earth, and on the pretext of sharing all with the poor they have, so to speak, reduced all to poverty.'' This is the exaggerated statement of a hostile observer, who had been an official of the treasury; but it testifies to the growing popularity, wealth, and power of monastic institutions.

p388
The ascetic spirit, which expressed itself in monasticism, affected the secular clergy also. The strict austerity of the Manichaean heretics was a certain challenge to the Church,\texthebrew{⁠}\footnote{Cp. Lea, \emph{History of Sacerdotal Celibacy}, I.33 \emph{sqq.}

} and in their extravagant praises of virginity some of the Christian fathers were barely able to stop short of the condemnation of marriage which was a tenet of the Manichees. The view that matrimony is a necessary evil naturally involved the question of the celibacy of the clergy. In this matter ecclesiastics were left free to follow the dictates of their own conscience, and no legislation was attempted, till a Roman council (about A.D. 384) summoned by Pope Siricius, forbade bishops, priests, and deacons to marry. ``Celibacy,'' it has been said, ``was but one of the many shapes in which the rapidly progressing sacerdotalism of Rome was overlaying religion with a multitude of formal observances.''\texthebrew{⁠}\footnote{Lea, \emph{ib.} 66.

} Against the encroachments of this sacerdotalism, a protestant movement was led in Gaul by Vigilantius, who denounced celibacy, fasting, prayers for the dead, relics, and the use of incense; but it did not survive his death. By degrees, the celibacy of the clergy became the rule in the west. In the eastern provinces, where Roman influence was not preponderant, it was otherwise. Marriage after ordination was forbidden, but compulsory separation of clergy who were already married was not imposed except in the case of bishops.\footnote{Cp. Socrates, V.22;

\emph{C. J.} I.3.45 ; Justinian, \emph{Nov.} CXXIII.12; \emph{Concilium Trullianum}, Canons 13, 30 (Mansi, XI.948, 956). Men who had been twice married were strictly excluded from holy orders.}

\xchapter{Chapter XII}

Short URL for this page:

bit.ly/BURLAT12

The new Emperor, Leo II, was a child of seven years, and the regency naturally devolved on his father Zeno. But with the consent of the Senate and the concurrence of the Empress Verina, the child conferred the Imperial dignity on his father, in the Hippodrome (February 9, A.D. 474) and died in the same year, leaving to Zeno nominally as well as actually the sole power (November 17).\footnote{See Candidus, p136; John Mal. XIV p376; Theodore Lector, I.24, 27; Theophanes A.M. 5966, 5967. The coronation in the Hippodrome (instead of the Hebdomon) was an innovation. We have coins of the joint reign with Dn Leo et Zeno Pp Aug and on the reverse the two Emperors seated, Zeno on Leo's left; and others with different reverses. There are also tremisses with Dn Zeno et Leo Caes on the obverse. See Sabatier I. Pl. VII.15, 16, 17; Pl. VIII.13.

}

Zeno was not beloved.\texthebrew{⁠}\footnote{He was married to Arcadia before he married Ariadne, and by her had a son, Zeno, of whom something more will be heard. Zeno was a very fast runner, according to a Ravenna chronicler known as the Anonymus Valesii (see below,

p423, n1 ), who had a marked liking for him. His speed of foot was ascribed to a peculiarity in his knee-caps; perhibent de eo quod patellas in genucula non habuisset, sed mobiles fuissent

(Anon. Val. IX.40) . Fast running was an Isaurian characteristic; compare the marvellous speed of Indacus (Suidas, \emph{sub} \textgreek{\texthebrew{Ἰ}νδακός}; John Ant. in \emph{F. H. G.} IV.617).

} He was unpopular both with the Byzantine populace and in senatorial circles.\texthebrew{⁠}\footnote{Cp. Joshua Stylites, 12.

} He was hated as an Isaurian. If we remember the depredations of the Isaurians in the reign of Arcadius, it is not surprising that they had an evil name, and it is more than probable that the soldiers introduced into the capital by Leo had not belied their reputation for rudeness and violence. Zeno's accession meant Isaurian ascendancy,

p390
high places for the Emperor's fellow-countrymen , and more rude mountaineers in the capital. Historians of the time vent their feelings by describing him as physically horrible and morally abominable,\texthebrew{⁠}\footnote{Cp.
Evagrius III.1 ; Malchus, \emph{fr.} 16. The prejudice of Malchus, who wrote under Zeno's successor, is undisguised.

} and he was said to be a coward.\texthebrew{⁠}\footnote{John Lydus, \emph{De mag.} III.45 , ``he could not bear even the picture of a battle.''

} His most trusted counsellor was the Isaurian Illus, who was, however, to prove a thorn in his side, and Trocundes, the brother of Illus, also rose into prominence.

The first year of the reign was crowded with anxieties. Vandals, Ostrogoths, Huns, and Arabs were all in arms against the Empire. King Gaiseric must have been deeply displeased by the murder of the Arian Aspar, with whom he is said to have been on friendly terms. After Leo's death, the Vandals descended on the western shores of Greece and captured Nicopolis. Zeno was not prepared for war. He sent to Carthage Severus, a man of high repute, who made a favourable impression on Gaiseric by refusing all his gifts. The king made him a present of all the captives who had fallen to the share of the royal family and allowed him to redeem others from any Vandals who willing to sell. A perpetual peace was then concluded between the two powers ( A.D. 474),\texthebrew{⁠}\footnote{Malchus, \emph{fr.} 3, Procopius, \emph{B. V.} I.7.

} and was maintained for nearly sixty years. Meanwhile Zeno's coronation had provoked Aspar's Ostrogothic relative Theoderic Strabo to new hostilities in Thrace. The Master of Soldiers in the Thracian provinces was captured and slain; but Illus took the field and terminated the war.

If the Emperor was able to cope with foreign foes by negotiation or arms, his position amid a hostile court and people was highly precarious. A formidable conspiracy was formed against him, of which the leading spirit was his mother-in\texthebrew{‑}law , the Augusta Verina.\texthebrew{⁠}\footnote{The fullest sources for this conspiracy are Candidus, p136, and John Ant. \emph{ib.}

} She had concurred in Zeno's elevation, but she did not like him, and being a woman of energy and ambition she found it distasteful to fall into the background, overshadowed by her daughter, the Augusta Ariadne. Her scheme was to raise to the throne and marry her paramour Patricius, who had formerly held the post of Master of Offices. She engaged the co-operation of her brother Basiliscus, who had been living in retirement at

p391
Heraclea on the Propontis, and Basiliscus succeeded in sedu­cing the Isaurian brothers Illus and Trocundes to abandon their loyalty to Zeno.\texthebrew{⁠}\footnote{According to the text in John Ant. Illus persuades Basiliscus, but it seems probable that this is a textual error, and that Basiliscus is really intended to be the subject of \textgreek{ποιε\texthebrew{ῖ}ται.}} When all the preparations were complete, the queen-mother , with consummate skill, persuaded Zeno that his life was in danger and that his only safety was in flight. Taking with him a large company of Isaurians, and supplying himself with treasure, he crossed over to Chalcedon (January 9, A.D. 475) and fled to Isauria.\texthebrew{⁠}\footnote{See Brooks, \emph{Emperor Zeno and the Isaurians}, p217, n19 .

} Those who accompanied him were fortunate, for, when the Emperor's flight was known, the populace indulged in their inveterate hatred of the Isaurians by a colossal massacre. Verina now hoped to reign as mistress of the palace, but she was outwitted by her brother, who was himself ambitious of the purple. The choice of the ministers and Senate fell not on Patricius but on Basiliscus, who was proclaimed and crowned Emperor at the Hebdomon palace. He immediately crowned his wife Zenonis as Augusta, and conferred the rank of Caesar upon his youthful son Marcus, whom he afterwards crowned Augustus.\texthebrew{⁠}\footnote{Theodore Lector, I.29 and Candidus, p136. That Marcus was successively Caesar and Augustus is borne out by the superscriptions of the Encyclicals (in Zacharias Myt. V.2, cp. 3 \emph{ad init.}; Evagrius III. 4 ,

5 , 7 ).

} The circumstances of his elevation naturally led to a breach with Verina, and, having good reason to fear her capacity for intrigue, he took the precaution of putting Patricius to death.\footnote{Candidus, \emph{ib.} Verina then intrigued to bring back Zeno; Basiliscus discovered her plots; and it might have gone hard for her, if Armatus had not contrived to conceal her.

}

Basiliscus reigned for twenty months and in that time he made himself extremely unpopular, chiefly by his ecclesiastical policy. He favoured the heresy of Monophysitism and issued a decree against the Council of Chalcedon. He and his wife had fallen under the influence of Timothy Aelurus, the bishop of Alexandria, who had come to Constantinople, and he went so far as to withdraw the Asiatic sees from the control of the bishop of Constantinople.\texthebrew{⁠}\footnote{On ecclesiastical affairs see below,

§ 3 .

} Acacius, the Patriarch, was roused by this injury to the rights of his see. He draped St. Sophia in black and appeared in mourning before a large sympathetic congregation. Basiliscus left the city.

The Emperor had made another enemy in the Ostrogothic

p392
Theoderic Strabo, who, as the enemy of Zeno, had supported his elevation, by bestowing a Mastership of Soldiers\texthebrew{⁠}\footnote{According to John Mal. XXV.378, 379 (comp. \emph{Chron. Pasch., sub a.} 478 ) he was mag. mil. in praesenti in 476. Otherwise Suidas (\emph{sub} \textgreek{\texthebrew{Ἀ}ρμάτιος}) \textgreek{στρατηγ\texthebrew{ὸ}ν \texthebrew{Ἰ}λλυρι\texthebrew{ῶ}ν}, and otherwise again Theoph. A.M. 5969 \textgreek{στρατηγ\texthebrew{ὸ}ν \texthebrew{ὄ}ντα τ\texthebrew{ῆ}ς Θρ\texthebrew{ᾴ}κης}. As Suidas is probably copying either Malchus or Candidus, perhaps Armatus was at first mag. mil. in Illyricum and afterwards in praesenti. Shestakov has made it probable that the articles of Suidas \emph{\textgreek{\texthebrew{Ἀ}ρμάτος}} (and \emph{\textgreek{Βασιλίσκος}})

which Müller (\emph{F. H. G.} IV) assigned to Malchus, \emph{fr.} 7, 8 , come partly from Candidus (see his \emph{Kandid Isauriski}, cp. \emph{Bibl.} II.2B ).

} on his relation Armatus, a young fop, who was the lover of the Empress Zenonis. Their love is described by a historian in a passage worthy of a romance.\footnote{The passage is in Suidas and in \emph{F. H. G.} IV. p117 is printed with the fragments of Malchus. But it is more probable that it comes from Candidus.

}

Basiliscus permitted Armatus, inasmuch as he was a kinsman, to associate freely with the Empress Zenonis. Their intercourse became intimate, and as they were both persons of no ordinary beauty they became extravagantly enamoured of each other. They used to exchange glances of the eyes, they used constantly to turn their faces and smile at each other; and the passion which they were obliged to conceal was the cause of dule and teen. They confided their trouble to Daniel a eunuch and to Maria a midwife, who hardly healed their malady by the remedy of bringing them together. Then Zenonis coaxed Basiliscus to grant her lover the highest office in the city.

The preferment which Armatus received from his uncle elated him beyond measure. He was naturally effeminate and cruel. Theoderic Strabo despised him as a dandy who only care for his toilet and the care of his body; and it was said that in the days of Leo he had punished a number of Thracian rebels by cutting off their hands. When he was exalted by his mistress's husband, he imagined that he was a man of valour, and dressed himself as Achilles, in which guise he used to ride about and astonish or amuse the people in the Hippodrome. The populace nicknamed him Pyrrhus, on account of his pink cheeks, but he took it as a compliment to his valour, and became still more inflated with vanity. ``He did not,'' says the historian, ``slay heroes like Pyrrhus, but he was a chamberer and wanton like Paris.''

Basiliscus, perhaps soon after his elevation, had despatched Illus and Trocundes against Zeno, who, now in his native fortresses,\texthebrew{⁠}\footnote{The strongholds called Salmon (locality unknown), Zachariah Myt. V.1.

} had resumed the life of an Isaurian chieftain. Basiliscus, however, failed to fulfil what he had promised to the two generals; and they received letters from some of the leading ministers at

p393
the court, urging them to secure the return of Zeno. For the city was now prepared to welcome the restoration of the Isaurian, to replace the Monophysite, whose unpopularity was increased by the fiscal rapacity of his ministers.\texthebrew{⁠}\footnote{Especially of Epinicus who, then a favourite of Verina, had in Leo's reign filled the highest financial offices; and was appointed, apparently by Basiliscus, praetorian prefect. (Suidas, \emph{sub nomine}, calls him \textgreek{\texthebrew{ὑ}παρχος τ\texthebrew{ῆ}ς πόλεως}, but this seems inconsistent with what is said about his oppression of the provinces, \textgreek{κα\texthebrew{ὶ} τ\texthebrew{ὰ ἔ}θνη κα\texthebrew{ὶ} τ\texthebrew{ὰ}ς τ\texthebrew{ῆ}ς πόλεις καπηλεύων κτλ}., which are only appropriate to a praetorian prefect. The notice must come either from Malchus or from Candidus.) Cp. also Suidas, \emph{sub} \textgreek{Βασιλίσκος} (which appears as Malchus, \emph{fr.} 7, in \emph{F. H. G.} IV. p116).

} Illus decided to change sides, and his resolution may have been reinforced by the fact that he had a certain hold over Zeno, having got into his power Longinus, Zeno's brother, whom he kept a prisoner in an Isaurian fortress. Accordingly, Zeno and Illus joined forces and started for Constantinople. When Basiliscus received news of this danger, he hastened to recall his ecclesiastical edicts and to conciliate the Patriarch and the people.\texthebrew{⁠}\footnote{See below,

§ 2 .

} But it was too late. Armatus, the Master of Soldiers, was sent with all available forces to oppose the advancing army of the Isaurians, but secret messages from Zeno, who promised to give him the Mastership of Soldiers for life and to confer the rank of Caesar on his son, induced him to betray his master. He avoided the road by which Zeno was advancing and marched into Isauria by another way. This betrayal decided the fate of Basiliscus. Zeno entered the capital without resistance in August 476. Basiliscus was sent to Cucusus in Cappadocia and there beheaded; his wife and children shared his fate. The promise which had been made to Armatus was kept to the letter. His son was created Caesar at Nicaea. But immediately afterwards the Emperor, by the advice of Illus, caused him to be assassinated, and the Caesar was stripped of his rank and compelled to take orders.\footnote{Evagr. III.24 . Malchus, \emph{ib.} The assassin was Onoulf, Mag. Mil. per Illyr., a brother of Odovacar, who at this time was establishing his power in Italy. We have coins of Basiliscus, and of Basiliscus and Marcus together, and of Zenonis (Sabatier, I. Pl. VIII.14\texthebrew{‑}20); the faces are all conventional.

}

A deplorable misfortune, which occurred in the reign of Basiliscus, is said to have helped, as accidents in superstitious ages always help, to render his government unpopular. This was an immense conflagration,\texthebrew{⁠}\footnote{Cedrenus I.618 = Zonaras XIV.22\texthebrew{‑}24. The ultimate source is evidently Malchus, see Suidas \emph{s.v.} \textgreek{Μάλχος.}} which, beginning in the quarter of Chalkoprateia, spread far and wide, redu­cing to ashes the

p394
adjacent colonnades and houses. But more serious was the destruction of the Basilica, the library founded by Julian, which contained no fewer than 120,000 books. Among these rolls, the intestine of a serpent, 120 feet long, on which the \emph{Iliad} and \emph{Odyssey} were written in golden characters, is specially mentioned. The fire spread along Middle Street and destroyed the palace of Lausus, which contained among its splendours some of the most beautiful works of Greek plastic art, the Cnidian Aphrodite, the Lindian Athene, and the Samian Here.

For the first few years after the restoration of Zeno, Illus was all-powerful . He was consul in A.D. 478; he was appointed Master of Offices, and created a patrician. But he was bitterly detested by the two Empresses, Verina and Ariadne, who resented his influence with Zeno. Attempts on his life were made at Verina's instigation. Her favourite, the Prefect Epinicus, suborned a barbarian to assassinate him. The attempt failed; the criminal confessed that the prefect had inspired his act; and Zeno, having deprived Epinicus of his office, handed him over to Illus who sent him to a castle in Isauria.\texthebrew{⁠}\footnote{John Ant. \emph{fr.} 95 (\emph{De ins.}). Cp. Brooks, \emph{op. cit.} 218, n56

for date.

} Some time elapsed, and then, leaving the capital on a pretext, Illus visited Epinicus in his prison and elicited a confession that he had been instigated by the queen-mother . He then refused (towards the end of A.D. 479) to return to Constantinople unless Verina were surrendered to him. Zeno, to whom Illus was indispensable, complied; she was sent to Tarsus where she was forced to become a nun and was confined by Illus in the castle of Dalisandus.\texthebrew{⁠}\footnote{Dalisandus, in the Decapolis of Isauria, is to be distinguished from Dalisandus in Lycaonia, see Ramsay, \emph{Hist. Geog.} pp335, 366, with the map opp. p330.

} The presence of Illus was sorely needed, on account of Ostrogothic hostilities in Illyricum and Thrace,\texthebrew{⁠}\footnote{An earthquake on Sept. 25, 479, had done terrible damage to the walls of the city, and an Ostrogothic assault would have been a serious danger. Cp. Marcellinus, \emph{sub} 480. Theoph. A.M. 5971, Brooks, \emph{C. Med. H.} I p476.

} and there was still a Gothic faction in the city. In his absence, Zeno had talked of taking the field himself, and there was much dissatisfaction at his failing to do so. He was

p395
accused of cowardice, but the true reason probably was that he feared not the enemy but his own army.\footnote{Cp. Brooks, \emph{Emp. Zeno and the Isaurians}, p219 . This article, to which I am under considerable obligations, has cleared up many difficulties in the chronology, and elucidated the whole story.

}

The treatment of Verina supplied a pretext to her son-in\texthebrew{‑}law, Marcian, to attempt to overthrow Zeno (end of A.D. 479).\texthebrew{⁠}\footnote{Date in John Ant. \emph{ib.} 3. This author and Eustathius (in

Evagr. III.26 ) are the fullest sources.

} Marcian, who was son of Anthemius, the western Emperor, had married Leontia, Leo's younger daughter, and claimed that he had a better right to the throne than Zeno, because his wife had been born in the purple. This claim, according to the theory of the Imperial succession, was entirely futile, but it illustrates how the idea that children born in the purple had a natural title to the throne was beginning to grow. The barbarians in the city rallied round Marcian and his brother Procopius,\texthebrew{⁠}\footnote{There was also a third brother Romulus (Theodore Lect. I.37). The chief barbarian associated was Busalbos, an officer \texthebrew{—} perhaps commander of one of the legions of the troops in praesenti.

} and the citizens were on their side. The brothers united their forces near the house of Caesarius, to the south of the Forum of Theodosius;\texthebrew{⁠}\footnote{\textgreek{Τ\texthebrew{ὴ}ν Καισαρίου ο\texthebrew{ἰ}κίαν}, to be identified with \textgreek{τ\texthebrew{ὰ} Καισαρίου} (of which there was a curator), in Theophanes A.M. 6054 (A.D. 561\texthebrew{‑}562). Evidently near the harbour of Caesarius, and this is confirmed by its proximity to the Forum of Bous, which we can infer from the passage in Theoph.

} and then one of them marched upon the palace, while the other attacked the house of Illus.\texthebrew{⁠}\footnote{\textgreek{Κατ\texthebrew{ὰ Ἰ}λλο\texthebrew{ῦ ἐ}ν το\texthebrew{ῖ}ς λεγομένοις Ο\texthebrew{ὐ}αράνου} (John Ant.). I can find no trace of this locality elsewhere.

} The Emperor nearly fell into their hands,\texthebrew{⁠}\footnote{The Stoa of the Delphax was attacked. This was evidently in the Palace, as indeed is expressly stated by Victor Tonn. \emph{sub} 523 intra palatium loco quod \textgreek{δελφακα} graeco vocabulo dicunt. Cp. Procopius, \emph{B. V.} 1.21 \textgreek{\texthebrew{ὅ}πη βασιλέως ε\texthebrew{ἶ}ναι στιβάδα ξυμβαίνει Δέλφικα το\texthebrew{ῦ}το καλο\texthebrew{ῦ}σι τ\texthebrew{ὸ} ο\texthebrew{ἴ}κημα.}} and during the day the rebels were victorious against the Imperial soldiers, on whose heads the citizens showered missiles from the roofs. But under the cover of night, Illus introduced into the city an Isaurian force from Chalcedon, and the next day Marcian's party was defeated. Marcian was ordained a priest and banished to Cappadocia; Leontia fled to a convent.\texthebrew{⁠}\footnote{Marcian escaped, and attacked Ancyra, but was defeated by Trocundes and imprisoned along with his wife in an Isaurian fortress.

} Theoderic Strabo was in league with Marcian, but did not reach the city in time to help him.

It was perhaps not long after this that the Empress Ariadne entreated Zeno to recall her mother. Zeno told her to ask Illus. The Empress sent for Illus and implored him with tears to

p396
release her mother. And Illus said, ``Why do you want her? Is it that she may set up another Emperor against your husband'' Then Ariadne said to Zeno, ``Is Illus to be in the Palace or I?'' and he replied, ``Do what you can. I prefer you.'' She suborned Sporacius, one of the Scholarian guards, to assassinate Illus, and the attempt was made, on the occasion of a spectacle in the Hippodrome, as Illus was walking through The Pulpita behind the Kathisma. The assassin's sword, aimed at the head, cut off the minister's right ear, and he was hewn to pieces on the spot.\texthebrew{⁠}\footnote{John Mal. XV.387 \emph{sqq.} For the position of the Pulpita cp. Theoph. A.M. 6024, p185, 10.

} Illus did not believe Zeno's asseverations that he was ignorant of the plot, and when the wound was healed he requested the Emperor to allow him to go to the East for change of air. Zeno relieved him of the duties of Master of Offices and appointed him Master of Soldiers in the East. Illus proceeded to Antioch, taking with him a considerable number of friends and adherents (481\texthebrew{‑}482), including Marsus and the pagan quaestor Pamprepius.\texthebrew{⁠}\footnote{Evagr. III.27 . Brooks, 225\texthebrew{‑}226 , gives reasons for thinking that Evagr. (Eustathius) and John Mal. are mistaken in supposing that Leontius also accompanied Illus. Zeno granted to Illus the special power of appointing dukes in the eastern provinces.

} Soon afterwards the patrician Leontius seems to have been sent to Antioch demanding the release of Verina, but Illus won him over to his interests and he did not return to Constantinople.\texthebrew{⁠}\footnote{See Joshua Styl. c14.

} The estrangement of the Emperor from his general was now complete, and a contest between the two Isaurians was inevitable. Illus and his party hoped to secure Egypt for their cause, and attempted, but without success, to take advantage of the ecclesiastical disputes which were at this time dividing Alexandria.\texthebrew{⁠}\footnote{See Asmus, ``Pamprepios,'' in \emph{B. Z.} XXII.332 \emph{sqq.} Pamprepius was sent to Alexandria, to combine measures with John Talaias who ascended the Patriarchal throne in 482, but in the same year (June) was deposed and succeeded by Peter Mongus, who was supported by Zeno. Peter organised an anti-pagan demonstration, and Pamprepius had to flee.

} The hostilities of the Ostrogoths prevented Zeno from taking any measures before the end of A.D. 483, or the spring of 484. When his hands were at last free, he commanded Illus to surrender Longinus (Zeno's brother) who had been a prisoner for many years. Illus refused, and Zeno deposed him from his command of the eastern army and appointed John the Scythian in his stead. At the same time he expelled the friends of Illus from Constantinople, confiscated

p397
their property, and bestowed it upon the cities of Isauria. War ensued and lasted for about four years.

Illus had employed the two years which he spent at Antioch (482\texthebrew{‑}484) in making himself popular and gaining friends. He counted, for the coming struggle, on the support of the orthodox adherents of the Council of Chalcedon, who had been displeased be an ecclesiastical decree (the Henotikon ) in which Zeno had expressly declined to maintain the dogmas of that assembly ( A.D. 481). He may also have hoped for some help from pagans. He was very intimate with the pagan philosopher Pamprepius, who had been appointed Quaestor through his influence, and had accompanied him to Antioch. Deciding not to assume the purple himself, Illus drew from his Isaurian prison the ex-tyrant Marcian, and proclaimed him Emperor. He had sought the assistance of the Patrician and king Odovacar in Italy; he had written to the Persian monarch Piroz and to some of the satraps of Roman Armenia. Odovacar refused; the Persians and Armenians promised help when the time came. A great defeat which the Persians suffered at the hands of the Ephthalites (January, A.D. 484; Piroz was slain) rendered it impossible for them to fulfill their promise.

Zeno sent an Isaurian force against the rebels.\texthebrew{⁠}\footnote{Under Conon, a fighting parson (he was bishop of Apamea), and Linges, a bastard brother of Leo.

} About the same time Illus changed his plans, and entered into an alliance with his old enemy the Emperor Verina who was still languishing in an Isaurian fortress.\texthebrew{⁠}\footnote{The castle of Papirios, to which she had been removed (cp. Theodore Lector, I.37). It seems to be the same as the fortress of Cherris (Brooks, \emph{ib.} 228) .

} He brought her to Tarsus, arrayed her in imperial robes; and it was decided to set aside Marcian,\texthebrew{⁠}\footnote{We are not told why Illus desired the co-operation of Verina to invest the rebellion with the prestige of legitimacy, and we may conjecture that he thought the association of the Empress with her son-in\texthebrew{‑}law Marcian would be too dangerous a combination.

} and to proclaim as Emperor the patrician Leontius. Verina crowned him Emperor, and a proclamation in her name was sent through the provinces of the East and Egypt. In this document she claims that the Empire belongs to her, that it was she who conferred it upon Zeno, and that now, since his avarice is ruining the state, she has determined to transfer it to the pious Leontius.\footnote{See John Mal. \emph{fr.} 35, \emph{De ins.} p165. Brooks well notices that the insistence in this document on the piety of Leontius alludes to Zeno's \emph{Henotikon} (\emph{op. cit.} 227) . Theophanes

(p398)
gives the date of the entry of Leontius into Antioch as June 27 ind. 7 (= 484). But a contemporary, Palchus the astrologer, gives the day of the coronation as July 19. See Cumont, ``L'Astrologue Palchos,'' in \emph{Revue de l'instruction publique en Belgique}, XL p1206, and though Palchus was mistaken in pla­cing the coronation at Antioch, his date must be accepted. If we correct \emph{June} in Theoph. to \emph{July}, Leontius entered Antioch only a week after the proclamation at Tarsus. The horoscope of Leontius given by Palchus was drawn incompletely by two astrologers, of whom one no doubt was Pamprepius. They inferred his success. Palchus shows that they over­looked certain data, which would have led to a true prognostication. Leontius appointed Pamprepius as his Master of Offices. For this and other appointments see John Mal. in \emph{De insidiis}, p165. For coins of Leontius minted at Antioch see Sabatier, I. Pl. VIII.22, 23.

}

p398
The new Emperor was received at Antioch,\texthebrew{⁠}\footnote{He was rejected at Chalcis, and at Edessa.

} and the rebellion spread. The Isaurian troops which Zeno had sent were obviously unable to cope with it, and Zeno sought the hope of Theoderic the Amal and his Ostrogoths. Theoderic, as Master of Soldiers in praesenti , joined the army of John the Scythian, and though he was recalled almost immediately, his followers seem to have remained and taken part in the campaign.\texthebrew{⁠}\footnote{The part played by Theoderic and the Ostrogoths is uncertain. Cp. John Ant. \emph{fr.} 214.4 and 6; Brooks, 228 .

} Rugian auxiliaries were also sent under the command of Aspar's son Ermenric. A battle was fought, the forces of Zeno were victorious, and Illus, Leontius, and Verina, with all their chief partisans, fled to the strong fortress of Cherris\texthebrew{⁠}\footnote{The fortress had been well supplied and strengthened by Zeno, as a place of refuge for himself in case of eventualities (Joshua Styl. c12). Art had assisted its natural strength. There was no path leading up to it save one so narrow that not even two persons could ascend at once (\emph{ib.} c17).

} in the Isaurian mountains (autumn, A.D. 484). The Empress died in a few days. The cause of Illus was now hopeless, but the fortress held out for nearly four years. It was taken by treachery (488), and Illus and Leontius were beheaded.\footnote{The protraction of the siege was partly due to the distraction of Theoderic's rebellion in 486, partly to the strength of the fortress. Illus made some proposals for peace about this time.ime. But he had fallen into despondency, and occupied himself with reading, committing the command of the garrison to Indacus. It was Indacus who betrayed the fortress. (According to Theoph., the husband of the widow of Trocundes was sent by Zeno to the fortress and arranged the treachery. Source, Theodore Lector(?).) Pamprepius was put to death by his friends before the end of the reign, because he had falsely foretold success.

}

The struggle between Illus and Zeno derives particular interest from the association of Illus with the prominent pagans who still flourished at Athens, Constantinople, and Alexandria. These men seem to have hoped that Illus, if victorious, would be able to secure public toleration for paganism.\texthebrew{⁠}\footnote{The full significance of this element in the rebellion of Illus has been brought out by J. R. Asmus, in

(p399)
his article on ``Pamprepios'' (\emph{B. Z.} XXII.320 \emph{sqq.}). The principal evidence is in the fragmentary \emph{Vita Isidori}, of Damascius (on which see Asmus, \emph{ib.} XVIII.424 \emph{sqq.}, and XIX.265 \emph{sqq.}), and the art. of Suidas, \textgreek{Παμπρέπιος}. There is an interesting statement in Zacharias, \emph{Vita Severi}, p40, that pagans in Caria (at Aphrodisias) offered sacrifices to the gods and inquired of the entrails of the victims whether Leontius, Illus, and Pamprepius would be victorious over Zeno. One of these Carians was the distinguished physician and philosopher, Asclepiodotus, a pupil of Proclus.

} It was impossible,

p399
of course, to stamp the movement with a pagan character. If Illus had come forward as a new Julian, he would have had no following. But there is little doubt that he was personally in sympathy with the ``Hellenes''; he was a man with intellectual interests and was inclined to the Neoplatonic philosophy. His close intimacy with the pagan savant, Pamprepius of Panopolis, who shared his fortunes, proves this. Pamprepius, who is described as swarthy and ugly, went in his youth from Egypt to the university of Athens, where he studied under the philosopher Proclus and was appointed professor of grammar (literature and philology). A quarrel with a magistrate forced him to leave Athens, and he betook himself to Constantinople, where pagans of talent, if they behaved discreetly, could still find a place.\texthebrew{⁠}\footnote{Ammonius of Alexandria seems to have taught philosophy at Constantinople in the reign of Zeno (cp. Asmus, ``Pamprepios,'' p326).

} At the request of Illus he delivered a lecture, probably explaining the doctrines of Neoplatonism, and Illus procured his appointment as professor of grammar at the university. He established himself in the favour of Illus by public recitation of a poem,\texthebrew{⁠}\footnote{Perhaps the \emph{\textgreek{\texthebrew{Ἰ}σαυρικά}}, mentioned by Suidas.

} in reward for which he received a pension. But when Illus was absent in Isauria ( A.D. 478), his enemies seized the opportunity to attack Pamprepius as a pagan and a sorcerer. He was banished from the city and retired to Pergamum; but Illus summoned him to Isauria, and then brought him back in triumph, and procured his appointment to the high post of Quaestorship. Henceforward his fortunes were bound up with those of Illus, to whom he acted as confidant and adviser throughout the struggle for the throne. The pagans blamed Pamprepius for the failure of the movement, and represented him as a traitor to the cause of his chief. But we may take it as certain that this charge was false, and that he was slain not because he was suspected of treachery, but because his prophecies had not come true and he had proved himself a blind guide.\footnote{See the conclusions of Asmus, \emph{op. cit.}

}

p400
The greater part of Zeno's reign had been troubled on the one hand by the hostile risings of the Ostrogoths, which have still to be described, and on the other by rebellion. In 488 both these troubles were terminated by the departure of the Goths from Italy and by the final suppression of Illus. The Emperor persisted in his policy of firmly establishing Isaurian predominance. His brother Longinus, who had managed to escape from his prison,\texthebrew{⁠}\footnote{In A.D. 485. Perhaps he had been set free by Illus, with the design of conciliating Zeno.

} was consul twice and princeps of the Senate.\texthebrew{⁠}\footnote{It is possible that he was also created mag. mil. in praes., and continued to hold this office in the first year of Anastasius; see

\emph{C. J.} XII.37.16 . Cp. John Mal. XV p386. The dates of his consul­ships are 486 and 490.

} Kottomenes had been appointed Master of Soldiers in praesenti , instead of Theoderic, in 484, and Longinus of Kardala at the same time became Master of Offices; both these men were Isaurians.\footnote{John Ant. \emph{fr.} 98, \textgreek{Καρδάλων} is the reading of the cod. Scorialensis, \textgreek{Καρδάμων} of the Parisinus.

}

A modern historian who was perhaps the first to say a good word for Zeno, observes that ``the great work of his reign was the formation of an army of native troops to serve as a counterpoise to the barbarian mercenaries''; and goes on to remark that the man who successfully resisted the schemes and forces of the great Theoderic cannot have been contemptible.\texthebrew{⁠}\footnote{Finlay, \emph{History of Greece}, vol. I p180.

} And even from the pages of a hostile contemporary writer\texthebrew{⁠}\footnote{Preserved in Suidas \emph{s.v.}\textgreek{Ζήνων} (probably from Malchus).

} we can see that he was not so bad as he was painted. He is said to have been in some respects superior to Leo, less relentless and less greedy. He was not popular,\texthebrew{⁠}\footnote{Yet the Ravenna chronicler known as the Anonymus Valesii represents him as very popular: Zeno recordatus est amore senatus et populi, munificus omnibus se ostendit, ita ut omnes ei gratias agerent. Senato Romano et populo tuitus est, ut etiam ei imagines per diversa loca in urbe Roma levarentur. Cuius tempora pacifica fuerunt.

(9.44) . One would think that the writer was an Isaurian. Compare also

9.40 : in republica omnino providentissimus, favens gentis suae.

} for his ecclesiastical policy of conciliation did not find general favour, and he was an Isaurian. But he was inclined to be mild; he desired to abstain from employing capital punishment. In the first year of his reign, Erythrius was Praetorian Prefect, a very humane man, who, when he saw that sufficient revenue could not be raised without severe oppression, resigned his office.\texthebrew{⁠}\footnote{Suidas \emph{s.v.}\textgreek{\texthebrew{Ἐ}ρύθριος} = Malchus \emph{fr.} 6. Erythrius seems to have succeeded Epinicus in 475; his tenure of office must have been very short. No extant constitutions are addressed to him. It is also possible that he was prefect in the last months of 476 after Zeno's restoration.

} In fiscal administration Zeno

p401
was less success ­ful than his predecessors and his successor Anastasius. We are told that he wasted all that Leo left in the treasury by donatives to his friends and inaccuracy in checking his accounts. In A.D. 477 the funds were very low, hardly sufficient to supply pay for the army. But the blame of this may rather rest with Basiliscus, who, reigning precariously for twenty months, must have been obliged to incur large expenses, to supply which he was driven to extortion, and in the following years the Ostrogoths were an incubus on the exchequer; while we must further remember that since the enormous outlay incurred by Leo's naval expedition the treasury had been in financial difficulties, which only a ruler of strict economy and business habits, like the succeeding Emperor Anastasius, could have remedied. Zeno was not a man of business, he was indolent and in many respects weak. Yet it is said that his reign would have been a good one but for the influence of the Praetorian Prefect Sebastian, who succeeded Erythrius, and introduced a system of selling offices.\texthebrew{⁠}\footnote{It is said that Sebastian used to buy for a small amount an office which Zeno bestowed on a friend, and then sell it to some one else for a much higher price, Zeno receiving the profit. He was Praetorian Prefect from 477 to 484. The decline of the Scholarian guards is attributed by Agathias

(V.15)

to Zeno, who bestowed appointments on Isaurian relatives of no valour.

} Of Sebastian we otherwise hear very little.

By his first wife Arcadia, Zeno had a son,\texthebrew{⁠}\footnote{Suidas, \emph{s.v.}\textgreek{Ζήνων}, probably from Malchus, see

\emph{F. H. G.} IV.118 .

} of the same name, whose brief and strangely disreputable career must have been one of the chief scandals at the court. His father desired that he should be carefully trained in manly exercises, but unscrupulous young courtiers, who wished to profit by the abundant supplies of money which the boy could command, instructed him in all the vulgar excesses of luxury and voluptuousness. They introduced him to boys of his own age, who did not refuse to satisfy his desires, while their adulation flattered his vanity to such a degree that he treated all who came in contact with him as if they were servants. His excesses brought on an internal disease, and he died prematurely, after lying for many days in a senseless condition. After his death, Zeno seems to have intended to devolve the succession upon his brother Longinus, who enjoyed a vile reputation for debauchery.\texthebrew{⁠}\footnote{See Suidas \emph{s.v.}\textgreek{Λογγ\texthebrew{ῖ}νος} (perhaps from Malchus).

} We have already

p402
seen how he was advanced to high posts of dignity. It is related that Zeno consulted a certain Maurianus, skilled in occult learning, who informed him that a silentiarius would be the next Emperor and would marry Ariadne. This prophecy was unfortunate for a distinguished patrician of high fame named Pelagius, who had once belonged to the silentiarii, for Zeno, seized with alarm and suspicion, put him to death.\texthebrew{⁠}\footnote{John Mal. XV p390. Arcadius, the Praetorian Prefect, expressed such indignation at this that Zeno sought to slay him, but Arcadius sought refuge in St. Sophia and escaped with the confiscation of his property.

} The Emperor's unpopularity naturally made him suspicious, and he was in bad health. An attack of epilepsy carried him off on April 9, A.D. 491.

The doctrinal decrees of Chalcedon were the beginning of many evils for the eastern provinces of the Empire. Theological discord, often accompanied by violence, rent the Church, and the Emperors found it utterly impossible to suppress the Monophysite, as they had suppressed the Arian, faith. In Alexandria, the monks and the majority of the population were devoted to the doctrine of One Nature, and on the death of Marcian the smouldering fire of dissatisfaction burst into flame. Timothy Aelurus,\texthebrew{⁠}\footnote{Said to have been called Aelurus or Cat, because he used to creep at night into the cells of the monks at Alexandria to incite them against Proterius (Theodore Lector, I.1). The view that he was a Herul (\textgreek{α\texthebrew{ἰ}λουρός} being a corruption of \textgreek{\texthebrew{Ἕ}ρουλος}) is not probable.

} an energetic Monophysite, was set up as rival Patriarch; Proterius was murdered in the baptistery ( A.D. 457, Easter) and his corpse was dragged through the city. Timothy sent a memorial to the Emperor Leo demanding a new Council, and Leo formally asked for the opinion of the bishops of Rome, Constantinople, Antioch, and Jerusalem, and other leading dignitaries of the Church.\texthebrew{⁠}\footnote{Fifty-five bishops, Simeon Stylites the younger, and two other monks.

} They condemned the conduct of Timothy and he was banished to the Chersonese.\texthebrew{⁠}\footnote{A.D. 460. He was succeeded by Timothy Salophaciolus (said to mean \emph{white-capped}), who retired to a monastery in 475, when the other Timothy returned, and on his death was reinstated in 477. He died in 482.

} At Antioch, the part of Timothy was played by Peter the Fuller, who during the reign of Leo was twice raised to the Patriarchal throne and twice ejected.\footnote{Theodore Lector, I.20; Liberatus, \emph{Brev.} c18.

}

p403
When Basiliscus ascended the throne, the Monophysite cause looked bright for a few months. Peter and Timothy were reinstated, and Basiliscus issued an Encyclical letter\texthebrew{⁠}\footnote{A.D. 476. This Encyclical will be found in

Evagrius, III.4 ; the Anti-encyclical, \emph{ib.} 7 (cp. Zachariah Myt. V.5); Zeno's Henotikon, \emph{ib.} 14.

} in which he condemned the Council of Chalcedon and the Tome of Leo. But this declaration raised a storm in Constantinople which he was unable to resist. The monks were up in arms, and the Patriarch Acacius,\texthebrew{⁠}\footnote{Elected 471, as successor to Gennadius, who succeeded Anatolius in 458.

} who was not a man of extreme views, found himself forced to oppose the Emperor's policy. Basiliscus hastened to retract, and he issued another letter, which was known as the anti-encyclical . But the settlement of the ecclesiastical struggle did not lie with him. Zeno returned, and a new policy was devised for restoring peace to the Church. His chief advisers here were Acacius and Peter Mongus, who had been the right-hand man of Timothy Aelurus. The policy was to ignore the Council of Chalcedon, but not to affirm anything contrary to its doctrine; and the hope was that the Monophysites and their antagonists would agree to differ, and would recognise that a common recognition of the great Councils of Nicaea and Constantinople was a sufficient bond of communion.

The \emph{Henotikon}, a letter addressed by the Emperor to the Church of Egypt, embodied this policy ( A.D. 481). It anathematises both Nestorius and Eutyches; declares the truth, and asserts the sufficiency, of the doctrine of Nicaea and Constantinople; and anathematises any who teach divergent doctrine ``at Chalcedon or elsewhere.'' As the document was intended to conciliate all parties, it was a blunder to mention Chalcedon; for this betrayed that the theological leanings of those who framed it were not favourable to the Chalcedonian dogma. The Monophysites gladly accepted it;\texthebrew{⁠}\footnote{Except an extreme party who were known as Akephaloi or ``Headless.''

} interpreting it as giving them full liberty to denounce Chalcedon and the \emph{Tome} of Pope Leo.

It is to be noted that Basiliscus by his \emph{Encyclical} and Zeno by his \emph{Henotikon} asserted the right of the Emperor to dictate to the Church and pronounce on questions of theological doctrine. They virtually assumed the functions of an Ecumenical Council. This was a claim which the see of Rome was not ready to admit except for itself. After the interchange of angry letters between

p404
Pope Simplicius and Acacius, a synod was held at Rome,\texthebrew{⁠}\footnote{A.D. 484 under Felix II, successor of Simplicius. One of the Sleepless monks of Studion pinned the sentence of excommunication on the back of Acacius as he was officiating in St. Sophia. Acacius retorted the sentence on the Pope.

} and Acacius and Peter Mongus, who was now Patriarch of Alexandria, were excommunicated.\footnote{After the death of Timothy Salophaciolus in 482, there was a struggle for the Patriarchal throne between Peter and John Talaias. Peter was supported by Zeno, and John, who was actually consecrated, betook himself to Rome and appealed to Simplicius.

}

The general result of the \emph{Henotikon} was to reconcile moderate Monophysites in Egypt and Syria, and to secure a certain measure of ecclesiastical peace in the East for thirty years\texthebrew{⁠}\footnote{Nominally till A.D. 518, but after A.D. 512 the spirit of the Henotikon did not prevail in the East (see below,

Chap. XIII § 2 ). Various views are held by modern writers of the Henotikon. Gelzer praises it unreservedly; Harnack considers it unfortunate, but admits that Zeno ``simply did his duty'' in issuing it (\emph{op. cit.} p228).

} at the cost of a schism with the West. But the extreme Monophysites were not reconciled to the policy of Acacius and Peter.

After the death of Olybrius, Leo was sole Roman Emperor for more than four months, and the Burgundian Gundobad, who had succeeded his uncle Ricimer as Master of Soldiers, directed the conduct of affairs in Italy. On March 5, A.D. 473, Glycerius, Count of the Domestics, was proclaimed Emperor at Ravenna ``by the advice of Gundobad,''\texthebrew{⁠}\footnote{Cassiodorus, \emph{Chron.}, Gundibato hortante. Marcellinus, \emph{Chron.}, Glycerius apud Ravennam plus praesumptione quam electione Caesar factus est (this was the view at Constantinople). John Ant. \emph{fr.} 92. For date see Anon. Cusp.

} just as Severus had been proclaimed in the same city by the advice of Ricimer. Of this Augustus, whose reign was to be brief, one important public act is recorded. Italy was threatened by an invasion of Ostrogoths who, under the leader ­ship of Widemir, began to move from Pannonia, but the diplomacy of Glycerius averted the storm, so that it fell on Gaul.

The election of Glycerius was not approved at Constantinople, and Leo selected another as the successor of Anthemius.\texthebrew{⁠}\footnote{John Ant. \emph{ib.}

} His choice was Julius Nepos, husband of the niece of the Empress, and military governor of Dalmatia, where he had succeeded his

p405
uncle, count Marcellinus.\texthebrew{⁠}\footnote{His parents were Nepotianus and a sister of Marcellinus.

} We do not hear that any resistance was offered to Nepos, who arrived in Italy, probably escorted by eastern troops; and it was not long before Gundobad, whether perforce or voluntarily, retired to Burgundy where, in the following year, he succeeded his father as one of the Burgundian kings.\texthebrew{⁠}\footnote{Cp. Schmidt, \emph{op. cit.} I.380\texthebrew{‑}381.

} Glycerius was deposed, and at Portus, the town at the mouth of the Tiber, he was ordained bishop of Salona.\texthebrew{⁠}\footnote{Anon. Val.

factus est episcopus; Marcellinus, \emph{Chron.}, in portu urbis Romae ex Caesare episcopus ordinatus est et obiit, where the form of expression suggests a doubt whether Glycerius ever reached Salona.

} Nepos was proclaimed Emperor and ruled at Rome (June 24, A.D. 474). Once more two Augusti reigned in unison.

To the vacant post of Master of Soldiers, which carried with it almost as a matter of course the title of Patrician, Orestes was appointed. This was that Orestes who had been the secretary of Attila, and he had married the daughter of a certain count Romulus. Possessing the confidence of the German troops he determined to raise his son to the Imperial throne.

We are told that Nepos, driven from Rome, went to Ravenna and, fearing the coming of Orestes, crossed over to Salona. This was on August 28, A.D. 475. The same year that saw the flight of Zeno from Constantinople saw the flight of Nepos from Ravenna. At Salona he lived for five years, and his Imperial authority was still recognised in the East and in Gaul. But in Italy the Caesar Julius was succeeded by the Caesar Augustulus, for so the young Romulus was mockingly nicknamed, whom his father Orestes invested with the Imperial insignia on October 31. These names, Julius, Augustulus, Romulus, in the pages of the chroniclers, meet us like ghosts re-arisen from past days of Roman history.\footnote{Am. Thierry made a similar remark. ``Ces rapprochements fortuits présentaient dans leur bizarrerie je ne sais quoi de surnaturel qui justifiait la crédulité et troublait jusqu'aux plus fermes esprits : on baissa la tête et on se tut.'' (\emph{Les Derniers Temps de l'empire d'occident}, p258).

}

It is important to remember that the position of Romulus was not constitutional inasmuch as he had not been recognised by the Emperor at Constantinople, in whose eyes Nepos was still the Augustus of the West. For twelve months Orestes ruled Italy in the name of his son. His fall was brought about by a mutiny of the troops. The army, which the Master of Soldiers commanded, seems to have consisted under Ricimer and

p406
his successors almost exclusively of East Germans, chiefly Heruls, also Rugians and Scirians. According to the usual custom,\texthebrew{⁠}\footnote{See above,

p206 .

} they were quartered on the Italians. But they were weary of this life. They desired to have roof-trees and lands of their own, and they petitioned Orestes to reward them for their services, by granting them lands and settling them permanently in Italy on the same principle on which various German peoples had been settled in other provinces. They did not demand the exceptionally large concession of two-thirds of the soil which had been granted by Honorius to the Visigoths; they asked for the only grant of one-third which had been assigned, for instance, to the Burgundians. But such a settlement in Italy was a very different thing from settlement in Gaul or Spain, and Orestes, notwithstanding his long association with Germans and Huns, was sufficiently Roman to be determined to keep the soil of Italy inviolate. He rejected the demand. The discontented soldiers found a leader in the Scirian Odovacar, one of the chief officers of Orestes.\texthebrew{⁠}\footnote{For the nationality of Odovacar see John Ant. 93,

Anon. Val. 45 . He was son of Edica, probably identical with Edeco, who acted as Attila's envoy to Byzantium in 448. His brother was Onoulf (Malchus, \emph{fr.} 8, John Ant. \emph{ib.}).

}

Ticinum

to which Orestes retired was easily taken, and the Patrician was slain at Placentia (August 28, A.D. 476). ``Entering Ravenna, Odovacar deposed Augustulus but granted him his life, pitying his infancy and because he was comely, and he gave him an income of six thousand solidi and sent him to live in Campania with his relatives.''\footnote{Anon. Val. VIII.38 .

}

The soldiers had proclaimed Odovacar king.\texthebrew{⁠}\footnote{He is styled rex Herulorum in \emph{Cons. Ital.} (\emph{Chron. Min.} I p313, cp. p309).

} But it was not as king over a mixed host of various German nationalities that Odovacar thought he could maintain his position in Italy. The movement which had raised him had no national significance, and if he retained the royal title of an East German potentate, it was as a successor of Ricimer, Gundobar, and Orestes that he hoped to govern the Italians. In other words, he had no idea of detaching Italy from the Empire, as Africa and much of Gaul and Spain had come to be detached. The legal position was to continue as before.\texthebrew{⁠}\footnote{He issued silver and bronze coins in his own name at Ravenna, without the title rex. The inscription was FLavius ODOVAC. See Wroth, \emph{Coins of Vandals}, p30.

} But the system of Ricimer was to

p407
be abandoned. There were be no more puppet Emperors in the West; Italy was to be under the sovranty of the Emperor at Constantinople, and its actual government was to be in the hands of Odovacar, who as Master of Soldiers was to be a minister of the Emperor, while he happened at the same time to be king of the East Germans who formed the army.

With this purpose in view Odovacar made the deposition of Romulus take the form of an abdication, and induced the Roman Senate to endorse formally the permanent institution of a state of things which had repeatedly existed in the days of Ricimer. A deputation of senators, in the name of Romulus, was sent to the Augustus at Constantinople to announce the new order of things. Zeno had already recovered the throne, from which Basiliscus had driven him, when the ambassadors arrived and informed him that they no longer needed a separate Emperor but that his sole supremacy would be sufficient; that they had selected Odovacar as a man capable of protecting Italy, being both a tried soldier and endowed with political intelligence. They asked Zeno to confer upon him the rank of Patrician and entrust him with the administration of Italy. They bore with them the Imperial insignia which Romulus had worn ( A.D. 477).\footnote{These details are preserved in a valuable fragment of Malchus (\emph{fr.} 10). Candidus relates that after the death of Nepos the Gallo-Romans (\textgreek{τ\texthebrew{ῶ}ν δυσμικ\texthebrew{ῶ}ν Γαλατ\texthebrew{ῶ}ν}) rejected the rule of Odovacar and sent an embassy to Zeno, but Zeno rather inclined to Odovacar

(\emph{fr.} 1, p136, \emph{F. H. G.} IV) .

}

At the same time messengers arrived from Nepos to congratulate Zeno on his restoration, to ask for his sympathy with one who had suffered the same misfortune as he, and to crave his aid in men and money to recover the throne. But for the existence of Nepos, the situation would have been simple. Zeno could not ignore his legal right, but was not prepared to support it with an army. He told the representatives of the Senate that of the two Emperors they had received from the East, they had slain Anthemius and banished Nepos; let them now take Nepos back. But he granted the other request. He sent to Odovacar a diploma conferring the Patriciate, and wrote to him, praising the respect for Rome and the observance of order which had marked his conduct, and bidding him crown his goodness by acknowledging the exiled Emperor. The fact that Verina was the aunt of the wife of Nepos was a consideration which helped to hinder Zeno from disowning him. Odovacar

p408
did not acknowledge the claim of Nepos, and Zeno cannot have expected that he would.

The events of A.D. 476 have been habitually designated as the ``Fall of the Western Empire.'' The phrase is inaccurate and unfortunate, and sets the changes which befell in a false light. No Empire fell in A.D. 476; there was no ``West Empire'' to fall. There was only one Roman Empire, which sometimes was governed by two or more Augusti. If it is replied that expression is merely a convenient one to signify what contemporary writers sometimes called the Hesperian realm ( Hesperium regnum ), the provinces which had been, since the death of Theodosius I, generally under the separate government of an Emperor residing in Italy, and that all that is meant is the termination of this line of western Emperors, it may be pointed out that A.D. 480 is in that case the significant date. For Julius Nepos, who died in that year, was the last legitimate Emperor in the West; Romulus Augustulus was only a usurper. The important point to seize is that, from the constitutional point of view, Odovacar was the successor of Ricimer, and that the situation created by the events of A.D. 476 was in this respect similar to the situation in the intervals between the reigns of the Emperors set up by Ricimer. If, on the death of Honorius, there had been no Valentinian to succeed him, and if Theodosius II had exercised the sovranty over the western provinces, and if no second Augustus had been created again before the western provinces had passed under the sway of Teutonic rulers, no one would have spoken of the ``Fall of the Western Empire.'' Yet this hypothetical case would be formally the same as the actual event of A.D. 476 or rather of A.D. 480. The West came finally, as it had more than once temporarily, under the sole sovranty of the Emperor reigning at East Rome.

The Italian revolution of A.D. 476 was, however, a most memorable event, though it has been wrongly described. It stands out prominently as an important stage in the process of the dismemberment of the Empire. It belongs to the same catalogue of chronological dates which includes A.D. 418, when Honorius settled the Goths in Aquitaine, and A.D. 435, when Valentinian ceded African lands to the Vandals. In A.D. 476 the same principle of disintegration was first applied to Italy. The settlement of Odovacar's East Germans, with Zeno's acquiescence,

p409
began the process by which Italian soil was to pass into the hands of Ostrogoths and Lombards, Franks and Normans. And Odovacar's title of king emphasised the significance of the change.

It is highly important to observe that Odovacar established his political power with the co-operation of the Roman Senate, and this body seems to have given him their loyal support throughout his reign, so far as our meagre sources permit us to draw inferences. At this time the senators who counted politically belonged to a few old and distinguished clans, possessing large estates and great wealth, particularly the Decii and the Anicii.\texthebrew{⁠}\footnote{There is a useful genealogical tree in Sundwall, \emph{Abh. zur Gesch. d. ausg. Römerthums}, p131, showing the relation­ships of the Decii who played a public part from 450 to 540.

} The leading men of these families received high honours and posts under Odovacar. Basilius, Decius, Venantius, and Manlius Boethius held the consul ­ship and were either Prefects of Rome or Praetorian Prefects;\texthebrew{⁠}\footnote{Flavius Caecina Decius Maximus \emph{Basilius} iunior was consul in 480, Praet. Pref. 483; Caecina Mavortius Basilius \emph{Decius} iunior, consul 486, Prefect of Rome, and then Praet. Pref. between 486 and 493; Fl. Decius Marius Basilius \emph{Venantius}, consul and Prefect of Rome 484; Flavius Manlius Boethius, consul and Prefect of Rome for the second time in 487, and Praet. Pref. earlier. Cp. CIL V.8120 .

} Symmachus and Sividius were consuls and Prefects of Rome;\texthebrew{⁠}\footnote{Quintus Aurelius Memmius Symmachus iunior, consul 485, Prefect of Rome probably in same year. Rufius Achillius Sividius, consul 488, and twice Prefect of Rome; cp. CIL XII.133 . A bronze tablet of Symmachus (Dessau, 8955) combines the names of Zeno and Odovacar: salvo d.n. Zenone et domno Odovacre.

} another senator of old family, Cassiodorus, was appointed a minister of finance.\texthebrew{⁠}\footnote{First com. r. pr. afterwards com. s. larg. Petrus Marcellinus Felix Liberius began, under Odovacar, a career which was to be long and distinguished, but we do not know what posts he held (cp. Cass. \emph{Var.} II.16). Those parts of the Imperial domains which were appropriated to the Emperor's private purse and were taken over by Odovacar, were placed under an official entitled comes et vicedominus noster; and this post might be held by a German. See Marini, \emph{Pap.} n82 (grant of land to count Pierius, A.D. 489). These patrimonial lands were chiefly in Sicily.

} The evidence indicates that while it was Odovacar's policy to appoint only men of Roman families to the Prefecture of the City, he allowed the Prefect to hold office only for a year, so that no man might win a dangerous political importance.\footnote{Sundwall, \emph{op. cit.} 181.

}

Yet the Roman nobility were now compelled to contribute more largely to the maintenance of the military forces which defended Italy. The greater part of the land belonged to them, and by the new settlement one-third of their estates was taken

p410
from the proprietors, and Odovacar's barbarian soldiers and their families were settled on them. It is not probable that the number of these soldiers exceeded 20,000 at the most, and it has been reasonably doubted whether this measure was actually carried out throughout the length and breadth of the peninsula.\texthebrew{⁠}\footnote{Cp. the remarks of Sundwall, 178, and 183, n4.

} We may suspect that the needs of the army were satisfied without a drastic application of the principle of partition. If the illustrious landowners had been mulcted on a large scale, it is hardly credible that they would have co-operated with the king as loyally as they seem to have done.

Soon after the government of Italy had passed into his hands, Odovacar's diplomacy achieved a solid success by indu­cing Gaiseric, who died in January, A.D. 477, to cede to him the island of Sicily. He undertook indeed to pay for it a yearly tribute, and the Vandal king reserved a foothold in the island, doubtless the western fortress of Lilybaeum.\texthebrew{⁠}\footnote{Victor Vit. I.4.

} The death of Julius Nepos has been mentioned. He was murdered by two of his retainers in his country house near Salona in May, A.D. 480. Odovacar assumed the duty of pursuing and executing the assassins, and at the same time established his own rule in Dalmatia.\texthebrew{⁠}\footnote{Dalmatia had been under Constantinople since the reign of Valentinian III (see above,

p125 ), and we must suppose that when Nepos was created Augustus in 474, Leo handed it over to him.

} The claims of Nepos, so long as he lived, had embarrassed the relations between Zeno and Odovacar; Zeno's acquiescence in Odovacar's position and the wishes of the Senate had been ambiguous and reserved. The death of Nepos relieved the situation, and there was no longer any difficulty at Constantinople about acknowledging the western consuls whom Odovacar chose. But the relations between the Emperor and his Master of Soldiers in Italy were always strained, and in A.D. 486 there was an open breach.\texthebrew{⁠}\footnote{In this year and 487 the names of the western consuls were not published in the East.

} Though Odovacar did not help the rebel Illus in his revolt, there were negotiations, and Zeno may have been suspicious and alarmed. Odovacar prepared an expedition into the Illyrian provinces, then pressed hard by the Ostrogoths, and Zeno averted it by instigating the Rugians to invade Italy.\texthebrew{⁠}\footnote{John Ant. \emph{fr.} 98, \emph{Exc. de Ins.} p138.

} Odovacar anticipated their attack by marching through Noricum and surprising them in the winter

p411
season (end of A.D. 487) in their territory beyond the Danube. Their king Feletheus and his queen were taken to Italy and beheaded, and with the death of his son, against whom a second expedition was sent, the Rugian power was destroyed.\footnote{Eugippius, \emph{Vita Severini}, c44. This source throws interesting light on the derelict provinces of Noricum, which for thirty years were exposed to the depredations of the Rugians, left unprotected by the Italian government, and virtually governed by St. Severinus.

}

Of the internal government we know little. The Church was unaffected by his rule;\texthebrew{⁠}\footnote{Only once does he seem to have intervened. When the clergy met to elect a Pope in succession to Simplicius in March 483, the Praetorian Prefect appeared on Odovacar's behalf, because Simplicius had urgently implored the king not to allow a new Pope to be elected without his consent.

} as an Arian he held aloof from ecclesiastical affairs. As to the working of the Roman administration under a German ruler, acting as an independent viceroy, and the limitations imposed on his power, we have abundant evidence regarding Odovacar's successor, Theoderic, and when we come to his reign the details will claim our attention.

In the reign of Arcadius the Visigoths had seemed likely to form a kingdom within the Illyrian peninsula, before they invaded Italy and established their home in the west. We shall now see how history repeated itself in the case of the Ostrogoths, how they too almost settled in the lands of the Balkans before they went westward to found a kingdom in Italy.\footnote{The chief sources for the events of this section are fragments of Malchus and John of Antioch, and the \emph{Getica} of Jordanes.

}

It will be remembered that after the collapse of the Hunnic power in A.D. 454 the Ostrogoths, over whom three brothers ruled, Walamir, Theodemir, and Widemir, were allowed by the Emperor Marcian to occupy northern Pannonia, as foederati .\texthebrew{⁠}\footnote{Jordanes, \emph{ib.} 268 , knew how it was apportioned among the three brothers. Theodemir's people were on the Plattensee and eastward towards the Danube.

} After some years they were provoked by the Emperor Leo, who refused to pay an annual sum of 100 pounds of gold which Marcian had granted them; and they ravaged the Illyrian provinces and seized Dyrrhachium. Peace was made in A.D. 461, the money grant was continued, and Theoderic,\texthebrew{⁠}\footnote{Theoderic may have been born about 454\texthebrew{‑}455. He is said to have been eight years old when he was sent to Byzantium. His mother seems to have been a concubine treated with the honours of a wife. Her name in

Anon. Val. XIV.58 , is Ereriliva, but she was a Catholic and took the christian name of Eusebia.

} the son

p412
of Theodemir, was sent as a hostage to Constantinople where he had the advantage of a Roman training. His education, however, in letters appears not to have advanced very fast, for it is said that he was never able to write. During those years his nation was engaged in wars with neighboring German peoples.\texthebrew{⁠}\footnote{As Gasquet (\emph{L'Empire byz.} p67) observes, ``what the barbarians hated most cordially was [not Romans but] other barbarians.'' Jordanes put it otherwise: the Ostrogoths made war cupientes ostentare virtutem (\emph{ib.} 52).

} They won a decisive victory over the Scirians which cost Walamir his life. His section of the Goths passed then under the rule of Theodemir, who had soon to resist a large combination of Scirians, Rugians, Gepids, and others. Both parties applied to the Emperor for support, and Leo, acting against the advice of Aspar who was friendly to the Ostrogoths, sent troops to help the Scirian league. In a sanguinary battle the Goths were victors ( A.D. 469), and their predominance on the Middle Danube was established.\texthebrew{⁠}\footnote{Priscus, \emph{fr.} 17, \emph{De leg. gent.}; John Ant. \emph{fr.} 90, \emph{De ins.};

Jordanes, \emph{ib.} 278 .

} Leo then considered it politic to cultivate their friendship and he allowed Theoderic to return to his people. The young prince at once distinguished himself in a campaign against the Sarmatians who had recently occupied Singidunum, and the Goths appropriated the city.

The last act of Theodemir seems to have been an invasion of the provinces of Dacia and Dardania, in which his army advanced as far as Naissus.\texthebrew{⁠}\footnote{Jordanes, \emph{Get.} 282\texthebrew{‑}286 , where events belonging to later incursions of Theoderic are mixed up with this invasion of Theodemir.

} Death befell him soon afterwards and Theoderic was elected as his successor in 471.\texthebrew{⁠}\footnote{Anon. Val. XVII.67 . Theoderic celebrated his tricennalia in A.D. 500.

} Soon after his accession (before 475) he seems to have led his people from their Pannonian homes to a new settlement in Lower Moesia, the same regions which had once been occupied by the Visigoths of Alaric.\texthebrew{⁠}\footnote{He seems to have resided in Novae.

Anon. Val. VI.42 .

} There is no evidence that this change of habitation was sanctioned by the Roman Emperor; but it does not seem to have been opposed at the time.

After the collapse of the Hunnic empire a large number of Ostrogoths had taken service in the Roman army, and formed the most important part of the German forces on whose support Aspar had maintained his power. We have already met their commander Theoderic (son of Triarius), called Strabo, ``squinter,'' who was not of very distinguished descent but was related through

p413
marriage to the family of Theodemir.\texthebrew{⁠}\footnote{John Ant. \emph{fr.} 98, Theoderic is said to be \textgreek{\texthebrew{ἀ}νεψιός} of Recitach son of Strabo. Schmidt (\emph{ib.} 127, n3) conjectures that Theodemir's sister had married Strabo's brother.

} We may call him Strabo to distinguish him from his more famous namesake. We saw the hostile attitude which he assumed towards Leo after the death of Aspar. The German troops gathered round him and proclaimed him king. He then sent an embassy to Leo, demanding for himself the post of Master of Soldiers in praesenti which Aspar had held, and the inheritance of Aspar, and for his troops grants of land in Thrace. The Emperor was willing to appoint him to the general ­ship, but refused the other demands. Then Strabo ravaged the territory of Philippopolis and reduced Arcadiopolis by starvation. These energetic proceedings extorted concessions from Leo; he agreed to pay a yearly stipend of 2000 lbs. in gold (= £90,000) to the Goths and to allot them a district in Thrace, and he conferred the post of Master of Soldiers in praesenti on Strabo, who was to fight for the Emperor against all enemies except the Vandals, and ``enemies'' doubtless included the Goths of Theoderic.\texthebrew{⁠}\footnote{Cp. Schmidt, \emph{op. cit.} I.136.

} He was, moreover, to be recognised as king of the Goths.\footnote{His wish to be recognised as king by the Emperor shows that he was not of royal descent. Dahn, \emph{Kön. der Germanen}, II.69.

}

In the troubles that followed Leo's death, Strabo naturally took the part of Basiliscus against his old foe, while Zeno was supported by Theoderic. After his restoration Zeno deprived Strabo of his military post and bestowed it on Theoderic, whom he also created a Patrician, confirming him in possession of the lands which his people had seized in Lower Moesia and promising him an annual stipend. He even adopted him as a son, according to the German right of adoption.

But there were no sincere feelings behind this favour and friendliness. The policy of the Emperor was to play off one Goth against the other. In the three following years ( A.D. 477\texthebrew{‑}479) the relations between him and the two rivals shifted rapidly through all the stages of possible combinations. In the first stage Zeno and Theoderic are combined against Strabo; in the second the two Theoderics join forces against Zeno; in the third Strabo and Zeno co-operate against Theoderic.

The drama began with an embassy from Strabo desiring reconciliation. The ambassadors reminded Zeno of the injuries

p414
which Theoderic had inflicted on the Empire, though he was called a Roman ``general'' and a friend. Zeno convoked the Senate, and it was concluded to be impossible to support the two generals and their armies, for the public resources were hardly sufficient to pay the Roman troops. The exchequer, it must not be forgotten, had not yet recovered from the failure of the Vandal expedition of the previous reign. As Strabo had always shown himself hostile at heart, was unpopular on account of his cruelty, and had assisted Basiliscus ``the tyrant,'' it was determined to reject his offer. Yet, as Zeno for a time withheld a reply, three friends of Strabo in Constantinople, Anthimus a physician, and two others, wrote him an account of the course which matters were taking; but the letters were discovered, the affair was examined by a senatorial commission of three persons, in the presence of the Master of Soldiers, and the three friends of the Goths were punished by flogging and exile.

Soon after this, probably in A.D. 478, the Emperor, perceiving that Strabo was becoming stronger and consolidating forces, and that Theoderic was hardly in a position to cope with him, deemed it wise to come to terms. He therefore sent an embassy proposing that the son of the chief should be sent to Byzantium as a hostage, and that Strabo himself should live as a private individual in Thrace, retaining what he had already secured by plunder, but binding himself to plunder no more. The chief refused, representing that it was impossible for him to withdraw now without paying the troops whom he had collected. Accordingly Zeno decided on war; troops were summoned from the dioceses of Pontus, Asia, and the East, and it was expected that Illus would assume the command. It seems, however, that Illus did not take the field, for we find Martinianus, his brother-in\texthebrew{‑}law , conducting a campaign against the Goths in the same year, and proving himself incompetent to maintain discipline in his own army. Then Zeno sent an embassy to Theoderic calling upon him to fulfil the duties of a Roman general and advance against the enemy. He replied that the Emperor and Senate must first swear that they will never make terms with the other Ostrogothic king. The senators took an oath that they would not do so unless the Emperor wished it, and the Emperor swore that he would not break the contract if it were not first violated by Theoderic himself.

p415
Theoderic then moved southwards. The Master of Soldiers of Thrace was to meet him with two thousand cavalry and ten thousand hoplites at a pass of Mount Haemus; when he had crossed into Thrace another force was to join him at Hadrianople, consisting of twenty thousand foot and six thousand horse; and, if necessary, Heraclea (on the Propontis) and the cities in the neighbourhood were prepared to send additional troops. But the Master of Soldiers was not at the pass of Mount Sondis, and the Goths when they advanced farther fell in with the army of Strabo, and the antagonists plundered one another's flocks and horses . Then Strabo, riding near his rival's camp, reviled him as a traitor to desert his own countrymen, and as a fool not to see through the plan of the Romans, who wished to rid themselves of the Goths, without trouble on their own part, by instigating them to mutual destruction, and were quite indifferent which party won. These arguments produced a power ­ful effect upon Theoderic's followers, and the two leaders made peace (478). This is the second stage of alliance, which we noted above. It was not to last long.

The reconciled Ostrogothic chieftains then sent ambassadors to Byzantium. Theoderic, upbraiding Zeno for having deceived him with false promises, demanded the concession of territory to his people, a supply of corn º to support his army till harvest time, and urged that if these demands were not satisfied, he would be unable to restrain his soldiers from plundering, in order to support themselves. Strabo demanded that the arrangements he had made with Leo (in A.D. 473) should be carried out, that the payment he had been accustomed to receive in former years should be continued, and that certain kinsmen of his, who had been committed to the care of Illus and the Isaurians, should be restored. We are not informed what answer Zeno made to the elder Theoderic, or whether he made any; to the son of Theodemir he replied, that if he consented to break with his namesake and make war upon him he would give him 2000 lbs. of gold and 10,000 lbs. of silver immediately, besides a yearly revenue of 10,000 nomismata, and the hand of a daughter of Placidia and Olybrius\texthebrew{⁠}\footnote{Probably Juliana, whom we afterwards find married to Areobindus.

} or of some other noble lady. But his promises did not avail, and Zeno prepared for war, notifying his intention to accompany the army in person. This intention created great

p416
enthusiasm among the soldiers, but at the last moment Zeno drew back, and they threatened a revolt, to prevent which the army was broken up and the regiments sent to their winter quarters.

When the army was disbanded, Zeno's only resort was to make peace on any terms with Strabo. In the meantime Theoderic, the son of Theodemir, was engaged in ravaging the fairest parts of Thrace in the neighbourhood of Mount Rhodope, which divides Thrace from Macedonia; he not only ruined the crops, but oppressed the farmers or slew them. Strabo, when he received Zeno's message, \texthebrew{—} remarking that he was sorry that the innocent husbandmen, for whose welfare Zeno\texthebrew{⁠}\footnote{``Zeno \emph{or Verina}'' (Malchus, \emph{fr.} 9, \emph{De leg. Rom.}). This seems to show that Verina had a preponderant influence at this time.

} did not care in the least, suffered from the ravages of his rival \texthebrew{—} concluded a peace on the conditions that Zeno was to supply a yearly payment sufficient to support thirteen thousand men; that he was to be appointed to the command of two scholae and to the post of Master of Soldiers in praesenti , and receive all the dignities which Basiliscus had bestowed upon him; that his kinsmen were to inhabit a city assigned by Zeno. The Emperor did not delay to execute this agreement; Theoderic was deposed from the office of Master of Soldiers, and Strabo appointed in his stead (before end of 478). This marks the third stage in these changeful relations.

Theoderic, now threatened by the superior forces of Strabo, was in a difficult position. But he managed to escape across Mount Rhodope into Macedonia (perhaps with the Emperor's collusion), and the town of Stobi felt the full brunt of his wrath. Thence he turned his steps toward Thessalonica, and the inhabitants felt so little confidence in Zeno that they actually believed that the Emperor wished to hand their city over to the barbarians. A sedition broke out which ended in the transference of the keys of the city from the Praetorian Prefect of Illyricum to the archbishop, a remarkable evidence of the fact that the people looked on the ministers of the Church as defenders against Imperial oppression. These suspicions of the Emperor's intentions were undoubtedly unjust. Zeno sent Artemidorus and Phocas to Theoderic, who was persuaded by their representations to stay his army and send an embassy to Byzantium. Theoderic

p417
demanded that a plenipotentiary envoy should be sent to treat with him. Zeno sent Adamantius, directing him to offer the Goths land in Pautalia (about Küstendil), and 200 lbs. of gold to supply food for that year, as no corn had been sown in the designated region. The motive of Zeno in choosing Pautalia was that if the Goths accepted it they would occupy a position between the Illyrian and Thracian armies, in which they might be more easily controlled.

Meanwhile Theoderic had proceeded by the Egnatian Way to Heraclea (Monastir), and had sent a message to one Sidimund,\texthebrew{⁠}\footnote{He was cousin of Aidoing, Count of the Domestics, and a friend of Verina; and he belonged to the Amal family.

} an Ostrogoth who had been in the service of Leo and had inherited an estate near Dyrrhachium, where he was living peaceably. Theoderic induced him to make an attempt to take possession of that important city of New Epirus, and for this purpose Sidismund employed an ingenious device. He visited the citizens individually, informing each that the Ostrogoths were coming with Zeno's consent to take possession of the city, and advising him to move his property with all haste to some other secure town or to one of the coast islands. The fact that his representations were listened to and that he effected the removal of a garrison of two thousand men proves that he possessed considerable influence. Theoderic was at Heraclea\texthebrew{⁠}\footnote{It is worth noti­cing that a sister of Theoderic, as well as his mother and brother, accompanied him on his march; she died at Heraclea and was buried there.

} when the messenger of Sidimund arrived with the news that the plan had been successfully carried out; and having burnt a large portion of the town because its inhabitants could not supply him with provisions, he set out for Epirus. He proceeded along the Egnatian Way, crossing the range of the Scardus mountains, and arrived at Lychnidus, which is now Ochrida. Built in a strong situation on the shore of Lake Ochrida, and well provided with water and victuals, Lychnidus defied the assault of the barbarians, who, unwilling to delay, hastened onwards, and having seized Scampae, the most important town between Lychnidus and Dyrrhachium, arrived at the goal of their journey.

It may be wondered whether at Dyrrhachium it entered the mind of Theoderic to ship his people across to the western peninsula and attack the Italian kingdom of Odovacar in the

p418
south. Adamantius, the ambassador who had been sent by Zeno to treat with him, seems to have thought it more likely that the Ostrogoths would employ vessels for the purpose of plundering the Epeirot or Dalmatian coasts, for he sent a post messenger to Dyrrhachium, to blame Theoderic for his hostile advance while negotiations were pending, and to exhort him to remain quiet and not to seize ships until he arrived himself.

Starting from Thessalonica, and passing Pella on the Via Egnatia, Adamantius came to Edessa, the modern Vodena, where he found Sabinian Magnus, and informed him that he had been appointed Master of Soldiers in Illyricum. The messenger, who had been sent to Dyrrhachium, returned in the company of a priest, to assure Adamantius that he might proceed confidently to the camp of Theoderic; and, having issued a mandate to collect all the troops available, the general and the ambassador moved forward to Lychnidus. Here Sabinian\texthebrew{⁠}\footnote{Sabinian was a strict disciplinarian, see Marcellinus, \emph{sub a.} 479: disciplinae praeterea militaris ita optimus institutor coercitorque fuit ut priscis Romanorum ductoribus comparetur.

} made difficulties about binding himself by oath to restore the hostages whom Theoderic was willing to deliver as a gage for the personal safety of Adamantius. This produced a deadlock; Theoderic naturally refused to give the hostages. Adamantius naturally refused to visit Theoderic.

Adamantius invented a simple solution of the difficulty, which led to a striking scene. Taking with him a body of two hundred soldiers he climbed by an obscure and narrow path, where horses had never set hoof before, and reached by a circuitous route an impregnable fort, built on a high cliff, close to the city of Dyrrhachium. At the foot of the cliff yawned a deep ravine, through which a river flowed. A messenger was sent to inform Theoderic that the Roman ambassador awaited him, and, attended by a few horse-soldiers , the Ostrogoth rode to the bank of the river. The physical features, the cliff, the chasm, and the river, are sufficiently simple and definite to enable us to call up vividly this strange scene. The attendants of both Adamantius and Theoderic had retired beyond range of earshot; and standing on the edges of the ravine the Ostrogothic king and the ambassador of the Roman Empire conversed together.

p419
``I elected to live,'' complained Theoderic, ``beyond the borders of Thrace, far away Scythia-ward , deeming that if I abode there I should trouble no man, and should be able to obey all the behests of the Emperor. But ye summoned me as to war against Theoderic, and promised, firstly, that the Master of Soldiers in Thrace would meet me with his army, yet he never appeared; secondly, ye promised that Claudius, the steward of the Gothic contingent, would come with the pay for my troops ( \textgreek{ξενικ\texthebrew{ῷ}} ), yet I never saw him; thirdly, ye gave me guides who, leaving the better roads that would have taken me to the quarters of the foe, led me by steep and precipitous rocky paths, where I wellnigh perished with all my train, advancing as I was with cavalry, waggons, and all the furniture of camp, and exposed to the attacks of the enemy. I was therefore constrained to come to terms with them, and owe them a debt of gratitude that they did not annihilate me, betrayed as I was by you and in their power.''

``The Emperor,'' replied Adamantius, ``bestowed upon you the title of Patrician, and created you a Master of Soldiers. These are the highest honours that crown the labours of the most deserving Roman officers, and nothing should induce you to cherish towards their bestower other than filial sentiments.'' Having endeavoured to defend or extenuate the treatment of which Theoderic complained, the envoy proceeded thus: ``You are acting intolerably in seizing Roman cities, while you are expecting an embassy; and remember that the Romans held you at their mercy, a prisoner, surrounded by their armies, amid the mountains and rivers of Thrace, whence you could never have extricated yourself, if they had not permitted you to withdraw, not even were your forces tenfold as great as they are. Allow me to counsel you to assume a more moderate attitude towards the Emperor, for you cannot in the end overcome the Romans when they press on you from all sides. Leave Epirus and the cities of this region \texthebrew{—} we cannot allow such great cities to be occupied by you and their inhabitants to be expelled \texthebrew{—} and go to Dardania, where there is an extensive territory of rich soil, uninhabited, and sufficient to support your host in plenty.''

To this proposal Theoderic replied that he would readily consent, but that his followers, who had recently endured many hardships, would be unwilling to leave their quarters in Epirus,

p420
where they had fully expected to pass the winter. He proposed a compromise, and engaged that if he were permitted to winter at Dyrrhachium he would migrate to Dardania in the ensuing spring. He added that he was quite ready to leave the unwarlike mass of his Ostrogoths in any city named by Zeno, and giving up his mother and sister as hostages, to take the field against Strabo with six thousand of his most martial followers, in company with the Illyrian army; when he had conquered his rival he expected to succeed to the post of Master of Soldiers and to be received in New Rome as a Roman.\texthebrew{⁠}\footnote{\textgreek{Τ\texthebrew{ὸ}ν \texthebrew{Ῥ}ωμαικ\texthebrew{ὸ}ν πολιτεύσοντα τρόπον}. For Julius Nepos see above,

p404 .

} He also observed that he was prepared, if the Emperor wished, ``to go to Dalmatia and restore Julius Nepos.'' Adamantius was unable to promise so much; it was necessary to send a messenger to Byzantium to consult the Emperor. And thus the interview terminated.

Meanwhile the military forces, stationed in the Illyrian cities, had assembled at Lychnidus, around the standard of Sabinian. It was announced to the general that a band of the Ostrogoths led by Theodimund, the brother of Theoderic, was descending in secure negligence from Mount Candaira, which separates the valley of the Genusus (Skumbi) from that of the Drilo. This band had formed the rear of the Ostrogothic line of march, and had not yet reached Dyrrhachium. Sabinian sent a few infantry soldiers by a circuitous mountain route, with minute directions as to the hour and place at which they were to appear; and himself with the rest of the army proceeded thither, after the evening meal, by a more direct way. Marching during the night he assailed the company of Theodimund at dawn of day. Theodimund and his mother, who was with him, fled with all speed into the plain, and, having crossed a deep gully, destroyed the bridge which spanned it to cut off pursuit. This act, while it saved them, sacrificed their followers, who turned at bay upon the Romans. Two thousand waggons and more than five thousand captives were taken, and a great booty ( A.D. 479).

After this the Emperor received two messages, one from Adamantius announcing the proposals of Theoderic, the other from Sabinian exaggerating his victory and dissuading him from the conclusion of peace. War seemed more honourable to Zeno and the pacific offers were rejected, Sabinian was permitted to continue the war, and for about a year and a half he held the

p421
Gauls in check in Epirus. But the active general was murdered by an ungrateful master,\texthebrew{⁠}\footnote{For the fate of Sabinian see John Ant. \emph{fr.} 97; for the date, 481, Marcellinus, \emph{sub a.}

} and John the Scythian and Moschian were sent to succeed him.

The revolt of Marcian towards the end of A.D. 479 had given Strabo a pretext for approaching Constantinople to assist the government. Having extorted money from Zeno, he received two of the conspirators in his camp and refused to surrender them. He was then once more deprived of his dignities and declared an enemy of the republic. He entered again into alliance with Theoderic and devastated Thrace. Zeno invoked the aid of the Bulgarians of the Lower Danube, but they were defeated by Strabo, who then advanced on Constantinople ( A.D. 481).

It was a surprise, and we are told that he would easily have captured the city if Illus had not set guards at the gates just in time. He attempted to cross over to Bithynia, but was defeated in a battle on the water, and departed to Thrace. Thence he set forth for Greece, with his son Recitach, his wife, and about 30,000 followers. At a place called the Stable of Diomede, on the Egnatian Road, his horse threw him one morning on a spear which was standing point upwards, close to his tent. The accident was fatal ( A.D. 481). Recitach succeeded him, and ruled in Thrace, ``performing more outrageous acts than his father had performed.''\texthebrew{⁠}\footnote{John Ant. \emph{fr.} 95. Another account will be found in Eustathius, \emph{fr.} 3

(\emph{apud Evagrium}, III.25) .

} Three years later Recitach was slain by Theoderic, son of Theodemir, whom Zeno instigated to the deed.\footnote{Recitach had murdered his uncles, so that the act of Theoderic (who was related to Strabo) was an act of blood-vengeance. John Ant. \emph{fr.} 98.

}

In 482 we find Theoderic \texthebrew{—} the name is no longer ambiguous \texthebrew{—} ravaging the provinces of Macedonia, and Thessaly, and capturing the town of Larissa. He was no longer held in check by the able general Sabinian who had been murdered the year before. The Emperor decided to make a new agreement. Parts of Moesia and Dacia Ripensis were conceded to the Ostrogoths, and Theoderic was appointed Master of Soldiers ( A.D. 483).\texthebrew{⁠}\footnote{Marcellinus, \emph{Chron., sub a.}

} In A.D. 484 he enjoyed the coveted distinction of giving his name to the year as consul, and he assisted Zeno against the rebel Illus. But a new breach soon followed. He devastated Thrace ( A.D. 486) and marched on Constantinople ( A.D. 487). Rhegium was occupied, Melantias was taken, and the capital once more

p422
threatened. But the intervention of his sister,\texthebrew{⁠}\footnote{Perhaps Amalafrida (Schmidt, \emph{op. cit.} I.147, n4).

} who was at Zeno's court, induced him to retire to his headquarters in Moesia, which he was soon to abandon for ever. The days of the Thracian period of Theoderic's career were numbered.

We have seen that there had been friction between the Emperor and his Viceroy in Italy, and that Odovacar had thoroughly defeated the Rugians whom Zeno had stirred up against him. The thought now occurred to Zeno or his advisers that he might at once punish Odovacar and deliver the Illyrian provinces from the mena­cing presence of the Ostrogoths by giving Theoderic a commission to supersede the ruler of Italy. Theoderic accepted the charge. A compact was made that (in the words of the chronicler) ``in case Odovacar were conquered, Theoderic should, as a reward of his labours, rule in place of Odovacar, until Zeno came himself.''\texthebrew{⁠}\footnote{Anon. Val. II.49 . We may conjecture that Theoderic, who had been mag. mil. in praes. in the East since 483, was now appointed mag. mil. in praes. in Italy, to replace Odovacar.

} The last condition is simply a way of saying that Zeno reserved all the Imperial rights of sovranty.

At the head of his people, numbering perhaps about 100,000,\texthebrew{⁠}\footnote{Schmidt, \emph{op. cit.} 1.152.

} Theoderic set forth from Moesia in the autumn of A.D. 488. Following the direct road to Italy, past Viminacium and Singidunum, he approached Sirmium, and here he was confronted by a formidable obstacle. This town was in the possession of the Gepids, who now blocked Theoderic's path. The place was taken after fierce fighting, but the Goths passed on with their booty and the Gepids reoccupied it. The winter, spring, and summer of the following year were spent somewhere between Sirmium and the Italian borders, and the causes of this delay are unknown.

It was not till the end of August ( A.D. 489) that, having crossed the Julian Alps, the Ostrogoths reached the river Sontius (Isonzo) and the struggle for Italy began. Of this memorable war we have only the most meagre outline. The result was decided within twelve months, but three and a half years were to elapse

p423
before the last resistance of Odovacar was broken down and Theoderic was completely master of Italy.\footnote{The chief sources are Ennodius (\emph{Panegyricus} and \emph{Vita Epiphanii}) and

the chronicle known as Anonymus Valesianus, Part 2 . The most recent editor, Cessi, has shown (correctly, I think) that it falls into two sections of different author­ship (1 = § 36-§ 77 ; 2 = § 78-§ 96 ). They are contrasted by the fact that the first is highly favourable to Theoderic, and the second undisguisedly censorious. The first was written before his death, the second probably between 527 and 534 (Cessi, CLXVI \emph{sqq.}). The conjecture of some that Maximian, archbishop of Ravenna, was the author, will not hold.

}

It was perhaps where the Sontius and the Frigidus meet that Theoderic found Odovacar in a carefully fortified camp, prepared to oppose his entry into Venetia. He had considerable forces, for besides his own army he had succeeded in enlisting foreign help.\texthebrew{⁠}\footnote{Ennodius, \emph{Pan.} p271, says rhetorically universas nationes, and tot reges quot sustinere generalitas milites vix valeret.

} We are not told who his allies were; we can only guess that among them may have been the Burgundians, who, as we know, helped him at a later stage. The battle was fought on August 28; Odovacar was defeated and compelled to retreat. His next line of defence was on the Athesis (Adige), and he fortified himself in a camp close to Verona, with the river behind him.\texthebrew{⁠}\footnote{Ennodius (\emph{ib.} 272) suggests that Odovacar chose the position to render flight impossible for his army.

} Here the second battle of the war was fought a month later (about Sept. 29)\texthebrew{⁠}\footnote{Sept. 29 or 30 (Hodgkin) seems implied by

Anon. Val. 50 .

} and resulted in a decisive victory for Theoderic. The carnage of Odovacar's men is said to have been immense; but they fought desperately and the Ostrogothic losses were severe;\texthebrew{⁠}\footnote{Caedis enormitas, Ennod. p273; ceciderunt populi ab utraque parte,

Anon. Val. \emph{ib.}

} the river was fed with corpses. The king himself fled to Ravenna. The greater part of the army, with Tufa who held the highest command, surrendered to Theoderic, who immediately proceeded to Milan.\footnote{Anon. Val. 51 , where it is said Tufa was appointed mag. mil. by Odovacar and his chief men. If so, Odovacar had usurped a right which belonged to the Emperor.

}

Northern Italy was now at the feet of the Goth; Rome and Sicily were prepared to submit, and it looked as if nothing remained to complete the conquest but the capture of Ravenna. But the treachery of Tufa changed the situation. Theoderic imprudently trusted him, and sent him with his own troops and a few distinguished Ostrogoths against Odovacar. At Faventia (Faenza) he espoused again the cause of his old master and handed over to him the Goths, who were put in irons.

p424
Theoderic made Ticinum (Pavia) his headquarters during the winter, and it is said that one of his motives for choosing this city was to cultivate the friendship of the old bishop Epiphanius, who had great influence with Odovacar. In the following year Odovacar was able to take the field again, to seize Cremona and Milan, and to blockade his adversary in Ticinum. At this juncture the Visigoths came to the help of the Ostrogoths and sent an army into Italy. The siege was raised and the decisive battle of the war was fought on the river Addua (Adda), in which Odovacar was utterly defeated (Aug. 11, A.D. 490). He fled for the second time to Ravenna. It was probably this victory that decided the Roman Senate to abandon the cause of Odovacar, and accept Theoderic. It made him master of Rome, southern Italy, and Sicily.

The agreement that Zeno made with Theoderic had been secret and unofficial. The Emperor did nothing directly to break off his relations with Odovacar.\texthebrew{⁠}\footnote{This is shown by the fact that the western consul of 490, Flavius Probus Faustus, assuredly nominated by Odovacar, was acknowledged in the East. Sundwall (\emph{Abh.} 187 \emph{sqq.}) is right, I think, in his treatment of the political situation in these years.

} But Odovacar seems some time before the battle of the Addua to have courted a formal rupture. He created his son Thela a Caesar, and this was equivalent to denouncing his subordination to the Emperor and declaring Italy independent.\texthebrew{⁠}\footnote{Sundwall would place the elevation of Thela at the beginning of 490. The fact is recorded by John Ant. \emph{fr.} 99, \emph{De ins.}

} He probably calculated that in the strained relations which then existed between the Italian Catholics and the East, on account of the ecclesiastical schism, the policy of cutting the rope which bound Italy to Constantinople would be welcomed at Rome and throughout the provinces. The senators may have been divided on this issue, but the battle of the Addua decided them as a body to ``betray'' Odovacar,\texthebrew{⁠}\footnote{John Mal. XV p383 \textgreek{πολεμήσας} (Th.) \textgreek{α\texthebrew{ὐ}τ\texthebrew{ῷ}} (Od.) \textgreek{κατ\texthebrew{ὰ} γνώμην προδοσίαν τ\texthebrew{ῆ}ς συγκλήτου \texthebrew{Ῥ}ώμης.}} and before the end of the year Festus, the princeps of the Senate, went to Constantinople to announce the success of Theoderic, and to arrange the conditions of the new Italian government.

Theoderic confidently believed that his task was now virtually finished. But the cause of his thrice-defeated enemy was not yet hopelessly lost. Tufa was still at large with troops at his command; and other unexpected difficulties beset the conqueror. The Burgundian king Gundobad sent an army into

p425
North Italy and laid waste the country.\texthebrew{⁠}\footnote{This episode is very obscure. The sources are Ennodius, \emph{Pan.} p276, \emph{Vit. Epiph.} 369 \emph{sqq.}; \emph{Hist. Misc.} XV.16 (cp. Cassiodorus, \emph{Var.} 12, 28 ). Ennodius gives no clear chronological indications.

Hodgkin

places the event in 490, before the battle of the Addua; but the circumstances seem to point to a later date, for Theoderic was apparently besieged in Ticinum till the arrival of the Visigoths and could not have dealt with the Burgundians. Schmidt's chronology is preferable (\emph{op. cit.} I.156). The motive for Gundobad's interference is intelligible: he may well have feared the enclosure of his kingdom between a Visigothic power on one side and an Ostrogothic on the other.

} Theoderic had not only to drive the invaders out, but he had also to protect Sicily against the Vandals, who seized the opportunity of the war to attempt to recover it. Their attempt was frustrated and they were forced to surrender the fortress of Lilybaeum as well as all their claims to the island.\footnote{Cassiodorus, \emph{Chron., sub a.} 491. Theoderic had also trouble with the Rugians who had joined his expedition. Having plundered Ticinum they went over to Tufa, but then quarrelled with him and returned to Theoderic. Cp. Ennod. \emph{Pan.} \emph{ib.}, \emph{Vit. Epiph.} 361 \emph{sqq.}

}

It seems to have been in the same year that Theoderic resorted to a terrible measure for destroying the military garrisons which held Italian towns for Odovacar. The Italian population was generally favourable to the cause of Theoderic, and secret orders were given to the citizens to slaughter the soldiers on a pre-arranged day. The pious panegyrist, who exultantly, but briefly, describes this measure and claims Providence as an accomplice, designates it as a sacrificial massacre'';\texthebrew{⁠}\footnote{Nex votiva, Ennodius, \emph{Pan.} p275. This atrocious act is not mentioned by Anon. Val. It is discussed by Dahn, \emph{Kön. der Germ.} II.80; Hodgkin, III.226 .

} and Theoderic doubtless considered that the treachery of his enemy's army in surrendering and then deserting justified an unusual act of vengeance. The secret of the plot was well kept, and it seems to have been punctually executed. The result was equivalent to another victory in the field; and nothing now remained for Theoderic but to capture the last stronghold of his adversary, the marsh city of Honorius.

The siege of Ravenna lasted for two years and a half. The Gothic forces entrenched themselves in a camp in the Pine-woods east of the city, but were not able entirely to prevent provisions from reaching the city by sea. Yet the blockade was not ineffective, for corn rose to a famine price. One attempt was made by Odovacar to disperse the besiegers. He made a sortie at night (July 10, A.D. 491) with a band of Herul warriors and

p426
attacked the Gothic trenches. The conflict was obstinate, but he was defeated.\texthebrew{⁠}\footnote{\emph{Consularia Italica}, p318 .

} Another year wore on, and it appeared that the siege might last for ever unless the food of the garrison could be completely cut off. Theoderic managed to procure a fleet of warships \texthebrew{—} we are not told whether they were built for the occasion, \texthebrew{—} and, making the Portus Leonis, \texthebrew{•} about six miles from Ravenna, his naval base, he was able to blockade the two harbours of the city (August, A.D. 492).\texthebrew{⁠}\footnote{These harbours are now dry land, and are marked, one by the Church of S. Apollinare in Classe, the other by that of S. Maria in Porto fuori.

} Odovacar held out for six months longer, but early in A.D. 493 negotiations, conducted by the bishop of Ravenna, issued in a compact between the two antagonists (February 25) that they should rule Italy jointly.\texthebrew{⁠}\footnote{Procopius, \emph{B. G.} I.1; John Ant. \emph{fr.} 99 (\emph{De ins.} p140).

} Theoderic entered the city a week later (March 5).

The only way in which the compact could have been carried out would have been by a territorial division. But Theoderic had no mind to share the peninsula with another king, and there can hardly be a doubt that, when he swore to the treaty, he had the full intention of breaking his oath. Odovacar's days were numbered. Theoderic, a few days after his entry into Ravenna, slew him with his own hand in the palace of Lauretum (March 15). He alleged that his defeated rival was plotting against him, but this probably was a mere pretext.\texthebrew{⁠}\footnote{Anon. Val. 54

dum ei Odoachar insidiaretur. In the other sources which depend on the Ravennate Annals (Anon. Cuspin., Cont. Prosperi Havn. and Agnellus) there is no mention of a plot, nor in John Ant.; but see Cassiodorus, \emph{Chron.} (Odoacrem molientem sibi insidias), and Procopius \emph{B. G.} I.1 \textgreek{\texthebrew{ἐ}πιβουλ\texthebrew{ῇ ἐ}ς α\texthebrew{ὐ}τ\texthebrew{ὸ}ν χρώμενον}. The plot was evidently part of the official Ostrogothic account.

} ``On the same day,'' adds the chronicler, ``all Odovacar's soldiers were slain wherever they could be found, and all his kin.''\footnote{``On the same day'' is not quite accurate. See John Ant. \emph{ib.}, who records that Odovacar's son, ``whom he had proclaimed Caesar,'' was exiled to Gaul, but returning to Italy was put to death. Sunigilda, Odovacar's wife, was starved to death. It is true that his brother was slain on the same day. The name of the son was Thela (Anon. Val.), and \textgreek{\texthebrew{Ὀ}κλάν} in John Ant. is probably an error for \textgreek{Θήλαν} (as Mommsen conjectured). The statement that all Odovacar's soldiers were killed is doubtless an exaggeration.

}

In three years and a half Theoderic had accomplished his task. The reduction of Italy cost him four battles, a massacre, and a long siege. His capital blunder had been to trust Tufa

p427
after the victory of Verona. We may be sure that throughout the struggle he spared no pains to ingratiate himself in the confidence of the Italian population. But when his rival had fallen, and when he was at last securely established, Theoderic's first measure was to issue an edict depriving of their civil rights all those Italians who had not adhered to his cause. This harsh and stupid policy, however, was not carried out, for the bishop Epiphanius persuaded the king to revoke it and to promise that there would be no executions.\footnote{Ennodius, \emph{ib.} 362 \emph{sqq.}

}

Two more services would be rendered to his country by Epiphanius before his death. The war had a disastrous effect on Italian agriculture.\texthebrew{⁠}\footnote{\emph{Ib.} 366, uides universa Italiae loca originariis uiduata cultoribus.

} Liguria had been devastated by the Burgundians; King Gundobad had carried thousands into captivity, and no husbandmen were left to till the soil and tend the vineyards. Theoderic was prepared to ransom the captives, and he charged Epiphanius with the office of persuading the Burgundian king to release them. The bishop, notwithstanding his infirm age, undertook the cold and difficult journey over the Alps in March ( A.D. 494), and was received by Gundobad at Lyons. To the arguments and prayers of the envoy, Gundobad, who was an excellent speaker, replied with the frank and cynical assertion that war permits and justifies everything which is unlawful in peace. ``War ignores the bridle of moderation which you, as a Christian luminary, teach. It is a fixed principle with belligerents that whatever is not lawful is lawful when they are fighting. The object of war is to cut up your opponent's strength at the roots.'' He went on to say that a peace had now been concluded \texthebrew{—} it had been sealed by the betrothal of a daughter of Theoderic to Gundobad's son Sigismund, \texthebrew{—} and that if the bishop and his companions would return to their homes he would consider what it were best to do in the interests of his soul and his kingdom. Epiphanius had gained his cause. Gundobad set free all prisoners who were in his own hands, without charge, and those who were the slaves of private persons were ransomed. More than six thousand were restored to Italy.\footnote{\emph{Ib.} 370 \emph{sqq.} Ennodius accompanied Epiphanius on this embassy.

}

The last public act of Epiphanius was to induce Theoderic

p428
to grant a reduction of the taxation of Liguria. ``The wealth,'' he urged, ``of a landed proprietor is the wealth of a good ruler.''\texthebrew{⁠}\footnote{Boni imperatoris est possessoris opulentia, \emph{ib.} 379.

} Theoderic remitted two-thirds of the taxes for A.D. 497. Epiphanius caught a chill in the cold marsh air of Ravenna and died on his return home.\texthebrew{⁠}\footnote{A.D. 497, at the age of 58.

\texthebrew{◂} previous

next \texthebrew{▸}Images with borders lead to more information.

The thicker the border, the more information.

(Details here.)

UP TO:

J. B. Bury
Homepage

Latin \& Greek

Texts

LacusCurtius

Home} He had played a considerable and beneficent part in Italian politics for nearly thirty years.

\xchapter{Chapter XIII, Part A}

Short URL for this page:

bit.ly/BURLAT13A

\xsubsection{(Part 1 of 2)}

On the evening of the day after Zeno's death, the Senate, the ministers, and Euphemius the Patriarch assembled in the palace, and a crowd of citizens and soldiers gathered in the Hippodrome (December 10, 491).\texthebrew{⁠}\footnote{The following description is taken from the contemporary document preserved in Constantine Porph. \emph{De cer.} I.92. Cp. above,

p316, n2 .

} Ariadne,\texthebrew{⁠}\footnote{Ariadne is represented on five diptychs belonging to the later part of the reign: namely, those of (1) Clementinus, cons. 513, at Liverpool; (2) Anthemius, cons. 515, one leaf, at Limoges; (3) Anastasius, cons. 517, in Bibl. nationale; (4) same, one leaf, at Verona; (5) same, one leaf, at Berlin; the other at South Kensington. With the help of these, by comparing the character of the head-dress, Delbrück has identified three female marble heads found in Italy as Ariadne's: (1) the head in the Lateran Museum, vulgarly known as St. Helena, on a bust which does not belong to it; (2) a head in the Palazzo dei Conservatori at Rome, found in 1887 near S. Maria dei Monti; (3) a head found at Rome but now in the Louvre. He considers them as probably Byzantine work. See his \emph{Porträts byz. Kais.}

} wearing the Imperial cloak, and accompanied by the Grand Chamberlain Urbicius, the Master of Offices, the Castrensis, the Quaestor, and others, but not by the Patriarch, then entered the Kathisma of the Hippodrome to address the people. She was warmly acclaimed. ``Long live the Augusta! Give the world an orthodox Emperor.'' Her speech was delivered by the Magister a libellis, who stood on the steps in front of the Kathisma. ``Anticipating your request, we have commanded the illustrious ministers, the sacred Senate, with the approval of the brave armies, to select a Christian and Roman Emperor, endowed with every royal virtue, who is not the slave of money, and who is, so far as a man may be, free

p430
from every human vice.'' \emph{People:} ``Ariadne Augusta, thou conquerest! O heavenly king, give the world a Basileus who is not avaricious!'' \emph{Ariadne:} ``In order that the choice may be pure and pleasing to God, we have commanded the ministers and the Senate, the vote of the army concurring, to make the election, in the presence of the Gospels, and in the presence of the Patriarch, so that no one may be influenced by friendship or enmity, or kinship, or any other private motive, but may vote with his conscience clear. Therefore, as the matter is weighty and concerns the welfare of the world, you must acquiesce in a short delay, till the obsequies of Zeno, of pious memory, have been duly performed, so that the election may not be precipitate.'' \emph{People:} ``Long live the Augusta! Cast out the thieving Prefect of the City! May all be well in thy time, Augusta, if no foreigner is imposed on the Romans!''\texthebrew{⁠}\footnote{\textgreek{Ε\texthebrew{ἰ} ο\texthebrew{ὐ}δ\texthebrew{ὲ}ν ξένον α\texthebrew{ὔ}ξει τ\texthebrew{ὸ} γένος τ\texthebrew{ῶ}ν \texthebrew{Ῥ}ωμαίων}. Probably the unpopular Prefect of the City was an Isaurian.

} \emph{Ariadne:} ``We have already anticipated your wishes. Before we came in, we appointed the illustrious Julian to the office of Prefect.'' \emph{People:} ``A good appointment! Long live the Augusta.'' After a few more words, Ariadne withdrew to the palace,\texthebrew{⁠}\footnote{\textgreek{Ε\texthebrew{ἰ}ς τ\texthebrew{ὸ}ν Α\texthebrew{ὐ}γουστέα} (so read for \textgreek{α\texthebrew{ὐ}γουστα\texthebrew{ῖ}ον}), \emph{De cer.} p421, l. 7.\texthebrew{—} The delphax seems to have been in the palace, but adjoining the Hippodrome. Ebersolt (\emph{Grand Palais}, p66) thinks it was an isolated building. See above,

p395 .

} and the ministers held a council in front of the Delphax to consult about the election. Urbicius proposed that the choice should be left to Ariadne, and the Patriarch, who was present, was sent to summon her. She chose Anastasius, a silentiary, and the Master of Offices sent the Counts of the Domestics and Protectors to fetch Anastasius from his house. He was kept that night in the Consistorium; notices were issued for a silentium\texthebrew{⁠}\footnote{\textgreek{Σιλέντιον κα\texthebrew{ὶ} κοβέντον} (= conventus), see above,

p24, n2 .

} to be held on the morrow; and the funeral of Zeno was performed.

Anastasius was a remarkable and well-known figure in Constantinople. He held unorthodox opinions, partly due, perhaps, to an Arian mother and a Manichaean uncle,\texthebrew{⁠}\footnote{Theodore Lector, II.7; Theoph. A.M. 5983.

} and he was possessed by religious enthusiasm, which led him to attempt to convert others to his own opinions. He did this in a curiously public manner. Having placed a chair in the church of St. Sophia, he used to attend the services with unfailing regularity

p431
and give private heterodox instruction to a select audience from his cathedra. By this conduct he offended the Patriarch Euphemius, who by Zeno's permission expelled him from the church and removed his chair of instruction;\texthebrew{⁠}\footnote{See Theophanes, A.M. 5982.

} but he was well thought of by the general public on account of his piety and liberality. It even appears that he may have at one time dreamt of an ecclesiastical career, for he was proposed for the vacant see of Antioch.\texthebrew{⁠}\footnote{In 488, when Palladius was elected. Compare A. Rose, \emph{Kaiser Anastasius I} (p13), who translates \textgreek{συνεψηφίσθη} rightly.

} The Patriarch was highly displeased at the Empress's choice of Anastasius, whom he stigmatised as unworthy to reign over Christians. His objections were overruled by the Senate and the Empress, but before he consented to take part in the coronation ceremony he insisted that the new Emperor should be required to sign a written declaration of orthodoxy. This was agreed to.

The officials dressed in white gathered in the Consistorium\texthebrew{⁠}\footnote{Not, it is expressly noticed, in the Arma (p422). \textgreek{Α\texthebrew{ἱ} πύλαι το\texthebrew{ῦ ἄ}ρματος} are mentioned in a seventh-century document, Const. \emph{De cer.} p628. Ebersolt (\emph{Le Grand Palais}, 63) thinks that the Arma was a dépôt of arms, near the Tribunal of the 19 Akkubita.

} on the following day (April 11), and were received ceremonially by Anastasius. The Patriarch was present, and now, if not before, he must have obtained the Emperor's signature to the declaration, which was lodged in the archives of St. Sophia under the care of the treasurer. Anastasius then left the Consistorium and ascended the steps of the portico\texthebrew{⁠}\footnote{The space in front of the Portico was the Tribunal of the 19 Akkubita. For details see Ebersolt, 62.

} of the triklinos of the Nineteen Akkubita. Here at the request of the senators he took a public oath that he would distress no person against whom he had a grudge, and that he would govern conscientiously. Then he proceeded to the triklinos of the Hippodrome, put on the Imperial tunic, girdle, leggings, and red boots,\texthebrew{⁠}\footnote{\textgreek{Στιχάριν διβητήσιν α\texthebrew{ὐ}ρόκλαβον, ζωνάριν, τουβία, καμπάγια βασιλικά} (p423).

} and entered the Kathisma, in front of which stood the troops, the standards lying on the ground. When he had been raised on a shield, and the torc placed on his head, the standards were raised, and he was acclaimed. Then he returned to the triklinos, when the Patriarch covered him with the Imperial cloak and crowned him. Reappearing in the Kathisma, he addressed the people, promising a donation of 5 nomismata and a pound of silver to

p432
each soldier \texthebrew{—} the same amount which had been given by Leo I. Among the enthusiastic acclamations with which he was greeted we may notice, ``Reign as thou hast lived! Thou hast lived piously! reign piously! Restore the army! Reign like Marcian!'' and ``Cast out the informers!''

A few weeks later Anastasius married Ariadne (May 20). His accession was undoubtedly a welcome change to Byzantium. He was a man of tall stature and remarkable for his fine eyes, which differed in hue.\texthebrew{⁠}\footnote{Hence called Dikoros. John Mal. XVI p392 describes his personal appearance.

} He is described as intelligent, well-educated , gentle, and yet energetic, able to command his temper, and generous in bestowing gifts.\texthebrew{⁠}\footnote{John Lydus, \emph{De mag.} I.47 . Zacharias Myt., well disposed to him as a Monophysite, says (VII.1) ``he was power­ful in aspect, vigorous in mind, and a believer.''

} A bishop of Rome wrote to him, ``I know that in private life you always strove after piety.''\footnote{Gelasius, in Mansi, XIII.30.

}

The first task imposed upon the new Emperor was to put an end to the unpopular predominance of the Isaurians, which had lasted for over twenty years. The choice of Anastasius had disappointed and alarmed the Isaurians, who had looked forward to the succession of Longinus. A riot in the Hippodrome soon gave Anastasius a pretext for driving them out of the city. During a spectacle at which the Emperor was present, the people clamoured against Julian, the Prefect of the City, who had done something which public opinion disapproved. Anastasius ordered his guards to intimidate the rioters, who then set fire to the Hippodrome, and pulled down and insulted the bronze statues of the Emperors. Not a few were slain in the tumult.\texthebrew{⁠}\footnote{John Ant. \emph{fr.} 100 (\emph{De ins.} p141). For date, Marcellinus, \emph{sub} 491.

} The Emperor found it politic to replace Julian by his own brother-in\texthebrew{‑}law Secundinus, but he attributed the disturbance to the machinations of the Isaurians. He expelled them all from the city. He forced Zeno's brother Longinus to take orders and banished him to the Thebaid. He confiscated Zeno's property, even selling his Imperial robes. He naturally withdrew the large allowances which Zeno had made to his fellow-countrymen , amounting to 1400 lbs. of gold.\texthebrew{⁠}\footnote{John Ant. \emph{ib.} p142. Evagrius

(III.35)

says 5000 pounds. His account of the war (from Eustathius?) is very inaccurate.

} A revolt had already broken out in Isauria,\texthebrew{⁠}\footnote{John Ant. \emph{ib.} p141 \textgreek{\texthebrew{ἤ}δη \texthebrew{ἀ}γγελθείσης τ\texthebrew{ῆ}ς κατ\texthebrew{ὰ} τ\texthebrew{ὴ}ν χώραν α\texthebrew{ὐ}τ\texthebrew{ῶ}ν \texthebrew{ἀ}ποστάσεως.}} and the rebels were now reinforced by the exiles

p433
from Constantinople, among them Longinus of Kardala.\texthebrew{⁠}\footnote{The ex-Master of Offices.

} Their total force is said to have numbered 100,000, and included Romans as well as Isaurians. The leaders in command were Linginines and Athenodorus.\texthebrew{⁠}\footnote{A man of wealth. (In Evagrius he is called Theodorus.) There was also a second leader of the same name. Linginenes (John Ant.) = \textgreek{Λονγινίνης \texthebrew{ὁ} χωλός} (John Mal. p393) = Libingis (Marcellinus). He was the comes Isauriae. Other prominent leaders were Conon, the fighting bishop, and Longinus of Selinus. The number of their forces is probably much exaggerated.

} They were met at Cotyaeum in Phrygia by an Imperial army under John the Scythian and John the Hunchback,\texthebrew{⁠}\footnote{John the Hunchback (\textgreek{\texthebrew{ὁ} κυρτός}) was Master of Soldiers in praesenti (John Mal. p393) and we may suppose that John the Scythian was still Master of Soldiers in the East. (Otherwise Theophanes A.M. 5985). Another general was the patrician Diogenianus, kinsman of the Empress (John Mal. \emph{ib.}). Justin (afterwards Emperor) took part in this battle. The number of the army given by John Ant. (2000) is corrupt. There were both Hunnic and Gothic auxiliaries.

} and were completely defeated, Linginines being slain. This battle shattered the power of the Isaurians irretrievably. But the defeated leaders did not submit and, just as in the case of the struggle between Illus and Zeno, warfare was carried on in the Isaurian mountains for several years before all the rebels were captured and killed.\texthebrew{⁠}\footnote{The chronology has been elucidated by Brooks, \emph{op. cit.} 235 \emph{sqq.} :\texthebrew{—}A.D. 493. Claudiopolis besieged by Diogenianus; his army blockaded by the Isaurians, and relieved by John the Hunchback; bishop Conon slain.

494\texthebrew{‑}497. Isaurians hold out in their fortresses, and are furnished with provisions by Longinus of Selinus, from the seaport of Antioch (not far from Selinus).

497. Longinus of Kardala and Athenodorus captured by John the Scythian, and their heads exposed on poles at Constantinople (cp. Evagrius, \emph{ib.} and Marcellinus, \emph{sub} 497).

498. Longinus of Selinus and two others who were holding out at Antioch captured. (Evagrius, \emph{ib.} and Marcellinus, \emph{sub} 498).

The year 497 was reckoned as the last of the war (cp. Marcellinus, and Theodore Lector, II.9).

The two Johns who conduct the war were rewarded by the consul­ship (498 and 499).

} It was not till A.D. 498 that the last of them, Longinus of Selinus, was taken and done to death by torture at Nicaea.

The Emperor settled large colonies of Isaurians in Thrace.\texthebrew{⁠}\footnote{Theoph. A.M. 5988. Cp. Procopius Gaz. \emph{Panegyr.} c10.

} The brief ascendancy of this people was now over for ever, but it was not to be regretted, for it had served the purpose of averting the far more serious peril of a German ascendancy, which might have brought upon the East the fate of Italy. Henceforward the foreign elements in the army were kept well in control by a preponderance of native troops.

It was fortunate for the Empire that the Isaurian struggle was over before a serious war broke out with Persia, which will

p434
be described in another chapter. But there was fighting from time to time with other enemies. The Blemyes troubled Egypt,\texthebrew{⁠}\footnote{See Joshua Styl. c20.

} the Mazices attacked Libya,\texthebrew{⁠}\footnote{John Ant. \emph{fr.} 74 (\emph{Exc. de virt. et vit.} p205). Probably during the Prefecture of Marinus, which seems to have begun in 512.

} the Tzani overran Pontus.\texthebrew{⁠}\footnote{Theodore Lector, II.19, perhaps in 505 or 506.

} The Saracens of the desert invaded Euphratesia, Syria, and Palestine in 498, but were thoroughly defeated. Another raid four years later was followed by a treaty of peace.\texthebrew{⁠}\footnote{(1) The Saracens of Hira, under Naman, who were vassals of Persia, overran the Euphratesian province and were defeated at Bithrapsas by Eugenius, the military commander in that province. Theoph. A.M. 5990. (2) The Saracens of Ghassan, of whom Harith was chief, overran Palestine and were defeated by Romanus, Dux of Palestine, \emph{ib.} Cp. Evagr. III.36 ; John of Nikiu, c89. (3) In 502 Phoenicia, Syria, and Palestine were overrun again by the bands of Harith, who retreated so quickly that Romanus could not reach them. Theoph. A.M. 5994. For the treaty see \emph{id.} A.M. 5995 and Nonnosus in

\emph{F. H. G.} IV p179 . For these Saracens see above,
Chap. IV § 1 .

} In A.D. 515 Cappadocia was laid waste by an irruption of the Sabeiroi who came down from the region of the Caucasus.\texthebrew{⁠}\footnote{Marcellinus, \emph{sub a.}; John Ant. 103, \emph{Exc. de ins.} p146 (from which it appears that this was a second incursion); John Mal. XVI p406. The Emperor then fortified the larger villages of Cappadocia, and remitted the taxes of the provinces which had suffered, for three years. The Sabeiroi were a Hunnic people (\textgreek{Ο\texthebrew{ὗ}ννοι Σαβήρ}) who lived north of the Caucasus, near the Caspian. Cp. Procopius, \emph{B. P.} II.29, \emph{B. G.} IV.3 and 11; above
p115 .

} But a more dangerous foe than any of these were the Bulgarians beyond the Danube.

After the disruption of the Hunnic empire in A.D. 454, a portion of the Huns had occupied the regions between the mouths of the Danube and the Dniester, where they were ruled by two of the sons of Attila. During the reign of Leo and Zeno, they sometimes raided the Roman provinces and sometimes supplied auxiliaries to the Roman armies.\texthebrew{⁠}\footnote{For the relations of the Empire to the Huns see Priscus, \emph{fr.} 18 \emph{Exc. de leg. gent.} p587 (where we learn that they were ruled by two of Attila's sons, Dengisch and Ernach), and \emph{fr.} 20; Marcellinus and \emph{Chron. Pasch., sub a.} 469 , where the defeat of the Huns and the slaying of Dengisch, whose head was brought to Constantinople, by Anagastus, mag. mil. of Thrace, is recorded (cp. John Ant. \emph{fr.} 89, \emph{Exc. de ins.} 205, where the date is 468, but perhaps the same event is not referred to. The text seems corrupt). In 480 Zeno called on the Huns (Bulgarians) to support him against the Ostrogoths, John Ant. \emph{fr.} 211.4. We have seen that Huns were employed by Anastasius against the Isaurians
(p433, n3) .

} They were kept in check by the Ostrogothic federates, but the departure of Theoderic from Italy had left the field clear for their devastations in Thrace and Illyricum, which throughout the reign of Anastasius suffered severely. These Huns now come to be known under the

p435
name of Bulgarians.\texthebrew{⁠}\footnote{See Marquart, \emph{Die Chronologie der alttürkischen Inschriften}, p77. (Cp. also Zeuss, \emph{Die Deutschen}, etc., 710 \emph{sq.}) The national Bulgarian tradition began the series of their kings with Avitochol, who may well be identical with Attila, and the second is Irnik, in whom we can hardly refuse to recognise Ernach (Attila's favourite son). Cp. Bury, \emph{The Chronological Cycle of the Bulgarians}, \emph{B. Z.} XIX p135.

} But we must distinguish these Bulgarians, who were also known as Unogundurs, from two other great Hunnic hordes who will presently come upon the scene of history: the Kotrigurs who lived between the Dnieper and the Don, and the Utigurs who lived to the south of the Don. These latter peoples were to disappear in the course of time; the Unogundurs were to be the founders of Bulgaria.

The Bulgarians were undoubtedly the foes who invaded the Empire in A.D. 493, defeated a Roman army, and killed Julian, Master of Soldiers.\texthebrew{⁠}\footnote{Marcellinus, \emph{sub a.} (Scythico ferro).

} The next recorded incursion was in A.D. 499, when Aristus, Master of Soldiers in Illyricum, lost more than a quarter of his army of 15,000 men in a battle against the Bulgarians.\texthebrew{⁠}\footnote{\emph{Id. sub a.} The scene of the battle was iuxta Tzurtam fluvium. The Roman army was accompanied by 520 wagons laden with arms. In the following year Anastasius encouraged the Illyrian troops by sending them a donative (\emph{id. sub} 500).

} Their depredations were repeated three years later ( A.D. 502), and on this occasion their progress was unopposed.\texthebrew{⁠}\footnote{\emph{Id. sub a.}, Theoph. A.M. 5994.

} Anastasius had determined to secure at least the immediate neighbourhood of the capital against the raids of the barbarians, and for this purpose he built a Long Wall,\texthebrew{⁠}\footnote{The building is recorded in \emph{Chron. Pasch.}, apparently under Indiction 15 =3rd consul­ship of Anastasius, that is A.D. 507. There are two entries under this year, (1) a demonstration in the circus, in favour of raising Areobindus to the throne; (2) the building of the Wall. Now (1) is recorded much more fully by Marcellinus under Ind. 5 = cons. of Paulus and Muscianus = 512; and all at dates between Ind. 15 and Ind. 6 have fallen out of our text. Hence it was inferred by Ducange that these two entries belong to Ind. 5. Rightly, but there is a deeper error, due not to the scribe but to the chronicler. The building of the Wall is lauded in the \emph{Panegyric} of Procopius (c21), and that oration cannot be dated later than 502 (as C. Kempen has shown in the Preface to his text). My view is that the date of the Wall is 497, which corresponded to an Ind. 5; and that the mistake arose through entering the notice under the Ind. 5 of the following cycle. Cp. above,
p289, n2 .

} the line of which can still be traced, from the Propontis to the Black Sea, at a distance of \texthebrew{•} about 40 miles west of Constantinople. The southern extremity was just to the west of Selymbria, and the northern between Podima and Lake Derkos. The fortification consisted of a stone wall \texthebrew{•} about 11 feet thick, without earthworks or ditch, and traces of round towers projecting \texthebrew{•} about 31 feet in front

p436
have been found. The length of the wall was 41 miles, and it corresponds roughly to the modern Turkish fortifications known as the Chatalja Lines, though the extreme points were further west.\texthebrew{⁠}\footnote{This account of the Wall is taken from C. Schuchhardt, \emph{Die Anastasius Mauer bei Constantinopel und die Dobrudscha-Wälle}, in the \emph{Jahrbuch des k. d. arch. Instituts}, XVI.107 \emph{sqq.} (1901). The dimensions given by

Evagrius, III.38 , and in Suidas (whose source is doubtless John Ant.), \emph{sub}\textgreek{\texthebrew{Ἀ}ναστάσιος} and \emph{sub}\textgreek{Τε\texthebrew{ῖ}χος} disagree with each other, and are all inaccurate. The settlements of Heruls ``in the lands and cities of the Romans'' recorded by Marcellinus, \emph{sub} 512, were evidently designed to strengthen the depopulated lands of the Illyrian peninsula.

} We do not hear of another invasion till A.D. 517, when a host of barbarian cavalry laid waste Macedonia, Epirus, and Thessaly, penetrating as far as Thermopylae.\texthebrew{⁠}\footnote{Marcellinus, \emph{sub a.}, Getae equites. I suspect that these are Bulgarians, whom elsewhere this chronicler calls Bulgares. Otherwise they must be Slavs, who are often designated as Getae. A thousand pounds of gold was sent to redeem the captives, but it was insufficient, and many were put to death.

} The consequences of the devastations of Germans and Huns for more than a hundred years was the depopulation of the Balkan provinces, the decline of its agricultural produce, and a considerable diminution of the Imperial revenue.\footnote{Cp. the undated law of Anastasius in

\emph{C. J.} X.27.2, 10

\textgreek{\texthebrew{ἐ}ν Θρ\texthebrew{ᾴ}κ\texthebrew{ῃ} γ\texthebrew{ὰ}ρ \texthebrew{ἐ}πειδ\texthebrew{ὴ} ο\texthebrew{ὐ}κ ε\texthebrew{ἰ}ς \texthebrew{ὁ}λόκληρον ε\texthebrew{ἰ}σφέρεται τ\texthebrew{ὰ} δημοσία δι\texthebrew{ὰ} τ\texthebrew{ὸ} προφάσει τ\texthebrew{ῶ}ν βαρβαρικ\texthebrew{ῶ}ν \texthebrew{ἐ}φόδων \texthebrew{ἐ}λαττωθ\texthebrew{ῆ}ναι το\texthebrew{ὺ}ς γεωργούς, κα\texthebrew{ὶ} μ\texthebrew{ὴ ἀ}ρκε\texthebrew{ῖ}ν τ\texthebrew{ὴ}ν \texthebrew{ἐ}ν ε\texthebrew{ἴ}δεσι συντέλειαν το\texthebrew{ῖ}ς κατ\texthebrew{}’ α\texthebrew{ὐ}τ\texthebrew{ὴ}ν \texthebrew{ἱ}δρυμένοις στρατιώταις.}}

If the elevation of Anastasius had been popular, his popularity did not continue. His reign was frequently troubled by seditions in Constantinople, which were in many cases provoked by his ecclesiastical policy. His purpose was to maintain the Henotikon of Zeno; his personal predilections were Monophysitic. We are ignorant of the cause of the sedition which broke out in A.D. 493, but it was evidently serious, as the statues of the Emperor and Empress were dragged through the city.\texthebrew{⁠}\footnote{Marcellinus, \emph{sub a.}

} The relations between Anastasius and the Patriarch Euphemius, who had been opposed to his elevation, were strained. Euphemius was devoted to the doctrine of Chalcedon, and had been planning a campaign against the Patriarch of Alexandria, first Peter, and then his successor Athanasius, both of whom anathematised the Council of Chalcedon and the Tome of Leo. Without the Emperor's knowledge he wrote a letter to Felix, the bishop of Rome, invoking his aid. The Patriarchs of Alexandria and Jerusalem informed

p437
the Emperor that Euphemius was a heretic;\texthebrew{⁠}\footnote{That is, a Nestorian.

} and a council was held at Constantinople which confirmed the Henotikon and deposed Euphemius ( A.D. 496).\texthebrew{⁠}\footnote{Zacharias Myt. VII.1. A copy of the letter to Felix was procured by Anastasius through his apocrisiarius at Rome and was sent to the Emperor. Other than purely ecclesiastical reasons entered into the quarrel between Anastasius and Euphemius. Anastasius suspected the Patriarch of secret intrigues with the Isaurian leaders. See Theodore Lector, II.9\texthebrew{‑}12 (who records an attempt on the Patriarch's life in St. Sophia). The same writer says that the Emperor endeavoured to recover from Euphemius the signed declaration of orthodoxy which he had made at his coronation, \emph{ib.} 8. Euphemius was banished to Euchaita.

} This led to a disturbance, and the people, rushing to the Hippodrome, supplicated the Emperor in vain to restore the Patriarch. Macedonius was appointed to the Patriarchal throne. He seems to have held much the same opinions as Euphemius, but he did not scruple to sign the Henotikon.\footnote{Theodore, II.13. The Monophysite Zacharias (VII.7) says that Macedonius ``omitted no intrigue of heart [\emph{sic}] to conceal his opinions.'' He had been a monk of the Akoimetoi, ``of whom there were about one thousand and who lived luxuriously in baths and in other bodily indulgences . . . and were adorned with the semblance of chastity, but were inwardly like whited sepulchres, full of all uncleanness. . . . And he used to celebrate the memory of Nestorius every year, and they used to celebrate with him.'' Perhaps there was some foundation for this attack; the Akoimetoi may have made a habit of personal cleanliness. Orthodox writers describe Macedonius as an ascetic.

}

A serious riot in the Hippodrome occurred in A.D. 498. The Prefect of the City had thrown into prison some members of the Green faction for the not uncommon offence of stone-throwing . The Greens demanded their release, and when the Emperor summoned the Excubitors to suppress them, there was a great uproar. Stones were thrown at the Kathisma, and one of these nearly hit Anastasius. The man who had thrown it was hewn in pieces by the Excubitors, and then the Greens set fire to the Bronze Gate of the Hippodrome. The fire spread not only to the Kathisma but also, in the other direction, to the Forum of Constantinople. Many offenders were punished, but a new Prefect, Plato, was appointed.\footnote{John Mal. XVI p394 (and \emph{Excerpta de ins.} p168) = \emph{Chron. Pasch. sub a.} The deposed Prefect was perhaps Secundinus. The succession of Prefects of the City in this reign seems to have been: Julianus 491; Secundinus 491; Plato 498; Helias (John Ant. \emph{fr.} 103, p142); Constantius Tzzurukkas (already in 501, Marcellinus, \emph{sub a.}); Plato, 512 [Marcellinus, \emph{sub a.}].

}

The pagan festival of the Brytae, which was celebrated with dancing,\texthebrew{⁠}\footnote{John Ant. \emph{ib.} Suidas \emph{sub \textgreek{Μαϊουμ\texthebrew{ᾶ}ς}}, a passage which does not prove that the Maïumas (in May) was identical with the Brytae. Combining Marcellinus with Joshua Styl. p35,

(p438)
we may infer that the date of the second riot, when Constantius was Prefect, was in 501, and the previous riot under Helias (James Ant.) in 500 (or 499). The edict prohibiting the feast was in 502 (Joshua). See further, John Mal. \emph{ib.} The date of Theophanes, A.D. 504\texthebrew{‑}505, must be rejected. More than 3000 were killed, acc. to Marcellinus, \emph{sub} 501. Procopius Gaz. \emph{Pan.} 16 probably refers to the licentiousness of the Brytae.\texthebrew{—} On the celebration of the festival of Brumalia (Nov. 24-Dec. 17) in the fifth and sixth centuries \texthebrew{—} notwithstanding its condemnation by the Church,

John Lydus, \emph{De mens.} IV § 158 \texthebrew{—} see J. R. Crawford, \emph{De Bruma et Brumalibus festis}, in \emph{B. Z.} XXIII.375 \emph{sqq.}

} repeatedly caused sanguinary riots among the demes,

p438
and in one of these disturbances ( A.D. 501) a bastard son of the Emperor was killed, and the Emperor forbade its celebration for the future throughout the Empire, thereby ``depriving the cities of the most beautiful dancing.'' He had already abolished the practice of contests with wild beasts ( A.D. 499).\footnote{Priscian, \emph{Pan.} 223 \emph{sqq.}; Procopius Gaz. \emph{Pan.} 15.

}

In A.D. 511 the Patriarch Macedonius, who no longer concealed his adhesion to the Council of Chalcedon, met the same fate as his predecessors. The Monophysites represented him as plotting against the Emperor, while the orthodox asserted that he was deposed because he declined to give up the profession of orthodoxy signed by the Emperor at his coronation. In any case, Anastasius had begun to move in the Monophysitic direction so far as to abandon the neutral spirit of the Henotikon. The position of Macedonius was not strong, because by signing the Henotikon he had alienated the orthodox monks of the capital. Seeking to win back their confidence he did not scruple to denounce Anastasius as a Manichaean. He was deposed by a local council in August, A.D. 511, was forced to surrender the document with the Emperor's signature, and was banished to Euchaita. Timothy, an undisguised Monophysite, was elected in his stead.

A distinguished Monophysite monk, Severus of Sozopolis, had, a few years before, arrived at Constantinople with a company of two hundred fellow-heretics and had been received with honour by Anastasius.\texthebrew{⁠}\footnote{In A.D. 508. Severus was brought up as a pagan, studied rhetoric at Alexandria and law at Berytus. He was baptized shortly before 490, and soon afterwards became a monk in the monastery of Peter the Iberian not far from Gaza. The cause of his visit to Constantinople was the persecution of Monophysite monks in Palestine by one Nephalius. He remained in the capital for three years. For his life we have two Syriac biographies by Zacharias and John; and some of his letters have been edited and translated by Brooks (see

Bibliography ).

} He caused scandal and disturbances by holding services in which the Trisagion (``Holy, holy, holy, Lord God of Hosts'') was chanted with the Monophysitic addition

p439
``Who wast crucified for us,'' which had been introduced at Antioch fifty years before. The new Patriarch Timothy interpolated this heretical phrase into the liturgy in St. Sophia. Anastasius, supported by the counsels of Marinus, Praetorian Prefect of the East,\texthebrew{⁠}\footnote{See Zacharias Myt. VII.9.

} determined to defy the religious sentiment of the people of Byzantium. On Sunday, Nov. 4 ( A.D. 512),\texthebrew{⁠}\footnote{The date depends on Marcellinus, \emph{sub a.}, whose account is the fullest. It is to be supplemented by \emph{Chron. Pasch.}, under the wrong year 507 , and

Evagrius, III.44 .

} the orthodox multitude in the Church drowned with their shouts the chanting of the heretical priests, and there was such a disturbance that Marinus and Plato, the Prefect of the City, interfered with armed force. Some were slain and others imprisoned. On the following day there was a more sanguinary conflict in the court of a church, and on Tuesday (Nov. 6) the orthodox congregated and formed a camp in the Forum of Constantine. The rioting now assumed the dimensions of a revolt. The general Areobindus was the husband of Juliana Anicia, who was the granddaughter of Valentinian III,\texthebrew{⁠}\footnote{Daughter of Placidia and Olybrius. There is a remarkable portrait of the princess (who died in 527) in the Vienna MS. of the work of Dioscorides on plants, which was written for her. She sits on a throne between the figures of Megalopsychia and Phronesis. The desire of the foundress (\textgreek{Πόθος τ\texthebrew{ῆ}ς φιλοκτίστου}) offers her the book, and the gratitude of the Arts kneels below her. See Kraus, \emph{Gesch. d. christl. Kunst}, I.460; Dalton, \emph{Byz. Art}, 460.

} and thus a member of the Theodosian house. The people proclaimed him Emperor and pulled down the statues of Anastasius. Celer, the Master of Offices, and Patricius, Master of Soldiers in praesenti , who were sent to pacify them, were driven off with showers of stones; the house of Marinus was burnt. On the next day the Emperor sent heralds to the people proclaiming that he was ready to abdicate, and appeared in the Kathisma of the Hippodrome without his crown. He was greeted with demands that Marinus and Plato should be thrown to the beasts. But in some extraordinary way he succeeded in calming the tumult. The crowd begged him to put on his crown and promised good behaviour.

It was unfortunate for the peace of the East that Anastasius was not indifferent in questions of religious doctrine. His reason prompted him to enforce the Henotikon and to lean to neither party in his ecclesiastical measures. He honestly endeavoured to carry out this policy up to the year A.D. 511\texthebrew{‑}512, but he was growing old, and, despairing of maintaining peace

p440
between the extreme parties, he threw himself into the arms of his Monophysite friends. It is to be observed that neither all the orthodox nor all the Monophysites demanded at this time a repudiation of the Henotikon; for the Monophysites could argue that it condemned the doctrines of Chalcedon, the orthodox that it did not.\texthebrew{⁠}\footnote{Those Monophysites who would not accept the Henotikon were known as the Akephaloi.

} The middle party, of whom Flavian of Antioch was the most prominent, sought to act more or less in the true spirit of the act of Zeno and leave the doctrine of Chalcedon severely alone. In the capital the difficulty of preserving peace was aggravated by the agitation of the Sleepless Monks of the monastery of Studion, who were uncompromising opponents of the Henotikon, and remained in communion with the Church of Rome. Some vain attempts had been made to end the schism. Pope Anastasius II, in his brief pontificate,\texthebrew{⁠}\footnote{A.D. 496\texthebrew{‑}498.

} desired to conclude it by a concession which was almost equivalent to a partial acceptance of the Henotikon. He sent to Constantinople two bishops proposing to withdraw the demand of his predecessors that the name of Acacius should be expunged from the roll of Patriarchs. On account of this policy he is one of the Popes for whom the Catholic Church has little good to say, and Dante found for him a suitable place in hell.\texthebrew{⁠}\footnote{\emph{Inferno}, XI.8 .

} His successors obstinately refused to heal the breach.\footnote{See below,

p464 .

}

Far more significant than the deposition of Macedonius, who had never approved of the Imperial policy, was the deposition of the Patriarch of Antioch, the moderate Flavian,\texthebrew{⁠}\footnote{At a synod held at Antioch, Zach. Myt. VII.10.

} and the election of the Pisidian Severus, whom we have already met as the leading theologian of the Monophysites and bitter foe of Chalcedon ( A.D. 512).\texthebrew{⁠}\footnote{The other most prominent Monophysite leader was Xenaias, bishop of Hierapolis, at whose instance a synod was held at Sidon in A.D. 512. Flavian's moderate policy at this synod enabled Xenaias to report to Anastasius that he was a heretic, and his ejection (not without violence) followed. Zach. Myt. VII.10. Severus was ``a confidant and friend'' of Probus, the nephew of Anastasius (\emph{ib.}). The influence of his nephew may well have counted for something in the old Emperor's change of policy; and the influence of Marinus counted too. If (see below,
p470 ) I am right in pla­cing the elevation of Marinus to the Pr. Prefecture in A.D. 512, this too may have some significance.

} On the occasion of his enthronement at Antioch, Severus anathematised the doctrinal decisions of that Council, and he determined to make his own Patriarchate as

p441
Monophysitic as that of Egypt. A synod at Tyre ( A.D. 513)\texthebrew{⁠}\footnote{\emph{Op. cit.} VII.12.

} condemned Chalcedon and confirmed the Henotikon, which was interpreted in the Monophysite sense. The triumphant party were ready for extreme measures, and the Emperor had to warn the Duke of Phoenicia Libanensis that he would countenance no bloodshed in dealing with recalcitrant bishops. But the general proceedings of the Monophysites, under the guidance of Severus, during the next few years, seem to have amounted to a persecution.

The reply to the revolution in the Emperor's policy was soon to come in the shape of a rebellion in Thrace.\footnote{See below,
§ 4 .

}

Anastasius was a conscientious ruler, and one of the great merits of his government was the personal attention which he paid to the control of the finances. A civil servant, who belonged to the bureau of the Praetorian Prefect, and began his career in this reign, asserts that the careful economy of Anastasius and his strictures in supervising the details of the budget saved the State, which ever since the costly expedition of Leo I against the Vandals had been on the brink of financial ruin.\footnote{John Lydus, \emph{De mag.} III.45 .

}

The economy of the Emperor enabled him to abolish the tax on receipts, known as the Chrysargyron, which weighed heavily on the poorest classes of the population.\texthebrew{⁠}\footnote{\emph{C. J.} XI.1 ; Procopius Gaz. \emph{Paneg.} 13; Priscian, \emph{Pan.} 149 \emph{sqq.} (argenti relevans atque auri pondere mundum); Theodore Lector, II.53. The hardship of the tax is described by Zosimus, II.38 . For date cp. Brooks, \emph{C. Med. H.} I.484; Stein, in \emph{Hermes}, LII.578.

} This act (May, A.D. 498) earned for him particular glory and popularity. The reception of the edict in the city of Edessa illustrates the universal joy which the measure evoked. ``The whole city rejoiced, and they all put on white garments, both small and great, and carried lighted tapers and censers full of burning incense,'' and praising the Emperor went to a church and celebrated the eucharist. They kept a merry festival during the whole week and resolved to celebrate this festival every year.\footnote{Joshua Styl. XXXI p22. The amount raised by the tax at Edessa (every four years) was 140 lbs. of gold. Anastasius is said to have burned all the documents relating to the collection of this tax, so as to place a difficulty in the way of its revival. See Procopius and Priscian, \emph{locc. citt.}, and

Evagrius, III.39 , where it is

(p442)
mentioned that the Emperor consulted the Senate. According to Cedrenus (that is, John Skylitzes, whom he transcribed), the hardships of the tax were brought to the attention of Anastasius by a deputation of monks from Jerusalem, and by a tragedy composed by Timotheus of Gaza

(I p627) . Timotheus was a grammarian, and he wrote zoological books on Indian animals, of which excerpts are preserved (see Krumbacher, \emph{Gesch. d. byz. litt.} pp631, 633, 582). It is to be noted that the abolition of the Chrysargyron gave special satisfaction to the Church, because the tax, which fell on the earnings of prostitutes, implicitly gave a legal recognition to vice (see Evagr. \emph{ib.}).

}

p442
The consequent loss of revenue suffered by the fisc was made good by an equivalent contribution from the revenue of the Private Estates.\texthebrew{⁠}\footnote{\textgreek{\texthebrew{Ἐ}κ τ\texthebrew{ῶ}ν \texthebrew{ἰ}δίων α\texthebrew{ὐ}το\texthebrew{ῦ}}, John Mal. XVI p398.

} The Imperial Estates seem to have received considerable additions in this reign, principally from the confiscation of the property of Zeno and the Isaurian rebels. In consequence of this increase, Anastasius found it expedient to institute a new finance minister, with similar functions to those of the Count of the Private Estates, who was to administer the recently acquired domains and all that should in future be acquired by the crown. This minister was designed by the title of Count of the Patrimony.\footnote{\emph{C. J.} I.34.1 . The Greek title was \textgreek{κόμης τ\texthebrew{ῆ}ς \texthebrew{ἰ}δκι\texthebrew{ῆ}ς κτήσεως} (in this constitution the Comes rer. priv. is called \textgreek{κόμης τ\texthebrew{ῆ}ς \texthebrew{ἰ}δκι\texthebrew{ῆ}ς περιουσίας}).

}

Perhaps the most important financial innovation introduced by Anastasius was in the method of collecting the annona . He relieved the town corporations of the responsibility for this troublesome task,\texthebrew{⁠}\footnote{See above,
p59 . The chief source is John Lydus, III.46 , 49 . Cp. John Mal. p400;

Evagr. III.42 .

} and assigned it to officials named vindices , who were probably appointed by the Praetorian Prefect. The appointments seem to have been given by auction to those who promised most,\texthebrew{⁠}\footnote{Lydus, III.49 .

} so that this form was equivalent to a revival of the old system of farming the revenue. Opinion was divided as to the effects of this change. On one hand it was said that the result was to impoverish the provinces;\texthebrew{⁠}\footnote{So Lydus (\emph{ib.}), who belonged to the anti-Marinus faction. Evagrius, \emph{ib.}, says \textgreek{\texthebrew{ὄ}θεν κατ\texthebrew{ὰ} πολ\texthebrew{ὺ} ο\texthebrew{ἵ} τε φόροι διερρύησαν τά τε \texthebrew{ἄ}νθη τ\texthebrew{ῶ}ν πόλεων διέπεσεν.}} on the other, that it was a great relief to the farmers.\texthebrew{⁠}\footnote{Priscian, \emph{Pan.} 193\texthebrew{‑}195:

Agricolas miserans dispendia saeva relaxas;

curia perversis nam cessat moribus omnis,

nec licet iniustis solito contemnere leges.

} One of the abuses which the measure may have been intended to remove was the unfair advantage enjoyed by the richer and more influential landowners, whom the curial bodies were afraid to offend. Under the new system, however, inequality of treatment could

p443
be secured in another way, by bribing the vindices . Anastasius hoped perhaps to mitigate this danger by strengthening the hands of the defensores and bishops, who were expected to protect the rights of subjects against official oppression. Those who condemned the new policy said that the vindices treated the cities like hostile communities.\footnote{John Lydus, \emph{ib.}

}

The originator of this revolutionary measure was an able financier of Syrian birth, named Marinus, who seems to have been the most trusted adviser of Anastasius, throughout the latter part of his reign. He began his career as a financial clerk under the Count of the East,\texthebrew{⁠}\footnote{John Lydus, III.36 . He was a scriniarius (or logothete, John Mal. XVI.400; cp. Stein, \emph{Studien}, p149). The scriniarii were clerks who kept the tax accounts. Originally, according to Lydus, they had no recognised place in the hierarchy of the civil service. They were incorporated in it by Theodosius the Great, and towards the end of the fifth century they became a very important body. The rationales of the financial ministries were recruited from them; and scriniarii sometimes rose to be Praetorian Prefects. John Lydus looked down upon them as mere accountants. They had not the liberal education of the Scholastici. Marinus is highly praised by his fellow heretic Zacharias of Mytilene (VII.9).

} and attained to the post of head of the tax department of the Praetorian Prefect.\texthebrew{⁠}\footnote{John Mal. XVI.400.

} In this capacity he gained the ear of the Emperor, and ultimately was elevated to the Praetorian Prefecture. The reform was probably carried out during his tenure of that post, but the date and duration of his Prefecture are a little uncertain.\texthebrew{⁠}\footnote{See

Appendix to this chapter .

} The immediate result of the new method of collecting the taxes was a considerable increase of the revenue and also of the private income of the Praetorian Prefect.\footnote{John Lydus, III.49 \textgreek{γίνεται μ\texthebrew{ὲ}ν πολύχρυσος ε\texthebrew{ἴ}περ τις \texthebrew{ἄ}λλος \texthebrew{ὁ} βασιλε\texthebrew{ὺ}ς κα\texthebrew{ὶ} μετ\texthebrew{}’ α\texthebrew{ὐ}τ\texthebrew{ὸ}ν \texthebrew{ὁ} Μαρ\texthebrew{ῖ}νος κα\texthebrew{ὶ ὅ}σοι Μαρινι\texthebrew{ῶ}ντες \texthebrew{ἁ}πλ\texthebrew{ῶ}ς.}}

It is not clear whether the reform of Marinus meant that the actual tax-collectors , who had hitherto been members of the town communities, were replaced by government officials. It seems more probable that the change consisted in pla­cing the local collectors under direct government control. They received their instructions from the vindex, and the provincial governor, who remained responsible for the taxation of the province, communicated with the vindex and not with the corporation of decurions. The new system was not permanent. Though it was not completely done away with, it was considerably modified in the following reigns. In some places the vindex survived,

p444
but in most of the provinces he disappeared, and there was probably a return to the old methods.\footnote{The local survival of the vindex is shown by Justinian, \emph{Nov.} 128 §§ 5 and 8; 38 \textgreek{το\texthebrew{ὺ}ς \texthebrew{ὀ}λεθρίους μισθωτ\texthebrew{ὰ}ς ο\texthebrew{ὓ}ς δ\texthebrew{ὴ} βίνδικας καλο\texthebrew{ῦ}σι}. There is clear evidence in the \emph{Novels} of Justinian that the local authorities shared in the collection of the taxes. Probably the system differed in different provinces. The term \textgreek{πολιτευόμενοι} refers to municipal authorities. Cp. \emph{Nov.} 130 § 3, p263; 128 § 5 \textgreek{ε\texthebrew{ἴ}τε \texthebrew{ἄ}ρχοντες ε\texthebrew{ἴ}τε πολιτευόμενοι, ε\texthebrew{ἴ}τε \texthebrew{ἐ}ξάκτορες ε\texthebrew{ἴ}τε βίνδικες ε\texthebrew{ἴ}τε κανονικάριοι} (special emissaries sent by the Praetorian Prefect); \emph{ib.} § 16 \textgreek{σιτ\texthebrew{ῶ}ναι} and \textgreek{διοικηταί} are appointed by the municipalities.

}

Other revenue questions occupied the anxious attention of the government at this period. The practice of converting the annona into money payments seems to have been considerably enlarged.\texthebrew{⁠}\footnote{This seems to be the meaning of the \textgreek{χρυσοτέλεια τ\texthebrew{ῶ}ν \texthebrew{ἰ}ούγων} introduced by Anastasius, John Mal. XVI.394;

Evagrius, III.42 . Cp. \emph{C. J.} X.27.2 .

} But the problem of sterile lands appears now to have become more acute than ever. This grave difficulty perpetually solicited the care and defied the statesman ­ship of the Imperial government. Farms were constantly falling out of cultivation through the impoverishment of their owners or the deficiency of labour. The heavy public burdens, aggravated by the oppression of officials, reduced many of the small struggling farmers to bankruptcy. This would have meant a considerable loss to the revenue, in the natural course of things, and the problem for the government was to avoid this loss by making others suffer for the unfortunate defaulters. For this purpose the small properties of the free farmers of a commune were regarded as a fiscal unity, liable for the total sum of the fiscal assessments of its members;\texthebrew{⁠}\footnote{These lands were hence called \textgreek{\texthebrew{ὁ}μόκηνσα.}} and when for any cause one property ceased to be solvent, the others were required to make good the deficiency. This addition to their proper contributions was known as an epibole .\texthebrew{⁠}\footnote{Adiectio sterilium.

} In the case of larger estates, which were not included in a commune, if one part became unproductive, the whole estate remained liable for the tax as originally estimated.\texthebrew{⁠}\footnote{Called \textgreek{\texthebrew{ὁ}μόδουλα.}} But a difficulty rose when parts of such an estate were sold or when it was divided among several heirs. Notwithstanding the division it was still treated as a fiscal unity, and if one of the proprietors became insolvent the government was determined that the deficiency should be made good by other portions of the original estate.\texthebrew{⁠}\footnote{\emph{C. Th.} XIII.11.9 .

} But there was a considerable difference of opinion as to the apportionment of the epibole in such a

p445
case. Should the whole estate be liable, or should the sterile property be annexed, along with its obligations, to the productive land in its immediate neighbourhood? The former solution would have assimilated the treatment of these estates to the lands of the communes. It is not clear what method was applied before the sixth century. We only know that the epibole in the two cases was not the same. In the reign of Anastasius an attempt seems to have been made to break down the distinction, and to have been successfully opposed by the Praetorian Prefect Zoticus ( A.D. 511\texthebrew{‑}512).\texthebrew{⁠}\footnote{Justinian, \emph{Nov.} 168, seems to be a fragment of a praetorian edict of Zoticus. It lays down that the \textgreek{\texthebrew{ἐ}πιβολή} only concerns property included in the census, and therefore does not apply to houses (in towns) as only farms and agricultural lands (\textgreek{χωρία}) are included in the census. For the tendency to assimilate \textgreek{\texthebrew{ὀ}μόκηνσα} and \textgreek{\texthebrew{ὀ}μόδουλα} see an additional fragment in Kroll's note \emph{ad loc.} A law of Anastasius lays down that the lands of the Imperial patrimony are not to be treated on the same principle as \textgreek{\texthebrew{ὀ}μόκηνσα}, which must mean that they are to be treated as \textgreek{\texthebrew{ὀ}μόδουλα}(\emph{C. J.} I.34.2) .

} Perhaps he defined the general method of dealing with sterile lands which was developed in the following reign by the Praetorian Prefect Demosthenes ( A.D. 520\texthebrew{‑}524).\texthebrew{⁠}\footnote{The edict of Demosthenes, addressed to the governor of Lydia, \textgreek{περ\texthebrew{ὶ ἀ}πόρων \texthebrew{ἐ}πιβολ\texthebrew{ῆ}ς}, is extant in the collection of Justinian's \emph{Novels} (166). The general tenor of the edict is: If a farm or a whole complex of property is sold by its proprietor (A), or on his death passes either to his children or to heirs who are outsiders (B); and if the purchasers or heirs should similarly alienate; and if the alienated land should become unproductive, then the \textgreek{\texthebrew{ἐ}πιβολή} is to fall on the property of the last purchaser or inheritor (C), not on all those who formerly possessed it. But if the last acquirer (C) is insolvent, then the burden must fall on those from whom he immediately acquired it (B). If they are insolvent, then the epibole shall be imposed on the original proprietor (A). Those on whom the epibole falls, whether few or many, shall bear it in proportion to the value of their fertile possessions. It seems evident that this edict was provoked by a particular case which the governor of Lydia referred to the Prefect. On the subject of the \textgreek{\texthebrew{ἐ}πιβολή}, see Zacharia , \emph{Gesch. d. gr.-röm. Retest}, ed. 3, 228 \emph{sqq.}; Monnier, \emph{Études de droit byzantin}, 345 \emph{sqq.}, 514 \emph{sqq.}, 642 \emph{sqq.}; Panchenko, \emph{O tainoi istorii Prokopiia}, 138 \emph{sqq.} I think we may fairly infer from the evidence (see last note) that the principle which governed the epibole in the case of \textgreek{\texthebrew{ὀ}μόκηνσα} was that of proximity. See further Justinian, \emph{Nov.} 128 §§ 7, 8. A remission of the epibole is mentioned in Joshua Stylites, c39, where Wright's translation has erroneously ``two folles.''

} The most important points in this ruling were, that the provincial governor was empowered to decide in each case on whom the epibole should fall; that the unproductive land, with all that appertained to it, including the colons, should be transferred to those who were made liable for its burdens; and that this liability should be determined not by proximity, but by the history of the property.

p446
The result of the economical policy of Anastasius and his financial reforms was that he not only saved the State from the bankruptcy which had threatened it, but, at his death, left in the treasury what in those days was a large reserve, amounting to 320,000 pounds of gold (about £14,590,000).\texthebrew{⁠}\footnote{Procopius, \emph{H. A.} 19 , where he is described as ``the most provident and economical of all Emperors.''

} His strict control of expenditure made him extremely unpopular with the official classes whose pockets suffered, and his saving policy, which probably included a great reduction of the expenses of the court, did not endear him to the nobles and ladies accustomed to the pageants and pleasures of Byzantine festivals. He was accused of avarice and stinginess, vices for which the men of Dyrrhachium, his native place, had a bad repute.\texthebrew{⁠}\footnote{John Lydus, \emph{De mag.} III.46 , quotes malicious verses which were placed on an iron statue of the Emperor in the Hippodrome.

} This accusation was unjust, and can be refuted by the admissions of one of the writers who report it.\texthebrew{⁠}\footnote{\emph{Ib.} 47 \textgreek{μεγαλόδωρος.}} Personally Anastasius was generous and open-handed ; he seldom sent any petitioners empty away; and several instances of his liberality to individuals are recorded. His ``parsimonious resource ­fulness,'' stigmatised by his successor Justin,\texthebrew{⁠}\footnote{\emph{C. J.} II.7.25

parca posterioris subtilitas principis.

} was entirely in the interests of the State; and the general tenor of his policy was to finance the Empire by economy in expenditure, and not to increase, but rather to reduce, the public burdens.\texthebrew{⁠}\footnote{His remission of arrears is recorded by John Ant. \emph{fr.} 100 (\emph{Exc. de ins.} p141), where it is also implied that confiscations of property were infrequent during his reign. The land taxes were remitted constantly in Mesopotamia during the Persian war (Joshua Styl. pp55, 63, 71, 75).

} This feature of his administration corresponded to his character. Though resolute and energetic, he was distinguished, like Nerva, by his mildness.

If he had not held heretical opinions, historians would have had little but praise for the Emperor Anastasius.

It remains to mention his useful monetary reform. For a long time past the general public had suffered great inconvenience through the bad quality of the copper money in circulation. It consisted of coins of very small denomination with no marks of value. Anastasius introduced a large copper follis, equivalent

p447
to forty sesterces, with smaller coins of the value of twenty, ten, and five sesterces, each clearly marked by a letter showing the value.\texthebrew{⁠}\footnote{M, K, I, and E. See Wroh, \emph{Imperial Byz. Coins}, I. XIII, XIV; LXXVIII\texthebrew{‑}IX. This type of bronze coinage remained current till the last quarter of the seventh century. The reform is noted in two texts, (1) the difficult and much discussed passage in Marcellinus, \emph{Chron., sub} 498, and (2) John Mal. XVI p400, which has been generally over­looked. From Malalas we learn that John the Paphlagonia, comes s. larg., carried out the reform: \textgreek{\texthebrew{ἅ}παν τ\texthebrew{ὸ} προχωρ\texthebrew{ὸ}ν κέρμα τ\texthebrew{ὸ} λεπτ\texthebrew{ὸ}ν \texthebrew{ἐ}ποίησε φολλερ\texthebrew{ὰ} προχωρε\texthebrew{ῖ}ν ε\texthebrew{ἰ}ς π\texthebrew{ᾶ}σαν τ\texthebrew{ὴ}ν \texthebrew{Ῥ}ωμαικ\texthebrew{ὴ}ν κατάστασιν \texthebrew{ἔ}κτοτε}. For \textgreek{προχωρ\texthebrew{ὸ}ν} we should, I think, read \textgreek{προχωρο\texthebrew{ῦ}ν} (an inspection of the unique Oxford MS. suggests that this was originally written. \textgreek{πρόχειρον} is another possibility). Perhaps a participle has fallen out. But the passage means, ``He converted all the small copper currency into follera which circulated henceforward in the Empire.'' Marcellinus says that the Romans called the new coins Terentiani, the Greeks follares \texthebrew{—} which corresponds to \textgreek{φολλερά}. The following table will show the relations of the chief gold, silver, and copper coins:

1 nomisma, or solidus (gold)

= 12 miliaresia (silver).

1 miliaresion º

= 2 keratia (siliquae) silver.

1 keration

= 6 M folles or follera (copper).

1 M follis

= 2 K coins (oboloi).

1 K coin

= 2 I coins (dekanumia).

1 I coin

= 2 E coins (pentanumia).

There were two small gold coins, the semissis = ½ nomisma and the tremissis = \texthebrew{1/3} nomisma. Roughly speaking the miliaresion corresponds to our shilling, the keration to sixpence, the follis to a penny.

} This mintage was a great practical benefit, and must have been highly appreciated by the poorer citizens.

He was always ready to spend money on useful public works. Besides the Long Wall of Thrace, he constructed a canal in Bithynia connecting the Gulf of Nicomedia with Lake Sophon, and thus realised an old project of the younger Pliny.\texthebrew{⁠} a Liberal sums were always forthcoming to repair injuries caused by war, to assist towns which were damaged by earthquake, to cleanse harbours, to build aqueducts or baths.\footnote{Cp. John Mal. XVI p409; John Lydus, III.47 ; Joshua Styl. p69. For the canal see Anna Comnena, X.5.

}

Partly through his religious policy and partly through his public economy Anastasius failed to secure the goodwill of various classes of his subjects; his unpopularity increased in the later years of his reign; and it was not surprising that an ambitious soldier should conceive the hope of dethroning him. Vitalian held the post of Count of the Federates, who were stationed in Thrace, and these troops now consisted chiefly of

p448
Bulgarians.\texthebrew{⁠}\footnote{His father Patriciolus also held the office of Count of the Federates (acc. to Theoph.), and he took part in the Persian war. He was a native of Zaldaba in Lower Moesia. It is possible that the family was of Gothic descent (Zacharias Myt VII.13). For the Federates see below,

vol. II chap. xvi § 1 .

} The immediate pretext for his revolt was the conduct of Hypatius, the Master of Soldiers in Thrace, whom the Federates regarded as responsible for depriving them of the provisions to whom they were entitled. But Vitalian claimed to be more than merely the leader of aggrieved soldiers.\texthebrew{⁠}\footnote{Acc. to Zach. Myt., he had a personal reason for hatred of Hypatius (\emph{ib.}).

} He pretended to represent the religious discontent, to voice orthodox indignation at the new form of the \emph{Trisagion}, and to champion the cause of the deposed Patriarch Flavian who was his personal friend, and the deposed Patriarch Macedonius. Vitalian was a man of exceptionally small stature and afflicted with a stammer; his enemies acknowledged his courage and cunning in war.

Hypatius seems to have been unpopular with the army. In A.D. 513\texthebrew{⁠}\footnote{The outbreak of the rising is generally placed in A.D. 514 (cp. Marcellinus). But the evidence in Wright, \emph{Catalogue Syn. MSS. Brit. Mus.} 333, adduced by Brooks (\emph{C. Med. H.} I.485) shows that the true date is 513, and there is nothing inconsistent with this in John Ant.

} Vitalian, by stratagem, compassed the death of two of the chief officers of the general's staff; gained over to his side the Duke of Lower Moesia; and then, capturing Carinus, a trusted friend of Hypatius, granted him his life on condition that he should help him to seize Odessus. Hypatius, unable to cope with the situation, withdrew to Constantinople. The rebel reinforced his Federate troops by a multitude of rustics, and, at the head of 50,000 men (it is said), advanced to Constantinople, hoping that the populace of the capital would rally to him as the champion of orthodoxy.

The Emperor commanded bronze crosses to be set up over the gates of the city, with inscriptions setting forth his own view of the cause of the rebellion.\texthebrew{⁠}\footnote{The object of the manifesto was doubtless to show that Vitalian's champion­ship of orthodoxy was only a pretext.

} He reduced by one-quarter the tax on the import of live stock for the inhabitants of Bithynia and Asia, in order to secure the loyalty of these provinces. The military authorities made what arrangements they could to meet the sudden crisis. When Vitalian occupied the suburbs and appeared before the walls, Patricius, Master of Soldiers in praesenti , who had won distinction in the Persian war and

p449
had considerably helped the advancement of Vitalian, was sent to confer with the rebel. Vitalian explained the purpose of his resort to arms. He was determined to rectify the injustices committed by Hypatius, and to obtain the ratification of the orthodox theological creed. He and his chief officers were invited into the city to discuss the matters at issue. He refused to accept the invitation himself, but his chief officers went on the following day and had an audience of the Emperor. Anastasius won them over by gifts and promises that the soldiers would receive all that was due, and by undertaking that the Church of Rome would be allowed to settle the religious questions in dispute. Vitalian had no option but to yield to the unanimous opinion of his officers, and he returned with his army to Lower Moesia to bide his time and mature new schemes.

The Emperor deposed the unpopular Hypatius and appointed in his stead Cyril, an officer of some experience, who immediately proceeded to Lower Moesia, perhaps with the purpose of capturing Vitalian by guile. But Vitalian was on the alert, and Cyril was assassinated. This act made it clear that the rebel was still a rebel, and a decree of the Senate was passed, in old Roman style, that Vitalian was an enemy of the republic. Alathar, a soldier of Hunnic origin, was appointed to succeed Cyril, but the supreme command of the Imperial army was assigned to another Hypatius, a nephew of the Emperor. This army, said to have been 80,000 strong, gained an inconsiderable victory (autumn, A.D. 513), which was soon followed by serious reverses.\texthebrew{⁠}\footnote{Julian, a clerk in the bureau of the Magister memoriae, was carried about in a cage until he was ransomed. We need not doubt that the numbers both of the army and of the losses are grossly exaggerated.

} Hypatius then fortified himself behind a rampart of wagons at Acris, on the Black Sea, near Odessus. In this entrenchment the barbarians attacked him, and, assisted by a sudden darkness, which a superstitious historian attributed to magic arts, gained a signal victory. The Romans, driven over precipices and into ravines, are said to have lost about 60,000 men. Hypatius himself ran into the sea, if perchance he might conceal himself in the waves, but his head betrayed him. Vitalian preserved him alive as a valuable hostage.\texthebrew{⁠}\footnote{Alathar was captured also and other officers. Vitalian paid ransoms to the Bulgarians who had taken them.

} This victory enabled him to pay his barbarian allies richly, and placed him in possession

p450
of all the cities and fortresses in Moesia and Scythia. The Emperor sent ambassadors with ten pounds of gold to ransom his nephew, but they were captured at Sozopolis (Sizeboli), which at the same time fell into the rebels' hands.

In the meantime a tumult, attended with loss of life, occurred at Constantinople, because Anastasius forbade the celebration of festivities in the evening on account of disorders in the Hippodrome. Among others the Prefect of the Watch was slain. This disturbance may have helped to dispose the Emperor to consider a compromise, when shortly afterwards ( A.D. 514) Vitalian, flushed with victory, appeared in the neighbourhood of the capital. He had collected in the Thracian ports a fleet of 200 vessels. These he sent to the Bosphorus, and marching himself along the coast occupied the European shores of the Straits. A certain John,\texthebrew{⁠}\footnote{He was known as son of Valeriane. This designation by the mother's name is very unusual. John Ant. \emph{ib.} p146.

} who seems to have been Master of Soldiers in praesenti , was sent to Sosthenion (Stenia) to treat with him. Conditions were arranged. Vitalian was appointed to the post of Master of Soldiers in Thrace, and Hypatius was liberated for a ransom of 9000 pounds of gold.

But the most important provision of the contract was that measures should be taken to establish peace in the Church by the convocation of a general Council, and it was agreed that a Council should be held at Heraclea in the following year.\texthebrew{⁠}\footnote{Victor Tonn. \emph{Chron., sub} 514. Theoph. A.M. 6006.

} Vitalian expressly insisted that Rome should be represented, and it was arranged that both he and the Emperor should communicate with Pope Hormisdas.\texthebrew{⁠}\footnote{We have the letters of Anastasius to Hormisdas: \emph{Coll. Avell., Ep.} 109 (Dec. 28, 514) and \emph{Ep.} 107 (Jan. 12, 515), of which the latter arrived at Rome first; the replies of Hormisdas, \emph{Ep.} 110 (July 8) and \emph{Ep.} 108 (April 4), his letter to Anastasius sent by the bishops who did not leave Rome till August, \emph{Ep.} 115 (Aug. 11), and the Indiculus of instructions to the bishops as to their behaviour, 116. In this document the Pope's correspondence with Vitalian is mentioned (p514), but it has not been preserved. The bishops returned to Rome, before the end of the year, with a letter from the Emperor to the Pope, containing a profession of faith and alleging that if he yielded on the question of Acacius, bloodshed would be the consequence, \emph{Ep.} 125. In July 516 he again wrote to Hormisdas, in the interests of unity, and at the same time to the Roman Senate, asking it to exert its influence with the pontiff, \emph{Epp.} 111, 113; he was told that the restoration of unity entirely rested with him, \emph{Epp.} 113, 114. In 517 there was a further interchange of letters, \emph{Epp.} 126, 127 (cp. 128, 129, 130), and finally Anastasius angrily broke off the correspondence, saying that he might put up with insult, but he would not tolerate being ordered, \emph{Ep.} 138 (July 11).

} The date of the Council was

p451
fixed for July 1, A.D. 515, but it never met. Delegates indeed were sent from Rome and arrived at Constantinople late in the year, but as the Pope adopted an uncompromising attitude in regard to the condemnation of the memory of Acacius, and as the Emperor held that it was unjust that living persons should be excluded from the Church on account of the dead,\texthebrew{⁠}\footnote{\emph{Ep.} 125, p539.

} no conciliation could be effected. A fruitless correspondence between Hormisdas and Anastasius ensued.

The Emperor appears to have also promised Vitalian that the bishops who had been driven from their sees should be restored,\texthebrew{⁠}\footnote{Victor and Theophanes, \emph{locc. citt.}

} but it is not clear whether this measure was intended to depend on the decisions of the Council. As the Council did not meet, and as the bishops were not restored, Vitalian was convinced that the Emperor had no intention of fulfilling his part of the bargain, and it was probably in the later months of the same year that he assembled his fleet anew, and reappeared with his army on the banks of the Bosphorus,\texthebrew{⁠}\footnote{The fullest account of the events of this year is given by John Malalas (and is summarised by Evagrius). His story does not completely tally with that of John Ant., who does not say a word about Marinus.

} whence he occupied Sycae, the region of the city, on the north side of the Golden Horn, which was in later times called Galata. It is surprising to find that the command of the Imperial forces was committed to Marinus, the Emperor's influential adviser, who had hitherto been employed only in civil affairs. This exceptional arrangement was due to the attitude of the two Masters of Soldiers in praesenti , Patricius and John, who were personal friends of Vitalian and his father. They hesitated to take command on the ground that if they were defeated they would be suspected of treason. The great financier, however, was equal to the crisis. The issue was decided by a naval battle at the mouth of the Golden Horn, in which the ships of the rebel were completely routed.\texthebrew{⁠}\footnote{The place is designated as opposite the Church of St. Thecla in Sycae, in the part of the Golden Horn \textgreek{\texthebrew{ὄ}που λέγεται τ\texthebrew{ὸ} Βυθάριν}, John Mal. 405 (\textgreek{περ\texthebrew{ὶ} τ\texthebrew{ὰ} καλούμενα Βυθάρια,}Evagrius, III.43 ). I know no other mention of the \textgreek{Βυθάρια}. John Ant.'s account is different. He says that a fast vessel commanded by Justin (Count of the Excubitors, afterwards Emperor) engaged with one of the enemy's ships off Chrysopolis and captured the crew, and that this success caused the flight of the other rebel ships. This is incredible as an account of the naval action; the exploit of Justin can only have been one incident.

} It is related that this victory was achieved by the use of a chemical compound, similar to the

p452
Greek fire of later days, which, projected upon the enemy's ships, set them on fire.\texthebrew{⁠}\footnote{This compound, according to John Mal., was supplied to Marinus by an Athens man of science named Proclus (not to be confounded with the famous Neoplatonist who had died in A.D. 485), and Proclus is said to have refused a reward of 400 lbs. of gold.

} Marinus then landed his forces at Sycae, slew the rebels whom he found there, and in the evening took up a position on the shores of the Bosphorus.\texthebrew{⁠}\footnote{The meaning of Anaplûs, which occurs in our sources (John Ant., John Mal., Evagr.), and has cause some difficulty, has been elucidated by Pargoire (\emph{Anaple et Sosthène}, in \emph{Izv. russk. arkh. Inst. v Kplie}, III.60 \emph{sqq.}). In these passages the \textgreek{\texthebrew{Ἀ}νάπλους} designates the whole European shore of the Bosphorus, or at all events the whole southern strip from Stenia southwards. But it is also found, in other texts, with two more restricted local meanings, designating points on the European shore corresponding to (1) Kuru Chesme and Arnaut Keui (see Marcellinus, \emph{Chron., sub} 481) and (2) Rumili Hissar. The first of these places was also called Hestiae, where there was a Church of St. Michael, built by Constantine, not to be confounded with that of Sosthenion (=Laosthenion), now Stenia, north of Rumili Hissar.

} In the night Vitalian fled with all the troops that were left to him and reached Anchialus, where he seems to have remained undisturbed during the next three years. The Emperor made a solemn procession to Sosthenion, which Vitalian had made his headquarters, and in the church of St. Michael, for which that place was noted, offered thanks to the archangel for the deliverance. All the rebels did not escape as easily as Vitalian. Tarrach, one of his henchmen, whom he had employed to assassinate Cyril, was burned at Chalcedon, and two others who happened to be taken were put to death.

The Empress Ariadne died in this year.\texthebrew{⁠}\footnote{Marcellinus, \emph{Chron., sub a.}

} Anastasius survived her by three years. He died at the age of eighty on the night of July 8\texthebrew{‑}9, A.D. 518.\texthebrew{⁠}\footnote{This date (which is given in Cyril, \emph{Vita S. Sabae}, p354) follows from the fact that Justin was elected on July 9 (John Mal. XVII p411 = \emph{Chron. Pasch., sub a.} ) and that Anastasius died during the previous night (Peter Patr. apud Const. Porph. \emph{De cer.} I.93 ). This agrees with the length of the reign of Anastasius given by Marcellinus, \emph{sub a.} Therefore the date of Theophanes (\emph{Chron., sub a.}), April 9, is false. See Tillemont, \emph{Hist. des Empereurs}, VI.586.

} He had no children and made no provision for the succession, though it was probably his intention to designate one of his three nephews, Probus, Pompeius, or Hypatius.\texthebrew{⁠}\footnote{Cp.
Anon. Val. 13 .

} His last months seem to have been troubled by new hostilities on the part of Vitalian, but the details are unknown to us.\footnote{Cyril, \emph{op. cit.} p340 \textgreek{συνεχόμενος \texthebrew{ὑ}π\texthebrew{ὸ} τ\texthebrew{ῶ}ν Βιταλιανο\texthebrew{ῦ} βαρβαρικ\texthebrew{ῶ}ν \texthebrew{ὀ}χλήσεων}. This was soon after the death of the Patriarch Timotheus, that is after April 5, A.D. 518. See Andreev, \emph{Konstantinopol'skie Patriarkhi}, p168.}

\xchapter{Chapter XIII, Part B}

\xsubsection{(Part 2 of 2)}

The rule of the Patrician Theoderic in Italy, if we date it from the battle of the Adda in A.D. 490, lasted thirty-six years. In its general constitutional and administrative principles it was a continuation of the rule of Odovacar. One of the first things Theoderic had to do was to settle his own people in the land, and this settlement was exactly similar to that which had been carried out by his predecessor. The Ostrogoths for the most part replaced Odovacar's Germans, who had been largely killed or driven out, though some of them who had submitted were permitted to retain their lands. The general principle was the assignment of one-third of the Roman estates to the Goths;\texthebrew{⁠}\footnote{A different view is maintained by Dumoulin (\emph{C. Med. H.} I.447). He thinks that the lands assigned to the Germans both by Odovacar and by Theoderic were one-third of the State lands (ager publicus). It may be doubted whether the number of the Ostrogothic army exceeded 25,000. Hodgkin (III.202) puts it at 40,000, and the number of the whole nation at 200,000. This figure seems too high.

} but the commission which carried out the division was under the presidency of a senator, Liberius, so that we may be sure the senatorial domains were spared so far as possible.

For six years the Emperor Anastasius hesitated to define his attitude to Theoderic,\texthebrew{⁠}\footnote{We saw that Theoderic, after his victory in 490, sent Flavius Festus, the chief of the Roman Senate, as an ambassador to Zeno (above,
p224 ). While Festus was still at Constantinople, Anastasius succeeded and refused to recognise Theoderic. A second embassy was sent in 492, led by another distinguished senator, Flavius Anicius Probus Faustus Niger (CIL VI.32 , 195), consul in 490, whom Theoderic had appointed Master of Offices. The result of the negotiations of Faustus was a partial recognition, as was shown by the fact that Anastasius permitted two western consuls to be nominated in 494. But Anastasius suspected Theoderic's intentions, and there was a breach. Faustus returned to Italy in 494. Then at the end of 496 (after the death of Pope Gelasius and the election of Anastasius II) Festus was again sent, and succeeded in concluding the definite arrangement of 497. The whole course of these negotiations has been ably examined by Sundwall (\emph{Abh.} 190 \emph{sqq.}), who makes it probable that they were closely affected by the ecclesiastical schism, and that the Synod of Rome held in May 495 by Gelasius and the intransigent attitude of the Italian bishops made it difficult for Anastasius to come to terms with the Senate, as the Senate itself was divided on the ecclesiastical question. That Theoderic depended mainly on the support of the Senate for regularising his position comes out very clearly in these transactions.

} but Theoderic carefully refrained from taking any measures that were incompatible with the position of a viceroy or that would render subsequent recognition difficult. At length they came to terms ( A.D. 497), and a definite arrangement was made which determined the position of Italy and the

p454
status of the Ostrogothic kingdom. Theoderic still held the office of Master of Soldiers which Zeno had conferred upon him. Anastasius confirmed him in this office and recognised him as Governor of Italy under certain conditions, which in their general scope must have corresponded to the arrangement which Zeno had made with Odovacar. These conditions determined the constitutional position of Theoderic.

Under this arrangement Italy remained part of the Empire, and was regarded as such officially both at Rome and at Constantinople. In one sense Theoderic was an independent ruler, but there were a number of limitations to his power, which implied the sovranty of the Emperor and which he loyally observed.\footnote{The following account is based on Mommsen's \emph{Ostgotische Studien} in \emph{Hist. Schr.} III p362 \emph{sqq.}

}

The position of the Ostrogothic king as a deputy comes out in the fact that he never used the years of his reign for the purpose of dating official documents. It comes out in the fact that he did not claim the right of coining money except in subordination to the Emperor.\texthebrew{⁠}\footnote{Under Theoderic, and under Odovacar before him, gold coins minted at Ravenna and Rome bore the name and types of the contemporary Emperors. Odovacar struck silver and bronze coins with his own name and portrait (the thick moustache is realistic). Theoderic's silver coins have the Emperor's bust on the obverse and his own monogram on the reverse. The bronze have the Imperial bust on the obverse. The only known coin on which Theoderic's bust appears is a large triple solidus, obviously struck for some particular occasion. Only one specimen is extant. It has been supposed that the bust, which is almost a half-length figure, was copied from an actual statue or mosaic picture of Theoderic. We know that such figures existed. See Wroth, \emph{Catalogue of the Coins of the Vandals, etc.} xxxi\texthebrew{‑}xxxii. Theoderic's coinage is ``singularly neat and even elegant'' (\emph{ib.}).

} It comes out, above all, in the fact that he did not make laws.\texthebrew{⁠}\footnote{In Procopius this is expressly asserted both of Theoderic and of his successors by representatives of the Goths. \emph{B. G.} II.6 p176.

} To make laws, leges in the full sense of the term, was reserved as the supreme prerogative of the Emperor. Ordinances of Theoderic exist, but they are not leges , they are only edicta ; and various high officials, especially the Praetorian Prefect, could issue an edictum . Nor was this difference between law and edict, in Theoderic's case, a mere difference in name. Theoderic did promulgate general edicts, that is, laws which did not apply only to special cases, but were of a general kind permanently valid, and which if they had been enacted by the Emperor would have been called laws. But the Praetorian Prefect had the right of issuing a general

p455
edict, \emph{provided it did not run counter to any existing law.} This meant that he could modify existing laws in particular points, whether in the direction of mildness or of severity, but could not originate any new principle or institution. The ordinances of Theoderic, which are collected in his code known as the
\emph{Edictum Theoderici}, exhibit conformity to this rule. They introduce no novelties, they alter no established principle. We are told that, when Theoderic first appeared in Rome, he addressed the people and promised that he would preserve inviolate all the ordinances of the Emperors in the past.\texthebrew{⁠}\footnote{Anon. Val. 66 . Compare Cassiodorus, \emph{Var.}
I.1 ;

XI.8 \emph{ad init.}

The general conservatism of Theoderic is emphasised in

\emph{Var.} III.9

propositi quidem nostri est nova construere sed amplius vetusta servare. Cp. I.25.

} Thus in legislation, Theoderic is neither nominally nor actually co-ordinate with the Emperor. His powers in this department are those of a high official, and though he employed them to a greater extent than any Praetorian Prefect could have done, on account of the circumstances of the case, yet his edicts are qualitatively on the same footing.

The right of naming one of the consuls of the year, which had belonged to the Emperor reigning in the West, was transferred by the Emperors Zeno and Anastasius to Odovacar and Theoderic.\texthebrew{⁠}\footnote{For the details of the arrangements as to the consulate see Mommsen, \emph{op. cit.} 226 \emph{sqq.}

} From A.D. 498 Theoderic nominated one of the consuls. On one occasion ( A.D. 522) the Emperor Justin waived his own nomination and allowed Theoderic to name both consuls \texthebrew{—} Symmachus and Boethius. But in exercising this right the Ostrogothic king was bound by one restriction. He could not nominate a Goth; only a Roman could fill the consul ­ship. The single exception corroborates the existence of the rule. In A.D. 519 Eutharic, the king's son-in\texthebrew{‑}law, was consul. But it is expressly recorded that the nomination was not made by Theoderic; it was made by the Emperor, as a special favour.\footnote{Cassiodorus, \emph{Var.} VIII.1 .

}

The capitulation which excluded Goths from the consul ­ship extended also to all the civil offices, which were maintained under Ostrogothic rule, as under that of Odovacar.\texthebrew{⁠}\footnote{It is to be noted that in most of his appointments to important offices Theoderic communicated his intentions to, or consulted with, the Senate.

} There was still the Praetorian Prefect of Italy, and when Theoderic acquired Provence, the office of Praetorian Prefect of Gaul was revived. There was the Vicarius of Rome; there were all the provincial governors, divided as before into the three ranks of consulars,

p456
correctors, and praesides . There was the Master of Offices. There were the two great finance ministries.\texthebrew{⁠}\footnote{A comes patrimonii was instituted, Odovacar's vicedominus (see above,
p409 ) under another name. Goths were eligible for this post.

} There was the Quaestorship of the Palace.\texthebrew{⁠}\footnote{All the officia or staffs of subordinate officials were maintained. In the State documents of Cassiodorus, officium nostrum means the staff of the Master of Offices. Both this minister and the Praetorian Prefect resided at Ravenna, but had representatives at Rome who, like themselves, were illustres.

} It may be added that Goths were also excluded from the honorary dignity of Patricius. Under Theoderic no Goth bore that title but Theoderic himself, who had received it from the Emperor.

The Roman Senate, to which Goths on the same principle could not belong, continued to meet and to perform much the same functions which it had performed throughout the fifth century. It was formally recognised by Theoderic as possessing an authority similar to his own.\footnote{Parem nobiscum reipublicae debetis adnisum,

Cassiod. \emph{Var.} II.24 .

}

If all the civil offices were reserved for the Romans, in the case of military posts it was exactly the reverse. Here it was the Romans who were excluded. The army was entirely Gothic; no Roman was liable to military service; and the officers were naturally Goths.\texthebrew{⁠}\footnote{The chief officers were called priors or counts.

} Theoderic was the commander of the army, as Master of Soldiers, for, though he did not designate himself by the title, he had retained the office, and no Master of Soldiers was appointed, subordinate to himself.\texthebrew{⁠}\footnote{Mommsen has illustrated this point by certain measures taken after Theoderic's death. His successor, Athalaric, was out of the question as commander of the forces, and the regent Amalasuntha appointed Tuluin, a Gothic warrior, and Liberius, a Roman, who was Praet. Prefect of Gaul, to be patricii praesentales. This involved two deviations from rule. Tuluin as a Goth was debarred from the dignity of patrician, and Liberius, as a Roman, from a military command. The office was simply that of mag. mil.; the moderation of the title illustrates the fact that the Mastership of Soldiers had become closely associated with the kingship through its long tenure by Theoderic. But I question whether Mommsen is right in assuming that Theoderic simply continued throughout his reign to hold the Mastership conferred on him by Zeno in 483. I conjecture that Zeno had appointed him mag. utriusque militiae in Italy before he set out (cp. above,
p422 ), and that this was confirmed by Anastasius.

} Though the old Roman troops and their organisation disappeared, it has been shown that the military arrangements were based in many respects on practices which had existed in Italy under Imperial rule.

The various disabilities of the Ostrogoths which have been described depended on the fact that they were not Roman

p457
citizens. They, like the Germans settled by Odovacar, had legally the same status as mercenaries or foreign travellers or hostages who dwelled in Roman territory, but might at any time return to their homes beyond the Roman frontier. The laws which applied only to Roman citizens, for instance those relating to marriage and inheritance, did not apply to them. But what may be called the ius commune , laws pertaining to criminal matters and to the general intercourse of life, applied to all foreigners who happened to be sojourning in Roman territory; and thus the Edict of Theoderic, which is based on Roman law, is addressed to Goths and Romans alike. The status of the Goths reminds us of a fundamental restriction of Theoderic's power. He could not turn a Goth into a Roman; he could not confer Roman citizen ­ship; that power was reserved to the Emperor.

Their quality, as foreign soldiers, determined the character of the courts in which the Ostrogoths were judged. The Roman rule was that the soldier must be tried by a military court, and military courts were instituted for the Goths. But here Theoderic interfered in a serious way with the rights of the Italians. All processes between Romans and Goths, to whichever race the accuser belonged, were brought before these military courts. A Roman lawyer was always present as an assessor, but probably no feature of the Gothic government was so unpopular as this. Like the Emperor, Theoderic had a supreme royal court, which could withdraw any case from a lower court or cancel its decision, and this tribunal seems to have been more active than the corresponding court of the Emperor. It is indeed in the domain of justice, in contrast with the domain of legislation, that the German kings in Italy sharply asserted their actual authority.

Besides being Master of Soldiers in regard to the Ostrogothic host, Theoderic was likewise the king of the people. He did not style himself rex Gotorum ; like Odovacar, he adopted the simple title of rex . This indefinite style was hardly due to the circumstance that the foreign settlers in Italy were not all Ostrogoths, that the remnant of Odovacar's Germans, and notably the Rugians,\texthebrew{⁠}\footnote{Procopius, \emph{B. G.} III.2.

} acknowledged his kingship. It was perhaps intended also to express his actual, as distinguished from his

p458
constitutional, relation to the Roman population. While the Roman citizens were constitutionally the subjects of the Emperor, of whom the Patrician Theoderic was himself a subject and official, they were actually in the hands of Theoderic, who was their real ruler. To designate this extra-constitutional relation, the word rex , which had no place in the constitutional vocabulary of Rome, was appropriate enough. It served the double purpose of expressing his regular relation to his German subjects, and his irregular relation, his quasi-kingship , to the Romans of Italy.\footnote{In an inscription commemorating his draining of the Pomptine marshes he is given the Imperial title of semper Augustus. He is there styled d. n. gloriosissimus adque inclytus rex, victor ac triumfator semper Aug., bono reipublicae natus, custos libertatis et propagator Romani nominis, domitor gentium, CIL X.6850 . In one inscription he is mentioned along with an Emperor, probably Anastasius: salvis domino . . . Augusto et gloriossimo rege Theoderico, CIL VI.1794 . He never wore the diadem.

}

The continuity of the administration of Odovacar with that of Theoderic was facilitated by the fact that some of the Roman ministers of Odovacar passed into the service of the Ostrogothic ruler, and probably the mass of subordinate officials remained unchanged. For instance, the first Praetorian Prefect of Italy under Theoderic was Liberius ( A.D. 493\texthebrew{‑}500), who had been one of the trusted ministers of Odovacar. Cassiodorus \texthebrew{—} father of the famous Cassiodorus whose writings are our chief authority for Theoderic's reign, \texthebrew{—} who had held both the great financial offices under Odovacar, continued to serve under Theoderic, and in the early years of the sixth century became Praetorian Prefect.\footnote{The thoroughly Roman character of the Italian kingdom is clear. There are one or two points in which Germanic influence has been suspected. (1) The saiones were marshals or messengers whom the king employed to intimate his commands. They might summon the Gothic soldiers to arms or recall a Roman official to a sense of duty. The office of saio may be a German institution, or there may be nothing German about it but the name. The functions of these officials correspond to those of the agentes in rebus, who also existed in Italy at this time, though, as Mommsen has shown, they were called comitiaci. They may have served as a model for the institution of the saiones. (2) By an Imperial law of A.D. 393 any person who considered his personal safety in danger might apply for special protection, tuitio, and a judge was bound to assign a civil officer (apparitor) to protect him. Tuitio is very prominent in Ostrogothic Italy; it was granted by the king himself, and was one of the methods by which he preserved peace and order among the two races. The quickening of this Roman custom, and its special association with the king, may have been partly due to the Germanic idea of the king's duty of protection (munt).

}

The constitutional system of administration which Theoderic accepted and observed was not a necessity to which he reluctantly

p459
or lukewarmly yielded. It was a system in which he seems to have been a convinced believer, and he threw his whole heart and best energies into working it. His object was to civilise his own people in the environment of Roman civilisation ( civilitas ). But he made no premature attempt to draw the two classes of his subjects closer, by breaking down lines of division. They were divided by religion and by legal status. So far as religion was concerned, the king was consistently tolerant, unlike the rulers of the Vandals and the Franks. His principle was: ``We cannot impose religion because no one can be compelled to believe against his will'' \texthebrew{—} a maxim which might well have been pondered on by Roman Emperors.\texthebrew{⁠}\footnote{Religionem imperare non possumus, quia nemo cogitur ut credat invitus.

Cassiodorus, II.27

} So extreme was his repugnance to influencing the creed of his fellow-creatures that an anecdote was invented that he put to death a Catholic deacon for embra­cing Arianism to please him. If there is any foundation for the story, there must have been other circumstances; but it is good evidence as to his religious attitude; if it was entirely invented, it proves his reputation.\footnote{To the Jews also he extended toleration and protection. Cassiodorus, \emph{loc. cit.}

}

And just as he accepted the duality of religion, he accepted the dual system by which Goths and Romans lived side by side as two distinct and separate peoples. He made no efforts to bring about fusion, his only aim was that the two nations should live together in amity. But little love was lost between them. The rude German barbarians despised the civilised Italians, and the Ostrogothic kingdom was overthrown before fusion could begin; but the development in Visigothic Spain, under similar conditions, makes it probable that fusion would have ensued, if the Ostrogothic power had endured. It says much for Theoderic's authority and tact that he was able to hold an equal balance between the two peoples, and to attain so nearly in practice to the difficult ideal which he set before him

After his death the concealed impatience of the Goths under his philo-Roman policy was soon to burst out and hurry them to disaster.

Although he aimed at maintaining peaceful relations with the Emperor throughout his long reign, this concord was

p460
threatened more than once, and there were even actual hostilities. A campaign which Theoderic undertook against the Gepids, in order to recover Sirmium and adjacent districts of the Prefecture of Italy which this people had occupied, led to a collision with the Imperial troops ( A.D. 504\texthebrew{‑}505). The events are obscure.\texthebrew{⁠}\footnote{The sources are Ennodius, \emph{Pan. Theod.} 277\texthebrew{‑}280; Cassiodorus, \emph{Chron., sub} 504 ; Marcellinus, \emph{Chron., sub} 505;

Jordanes, \emph{Get.} 300\texthebrew{‑}301 ; Cassiodorus,

\emph{Var.} VIII.10.4 ; \emph{Or.} p473. There are difficulties in reconciling them. Cp. Hodgkin, III.438 \emph{sqq.} ; Schmidt, \emph{Gesch. der deutschen Stämme}, I.310.

} It would seem that the Gepids yielded with little resistance, in consequence of internal dissensions. But the expedition which Theoderic sent against them aroused the suspicions of Anastasius. At this time the central provinces of the Balkan peninsula were exposed to the depredations of a Hun, named Mundo, who had organised a band of brigands. The government sent the Master of Soldiers, Sabinian, to capture him, and Sabinian was supported by a formidable force of allied Bulgarians. Mundo appealed for help to the Ostrogothic general Pitzias, who was engaged in completing the occupation of the territory which he had won from the Gepids. Our informants do not explain why he should have made the brigand's cause his own, or regarded Sabinian's movements as a threat to the Goths; but he marched into Dacia and won a decisive victory over the Bulgarians. Mundo also inflicted a severe defeat on Sabinian at Horrea Margi.\texthebrew{⁠}\footnote{The place is given by Marcellinus, who says nothing of Ostrogoths or Bulgarians, and by Jordanes, who does not mention Bulgarians. From Ennodius one might infer that the battle was fought in two sections; he passes lightly over Sabiniani ducis abitionem turpissimam.

} The key to this episode probably is that Anastasius viewed with alarm the Gothic occupation of the important frontier town of Sirmium; he preferred that it should be in the hands of the Gepids than in those of his viceroy.\texthebrew{⁠}\footnote{The words of Ennodius are important: per foederati Mundonis adtrectationem Graecia est professa discordiam, secum Bulgares suos in tutela deducendo. It is to be observed that Mundo is described as an ally of the Ostrogoths. We are told nothing of his subsequent fortunes.

} After the defeat of Sabinian, he must have acquiesced in Theoderic's restoration of the Prefecture of Italy to its old limits, for no further hostilities followed.\footnote{Ennodius, \emph{loc. cit.}, ad limitem suum Romana regna remearunt. For the organisation of Sirmian Pannonia by Theoderic see Cass. \emph{Var.} III. 23 , 24 .

}

These operations in the region of the Save were probably connected with an attempt to make his authority felt in the Pannonian province. Of the conditions in Noricum and Pannonia

p461
at this time we have no clear idea. But we know that about the year 507 Theoderic settled a portion of the Alamannic people in Pannonia, perhaps in Savia. The remnant of this people, after their defeat by Clovis (perhaps in A.D. 495), had wandered southward into Raetia to escape the sword or the yoke of the victor. Clovis requested Theoderic to surrender them, and we possess Theoderic's reply. He deprecated the Frank king's desire to push his victory further. ``Hear the counsel,'' he wrote, ``of one who is experienced in such matters. Those wars of mine have been profitable, the ending of which has been guided by moderation.'' He took the Alamanni under his protection and gave them a home within the borders of his kingdom.\footnote{See Cassiodorus, \emph{Var.} II.41 and cp. III.50 , where we see the Alamanni, on their way from Raetia to Pannonia, passing through Noricum. The date of both these letters is 507. Also Ennodius, \emph{Paneg.} c. xv (this work belongs to the same year). Cp. Mommsen, \emph{Preface} to Cassiod. pp. xxxiii\texthebrew{‑}xxxiv; Dahn, \emph{Kön. der Germ.} IX.1 p64. Parts of Noricum and Raetia were occupied about the year 500 by Marcomanni and Quadi coming from Bohemia, driven westward by the Slavs; they were now known under the name of Bajuvarii, Bavarians (cp. Jordanes, \emph{Get.} 280 ). Thus was founded Bavaria.

}

In his relations with foreign powers, Theoderic acted as an independent sovran. The four chief powers with which he had to reckon were the Visigoths, the Burgundians, the Franks, and the Vandals. It was natural that he should look for special co-operation from the Visigoths, who were a kindred folk. But his policy at first was not to draw the Visigoths into a close intimate alliance, which might seem a threat to the other powers. He sought to form bonds of friendship with all the reigning houses, by means of matrimonial alliances. If he wedded one of his daughters to the Visigothic king, Alaric II, the other married Sigismund ( A.D. 494), who became king of the Burgundians after his father Gundobad's death. Theoderic himself took as his second wife a Frankish princess, sister of Clovis. And his own sister married Thrasamund, king of the Vandals ( A.D. 500). Thus he formed close ties with all the chief powers of the West.\texthebrew{⁠}\footnote{His niece married Hermanfrid, king of the Thuringians. He adopted as a son the king of the Heruls

(Cassiodorus, \emph{Var.} IV.2) . He gave Lilybaeum to his sister Amalafrida when she married Thrasamund, acc. to Procopius, \emph{B. V.} I.8, 11, whose statement is illustrated by an inscription on a boundary stone near that town marking the fines inter Vandalos et Gothos, CIL X.7232 . Amalfrida Theodenanada (Dessau, 89990) seems to be a different person from Theoderic's sister, perhaps her daughter.

} One object of this policy was doubtless to maintain the existing order of things, to preserve peace in western Europe,

p462
and secure Italy against attack. But we can hardly be wrong in thinking that it was also the purpose of Theoderic to secure his own position in Italy, in relation to the Imperial power. He could hardly fail to foresee that the day might come when Anastasius or one of his successors might decide to bring Italy under his immediate government or to deal with himself as Zeno had dealt with Odovacar. To meet such a danger, it would be much to have behind him the support of the western powers. As the centre and head of the system, linking together the German royalties, he would be in a far stronger position in regard to his sovran at Constantinople than Odovacar had been standing alone.

The family alliances of Theoderic did not avail to hinder war. He could not avert the inevitable struggle between the Franks and the Visigoths in Gaul. No moment in his reign caused him perhaps more anxiety than when Clovis declared war upon Alaric. Theoderic did what he could. We have the three letters which he wrote at this crisis to Alaric, to Gundobad, and to Clovis himself.\texthebrew{⁠}\footnote{Cass. \emph{Var.} III.1, 2, 4. He also wrote a circular letter to the kings of the Thuringians, Heruls, and Varni, \emph{ib.} 3.

} It was in vain. Theoderic promised armed help to his son-in\texthebrew{‑}law. But for some reason he was unable to render it. It would seem that he had calculated that the Burgundians would not side with the Franks, and that they cut him off so that he could not reach Aquitaine in time to intervene in the struggle. On the field of Vouillé (near Poictiers) the Visigothic king fell and Aquitaine was annexed to the dominion of the Franks ( A.D. 507). But in the following years the generals of Theoderic conducted campaigns in Gaul. They succeeded in rescuing Arles and in saving Narbonensis for the Visigothic kingdom. They wrested Provence from Burgundy and annexed it to Italy. At the same time the personal power of Theoderic received another extension. The heir of Alaric was a child, and the government of his realm was consigned to Theoderic, who was his grandfather and most power ­ful protector. For the rest of his life Theoderic ruled Spain and Narbonensis. Thus no inconsiderable part of the western section of the old Roman Empire was under his sway: Spain, Narbonensis, and Provence, Italy and Sicily, the two provinces of Raetia, Noricum, part of Pannonia, and Dalmatia.

p463
Thus the war in Gaul involved Theoderic, in spite of his relations to the royal houses, in hostilities against both the Franks and the Burgundians. The Burgundian alliance does not seem to have led to any close intimacy. Gundobad remained an Arian till his death ( A.D. 516), but he took good care to remain on friendly terms with Anastasius. His son Sigismund, Theoderic's son-in\texthebrew{‑}law, who succeeded him, had been converted to Catholicism\texthebrew{⁠}\footnote{Avitus, \emph{C. Arianos}, p2. The letters and works of Avitus, and the \emph{Vita} of Caesarius, bishop of Arles, throw some general light on the history of Burgundy during the first quarter of the sixth century.

} by Avitus, the bishop of Vienne, and appears to have been completely in the hands of Avitus and the Catholic clergy. He looked to the Emperor as his over ­lord, and addressed him in almost servile terms.\texthebrew{⁠}\footnote{Avitus, \emph{Epp.} 93 and 94. He writes for instance: vester quidem est populus meus et plus me servire vobis quam illi praeesse delectat.

} Theoderic was alarmed at the prospect of political intimacy between Burgundy and Constantinople, and he would not allow Sigismund's messengers to travel through Italy to the East.\texthebrew{⁠}\footnote{\emph{Ep.} 94.

} The strained relations between the courts were shown by the circumstance that the consul ­ship of Eutharic was not accepted in Burgundy as the date of A.D. 519.\texthebrew{⁠}\footnote{CIL XII.1500 .

} Theoderic probably placed his hopes in his grandson Sigeric, who, though he had been converted to the Catholic creed, was not on good terms with his father. His mother was dead, and Sigismund had taken a second wife. We know nothing authentic of the breach between father and son, but the end was that Sigeric was put to death by his father's orders ( A.D. 522).\texthebrew{⁠}\footnote{Marius Avent., \emph{sub a.} A legend grew up that his stepmother, whom he had insulted, accused him of treason,

Gregory of Tours, \emph{H. F.} III.5 .

} Theoderic prepared for war to avenge his grandson, but it was the Franks, not the Ostrogoths, who were to punish Sigismund. It was not to their mind that Theoderic should have a free hand in Burgundy, and moving more quickly, they captured Sigismund and his family and subdued a part of the kingdom. An Ostrogothic force arrived afterwards and annexed the district between the Isère and the Durance to Theoderic's realm ( A.D. 523).\footnote{Cassiodorus, \emph{Var.} VIII.10.8 . Cp. \emph{Vita Caesarii}, I.60.

}

The war between the Franks and Visigoths seems to have led to friction between Theoderic and the Emperor. In that struggle Clovis posed as the champion of Catholic orthodoxy, going forth to drive the Arian heresy from the confines of Gaul,

p464 and all the sympathies of the Gallo-Roman Church were with the Franks. The Emperor afterwards showed his approbation of the Merovingian king by conferring upon him the honorary consul ­ship.\texthebrew{⁠}\footnote{Gregory of Tours, \emph{Hist. Fr.} II.38

ab Anastasio imperatore codecillos de consulato accepit . . . et ab ea die tamquam consul aut augustus est vocitatus. The expression tamquam consul seems to be equivalent here to ex consule, the official title of honorary consuls (augustus seems to be a mistake of Gregory; if Clovis did assume it, it certainly was not conferred on him). In the \emph{Lex Salica} (ed. Behrend, p125) Clovis is called proconsul. Mommsen has suggested that this is a mistake for praecelsus. It is possible that Gregory has confused consul­ship with proconsul­ship, and that Anastasius really conferred an honorary proconsul­ship. This, perhaps, is less likely.

} Theoderic meanwhile was supporting the Visigoths, and we may conjecture that his Gallic policy was disapproved by Anastasius, who ( A.D. 508) despatched a squadron of a hundred ships to ravage the coasts of Apulia.\footnote{Marcellinus, \emph{Chron., sub a.} The ships carried 8000 soldiers.

}

The ecclesiastical relations between Rome and Constantinople affected the political situation in Italy, more or less, throughout the reign of Theoderic.\texthebrew{⁠}\footnote{This has been best elucidated by Sundwall, \emph{op. cit.}, which I have used much in what follows. Pfeilschifter's \emph{Theod. und die Kathol. Kirche} is indispensable.

} This was partly due to the fact that the great Roman families were now all Christian, and many of the senators held strong opinions on the subject of the schism which the Henotikon of Zeno had provoked. Festus had taken advantage of his political mission to Constantinople in A.D. 497 to attempt to heal the schism. He told the Emperor that he had hopes of indu­cing the º Pope Anastasius to sign the Henotikon. But when he returned to Italy the Pope was dead.\texthebrew{⁠}\footnote{Died November 19, 498.

} Festus, however, only represented the opinion of part of the Senate. There was a marked division in the views of the senators, of whom an influential section were opposed to any compromise on the theological question. This difference of opinion led to a bitter struggle over the election of a new Pope. Two men were elected on the same day (November 22, A.D. 498), Laurentius, the candidate of Festus and the party of reconciliation, and Symmachus, supported by the orthodox, who were prepared to make no concessions. Two rival Popes were enthroned in Rome, each upheld by strong and determined partisans, and for years the city was disturbed by sanguinary tumults.\texthebrew{⁠}\footnote{Cp. Theodore Lector, II.17. The most prominent supporter of Symmachus was Faustus.

} An appeal was made to Theoderic to decide between

p465
the two claimants. It is a remarkable episode in the history of the Church that such a question should be referred to an Arian. As the tranquillity of Italy was in peril, the ruler could not stand aloof, and he consented to give a decision. He was conscious of his obligations to Festus, but the clergy, especially the clergy of North Italy, were as a body adherents of Symmachus, and it was in favour of Symmachus that Theoderic decided ( A.D. 499).

But the matter was not finally settled by the king's arbitrament. The behaviour of Symmachus was aggressive and uncompromising,\texthebrew{⁠}\footnote{He addressed a letter to the Emperor, to which the Emperor after some delay replied by a manifesto, and Symmachus rejoined in a rather violent Apologetic, which will be found in Thiel, \emph{Epp. R. Pont.} p700 \emph{sqq.}

} and charges were brought against him, which were submitted to a synod held two years later. He was acquitted and recognised as the legitimate bishop of Rome,\texthebrew{⁠}\footnote{The date of the Synodus Palmaris was probably early summer 502: Sundwall, \emph{op. cit.} p206. Its enemies called it the Synodus absolutionis incongruae, and it was defended in a pamphlet by Ennodius, \emph{Libellus Romano Synodo}, 287 \emph{sqq.}

} but his conduct alienated Theoderic, and no steps were taken to remove or suppress Laurentius, who continued to maintain his papal pretensions at Rome for the next few years. But in A.D. 505 there was a revulsion of feeling. The adherents of Laurentius were chiefly men who considered the maintenance of close relations with the Imperial court a fundamental interest of Italy. But their Italian sentiments were aroused by the incidents connected with Sirmium. Here their sympathy was with Theoderic, and it seems highly probable that the hostilities between the troops of Anastasius and those of his viceroy in Dacia were partly at least responsible for a general change of opinion in favour of Symmachus.\texthebrew{⁠}\footnote{Sundwall, p212.

} This made the position of Laurentius impossible, and he was obliged to retire before the end of A.D. 506.

Thus ten years after the settlement which had been arranged between Theoderic and the Emperor, the policy of the Gothic ruler had brought it about that Italy presented a united front, and the influence of Constantinople now reached its lowest point. The Church and the Senate were united against the East on the ecclesiastical question. In the spring of A.D. 507 Ennodius, one of the leading dignitaries of the Italian Church,

p466
pronounced his Panegyric on the Arian king.\texthebrew{⁠}\footnote{Sundwall, pp42\texthebrew{‑}43, has fixed the date.

} But this situation was only momentary. Hitherto Theoderic had followed the example of Odovacar in basing his government on close co-operation with the great Roman families, members of which were chosen to fill the highest civil posts, especially the Prefecture of Rome and the Praetorian Prefecture of Italy. But from this time forward we can mark the beginning of a new policy. Probus Faustus Niger, who had been the leading champion of Symmachus in the conflict over the Papal throne, is indeed Prefect of Italy from A.D. 507\texthebrew{‑}512, but we find new men, who do not belong to the senatorial circle, appointed Prefects of the City.\texthebrew{⁠}\footnote{Agapitus, 507\texthebrew{‑}509, followed by Artemidorus, and then Argolicus. Sundwall, p215.

} It was apparently the aim of Theoderic to diminish his dependence on the senate. At Ravenna he had gathered round him a circle of other ministers of provincial origin who were devoted to his interests. To such were entrusted the financial offices; from such were generally selected the Master of Offices and the Quaestor.

Of Theoderic's acts and policy throughout the rest of the reign of Anastasius we know very little. He looked with favour on the vain attempts of Vitalian to restore the unity of the Church, and was ready to co-operate with Pope Hormisdas to bring it about.\texthebrew{⁠}\footnote{Hormisdas succeeded Symmachus in 514. It may be mentioned here that it was in these two pontificates that the Scythian monk Dionysius Exiguus worked at Rome under the auspices of the Roman Church, translating into Latin the ``Apostolical Canons'' and the canons of the great Councils. This collection, to which he added the canons the Council of Sardica and the African Councils, became authoritative, and protected the chancery of the church of Rome (where Greek was little known) from being imposed upon by forgeries. Dionysius also established the custom of dating events from the Nativity, and introduced the cycle of 532 years (= 28 solar cycle \ensuremath{\times} 19 lunar cycle), invented by Victorius of Aquitaine for the computation of Easter. For his works see \emph{P. L.} 67. Maassen, \emph{Gesch. der Quellen u. der Litt. des Canonischen Rechtes}, vol. I.

} It would be a mistake to read into his \emph{Edict}, which was probably issued in A.D. 512, any design of diminishing the power or prestige of the senatorial classes.\texthebrew{⁠}\footnote{See Gaudenzi's article in \emph{Zeitschrift der Savigny-Stiftung}, VII.1, 1886.

} Throughout the provinces Romans and Goths alike were constantly attempting to encroach upon the lands of their neighbours; many acts of violence occurred;\texthebrew{⁠}\footnote{Cassiodorus, \emph{Var.} VIII.27 .

Lécrivain, \emph{Le Sénat}, 178 \emph{sqq.}

} and the principal object of the Edict seems to have been to put an end to these illegalities and disorders.

p467
The relations between Ravenna and Constantinople were never cordial. Italians who were banished from Italy by Theoderic were treated with marked favour at the Byzantine court, and received posts in the Imperial service. We learn this fact from Priscian, the distinguished African grammarian, who, leaving the realm of the Vandals, had settled in Constantinople and sympathised with the national feeling of the Italians against Gothic rule.\texthebrew{⁠}\footnote{Priscian, \emph{Paneg. in Anastas.} vv. 242 \emph{sqq.} V. 265 expresses the hope that Gothic rule will not last long:

utraque Roma tibi nam spero pareat uni.

Priscian was a friend of Symmachus, to whom he dedicated three minor works (Sandys, \emph{Hist. Class. Scholarship}, I p258). Cp. Usener, \emph{Anecd. Holders}, 26.

} The presence of these exiles, who, we may be certain, maintained a frequent correspondence with their friends in Rome, is a circumstance which must not be lost sight of in studying the relations of Theoderic with the Emperor and with the Roman Senate.

It is remarkable that Theoderic, who was educated at Constantinople and was imbued with sincere admiration for Greek and Roman civilisation, was illiterate. It is recorded that he was unable to write his own name. He caused a gold stencil plate to be pierced with the four letters legi (I have read), so that he could sign documents by drawing a pen through the holes.\footnote{Anon. Val. 79 , where, however, it seems to be implied that at the end of ten years he had learned to write the four letters. If there is any truth in this we must suppose that the letters were arranged in an elaborate monogram, which would explain the use of a stencil plate, without having recourse to the inference of the chronicler that Theoderic was unable to write at all. The same device was adopted by the illiterate Emperor Justin, according to Procopius

(\emph{H. A.} 6) . In Anon. Val. \emph{loc. cit.}, quattuor litteras legi habentem has manuscript authority and is read by the latest editor, Cessi; cp. the passage of Procopius (\textgreek{γραμμάτων τεττάρων, \texthebrew{ἄ}περ \texthebrew{ἀ}ναγν\texthebrew{ῶ}ναι τ\texthebrew{ῇ} Λατίνων φων\texthebrew{ῇ} δύναται}). The old reading was regis, and Valesius inserted Theod. after habentem. But the signature Theod or Theodoricus has more than four letters. A. J. Evans (\emph{Antiquarian Researches in Illyricum}, Part III pp22, 23), with the text of Valesius before him, thought that a monument of Theodoricus is meant such as is found on coins and on an engraved gem, apparently a seal of an official of Theoderic.

}

Theoderic chose Ravenna, the city of Honorius and Placidia and Valentinian, as his capital. The Emperors who reigned in the days of Ricimer had seldom resided in the palace of the Laurelwood (Lauretum), but Odovacar had made it his home. Theoderic built a new palace in another part of the city, and erected beside it a new church dedicated to St. Martin, in which his Arian Goths worshipped. Of the palace only a wall, if anything,

p468
remains. But the church, one of the fine works of the Ravennate school of architecture, still stands. It was afterwards dedicated to St. Apollinaris, and is known as San t' º Apollinare Nuovo.\texthebrew{⁠}\footnote{See Rivoira, \emph{Lombardic Architecture}, I.40 \emph{sqq.} The Corinthian capitals in the nave and the ambo are of Byzantine workman­ship. The palace of Theoderic is represented in the mosaics. The round campanile belongs to the ninth century (\emph{ib.} 45).

} Of the mosaic pictures which adorn the nave only those which are aloft near the roof, \texthebrew{—} scriptural scenes, \texthebrew{—} and the figures between the windows, belong to Theoderic's reign; the decoration of the church was not completed till thirty years after his death.\texthebrew{⁠}\footnote{Cp. Dalton, \emph{Byz. Art}, 350. See below,
vol. II p285 .

} We may assume that it was he who built the Arian baptistery which survives as S. Maria in Cosmedin. It is interesting to learn that near the State factories at the port of Classis he drained a portion of the marshes and planted an orchard.\footnote{Rex Theodericus . . . fabricis suis amoena coniugens, sterili palude siccata, hos hortos suavi pomorum fecunditate ditavit.
CIL XI.10 ;

Jordanes, \emph{Get.} 151 .

}

Ravenna has another famous memorial of Theoderic, the round mausoleum which he built for himself. It was ``covered by a cupola consisting of a single piece of Istrian limestone, the circumference of which is provided with twelve handles, intended, without doubt, to lift by means of ropes and drop into its place this wonder ­ful inverted basin.''\texthebrew{⁠}\footnote{Rivoira, \emph{ib.} 54. He remarks on the ability displayed in the construction of the building and its excellent proportions. Internally it is shaped like a cross with equal arms, must have been inspired by some Roman sepulchral edifice. Rivoira acknowledges the impulse given by Theoderic to art. Many public works were carried out by his direction, \emph{e.g.} the restoration of the aqueducts of Ravenna, of the walls of Rome, and of the Theatre of Pompey; the construction of baths at Verona. Literature as well as art flourished under Theoderic. Cassiodorus, Boethius, and Ennodius were the most distinguished writers, but they do not exhaust the list.

} We must suppose that the body of the king once lay in the sepulchre which was designed to receive it. What befell it is a matter for conjecture; we only know that three hundred years later the tomb had long been empty.\footnote{This is recorded by Agnellus (who wrote in the ninth century), \emph{Lib. Pont.} (in \emph{Scr. r. Lang.}) p304. The sarcophagus was a porphyry urn.

}

Under the rule of Theoderic, Italy is said to have enjoyed peace, prosperity, and plenty, such as she had not known for many a long year. His success was due not only to his political and military capacity, but also to his rigorous though humane ideal of justice. The praises of Italian panegyrists are borne

p469
by the verdict of one who was afterwards employed in active hostility against Theoderic's successors. If a Ravennate chronicler asserts that the king ``did nothing wrong'' ( nihil perperam gessit ),\texthebrew{⁠}\footnote{Anon. Val. 60 .

} the historian Procopius makes a statement, hardly less unqualified, in regard to the justice of the administration, and dwells on the deserved devotion which his subjects entertained towards him.\texthebrew{⁠}\footnote{The encomium of Procopius will be found in \emph{B. G.} I.1.

} The peace and plenty of his times are illustrated with vivid hyperboles in an Italian chronicle.\texthebrew{⁠}\footnote{Anon. Val. 72, 73 . The laudatory notices in this chronicler were perhaps inspired by the Panegyric of Ennodius. See

Dumoulin's article in \emph{Revue historique}, 1902 .

} ``Merchants from divers provinces used to throng to him. For so perfect was the public order that if a man wished to leave his silver or gold in his field, it was respected as much as if it were within the walls of a town. This was shown by the fact that he built no new gates for any town in all Italy, nor were the gates of any town ever closed. Any one could go about his business at any hour of the night just as if it were day. In his time sixty modii of wheat cost a solidus, and thirty amphorae of wine were sold for the same price.''\texthebrew{⁠}\footnote{Thus a modius of wheat (= about 2 gallons) cost 2½d., or a bushel cost 10d. and a quarter 6s. 8d. As the Roman amphora was nearly 6 gallons, a gallon of wine cost less than 1d. The price of wheat in Julian's time was between 1/10th and 1/15th of a solidus for a modius (\emph{Misopogon}, 369), \emph{i.e.} about a shilling. In the sixth century 16 artabae of Egyptian wheat were sold for two solidi (\emph{Pap. Cairo}, I.67062). An artaba is generally reckoned = 3\texthebrew{1/3} modii, but at this time it was equated with 3 modii (see tables in \emph{Pap. Cairo}, II.67138), so that the price was 1/24th solidus = about 6d. a modius. On the other hand, in the accounts of Ammonius (\emph{ib.}) we find 25 art. sold for 99½ keratia, or 1 art. for 1/6th solidus and 1 modius for 8\texthebrew{1/3} d. (or somewhat more if the solidus was equated with 22 instead of 24 keratia).

\texthebrew{◂} previous

next \texthebrew{▸}Images with borders lead to more information.

The thicker the border, the more information.

(Details here.)

UP TO:

J. B. Bury
Homepage

Latin \& Greek

Texts

LacusCurtius

Home} If this cheapness of provisions was normal, it would be one of the most convincing signs of the prosperity of Italy under Theoderic's government. But notwithstanding the improvement in their material conditions and in their general security, we can hardly believe that the Italians, with the barbarians settled in their midst, regarded themselves as steeped in felicity.

\xchapter{Chapter XIII, Appendix}

There are considerable difficulties as to the succession of the Praet. Prefects in this reign. The evidence will be found collected in Borghesi, \emph{Les Préfets de Prétoire}, I .370 \emph{sqq.}, but his results are not clear or satisfactory. The dates in \emph{C.J.} are our main guide. The following seem to be fairly certain: Matronianus, A.D. 491, July (\emph{C.J.} VII .39.4 ;
I .22.6 ); Hierius, A.D. 494 (John Mal. XVI p392)\texthebrew{‑}496, Feb. 13 (\emph{C.J.} VI .21.16); Euphemius, A.D. 496, April 1\texthebrew{‑}Aug. 21 (\emph{ib.} X .16.13;
X .19.9); Polycarpus, A.D. 498, April 1 (\emph{ib.} V .30.4); Constantine, A.D. 502, Feb. 15\texthebrew{‑}July 21 (\emph{ib.} III .13.7\texthebrew{‑}6, 20.18); Appion, A.D. 503 (John Mal. XVI p398); Leontius, A.D. 503\texthebrew{‑}504 (John Lyd. III .17) ; Constantine again, A.D. 505 (\emph{C.J.} II .7.22 , but the month Iul. is wrong; Krüger suggests Ian. ); Eustathius, A.D. 505, April 19\texthebrew{‑}506, Nov. 20 (\emph{ib.} I .4.19;
II .7.23); Zoticus, A.D. 511\texthebrew{‑}512 (Cyrillus, \emph{Vita S. Sabae}, pp290, 294; this agrees with the chronological indications in John Lyd. III .27 ; from whom we also learn that Zoticus held office for little more than a year); Sergius, A.D. 517, April 1\texthebrew{‑}Dec. 1 (\emph{C.J.} V .27.6 ;
II .7.24 ).

The Prefects of uncertain date are Armenius, Arcadius, Leontius (\emph{ib.} XII .50.23 ;
XII .37.7 ;
VII .39.6 ), and Marinus. As to Leontius, he held office after 500 (cp. \emph{ib.} VII .39.5 , and John Lyd. III .17 ). For the Prefecture of Marinus we have the limits 498 (John Lyd. III .36) and 515, in which year he was ex\texthebrew{‑}Pr. Pr. (John Mal. XVI pp403, 405, 407). He was influential with Anastasius in the Prefecture of Zoticus (Cyrillus, \emph{loc. cit.}), and it is to be noted that Zacharias of Mytilene ( VII .9), speaking of him as the Emperor's friend and confidant, describes him as a chartularius ( A.D. 511). The people of Constantinople held him as partly responsible for the ecclesiastical measures which caused the riot of Nov. 512, and his house was burnt down (Marcell. \emph{Chron.}, \emph{sub a.}). On the whole, I would conjecture that he became Prefect in that year, having succeeded Zoticus. It does not follow from John Lyd. \emph{loc. cit.} (as Borghesi supposes) that he immediately succeeded Polycarpus. In the latter part of his reign, Anastasius appointed only Scholastici ( \textgreek{\texthebrew{ῥ}ήτορες , λογικοί} ) to the Prefecture (John Lyd. III .50 ; Priscian, \emph{Pan.} 246\texthebrew{‑}251), in accordance with the old tradition of the civil service. For the training of the scholasticus cp. Macarius, \emph{Hom.} 15.42, in Migne, \emph{P.G.} XXXIV .604. \texthebrew{—} Marinus is meant by the \textgreek{Μαριανός} who is mentioned in Justinian, \emph{Nov.} 96 § 15, as is evident from the context. He was Praetorian Prefect again under Justin A.D. 519 (\emph{C.J.} V .27.7 ;
II .7.25 ). \texthebrew{—} There is a slight difficulty about Appion, though John Malalas (source: Eustathius of Epiphania) says expressly that the patrician

p471 Appion was appointed \textgreek{\texthebrew{ἔ}παρχος πραιτωρίων \texthebrew{ἀ}νατολ\texthebrew{ῆ}ς} and sent to the East on the outbreak of the Persian War. This seems to harmonise with Joshua Styl. LV p44, who states that Appion the hyparch was at Edessa in May 503. But it would be very strange for a Praet. Prefect to proceed himself to the seat of war to supervise the commissariat, and we should naturally take hyparch to mean the officer called prefect of the camp, \textgreek{\texthebrew{ὁ} το\texthebrew{ῦ} στρατοπέδου \texthebrew{ἔ}παρχος} (Procopius, \emph{B.V.} I .11), both here and \emph{ib.} LXX , where we learn that Calliopius became hyparch in May 504, º an office which he occupied till 506, \emph{ib.} XCIX . We cannot suppose Calliopius to have been Praet. Prefect, as the post was held by Constantine and Eustathius in 505\texthebrew{‑}506, and it is a little difficult to interpret hyparch differently in the two cases. But we have to take into consideration the statement of John Lyd. III .17 that Anastasius was ``moved with anger against Appion,'' \textgreek{\texthebrew{ἀ}νδρ\texthebrew{ὸ}ς \texthebrew{ἐ}ξοχωτάτου κα\texthebrew{ὶ} κοινωνήσαντος α\texthebrew{ὐ}τ\texthebrew{ῷ} τ\texthebrew{ῆ}ς βασιλείας \texthebrew{ὅ}τε Κωάδης \texthebrew{ὁ} Πέρσης \texthebrew{ἐ}φλέγμαινε, Λεοντίου τ\texthebrew{ὴ}ν \texthebrew{ἐ}παρχότητα διέποντος} . This seems to mean that Apion was Praet. Pref. at the outbreak of the Persian War, but fell into disfavour and was succeeded by Leontius, and establishes the Prefecture of Appion. I am inclined to think that Joshua's Appion was a different person.

A page or image on this site is in the public domain ONLY if its URL has a total of one * asterisk. If the URL has two ** asterisks, the item is copyright someone else, and used by permission or fair use. If the URL has none the item is \emph{©} Bill Thayer.

See

my copyright page

for details and contact information.

Page updated:
20 Jul 07

Accessibility

\xchapter{Chapter XIV}

Short URL for this page:

bit.ly/BURLAT14

The rulers of Constantinople would hardly have steered their section of the Empire with even such success as they achieved through the dangers which beset it in the fifth century, had it not been that from the reign of Arcadius to that of Anastasius their peaceful relations with the Sassanid kings of Persia were only twice interrupted by brief hostilities. The unusually long duration of this period of peace, notwithstanding the fact that the conditions in Armenia constantly supplied provocations or pretexts for war, was in a great measure due to the occupation of Persia with savage and dangerous enemies who threatened her north-eastern frontier, the Ephthalites or White Huns, but there was a contributory cause in the fact that the power of the Sassanid kings at this time was steadily declining. It is significant that when, at the end of the fifth century, a monarch arose who was able to hold his own against the encroachments of the Zoroastrian priesthood and the nobility, grave hostilities immediately ensued which were to last with few and uneasy intervals for a hundred and thirty years.

At the accession of Arcadius, Varahran IV was on the Persian throne, but was succeeded in A.D. 399 by Yezdegerd I. The policy of this sovran was favourable to his Christian subjects, who had been allowed to recover from the violent persecution which they had suffered at the hands of Sapor, the conqueror of Julian; and he was an object of veneration to Christian historians,\texthebrew{⁠}\footnote{Compare \emph{e.g.} Socrates, VII.8; \emph{Chron. Edess.} (ed. Guidi), p107. See Labourt, \emph{Le Christianisme dans l'empire perse}, 91\texthebrew{‑}93.

} while the Magi and the chroniclers of his own kingdom

p2
detested his name. After the death of Arcadius there were negotiations between the courts of Constantinople and Ctesiphon, but it is difficult to discover precisely what occurred. There is a record, which can hardly fail to have some foundation, that in his last illness Arcadius was fretted by the fear that the Persians might take advantage of his son's infancy to attack the Empire, and that he drew up a testament in which he requested the Great King to act as guardian of his son.\texthebrew{⁠}\footnote{Procopius, \emph{B. P.} I.2 ; Theophanes, A.M. 5900 (a notice evidently drawn from the same source as that in Michael Syrus, VIII.1. Haury's view (\emph{Zur Beurteilung des Procop. 21}) that Arcadius appointed Yezdegerd ``guardian'' in 402, when he crowned Theodosius, cannot be accepted. Agathias

(IV.26)

expresses scepticism about this statement of Procopius, and many modern writers (\emph{e.g.} Tillemont, Gibbon, Nöldeke) have rejected it. (See P. Sauerbrei, \emph{König Jazdegerd, der Sünder}, in \emph{Festschrift Albert v. Bamberg}, Gotha, 1905; on the other hand, Haury, \emph{B.Z.} XV.291 \emph{sqq.} Cp. Güterbock, \emph{Byzanz und Persien}, 28). But such a recommendation of a child heir to a foreign monarch is not without parallels. Heraclius, when he went forth against Persia, is said to have placed his son under the guardian­ship of the Chagan of the Avars. Kavad proposed that Justin I should adopt his son Chosroes (see below,

p79 ).

} There seems no reason not to accept this statement, provided we do not press the legal sense of guardian,\texthebrew{⁠}\footnote{Procopius uses the word \textgreek{\texthebrew{ἐ}πίτροπος} (= tutor), Theophanes (A.M. 5900), \textgreek{κουράτωρ.}} and take the act of Arcadius to have been simply a recommendation of Theodosius to the protection and goodwill of Yezdegerd. The communication of this request would naturally be entrusted to the embassy, which, according to the traditional etiquette, announced the accession of a new Emperor at the Persian court.\texthebrew{⁠}\footnote{According to Skylitzes

(Cedrenus, I.586)

Arcadius sent 1000 lbs. of gold to Yezdegerd. This is not improbable; the embassy announcing the Emperor's decease would in any case offer gifts.

} Yezdegerd took the wish of his ``brother'' as a compliment and declared that the enemies of Theodosius would have to deal with him.

Whatever be the truth about this record, which is not mentioned by contemporary writers,\texthebrew{⁠}\footnote{Sauerbrei (\emph{op. cit.}) seems to be right in his conclusion that the notice in Theophanes is not taken from Procopius but from a common source. If this is so, the record is not later than the fifth century. Skylitzes seems to have had access to this source or to an independent derivative.

} there is no doubt that there were transactions between the two governments at this juncture, and either a new treaty or some less formal arrangement seems to have been concluded, bearing chiefly on the position of Persian Christians and perhaps also on commerce. The Imperial Government employed the good offices of Maruthas, bishop of Martyropolis,\texthebrew{⁠}\footnote{Socrates, \emph{loc. cit.}

} who, partly on account of his medical

p3
knowledge, enjoyed much credit with Yezdegerd, to persuade the king to protect his Christian subjects. Yezdegerd inaugurated a new policy, and for the next twelve years the Christians of Persia possessed complete ecclesiastical freedom.\footnote{The important Council of Seleucia held in 410 was the immediate outcome of the new situation. It is stated in the Acts of this Council that Yezdegerd ordained that the churches destroyed by his predecessors should be rebuilt, that all who had been imprisoned for their faith should be set at liberty, and the clergy should be free to move about without fear, \emph{Synodicon orientale}, ed. Chapot, p254. See Labourt, \emph{op. cit.} 91 \emph{sqq.}

}

It is possible that at the same time the commercial relations between the two realms were under discussion. It was the policy of both powers alike to restrict the interchange of merchandise to a few places close to the frontier. Persian merchants never came to Constantinople, Roman merchants never went to Ctesiphon. The governments feared espionage under the guise of trade, and everything was done to discourage free intercourse between the two states. Before the treaty of Jovian, Nisibis was the only Roman town in which Persian merchants were allowed to trade.\texthebrew{⁠}\footnote{Peter Patric. \emph{fr.} 3 (\emph{Leg. Rom.} p4).

} After the loss of Nisibis, Callinicum seems to have become the Roman market for Persian merchandise, but we hear nothing of the new arrangements until the year 408\texthebrew{‑}409, when an Imperial edict was issued for the direction of the governors of the frontier provinces.\texthebrew{⁠}\footnote{\emph{C. J.} IV.63.4 . The motive of the restriction of trade to certain places is stated plainly: ne alieni regni, quod non convenit, scrutentur arcana. Artaxata was subsequently replaced by Dubios (Dovin) not far to the north-east; camp. Procopius, \emph{B. P.} II.25.

} From it we learn that the two governments had agreed that the Persian towns of Nisibis and Artaxata and the Imperial town of Callinicum should be the only places to which Persian and Roman traders might bring their wares and resort to transact business. Taken in connection with the fact that the two governments had been engaged in negotiations, this promulgation of the edict at this time suggests that if a new compact regarding commercial relations was not concluded, an old agreement, which may have been laxly executed, was confirmed.\footnote{Güterbock (\emph{op. cit.} 74\texthebrew{‑}75) refers the agreement to the treaty of 387, but why not to that of 363? The words of the edict are loca in quibus foederis tempore cum memorata natione nobis convenit. Sozomen makes the remarkable statement that the Persians prepared for war juncture, and then concluded a peace for 100 years (IX.4 \emph{ad init.}). It is curious that he should have confused the peace of 422 with the transactions of 408. Haury (\emph{loc. cit.} p294) suggests that there was actually a movement in Byzantium against the succession of Theodosius and that

(p4)
Yezdegerd threatened to intervene. It may be observed that the appointment of the Persian eunuch Antiochus to educate Theodosius had nothing to do with Yezdegerd.

}

p4
At the very end of Yezdegerd's reign the friendly understanding was clouded. All might have gone well if the Christian clergy had been content to be tolerated and to enjoy their religious liberty. But they engaged in an active campaign of proselytism and were so success ­ful in converting Persians to Christianity that the king became seriously alarmed.\texthebrew{⁠}\footnote{The incident which immediately provoked the persecution was the outrageous act of a priest who destroyed a fire-temple near his church. Theodoret, V.38; Labourt, \emph{op. cit.} 106 \emph{sq.}

} It was perfectly natural that he should not have been disposed to allow the Zoroastrian religion to be endangered by the propagation of a hostile creed. It is quite certain that if there had been fanatical Zoroastrians\texthebrew{⁠}\footnote{There were some old Zoroastrian communities in Cappadocia \texthebrew{—} settlers from Babylonia \texthebrew{—} in the time of the Achaemenids, which still existed in the fourth and fifth centuries (cp. Basil, \emph{Epp.} 258\texthebrew{‑}325); they were known as Magusaeans (\textgreek{Μαγουσα\texthebrew{ῖ}οι}). Strabo notices them, º

XV.3.15

\textgreek{\texthebrew{ἐ}ν δ\texthebrew{ὲ} τ\texthebrew{ῇ} Καππαδοκί\texthebrew{ᾲ}} (\textgreek{πολ\texthebrew{ὺ} γ\texthebrew{ὰ}ρ \texthebrew{ἐ}κε\texthebrew{ῖ} τ\texthebrew{ὸ} τ\texthebrew{ῶ}ν Μάγων φ\texthebrew{ῦ}λον, ο\texthebrew{ἳ} κα\texthebrew{ὶ} πύραιθοι καλο\texthebrew{ῦ}νται· πολλ\texthebrew{ὰ} δ\texthebrew{ὲ} κα\texthebrew{ὶ} τ\texthebrew{ῶ}ν Περσικ\texthebrew{ῶ}ν θε\texthebrew{ῶ}ν \texthebrew{ἱ}ερά}), \textgreek{κτλ}. See Cumont, \emph{Les Mystères de Mithra}, ed. 3, pp11, 12.

} in the Roman Empire and they had undertaken to convert Christians, the Christian government would have stopped at nothing to avert the danger. Given the ideas which then prevailed on the importance of State religions, we cannot be surprised that Yezdegerd should have permitted acts of persecution. Some of the Christians fled to Roman territory. The Imperial government refused to surrender them ( A.D. 420) and prepared for the event of war.\texthebrew{⁠}\footnote{A constitution authorising the inhabitants of the Eastern and Pontic provinces to build walls round their homes (May, 420) is interpreted as a measure taken in view of impending invasion.

\emph{C. J.} VIII.10.10 . Cp. Lebeau, V p493.

} Yezdegerd died at this juncture, and was succeeded by his son Varahran V, who was completely under the influence of the Zoroastrian priests, and began a general persecution.\texthebrew{⁠}\footnote{Labourt, 110 \emph{sqq.}

} Some outages were committed on Roman merchants. The war which resulted lasted for little more than a year, and the Roman armies were success ­ful.\texthebrew{⁠}\footnote{The general Ardaburius operated in Arzanene and gained a victory, autumn 421, which forced the Persians to retreat to Nisibis, which Ardaburius then besieged. He raised the siege on the arrival of an army under Varahran, who proceeded to attack Resaina. Meanwhile the Saracens of Hira, under Al\texthebrew{‑}Mundhir, were sent to invade Syria, and were defeated by Vitianus. During the peace negotiations the Persians attacked the Romans and were defeated by Procopius, son-in\texthebrew{‑}law of Anthemius (Socrates, VII.18, 20). The Empress Eudocia celebrated the war in a poem in heroic metre (\emph{ib.} 21).

} Then a treaty was negotiated by which peace was made

p5
for a hundred years ( A.D. 422). Varharan undertook to stay the persecution; and it was agreed that neither party should receive the Saracen subjects of the other.\footnote{Malchus (\emph{fr.} 1 in \emph{De Leg. gent.} p568) refers this provision to the peace concluding ``the greatest war'' in the time of Theodosius. This obviously means that of 422, not that of 442.

}

The attention of Varahran was soon occupied by the appearance of new enemies beyond the Oxus, who for more than a hundred years were constantly to distract Persian arms from the Roman frontier.\texthebrew{⁠}\footnote{The best study of the history of the Ephthalites is the memoir of Ed. Drouin in \emph{Le Muséon}, XIV (1895). See also A. Cunningham, \emph{Ephthalite or White Huns}, in \emph{Transactions} of Ninth International Oriental Congress, London, 1892.

Thayer's Note: Drouin's paper was published in four parts,

pp73\texthebrew{‑}84

\texthebrew{•}pp141\texthebrew{‑}161

\texthebrew{•}pp232\texthebrew{‑}247

\texthebrew{•}pp277\texthebrew{‑}288

} The lands between the Oxus and Jaxartes had for some centuries been in the hands of the Kushans. The Kushans were now conquered (\emph{c.} A.D. 425) by another Tartar people, who were known to the Chinese as the Ye\texthebrew{‑}tha, to Armenian and Arabic writers as the Haithal, and to the Greeks as the Ephthalites.\texthebrew{⁠}\footnote{Theophylactus Simocatta gives the alternative name of \textgreek{\texthebrew{Ἀ}βδελοί} (\emph{Hist.} VII.7.8).

} The Greek historians sometimes classify them as Huns, but add the qualification ``white,'' which refers to their fair complexion and distinguishes them from the true Huns ( Hiung\texthebrew{‑}nu ), who were dark and ugly.\texthebrew{⁠}\footnote{See \emph{e.g.} Proc. \emph{B.P.} I.3; Cosmas, \emph{Christ. Top.} XI.11. Procopius states that their habits were not nomad.

} The Ephthalites belonged in fact not to the Hiung\texthebrew{‑}nu , but to a different Turanian race, which was known to the Chinese as the Hoa. Their appearance on the Oxus marked a new epoch in the perennial warfare between Iran and Turan. They soon built up a considerable empire extending from the Caspian to the Indus, including Chorasmia, Sogdiana, and part of north-western India.\texthebrew{⁠}\footnote{Cosmas, \emph{l.c.}

} Their chief town was Balkh, and Gurgan\texthebrew{⁠}\footnote{\textgreek{Γοργώ}, Procop. \emph{l.c.}

} (on the river of the same name which flows into the Caspian) was their principal frontier fortress against Persia. The first hostilities against the Ephthalites broke out in A.D. 427 and resulted in a complete victory for Varahran.\footnote{The following is a chronological list of the Perso-Ephthalite wars (Drouin, \emph{op. cit.} p288):

A.D. 427 war under Varahran.

`` 442\texthebrew{‑}449 war under Yezdegerd II.

`` 450\texthebrew{‑}451 ``

`` 454 ``

`` 474\texthebrew{‑}476 `` Perozes.

`` 482\texthebrew{‑}484 `` ``

`` 485 war in interregnum.

`` 503\texthebrew{‑}513 war under Kavad.

`` 556\texthebrew{‑}557 `` Chosroes.

}

The reign of Theodosius II witnessed a second but less serious disturbance of the peace, soon after the accession of Yezdegerd II ( A.D. 438). The cause is uncertain. It has been conjectured, without sufficient evidence, that the Persian king was in league

p6
with Attila and Gaiseric for the destruction of the Empire.\texthebrew{⁠}\footnote{Güldenpenning, \emph{op. cit.} 340.

} It is possible that Persian suspicions had been provoked by the erection of a fortress at Erzerum in Roman Armenia, on the Persarmenian frontier, which was named Theodosiopolis.\texthebrew{⁠}\footnote{Moses of Chorene relates its foundation by Anatolius in his

\emph{Hist. Arm.} III.59 . As Book III ends in A.D. 433, this seems to be the lower limit for the date.

Procopius, \emph{Aed.} III.5 , p255 (cp. p210), ascribes the foundation to Theodosius I (and so Chapot, \emph{op. cit.} p361); but his confusion between the two Emperors of that name is quite clear in

III.1 , p210.

} This stronghold was to have a long history, reaching down to the present day, as one of the principal eastern defences of Asia Minor. Whatever motives may have instigated him to violate the peace, Yezdegerd raided Roman Armenia ( A.D. 440).\texthebrew{⁠}\footnote{And in Mesopotamia he advanced as far as Nisibis. See Elisha Vartabed, \emph{Hist. Arm.} c1, p184.

} Menaced, however, in his rear by an invasion of the Ephthalites he was easily bought off by Anatolius, the Master of Soldiers in the East, and Aspar. A new peace was then concluded ( A.D. 442), probably confirming the treaty of A.D. 422, with the additional stipulations that neither party should build a fortress within a certain distance of the frontier, and that the Romans should (as had been agreed by the treaty of A.D. 363) contribute a fixed sum to keep in repair the defences of the Caspian Gates against the barbarians beyond the Caucasus. ``Caspian Gates'' is a misleading name; for it was used to designate not, as one would expect, passes at the eastern extremity of the range, but passes in the centre, especially that of Dariel, north of Iberia. These danger-points were guarded by the Romans so long as they were over ­lords of Iberia, but now they abandoned Iberia to Persian influence and were therefore no longer in a position to keep garrisons in the mountain passes.\footnote{The Persians built the fortress of Biraparach (\textgreek{\texthebrew{Ἰ}ουροειπαάχ} Priscus, \emph{fr.} 15, \emph{De leg. gent.} p586; \textgreek{Βιραπαράχ} John Lydus, \emph{De mag.} III.52 ) probably in the pass of Dariel; and the fortress Korytzon (Menander, \emph{fr.} 3, \emph{De leg. Rom.} p180), which seems to be the Tzur of Procopius (\emph{B. G.} IV.3; De Boor conjectures \textgreek{χώρου Τζόν} in Menander), perhaps farther east. Cp. P.\texthebrew{‑}W. \emph{s.v.} Biraparach; Chapot, \emph{op. cit.} p369. See also Procopius, \emph{B. P.} I.10. Procopius (\emph{ib.} 2 \emph{ad fin.}) confounds the war of 420\texthebrew{‑}422 with that of 440\texthebrew{‑}441.

}

The greater part of Yezdegerd's reign was troubled by war with the Ephthalites. He made energetic efforts to convert Persian Armenia to the religion of Zoroaster, but the Armenians were tenacious of their Christianity and offered steady resistance to his armies. Since A.D. 428, when the last Arsacid king, Ardashir, had been deposed by the Persian monarch at the

p7
request of the Armenians themselves, the country had been ruled by Persian governors ( marzbans ).\texthebrew{⁠}\footnote{Cp. Lazarus, \emph{Hist. Arm.} c15, p272. Vramshapu had reigned from 392 to 414, then Chosroes III for a year, after whose death Yezdegerd appointed his own son Sapor. In 422 Varahran agreed to the accession of Ardashir, Vramshapu's son (Moses Chor. \emph{Hist. Arm.} III. c18).

} In A.D. 450 the Armenians sent a message to Constantinople imploring the Emperor to rescue them and their faith. Marcian, who had just come to the throne and was threatened by Attila, was not in a position to go to war with Persia for the sake of the Persarmenian Christians. He determined to be neutral, and Yezdegerd was informed that he need fear no hostilities from the Empire.\texthebrew{⁠}\footnote{Elisha Vartabed, \emph{Hist. Arm.} c3, pp206\texthebrew{‑}207; Lazarus, \emph{op. cit.} c36, p298. A full and tedious account of the wars in Armenia will be found in these writers who were contemporary. Elisha's history ends in 446, Lazarus comes down to the accession of Valakhesh.

} The war between the Armenians and their over ­lord continued after the death of Yezdegerd ( A.D. 453) during the reign of Firuz (Perozes), under the leader ­ship of Vahan the Mamigonian.

Firuz perished in a war with the Ephthalites, whose king had devised a cunning stratagem of covered ditches which were fatal to the Persian cavalry ( A.D. 484).\texthebrew{⁠}\footnote{Procopius, \emph{B. P.} I.4; Lazarus, \emph{op. cit.} c73.

} Valakhesh (Balas), perhaps his brother, followed him, and enjoyed a shorter but more peaceable reign. He made a treaty with the enemy, consenting to pay them a tribute for two years. He pacified Armenia by granting unreserved toleration; Vahan was appointed its governor; and Christianity was reinstated. Valakhesh died in A.D. 488.

During this period \texthebrew{—} the reigns of Marcian, Leo, and Zeno \texthebrew{—} there had been no hostilities between the two empires, but there had been diplomatic incidents. About A.D. 464 Perozes had demanded money from Leo for the defence of the Caucasian passes, had complained of the reception of Persian refugees, and of the persecution of the Zoroastrian communities which still existed on Roman territory.\texthebrew{⁠}\footnote{See above,

p4, n2 .

} Leo sent an ambassador who was received by the king, perhaps on the frontier of the Ephthalites, and the matters seem to have been amicably arranged.\texthebrew{⁠}\footnote{Priscus, \emph{fr.} 15, \emph{De Leg. gent.} p586, \emph{frs.} 11, 12, \emph{De leg. Rom.} It is difficult to reconcile the chronology with what is otherwise known of the first campaign of Perozes against the Ephthalites, whom Priscus apparently means by the Kidarites. The Kidarites proper seem to have been Huns who had settled in the trans-Caucasian

(p8)
country and threatened the pass of Dariel, and they are meant in another passage of Priscus (\emph{fr.} 22, \emph{De leg. gent.}) where Perozes announces to Leo that he has defeated them, \emph{c.} A.D. 468. For the Kidarites, and this assumed confusion, see Drouin, \emph{op. cit.} 143\texthebrew{‑}144.

} Ten years later an incident occurred which illustrates

p8
the danger of the extension of Persian influence to the Red Sea, although the Persian Government was in this case in no way responsible.\texthebrew{⁠}\footnote{The source is Malchus, \emph{fr.} 1, \emph{De leg. gent.} p568. Cp. Khvostov, \emph{Ist. vost. torgovli Egipta}, I. p199. Jotaba has been identified with Strabo's Dia

(XVI.4.18) , now Tiran. It was inhabited by a colony of Jews, once independent, according to Procopius, \emph{B. P.} I.19.4.

} A Persian adventurer, Amorkesos, who ``whether because he was not success ­ful in Persia or for some other reason preferred Roman territory,'' settled in the province of Arabia. There he lived as a brigand, making raids, not on the Romans but on the Saracens. His power grew and he seized Jotaba, one of the small islands in the mouth of the gulf of Akaba, the eastern inlet formed by the promontory of Sinai. Jotaba belonged to the Romans and was a commercial station of some importance. Driving out the Greek custom-house officers, Amorkesos took possession of it and soon amassed a fortune by collecting the dues. He made himself ruler of some other places in the neighbourhood, and conceived the desire of becoming a phylarch or satrap of the Saracens of Arabia Petraea, who were nominally dependent on the Roman Emperor. He sent an ecclesiastic to Leo to negotiate the matter, and Leo graciously signified his wish to have a personal interview with Amorkesos. When the Persian arrived, he shared the Imperial table, was admitted to assemblies of the Senate, and even honoured with precedence over the patricians. The Byzantines, it appears, were scandalised that these privileges should be accorded to a fire-worshipper , and Leo seems to have been obliged to pretend that his guest intended to become a Christian. On his departure Leo gave him a valuable picture, and compelled the members of the Senate to present him with gifts; and, what was more important, he transferred to him the possession of Jotaba, and added more villages to those which he already governed, granting him also the coveted title of phylarch.\texthebrew{⁠}\footnote{Leo was criticised for inviting Amorkesos to his court, and for permitting the foreigner to see the towns through which he help travel, unarmed and defenceless. Malchus, \emph{ib.}

} Jotaba, however, was not permanently lost. The Imperial authority there was re-established in the reign of Anastasius.\footnote{In A.D. 498 by Romanus (see above,

Chap. XIII § 1, p432 ). Theophanes, A.M. 5990. It was arranged that Roman traders should live in the island. Cp. Procopius, \emph{ib.}

}

p9
Valakhehs was succeeded on the Persian throne by Kavad, the son of Perozes. Kavad was in some ways the ablest of all the Sassanid sovrans. His great achievement was to restore the royal power, which had been gradually declining since the end of the fourth century, and was now well on its way towards the destiny which two hundred years later was to overtake the Merovingian kings of France. The kings had failed to retain their own authority over the Magian priesthood and the official or bureaucratic nobility, and the state was really managed by the principal minister whose title was wazurg-framadhar , and whose functions may be compared to those of a Praetorian Prefect.\texthebrew{⁠}\footnote{See Stein's important study of the reforms of Kavad and Chosroes, \emph{Ein Kapitel vom persischen und vom byzantinischen Staate} (\emph{Byz.-neugr. Jahrbücher}, I, 1920) p57.

} It was one of these ministers to whom Kavad owed his elevation.

Kavad might not have found it easy to emancipate the throne from the tutelage to which it had so long submitted, if there had not been a remarkable popular movement at the time of which he boldly took advantage.\texthebrew{⁠}\footnote{See Rawlinson, \emph{op. cit.} 342 \emph{sqq.}; Nöldeke, \emph{Tabari}, 455 \emph{sqq.} Cp. Tabari, 141 \emph{sq.},

Agathias, IV.27 ;

Procopius, \emph{B. P.} I.5 . The Mazdakites are designated as Manichaeans in John Mal. XVIII p444 , and the fuller account of Theophanes, A.M. 6016 . Both these notices are derived from Timotheus, a baptized Persian.

} A communist had arisen in the person of Mazdak, and was preaching successfully among the lower classes throughout Persia the doctrines that all men are equal, that the present state of society is contrary to nature, and that the acts condemned by society as crimes are, as merely tending to overthrow an unjustifiable institution, blameless. Community of property and wives was another deduction. Kavad embraced and actually helped to promulgate these anarchical doctrines. His conversion to Mazdakism was not, of course, sincere; his policy was to use the movement as a counterpoise to the power of the nobles and the Zoroastrian priests. There was a struggle for some years of which we do not know the details, but at length the nobles managed to immure the dangerous king in the Castle of ``Lethe'' ( A.D. 497).\texthebrew{⁠}\footnote{Giligerda, in Susiana.

} Mazdak was imprisoned, but forcibly released by his disciples. After a confinement of two or three years Kavad found means to escape, and with the help of the Ephthalites was reinstated on the throne ( A.D. 499).

p10
During his reign Kavad began a number of reforms in the organisation of the state which tended to establish and secure the royal authority. He did not do away with the high office of wazurg-framadhar , but he deprived it of its functions and it became little more than a honorific title.\texthebrew{⁠}\footnote{See Stein (\emph{ib.} p65), who suggests with much probability (p52) that the institution of the astabedh, a minister whose functions are compared by Greek and Syrian writers to those of the magister officiorum, was due to Kavad. The first mention of this official is in Joshua Styl. c59 (A.D. 502); see also Procopius, \emph{B. P.} I.11.25.

} He began a new survey of the land, for the purpose of instituting a system of sound finance.\texthebrew{⁠}\footnote{Tabari, p241.

} Towards the end of his reign his position was so strong that he was able to take measures to suppress the anti-social Mazdakite sect, which he had suffered only because the hostility between these enthusiasts and the nobles and priests helped him to secure and consolidate the royal power.

It was some time after the restoration of Kavad that hostilities broke out, after sixty years of peace between Persia and the Empire. In their financial embarrassments the Sassanid kings were accustomed to apply to Constantinople, and to receive payments which were nominally the bargained contribution to the defence of the Caucasian passes. The Emperors Leo and Zeno had extricated Perozes from difficulties by such payments.\texthebrew{⁠}\footnote{Joshua Styl. p7.

} But in A.D. 483 the Persians repudiated a treaty obligation. It had been agreed by the treaty of Jovian that Persia was to retain Nisibis for 120 years and then restore it to the Romans. This period now terminated and the Persians declined to surrender a fortress which was essential to their position in Mesopotamia. The Emperor Zeno did not go to war, but he refused to make any further payments for the defence of the Caucasus. When king Valakhesh applied to him he said: ``You have the taxes of Nisibis, which are due rightfully to us.''\texthebrew{⁠}\footnote{\emph{Ib.} pp7, 12.

} The Imperial Government cannot have seriously expected Persia to fulfil her obligation in regard to Nisibis, but her refusal to do so gave the Romans the legal right to decline to carry out their contract to supply money. Anastasius followed the policy of Zeno when Kavad renewed the demand with menaces in A.D. 491.\footnote{\emph{Ib.} p13. ``As Zeno did not send, so neither will I, until thou

(p11)
restorest to me Nisibis.'' Kavad applied again during the Isaurian War, and Anastasius offered to send him money as a loan, but not as a matter of custom (\emph{ib.} p15).

}

p11
After his restoration Kavad was in great straits for money. He owed the Ephthalites a large sum which he had undertaken to pay them for their services in restoring him to the throne, and he applied to Anastasius. The Emperor had no intention of helping him, as it appeared to be manifestly to the interest of the Empire to promote hostility and not friendship between the Ephthalites and the Persians. It is said that his refusal took the form of a demand for a written acknowledgment ( cautio ), as he knew that Kavad, unfamiliar with the usages of Roman law, would regard such a mercantile transaction as undignified and intolerable.\texthebrew{⁠}\footnote{Procopius, \emph{B. P.} I.7; Theodorus Lector, II.52; Theophanes, \emph{sub} A.M. 5996. John Lydus (\emph{De mag.} III.52) attributes the war to a demand for the costs of maintaining the castle of Biraparach, and doubtless the question of the Caucasian defences was mentioned in the negotiations. Kavad refers to the demand for money in his letter to Justinian quoted by John Mal. XVIII p450.

} Kavad resolved on war, and the Hundred Years' Peace was broken, not for the first time, after a duration of eighty years (August, A.D. 502).\footnote{Joshua Styl. p37.

}

The Persian monarch began operations with an invasion of Armenia, and Theodosiopolis fell into his hands by treachery. Then he marched southwards, attacked Martyropolis which surrendered, and laid siege to Amida. This city, after a long and laborious winter siege beginning in October, was surprised in January ( A.D. 503), chiefly through the negligence of some monks who had undertaken to guard one of the towers, and having drunk too much wine slumbered instead of watching.\texthebrew{⁠}\footnote{But whether the monks were to blame is doubtful (Haury, \emph{Zur Beurteilung des Proc.} 23).

} There was a hideous massacre which was stayed by the persuasions of a priest, the survivors were led away captive, and Amida was left with a garrison of 3000 men.\footnote{The siege of Amida is described by Joshua Styl. cc. l, liii; Zacharias Myt. VII.3; Procopius, \emph{B. P.} I.7. Eustathius of Epiphania described it in his lost history

(Evagrius, III.37) , and may have been the source of both Procopius and Zacharias; if not, Procopius must have used Zacharias (cp. Haury, \emph{Proleg.} to his ed. of Procopius, pp19\texthebrew{‑}20). The stories in the three sources are carefully compared by Merten, \emph{De bello Persico}, 164 \emph{sqq.}

}

On the first news of the invasion the Emperor had sent Rufinus as an ambassador to offer money and propose terms of peace.\texthebrew{⁠}\footnote{During the siege of Amida, Roman Mesopotamia was invaded and plundered by the Saracens of Hira under Naman (Joshua Styl. p39 \emph{sqq.}).

} Kavad detained him till Amida fell, and then

p12
despatched him to Constantinople with the news. Anastasius made military preparations, but the forces which he sent were perhaps not more than 15,000 men.\texthebrew{⁠}\footnote{So Marcellinus, \emph{sub a.} Joshua Styl. gives 40,000 men to Patricius and Hypatius and 12,000 to Areobindus.

} And, influenced by the traditions of the Isaurian campaigns, he committed the error of dividing the command, in the same theatre of war, among three generals. These were the Master of Soldiers in the East, Areobindus, the great-grandson of Aspar (on the mother's side) and son-in\texthebrew{‑}law of the Emperor Olybrius; and the two Masters of Soldiers in praesenti , Patricius, and the Emperor's nephew Hypatius, whose military inexperience did not deserve such a responsible post.\footnote{Priscian's Panegyric on Anastasius may perhaps be dated to this year. For he says of Hypatius quem vidit validum Parthus sensitque timendum (p300) and does not otherwise mention the war. Among the subordinate commanders were Justin (the future Emperor); Patriciolus and his son Vitalian; Romanus. Areobindus was Consul in 506, and his consular diptych is preserved at Zürich, with the inscription Flavius Areobindus Dagalaiphus Areobindus, V. I., Ex C. Sacri Stabuli et Magister Militum Per Orientem Ex Consule Consul Ordinarius. See CIL XIII.5245 ; Meyer, \emph{Zwei ant. Elfenb.} p65.

}

The campaign opened (May, A.D. 503) with a success for Areobindus, in the neighbourhood of Nisibis, but the enemy soon mustered superior forces and compelled him to withdraw to Constantia. The jealousy of Hypatius and Patricius, who with 40,000 men had encamped\texthebrew{⁠}\footnote{At Siphrios, 9 miles from Amida.

} against Amida, induced them to keep back the support which they ought to have sent to their colleague. Soon afterwards the Persians fell upon them, their vanguard was cut up, and they fled with the rest of their army across the Euphrates to Samosata (August).\footnote{John Lydus (\emph{De mag.} III.53) attributes the ill-success of the Romans to the incompetence of the generals, Areobindus, who was devoted to dancing and music, Patricius and Hypatius, who were cowardly and inexperienced. This seems borne out by the narratives of Procopius and Joshua. Cp. Haury, \emph{Zur Beurt. des Proc.} 24\texthebrew{‑}25.

}

Areobindus meanwhile had shut himself up in Edessa, and Kavad determined to attack it. The Christian legend of Edessa was in itself a certain challenge to the Persian kings. It was related that Abgar, prince of Edessa and friend of the Emperor Augustus, suffered in his old age from severe attacks of gout. Hearing of the miraculous cures which Jesus Christ was performing in Palestine, Abgar wrote to him, inviting him to leave a land of unbelievers and spend the rest of his life at Edessa. Jesus declined, but promised the prince recovery from his disease.

p13
The divine letter existed, and the Edessenes afterwards discovered a postscript, containing a pledge that their city would never be taken by an enemy. The text of the precious document was inscribed on one of the gates, as a sort of phylactery, and the inhabitants put implicit confidence in the sacred promise.\texthebrew{⁠}\footnote{Procopius, \emph{B. P.} II.12.

} It is said that the Saracen sheikh Naman urged on Kavad against Edessa, and threatened to do there worse things than had been done at Amida. Thereupon a wound which he had received in his head swelled, and he lingered in pain for two days and died.\texthebrew{⁠}\footnote{Joshua Styl. p47.

} But notwithstanding this sign Kavad persisted in his evil intention.

Constantia lay in his route, and almost fell into his hands. Here we have a signal example of a secret danger which constantly threatened Roman rule in the Eastern provinces, the disaffection of the Jews. The Jews of Constantia had conspired to deliver the city to the enemy, but the plot was discovered, and the enraged Greeks killed all the Jews they could find. Disappointed of his hope to surprise the fortress, Kavad did not stay to attack it, but moved on to Edessa. He blockaded this city for a few days without success (September 17), and Areobindus sent him a message: ``Now thou seest that the city is not thine, nor of Anastasius, but it is the city of Christ who blessed it, and it has withstood thy hosts.''\texthebrew{⁠}\footnote{This idea recurs in Procopius, who describes (\emph{B.P.} II.26 \emph{ad init.}) the Mesopotamian campaign of Chosroes, in which he besieged Edessa, as warfare ``not with Justinian nor with any other man, but with the God of the Christians.''

} But he deemed it prudent to induce the Persians to withdraw by agreeing to pay 2000 lbs. of gold at the end of twelve days and giving them hostages. Kavad withdrew, but demanded part of the payment before the appointed day. When this was refused he returned and renewed the blockade (September 24), but soon abandoned the enterprise in despair.

The operations of the following year were advantageous to the Empire. The evils of a divided command had been realised, Hypatius was recalled, and Celer, the Master of Offices, an Illyrian, was invested with the supreme command.\texthebrew{⁠}\footnote{I infer the superior authority of Celer from Joshua Styl. p55. He had arrived, early in 504, with a reinforcement of 2000 according to Marcellinus, but with a very large army according to Joshua.

} He invaded and devastated Arzanene; Areobindus invaded Persian Armenia;

p14
Patricius undertook the recovery of Amida. The siege of this place lasted throughout the winter till the following year ( A.D. 505). The garrison, reduced to the utmost straits by famine, finally surrendered on favourable terms. The sufferings of the inhabitants are illustrated by the unpleasant story that women ``used to go forth by stealth into the streets of the city in the evening or in the morning, and whomsoever they met, woman or child or man, for whom they were a match, they used to carry him by force into a house and kill and eat him either boiled or roasted.'' When this practice was betrayed by the smell of the roasting, the general put some of the women to death, but he gave leave to eat the dead.\footnote{Joshua Styl. p62. Cp. Procopius, \emph{B. P.} I.9 p44, from which it would appear that it was the few Roman inhabitants who were reduced to such straits.

}

The Romans paid the Persians 1000 lbs. of gold for the surrender of Amida. Meanwhile Kavad was at war with the Ephthalites, and he entered into negotiations with Celer, which ended in the conclusion of a truce for seven years ( A.D. 505).\texthebrew{⁠}\footnote{\emph{Ib.} p45. John Lydus, \emph{loc. cit.}

} It appears that the truce was not renewed at the end of that period, but the two empires remained actually at peace for more than twenty years.

It has been justly observed that in these oriental wars the Roman armies would hardly have held their own, but for the devoted loyalty of the civil population of the frontier provinces. It was through their heroic co-operation and patience of hunger that small besieged garrisons were able to hold out. Their labours are written in the remains of the stone fortresses in these regions.\texthebrew{⁠}\footnote{Chapot, \emph{op. cit.} p376.

} And they had to suffer sorely in time of war, not only from the enemy, but from their defenders. The government did what it could by remitting taxes; but the ill-used which they experienced from the foreign, especially the German, mercenaries in the Imperial armies was enough to drive them into the arms of the Persians. Here is the vivid description of their sufferings by one of themselves.\texthebrew{⁠}\footnote{Joshua Styl. p68. Cp. pp71\texthebrew{‑}73.

} º

``Those who came to our aid under the name of deliverers plundered us almost as much as our enemies. Many poor people they turned out of their beds and slept in them, whilst their owners lay on the ground in cold weather. Others they drove out of their own houses, and went in and dwelt in them. The

p15
cattle of some they carried off by force as if it were spoil of war; the clothes of others they stripped off their persons and took away. Some they beat violently for a mere trifle; with others they quarrelled in the streets and reviled them for a small cause. They openly plundered every one's little stock of provisions, and the stores that some had laid up in the villages and cities. Before the eyes of every one they ill-used the women in the streets and houses. From old women, widows, and the poor they took oil, wood, salt, and other things for their own expenses, and they kept them from their own work to wait upon them. In short they harassed every one both great and small. Even the nobles of the land, who were set to keep them in order and to give them their billets, stretched out their hands for bribes; and as they took them from every one they spared nobody, but after a few days sent other soldiers to those upon whom they had quartered them in the first instance.''

This war taught the Romans the existence of a capital defect in their Mesopotamian frontier. While the Persians had the strong fort of Nisibis against an advance to the Tigris, the Romans had no such defence on their own frontier commanding the high road to Constantia. After the conclusion of the treaty, Anastasius immediately prepared to remedy this weakness. At Daras, close to the frontier and a few miles from Nisibis, he built an imposing fortified town, provided with corn º -magazines, cisterns, and two public baths. He named it Anastasiopolis, and it was for the Empire what Nisibis was for Persia. Masons and workmen gathered from all Syria to complete the work while Kavad was still occupied by his Ephthalite war. He protested, for the building of a fort on the frontier was a breach of treaty engagements, but he was not in a position to do more than protest and he was persuaded to acquiesce by the diplomacy and bribes of the Emperor, who at the same time took the opportunity of strengthening the walls of Theodosiopolis.\footnote{Procopius, \emph{B. P.} II.10; Joshua Styl. p70. The fortifications of Daras will be described below,

Chap. XVI § 3 , in connection with the siege of Chosroes.

\texthebrew{◂} previous

next \texthebrew{▸}Images with borders lead to more information.

The thicker the border, the more information.

(Details here.)

UP TO:

J. B. Bury
Homepage

Latin \& Greek

Texts

LacusCurtius

Home}

\xchapter{Chapter XV, Part A}

Short URL for this page:

bit.ly/BURLAT15A

\xsubsection{(Part 1 of 3)}

Anastasius had made no provision for a successor to the throne, and there was no Augusta to influence the election. Everything turned out in a way that no one could have foreseen. The most natural solution might have seemed to be the choice of one of the late Emperor's three nephews, Probus, Pompeius, or Hypatius. They were men of average ability, and one of them, at least, Pompeius, did not share his uncle's sympathy with the Monophysitic creed. But they were not ambitious, and perhaps their claims were not seriously urged.\footnote{It is said, indeed, that there were many who wished that one of them should succeed

(Evagrius, \emph{H. E.} 4.2) . Anastasius had other relatives who were eligible (``numerous and very distinguished,'' Procopius, \emph{B. P.} I.11 ).

}

The High Chamberlain Amantius hoped to play the part which Urbicius had played on the death of Zeno, and he attempted to secure the throne for a certain Theocritus, otherwise unknown, who had probably no qualification but personal devotion to himself. As the attitude of the Palace guards would probably decide the election, he gave money to Justin, the Count of the Excubitors, to bribe the troops.\footnote{John Mal. XVII.410 (cp. \emph{Chr. Pasch., sub a.}; Cramer, \emph{Excerpta}, II.318); Marcellinus, \emph{sub} 519. Theocritus, described by Marcellinus as Amantii satelles, is designated as \textgreek{\texthebrew{ὁ} δομεστικός} in John Mal. \emph{fr.} 43, \emph{De ins.} p170. It means the ``domestic'' of Amantius, see Zach. Myt. IX.1.

}

In the morning (July 9) the people assembled in the Hippodrome and acclaimed the Senate. ``Long live the Senate! Senate of the Romans, tu vincas! We demand our Emperor, given by God, for the army; we demand our Emperor, given by God, for the world!'' The high officials, the senators, and

p17
the Patriarch had gathered in the Palace, clad most of them in mouse-coloured garments, and sat in the great hall, the Triklinos of the Nineteen Akkubita. Celer, the Master of Offices, urged them to decide quickly on a name and to act promptly before others (the army or the people) could wrest the initiative from their hands. But they were unable to agree, and in the meantime the Excubitors and the Scholarians were acting in the Hippodrome. The Excubitors proclaimed John, a tribune and a friend of Justin, and raised him on a shield. But the Blues would not have him; they threw stones and some of them were killed by the Excubitors. Then the Scholarians put forward an unnamed patrician and Master of Soldiers, but the Excubitors would not accept him and he was in danger of his life. He was rescued by the efforts of Justin's nephew, the candidatus Justinian. The Excubitors then wished to proclaim Justinian himself, but he refused to accept the diadem. As each of these persons was proposed, their advocates knocked at the Ivory Gate, which communicated between the Palace and the Hippodrome, and called upon the chamberlains to deliver the Imperial robes. But on the announcement of the name, the chamberlains refused.

At length, the Senate ended their deliberations by the election of Justin, and constrained him to accept the purple. He appeared in the Kathisma of the Hippodrome and was favourably received by the people; the Scholarians alone, jealous of the Excubitors, resented the choice. The coronation rite was immediately performed in the Kathisma. Arrayed in the Imperial robes, which the chamberlains at last delivered, he was crowned by the Patriarch John; he took the lance and shield, and was acclaimed Basileus by the assembly. To the troops he promised a donation of five nomismata (£3: 7: 6) and one pound of silver for each man.

Such is the official description of the circumstances of the election of Justin.\texthebrew{⁠}\footnote{Preserved in Constantine Porph. \emph{Cer.} I.93 (taken from the \emph{Katastasis} of Peter the Patrician).

} If it is true so far as it goes, it is easy to see that there was much behind that has been suppressed. The intrigue of Amantius is ignored. Not a word is said of the candidature of Theocritus which Justin had undertaken to support. If Justin had really used his influence with the

p18
Excubitors and the money which had been entrusted to him in the interest of Theocritus, it is hardly credible that the name of Theocritus would not have been proposed in the Hippodrome. If, on the other hand, he had worked in his own interest, as was naturally alleged after the event,\texthebrew{⁠}\footnote{Evagrius, \emph{loc. cit.}: Zach. Myt. \emph{loc. cit.}

} how was it that other names, but not his, were put forward by the Excubitors? The data seem to point to the conclusion that the whole \emph{mise en scène} was elaborately planned by Justin and his friends. They knew that he could not count on the support of the Scholarians, and, if he were proclaimed by his own troops alone, the success of his cause would be doubtful. The problem therefore was to manage that the initiation should proceed from the Senate, whose authority, supported by the Excubitors, would rally general consent and overpower the resistance of the Scholarian guards. It was therefore arranged that the Excubitors should propose candidates who had no chance of being chosen, with the design of working on the fears of the Senate. Justin's friends in the Senate could argue with force: ``Hasten to agree, or you will be forestalled, and some wholly unsuitable person will be thrust upon us. But you must choose one who will be acceptable to the Excubitors. Justin fulfils this condition. He may not be an ideal candidate for the throne, but he is old and moderate.'' But, however the affair may have been managed by the wirepullers, Justin ascended the throne with the prestige of having been regularly nominated by the Senate, and he could announce to the Pope that ``We have been elected to the Empire by the favour of the indivisible Trinity, by the choice of the highest ministers of the sacred Palace, and of the Senate, and finally by the election of the army.''\footnote{\emph{Coll. Avellana, Ep.} 141.

}

The new Emperor, who was about sixty-six years of age, was an Illyrian peasant. He was born in the village of Bederiana in the province of Dardania, not far from Scupi, of which the name survives in the town of Üsküb, and his native language was Latin.\texthebrew{⁠}\footnote{Born in 452, if he was 75 at his death (John Mal. XVII.424; but 77 acc. to \emph{Chr. Pasch.}). Bederiana is represented by the modern village of Bader. Justinian speaks of Scupi, which he renamed Justiniana Prima, as patria nostra (\emph{Nov.} 11). On the identification, see Evans, \emph{Arch. Researches}, II.141 \emph{sqq.}

} Like hundreds of other country youths,\texthebrew{⁠}\footnote{Cp. in the address of Germanus to his soldiers, in Procopius, \emph{B. V.} II.16 \textgreek{\texthebrew{ὑ}μ\texthebrew{ᾶ}ς \texthebrew{ἐ}ξ \texthebrew{ἀ}γρο\texthebrew{ῦ} ξύν τε τ\texthebrew{ῇ} πήρ\texthebrew{ᾳ} κα\texthebrew{ὶ} χιτωνίσκ\texthebrew{ῳ ἐ}νί.}} he set forth

p19
with a bag of bread on his back and walked to Constantinople to better his fortune by enlisting in the army. Two friends accompanied him, and all three, recommended by their physical qualities, were enrolled in the Palace guards.\texthebrew{⁠}\footnote{Procopius, \emph{H. A.} VI . John Mal. (XVII.410) describes Justin as good-looking, with a well-formed nose, and curly grey hair.

} Justin served in the Isaurian and Persian wars of Anastasius, rose to be Count of the Excubitors, distinguished himself in the repulse of Vitalian, and received senatorial rank.\texthebrew{⁠}\footnote{\emph{Ib.}; John Ant. \emph{De ins.}, \emph{fr.} 100 (p142) Theodorus Lector, II.37.

} He had no qualifications for the government of a province, not to say of an Empire; for he had no knowledge except of military matters, and he was uneducated.\texthebrew{⁠}\footnote{John Lydus, \emph{De mag.} III.51 , Procopius, \emph{ib.}, John Mal. \emph{ib.} \textgreek{\texthebrew{ἀ}γράμματος.}} It is even said that he could not write and was obliged, like Theoderic the Ostrogoth, to use a mechanical device for signing documents.

He had married a captive whom he had purchased and who was at first his concubine. Her name was Lupicina, but she was crowned Augusta under the more decorous name of Euphemia.\texthebrew{⁠}\footnote{Victor Tonn. \emph{s. a.} 518; Theodore Lector, II.37, cp. Cramer \emph{Anecd. Par.} II.108; Procopius, \emph{H. A.} VI , IX . On a coin supposed to represent Euphemia, see Wroth, \emph{Imp. Byzantine Coins}, I p. xiv n4. There are miniature representations of Justin and Euphemia on the two extremities of the horizontal bar of a silver cross preserved in the Treasury of St. Peter's at Rome. The cross bears the inscription:

ligno quo Christus humanum subdidit hostem

dat Romae Justinus opem et socia decorem.

From the style of the headcap of the Empress, Delbrück (\emph{Porträts byz. Kais.}) was able to infer that Justin I and Euphemia (not Justin II and Sophia) are in question.

} In his success ­ful career the peasant of Bederiana had not forgotten his humble relatives or his native place. His sister, wife of Sabbatius, lived at the neighbouring village of Tauresium\texthebrew{⁠}\footnote{Now Taor. Justinian built a rectangular wall round it, with a tower at each corner, and called it Tetrapyrgia.

Procopius, \emph{Aed.} IV.1.18 .

} and had two children, Petrus Sabbatius and Vigilantia. He adopted his elder nephew, brought him to Constantinople, and took care that he enjoyed the advantages of an excellent education. The young man discarded the un-Roman names of Peter and Sabbatius\texthebrew{⁠}\footnote{His name, however, appeared in full on his consular diptychs of 521. CIL V.8210 , 3: Fl. Petr. Sabbat. Iustinian., v.i., comes, mag. eqq. et p. praes., et consul ord.

Thayer's Note: The notion that Justinian was a Slav whose original name was Uprauda or Vpravda, whose mother's name was Bigleniza, and so on with the rest of his family, has been firmly exploded; the source of it is a fraud, very likely perpetrated by a 17c Slav. The details, and the detective work in unmasking the imposture, may be found in Bryce and Jireček,

``Life of Justinian by Theophilus'' . Such, however, is the persistence of error that 120 years after it was corrected, the nonsense is still plastered all over the Web.

} and was known by the adoptive name of Justinianus. He was enrolled among the candidati. Justin had other nephews and seems to have cared also for their fortunes. They were liberally educated and were destined to

p20
play parts of varying distinction and importance on the political scene.\footnote{Vigilantia, who married Dulcissimus, had three children, Justin afterwards Emperor), Marcellus, and Praejecta. A brother of Justin, or another sister, had three sons, Germanus, Boraides, and Justus (cp. Procopius, \emph{B. P.} I.24.53; \emph{B. G.} III.31.12). Germanus, who was to play a considerable part, was thus the cousin of Justinian. He married (1) Passara, by whom he had two sons, Justin and Justinian, and one daughter Justina; (2) the Ostrogothic princess Matasuntha, by whom he had one son Germanus. See the Genealogical Table . Another cousin of Justinian, named Marcian, is mentioned by John Malalas, XVIII.496. I conjecture that he may be identical with Justin, son of Germanus. For this Justin was consul in 540, and on his consular diptych his name runs: Fl. Mar. Petr. Theodor. Valent. Rust. Boraid. Germ. Iust. (Meyer, \emph{Zwei ant. Elf.} p10). I take Mar. to be for Marcianus. Germanus was the \textgreek{\texthebrew{ἀ}νεψιός} of Justinian (Proc. \emph{locc. citt.}), and this has generally been taken to mean nephew, so that Justinian would have had a brother or a second sister. But I agree with Kallenberg (\emph{Berl. phil. Wochenschrift}, XXXV.991) that in \emph{B. G.} IV.40.5 \textgreek{\texthebrew{Ἰ}ουστ\texthebrew{ῖ}νος \texthebrew{ὁ} Γερμανο\texthebrew{ῦ} θε\texthebrew{ῖ}ος} should be retained (all the editions print the emendation \textgreek{\texthebrew{Ἰ}ουστινιανός}).

}

The first care of Justin was to remove the disaffected; Amantius and Theocritus were executed, and three others were punished by death or exile.\texthebrew{⁠}\footnote{Marcellinus, \emph{s. a.} 519, Procopius, \emph{H. A.} 6 ; John Mal. XVII.410, and \emph{fr.} 43, \emph{De ins.} Marcellinus describes Amantius as a Manichean; Procopius says that there was no charge against him, except of using insulting language about the Patriarch; Malalas speaks of a demonstration against Amantius and Marinus in St. Sophia.

} His next was to call to Constantinople the influential leader who had shaken the throne of Anastasius. Before he came to the city, Vitalian must have been assured of the religious orthodoxy of the new Emperor, and he came prepared to take part in the reconciliation of Rome with the Eastern Churches. He was immediately created Master of Soldiers in praesenti ,\texthebrew{⁠}\footnote{John Mal. \emph{ib.} 411.

} and in A.D. 520 he was consul for the year. The throne of Justin seemed to be firmly established. The relatives of Anastasius were loyal; Pompeius co-operated with Justinian and Vitalian in the restoration of ecclesiastical unity. Marinus, the trusted counseller of the late sovran, was Praetorian Prefect of the East in A.D. 519.\footnote{\emph{C. J.}

V.27.7 ;

II.7.25 ; John Mal. \emph{ib.} records that Appion, who had been exiled by Anastasius, was recalled and made Pr. pr. Or. Perhaps he held the post in 518\texthebrew{‑}519 and was succeeded by Marinus.

}

The reunion with Rome, which involved the abandonment of the Henotikon of Zeno, the restoration of the prestige of the Council of Chalcedon, and the persecution of the Monophysites, was the great inaugural act of the new dynasty.\texthebrew{⁠}\footnote{See below,
Chap. XXII § 4 .

} The Emperor's nephew, Justinian, was deeply interested in theological questions, and was active in bringing about the ecclesiastical

p21
revolution. His intellectual powers and political capacity must have secured to him from the beginning a preponderant influence over his old uncle, and he would naturally regard himself as the destined successor to the throne. Immediately after Justin's election, he was appointed Count of the Domestics; and then he was invested with the rank of patrician, and was created a Master of Soldiers in praesenti .\texthebrew{⁠}\footnote{See \emph{Coll. Avell.} 162, 154, 230 (p696). He is mag. eqq. et p. praes. on his consular diptychs.

} His detractors said that he was unscrupulous in removing possible competitors for political influence. The execution of Amantius was attributed to his instigation.\texthebrew{⁠}\footnote{Procopius, \emph{H. A.} VI .

} Vitalian was a more formidable rival, and in the seventh month of his consul ­ship Vitalian was murdered in the Palace. For this crime, rightly or wrongly, Justinian was also held responsible.\texthebrew{⁠}\footnote{\emph{Ib.} Here Procopius is supported by Victor Tonn. \emph{s. a.} 523 (Iustiniani patrii factione). Loofs (\emph{Leontius}, 259) does not believe in the guilt of Justinian. John Malalas (\emph{De ins.}, \emph{fr.} 43) seems to connect the murder (which was committed in the Delphax in the Palace) with riotous demonstrations of the Blues and Greens in the Hippodrome and the streets.

} During the remaining seven years of the reign we may, without hesitation, regard him as the directing power of the Empire.\texthebrew{⁠}\footnote{In the Secret History Procopius treats the reign of Justin as virtually part of that of Justinian. That this view of Justinian's influence was generally accepted is shown by passages in the \emph{Public History} and the \emph{De aedificiis}. \emph{B. V.} I.9.5;

\emph{Aed.} I.3.3 .

} He held the consul ­ship in A.D. 521 and entertained the populace with magnificent spectacles.\texthebrew{⁠}\footnote{He spent 288,000 (£18,000) on the shows, exhibited 20 lions and 30 leopards. Marcellinus, \emph{s. a.}

} When he was afterwards elevated to the rank of nobilissimus ,\texthebrew{⁠}\footnote{Before 527; Marcellinus, \emph{s. a.} Victor Tonn. states that Justinian was created Caesar in 525 (\emph{s. a.}), but his authority is inferior. Cyril, \emph{Vit. Sabae}, p386, does not mention either title.

} it was a recognition of his position as the apparent heir to the throne. We may wonder why he did not receive the higher title of Caesar; perhaps Justin could not overcome some secret jealousy of the brilliant nephew whose fortune he had made.

Justinian's power behind the throne was sustained by the enthusiastic support of the orthodox ecclesiastics, but he is said to have sought another means of securing his position, by attracting the devotion of one of the Factions of the Hippodrome. Anastasius had shown favour to the Greens; and it followed almost as a matter of course that Justinian should patronise the Blues. In each party there was a turbulent section which was a standing menace to public order, known as the

p22
Partisans,\texthebrew{⁠}\footnote{\textgreek{\texthebrew{ῶ}ταιΟ\texthebrew{ἱ} στασι\texthebrew{ῶ}ται}. Procopius, \emph{H. A.} VII . He says that they affected a peculiar dress, wearing very wide sleeves drawn tight at the wrist, and imitating the costume of the Huns in trousers and shoes. They allowed their beards and moustaches to grow, shaved the head in front and wore the hair long behind. They used to go about in organised bands at night and rob the passers-by. For the connexion of Justinian with the Blues, which rests on the evidence of the Secret History, cp. Panchenko, \emph{O tain. ist. Prok.} 89 \emph{sqq.}

} and Justinian is alleged to have enlisted the Blue Partisans in his own interest. He procured official posts for them, gave money to those who needed it, and above all protected them against the consequences of their riots. It is certain that during the reign of Justin, both the capital and the cities of the East were frequently troubled by insurrections against the civil authorities and sanguinary fights; and it was the Blue Faction which bore the chief share of the guilt.\texthebrew{⁠}\footnote{John Mal. XVII.416 \textgreek{τ\texthebrew{ὸ} Βένετον μέρος \texthebrew{ἐ}ν πάσαις τα\texthebrew{ῖ}ς πόλεσιν \texthebrew{ἠ}τάκτει} (cp. \emph{H. A.} VIII. \emph{ad init.} ) \textgreek{κα\texthebrew{ὶ ἐ}τάρασσον τ\texthebrew{ὰ}ς πόλεις λιθασμο\texthebrew{ῖ}ς κα\texthebrew{ὶ} καταβασίαις κα\texthebrew{ὶ} φόνοις}. This refers to the first years of the reign. Cp. Mansi, VIII.1106 (relating to Syria Secunda).

} The culminating scandal occurred in A.D. 524.\texthebrew{⁠}\footnote{The date 524 may be inferred by combining John Mal. (\emph{ib.}), who gives indiction 3 (=524\texthebrew{‑}525) for the fall of Theodotus, with Theophanes (A.M. 6012), who mentions the sixth year of Justin; and it is confirmed by

\emph{C. J.} II.7.26 , which was addressed to Theodotus in 524. Theodore was Prefect of the City in 520 (John Mal. \emph{fr.} 43, \emph{De ins.}); Theodotus was appointed in 522\texthebrew{‑}523 (John Mal. XVII.416); and was succeeded by Theodore Têganistes (\emph{ib.}). Hence in

\emph{C. J.} IX.19.6

(A.D. 526) Theodoto is probably an error for Theodoro.

} On this occasion a man of some repute was murdered by the Partisans in St. Sophia. Justinian happened to be dangerously ill at the time, and the matter was laid before the Emperor. His advisers seized the opportunity to urge upon him the necessity of taking rigorous measures to suppress the intolerable licence of these enemies of society. Justin ordered the Prefect of the City, Theodotus Colocynthius, to deal out merciless justice to the malefactors.\texthebrew{⁠}\footnote{\emph{H. A.} IX . The other sources do not mention Justinian's illness.

} There were many executions, and good citizens rejoiced at the spectacle of assassins and plunderers being hanged, burned, or beheaded.\texthebrew{⁠}\footnote{Marcellinus, whose notice, though dated A.D. 523, must refer to this affair.

} Theodotus, however, was immediately afterwards deprived of his office and exiled to Jerusalem, and his disgrace has been attributed to the resentment of Justinian who had unexpectedly recovered from his disease.\texthebrew{⁠}\footnote{Procopius (\emph{ib.}) says that some of the friends of Theodotus were tortured, and confessed that he had spoken disloyally against Justinian, but that the Quaestor Proclus took his part and proved that he had done nothing to deserve death. John Malalas (\emph{ib.}) has a different story. He ascribes Justin's anger to the fact that Theodotus had executed a rich senator without consulting himself. Both accounts may be true. According to

(p23)
to John of Nikiu (p503) he arrested a nephew of Justin. Procopius and Malalas agree that at Jerusalem Theodotus remained in concealment, believing that his life was in danger.

} However this may have been, the

p23
Blues had received an effective lesson, and during the last years of the reign not only the capital but the provincial cities also enjoyed tranquillity.\footnote{John Mal. \emph{ib.}, where it is stated the public spectacles were generally prohibited, and that all professional dancers were banished from the East, but that an exception was made in the case of Alexandria.

}

There were few events of capital importance during the reign of Justin. Its chief significance lay in the new orientation of religious policy which was inaugurated at the very beginning, and in the long apprentice ­ship to statecraft which it imposed on Justinian before the full power and responsibility of government devolved on him. Next to him the most influential minister was Proclus the Quaestor, an incorruptible man who had the reputation of an Aristides.\texthebrew{⁠}\footnote{Procopius, \emph{B. P.} I.11.11;

\emph{H. A.} VI.13\texthebrew{‑}14 ; John Lydus, \emph{De mag.} III.20 \textgreek{Πρόκλος \texthebrew{ὁ} δικαιότατος.}} There was some danger of a breach with the Ostrogothic ruler of Italy in A.D. 525\texthebrew{‑}526, but this menace was averted by his death,\texthebrew{⁠}\footnote{See below,
Chap. XVIII § 1 .

} and the Empire enjoyed peace till the last year of the reign, when war broke out with Persia.

In the spring of A.D. 527 Justin was stricken down by a dangerous illness, and he yielded to the solicitations of the Senate to co-opt Justinian as his colleague. The act of coronation was performed in the great Triklinos in the Palace (on April 4), and it seems that the Patriarch, in the absence of the Emperor, placed the diadem on the head of the new Augustus. The subsequent ceremonies were carried out in the Delphax, where the Imperial guards were assembled, and not, as was usual, in the Hippodrome.\texthebrew{⁠}\footnote{Constantine Porph. \emph{Cer.} I.95 (from the \textgreek{κατάστασις} of Peter the Patrician).

} Justin recovered, but only to survive for a few months. He died on August 1, from an ulcer in the foot where, in one of his old campaigns, he had been wounded by an arrow.\footnote{John Mal. XVII.424.

}

The Emperor Justinian was about forty-five years old when he ascended the throne.\texthebrew{⁠}\footnote{Zonaras, XIV.5.40 (we do not know the source).

} Of his personal appearance we can

p24
form some idea from the description of contemporary writers\texthebrew{⁠}\footnote{Procopius, \emph{H. A.} VIII.12\texthebrew{‑}13 ; John Mal. XVIII.425.

} and from portraits on his coins and in mosaic pictures.\texthebrew{⁠}\footnote{There are two pictures at Ravenna, one in the apse of S. Vitale dating from A.D. 547, the other in the nave of S. Apollinare Nuovo, about ten years later. The former is a bad portrait; the face is oval, whereas all the other evidence both literary and monumental concurs in showing that it was round. He also wears a moustache; perhaps this was true in 547, though not in 538 nor 557 (John Mal. speaks of a beard, but if this is not simply an error, it must refer to the very end of his life). The other picture, truer to life, shows a round, smooth, shaven face, and conveys the impression of a man who is losing his old energy. The evidence of the coinage, admirably elucidated by Wroth (\emph{Byz. Coins}, I pp. xc\texthebrew{‑}xcii), is more important. The early coins of the reign display a purely conventional face, but in A.D. 538 changes were introduced. Bronze money was inscribed with the year of issue, and a new Imperial bust appeared both on gold and on bronze coins, and was not changed again. The private bust on the gold of A.D. 527 had a three-quarters face; the new bust showed a full face, shaven, round, plump, with a slight smile \texthebrew{—} unquestionably a genuine portrait of the man whom the picture in S. Apollinare shows when he was 20 years older. There is also a gold medallion (perhaps of 534; cp. Wroth, I p25, and

Cedrenus, I p649) , on which the Emperor's bust appears with round shaven face. The fifteenth-century drawing of Justinian's equestrian statue in the Augusteum (reproduced in Diehl, \emph{Justinien}, p27) does not help, nor the silver disk of Kerch which shows an Emperor on horseback (Diehl, p30, but the identity of the Emperor is doubtful). The Barberini ivory (Diehl, frontispiece) would be useful, if its date were certain, but some ascribe it to the age of Constantine the Great. Compare (as well as Wroth) Diehl's interesting appreciation of the mosaics.

} He was of middle height, neither thin nor fat; his smooth shaven face was round, he had a straight nose, a firm chin, curly hair which, as he aged, became thin in front. A slight smile seems to have been characteristic. The bust which appears on the coinage issued when he had reached the age of fifty-six , shows that there was some truth in the resemblance which a hostile writer detected between his countenance and that of the Emperor Domitian.

His intellectual talents were far above the ordinary standard of Roman Emperors, and if fortune had not called him to the throne, he would have attained eminence in some other career. For with his natural gifts he possessed an energy which nothing seemed to tire; he loved work, and it is not improbable that he was the most hardworking man in the Empire. Though his mind was of that order which enjoys occupying itself with details, it was capable of conceiving large ideas and embra­cing many interests. He permitted himself no self-indulgence ; and his temperance was ascetic. In Lent he used to fast entirely for two days, and during the rest of the season he abstained

p25
from wine and lived on wild herbs dressed with oil and vinegar. He slept little and worked far into the night.\texthebrew{⁠}\footnote{The statements of Procopius in

\emph{H. A.} XIII.28\texthebrew{‑}30 º
and in

\emph{Aed.} I.7.7\texthebrew{‑}11

are almost identical. In the latter passage it is said that these excessively abstemious practices caused a painful disease in the knee, which was miraculously healed by the relics of saints which had been discovered at Melitene.

} His manners were naturally affable. As Emperor he was easily accessible, and showed no offence if a bold or tactless subject spoke with a freedom which others would have resented as disrespectful. He was master of his temper, and seldom broke out into anger.\texthebrew{⁠}\footnote{\emph{H. A.} 13.1\texthebrew{‑}3 .

} He could exhibit, too, the quality of mercy. Probus, the nephew of Anastasius, accused of reviling him, was tried for treason. When the report of the trial was laid before the Emperor he tore it up and said to Probus, ``I pardon you for your offence against me. Pray that God also may pardon you.''\footnote{John Mal. XVIII.438. We shall meet another instance in the case of the conspirator Artabanes. Cp. Procopius, \emph{Aed.} I.1. 10

and

16 .

}

The reign of a ruler endowed with these estimable qualities, animated by a strong and unflagging sense of duty, devoting himself day and night to the interests of the State\texthebrew{⁠}\footnote{Cp. John Lydus, \emph{De mag.} III.55 \textgreek{τ\texthebrew{ὸ}ν πάντων βασιλέων \texthebrew{ἀ}γρυπνότατον}. Cp. \textgreek{βασιλ\texthebrew{ῆ}ος \texthebrew{ἀ}κοιμήτοιο} in the verses inscribed in the church of SS. Sergius and Bacchus, which he built (CIG IV.8639).

} for thirty-eight years, could not fail to be memorable. Memorable assuredly it was. Justinian wrought not only for his own time but for posterity. He enhanced the prestige of the Empire and enlarged its borders. He bequeathed, by his monumental work in Roman law, an enduring heritage to Europe; while the building of the Church of St. Sophia would in itself be an imperishable title to the gratitude of men. These achievements, however, are only one side of the picture. The successes and glories of his reign were to be purchased at a heavy cost, and the strain which he imposed on the resources of the State was followed by decline and disaster after his death. Perhaps no more scathing denunciation of the character, aims, and methods of a ruler has ever been written than the notorious indictment which the contemporary historian Procopius committed to the pages of a Secret History, wherein Justinian is represented as a malignant demon in human form.\texthebrew{⁠}\footnote{The credibility of the \emph{Secret History} is discussed below,

Chap. XXIV .

} Though the exaggerations of the writer are so gross and manifest that his venomous pen defeats its own object, there is sufficient evidence from other

p26
sources to show that the reign of Justinian was, in many ways, far from being a blessing to his subjects.

The capital error of Justinian's policy was due to a theory which, though not explicitly formulated till quite recent times, has misled many eminent and well-meaning sovrans and statesmen in all periods of history. It is the theory that the expansion of a state and the exaltation of its prestige and honour are ends in themselves, and valuable without any regard to the happiness of the men and women of whom the state consists. If this proposition had been presented nakedly either to Justinian or to Louis XIV, he would have indignantly repudiated it, but both these monarchs, like many another, acted on it, with most unhappy consequences for their subjects. Justinian possessed imagination. He had formed a high ideal of the might and majesty of the Empire of which he was the master. It humiliated him to contrast its moderate limits with the vast extent of territory over which the word of Constantine or Theodosius the Great had been law. He was dazzled by the idea of restoring the old boundaries of the Roman Empire. For though he only succeeded in recovering, as we shall see, Africa, Italy, and a small strip of Spain, his designs reached to Gaul, if not to Britain. After he had conquered the African provinces he announced his ambitious policy. ``We have good hopes that God will grant us to restore our authority over the remaining countries which the ancient Romans possessed to the limits of both oceans and lost by subsequent neglect.''\texthebrew{⁠}\footnote{\emph{Nov.} 30, § 11 , published just after the conquest of Sicily, in 536.

} In drawing up this magnificent programme, Justinian did not consider whether such an extension of his government would make his subjects, who had to bear the costs of his campaigns, happier or better. He assumed that whatever increased the power and glory of the state must also increase the well-being of its members. The resources of the state were not more than sufficient to protect the eastern frontier against the Persians and the Danubian against the barbarians of the north; and if the Emperor had been content to perform these duties more efficiently than his predecessors, he would unquestionably have deserved better of his subjects.

His conception of the greatness of the Empire was indissolubly associated with his conception of the greatness of its sovran,

p27
and he asserted the absolutism of the autocrat in a degree which no Emperor had hitherto attempted.\texthebrew{⁠}\footnote{Agathias, V.14

\textgreek{α\texthebrew{ὐ}τοκράτωρ \texthebrew{ὀ}νόματι τε κα\texthebrew{ὶ} πράγματι \texthebrew{ἀ}πεδέδεικτο.}} This was conspicuously shown in the dictator ­ship which he claimed over the Church. He was the first Emperor who studied dogmatic questions independently and systematically, and he had all the confidence of a professional theologian. A theologian on the throne is a public danger, and the principle of persecuting opinion, which had been fitfully and mildly pursued in the fifth century, was applied rigorously and systematically under Justinian. His determination to be supreme in all departments made him impatient of advice; he did not like his commands to be discussed, and he left to his ministers little latitude for decision. His passion for dealing personally with the minute details of government had the same unfortunate results as in the case of Philip II.\texthebrew{⁠}\footnote{Diehl has noted the resemblance. Diehl's judgment of the Emperor's character is that, with many high qualities, he had ``une âme de valeur plutôt médiocre'' (p21).

} Like other autocrats, he was jealous and suspicious, and ready to listen to calumnies against his most loyal servants. And there was a vein of weakness in his character. He faltered at one supremely critical moment of his reign, and his consort, Theodora, had an influence over him which no woman could have exercised over an Augustus or a Constantine.

It was probably before he had any prospect of the throne that Justinian formed a violent attachment to a girl of exceptional charms and talents, but of low birth and blemished reputation. Theodora had already borne at least one child to a lover\texthebrew{⁠}\footnote{A daughter, whose son Anastasius or Athanasius married Joannina, daughter of Belisarius. Procopius,

\emph{H. A.} IV.37 ; John Eph. Part III V.1. (Perhaps the same Athanasius is meant in John Philoponus, \emph{De opif. mundi}, I \emph{Prooem.}, as Reichardt thinks.) According to

\emph{H. A.} 17, 16 \emph{sqq.}

she had also a son John, who was taken as an infant by his unnamed father to Arabia because Theodora wished to destroy her offspring. When he had grown up, he was informed by his dying father of the secret of his birth. He went to Byzantium and revealed himself to the Emperor, who arranged that he should never be seen again.

} when she captured the heart of the future Emperor. According to a tradition \texthebrew{—} and perhaps she countenanced this story herself, for she could not deny the humility of her birth \texthebrew{—} she had come from Paphlagonia to the capital, where she was

p28
discovered by Justinian, making a scanty living by spinning wool.\texthebrew{⁠}\footnote{\textgreek{Πάτρια}, p248. To commemorate her old abode she founded the church of St. Panteleemon on the site.

} But contemporary rumours which were circulated by her enemies assigned to her a less respectable origin, and told a circumstantial story of a girlhood spent in singular infamy. She was said to be the daughter of Acacius, who was employed by the Green Faction at Constantinople as keeper of the wild beasts,\texthebrew{⁠}\footnote{\textgreek{\texthebrew{Ἀ}ρκτοτρόφος}, bear-keeper, was the term.

} which they exhibited at public spectacles. When Acacius died his widow married his successor, but this man was soon deprived of the office in favour of another who paid a bribe to obtain it.\texthebrew{⁠}\footnote{The official known as \textgreek{\texthebrew{ὀ}ρχηστής} had these appointments in his hands.

} The woman sent her three little daughters, Comito,\texthebrew{⁠}\footnote{We know from another source than \emph{H. A.} that Theodora had a sister Comito. She married Sittas, Master of Soldiers. John Mal. XVIII.430.

} Theodora, and Anastasia, in the guise of suppliants with fillets on their heads, to beg the Greens assembled in the Hippodrome to reinstate their stepfather who had been so unjustly treated. The Greens obdurately refused; but the Blues had compassion and appointed the man to be their own bear-keeper , as the post happened to be vacant. This incident of her childhood was said to be the explanation of the Empress Theodora's implacable hostility towards the Greens. The three sisters, when they were older, went on the stage, and in those days an actress was almost synonymous with a prostitute. According to the scandalous gossip, which is recorded with malicious relish in the \emph{Secret History} of Procopius, Theodora showed exceptional precocity and shamelessness in a career of vice. Her adventures were not confined to Constantinople. She went to the Libyan Pentapolis as the mistress of a new governor, but having quarrelled with him she betook herself to Alexandria, and worked her way back to the capital, where she entrapped Justinian.\footnote{The account of her career in

\emph{H. A.} IX

stands alone. Some have thought that it gains some support from a passage in John Eph. \emph{Comm.} p68, where the Empress is described as \textgreek{τ\texthebrew{ὴ}ν \texthebrew{ἐ}κ το\texthebrew{ῦ} πορνείου} (so Panchenko, \emph{op. cit.} 73), but these words are certainly an interpolation, for it is incredible that they were written by John, who was a devoted admirer of the Empress (cp. Diehl, \emph{op. cit.} 42). Was the interpolator acquainted with the \emph{Secret History}? Perhaps the expression is due to a tradition that Theodora had acted at a theatre at Constantinople which was in a street known by the suggestive name of \textgreek{Πόρναι}. See Justinian,

\emph{Nov.} 105, § 1

\textgreek{πρόοδον τ\texthebrew{ὴ}ν \texthebrew{ἐ}π\texthebrew{ὶ} τ\texthebrew{ὸ} θέατρον \texthebrew{ἄ}γουσαν \texthebrew{ἣ}ν δ\texthebrew{ὴ} Πόρνας καλο\texthebrew{ῦ}σιν.}}

This chapter of her biography, which reposes solely on the

p29
testimony of enemies, has more value as a picture of contemporary manners than as an indictment of the morals of Theodora. It is difficult to believe that if her girlhood had been so steeped in vice and infamy as this scandalous document asserts, she could have so completely changed as to develop into a matron whose conjugal chastity the same enemies could not seriously impugn, although they were ready to insinuate suspicions.\texthebrew{⁠}\footnote{Cp. \emph{H. A.} 16.11

\textgreek{\texthebrew{ὑ}ποψίας δ\texthebrew{ὲ} συμπεσούσης α\texthebrew{ὐ}τ\texthebrew{ῇ ἐ}ρωτολήπτ\texthebrew{ῳ} ε\texthebrew{ἶ}ναι ε\texthebrew{ἰ}ς τ\texthebrew{ῶ}ν ο\texthebrew{ἰ}κετ\texthebrew{ῶ}ν \texthebrew{ἕ}να \texthebrew{Ἀ}ρεόβινδον \texthebrew{ὄ}νομα}. The supplement is Haury's. Diehl observes that it is not recorded that any personal taunts were levelled at Theodora during the Nika revolt (p44); but this is not quite true (see \emph{Chron. Pasch., sub} 532 \textgreek{τ\texthebrew{ὰ}ς \texthebrew{ὑ}βριστικ\texthebrew{ὰ}ς φων\texthebrew{ὰ}ς \texthebrew{ἃ}ς \texthebrew{ἔ}λεγον . . . ε\texthebrew{ἰ}ς τ\texthebrew{ὴ}ν α\texthebrew{ὔ}γουσταν Θεοδώραν}).

} But it would be foolish to argue that the framework of the story is entirely fictitious. Theodora may have been the daughter of a bear-keeper , and she may have appeared on the stage. And her youth may have been stormy; we know that she was the mother of an illegitimate child.

After the rise in his fortunes through the accession of his uncle, Justinian seems to have secured for his mistress the rank of a patrician.\texthebrew{⁠}\footnote{\emph{H. A.} 9.30 . If this was so, the law of Justin relaxing the rule which forbade senators to marry actresses ( \emph{C. J.} V.4.23 ; 520\texthebrew{‑}523) was not required, as has been supposed, for the purpose of making Justinian's marriage possible. Cp. Panchenko, \emph{op. cit.} 74. John of Ephesus refers to Theodora's activity in the matter of a Monophysite deacon, while she was still a patrician, but this was probably after her marriage (\emph{Comm.} p68).

} He wished to marry her, but the Empress Euphemia resolutely opposed this step, and it was not till after her death\texthebrew{⁠}\footnote{\emph{H. A.} 9.47 . The year is not known, but she died before Justin.

} that Theodora became the wife of Justinian. When he was raised to the throne, she was, as a matter of course, crowned Augusta.

Her beauty and charm were generally acknowledged. We may imagine her as a small pale brunette, with a delicate oval face and a solemn intense expression in her large black eyes.\texthebrew{⁠}\footnote{\emph{H. A.} 10.11

\textgreek{ε\texthebrew{ὐ}πρόσωπος κα\texthebrew{ὶ} ε\texthebrew{ὔ}χαρις \texthebrew{ἄ}λλως}, cp. \emph{Aed.} I.11.8 . Procopius describes her as \textgreek{κολοβός}, an uncomplimentary way of saying that she was \emph{petite.}

} Portraits of her are preserved in marble, in mosaics, and on ivory. There is a life-size bust of her at Milan, which was originally coloured; the tip of the nose is broken off, but the rest is well preserved, and we can see the attractiveness of her face.\texthebrew{⁠}\footnote{The identification is due to Delbrück, \emph{Porträts byz. Kais.}, whose arguments have convinced me. He proved in the first place by a very complete examination of the headgear of Empresses in the fifth and sixth centuries that the bust belongs to the sixth; and that it is Theodora's is demonstrated by a comparison with the mosaics and the ivories. It is in the archaeological museum of Castel

(p30)
Sforzesco. It is probably eastern work, and must have been set up at Milan either in 538, during the few months in which the town was in Imperial hands, or before 535.

} Then

p30
we have two ivory tablets representing her in imperial robes.\texthebrew{⁠}\footnote{These tablets (of which one is at the Bargello in Florence, the other at Vienna) seem to be leaves of the same diptych. Gräven thought that the lady was Amalasuntha, but the diadem, which Gothic royalties never wore, disproves this. For comparison we have a small portrait of Theodora on the consular diptych of Justin (A.D. 540) which is preserved at Berlin. There can be little doubt that Delbrück is right in his identification. On the Bargello tablet the Empress is standing with a sceptre in her left hand and a cruciger globe in her right. On the segmentum of her chlamys is a male bust with a sceptre in his left and the handkerchief ( mappa ) in his right. This points to the consular games, so that the presumption is that the tablet was associated with a consul­ship of Justinian, and this would date it to 528 or 533 or 534. On the Vienna tablet the Empress is enthroned, and on the chlamys is a female bust, which Delbrück suggests might be that of her niece Sophia, afterwards Empress.

} These three portraits show her probably as she was from the age of thirty to thirty-five . She is visibly an older woman in the mosaic picture in the church of S. Vitale at Ravenna (\emph{c.} A.D. 547), but the resemblance to the bust can be discerned in the shape of the face, in the mouth, and in the eyes. But the dominion which she exercised over Justinian was due more to her mental qualities than to her physical charms. A contemporary writer praises her as ``superior in intelligence to all the world,''\texthebrew{⁠}\footnote{John Lydus, \emph{De mag.} III.69 .

} and all that we know of her conduct as Empress shows that she was a woman of exceptional brain and courage. Her influence in the Emperor's counsels was publicly acknowledged in a way which had no precedent in the past. In a law which aimed at suppressing corruption in the appointment of provincial governors, the Emperor declared that in framing it ``we have taken as partner in our counsels our most pious consort given to us by God.''\texthebrew{⁠}\footnote{\emph{Nov.} 8 , A.D. 535.

} At the end of the law an oath of allegiance is prescribed. The official is to swear loyalty to ``our divine and pious despots, Justinian, and Theodora, the consort of his throne.'' But although Justinian's devotion to his wife prompted him to increase her dignity and authority in the eyes of the Empire in unusual ways, it would be a mistake to suppose that legally she possessed powers which former Empresses had not enjoyed or that she was co-regent in the constitutional sense.\texthebrew{⁠}\footnote{Compare the remarks of Panchenko, \emph{op. cit.} 74\texthebrew{‑}76.

} Custom was strained to permit her unusual privileges. For instance, she is said to have received foreign envoys and presented them with gifts ``as if the Roman Empire were under

p31
her rule.''\texthebrew{⁠}\footnote{Procopius, \emph{H. A.} 30.24 .

} Chosroes was amazed when his minister Zabergan showed him a letter which he had received from Theodora urging him to press his master to make peace.\texthebrew{⁠}\footnote{\emph{Ib.} II.32 \emph{sqq.} Theodora had known Zabergan when he had come as envoy to Constantinople.

} Such incidents might well give the impression that the Empire was ruled by two co-equal sovrans, and some thought that Theodora had greater power than Justinian himself.\texthebrew{⁠}\footnote{The view is expressed by Zonaras, XIV.6.5\texthebrew{‑}6. It is highly remarkable that no coins were issued with Theodora's name and face, an honour which had been accorded to all Augustae until Justin's reign (there are no coins of Euphemia, see Wroth, \emph{Imp. Byz. Coins}, XIV note 4). In the reign following, Justin and Sophia appear together in many issues.

} Such power as she possessed she owed to her personal influence over her husband and to his toleration of her intervention in public affairs.

She was not indeed content to pursue her aims merely by the legitimate means of persuading the monarch. When she knew that he had resolutely determined on a line of policy which was not in accordance with her own wishes, she did not scruple to act independently. The most important matter in which their views diverged was ecclesiastical policy. Theodora was a devoted Monophysite, and one of her constant preoccupations was to promote the Monophysitic doctrine and to protect its adherents from the penal consequences which they incurred under Justinian's laws. Her husband must have been well aware that she had an intelligence department of her own and that secret intrigues were carried on of which he would not have approved. But she was clever enough to calculate just how far she could go.

Her power of engaging in independent political action was due to her economic independence. She had large financial resources at her disposal, for which apparently she had to render no account. The personal expenses of an Emperor's consort and the maintenance of her household were provided by estates in Asia Minor which were managed by a high steward known as the Curator of the House of Augusta,\texthebrew{⁠}\footnote{Curator divinae domus serenissimae Augustae,

\emph{C. J.} VII.37.3 .

} who was responsible to her. Justinian appears to have increased these estates considerably for the benefit of Theodora.\texthebrew{⁠}\footnote{For her estates in Cappadocia, yielding a revenue of 50 lbs. of gold (over £2250), see

\emph{Nov.} 30.6 ; in Helenopontus,

\emph{Nov.} 28.5 ; in Paphlagonia,

\emph{Nov.} 29.4 .

} He gave her large donations on the occasion of her marriage.\texthebrew{⁠}\footnote{\emph{C. J.} \emph{ib.}

Procopius says that he lavished on her large sums of money before his marriage,

\emph{H. A.} 9.31 .

} The house known

p32
as the palace of Hormisdas, in which Justinian had resided before his elevation to the throne, was enlarged and enclosed within the precincts of the Great Palace, and placed at the disposal of the Empress.\footnote{Procopius, \emph{Aed.} I. 4.1

and

10.4 ; John Eph. \emph{Comm.} c42.

}

Theodora did much to deserve the reputation of a beneficent queen, always ready to use her influence for redressing wrongs,\texthebrew{⁠}\footnote{John Lydus, \emph{De mag.} III.69 .

} and particularly solicitous to assist the unhappy of her own sex.\texthebrew{⁠}\footnote{Procopius, \emph{B. G.} III.32 \textgreek{\texthebrew{ἐ}πεφύκει γ\texthebrew{ὰ}ρ \texthebrew{ἀ}ε\texthebrew{ὶ} δυστυχούσαις γυναιξ\texthebrew{ὶ} προχρε\texthebrew{ῖ}ν.}} To her initiative are ascribed the stringent laws which were passed to suppress the traffic in young girls, which flourished as actively then as in modern Europe, and was conducted by similar methods, which the legislator graphically describes. Agents used to travel through the provinces to entice to the capital poor girls, sometimes under ten years of age, by the bait of fine clothes and an easy life. Indigent parents were easily persuaded by a few gold coins to consent to the ruin of their daughters. The victims, when they came to the city, were fed and clothed miserably, and kept shut up in the houses of ill-fame, and they were forced to sign written contracts with their infamous masters. Sometimes compassionate patrons of these establishments offered to deliver one of these slaves from her misery by marrying her, but the procurers generally refused to consent. The new edict forbade the trade and ordered that all procurers should be banished from Constantinople.\texthebrew{⁠}\footnote{\emph{Nov.} 14 (A.D. 535), addressed to the people of Constantinople.

} The principle of compensation, however, seems to have been applied. The patrons were allowed to state on oath how much money they had given to the parents of each girl; the average price was five nomismata, and Theodora paid the total out of her private purse.\texthebrew{⁠}\footnote{This is related by John Mal. XVIII.440\texthebrew{‑}441, as if it occurred in A.D. 529. There is no reason to question the date. The words \textgreek{κελεύσασα} (Theodora) \textgreek{το\texthebrew{ῦ} λοιπο\texthebrew{ῦ} μ\texthebrew{ὴ} ε\texthebrew{ἶ}ναι πορνοβοσκο\texthebrew{ὺ}ς} need not be an anticipatory reference to \emph{Nov.} 14, but may refer to measures taken by the Prefect of the City. A recrudescence of the forbidden practices was inevitable, and may have necessitated the legislation of 535. A law of Leo I against the traffic will be found in

\emph{C. J.} XI.41.7 .

} To receive unfortunate women who abandoned a life of shame, a palace on the Asiatic shore of the Bosphorus, not far from the Black Sea, was converted into a convent which was known as Metanoia or Repentance.\footnote{Procopius, \emph{Aed.} I.9.5\texthebrew{‑}10 , and

\emph{H. A.} 17.5 . In the latter passage the author represents the action of Theodora as tyrannical. ``She collected in the middle of the Forum more than 500 prostitutes who made

(p33)
a bare living, and sending them across to the convent called Metanoia she shut them up and forced them to change their way of life. Some of them threw themselves down from a height (the roof or a high window) and in this way escaped the compulsory change.''

}

p33
Theodora was perhaps too eager to interfere, as a sort of beneficent providence, in the private affairs of individual persons, and her offices were not always appreciated. She is said to have forced two sisters, who belonged to an old senatorial family and had lost their husbands, to marry against their will vulgar men who were utterly unworthy of them.\texthebrew{⁠}\footnote{\emph{H. A.} 7.7 \emph{sqq.}

Cp. the case of Saturninus, son of the mag. off. Hermogenes, \emph{ib.} 32 \emph{sqq.} Her interference in the domestic affairs of Belisarius, of which Procopius knows so much, is assuredly not entirely invented. In the case of Artabanes and Praejecta she had a locus standi, as Praejecta was Justinian's niece (\emph{B. G.} III.31, see below,
p146 ).

} And her enemies alleged that in her readiness to espouse the cause of women she committed grave acts of injustice and did considerable harm. Wives who were divorced for adultery used to appeal to the Empress and bring accusations against their husbands, and she always took their part and compelled the unfortunate men to pay double the dowry, if she did not cause them to be whipped and thrown into prison. The result was that men put up with the infidelity of their wives rather than run such risks.\texthebrew{⁠}\footnote{\emph{Ib.} 17.24 \emph{sqq.}

} It is impossible to decide how much truth there may be in these charges, but they illustrate Theodora's desire to be the protectress and champion of her own sex.

There can be little doubt that the Empress used her position to exercise a patronage in appointments to offices, which was not always in the public interest, and that she had few scruples in elevating her favourites and disgra­cing men who displeased her.\texthebrew{⁠}\footnote{The case of Peter Barsymes is a notable example,

\emph{H. A.} 22.22 \emph{sqq.}

She is said to have intended to create Theodosius, the lover of Antonina, a mag. mil., \emph{ib.} III.19. Peter the Patrician may have owed his promotion to the post of mag. off. to her favour (see below,
p166 ).

} It must, however, be confessed that in the two cases in which we have good evidence that she intervened to ruin officials, her intervention was beneficial. Thus she procured the disgrace of an Imperial secretary named Priscus, an unprincipled man who had grown rich at the public expense. He was alleged to have spoken against her, and as she could not prevail on Justinian to take action, she caused the man to be put on board a ship and transported to Cyzicus, where he was tonsured. Justinian acquiesced in the accomplished fact and confiscated his property.\footnote{The details are told in

\emph{H. A.} 16.7\texthebrew{‑}10 , but the main fact is confirmed by John Mal. XVIII.449 and \emph{De ins.}

(p34)
\emph{fr.} 45. The other more important case is that of John of Cappadocia, see below,
p57 . Haury is certainly right in supposing that in describing Priscus as \textgreek{Παφλαγών}, Procopius does not mean that he was a Paphlagonian, but is alluding to Cleon, ``the Paphlagonian'' of the \emph{Knights} of Aristophanes (\emph{B. Z.} IX.674).

}

p34
Procopius, in his \emph{Secret History}, has several stories to tell of cruel punishments which she inflicted privately on persons who had offended her. Lurid tales were whispered of the terrible secret dungeons of her palace in which men disappeared for ever,\texthebrew{⁠}\footnote{Private prisons were forbidden by law,

\emph{C. J.} IX.5.2

(529).

} and the known fact that she had the means of maintaining heretics in concealment for years made gossip of this kind appear credible. Whatever may be the truth about her alleged vengefulness and cruelty, it is certain that she was feared.

There was no disguise in the attitude which she assumed as head of the ecclesiastical opposition to Justinian's policy, and he must have been fully aware that secret intrigues were carried on which he would not have sanctioned. It seemed indeed difficult to believe that a man of his autocratic ideas would have tolerated an independent power beside his own; and the theory was put forward that this apparent discord between their aims and views was a political artifice deliberately planned to blind their subjects, and to facilitate the transactions which the emperors could not openly permit.\texthebrew{⁠}\footnote{\emph{H. A.} 10.13 \emph{sqq.}

Theodora's active partisan­ship for the Blue Faction, of which Justinian professed to disapprove, is given as an instance.

} This theory may contain a small measure of truth so far as ecclesiastical policy is concerned. It may have been convenient to the Emperor to allow the severities which his policy forced him to adopt against the Monophysites to be mitigated by the clandestine and illegal protection which the Empress afforded to them. But otherwise the theory can hardly be entertained seriously. We can only regard the latitude which was allowed to Theodora as due to Justinian's weakness. And she was clever enough to know how far she could venture.

Her habits presented contrast to the temperance and simplicity of Justinian. She spent a long time in her bath. At her meals she indulged in every kind of food and drink. She slept long both at night and in her daily siesta.\texthebrew{⁠}\footnote{\emph{H. A.} 15.6 \emph{sqq.}

} She spent many months of the year in the suburban palaces on the seashore, especially at Hêrion (on the coast of Bithynia, opposite the Islands of the Princes), which Justinian enlarged and

p35
improved.\texthebrew{⁠}\footnote{\emph{Ib.}, where it is said that the court suffered much discomfort at Hêrion. For the reconstruction of the palace see \emph{Aed.} I.11.16\texthebrew{‑}22.

} Sometimes she visited the hot springs of Pythia (in Bithynia), where Justinian also built an Imperial residence. On these occasions she was attended by an immense retinue of patricians and chamberlains.\texthebrew{⁠}\footnote{Four thousand on the occasion of her visit in 529, acc. to John Mal. XVIII.441. For the palace see \emph{Aed.} V.3, 16\texthebrew{‑}20.

} For Theodora had all a parvenue's love of pomp and show, and she was probably encouraged by the Emperor, who, though simple in his own tastes, thought much of public splendour and elaborate ceremonial as a means of enhancing the Imperial majesty. We are told that new and abasing forms of etiquette were introduced at court. When the Senate appeared in the presence of the Emperor, it had been the custom for one of the patricians to kiss the sovran on the right breast, and the sovran replied to the salutation by kissing the head of the patrician. No corresponding ceremony was practised in the case of the Empress. Under Justinian and Theodora it became obligatory that all persons, of whatever rank, should prostrate themselves on entering the presence of the Emperor and of the Empress alike. The spirit of oriental servility in the Palace was shown by the fact that officials and members of the court who, in talking among themselves used to speak of the sovrans as the ``Emperor'' and the ``Empress'' ( Basileus and Basilis ), now began to designate them as the ``Lord'' and the ``Lady'' ( Despotes and Despoina ) and described themselves as their slaves; and any one who did not adopt these forms was considered to have committed an unpardonable solecism.\footnote{\emph{H. A.} 30.21\texthebrew{‑}26 .

}

It is not improbable that, if Justinian had wedded a daughter of one of the senatorial families, many people would have been happier, and the atmosphere of the Palace would have been less dangerously charged with suspicion and intrigue. But, if Theodora was greedy of power and often unscrupulous in her methods, her energy and determination on one occasion rescued the throne, and on another rendered a signal service to the community. And there is no reason to suppose that in her conduct generally she was not honestly convinced that, if she employed irregular means, she was acting in the true interests of the State.

The brilliancy of Justinian's reign did not bring happiness or contentment to his subjects. His determination to increase the power of the throne and retain his government more completely in his own hands caused dissatisfaction in the senatorial circles and inevitably led to tyranny; and his ambitious plans of expansion involved expenses that could only be met by increasing the financial burdens which already weighed too heavily on the people.

The frugal policy of Anastasius had bequeathed to his successor a reserve of 320,000 lbs. of gold (about 14½ millions sterling). In the reign of Justin these savings were dissipated, as well as a further amount of 400,000 lbs. which had come into the treasury in addition to the regular revenue.\texthebrew{⁠}\footnote{John Lyd. \emph{De mag.} III.51 ;

Procopius, \emph{H. A.} 19.7\texthebrew{‑}8

(\textgreek{ο\texthebrew{ὐ}δεν\texthebrew{ὶ} νόμ\texthebrew{ῳ}} seems to mean ``irregularly'').

} A heavy tax on the exchequer was caused by the terrible earthquake of May A.D. 526, which laid the city of Antioch in ruins and destroyed, it is said, 250,000 people.\texthebrew{⁠}\footnote{John Lyd. \emph{ib.} 54 . The details of this disaster are vividly described by the Antiochene writer John Mal. XVII.419 \emph{sqq.} The Emperor showed his sorrow by appearing in St. Sophia without his diadem. Another serious earthquake befell Antioch at the end of 529, \emph{ib.} 442.

} In the following year war broke out with Persia, and when Justinian came to the throne, the financial position was not such as to justify any extraordinary enterprises. It is asserted by a civil servant who had a long career in the office of the Praetorian Prefect of the East, that the unfavourable financial situation was chiefly caused by the incompetence of those who had held the Prefecture in the reign of Justin.\texthebrew{⁠}\footnote{John Lyd. III.51 .

} Justinian after some time found a man for the post who knew how to fill the treasury.

John, a native of Caesarea in Cappadocia, began as a clerk in the office of a Master of Soldiers. In this capacity he became, by some chance, known to Justinian, and he was promoted to the post of logothete, a name which had now come into general use for those responsible officials who, under the Praetorian Prefect, controlled the operations of the subordinate assessors and collectors of taxes in the provinces.\texthebrew{⁠}\footnote{The office of logothete is discussed at length by Panchenko, \emph{op. cit.} 106 \emph{sqq.} Stein has shown that it is probably the Greek equivalent of scriniarius (cp. above,
Vol. I Chap. XIII § 3, p442 ).

} In the case of Marinus,

p37
this post had been a stepping-stone to the Prefecture itself, and John had the same luck. He was first raised to the rank of an illustris, and became Praetorian Prefect before A.D. 531.\texthebrew{⁠}\footnote{John Lyd. \emph{ib.} 57 . John is Pr. Pr. on April 30, 531

(\emph{C. J.} VI.27.5) , Julian on Feb. 20 (\emph{ib.} III.1.16).

} He had not the qualifications which might have been thought indispensable for the duties of this ministry, for he had not received a liberal education, and could barely read and write; but he had the qualification which was most essential in the eyes of the Emperor, talent and resource ­fulness in raising money. His physical strength and energy were enormous, and in difficulties he was never at a loss.\texthebrew{⁠}\footnote{The sources for John's character are Procopius, \emph{B. P.} I.24 and John Lyd. \emph{ib.} 57 \emph{sqq.} (where we get the details). The two pictures agree.

} He is described as the boldest and cleverest man of his time.\texthebrew{⁠}\footnote{Procopius, \emph{B. V.} I.10.7.

} But he was absolutely unscrupulous in his methods, and while he supplied the Emperor with the funds which he required, he also became himself enormously rich and spent his money on gluttony and debauchery. ``He did not fear God, nor regard man.'' The provinces of Lydia and Cilicia were a conspicuous scene of his operations. He procured the appointment of another Cappadocian, also named John, to the governor ­ship of Lydia \texthebrew{—} a man after his own heart, enormously fat and popularly known as Maxilloplumacius ( Flabby-jaw ). With the help of this lieutenant, the Prefect ruined Lydia and its capital, Philadelphia. He visited the province himself, and we are told that when he had done with it, he had left not a vessel in a house, nor a wife, a virgin, or a youth unviolated. The exaggeration is pardonable, for our informant was born at Philadelphia. The same writer gives particular instances \texthebrew{—} some of which had come under his own observation \texthebrew{—} of the violent means to which John the Cappadocian resorted to extort money from rich persons. He had dark dungeons in the Prefect's residence, and he made use of torture and painful fetters.\footnote{John Lydus can find no terms too strong to describe the Prefect's cruelty and luxury (Phalaris, Cyclops, Briareus, Sardanapalus, etc.). For his debauchery see cc62 , 65 . As John Lydus was an official in the Prefecture, his testimony is valuable, though on the other hand he may have had private reasons of animosity (cp. c96 ), and his respect for the traditions of the office made a rude uneducated Prefect, who introduced practical innovations in the conduct of business, repellent to him. It seems certain that Maxilloplumacius was governor of Lydia, for he must be meant by \textgreek{τ\texthebrew{ὸ}ν \texthebrew{ὁ}μόγνιον κα\texthebrew{ὶ ὁ}μόψυχον τ\texthebrew{ῆ}ς α\texthebrew{ὐ}το\texthebrew{ῦ} βδελύριας \texthebrew{ὕ}παρχον} (that is, \textgreek{\texthebrew{ἔ}παρχον}) c61.

}

While contemporary writers agree in painting John as a

p38
coarse monster, without a single redeeming quality, we must make some allowance for exaggeration. It is unlikely that he would have enjoyed so long the confidence of the Emperor if his sole recommendation had been skill in plundering the provinces. As a matter of fact, we shall see that during his second tenure of the Prefecture, which lasted about nine years, a series of provincial reforms was carried through which intimately concerned his own sphere of administration and in some respects diminished his power. This could not have been done without his co-operation , and we cannot fairly withhold from him part of whatever credit the legislation deserves. We may conjecture that he won and retained his influence over the Emperor, not only through his success in replenishing the treasury, but also partly through his independence, which was displayed when he openly opposed the project of conquering Africa, and partly through the fact that he was not hampered by conservative prejudices. It was chiefly his indifference to the traditions of the civil service that made him unpopular among the officials of the Prefecture.

Besides increasing the revenue by fair means and foul, John had recourse to economies which were stigmatised by contemporary opinion as injurious to the public interest. He cut off or reduced the service of the State post, with the exception of the main line to the Persian frontier. The post from Chalcedon to Dakibiza was abolished, and replaced by a service of boats to Helenopolis, while in southern Asia Minor and Syria asses were substituted for horses and the speed of travelling was diminished. The results were twofold. The news of disasters in the provinces, which demanded prompt action, was slow in reaching Constantinople. More serious was the consequence for the farmers in the inland provinces, who, deprived of the public means of transport, were obliged to provide for the transmission of their produce to the ports to be conveyed to the capital. Large quantities of corn º rotted in the granaries; the husbandmen were impoverished; and the Prefect's officials pressed for payment of the taxes in gold.\texthebrew{⁠}\footnote{Procopius, \emph{H. A.} 30.8\texthebrew{‑}11 . John Lyd. \emph{ib.} 61 .

} Multitudes of destitute people left their homes and went to Constantinople.\footnote{John Lyd. \emph{ib.} 70 . John had attempted to make a breach between Justinian and Theodora, Procopius, \emph{B. P.} I.25.4;

\emph{H. A.} 17.38 .

}

p39
Justinian was well satisfied with the fruits of John's administration, and only too ready to shut his eyes to the methods by which the funds he needed were procured. How far he was really innocent it is impossible to determine, but we are assured that the ministers and courtiers always praised the Prefect to the Emperor, even though they had personal grievances against him. At length Theodora, who disliked the Cappadocian and was well acquainted with his iniquities, endeavoured to open Justinian's eyes and to show him that, if the tyrannical administration were allowed to continue, his own position would be endangered. If her arguments produced any effect on his mind, he wavered and postponed action\texthebrew{⁠}\footnote{John Lydus (\emph{ib.} 69) , who says that Justinian was deterred from making a change because John had deliberately, in order to ensure his permanent occupation of the post, introduced such confusion into the book-keeping that a successor would have found it almost impossible to carry on the administration.

\texthebrew{◂} previous

next \texthebrew{▸}Images with borders lead to more information.

The thicker the border, the more information.

(Details here.)

UP TO:

J. B. Bury
Homepage

Latin \& Greek

Texts

LacusCurtius

Home} until action was suddenly forced on him by a revolutionary outbreak which well-nigh cost him his throne.

\xchapter{Chapter XV, Part B}

Short URL for this page:

bit.ly/BURLAT15B

\xsubsection{(Part 2 of 3)}

The famous rising at Constantinople, which occurred in the first month of A.D. 532 and wrecked the city, was the result of widely prevailing discontent with the administration, but it began with a riot of the Hippodrome factions which in ordinary circumstances would have been easily suppressed.\texthebrew{⁠}\footnote{The contemporary sources are

Procopius, \emph{B. P.} I.24 ; John Mal. XVIII.473 \emph{sqq.} and \emph{fr.} 46, \emph{De ins.}; John Lyd. \emph{De mag.} III.70 ; Marcellinus, \emph{Chron., sub a.} (also notices, not very important, in Zacharias Myt. IX.4, Victor Tonn. \emph{sub a.} Theodore Lector in Cramer, \emph{Anecd. Par.} II.112). The narrative of Malalas was originally much fuller than in the abbreviated text which we possess, but the missing parts can be largely supplied from \emph{Chron. Pasch.} and Theophanes, where we find many details for which it can be shown that Malalas was the source, as well as some derived from elsewhere (cp. the analysis in Bury, \emph{The Nika Riot}, 98 \emph{sqq.}). The short notice of Marcellinus is important. The Illyrian Marcellinus had been the cancellarius in the official staff of Justinian when he was Master of Soldiers. Before 527 he retired and became a priest. His chronicle, in its first form, reached the year 518, but afterwards he brought it down to 534. His personal relations to Justinian make it probable that his account of the revolt represents it in the light in which the court desired it to be viewed. The popular dissatisfaction which made the rising possible is entirely ignored, and the whole movement is explained as a conspiracy organised by the nephews of Anastasius, who appear as the prime movers. It must, of course, be remembered that when Marcellinus wrote, \emph{c.} 534, it was impossible to refer to the oppressions of John the Cappadocian, who was then in office. John Lydus ignores the conspiracy and the elevation of Hypatius, and represents the unpopularity of the Prefect as the sole cause of

(p40)
the revolt. For modern studies of the Nika see Bibliography, under

Schmidt, W. A. , and

Kalligas . Cp. Hodgkin, \emph{Italy}, III.618 \emph{sqq.} ; Ranke, \emph{Weltgeschichte}, II.2, pp23 \emph{sqq.}; Diehl, \emph{Justinien}, 462 \emph{sqq.}

} We saw

p40
how Justinian in his uncle's reign patronised the Blues and made use of them as a political support. But when he was safely seated on the throne, he resolved no longer to tolerate the licence of the factions, from the consequences of which he had formerly protected the Blues. Immediately after his accession he laid injunctions on the authorities in every city that the disorders and crimes of the factions should be punished impartially.\footnote{John Mal. XVII.416.

}

A number of persons belonging to both factions had been arrested for a riot in which there had been loss of life. Eudaemon, the Prefect of the City, held an inquiry, and finding seven of the prisoners guilty of murder, he condemned four to be beheaded, and three to be hanged. But in the case of two the hangman blundered and twice the bodies fell, still alive, to the ground. Then the monks of St. Conon, which was close to the place of execution, interfered, and taking up the two criminals, one of whom was a Blue, the other a Green, put them in a boat and rowed them across the Golden Horn to the asylum of the church of St. Laurentius.\texthebrew{⁠}\footnote{St. Conon was in the 13th Region, across the Golden Horn. St. Laurentius was close to Blachernae.

} The Prefect, on hearing what had occurred, sent soldiers to guard the church.\footnote{Jan. 11 (substitution). This was perhaps the day of the curious scene which occurred between the Emperor and the Greens, and is recorded by Theophanes (cp. \emph{Chron. Pasch.} ), if the chronicler is right in associating it with the Nika riot. But it is possible that the chronicler is mistaken, and in any case the subject of the conversation has no apparent connexion with the sedition. I have therefore transferred it to an Appendix (below,

p71 ).

}

The ides of January fell three days later (Tuesday), and, according to custom, horse races were held in the Hippodrome, and the Emperor was present. Both the Blues and the Greens importuned the Emperor with loud prayers to show mercy to the two culprits who had been rescued by accident from the gallows. No answer was accorded, and at the twenty-second race the spectators were amazed to hear the unexpected exclamation, ``Long live the humane Greens and Blues!'' The cry announced that the two parties would act in concert to force the government to grant a pardon, and it is probable that their leaders had previously arranged to co-operate . When the races were over, the factions agreed on a watchword, nika , ``conquer,'' and the rising which followed was known as the

p41
Nika Revolt. The united factions were known for the time as the Green-Blues ( Prasino-venetoi ).

In the evening the mob of rioters assembled at the Praetorium and demanded from the Prefect of the City what he intended to do with the refugees in St. Laurentius. No answer was given, and the rioters broke into the prison, released the criminals who were confined in it, killed some of the officials, and set fire to the building, which was partly burned. Elated by success they rushed eastward to the Augusteum and committed graver acts of incendiarism. They fired the Chalke, the entrance of the Great Palace, and not only was this consumed, but the flames spread northward to the Senate-house and the church of St. Sophia. These buildings were burned down.\footnote{On the order of the different conflagrations during the sedition see Bury, \emph{op. cit.} 114 \emph{sqq.}

}

On the following morning (Wednesday, January 14) the Emperor ordered the races to be renewed. But the Blues and Greens were not in the humour for witnessing races. They set on fire the buildings at the northern end of the Hippodrome, and the conflagration destroyed the neighbouring baths of Zeuxippus with the portico of the Augusteum. It is probable that on this occasion Justinian did not appear in the Kathisma, or face the multitudes who were now clamouring in the Hippodrome, no longer interceding for the lives of the two wretches who had escaped the hangman, but demanding that three unpopular ministers should be deprived of their offices. The demonstration was directed against Eudaemon, Prefect of the City, Tribonian the Quaestor, and John of Cappadocia, and the situation had become so serious that the Emperor decided to yield.\texthebrew{⁠}\footnote{The demands were reported to the Emperor by Basilides, Mundus, and Constantiolus, who had issued from the Palace and addressed the multitude, asking them what they wanted. \emph{Chron. Pasch.} ; cp. John Mal. 475 .

} Tryphon was appointed Prefect of the City, Basilides Quaestor, and Phocas, a man of the highest probity, was persuaded to undertake the office of Praetorian Prefect.

These concessions would probably have satisfied the factions and ended the trouble, like similar concessions in previous reigns, if the decision had depended solely on the leaders of the Blues and Greens. But the movement now wore an aspect totally different from that of the previous day. We saw how the city had been filled by throngs of miserable country folk

p42
from the provinces who had been ruined by the fiscal administration of the Praetorian Prefect and were naturally animated by bitter resentment against the Emperor and the government. It was inevitable that they should take part in the disturbances;\texthebrew{⁠}\footnote{This is brought out in the account of John Lydus. The agitation against John marks the change from the faction riot to the popular sedition. The quarrel of the Blues was with the Prefect of the City. John had always posed demonstratively as a lover of the Blue party (John Lyd. III.62 \emph{ad init.} ; cp. Zach. Myt. IX.14).

} it was at least a good opportunity to compass the fall of the detested Cappadocian; and the riot thus assured the character and proportions of a popular rising. But there were other forces in the background, forces which aimed not merely at a reform of the administration, but at a change in the dynasty. The policy of Justinian in seeking to make his power completely independent of the Senate and the Imperial Council had caused deep animosity in the senatorial class, and the disaffected senators seized the opportunity to direct the rising against the throne.\texthebrew{⁠}\footnote{Marcellinus, \emph{s. a.}: iam plerisque nobilium coniuratis, and the fact comes out in the narrative of Procopius. The conspirators took care that the people should be supplied with arms (\emph{ib.}).

} We must attribute to the secret agitation of these men and their agents the fact that the removal of the obnoxious ministers, especially of John, failed to pacify the people.

The plan was to set on the throne one of the nephews of Anastasius, unfortunate victims of their kinship to an Emperor. For we must acquit them of any ambitious designs of their own. They had been well treated by Justin and Justinian, and their only desire was to live in peace. Pompeius and Hypatius were out of the reach of the insurgents; they and many other senators were with Justinian in the Palace. It was therefore decided to proclaim Probus Emperor, and the mob rushed to his house. But they did not find him, for, fearing what might happen, he had left the city. In their angry disappointment they burned his house.

It was assuredly high time for the Emperor to employ military force to restore order. But the Palace guards, the Scholarians and Excubitors, were unwilling to do anything to defend the throne. They had no feeling of personal devotion to Justinian, and they decided to do nothing and await events. Fortunately for Justinian there happened to be troops of a more irregular kind in the city, and two loyal and experienced commanders.

p43
Belisarius, who as Master of Soldiers in the East had been conducting the war against Persia, had recently been recalled, and he had in his service a considerable body of armed retainers, chiefly of Gothic\texthebrew{⁠}\footnote{John Mal. 475 \textgreek{μετ\texthebrew{ὰ} πλήθους Γοτθικο\texthebrew{ῦ}} (cp. Procopius, \emph{B. P.} I.24.40 ). For the practice of keeping private bands of retainers see above,

Chap. ii § 2, p43 .

} race. Mundus, a general who had done good service in the defence of the Danube, was also in the capital with a force of Heruls. But all the soldiers on whom the Emperor could count can hardly have reached the number of 1500.

It was perhaps on Thursday (January 15)\texthebrew{⁠}\footnote{Or perhaps on Wednesday. See Bury, \emph{op. cit.} 107.

} that Belisarius rode forth at the head of Goths and Heruls to suppress the revolution. There was a battle, possibly in the Augusteum; many were killed; but the soldiers were too few to win a decisive victory, and the attack only exasperated the people. The clergy, it may be noted, seem to have made some vain attempts to restore order.\footnote{This comes from Zonaras, XIV.6.14, who also mentions that women took part in the disorders (\emph{ib.} 16). On his source cp. Sauerbrei, \emph{De font. Zon. quaest.} 77.

}

During the two following days there was desultory street fighting, and another series of conflagrations. On Friday the mob again set fire to the Praetorium, which had only been partly damaged, but there was a strong north wind which blew the flames away from the building. They also set fire to the baths of Alexander, and the same wind carried the conflagration to the neighbouring hospice of Eubulus and hence to the church of St. Irene and the hospice of Sampson. On Saturday there was a conflict between the soldiers and the insurgents in the street which led northward from Middle Street to the Basilica and the quarter of Chalkoprateia.\texthebrew{⁠}\footnote{On the topography see Bury, \emph{op. cit.} 11 \emph{sqq.} Beliaev, \emph{Khram Bog. Khalk.}; Krasnosel'tsev, \emph{Zamietka po voprosu o miestopol. Khalk. Khrama} (Lietopis ist.-phil. obshchestva of Odessa, Viz. otd. II. 1894), 309 \emph{sqq.}

} It would appear that some of the mob had occupied the Octagon, a building close to the Basilica, and the soldiers set it on fire. The same fatal north wind was blowing, and the flames, wafted southwards, spread to the church of St. Theodore Sphoracius and to the palace of Lausus, which was consumed with all its treasures, and thence raged along Middle Street, in the direction of the Forum of Constantine, destroying the colonnades and the church of St. Aquilina.\texthebrew{⁠}\footnote{Theophanes connects the burning of the Praetorium with this conflagration, but see Bury, \emph{op. cit.} 116\texthebrew{‑}117.

} We can imagine how great must have been the alarm

p44
in the Palace, which was almost in a state of siege.\texthebrew{⁠}\footnote{It may be conjectured that a shortage of food and water was feared in the Palace; for after the suppression of the revolt Justinian constructed cisterns and a granary close to the Palace, in order to have supplies in emergencies, John Mal. 477.

} Justinian could not trust his guards, and he had strong and not unjustified suspicions that many of the senators who surrounded him were traitors. Fearing their treachery, he ordered them all to leave the Palace on Saturday at nightfall,\texthebrew{⁠}\footnote{Procopius, \emph{ib.} 40 . \emph{Chron. Pasch.} places the dismissal of the senators \emph{after} Justinian's appearance in the Hippodrome on Sunday. See Bury, 108.

} except a few like the Cappadocian, whose loyalty was certain or whose interests were bound up with his own. He particularly suspected Hypatius and Pompeius, and when they protested against deserting him, his suspicions only grew stronger, and he committed the blunder of dismissing them.

On Sunday morning (January 18) the Emperor made an effort in person to pacify the people. He appeared in the Kathisma of the Hippodrome with a copy of the gospels in his hands, and a large crowd assembled. He swore on the holy book that he would grant an amnesty without any reservations and comply with the demands of his subjects. But the great part of the crowd was bitterly hostile. They cried, ``You are perjuring yourself,''\texthebrew{⁠}\footnote{\emph{Chron. Pasch.} \textgreek{\texthebrew{ἐ}πιορκε\texthebrew{ῖ}ς, σγαύδαρι}. It has been proposed to read \textgreek{γαύδαρι} and explain it as = \textgreek{γάδαρε}, ``ass.''

} and ``You would keep this oath to us as you kept your oath to vitalian.'' And there were shouts of ``Long live Hypatius!'' Meanwhile it had become known that the nephews of Anastasius had left the Palace. The people thronged to the house of Hypatius, and in spite of his own reluctance and the entreaties of his wife Maria, who cried that he was being taken to his death, carried him to the Forum of Constantine, where he was crowned with a golden chain wreathed like a diadem.

A council was then held by Hypatius and the senators who were supporting his cause.\texthebrew{⁠}\footnote{Procopius, \emph{ib.} § 25 , speaks as if all the senators, who were not in the Palace, were present.

} Here we can see clearly that the insurrection was guided and fomented by men of high position who were determined to overthrow Justinian. The question was debated whether the Palace should be attacked immediately. One of the senators, Origen, advised delay. He proposed that the new Emperor should occupy for the moment one of the

p45
smaller Imperial palaces\texthebrew{⁠}\footnote{Two are mentioned (\emph{ib.} § 30), Placillianae and Helena.

} and prosecute the war against his rival with deliberation, leaving nothing to chance. But his advice did not prevail, and Hypatius, who was himself in favour of prompt action, proceeded to the Hippodrome and was installed in the Kathisma. The insurgents crowded the huge building in dense masses, and reviled Justinian and Theodora.\footnote{\emph{Chron. Pasch.} ; cp. Cramer, \emph{Anecd. Par.} II.320.

}

In the meantime, another council was being held in the Palace. The situation seemed desperate. To many, including the Emperor himself, there seemed no resource but escape by sea. John the Cappadocian recommended flight to Heraclea, and Belisarius agreed. This course would have been adopted had it not been for the intervention of the Empress Theodora, whose indomitable courage mastered the wavering spirits of her husband and his councillors. A writer, who may well have heard the scene described by Belisarius himself, professes to reproduce her short speech, and even his sophisticated style hardly spoils the effect of her vigorous words:

The present occasion is, I think, too grave to take regard of the convention that it is not meet for a woman to speak among men. Those whose dearest interests are exposed to extreme danger are justified in thinking only of the wisest course of action. Now in my opinion, on the present occasion, if ever, flight is inexpedient even if it should bring us safety. It is impossible for a man, when he has come into the world, not to die; but for one who has reigned it is intolerable to be an exile. May I never exist without this purple robe and may I never live to see the day on which those who meet me shall not address me as ``Queen.''\texthebrew{⁠}\footnote{\textgreek{Δέσποινα.}} If you wish, O Emperor, to save yourself, there is no difficulty; we have ample funds. Yonder is the sea, and there are the ships. Yet reflect whether, when you have once escaped to a place of security, you will not prefer death to safety. \emph{I} agree with an old saying that ``Empire is a fair winding-sheet .''\texthebrew{⁠}\footnote{\textgreek{Καλ\texthebrew{ὸ}ν \texthebrew{ἐ}ντάφιον \texthebrew{ἡ} βασιλεία \texthebrew{ἐ}στί}, \emph{B. P.} \emph{ib.} 33\texthebrew{‑}38. The phrase comes from Isocrates, \emph{Archidamos}, § 45 \textgreek{καλόν \texthebrew{ἐ}στιν \texthebrew{ἐ}ντάφιον \texthebrew{ἡ} τυραννίς.}}

Theodora's dauntless energy communicated itself to her hearers, and they resolved to remain and fight.

In the Hippodrome it was believed that they had already fled. Hypatius, we are told, still doubtful of his chances of success, had secretly sent a message to the Palace, advising Justinian to attack the people crowded in the Hippodrome. Ephraem, the messenger, gave the message to Thomas, an Imperial

p46
secretary, who, ignorantly or designedly, informed him that Justinian had taken to flight.\texthebrew{⁠}\footnote{Thomas was a pagan and was possibly disloyal. The episode is related in \emph{Chron. Pasch.} and came from Malalas .

} Ephraem proclaimed the news in the Hippodrome, and Hypatius now played the Imperial part with confidence, but the people were soon undeceived.

Justinian sent out a trusted eunuch, named Narses, with a well-filled purse to sow dissensions and attempt to detach the Blue faction from the rebellion.\texthebrew{⁠}\footnote{John Mal. 476. Procopius does not mention this incident. Narses, a Persarmenian, was \textgreek{ταμίας τ\texthebrew{ῶ}ν βασιλικ\texthebrew{ῶ}ν χρημάτων} since A.D. 530 (Procopius,

\emph{B. P.} I.15.31 , cp. \emph{B. G.} II.13.16 ). º At this time he was not Praepositus s. cub., a post which he afterwards filled (between 540 and 552), see CIL VI.1199 . The same financial post was held by Rusticus in 554 (see

Agathias, III.2) , but I cannot agree with Stein that this was an extraordinary post, existent only in time of war, and officially named comes s. largitionum (\emph{Studien}, 163 \emph{sqq.}).

} He could insinuate that Hypatius, like his uncle, would be sure to protect their rivals the Greens, and remind them of the favour which Justinian had shown them in time past and of the unwavering goodwill of Theodora. While Narses fulfilled this mission, Belisarius and Mundus prepared to attack. At first Belisarius thought it would be feasible to reach the Kathisma directly from the Palace and pluck the tyrant from his throne. But the way lay through a building occupied by a portion of the guards, and they refused to let him pass.\texthebrew{⁠}\footnote{We have no clear knowledge as to the communications between the Palace and the Hippodrome in the sixth century. We have more information about the arrangements existing in the ninth and tenth, but there must have been considerable changes in the meantime. The Emperors always reached the Kathisma from the Daphne portion of the Palace by a winding stair, Kochlias (which must, however be distinguished from the Kochlias leading to a gate by which Mundus left the Palace, Procop. \emph{ib.} § 43 ). If Belisarius hoped to reach Hypatius by this way, we must suppose that there was a guard-room at the top of the staircase, in the Kathisma structure, for we know that the Daphne buildings, in which there was no room for guards, adjoined the walls of the Hippodrome. But there may have been at this time a direct communication between the part of the Hippodrome north of the Kathisma and the quarters of the guards north of Daphne; and this seems the most probable explanation.

} The Emperor then ordered him to lead his troops, as best he could, through the ruins of the Chalke into the Augusteum. With great difficulty, climbing through the debris of half-burnt buildings, they made their way round to the western entrance of the Hippodrome and stationed themselves just inside, at the portico of the Blues, which was immediately to the right of the Kathisma.\texthebrew{⁠}\footnote{Procopius, \emph{ib.} § 49 . See above,

Vol. I, Chap. III, p81 . This passage shows conclusively that there was a western entrance to the Hippodrome, close to the Kathisma. It is possible that the Nekra gate by which Mundus

entered was on the same side, near the south end, rather than on the eastern side, as has been generally supposed. We might infer from Procopius that when Belisarius and his troops forced their way into the Hippodrome, Mundus and his Heruls stationed themselves just outside the same (main) entrance (\textgreek{πλησίον που \texthebrew{ἐ}στηκ\texthebrew{ὼ}ς Μο\texthebrew{ῦ}νδος}), and when the fighting began went to the Nekra. On the other hand, Mundus seems to have been awaiting the action of Belisarius on the eastern side, for he left the Palace by a gate where there was a winding descent (Procop. \emph{ib.} § 43) \texthebrew{—} evidently an issue on the south of the Daphne, where there was a fall in the ground. It is therefore more natural to suppose that the Nekra, by which he entered the Hippodrome, was on the same side.

} In order

p47
to gain access to the Kathisma itself, it would have been necessary to pass though a small gate on the left, which was shut and guarded. If Belisarius attempted to force this gate, his men would have been exposed to an attack from the crowd in the rear. He therefore determined to charge the people. He drew his sword and gave the word. Though many of the populace had arms, there was no room in the dense throng to attempt an orderly resistance, and confronted by the band of disciplined soldiers the mob was intimidated and gave way. Moreover there were dissensions among them, for the bribes of Narses had not been fruitless. They were cut down without mercy, and then Mundus appeared with his Heruls to help Belisarius in the work of slaughter. Mundus had left the Palace by another way, and he now entered the Hippodrome by a gate known as Nekra. The insurgents were between two fires, and there was a great carnage. It was said that the number of the slain exceeded 30,000.\footnote{Procopius, ``more than 30,000,'' and John Mal. ``35,000 more or less,'' roughly agree. John Lyd. gives an exaggerated figure, 50,000; Zonaras, 40,000.

}

Two nephews of Justinian, Boraides and Justus, then entered the Kathisma without meeting resistance.\texthebrew{⁠}\footnote{Procopius, § 53 . Acc. to John Mal. 476 , it was Belisarius who arrested the two brothers.

} They seized Hypatius, who had witnessed the battle from his throne, and secured Pompeius, who was with him. The brothers were taken into the Palace, and, notwithstanding the tears of Pompeius and the pleadings of Hypatius that he had acted under compulsion,\texthebrew{⁠}\footnote{Procopius contrasts the weakness of Pompeius with the more dignified tone of Hypatius. Acc. to John Mal. above, they urged that they had deliberately collected the crowd in the Hippodrome, in order to facilitate the suppression of the rising. Justinian ironically observed, ``If you had such influence with the insurgents, why did you not hinder them from burning down the city?''

} they were executed on the following day and their bodies were cast into the sea. The Emperor, suspicious though he was, probably believed that they were not morally guilty, but feared

p48
that they would be used as tools in future conspiracies. They were too dangerous to be allowed to live, but their children were spared.

The throne of Justinian was saved through the moral energy of Theodora and the loyal efforts of Belisarius. It was not only saved, but it rested now on firmer foundations, for it gave the Emperor the opportunity of taking vigorous measures to break down the opposition of the senatorial nobles to his autocracy. There were no more executions, but eighteen senators who had taken a leading part in the conspiracy were punished by the confiscation of their property and banishment.\texthebrew{⁠}\footnote{The number 18 is in John Mal. \emph{fr.} 46 , \emph{De ins.} and \emph{Chron. Pasch.} ; Procopius says generally \textgreek{\texthebrew{ἁ}πάντων}(\emph{ib.} § 57) >. If 18 is correct, Procopius must have exaggerated either here or

\emph{ib.} § 25 ; but the latter passage is borne out by Marcellinus (plerisque nobilium).

} At a later time, when he felt quite secure, Justinian pardoned them and restored to them any of their possessions which he had not already bestowed on others, and a similar restitution was even made to the children of Hypatius and Pompeius.

The news of the Emperor's victory over his enemies and the execution of the usurper was proclaimed in the cities throughout the Empire. For a long time after this event the factions of the Hippodrome seem to have been on their good behaviour, if we may judge by the silence of the chroniclers. During the last twenty years of the reign riots and faction fights occurred from time to time, but the rival parties did not combine again and the disorders were easily put down.\footnote{In 547 the factions quarrelled, and the disturbance was suppressed by the Excubitors with some blood-shed (John Mal. XVIII.483). Riots are recorded in 548 and 551 (\emph{ib.} 484). In May 556 there was a more serious popular demonstration against the Prefect of the City, due to a famine which lasted for three months; the leaders of the Blues were seized and punished (\emph{ib.} 488). In May 559 there were conflicts between the Blues and Greens for two days, and the disorder was accompanied by incendiarism; the house of Peter Barsymes the Praet. Prefect was burnt, and the intervention of the Excubitors was necessary (\emph{ib.} 490). Other disturbances are recorded in 562 (\emph{ib.} 492) and 563 (Theophanes, A.M. 6055), and (without date) in John Mal. \emph{frs.} 50 and 51, \emph{De ins.}

}

After the suppression of this formidable rebellion, one of the first anxieties of the Emperor was to set about rebuilding the edifices which had been destroyed by fire, above all the church of St. Sophia. He was sitting amidst ruin and devastation,

p49
and it would be natural if he had thought of nothing but restoring the wrecked buildings as rapidly as possible; but he saw in the calamity an opportunity for making his capital more magnificent, and constructing a church which would be the wonder of the world. The damage might well have been made good in two years if he had been content to rebuild on the same scale; the work he designed took five years, and considering what was accomplished the time seems incredibly short.

Forty days after the tumult had subsided, the ruins of the church were cleared away, neighbouring houses were bought up, and space was provided for a new temple of the Divine Wisdom. The plans were prepared by Anthemius of Tralles, an architect and engineer who possessed imagination as well as mastery of his craft, and to him was entrusted the direction of the work, with Isidore of Miletus as his assistant. It is to be noted that both these architects were natives of Asia Minor. We cannot doubt that Anthemius had already given proofs of his skill as a builder, and it is not bold to conjecture that he was the architect of the church of SS. Sergius and Bacchus, which Justinian and Theodora had caused to be erected at the beginning of their reign. Justinian had extended the precincts of the Great Palace to take in the house of Hormisdas \texthebrew{—} on the seashore, south of the Hippodrome \texthebrew{—} which had been his residence before he ascended the throne; and close to it he built two churches side by side with a common court, a basilica of SS. Peter and Paul, which has disappeared,\texthebrew{⁠}\footnote{Van Millingen, \emph{Byz. Churches}, 62 \emph{sqq.} The church was founded in 527.

} and the octagonal domed church of SS. Sergius and Bacchus, which has survived, converted by the Turks into a mosque which they call the Little St. Sophia. The names of the Emperor and Empress are associated in the metrical inscription which is still to be seen on the frieze and their monograms can be read on the beautiful ``melon'' capitals. Modern architects have paid tribute to the remarkable skill with which the dome has been buttressed and weighted, and we may divine that it was the skill of Anthemius, of whom a contemporary said that he ``designed wonder ­ful works both in the city and in many other places which would suffice to win him everlasting glory in the memory of men so long as they stand and endure.''\footnote{Agathias, \emph{Hist.} V.8.

}

His plan of St. Sophia was different. It is a Greek cross

p50
\texthebrew{•} (about 250 by 225 feet) with a dome rising above the quadrilateral space between the arms to the height of 180 feet. He undertook to solve the problem of pla­cing a great aerial cupola, 100 feet in diameter, over this space which was 100 feet square. Hitherto cupolas had been set over round spaces. At each angle of the square Anthemius erected a massive pier, in which the settings of the stones were strengthened by special methods. These piers supported the four arches and pendentives on which the ribbed dome rested, and he calculated on securing stability by the semi-domes on the east and west and buttresses on the north and south. To diminish the weight of the dome very light materials were used, tiles of a white spongy earth manufactured at Rhodes.\footnote{It has been a subject of debate whether the architects of the sixth century, in their dome constructions, derived their inspiration from the East, or were simply working along the lines of Roman architectural tradition and adopting suggestions from Roman models. Rivoira has conjectured that Anthemius (whose brother Alexander was a physician at Rome, \emph{loc. cit.}) had visited Rome and picked up ideas there. He calls attention (\emph{Lombardic Arch.} I p79) to the family likeness between the plan of St. Sophia ``and the two halls of the Baths of Agrippa and Nero as well as the Basilica Nova of Maxentius and Constantine.'' He also notes (with Choisy) the resemblance of SS. Sergius and Bacchus with the Licinian Nymphaeum; and he thinks the designer of that church derived from the Serapeum of Hadrian's Villa the idea of a dome of which the surface is a rhythmic sequence of flat and concave sections unsupported by pendentives, simply flush with the course of the drum from which they start (p81). Further, he holds that the visible radiating ribs with which Isidore provided the dome of St. Sophia, when he restored it, were suggested by the Mausoleum Augustorum, a building (of fifth century) consisting of two rotundas, of one of which (it survived till the sixteenth century) a sketch has been preserved (pp82\texthebrew{‑}83). Strzygowski, on the other hand, contends that the origin of domed churches was oriental, and in one of his latest works he particularly connects it with Armenia (\emph{Die Baukunst der Armenier und Europa}, 2 vols., 1918).

}

The material of St. Sophia, as of most Byzantine churches, was brick. Its exterior appearance, seen from below, does not give a true impression of its dimensions. The soaring cupola is lost and buried amid the surrounding buttresses that were added to secure it in later ages. From afar one can realise its proportions, lifted high above all the other buildings and dominating the whole city like a watch-tower , as Procopius described it. But in it, as in other Byzantine churches, the contrast between the plainness of the exterior and the richness of the interior decoration is striking. Although the mosaic pictures, including the great cross on a starry heaven at the summit of the dome,

p51
are now concealed from the eyes of faithful Moslems by whitewash, the marbles of the floor, the walls, and the pillars show us that the rapturous enthusiasm of Justinian's contemporaries as to the total effect was not excessive. The roof was covered with pure gold, but the beauty of the effect lay, it was observed, rather in the answering reflexions from the marbles than from the gold itself. The marbles from which were hewn the pillars and the slabs that covered the walls and floor were brought from all quarters of the world. There was the white stone from the quarries in the Proconnesian islands near at hand, green cipollino from Carystus in Euboea, verde antico from Laconia and Thessaly, Numidian marble glinting with the gold of yellow crocuses, red and white from Caria, white-misted rose from Phrygia, porphyry from Upper Egypt. To Procopius the building gave the impression of a flowering meadow.

While the artists of the time showed skill and study in blending and harmonising colours, the sculptured decoration of the curves of the arches with acanthus and vine tendrils, and the beauty of the capitals of white Proconnesian marble, are not less wonder ­ful. The manufacture of capitals for export had long been an industry at Constantinople, and we can trace the evolution of their forms. The old Corinthian capital, altered by the substitution of the thorny for the soft acanthus, had become what is known as the ``Theodosian'' capital.\texthebrew{⁠}\footnote{An example is preserved in the fifth-century church of St. John of Studion.

} But it was found that this was not suitable for receiving and supporting the arch, and the device was introduced of pla­cing above it an intermediate ``impost,'' in the form of a truncated and reversed pyramid, which was usually ornamented with vine or acanthus, a cross or a monogram. Then, apparently early in the sixth century, the Theodosian capital and the impost were combined into a single block, the ``capital impost,'' which assumed many varieties of form.\footnote{See Lethaby and Swainson, \emph{S. Sophia}, 247 \emph{sqq.}; Diehl, \emph{L'Art byz.} 128; and cp. Rivoira, \emph{op. cit.} I.128.

}

The building was completed in A.D. 537, and on December 26 the Emperor and the Patriarch Menas drove together from St. Anastasia to celebrate the inaugural ceremonies.\texthebrew{⁠}\footnote{Theophanes, A.M. 6030.

} But Anthemius had been overbold in the execution of his architectural design, and had not allowed a sufficient margin of safety for the

p52
support of the dome. Twenty years later the dome came crashing down, destroying in its fall the ambo and the altar (May, A.D. 558).\texthebrew{⁠}\footnote{Agathias, V.9 . John Mal. XVIII p489 , ``in the 6th indiction.'' Theophanes follows Malalas, but puts the notice under a wrong year, A.M. 6051, which would mean A.D. 559. The whole dome does not seem to have collapsed; as it consisted of independent sections this was possible. See Jackson, \emph{Byz. Architecture}, I.87.

} Anthemius was dead, and the restoration was undertaken by Isidore the Younger. He left the semi-domes on the east and west as they were, but widened the arches on the north and south, making ``the equilateral symmetry'' more perfect, and raised the height of the dome by \texthebrew{•} more than twenty feet. The work was finished in A.D. 562, and on Christmas Eve the Emperor solemnly entered it. The poet Paul, the silentiary, was commanded to celebrate the event in verse, and a few days later\texthebrew{⁠}\footnote{Friedländer (p110, Preface to ed. of Paul) conjectures on Epiphany, Jan. 6. Paul, son of Cyrus, was of good family and great wealth.

} he recited in the Palace the proem of his long poem describing the beauties of the church. Justinian then proceeded in solemn procession to St. Sophia, and in the Patriarch's palace, which adjoined the church, he recited the rest. It was a second inauguration, and the effort of Paul was not unworthy of the occasion.

Terrible, thought a writer of the day, as well as marvellous, the dome of St. Sophia ``seems to float in the air.'' It was pierced by forty windows, the half-domes by five, and men were impressed by the light which flooded the church. ``You would say that sunlight grew in it.'' Lavish arrangements were made for artificial illumination for the evening services. A central chandelier was suspended by chains from the cornice round the dome over the ambo; the poet compared it to a circular dance of lights:

And in other parts of the building there were rows of lamps in the form of silver bowls and boats.

Justinian did not regard expense in decorating with gold and precious stones the ambo which stood in the centre under the dome. Similar sumptuousness distinguished the sanctuary of the apse \texthebrew{—} the iconostasis and the altar which was of solid gold. The Patriarch's throne was of gilded silver and weighed

p53
40,000 lbs. A late record states that the total cost of the building and furnishing of St. Sophia amounted to 320,000 lbs. of gold, which sent to our mint to\texthebrew{‑}day would mean nearly fourteen and a half million sterling,\texthebrew{⁠}\footnote{The figure is given in the \emph{\textgreek{Διήγησις περ\texthebrew{ὶ} τ\texthebrew{ῆ}ς \texthebrew{ἁ}γ}}. \textgreek{Σοφίας}, c25, p102; it excludes the cost of the sacred vessels and the numerous private gifts. It is a curious coincidence that the sum is identical with the reserve saved by Anastasius (see above,

§ 4, p36 ). Three hundred and sixty-five lbs. of gold (\emph{c.} £1,642,500) were said to have been spent on the ambo, \emph{ib.} This narrative is marked not only by miracles and obvious fictions but by curious errors, such as dating the beginning of the building of the church to A.D. 538. The fabulous sum which the building is said to have cost might have been reached if 10,000 workmen had been continuously employed and received the wages \emph{de luxe} which Justinian is said to have lavished on them (\emph{\textgreek{Διήγησις}}, § 9 and § 20). It is remarkable, however, that the compiler believed that his figures were derived from an account of the expenses kept by Strategius, who was Count of the Sacred Largesses (as we otherwise know) in 536\texthebrew{‑}537 (on the credibility of the document cp. Preger, \emph{B. Z.} X.455 \emph{sqq.}). Uspenski, \emph{Ist. viz. imp.} I.532, accepts the figures. For a more trustworthy figure in connexion with the building see below,

p55, n3 . I should be surprised if the total expenses amounted to a million sterling. If they were partly paid out of the Private Estate, the confiscated property of the rich senators concerned in the rebellion would have gone a long way.

} a figure which is plainly incredible.

But this, though it was the greatest item in the Emperor's expenditure on restoring and beautifying the city, was only one. The neighbouring church of St. Irene also rose from its ashes, as a great domed basilica, the largest church in Constantinople except St. Sophia itself.\texthebrew{⁠}\footnote{See George, \emph{Church of St. Eirene.}

} The monograms of Justinian and Theodora are still to be read on the capitals of its pillars. More important as a public and Imperial monument was the Church of the Holy Apostles in the centre of the city, which had not been injured by fire, but had suffered from earthquakes and was considered structurally unstable. Justinian pulled it down and rebuilt it larger and more splendid, as a cruciform church with four equal arms and five domes. Though it was destroyed by the Turks to make room for the mosque of Mohammed the Conqueror, descriptions are preserved which enable us to restore its plan.\texthebrew{⁠}\footnote{The history, plan, and mosaics of the church are fully treated in Heisenberg's work, \emph{Die Apostelkirche.}

} San Marco at Venice was built on a very similar design and gives the best idea of what it was like. It may have been begun after the completion of St. Sophia, for it was dedicated in A.D. 546; but the mosaic decoration, of which full accounts have come down to us, was not executed till after Justinian's death, and it has been shown that these pictures, which may

p54
belong to the time of his immediate successors, were designed and selected with a dogmatic motif. ``The two natures of Christ in one person are the theme of the whole cycle.''\texthebrew{⁠}\footnote{Heisenberg, \emph{op. cit.} p168.

} The use of pictures for propagating theological doctrine was understood in the sixth century; we shall see another example at Ravenna.\footnote{Below,

Chap. XVIII § 12 . º

}

The principal secular buildings which had been destroyed by the fires of the Nika riot and were immediately rebuilt were the Senate-house , the baths of Zeuxippus, the porticoes of the Augusteum, and the adjacent parts of the Palace. The Chalke had been burnt down, and the contiguous quarters behind it \texthebrew{—} the portico of the Scholarian guards and the porticoes of the Protectors and Candidates. All these had to be rebuilt.\texthebrew{⁠}\footnote{Procopius, \emph{Aed.} 1.10 . Cp. \emph{Chron. Pasch., sub} 532 .

} But at the same time Justinian seems to have made extensive changes and improvements throughout the Palace; we are told that he renovated it altogether.\texthebrew{⁠}\footnote{Procopius, \emph{ib.} 10 \textgreek{νέα μ\texthebrew{ὲ}ν τ\texthebrew{ὰ} βασίλεια σχεδόν τι πάντα}. I conjecture that Justinian may have designed the Chrysotriclinos buildings, and that Justin II, to whom they are attributed in our sources, may have only completed them.

} Of the details we hear nothing, except as to the Chalke itself. You go through the great gate of the Chalke from the Augusteum, and then through an inner bronze gate into a domed rectangular room, decorated by mosaic pictures showing the Vandal and Italian conquests, with Justinian and Theodora in the centre, triumphing and surrounded by the Senate.\footnote{Procopius, \emph{ib.} The central dome was sustained by four piers, and two arches on the south and two on the north, springing from the walls, formed the vaulted ceilings north and south of the dome.

}

If the Emperor spent much on the restoration and improvement of the great Palace, he appears to have been no less lavish in enlarging and embellishing his palatial villa at Hêrion, on the peninsula which to\texthebrew{‑}day bears the name of Pharanaki, to the south-east of Chalcedon.\texthebrew{⁠}\footnote{The site has been fixed beyond dispute by Pargoire, \emph{Hiéria.} The name of the peninsula took many forms. In the time of Justinian it was generally called \textgreek{\texthebrew{Ἱ}ερόν} or \textgreek{\texthebrew{Ἤ}εριον}, in the time of Heraclius \textgreek{\texthebrew{Ἱ}έρεια.}} It was the favourite resort of Theodora in summer; she used to transport her court there every year.\texthebrew{⁠}\footnote{Her courtiers did not all like the change, if we can draw any inference from the complaint of Procopius that Theodora inhumanly exposed them to the perils of the sea and that they suffered from the lack of the comforts they enjoyed in the city.

\emph{H. A.} 15.36\texthebrew{‑}37 .

} Here Justinian created a small town, with a splendid church

p55
dedicated to the Mother of God, baths, market-places , and porticoes; and constructed a sheltered landing-place by building two large moles into the sea.\footnote{Procopius, \emph{Aed.} I.11 . The historian was perhaps thinking particularly of the construction of this port when he was writing,

\emph{H. A.} 8.7\texthebrew{‑}8 , as Pargoire suggests (\emph{op. cit.} p58).

\texthebrew{◂} previous

next \texthebrew{▸}Images with borders lead to more information.

The thicker the border, the more information.

(Details here.)

UP TO:

J. B. Bury
Homepage

Latin \& Greek

Texts

LacusCurtius

Home}

\xchapter{Chapter XV, Part C}

Short URL for this page:

bit.ly/BURLAT15C

\xsubsection{(Part 3 of 3)}

The nine or ten years following the suppression of the Nika revolt were the most glorious period of Justinian's reign. He was at peace with Persia; Africa and Italy were restored to his dominion. The great legal works which he had undertaken were brought to a success ­ful conclusion; and Constantinople, as we have just seen, arose from its ashes more magnificent than ever.

But the period was hardly as happy for the subjects as it was satisfactory to their ruler. For a short time the fiscal exactions under which they had groaned may have been alleviated under the milder administration of the popular Phocas, who had succeeded the Cappadocian, and who at least had no thought of using his office to enrich himself.\texthebrew{⁠}\footnote{Phocas, son of Craterus, was of good family and well off. He began as a silentiary in the Palace. He had been prosecuted as a pagan in 529 (John Mal. XVIII.449, where it is falsely said that he was put to death). See the panegyric of John Lydus (III.70 \emph{sqq.}) , where his liberality and personal frugality are praised. When he became Prefect he set himself to learn Latin, and John Lydus procured him the services of an instructor. He is highly spoken of by the Emperor in \emph{Nov.} 82, § 1 (539).

} But Phocas was soon removed, probably because his methods failed to meet the financial needs of the Emperor, engaged in preparations for the African expedition and in plans for rebuilding the city on a more splendid scale. In less than a twelvemonth John the Cappadocian was once more installed in the Prefecture, and was permitted for eight or nine years to oppress the provinces of the East.\texthebrew{⁠}\footnote{John Mal., \emph{ib.} 477, dates John's restoration to 532. We know from Procopius, \emph{B. V.} 1.10 and 13, that he was Prefect before June 533. He fell in the tenth year of his Prefecture (\emph{B. P.} I.25.3), and as his fall can be fixed to May 541 (see below), Procopius seems to count as if his two tenures of office had been continuous. Phocas held the Prefecture long enough to furnish 4000 lbs. of gold (£288,000) towards the building of St. Sophia (John Lyd. \emph{ib.} 76) , but for less than a year

(Proc. \emph{H. A.} 21.7) .

} Justinian did not feel himself bound by the promises he had made to the insurgents, seeing that they had been made in vain. He also restored to the post of Quaestor

p56
the great jurist Tribonian, who, otherwise most fitted to adorn the office, seems to have been somewhat unscrupulous in indulging his leading passion, a love of money.

The only person whom John the Cappadocian feared was the Empress. He knew that she was determined to ruin him. He was unable to undermine her influence with Justinian, but that influence did not go far enough to shake Justinian's confidence in him. He dreaded that her emissaries might attempt to assassinate him, and he kept around him a large band of armed retainers, a measure to which no Praetorian Prefect except Rufinus had resorted before. He was exceedingly superstitious, and impostors who professed to foretell the future encouraged him in the hope that he would one day sit on the Imperial throne. In A.D. 538 he enjoyed the expensive honour of the consul ­ship.\footnote{\emph{Fasti cons.} (p56): Flavius Johannes. Cp. CIL VI.32 , 042.

}

If there was one man whom John detested and envied it was the general Belisarius, who in A.D. 540 arrived at Constantinople, bringing the king of the Ostrogoths as a captive in his train. If any man was likely to be a dangerous rival in a contest for the throne, it was the conqueror of Africa and Italy, who was as popular and highly respected as John himself was unpopular and hated. As a matter of fact, thoughts of disloyalty were far from the heart of Belisarius, but he was not always credited with unswerving fidelity to Justinian, even by Justinian himself.

Belisarius, like his master, was born in an Illyrian town,\texthebrew{⁠}\footnote{Germania, on the borders of Illyricum and Thrace, Procop. \emph{B. V.} I.11.21 . Cp. the town of Germae (\textgreek{Γερμαή}) in Dardania

(\emph{Aed.} IV.1.31) ; Hierocles, \emph{Synecd.} 654.5).

} and, like his master, he had married a woman whose parents were associated with the circus and the theatre and who was the mother of children before she married the soldier.\texthebrew{⁠}\footnote{Her father and grandfather were chariot-drivers, her mother a disreputable actress (\textgreek{τ\texthebrew{ῶ}ν τινος \texthebrew{ἐ}ν Θυμέλ\texthebrew{ῃ} πεπορνευμένων}),

Procop. \emph{H. A.} 1.11 . She is said to have been 60 years old in 544 (\emph{ib.} IV.41). She had borne a daughter to Belisarius, Joannina. An illegitimate daughter married Ildiger, an officer who was prominent in the wars of the time; and by another illegitimate child she had a granddaughter who married Sergius, a nephew of Solomon the eunuch (see below,

p145 ). For his son Photius see below,

p60 .

} Unlike Theodora, she did not mend her morals after her marriage, and her amours led to breaches with her husband. But notwithstanding temporary estrangements she preserved the affection

p57
of her husband, who had a weak side to his character, and she faithfully accompanied him on his campaigns and worked energetically in his interests. She often protected him, when he was out of favour with the Emperor, through her influence with Theodora, who found her a useful ally and resource ­ful agent.

The cunning of this unscrupulous woman compassed the fall of John of Cappadocia. She was interested in destroying him as the enemy both of her husband and of her Imperial mistress. The only hope of damaging him irretrievably in the eyes of the Emperor was to produce clear evidence that he entertained treasonable designs, and for this purpose Antonina resorted to the vile arts of an \emph{agent provocateur.} The Prefect had a daughter, his only child, whom he loved passionately; it was the one amiable trait in his repulsive character. His enemies could cast no reproach on the virtue of Euphemia, but she was very young and fell an easy victim to the craft of Antonina. It was in April or May A.D. 541 that the treacherous scheme was executed. Belisarius had set out in the spring to take command in the Persian war, and his wife had remained for a short time at Constantinople before she followed him to the East. She employed herself in cultivating the acquaintance of Euphemia, and having fully won her friendship she persuaded the inexperienced girl that Belisarius was secretly disaffected towards Justinian. It is Belisarius, she said, who has extended the borders of the Empire, and taken captive two kings, and the Emperor has shown little gratitude for his services. Euphemia, who, taught to see things through her father's eyes, feared Theodora and distrusted the government, listened sympathetically to the confidences of her friend. ``Why,'' she asked, ``does Belisarius not use his power with the army to set things right?'' ``It would be useless,'' said Antonina, ``to attempt a revolution in the camp without the support of civilian ministers in the capital. If your father were willing to help, it would be different.'' Euphemia eagerly undertook to broach the matter to her father. John, when he heard his daughter's communication, thought that a way was opened for realising the vague dreams of power which he had been cherishing. It was arranged that he should meet Antonina secretly. She was about to start for the East, and she would halt for a night near Chalcedon at the palace

p58
of Rufinianae, which belonged to her husband.\texthebrew{⁠}\footnote{At Jadi Bostan (see above,

Vol. I p87 ) We do not know how this Imperial residence passed into the possession of Belisarius. At \textgreek{Παντείχιων}, near Chalcedon, evidently to be identified with Pendik, Belisarius had also a villa. See Pargoire, \emph{Rufinianes}, pp459, 477.

} Hither John agreed to come secretly, and the day and hour were arranged. Antonina then informed the Empress of all she had done and the details of the scheme. It was essential that the treasonable conversation should be overheard by witnesses, whose testimony would convince Justinian. Theodora, who entered eagerly into the plot, chose for this part the eunuch Narses, and Marcellus, commander of the Palace guards, a man of the highest integrity, who stood aloof from all political parties, and never, throughout a long tenure of his command, forfeited the Emperor's respect.\texthebrew{⁠}\footnote{It is uncertain whether he was Count of the Excubitors or Commander of the Scholarians. For his character see Procopius, \emph{B. G.} III.32.22.

} Theodora did not wait for the execution of the scheme to tell Justinian of what was on feet, and it was said that he warned John secretly not to keep the appointment. This may not be true. In any case, John arrived at Rufinianae at midnight, only taking the precaution of bringing some of his armed retainers. Antonina met him outside the house near a wall behind which she had posted Marcellus and Narses. He spoke, without any reserve, of plans to attempt the Emperor's life. When he had fully committed himself, Narses and Marcellus emerged from their hiding-place to seize him. His men, who were not far off, rushed up and one of them wounded Marcellus. In the fray John succeeded in escaping, and reaching the city he sought refuge in a sanctuary.\texthebrew{⁠}\footnote{\textgreek{\texthebrew{Ἐ}ς τ\texthebrew{ὸ ἱ}ερόν} (\emph{B. P.} I.25.30), where Haury conjectures that the words \textgreek{το\texthebrew{ῦ} Λαυρεντίου} have fallen out.

} The historian who tells the tale thought that if he had gone boldly to the Palace he would have been pardoned by Justinian. But the Empress now had the Emperor's ear. John was deprived of the office which he had so terribly abused and banished to Cyzicus, where he was ordained a deacon against his will. His large ill-gained possessions were forfeited as a matter of course, but the Emperor showed his weakness for the man by letting him retain a considerable portion, which enabled him to live in great luxury in his retirement.\footnote{All the details are derived from Procopius, \emph{B. P.} I.25. In the summary notice of John Mal. XVIII.480 (cp. \emph{De ins.} p172), the date is given as August. But Theodotus had succeeded John as Prefect between May 7 and June 1, 541 (\emph{Nov.} 109 and 111).

}

But he was not long suffered to enjoy his exile in peace. The

p59
bishop of Cyzicus, Eusebius, was hated by the inhabitants. They had preferred charges against him at Constantinople, but his influence there was so great that he was able to defy Cyzicus. At last some young men, who belonged to the local circus factions, murdered him in the market-place . As it happened that John and Eusebius were enemies, it was suspected that John was accessory to the crime, and, considering his reputation, the suspicion was not unnatural. Senators were sent to Cyzicus to investigate the murder. John's guilt was not proved;\texthebrew{⁠}\footnote{\textgreek{Ο\texthebrew{ὐ} λίαν \texthebrew{ἐ}ξελήλεγκτο}, Procop. \emph{ib.} Theodora's name is not mentioned in the episode as narrated here, but the part which she played comes out in the supplementary story in

\emph{H. A.} 17.40 \emph{sqq.}

Procopius wrote \emph{B. P.} I.25 in the third year of John's imprisonment in Egypt (\textgreek{τρίτον το\texthebrew{ῦ}το \texthebrew{ἔ}τος α\texthebrew{ὐ}τ\texthebrew{ὸ}ν \texthebrew{ἐ}ντα\texthebrew{ῦ}θα καθείρξαντες τηρο\texthebrew{ῦ}σιν}, § 43), probably in 544\texthebrew{‑}545. For when he says (§ 44) that the retribution for his acts overtook John ten years later, he seems to mean that ten years elapsed between John's reappointment after the Nika, in 532\texthebrew{‑}533, and the affair at Cyzicus, which would thus have occurred in 542\texthebrew{‑}543.

} but the commission of inquiry must have received secret orders to punish him rightly or wrongly, for he was stripped and scourged like a common highwayman, and then put on board ship, clad in a rough cloak. The ship bore him to Egypt, and on the voyage he was obliged to support life by begging in the seaports at which it called. When he reached Egypt he was imprisoned at Antinoopolis. For these illegal proceedings the Emperor, we may be sure, was not responsible, and no private enemy could have ventured to resort to them. The hand of Theodora could plainly be discerned. But she was not yet satisfied with his punishment; she desired to have him legally done to death.\texthebrew{⁠}\footnote{What follows comes from \emph{H. A.} 17 , where the date given is \textgreek{τέτρασι \texthebrew{ἐ}νιαυτο\texthebrew{ῖ}ς \texthebrew{ὑ}στερον} (\emph{i.e.} 546\texthebrew{‑}547, see last note).

} Some years later she got into her power two young men of the Green faction who were said to have been concerned in the murder of the bishop. By promises and threats she sought to extract a confession implicating John the Cappadocian. One of them yielded, but the other, even under torture, refused. Baffled in her design she is said to have cut off the hands of both the youths.\texthebrew{⁠}\footnote{This was done publicly in the Forum of Constantine. Justinian, according to Haury's very probable emendation of the text (\emph{ib.}, \emph{ad fin.}), pretended to know nothing about it.

} John remained in prison till her death, after which he was allowed by the Emperor to return to Constantinople, a free man, but a priest. Yet it was said that he still dreamed of ascending the throne.\footnote{Procop. \emph{B. P.} II.30.

}

p60
It is incontestable that Theodora performed a public service by delivering the eastern provinces from the government of an exceptionally unscrupulous oppressor, and that his sufferings, although they were illegally inflicted, were richly deserved. But the revolting means imagined by her unprincipled satellite Antonina and approved by herself, the employment of the innocent girl to entrap her father,\texthebrew{⁠}\footnote{We are not told what became of Euphemia.

} do not raise her high in our estimation. It must be observed, however, that the public opinion of that time found nothing repulsive in a stratagem which to the more delicate feelings of the present age seems unspeakably base and cruel. For the story is told openly in a work which the author could not have ventured to publish if it had contained anything reflecting injuriously on the character of the Empress.\footnote{One point was indeed omitted, the shameless perjury of Antonina, and in the \emph{Secret History} Procopius holds it up to censure. She had sworn both to John and to his daughter by the most solemn oaths that a Christian can take that no treachery was intended.

\emph{H. A.} 2.61 .

}

It was not long before the Empress had an opportunity of repaying her friend for her dexterous service. Belisarius and Antonina had adopted a youth named Theodosius, for whom Antonina conceived an ungovernable passion. Their guilty intrigue was discovered by Belisarius in Sicily ( A.D. 535\texthebrew{‑}536), and he sent some of his retainers to slay the paramour, who, however, escaped to Ephesus.\texthebrew{⁠}\footnote{The story is told in

Procopius, \emph{H. A.} 1.15 \emph{sqq.}

Theodosius was a Thracian; his parents belonged to the Eunomian sect, and he was baptized into the true faith at his adoption just before the Vandal expedition started. Perhaps he is referred to in \emph{B. V.} I.12.2 as ``one of the soldiers recently baptized.''

} But Antonina persuaded her uxorious husband that she was not guilty, regained his affection, and induced him to hand over to her the servants who had betrayed her amour. It was reported that having cut out their tongues she chopped their bodies in small pieces and threw them into the sea. Her desire for Theodosius was not cooled by an absence of five years, and while she was preparing her intrigue against John the Cappadocian she was planning to recall her lover to her side when Belisarius departed for the East. But her son Photius, who had always been jealous of the favourite preferred to himself, penetrated her design and revealed the matter to his stepfather, and they bound themselves by solemn oaths to punish Theodosius. They decided, however, that

p61
nothing could be done immediately; they must wait till Antonia followed her husband to the East. Photius accompanied Belisarius in the campaign, and for some months Antonina enjoyed the society of her paramour at Constantinople. When, in the summer of the year, she set out for Persia, Theodosius returned to Ephesus. The general met his wife, showed his anger, but had not the heart to slay her. Photius hastened to Ephesus, seized Theodosius, and sent him under a guard of retainers to be imprisoned in a secret place in Cilicia. He proceeded himself to Constantinople, in possession of the wealth which Theodosius had been allowed to appropriate from the spoils of Carthage.\texthebrew{⁠}\footnote{\emph{H. A.} 1.17 , ``it is said that he had got by plunder 10,000 nomismata from the palaces of Carthage and Ravenna.'' He was with Belisarius at Carthage, but was at Ephesus when Ravenna was taken; so that, if the report is true, Belisarius or Antoninus must have sent him a share of the Italian plunder.

} But the danger of her favourite had come to the ears of Theodora. She caused Belisarius and his wife to be summoned to the capital and she forced a reconciliation upon the reluctant husband. Then she seized Photius and sought by torture to make him reveal the place where he had concealed Theodosius. The secret was disclosed, however, through another channel and Theodosius was rescued; the Empress concealed him in the Palace, and presented him to Antonina as a delight ­ful surprise.

The unhappy Photius, who showed greater force of character than Belisarius, was kept a captive in the dungeons of Theodora. He escaped twice, but was dragged back from the sanctuaries in which he had sought refuge. His third attempt at the end of three years was success ­ful; he reached Jerusalem, became a monk, and escaped the vengeance of the Empress. He survived Justinian, and in the following reign was appointed, notwithstanding his religious quality, to suppress a revolt of the Samaritans, a task which he carried out, we are told, with the utmost cruelty, taking advantage of his powers to extort money from all the Syrian provinces.\footnote{John Eph. \emph{Hist. ecc.} part 3, book I cc31, 32. This notice gives a very unfavourable impression of Photius, with whom Procopius is sympathetic. They agree in the statement that he became a monk.

}

Justinian had been fourteen years on the throne when the Empire was visited by one of those immense but rare calamities in the presence of which human beings could only succumb helpless and resourceless until the science of the nineteenth century began to probe the causes and supply the means of preventing and checking them. The devastating plague, which began its course in the summer of A.D. 542 and seems to have invaded and ransacked nearly every corner of the Empire, was, if not more malignant, far more destructive, through the vast range of its ravages, than the pestilences which visited ancient Athens in the days of Pericles and London in the reign of Charles II; and perhaps even than the plague which travelled from the East to Rome in the reign of Marcus Aurelius. It probably caused as large a mortality in the Empire as the Black Death of the fourteenth century in the same countries.\footnote{The best and principal authority is Procopius (\emph{B. P.} II.22\texthebrew{‑}23), who was living in Constantinople during the visitation. But we have a second first-hand source in John of Ephesus, who was in Palestine when the plague broke out there, and then, travelling to Mesopotamia and returning to Constantinople through Asia Minor, observed its ravages in Cilicia, Cappadocia, and Bithynia (extracts from his History in Land's ed. of the \emph{Commentarii}, p227 \emph{sqq.}). The date of the outbreak is fixed to 542 by the order of events in Procopius, and this agrees with the specific date of another contemporary, John Malalas (XVIII p48), who places it in the 5th indiction (541\texthebrew{‑}542), and with that of Evagrius

(IV.29) , who places it two years after the Persian capture of Antioch (June 540). The date of John of Ephesus (year of Alexandria 855 = A.D. 543\texthebrew{‑}544) may be explained by supposing that the year of the journey was 543. The plague was also noticed by Zachariah of Mitylene in the lost portion of book X, but a short extract is preserved in the chronicle of Michael Syrus (IX.28; see Zachariah, p313).

}

The infection first attacked Pelusium, on the borders of Egypt, with deadly effect, and spread thence to Alexandria and throughout Egypt, and northward to Palestine and Syria. In the following year it reached Constantinople, in the middle of spring,\texthebrew{⁠}\footnote{Procopius, \emph{op. cit.} 22, 9.

} and spread over Asia Minor and through Mesopotamia into the kingdom of Persia.\texthebrew{⁠}\footnote{Zachariah mentions Nubia, Ethiopia, Armenia, Arzanene, and Mesopotamia. For Persia cp. also Procopius, \emph{op. cit.} 25, 12.

} Travelling by sea, whether from Africa or across the Adriatic, it invaded Italy and Sicily.\footnote{\emph{Cont. Marcellini, sub} 453, mortalitas magna Italiae solum devastat, Orientem iam et Illyricum peraeque attritos. Zachariah's list of the visited countries includes Italy, Sicily, Africa, and Gaul.

}

It was observed that the infection always started from the coast and went up to the interior, and that those who survived

p63
it had become immune. The historian Procopius, who witnessed its course at Constantinople, as Thucydides had studied the plague at Athens, has detailed the nature and effects of the bubonic disease, as it might be called, for the most striking feature was a swelling in the groin or in the armpit, sometimes behind the ear or on the thighs. Hallucinations occasionally preceded the attack. The victims were seized by a sudden fever, which did not affect the colour of the skin nor make it as hot as might be expected.

The fever was of such a languid sort from its commencement and up till evening that neither to the sick themselves nor a physician who touched them would it afford any suspicion of danger. . . . But on the same day in some cases, in others on the following day, and in the rest not many days later a bubonic swelling developed. . . . Up to this point everything went in about the same way with all who had taken the disease. But from then on very marked differences developed. . . . There ensued with some a deep coma, with others a violent delirium, and in either case they suffered the characteristic symptoms of the disease. For those who were under the spell of the coma forgot all those who were familiar to them and seemed to be sleeping constantly. And if any one cared for them, they would eat without waking, but some also were neglected and these would die directly through lack of sustenance. But those who were seized with delirium suffered from insomnia and were victims of a distorted imagination; for they suspected that men were coming upon them to destroy them, and they would become excited and rush off in flight, crying out at the top of their voices. And those who were attending them were in a state of constant exhaustion and had a most difficult time. . . . Neither the physicians nor other persons were found to contract this malady through contact with the sick or with the dead, for many who were constantly engaged either in burying or in attending those in no way connected with them held out in the performance of this service beyond all expectation. . . . [The patient] had great difficulty in the matter of eating, for they could not easily take food. And many perished through lack of any man to care for them, for they were either overcome with hunger, or threw themselves from a height.

And in those cases where neither coma nor delirium came on, the bubonic swelling became mortified and the sufferer, no longer able to endure the pain, died. And we would suppose that in all cases the same thing would have been true, but since they were not at all in their senses, some were quite unable to feel the pain; for owing to the troubled condition of their minds they lost all sense of feeling.

Now some of the physicians who were at a loss because the symptoms were not understood, supposing that the disease centred in the bubonic swellings, decided to investigate the bodies of the dead. And upon opening some of the swellings they found a strange sort of carbuncle [ \textgreek{\texthebrew{ἄ}νθραξ} ] that had grown inside them.

p64
Death came in some cases immediately, in others after many days; and with some the body broke out with black pustules about as large as a lentil, and these did not survive even one day, but all succumbed immediately. With many also a vomiting of blood ensued without visible cause and straightway brought death. Moreover I am able to declare this, that the most illustrious physicians predicted that many would die, who unexpectedly escaped entirely from suffering shortly afterwards, and that they declared that many would be saved who were destined to be carried off almost immediately. . . . While some were helped by bathing others were harmed in no less degree. And of those who received no care many died, but others, contrary to reason, were saved. And again, methods of treatment showed different results with different patients. . . . And in the case of women who were pregnant death could be certainly foreseen if they were taken with the disease. For some died through miscarriage, but others perished immediately at the time of birth with the infants they bore. However they say that three women survived though their children perished, and that one woman died at the very time of child-birth but that the child was born and survived.

Now in those cases where the swelling rose to an unusual size and a discharge of pus had set in, it came about that they escaped from the disease and survived, for clearly the acute condition of the carbuncle had found relief in this direction, and this proved to be in general an indication of returning health. . . . And with some of them it came about that the thigh was withered, in which case, though the swelling was there, it did not develop the least suppuration. With others who survived the tongue did not remain unaffected, and they lived on either lisping or speaking incoherently and with difficulty.\texthebrew{⁠}\footnote{I have borrowed Dewey's translation.

}

This description\texthebrew{⁠}\footnote{It agrees generally with the less accurate description of John of Ephesus.

} shows that the disease closely resembled in character the terrible oriental plague which devastated Europe and parts of Asia in the fourteenth century. In the case of the Black Death too the chief symptom was the pestboils, but the malady was generally accompanied by inflammation of the lungs and the spitting of blood, which Procopius does not mention.\footnote{See Hecker, \emph{The Epidemics of the Middle Ages} (transl. Babington, ed. 3), p3 \emph{sqq.} The coma in some cases, the sleeplessness in others, are mentioned.

}

In Constantinople the visitation lasted for four months altogether, and during three of these the mortality was enormous. At first the deaths were only a little above the usual number, but as the infection spread 5000 died daily, and when it was at its worst 10,000 or upward.\texthebrew{⁠}\footnote{So Procopius. John of Ephesus (\emph{op. cit.} p234) says that 5000, 7000, 12,000, even 16,000 corpses of the poor were removed daily from the streets in the early stage of the plague; and they were counted by men stationed at the harbours, the ferries, and the gates.

} This figures are too

p65
vague to enable us to conjecture how many of the population were swept away; but we may feel sceptical when another writer who witnessed the plague assures us that the number of those who died in the streets and public places exceeded 300,000.\texthebrew{⁠}\footnote{John Eph. \emph{ib.}

} If we could trust the recorded statistics of the mortality in some of the large cities which were stricken by the Black Death \texthebrew{—} in London, for instance, 100,000, in Venice 100,000, in Avignon 60,000 \texthebrew{—} then, considering the much larger population of Constantinople, we might regard 300,000 as not an excessive figure for the total destruction. For the general mortality throughout the Empire we have no data for conjecture; but it is interesting to note that a physician who made a careful study of all the accounts of the Black Death came to the conclusion that, without exaggeration, Europe (including Russia) lost twenty-five millions of her inhabitants through that calamity.\footnote{Hecker, \emph{op. cit.} p29.

}

At first, relatives and domestics attended to the burial of the dead, but as the violence of the plague increased this duty was neglected, and corpses lay forlorn not only in the streets, but even in the houses of notable men whose servants were sick or dead. Aware of this, Justinian placed considerable sums at the disposal of Theodore, one of his private secretaries,\texthebrew{⁠}\footnote{Referendarii; there were fourteen (Justinian,

\emph{Nov.} 10 ). For these officials, not to be confounded with the magistri, see Bury, \emph{Magistri scriniorum}, \textgreek{\texthebrew{ἀ}ντιγραφ\texthebrew{ῆ}ς}\emph{, and} \textgreek{\texthebrew{ῥ}εφερενδάριοι.}} to take measures for the disposal of the dead. Huge pits were dug at Sycae, on the other side of the Golden Horn, in which the bodies were laid in rows and tramped down tightly; but the men who were engaged on this work, unable to keep up with the number of the dying, mounted the towers of the wall of the suburb, tore off their roofs, and threw the bodies in. Virtually all the towers were filled with corpses, and as a result ``an evil stench pervaded the city and distressed the inhabitants still more, and especially whenever the wind blew fresh from that quarter.''\texthebrew{⁠}\footnote{Procopius \emph{B. P.} II.23. His account of the measures for the disposal of the corpses agrees with that of John Eph., who, however, does not mention the towers, and says that each of the pits could contain 70,000 bodies. One would have thought that all the arrangements would have been made by the Prefect of the City, but that functionary is not mentioned.

} It is particularly noted that members of the Blue and Green parties laid aside their mutual enmity and co-operated in the labour of burying the dead.

p66
During these months all work ceased; the artisans abandoned their trades. ``Indeed in a city which was simply abounding in all good things starvation almost absolute was running riot. Certainly it seemed a difficult and very notable thing to have a sufficiency of bread or of anything else.''\texthebrew{⁠}\footnote{\emph{Ib.} (Dewey's rendering).

} All court functions were discontinued, and no one was to be seen in official dress, especially when the Emperor fell ill. For he, too, was stricken by the plague, though the attack did not prove fatal.\footnote{\textgreek{Κα\texthebrew{ὶ} α\texthebrew{ὐ}τ\texthebrew{ῷ} γ\texthebrew{ὰ}ρ ξυνέπεσε βουβ\texthebrew{ῶ}να \texthebrew{ἐ}π\texthebrew{ῆ}ρθαι}. Gibbon suggests that Justinian may have owed his safety to his abstemious habits. Another serious illness of Justinian is recorded by

Procopius, \emph{Aed.} I.7.6 \emph{sqq.}

}

Our historian observed the moral effects of the visitation. Men whose lives had been base and dissolute changed their habits and punctiliously practised the duties of religion,\texthebrew{⁠}\footnote{Cp. Hecker, \emph{op. cit.} p31.

} not from any real change of heart, but from terror and because they supposed they were to die immediately. But their conversion to respectability was only transient. When the pestilence abated and they thought themselves safe they recurred to their old evil ways of life. It may be confidently asserted, adds the cynical writer, that the disease selected precisely the worst men and let them go free.

Fifteen years later there was a second outbreak of the plague in Constantinople (spring A.D. 558), but evidently much less virulent and destructive. It was noticed in the case of this visitation that females suffered less than males.\footnote{Agathias, V.10 . He identifies this disease with that of 543. Of the plagues of the fourteenth century, which were frequent after 1350 until 1383, Hecker says that he does not consider them as the same as the Black Death. ``They were rather common pestilences without inflammation of the lungs, such as in former times, and in the following centuries, were excited by the matter of contagion everywhere existing'' (p27). The plague in Italy and Gaul which Marius Avent. notices, \emph{sub a.} 570, as morbus validus cum profluvio ventris et variola seems to be the same as that whose ravages in Liguria Paulus Diac. (\emph{Hist. Lang.} II.4) describes. Paul mentions the pestboils.

}

The Empress Theodora died of cancer of June 28, A.D. 548.\texthebrew{⁠}\footnote{Date: John Mal. XVIII p484, cp. Procopius, \emph{B. G.} III.30.4 (in \emph{Consularia Ital., sub a.}, the day is given as June 27). The disease is named by Victor Tonn. \emph{sub} 549.

} Her death was a relief to her numerous enemies, but to Justinian it must have been a severe blow. We would give much to have a glimpse into their private life or a record of one of their intimate

p67
conversations. We have no means of lifting even a corner of the veil. But it is a significant fact that, though they disagreed on various questions of policy, scandal, which had many evil things to tell of them both, never found any pretext to suggest that they quarrelled or were living on bad terms.

Soon after this event a conspiracy was formed against the Emperor's life, which had little political significance but created a great sensation because men of his own family were indirectly involved.\texthebrew{⁠}\footnote{The source is Procopius, \emph{B. G.} III.32. The date seems to be the latter part of 548.

} A general named Artabanes, of Armenian race, whom we shall meet as a commander in Africa, had conceived the ambition of marrying the Emperor's niece Praejecta, but the plan had been thwarted by Theodora, who compelled him to live again with the wife whom he had put away.\texthebrew{⁠}\footnote{See below,

p146 .

} After her death he repudiated his wife for the second time, but Praejecta, who had been given to another, was lost to him, and he bore no goodwill towards the Emperor. His disaffected feelings would not have prompted him to initiate any sinister design, but a kinsman of his, one Arsaces, was animated by a bitter desire of revenge upon Justinian, who, when he was found guilty of a treacherous correspondence with the king of Persia, had ordered him to be scourged lightly and paraded through the streets on the back of a camel. Arsaces fanned into flame the smouldering resentment of Artabanes, and showed him how easy it would be to kill the Emperor, ``who is accustomed to sit without guards till late hours in the night, in the company of old priests, deep in the study of the holy books of the Christians.'' But perhaps what did most to secure the adhesion of Artabanes was the prospect that Germanus, Justinian's cousin, and his two sons\texthebrew{⁠}\footnote{See above,

p20 .

} would sanction, if they did not take an active part in, the design.

For Germanus, at this time, had a personal grievance against the Emperor. His brother Boraides had died, leaving almost all his property to Germanus, allowing his daughter to receive only so much as was required by the law. But Justinian, deeming the arrangement unfair, overrode the will in the daughter's favour.\texthebrew{⁠}\footnote{\emph{B. G.} III.31.17\texthebrew{‑}18.

} relying on the indignation which this arbitrary act had aroused in the family, Arsaces opened communications with Justin, the elder son of Germanus. Having bound him by oath not to reveal

p68
the conversation to any person except his father, he enlarged on the manner in which the Emperor ill-treated and passed over his relatives, and expressed his conviction that it would go still harder with them when Belisarius returned from Italy. He then revealed the plan of assassination which he had formed in conjunction with Artabanes and Chanaranges, a young and frivolous Armenian who had been admitted to their counsels.

Justin, terrified at this revelation, laid it before his father, who immediately consulted with Marcellus, the Count of the Excubitors, whether it would be wise to inform the Emperor immediately. Marcellus, an honourable, austere, and wary man, dissuaded Germanus from taking that course, on the ground that such a communication, necessitating a private interview with the Emperor, would inevitably become known to the conspirators and lead to the escape of Arsaces. He proposed to investigate the matter himself, and it was arranged that one of the conspirators should be lured to speak in the presence of a concealed witness. Justin appointed a day and hour for an interview between Germanus and Chanaranges, and the compromising revelations were overheard by Leontius, a friend of Marcellus, who was hidden behind a curtain. The programme of the matured plot was to wait for the arrival of Belisarius and slay the Emperor and his general at the same time; for if Justinian were slain beforehand, the conspirators might not be able to contend against the soldiers of Belisarius. When the deed was done, Germanus was to be proclaimed Emperor.

Marcellus still hesitated to reveal the plot to the Emperor, through friendship or pity for Artabanes. But when Belisarius was drawing nigh to the capital he could hesitate no longer, and Justinian ordered the conspirators to be arrested. Germanus and Justin were at first not exempted from suspicion, but when the Senate inquired into the case, the testimony of Marcellus and Leontius, and two other officers to whom Germanus had prudently disclosed the affair, completely cleared them. Even then Justinian was still indignant that they had concealed the treason so long, and was not mollified until the candid Marcellus took all the blame of the delay upon himself. The conspirators were treated with clemency, being confined in the Palace and not in the public prison. Artabanes was not only soon pardoned

p69
but was created Master of Soldiers in Thrace and sent to take part in the Ostrogothic war.\footnote{\emph{B. G.} III.39.8.

}

Another plot to assassinate Justinian was organised by a number of obscure persons in November A.D. 562,\texthebrew{⁠}\footnote{John Malalas, XVIII.493, and \emph{De ins.} \emph{fr.} 49; Theophanes, A.M. 6055. The conspirators were: Ablabius, son of Miltiades; Marcellus, Vitus, and Eusebius, bankers; Sergius, now of Aetherius, the curator domus divinae; Isaac, \textgreek{\texthebrew{ὁ ἀ}ργυροπράτης \texthebrew{ὁ} κατ\texthebrew{ὰ} Βελισάριον}; and Paul, a retainer of Belisarius. The workshop of Marcellus was near St. Irene, and he is described as \textgreek{\texthebrew{ὁ} κατ\texthebrew{ὰ} Α\texthebrew{ἰ}θέριον τ\texthebrew{ὸ}ν ξουράτωρα}. Marcellus, who was arrested as he entered the Palace with a dagger (\textgreek{βο\texthebrew{ῦ}γλιν} = pugio), killed himself on the spot. Paul the Silentiary refers to this conspiracy in \emph{S. Sophia}, 22 \emph{sqq.}

} and would hardly merit to be recorded if it had not injured Belisarius. One of the conspirators talked indiscreetly to Eusebius, Count of the Federates, and they were all arrested. Their confessions involved two followers of Belisarius, who, seized and examined by the Prefect of the City and the quaestor, asserted that Belisarius was privy to the plot. The Emperor convoked a meeting of the Senate and Imperial Council; the depositions of the prisoner were read; and suspicion weighed heavily on the veteran general. He made no resistance when he was ordered to dismiss all his armed retainers, and he remained in disgrace till July A.D. 563, when he was restored to favour.\texthebrew{⁠}\footnote{Theophanes, \emph{ib.} The technical term for political disgrace is \textgreek{\texthebrew{ἀ}γανάκτησις.}} His character and the whole record of his life make it highly improbable that he was guilty of disloyalty in his old age. He died in March A.D. 565.\texthebrew{⁠}\footnote{Theophanes, A.M. 6057. His property went to the Imperial house of Marina. Antonina survived him according to \emph{\textgreek{Πάτρια}}, p254.

} His disgrace, though it was brief, made such an impression on popular imagination in later times that a Belisarius legend was formed, which represented the conqueror of Africa and Italy as ending his days as a blind beggar in the streets of Constantinople.\footnote{The earliest mention of this legend is in \textgreek{Πάτρια}, p160 (tenth century), cp. Tzetzes, \emph{Chil.} III.339 \emph{sqq.} The extant medieval romance on the subject took shape in the age of the Palaeologi. Belisarius, slandered by jealous nobles, is shut up in a tower. The Emperor releases him on the insistent demand of the people that he should be the leader of a military expedition against England (\textgreek{νησ\texthebrew{ὶ}ν τ\texthebrew{ῆ}ς \texthebrew{Ἐ}γγλητέρας}). Belisarius sails thither, takes the English \textgreek{κάστρον}, and returns to Constantinople with the king of England as prisoner. He is again accused of treason, and is blinded. The three known versions of the story will be found in Wagner, \emph{Carmina Graeca medii aevi}, p304 \emph{sqq.} There is a historical case of disgraced generals (Peganes and Symbatios) being blinded and set to beg in the streets in the ninth century (see Bury, \emph{Eastern Roman Empire}, p176), and perhaps this suggested the mythical fate of Belisarius.

}

As Justinian had no children of his own, it was incumbent on him to avert the possibility of a struggle for the throne after his death by designating a successor. So long as Theodora was alive the importance of providing for the future was not so serious, as it might be reasonably supposed that she would be able to control the situation as successfully as Pulcheria and Ariadne. But after her death it was a dereliction of duty on the part of Justinian, as it had been on the part of Anastasius, not to arrange definitely the question of the succession. His failure to do so was probably due partly to his suspicious and jealous temper, and partly to an inability to decide between the two obvious choices.

Of his three cousins, Germanus, Boraides, and Justus,\texthebrew{⁠}\footnote{See above,

p20 . Germanus died in 550

(Procopius, \emph{B. G.} III.40.9) ; Boraides \emph{c.} 547\texthebrew{‑}548 (see above,

p67 ); Justus in 545 (\emph{B. P.} II.28.1. For the marriage of Germanus see above,

p20 , and for his death, below,

p254 .

} only Germanus survived Theodora, but he, who was an able man and whom the popular wish would have called to the throne, died two years later. His two sons, Justin and Justinian, were competent officers. We have seen them occupying important military posts, and if they were not trusted with the highest commands, it is probable that they did not display ability of the first rank. They were both unreservedly loyal to the sovran, and Justin seems, like his father, to have enjoyed general respect and popularity.\texthebrew{⁠}\footnote{Cp. Evagrius, V.1 .

} If Justinian had decided to create him Caesar or Augustus, the act would have been universally applauded.

The influence of Theodora had rendered it impossible for the Emperor, in her lifetime, to show any special preference for this branch of his kin. Germanus, whose amiable qualities and sense of justice endeared him to others, was hated and suspected by her. She resolved that his family should not multiply. He had children, but he should have no grandchildren. In this design she so far succeeded that neither of his sons married till after her death. All her efforts, however, did not prevent his daughter Justina from espousing the general John, nephew of Vitalian, but she threatened that she would destroy John and he went in fear of his life.\footnote{Procopius, \emph{H. A.} 5.8\texthebrew{‑}15 . John, who was in Italy, avoided associating with Belisarius, through fear of Antonina. If we are to connect this with the statement in

\emph{B. G.} III.18.25

(p71)
\textgreek{παρ\texthebrew{ὰ} Βελισάριον ο\texthebrew{ὐ}κέτι \texthebrew{ἤ}ει} (though a different motive is assigned), the date of the marriage of John and Justina would be A.D. 546.

}

p71
Justinian had nephews, sons of his sister Vigilantia,\texthebrew{⁠}\footnote{See above,

p20 .

} and on the eldest of these, Justin, the Empress bestowed her favour. Her desire was that her own blood should be perpetuated in the dynasty, and she married her niece Sophia, a woman who possessed qualities resembling her own, to Justin. After her death, Justinian seems to have been convinced that the conspicuous merits of Germanus entitled him to the succession, but he was unable to bring himself to take a definite decision. When Germanus died the choice lay between the two Justins, the nephew and the cousin, and we may divine that there was a constant conflict between their interests at court. The Emperor's preference inclined, on the whole, to Justin, the husband of Sophia. He created him Curopalates, a new title of rank which raised him above the other Patricians, yet did not give him the status of an heir apparent which would have been conferred by the title of Caesar or even Nobilissimus.\texthebrew{⁠}\footnote{He was Curopalates in 559, when he repressed a Hippodrome riot, John Mal. XVIII.491. See

Evagrius, V.1 . Sophia was perhaps a daughter of Comito and Sittas.

} But Justin enjoyed the great advantage of living in the Palace and having every opportunity to prepare his way to the throne; while the services of his rival and namesake were employed in distant Colchis.

Not the least of Theodora's triumphs was the posthumous realisation of her plan for the succession. Justinian died on Nov. 14, A.D. 565,\texthebrew{⁠}\footnote{\emph{Chron. Pasch., sub a.} Theophanes gives Nov. 11 as the day. His funeral is described by Corippus in \emph{Laud. Iust.} III.4 \emph{sqq.} The Empress Sophia laid over his bier a purple cloth on which were embroidered in gold pictures of his achievements, Iustinianorum series tota laborum, \emph{ib.} II.276 \emph{sqq.}

} and Justin, the son of Vigilantia, supported by the Senate and the Excubitors, secured the throne without a struggle.

The chronicle of Theophanes contains a remarkable record of a conversation between Justinian and the Green party in the Hippodrome. It is apparently an official record (preserved in the archives of the Greens?), under the title \textgreek{\texthebrew{Ἄ}κτα δι\texthebrew{ὰ} Καλοπόδιον τ\texthebrew{ὸ}ν κουβικουλάριον}p72
\textgreek{κα\texthebrew{ὶ} σπαθάριον} , and is inserted after the short summary of the Nika riot which the chronicler has prefixed to his detailed narrative. But it exhibits no connexion whatever with the causes of that event, and may record an incident which occurred at some other period of the reign.\texthebrew{⁠}\footnote{See P. Maas, \emph{Metrische Akklamationen der Byzantiner}, \emph{B. Z.} XXI.49\texthebrew{‑}50. He reprints the Greek text so as to bring out the rhythmical character of the conversation. The \emph{Acta} were known to the compiler of the \emph{Chron. Pasch.} (\emph{c.} A.D. 630) , who reproduces the opening words of the Greens. He substitutes plural verbs for singular, but otherwise agrees closely with Theophanes. Both writers probably copied from a sixth-century chronicle.

} It seems likely that Calopodius who had offended the Greens is the same as Calopodius who was praepositus s. cub. in A.D. 558 (John Mal. XVIII p490).

As we are totally ignorant of the circumstances, a great part of this allusive dialogue is very obscure. Some act on the part of the chamberlain Calopodius had excited the anger of the Greens; they begin by complaining of this in respect ­ful tones, and obtaining no satisfaction go on to air their grievances as an oppressed party, with violent invective. A mandator or herald speaks for the Emperor, standing in front of the kathisma, and the Greeks evidently have a single spokesman.

\emph{Greens.} Long may you live, Justinian Augustus! Tu vincas. I am oppressed, O best of sovrans, and my grievances, God knows, have become intolerable. I fear to name the oppressor, lest he prosper the more and I endanger my own safety.

\emph{Mandator.} Who is he? I know him not.

\emph{Greens.} My oppressor, O thrice august! is to be found in the quarter of the shoemakers.\footnote{\textgreek{Ε\texthebrew{ἰ}ς τ\texthebrew{ὰ} τζαγγαρε\texthebrew{ῖ}α ε\texthebrew{ὑ}ρίσκεται} \texthebrew{—} an allusion to the name of Calo-podius (as Maas points out).

}

\emph{Mandator.} No one does you wrong.

\emph{Greens.} One man and one only does me wrong. Mother of God, may he be humbled ( \textgreek{μ\texthebrew{ὴ ἀ}νακεφαλίσ\texthebrew{ῃ}} )!

\emph{Mandator.} Who is he? We know him not.

\emph{Greens.} Nay, you know well, O thrice august! I am oppressed this day.

\emph{Mandator.} We know not that anyone oppresses you.

\emph{Greens.} It is Calopodius, the spathar, who wrongs me, O lord of all!

\emph{Mandator.} Calopodius has no concern with you.\footnote{\textgreek{Ο\texthebrew{ὐ}κ \texthebrew{ἔ}χει πρ\texthebrew{ᾶ}γμα.}}

\emph{Greens.} My oppressor will perish like Judas; God will requite him quickly.

\emph{Mandator.} You come, not to see the games, but to insult your rulers.

\emph{Greens.} If anyone wrongs me, he will perish like Judas.

\emph{Mandator.} Silence, Jews, Manichaeans, and Samaritans!

\emph{Greens.} Do you disparage us with the name of Jews and Samaritans? The Mother of God is with all of us.

\emph{Mandator.} When will ye cease cursing yourselves?

\emph{Greens.} If anyone denies that our lord the Emperor is orthodox, let him be anathema, as Judas.

p73
\emph{Mandator.} I would have you all baptized in the name of one God.

\emph{Greens.} I \emph{am} baptized in One God.\footnote{The Greens apparently take up the words of the mandator, \textgreek{ε\texthebrew{ἰ}ς \texthebrew{ἕ}να βαπτίζεσθαι}, in a Monophysitic sense. The preceding words, \textgreek{ο\texthebrew{ἱ} δ\texthebrew{ὲ} πράσινοι} º \textgreek{\texthebrew{ἐ}βόησαν \texthebrew{ἐ}ράνω \texthebrew{ἀ}λλήλων κα\texthebrew{ὶ ἔ}κραζον \texthebrew{ὡ}ς \texthebrew{ἐ}κέλευσεν \texthebrew{Ἄ}ντλας}, seem to imply that while the conversation was throughout conducted by a spokesman (Antlas?), here the whole party shouted together.

}

\emph{Mandator.} Verily, if you refuse to be silent, I shall have you beheaded.

\emph{Greens.} Every person seeks a post of authority, to secure his personal safety. Your Majesty must not be indignant at what I say in my tribulation, for the Deity listens to all complaints. We have good reason, O Emperor! to mention all things \emph{now.}\texthebrew{⁠}\footnote{\textgreek{\texthebrew{Ὀ}νομάζομεν \texthebrew{ἄ}ρτι πάντα}. The sense demands that \textgreek{\texthebrew{ἄ}ρτι} should be the emphatic word.

} For we do not even know where the palace is, nor what is the government. If I come into the city once, it is sitting on a mule;\texthebrew{⁠}\footnote{\textgreek{\texthebrew{Ὄ}ταν ε\texthebrew{ἰ}ς βορδώνιν καθέζομαι}. Prisoners were drawn by mules to execution or punishment.

} and I wish I had not to come then, your Majesty.\footnote{\textgreek{Ε\texthebrew{ἴ}θοις μηδ\texthebrew{ὲ} τότε, τρισαύγουστε.}}

\emph{Mandator.} Every one is free to move in public, where he wishes, without danger.

\emph{Greens.} I am told I am free, yet I am not allowed to use my freedom. If a man is free but is suspected as a Green, he is sure to be publicly punished.

\emph{Mandator.} Have ye no care for your lives that ye thus brave death?

\emph{Greens.} Let this (green) colour be once uplifted\texthebrew{⁠}\footnote{\textgreek{\texthebrew{Ἐ}παρθ\texthebrew{ῇ} τ\texthebrew{ὸ} χρ\texthebrew{ῶ}μα το\texthebrew{ῦ}το κα\texthebrew{ὶ ἡ} δίκη ο\texthebrew{ὐ} χρηματίζει.}} \texthebrew{—} then justice disappears. Put an end to the scenes of murder, and let us be lawfully punished. Behold, an abundant fountain; punish as many as you like. Verily, human nature cannot tolerate these two (contradictory) things. Would that Sabbatis had never been born, to have a son who is a murderer. It is the twenty-sixth murder that has been committed in the Zeugma;\texthebrew{⁠}\footnote{De Boor prints \textgreek{ε\texthebrew{ἰ}κότως \texthebrew{ἕ}κτος}. Sabbatius, it will be remembered, was Justinian's father.

} the victim was a spectator in the morning, in the afternoon, O lord of all! he was butchered.

\emph{Blues.} Yourselves are the only party in the hippodrome that has murderers among their number.

\emph{Greens.} When ye commit murder ye leave the city in flight.

\emph{Blues.} Ye shed blood, and debate. Ye are the only party here with murderers among them.

\emph{Greens.} O lord Justinian! they challenge us and yet no one slays them. Truth will compel assent.\texthebrew{⁠}\footnote{\textgreek{Νοήσει \texthebrew{ὁ} μ\texthebrew{ὴ} θέλων.}} Who slew the woodseller in the Zeugma, O Emperor?

\emph{Mandator.} Ye slew him.

\emph{Greens.} Who slew the son of Epagathus, Emperor?

\emph{Mandator.} Ye slew him too, and ye slander the Blues.

\emph{Greens.} Now have pity, O Lord God! The truth is suppressed. I should like to argue with them who say that affairs are managed by God. Whence comes this misery?

p74
\emph{Mandator.} God cannot be tempted with evil.

\emph{Greens.} God, you say, cannot be tempted with evil? Who is it then who wrongs me? Let some philosopher or hermit explain the distinction.

\emph{Mandator.} Accursed blasphemers, when will ye hold your peace?

\emph{Greens.} If it is the pleasure of your Majesty, I hold my peace, albeit unwillingly. I know all \texthebrew{—} all, but I say nothing. Good\texthebrew{‑}bye, Justice! you are no longer in fashion. I shall turn and become a Jew. Better to be a ``Greek'' (pagan) than a Blue, God knows.

\emph{Blues.} You are detestable, I cannot abide the sight of you. Your enmity dismays me.

\emph{Greens.} Let the bones of the spectators be exhumed.\texthebrew{⁠}\footnote{\emph{I.e.} let them be murdered. A customary form of curse in the Hippodrome, cp. Theophanes, A.M. 6187 \textgreek{\texthebrew{ἀ}νασκαφ\texthebrew{ῇ} τ\texthebrew{ὰ ὀ}στέα \texthebrew{Ἰ}ουστινιανο\texthebrew{ῦ}} (Justinian II).

}

The language of this astonishing dialogue obeys metrical laws, which concern not quantity but the number of the syllables and the accentuation of the last word in each clause. The most frequently occurring form is five syllables with the penultimate accented + four with the antepenultimate (or ultimate?) accented, \emph{e.g.}:

It is evident that to converse in metrical chant both the Imperial mandator and the spokesmen of the demes must have had a special training in the art of improvising.\footnote{An earlier example of metrical cries is preserved in inscriptions on the monument of the famous charioteer Porphyrius (reign of Anastasius). See Woodward's publication in the Appendix to George, \emph{Church of Saint Eirene}; and his paper in the \emph{Annual of the British School at Athens}, XVII.88 \emph{sqq.} The dialogue has considerable interest as a sample of the spoken Greek of the sixth century.

\texthebrew{◂} previous

next \texthebrew{▸}Images with borders lead to more information.

The thicker the border, the more information.

(Details here.)

UP TO:

J. B. Bury
Homepage

Latin \& Greek

Texts

LacusCurtius

Home}

\xchapter{Chapter XVI}

Short URL for this page:

bit.ly/BURLAT16

Our records of the Persian war conducted by the generals of Anastasius, which was described in a former chapter, give us little information as to the character and composition of the Imperial army. But we may take it as probable that the military establishment was already of much the same kind as we find it a quarter of a century later in the reign of Justinian. In the course of the fifth century the organisation of the army underwent considerable changes which our meagre sources of information do not enable us to trace. During that period, since the early years of Theodosius II, we have no catalogue of the military establishment, no military treatises,\texthebrew{⁠}\footnote{With the exception of that of Vegetius, which does not help much. See above,

Vol. I, p250 . º

} no military narratives. When we come to the reign of Justinian, for which we have abundant evidence,\texthebrew{⁠}\footnote{Besides Procopius, the chief source, we have four tactical documents which supplement and illustrate his information. (1) Fragments of a tactical work by Urbicius, who wrote in the reign of Anastasius. (2) Anonymous \textgreek{Βυζαντιος Περ\texthebrew{ὶ} στρατηγικ\texthebrew{ῆ}ς}, and (3) a \textgreek{\texthebrew{ἑ}ρμήνεια} or glossary of military terms, from the reign of Justinian. (4) Pseudo-Mauricius, \emph{Stratêgikon}, from the end of the sixth century. For editions of these works see

Bibliography . In regard to the date of the \emph{Stratêgikon} (falsely ascribed to the emperor Maurice), it is quite clear that it was composed \emph{after} the reign of Justinian, and \emph{before} the institution of the system of Themes, which is probably to be ascribed to Heraclius. Thus we get as outside limits A.D. 585-\emph{c.} 615. It is º quite perverse to date it (with Vári and others) to the eighth century. For modern studies of the sixth-century armies see

Bibliography II.2, C

under Benjamin, Maspéro, Aussaresses, Grosse, Müller.

} we find that the old system of the fourth century has been changed in some important respects.

The great commands of the Masters of Soldiers, and the distinction between the comitatenses and the limitanei , have not

p76
been altered; but the legions, the cohorts, and the alae , the familiar units of the old Roman armies, have disappeared both in name and in fact, and to the comitatenses and limitanei has been added a new organisation, the foederati , a term which has acquired a different meaning from that which it bore in the fourth century.

The independent military unit is now the numerus , a company generally from 200 to 400 soldiers, but sometimes varying below or above these figures. In old days it was necessary to divide the legion for the purpose of garrisoning towns; on the new system each town could have a complete, or more than one complete unit. These companies were under the command of tribunes.\footnote{The Greek name of the numerus is \textgreek{\texthebrew{ἀ}ριθμός} or \textgreek{τάγμα} (Sozomen, \emph{H. E.} I.8 \textgreek{τ\texthebrew{ὰ Ῥ}ωμαίων τάγματα \texthebrew{ἃ} ν\texthebrew{ῦ}ν \texthebrew{ἀ}ριθμο\texthebrew{ὺ}ς καλο\texthebrew{ῦ}σι}); \textgreek{κατάλογος} is used in the same sense, \emph{e.g.} Procopius, \emph{B. P.} I.15. For the evidence as to its strength cp. Maspéro, \emph{op. cit.} 116 \emph{sq.}, who remarks that it was a tactical principle to vary the strength of the numeri in order to deceive the enemy (cp. Pseudo-Maurice, \emph{Strat.} I.4 \emph{ad fin.} \textgreek{χρ\texthebrew{ὴ} μηδ\texthebrew{ὲ} πάντα τ\texthebrew{ὰ} τάγματα \texthebrew{ἐ}πιτηδεύειν πάντως \texthebrew{ἴ}σα ποιε\texthebrew{ῖ}ν κτλ}.). But the theoretical strength of the infantry numerus which Urbicius and the tacticians of Justinian's reign call \textgreek{σύνταγμα} was 256 (\emph{\textgreek{\texthebrew{Ἑ}ρμήνεια}} 12, cp. Pseudo-Maurice, XII.8; Urbicius says 250). These authorities nearly agree as to the tactical divisions of an army. The chief division, according to the \emph{\textgreek{\texthebrew{Ἑ}ρμήνεια}}, are: phalanx = 4096, meros = 2048, chiliarchia = 1024, º pentakosarchia =512, syntagma (tagma) = 256, taxis = 128, tetrarchia = 64, lochos = 16 (sometimes 8 or 12; Urbicius says 25). Pseudo-Maurice contemplates rather higher figures: the tagma should vary from 300 to 400 as a maximum; the chiliarchy, from 2000 to 3000; the meros, which consists of \textgreek{μο\texthebrew{ῖ}ραι}, should not exceed 6000 or 7000 (\emph{Strat.} 1.4).

}

Apart from the guard-troops stationed in the capital, the armed forces of the Empire fall into five principal categories. (1) The technical name comitatenses is little used. These troops, who are recruited almost exclusively among subjects of the Empire chiefly in the highlands of Thrace, Illyricum and Isauria, are now generally distinguished as stratiotai , regular Roman soldiers, from the other sections of the army.\texthebrew{⁠}\footnote{It is notable that Procopius sometimes speaks of the Isaurian regiments if they were distinct from the other Roman troops (\textgreek{κατάλογοι}), as in \emph{B. G.} I.5.2; but they were included among the stratiôtai.

}

(2) The limitanei perform the same duty of protecting exposed frontiers, and on the same conditions as before.

(3) The foederati , who must have been organised in the fifth century, are the new and striking feature which is revealed to us by the history of the campaigns of Belisarius. They are the most useful part of the field army, and they consist entirely of

p77
cavalry. They were originally recruited exclusively from barbarians, who volunteered for Imperial service, and were organised as Roman troops under Roman officers;\texthebrew{⁠}\footnote{The position of the Foederati was misconceived by Mommsen and by Benjamin (who held that they were recruited by Roman officers as a private speculation) and has been elucidated by J. Maspéro (\emph{Organ. mil.} and \emph{\textgreek{Φοιδεράτοι}}). His arguments seem to me convincing. The growth of the Federate troops was gradual, and appears to have begun in the reign of Honorius (Olympiodorus, \emph{fr.} 7). Areobindus is a Count of the Federates under Theodosius II (John Mal. XIV.364); in the time of Anastasius, Patriciolus (Theophanes, A.M. 6005) and probably Vitalian held the same post. There was a special bureau of \textgreek{χαρτουλάπιοι φοιδερ\texthebrew{ὰ}των} to deal with the payment of these troops ( \emph{C. J.} XII.37.19 , probably a law of Anastasius), who seem to have been considered more honourable and doubtless received higher pay than the comitatenses. For the technical use of Stratiôtai see Justinian, \emph{Nov.} 116 \textgreek{στρατι\texthebrew{ῶ}ται κα\texthebrew{ὶ} φοιδερ\texthebrew{ᾶ}τοι}, \emph{Nov.} 117.11; Procopius, \emph{B. P.} I.17.46 \textgreek{\texthebrew{Ῥ}ωμα\texthebrew{ῖ}οι στρατι\texthebrew{ῶ}ται}, \emph{B. V.} I.11.2; \emph{B. G.} IV.26.10.

} but in the sixth century Roman subjects were not debarred from enlisting in their companies.\texthebrew{⁠}\footnote{Procopius, \emph{B. V.} I.11.

} The degradation of the term Federates to designate these forces was not very happy, and it has naturally misled modern historians into confusing them with (4) the troops to whom the name was properly applied in the fourth century, and who are now distinguished as \emph{Allies}:\texthebrew{⁠}\footnote{\textgreek{Σύμμαχοι.}} the bands of barbarians, Huns for instance, or Heruls, who, bound by a treaty with the Empire, furnished, in return for land or annual subsidies, armed forces which were led by their native chiefs.

To these we must add (5) another class of fighting men, who were not in the employment of the government, the private retainers of the military commanders. The rise of the custom of keeping bands of armed followers has already been noticed.\texthebrew{⁠}\footnote{See above,

Chap. II § 2 .

} It was adopted not only by generals and Praetorian Prefects, but by officers of subordinate rank and wealthy private persons.\texthebrew{⁠}\footnote{Benjamin (\emph{op. cit.} 24 \emph{sqq.}) has collected instances from Procopius and Agathias. Egyptian papyri supply evidence for the employment of these Bucellarians in Egypt by large landowners. See the instances cited by Maspéro, \emph{Organ. milit.} 66 \emph{sqq.}

} The size of the retinues depended upon the wealth of the employer. Belisarius, who was a rich man, kept at one time as many as 7000.\footnote{Procopius, \emph{B. G.} III.1.20. Valerian, Mag. mil. of Armenia, had more than 1000 retainers (\emph{ib.} XXVII.3). Narses had less than 400 (Agathias, I.19).

}

There were two distinct classes of retainers, the hypaspistai , shield-bearers , who were the rank and file, and the doryphoroi , spear-bearers , who were superior in rank, fewer in number, and corresponded to officers. Belisarius himself and Sittas had been

p78
doryphoroi in the retinue of Justinian before he ascended the throne. The doryphoroi on accepting service were obliged to take a solemn oath not only of fidelity to their employer, but also of loyalty to the Emperor,\texthebrew{⁠}\footnote{\emph{B. V.} II.18.6. The superior position of the doryphoroi is illustrated by the fact that individual hypaspistai are very seldom named by Procopius, whereas he mentions by name 47 doryphoroi. Benjamin, \emph{op. cit.} 32\texthebrew{‑}33.

} a circumstance which implies an official recognition by the government. They were often employed on confidential missions, they stood in the presence of their master at meals, and attended him closely in battle. Both the doryphoroi and the hypaspistai seem to have been entirely mounted troops. The majority of them were foreigners (Huns and Goths), or mountaineers of Thrace and Asia Minor.

As a rule, in the campaigns of the sixth century, we find the armies composed mainly of comitatenses and foederati, but always reinforced by private retainers and barbarian allies. A single army in the field generally numbered from 15,000 to 25,000 men, a figure which probably it seldom exceeded; 40,000 was exceptionally large. The total strength of the Imperial army under Justinian was reckoned at 150,000.\footnote{Agathias, V.13 \emph{ad fin.} The figure is probably very close to the truth.

}

The tactics and equipment of the Imperial armies had been considerably altered by the necessity of adapting them to the military habits of their oriental foes. At this time, in establishment and equipments, the Persians differed so little from the Romans that a Roman corps might have appeared in a Persian, or a Persian in a Roman army, with little sense of discrepancy. The long eastern warfare of the third and fourth centuries had been a school in which the Romans transformed in many ways their own military traditions and methods. They adopted from their adversaries elaborate defensive armour, cuirasses, coats of mail, casques and greaves of metal. At the end of the fourth century there were cuirassiers forming \emph{corps d'élite}, and in the sixth these heavily armed ``iron cavalry''\texthebrew{⁠}\footnote{Ferreus equitatus,

Amm. Marc. XIX.1.2 .

} ( catafractarii ) have become a still larger and more important section of the army. Another result of the eastern wars was the universal practice of archery, which the old Roman legions despised. The heavy cavalry were armed with bow and arrows as well as with lance and sword.

In his old age king Kavad was troubled and anxious about the succession to his throne, which he desired to secure to Chosroes his favourite son. But Chosroes was not the eldest, and his father feared that when he died the Persian nobles would prefer one of the elder brothers and put Chosroes to death. Accordingly he conceived the idea of pla­cing his favourite under the protection of the Roman Emperor, as Arcadius had recommended Theodosius to the protection of Yezdegerd. But his proposal took a strange form. He asked Justin to adopt Chosroes. Both Justin and Justinian were at first attracted by the proposal, but the influence of the quaestor Proclus induced them to refuse. Proclus, who viewed the matter as a lawyer, represented the request as insidious; for the adopted son might assert a claim to the father's inheritance; the Persian king might claim the Roman Empire.

The refusal of his request was deeply resented by Kavad, and there were causes of friction in the Caucasian regions which led to a new breach between the two great powers.\texthebrew{⁠}\footnote{Legally the two powers seem to have been in a state of war, for the armistice of seven years (A.D. 505) had not been renewed. This may be inferred from the statement of John Malalas (XVIII p478) that the peace of 532 terminated in a war which had lasted for 31 years, \emph{i.e.} since 502.

} Both governments were actively pushing their interests in that part of the world.

The Pontic provinces, as well as Roman Armenia, constantly suffered from the depredations of the Tzani, a heathen people who maintained their independence in an inland district on the borders of Colchis and Armenia, and lived by brigandage. The Imperial government was in the habit of giving them a yearly allowance to purchase immunity, but they paid little regard to the contract. One of the achievements of Justin's peaceful reign was partially to civilise these wild mountaineers. Sittas, the brother-in\texthebrew{‑}law of Theodora, was sent against them. He subdued them, enrolled them in the Roman armies, and they were induced to embrace Christianity.\footnote{Justinian,

\emph{Nov.} 28 . Proc. \emph{B. P.} 1.15. They returned to their old marauding habits and had to be reduced again in A.D. 558. Agathias, V.1.2.

}

The reduction of the Tzani proved to be a preliminary to a more active policy in Caucasian countries. South of the

p80
great range, between the Euxine and the Caspian, lay three kingdoms: in the west, Colchis, the land of the Lazi, whose name is still preserved in Lazistan; in the centre, Iberia or Georgia; and in the east, almost beyond Roman vision, Albania,

The importance of Lazica, in Roman eyes, was twofold. It was a barrier against a Persian advance through Iberia to the coasts of the Black Sea.\texthebrew{⁠}\footnote{For Roman interference in the domestic affairs of Colchis in the reign of Marcian see Priscus, \emph{fr.} 8 \emph{De leg. Rom., fr.} 12 \emph{De leg. gent.} (cp. also \emph{frs.} 16 and 22).

} In the reign of Justin, Tzath, the king of the Lazi, who had hitherto been friendly to Persia, visited Constantinople and became a client of the Emperor.\texthebrew{⁠}\footnote{John Mal. XVIII p412. He was baptized a Christian and married a Roman lady, Valeriana, daughter of Nomos a patrician. Justin crowned him, and the chronicler describes his royal robes at some length.

} Perhaps this change of policy was caused by the development of Persian designs in Iberia. This country had long been a client state of Persia, but it was devoted to the Christian faith. Kavad either resolved to assimilate it to Persian civilisation or sought a pretext for invading it, and he issued a command to the Iberians to abandon the custom of burying their dead. Gurgenes, the Iberian king, turned to the Roman Emperor for protection.\texthebrew{⁠}\footnote{Justin sent Probus, the nephew of Anastasius, with a large sum of money, to Bosporus, to induce the Huns of the Crimea to help the Iberians; but he was unsuccessful (Proc. \emph{B. P.} 1.12).

} A force was sent to Lazica, while a Persian army invaded Iberia, and Gurgenes, with his family, fled within the Lazic borders and proceeded to Constantinople. Roman garrisons were placed in the Lazic forts on the Iberian frontier,\texthebrew{⁠}\footnote{But they soon departed, and the natives were unable to defend the forts against the Persians. Proc. \emph{B. P.} I.13, p58. Sittas was a mag. mil. in praes., and he was now appointed to the newly created post of Mag. mil. per Armeniam. He seems to have held the two posts concurrently. During peace his headquarters were at Constantinople. See Proc. \emph{B. P.} I.15, p74, II.3, p154; John Mal. XVIII p429.

} and Sittas with Belisarius, who now first appears upon the scene, made a success ­ful incursion into Persarmenia. In a second expedition the Romans were defeated by two able commanders, Narses and Aratius, who afterwards deserted and entered Roman service.\footnote{Procopius, \emph{ib.} 1.12. This is probably the incursion (noticed by John Mal. p427) under Gilderic and others.

}

Thus the war began before the death of Justin. Perhaps it might have been averted if his successor had not determined to

p81
build a new fortress near Daras. Belisarius, who had been appointed commandant of Daras, was directed to build the work, and as the building operations were progressing, a Persian army, 30,000 strong, under the prince Xerxes, invaded Mesopotamia ( A.D. 528).\texthebrew{⁠}\footnote{John Mal. p441. For the events of 528 we have to combine Procopius (\emph{B. P.} I.13) and Malalas. The two narratives are carefully compared by Sotiriadis in \emph{Zur Kr. v. Joh. v. Ant.} p114 \emph{sq.} It is to be noted that Belisarius held only a subordinate position and was in no way responsible for the defeat. The operations of 529 are entirely omitted by Procopius. For the fortress at Minduos, which the Romans tried to build, see Proc. \emph{ib.} and Zacharias Myt. IX.2. º

} The Romans, under several leaders who had joined forces, were defeated in a disastrous battle; two of the commanders were slain and three captured. Belisarius luckily escaped. The foundations of the new fortress were left in the hands of the enemy. But the victors had lost heavily and soon retreated beyond the frontier. Justinian sent more troops and new captains to the fortresses of Amida, Constantia, Edessa, Sura, and Beroea; and formed a new army (of Illyrians and Thracians, Scythians and Isaurians) which he entrusted to Pompeius, probably the nephew of Anastasius.\texthebrew{⁠}\footnote{John Mal. p442. Pompeius was a patrician, and it is not very likely that there were two patricians of this name.

} But no further operations are recorded in this year, which closed with a severe winter.

The hostilities of A.D. 529 began in March with a combined raid of Persian and Saracen forces, under the guidance of Mundhir, king of Hira, who penetrated into Syria, almost to the walls of Antioch, and retreated so swiftly that the Romans could not intercept him. Reprisals were made by a body of Phrygians who plundered Persian and Saracen territory (April). Pompeius seems to have accomplished nothing, and Belisarius was appointed Master of Soldiers in the East.\texthebrew{⁠}\footnote{In succession to Hypatius (before June) acc. to John Mal. p445. Hypatius (\emph{ib.} 423) had been created mag. mil. Or. between April and August 527. It is difficult to reconcile this with the statements of Procopius, who places both the appointment and the deposition of Hypatius before April 527 º (\emph{B. P.} I.11 p53 and p55, compared with I.13 p59). It is possible that the notice of the deposition is an anticipation; the whole section beginning \textgreek{μετ\texthebrew{ὰ} δέ}, p54, to end of chap. II may be a chronological digression. But Zacharias, IX.1, states that Timus (otherwise unknown, perhaps an error for Timostratus) was mag. mil. when Justin died, and that Belisarius succeeded him (IX.2). If this is right, Malalas is wrong.

} The rest of the year was occupied with ineffectual negotiations.\footnote{It is remarkable that in the summer of 529 Justinian should have sent the customary friendly embassy to announce his accession. Hermogenes was the envoy (John Mal. pp447\texthebrew{‑}448). He returned with a letter from Kavad, of which the text is given (\emph{ib.} p449).

}

p82
Belisarius was now to win his military laurels at the early age of twenty-five . There was still talk of peace, but Kavad seems not to have really desired it, and the ambassador, Rufinus, waited idle at Hierapolis. Hermogenes, the Master of Offices, was sent out to help the young general with his experience, and they concentrated at Daras an army of 25,000 mixed and undisciplined troops. Perozes, who had been appointed mihran or commander-in\texthebrew{‑}chief of the Persian army, arrived at Nisibis in June\texthebrew{⁠}\footnote{Theophanes supplies the date.

} ( A.D. 53), at the head of 40,000 troops, confident of victory. They advanced within two miles of Daras, and the mihran sent to Belisarius a characteristically oriental message, that, as he intended to bathe in the city on the morrow, a bath should be prepared for his pleasure.

The Romans made preparations for battle, just outside the walls of the town. The Persians arrived punctually as their general signified, and stood for a whole day in line of battle without venturing to attack the Romans, who were drawn up in carefully arranged positions. In the evening they retired to their camp,\texthebrew{⁠}\footnote{During the afternoon the armies were diverted by two single combats, in which a Byzantine professor of gymnastics, who had accompanied the army unofficially, slew two Persian champions.

} but returned next morning, resolved not to let another day pass without a decisive action, and found their enemy occupying the same positions as on the previous day. They were themselves now reinforced by a body of 10,000, which arrived from Nisibis. The Roman dispositions were as follows:

About a stone's throw from the gate of Daras that looks towards Nisibis a deep trench was dug, interrupted by frequent ways for crossing. This trench, however, was not in a continuous right line; it consisted of five sections. At each end of a short central trench, which was parallel to the opposite wall of the city, a trench ran outwards almost at right angles; and where each of these perpendicular trenches or ``horns'' terminated, two long ones were dug in opposite directions at right angles, and consequently almost parallel to the first trench. Between the trenches and the town Belisarius and Hermogenes were posted with the infantry. On the left, behind the main ditch and near the left ``horn,'' was a regiment of cavalry under Buzes, and 300 Heruls under Pharas were stationed on a rising ground, which the Heruls occupied in the morning, at the

p83
suggestion of Pharas and with the approval of Belisarius. Outside the angle made by the outermost ditch and the horn were placed 600 Hunnic cavalry, under the Huns Sunicas and Aigan. The disposition on the right wing was exactly symmetrical. Cavalry under John (the son of Nicetas), Cyril, and Marcellus occupied the position corresponding to that occupied by Buzes on the left, while other squadrons of Hunnic horse, led by Simas and Ascan, were posted in the angle.

Half of the Persian forces stood in a long line opposite to the Roman dispositions, the other half was kept in reserve at some distance in the rear. The mihran commanded the centre, Baresmanas the left wing, and Pityaxes the right. The corps of Immortals, the flower of the army, was reserved for a supreme occasion. The details of the battle have been described by a competent eye-witness .\footnote{Procopius, \emph{B. P.} I.14. A diagram will make the arrangement of the forces clear.

IMAGE ZZZ

}

As soon as noon was past the barbarians began the action. They had reserved the engagement for this hour of the day because they are themselves in the habit of eating only in the evening, while the Romans eat at noontide, so that they counted on their offering a less vigorous resistance if they were attacked fasting. At first each side discharged volleys of arrows and the air was obscured with them; the barbarians shot more darts, but many fell on both sides. Fresh relays of the barbarians were always coming up to the front, unperceived by their adversaries; yet the

p84
Romans had by no means the worst of it. For a wind blew in the faces of the Persians and hindered to a considerable degree their missiles from operating with effect. When both sides had expended all their arrows, they used their spears, hand to hand. The left wing of the Romans was pressed most hardly. For the Cadisenes, who fought at this point with Pityaxes, had advanced suddenly in large numbers, and having routed their opponents, pressed them hard as they fled, and slew many. When Sunicas and Aigan with their Huns saw this they rushed on the Cadisenes at full gallop. But Pharas and his Heruls, who were posted on the hill, were before them (the Huns) in falling on the rear of the enemy and performing marvellous exploits. But when the Cadisenes saw the cavalry of Sunicas also coming against them from the side, they turned and fled. The rout was conspicuous when the Romans joined together and great slaughter was inflicted on the enemy.

The mihran [meanwhile] secretly sent the Immortals with other regiments to the left wing. When Belisarius and Hermogenes saw them, they commanded Sunicas, Aigan, and their Huns, to go to the angle on the right where Simas and Ascan were stationed, and placed behind them many of the retainers of Belisarius. Then the left wing of the Persians, led by Baresmanas, along with the Immortals, attacked the Roman right wing at full speed. And the Romans, unable to withstand the onset, fled. Then those who were stationed in the angle (the Huns, etc.) attacked the pursuers with great ardour. And coming athwart the side of the Persians they cleft their line in two unequal portions, the larger number on the right and a few on the left. Among the latter was the standard-bearer of Baresmanas, whom Sunicas killed with his lance. The foremost of the Persian pursuers, apprehending their danger, turned from their pursuit of the fugitives to oppose the attackers. But this movement placed them between enemies on both sides, for the fugitive party perceived what was occurring and rallied. Then the other Persians and the corps of the Immortals, seeing the standard lowered and on the ground, rushed with Baresmanas against the Romans in that quarter. The Romans met them, and Sunicas slew Baresmanas, hurling him to earth from his horse. Then the barbarians fell into great panic, and forgot their valour and fled in utter disorder. And the Romans closed them in and slew about five thousand. And thus both armies were entirely set in motion; that of the Persians for retreat and that of the Romans for pursuit. All the infantry of the defeated army threw away their shields, and were caught and slain pell-mell. Yet the Romans pursued only for a short distance, for Belisarius and Hermogenes would not permit them to go further, lest the Persians, compelled by necessity, should turn and rout them if they followed rashly; and they deemed it sufficient the keep the victory untarnished, this being the first defeat experienced by the Persians for a long time past.\texthebrew{⁠}\footnote{It is curious that Zacharias, IX.3, in his notice of the battle, does not mention Belisarius. He names Sunicas, Buzes, and Simuth (Simas?).

}

It will be observed that this battle \texthebrew{—} the first of which we have any full description since the fourth century \texthebrew{—} was fought and

p85
won entirely by cavalry. It has been pointed out that the dispositions of Belisarius show his ``deliberate purpose to keep his infantry out of the stress of the fight.''\texthebrew{⁠}\footnote{Oman (\emph{Art of War}, p29), who has well elucidated the battle. In one point I disagree with his plan. The central trench (C) was evidently, from the description of Procopius, much shorter than the wing trenches (A, A´), and the lines of infantry must have extended considerably beyond it on either side. But this only brings out and confirms his interpretation of the tactical plan of Belisarius, to force the enemy to attack the wings.

} This was done by throwing forward the wings, and leaving only a comparatively short space between them, so that they drew upon themselves the chief attack of the enemy. We are not told how the Persians disposed their horse and foot. The foot may have been in the centre. But the fighting was evidently done by the cavalry, for the infantry was not efficient. Belisarius, addressing his soldiers before the battle, described the Persian infantry as ``a crowd of miserable peasants who only come into battle to dig through walls and strip the slain and generally to act as servants to the soldiers (that is, the cavalry).'' We may conjecture that while in mere numbers the Romans were fighting one to two, the great excess of the Persian forces was chiefly in the infantry, and that otherwise they were not so unevenly matched.

About the same time the Roman arms were also success ­ful in Persarmenia, where a victory was gained over an army of Persarmenians and Sabir auxiliaries, which, if it had not been overshadowed by the victory of Daras, would have probably been made more of by the Greek historians.\footnote{About this time Narses the Persarmenian, with his two brothers, deserted to Rome (Proc. \emph{B. P.} I.15).

}

After the conspicuous defeat which his army had experienced, Kavad was not disinclined to resume negotiations, and embassies passed between the Persian and Roman courts;\texthebrew{⁠}\footnote{See Proc. \emph{B. P.} I.16; John Mal. p454. Sotiriadis (\emph{op. cit.} p119) points out the difficulties in the text and gives a probable solution. For the Samaritan rising, \emph{ib.} 445.

} but at the last moment the persuasions and promises of fifty thousand Samaritans induced him to break off negotiations on a trifling pretext. The Samaritans had revolted in A.D. 529, and the fifty thousand, who had escaped the massacre which attended the suppression of the rebellion, actuated by the desire of revenge, engaged to betray Jerusalem and Palestine to the foe of the Empire. The plot, however, was discovered and forestalled.

In the following spring ( A.D. 531), at the instigation of

p86
Mundhir, in whose advice Kavad had great confidence, fifteen thousand Persian cavalry under Azareth crossed the Euphrates at Circesium with the intention of invading Syria. They marched along the banks of the river to Callinicum, thence by Sura to Barbalissus, whence taking the western road they pitched their camp at Gabbula, \texthebrew{•} twelve miles from Chalcis, and harried the neighbourhood. Meanwhile Belisarius arrived at Chalcis, where he was joined by Saracen auxiliaries under Harith. His army was 22,000 strong, but he did not venture to attack the enemy, who numbered 30,000, and his inactivity aroused considerable discontent among both officers and soldiers.\texthebrew{⁠}\footnote{Procopius, \emph{B. P.} p92.

} The Hun captain Sunicas set at naught the general's orders, and attacking a party of Persians not only defeated them, but learned from the prisoners whom he took the Persian plan of campaign, and the intention of the foe to strike a blow at Antioch itself. Yet the success of Sunicas did not in the eyes of Belisarius atone for his disobedience, and Hermogenes, who arrived at this moment on the scene of action from Constantinople, arranged with difficulty the quarrel between the general and the captain. At length Belisarius ordered an advance against the enemy, who had meanwhile by their siege engines taken the fortress of Gabbula (near Chalcis) and other places in the neighbourhood. Laden with booty, the Persians retreated and reached the point of the right Euphrates bank opposite to the city of Callinicum, where they were overtaken by the Romans. A battle was unavoidable, and on the 19th of April the armies engaged. What really happened on this unfortunate day was a matter of doubt even for contemporaries; some cast the blame on Belisarius, others accused the subordinate commanders of cowardice.\footnote{Compare the conflicting accounts of Procopius (\emph{B. P.} I.18), the secretary of Belisarius, and Malalas. We have no means of determining the source of the latter, but in many cases he furnishes details omitted by the former. The account of Zacharias, IX.4 throws no light, but he mentions that the wind was blowing in the face of the Romans.

}

At Callinicum the course of the Euphrates is from west to east. The battle was fought on the bank of the river, and as the Persians were stationed to the east of the Romans, their right wing and the Roman left were on the river. Belisarius and his cavalry occupied the centre; on the left were the infantry and the Hunnic cavalry under Sunicas and Simas; on the right were Phrygians and Isaurians and the Saracen auxiliaries under

p87
their king Harith.\texthebrew{⁠}\footnote{I cannot agree with the plan of the battle implied by Sotiriadis (p123), which would place the Persians \emph{west} of the Romans. I adopt the reverse position, and thus bring the statements of Malalas into accordance with those of Procopius. In the mere fact of the position of troops there is no reason why the two accounts should differ. According to Sotiriadis, ``the northern part'' (\textgreek{τ\texthebrew{ὸ ἀ}ρκ\texthebrew{ῷ}ον μέρος}) of the Roman army was the right wing; according to my explanation, it was the left.

} The Persians began the action by a feigned retreat, which had the effect of drawing from their position the Huns on the left wing; they then attacked the Roman infantry, left unprotected, and tried to ride them down and press them into the river. But they were not as success ­ful as they hoped, and on this side the battle was drawn. On the Roman right wing the fall of Apscal, the captain of the Phrygian troops, was followed by the flight of his soldiers; a panic ensued, and the Saracens acted like the Phrygians; then the Isaurians made for the river and swam over to an island. How Belisarius acted, and what the Hun captains were doing in the meantime, we cannot determine. It was said that Belisarius dismounted, rallied his men, and made a long brave stand against the charges of the Persian cavalry. On the other hand, this valiant behaviour was attributed to Sunicas and Simas, and the general himself was accused of fleeing with the cowards and crossing to Callinicum. There is no clear evidence to prove that the defeat was the fault of Belisarius; though perhaps an over-confident spirit in his army prevailed on him to risk a battle against his better judgment.

The Persians retreated, and the remnant of the Roman army was conveyed across the river to Callinicum. Hermogenes\texthebrew{⁠}\footnote{It may be suspected that Hermogenes presented the behaviour of Belisarius in a suspicious light. He was a Hun, and sympathised doubtless with Sunicas and Simas.

} sent the news of the defeat to Justinian without delay, and the Emperor despatched Constantiolus to investigate the circumstances of the battle and discover on whom the blame, if any, rested. The conclusions at which Constantiolus arrived resulted in the recall of Belisarius and the appointment of Mundus to the command of the eastern armies.\texthebrew{⁠}\footnote{We cannot, I think, infer from the recall of Belisarius that the verdict of Constantiolus was adverse to him; on the contrary, if it had been adverse to him, the informant who furnished Malalas with his narrative, and who was evidently unfriendly to Belisarius, would have certainly stated the fact in distinct terms. Probably the reason of his recall was the circumstance that a bad feeling prevailed between him and the subordinate commanders; and Justinian saw that this feeling was a sure obstacle to success. The investigation of Constantiolus would naturally have shown up these jealousies and quarrels in the clearest light.

} It is significant of the difference

p88
between the spirit of the Persian and of the Roman governments that while Belisarius was recalled, with honour, after his defeat, the victorious Azareth was disgraced. He had been sent against Antioch and he had not approached it, and his victory had been bought with great losses.

The arms of Mundus were attended with success. Two attempts of the Persians to take Martyropolis were thwarted, and they experienced a considerable defeat. But the death of the old king Kavad and the accession of his son Chosroes (September 13, 531) led to the conclusion of a treaty which was known as ``the Endless Peace.'' The negotiations were conducted on the Roman side by Hermogenes and Rufinus, who was a \emph{grata persona} with Chosroes, and were protracted during the winter, because the Persians were unwilling to restore the forts they had taken in Lazica. They finally yielded and the treaty was ratified in spring A.D. 532.\texthebrew{⁠}\footnote{In the sixth year of Justinian, therefore after April 1. \emph{B. P.} p117.

} On their part the Romans restored two important fortresses in Persarmenia.\texthebrew{⁠}\footnote{Pharangion and Bôlon.

} The other conditions were that the Emperor would pay 11,000 lbs. of gold for the defence of the Caucasian passes, that the headquarters of the duke of Mesopotamia were no longer to be at Daras but at Constantia, and that the Iberian refugees at Constantinople might, as they chose, either remain there or returned to their own country.\footnote{Procopius, \emph{B. P.} I.22. John Mal. XVIII p477 states that the two monarchs agreed, as brothers, to supply each other with money or men in case of need. This may seem improbable, but such an agreement seems to have been made in a private treaty, see Joshua Styl. c. VIII. I conjecture that this refers to the treaty of 442: the stipulated help consisted of 300 able-bodied men \emph{or} 300 staters.

}

This treaty made no change in the frontiers between Roman and Persian Armenia. In the early years of Chosroes Persian Armenia was peaceful and contented under a native vassal prince and the Christians enjoyed full toleration. But at the same time the Armenian Church was drifting apart from Constantinople and Rome. The decisions of Chalcedon had been indeed accepted, but the Armenian theologians viewed them with some suspicion from the first; the ecclesiastical policy of Zeno and Anastasius confirmed them in their doubts; and the Henotikon of Zeno had been approved in a council held in A.D. 491. On the restoration of the doctrine of Chalcedon by Justin

p89
the Armenians displayed their Monophysitic leanings, and a definite and permanent schism between the Armenian and Greek Churches was the result. This separation was the work of the patriarch Narses, who secured the condemnation of the dogma of the Two Natures,\texthebrew{⁠}\footnote{Perhaps at a synod, c. 527. See Tournebize, \emph{Hist. de L'Arménie}, I.90\texthebrew{‑}91. Le Quien, \emph{Or. Christ.} I. pp1381\texthebrew{‑}1384.

} and at the Synod of Duin held just after his death, in A.D. 551, the independence of the Armenian Church was confirmed and a reform of the calendar was inaugurated. The Armenian era began on July 11, A.D. 552. The schism had its political consequences. Chosroes could profit by the fact that Greek influence declined in Persarmenia and Greek political agents were less favourably received.

The reign of Chosroes Nushirvan\texthebrew{⁠}\footnote{Khosru = Hu-srava (fair glory) is etymologically identical with \textgreek{ε\texthebrew{ὔ}}-\textgreek{κλεια}. The proper form of Nushirvan is Anosha-revan = of immortal soul.

} extended over nearly half of the sixth century, and may be called the golden or at least the gilded period of the monarchy of the Sassanids. His father Kavad had prepared the way for his brilliant son, as Philip of Macedon had prepared the way for Alexander. It was a period of energetic reforms, in some of which, as in the working out of a new land system, Chosroes was only continuing what his father had begun. This system was found to work so well that after their conquest of Persia the Saracen caliphs adopted it unaltered. In the general organisation some changes were made. The Persian empire was divided into four great circumscriptions each of which was governed by a marzban who had the title of ``king.'' The military government of these districts was now transferred to four spahbedhs , the civil government to four pādhospans , and the marzbans, though allowed to retain the honourable title, were reduced to second-class rank and were subordinate to the spahbedhs.\texthebrew{⁠}\footnote{For these changes see Stein, \emph{Ein Kapitel vom pers. Staate}, who (p66) would attribute the institution of the 4 pādhospans to Kavad.

} The most anxious pains of Chosroes were spent on the army, and it is said that when he reviewed it he used to inspect each individual soldier. He reduced its cost and increased its efficiency. But he also encouraged literature and patronised the study of Persian history. Of his personal culture the envy or impartiality of a Greek historian speaks with

p90
contempt as narrow and superficial;\texthebrew{⁠}\footnote{Agathias (II.28), who asks how one brought up in the luxury of an oriental barbarian could be a philosopher or a scholar. For the reception of Greek philosophers at the Persian court see

p370 .

} on the other hand, he has received the praises of an ecclesiastical writer. ``He was a prudent rather than a wise man, and all his lifetime he assiduously devoted himself to the perusal of philosophical work. And, as was said, he took pains to collect the religious books of all creeds, and read and studied them, that he might learn which were true and wise and which were foolish. . . . He praised the books of the Christians above all others, and said, 'These are true and wise above those of any other religion.' ``\texthebrew{⁠}\footnote{John Eph. VI.20. John apologises for thus eulogising a Magian and an enemy. What he says about the king's Christian proclivities is more edifying than probable. But Chosroes was not fanatical. He allowed one of his wives and her son to profess Christianity. In the eyes of Procopius, Chosroes was the typical oriental tyrant, cruel and perfidious.

} As a success ­ful and, judged by the standards of his age and country, enlightened ruler, Chosroes stands out in the succession of Sassanid sovrans much as Justinian stands out in the succession of the later Roman emperors.

The Emperor Justinian had, with the energy and thoroughness which distinguished the first half of his long reign, made use of the years of peace to strengthen the defences of the eastern provinces. Sieges were the characteristic feature of the wars on the oriental frontier, and walls were wellnigh as important as men. The fortifications of many of the most important cities and strongholds had fallen into decay, many had weak points, some were ill furnished with water. All the important towns in Mesopotamia and Osrhoene, and not a few of those in northern Syria were restored, repaired, or partly rebuilt in the reign of Justinian under the supervision of expert engineers. An account of these works has been preserved,\texthebrew{⁠}\footnote{Procopius, \emph{Aed.} book II .

} and most of them were probably executed between A.D. 532 and 539. The fortresses on the Pontic or Armenian border were similarly strengthened.\texthebrew{⁠}\footnote{\emph{Ib.} book III.

} Here, too, an important administrative change was made. Roman Armenia beyond the Euphrates, which had hitherto been governed by native satraps,\texthebrew{⁠}\footnote{Zeno abolished hereditary succession to the satrapies (except in the case of Belabitene), and vested the nomination in the Emperor (Procop. \emph{ib.} III.1) .

} under the general control of a military officer,\texthebrew{⁠}\footnote{The Comes Armeniae, who had been abolished in 528, when the mag. mil. per Arm. was created. See above,

p80 .

} was organised as a regular province

p91
under a governor of consular rank, and was officially designated as the Fourth Armenia. The satraps were abolished. Martyropolis was the chief town and residence of the governor.\footnote{A.D. 536. Justinian,

\emph{Nov.} 31, § 3 . At the same time considerable changes were made in the East Pontic provinces of 1st, 2nd, and 3rd Armenia, which will be noticed in another place (below,

p344 ).

}

When Chosroes concluded the ``Endless Peace'' with Justinian, he had little idea that the new Emperor was about to embark on great enterprises of conquest. Within seven years from that time ( A.D. 532\texthebrew{‑}539) Justinian had overthrown the Vandal kingdom of Africa, and had reduced the Moors; the subjection of the Ostrogothic lords of Italy was in prospect, Bosporus º and the Crimean Goths were included in the circle of Roman sway, while the Homerites of southern Arabia acknowledged the supremacy of New Rome. Both his friends and his enemies said, with hate or admiration, ``The whole earth cannot contain him; he is already scrutinising the aether and the remote places beyond the ocean, if he may win some new world.''\texthebrew{⁠}\footnote{Procopius, \emph{B. P.} II.3 (p160).

} The eastern potentate might well apprehend danger to his own kingdom in the expansion of the Roman Empire by the reconquest of its lost provinces. We may consider it natural enough that Chosroes should have seized or invented a pretext to renew hostilities, when it seemed but too possible that if Justinian were allowed to continue his career of conquest undisturbed the Romans might come with larger armies and increased might to extend their dominions in the East at the expense of the Sassanid empire.

Hostilities between the Saracens of Hira and their enemies of Ghassan supplied Chosroes with the pretext he desired. The Roman provinces had constantly suffered from the inroads of the Ghassanid tribes who obeyed no common ruler, and one of the early achievements of Justinian's reign was the creation of a Ghassanid state under the government of supreme phylarch, nominated by the Emperor. This client state formed a counterpoise to the Lakhmids of Hira, who were clients of Persia. Harith was appointed phylarch, and received the title of king and the dignity of patrician.\texthebrew{⁠}\footnote{See

Procopius, \emph{B. P.} I.17 ; Theophanes, A.M. 6056 . Harith reigned \emph{c.} A.D. 528\texthebrew{‑}570. Mundhir, the veteran chief of Hiras, was similarly allowed by the Persians to bear the title of king

(Procop. \emph{ib.}) . He reigned for about 50 years (A.D. 508\texthebrew{‑}554); see Tabari, p170 (Noeldeke's note). He was exceptionally

(p92)
barbarous. He sacrificed the son of his enemy Harith

(Procop. \emph{B. P.} II.28)

and on another occasion 400 nuns (cp. Noeldeke, \emph{l.c.}) to the goddess Uzza. For these kingdoms see Huart, \emph{Hist. des Arabes}, chap. IV.

} The cause of contention at this

p92
juncture between the two Saracen powers was a tract of waste land called Strata, to the south of Palmyra, a region barren of trees and fruit, scorched dry by the sun, and used as a pasture for sheep. Harith the Ghassanid could appeal to the fact that the name Strata was Latin, and could adduce the testimony of the most venerable elders that the sheep-walk belonged to his tribe. Mundhir, the rival sheikh, contented himself with the more practical argument that for years back the shepherds had paid him tribute. Two arbitrators were sent by the Emperor, Strategius, Count of the Sacred Largesses, and Summus, the duke of Palestine. This arbitration supplied Chosroes with a pretext for breaking the peace. He alleged that Summus made treasonable offers to Mundhir, attempting to shake his allegiance to Persia; and he professed to have in his possession a letter of Justinian to the Ephthalites, urging them to invade his dominions.\footnote{Procopius says that he does not know whether these allegations were true or false (\emph{B. P.} II.1). The second Book of his \emph{De bello Persico} is our main source for the war which ensued. It comes down to the end of A.D. 549.

}

About the same time suggestions from without urged the thoughts of Chosroes in the direction which they had already taken. An embassy arrived from Witigis, king of the Ostrogoths, now hard pressed by Belisarius, and pleaded with Chosroes to act against the common enemy ( A.D. 539).\texthebrew{⁠}\footnote{See below,
chap. XVIII § 9. The reader may ask how the details of this embassy were known. Procopius tells us in another place (\emph{B. P.} II.14) that the interpreter, returning from Persia, was captured near Constantia by John, duke of Mesopotamia, and gave an account of the embassy. The pseudo-bishop and his attendant remained in Persia.

} Another embassy arrived from the Armenians making similar representations, deploring and execrating the Endless Peace, and denouncing the tyranny and exactions of Justinian, against whom they had revolted. The history of Armenia, now a Roman province,\texthebrew{⁠}\footnote{See above,

p90 .

} had been unfortunate during the years that followed the peace. The first governor, Amazaspes, was accused by one Acacius of treachery, and, with the Emperor's consent, was slain by the accuser, who was himself appointed to succeed his victim. Acacius was relentless in exacting a tribute of unprecedented magnitude (£18,000); and some Armenians, intolerant of his cruelty, slew him and fled. The Emperor immediately despatched

p93
Sittas, the Master of Soldiers per Armeniam , to recall the people to a sense of obedience, and, when Sittas showed himself inclined to use the softer methods of persuasion, insisted that he should act with sterner vigour. The rebellion became general. Sittas was accidentally killed soon afterwards, but the rebels found themselves unequal to coping with the Roman forces, which were then placed under the command of Buzes, and they decided to appeal to the Persian monarch. The servitude of their neighbours the Tzani and the imposition of a Roman duke over the Lazi of Colchis confirmed them in their fear and detestation of Roman policy.

Accordingly Chosroes, in the autumn of A.D. 539, decided to begin hostilities in the following spring, and did not deign to answer a pacific letter from the Roman Emperor, conveyed by Anastasius, whom he retained an unwilling guest at the Persian court.\texthebrew{⁠}\footnote{Theodora also wrote a letter to Zabergan, whom she knew personally as he had come to Constantinople as an envoy, requesting him to urge Chosroes to preserve peace. But this letter may have been sent later, in 540 or even 541. Chosroes made use of it to quell discontent among his troops, arguing that a state must be weak in which women intervened in public affairs.

Procopius, \emph{H. A.} II.32\texthebrew{‑}36 . º

} The war which thus began lasted five years, and in each year the king himself took the field. He invaded Syria, Colchis, and Commagene in successive campaigns; in A.D. 543 he began but did not carry out an expedition against the northern provinces; in next year he invaded Mesopotamia; and in A.D. 545 a peace was concluded.

\xsubsection{I. Invasion of Syria ( A.D. 540)\texthebrew{⁠}\footnote{Procopius, \emph{B. P.} II.5\texthebrew{‑}14.

}}

Avoiding Mesopotamia, Chosroes advanced northwards with a large army along the left bank of the Euphrates. He passed the triangular city of Circesium, but did not care to assault it, because its walls, built by Diocletian, were too strong; while he disdained to delay at the town of Zenobia (Halebiya), named after the queen of Palmyra, because it was too insignificant. But when he approached Sura his horse neighed and stamped the ground; and the magi, who attended the king, seized the incident as an omen that the city would be taken. On the first day of the siege the governor was slain, and on the second the bishop of the place visited the Persian camp in the name of

p94
the dispirited inhabitants, and implored Chosroes with tears to spare the town. He tried to appease the implacable foe with an offering of birds, wine, and bread, and engaged that the men of Sura would pay a sufficient ransom. Chosroes dissembled the wrath he felt against the Surenes because they had not submitted immediately; he received the gifts and said that he would consult the Persian nobles regarding the ransom; and he dismissed the bishop, who was well pleased with the interview, under the honourable escort of Persian notables, to whom the monarch had given secret instructions.

``Having given his directions to the escort, Chosroes ordered

p95
his army to stand in readiness, and to run at full speed to the city when he gave the signal. When they reached the walls the Persians saluted the bishop and stood outside; but the men of Sura, seeing him in high spirits and observing how he was escorted by the Persians, put aside all thoughts of suspicion, and, opening the gate wide, received their priest with clapping of hands and acclamation. And when all had passed within, the porters pushed the gate to shut it, but the Persians placed a stone, which they had provided, between the threshold and the gate. The porters pushed harder, but for all their violent exertions they could not succeed in forcing the gate into the threshold-groove . And they did not venture to throw it open again, as they apprehended that it was held by the enemy. Some say that it was a log of wood, not a stone, that was inserted by the Persians. The men of Sura had hardly discovered the guile, ere Chosroes had come with all his army and the Persians had forced open the gate. In a few moments the city was in the power of the enemy.''\texthebrew{⁠}\footnote{Procopius, \emph{B. P.} II.5.

} The houses were plundered; many of the inhabitants were slain, the rest were carried into slavery, and the city was burnt down to the ground. Then the Persian king dismissed Anastasius, bidding him inform the Emperor in what place he had left Chosroes the son of Kavad.

Perhaps it was merely avarice, perhaps it was the prayers of a captive named Euphemia, whose beauty attracted the desired of the conqueror, that induced Chosroes to treat with unexpected leniency the princes of Sura. He sent a message to Candidus, the bishop of Sergiopolis, suggesting that he should ransom the 12,000 captives for 200 lbs. of gold (15s. a head). As Candidus had not, and could not immediately obtain, the sum, he was allowed to stipulate in writing that he would pay it within a year's time, under penalty of paying double and resigning his bishopric. Few of the redeemed prisoners survived long the agitations and tortures they had undergone.

Meanwhile the Roman general Buzes was at Hierapolis. Nominally the command in the East was divided between Buzes and Belisarius: the provinces beyond the Euphrates being assigned to the former, Syria and Asia Minor to the latter. But as Belisarius hadn't yet returned from Italy, the entire army was under the orders of Buzes.

p96
Informed of the presence of Chosroes in the Roman provinces, Justinian despatched his cousin Germanus to Antioch, with a small body of three hundred soldiers.\texthebrew{⁠}\footnote{Cp. John Mal. bk. XVIII p480. \emph{Contin. Marcell., s. a.}

} The fortifications of the ``Queen of the East'' did not satisfy the careful inspection of Germanus, for although the lower parts of the city were adequately protected by the Orontes, which washed the bases of the houses, and the higher regions seemed secure on impregnable heights, there rose outside the walls adjacent to the citadel\texthebrew{⁠}\footnote{The citadel was called Orocasias.

} a broad rock, almost as lofty as the wall, which would inevitably present to the besiegers a fatal point of vantage. Competent engineers said that there would not be sufficient time before the arrival of Chosroes to remedy this defect by removing the rock or enclosing it within the walls. Accordingly Germanus, despairing of resistance, sent Megas, the bishop of Beroea, to divert the Persian advance from Antioch by the influence of money or entreaties. The army had already crossed the Euphrates, and Megas arrived as it was approaching Hierapolis, from which Buzes had withdrawn a large part of the garrison. He was informed by the great king that it was his unalterable intention to subdue Syria and Cilicia. The bishop was constrained or induced to accompany the army to Hierapolis, which was strong enough to defy a siege, and was content to purchase immunity by a payment of 2000 lbs. of silver. Chosroes then consented to retire without assaulting Antioch on the receipt of 1000 lbs. of gold (£45,000), and Megas returned speedily to Beroea.\texthebrew{⁠}\footnote{Probably \emph{via} Batnae.

} From this city the avarice of the Sassanid demanded double the amount he had exacted at Hierapolis; the Beroeans gave him half the sum, affirming it was all they had; but the extortioner refused to be satisfied, and proceeded to demolish the city.

From Beroea he advanced to Antioch, and demanded the 1000 lbs. with which Megas had undertaken to redeem it; and it is said that he would have been contented to receive a smaller sum. Germanus and the Patriarch had already departed to Cilicia, and the Antiochenes would probably have paid the money had not the arrival of six thousand soldiers from Phoenicia Libanensis, led by Theoctistus and Molatzes, infused into their

p97
hearts a rash and unfortunate confidence. Julian, an Imperial secretary, who had arrived at Antioch as an ambassador, bade the inhabitants resist the extortion; and Paul, the interpreter of Chosroes, who approached the walls and counselled them to pay the money, was almost slain. Not content with defying the enemy by a refusal, the men of Antioch stood on their walls and loaded Chosroes with torrents of scurrilous abuse, which might have inflamed a milder monarch.

The siege which ensued was short. It seems not to have occurred to the besieged that they should themselves occupy the dangerous rock outside the citadel, and it was seized by the enemy. The defence at first was brave. Between the towers, which crowned the wall at intervals, platforms of wooden beams were suspended by ropes attached to the towers, that a greater number of defenders might man the walls at once. But during the fighting the ropes gave way and the suspended soldiers were precipitated, some without, some within the walls; the men in the towers were seized with panic and left their posts. The confusion was increased by a report that Buzes was coming to the rescue; and a multitude of women and children were crushed or trampled to death. But the gate leading to the remote suburb of Daphne was purposely left unblocked by the Persians; Chosroes seems to have desired that the Roman soldiers and their officers should be allowed to leave the city unmolested; and some of the inhabitants escaped with the departing army. But the young men of the Hippodrome factions made a valiant and hopeless stand against superior numbers; and the city was not entered without a considerable loss of life, which Chosroes pretended to deplore. It is said that two illustrious ladies cast themselves into the Orontes, to escape the cruelties of oriental licentiousness.

It was nearly three hundred years since Antioch had experienced the presence of a human foe, though it suffered frequently and grievously from the malignity of nature. The Sassanid Sapor had taken the city in the ill-starred reign of Valerian, but it was kindly dealt with then in comparison with its treatment by Chosroes. The cathedral was stripped of its wealth in gold and silver and its splendid marbles. Orders were given that the whole town should be burnt, except the

p98
cathedral, and the sentence of the relentless conqueror was executed as far as was practicable.

While the work of demolition was being carried out, Chosroes was treating with the ambassadors\texthebrew{⁠}\footnote{Julian mentioned above, and John, son of Rufinus (doubtless the same Rufinus who had been employed by Anastasius as ambassador to Kavad).

} of Justinian, and expressed himself ready to make peace, on condition that he received 5000 lbs. of gold, paid immediately, and an annual sum of 500 lb. nominally for the defence of the Caspian Gates. While the ambassadors returned with this answer to Byzantium, Chosroes advanced to Seleucia, the port of Antioch, and looked upon the waters of the Mediterranean; it is related that he took a solitary bath in the sea and sacrificed to the sun. In returning he visited Daphne, which was not included in the fate of Antioch, and thence proceeded to Apamea, whose gates he was invited to enter with a guard of 200 soldiers. All the gold and silver in the town was collected to satisfy his greed, even to the jewelled case in which a piece of the true cross was reverently preserved. He spared the precious relic itself, which for him was devoid of value. The city of Chalcis purchased its safety by a sum of 200 lbs. of gold; and having exhausted the provinces to the west of the Euphrates, Chosroes decided to continue his campaign of extortion in Mesopotamia, and cross the river at Obbane, near Barbalissus, by a bridge of boats. Edessa, the great stronghold of western Mesopotamia, was too strong itself to fear a siege, but paid 200 lbs. of gold for the immunity of the surrounding territory from devastation.\texthebrew{⁠}\footnote{The people of Edessa were generous enough to subscribe to ransom the Antiochene captives; farmers who had no money gave a sheep or an ass, prostitutes stripped off their ornaments. But, according to Procopius (\emph{B. P.} II.13), Buzes, who happened to be there, seized the money that was collected and allowed the captives to be carried off to Persia.

} At Edessa, ambassadors arrived from Justinian, bearing his consent to the terms proposed by Chosroes; but in spite of this the Persian did not shrink from making an attempt to take Daras on his homeward march.

The fortress of Daras, which Anastasius had erected to replace the long-lost Nisibis as an outpost in eastern Mesopotamia, was built on three hills, on the highest of which stood the citadel. One of the other heights projected from higher hills behind and could not be surrounded by the walls, which

p99
were built across it. There were two walls between which stretched a space of fifty feet, used by the inhabitants for the pasture of domestic animals. The climate of Mesopotamia, the severe snows of winter followed by the burning heats of summer, tried the strength of masonry, and Justinian found it necessary to repair the fortress. He did far more than repair it. He raised the inner wall by a new story, so that it reached the unusual height of sixty feet, and he secured the supply of water by diverting the river, which flowed outside the walls, into the town by means of a channel worked between the rocks. He also built barracks for the soldiers, so that the inhabitants were spared the burden of quartering them.\footnote{See Procopius,

\emph{De aed.} II.1 ; \emph{B. P.} II.13. The towers were 100 feet high. The details of the description of Procopius have been verified by the discoveries of Sachan on the site (\emph{Reise in Syrien und Mesopotamien}, 395 \emph{sqq.}). Cp. Chapot, \emph{op. cit.} 313 \emph{sqq.}

}

Chosroes attacked the city on the western side, and burned the western gates of the outer wall, but no Persian was bold enough to enter the interspace. He then began operations on the eastern, the only side of the rock-bound city where digging was possible, and ran a mine under the outer wall. The vigilance of the besieged was baffled until the subterranean passage had reached the foundations of the outer wall; but then, according to the story, a human or superhuman form in the guise of a Persian soldier advanced near the wall under the pretext of collecting discharged missiles, and while to the besiegers he seemed to be mocking the men on the battlements, he was really informing the besieged of the danger that was creeping upon them unawares. The Romans then, by the counsel of Theodore, a clever engineer, dug a deep transverse trench between the two walls so as to intersect the line of the enemy's excavation; the Persian burrowers suddenly ran or fell into the Roman pit; those in front were slain, and the rest fled back unpursued through the dark passage. Disgusted at this failure, Chosroes raised the siege on receiving from the men of Daras 1000 lbs. of silver. Justinian, indignant at his enemy's breach of faith, broke off the negotiations for peace.

When he returned to Ctesiphon the victorious monarch built a new city near his capital, on the model of Antioch, with whose spoils it was beautified, and settled therein the captive inhabitants of the original city, the remainder of whose days

p100
was perhaps more happily spent than if the generosity of the Edessenes had achieved its intention. The name of the new town, according to Persian writers,\texthebrew{⁠}\footnote{See Tabari, pp341\texthebrew{‑}342; Rawlinson, \emph{Seventh Oriental Monarchy}, p395. The new Antioch had one remarkable privilege; slaves who fled thither, if acknowledged by its citizens as kinsmen, were exempted from the pursuit of their Persian masters.

} was Rumia (Rome); according to Procopius it was called by the joint names of Chosroes and Antioch ( Chosro-Antiocheia ).\footnote{Procopius, \emph{B. P.} II.15\texthebrew{‑}19. Antioch itself was rebuilt by Justinian. The circuit of the wall was contracted, and the hill cliffs of Orocasias were not included within the line. The course of the Orontes was diverted so that it should flow by the new walls.

Procop. \emph{De aed.} II.10 .

}

\xsubsection{II. The Persian Invasion of Colchis, and the campaign of Belisarius in Mesopotamia ( A.D. 541)}

From this time forth the kingdom of Lazica or Colchis began to play a more important part in the wars between the Romans and Persians. This country seems to have been then far poorer than it is to\texthebrew{‑}day; the Lazi depended for corn, º salt, and other necessary articles of consumption on Roman merchants, and gave in exchange skins and slaves; while ``at present Mingrelia, though wretchedly cultivated, produces maize, millet, and barley in abundance; the trees are everywhere festooned with vines, which grow naturally, and yield a very tolerable wine; while salt is one of the main products of the neighbouring Georgia.''\texthebrew{⁠}\footnote{Rawlinson, \emph{op. cit.} p406, where the facts are quoted from Haxthausen's \emph{Transcaucasia.} Procopius himself mentions (\emph{B. G.} IV.14) that the district of Muchiresis in Colchis was very fertile, produ­cing wine and various kinds of corn.

} The Lazi were dependent on the Roman Empire, but the dependence consisted not in paying tribute but in committing the choice of their kings to the wisdom of the Roman Emperor. The nobles were in the habit of choosing wives among the Romans; Gubazes, the king who invited Chosroes to enter his country, was the son of a Roman lady, and had served as a silentiary in the Byzantine palace.\texthebrew{⁠}\footnote{See Procop.

\emph{B. P.} II.29 ;

\emph{B. G.} IV.9 . Previous Lazic kings had married Roman ladies of senatorial family.

} The Lazic kingdom was a useful barrier against the trans-Caucasian Scythian races, and the inhabitants defended the mountain passes without causing any outlay of men or money to the Empire.

But when the Persians seized Iberia it was considered necessary to secure the country which barred them from the sea by the

p101
protection of Roman soldiers, and the unpopular general Peter, originally a Persian captive, was not one to make the natives rejoice at the presence of their defenders. Peter's successor was John Tzibus, a man of obscure station, whose unscrupulous skill in raising money made him a useful tool to the Emperor. He was an able man, for it was by his advice that Justinian built the town of Petra, to the south of the Phasis.\texthebrew{⁠}\footnote{The site of Petra is uncertain. It has been identified with Ujenar (by Dubois de Montpéreux, \emph{Voyage autour du Caucase}, III.86), 15 miles SE of the mouth of the Phasis and 12 miles from the coast. But the description of Procopius, \emph{B. G.} II.17 , suggests that it was quite close to the sea.

} Here he established a monopoly and oppressed the natives. It was no longer possible for the Lazi to deal directly with the traders and buy their corn and salt at a reasonable price; John Tzibus, perched in the fortress of Petra, acted as a middleman, to whom both buyers and sellers were obliged to resort, and pay the highest or receive the lowest prices. In justification of this monopoly it may be remarked that it was the only practicable way of imposing a tax on the Lazi; and the imposition of a tax might have been deemed a necessary and just compensation for the defence of the country, notwithstanding the facts that it was garrisoned solely in Roman interests, and that the garrison itself was unwelcome to the natives.

Exasperated by these grievances, Gubazes, the king of Lazica, sent an embassy to Chosroes, inviting him to recover a venerable kingdom, and pointing out that if he expelled the Romans from Lazica he would have access to the Euxine, whose waters could convey his forces against Byzantium, while he would have an opportunity of establishing a connexion with those other enemies of Rome, the barbarians north of the Caucasus.\texthebrew{⁠}\footnote{Another element in the Colchian policy of Chosroes was the circumstance that if Lazica were Persian, the Iberians would have no power in the rear to support them if they revolted. Compare Procopius, \emph{B. P.} II.28.

} Chosroes consented to the proposals of the ambassadors; and keeping his real intention secret, pretended that pressing affairs required his presence in Iberia.

Under the guidance of the envoys, Chosroes and his army passed into the thick woods and difficult hill-passes of Colchis, cutting down as they went lofty and leafy trees, which hung in dense array on the steep acclivities, and using the trunks to smooth or render passable rugged or dangerous places. When

p102
they had penetrated to the middle of the country, they were met by Gubazes, who paid oriental homage to the great king. The chief object was to capture Petra, the stronghold of Roman power, and dislodge the tradesman, as Chosroes contemptuously termed the monopolist, John Tzibus. A detachment of the army under Aniabedes was sent on in advance to attack the fortress; and when this officer arrived before the walls he found the gates shut, yet the place seemed totally deserted, and not a trace of an inhabitant was visible. A messenger was sent to inform Chosroes of this surprise; the rest of the army hastened to the spot; a battering-ram was applied to the gate, while the monarch watched the proceedings from the top of an adjacent hill. Suddenly the gate flew open, and a multitude of Roman soldiers rushing forth overwhelmed those Persians who were applying the engine, and, having killed many others who were drawn up hard by, speedily retreated and closed the gate. The unfortunate Aniabedes (according to others, the officer who was charged with the operation of the battering-ram ) was impaled for the crime of being vanquished by a huckster.

A regular siege now began. It was inevitable that Petra should be captured, says our historian Procopius, in the vein of Herodotus,\texthebrew{⁠}\footnote{\textgreek{Κα\texthebrew{ὶ} φ\texthebrew{ὰ}ρ \texthebrew{ἔ}δει Πέτραν Χόσρωι \texthebrew{ἀ}λ\texthebrew{ῶ}ναι.}} and therefore John, the governor, was slain by an accidental missile, and the garrison, deprived of their commander, became careless and lax. On one side Petra was protected by the sea, landwards inaccessible cliffs defied the skill or bravery of an assailant, save only where one narrow entrance divided the line of steep cliffs and admitted of access from the plain. This gap between the rocks was filled by a long wall, the ends of which were commanded by towers constructed in an unusual manner, for instead of being hollow all the way up, they were made of solid stone to a considerable height, so that they could not be shaken by the most power ­ful engine. But oriental inventiveness undermined these wonders of solidity. A mine was bored under the base of one of the towers, the lower stones were removed and replaced by wood, the demolishing force of fire loosened the upper layers of stones, and the tower fell. This success was decisive, as the besieged recognised; they readily capitulated, and the victors did not lay hands on any property in the fortress save the possessions of the defunct governor.

p103
Having placed a Persian garrison in Petra, Chosroes remained no longer in Lazica, for the news had reached him that Belisarius was about to invade Assyria, and he hurried back to defend his dominions.

Belisarius, accompanied by all the Goths whom he had led in triumph from Italy, except the Gothic king himself, had proceeded in the spring to take command of the eastern army in Mesopotamia.\texthebrew{⁠}\footnote{The Italian generals accompanied Belisarius. One of them, Valerian, succeeded Martin as general in Armenia; Martin had been transferred to Mesopotamia.

} Having found out by spies that no invasion was meditated by Chosroes, whose presence was demanded in Iberia \texthebrew{—} the design on Lazica was kept effectually concealed \texthebrew{—} the Roman general determined to lead the whole army, along with the auxiliary Saracens of Harith into Persian territory. It is remarkable that in this campaign although Belisarius was chief in command he never seems to have ventured or cared to execute his strategic plans without consulting the advice of the other officers. It is difficult to say whether this was due to distrust of his own judgment and the reflexion that many of the subordinate generals had more recent experience of Persian warfare than himself,\texthebrew{⁠}\footnote{This is dwelt on in one of the speeches which Procopius places in the mouth of Belisarius (\emph{B. P.} II.16).

} or to a fear that some of the leaders in an army composed of soldiers of many races might prove refractory and impatient of too peremptory orders. At Daras a council of war decided on an immediate advance.

The army marched towards Nisibis, which was too strong to be attacked, and moved forward to the fortress of Sisaurana , where an assault was at first repulsed with loss.\texthebrew{⁠}\footnote{Between Nisibis and the Tigris (the same as Sisar in

Amm. Marc. XVIII.6.9 ). º

} Belisarius decided to invest the place, but as the Saracens were useless for siege warfare, he sent Harith and his troops, accompanied by 1200 of his own retainers, to invade and harry Assyria, intending to cross the Tigris himself when he had taken the fort. The garrison was not supplied with provisions, and soon consented to surrender; all the Christians were dismissed free, the fire-worshippers were sent to Byzantium\texthebrew{⁠}\footnote{These Persians, with their leader Bleschanes, were afterwards sent to Italy against the Goths. It was Roman policy to employ Persian captives against the Goths, Gothic captives against the Persians.

} to await the Emperor's pleasure, and the fort was levelled to the ground.

Meanwhile the plundering expedition of Harith was success ­ful,

p104
but he played his allies false. Desiring to retain all the spoils for himself, he invented a story to rid himself of the Romans who accompanied him,\texthebrew{⁠}\footnote{Trajan and John the Glutton were in command of these 1200 \textgreek{\texthebrew{ὑ}πασπισταί}. When they separated from Harith they proceeded to Theodosiopolis, in order to avoid a hostile army which did not exist.

} and he sent no information to Belisarius. This was not the only cause of anxiety that vexed the general's mind. The Roman, especially the Thracian, soldiers were not inured to the intense heat of the dry Mesopotamian climate in midsummer, and disease broke out in the army, demoralised by physical exhaustion. All the soldiers were anxious to return to more clement districts. There was nothing to be done but yield to the prevailing wish, which was shared by all the generals. It cannot be claimed that the campaign of Belisarius accomplished much to set off against the acquisition of Petra by the Persians.

It was indeed whispered by the general's enemies that he had culpably missed a great opportunity. They insinuated that if, after the capture of Sisaurana, he had advanced beyond the Tigris he might have carried the war up to the walls of Ctesiphon. But he sacrificed the interests of the Empire to private motives, and retreated in order to meet his wife who had just arrived in the East and punish her for her infidelity.\texthebrew{⁠}\footnote{Procopius, \emph{H. A.} 17.25 . For the story see above,

p60 .

} The scandals may be true, but it is impossible to say how far they affected the military conduct of Belisarius.

\xsubsection{III. The Persian Invasion of Commagene ( A.D. 542)\texthebrew{⁠}\footnote{Procopius, \emph{B. P.} II.20, 21.

}}

The first act of Chosroes when he crossed the Euphrates in spring was to send 6000 soldiers to besiege the town of Sergiopolis because the bishop Candidus, who had undertaken to pay the ransom of the Surene captives two years before, was unable to collect the amount, and found Justinian deaf to his appeals for aid. But the town lay in a desert, and the besiegers were soon obliged to abandon their design in consequence of the drought. It was not the Persian's intention to waste his time in despoiling the province of Euphratensis; he purposed to invade Palestine and plunder the treasures of Jerusalem. But this exploit was reserved for his grandson of the same name, and the invader returned to his kingdom having accomplished

p105
almost nothing. This speedy retreat was probably due to the outbreak of the Plague in Persia, though the Roman historian attributes it to the address of Belisarius.

Belisarius travelled by post-horses ( veredi ) from Constantinople to the Euphratensian province, and taking up his quarters at Europus\texthebrew{⁠}\footnote{Yerabus. Cp. Chapot, \emph{op. cit.} p280.

} on the Euphrates, he collected there the bulk of the troops who were dispersed throughout the province in its various cities. Chosroes was curious about the personality of Belisarius, of whom he had heard so much, \texthebrew{—} the conqueror of the Vandals, the conqueror of the Goths, who had led two fallen monarchs in triumph to the feet of Justinian. Accordingly he sent Abandanes\texthebrew{⁠}\footnote{One of the \textgreek{βασιλικο\texthebrew{ὶ} γραμματε\texthebrew{ῖ}ς} (notarii).

} as an envoy to the Roman general on the pretext of learning why Justinian had not sent ambassadors to negotiate a peace.

Belisarius did not mistake the true nature of this mission, and determined to make an impression. Having sent a body of one thousand cavalry to the left bank of the river, to harass the enemy if they attempted to cross, he selected six thousand tall and comely men from his army and proceeded with them to a place at some distance from his camp, as if on a hunting expedition. He had constructed for himself a pavilion\texthebrew{⁠}\footnote{\textgreek{Παπυλεών}, which Procopius introduces with one of his usual apologetic formulae for words that are not Greek.

} of thick canvas, which he set up, as in a desert spot, and when he knew that the ambassador was approaching, he arranged his soldiers with careful negligence. On either side of him stood Thracians and Illyrians, a little farther off the Goths, then Heruls, Vandals, and Moors; all were arrayed in close-fitting linen tunics and drawers, without a cloak or epomis to disguise the symmetry of their forms, and, like hunters, each carried a whip as well as some weapon, a sword, an axe, or a bow. They did not stand still, as men on duty, but moved carelessly about, glancing idly and indifferently at the Persian envoy, who soon arrived and marvelled.

To the envoy's complaint that the Emperor had not sent an embassy to his master, Belisarius answered, with an air of amusement, ``It is not the habit of men to transact their affairs as Chosroes has transacted his. Others, when aggrieved, send an embassy first, and if they fail in obtaining satisfaction, resort

p106
to war; but he attacks and then talks of peace.'' The presence and bearing of the Roman general, and the appearance of his followers, hunting indifferently at a short distance from the Persian camp without any precautions, made a profound impression on Abandanes, and he persuaded his master to abandon the proposed expedition. Chosroes may have reflected that the triumph of a king over a general would be no humiliation for the general, while the triumph of a mere general over a king would be very humiliating for the king; such at least is the colouring that general's historian puts on the king's retreat. According to the same authority, Chosroes hesitated to risk the passage of the Euphrates while the enemy was so near, but Belisarius, with his smaller numbers, did not attempt to oppose him.\texthebrew{⁠}\footnote{A Persian army always carried with it materials for constructing pontoons (Proc. \emph{B. P.} II.21), and they crossed by such a bridge on this occasion.

} A truce was made, and a rich citizen of Edessa was delivered, an unwilling hostage, to Chosroes. In their retreat, the Persians turned aside to take and demolish Callinicum, the Coblenz of the Euphrates, which fell an easy prey to their assault, as the walls were in process of renovation at the time. This retirement of Chosroes, according to Procopius, procured for Belisarius greater glory than he had won by his victories in the West. But Belisarius was now recalled to conduct the war in Italy.

The account of Procopius, which coming from a less able historian would be rejected on account of internal improbability, cannot be accepted with confidence. It displays such a marked tendency to glorify Belisarius, that it can hardly be received as a candid story of the actual transactions. Besides, there is a certain inconsistency. If Chosroes retired \emph{for fear of} Belisarius, as Procopius would have us believe, why was it he who received the hostage, and how did he venture to take Callinicum? As there actually existed a sufficient cause, unconnected with the Romans, to induce his return to Persia, namely the outbreak of the Plague, we may suspect that this was its true motive.\footnote{So Rawlinson (\emph{op. cit.} p461), who perhaps is more generous to Procopius than he deserves. The Plague broke out in Persia in the summer of 542.

}

p107

\xsubsection{IV. The Roman Invasion of Persarmenia ( A.D. 543)\texthebrew{⁠}\footnote{Proc. \emph{ib.} 24, 25.

}}

In spite of the Plague Chosroes set forth in the following spring to invade Roman Armenia. He advanced into the district of Azerbiyan (Atropatene), and halted at the great shrine of Persian fire-worship , where the Magi kept alive an eternal flame, which Procopius wished to identify with the fire of Roman Vesta. Here the Persian monarch waited for some time, having received a message that two Imperial ambassadors\texthebrew{⁠}\footnote{Constantianus, an Illyrian, and Sergius of Edessa, both rhetors and men of parts.

} were on their way to him. But the ambassadors did not arrive, because one of them fell ill by the road; and Chosroes did not pursue his northward journey, because the Plague broke out in his army. His general Nabedes sent the bishop of Dubios to Valerian, the general in Armenia, with complaints that the expected embassy had not appeared. The bishop was accompanied by his brother, who secretly communicated to Valerian the valuable information that Chosroes was just then encompassed by perplexities, the spread of the Plague, and the revolt of one of his sons. It was a favourable opportunity for the Romans, and Justinian directed all the generals stationed in the East to join forces to invade Persarmenia.

Martin was now Master of Soldiers in the East. He does not appear to have possessed much actual authority over the other commanders. They at first encamped in the same district, but did not unite their forces, which in all amounted to about thirty thousand men. Martin himself, with Ildiger and Theoctistus, encamped at Kitharizon, a fort about four days' march from Theodosiopolis; the troops of Peter and Adolius took up their quarters in the vicinity; while Valerian stationed himself close to Theodosiopolis and was joined there by Narses with a body of Heruls and Armenians. The Emperor's cousin Justus and some other commanders remained during the campaign far to the south in the neighbourhood of Martyropolis, where they made incursions of no great importance.

At first the various generals made separate inroads, but they ultimately united their regiments in the spacious plain of Dubios, eight days from Theodosiopolis. This plain, well suited for equestrian exercise, and richly populated, was a famous rendez-vous

p108
for traders of all nations, Indian, Iberian, Persian, and Roman.\texthebrew{⁠}\footnote{Dubios corresponds to Duin.

} \texthebrew{•} About thirteen miles from Dubios there was a steep mountain, on the side of which was perched a village called Anglôn, protected by a strong fortress. Here the Persian general Nabedes, with four thousand soldiers, had taken up an almost impregnable position, blocking the precipitous streets of the village with stones and wagons. The ranks of the Roman army, as it marched to Anglôn, fell into disorder; the want of union among the generals, who acknowledged no supreme leader, led to confusion in the line of march; mixed bodies of soldiers and sutlers turned aside to plunder; and the security which they displayed might have warranted a spectator in prophesying a speedy reverse. As they drew near to the fortress, an attempt was made to marshal the somewhat demoralised troops in the form of two wings and a centre. The centre was commanded by Martin, the right wing by Peter, the left by Valerian; and all advanced in irregular and wavering line, on account of the roughness of the ground.\texthebrew{⁠}\footnote{Procopius assigns as an additional cause the want of discipline or previous marshalling of the troops; but I feel some suspicions of the whole account of this campaign.

} The best course for the Persians was obviously to act on the defensive. Narses and his Heruls, who were probably on the left wing with Valerian, were the first to attack the foes and to press them back into the fort. Drawn on by the retreating enemy through the narrow village streets, they were suddenly taken in the flank and in the rear by an ambush of Persians who had concealed themselves in the houses. The valiant Narses was wounded in the temple; his brother succeeded in carrying him from the fray, but the wound proved mortal. This repulse of the foremost spread the alarm to the regiments that were coming up behind; Nabedes comprehended that the moment had arrived to take the offensive and let loose his soldiers on the panic-stricken ranks of the assailants; and all the Heruls, who fought according to their wont without helmets or breastplates,\texthebrew{⁠}\footnote{The Herul's only armour was a shield and a cloak of thick stuff.

} fell before the charge of the Persians. The Romans did not tarry; they cast their arms away and fled in wild confusion, and the mounted soldiers galloped so fast that few horses survived the flight; but the Persians, apprehensive of an ambush, did not pursue.

Never, says Procopius, did the Romans experience such a

p109
great disaster. This exaggeration inclines us to be sceptical. We can hardly avoid detecting in his narrative a desire to place the generals in as bad a light as possible, just as in his description of the hostilities of the preceding year we saw reason to suspect him unduly magnifying the behaviour of his hero Belisarius. In fact his aim seems to be to draw a strong and striking contrast between a brilliant campaign and a miserable failure. We have seen reason to doubt the exceptional brilliancy of the achievement of Belisarius; and we may wonder whether the defeat at Anglôn was really overwhelming.

\xsubsection{V. The Persian Invasion of Mesopotamia; Siege of Edessa ( A.D. 544)\texthebrew{⁠}\footnote{Procopius, \emph{B. P.} II.26\texthebrew{‑}28.

}}

His failure at Edessa in the first year of the war had rankled in the mind of the Sassanid monarch. The confidence of the inhabitants that they enjoyed a special divine protection in virtue of the letter of Jesus to Abgar was a challenge to the superstition of the Fire-worshippers , and the Magi could not bear the thought that they had been defeated by the God of the Christians. Chosroes comforted himself by threatening to enslave the Edessenes, and make the site of their city a pasture for sheep. But the place was strong. Its walls had been ruined again and again by earthquakes, against which the divine promise did not secure it, and again and again rebuilt. It had suffered this calamity recently ( A.D. 525) and had been restored by Justin, who honoured it by his own name. But Justinopolis had as little power over the tongues of men as Anastasiopolis or Theupolis. Edessa, the city of Abgar, remained Edessa, as Daras remained Daras and Antioch Antioch. Justinian had reconstructed the fortifications and made it stronger than ever, and installed hydraulic arrangements to prevent the inundations of the river Scyrtus which flowed through the town.\footnote{Evagrius, \emph{H. E.} IV.8 ; Procopius,

\emph{De aed.} II.7 ;

\emph{H. A.} 18 . The chief feature of the fortifications of Justinian were the new walls which he built to the crest of a hill overtopping the citadel. For the plan of the castle see Texier and Pullan, \emph{Byzantine Architecture}, p183.

}

Realising the strength of the place, Chosroes would have been glad to avoid the risk of a second failure, and he proposed to

p110
the inhabitants that they should pay him an immense sum or allow him to take all the riches in the city. His proposal was refused, though if he had made a reasonable demand it would have been agreed to; and the Persian army encamped at somewhat less than a mile from the walls. Three experienced generals, Martin, Peter, and Peranius, were stationed in Edessa at this time.

On the eighth day from the beginning of the siege, Chosroes caused a large number of hewn trees to be strewn on the ground in the shape of an immense square, at about a stone's throw from the city; earth was heaped over the trees, so as to form a flat mound, and stones, not cut smooth and regular as for building, but rough hewn, were piled on the top, additional strength being secured by a layer of wooden beams placed between the stones and the earth. It required many days to raise this mound to a height sufficient to overtop the walls. At first the workmen were harassed by a sally of Huns, one of whom, named Argek, slew twenty-seven with his own hand. This could not be repeated, as henceforward a guard of Persians stood by to protect the builders. As the work went on, the mound seems to have been extended in breadth as well as in height, and to have approached closer to the walls, so that the workmen came within range of the archers who manned the battlements, but they protected themselves by thick and long strips of canvas, woven of goat hair, which were hung on poles, and proved an adequate shield. Foiled in their attempts to obstruct the progress of the threatening pile, which they saw rising daily higher and higher, the besieged sent an embassy to Chosroes. The spokesman of the ambassadors was the physician Stephen, a native of Edessa, who had enjoyed the friendship and favour of Kavad, whom he had healed of a disease, and had superintended the education of Chosroes himself. But even he, influential though he was, could not obtain more than the choice of three alternatives \texthebrew{—} the surrender of Peter and Peranius, who, originally Persian subjects, had presumed to make war against their master's son; the payment of 50,000 lbs. of gold (two million and a quarter pounds sterling); or the reception of Persian deputies, who should ransack the city for treasures and bring all to the Persian camp. All these proposals were too extravagant to be entertained for an instant;

p111
the ambassadors returned in dejection, and the erection of the mound advanced. A new embassy was sent, but was not even admitted to an audience; and when the plan of raising the city wall was tried, the besiegers found no difficulty in elevating their structure also.

At length the Romans resorted to the plan of undermining the mound, but when their excavation had reached the middle of the pile the noise of the subterranean digging was heard by the Persian builders, who immediately dug or hewed a hole in their own structure in order to discover the miners. These, knowing that they were detected, filled up the remotest part of the excavated passage and adopted a new device. Beneath the end of the mound nearest to the city they formed a small subterranean chamber with stones, boards, and earth. Into this room they threw piles of wood of the most inflammable kind, which had been smeared over with sulphur, bitumen, and oil of cedar. As soon as the mound was completed,\texthebrew{⁠}\footnote{Just before its completion, Martin made proposals for peace, but the Persians were unwilling to treat.

} they kindled the logs, and kept the fire replenished with fresh fuel. A considerable time was required for the fire to penetrate the entire extent of the mound, and smoke began to issue prematurely from that part where the foundations were first inflamed. The besieged adopted an obvious device to mislead the besiegers. They cast burning arrows and hurled vessels filled with burning embers on various parts of the mound; the Persian soldiers ran to and fro to extinguish them, believing that the smoke, which really came from beneath, was caused by the flaming missiles; and some thus employed were pierced by arrows from the walls. Next morning Chosroes himself visited the mound and was the first to discover the true cause of the smoke, which now issued in denser volume. The whole army was summoned to the scene amid jeers of the Romans, who surveyed from the walls the consternation of their foes. The torrents of water with which the stones were flooded increased the vapour instead of quenching it and caused the sulphurous flames to operate more violently. In the evening the volume of smoke was so great that it could be seen as far away to the south as at the city of Carrhae;\texthebrew{⁠}\footnote{The distance of Carrhae from Edessa was \texthebrew{•} about thirty miles.

} and the fire,

p112
which had been gradually working upwards as well as spreading beneath, at length gained the air and overtopped the surface. Then the Persians desisted from their futile endeavours.

Six days later an attack was made on the walls at early dawn, and but for a farmer who chanced to be awake and gave the alarm, the garrison might have been surprised. The assailants were repulsed; and another assault on the great gate at mid-day likewise failed.\texthebrew{⁠}\footnote{At this juncture the Persians desired to treat, and informed the garrison that a Roman ambassador from Constantinople had arrived in their camp. They allowed the ambassador to enter Edessa, but Martin was suspicious of their intentions, and feigning to be ill said that he would send envoys in three days.

} One final effort was made by the baffled enemy. The ruins of the half-demolished mound were covered with a floor of bricks, and from this elevation a grand attack was made. At first the Persians seemed to be superior, but the enthusiasm which prevailed in the city was ultimately crowned with victory. The peasants, even the women and the children, ascended the walls and took a part in the combat; cauldrons of oil were kept continually boiling, that the burning liquid might be poured on the heads of the assailants; and the Persians, unable to endure the fury of their enemies, fell back and confessed to Chosroes that they were vanquished. The enraged despot drove them back to the encounter; they made yet one supreme effort, and were yet once more discomfited. Edessa was saved, and the siege unwillingly abandoned by the disappointed king, who, however, had the satisfaction of receiving 500 lbs. of gold from the weary though victorious Edessenes.

In the following year, A.D. 545, a truce\texthebrew{⁠}\footnote{Procopius, \emph{B. P.} II.28. The 2000 lbs. were calculated at the rate of 400 a year.

} was concluded for five years, Justinian consenting to pay 2000 lbs. of gold. But Chosroes refused to assent to the Emperor's demand that this truce should apply to operations in Lazica, where he believed that he held a strong position. Hence during the duration of the truce, there was an ``imperfect'' war between the two powers in Colchis. Justinian readily acceded to a request of the king to permit a certain Greek physician, named Tribunus,\texthebrew{⁠}\footnote{Cp. Zacharias Myt. XII.7, where he is called Tribonian.

} to remain at the Persian court for a year. Tribunus of Palestine, the best medical authority of the age, was, we are told, a man of distinguished virtue and piety, and highly valued by Chosroes,

p113
whose constitution was delicate and constantly required the services of a physician. At the end of the year the king permitted him to ask a boon, and instead of proposing remuneration for himself he begged for the freedom of some Roman prisoners. Chosroes not only liberated those whom he named, but others also to the number of three thousand.

The Lazi soon found that the despotism of the Persian fire-worshipper was less tolerable than the oppression of the Christian monopolists, and repented that they had taught the armies of the great king to penetrate the defiles of Colchis. It was not long before the Magi attempted to convert the new province to a faith which was odious to the christianised natives, and it became known that Chosroes entertained the intention of removing the inhabitants and colonising the land with Persians. Gubazes, who learned that Chosroes was plotting against his life, hastened to seek the pardon and the protection of Justinian. In A.D. 549, 7000 Romans were sent to Lazica, under the command of Dagisthaeus, to recover the fortress of Petra. Their forces were strengthened by the addition of a thousand Tzanic auxiliaries.

The acquisition of Colchis pleased Chosroes so highly, and the province appeared to him of such eminent importance, that he took every precaution to secure it.\texthebrew{⁠}\footnote{He tried to build a fleet in the Euxine, but the material was destroyed by lightning.

} A highway was constructed from the Iberian confines through the country's hilly and woody passes, so that not only cavalry but elephants could traverse it. The fortress of Petra was supplied with sufficient stores of provisions, consisting of salted meat and corn, to last for five years; no wine was provided, but vinegar and a sort of grain from which a spirituous liquor could be distilled. The armour and weapons which were stored in the magazines would, as was afterwards found, have accoutred five times the number of the besiegers; and a cunning device was adopted to supply

p114
the city with water, while the enemy should delude themselves with the idea that they had cut off the supply.

When Dagisthaeus laid siege to the town the garrison consisted of 1500 Persians. He committed the mistake of not occupying the clisurae or passes from Iberia into Colchis, so as to prevent the arrival of Persian reinforcements. The siege was protracted for a long time, and the small garrison suffered heavy losses. At last Mermeroes, allowed to enter Colchis unopposed with large forces of cavalry and infantry, arrived at the pass which commands the plain of Petra. Here his progress was withstood by a hundred Romans, but after a long and bloody battle the weary guards gave way, and the Persians reached the summit. When Dagisthaeus learned this he raised the siege.

Mermeroes left 3000 men in Petra and provisioned it for a short time. Leaving 5000 men under Phabrigus in Colchis, and instructing them to keep Petra supplied with food, he withdrew to Persarmenia. Disaster soon befell these troops; they were surprised in their camp by Dagisthaeus and Gubazes in the early morning, and but few escaped. All the provisions brought from Iberia for the use of Petra were destroyed, and the eastern passes of Colchis were garrisoned.\footnote{At this point the two books of Procopius known as \emph{De bello Persico} come to an end, but the thread of the narrative is resumed in the \emph{De bello Gothico}, Bk. IV, which was written after the other books had been given to the world. Procopius apologises for the necessity which compels him to abandon his method of \emph{geographical} divisions. (\emph{B. G.} IV.1).

}

In the spring of A.D. 550 Chorianes entered Colchis with a Persian army, and encamped by the river Hippis, where a battle was fought in which Dagisthaeus was victorious, and Chorianes lost his life. Dagisthaeus, however, was accused of misconducting the siege of Petra, through disloyalty or culpable negligence. Justinian ordered his arrest, and appointed Bessas, who had recently returned from Italy, in his stead. Men wondered at this appointment, and thought that the Emperor was foolish to entrust the command to a general who was far advanced in years, and whose career in the West had been inglorious; but the choice, as we shall see, was justified by the result.

The first labour that devolved on Bessas was to suppress a revolt of the Abasgians. The territory of this nation extended along the lunated eastern coast of the Euxine, and was separated from Colchis by the country of the Apsilians, who inhabited

p115
the district between the western spurs of Caucasus and the sea. The Apsilians had long been Christians, and submitted to the lordship of their Lazic neighbours, who had at one time held sway over the Abasgians. Abasgia was governed by two princes, of whom one ruled in the west and the other in the east. These potentates increased their revenue by the sale of beautiful boys, whom they tore in early childhood from the arms of their reluctant parents and made eunuchs; for in the Roman Empire these comely and useful slaves were in constant demand, and secured a high price from the opulent nobles. It was the glory of Justinian to bring about the abolition of this unnatural practice; the people supported the remonstrances which the Emperor's envoy, himself an Abasgian eunuch, made to their kings; the royal tyranny was abolished, and a people which had worshipped trees embraced Christianity, to enjoy, as they thought, a long period of freedom under the protection of the Roman Augustus. But the mildest protectorate tends insensibly to become domination. Roman soldiers entered the country, and taxes were imposed on the new friends of the Emperor. The Abasgi preferred the despotism of men of their own blood to servitude to a foreign master, and they elected two new kings, Opsites in the east and Sceparnas in the west. But it would have been rash to brave the jealous anger of Justinian without the support of some stronger power, and when Nabedes, after the great defeat of the Persians on the Hippis, visited Lazica, he received sixty noble hostages from the Abasgians, who craved the protection of Chosroes. They had not taken warning from the repentance of the Lazi, that it was a hazardous measure to invoke the Persian. The king, Sceparnas, was soon afterwards summoned to the Sassanid court, and his colleague Opsites prepared to resist the Roman forces which Bessas despatched against him under the command of Wilgang (a Herul) and John the Armenian.

In the southern borders of Abasgia, close to the Apsilian frontier, an extreme mountain of the Caucasian chain descends in the form of a staircase to the waters of the Euxine. Here, on one of the lower spurs, the Abasgi had built a strong and roomy fastness in which they hoped to defy the pursuit of an invader. A rough and difficult glen separated it from the sea, while the ingress was so narrow that two persons could not enter

p116 abreast, and so low that it was necessary to crawl. The Romans, who had sailed from the Phasis, or perhaps from Trapezus, landed on the Apsilian borders, and proceeded by land to this glen, where they found the whole Abasgian nation arrayed to defend a pass which it would have been easy to hold against far larger numbers. Wilgang remained with half the army at the foot of the glen, while John and the other half embarked in boats which had accompanied the coast march of the soldiers. They landed at no great distance, and by a circuitous route were able to approach the unsuspecting foe in the rear. The Abasgians fled in consternation towards their fortress; fugitives and pursuers, mingled together, strove to penetrate the narrow aperture, and those inside could not prevent enemies from entering with friends. But the Romans when they were within the walls found a new labour awaiting them. The Abasgi fortified themselves in their houses, and vexed their adversaries by showering missiles from above. At length the Romans employed the aid of fire, and the dwellings were soon reduced to ashes. Some of the people were burnt, others, including the wives of the kings, were taken alive, while Opsites escaped to the neighbouring Sabirs.

The truce of five years had now elapsed (April, A.D. 550), and while new negotiations began between the courts of Constantinople and Ctesiphon, Bessas addressed himself to the enterprise in which Dagisthaeus had failed, the capture of Petra. The garrison was brave and resolute, and the siege was long. But the persistency of Bessas achieved success and the stronghold fell in the early spring of A.D. 551. The gallant soldier, John the Armenian, was slain in the final assault. When Mermeroes, who was approaching to relieve Petra, heard the news, he retraced his steps, in order to attack Archaeopolis and other fortresses on the right bank of the Phasis.\texthebrew{⁠}\footnote{At this time the total number of Roman soldiers in Lazica amounted to 12,000. Of these 3000 were stationed at Archaeopolis, the remaining 9000, with an auxiliary force of 800 Tzani, were entrenched in a camp near the mouth of the Phasis. A year later the forces amounted to 50,000 (Agathias, III.8). A comparison of these numbers with those of the expeditions to Africa and Italy (see the following chapters) shows the importance of the occupation of Lazica in the eyes of the Imperial government.

} His siege of Archaeopolis\footnote{Dubois de Montpéreux (\emph{op. cit.} III.51) finds Archaeopolis at Nakolevi, on the Chobos. Mermeroes made another attack on it in 552.

}

p117
was a failure. He suffered a considerable defeat and was forced to retire. He succeeded in taking some minor fortresses in the course of the following campaigns ( A.D. 552\texthebrew{‑}554).\texthebrew{⁠}\footnote{There has been some difficulty about the chronology of the last years of the Lazic war. The narrative of Procopius ends \emph{B. G.} IV.17. He marks the winter 551\texthebrew{‑}552 in c16, the spring of 552 in c17, and the failure of Mermeroes in that year. The story is continued by Agathias, II.18, who refers briefly to the futile attacks of Mermeroes on Archaeopolis, mentioned by Procopius, and then describes the continuation of hostilities, without mentioning that a winter had intervened. In II.22 he notices the death of Mermeroes, and II.27 places that event in the 28th year of Justinian and 25th of Chosroes (but the 24th of Chosroes corresponds to 28 Justin.) = A.D. 554\texthebrew{‑}555. this means that the events related in II.19\texthebrew{‑}22 occurred in 553\texthebrew{‑}554, and that the author has omitted to distinguish the years. After this point he invariably marks the years (III.15, spring 555; 28 and IV.12, winter 555\texthebrew{‑}556; IV.13, spring and summer 556; 15, winter 556\texthebrew{‑}557). This chronology (so Clinton, \emph{fr., sub annis}) is borne out by the notice of the earthquake in V.3, which is dated by John Mal. XVIII p488.

} His death, which occurred in the autumn of A.D. 554, was a serious loss to Chosroes, for, though old and lame, and unable even to ride, he was not only brave and experienced, but as unwearying and energetic as a youth. Nachoragan was sent to succeed him.

Although the operations of the Persians in these years had been attended with no conspicuous success, they had gained one considerable advantage without loss to themselves. The small inland district of Suania, in the hills to the north of Lazica, had hitherto been a dependency on that kingdom. Its princes were nominated by the Lazic kings. The Suanians now ( A.D. 552) repudiated this connexion and went over to the Persians, who sent troops to occupy the territory.\footnote{We learn this from the negotiations of A.D. 562; Menander, \emph{De leg. Rom.} \emph{fr.} 3, pp178, 186\texthebrew{‑}187. In the reign of Leo, some Suanian forts had been seized by the Persians, and the Suanians had sought help from the Emperor, \emph{c.} A.D. 468. Priscus, \emph{De leg. gent. fr.} 22.

}

In the meantime the question of the renewal of the five years' truce had been engaging the attention of the Roman and Persian courts, and the negotiations had continued for eighteen months. At length it was renewed ( A.D. 551, autumn) for another period of five years, the Romans agreeing to pay 2600 lbs. of gold,\texthebrew{⁠}\footnote{At 400 lbs. annually, the rate agreed on in 445. The extra 5600 were for the year and a half spent in negotiation. See Proc. \emph{B. G.} IV.15.

} and, as before, it was not to affect the hostilities in Colchis. A contemporary states that there was much popular indignation that Chosroes should have extorted from the Empire 4600 lbs. of gold in eleven and a half years, and the people of Constantinople murmured at the excessive consideration which the Emperor

p118
displayed towards the Persian ambassador Isdigunas\texthebrew{⁠}\footnote{Izedh-Gushnasp (\textgreek{\texthebrew{Ἰ}εσδεγουσνάφ} in Menander). The solemnities observed in the reception and treatment of this embassy were recorded in Peter the Patrician, and are preserved in Constantine Porph. \emph{Cer.} I.89 and 90. The ambassador is here called \textgreek{\texthebrew{Ἰ}έσδεκος} (p405). He returned to Persia in spring A.D. 552 and the treaty received the seal of Chosroes (Proc. \emph{B. G.} I.17).

} and his retinue, who were permitted to move about in the city, without a Roman escort, as if it belonged to them.

Meanwhile king Gubazes, who had been engaged in frequent quarrels with the Roman commanders, sent a complaint to Justinian accusing them of negligence in conducting the war. Bessas, Martin, and Rusticus were specially named. The Emperor deposed Bessas from his post, but assigned the chief command to Martin and did not recall Rusticus. This Rusticus was the Emperor's pursebearer who had been sent to bestow rewards on soldiers for special merit. He and Martin determined to remove Gubazes. To secure themselves from blame, they despatched John, brother of Rusticus, to Justinian with the false message that Gubazes was secretly favouring the Persians. Justinian was surprised, and determined to summon the king to Constantinople. ``What,'' asked John, ``is to be done, if he refuses?'' ``Compel him,'' said the Emperor; ``he is our subject.'' ``But if he resist?'' urged the conspirator. ``Then treat him as a \emph{tyrant.}'' ``And will he who should slay him have naught to fear?'' ``Naught, if he act disobediently and be slain as an enemy.'' Justinian signed a letter to this effect, and armed with it John returned to Colchis. The conspirators hastened to execute their treacherous design. Gubazes was invited to assist in an attack on the fortress of Onoguris, and with a few attendants he met the Roman army on the banks of the Chobus. An altercation arose between the king and Rusticus, and on the pretext that the gainsayer of a Roman general must necessarily be a friend of the enemy, John drew his dagger and plunged it in the royal breast. The wound was not mortal but it unhorsed the king, and when he attempted to rise from the ground, a blow from the squire of Rusticus killed him outright.\footnote{Agathias, III.2\texthebrew{‑}4. These events belong to the autumn and winter 554\texthebrew{‑}555.

}

The Lazi silently buried their king according to their customs, and turned away in mute reproach from their Roman protectors. They no longer took part in the military operations, but hid

p119
themselves away as men who had lost their hereditary glory. The other commanders, Buzes and Justin the son of Germanus, concealed the indignation which they felt, supposing that the outrage had the Emperor's authority. Some months later, when winter had begun, the Lazi met in secret council in some remote Caucasian ravine, and debated whether they should throw themselves on the protection of Chosroes. But their attachment to the Christian religion as well as their memory of Persian oppression forbade them to take this step, and they decided to appeal for justice and satisfaction to the Emperor, and at the same time to supplicate him to nominate Tzath, the younger brother of Gubazes, as their new king. Justinian promptly complied with both demands. Athanasius, a senator of high repute, was sent to investigate the circumstances of the assassination, and on his arrival he incarcerated Rusticus and John, pending a trial. In the spring ( A.D. 555) Tzath arrived in royal state, and when the Lazi beheld the Roman army saluting him as he rode in royal apparel, a tunic embroidered with gold reaching to his feet, a white mantle with a gold stripe, red shoes, a turban adorned with gold and gems, and a crown, they forgot their sorrow and escorted him in a gay and brilliant procession. It was not till the ensuing autumn\texthebrew{⁠}\footnote{Cp. Agathias, IV.12 \emph{ad init.}

} that the authors of the death of the late king were brought to justice, and the natives witnessed the solemn procedure of a Roman trial. Rusticus and John were executed. Martin's complicity was not so clear, and the Emperor, to whom his case was referred, deposed him from his command in favour of his own cousin Justin, the son of Germanus.\texthebrew{⁠}\footnote{Justin was created \textgreek{στρατηγ\texthebrew{ὸ}ς α\texthebrew{ὐ}τοκράτωρ}, Agathias, IV.21.

} Martin perhaps would not have been acquitted if he had not been popular with the army and a highly competent general.

Immediately after the assassination of Gubazes, the Romans who had assembled in full force before the fortress of Onoguris sustained a severe and inglorious defeat at the hands of 3000 Persians ( A.D. 554). In the following spring, Phasis (Poti), at the mouth of the like-named river, was attacked by Nachoragan, and an irregular battle before this town resulted in a victory for Martin which wiped out the disgrace of Onoguris.\texthebrew{⁠}\footnote{In consequence of this failure, Nachoragan was flayed alive by the order of Chosroes.

} In the

p120
same year, the Misimians, a people who lived to the north-east of the Apsilians and like these and the Suanians were dependent on Lazica, slew a Roman envoy who was travelling through their country and had treated them with insolence. Knowing that this outrage would be avenged they went over to Persia. This incident determined the nature of the unimportant operations of A.D. 556. A Persian army prevented the Romans from invading the land of the Misimians. But a punitive expedition was sent in the ensuing winter and was attended with an inhuman massacre of the Misimians, who finally yielded and were pardoned. This expedition was the last episode of the Lazic War.

The truce of five years expired in the autumn of A.D. 556. Both powers were weary of the war, and the course of the campaigns had not been encouraging to Chosroes. It is probable, too, that he was preparing for a final effort to destroy, in conjunction with the Turks, the kingdom of the Ephthalites. Early in the year he had sent his ambassador Isdigunas to Constantinople\texthebrew{⁠}\footnote{Malalas, XVIII p488, notes the presence of the Persian ambassador in May.

} to negotiate a renewal of the truce which would soon expire. It was intended that the arrangement should be a preliminary to a treaty of permanent peace, and this time it was not to be imperfect, it was to extend to Lazica as well as to Armenia and the East. The truce was concluded ( A.D. 557) on the terms of the \emph{status quo} in Lazica, each power retaining the forts which were in its possession; there was no limit of time and there were no money payments.\footnote{Agathias, IV.30.

}

The historical importance of the Lazic War lay in the fact that if the Romans had not succeeded in holding the country and thwarting the design of Chosroes, the great Asiatic power would have had access to the Euxine and the Empire would have had a rival on the waters of that sea. The serious menace involved in this possibility was fully realised by the Imperial government and explains the comparative magnitude of the forces which were sent to the defence of the Lazic kingdom.

It is not clear why five years were allowed to lapse before this truce of A.D. 557 was converted into a more permanent

p121
agreement. Perhaps Chosroes could not bring himself to abandon his positions in Lazica, and he knew that the complete evacuation of that country would be insisted on as an indispensable condition by the Emperor. At length, in A.D. 562, Peter the Master of Offices, as the delegate of Justinian, and Isdigunas, as the delegate of Chosroes, met on the frontiers to arrange conditions of peace.\texthebrew{⁠}\footnote{Our source for these transactions is Menander Protector, \emph{fr.} 3, \emph{De leg. Rom.} The provisions have been commented on at length by Güterbock, \emph{Byzanz und Persien}, 57 \emph{sqq.}

} The Persian monarch desired that the term of its duration should be long, and that, in return for the surrender of Lazica, the Romans should pay at once a sum of money equivalent to the total amount of large annual payments for thirty or forty years; the Romans, on the other hand, wished to fix a shorter term. The result of the negotiations was a compromise. A treaty was made for fifty years, the Roman government undertaking to pay the Persians at the rate of 30,000 gold pieces (£18,750) annually. The total amount due during the first seven years was to be paid at once, and at the beginning of the eighth year the Persian claim for the three ensuing years was to be satisfied. The inscription of the Persian document, which ratified the compact, was as follows:

``The divine, good, pacific, ancient Chosroes, king of kings, fortunate, pious, beneficent, to whom gods have given great fortune and great empire, the giant of giants, who is formed in the image of the gods, to Justinian Caesar our brother.''

The most important provision of the treaty was that Persia agreed to resign Lazica to the Romans. The other articles were as follows:

(1) The Persians were bound to prevent Huns, Alans, and other barbarians from traversing the central passes of the Caucasus with a view to depredation in Roman territory; while the Romans were bound not to send an army to those regions or to any other parts of the Persian territory. (2) The Saracen allies of both States were included in this peace. (3) Roman and Persian merchants, whatever their wares, were to carry on their traffic at certain prescribed places,\texthebrew{⁠}\footnote{Doubtless Nisibis, Dubios, and Callinicum. Cp. Güterbock, \emph{op. cit.} 78.

} where custom-houses were stationed, and at no others. (4) Ambassadors between the two

p122
States were to have the privilege of making use of the public posts, and their baggage was not to be liable to custom duties. (5) Provision was made that Saracen or other traders should not smuggle goods into either Empire by out\texthebrew{‑}of-the\texthebrew{‑}way roads; Daras and Nisibis were named as the two great emporia where these barbarians were to sell their wares.\texthebrew{⁠}\footnote{The word for smuggling is \textgreek{κλεπτοτελωνε\texthebrew{ῖ}ν.}} (6) Henceforward the migration of individuals from the territory of one State into that of the other was not to be permitted; but any who had deserted during the war were allowed to return if they wished. (7) Disputes between Romans and Persians were to be settled \texthebrew{—} if the accused failed to satisfy the claim of the plaintiff \texthebrew{—} by a committee of men who were to meet on the frontiers in the presence of both a Roman and a Persian governor. (8) To prevent dissension, both States bound themselves to refrain from fortifying towns in proximity to the frontier. (9) Neither State was to harry or attack any of the subject tribes or nations of its neighbour. (10) The Romans engaged not to place a large garrison in Daras, and also that the magister militum of the East\texthebrew{⁠}\footnote{In both these cases the same expression is used, \textgreek{τ\texthebrew{ὸ}ν τ\texthebrew{ῆ}ς \texthebrew{ἕ}ω στρατηγόν}, and must refer to the same officer. The Latin translation in Müller's edition is misleading, if not positively erroneous; in the first place the words are rendered dux orientis, in the second place praefectum orientis, which would naturally mean the Praetorian Prefect of the East. The reference of legal disputes to the Master of Soldiers is noteworthy.

\texthebrew{◂} previous

next \texthebrew{▸}Images with borders lead to more information.

The thicker the border, the more information.

(Details here.)

UP TO:

J. B. Bury
Homepage

Latin \& Greek

Texts

LacusCurtius

Home} should not be stationed there; if any injury in the neighbourhood of that city were inflicted on Persian soil, the governor of Daras was to pay the costs. (11) In the case of any treacherous dealing, as distinct from open violence, which threatened to disturb the peace, the judges on the frontier were to investigate the matter, and if their decision was insufficient, it was to be referred to the Master of Soldiers in the East; the final appeal was to be made to the sovran of the injured person. (12) Curses were imprecated on the party that should violate the peace.

A separate agreement provided for the toleration of the Christians and their rites of burial in the Persian kingdom. They were to enjoy immunity from persecution by the Magi and, on the other hand, they were to refrain from proselytising.

When the sovrans had learned and signified their approbation of the terms on which their representatives had agreed, the two ambassadors drafted the treaty each in his own language. The

p123
Greek draft was then translated into Persian, and the Persian into Greek, and the two versions were carefully collated. A copy was then made of each. The original versions were sealed by the ambassadors and their interpreters, and Peter took possession of the Persian, and Isdigunas of the Greek, while of the unsealed copies Peter took the Greek and Isdigunas the Persian. It is rarely that we get a glimpse like this into the formal diplomatic procedure of ancient times.

One question remained undecided. The Romans demanded that with the resignation of their pretensions to Lazica the Persians should also evacuate the small adjacent region of Suania. No agreement was reached by the plenipotentiaries, but the question was not allowed to interfere with the conclusion of the treaty, and was reserved for further negotiation. For this purpose Peter went in the following year ( A.D. 563) to the court of Chosroes, but Chosroes refused to agree to his argument that Suania was a part of Lazica. In the course of the conversations, the king made the remarkable proposal that the matter should be left to the Suanians themselves to decide. Peter would not entertain this, as Chosroes probably anticipated, and the negotiations fell through.

\xchapter{Chapter XVII}

Short URL for this page:

bit.ly/BURLAT17

It was the claim of the Roman Empire, from its foundation, to be potentially conterminous with the inhabited world and to embrace under its benignant sway the human race. Roman poets often spoke of it simply as the world ( orbis ). This pretentious idea, which was inherited by the Church, might well have been extinguished by the losses which Rome had sustained. Her territory had not been extended since the days of Trajan, and since the beginning of the fourth century her borders had been gradually retreating. All the western provinces were barbarian kingdoms; Italy itself, with Rome, was no more than a nominal dependency. The idea of restoring the Empire to its ancient limits seems to have floated before the mind of Justinian, but it is difficult to say whether he conceived it from the first as a definite aim of policy. He seized so promptly the opportunities which chance presented to him of recovering lost provinces in the lands of the Mediterranean, that we may suspect that he would have created pretexts, if they had not occurred.

His ambition found its first theatre in Africa. A revolution at Carthage in A.D. 531 gave the desired opportunity for intervention. The perpetual peace which Gaiseric had concluded with the Roman government ( A.D. 476) had, under his successors, been faithfully observed on both sides. There appear to have been no hostilities except during the war with Odovacar and Theoderic, when king Gunthamund took advantage of the situation to make descents on Sicily and inflicted a defeat upon the Goths.\texthebrew{⁠}\footnote{Cassiodorus, \emph{Chron., sub} 491; Dracontius, \emph{Satisfactio}, vv. 213\texthebrew{‑}214.

} The Catholic Christians endured more or less cruel

p125
persecutions at the hands of Huneric, Gunthamund, and Trasamund,\texthebrew{⁠}\footnote{The contemporary bishop of Vita wrote the story of these persecutions in his \emph{Hist. Pers. Afr. Prov.}

} and the Emperors occasionally protested.\texthebrew{⁠}\footnote{Huneric allowed the Church of Carthage to ordain a bishop at the request of Zeno and his sister-in\texthebrew{‑}law Placidia. Victor, II.2, cp. I.51.

} These kings pursued the policy of Gaiseric and looked with suspicion and jealousy on any relations between their African subjects and Constantinople. The poet Dracontius was thrown into prison by Gunthamund for celebrating the praises of a foreign potentate, and wrote a recantation and apology for his fault. The potentate was undoubtedly Zeno.\texthebrew{⁠}\footnote{Dracontius, \emph{ib.} 93:

culpa mihi fuerat dominos reticere modestos

ignotumque mihi scribere vel dominum.

Dracontius was the most considerable of the obscure Latin poets between Sidonius and Corippus. His most ambitious work was the \emph{De laudibus Dei}, but his pagan poems, the \emph{Orestes}, and the ten short pieces collected under the title of \emph{Romulea} are more interesting. In the reigns of Trasamund, who seems to have encouraged letters, and his successors there was a good deal of literary activity in Africa. We have a verse panegyric on Trasamund by Florentinus (\emph{Anthol. Lat.} No. 376), poems of Felix on the public Thermae which the same king built at Alianae, near Carthage, where the kings had a palace (\emph{ib.} 210\texthebrew{‑}214). We have also the \emph{Book of Epigrams} of Luxorius (\emph{ib.} 287\texthebrew{‑}375), of which the most interesting is that on the death of Damira, the infant daughter of Oageis, a kinsman of king Hilderic.

} But there was no breach and the relations between Trasamund and Anastasius were rather friendly.\texthebrew{⁠}\footnote{Procopius, \emph{B. V.} I.8.14.

} Then Hilderic, the son of Huneric, came to the throne ( A.D. 523).\texthebrew{⁠}\footnote{zzz GENEALOGICAL TABLE

} The fact that he was the grandson of Valentinian III was calculated to promote closer intimacy with Constantinople,\texthebrew{⁠}\footnote{The poet Florentinus (\emph{Anth. Lat.} 215) hailed him as

Vandalirice potens, gemini diadematis heres,

and reminded him of the victories of his Roman ancestors, Theodosius and Valentinian III:

ampla Valentiniani virtus cognita mundo

hostibus addictis ostenditur arce nepotis.

} and under his mild rule persecution ceased. He was the guest-friend of Justinian, and that astute prince probably aimed at making the Vandal state a dependency of the Empire, through his influence on the unwarlike king.\texthebrew{⁠}\footnote{Procopius, \emph{B. V.} I.9.8 , cp. 19\texthebrew{‑}23 .

} Hilderic's complaisance

p126
to Constantinople aroused dissatisfaction; the opposition was headed by his cousin Gelimer, who usurped the throne in A.D. 530 and threw Hilderic into prison. Justinian at once intervened. He addressed to the usurper a letter of remonstrance, appealing to the testament of Gaiseric and demanding the restoration of the rightful king. Gelimer replied by pla­cing Hilderic under a stricter guard. The Emperor then despatched an ultimatum requiring Gelimer to send the deposed sovran to Constantinople, otherwise he would regard the treaty with Gaiseric as terminated. Gelimer replied defiantly that the matter concerned the Vandals themselves, and that it was not Justinian's business. He probably saw through Justinian's designs and knew that if he yielded he might postpone but would not avert war.\footnote{Cp. Diehl, \emph{L'Afrique byzantine}, p6. The true form of Gelimer's name is Geilamir (so his coins and CIL VIII.17 .412). The date of his usurpation 530 follows from the length of Hilderic's reign given as 7 years by Procopius, 7 years 3 months by Victor Tonn., and in the shorter edition of the Vandal \emph{Laterculus Regum} as 7 years 14 days. If we take the last figure we get May 19, 530, as the day of Hilderic's defeat. Victor Tonn. places it in 531, and John Mal. places the application of Hilderic to Justinian in the same year (XVIII.459). 531 is the date usually assigned by modern writers (Clinton, Diehl, etc.); but Schmidt is right in deciding for 530 (\emph{Gesch. der Wand.} p124).

}

The Emperor decided that the time had come to attempt the conquest of Africa, and as soon as peace had been concluded with Persia in spring of A.D. 532, the preparations were hurried forward. In his eyes it was no war of aggression; it was the suppression of tyrants in provinces over which the Emperors had always tacit ­ly reserved their rights ( iura imperii ). The ecclesiastics were ardently in favour of an enterprise which would rescue their fellow-Catholics in Africa from the oppression of Arian despots.\texthebrew{⁠}\footnote{The war was also welcomed by the eastern traders residing at Carthage, who saw in the reunion of Africa with the Empire advantage to their commercial interests. Procopius, \emph{B. V.} I.20.5.

} But from his counsellors and ministers Justinian received no encouragement. The disaster of the great expedition of the Emperor Leo was not forgotten. Their minds were still possessed by the formidable prestige which the Vandal power had attained under Gaiseric both by land and sea. The Empire had not kept up a power ­ful navy, and without command of the sea the hazard of attempting to transport an army and land it on a hostile coast could not be denied. The Praetorian Prefect, John of Cappadocia, explained to the Emperor the difficulties and risks of the undertaking in the plainest words,

p127
and earnestly endeavoured to dissuade him from an adventure which the opinion of experts unreservedly condemned. And this view was justified, although its advocates probably had not realised how far the military strength of the Vandals had decayed since the days of Gaiseric. But notwithstanding this decline, the events of the campaign show that if Gelimer had not committed the most amazing mistakes, which his enemies could not have foreseen, the Roman army would probably have suffered an inglorious defeat. Justinian turned deaf ears to the gloomy anticipations of his counsellors, he believed in the justice of his cause, he believed that Heaven was on his side,\texthebrew{⁠}\footnote{For the religious motives cp. Procopius, \emph{ib.} 10.10.19\texthebrew{‑}20 ; Diehl, \emph{op. cit.} 7\texthebrew{‑}8.

} and he had confidence in the talents of his general Belisarius, whom he destined to the command of the expedition and invested with the fullest powers, giving him a new title equivalent to imperator , which had long been restricted to the Emperors themselves.\footnote{Procopius (\emph{ib.} 11.20) does not actually say that he was designated as \textgreek{στρατηγ\texthebrew{ὸ}ς α\texthebrew{ὐ}τοκράτωρ}, but that all his acts were to be valid \textgreek{\texthebrew{ἄ}τε α\texthebrew{ὐ}το\texthebrew{ῦ} βασιλέως α\texthebrew{ὐ}τ\texthebrew{ὰ} διαπεπραγμένου}. But as he had ceased to be mag. mil. per Orientem, and nothing is said of his appointment to another of the regular military commands, we may infer that it was on this occasion that Justinian introduced the new and exceptional post of \textgreek{στρατ. α\texthebrew{ὐ}τοκράτωρ}. It is to be remembered that \textgreek{α\texthebrew{ὐ}τοκράτωρ} is the \emph{official} equivalent of imperator. We shall hereafter meet other commanders bearing the same title and authority (Germanus, Narses, Justin).

}

The small numbers of the army, deemed sufficient for the conquest of a people who had the military reputation of the Vandals, is surprising. It consisted of not more than 16,000 men. Perhaps this was as much as it was considered possible to transport with safety; and if it were annihilated, the loss would not be irreparable. There were 10,000 infantry, which were drawn partly from the Comitatenses and partly from the Federates. There were 5000 excellent cavalry, of whom more than 3000 were similarly composed, and the remainder were private retainers of Belisarius.\texthebrew{⁠}\footnote{Procopius, \emph{ib.} 11.2 \emph{sqq.} The bucellarii (\textgreek{δορυφόροι κα\texthebrew{ὶ ὑ}πασπισταί}) of Belisarius were probably at least 1400 or 1500 (cp. Diehl, \emph{ib.} 17 note. It seems clear from the whole context in Procopius that the Heruls and Huns were not included in the 5000 cavalry (though Diehl hesitates).

} There were two additional bodies of allied troops, both mounted archers, 600 Huns and 400 Heruls. The whole force was transported on 500 vessels, guarded by ninety-two dromons or ships of war.

The hundred years of their rule in Africa had changed the spirit and manners of the Vandals. They had become less

p128
warlike; they had adopted the material civilisation and luxuries of the conquered provincials; and their military efficiency had declined since Gaiseric's death. It may be doubted whether their army numbered more than 30,000 men.\texthebrew{⁠}\footnote{Diehl, \emph{op. cit.} p9, says less than 40,000; cp. Pflugk Harttung, \emph{Hist. Zeitschrift}, LXI p70 (1889). Note that the figure of 80,000 warriors given in

Procopius, \emph{H. A.} 18.6

is merely a repetition of his mistake in \emph{B. V.} I.5.18 (see above,

Vol. I, p246 ).

} It consisted entirely of cavalry, arrayed in inferior armour, who fought with lance and sword, and were, like other German peoples, unskilled in archery and the use of the javelin. Their king, although he was more martial than his predecessor, was a man of sentimental temperament, who had no military or political talents. The situation required a leader of exceptional ability. For the kingdom was divided against itself. Gelimer's Roman subjects longed for restoration to the Empire and would do all they could to assist the invaders. Even among the Vandals there were the adherents of Hilderic. The Moorish tribes of the interior could not be trusted to remain friendly or neutral if fortune seemed to incline to the Roman cause.

Before the Imperial army set sail from the Bosphorus, two events happened, and Gelimer committed two astounding blunders. The inhabitants of Tripolitana\texthebrew{⁠}\footnote{Led by a certain Pudentius, who was in correspondence with Justinian and was assisted by a small body of troops sent from Constantinople. Procopius, \emph{ib.} 10.5.

} revolted from the Vandals, and Gelimer made no attempt to recover it. This was a fatal policy, for it would enable the Roman army, if it reached the coast of Africa in safety, to land on a friendly soil. Shortly before this the Vandal governor of Sardinia\texthebrew{⁠}\footnote{He was a Goth, named Godas.

} had proclaimed himself independent of Carthage, and when he heard of Justinian's project he offered his submission to the Emperor. Gelimer despatched a force of 5000 men and 120 ships to recover the island. He thus deprived himself of a considerable fraction of his army and virtually of his whole effective naval strength.\texthebrew{⁠}\footnote{Procopius describes the ships as ``the best sailers'' (\emph{ib.} 11.24) . If they were only part of the fleet, the rest was not strong enough to attempt any action. No inference can be drawn from \textgreek{παντ\texthebrew{ὶ} τ\texthebrew{ῷ} στόλ\texthebrew{ῷ}} (\emph{ib.} 25.17 and 21 ), which means the whole Sardinian squadron.

} The Vandal fleet which was reputed so formidable played no part in the war. This curious perversity of Gelimer, in wasting his strength on the recovery of a distant island whose disaffection could hardly have affected the course of events,\texthebrew{⁠}\footnote{Gelimer, no doubt, believed that the Sardinian expedition would return before the enemy landed in Africa.

} and

p129
neglecting to suppress the movement in Tripolitana, whose possession was of the first importance, was perhaps decisive for the whole issue of the war.

If the Sardinian revolt was a piece of luck for Justinian, the attitude of Italy was hardly less fortunate. After the death of Trasamund, his Ostrogothic wife Amalafrida had been imprisoned and afterwards murdered,\texthebrew{⁠}\footnote{See

next chapter, § 1, p158 .

} and this led to an irreconcilable breach between the courts of Carthage and Ravenna. The Ostrogothic government willingly supported the Imperial expedition by pla­cing the harbours of Sicily at its disposal.

The Roman forces set sail from Constantinople in June A.D. 533. Before their departure the ship of the general moored in front of the Imperial palace, and the Patriarch offered prayers for the success of the expedition. Among those who witnessed their sailing perhaps most who were competent to judge believed that they would never return. Belisarius was accompanied by his wife Antonina, and by the historian Procopius, who again acted as his legal assessor, and to whom we owe the story of the war. The domesticus , or chief of the general's staff, was the eunuch Solomon, a native of Mesopotamia, one of those able eunuchs whom we frequently meet on the stage of Byzantine history.

The voyage from the Bosphorus to Sicily was marked by many halts,\texthebrew{⁠}\footnote{Nine days were spent at Heraclea and Abydus, so, as the expedition sailed ``about the summer solstice'' (June 21), it left the Dardanelles about July 1. There was a long delay at Methone (Modon, in Messenia) where the army suffered (just as modern armies so often suffer from the dishonesty of contractors) from the greed of the Praetorian Prefect, on whom it devolved to provide the soldiers with the bread necessary for the voyage. It was found that the bread had gone bad, because it had been baked only once, instead of twice. Five hundred soldiers fell victims to dysentery caused by the putrefying dough. The Prefect had saved both fuel and flour. Belisarius was praised for complaining to the Emperor, but no punishment was inflicted on the guilty.

} and the shore of Africa was not reached till the beginning of September. Procopius commemorates the practical foresight of Antonina in storing a large number of jars of water, covered with sand, in the hold of the general's ship, and tells how this provision stood them in good stead in the long run from Zacynthus to Catane. Belisarius had been full of misgivings about the voyage from Sicily to Africa, expecting that the enemy would attack him by sea. He now learned for the

p130
first time (from a man who had just arrived from Carthage) that the Vandal fleet had been sent to Sardinia; and equally welcome was the news that Gelimer was unaware that the Roman expedition was on its way and had made no preparation to meet it, at Carthage or elsewhere.

The fleet made land at Caputvada (Ras Kapudia) on the African coast, and the army disembarked and fortified a camp. Before landing Belisarius had held a council of war, and some of his generals argued that it would be the better plan to sail straight for Carthage and surprise it, but Belisarius overruled this view; there was the chance of a hostile fleet appearing, and he knew that the soldiers were afraid of a naval attack. Caputvada is \texthebrew{•} sixty-six Roman miles south of Hadrumetum (Sousse) and one hundred and sixty-two from Carthage,\texthebrew{⁠}\footnote{Procopius reckons the distance as five days' journey for an unencumbered man, \emph{B. V.} I.14.17 .

} so that if his army marched \texthebrew{•} slightly over eleven miles a day, he was fourteen days' journey from his goal. The road ran close to the coast, and the fleet was instructed to sail slowly and keep within hail of the army. A squadron of 300 horse, under John the Armenian, was sent ahead as an advance guard at a distance of \texthebrew{•} three miles, and the corps of 600 Huns was ordered to march at the same distance to the left of the road, to protect the army from a flank attack. The first town on their route was Syllectum (Selketa), which was seized quietly by a ruse. The overseer of the public post deserted and delivered all the horses to Belisarius, who rewarded him with gold and gave him a copy of a letter addressed by the Emperor to the leading men\texthebrew{⁠}\footnote{\textgreek{\texthebrew{Ἄ}ρχοντες.}} of the Vandals, to make public. It ran thus:

``It is not our purpose to go to war with the Vandals, nor are we breaking our treaty with Gaiseric. We are only attempting to overthrow your tyrant, who making light of Gaiseric's testament keeps your king a prisoner, and killed those of his kinsmen whom he hated, and having blinded the rest keeps them in prison, not allowing them to end their sufferings by death. Therefore join us in freeing yourselves from a tyranny so wicked, that you may enjoy peace and liberty. We give you pledges in the name of God that we will give you these blessings.

As the man did not venture to publish the letter openly but

p131
only showed it secretly to his friends, it produced no effect. During their march northward the friendliness of the inhabitants supplied the invaders with provisions, and Belisarius took the strictest measures to prevent his soldiers from alienating the sympathies of the population by marauding and looting. It will be remembered how in England's war with her American colonies the shameless pillaging of the property of the colonial loyalists, by the Hessian mercenaries whom George III had hired, drove them into the ranks of the rebels, and the English generals were incapable of keeping a firm hand on their auxiliaries. Belisarius had a more difficult task. Want of discipline, as we shall see, was the weak point in his mixed army. But for the present he succeeded in restraining the appetites of his barbarian troops, and advanced comfortably towards the Vandal capital.

Passing Thapsus, Leptis, and Hadrumetum, the army reached Grasse, where the Vandal kings had a villa and a beautiful park, full of fruit trees, and as the fruit was ripe the soldiers ate their fill. This place, now Sidi-Khalifa , is still famous for its fruit gardens.\texthebrew{⁠}\footnote{See Tissot, \emph{Géographie}, II p116. It is close to Fradiz, the ancient Aphrodisium.

} During the night of the halt at Grasse some of the Roman scouts met enemy scouts and after exchanging blows both parties retired to their camps. Thus Belisarius learned for the first time that the enemy was not far away. It was, in fact, the king who was following them but keeping out of sight. Gelimer was at Hermiane\texthebrew{⁠}\footnote{In Byzacena.

} when he learned of the Roman disembarkation. He sent orders immediately to his brother Ammatas at Carthage to kill Hilderic and the other prisoners, and, collecting all the troops in the city, to be ready to attack the Roman army at a given time and place. He marched southward himself at the head of his army to follow and observe the advance of the invaders without being seen himself. His plan was to surprise and surround the enemy at a spot near Tunis and \texthebrew{•} ten miles from Carthage.

Not far from Grasse the high road to Carthage left the coast and crossed the promontory which runs out into Cape Bon. Here the army and the ships parted company, and the naval commander was instructed not to put in at Carthage but to

p132
remain \texthebrew{•} about three miles out at sea until he should be summoned. The road rejoined the coast at Ad Aquas, which is now Hammam el\texthebrew{‑}Enf, \texthebrew{•} twenty-three miles from Carthage. By the fourth day\texthebrew{⁠}\footnote{From Grasse. The distance to Ad Aquas is \texthebrew{•} about 50 miles. The date was the eve of St. Cyprian's day, Sept. 13 (\emph{B. V.} I.21.23) . If the army landed at Caputvada on Sept. 2, they reached Syllectum (a long day's march of 19 miles) on Sept. 3, Hadrumetum Sept. 5, Grasse Sept. 9, Ad Aquas Sept. 12. This would mean much longer marches between Grasse and Ad Aquas than the \texthebrew{•} 80 stades (11½ miles) which Procopius says was the average day's march.

} (September 13) the army was approaching Tunis, and it was perhaps at the northern extremity of the defile of Hammam el\texthebrew{‑}Enf, on a rocky spur of the Jebel Bu-Kornin \texthebrew{—} the two-horned hill \texthebrew{—} that Belisarius, neglecting no precautions and hesitating to risk an engagement with his whole army, made a stockaded camp in which he ordered his infantry to remain while he rode down into the plain with the cavalry.\texthebrew{⁠}\footnote{The position of the camp, at Darbet es-Sif, is Tissot's plausible conjecture, \emph{op. cit.} p121. Procopius notes that the place was \texthebrew{•} 7 miles from Decimum (\emph{ib.} 19.1) .

} John the Armenian had ridden on in advance, as usual, while the Huns were some miles to the left, west of the Bu-Kornin hills. Belisarius had no idea of the excellent strategic plan which the enemy had devised to destroy him.

If we walk out of the modern town of Tunis by the south-eastern gate, Bab Alleona, we soon reach the railway station of Jebel Jellud, and near it was the Roman station Ad Decimum, at the tenth milestone from Carthage. On the left are a number of little eminences of which the highest is named Megrin, on the right the hill of Sidi Fathalla, behind which extends to the west the Sebkhaes-Sejuni or Salt-plain, an arid treeless tract then as now.\texthebrew{⁠}\footnote{\textgreek{Πεδίον \texthebrew{Ἀ}λ\texthebrew{ῶ}ν}, \emph{B. V.} I.18.12 This indication and the reference to the hills on either side of the road, \emph{ib.} 19.19, are the important determinants in the identification of Ad Decimum, which is due to Tissot. The fact that the place was a mutatio \texthebrew{•} ten Roman miles from Carthage is not enough as we do not know how far the city of Carthage extended southward and from what point the distance was measured. Tissot (\emph{ib.} 114 \emph{sqq.}) has thrown much light on the topography of the battle.

} This was the place in which Gelimer had planned to surround the Romans. Ammatas coming from Carthage was to confront them in the defile; when they were engaged with him, Gibamund, the king's nephew,\texthebrew{⁠}\footnote{Gibamund is mentioned as the builder of Thermae in a metrical inscription found at Tunis (CIL VIII.25362 ):

Gaude operi, Gebamunde, tuo, regalis origo,

deliciis sospes utere cum populo.

} with 2000 men, advancing across the Salt-plain, was to descend from the hill on their left,

p133
while Gelimer himself with the main army was to come upon them in the rear. The time at which the Romans might be expected to reach Ad Decimum was nicely calculated, and the plan all but succeeded.

Ammatas committed the error of appearing with a few men at Ad Decimum some hours before the appointed time, probably for the purpose of surveying the ground. He arrived at noon and came face to face with the troops of John. He was a brave warrior and he killed with his own hand twelve of John's best men before he fell himself. His followers fled and swept back in a hot-foot race to the shelter of Carthage the other troops who were marching negligently in bands of twenty or thirty to the appointed place. John and his riders pursued and slew as far as the city gates.

While this action was in progress, the Huns had reached the Plain of Salt and fell in with the forces of Gibamund who were moving eastward to Sidi Fathalla, and, although in numbers they were less than one to three, utterly annihilated them.

p134
The Huns enjoyed the battle; the Vandals, they thought, were a feast which God had prepared for them.\footnote{\emph{B. V.} I.18.18. From the statement (\emph{ib.} 12) that the Salt-plain is 40 stades from Decimum, we may infer that the engagement occurred at that distance \texthebrew{•} (5 to 6 miles). The eastern edge of the Salt-plain is much nearer Decimum.

}

Of these two events Belisarius knew nothing as he descended from Hammam el\texthebrew{‑}Enf into the plain of Mornag. His Federate cavalry rode in advance, the regular cavalry and his own retainers at some distance in the rear. Crossing the stream Oued Miliane, the road to Tunis passes Maxula (Rades), which lies between the sea and the southern shore of the lake of Tunis.\texthebrew{⁠}\footnote{From Maxula the shortest road to Carthage was along the shore, but this way was impracticable on account of the canal connecting the sea with the lake.

} The Federates, when they reached Ad Decimum, saw the corpses of their comrades and those of Ammatas and some Vandals. The people of the place told them what had happened and they climbed the hills to reconnoitre. Presently they discerned a large cloud of dust to the south and then a large force of Vandal cavalry. They sent, at once, a message to Belisarius urging him to hasten. It was Gelimer's army that was coming. Having followed Belisarius at a safe distance along the main road he had doubtless left it at Grombalia, and keeping to the west of the Jebel Bu-Kornin proceeded along a road which is still used by the natives for travelling between Grombalia and Tunis. The hilly nature of the ground did not permit him to see either the movements of Belisarius on his right or the disaster of his nephew on his left. When his vanguard reached Ad Decimum there was a contest with the Roman Federates to win possession of an eminence (possibly Megrin), in which the Vandals were success ­ful. The Federates then fled for a mile along the road to rejoin their own army and met Uliaris with 800 guardsmen, who seeing them galloping in disorder turned themselves and galloped back to Belisarius.

Gelimer now had the victory in his hands, but the gods were determined to destroy him. The historian who tells the tale and who witnessed the cavalry riding back in terror to the commander-in\texthebrew{‑}chief , declares that ``Had Gelimer pursued immediately I do not think that even Belisarius would have withstood him, but our cause would have been utterly ruined, so large appeared the multitude of the Vandals and so great the

p135
fear they inspired; or if he had made straight for Carthage he would have slain easily all the men with John, and would have preserved the city and its treasures, and would have taken our ships which had approached near, and deprived us not only of victory but of the means of escape.''\footnote{\emph{B. V.} I.19.25 \emph{sqq.}

}

Gelimer was a man of sentimental temperament. When he reached Ad Decimum and saw the dead body of his brother he was completely unmanned. He set up loud lamentations and could think of nothing but burying the corpse; and so, as the historian remarks, ``he blunted the edge of opportunity,'' and such an opportunity did not recur.

Meanwhile Belisarius had rallied the fugitives and administered a solemn rebuke. On learning exactly what had happened, he rode at full speed to Decimum and found the barbarians in complete disorder. They did not wait for his attack but fled as fast as they could, not towards Carthage but westward towards Numidia. They lost many, and the fighting ended at night, when John's troops and the Huns arrived on the scene. A considerable victory had been gained, but it was a victory which Gelimer had presented to Belisarius; it ought to have been a defeat.

The night was passed at Decimum, and on the following day Antonina arrived with the infantry and the whole army marched to Carthage, arriving at nightfall. Its inhabitants opened the gates and welcomed the victor with a brilliant illumination. But Belisarius was cautious, and he would not enter that night, partly because he feared an ambuscade and partly because he was resolved that his soldiers should not plunder the city. The next day (September 15) the army marched in, in formation of battle. Belisarius need not have been afraid; no snare was set.

He seated himself on the king's throne, and consumed the dinner which Gelimer had confidently ordered to be ready for his own victorious return. The inhabitants welcomed the deliverer, and the Imperial fleet sailed into the lake of Tunis. Belisarius lost no time in repairing the walls of the city and rendering it capable of sustaining a siege. Meanwhile the Moorish tribes of Numidia and Byzacium, learning the issue of the battle, hastened to send friendly embassies to the conqueror.\footnote{Before the Vandal occupation it had been the custom of the client Moorish chiefs to receive as tokens of office from the Emperor a gilded silver staff and a silver cap in the form of

(p136)
a crown, a white cloak, a white tunic, and a gilded boot. Belisarius sent these to them now and gave them presents of money. \emph{B. V.} I.25.7.

}

p136
Gelimer and his vanquished army had fled to the plain of Bulla Regia.\texthebrew{⁠}\footnote{Hammam Daraji, on the border of the Proconsular province and Numidia.

} His first care was to send the bad news to his brother Tzazo, who commanded the Sardinian expedition, imperatively recalling him. Tzazo, who had succeeded in re-establishing the Vandal authority in Sardinia, returned with his troops, and Gelimer thus reinforced marched towards Carthage. He cut the aqueduct, and he attempted to prevent provisions from arriving in the city, which he hoped to reduce by blockade. He sent secret agents to undermine the loyalty of the inhabitants and the Imperial army. In this he had some success. The auxiliary Huns seem to have determined to stand aloof in the approaching struggle and then rally to the aid of the victorious party.\footnote{For the discontent of the Huns, who feared that if the Romans were victorious they would be kept in Africa, see \emph{B. V.} II.1.6. Belisarius swore to them that when the Vandals were defeated he would send them home with all their booty, but notwithstanding this they played a double game at Tricamaron (\emph{ib.} 2.3).

}

About the middle of December Belisarius judged that the time had come to bring matters to an issue. Gelimer had pitched his camp at Tricamaron,\texthebrew{⁠}\footnote{The place has not been identified.

} on the banks of the Mejerda, \texthebrew{•} about twenty miles west of Carthage. Here were collected not only his soldiers but their wives and children and property. The battle of Tricameron was in some respects a repetition of the battle of Ad Decimum. It was a battle of cavalry. The Roman infantry was again far behind and did not come up till the late afternoon when the issue was virtually decided. It was only after repeated charges that the mailed Roman horsemen succeeded in breaking the enemy's lines. Tzazo and many others of the bravest officers fell. The Vandals fled to their camp, and the Huns who had hitherto refused to join in the combat now joined in the pursuit. As soon as the infantry arrived, the victors fell upon the camp, and Gelimer, seeing that all was lost, fled with a few attendants into the wilds of Numidia. All his soldiers who could escape sought refuge in the churches of the surrounding district. There was no pursuit. The Roman troops thought of nothing but of seizing the rich spoil, women and treasures, which awaited them in the camp.

p137
The general was utterly powerless to restore discipline, and he passed an anxious night. He feared that some of the enemy, realising the situation, would attack his disorderly troops; and ``if any thing of the kind had happened,'' says Procopius, ``I think that not a Roman would have escaped to enjoy his booty.'' The victory of Tricamaron (middle of December, A.D. 533) destroyed the Vandal kingdom. But it was due to the weakness and incompetence of the king. He had no idea of using to advantage his great numerical preponderance in cavalry. Even after the defeat, if he had not run away, he might have annihilated the enemy busy with their loot.

It is to be observed that both the actions of the short campaign were fought and won by the Roman cavalry, as in the battle of Daras. The more numerous infantry might almost as well not have been in Africa. There is room for wonder whether if Belisarius had been opposed to a commander of some ability and experience in warfare, he would not have been hopelessly defeated. His secretary, Procopius, expresses amazement at the issue of the war, and does not hesitate to regard it not as a feat of superior strategy but as a paradox of fortune.\texthebrew{⁠}\footnote{\emph{Ib.} 7.18 \emph{sqq.} See Diehl, \emph{op. cit.} 30 \emph{sqq.}

} But if in this campaign Belisarius did not display signal military talent, there can be no question as to his skill in holding together the undisciplined and heterogeneous troops which he commanded. The Federates thought of nothing but securing booty; they were inclined to regard themselves as independent allies; again and again, but for the general's firmness and tact, their insubordinate spirit might have been disastrous.

The Vandal warriors who had fled to the asylum of sanctuaries surrendered to the Roman general, who promised that they would be well treated and sent to Constantinople in spring. All the treasures belonging to Gelimer were seized in Hippo Regius.\texthebrew{⁠}\footnote{A silver basin, found in 1875 not far from Feltre, with the inscription Geilamir Vandalorum et Alanorum rex (CIL VIII.17 .412), must have come from this treasure. Mommsen (\emph{Hist. Schr.} I p566) conjectures it may have been a gift from Belisarius to a Herul officer, who might have taken it to Italy.

} Belisarius then made arrangements to assert the Imperial authority throughout the Vandal dominions, of which he had yet occupied but a small part. He sent detachments by sea to take possession of Sardinia and Corsica, the Balearic

p138
Islands, the fortress of Septum in Tingitana, on the straits of Gibraltar, and Caesarea (Cherchel) on the coast of Mauretania. But the task of establishing Roman administration throughout the African provinces, and especially in the three Mauretanias, was to require several years and far more strenuous military exertions than were needed to destroy the power of the Vandals.

Gelimer had fled to Mount Papua in the wilds of Numidia, where he found among the Moors a miserable but impregnable refuge. Here for three months he and the friends who were with him endured hunger and cold, blockaded by the Herul leader Pharas, whose followers watched the paths at the foot of the mountain. It was a tedious watch during the cold winter months. Pharas sent a friendly message to the king counselling him to surrender. The pride of Gelimer could not yet brook the thought, but he besought Pharas to send him a loaf, a sponge, and a lyre. He had not tasted baked bread since he had come to the mountain; he wanted a sponge to dry his tears; and a lyre that he might sing a song which he had composed on his misfortunes. The curious request, which was readily granted, illustrates the temperament of Gelimer who loved the luxury of grief.\texthebrew{⁠}\footnote{We saw how he indulged it at an inopportune moment, at the battle of Ad Decimum. His meeting with his brother Tzazo gave him another opportunity.

} At length (in March) pitying the sufferings of his faithful attendants, he surrendered, assured of honourable treatment. He was taken to Constantinople, where he adorned the triumph of Belisarius. When he saw the Emperor sitting in all his splendour in the Kathisma of the Hippodrome, he repeated to himself, ``Vanity of vanities, all is vanity.'' An ample estate in Galatia was granted to him, and the dignity of Patrician would have been conferred on him, if he had not resolutely refused to abandon his Arian religion.\footnote{The progress of bigotry is to be noted. This new condition for the Patriciate was evidently laid down by Justinian. In the fifth century Aspar and Theoderic, both Arians, had been created patricians.

}

The difficulties of the command of Belisarius were illustrated by the intrigues which the subordinate generals began to spin against him after his final success. They wrote secretly to Constantinople insinuating that he was aiming at the throne. Justinian doubtless knew what these charges were worth. He gave Belisarius the choice of returning to Constantinople or of

p139
remaining in Africa. Belisarius prudently chose to return, and was rewarded by a triumph, which at this time was an exceptional honour for a private person ( A.D. 534). He brought back with him a captive king with the choicest of the Vandal warriors;\texthebrew{⁠}\footnote{Most of them were formed into five cavalry regiments, known as Vandali Iustiniani, and stationed on the Persian frontier, Procopius, \emph{ib.} 14.17 . Some entered the private service of Belisarius

(\emph{B. G.} III.1) . Perhaps there were about 3000 in all.

} an immense treasure; and what above all appealed to the piety of the Emperor and to the sentiment of orthodox Christians, King Solomon's golden vessels of which Gaiseric had robbed Rome and of which Titus had despoiled Jerusalem.\texthebrew{⁠}\footnote{Also the Imperial ornaments which Gaiseric had taken from Rome,

\emph{C. J.} I.27.1, § 6 . For the mosaics on the walls of the Chalce see above,

p54 .

} He was soon to be entrusted with the conduct of a longer and more arduous enterprise.

The general idea of the Emperor's scheme for the administration of the African provinces was to wipe out all traces of the Vandal conquest, as if it had never been, and to restore the conditions which had existed before the coming of Gaiseric. The ecclesiastical settlement, which lay near Justinian's heart, was easy and drastic.\texthebrew{⁠}\footnote{Justinian, \emph{Nov.} 37.

} All the churches which the conquered Arians had taken for their own worship were restored to the Catholics, and heretics were treated with the utmost intolerance. Vandals, even those who were converted from their religious errors, were excluded from public offices. The rank and file of the Vandal fighting men became the slaves of the Roman soldiers who married the women. All the estates which had passed into the hands of the barbarians were to be restored to the descendants of the original owners who could establish their claims, \texthebrew{—} a measure which led to the forgery of titles and endless lawsuits.\texthebrew{⁠}\footnote{\emph{Nov.} 36 (A.D. 535) makes provisions to remedy these evils.

} The ultimate result of the whole policy was the disappearance of the Vandal population in Africa.

When he received the news of the victory of Tricamaron, Justinian must have proceeded immediately, if he had not already begun, to prepare the details of the future government of Africa; for the whole scheme was published in April A.D. p140 534.\texthebrew{⁠}\footnote{\emph{C. J.} I.27.1 and 2 . The total cost of the administration was less than £30,000. Cp. above,

Vol. I, p33, n1 .

} Its general character was modelled on the system which was in force before the Vandal conquest, but the changed circumstances required some modifications. Formerly Africa had been a diocese of the Prefecture of Italy. This arrangement could not be maintained as Italy was in the hands of the Ostrogoths. Hence the civil governor was invested with the title of Praetorian Prefect of Africa, and enjoyed the corresponding dignity and emoluments. Under him were the governors of the seven provinces: Proconsularis, Byzacena, Tripolitana, Numidia, the two Mauretanias, and Sardinia.\texthebrew{⁠}\footnote{\emph{C. J.} \emph{ib.} I § 12 . In the text we must read Zeugi for Tingi, and Mauritaniae for Mauritania. See Diehl, \emph{op. cit.} 107 \emph{sqq.} Proconsularis Carthage (=Zeugi), Byzacena and Trip. were under Consulares; the other four under praesides. On the old system, Numidia has a consularis, Trip. a praeses. This change may be explained, as Diehl suggests, by the fact that in 534 Tripolitana was regarded as entirely conquered, while most Numidia had still to be occupied.

} But the compass of the Second or Western Mauretania (Caesariensis) was extended so as to include Tingitana, which in old days had belonged to the diocese of Spain.

The military establishment was placed under a Master of Soldiers,\texthebrew{⁠}\footnote{Magister militum Africae. Under him was a magister peditum (Proc. \emph{B. V.} II.16.2).

} a new creation, since in old days the armies of Africa had been under the supreme command of the Master of Soldiers in Italy. The fundamental distinction between the mobile army and the frontier troops was retained. The mobile army consisted of the divisions of the comitatenses who had been sent with Belisarius, of foederati , and of native African troops ( gentiles ).\texthebrew{⁠}\footnote{It included also a small body of guard troops (Excubitores), who were sent with Solomon in 534.

} The frontier troops were distributed in four districts, under dukes, who had authority also over mobile troops stationed in these military provinces.\texthebrew{⁠}\footnote{Tripolitana, Byzacena, Numidia, and Mauretania. The chief stations, where the dukes resided, were respectively Leptis Magna, Capsa or Thelepte, Cirta, and Caesarea. A commander of subordinate dignity, with the title of tribunus, was stationed at Septum. The military received larger salaries than the civil governors.

} The establishment of this organisation throughout Africa was retarded for some years by wars and mutinies, but it was begun by Belisarius before he departed, and it was gradually carried out, along with an elaborate scheme of fortification against the inroads of the Moorish tribes.

The Moors began hostilities before the Romans had time to make provision for the defence of the country or to organise

p141
the new civil administration. The situation was so grave that Justinian, when he sent Solomon in autumn ( A.D. 534) to replace Belisarius, united in his hands the supreme civil as well as military authority. Solomon was Praetorian Prefect as well as Master of Soldiers.\texthebrew{⁠}\footnote{The first Pr. Pr. of Africa had been Archelaus

(\emph{C. J.} I.27.1) . Solomon was the first mag. mil. Afr.

} This appointment struck the note of a change in the principles of provincial administration which had prevailed since Diocletian. We shall see how elsewhere Justinian departed from the general rule of a strict separation of the civil and military powers. In Africa, although the two offices were seldom united, perhaps only on three occasions,\texthebrew{⁠}\footnote{Solomon (534\texthebrew{‑}536); Solomon (539\texthebrew{‑}543); Theodore (569; John Bicl. \emph{sub a.}, cp. Diehl, \emph{op. cit.} 599). Perhaps Sergius (544) should be added (Marcellinus, \emph{sub} 541).

} there is a tendency from the beginning to subordinate the Praetorian Prefect to the Master of Soldiers,\texthebrew{⁠}\footnote{This seems to be the case with Symmachus under Germanus (536) and Athanasius under Areobindus (546). Diehl, \emph{op. cit.} 117.

} and before the end of the century the Master of Soldiers will become a real viceroy with the title of Exarch.

The leading feature of the history of North Africa from the Roman reconquest to the Arab invasion in the middle of the seventh century is a continuous struggle with the Moors, broken by short periods of tranquillity. Each province had its own enemies. Tripolitana was always threatened by the Louata, Byzacena by the Frexi;\texthebrew{⁠}\footnote{There were other tribes besides the Louata and Frexi, but these were the most prominent. In Justinian's time, Antalas was the chieftain of the Frexi and their confederate tribes; Cutsina was the chief of other tribes in the same region; Iabdas was king of the Aurasian Moors; Mastigas and Masuna were the leading princes of the Mauretanian Moors.

} the townspeople of Numidia lived in dread of the Moors of the Aurasian hills. Mauretania was largely occupied by Berber tribes. The Roman government never succeeded in effecting a complete subjugation of the autochthonous peoples. It was not an impossible task, if the right means had been taken. But the Roman army was hardly sufficient in numbers to maintain effectively the defence of a long frontier, against enemies whose forces consisted of light cavalry, immensely more numerous. This numerical inferiority might have mattered little if the troops had been trustworthy. But they were always ready to revolt against discipline, and in war their thoughts were not on protecting the provinces but on

p142
securing booty. They could do work under a commander who knew how to handle them, but such commanders were rare. Most of the military governors found their relations with their own soldiers as difficult a problem as their relations with the Moors. Here we touch on a second cause of the failure of the Romans to secure a lasting peace in Africa \texthebrew{—} the unfitness of so many of their military governors. A succession of men like Belisarius, Solomon, and John Troglita would probably have succeeded, if not in establishing permanent and complete tranquillity, at least in defending the frontiers efficiently. But when a commander of this type had weathered a crisis or retrieved a disaster, he was too often succeeded by an incompetent man, who had no control over the soldiers, no skill in dealing with the Moors, and who undid by his inexperience all that his predecessor had accomplished. And apart from these weaknesses, it has been remarked with justice that the general military policy was not calculated to pacify the restless barbarians beyond the frontier. It was a policy of strict defence. The elaborate system of fortresses which were speedily erected throughout the provinces stood the inhabitants in good stead, but they did not prevent raids, and the Romans only opposed raids on Roman soil. Far more would have been effected if the Romans had taken the offensive whenever there was a sign of restlessness and sent flying columns beyond the frontier to attack the Moors on their own ground. Finally the want of success in dealing with the Moorish danger may have been partly due to defective and inconsistent diplomacy.\footnote{Illustrations of all these points will be found in Diehl, \emph{op. cit.}

}

The one fact in the situation which enabled the Romans to maintain their grip on Africa was the disunion among the Moors. On more than one occasion they suffered such crushing disasters that if the Moors had made a determined and united effort the Imperial armies would easily have been driven into the sea. But the jealousies and quarrels among the chieftains º hindered common action; and if one began a hostile movement, the Romans could generally depend on the quiescence or assistance of his neighbour.

On his arrival in Africa ( A.D. 534) Solomon\texthebrew{⁠}\footnote{Belisarius left most of his cavalry behind; Justinian sent new forces; and Solomon seems to have disposed of about 18,000 men (Diehl, 67, note 4).

} had immediately

p143
to take the field against Cutsina and other Moorish leaders who descended upon Byzacena, while Iabdas was devastating Numidia. He defeated the former at Mamma, but not decisively; they returned with reinforcements, and were thoroughly beaten in the important battle of Mount Burgaon (early in A.D. 535).\texthebrew{⁠}\footnote{These localities have not been certainly identified.

} An expedition against the Numidian Moors in the following summer was unsuccessful, but Solomon lost no time in setting about the erection of fortified posts along the main roads in Numidia and Byzacena. In A.D. 536 the Emperor regarded peace as established and the Moors as conquered.\footnote{\emph{Nov.} XXX.11.2.

}

The task of keeping the natives in check had at least been well begun; but it was interrupted by a dangerous military revolt.

Various causes contributed to the mutiny. The pay of the soldiers had fallen into arrears, because the taxes from which it should have been defrayed had not been paid up. There was dissatisfaction about the division of booty. There were many Arians among the barbarian federates in the army who were ill-pleased at the intolerant religious policy which had been set in motion.\texthebrew{⁠}\footnote{Procopius says there were 1000, some of them Heruls; and they were instigated by the Arian clergy of the Vandals who had lost their churches and their incomes. \emph{B. V.} II.14.12\texthebrew{‑}13.

} Men who had married Vandal women claimed the lands which had belonged to their fathers or husbands and had been confiscated by the State. Above all, Solomon did not understand the art of tempering discipline by indulgence and was not a favourite with either officers or men. A conspiracy was formed to murder him at Easter ( A.D. 536). It miscarried because the courage of those who were chosen to do the deed failed them, and then a great number of the disaffected, fearing discovery, left Carthage and assembled in the plain of Bulla Regia. Those who were left behind soon threw off the pretence of innocence and the city was a scene of massacre and pillage. Solomon, having charged his lieutenants Theodore and Martin to do what they could in his absence, escaped by night, along with his assessor, the historian Procopius, and sailed for Sicily, to invoke the aid of Belisarius, who had just completed the conquest of the island. Belisarius did not lose a moment in setting sail for Carthage, in which he found Theodore beleaguered by the

p144
rebels. They were about 9000 strong\texthebrew{⁠}\footnote{Including about 1000 Vandals, of whom 400 had returned from the east. On the way from Constantinople to Syria \texthebrew{—} where they were to form part of the frontier forces \texthebrew{—} they succeeded in seizing a ship at Lesbos and landing in Africa.

} and under the command of Stotzas, who was one of the private retainers of Martin. The design of this upstart was to form an independent kingdom in Africa for himself.

Theodore was on the point of capitulating when Belisarius arrived, and on the news of his appearance the rebels hastily raised the siege and took the road for Numidia. It was a high compliment to the prestige of the conqueror of the Vandals. With the few troops who had remained loyal in Carthage, and a hundred picked men whom he had brought with him, Belisarius overtook Stotzas at Membressa\texthebrew{⁠}\footnote{Mejez el\texthebrew{‑}Bab, on the river Bagradas (Mejerda).

Thayer's Note: Hodgkin gives details of the battle

(\emph{Italy and Her Invaders}2, IV.33 ff.) .

} and defeated him. The rebels fled, but they did not submit. Belisarius could not remain: news from Sicily imperatively recalled him. He arranged that Solomon should withdraw from the scene, and that two officers, Theodore and Ildiger, should assume responsibility until the Emperor appointed Solomon's successor. Soon after his departure the situation became worse, for the troops stationed in Numidia, who had been moved to cut off the retreat of Stotzas, declared in his favour. Two-thirds of the army were now in rebellion.\footnote{This was found to be the case by Germanus, who investigated the military register on his arrival. Proc. \emph{B. V.} II.16.3.

}

Justinian was happily inspired at this grave crisis. He sent the right man to deal with it, his cousin Germanus, the patrician, who already had had experience of warfare on the Danube, as Master of Soldiers in Thrace. He was appointed Master of Soldiers of Africa, with extraordinary powers, and it was hoped that his prestige as a member of the Imperial family would have its influence in recalling the rebels to a sense of loyalty. His first act was to proclaim that he had come not to punish the mutineers, but to examine and rectify their grievances. This announcement was at once effective. Many of the soldiers left the camp of the rebels and reported themselves at Carthage. When it was known that they were handsomely treated and that they received arrears of pay even for the weeks during which they were in rebellion, large numbers deserted the cause of Stotzas, and Germanus found himself equal in strength to the

p145
insurgents. Stotzas, seeing that his only chance was to strike quickly, advanced on Carthage. A desperate battle was fought at Scalas Veteres (Cellas Vatari) in the spring ( A.D. 537), and the rebels were defeated. Moorish forces, under Iabdas and other chiefs, who had promised to support Germanus, were spectators of the combat, but according to their usual practice they took no part till the victory was decided, and then they joined in the pursuit, instead of falling on the exhausted victors.\footnote{Diehl, \emph{op. cit.} 87.

}

Germanus remained in Africa for two years and succeeded in re-establishing discipline in the army. Then the experienced Solomon was sent out to replace him A.D. 539) and to complete the military organisation of the provinces and the system of defence, in which Justinian took a keen personal interest. He began by weeding out of the army all those whom he suspected as doubtful or dangerous, sending them to Italy or the East, and he expelled from Africa the Vandal females who had done much to instigate the mutiny. After success ­ful campaigns against the Aurasian Moors, he established his power solidly in Numidia and Mauretania Sitifensis, and carried out the vast work of strengthening the defences of the towns and build hundreds of forts. Africa enjoyed a brief period of peace to which, amid subsequent troubles, the provincials looked back with regret.

The great pestilence which devastated the Empire in A.D. 542 and 543 visited Africa and took a large toll from the army. At the same time new troubles threatened from the Moors. The Emperor, who gratefully recognised the services and abilities of Solomon, appointed his nephew Sergius\texthebrew{⁠}\footnote{Sergius was not only nephew of Solomon; he was also son-in\texthebrew{‑}law of Antonina.

} duke of Tripolitana. It was a thoroughly bad appointment. Sergius was incompetent, arrogant, and debauched; he was not even a brave soldier; and he proved a governor of the well-known type who cannot avoid offending the natives. An insolent outrage committed against a deputation of the Louata provoked that people to arms; and by an unfortunate coincidence Solomon at the same time succeeded in offending the power ­ful chief Antalas, who had hitherto been friendly. The Moors joined forces, and in the battle of Cillium\texthebrew{⁠}\footnote{Kasrin, west of Sbeitla. Victor Tonn. \emph{sub} 543; Procopius, \emph{B. V.} II.21.

} ( A.D. 544) the Romans were utterly defeated and Solomon was slain.

p146
The Imperial rule in Africa was again in grave danger. The news of the defeat stirred the Berber tribes all along the frontier; even the Visigoths seized the occasion to send forces across the straits, and unsuccessfully besieged Septum.\texthebrew{⁠}\footnote{Isidore, \emph{Hist. Goth.} 42 , p284;

Procop. \emph{B. G.} II.30.15 .

} The Emperor made the fatal mistake of appointing Sergius, who was at once incapable and unpopular, as Solomon's successor. Stotzas, who since his defeat by the Germans had lived with a handful of followers in the wilds of Mauretania, now reappeared upon the scene and joined the Moors of Antalas, while Sergius quarrelled with his officers. Instead of superseding him, he despatched a second incompetent commander, the patrician Areobindus, who had married his own niece Praejecta. He made Areobindus co-ordinate with Sergius, but he was to command the army of Byzacena, Sergius that of Numidia. The two generals did not agree, and misfortune ensued. The Byzacene forces, relying on the support of Sergius, who left them in the lurch, were severely defeated at Thacia, between Sicca Veneria (el\texthebrew{‑}Kef) and Carthage (end of A.D. 545).\texthebrew{⁠}\footnote{Thacia has been identified with Borj-Messaudi. In this battle, John, son of Sisinniolus, one of the best officers in the army, fell.

} After this disaster Sergius was relieved of his post and Areobindus replaced him. He was a man of little merit, and in a few months he was removed by a conspiracy. Guntarith, the duke of Numidia, aspired to play the part of Stotzas, and having come to an understanding with some of the Moorish chiefs, he suddenly seized the palace at Carthage, and Areobindus was assassinated (March A.D. 546). º Praejecta fell into the hands of Guntarith, who formed the plan of marrying her. But Guntarith's supremacy lasted little over a month. A portion of the army remained loyal and found a leader in an Armenian officer, Artabanes, who brought about the murder of the rebel at a banquet (May). Justinian appointed Artabanes Master of Soldiers of Africa, and Praejecta offered her hand to her deliverer.\texthebrew{⁠}\footnote{See above,

p33 .

} But Artabanes was already married and Theodora refused to permit a divorce. He followed Praejecta to Constantinople, and the Emperor tried to console him by creating him Master of Soldiers in praesenti and Count of the Federates.

The situation was deplorable. The ravages of the Moors

p147
during the last three years had exhausted and depopulated the provinces. At last Justinian made a happy appointment. John Troglita, who had served with distinction under Belisarius and Solomon and was thoroughly acquainted with the conditions of the country, was recalled from the East, where he had given new proofs of military talent, and sent to take command of the armies of Africa (end of A.D. 546). Happily the Moors were divided, and John was a diplomatist as well as a general. He was able to secure the help of Moorish contingents in his campaigns. Early in A.D. 547 he inflicted a decisive defeat on the most dangerous of his opponents, Antalas.\texthebrew{⁠}\footnote{The scene of this battle is unknown; probably somewhere south of Sufetula, see Diehl, \emph{op. cit.} 370.

} But the troubles of Africa were not yet over. A few months later, the Berbers of Tripolitana rose under Carcasan, and won a crushing victory over the Imperial troops in the plain of Gallica.\texthebrew{⁠}\footnote{Now Maret, south-east of Gabes. \emph{Ib.} 374.

} Antalas took the field again and joined his triumphant neighbours. But the Roman cause was retrieved in the great battle of the Fields of Cato,\texthebrew{⁠}\footnote{Unknown locality in Byzacena. Corippus \emph{Johannis}, VIII.165. On this occasion many Moors, especially the faithful Cutsina, fought for the Romans.

} where seventeen Moorish leaders fell, among them Carcasan (early in A.D. 548). This victory secured for Africa complete tranquillity for nearly fourteen years. The relations between the Empire and the dependent Moorish princes were renewed and revised. The administration of the provinces was placed on a normal footing. The inhabitants and the wasted lands had time to recover from the devastations. The military defences of the frontier were re-established and improved.\texthebrew{⁠}\footnote{Cp. Diehl, 380.

} John Troglita, who seems to have governed Africa for about four years after his great victory, stands out, with Belisarius and Solomon, as the third hero of the Imperial reoccupation of Africa. His deeds inspired the African poet Corippus, whose \emph{Johannis} tells us nearly all we know of his campaigns.\footnote{The full narrative of Procopius, \emph{B. V.}, stops with the arrival of John. But he mentions briefly the three battles of A.D. 547\texthebrew{‑}548, and grimly concludes with words which sum up the terrible sufferings which the provinces had endured: ``Thus, at long last, the Libyans who survived, few in number and very poor, won some rest.''

}

Justinian was to have one more war in Africa, and it appears to have been entirely due to the stupid treachery of the military governor. The loyalty of the aged chief Cutsina was secured by an annual pension. In A.D. 563, when he came to Carthage to

p148
receive the money, he was assassinated by order of John Rogathinus, the Master of Soldiers.\texthebrew{⁠}\footnote{The source is John Mal. XVIII.495, transcribed and completed in Theophanes, A.M. 6055. John is called simply \textgreek{\texthebrew{ἄ}ρχων}; but this, as Diehl points out (456), certainly means the mag. mil., not the Praet. Pref., who at this time was probably Areobindus.

} The motive of the crime is unknown, but the sons of the murdered Moor immediately raised Numidia in revolt. The forces in the province were insufficient to cope with the insurrection, and the Emperor was compelled to send an army under his nephew Marcian, who succeeded, perhaps by diplomatic means, in re-establishing peace.

While Solomon was fighting with the Moors, he was at the same time engaged in carrying out a large scheme of defensive fortification to protect the African provinces against the incursions of the barbarians in the future; he was fortifying and rebuilding old towns and constructing new fortresses. The building of fortresses was one of the notable features of Justinian's policy. All the provinces exposed to foes in the East, in the Balkan peninsula, and in Africa were protected by forts, constructed on principles carefully thought out; but it is in Africa, where the soil is covered with their ruins, that the system of defence which was employed can best be studied. The numerous walls and citadels dating from the days of Solomon, which are still to be seen, are the best commentary on the principles and rules laid down in contemporary military handbooks.\footnote{The whole system of the African defences has been explained and illustrated at length by Diehl in \emph{L'Afrique byzantine}, to which I may refer the reader who is interested in the subject. He has written his admirable description with the work of the Anonymus Tacticus beside him, and refers throughout to its pages. Here I can only indicate briefly the general character of the defensive system. Details would be useless without illustrations.

}

Fortified towns, connected by a chain of small forts, formed the first frontier defence. Behind this there was a second barrier, larger towns with larger garrisons, which were all to afford a refuge to the inhabitants of the neighbourhood in case of an invasion. When the watchmen in the frontier stations discerned mena­cing movements of the tribes, they transmitted the alarm by the old system of fire signals by night or smoke signals by day,\texthebrew{⁠}\footnote{See Anon. Tact. VIII, and compare the notes of the editors, p315.

} so that the people of the villages might have time to find

p149
refuge in the walled towns and the garrisons of the inland places might be prepared.

In many cases the towns were entirely surrounded by walls, and in some had the additional defence of detached forts. In other cases they were open, and protected by the citadel. The neighbouring strongholds of Theveste, Thelepte, and Ammaedera on the frontier of Byzacena present good examples of the three types. The features of a fully fortified town were a wall with towers, an outer wall, and a fosse; the space between the two walls being large enough to accommodate the refugees who flocked in from the open country in a time of danger. But this scheme is not invariably found; sometimes there was no outer wall, sometimes there was no ditch. These variations depended upon local circumstances, as the form of the fortress depended on the nature of the ground. A rectangular shape was adopted when it was possible, but very irregular forms were sometimes required by the site. Theveste is a well-preserved example of the large fortress, rectangular, measuring about 350 by 305 yards, with three gates, and frontier towers; Thamugadi of the smaller castle (about 122 by 75 yards), with a tower at each corner and in the centre of each side. Small forts, like Lemsa, had a tower at each of the four angles.

From Capsa (Gafsa) in the Byzacene province to Sabi Justiniana and Thamalla in Mauretania Sitifensis the long line of fortresses can be traced round the north foothills of the Aurasian mountains. Thelepte, Theveste, with Ammaedera behind it to the north, Mascula and Bagai, Thamugadi, Lambaesis, Lambiridi, Cellae, and Tubunae\texthebrew{⁠}\footnote{These places are now known as Medinet el\texthebrew{‑}Kedima, Tébessa, Haïdra, Khenchela, Ksar Bagai, Timgad, Lambèse, Ouled Arif, Zerga, Tobna.

} were the principal advanced military stations, which were connected and flanked by small castles and redoubts. When invaders from the south had penetrated this line, the inhabitants might seek shelter in Sufes (Sbiba) and Chusira (Kessera) in Byzacena; in Laribus (Lorbeus), Sicca Veneria (Kef), Tubursicum Bue (Tebursuk), Thignica (Aïn Tunga) in the Proconsular Province; Madaura (Mdaurech), Tipasa (Tifech), Calama (Guelma), Tigisis (Aïn el\texthebrew{‑}Borj) in Numidia, to mention a few of the military posts in the interior.

The Mauretanian provinces were more lightly held. It is interesting to observe that Justinian took special care to

p150
strengthen by impregnable walls the fortress of Septum on the straits of Gades. This ultimate outpost of the Empire was to be a post of observation. He gave express directions that it should be entrusted to a loyal and judicious commander, who was to watch the straits, gather information as to political events in Spain and Gaul, and send reports to his superior the duke of Mauretania.\footnote{\emph{C. J.} I.27.2, 2 .

Procopius, \emph{Aed.} VI.7.14\texthebrew{‑}16 .

\texthebrew{◂} previous

next \texthebrew{▸}Images with borders lead to more information.

The thicker the border, the more information.

(Details here.)

UP TO:

J. B. Bury
Homepage

Latin \& Greek

Texts

LacusCurtius

Home}

\xchapter{Chapter XVIII, Part A}

Short URL for this page:

bit.ly/BURLAT18A

\xsubsection{(Part 1 of 5)}

The ecclesiastical reunion of Rome with the East, accomplished by Justinian and Pope Hormisdas, soon produced political effects. It would be rash to suppose that the idea of abolishing the Gothic viceroyalty and reasserting the immediate power of the Emperor in Italy had assumed a definite shape in the mind of Justinian in the early years of his uncle's reign. His own strong theological convictions may suffice to account for his policy. But the restoration of ecclesiastical unity was evidently the first step that would have been taken by a statesman who nursed the design of overthrowing the Gothic power. The existence of the schism did not indeed reconcile the Italian Catholics to the administration of the Goths, but it tended to render many of them less eager for a close political bond with Constantinople.

The death of Anastasius, with whom Theoderic never had been on terms of amity, was an important event for the Italian government. It can hardly be a coincidence that it was after Justin's succession that arrangements were made for the succession to the Ostrogothic throne. Theoderic had no male children. His daughter Amalasuntha had received a Roman education, and he had selected as her husband Eutharic, an Ostrogoth of royal lineage who was living obscurely in Spain.\footnote{He was descended from the famous king Hermanric (Jordanes gives the genealogy,

\emph{Get.} 81 ), and was discovered by Theoderic when he assumed the regency of Spain. His full name was Eutharic Cilliga (CIL VI.32003 ; Cassiodorus, \emph{Chron., sub} 518). Rome was surprised and delighted by the magnificent shows of wild beasts procured from Africa and the lavish largesses which signalised his assumption of the consul­ship in January 519 (\emph{ib., sub a.}). It was probably on the occasion of his consul­ship

(p152)
that Cassiodorus eulogised him in the Senate house (\emph{Var.} IX.25) in an oration of which a fragment is preserved (\emph{Paneg.} pp465 \emph{sqq.}, cp. p470).

}

p152
The marriage was celebrated in A.D. 515, and a son, Athalaric, was born three years later. This infant Theoderic destined to be his successor. It was the right of the Goths to choose their own king, but the choice could hardly be made without an understanding with the Emperor if the future king was to be also the Emperor's viceroy and Master of Soldiers in Italy. That Justin was consulted, and that he agreed to Theoderic's plan, seems to be clearly shown by the fact that Eutharic was nominated consul for A.D. 519. As Goths were strictly excluded from the consul ­ship, this could only be done by the personal motion of the Emperor, who thus signified his approbation of the settlement of the succession to the Italian throne.

When the reunion of the Churches was accomplished, Justin paid a marked compliment both to Theoderic and to the Senate by resigning the nomination of an eastern consul for A.D. 522 in order that the two sons of the distinguished Roman senator Boethius might fill the consul ­ship as colleagues.\texthebrew{⁠}\footnote{On this occasion Boethius pronounced a panegyric on the king.

} It seemed as if cordial relations between Ravenna and Constantinople might now be firmly established, yet within a year the situation became more difficult and dangerous than ever.

We have no precise information as to the views of Eutharic.\texthebrew{⁠}\footnote{See

Anon. Val.

(the writer hostile to Theoderic)

80 .

} It appears that he entertained strong national feelings and was devoted to his Arian faith; and he may have been somewhat impatient of the moderate policy of his father-in\texthebrew{‑}law and the compromises to which it led. We do not know whether he would have been prepared to denounce the capitulations and cut Italy off from the Empire as an independent Gothic state. But he was suspicious of the intentions of the Emperor and of the loyalty of the Roman Senate. He died in the course of A.D. 522, but he may have influenced the situation by propagating these suspicions in Gothic circles. And the suspicions seemed to be confirmed by the edicts which Justin issued against the Arians. The Goths connected these efforts for the extinction of Arianism with the reunion of the Church; they feared that the Imperial policy would provoke an anti-Arian movement

p153
in Italy; and the consequence was a growing mistrust of the Senate, and especially of these senators who had taken a prominent part in terminating the schism. Pope Hormisdas was trusted by Theoderic, but he died in August A.D. 523, and his successor, John I, was associated with those who desired a closer dependency of Italy on the Imperial government, as a means of attaining greater power and freedom for the Roman Senate.

It had been a token of Theoderic's goodwill when in autumn, A.D. 522, he appointed Boethius to the post of Master of Offices. Anicius Manlius Torquatus Severinus Boethius was a man of illustrious birth and ample fortune, whose life was dedicated to philosophy and science.\texthebrew{⁠}\footnote{He had been consul in 510. He constructed a sun-dial, a water-clock, and a celestial globe at Theoderic's request, to be sent as gifts to the Burgundian king. For his writings see below,

§ 11 .

} Translated from the society of his kinsmen and friends at Rome into the court circles of Ravenna, he did not find himself at home and could not make himself popular. His severe ethical standards repelled the pliant and opportune palatine officials who surrounded the king, and probably he was not very tactful.\texthebrew{⁠}\footnote{He seems to have opposed and prevented the appointment of a certain Decoratus to the Quaestorship, because he considered him as having mentem nequissimi scurrae delatorisque (\emph{De consol. Phil.} iii.4). Decoratus became Quaestor after the fall of Boethius. The epitaph in Rossi, \emph{Inscr. Chr.} ii p113, may be his.

} He had held office for about a year when a storm suddenly burst over his head.

An official seized letters which had been despatched by some Roman senators to the Emperor.\texthebrew{⁠}\footnote{Severus was the name of the official (Suidas, \emph{s.v.}\textgreek{Σεβ\texthebrew{ῆ}ρος}). The sources for the following events are Boethius, \emph{De cons. Phil.} i.4;
Anon. Val. 85\texthebrew{‑}87 ; \emph{Lib. pont., Vita Johannis}, pp275\texthebrew{‑}276; Procopius, \emph{B. G.} I.1. The questions relating to the legal procedure and the exact nature of the charges have been much discussed, most recently by Cessi in his introduction to Anon. Val., where the literature of the subject will be found, p. cxxv.

} In this correspondence passages occurred which could be interpreted as disloyal to the government of Theoderic,\texthebrew{⁠}\footnote{Adversus regnum regis. The passage in Suidas suggests that Albinus himself was not the writer of any of the letters; they were written by his friends. He had been consul in 493 and Praet. Pref. of Italy in 513. He seems to have belonged to the Decian family (cp. Sundwall, \emph{Abh.} p87).

} and the patrician Faustus Albinus junior was particularly compromised. The matter passed into the hands of Cyprian, a referendarius whose duty it was to prepare the case for the king's Consistorium, which was the

p154
tribunal for cases of treason.\texthebrew{⁠}\footnote{For the duties of the referendarii see

Cass. \emph{Var.} VI.17

(per eum nobis causarum ordines exponuntur). Boethius (\emph{loc. cit.}) speaks of Cyprian as a delator. Severus was the delator, not Cyprian, who only handled the delatio; and Cessi defends him as having simply performed his legal duty. But he may have shown a zeal and partiality in the prosecution which would explain the strong language of Boethius. Cp.

Anon. Val. 86 .

} It is important to note that Cyprian was a man of unusual parts, and enjoyed the confidence of Theoderic, whom he used often to accompany on his rides.\texthebrew{⁠}\footnote{See the panegyrics of Cyprian's qualities in the letters conferring on him the office of Comes s. larg. and recommending him to the Senate of which this appointment made him a member, A.D. 524. Cassiodorus, \emph{Var.} V.40 and 41. He knew Gothic as well as Greek and Latin (trifarii linguis).

} The intercepted letters of the friends of Albinus justified an investigation. Boethius was a member of the Consistory \emph{ex officio}, and he spoke in defence of Albinus.\texthebrew{⁠}\footnote{I cannot agree with Cessi (p. cxlix) that it can be inferred from Cass. \emph{ib.} VI.6.2 that it was a duty of the Master of Offices to defend accused senators. But any member of the Consistory could express his opinion on a case.

} It was impossible to deny the material facts, and Boethius took the line that Albinus was acting not in his private capacity but as a senator, and therefore was not alone responsible for his act. ``The whole Senate, including myself, is responsible; there can be no action against Albinus as an individual.'' This defence was construed as a confession, and made the ground of a charge of treason against Boethius himself, and three men who belonged to ministerial circles but were under a cloud came forward to support the charge. He was arrested, and, as a matter of course, deprived of his office. Cassiodorus was appointed in his stead, and it may be ascribed to his influence that no attempt was made to involve other members of the Senate in the crime.\footnote{Cp. Sundwall, \emph{Abh.} p246. Sundwall thinks that the guilt of Boethius lay not in his defence of Albinus, but in trying to suppress the accusation (p243). Cyprian was the subordinate of Boethius, and Boethius appears, on grounds of procedure, to have raised objections to the denunciation of Severus being received (delatorem ne documenta deferret quibus senatum maiestatis reum faceret impedire criminamur (\emph{De cons. Phil.} i.4)). But surely this was only an incidental point, not the serious charge. The three witnesses were Basilius, Opilio, and Gaudentius, of whom Opilio was Cyprian's brother and a relative, perhaps son-in\texthebrew{‑}law, of Basilius. Boethius says that Basilius was in debt, and the other two had been condemned to exile ob multiplices fraudes. A year or two after Theoderic's death Opilio was appointed com. sacr. larg. and eulogised by Cassiodorus (\emph{Var.} VIII.16 and 17) in the usual way.

}

Up to this point there is no reason for thinking that there was anything illegal in the procedure; but now, instead of completing the process of Albinus and trying Boethius before the

p155
Consistory, the matter was taken out of the hands of that body, and the two men were thrown into prison at

Ticinum

(late autumn, A.D. 523). Thither the Prefect of Rome was summoned, and with him the king proceeded with the investigation of the case.\texthebrew{⁠}\footnote{Boethius speaks of a forged letter which was used against him, \emph{loc. cit.} In ordinary criminal trials of senators the tribunal consisted of the Prefect and five senators (Mommsen, \emph{Strafrecht}, 287), and this procedure may have been adopted as Sundwall suggests (\emph{op.cit.}, 248), although, as it was a case of treason, the proper tribunal was the Consistorium. In any case, I am sure that the ordinary view that the Senate tried the case and sentenced Boethius is mistaken (so Cessi, cxlviii).

} Boethius was found guilty and condemned to death. Albinus drops out of the story, his fate is not recorded. Theoderic was determined to teach the Senate a lesson, but perhaps he thought it better to let the course of political events guide him to an ultimate decision as to the fate of the distinguished philosopher. In his dungeon\texthebrew{⁠}\footnote{Probably at Ticinum. But Boethius himself had been transferred to Calvenzano
(Anon. Val. 87)

near Melegnano, \texthebrew{•} about seven miles south of Milan.



14a
14b
We do not know when the sentence was passed. Nine months or more seem to have elapsed between his arrest and execution. October 23 was the date accepted in ecclesiastical tradition for his death, but this tradition only emerges three centuries later and has been questioned (cp. the article on Boethius in \emph{D. Chr. B.}). See Pfeilschifter, \emph{Theod. der Grosse}, p164.

} Boethius composed his famous book on the \emph{Consolation of Philosophy}, and probably expected that his sentence would be mitigated. But he was put to death (in the late summer or autumn of A.D. 524),\texthebrew{⁠} and, it was said, in a cruel manner. A cord was tightened round his head, and he was despatched with a club.\texthebrew{⁠}While Boethius was awaiting his trial, the senators met and debated. They were thoroughly alarmed, and passed decrees designed to exculpate themselves, and therefore repudiating Boethius and Albinus. The only man perhaps who stood by Boethius was his father-in\texthebrew{‑}law Symmachus, the head of the Senate. He may have used strong language; he declined at least to associate himself with the subservient decrees. Thereby he laid himself open to the charge that he defended treason and sympathised with traitors. He was arrested, taken to Ravenna, and executed. It was a foolish act, the precaution of a tyrant.\footnote{The passage in Boethius, i.4, an optasse illius ordinis salutem nefas vocabo? ille quidem suis de me decretis uti hoc nefas esset efficeret, is important. In my opinion it supplies the key to the arrest of Symmachus, which historians have not explained. It is evident that he did not subscribe to the decree repudiating his son-in\texthebrew{‑}law, for whose wrongs he was mourning

(p156)
(Boethius, \emph{op.cit.}, ii.4), and he can hardly have failed to speak in his defence. This attitude in the Senate furnished Theoderic with an excuse for arresting a man whom he had reason to fear as a near and dear relative of Boethius. The sources say nothing of a trial, but it seems unlikely that this formality was dispensed with. The date of the death of Symmachus fell in 525 (Marius Avent., \emph{sub a.}), probably in the first half. I cannot agree with the transposition in the text of Anon. Val. proposed by Cessi (\emph{op.cit.}, cxxvii), which, by removing § 88\texthebrew{‑}91 so as to follow § 84, would make the notice of the death of Symmachus, § 92, follow immediately that of the death of Boethius.

}

p156
It is probable that these events had some connexion with an Imperial edict which was issued about this time, threatening Arians with severe penalties, excluding them from public offices and from service in the army, and closing their churches. Theoderic was alarmed. He resented the revival of pains and penalties against his fellow-religionists in the East, and he saw in the edict an encouragement to the Italians to turn against their Arian fellow-subjects . But the edict is not preserved, and we do not know the exact date of its promulgation; so that we cannot decide whether it influenced Theoderic's policy before the execution of Boethius. It may not have been issued till after his death.\texthebrew{⁠}\footnote{We possess the edict against heretics of Justin and Justinian (the beginning of it is lost) issued in 527 ( \emph{C. J.} I.5.12 ; we are able to date it through the reference in

I.5.18, §4 ), which contains the exception in favour of the Goths, made after the negotiations with Pope John. The date of 523\texthebrew{‑}524 for a measure against the Arians depends on Theophanes, A.M. 6016, where the mission of Pope John is misdated. I cannot agree with Pfeilschifter (\emph{op.cit.}, 168) that

\emph{C. J.} I.5.12

was issued in 523, with its reserve in favour of the Goths, and that, notwithstanding the reserve, severe measures were taken in winter of 523\texthebrew{‑}524 as a reprisal for the proceedings against Albinus and Boethius. It seems more probable that there was special legislation against the Arians in 524, provoked perhaps by the wealth of the Arian churches (\textgreek{πλο\texthebrew{ῦ}τόν τινα \texthebrew{ἀ}κο\texthebrew{ῆ}ς κρείττω,}Procopius, \emph{H. A.} 11.16 ), and that the persecution began without any reference to Italian politics. It must be emphasised that in the prosecution of Boethius there was no anti-Catholic tendency; his opponents (Cyprian, etc.) were Catholics.

} We can only say that severe measures against the Arians had been adopted, and reported in Italy, before the autumn of A.D. 525. Theoderic determined to bring matters to an issue at Constantinople by coming forward as the protector of his fellow-heretics in the East. He selected as his ambassador John, the bishop of Rome, who was induced to undertake the distasteful commission of urging the Emperor to relax his policy and of conveying to him the royal threat that, if he persisted, reprisals on Italian Catholics would be the consequence. The Pope set forth, accompanied by several bishops and prominent senators, some time between the beginning of September and

p157
the end of November, A.D. 525.\texthebrew{⁠}\footnote{The determination of this date, which is due to Pfeilschifter, depends (1) on the fact that John was in Constantinople at Christmas (so that he could not have started later than the end of November); this is known from the statement of a contemporary priest Procopius, who, when John arrived, was translating a Latin work into Greek (see

note to Bonn ed. of \emph{Chron. Pasch.} II.120) ; (2) on a letter of Boniface, primicerius notariorum, addressed to the Pope before he left Italy in the 4th indiction which began on September 1 (so that he could not have started before September),

Pitra, \emph{Analecta novissima}, I.466 , where the addressee is wrongly said to be John IV. Of the senators who were with John, Importunus and Theodorus, members of the great Decian family, and Agapetus were ex-consuls.

} He was received in the eastern capital with an honourable welcome, and remained there at least five months. He celebrated Christmas and Easter in St. Sophia, and successfully vindicated his right to sit on a higher throne than the Patriarch's. It is recorded, and perhaps we have no right to question the statement, that Justin, though long since duly crowned, caused the Pope to crown him again.\texthebrew{⁠}\footnote{This is mentioned only in \emph{Lib. pont.}

} The mission succeeded in its principal object. The Emperor agreed to restore their churches to the Arians and permit them to hold their services. He refused to allow converted Arians to return to their old faith, but the main demand of Theoderic was conceded. Yet when the Pope and his companions returned to Ravenna in the middle of May\texthebrew{⁠}\footnote{They must have started immediately after Easter day, which fell in 526 on April 19.

} their reception was the reverse of that which success ­ful envoys might expect. They were arrested and thrown into prison.\texthebrew{⁠}\footnote{\emph{Lib. pont.}, in custodia eos omnes adflictos cremauit (tortured).
Anon. Val. 93

says only in offensa sua eum esse iubet, which Cessi interprets (p. clix) as pointing to a rigorous surveillance.

} John, who had been ailing when he started for the East, died a few days later (May 18, A.D. 526);\texthebrew{⁠}\footnote{\emph{Lib. Pont.} XV kal. Iun.

} his body was taken to Rome and interred in St. Peter's; there was a popular demonstration at his funeral and he was regarded as a martyr.

There was a contested election for the succession to the vacant see. It was probably a contest of strength between the Italians who were friendly to the Ostrogothic regime and those who were not. The former succeeded in securing the victory of their candidate after a struggle of two months, and the election of Felix IV (July 12) was a satisfaction to Theoderic, who had expressly signalised his wishes in the matter to the members of the Senate.\footnote{This is clear from the letter of Athalaric to the Senate commending that body for having accepted his grandfather's choice. Cass. \emph{Var.} VIII.15.

}

p158
But the days of Theoderic were numbered. Seven weeks later he was seized by dysentery, and died on August 30. Before his death he called together the Goths of his entourage and, presenting to them his grandson Athalaric as their future king, enjoined upon them to keep on good terms with the Senate and the Roman people, and always to show the becoming respect to the Emperor.\texthebrew{⁠}\footnote{Jordanes, \emph{Get.} 304 . As Justin had agreed to the concession which he demanded, and as he had secured the man he desired as head of the Catholic Church, it is perfectly incredible that four days before his death Theoderic should have drawn up a decree empowering the Arians to take possession of Catholic churches, as
Anon. Val. 94

asserts. The statement, which stands alone, his generally been accepted, but Pfeilschifter is assuredly right in rejecting it.

} Popular legend did not fail to connect his end with his recent acts of tyranny. It was said that a huge fish had been served at the royal table, and that to the king's imagination, tortured by conscience, its head, with long teeth and wild eyes, assumed the appearance of Symmachus. Theoderic took to his bed in terror, and declared to his physician his remorse for the slaughter of the illustrious senators.\footnote{Procopius, \emph{ib.};

Jordanes, \emph{Get.} 69 .

}

During the last year of his life he had been distressed by the fate of his sister Amalafrida, the widow of king Trasamund.\texthebrew{⁠}\footnote{In his last years it may be noted that he had experienced a succession of losses, which he must have keenly felt, by the deaths of Ennodius, a trusted friend (in 521), his son-in\texthebrew{‑}law Eutharic (523), his grandson Sigeric of Burgundy (522; see above,

chap. xiii § 5 ); and the death of Pope Hormisdas, with whom he was always on cordial terms, was also a blow.

} She had remained in Africa after her husband's death, and was probably useful to her brother in maintaining the good relations between the courts of Ravenna and Carthage which her marriage had inaugurated. But as king Hilderic leaned more and more towards Constantinople, and fell under the influence of Justinian, he drew away from the Goths, and his friendship with Theoderic cooled. Amalafrida, who had her own Gothic entourage in her adopted country, was accused, rightly or wrongly, of conspiring against the king, and was thrown into prison, where she died, from natural causes it was given out, but it was suspected that her death was violent. All her Goths were killed. Theoderic, if he had lived, would doubtless have attempted to wreak vengeance on Hilderic. After his death his daughter was not in a position to do more than address to the king of the Vandals a strong remonstrance.\footnote{The letter was written by Cassiodorus and is preserved (\emph{Var.} IX.1) . The other sources are Procopius, \emph{B. V.} III.9.3\texthebrew{‑}5, and Victor Tonn.

(p159)
\emph{sub} 523 (the year of Trasamund's death). This date has been reasonably questioned by Schmidt (\emph{Gesch. der Wand.} 122). For the letter of Cassiodorus strongly suggests that the queen's death was quite recent. I should be inclined to date it early in 526, and to connect with it (as Schmidt does) Theoderic's urgency in completing his naval armament, and collecting it at Classis on June 13. See \emph{Var.} V.17.3 non habet quod nobis Graecus imputet aut Afer insultet (also letters 16, 18, and 19).

}

Theoderic was succeeded by a child, his grandson Athalaric, whom his daughter Amalasuntha had borne to Eutharic, and Amalasuntha held the reins of government as regent during her son's minority. She had received a Roman education at Ravenna; she was brave and intelligent,\texthebrew{⁠}\footnote{Compare the laudatory remarks of Procopius, \emph{B. G.} I.2 p10; 4 p24. There is a miniature representation of Amalasuntha on the consular diptych of Orestes (A.D. 530) which is preserved at South Kensington. She wears earrings and pendants from her head-dress (\textgreek{πρεπενδούλια}, like those worn by Roman Empresses of the time), but no diadem.

} and perhaps sincerely believed in the ideal of blending the Italians and Goths into a united nation. Even if her convictions and sentiments had been different, the inherent weakness of a regency would have forced her to follow her father's last advice, to keep on good terms with the Emperor and to conciliate the Senate. The restoration of the confiscated properties of Boethius and Symmachus to their children was a pledge of the change. The Roman people was assured that no difference would be made in the treatment of Romans and Goths,\texthebrew{⁠}\footnote{Cassiodorus, \emph{Var.} VIII.3.

} and when the Senate and people swore loyalty to the young king, he also took an oath of good government to them. The Senate was invited to express its demands and desires.\texthebrew{⁠}\footnote{\emph{Ibid.} 2. A letter was also addressed to the Catholic clergy, requesting their prayers for Athalaric (\emph{ib.} 8).

} Ambassadors were sent to the Emperor bearing a letter\texthebrew{⁠}\footnote{\emph{Ibid.} 1. This letter is addressed in the MSS. to Justinian. There is a reference in the text to the Emperor's old age, and it is simpler to adopt Mommsen's correction Iustino than the explanation attempted by Martroye (p158\texthebrew{‑}159). It is to be noted that the Gothic general Tuluin (who had commanded in all Theoderic's campaigns since 504) was created a Patrician soon after Theoderic's death, and thus became a member of the Senate (Cass. \emph{Var.} VIII.9 and 10). This act must have been sanctioned by Justin.

} in which he was requested to aid the youth of Athalaric, and it was suggested that the tomb should be allowed to bury old hatreds: Claudantur odia cum sepultis.

Amalasuntha determined to give her son the education of Roman princes, and she confided him to the care of three civilised

p160
Goths, who shared her own views. But the Goths, as a whole, had no comprehension of the ideal of Italian civilisation at which she, like her father, aimed; they believed only in the art of war; and they regarded themselves as victors living in the midst of a vanquished population. It outraged their barbarian sentiments that their king should receive an education in the humanities. Their indignation was aroused when Athalaric, chastised by his mother for some fault, was found in tears. They whispered that the queen wished to do away with her son and marry again. Some of the leaders of this faction then sought an audience of Amalasuntha, and protested against the system of training which she had chosen for the king. A literary education, they urged, promotes effeminacy and cowardice; children who fear the whip cannot face the sword and spear; look at Theoderic, who had no idea of letters; let Athalaric be brought up in manly exercises with companions of his own age. Amalasuntha feigned to be persuaded by arguments with which she profoundly disagreed. She feared that, if she refused, she would be deposed from the regency, for there were but few among the Goths who sympathised with her ideas and policy. Athalaric was released from the discipline of pedagogues, but even the enemies of a liberal education would hardly have contended that the new system was a success. He was of a weak and degenerate nature, and the Gothic youths with whom he associated soon led him into precocious debauchery which ruined his health.

As time went on, the dissatisfaction of the Goths with the rule of Amalasuntha increased, and she became aware that a plot was on foot to overthrow her. She sent three of the most dangerous men who were engaged in the agitation against her to different places on the northern frontier, on the pretext of military duty. Finding that they still carried on their intrigues, she decided on stronger measures. Fully estimating the hazards of her position, she took the precaution of providing herself with a retreat. She wrote to Justinian, asking if he would receive her in case of need. The Emperor, who probably did not view with dissatisfaction the situation in Italy, cordially agreed, and prepared a mansion at Dyrrhachium for the queen's reception on her journey to Constantinople. Thus secured, Amalasuntha proceeded to the commission of murders, which it is common to palliate or justify by the plea of political necessity.

p161
She sent some devoted Goths to assassinate the three arch-conspirators . She stowed 40,000 gold pieces in a vessel, which she sent to Dyrrhachium, directing that it should not be unloaded before her arrival. When she learned that the murders had been duly accomplished she recalled the ship and remained at Ravenna.

It is important to realise that the Ostrogothic kingdom was now politically isolated. The system of friendly understandings, cemented by family alliances, which Theoderic had laboured to build up among the western Teutonic powers was at the best a weak guarantee of peace; after his death it completely broke down. We have seen how the alliance with the Vandals was ruptured, and how Amalafrida, Theoderic's sister, was put to death by Hilderic,\texthebrew{⁠}\footnote{See above,

p129 .

} an injury which Amalasuntha was not in a position to avenge. The Thuringians, whose queen was her cousin, were attacked and conquered by the Franks.\texthebrew{⁠}\footnote{In 529 and following years. See below,

Chap. XVIII § 13 .

} The Franks were also intent on driving the Visigoths from the corner of Gaul which they still retained; the young king Amalaric, the grandson of Theoderic, was killed ( A.D. 531), and Theudis, who succeeded him, had enough to do to maintain the possession of Septimania. From that quarter the Ostrogoths could look for no support. The power of the Franks became more formidable by their conquest of Burgundy ( A.D. 532\texthebrew{‑}534),\texthebrew{⁠}\footnote{Between 529 and 532 Amalasuntha restored to Burgundy districts between the Isère and the Durance, which had been taken by Theoderic in 523. Cass. \emph{Var.} XI.1.13; cp. VIII.10.8.

} and there was also the danger that the Ostrogothic provinces in Gaul might be attacked by their insatiable ambition. Thus the Italian regency would have been forced, even if there had been no internal difficulties, to conduct itself demurely and respectfully towards the Imperial power to which constitutionally it owed allegiance.

Amalasuntha had one near relative in Italy, her cousin Theodahad, the son of Amalafrida, queen of the Vandals, by a first marriage. He was the last person to whom she could look for help in her difficulties. Theodahad had none of the soldierly instincts of his race. He had enjoyed a liberal education and was devoted to the study of the philosophy of Plato. But he was far from being free from the passions which philosophy condemns. The ruling trait of his character was cupidity.

p162
He had estates in Tuscany, and by encroachments on the properties of his neighbours he had gradually acquired a great part of that province.\texthebrew{⁠}\footnote{Theoderic had on more than one occasion to check and reprove his nephew's cupidity. Cass. \emph{Var.} IV.39; v.12.

} ``He considered it a misfortune to have a neighbour.'' The Tuscans had complained of his rapacity, and Amalasuntha had forced him to make some restitutions, earning his undying hatred. He was not, however, naturally ambitious of power. His ideal was to spend the last years of his life in the luxury and society of Constantinople. When he first appears on the stage of history he takes a step to realise this desire. Two eastern bishops had come to Rome on business connected with theological doctrine.\texthebrew{⁠}\footnote{Hypatius of Ephesus and Demetrius of Philippi. Cp. Liberatus, \emph{Brev.} xx. They bore a letter from the Emperor to the Pope, which, with the Pope's reply, is preserved in

\emph{C. J.} I.1.8 . The former is dated June 6, 533, the latter March 25, 534. Thus the negotiations with Amalasuntha and Theodahad are dated to 533\texthebrew{‑}534, and we can infer that Alexander and the bishops left Italy at the end of March 534.

} Theodahad entrusted them with a message to Justinian, proposing to hand over to him his Tuscan estates in return for a large sum of money, the rank of senator, and permission to live at Constantinople.

Along with these two bishops, Alexander, an Imperial agent, had arrived in Italy. His ostensible business was to present to the regent some trifling complaints of unfriendly conduct.\texthebrew{⁠}\footnote{The three complaints were that she retained the fortress of Lilybaeum, which belonged to the Emperor; that she gave refuge to ten Huns who had deserted from the Imperial army; and that in a campaign against the Gepids in the neighbourhood of Sirmium, the Goths had committed hostile acts against Gratiana, a town in Moesia. The claim of the Emperor to Lilybaeum was founded on the circumstance that Theoderic gave it to his sister when she married Hilderic (see above,

Vol. I p461 ). Belisarius had demanded its surrender at the end of the Vandalic war.

} At a public audience, Amalasuntha replied to the charges, dwelt on their triviality, and alleged her services to the Emperor in allowing his fleet to make use of Sicily in the expedition against the Vandals. But this performance was only intended to deceive the Goths. Justinian had followed closely events in Italy, and the real purpose of Alexander's visit was to conclude a secret arrangement with the regent. Her position was now more critical than ever. The premature indulgences of Athalaric had brought on a decline, and he was not expected to live. On his death her position, unpopular as she was with the Goths, would hardly be tenable, and she thought of resigning her power into the

p163
hands of the Emperor.\texthebrew{⁠}\footnote{If Amalasuntha seriously contemplated this step, it was, though defensible theoretically on constitutional grounds, an act of gross treachery towards her own people.

} She communicated her intention to Alexander, who then returned to Constantinople with the bishops. On receiving the messages of Amalasuntha and of Theodahad, Justinian sent a new agent to Italy, Peter of Thessalonica, an able and persuasive diplomatist.

Meanwhile Athalaric died.\texthebrew{⁠}\footnote{On October 2, 534, according to \emph{Con. Ital.} (Agnellus), in Chron. min. i.333. This accords with the statements of Procopius (\emph{B. G.} I.4) and Jordanes (\emph{Rom.} 367) that Athalaric reigned eight years. For the coins of his reign see Wroth, \emph{Coins of the Vandals}, p. xxxiii. An inscription records that he constructed seats in the amphitheatre of Ticinum in 528\texthebrew{‑}529 (d. n. Atalaricus rex gloriosissimus, CIL V.6418 ).

} But now that the critical moment had come, Amalasuntha, who enjoyed power, could not bear to part with it, and she committed a fatal blunder. She sent for her cousin Theodahad, assured him that, in attempting to curb his rapacity, her intention had been to prevent him from making himself unpopular, and offered him the title of king, on condition that she should retain in her own hands the exercise of government. Dissembling the bitter animosity which he felt towards her and of which she can have had little conception, he consented to her terms, and took a solemn oath to fulfil all she demanded. As soon as he was proclaimed king,\texthebrew{⁠}\footnote{\emph{Cons. Ital. loc. cit.}: alia die, which seems to mean October 3, the day after Athalaric's death.

} formal letters were addressed to the Senate, in which Amalasuntha dwelled upon Theodahad's literary tastes, and Theodahad enlarged on Amalasuntha's wisdom, professing his resolve to imitate her.\texthebrew{⁠}\footnote{Cassiodorus, \emph{Var.} X.3 and 4. Cassiodorus had been appointed Praet. Prefect before the end of Athalaric's reign. Cp. \emph{Var.} XI.7.

} Letters were also despatched to Justinian, informing him of what had happened.\footnote{\emph{Ib.} 1 and 2.

}

But after the first hypocritical formalities, Theodahad lost little time in throwing off the mask. He gathered together the relatives of the three Goths who had been murdered by Amalasuntha's orders; the Gothic notables who were faithful to her were slain, and she was herself seized and imprisoned in an island in Lake Bolsena in Tuscany, which probably belonged to the king.\texthebrew{⁠}\footnote{There are two islands in the lake, Bisentina and Martana; the latter is supposed to be the scene of the tragedy.

} She was then forced to write a letter to Justinian, assuring him that she had suffered no wrong. Theodahad wrote himself

p164
to the same effect, and committed the letters to two senators \texthebrew{—} Liberius, the Praetorian Prefect of Gaul, and Opilio \texthebrew{—} to bear to Constantinople.

In the meantime, Peter, the new agent whom the Emperor had selected to continue the secret negotiations, had started. Travelling by the Egnatian Road, Peter met on his way the Goths who bore the news of Athalaric's death and Theodahad's elevation to the throne; and on reaching the port of Aulon (Valona), he met Liberius and Opilio, who informed him of the queen's captivity. Peter sent a fast messenger to Constantinople and awaited further orders.\texthebrew{⁠}\footnote{That Peter's messengers outstripped the Italian ambassadors is clear from the narrative in Procopius. They arrived after Justinian had forwarded his instructions. Liberius told the whole truth, but Opilio sought to defend Theodahad. Historians have not observed that there is an interesting notice of the Emperor's reception of Liberius in Constantine Porph. \emph{De cerim.} i.87 (taken doubtless from a work of this same Peter): \textgreek{Λίβερ \texthebrew{ὁ} πατρίκιος κα\texthebrew{ὶ ἔ}παρχος Γαλλι\texthebrew{ῶ}ν \texthebrew{ἐ}πεμφθη \texthebrew{ἐ}ντα\texthebrew{ῦ}θα παρ\texthebrew{ὰ} Θευδ\texthebrew{ᾶ} το\texthebrew{ῦ ῥ}ηγ\texthebrew{ὸ}ς Γότθων κα\texthebrew{ὶ} τ\texthebrew{ῆ}ς συγκλήτου \texthebrew{Ῥ}ωμαίων}. Justinian accorded him the same honours as were due to a Praet. Pref. of the East. He afterwards passed into Justinian's service, and was Augustal Prefect of Egypt. We shall meet him again in Sicily and Spain.

} Justinian immediately wrote a letter to Amalasuntha, assuring her of his protection, and instructed Peter to make it clear to Theodahad and the Goths that he was prepared to support the queen. But the Emperor's authority and his envoy's representations did not avail to save Amalasuntha.\texthebrew{⁠}\footnote{Jordanes, \emph{Get.} 306 .

} She was killed \texthebrew{—} strangled, it was said, in a bath \texthebrew{—} in the lonely island by the relatives of the Goths whom she had slain, and who had persuaded Theodahad that her death was necessary to his own safety. Goths and Romans were alike shocked by the fate of Theoderic's daughter, whose private virtues were acknowledged by all. Peter told Theodahad, in the name of Justinian, that the criticism which had been perpetrated meant ``a war without truce.'' The king pleaded that it had been committed against his will, but he continued to hold the assassins in honour.\footnote{The chronology is beset by serious difficulties. The fact that Peter met the envoys of Amalasuntha at Thessalonica and those of Theodahad at Aulon, shows that her captivity began very soon after Theodahad's accession (we cannot infer anything from the vague post aliquantum tempus of

Jordanes, \emph{Get.} 306 ). The only evidence for the date of her murder is \emph{Cons. Ital., loc. cit.}, where she is said to have been imprisoned in Lake Bolsena on April 30, 535. It seems probable that this may really be the date of her murder; it cannot be that of her imprisonment, for she was already confined \emph{in the island} when Theodahad sent his envoys to Justinian (\emph{B. G.} I.4). In this view, I think, Leuthold (\emph{Untersuchungen zur ostgotischen Gesch.} 26) is right. Peter, who seems to have left Constantinople

(p165)
in October, and was detained for some time at Aulon to receive instructions, must have reached Italy before the end of the year, or at latest in January. Here the difficulty arises. Procopius says (\emph{ib.}) \textgreek{Πέτρου δ\texthebrew{ὲ ἀ}φικομένου \texthebrew{ἐ}ς \texthebrew{Ἰ}ταλίαν \texthebrew{Ἀ}μαλασούνθ\texthebrew{ῃ} ξυνέβη \texthebrew{ἐ}ξ \texthebrew{ἀ}νθρώπων \texthebrew{ἀ}φανισθ\texthebrew{ῆ}ναι}. This has been generally interpreted to mean, Peter on his arrival in Italy found that Amalasuntha had already been done away with. We should thus be faced with three alternatives: (1) the rejection of the date April 30; (2) the extremely unlikely hypothesis that Peter remained at Aulon till May; (3) the hypothesis, put forward by Leuthold (\emph{ib.} 24), that Procopius deliberately falsified facts. But the words of Procopius admit a different interpretation. In fact, they properly mean: After the arrival of Peter in Italy, it occurred that Amalasuntha was done away with (\textgreek{\texthebrew{ἀ}φανισθ\texthebrew{ῆ}ναι} is aorist, not pluperfect). My own view is that Procopius designedly made his statement ambiguous. He was treading on delicate ground, and he was afraid to force on the reader's attention the fact that Peter was some time (about four months) in Italy and was unable (or unwilling) to save the queen's life.

}

p165
This brief story of Amalasuntha's tragic end, told by Procopius in his \emph{History of the Wars}, raises some perplexing questions, which might compel us, even if we had no other evidence, to suspect the presence of unexplained circumstances in the background. It is difficult to understand Theodahad's motive in permitting the murder, knowing, as well he knew, that such an act would cause the highest displeasure to Justinian and might lead to war, which, as his subsequent policy shows, he desired, almost at any cost, to avoid. Peter was in Italy at the time, and had been there for some months before the event. He had been instructed by the Emperor to champion the cause of Amalasuntha. How was it that he was not only unable to restore her to liberty but could not even save her life? When we find that Procopius is silent as to any efforts of the ambassador in the queen's behalf, and even, by an ambiguous sentence, allows his readers to believe that Peter arrived too late to interfere, there is ground for suspecting that the tale is only half told.

An explanation is forthcoming from the pen of Procopius himself. In his \emph{Secret History} he had added a sinister supplement, which, he says, ``it was impossible for me to publish through fear of the Empress.''\texthebrew{⁠}\footnote{\emph{H.A.} 16 \emph{ad init.} Cp. above,

p34 .

} According to this story, Theodora viewed with alarm the prospect of Amalasuntha seeking refuge at Constantinople. She feared that this handsome and strong-minded woman might gain an influence over the Emperor, and she suborned Peter, by promises of money and office, to procure the death of the queen of the Goths. ``On arriving in Italy, Peter persuaded Theodahad to despatch Amalasuntha. And

p166
in consequence of this he was promoted to the dignity of Master of Offices, and won great power and general detestation.''\texthebrew{⁠}\footnote{\textgreek{Μάλιστα πάντων \texthebrew{ἒ}χθους}, Haury's correction of \textgreek{\texthebrew{ἐ}χθρ\texthebrew{ῶ}ν.}} The credibility of this story has been doubted,\texthebrew{⁠}\footnote{Gibbon (c. xli) accepts it in his text, but throws some doubt about it in his note. Hodgkin rejects it as ``a malicious after-thought of the revenge­ful old age of Procopius'' (iii.720) , but does not discuss the evidence. Diehl (\emph{Justinien}, 181) observes: ``La chose semble bien douteuse, quoique Théodora entretînt à ce moment même avec Théodat et sa femme une assez mystérieuse correspondance et que Pierre lui fût tout dévoué.'' Dahn thought that the story may perhaps be a pure invention (\emph{Prokopius}, 379).

} but the evidence in its favour is considerably stronger than has been realised.

It may be observed, in the first place, that it supplies an adequate explanation of the conduct of Theodahad in consenting to the crime. Relying on the influence of the power ­ful Empress, he might feel himself safe in complying with the wishes of the Gothic enemies of Amalasuntha and ignoring the Emperor's threats. And, in the second place, there is nothing incredible in Theodora's complicity. There is nothing in her record to make us suppose that she was incapable of such a crime, and the motive was surely sufficient. It must be remembered that, on the scene of public affairs, Amalasuntha was, next to Theodora herself, the most remarkable living woman. She possessed advantages of person and education, which report might magnify, and in her eight years of government she had shown strength of mind and even unscrupulousness. But if in these respects she might compete with the Empress, her unblemished private character and her royal birth were advantages which Theodora could perhaps be hardly expected to forgive. Whatever be the truth about Theodora's early career, her origin was of the lowest, and report, rightly or wrongly, was busy with the licentiousness of her youth. We can well understand that Theodora would have been ready to go far in order to prevent the arrival at Constantinople of a king's daughter who might gain an influence over the Emperor and would in any case inevitably challenge comparisons unfavourable to herself.

The statement of Procopius respecting Theodora's part in the drama must be admitted to be perfectly credible, but, in the absence of corroborative evidence, it would be open to us to dismiss it as the specious invention of malice. We have,

p167
however, independent evidence which corroborates Procopius in one important particular. It is an essential point in his story that Peter was the devoted agent of Theodora, and that she procured his appointment as ambassador to Ravenna. This is fully borne out by letters which Theodahad addressed to the Empress, when Peter returned to Constantinople after the murder. In these letters the ambassador is unambiguously described as her confidential envoy.\texthebrew{⁠}\footnote{Cassiodorus, \emph{Var.} X.20 additum est etiam gaudio meo quod talem virum vestra serenitas destinavit; iv 23 legatum vestrum . . . Petrum . . . vestris obsequiis inhaerentem.

} Here too we learn the significant fact that she enjoined on Theodahad that, if he made any request to the Emperor, he should first submit it to her.\texthebrew{⁠}\footnote{\emph{Ib.} 20 hortamini eum ut quicquid expetendum a triumphali principe domno iugali vestro credimus, vestris ante sensibus ingeramus.

} Moreover, in a letter of Theodahad's wife Gudeliva to Theodora, there is a mysterious passage which, in the light of what Procopius tells us, can be most easily explained as a veiled reference to the crime. ``While it is not seemly,'' wrote Gudeliva, ``that there should be any discord between the Roman realms, an affair has occurred of such a kind as fitly to render us dearer to you.''\texthebrew{⁠}\footnote{\emph{Ib.} 21 emersit tamen et qualitas rei, quae nos efficere cariores vestrae debeat aequitati. It must be remembered that the interests of Gudeliva were also involved. Her position as wife of the king, with Amalasuntha as the queen, would have been intolerable.

\texthebrew{◂} previous

next \texthebrew{▸}Images with borders lead to more information.

The thicker the border, the more information.

(Details here.)

UP TO:

J. B. Bury
Homepage

Latin \& Greek

Texts

LacusCurtius

Home} In a letter despatched immediately after the murder, this sentence bears an ominous significance.

The story of Procopius implies that the secret intrigues were known to a wide circle. Even if that were not so, he might have received information from Antonina, who was in Theodora's confidence, or from Peter himself. We must remember too that Theodahad, when he abandoned all thoughts of peace, had no motive to conceal the guilty intervention of Theodora. The conclusion that she did intervene and that Peter, acting by her orders, promoted the murder of Amalasuntha by hints and indirections, while he was ostensibly, in obedience to Justinian, acting in the interests of the queen, seems to be warranted by the evidence considered as a whole. This evidence would, of course, be far from sufficient to procure her conviction in a legal court. No public prosecutor could act on it. But where a jury would not be justified in convicting, public opinion is frequently justified in judging that a charge is true.

\xchapter{Chapter XVIII, Part B}

\xsubsection{(Part 2 of 5)}

Soon after the crime Peter returned to Constantinople. He bore letters from Theodahad and his wife to Justinian and Theodora; and he was to be followed presently by an Italian ecclesiastical, perhaps Pope Agapetus himself.\texthebrew{⁠}\footnote{If we read the six letters in Cass. \emph{Var.} X.19\texthebrew{‑}24 together, it appears unquestionable that they were despatched about the same time. 19 (Theod. to Justinian), 20 (Theod. to Theodora), 21 (Gudeliva to Theodora) were entrusted to Peter; 22 (Theod. to Justinian), 23 (Theod. to Theodora), 24 (Gudeliva to Theodora) were to be committed to the unnamed ecclesiastic, who is referred to in the same terms, in five of the letters, as illum virum venerabilem. It is obvious from the tenor that hostilities had not yet begun, and that Theodahad's object in writing was to avert them. Further, 19 was the \emph{first} communication addressed by him to the Emperor since Peter's arrival early in 535, as is clear from the first sentence, in which he thanks Justinian for congratulating him on his elevation (gratias divinitati referimus quod provectum nostrum clementiae vestrae gratissimum esse declarastis). We may therefore confidently date these documents to summer 535. Another document also belongs here, xi.13, an appeal of the Roman Senate to preserve peace. The tenor seems to point to a date before the outbreak of war, and the statement that the letter would be delivered per illum virum venerabilem legatum piissimi regis nostri seems conclusive. We may naturally connect this letter with the record found only in Liberatus, \emph{Brev.} 21, that Theodahad compelled the Senate by threats to make such an appeal. This vir venerabilis is supposed by some to have been Rusticus, who afterwards accompanied Peter to Constantinople in 536 (see below,

p173 ; cp. Dahn, \emph{Kön. der Germ.} II.203). Leuthold, Körbs, and others have attempted to prove that Pope Agapetus is meant, and that the letters were written and sent in February 536. But the tenor of the letters is quite inappropriate to the situation then. The important passage is 19.4, 5, and I am inclined to think that Theodahad means he will send the Pope, not with Peter, but later on, as his own envoy. The Pope, however, did not start till a much later date (see below,

p172 ).

} The object of Theodahad was to avert hostilities, and it is clear that he relied, above all, on the influence of Theodora. It is said that he forced the Roman senators to address Justinian in behalf of peace, by threatening to slay them, with their wives and children, if they refused. And we possess a letter of the Senate, drawn up by the Praetorian Prefect Cassiodorus, professing deep affection for the Amal ruler, nourished at the breasts of Rome, and imploring the Emperor to keep the peace. But the king's hopes of a peaceful settlement were vain. The Emperor immediately prepared for war. The idea of restoring the Imperial power in Italy had probably and long in his mind, his diplomacy had been occupied with it during the past year, and in a law issued six weeks before the murder of the queen he seems to allude to the

p169
Italian enterprise.\texthebrew{⁠}\footnote{Justinian, \emph{Nov.} 6 (March 6, 535) \emph{ad init.}: ea quae sunt firma habebimus et quae nondum hactenus venerunt adquirimus.

} If Theodahad were willing to abdicate and give the Emperor peaceful possession, well and good, but the only alternative was war. On this Justinian was fully resolved, and Peter, who returned to Italy during the summer ( A.D. 535), must have been the bearer of this ultimatum. In the meantime Justinian pushed on the preparations for war.\footnote{Procopius, \emph{B. G.} I.5.1, ``as soon as he learned what had happened to Amalasuntha, being in the ninth year of his reign, he entered upon war'' (the ninth regnal year began April 1, 535). Justinian must have heard the news by the end of May, and Belisarius may have sailed at the end of June. His sailing marked the beginning of the war, and thus we can understand (cp. Körbs, \emph{Untersuchungen}, p60) why the years of the war as reckoned by Procopius run from summer solstice to summer solstice. This reckoning is quite clear; it was pointed out by Eckhardt and Leuthold and has been established by Körbs. Thus year 1 = end of June 535 to end of June 536. The system has been misunderstood because, in noting the end of each year, Procopius uses the formula ``the winter was over and the 1st (2nd, etc.) year of the war came to an end,'' in imitation of Thucydides. The inference that the new year began with the spring equinox was natural, but is inconsistent with the narrative. The end of the winter and the end of the year of war are not coincident, and the former is only introduced to remind the reader of Thucydides (Körbs, 56 \emph{sqq.}).

}

The war against the Goths was begun in a very different way from the war against the Vandals. The Emperor had taken his subjects into his confidence when he prepared the African expedition; all the world knew that he was committed to the subjugation of Africa. But the outbreak of hostilities, which was to lead to the subjugation of Italy, was carefully concealed so long as concealment was possible; and the first steps were so contrived as not to commit the government to immediate operations on an extensive scale, if the task should appear too formidable. It is probable that Justinian was still waiting on events in Italy, and calculating that Theodahad, who was devoid of military spirit and capacity, would on the first symptoms of danger yield to all his demands. It was a calculation in which too little account was taken of the feelings of the Ostrogothic people.

The first operations in the war would indeed have been dictated, in any case, by geographical considerations. To occupy the Gothic province of Dalmatia, which was accessible by land, and that of Sicily, which was the most easily accessible by sea, were obviously, for a power which commanded the sea, the first

p170
things to be done. The possession of these two provinces would provide the bases for the conquest of Italy. Mundus, the loyal Gepid, Master of Soldiers in Illyricum, led the forces against Dalmatia. The resistance there seems to have been weak. He defeated the Goths and occupied Salona.\footnote{Perhaps in August or September 535.

}

The conqueror of Africa was marked out for the command of the overseas expedition, and the full powers of an imperator were again conferred upon him.\texthebrew{⁠}\footnote{He is named a \textgreek{στρατηγ\texthebrew{ὸ}ς α\texthebrew{ὐ}τοκράτωρ} (Procopius, \emph{B. G.} I.5.4. See above,

p127 ).

} But the army which was entrusted to him was hardly half as strong as that which he had led against the Vandals. It consisted of 4000 legionaries and Federates; a special division of 3000 Isaurians under Eunes; 200 Huns, 300 Moors; and the armed retainers of Belisarius, who may have amounted to several hundreds. Thus the total strength was about 8000. The principal generals were Constantine and Bessas, both Thracians; and the Iberian prince Peranius.\texthebrew{⁠}\footnote{Son of king Gurgen.

} Belisarius was accompanied by his stepson Photius, still a stripling, but strong and intelligent beyond his age.

The purpose of the expedition was kept secret. It was given out that destination of the fleet was Carthage, and no one had any idea that its sailing was the first step in a new enterprise. Belisarius was instructed that on landing in Sicily he should still pretend to be on his way to Africa, and should do nothing until he had discovered whether the island could be subjugated without trouble. This would evidently depend on the disposition of the Sicilians and the strength of the Gothic garrisons. If it appeared that he was likely to meet with a serious resistance, he was to proceed to Africa as if no other intention had been entertained. He was to run no risks with his small army. This cautious plan of action shows that the Emperor was not yet prepared to commit himself to an Italian campaign. The operations of Belisarius and Mundus were designed, in the first instance, as auxiliary to the Imperial diplomacy.

If war could not be avoided, Justinian calculated upon obtaining some aid from beyond the Alps. He sent an embassy to the kings of the Franks, urging that it was their interest as a Catholic power to co-operate with him against the Arian Goths,

p171
and as he supported his arguments by gold, he secured unreserved promises of assistance.\footnote{The Merovingian kingdom was now divided among three rulers, Choltachar and Childebert, sons of Chlodwig, and their nephew Theodebert. It was Theodebert who was the leading spirit in the Italian policy of the Franks (see below,

§ 8 ).

}

Belisarius disembarked at Catane and he found his work easier than he could well have anticipated. Having seized Catane, he occupied Syracuse, and from the summary statement which has come down it would almost seem that no resistance was offered anywhere and that no military operations were necessary, except at Panormus. Here the fortifications were strong, and the Gothic garrison, which was probably larger than in the other cities, refused to surrender. The Imperial fleet sailed into the harbour, which was unfortified. The masts of the ships overtopped the walls of the town, and Belisarius conceived the device of hoisting boats, full of soldiers, to the tops of the masts, so that they could shoot down upon the defenders. To this menace the Goths, who must have been half-hearted in their resistance, immediately yielded. The restoration of Roman rule in the island was completed before the end of December. Belisarius was one of the consuls of the year, and on December 31 he was able to enter Syracuse and formally lay down his office. The coincidence seemed to his contemporaries a signal favour of fortune.

The ease with which Sicily was reduced shows that the Sicilians were ready to exchange the yoke of Ravenna for that of New Rome, and that there were not large Gothic forces in the island. It may be observed that it would have been far more difficult for a small garrison, in those days, to hold a town of considerable size against a foe, in spite of the wishes of the inhabitants, than in modern times. A slender force, armed with sword and spear, could not defy a numerous populace, as they might if they were supplied with firearms.

In the meantime communications had passed between Rome and Constantinople.\texthebrew{⁠}\footnote{It is highly probable that Theodahad was at Rome during the whole winter 535\texthebrew{‑}536, and directed the negotiations from there (Cassiodorus, \emph{Var.} X.14 and 18, cp. xii.18 and 19)., and that it was then that (as Wroth saw) he issued the fine bronze coins which bear his bust, with Dn Theodahatus rex on the obverse, and Victoria Principum on the reverse. On these coins the head is ``neither bare nor bound with a diadem, but wearing a closed crown ornamented with jewels and two stars. His robe is also richly ornamented with jewels

(p172)
and a cross. The hair is cut short; the face beardless, but with a moustache such as has been seen already on the portrait coins of Odovacar and Theoderic'': Wroth (\emph{Coins of Vandals, etc.}, p. xxxiv), who sees no reason not to regard it as a true portrait, and suggests that it should be connected with Cassiodorus, \emph{Var.} VI.7 ut figura vultus nostri metallis usualibus imprimatur, etc. The other (Ravennate) coinage of Theodahad is of the ordinary type.

} Alarmed by the operations in

p172
Dalmatia and Sicily, king Theodahad made a new effort to persuade the Emperor to desist from his purpose. He induced Pope Agapetus to undertake the office of ambassador to Constantinople (early winter, A.D. 535).\texthebrew{⁠}\footnote{Liberatus, \emph{Brev.} 21, ipsi papae et senatui Romano interminatur non solum senatores sed et uxores et filios filiasque eorum gladio se interemturum nisi egissent apud imperatorem ut destinatum exercitum suum de Italia submoveret (Italia is used in its political sense, including Sicily). The date 535 is supplied by \emph{Cont. Marcellini}, \emph{s.a.} It has generally been supposed that Agapetus was sent in 536, on account of a passage in the life of this Pope in \emph{Lib. pont.} p287 ambulavit Constantinopolim X kal. Mai., where Clinton read X kal. Mart. = Feb. 21, which is taken to be the date of his arrival at Constantinople, though it would naturally be that of his departure from Rome. Now as X kal. Mai., April 22, is the date of the Pope's death, Duchesne is certainly right in regarding the date as an interpolation. Against the emendation it may further be urged that Feb. 21-March 13, on which day Agapetus consecrated the Patriarch Menas, is too short a time for the proceedings preliminary to this ecclesiastical victory. Finally, the mission of the Pope \emph{after} the despatch of Peter with the two letters is unintelligible. These objections are conclusive. The view that Agapetus accompanied Peter is quite untenable. Agapetus was in Italy on Sept. 9 (see his letters in Mansi, viii.848, 850) and as late as Oct. 15, if the date of the letter to Justinian (\emph{ib.} 850) is correct. He may have started in November or December.

} The appeal did not avail. We are not told how the Pope discharged the duties of a mission which he seems to have undertaken reluctantly, but he soon became absorbed in the ecclesiastical controversies of Constantinople, where he remained till his death (April 22, 536). Meanwhile the successes of Mundus and Belisarius increased the fears of Theodahad, and the fall of Panormus seems to have been decisive. The Imperial envoy Peter, who had returned from Constantinople to Rome, was able to take advantage of the completion of the conquest of Sicily to persuade the vacillating king to attempt to come to terms with his master.\texthebrew{⁠}\footnote{From the order of the narrative in Procopius, i.6 \emph{ad init.}, it is natural to infer that Peter's interviews with the king were later than Dec. 31, 535. The inference, however, is not quite certain, for \textgreek{τα\texthebrew{ῦ}τα} (\textgreek{\texthebrew{ἐ}πε\texthebrew{ὶ} δ\texthebrew{ὲ} τα\texthebrew{ῦ}τα Πέτρος \texthebrew{ἔ}μαθεν}) \emph{might} mean loosely the progress of Belisarius in Sicily.

} Theodahad's fears made him amenable, and he handed to Peter a letter in which he offered to resign Sicily and to submit to a number of capitulations, which would clearly establish and confirm the Emperor's overlord ­ship.\texthebrew{⁠}\footnote{Theodahad undertook (1) to send to the Emperor yearly a gold crown weighing 300 lbs.; (2) to furnish

300 Gothic soldiers, when required; (3) not to put to death, or confiscate the property of, any senator or cleric, without the Emperor's consent; (4) not to confer the patrician or the senatorial dignity, except with the same consent; (5) not to suffer his name to be acclaimed before that of the Emperor, in the circus or theatre; (6) to allow no statue to be set up to himself alone, without a statue of the Emperor standing on the right.

} Peter set out, but he had only reached

p173
Albano when he was recalled. Theodahad's craven spirit was tortured by the fear that his terms would be rejected, and he had decided to seek Peter's advice. The historian Procopius records a curious conversation between the king and the ambassador.\texthebrew{⁠}\footnote{There seems no reason to suppose that this conversation was invented by the historian. I take it as an important indication that Procopius was personally acquainted with Peter, and that Peter was his source for diplomatic history of the years 534\texthebrew{‑}536.

} ``Suppose my terms do not satisfy Justinian, what will happen?'' asked the king. ``You will have to fight,'' said Peter. ``Is that fair, my dear ambassador?'' ``Why not?'' replied Peter; ``it is fair that every man should be true to his own character.'' ``What do you mean?'' ``Your interest is philosophy,'' said Peter, ``while Justinian's is to be a good Roman Emperor. Observe the difference. It could never be seemly for a philosopher to cause death to men, and in such numbers; especially for a Platonist, whose hands should be pure of blood. Whereas it is natural that an Emperor should seek to recover territory which of old belongs to his dominion.'' Theodahad then swore, in Peter's presence, and caused queen Gudeliva to swear likewise, that he would deliver Italy over to Justinian, in case his first proposals were rejected. He wrote a letter to this effect, stipulating only that lands produ­cing a yearly revenue of 1200 lbs. of gold should be secured to him; but he made Peter promise by oath that he would first deliver the previous letter, and only produce the second in case the first proved unacceptable. In agreeing to this arrangement, Peter may seem to have had a strange idea of the duties of an ambassador, but we may take it for granted that he was perfectly certain that the compromise offered in the first communication would be rejected.\texthebrew{⁠}\footnote{Along with Peter, Theodahad sent Rusticus, an ecclesiastic, probably a bishop (cp. Leuthold, \emph{Untersuchungen}, p40).

} Rejected it was; the second letter was presented, and the Emperor was highly pleased. Peter was sent once more to Italy, along with another agent, to confirm the agreement, and to arrange that the estates of the patrimonium should be assigned to Theodahad's use.\texthebrew{⁠}\footnote{Probably in Tuscany. \emph{B. G.} I.6.26, cp. 4.1.

} Instructions

p174
were sent to Belisarius, who was still in Sicily, to be prepared to take possession of the royal palaces and assume the control of Italy.

When the ambassadors arrived they found Theodahad no longer in the same mood. Things in the meantime had been occurring in Dalmatia, where a considerable Gothic army had arrived to recover the province. Maurice, the son of Mundus, went out with a small force to reconnoitre, and fell in a sanguinary skirmish. His father, excited by grief and anger, immediately marched against the Goth, and almost annihilated their forces, but in the heat of a rash pursuit was mortally wounded.\texthebrew{⁠}\footnote{The deaths of Mundus and his son seemed to the Italians to supply the interpretation of a Sibylline oracle, which had evidently been in their minds since the conquest of the Vandals. The Latin words, as handed down in the MSS. of Procopius (\emph{B. G.} I.7, p33) present some difficulty. Africa capta mundus cum nat is quite clear, but is followed by uperistal in one MS., \textgreek{ζρεριστασι} in another. The Greek version, given by Procopius, \textgreek{\texthebrew{ἡ}νίκα \texthebrew{ἂ}ν \texthebrew{Ἀ}φρικ\texthebrew{ὴ ἔ}χεται, \texthebrew{ὁ} κόσμος ξ\texthebrew{ὺ}ν τ\texthebrew{ῷ} γόν\texthebrew{ῷ ὀ}λε\texthebrew{ῖ}ται}, points to the future of pereo. Braun and Haury read cum nato peribit, Comparetti peribunt. Periet, a form which occurs in Corippus, would be nearer (cp. Bury, \emph{B.Z.} xv.46).

} His death rendered the victory equivalent to a defeat. The Imperial army, in which it seems that there was none competent to take his place, withdrew from Dalmatia. The field was left to the Goths, but they too had lost their commander, and they did not at first venture to occupy Salona, where the Roman population was not friendly.

The news of these events elated Theodahad, whose unstable mind was vacillating between fear of war and the pleasures of royalty. When the Imperial ambassadors arrived, full of confidence and disregardful of his oath, he refused to fulfill his contract. The Gothic notables, to whom Justinian had sent a conciliatory letter, supported him in his refusal, and he went so far as to detain the ambassadors in close confinement.

On learning what had occurred the Emperor appointed Constantian, his Count of the Stable, to lead the Illyrian army to recover Dalmatia, and sent orders to Belisarius to invade Italy. The task of Constantian was easily enough accomplished. He transported his troops by sea from Dyrrhachium to Epidaurus (Ragusa), and the Goths, who had meanwhile seized Salona, believing that they could not defend it, withdrew towards Scardona. Marching to Salona, Constantian rebuilt parts of

p175
the walls, which were in disrepair, and the Gothic army then retired to Ravenna.\footnote{The loss of Dalmatia may be placed in March 536; the recovery in May-June. Peter must have returned to Italy by the beginning of April. Procopius notes the end of the first year of the war, after the recovery of Dalmatia (i.7 \emph{ad fin.}).

}

Belisarius was preparing to transport his army to Italy when he was summoned to Africa to suppress the military mutiny, with which Solomon was unable to cope (last days of March).\texthebrew{⁠}\footnote{See above,

p143 . Belisarius must have started for Carthage in the first days of April. He cannot have been back in Sicily before May, and he did not cross to Italy before the end of June (the second year of the war had already begun, Procopius \emph{B. G.} I.7.37 and 8.1). He did not reach Naples till October (winter was approaching, \emph{ib.} 9. 9).

} On his return, leaving garrisons in Syracuse and Panormus, he crossed the straits and landed at Rhegium. The defence of the straits was in the hands of Evermud, son-in\texthebrew{‑}law of the king. His forces were probably insignificant; he deserted to Belisarius, was sent to Constantinople, and rewarded by the patrician dignity. The general advanced by the coast road to Naples, accompanied by the fleet, and he met with no opposition.

He encamped before Naples, and received a deputation of citizens, who implored him not to press them to surrender; Naples is a place of no importance, they said, let him pass on and take Rome. The general, observing that he had not asked them for advice, promised that the Gothic garrison would be allowed to depart unharmed, and he privately promised large rewards to Stephen, the head of the deputation, if he could prevail upon the citizens to surrender. A meeting was held, and two influential orators, Pastor and Asclepiodotus, who were loyal to the Goth interest, induced the citizens to put forward demands which they were sure would not be granted. But Belisarius agreed to everything. Then Pastor and his fellow in public harangues urged that the general was not in a position to guarantee their security, and that the city was too strong to be taken. This view was supported by the Jews, who, favoured by Theoderic's policy, were deeply attached to Gothic rule, and it carried the day.

p176
Belisarius decided to besiege the place, but it proved a more difficult operation than he had expected. He cut the aqueduct, but this caused little inconvenience, as the town had good wells. The besiegers had no points of vantage from which they could conduct the attack. Ancient Naples included within its walls only a small portion of the modern city. It corresponded to a rectangular area of about 1000 by 800 yards, in which the church of San Lorenzo would be close to the centre. But the ground must have been distinctly higher than the modern level, to give the besieged the advantages which they possessed.\texthebrew{⁠}\footnote{Hodgkin, iv.49 \emph{sqq.}

} Having wasted some weeks and incurred serious losses in men, Belisarius, impatient to advance against Rome and meet Theodahad, determined to abandon the siege. But the luck which had signally favoured him hitherto was again with him. He had given orders to the army to prepare for departure, when a curious Isaurian, climbing into the broken aqueduct in order to inspect its construction, discovered that, near the walls, the channel had been pierced through solid rock, and that the aperture was still open, too narrow to admit a man in armour, but capable of being enlarged. Belisarius acted promptly. Files were employed to enlarge the opening, so as to make no noise. But before making use of this means of entering the city, the general gave the Neapolitans another chance to avoid bloodshed and the horrors of a sack. He summoned Stephen to his camp, assured him that it was now impossible that the city should not fall into his hands, and implored him to persuade his fellow-citizens to capitulate and avoid the miseries which would befall them. Stephen returned in tears, but the people refused to listen. They were convinced that the appeal of Belisarius was merely a ruse.

Six hundred men crept through the aqueduct at night, slew the sentinels on the northern wall, and enabled the Roman troops who were waiting below with scaling-ladders to ascend on the battlements. The horrors which Belisarius had anticipated ensued, and the Huns particularly distinguished themselves in the work of murder and plunder. At length the general succeeded in gathering the troops together and staying the carnage. Swords were sheathed and captives were released. Eight hundred Goths who were taken were well treated. The Neapolitans turned with anger against the two demagogues

p177
whom they held responsible for all that had befallen them. They slew Asclepiodotus; they found Pastor already dead, stricken by apoplexy when he knew that the city was taken.\footnote{Naples probably fell early in November. The siege lasted twenty days. We must allow some time for the march to Rome, which was reached on Dec. 9.

}

The people of Naples had confidently expected that king Theodahad would have sent an army to relieve their city. He seems to have been paralysed by fear; he took no measures for defence of his kingdom or of any part of it. Disgusted with his inactivity, the Goths of Rome and the province of Campania decided, after the fall of Naples, to depose him and elect a leader of military experience. They met at Regata in the Pomptine marshes,\texthebrew{⁠}\footnote{Regata (not Regeta) was 280 stades from Rome, near the Decennovian canal, which drained the marshes and reached the sea at Terracina. The text of Procopius (i.11) in this passage equates 19 miles with 113 stades (\textgreek{τρισκαίδεκα και \texthebrew{ἐ}κατόν}). The reading cannot be sound. Procopius reckoned 7 stades to a mile (see above,

p132 ). We should probably read \textgreek{τρε\texthebrew{ῖ}ς κα\texthebrew{ὶ} τριάκοντα} (\textgreek{λγ}´ instead of \textgreek{ιγ}´), 133\texthebrew{⁄}7 = 19. See Haury, \emph{B.Z.} xv.297, who conjectures that Regata is a later name for Forum Appii.

} and, as there was no suitable member of the royal family of the Amals, their choice fell on Witigis, a man of undistinguished birth, who had earned some repute in the campaigns against the Gepids. He was acclaimed king (November, A.D. 536),\texthebrew{⁠}\footnote{Procopius, \emph{ib.} Jordanes, \emph{Rom.} 372,

\emph{Get.} 310 . The date of the accession of Witigis cannot have been prior to October. For Theodahad, who was put to death a few days afterwards, was in the third year of his reign (\emph{B. G.} I.11), which began not earlier than October 3. There is no reason for rejecting the record in \emph{Cons. Ital.} (\emph{loc. cit.}) that he was slain mense Decembris, and thus we obtain the end of November for the elevation of Witigis. This accords with the probable date of the siege of Naples (see above,

p177, n1 ). The date usually assigned is August (so Clinton, Hodgkin , Martroye, etc.) on the testimony of a passage in \emph{Lib. pont., Silverius}, p290), where it is stated that two months after the ordination of Pope silverius (June 20) Theodahad was killed and Witigis elected. As this contradicts all the other evidence, we must reject it. I have little doubt that the text (post menses vero ii.) offers an instance of the very common numerical confusion of ii with v. Leuthold (\emph{op. cit.} 48), who is right in his view of the chronology, makes a subtle and unsatisfactory attempt to explain the wrong date in this text.

} and Cassiodorus, whose impartial pen was prepared to serve him, as it had served Theoderic and Amalasuntha, and as it had served Theodahad, announced to all the Goths the election of one not chosen, like Theodahad, ``in the recesses of a royal bedchamber, but in the expanse of the boundless Campagna;\texthebrew{⁠}\footnote{In campis late patentibus.

} of one who owed his dignity first to Divine grace, but secondly to the free judgment of the people; of one who knew the brave men in his army by comrade ­ship, having

p178
stood shoulder to shoulder with them in the day of battle.''\texthebrew{⁠}\footnote{Cassiodorus, x.31 (Hodgkin's paraphrase, iii.74 ).

} The event proved that the choice of the Goths was undiscerning. Witigis was a responsible soldier, and would have been a valuable leader of a division under an able commander, but he possessed none of the higher qualities demanded in one who was called to lead a nation against a formidable invader.

Theodahad, who had hitherto been residing at Rome, fled incontinently to Ravenna. Witigis decided that he must die, and sent a certain Optaris to bring him alive or dead. Optaris was selected because he had a personal grudge against Theodahad. Travelling night and day without a pause, he overtook the fugitive, flung him on the ground, and butchered him like a sacrificial victim.

The new king immediately marched to Rome and held a council. Everything depended on the plan of campaign that was now formed. The Goths were menaced by two dangers, the imminent advance of Belisarius from the south, and the hostile attitude of the Franks in the north. The main forces of the Goths were stationed in the northern frontier provinces, in Provence and Venetia. Witigis proposed, and the proposal was accepted, first of all to deal with the Franks, and then to take the field against Belisarius with all the forces of the kingdom. It is safe to say that this plan of postponing the encounter with the most dangerous enemy was unwise.\texthebrew{⁠}\footnote{Compare the remarks of Hodgkin, iv.76 .

} The best chance of the Goths would have been to hurry the main part of their troops from the north, and either join battle with the Imperial army before it reached Rome, or else hold Rome strongly and force Belisarius to undertake a siege which would be long and difficult. In the meantime an envoy could be sent to negotiate with the Franks. The place of Witigis himself was at Rome, the threatened point, and he committed a fatal blunder when he started for Ravenna ``to make arrangements for the war.'' He left a garrison of 4000 men in Rome, under Leuderis, extracted an oath of fidelity from Silverius, the Pope,\texthebrew{⁠}\footnote{Silverius (son of a previous Pope, Hormisdas) succeeded Agapetus in June 536. It is said that he was imposed on the electors by Theodahad (\emph{Lib. pont., loc. cit.}), but this may be only an invention. Cp. Liberatus, \emph{Brev.} c22; Hodgkin, iv.92 .

} and from

p179
the Senate and people, and took a number of senators with him as hostages.

At Ravenna, Witigis married, against her will, Matasuntha, the sister of Athalaric, in order to link himself with the dynasty of Theoderic; and the wedding was celebrated in a florid oration by Cassiodorus.\texthebrew{⁠}\footnote{Fragments have been preserved and, edited by Traube, are included in Mommsen's ed. of the \emph{Variae}, p473 \emph{sqq.} The orator refers contemptuously to the effeminate training of Athalaric by his mother, and enlarges on the military career and prowess of Witigis. An elaborate description of the sumptuous pomp of the nuptials follows. There are coins with the monogram of Matasuntha, and they have been generally attributed to the reign of Witigis, but for another view see below,

p254 . The coinage of Witigis is of the ordinary type.

} He then proceeded to negotiate with the Franks. We saw how they had been induced by Justinian to promise their co-operation . But Theodahad had made them an attractive offer. He proposed to hand over to them the Ostrogothic territory in Gaul, along with 2000 lbs. of gold, in return for their engagement to assist him in the war. He died before the transaction was concluded. Witigis saw that the best thing to be done was to carry out this arrangement. The Frank kings consented, but, as they did not wish openly to break their compact with Justinian, they promised secretly to send as auxiliaries ``not Franks, but men of their tributary peoples.''\footnote{Procopius says that the three Frank kings divided the money and the land (\emph{B. G.} I.13.27). But they did not consider their title secure until the land was formally ceded by the Emperor, and Justinian deemed it wise to agree (III.33.3).

}

At the same time a last attempt was made to come to terms with the Emperor. It was plausible to argue that, as the murder of Amalasuntha had been the alleged reason for invading Italy, the cause for war was removed by the punishment of Theodahad and the elevation of Matasuntha to the throne. What more could the Goths do? Witigis wrote to Justinian to this effect,\texthebrew{⁠}\footnote{Cassiodorus, \emph{Var.} X.32. Witigis at the same time appealed to the Catholic bishops of Italy to pray for peace, and wrote to the Praet. Pref. of Illyricum asking him to help his ambassadors at Thessalonica on their journey. \emph{Ib.} 34, 35.

} and likewise to the Master of Offices urging him to work for peace.\texthebrew{⁠}\footnote{\emph{Ib.} 33.

} As to these negotiations we possess only the documents drawn up at Ravenna, and have no information as to the Emperor's reply. We may conjecture that he offered Witigis the simple alternative between war and submission.

In the meantime Belisarius had left Naples and was marching

p180
northward. The Romans, warned by the experiences of Naples, and urged by the Pope, who had no scruples in breaking his oath to Witigis, sent a messenger inviting him to come. He had placed small garrisons in Naples and Cumae, the only forts in Campania, and marching by the Via Latina he entered Rome on December 9, A.D. 536,\texthebrew{⁠}\footnote{There is a lacuna in the text of Procopius (I.14.14) , but it can be supplied from

Evagrius, \emph{H. E.} IV.19 . See Haury, \emph{ad loc.}

} by the

Porta Asinaria , close to the Basilica of the Lateran.\texthebrew{⁠}\footnote{The old gate, which is walled up, stands beside the Porta San Giovanni.

\texthebrew{◂} previous

next \texthebrew{▸}Images with borders lead to more information.

The thicker the border, the more information.

(Details here.)

UP TO:

J. B. Bury
Homepage

Latin \& Greek

Texts

LacusCurtius

Home} On the same day the Gothic garrison discreetly withdrew by the

Porta Flaminia . Their leader, Leuderis, remained, and was sent to the Emperor with the keys of the city gates.

\xchapter{Chapter XVIII, Part C}

\xsubsection{(Part 3 of 5)}

The Romans soon learned to their deep chagrin that it was the intention of Belisarius to remain in their city and expose it to the hardships of a siege. With the small forces at his disposal, this was the only prudent course open to him. Taking up his quarters in the

Domus Pinciana , on the Pincian Hill, in the extreme north of the city, the general immediately set about strengthening the fortifications. The great walls of Aurelian, which encompassed the city in a circuit of \texthebrew{•} about twelve miles, had been repaired more than a hundred years ago, in the reign of Honorius, and recently by Theoderic.\texthebrew{⁠} But Belisarius found many dilapidations to make good, and he added some new fortifications. A wide ditch was dug on the outer side. The wall, as originally constructed, was well adapted for defence. A special feature was a covered way running round the inside of the wall to facilitate the passage of troops from one point to another. Some portions of this arched gallery still remain.\texthebrew{⁠}\footnote{Near the Porta Asinaria.

} Considering the vicissitudes through which Rome subsequently passed in a period of thirteen hundred years, the walls which the army of Belisarius defended are wonderfully preserved.\footnote{A summary enumeration of the towers, battlements, and loopholes, from the Porta Flaminia to the Porta Metrovia, compiled (copied from an earlier document?) in the eighth century, is extant (text in Jordan, \emph{Top. der Stadt Rom}, II.578).

}

At the same time measures were taken to supply the city with stores of grain imported from Sicily. But Belisarius appears not to have expected that Rome would be attacked by

p181
a formidable army. He diminished his garrison by flinging out forces northward to seize commanding positions along the Flaminian Way \texthebrew{—} Narni, Spoleto, and Perugia, and some lesser strongholds. In the meantime Witigis had sent a considerable detachment to Dalmatia. Salona was besieged by land and sea, but the diversion ended in failure, and the province remained in Imperial hands.\texthebrew{⁠}\footnote{The end of this expedition is not related by Procopius.

} An attempt to recover Perugia was also defeated. But the confidence of the Goths rose when they realised the weakness of the forces with which Rome was held, and heard rumours of the discontent of the inhabitants at the military occupation of their city. The king decided to throw all his strength into the recovery of Rome, and he marched southward at the head of an army, which is thought by some to have numbered 150,000 warriors, most of them heavily armed, with horses protected by mail. The figure must be far in excess of the truth,\texthebrew{⁠}\footnote{The figure, nevertheless, was given by Belisarius in his letter to Justinian (\emph{B. G.} I.24); see below,
p186 . 150,000 Gothic warriors would mean a Gothic operation approaching 700,000. When they entered Italy the number of the Ostrogoths perhaps hardly reached 100,000, and they cannot have multiplied seven times in forty-five years. Moreover, if they had been so strong, it would have been out of the question to attempt to conquer them with the small forces of Belisarius. The figure is also disproved by the circumstances of the siege (see below,
pp183 \emph{sq.} ). I am inclined to believe that the number mentioned by Belisarius represents an estimate of the total Gothic population of Italy.

} but there can be no doubt that the Gothic host was large compared with the army of 5000 against which it was advancing. Belisarius was now dealing with a very different problem from that which had faced him in his campaign against the Vandals. He hastily recalled the generals, Bessas and Constantine, whom he had sent into Tuscany, bidding them abandon all their positions except Perugia, Spoleto, and Narni, in which they were to leave small garrisons.

Witigis did not delay to reduce these three places. The occupation of Narni was important.\texthebrew{⁠}\footnote{Narnia is 54 Roman miles from Rome. Its situation is described in some detail by Procopius (I.17.8\texthebrew{‑}11) , and the description is not irrelevant, as showing that even with a small garrison it could bar the progress of an army. One arch and some piers of the bridge of Augustus which Procopius mentions still remain.

} It forced the Gothic army, just as, more than a hundred years before, the army of Alaric had been forced, to diverge from the Flaminian Road to the east, to march through the Sabine country, and approach Rome by the Via Salaria, instead of marching by the Via p182 Flaminia.\texthebrew{⁠}\footnote{The usual view has been that the Goths advanced by the Via Flaminia (regaining it somewhere presumably between Narnia and Ad Tiberim, now Magliano, where there was a bridge), and that the bridge where Belisarius placed the garrison was the Pons Milvius, now Ponte Molle, 2 miles from Rome. This view was held by Gibbon and maintained by Hodgkin . But it is certainly erroneous and inconsistent with the story. If the fighting had been at the Milvian Bridge the Roman fugitives would have returned to the Porta Flaminia, not to the Porta Salaria. The cause of the error is that Procopius does not name the bridge, but calls it simply \textgreek{Τιβέριδος το\texthebrew{ῦ} ποταμο\texthebrew{ῦ} γεφύρα}, ``a bridge of the Tiber'' (\emph{ib.} xiii) . Hence, as the Milvian was the only bridge which spanned the Tiber north of the city, it was naturally supposed to be meant. But \textgreek{Τίβερις} is ambiguous in Procopius; it means (1) the Tiber, (2) the Anio. That it means the Anio in this passage is shown by the statement in the context (\emph{ib.} xiv) that there are bridges over the river in other places (\textgreek{πολλαχόσε το\texthebrew{ῦ} ποταμο\texthebrew{ῦ}}), meaning, of course, in the neighbourhood of Rome. This is not true of the Tiber, which had only the one bridge outside the city; but it is true of the Anio, which is crossed, near Rome, by the Via Nomentana and the Via Tiburtina, as well as by the Via Salaria. In two other passages (\emph{B. G.} III.10.23, and 24.31) \textgreek{Τίβερις} clearly means the Anio. This was the view of Gregorovius, \emph{Rome in the Middle Ages}, I.372; is accepted by Hartmann, \emph{Gesch. Italiens}, I.295, n19; and has been defended in a special monograph by L. Fink, \emph{Das Verhältnis der Aniobrücken zur Mulvischen Brücke in Prokops Gothenkrieg}, 1907. Procopius knew the localities, and the ambiguous use of \textgreek{Τίβερις} cannot be due to ignorance. The explanation may be found in the modern name of the Anio, Teverone, and the use in Procopius be taken to show that the old name had passed out of common speech before his time.

} When Witigis reached the Ponte Salario, where the road crosses the Anio, a few miles from the city, he found himself arrested by a fort which Belisarius had built on the bridge with the object of gaining time in order to procure more provisions.

But the garrison of the fort failed him. On the arrival of the Goths they decamped by night, and the enemy secured the bridge. Next day the general, ignorant of the cowardice of his men, rode towards the bridge with a thousand horsemen, and found that the Goths had crossed. A cavalry engagement ensued, in which Belisarius, carried away by the excitement of battle, indiscreetly exposed himself. Deserters knew his dark-grey horse with a white head, and urged the Goths to aim at him. But he escaped unwounded. There were severe losses on both sides, and the small Roman band was in the end forced to flee. They reached the Salarian Gate about sunset, and the sentinels, not recognising the general begrimed with dust of battle, and already informed by fugitives that he was slain, refused to open. Belisarius turned and charged the pursuers, who retreated, thinking that a new army had issued from the gate. He then succeeded in obtaining admission, and spent the

p183
night in making arrangements for the defence of the city. Each gate was assigned to the charge of a different leader. One more incident occurred before the night was over. Witigis sent an officer to make a speech outside the Salarian Gate. This man, whose name was Wacis, reproached the Romans for their treachery to the Goths and for preferring the protection of Greeks,\texthebrew{⁠}\footnote{Procopius reproduces the Latin name \texthebrew{—} \textgreek{Γραικοί} (I.18) .

} people, he said, who had never visited Italy before except in the capacity of actors or thieving mariners. No one made any reply to his outburst and he retired.

On the following day the siege began.\texthebrew{⁠}\footnote{For the chronology of the siege Procopius supplies the following data. It ended ``about the spring equinox'' in 538, and lasted 1 year and 9 days (II.10 , p194). It began at the beginning of March (\textgreek{Μαρτίου \texthebrew{ἰ}σταμένου}, II.24 , p122). In the \emph{Lib. pont. (Silverius)}, Feb. 21 is mentioned as the first day of the siege. It is not easy to reconcile this difference.

} It was to last a year and nine days, far longer than either of the belligerents anticipated. The Goths did not attempt to surround the whole circuit of the city. They constructed seven camps, one on the west side of the river, in the region of the Vatican, then known as the Campus Neronis. The other six were east of the Tiber, on the northern and eastern sides of the city.\texthebrew{⁠}\footnote{They can be located as follows: (1) just north of the Flaminian Gate, Porta del Popolo; (2) in the grounds of the Villa Borghese, to command the Salarian and Pincian Gates; (3) on the Via Nomentana, to command the Porta Nomentana; (4) and (5) on the Via Tiburtina, near San Lorenzo, and (6) farther south, on the Via Praenestina, to command the corresponding gates (P. Tiburtina and P. Labicana). It is possible that the Porta Pinciana had been closed and was reopened by Belisarius (see Jordan, \emph{Topogr.} I.1., p352, cp. II. p578).

} One of them was under the command of Witigis himself.\texthebrew{⁠}\footnote{\textgreek{Τ\texthebrew{ῶ}ν δ\texthebrew{ὲ} δ\texthebrew{ὴ ἄ}λλων Ο\texthebrew{ὔ}ττιγις \texthebrew{ἡ}γε\texthebrew{ῖ}το \texthebrew{ἔ}κτος α\texthebrew{ὐ}τός. \texthebrew{ἄ}ρχων γ\texthebrew{ὰ}ρ \texthebrew{ἦ}ν \texthebrew{ἔ}ς κατ\texthebrew{ὰ} χαράκωμα \texthebrew{ἔ}καστον} (\emph{B. G.} I.19). Hodgkin (IV.148) says that Procopius is ``rather vague here.'' Could he have been more explicit? During the siege the Goths profaned and damaged many tombs of Christian martyrs outside the walls. We know this from verses which were afterwards inscribed on the sepulchres when Pope Vigilius restored them. See \emph{Anth. Lat. Supp.} I. Nos. 83, 87, 89, 99.

} Thus from the Porta Maggiore to the Porta S. Paolo and the river there was no leaguer. The whole circuit of the Aurelian Wall, including the Transtiberine region, was \texthebrew{•} less than thirteen miles,\texthebrew{⁠}\footnote{This includes the wall along the river from the P. Flaminia to the Pons Aurelius (Ponte Sisto). See Jordan, I.1.343\texthebrew{‑}344. It must be remembered that the Transtiberine portion of the Aurelian Wall enclosed only the Janiculum and the southern part of the modern Transtiberine town, reaching the river just north of the Ponte Sisto. The Gothic camp on this side of the river was far from the walls, and its principal purpose was to prevent the Romans from destroying the Milvian Bridge (so Procopius). It was under the command of Marcias, who had been the commander of an army in Provence, and arrived on the scene after the siege had begun.

} so if Witigis had the

p184
huge host which he is supposed to have led against Rome he would have had a man to every foot of the wall and an army of more than 10,000 to spare. He could not have decided that he had too few to blockade the city completely.\footnote{\textgreek{Ο\texthebrew{ὐ}χ ο\texthebrew{ἶ}οί τε \texthebrew{ὄ}ντες \texthebrew{ὅ}λ\texthebrew{ῳ} τ\texthebrew{ῷ} στρατοπέδ\texthebrew{ῳ} τ\texthebrew{ὸ} τε\texthebrew{ῖ}χος περιλαβέσθαι κύκλ\texthebrew{ῳ}}, \emph{B. G.} I.19.2.

}

The first operation of the Goths was to cut the numerous aqueducts which traversed the Campagna and supplied Rome with water from the Latin hills. The destruction of these magnificent works, although it caused some inconvenience, hardly affected the fortunes of the siege; but it had far-reaching consequences for the future of Rome.\texthebrew{⁠}\footnote{Compare Hodgkin (IV. ch. vi) , who gives an interesting account of the aqueducts, based on Lanciani's monograph \emph{Le acque e gli acquedotti di Roma Antica.} There were eleven principal aqueducts, including that of Alexander Severus. Procopius says that there were fourteen (I.19) .

} Since the third century B.C. the city had been excellently supplied with pure water, and new conduits had constantly been built to meet the growing needs of the inhabitants. For a thousand years after the act of demolition wrought by the Goths, the Romans were again, as in the early Republic, compelled to draw their water from the Tiber and the wells. The time-honoured habits of luxurious bathing, which had been such a conspicuous feature of their civilisation, came to an end. The aqueducts might easily have been restored at the end of the war, and doubtless this would have been done if Rome had again become an Imperial residence, but the comfort and cleanliness of the people were no object of care to the medieval popes, who regarded the ancient Thermae as part of the unregenerate life of paganism. The long lines of arcades which crossed the Campagna were allowed to fall into ruin.

The cutting of the aqueducts caused an immediate difficulty. There was no water to turn the corn º mills which supplied the Romans with bread. The inventive brain of Belisarius devised an expedient. Close to a bridge (probably the Pons Ælius) through whose arch the stream of the Tiber bore down with considerable force, he stretched from bank to bank tense ropes to which he attached two boats, separated by a space of two feet. Two mills were placed on each boat, and between the skiffs was suspended the water-wheel , which the current easily turned. A line of boats was formed, and a series of mills in the bed of the river ground all the corn that was required. The efforts of the

p185
enemy to disconcert this ingenious device and break the machines by throwing trees and corpses into the water were easily thwarted by Belisarius; he stretched across the stream iron chains which formed an impassable barrier against all dangerous obstacles that might harm his boats or wheels.

The Romans chafed under the hardships which the first days of the siege brought upon them and which seemed likely to increase. Witigis, informed of their discontent by deserters, thought that Belisarius, under the influence of public opinion, might be induced to relinquish his plan of defending Rome if a favourable proposal were made. He sent envoys, whom Belisarius received in the presence of his generals and the senators. The Gothic spokesman enlarged on the miseries which the siege must inflict on the Romans, and offered to permit the Imperial army to leave the city unharmed and with all their property. The reply of Belisarius was a stern refusal. ``I tell you,'' he said, ``the time will come when you shall be glad to hide your heads under the thorn bushes and shall be unable to do so. Rome belongs to us of old. You have no right to it. It is impossible for Belisarius to surrender it, while he is alive.''

A grand attempt to take the city by assault soon followed. The walls were attacked in various places, but everywhere the besiegers were repelled. The fighting was particularly severe near the Aurelian Gate, west of the Tiber, where the Goths attacked the great quadrangular Mausoleum of Hadrian, and the defenders, hard pressed, hurled statues down upon the enemy.\footnote{It was supposed that the Gothic losses in dead on this day were 30,000.

}

Belisarius, though he openly expressed complete confidence, was well aware of the dangers and difficulties of his situation, and knew that success was hardly possible unless new troops came at once his aid. He wrote a letter to Justinian, in which he reported his operations and urged in the strongest language his need of reinforcements. ``So far,'' he wrote, ``all has gone well, whether our success be due to valour or to fortune, but in order that this success may continue, it behoves me to declare plainly what it behoves you to do. Though God orders all things as He wills, yet men are praised or blamed according to their success or failure. Let arms and soldiers be sent to us in such numbers that henceforward we may wage the war on terms of equality. Let the conviction penetrate your mind, O Emperor, that if

p186
the barbarians overcome us now, we shall lose not only your dominion of Italy but the army also, and besides this we shall suffer the immense disgrace of failure, not to speak of the shame of bringing ruin on the Romans who preferred loyalty to your throne to their own safety. Understand that it is not possible to hold Rome long with ever so large a host. It is surrounded by open country and, not being a seaport, it is cut off from supplies. The Romans are now friendly, but if their hardships are protracted the pinch of famine will force them to do many things against their own wishes. For myself, I know that my life belongs to your Majesty, and I shall not be forced out of this place while I live. But consider how such an end to the life of Belisarius would affect your reputation.''\footnote{I have reproduced a good part of this document because, if not a literal copy of the original letter, there is every reason to believe that Procopius, who may have written it himself, reproduced its actual tenor (\emph{B. G.} I.24).

}

The Emperor had despatched reinforcements in December under Valerian and Martin, but they spent the winter months in Greece and had not yet arrived. On receiving the urgent appeal of his general, Justinian ordered them to proceed without delay, and prepared to raise a new armament. Meanwhile, on the day following the Gothic assault, Belisarius sent the women and children and the slaves who were not employed in garrison duties out of the city. Some travelled by boat down the Tiber, others departed by the Appian Way. The enemy made no attempt to hinder their departure. The artisans and tradespeople, whose occupation was almost gone, were drafted into the garrison, mixed with the regular soldiers, and paid a small wage for their services.\footnote{A scandal was created at this time by some persons who attempted to open the gates of the temple of Janus in the Forum. Since Rome had become Christian the temple had been kept shut in war as well as in peace. The efforts of the secret pagans (who remained undiscovered) failed to open the bronze doors, but damaged the bolts or hinges so that they would not shut tight.

}

Enraged, perhaps, at the failure of his attack, Witigis put to death the senators whom he kept at Ravenna as hostages, except a few who managed to escape. It was an act of barbarity which was seldom practised, and was as useless as it was cruel. At the same time he occupied Portus, at the mouth of the Tiber. This was a serious blow to the besieged, for Portus had for centuries been the port of Rome, with which it was connected by an excellent road and a towpath along the right bank of the

p187
river, so that heavy barges laden with supplies could be towed up by oxen without the aid of oars or sails. The older harbour of Ostia, over against Portus, remained in the hands of the Romans, but there was no towpath, so that the river traffic from here depended on the wind. Moreover, when the Goths threw a garrison of a thousand men into Portus, boats could not anchor at Ostia, and were forced to put in at Antium, a day's journey distant.\texthebrew{⁠}\footnote{On the decline of Ostia and rise of Portus see G. Calza, \emph{Gli scavi recenti nell' abitato di Ostia}, in the \emph{Mon. antich.} of the \emph{Acc. dei Lincei}, XXVI 1920. Ostia had been a rich and prosperous city, as the excavated ruins show, till the beginning of the fourth century, for the Portus of Claudius across the Tiber had been simply the port of Ostia where all the business was transacted. The policy of Constantine brought Portus directly under the central administration at Rome, and the gradual decline of the old port set in. Cassiodorus (\emph{Var.} VII.9), in describing the two towns as duo lumina, has in mind Ostia's ancient importance.

} The secretary of Belisarius regrets that 300 men could not have been spared to secure Portus, which was so strong that even so few could have held it.

About three weeks later Martin and Valerian arrived with a force of 1600 cavalry, mostly Huns and Slavs, and they succeeded in eluding the Goths and entering Rome. Sorties were carried out after their arrival with uniform success, which Belisarius ascribed to the superiority of his well-trained mounted archers; and, if he could have had his way, he would have continued to wear down the enemy by constant small sallies, in which little was risked. But the army, rendered confident through their successes and convinced of their superiority to the barbarians, clamoured for a pitched battle, and their leader, wearied by their importunities, reluctantly yielded. A general action was fought in the north of the city, on both sides of the river, and the Romans, routed by sheer weight of numbers, were driven back within the walls.

Towards the end of June the besieged began to feel the pinch of hunger and disease. There was only enough corn to feed the soldiers, and the Goths tightened the blockade, hitherto conducted with remarkable negligence, by constructing a fortress at the junction of two aqueducts commanding the Appian and Latin Ways.\texthebrew{⁠}\footnote{The medieval tower, known as the Torre Fiscali, about 3½ miles from Rome, marks the place.

} The citizens\texthebrew{⁠}\footnote{Many had looked forward to the month of July as the term of their sufferings, on the faith of a Sibylline oracle which predicted that in that month the Romans would have a new king, and Rome would have nothing more to fear from the Getae.

(p188)
Procopius (I.24) reproduced the Latin words of the oracle, but they have been corrupted in the MSS. For attempts to restore the original see Bury (\emph{B. Z.} XV.45), and H. Jackson (\emph{Journal of Philology}, XXXIII.142) , who rightly points out that the word before mense (which is quite clear in the MSS.) must have been quinto, not quintili, as has generally been assumed.

} urged Belisarius to risk a battle.

p188
He refused, but held out promises that large reinforcements and supplies would soon arrive. The prospect of approaching relief was based only on rumour, and he sent his secretary Procopius to Campania to discover whether the report was true, to collect provision ships, and to send to Rome all the troops that could be spared from the garrisons of the Campanian towns. Procopius left Rome at night by the southern gate of St. Paul, and, eluding the Goths, reached Naples and executed his orders. Some time afterwards Belisarius sent Antonina to Naples, where, in a place of safety, she might help, with her considerable capacity for organisation, in the task of sending relief to Rome. She found that Procopius had already raised 500 soldiers and had loaded a large number of vessels with corn. But the reinforcements, so anxiously awaited, had not yet come, though they were on their way. They seem to have arrived in the month of November.\texthebrew{⁠}\footnote{They cannot have arrived sooner, for they must have reached Rome in December, since after their arrival there the Goths \emph{immediately} despaired and sent envoys to Belisarius (\textgreek{ε\texthebrew{ὐ}θ\texthebrew{ὺ}ς μ\texthebrew{ὲ}ν \texthebrew{ἀ}πεγίνωσκον τ\texthebrew{ὸ}ν πόλεμον} (\emph{B. G.} II.6 \emph{ad init.}), and this was shortly before the winter solstice (II.7, p181). It seems to follow that Procopius cannot have been sent to Naples before September or October, though one would naturally infer from the narrative that he was sent in July (so Hodgkin, IV.246) . But it cannot be supposed that he would have kept back for four months the 500 men whom he collected and who were sorely needed at Rome. Antonina appears to have been back in Rome before November 18, for she took part in the deposition of Pope Silverius (\emph{Lib. pont.} lx \emph{Silv.} p292): cp. below,
p379 .

} 3000 Isaurians disembarked at Naples and 1800 cavalry at Otranto. Of their commanders the most distinguished was John, the nephew of Vitalian, one of the bravest and most skilful officers who served under Belisarius.\footnote{He is called the ``Sanguinary'' in the \emph{Lib. pont.} (\emph{ib.}).

}

In the meantime the army of Witigis was suffering, as well as the Romans, from famine and disease. It was steadily declining in numbers when the discouraging tidings came that new forces were on their way to the relief of Rome. The 3000 Isaurians were sent by sea to Ostia, but John, the nephew of Vitalian, with his 1800 cavalry and the 500 who had been raised by Procopius, marched by the Appian Way, followed by a train of waggons laden with food. To prevent the Goths from intercepting them in force, Belisarius arranged a strong sortie on the

p189
camp near the Flaminian Gate. It was completely success ­ful; the Goths were utterly routed. This was the turning-point in the siege. Witigis despaired of taking Rome and sent envoys to Belisarius, the chief of whom was a distinguished but unnamed Italian.

The conversation between the general and the spokesman of the Goths is reported by Procopius, and, as we may safely assume that he had returned to Rome and was present at the interview, it is possible that he has given, at least partly, the tenor of the dialogue.

We may take the later part of this conversation as a genuine report. Nor is it improbable that the Italian delegate of the Goths raised the question of the constitutional position of Italy and the legitimacy of the Ostrogothic government. If so, it is interesting to observe that both his argument and the reply of Belisarius misrepresented historical facts. On the Gothic side it was stated that Odovacar's offence, in the eyes of Zeno, lay in the dethronement of Romulus Augustulus, whereas Zenos regarded Augustulus as a usurper, and it was out of respect for the rights of Julius Nepos that he at first refused to recognise Odovacar. But he did recognise him subsequently, so that Odovacar, at least during the later years of his reign, was as little a ``tyrant'' as Theoderic himself. Belisarius distorted facts more seriously. He completely ignored the definite agreement concluded between Theoderic and the Emperor Anastasius. It was on this agreement that the legitimacy of Ostrogothic rule rested, and its existence invalidated the argument of Belisarius. It is not too much to read between the lines that Procopius himself considered that legally the Goths had a good case.

While Belisarius was receiving the envoys the reinforcements were arriving at Ostia. The same night he rode down to the port and arranged that the provisions should be transported up

p191
the river and that the troops should march to Rome without delay. His confidence that the enemy would not interfere with the operations was justified by the event. The arrangements for an armistice of three months were then completed. Hostages were interchanged, and a guarantee was given that even if the truce were violated in Italy, the envoys should be allowed to return unharmed from Constantinople.

Rome was revictualled, but the Goths in their camps and fortresses were suffering from want of food. The secretary of Belisarius observes that the cause of this scarcity was the Imperial sea-power , which prevented them from receiving the imports on which Italy depended. The shortage of food decided Witigis to remove his garrisons from Portus, Centumcellae (Cività º Vecchia), and Album, and these places were promptly occupied by Imperial troops. The Goths complained of this action as a breach of the truce, but Belisarius laughed at them. He certainly put a free interpretation on the meaning of an armistice. He sent John, in command of 2000 troops, to spend the rest of the winter on the borders of Picenum, with instructions that, in case the enemy should break the truce, he was to swoop down on the Picentine territory, plunder it, and make slaves of the Gothic women and children.

About this time the attention of Belisarius was directed to the situation of northern Italy, where the inhabitants were watching the struggle with lively interest. Prominent citizens of Milan, along with Datius the archbishop, succeeded in reaching Rome, and begged him to send a small force to the north, assuring him that it would be an easy matter not only to hold Milan but also to procure the revolt of the whole province of Liguria. Belisarius consented to the plan, but he could not execute it during the truce, and the Milanese emissaries remained at Rome for the winter.

Soon after this a tragic incident occurred, which, if we may believe the secretary of Belisarius, was connected with domestic scandals in the general's household. When Witigis was preparing to march on Rome, Praesidius, a distinguished citizen of Ravenna, rode with a few servants to Spoletium with the purpose of joining the Imperialist cause. The only valuables he carried with him were two daggers with sheaths richly adorned with gold and gems. He halted at a church outside Spoletium,

p192
which was then held by Constantine. This general heard about the precious daggers, and sent one of his followers to the church, who forced Praesidius to surrender his treasure. Praesidius went on to Rome, intent on complaining to Belisarius, but the emergencies and dangers of the siege hindered him from troubling the commander with his private grievance. As soon as the truce had been arranged he made his complaint and demanded redress. Belisarius urged Constantine to restore the weapons, but in vain. Then one day, as he was riding in the Forum, Praesidius seized his bridle, and loudly demanded whether it was permitted by the Imperial laws that when a suppliant arrived from the camp of the enemy he should be robbed of his property. Belisarius was compelled to promise that the daggers should be restored, and summoning Constantine to a private room, in the presence of other generals, told him that he must give up the daggers. Constantine replied that he would rather throw them into the Tiber. Belisarius called his guards. ``I suppose they are to slay me,'' said Constantine. ``Certainly not,'' said Belisarius, ``but to force your armour-bearer to restore the daggers.'' But Constantine, believing that he was to die, drew his dagger and tried to stab Belisarius in the belly. Starting back, Belisarius seized Bessas and sheltered himself behind him, while Valerian and Ildiger dragged Constantine back. Then the guards came in, wrested the weapon from Constantine, and removed him. Some time afterwards he was put to death.

His execution was severely condemned by Procopius, who denounces it as the only impious act ever committed by Belisarius, and as an act out of keeping with his character, which was distinguished by fairness and leniency. This verdict is remarkable, for at no time would the capital penalty be considered an unjust severity in the case of an officer who attempted the life of his superior. But in his \emph{Secret History} Procopius supplements the story and thereby explains his condemnation of the act. If we may believe what he there relates, Constantine was sacrificed to the hatred of Antonina. The scandalous anecdote is that when Belisarius had discovered in Sicily his wife's disgraceful intrigue with Theodosius,\texthebrew{⁠}\footnote{\emph{H. A.} 1, p10 . It is to be noted that Procopius was not in Sicily when this scandal occurred. He was in Africa. See above,
p143 .

} Constantine expressed his

p193
sympathy with the injured husband, and observed, ``If it were my case, I would have slain the woman and not the young man.'' The words were reported to Antonina, who bided her time for revenge. The affair of Praesidius brought her the opportunity to punish Constantine for his offensive words. Her persuasions induced Belisarius to order the execution, and, according to Procopius, the Emperor was seriously displeased at the death of such a capable general.

Soon after this incident the truce was unequivocally broken by repeated endeavours of the Goths to steal secretly into Rome. They planned to gain an entrance through the aqueduct known as the

Aqua Virgo , near the Pincian Gate, but their explorations in the tunnel were revealed by the light of their torches. Another device was to drug the guards of a low section of the wall, on the north-western side of the city, with the help of two Romans, who were bribed. But one of them informed Belisarius and the scheme was frustrated. On another occasion the Goths openly attacked and were repelled. In retaliation for these acts Belisarius sent orders to John to descend upon the Picentine provinces. Some preparations for this eventuality had been made by the Goths. John was opposed by a force under Ulitheus, an uncle of the king, but the Romans were victorious, and Ulitheus was slain. This battle must have been fought somewhere in the southern province of Picenum,\texthebrew{⁠}\footnote{Picenum suburbicarium, of which the chief towns were Ancona, Auximum, Firmum, and Asculum.

} for John then marched to

Auximum

(Osimo). Finding that it had strong natural defences, he made no attempt to take it, but marched forward into the northern Picenum and reached

Urbinum . He judged that Urbinum, like Auximum, might be difficult to capture, and went on to

Ariminum . In leaving two fortresses held by the enemy in his rear, John disobeyed the express injunctions of his commander-in\texthebrew{‑}chief .\texthebrew{⁠}\footnote{Belisarius had ordered him to lay siege to any fortress that lay on his route, and if he failed to take it, not to advance farther. Procopius (II.10) justifies John's disobedience.

} But his disobedience had a useful result. He shrewdly foresaw that the seizure of Ariminum, which is only a day's march from Ravenna, would compel Witigis, fearing for the safety of the Gothic capital, to raise the siege of Rome. Ariminum offered no resistance, the garrison fled to Ravenna. John presently received a message from the Gothic queen. Matasuntha hated the husband to

p194
whom she had been united against her will, and now she impetuously proposed to betray Ravenna and to marry John, though he must have been completely a stranger to her.

When the news of the fall of Ariminum reached Rome, the Goths immediately burned the palisades of their camps and prepared to depart. Belisarius did not allow them to go unharmed. He waited till about half their host had crossed the Milvian Bridge and then attacked them with all his forces. Their losses were considerable. Besides those who were slain in combat many were drowned in the Tiber. Thus the siege of Rome, which had lasted for a year and nine days, came to an end about the middle of March, A.D. 538. It had furnished Witigis with an opportunity to demonstrate his incompetence, and Belisarius to display his resource ­fulness.

Small as his forces were, Belisarius seems throughout to have been sanguine that he would be able to overcome the resistance of the Goths. It had been, and was to be, a war of sieges; if the enemy had met him in the open field, after the arrival of the reinforcements, it is possible that he would have won a decisive victory, and the conquest of Italy might have been achieved almost as rapidly as the conquest of Africa. He was asked during the siege of Rome how it was that he was so confident, seeing the disparity in strength between the army of the enemy and his own. His reply was that he relied on the superiority of his tactic. ``Ever since we first met the Goths,'' he said,\texthebrew{⁠}\footnote{Procopius, \emph{B. G.} I.27.26.} ``in small engagements, I studied the differences in our tactical methods for the purpose of adapting my tactic so as to make up for the inferiority of my numbers. I found that the chief difference is that almost all our Roman troops and our Hunnic allies are excellent horse-archers , whereas the Goths are totally unpractised in this force of warfare. Their cavalry are accustomed to use only lances and swords, while their bowmen are unmounted and go into battle under the cover of their heavy armed cavalry. And so, except in hand-to\texthebrew{‑}hand fighting, their cavalry have no means of protecting themselves against the missiles of the enemy and can easily be cut up, and their infantry are ineffectual against mounted forces.'' But no tactic, however able, would have succeeded against the Goths,

p195
who were brave and well disciplined, if their army had been as vast as that which the historian alleges Witigis led against Rome.

\xchapter{Chapter XVIII, Part D}

\xsubsection{(Part 4 of 5)}

After the raising of the siege of Rome the scene of war shifts northward, to the fortresses along the Flaminian Way, in the lands of Umbria and Picenum, and to the provinces beyond the Po, where fighting was still to go on for two years before Belisarius succeeded in capturing the Gothic capital.

The Flaminian Way, which, traversing the Apennines, connected Rome with Ravenna, reached the Hadriatic at Fanum Fortunae (Fano), whence, following the coast, it led to Ariminum, and was continued to Ravenna. The general disposition of the belligerent forces in these districts is easy to grasp. The principal fortified hill towns to the west of the Flaminian Way, with the exception of Perusia, were held by the Goths, and those to the east, with the exception of Auximum (Osimo), by the Romans.\texthebrew{⁠}\footnote{Cp. Hodgkin, \emph{op. cit.} iv.288 . Urbs Vetus (Orvieto) was held by 1000 Goths, Tuder (Todi) by 400, Clusium (Chiusi) by 1000, Urbinum by 2000, Mona Feletris (Montefeltro) by 500, Caesena by 500, Auximum by 4000. Narnia, Spoletium, and Firmum were in the hands of the Romans. The distance from Rome to Ravenna by the Via Flaminia is 370 kils.; from Fanum to Ravenna about 96.

} Ariminum, as we saw, had been somewhat audaciously occupied by John, the nephew of Vitalian, with 2000 Isaurians, and

Ancona

was securely held by Conon.

It appeared to Belisarius that it would be a serious error to keep 2000 excellent cavalry, who would be invaluable in open warfare, shut up in Ariminum, and only tempting the Goths to besiege it. Accordingly, as soon as the enemy retired from Rome, his first care was to send forward Martin and Ildiger at the head of 1000 horsemen to order John to withdraw from Ariminum, and replace his Isaurians by a small force of infantry taken from the garrison at Ancona, which could easily spare them. As the retreating army of Witigis had diverged from the Flaminian highroad in order to avoid the forts of Narnia and Spoletium, no obstacle opposed the advance of Martin and Ildiger until they reached Petra Pertusa, ``The tunnelled Rock,'' a pass between Cales (Cagli) and Forum Sempronii (Fossombrone),

p196
about twenty-five miles from the Hadriatic Sea.\texthebrew{⁠}\footnote{It is called Intercisa in the Itineraries, and is a little to the east of the modern Acqualagna. The pass is 125 feet long, over 17 feet high, and broad. On either side are the halves of a mountain: on right Monte Paganuccio, 3259 feet high, on left Monte Pietralata, 2960 feet high.

}
This pass, now known as the Passo di Furlo, is accurately described by Procopius. The Flaminian Road comes up against a high wall of rock, on the right of which a river descends with such a rapid current that it would be death to attempt to cross it, and on the left the precipitous cliff to which the rock belongs rises so high that men standing on its summit would appear to those below like the smallest birds. The Emperor Vespasian bored a tunnel through this rock, as an inscription on the spot records. It was a natural fortress, well adapted for defence. The Roman troops who now advanced found it held by a Gothic garrison and closed by doors at either end. The Goths, who had their women and children with them, lived in houses outside the tunnel, apparently on the Hadriatic side. When it was found impossible to make any impression on the well-fortified entrance to the passage, some men were sent to the top of the cliff, and dislodging huge fragments of rock they rolled them down on the Gothic block-houses below. The enemy immediately surrendered, and Martin and Ildiger, leaving a small garrison behind them, continued their journey to Fanum. From here they had to ride southward to Ancona to pick up a detachment of foot-soldiers to replace the Isaurians at Ariminum. Then retra­cing their steps to Fanum they arrived safely at their destination, and delivered the commands of Belisarius to John. But John declined to obey, and leaving the foot-soldiers with him Martin and Ildiger departed to report the issue of their errand to the commander-in\texthebrew{‑}chief .

The insubordination of John strikes the note of the subsequent course of the Roman conduct of the war. Counsels were divided, and the commander-in\texthebrew{‑}chief could no longer depend on his generals to conform to his plans. Belisarius was slow and cautious, but it is probable that, if he had been able to have his own way and secure the punctual obedience of his subordinates, the war would have been shortened. John was an excellent but sometimes over-confident soldier. He was impatient of the cautious deliberation of Belisarius, and doubtless thought that he was himself more worthy of the post of supreme commander.

p197
In the present instance, the event speedily showed that Belisarius was right. Witigis had no sooner crossed the Apennines than he addressed himself to the siege of Ariminum. Failing in his assaults, he sat down to take it by hunger, and the besieged were presently reduced to extreme distress (April, A.D. 538).

Belisarius meanwhile had begun to advance northward from Rome, to carry out methodically his plan of redu­cing, first of all, the Gothic fortresses west of the Apennines. It was about the middle of the year. Clusium and Tuder surrendered on his approach. His next object would have been the reduction of Urbs Vetus, but the execution of his plan was disarranged by the arrival of reinforcements from the East which now reached Picenum under the command of the eunuch Narses, keeper of the Emperor's privy purse.\texthebrew{⁠}\footnote{They probably landed at Ancona. Shortly before this Witigis had sent a force against it. Conon, the commander, gave the Goths battle outside the walls and was severely defeated; but the fortress was not taken.

} The new army was 7000 strong, consisting of 5000 Roman troops under another Narses and Justin the Master of Soldiers in Illyricum, and 2000 Herul auxiliaries under their own leaders. Such an important addition to the Imperial fighting forces modified the situation, and Belisarius, leaving Urbs Vetus unreduced, marched to Picenum to confer with Narses and arrange the future conduct of the war. They met at Firmum and a council was held which had weighty consequences. The urgent question was the relief of Ariminum, which was hard pressed and might be forced to surrender through hunger. Should the army march to its relief immediately? Belisarius was opposed to this course on military grounds. So long as Auximum was held by the enemy, an advance against the Goths at Ariminum would expose his rear to an attack from the garrison of that fortress. The majority of the generals present agreed, and held that no risks should be taken to save John, whose predicament was due to his own rashness and insubordination. Narses, who was a personal friend of John, opposed this view. He pointed out that the disobedience of John was a side-issue which ought not to affect their decision. After the relief of Ariminum John could be punished for defying the commands of Belisarius. But it would be highly inexpedient, he argued, considering not only the material loss, but also the moral consequences, to allow an important city and a large

p198
body of troops \texthebrew{—} not to speak of a vigorous general \texthebrew{—} to fall into the hands of the foe.

While the council was sitting, a soldier from Ariminum who had eluded the blockade arrived in the camp with a letter from John. Its purport was: ``All our supplies have long since failed us. Unable to resist the enemy, we cannot hold out against the pressure of the inhabitants, and within seven days we shall have reluctantly to surrender ourselves and the city. Our extreme necessity is, I think, an adequate excuse for an act which may appear unbecoming.'' This message, simply announcing a fact and making no demand for succour, strengthened whatever effect may have been produced by the arguments of Narses. Belisarius decided to do all that could be done to save Ariminum, though he still felt grave scruples whether it was a wise thing to do.

It would be bold for a modern critic, with the meagre evidence at his disposal, to assert that the hesitations of the commander-in\texthebrew{‑}chief were unjustified, but it is difficult to resist the impression that the course recommended by Narses was the right one. It required military skill, but when Belisarius set his mind to the problem he solved it triumphantly. In order to mitigate the danger from Auximum, he posted a thousand men to the east of it near the coast. A large force was sent by sea to Ariminum under the command of Ildiger, who was instructed not to disembark until a second army, which, led by Martin, was to march along the coast road, approached the city. Martin, when he arrived, was to light many more fires than were required, in order to deceive the enemy as to the number of his troops. Belisarius, accompanied by Narses, led the rest of the army by an inland mountainous route with the purpose of descending on Ariminum from the north-west.\texthebrew{⁠}\footnote{The only indication that Procopius gives of the route is that they passed Urbs Salvia (Urbesaglia), º which is half-way between Fermo and Nocera. They must of course have crossed the Flaminian Way. Hodgkin's view that they followed the Flaminian Way from Nocera, so that from Fano onward their route would have coincided with that of Martin, is entirely incompatible with the narrative of Procopius. There can, I think, be little doubt that they left the Flaminian Way at Scheggia or at Acqualagna and followed one of the routes noticed below, pp288-289. Compare the retreat of Garibaldi in 1849, via San Marino to Musano (due west of Rimini). See G. M. Trevelyan, \emph{Garibaldi's Defence of the Roman Republic}, chaps. xiii, xiv.

} For the full success of the plan it was necessary that the arrivals of the three armies on the

p199
scene should be timed to coincide. At a day's journey from Ariminum a few Goths fell in with the army of Belisarius, and hardly realised that they were in the presence of an enemy till Roman arrows began to work havoc among them. Some fell, others crawled wounded behind the shelter of rocks. From their concealment they could see the standards of Belisarius, and they received the impression of an army far in excess of its actual numbers. In the night they made their way to the camp of Witigis at Ariminum, and arriving at mid-day reported the approach of Belisarius with an innumerable host. The Goths immediately formed in battle order on the northern side of the city and spent the afternoon looking towards the hills. When night fell and they were composing themselves to rest, they suddenly saw to the south-east the blaze of the fires which had been kindled by the troops of Martin. They realised that they were in danger of being surrounded, and passed the night in terror. When morning came and they looked out to sea, they beheld a great armament of hostile ships approaching. In fear and confusion they broke up their camp, and no man thought of anything but reaching the shelter of Ravenna. If the garrison of the city had rushed out and dealt death among the panic-stricken fugitives, Procopius thought that the war might have ended there and then. But the soldiers of John were too exhausted by their privations to seize the moment.

Ildiger and the troops who had come by sea were the first to arrive in the abandoned camp of the barbarians. Belisarius arrived at mid-day. When he met John, pale and gaunt with hunger, he could not forbear remarking that he ought to thank Ildiger. John dryly replied that his gratitude was due not to Ildiger but to Narses.

The relief of Ariminum, accomplished without the loss of a single life, was a new proof of the military capacity of Belisarius, but it was a moral triumph for Narses, since but for his influence it would never have been undertaken. Distrust and division ensued between the commander-in\texthebrew{‑}chief and the chamberlain, and the bloodless victory hardly compensated for the injuries which this dissension inflicted on the Imperial cause. Narses

p200
felt, and his friends convinced him, that it was beneath the dignity of his office to act in subordination to a general, and he determined to use the forces which he had brought to Italy according to his own discretion. In accordance with this resolution he excused himself repeatedly from complying with requests or orders from Belisarius, who at length convoked a military council to clear up the situation.

At this council Belisarius did not at first insist upon his rights as commander-in\texthebrew{‑}chief or rebuke Narses for disobedience. He pointed out that the enemy were far from being defeated; Witigis had still an army of tens of thousands at Ravenna; the situation in Liguria was serious; Auximum with its large and valiant garrison was still uncaptured, as well as other strong places like Urbs Vetus. He proposed that a portion of the army should be sent to Liguria, to the rescue of Milan, which was in grave peril, and that the remaining forces should be employed against the Gothic fortresses south and west of the Flaminian Way, and first of all against Auximum. Narses replied. He contended that it was inexpedient that all the Imperial forces should be concentrated on the two objects of Auximum and Milan. Let Belisarius undertake these enterprises, but he would attempt the conquest of the Aemilian province. This would have the probable advantage of retaining the main army of the Goths at Ravenna, so that they would be unable to send aid to the places attacked by Belisarius. But Belisarius was opposed to any plan which involved a dissipation of forces, and he decided to assert his authority. He produced a letter which the Emperor had recently addressed to the commanders of the troops in Italy. It was conceived in these terms:

``In sending Narses our purser to Italy we do not invest him with the command of the army. It is our wish that Belisarius alone shall lead the whole army as seems good to him, and it behoves you all to obey him \emph{in the interest of our State.}''\footnote{\emph{B. G.} II.18.28 \textgreek{α\texthebrew{ὐ}τ\texthebrew{ῷ} τε \texthebrew{ὑ}μ\texthebrew{ᾶ}ς \texthebrew{ἔ}πρεσθαι \texthebrew{ἅ}παντας \texthebrew{ἐ}π\texthebrew{ὶ} τ\texthebrew{ῷ} συμφέροντι τ\texthebrew{ῇ ἡ}μετέρ\texthebrew{ᾳ} πολιτεί\texthebrew{ᾳ} προσήκει.}}

In the last phrase there was a possible ambiguity of which Narses at once took advantage, interpreting it as a reservation limiting the duty of obedience. The plan of Belisarius, he said, is not in the interest of the State, and therefore we are not

p201
bound to obey him. It may seem difficult to suppose that Justinian intended to lay down a principle which logically led to military anarchy, since it was open to every commander to take a different view of the wisdom of a strategy plan. Yet we cannot consider it impossible that the insertion of the words ``in the interest of the State'' was designed as a check on the authority of the commander-in\texthebrew{‑}chief . For if the Emperor had really meant to enjoin unconditional obedience, the phrase in question was entirely unnecessary. The fact that the trusted keeper of his privy purse should have been chosen for a military mission lends colour to the suspicion that Justinian was dissatisfied with the progress of the war, and doubtful whether Belisarius was conducting it with the necessary energy. It would be going too far to suggest that he wished to deprive Belisarius of the undivided glory of conquering Italy, though we are told that this was the personal object of Narses.

Belisarius was not in a position to enforce his claims, and he had sufficient self-restraint to avoid an actual breach. Matters were smoothed over for the time, and the co-operation of the commanders, though it was far from cordial, continued. A large force was despatched against Urbs Vetus, and Belisarius, again postponing his intention of redu­cing Auximum, marched to the siege of Urbinum, accompanied by Narses and John. But the forces of the rival commanders did not mingle; they encamped separately on the eastern and western sides of the city. The garrison of Urbinum, which is situated on a high hill at a strenuous day's journey from Ariminum, refused an invitation of Belisarius to surrender; they had abundance of provisions and trusted in the strength of the city. Narses, deeming the place impregnable, considered it waste of time to remain, and, withdrawing to Ariminum, sent John, at the head of all his forces, against Caesena. Failing to take this place, John, who was impatient of sieges, advanced against Forum Cornelii (Imola), which he captured by surprise, and then easily subjugated the whole Aemilian province.

Meanwhile fortune played into the hands of Belisarius. Urbinum was supplied by a single spring. It suddenly ran dry, and deprived of water the Goths could only capitulate. Narses is said to have received the news of this success with deep chagrin.

It was now December ( A.D. 538)\texthebrew{⁠}\footnote{\emph{B. G.} II.20.1.

} and Belisarius decided that it was inopportune then to attempt the siege of Auximum, which promised to prove a difficult enterprise. He left a large force in Firmum to protect the country against the ravages of the garrison, and marched himself to Urbs Vetus, where provisions were already running short. The place could hardly have been taken by assault. It is a natural stronghold, requiring no artificial fortifications, \texthebrew{—} built on an isolated hill rising out of hollow country. This hill, level at the top, is precipitous below, and is surrounded by cliffs of the same height, between which and the hill itself flows a large and impassable river, according to Procopius, entirely encircling the hill except at one point where the city could be approached from the cliffs. At the present day, Orvieto is not surrounded by water. The river Paglia flows round the northern and eastern sides of the hill, to join the river Chiana, but on the south and west there is no such natural moat. It is supposed that the Paglia may have changed its course.\texthebrew{⁠}\footnote{Cp. Hodgkin, iv p338 .

} Hunger was the only weapon which could avail against a brave garrison, and the Goths, when they had been reduced to consuming hides softened in water, surrendered at last to Belisarius (spring, A.D. 539).

In the meantime important events had been happening beyond the Po. Immediately after the Goths had raised the siege of Rome, Belisarius, in fulfilment of his promise to Datius, the archbishop of Milan, had sent 1000 Isaurians and Thracians under the command of Mundilas to Liguria (April, A.D. 538). They went by sea from Porto to Genoa, and, crossing the Po, they succeeded in occupying Milan, Bergamum, Comum, Novaria, and all the strong places of inland Liguria except

Ticinum

(Pavia). On hearing the news Witigis sent his nephew Uraias to recover Milan, and he received power ­ful aid from abroad.

Theodebert, grandson of Chlodwig, had succeeded his father Theoderic as king of Austrasia in A.D. 533. Besides the Austrasian dominion on both sides of the Rhine, with its capital at Metz, he ruled over a portion of Aquitania and a portion of Burgundy which had recently been conquered by his uncles. We

p203
possess a letter which he wrote to Justinian, probably in an early stage of the war, offering excuses for his failure to send to Italy a force of 3000 men which he had promised. As he styles Justinian ``father,''\texthebrew{⁠}\footnote{D. illustro et praecellentissimo domno et patri Iustiniano imperatore Theodebertus rex (\emph{Epp. Merow. et Kar. aevi} i \emph{Epp. Austrasicae}, 19). The soldiers were to have been sent to the help of the patrician Bregantinus (perhaps = Vergentinus, below, p205, n1 )

} it may be inferred that the Emperor formally adopted him as a son when he sought an assurance of the co-operation of the Franks before the outbreak of the war. But Theodebert was ambitious and treacherous, and his filial relation to Justinian was no obstacle to his policy of playing fast and loose between the two belligerents. At this crisis he resolved to assist the Goths, and 10,000 Burgundians crossed the Alps to co-operate with Uraias. He sophistically professed that he was not violating his convention with the Emperor, because no Franks were in the army; the Burgundians, forsooth, were acting as an independent people, without his authority.\texthebrew{⁠}\footnote{Procopius, \emph{B. G.} II.12.38.

} The Gothic and Burgundian forces blockaded Milan, which Mundilas held with only 300 soldiers as the rest of his force had been distributed in the other Ligurian fortresses.\texthebrew{⁠}\footnote{Marcellinus, \emph{sub} 539, records that Theodebert devastated the Aemilian province and plundered Genoa.

} The able-bodied civilian inhabitants were therefore called upon to take part in the defence.

After the relief of Ariminum, Belisarius despatched a large army under Martin and Uliaris to the relief of Milan. These commanders encamped on the southern bank of the Po; they were afraid to face the host of barbarians who were besieging the city. Mundilas despatched a messenger, who managed to evade the sentinels of the enemy, to plead the urgent need of the besieged, and was sent back with promises of speedy aid, which Martin and Uliaris made no effort to fulfil. At last, after a delay so long that it amounted to treason to the Imperial cause, they wrote to Belisarius, representing their forces as hopelessly inadequate to cope with the enemy and requesting him to send John and Justin, who were in the neighbouring province of Aemilia, to reinforce them. Belisarius complied, but John and Justin refused to move without the authority of Narses. Belisarius wrote to Narses, who gave the requisite order. John proceeded to collect ships for the purpose of crossing the Po,

p204
but before his preparations were completed he fell ill. Thus delay ensued upon delay, and meanwhile the unhappy inhabitants of Milan were starving. When they were reduced to feeding on dogs and mice, Gothic envoys waited on Mundilas, inviting him to capitulate on the condition that he and all his soldiers should have their lives spared. He was ready to accept these terms if they would agree to spare the inhabitants. But the Goths, who were infuriated against the disloyal Ligurians, did not conceal their determination to wreak a bloody vengeance. Mundilas therefore refused, but his hands were soon forced. He attempted to induce the soldiers to make a desperate sally against the foe, but, worn as they were by the sufferings of the siege, they had not the courage to embrace so forlorn a hope. They compelled their leader to agree to the terms which the Goths had proposed.

Mundilas and the soldiers were placed in honourable captivity, in accordance with the agreement. Milan and its inhabitants felt the full fury of a host of savages. All the adult males, who according to Procopius numbered 300,000, were massacred; all the women were presented as slaves to the Burgundians. The city itself was razed to the ground. It was the wealthiest and most populous town in Italy then, as now, and if Procopius is near the truth in his estimate of the number of males who were slain, it must have been nearly as populous as it is to\texthebrew{‑}day.\footnote{\emph{B. G.} II.21.39. The population of modern Milan is between 600,000 and 700,000 (that of Rome is over 500,000). Procopius describes it (\emph{B. G.} II.7.38) as the most populous Italian city next to Rome. It seems probable that he has immensely exaggerated the number of the slain. For the massacre see also \emph{Cont. Marcell.}, \emph{s.a.}, and Marius of Aventicum, \emph{Chron.}, \emph{s.a.} 538 (senatores et sacerdotes cum reliquis populis etiam in ipsa sacrosancta loca interfecti sunt). Datius the archbishop escaped.

}

In the long series of deliberate inhumanities recorded in the annals of mankind, the colossal massacre of Milan is one of the most flagrant. Historians have passed it over somewhat lightly. But the career of Attila offers no act of war so savage as this vengeance, carried out by the orders of the nephew of the Gothic king. It gives us the true measure of the instincts of the Ostrogoths, claimed by some to have been the most promising of the German invaders of the Empire.

Reparatus, the Praetorian Prefect of Italy, was found in the city. He was the brother of Pope Vigilius, but this did not save him. He was cut in pieces and thrown to the dogs.

p205
Cerventinus,\texthebrew{⁠}\footnote{Procopius calls him Vergentinos, \emph{B. G.} II.21.41.

} another brother, escaped to Dalmatia, and went on to Constantinople to announce the calamity to Justinian. The fall of Milan, which happened towards the end of March\texthebrew{⁠}\footnote{\emph{B. G.} II.22.1 \textgreek{\texthebrew{ὁ} χειμ\texthebrew{ὼ}ν \texthebrew{ἐ}τελεύτα \texthebrew{ἤ}δη.}} ( A.D. 539), led to the immediate recovery of all Liguria by the Goths. The news came as a heavy blow to Belisarius, but it was an irresistible proof of the unwisdom of divided military authority by which the Emperor himself could not fail to be impressed. Belisarius wrote to him explaining all the circumstances and showing where the blame rested. Justinian inflicted no punishment on those who were in fault,\texthebrew{⁠}\footnote{Belisarius would not allow Uliaris to appear in his presence, \emph{ib.} ii.

} but he immediately recalled Narses, and in language which was not ambiguous confirmed the supreme authority of Belisarius.\footnote{The Heruls, who had come to Italy with Narses, refused to remain when he was recalled.

}

In the meanwhile Witigis, while the fate of Liguria still remained undecided, was seriously alarmed for the safety of Ravenna. The Romans were firmly established at Ariminum and urbinum, and he expected that at any moment Belisarius might advance against his capital. Early in the year he resolved to seek foreign help. He first applied to Wacho, king of the Langobardi, who dwelled beyond the Danube. But no succour was forthcoming from this quarter. Wacho, who was an ally of the Emperor, did not consider it expedient to imitate the double-dealing of the treacherous Franks. The Goths then conceived the idea of appealing to a greater power, the king of Persia himself. They argued, with truth, that Justinian would not have embarked on his enterprises in the West if he had not been secured in the East by the peace which he had concluded with Chosroes. If they could succeed in embroiling him in a war with Persia, it would be impossible for him to continue the war in Italy. As the only practicable route to Persia lay through Imperial territory, it would have been difficult to sent Gothic ambassadors. By large bribes two Ligurian priests, who could travel without exciting suspicion, were induced to undertake the mission, and they succeeded in reaching the court of

p206
Chosroes and delivering a letter from Witigis. This appeal was hardly the chief motive which determined Chosroes to reopen hostilities, but undoubtedly it produced its effect. It must have impressed upon him that in considering his foreign policy it would be wise to take account of the situation in the western Mediterranean. He resolved on war, as Witigis hoped, but his operations began too late to rescue Witigis from disaster.

The report that negotiations were passing between Ravenna and Ctesiphon reached Justinian (in June), and inclined him to the idea of ending the Italian war by a compromise as soon as possible, so as to set Belisarius free to take command on the eastern frontier. He accordingly released the Gothic envoys whom he had detained for more than a year,\texthebrew{⁠}\footnote{They had been sent to Constantinople during the siege of Rome. See above,
p191 .

} and promised to send ambassadors of his own to discuss peace. When these Goths arrived in Italy, Belisarius would not allow them to proceed to Ravenna till Witigis surrendered the Roman envoys, Peter and Athanasius, who had been held prisoners for four years.\texthebrew{⁠}\footnote{We may infer from \emph{B. G.} II.22.25 that this happened at the end of June or beginning of July 539. Körbs (\emph{op. cit.} 31 \emph{sqq.}) calculates that the Gothic embassy to Ctesiphon started in March from Ravenna. They could reach the Bosphorus in about four weeks, if they travelled quickly at the rate of 75 kils. a day. To reach Ctesiphon meant five or six weeks more, so that they might have arrived about the middle of May. Thus Justinian could have heard of their arrival and released the Gothic envoys before the end of June. For the distances of Ravenna from Constantinople via Aquileia see below,
p225 .

} The Emperor rewarded these men for their services by creating Peter Master of Offices, and Athanasius Praetorian Prefect of Italy.

Italy indeed needed peace. Agriculture had ceased in the provinces devastated by war, and in Liguria and Aemilia, in Etruria, Umbria, and Picenum the inhabitants were dying of hunger and disease. It was said that in Picenum alone 50,000 tillers of the soil perished. Procopius noted the emaciation, the livid colour, and the wild eyes of the people, suffering either from want of food, or from a surfeit of indigestible substitutes like acorn bread. Cannibalism occurred, and a ghastly story was told of two women who lived in a lonely house near Ariminum where they offered a night's lodging to passers-by . They killed seventeen of these guests in their sleep and devoured their flesh. The eighteenth woke up as the cannibals were about to despatch him; he forced them to confess, and slew them. Scattered

p207
over the country-sides were the unburied corpses of those who had died while they sought with feeble hands to tear blades of grass from the ground. The Imperial armies suffered little, for they received a constant supply of provisions by sea from Calabria and Sicily.

Belisarius in the meantime prosecuted his plans. He considered it essential to capture Auximum and Faesulae before he advanced upon Ravenna. Pla­cing Martin and John at Dertona (Tortona) to defend the line of the Po against Uraias, he sent Justin and Cyprian to blockade Faesulae, and undertook himself the most important of his tasks, the siege of Auximum. These two sieges occupied more than six months (April to October or November, A.D. 539).\footnote{\emph{Cont. Marcell.}, \emph{s.a.}, states that these cities were taken in 539 septimo mense, but erroneously records their capture before the fall of Milan. The sieges must have begun in April or May, cp. \emph{B. G.} XXII.1, and so the seventh month would be October-November.

}

The army on the Po succeeded, as Belisarius anticipated, in hindering Uraias from marching to the aid of Faesulae. The two hosts, reluctant to risk a trial of strength, remained immobile on the banks of the river, till a new enemy appeared upon the scene. The Franks regarded the calamities of Italy as an opportunity for themselves and were as perfidious towards the Goths as towards the Empire; and Theodebert himself, at the head (it is said) of 100,000 men, descended from the Alps for the plunder and destruction alike of Goths and Imperialists, with both of whom they had recently sworn alliance. Procopius describes their equipment. There were a few mounted spearmen in attendance on the king, the rest were infantry armed with a sword, a shield, and an axe. At the first onset of battle a shower of axes fell upon the foes, shattering shields and killing men.

The Goths, fondly imagining that Theodebert was coming to their aid in fulfilment of his promises, rejoiced to hear of his approach. At Ticinum, where a bridge spanned the Po, at its confluence with the stream of Ticinus, the Goths who guarded it gave the Franks every assistance to cross the river. As soon as they held the bridge, the invaders threw off the mask. They seized the women and children of the Goths, slaughtered them, and threw their bodies into the river. Procopius saw a religious

p208
significance in this act. ``These barbarians,'' he says, ``though converted to Christianity retain most of their old beliefs and still practise human sacrifices.'' Having crossed the Po, the Franks advanced southward towards Dertona, near which Uraias and his army were encamped not far from the Roman camp. The Goths went forth to welcome their allies, and were received by a shower of axes. They turned in headlong flight, rushed wildly through the camp of the astonished Romans, and pursued the road to Ravenna. The Romans imagined that Belisarius must have suddenly arrived and surprised the Gothic camp; and issuing forth to meet him they found themselves confronted by the immense army of the Franks. They were forced to fight, but were easily routed and retired to Tuscany.

The victors were in possession of two deserted camps, supplied, however, with provisions. The food did not go far among so many, and in the desolated country they found no subsistence but oxen and the water of the Po. It is satisfactory to know that they paid a heavy price for their rapacity. Dysentery broke out, and large numbers \texthebrew{—} a third of the host, it was reported \texthebrew{—} died. The survivors were bitter against their king for leading them into a place of desolation to perish of hunger and disease. Then a letter arrived from Belisarius, reproaching Theodebert for his treachery, mena­cing him with the anger of the Emperor, and advising him to attend to his domestic affairs instead of running into danger by interfering in matters which did not concern him. The barbarians retreated ingloriously across the Alps.\footnote{The retreat gave Justinian the pretext for assuming the title Francicus.

}

This episode had little influence on the course of the war. All the efforts of the Imperial forces were concentrated on the blockades of Auximum and Faesulae. The flower of the Gothic army was holding Auximum and was resolved to hold it to the end. When the provisions began to give out, the commander of the garrison sent an urgent message to Ravenna, imploring Witigis to send an army to relieve them. Immediate help was promised, but nothing was done. Time wore on, the garrison was sorely pressed by hunger, and too careful a watch was kept to allow any one to steal out of the town. But the Goths managed to bribe a soldier named Burcentius, who was keeping guard at mid-day in an isolated spot near the walls, to carry a

p209
letter to Ravenna. Burcentius executed the errand, and Witigis again sent back good words which were read aloud to the garrison, and encouraged them to hold out. As no help came, they again employed the services of the traitor and informed the king that they would be compelled to surrender within five days.

Belisarius meanwhile had repeatedly urged them to surrender on favourable terms, and knowing that they were starving he was puzzled at their refusal to comply. A Slavonic soldier, hidden in a bush for the purpose, succeeded in capturing alive a Goth who had crept out of the city at dawn to gather grass. The prisoner disclosed the treachery of Burcentius, and Belisarius delivered him to his comrades to do with him what they would. They burned him alive in sight of the walls.

The chief water-supply of Auximum was derived from a huge cistern, built in a rocky place outside the walls, so that men had to come out of the city in order to fill their water jars. Belisarius sent some Isaurians to attempt to destroy the cistern, but the masonry resisted all their efforts. Then he poisoned the spring by throwing in quicklime with dead animals and noxious herbs. But there was another small well inside the city, and, though sadly insufficient for their needs, it enabled the loyal Goths to postpone surrender.

The end was brought about by the capitulation of the starving defenders of Faesulae. The captives were brought to Auximum and paraded in front of the walls, and this sight determined the garrison, convinced at last that they had nothing to hope from Ravenna, to follow the example of Faesulae. The terms arranged were that they should give up half of their possessions to be divided among the besiegers, and should pass into the service of the Emperor. They cannot be reproached for having accepted these conditions. Their king had basely left them to their fate. He had professed to regard Auximum as the key of Ravenna, but such was his cowardice that he could not bring himself to send to its relief any portion of the considerable army which was idly protecting his capital.

Auximum fell in October or November and Belisarius lost no time in preparing an advance upon Ravenna. New forces

p210
had just arrived from Dalmatia, and these he ordered to guard the northern bank of the Po, while another contingent was sent to patrol the southern. The purpose of these dispositions was to prevent food stores from being sent down the river from Liguria. The Imperial command of the sea effectively hindered any attempts to supply the city from elsewhere.

The one thing that Belisarius had now to fear was that the Franks might again descend into Italy and again aid the Goths as they had aided them at Milan. When he learned that a Frank embassy was coming to Ravenna, he sent ambassadors to Witigis. The Frank proposal was that Goths and Franks should make common cause, and, when they had driven the Roman invaders from Italy, should divide the peninsula between them. The Imperial envoys warned the Goths against entertaining the insidious offer of a people whose rapacity was only equalled by their treachery. Their rapacity was proved by the way they had dealt with the Burgundians and Thuringians; their treachery the Goths knew to their own cost by the events of a few months ago. Witigis and his counsellors decided that it would be wiser to come to terms with the Emperor than to trust such a dangerous ally as Theodebert, and the Frank envoys were sent empty away. Hostilities were suspended, and negotiations opened with Belisarius, who, however, did not relax his precautions against the introduction of provisions into Ravenna. He even bribed some one to set fire to the public corn º store in the city \texthebrew{—} at the secret suggestion, it was said, of the queen Matasuntha. Some of the Goths ascribed the conflagration to treachery, others to lightning; the one theory suggested enemies among themselves, the other an enemy in heaven.

Uraias in the meantime was preparing to come to the aid of his uncle with 4000 men, most of whom he had taken from the garrisons which held the forts of the Cottian Alps. But John and Martin hurried westward, seized the forts, and captured the wives and children of the Goths. On hearing that their families were in the hands of the Romans, the soldiers of Uraias deserted him and went over to John; and Uraias was forced to remain inactive in Liguria.

Two senators, Domnicus and Maximin, now arrived from Constantinople, bearing the Emperor's instructions for the

p211
conclusion of peace. The menace of Persia inclined Justinian to grant more lenient terms than the military situation seemed to warrant. He proposed a territorial division of Italy. All the lands north of the Po should be retained by Witigis, all the lands south of the Po should be retained by the Emperor. The royal treasury of Ravenna should be divided equally between the two contracting powers. Witigis and the Goths were surprised by a proposal which was far more favourable than they had looked for, and they accepted it without hesitation. But it did not please Belisarius. He saw within his reach a complete victory to compensate for the toils and anxieties of five weary years. He had dethroned and led captive the king of the Vandals; he was determined to dethrone and lead captive the king of the Ostrogoths. When the ambassadors returned from Ravenna to his camp and asked him to ratify by his signature the treaty of peace, he declined. As his refusal was severely criticised by some of the generals as an act of disobedience to the emperors's decision, he summoned a military council and asked those present whether they approved of the division of Italy, or whether they deemed it practicable to conquer it entirely. All the officers were unanimous in approving the terms which the Emperor had dictated, and Belisarius required them to put in writing their opinion that nothing would be gained by continuing the war, so that he should be exonerated from blame if future events were to prove that it would have been wiser to carry to completion the overthrow of the Gothic kingdom.\footnote{Procopius curiously says that Belisarius was pleased with the opinion of the generals (\emph{B. G.} II.29.16 \textgreek{\texthebrew{ἡ}σθείς}), though he had just recorded (\emph{ib.} iv) his vexation at the conditions prescribed by Justinian.

}

But the refusal of the commander-in\texthebrew{‑}chief to sign the treaty had already produced an unfavourable impression at Ravenna. Witigis suspected that the negotiations were a trap, and refused to execute the agreement unless Belisarius signed it and gave them a sworn guarantee of good faith. Famine meanwhile was doing its work, and the discontent of the Goths with their incompetent king reached a climax.

Then a remarkable idea occurred \texthebrew{—} we do not know from what quarter the suggestion came \texthebrew{—} to the men of weight among the Goths. Why not revert to the political condition of Italy

p212
as it existed before the days of Theoderic, before the days of Odovacar? The regime of Witigis had discredited Ostrogothic royalty, and they would feel no repugnance to submitting to the direct authority of a western Emperor\texthebrew{⁠}\footnote{\emph{B. G.} II.29.18 \textgreek{βασιλέα τ\texthebrew{ῆ}ς \texthebrew{ἐ}σπερίας Βελισάριον \texthebrew{ἀ}νειπε\texthebrew{ῖ}ν \texthebrew{ἔ}γνωσαν.}} residing at Rome or Ravenna, if that Emperor were Belisarius, whom they deeply respected both as a soldier and as a just man. They entertained no doubts that he would eagerly accept the offer of a throne. They did not know his uncompromising loyalty or suspect that there was no rôle that seemed more thoroughly detestable to him than the rôle of a usurper. He had once taken an express and solemn oath that he would never aspire to the throne so long as Justinian was alive. But when messengers of the Goths privately sounded him on the plan he professed to welcome it with pleasure. For he saw in it a means of bringing his work to a speedy and triumphant conclusion. When these clandestine negotiations came to the knowledge of Witigis, he resigned himself to the situation and sent a secret message to Belisarius urging him to accept the offer.\footnote{Martroye (\emph{L'Occident}, p401) argues that this intrigue had been arranged in consultation with Chosroes. His reason is that in autumn 539 the Armenian envoys who urged Chosroes to declare war pointed out that Justinian had lost his two best generals, Sittas who had been killed, and Belisarius who had left his service to be the sovran of Italy. Martroye is mistaken in supposing that the Gothic proposal was made to Belisarius at about the same time; this cannot have been earlier than Jan. or Feb. 540.

}

Belisarius then summoned a meeting of the generals and invited the presence of the two Imperial ambassadors. He asked them whether they would approve if, without striking another blow, he should succeed in recovering the whole of Italy, in taking captive Witigis, and seizing all his treasure. The assembly agreed that it would be a magnificent achievement, and urged him to accomplish it if he could. Having in this way protected himself against misinterpretation of his motive in pretending to yield to the Gothic proposal, he sent confidential messengers to Ravenna to announce his definite acceptance. Official envoys were sent back to the camp, nominally to continue the discussion of peace terms, but privately to receive from the commander pledges of his good faith. He gave them sworn pledges on all matters save his willingness to accept the purple; on that point he deferred his oath till he should stand in the presence of Witigis and the Gothic magnates. The envoys

p213
were satisfied; they could not imagine that he would reject the Imperial diadem.

He then made his arrangements for entering Ravenna. He dispersed a part of his army, under the command of those leaders who were ill-disposed towards himself \texthebrew{—} John, Aratius, his brother Narses, and Bessas \texthebrew{—} to various destinations, on the pretext that it was difficult to provide the requisite commissariat for the whole army in one place. He sent his fleet laden with corn and other foods to the port of Classis, to fill the starving mouths of Ravenna. Then he advanced with his army and entered the city in May A.D. 540.\texthebrew{⁠}\footnote{Agnillus, \emph{Lib. Pont.} p101.

} It is disappointing that the historian does not describe the scene in which Belisarius undeceived the Gothic king and nobles as to his intentions. We are only told that he kept Witigis in honourable captivity, and that he allowed all the Goths who lived in the cis-Padane provinces to return to their homes. He seized the treasures of the palace, but the Goths were allowed to retain all their private property, and plundering was strictly forbidden.

Most of the garrisons of the strong places north of the Po voluntarily surrendered,\texthebrew{⁠}\footnote{Tervisium (Treviso) and the forts of Venetia are expressly mentioned. Caesena, west of Ariminum, had held out till the fall of Ravenna.

} apparently under the impression that Italy was to be ruled by Belisarius. Ticinum, which was the headquarters of Uraias, and Verona, which was held by Ildibad, were the chief exceptions. When the Gothic notables of the northern provinces realised that Belisarius had made ``the great refusal'' and was about to return to Constantinople, they proceeded to Ticinum and urged Uraias to assume the royal insignia and place himself at their head to fight a desperate battle for freedom. Uraias was ready to fight, but he declined to step into the place of Witigis; the nephew of such an unlucky ruler would not, he declared, have the necessary prestige. He advised them to choose as their king Ildibad, a man of conspicuous energy and valour, and a nephew of Theudis the king of the Visigoths. Accordingly Ildibad, at the request of the Gothic leaders, came from Verona and suffered himself to be proclaimed king. But he persuaded his followers to make one more effort to induce Belisarius to recall his decision. A deputation waited on the commander, who was making his preparations to leave Ravenna.

p214
They upbraided him, with justice, for having broken faith. But reproaches and enticements produced no effect. Belisarius told them definitely that he would never assume the Imperial name in Justinian's lifetime. Soon afterwards he left the shores of Italy, taking with him the dethroned king and queen, many leading Goths, and the royal treasure.\footnote{Probably in June.

}

The impregnable fidelity of Belisarius to Justinian's throne, under a temptation which few men in his position would have resisted, is the fact which has been chiefly emphasised by historians in describing these tortuous transactions. But his innocence of criminal disloyalty in thought or deed does not excuse his conduct. He was guilty of a flagrant violation of his promises to the Goths, and he was guilty of gross disobedience to the Emperor's orders. It was not the business of the commander-in\texthebrew{‑}chief to decide the terms of peace; that was entirely a question for the Emperor. We can understand his unwillingness to allow the complete victory, which seemed within his grasp, to escape him; but it would be difficult to justify the chicanery which he employed at first in protracting the negotiations, and then in deceiving the enemy by pretended disloyalty to his master. Nor was his policy justified by success. It did not lead automatically to the complete conquest of Italy and the extension of Imperial authority to the Alps. When he sailed for Constantinople, he left behind him in the provinces north of the Po, enemies who had not submitted and a new Ostrogothic king who was bound by no covenant. A resumption of hostilities could not fail to ensue. If the peace which Justinian offered to the Goths had been concluded, and Witigis had remained as the recognised ruler of trans-Padane Italy, bound to the Empire by treaty, the arrangement could not indeed have been final, but the Emperor was justified in calculating that it would ensure for some years to come the tranquillity of Italy, and enable him to throw all his forces into the imminent struggle with Persia.

It is as little surprising then that when the victorious general disembarked at Constantinople with a captive king in his train, the Emperor should have given him a cold reception and denied him the honours of a triumph, as that the people, dazzled by the distinction of his captives and the richness of his spoil,

p215
and measuring his deserts by these spectacular results,\texthebrew{⁠}\footnote{The populace indeed were not permitted to see the treasures. They were exhibited to the senators in the Palace. \emph{B. G.} III.1.3. It is possible that the victory over Witigis was celebrated in the following year when Justinian made a triumphal entry into Constantinople on August 12, through the Charisian Gate. The ceremony is briefly described in Constantine Porphyrog., \emph{App. ad libr. prim. de Cer.} p497, a passage evidently taken or abridged from Peter the Patrician. The immediate motive of the triumph may have been the success of Belisarius in Mesopotamia in capturing the fort of Sisaurana. See Serruys, in \emph{Revue des études grecques}, XX.240 \emph{sqq.} 1907. The question has also been discussed by Martroye, who rejects the explanation of Serruys, thinks that the date meant is 540, not 541, and explains the ceremony as a solemn entry, distinct from a triumph, in which Belisarius took part. See \emph{Mém. de la Soc. des Antiquaires de France}, 1912, p25 (the article is only known to me through Bréhier's notice in \emph{Rev. Hist.} CI, Sept.-Déc. 1912, p325).

} should have attributed the Imperial attitude to jealousy. Though the enemies of Belisarius did all they could to poison Justinian's mind with suspicions, he can hardly have had serious doubts of his general's loyalty, yet it must have been far from agreeable to him to know that a subject had been given the opportunity of rejecting the offer of a throne. But, apart from this, it must be admitted that he was justified in refusing a triumph to a general who, whatever his services had been, had deliberately frustrated his master's policy. That the anxiety of the Emperor to hasten the departure of Belisarius from Italy was not entirely due to the urgent need of his services in the East, may be inferred from the fact that he was not sent against the Persians until the ensuing spring.

The Gothic prisoners were honourably treated. Witigis received the title of patrician and an estate on the confines of Persia.\texthebrew{⁠}\footnote{Jordanes, \emph{Get.} § 313 ; \emph{Hist. Misc.} xvi p107 administrationem illi Persarum tribuit terminos. It is not quite clear what this means.

} He survived his dethronement for two years.

It is naturally to be assumed that, as the provinces of Italy were gradually recovered, measures were taken for securing the civil administration. In some cases probably the Italians who served under the Goths were allowed to continue in their posts as governors of provinces, in others new men must have been appointed. But it was also necessary, perhaps even before the capture of Rome, to set up a central financial administration.\texthebrew{⁠}\footnote{\emph{Nov.} 75 (Dec. 537) seems to imply that a comes s. patrimonii per Italiam had already been appointed. He dealt with the patrimonial estates in Sicily as well as in Italy.

} Sicily had been reorganised after its submission to Belisarius and committed to the government of a Praetor, who had the

p216
responsibility for military as well as for civil affairs.\texthebrew{⁠}\footnote{\emph{Nov.} 75 (Dec. 537). This law provides that the Quaestor at Constantinople should be the court of appeal for Sicilian lawsuits.

\texthebrew{◂} previous

next \texthebrew{▸}Images with borders lead to more information.

The thicker the border, the more information.

(Details here.)

UP TO:

J. B. Bury
Homepage

Latin \& Greek

Texts

LacusCurtius

Home} Amid the din of arms these administrative measures occupied little attention, and they were soon to be upset or endangered by the renewal of war throughout the whole peninsula.

It may be convenient to the reader to have before him a table showing some of the distances on the routes between Italy and the East.

The usual rate of travelling, on horseback, varied from 60 to 75 kilometres daily. The regular rate of marching for an army was from 15 to 17 kilometres, but might be considerably more if the army was small or when there was a special need for haste.

The average rate of sailing was from 100 to 150 sea miles in 24 hours.

\xchapter{Chapter XVIII, Part E}

Short URL for this page:

bit.ly/BURLAT18E

\xsubsection{(Part 5 of 5)}

The power of the Ostrogoths was not yet broken. They were soon to regain much that they had lost, and under a new warrior king to wage a war which was well-nigh fatal to the ambitions of Justinian. But before we proceed to the second chapter of the reconquest of Italy, we may glance at the peaceful work of three eminent Italians who shed lustre on the Ostrogothic period, and secured a higher place in the eyes of posterity than the kings and warrior who in their own lifetime possessed the stage.

It is hardly too much to say that Boethius had a more genuine literary talent than any of his contemporaries, either Latin or Greek. We have seen how he composed in prison the ``golden volume''\texthebrew{⁠}\footnote{``Not unworthy of the leisure of Plato or Tully,'' Gibbon, iv p415. On Boethius cp. Ebert, \emph{Gesch. der Litt. des Mittelalters}, º i.485 \emph{sqq.}; H. F. Stewart, \emph{Boethius.}

} which has immortalised him, the \emph{Consolation of Philosophy.} It was one of the best known and most widely read books throughout the Middle Ages, notwithstanding the fact that it ignores Christianity,\texthebrew{⁠}\footnote{Stewart (\emph{op. cit.} 106 \emph{sqq.}) would explain this by the view that the work ``is intensely artificial.'' The verses, he says, are smooth and cold. There is ``nothing that suggests a heart beating itself out against the bars of its prison.'' Thus the book does not express the personal beliefs of the author, who composed simply as a diversion, to pass the time.

} though its Platonism has a Christian colouring. It was translated into Anglo-Saxon by King Alfred, and into English by Chaucer.\footnote{And after the Middle Ages by Queen Elizabeth. There were early translations in the principal European languages.

}

The \emph{Consolation of Philosophy} has indeed a considerable charm, which is increased by the recollection of the circumstances in which it was composed. A student who, maintaining indeed a lukewarm connection with politics, had spent most of his days in the calm atmosphere of his library, where he expected to end his life, suddenly found himself in the confinement of a dismal

p217
prison with death impending over him. There is thus in his philosophical meditations an earnestness born of a real need of consolation, while at the same time there is a pervading serenity. Poems, sometimes lyrical, sometimes elegiac, break the discussion at intervals,\texthebrew{⁠}\footnote{This form of mixed verse and prose was originated by the Greek Cynic Menippus, and in later literature had been employed by Varro, Petronius, Seneca, and more recently by Martianus Capella the Neoplatonist (fifth century) in his allegorical work ``On the nuptials of Philology and Mercury, and the seven liberal arts.'' It was probably this work that suggested the form to Boethius.

} like organ chants in a religious service.

The problem of the treatise\texthebrew{⁠}\footnote{Book I contains the story of the writer's personal wrongs, which he relates to Philosophia; Bk. II a discussion on Fortune; Bk. III passes to the Summum Bonum; in Bk. IV Philosophia justifies God's government; Bk. V deals with free will.

} is to explain the ``unjust confusion'' which exists in the world, the eternal question how the fact that the evil win often the rewards of virtue ( pretium sceleris \texthebrew{—} diadema ) and the good suffer the penalties of crime, can be reconciled with a `` deus, rector mundi .'' If I could believe, says Boethius, that all things were determined by chance and hazard, I should not be so perplexed. In one place he defines the relation of fate to the Deity in the sense that fate is a sort of instrument by which God regulates the world according to fixed rules. In other words, fate is the law of phenomena or nature, under the control of the Supreme Being, which he identifies with the Summum Bonum or highest good. His discussion of the subject is not very illuminating \texthebrew{—} did it really satisfy him?

But the metaphysical discussion does not interest the student of literature so much as the setting of the piece and things said incidentally.\texthebrew{⁠}\footnote{It is interesting to notice that Dante's famous verses

Nessun maggior dolore

Che ricordarsi del tempo felice

Nella miseria,

come from Boethius (II.IV). It has been pointed out that the idea occurs in

Synesius, \emph{Ep.} 57

( Sandys, \emph{Hist. of Classical Scholarship}, I.243, n2 ). Dante assigned to Boethius a place in the Fourth Heaven.

} Boethius imagines his couch surrounded by the Muses of poetry, who suggest to him accents of lamentation. Suddenly there appears at his head a strange lady of lofty visage. There was marvellous fluidity in her stature; she seemed sometimes of ordinary human height, and at the next moment her head touched heaven, or penetrated so far into its recesses that her face was lost to the vision. Her eyes too were unnatural, brilliant and transparent beyond the power of human eyes,

p218
of fresh colour and unquenchable vigour. And yet at the same time she seemed so ancient of days ``that she could not be taken for a woman of our age.'' Her garments were of the finest threads, woven by some secret art into an indissoluble texture, woven, as she told Boethius, by her own hands. And on this robe there was a certain mist of neglected antiquity, the sort of colour that statues have which have been exposed to smoke. On the lower edge of the robe was the Greek letter \textgreek{Π} (the initial of \textgreek{Πρακτική} , Practical Philosophy), from which stairs were worked leading upwards to the letter \textgreek{Θ} ( \textgreek{Θεωρητική} , Pure Philosophy). And her garment had the marks of violent usage, as though rough persons had tried to rend it from her and carried away shreds in their hands. The lady was Philosophia; she bore a sceptre and parchment rolls. She afterwards explained that the violent persons who had rent her robe were the Epicureans, Stoics, and other late schools; they succeeded in tearing away patches of her dress, fancying severally that they had obtained the whole garment. Philosophia's first act is to drive out the Muses, whom she disdainfully terms ``theatrical strumpets,'' and she remarks that poetry ``accustoms the minds of men to the disease but does not set them free.''\footnote{Ed. Peiper, p5: hominumque mentes musae assuefaciunt morbo, non liberant.

}

A striking feature of the \emph{Consolatio} is the interspersion of the prose dialogue with poems at certain intervals, which, like choruses in Greek tragedy, appertain to the proceeding argument. Thus the work resembles in form Dante's \emph{Vita Nuova}, where the sonnets gather up in music the feelings occasioned by the narrated events. These poems, which betray the influence of Seneca's plays,\texthebrew{⁠}\footnote{Peiper in his edition gives a list of passages which contain excerpts from or echoes of Seneca's tragedies.

} have all a charm of their own, and metres of various kinds are gracefully employed.

One poem, constructed with as much care as a sonnet,\texthebrew{⁠}\footnote{ii. viii. p48; it consists of thirty lines thus arranged, 4+4+4+3 = 4+4+4+3.

} sings of the ``love that moves the sun and stars,''

an idea familiar to modern readers from the last line of Dante's \emph{Divina Commedia}, but which is as old as Empedocles. As

p219
an example of his metrical devices take two lines of a stanza, where the author is illustrating the return of nature to itself by a caged bird, which, when it beholds the green wood once more, spurns the sprinkled crumbs \texthebrew{—}Immediately after this poem Boethius proceeds, ``Ye too, O creatures of earth! albeit in a vague image, yet do ye dream of your origin'' ( vos quoque, O terrena animalia! tenui licet imagine vestrum tamen principium somniatis ).

The delicate feeling of Boethius for metrical effect may be illustrated by the poem on the protracted toils of the siege of Troy and the labours of Hercules. It is written in Sapphic metre, but the short fourth lines are omitted until the end. The effect of this device is that the mind and voice of the reader continue to travel without relief or metrical resting-place until all the labours are over and heavenly rest succeeds in the stars of the concluding and only Adonius \texthebrew{—}If the \emph{Consolation} had never been written, Boethius would still have had his place in the list of men who have done service to humanity. Possessing the multifarious learning characteristic of the time, he devoted himself especially to the philosophy of the great masters, Plato and Aristotle, and at an early age he conceived the ambitious idea of translating into Latin, and writing commentaries on, all their works.\texthebrew{⁠}\footnote{\textgreek{Περ\texthebrew{ὶ ἑ}ρμηνείας}, ii.2.

} Of this task of a lifetime he succeeded only in completing the logical works of Aristotle, but these translations were of capital importance, in keeping alive the study of logic throughout the Middle Ages,\texthebrew{⁠}\footnote{We have also his version and exegesis of the \emph{\textgreek{Περ\texthebrew{ὶ ἑ}ρμηνείας}}, in two forms (one elementary, the other for advanced students). His translation of the \emph{Isagoge} of Porphyry to Aristotle's \emph{Categories} was a vulgar handbook in the Middle Ages. He also wrote a commentary on Cicero's \emph{Topica.}

} and he raised the question as to the nature of genera and species , which was to be fought out towards the end of that period in the debate between the Nominalists and the Realists. His polymathy carried him into other fields. He translated (perhaps)

p220
the Geometry of Euclid,\texthebrew{⁠}\footnote{He certainly wrote an original treatise \emph{De geometria.} It is doubtful whether the extant translation of Euclid ascribed to him is really his.

} wrote treatises on arithmetic and music, and even ventured into the region of theological doctrine.\texthebrew{⁠}\footnote{\emph{De trinitate}, and one or two other tracts. It has been supposed that they were written for a literary society which met every week to read and discuss papers; see Stewart, \emph{op. cit.} p132.

} Though he was a professing Christian, he did not yield to Symmachus, the illustrious pagan ancestor of his wife, in enthusiasm for the ancients, and his aim was to keep alive in Italy the quickening influence of Greek science.\texthebrew{⁠}\footnote{Cp. Cass. \emph{Var.} I.45 (letter of Theoderic, A.D. 509\texthebrew{‑}510) quascumque disciplinas vel artes fecunda Graecia per singulos viros edidit, te uno auctore patrio sermone Roma suscepit. Here his various literary activities are enumerated. Ennodius addressed Boethius as an indefatigable student quem in annis puerilibus . . . industria fecit antiquum (\emph{Epp.} vii.13).

} Writing in the year of his consul ­ship ( A.D. 510), he said, ``The cares of office hinder me from devoting all my time to these studies (in logic), but I think it may be considered of some public utility to instruct my fellow-citizens in the subject.''\footnote{\emph{Comm. in Arist. Cat.} ii. \emph{Praef.} (Migne, \emph{P. L.} lxiv.201).

}

The other eminent man of letters, who shed a certain lustre on Ostrogothic Italy, Magnus Aurelius Cassiodorus Senator,\texthebrew{⁠}\footnote{His contemporaries called him Senator. The facts known about the family and the dates of his career are summarised by Mommsen in the \emph{Prooemium} to his edition. On his literary work see

Hodgkin, \emph{Letters of Cassiodorus} ;

Sandys, \emph{History of Classical Scholarship}, I.244 \emph{sqq.}

} was of inferior fibre to Boethius in literary taste as well as in personal character, but he was no less genuinely interested in intellectual pursuits, and posterity owes him an even greater debt. The Cassiodori, who seem originally to have come from Syria, acquired an estate at Scyllacium (Squillace) on the eastern coast of Bruttii. The great-grandfather of Cassiodorus successfully defended this province and Sicily against raids of the Vandals.\texthebrew{⁠}\footnote{\emph{Var.} I.4. Above,

vol. I p258 .

} His grandfather, a friend of Aetius, was employed on an embassy to the Huns,\texthebrew{⁠}\footnote{\emph{Ib.}

} and we have seen

how his father filled high posts under Odovacar and Theoderic. Born himself not long before Theoderic's invasion, he was a boy when his father became Praetorian Prefect\texthebrew{⁠}\footnote{About A.D. 501.

} and employed him as a legal assistant in his bureau. He won the king's notice by a panegyric which he pronounced on some public occasion, and was appointed to the high office of Quaestor of the Palace at an unusually early age.\texthebrew{⁠}\footnote{In his early twenties. His Quaestorship falls between the years 507 and 511.

} In this post he conducted the official correspondence of the

p221
king, and in the composition of State documents he found congenial employment for his rhetorical talent. After he laid down this office ( A.D. 511), he seems to have taken no part in public affairs (except in the year of his consul ­ship, A.D. 514) till the close of Theoderic's reign, when he was appointed Master of Offices.\texthebrew{⁠}\footnote{About A.D. 523\texthebrew{‑}524. In this post he also fulfilled many of the duties of Quaestor (cp. ix.25).

} He continued to hold this dignity in the first years of the following reign, and after an interval of retirement\texthebrew{⁠}\footnote{A.D. 527\texthebrew{‑}533.

} he became Praetorian Prefect, and remained in that post during the stormy years which followed, content to play the ignoble rôle of a time-server , apparently as loyal to Theodahad as he had been to Amalasuntha,\texthebrew{⁠}\footnote{In his Oration on the marriage of Witigis, he condemns Amalasuntha for her education of Athalaric (p473): fas fuit illam sub pietatis excusationis peccare.

} and on Theodahad's fall turning without hesitation to the rising sun of Witigis. But we have every reason to believe that throughout his career he did not waver in a sincere conviction that Italy was better off under Ostrogothic government than she would have been under the control of Constantinople. It is possible that he retired from public life before the capture of Ravenna, but while he was still Prefect, he published ( A.D. 537) a collection of the official letters and State papers, which he had composed during his three ministries.\texthebrew{⁠}\footnote{The \emph{Variae} thus fall into three chronological groups. (1) quaestoriae, 507\texthebrew{‑}511 = Books i\texthebrew{‑}iv; (2) magisteriae, 523\texthebrew{‑}527 = Books v (except last two letters, A.D. 511), viii, ix.1\texthebrew{‑}14; (3) praefectoriae, 533\texthebrew{‑}537 = Books ix.15\texthebrew{‑}25, x, xi, xii. Books vi and vii contain \emph{Formulae} for admission to various offices of state. In 540 he added to this collection as Book xiii a treatise \emph{De anima.}

} This collection is a mine of information for the administration and condition of Ostrogothic Italy, and we have to thank perhaps the literary vanity of Cassiodorus for the ample knowledge that we possess of Theoderic's policy; but it bears all the signs of having been carefully expurgated. As the work was published when the issue of the war was uncertain, he consulted his own interests by cutting out anything that could offend either the Emperor or the Goths,\texthebrew{⁠}\footnote{Cp. Mommsen, \emph{loc. cit.} p. xxii.

} and it is probable that many documents which would clear up some of our uncertainties as to the relations between Ravenna and Constantinople have been omitted altogether.

Few rhetorical compositions, and perhaps no public documents, offer greater difficulties to the reader when he attempts to arrive

p222
at the plain fact which the author intends to convey. ``It is ornament alone,'' he says in his Preface, ``that distinguishes the learned from the unlearned,'' and, true to this maxim of decadent rhetoric, he obscures the simplest and most trivial statements in a cloud of embellishments. But to appreciate his inflated style we must remember that he was, after all, only improving upon what had been, since Diocletian, the traditional style of the Imperial chancery. We have innumerable constitutions of the fourth and fifth centuries, in which the vices of adornment and contorted phraseology make it a laborious task to discover the meaning. Cassiodorus exerted his ingenuity and command of language in elaborating the sublime style, always frigid, but ludicrously inappropriate to legal documents and State papers.

In his later years Cassiodorus betook himself to his ancestral estate at Squillace,\texthebrew{⁠}\footnote{There is a charming description of the situation and amenities of the town of Squillace in

\emph{Var.} XII.15 . After the capture of Ravenna in 539 he seems to have gone to Constantinople and lived there for about fifteen years, returning to Italy after the final conquest of Italy by Narses. At least the reference to him in a letter of Vigilius excommunicating two Italian deacons (Mansi, ix.357) seems to imply that he was at Constantinople about 550. Cp. Sundwall, \emph{Abh.} p156.

Thayer's Note: For those that Latin disagrees with, see

this Italian translation .

} and devoted the rest of a long life to religion and literature. He became a monk and founded two monasteries, one, up in the hills at Castellum, a hermitage for those who desired solitary austerity, the other, built inside the fish-ponds of his own domain and hence called Vivarium , for monks who were content to live in the less strict conditions of a monastic society. At Vivarium, where he lived himself, Cassiodorus introduced a novelty\texthebrew{⁠}\footnote{It had indeed been anticipated in the monastery of St. Martin of Tours, where young monks had practised the copying of MSS. (Sulpicius Severus, \emph{Vita S. Mart.} vii).

} which led to fruitful results for posterity. He conceived the idea of occupying the abundant leisure of the brethren with the task of multiplying copies of Latin texts. There was a chamber known as the scriptorium or `` writing-room `` in the monastery, in which those monks who had a capacity for intellectual labour, used to copy both pagan and Christian books, working at night by the light of self-filling ``mechanical lamps.'' It is well known that the preservation of our heritage of Latin literature is mainly due to the labours of monastic copyists. The originator of the idea was Cassiodorus. His example was adopted in other religious establishments, and monastic libraries came to be a regular institution.

p223
Most of the works of Cassiodorus have come down to us. The great exception is his \emph{History of the Ostrogoths}, in which he attempted to reconstruct a historical past for the Gothic race.\texthebrew{⁠}\footnote{On this work (in 12 books), which belongs to the period of leisure between 527 and 533, see his \emph{Praefatio} to the \emph{Variae}, and \emph{Var.} IX.25 (originem Gothicam historiam fecit esse Romanam). For the later history of the Goths he drew on their own traditions, and he used the work of Ablabius, who wrote a Roman history in Greek \emph{c.} A.D. 400 (he is \emph{possibly} referred to in \emph{Var.} X.22 ). The work of Cassiodorus was liberally used by Jordanes, so that we can form an idea of its scope and contents.

} Starting with the two false assumptions that the Goths were identical on the one hand with the Getae and on the other hand with the Scythians, he was able to produce, for the records of Greek and Roman antiquity, a narrative which represented them as playing a great part on the stage of history at a time when they were really living in obscurity on the Lower Vistula, utterly beyond the horizon of Mediterranean civilisation.

The principal works of his later years were intended for the instruction of his monks \texthebrew{—} the \emph{Institutions} and a treatise on \emph{Orthography.} The \emph{Institutions} consisted of two independent parts, of which the first, \emph{De institutione divinarum litterarum}, was intended as an introduction to the study of the manuscripts of the Bible,\texthebrew{⁠}\footnote{Cassiodorus went to much expense in procuring MSS. (\emph{De ins.} I.8). He mentions a large codex containing Jerome's version of the Scriptures, and ``it has been conjectured that part of it survives in the first and oldest quaternion of the \emph{codex Amiatinus} of the Vulgate, now in the Laurentian Library in Florence'' (Sandys, \emph{op. cit.} i.251).

} and contains an interesting disquisition on the question of correcting the text. The second part is a handbook on the Seven Liberal Arts. The two together offered a general survey of sacred and secular learning.\texthebrew{⁠}\footnote{These books were written before A.D. 555. Of his other works need only be mentioned his Commentary on the Psalms and his Ecclesiastical History (\emph{Hist. Tripartita}).

} The manual on spelling was composed, for the guidance of copyists, in the ninety-third year of his age (\emph{c.} A.D. 580).\texthebrew{⁠}\footnote{For this work he had the advantage of using, besides older treatises, that of his elder contemporary Priscian (an African provincial), whose Panegyric on Anastasius has been referred to above,

vol. I p467 . Priscian's Grammar was a standard text-book in the Middle Ages. It was transcribed at Constantinople by a pupil in 526\texthebrew{‑}527. Sandys, \emph{op. cit.} 258\texthebrew{‑}259.

} Thus he had lived to see great changes. He had witnessed the complete subjugation of Italy by Justinian, and when, at the age of eighty, he saw many of its provinces pass under the yoke of the Lombard barbarians, it may well have occurred to him that if the Ostrogothic rule had

p224
been allowed to continue this calamity would have been spared to his fellow-countrymen .

While Boethius was immersed in the study of philosophy in his library, with its walls decorated with ivory and glass, and while Cassiodorus was engaged in his political and rhetorical labours in the Palace at Ravenna, another young man, of about the same age as they, who was destined to exert a greater influence over western Europe than any of his Italian contemporaries, was spending his days in austere religious practices in the wild valleys of the upper Anio. St. Benedict, who belonged to the same Anician gens as Boethius, was born at

Nursia , in an Apennine valley, \texthebrew{•} about twenty miles east of Spoleto.\texthebrew{⁠}\footnote{A.D. 480.

} Sent to Rome to study, he was so deeply disgusted by the corruption of his school companions and by the vice of the great city, that at the age of fourteen he set out with a faithful nurse for the ``desert,'' and at length took up his abode in a cave at Sublaqueum (Subiaco), near the sources of the Anio, where he lived as a hermit. The temptations which he resisted, the perils which he escaped, and the legends which rapidly gathered round him, may be read in the biography written by his admirer, Pope Gregory the Great.\texthebrew{⁠}\footnote{Or in

Hodgkin, \emph{Italy and her Invaders}, IV chap. XVI .

} After his fame had gone abroad, his solitude was interrupted, for men who desired to embrace the monastic life flocked to him from all parts. He founded twelve monasteries in the neighbourhood of Subiaco. In A.D. 528 he left the peaceful region and went into Campania, where, at Monte Cassino, halfway between Rome and Naples, he found a congenial task awaiting him. Here, notwithstanding all the efforts of Christian Emperors and priests to extirpate the old religions, there still stood an altar and statue of Apollo in a sacred grove, and the surrounding inhabitants practised the rites of pagan superstition. Benedict induced them to burn the grove and demolish the altar and image, and on the height above he founded the great monastery where he lived till his death. The Rule which he drew up for his monks avoids the austerities of Egyptian monasticism, and he expressly says that he wished to ordain nothing hard or burdensome.\texthebrew{⁠}\footnote{See \emph{Prologue.}

} Within three hundred years this code of laws had superseded all others in western Europe, where it held much the same position as St. Basil's in the East.\footnote{Benedict knew Basil's Rule and refers to it (chap. lxxiii).

}

p225
Benedict himself did not anticipate that the order which he founded would ultimately become the learned order in the Church. He ordained that ``because idleness is an enemy of the soul,'' the brethren should occupy themselves ``at specified times in manual labour, and at other fixed hours in holy reading,''\texthebrew{⁠}\footnote{\emph{Reg. Ben.} chap. xlviii (Gasquet's translation).

} but there is no indication that he included in manual labour the transcription of MSS. \texthebrew{⁠}\footnote{The rule that no brother is to keep as his own, without the abbot's leave, books or tablets or pens (chap. xxxiii) proves nothing.

\texthebrew{◂} previous

next \texthebrew{▸}Images with borders lead to more information.

The thicker the border, the more information.

(Details here.)

UP TO:

J. B. Bury
Homepage

Latin \& Greek

Texts

LacusCurtius

Home} Probably this was not introduced at Monte Cassino till after his death, under the influence of Cassiodorus.

It may be convenient to the reader to have before him a table showing some of the distances on the routes between Italy and the East.

The usual rate of travelling, on horseback, varied from 60 to 75 kilometres daily. The regular rate of marching for an army was from 15 to 17 kilometres, but might be considerably more if the army was small or when there was a special need for haste.

The average rate of sailing was from 100 to 150 sea miles in 24 hours.

\xchapter{Chapter XIX, Part A}

\xsubsection{(Part 1 of 4)}

The policy of Belisarius had frustrated the conclusion of a peace which would have left the Goths in peaceful possession of Italy north of the Po. Such a peace could hardly have been final, but it would have secured for the Empire a respite of some years from warfare in the west at a time when all its resources were needed against the great enemy in the east. If Belisarius had not been recalled, he would probably have completed the conquest of the peninsula within a few months. This, which would have been the best solution, was defeated by the jealousy of Justinian; and the peace proposed by the Emperor, which was the next best course, was defeated by the disobedience of his general. Between them they bear the responsibility of inflicting upon Italy twelve more years of war.

The greater blame must be attached to Justinian. He had indeed every reason to be displeased with the behaviour of Belisarius, but the plainest common sense dictated that, if he could no longer trust Belisarius, he should replace him by another commander-in\texthebrew{‑}chief . Of the generals who remained in Italy the most distinguished was John, the nephew of Vitalian. But instead of appointing him or another to the supreme command, the Emperor allowed the generals to exercise co-equal and independent authority each over his own troops. In consequence of this unwise policy there was no effective co-operation ; each commander thought only of his own interests. They plundered the Italians, and allowed the soldiers to follow their example, so that discipline was undermined. In a few months so many blunders were committed that the work accomplished by Belisarius

p227
in five arduous years was almost undone, the Goths had to be conquered over again, and it took twelve years to do it.\footnote{For the second period of the Italian war, our source is still Procopius. But the historian is no longer writing from personal knowledge, for he probably never returned to Italy. Haury indeed is confident that he was in Italy during the twelfth year of the war, 546\texthebrew{‑}47, because he wrote about this year as much as he wrote about the first years of the war (\emph{Procopiana}, I p9), but this is far from decisive. Chance may have supplied him with fuller information for the events connected with the siege of Rome. And as a matter of fact, the events of the last six months of 552 run to as great a length as those of the whole year 546\texthebrew{‑}547.

}

The situation was aggravated by the prompt introduction of the Imperial financial machinery in the conquered provinces. The logothete\texthebrew{⁠}\footnote{For the logothetes see below,

p358 .

} Alexander, an expert in all the cruel methods of enriching the treasury and the tax-collector at the expense of the provincials, arrived, and soon succeeded in making both the Italians and the soldiers thoroughly discontented. Having established his quarters at Ravenna, he required the surviving Italian officials of the Gothic kings to account for all money that had passed through their hands during their years of service, and compelled them to make good deficits out of their own pockets. It cannot be doubted that many of these officials had made illegitimate profits and we need not waste much pity on them; but Alexander extended his retrospective policy to all private persons who had any dealings with the fisc of Ravenna. In an inquiry into transactions of twenty, thirty, or forty years ago, conducted by a man like Alexander, it is certain that grave injustices were done.

He was acting on the constitutional principle that Italy was, throughout the Gothic régime, subject to the Imperial authority, and that the kings and their servants were responsible to the Emperor for all their acts. But his proceedings were calculated to alienate the sympathies of the Italians and render the government of Justinian unpopular. At the same time, by curtailing the pay of the soldiers on various pretexts, he caused a deep sense of injustice in the army.

After the departure of Belisarius, Vitalius was stationed in Venetia, Constantian commanded the troops in Ravenna, Justin held Florence, Conon Naples, Cyprian Perusia, and Bessas perhaps had his quarters in Spoletium.\texthebrew{⁠}\footnote{See \emph{B. G.} III.1 § 34; 3 § 2; 5 § 1; 6 § 2; 6 § 8 and 5 § 4; 6 § 8.

} North of the Po, the only important places still held by the Goths were

Ticinum ,

p228
which king Ildibad made his residence, and Verona.\texthebrew{⁠}\footnote{Procopius, \emph{B. G.} 1 § 27, says only Ticinum; but it is clear from the narrative that Verona had remained in their hands throughout.

} The army of Ildibad amounted at first to little more than a thousand men, but he gradually extended his authority over Liguria and Venetia. The Roman generals did nothing to prevent this revival of the enemy's strength, and it was not till he approached Treviso, which appears to have been the headquarters of Vitalius, that Ildibad met any opposition. Vitalius, whose forces included a considerable body of Heruls, gave him battle and was decisively defeated, Vitalius barely escaping, while the Herul leader was slain.

Ildibad did not live long enough to profit by the prestige which his victory procured him. His death was indirectly due to a quarrel with Uraias, to whose influence he had owed his crown. The wife of Uraias was beautiful and wealthy, and one day when she went to the public baths, in rich apparel and attended by a long train of servants, she met the queen, who was clad in a plan dress (for the royal purse was ill-furnished ), and treated her with disrespect. The queen implored Ildibad to avenge her outraged dignity, and soon afterwards Uraias was treacherously put to death. This act caused bitter indignation among the Goths, yet none of them was willing to avenge the nephew of Witigis. But a Gepid belonging to the royal guard, who had a personal grudge against the king, murdered Ildibad at a banquet in the palace ( A.D. 541, about May). He would not have ventured on the crime if he had not known that it would please the Goths, as a just retribution for the murder of Uraias.

The event came as a surprise, and the Goths could not immediately agree on the choice of a successor to the throne. The matter was decided in an unexpected way. The Rugian subjects of Odovacar, who had submitted after his fall to the rule of Theoderic, had never merged themselves in the Gothic nationality, but had maintained their identity as a separate people in northern Italy. They seized the occasion to proclaim as king Eraric, the most distinguished of their number. The Goths were vexed at the presumption of the Rugians, but nevertheless they recognised Eraric, and endured his rule for five months, presumably because there was none among themselves on whose fitness for the throne they could agree.

p229
Eraric summoned a council and persuaded the Goths to consent to his sending an embassy to Constantinople for the purpose of proposing peace on the same terms which the Emperor had offered to Witigis. But the Rugian was a traitor. He selected as ambassadors creatures of his own, and gave them secret instructions to inform Justinian privately that he was prepared, in return for the Patriciate and a large sum of money, to abdicate and hand over northern Italy to the Empire.

In the meantime he made no pretence of carrying on the war, and the Goths regretted the energy of Ildibad. Looking about for a worthy successor, they bethought them of Totila,\texthebrew{⁠}\footnote{So he is always called by Procopius, but on the coins of his reign, both silver and copper, his name is invariably Baduila, which is also found in the \emph{Hist. Miscella}, B. xvi. p107. Jordanes (\emph{Rom.} 380) uses both names. The reason of the double designation has not been cleared up. Totila issued at first coins with the head of Justinian, but this recognition of the Emperor was soon abandoned, and the bust of Anastasius was substituted. Finally, in his last years, he issued silver and bronze coins with his own bust (imitated from that of Anastasius), A.D. 549\texthebrew{‑}552. The regal mint was at Ticinum, but coins were afterwards struck at Rome. See Wroth, \emph{Coins of the Vandals}, etc., xxxvii\texthebrew{‑}xxxviii.

} Ildibad's nephew, a young man who had not yet reached his thirtieth year and had acquired some repute for energy and intelligence. He had been appointed commander of the garrison of Treviso, and after his uncle's assassination, despairing of the Gothic cause, he had secretly opened negotiations with Ravenna, offering to hand over the town. A day for the surrender was fixed when he received a message from the Gothic nobles who were conspiring against Eraric, inviting him to become their king. Concealing his treacherous intrigue with the enemy, he accepted the proposal on condition that Eraric should be slain before a certain day, and he named the day on which he had undertaken to admit the Romans into the town. Eraric was duly put to death by the conspirators and Totila ascended the throne ( A.D. 541, September or October).

Eraric's ambassadors seem to have been still at Constantinople when the news of his murder and Totila's accession arrived. Justinian was incensed at the supine conduct of his generals who had failed to take advantage of Eraric's incapacity, and his indignant messages at last forced them to plan a common

p230
enterprise. They met at Ravenna and decided that Constantian and Alexander should advance upon Verona with 12,000 men. One of the Gothic sentinels was bribed to open a gate, and when the army approached the city, a picked band led by an Armenian, Artabazes, was sent forward at night to enter and take possession. Artabazes did his part, and Verona would have been captured if the commanders had not wasted the night in quarrelling over the division of the expected booty. When they arrived at last, the Gothic garrison had regained possession of the place and barred the gates, and the little band of Artabazes, having no other means of escape, leaped from the walls and all but a few were killed by the fall.\footnote{The attempt on Verona and the battles of Faventia and Mugello are probably to be placed in the spring of 542. See \emph{Cont. Marcell.} §§ 2, 3, \emph{sub a.}

}

The army retreated across the Po and encamped on the stream of Lamone,\texthebrew{⁠}\footnote{Cp. Hodgkin, \emph{op. cit.} IV.444 . The Lamone is the ancient Anemo.

} near Faventia. Totila marched against them at the head of 5000 men, and in the battle which ensued gained a brilliant victory, all the Imperial standards falling into his hands.\texthebrew{⁠}\footnote{The most striking incident in the battle was the single combat between Valaris, a gigantic Goth, and Artabazes, in which the Goth was slain and Artabazes mortally wounded.

} Verona and Faventia exhibited the evil of a divided command.

Totila was encouraged by this success to take the offensive in Tuscany. He sent a force against Florence, where Justin, who had helped to capture it three years before, was in command. John, Bessas, and Cyprian hastened to its relief, and on the appearance of their superior forces, the Goths raised the siege and moved up the valley of the Sieve. This locality was then known as Mucellium, and the name survives as Mugello. The Roman army pursued them, and John with a chosen band pushed on to engage the enemy while the rest followed more slowly. The Goths, who had occupied a hill, rushed down upon John's troops. In the hot action which ensued, a false rumour spread that John had fallen, and the Romans retired to join the main army, which had not yet been drawn up in order of battle, and was easily infected with their panic. All the troops fled disgracefully, and the Goths pursued their advantage. The prisoners were well treated by Totila and induced to served under his banner. The defeated generals abandoned all thought of

p231
further co-operation and hastily retreated, Bessas to Spoletium, Cyprian to Perusia, and John to Rome.

The victory of Mugello, however, did not lead to the defection of Tuscany, and Justin remained safely in Florence. Totila captured some places in Umbria \texthebrew{—} Caesena and Petra Pertusa,\texthebrew{⁠}\footnote{The Continuer of Marcellinus, \emph{sub} 542, adds Urbinum and Mons Feletris.

}\texthebrew{—} but then instead of pursuing steadily the conquest of central Italy, where the Imperialist forces, concentrated in strong cities, were too formidable for his small army, he decided to transfer his operations to south of the peninsula. There the success of his arms was swift and sweeping. Avoiding Rome, he marched to Beneventum, which was an easy prey, and razed its walls to the ground. The provinces of Lucania and Bruttii, Apulia, and Calabria acknowledged his authority and paid him the taxes which would otherwise have gone to satisfy the demands of the Imperial soldiers, to whom long arrears were owed. Totila had meanwhile laid siege to Naples, which Conon was holding with a garrison of 1000 Isaurians. He collected considerable treasure from Cumae and other fortresses in the neighbourhood, and created a good impression by his courteous treatment of the wives and daughters of Roman senators whom he found in these places and allowed to go free. This is one instance, and we shall meet others, of the policy which he often followed of winning the sympathy of the Italians by a more generous treatment than they were prepared to expect from an enemy.

The news of the revival of the Gothic power and the danger of Naples alarmed the Emperor, and he took some measures to meet the crisis, but they were far from sufficient. Instead of confiding the supreme command to an experienced general, he appointed a civilian, Maximin, to be Praetorian Prefect of Italy, and gave him powers of general supervision over the conduct of the war, sending with him Thracian and Armenian troops and a few Huns. Maximin, who seems to have been one of the worst choices the Emperor could have made, sailed to Epirus and remained there unable to decide what to do. Soon afterwards Demetrius, an officer who had formerly served under Belisarius, was sent to the west. He appears to have been invested with the office of Master of Soldiers, but we find him acting under the orders of Maximin.\texthebrew{⁠}\footnote{\emph{B. G.} III.6.13 \textgreek{Δημήτριον στρατήγον}; cp. 7.3.

} He sailed straight to

p232
Sicily, where he learned how severely Naples was suffering from lack of food, and he made prompt preparations to bring help. He had only a handful of men, but collecting as many vessels as he could find in the Sicilian harbours, he loaded them with provisions and set sail in the hope that the enemy would believe that they were conveying a large army. It is thought that if this bold design had been executed the Goths would have withdrawn from Naples and the city might have been saved. But before Demetrius reached his destination, he revised his plan and made for Porto, hoping to obtain some reinforcements from Rome. But the Roman garrison was demoralised and refused to join in an expedition which seemed full of danger. Demetrius then sailed for the bay of Naples. Totila meanwhile had been fully informed of the facts and had a number of war vessels ready to attack the transports when they were close to the shore. Most of the crews were slain or made prisoners; Demetrius was one of the few who escaped in boats.

Another attempt to relieve Naples was another failure. Maximin and the forces which accompanied him had at last left Epirus and reached Syracuse. Moved by the importunate messages of Conon for help, he consented, although it was now midwinter, to send these troops to Naples, and Demetrius, who had made his way back to Sicily, accompanied this second expedition. It reached the bay of Naples safely, but there a violent gale arose which drove the ships ashore close to the Gothic camp. The crews were easily slain or captured, and Demetrius fell into the hands of Totila.\footnote{Another Demetrius (originally a Cephallenian sailor who had distinguished himself in the campaigns of Belisarius and had been appointed an overseer of some kind in Naples) had succeeded in leaving the city to community with the relief expedition. He, too, fell into the hands of the Goths. He had given dire offence to Totila, whom, on his first appearance before the walls of Naples, he had overwhelmed with insolent abuse. The king now punished him by cutting out his tongue and cutting off his hands.

}

The Neapolitans were starving, and Totila proposed generous terms. ``Surrender,'' he said, ``and I will allow Conon and all his soldiers to depart unhurt and take all their property with them.'' Still hoping that help might come, Conon promised to surrender on these terms in thirty days. Confident that there was no chance of relief forthcoming, Totila replied, ``I will give you three months, and in the meantime will make no attempt to

p233
take the city.'' But before the term had run out, the exhausted garrison and citizens abandoned hope and opened the gates ( A.D. 543, March or April).

On this occasion Totila exhibited a considerable humanity which was not to be expected, as the historian Procopius remarks, from an enemy or a barbarian. He knew that if an abundance of food were at once supplied, the famished inhabitants would gorge themselves to death. He posted sentinels at the gates and in the harbour and allowed no one to leave the city. Then he dealt out small rations, gradually increasing the quantity every day until the people had recovered their strength. The terms of the capitulation were more than faithfully observed. Conon and his followers were embarked in ships with which the Goths provided them, and when, deciding to sail for Rome, they were hindered by contrary winds, Totila furnished horses , provisions, and guides so that they could make the journey by land.

The fortifications of Naples were partly razed to the ground.\footnote{If there is any foundation for the tradition, preserved in Gregory I, \emph{Dial.} ii cc14, 15, that Totila visited St. Benedict at Monte Cassino, the incident must have occurred either before or soon after the siege of Naples. The year of Benedict's death is uncertain, perhaps 544 (so Pagi and Martroye, \emph{L'Occident}, p436); the traditional date, accepted by the Order, is March 21, 543.

}

In the meantime the generals of Justinian were making no efforts to stem the tide of Gothic success. They plundered the Italians and spent their time in riotous living. Then Constantian wrote to the Emperor, stating bluntly that it was impossible to cope with the enemy.\texthebrew{⁠}\footnote{\emph{B. G.} III.9.5; the other leaders seem either to have enclosed separate letters to the same effect or to have attached their signatures to Constantian's communication.

} These messages did not arouse Justinian to action till they were reinforced by news of Totila's next movements.

Totila felt that he was now in a position to attack Rome itself. He began his operations by writing a letter to the Senate, in which he contrasted Gothic with ``Greek'' rule and attempted to show that it was the interest of the Italians that the old régime of the days of Theoderic and Amalasuntha should be restored. The letter was conveyed to Rome by Italian prisoners,

p234
but John, who was in command of the garrison, forbade the senators to reply. Totila then contrived that a number of placards, announcing that he bound himself by the most solemn oaths not to harm the Romans, should be smuggled into Rome and posted up. John suspected that the Arian clergy were his agents and expelled them all from the city.

Totila then sent part of his army to besiege Otranto, and with the rest advanced upon Rome (spring, A.D. 544). Thereupon Justinian at last decided to recall Belisarius from Persia and send him to Italy to assume the supreme command, as the only means of retrieving the situation.\footnote{It was said that Belisarius persuaded Justinian to send him to Italy, by a promise that the war would be self-supporting and that he would never ask for money ( Procopius, \emph{H. A.} 4.39

\textgreek{\texthebrew{ὥ}ς φασι}), but we find him writing for money at the end of the first year (\emph{B. G.} III.12.10; see below,

p235 ). He held the office of Comes stabuli.

}

The first thing Belisarius had to do was to collect some troops in Europe, for it was impossible to weaken the eastern front by bringing any regiments with him from Asia. At his own cost and with the assistance of Vitalius, who had recently been appointed Master of Soldiers in Illyricum, he recruited 4000 men in the Thracian and Illyrian provinces, and proceeded to Salona. His first care was to send a relief expedition to Otranto (summer, A.D. 544), and this enterprise was completely success ­ful.\texthebrew{⁠}\footnote{It was carried out by Valentine, who had won distinction in the siege of Rome by Witigis.

} The siege was raised and the town supplied with provisions for a year. This was a good beginning, but Belisarius then, persuaded by Vitalius,\texthebrew{⁠}\footnote{\emph{B. G.} III.13.14.

} committed a serious mistake. He made Ravenna his base, and he could hardly have chosen a less suitable place for offensive operations of which the most important and pressing objects were to succour Rome and recover Naples and southern Italy.

Some of the fortresses in the province of Aemilia, including Bononia, were occupied, but the Illyrian troops who won these successes, having suddenly received the news that their homes were being devastated by an army of Huns, stole away and marched back to their own country. Bononia could no longer be held, and soon afterwards Auximum surrendered to the Goths, who inflicted a severe defeat on a small force which Belisarius had sent to its relief. At the end of the first year of his command the general had little to show but the saving of

p235
Otranto.\texthebrew{⁠}\footnote{He rebuilt indeed the walls and defences of Pisaurum (Pesaro), which had been, like those of Fano, dismantled by Witigis.

} Meanwhile Totila was blockading Rome, now under the command of Bessas, and he had taken Tibur. The fall of this place was due to a dispute between the inhabitants and the Isaurian garrison. The Isaurians betrayed it to the enemy, and all the inhabitants, including the bishop, were put to death in a way which the historian declines to describe on the ground that he is unwilling to ``leave to future times memorials of atrocity.''

Belisarius saw that the Imperial cause in Italy was lost unless he received power ­ful reinforcements and money to pay them. In the early summer of A.D. 545 he wrote to the Emperor setting forth the difficulties of the war. ``I arrived in Italy without men, horses , arms, or money. The provinces cannot supply me with revenue, for they are occupied by the enemy; and the numbers of our troops have been reduced by large desertions to the Goths. No general could succeed in these circumstances. Send me my own armed retainers and a large host of Huns and other barbarians, and send me money.'' With a letter to this effect, he sent John to Constantinople under a solemn pledge that he would return immediately. But John, instead of pressing the urgent needs of his commander, delayed in the capital and advanced his own fortunes by marrying the daughter of Germanus, the Emperor's cousin.

It was probably late in the year that John came at last with a new army. Belisarius had gone over to Dyrrhachium to await his arrival and had sent another importunate message to the Emperor. Isaac the Armenian accompanied John, and the Emperor had sent Narses to the land of the Heruls to secure a host of those barbarians\texthebrew{⁠}\footnote{Procopius does not record their numbers; evidently he had no accurate information (\emph{B. G.} III.13.21).

} to take part in the operations of the following spring.

Totila, in the meantime, had been taking town after town in Picenum and Tuscany. Fermo and Ascoli, Spoleto and Assisi, were compelled to capitulate.\texthebrew{⁠}\footnote{Cp. \emph{Contin. Marcell., sub} 545.

} He offered large bribes to Cyprian to surrender Perusia, and, finding him incorruptible, suborned one of his retainers to assassinate him. But the foul murder did not effect its purpose, as the garrison remained loyal to the Emperor. The Goths had now secured effective command of the Flaminian Way, and it was impossible for Imperial troops

p236
to march from Ravenna overland to the relief of Rome. The only place which the Imperialists still held in the Aemilian province was Placentia, an important fortress, because here the Aemilian Way crossed the Po. Totila presently sent an army against it, and captured it at the end of a year, when the inhabitants were so pressed by hunger that they were driven to cannibalism (May A.D. 545 to May 546).\footnote{\emph{B. G.} III.13.8 and 16.2.

\texthebrew{◂} previous

next \texthebrew{▸}Images with borders lead to more information.

The thicker the border, the more information.

(Details here.)

UP TO:

J. B. Bury
Homepage

Latin \& Greek

Texts

LacusCurtius

Home}

\xchapter{Chapter XIX, Part B}

\xsubsection{(Part 2 of 4)}

It was towards the end of A.D. 545 or early in A.D. 546 that Totila began to besiege Rome in person and with vigour. He had already cut off sea-borne supplies by a considerable fleet of light ships stationed at Naples and in the Liparaean Islands. The whole province of Campania seems to have been subject to the Goths, who, we are told, both here and in the rest of Italy, left the land to the Italians to till peaceably, only requiring them to pay the taxes which would otherwise have been exacted by the Emperor.\texthebrew{⁠}\footnote{\emph{Ib.} 131; this statement, however, is in direct contradiction with 9.3, where the provincials are said to have been robbed of their lands by the Goths and of their movable property by the Imperial soldiers.

} Of the two ports at the mouth of the Tiber, Ostia was in the possession of the Goths, while Portus was held for the Emperor by Innocent. It will be remembered that during the former siege by Witigis, the position was just the reverse; the Romans were in Ostia, the Goths in Portus.

Belisarius despatched Valentine and Phocas, one of his guards, with 500 men by sea to reinforce the garrison of Portus. The troops in Rome numbered 3000, and if Bessus, their commander, had co-operated actively with the leaders at Portus, it might have been possible to secure the passage of foodships up the Tiber. But he refused to allow any of his men to hazard a sortie. Valentine and Phocas, with their small forces, attempted a surprise attack on the Gothic camp, but they fell into an ambuscade which Totila, informed of their plan by a deserter, had set for them. Most of the Romans, including the two leaders, perished.

Not long afterwards Pope Vigilius, who was staying at Syracuse on his way to Constantinople, sent a flotilla of corn-ships to feed the starving city. The Goths saw them approaching

p237
and posted an ambush.\texthebrew{⁠}\footnote{It is not clear what the \textgreek{\texthebrew{Ῥ}ωμαίων λιμήν} was, behind the walls of which (\textgreek{τ\texthebrew{ῶ}ν τειχ\texthebrew{ῶ}ν \texthebrew{ἐ}ντος}) the Goths concealed themselves (\emph{ib.} 15.10), and into which the foodships sailed, according to Procopius. Probably he misconceived the topography. The ships were making for the harbour of Porto, and were captured near the mouth of the channel which divides Porto from the Isola Sacra. This appears to be the interpretation of Hodgkin, \emph{op. cit.} IV.526 . It is difficult to see how the Goths could have hidden themselves in the harbour of Porto (which was presumably under the control of the garrison) as Martroye assumes (\emph{op. cit.} 446). As to the presence of Vigilius in Sicily see below,

p385 .

} The garrison of Portus, who could see the movements of the enemy from the walls, waved garments and signalled to the ships to keep away from the harbour and land elsewhere, but the crews mistook the signals for demonstrations of welcome, and sailing into the trap which had been laid for them were easily captured and slain. A bishop who accompanied the convoy was seized and interrogated by the king. His replies were unsatisfactory, and Totila, convinced that he was lying, punished him by cutting off his hands.

The pressure of hunger in Rome was now so severe that it was decided to ask Totila for a truce of a few days, on the understanding that, if no help arrived before it expired, the city would be surrendered. One of the Roman clergy, the deacon Pelagius, who was afterwards to fill the chair of St. Peter, undertook the mission. As representative of the Roman see at Constantinople he had ingratiated himself in the favour of Justinian, he enjoyed a high reputation in Italy, and had won popularity by employing his considerable wealth to relieve the sufferings of the siege. Totila received him with the courtesy due to a man of his character and influence, but made a speech, if we can trust the historian, which had the effect of preventing any attempt at negotiation.

``The highest compliment I can pay to an ambassador,'' such was the drift of the king's statement, ``is candour. And so I will tell you plainly at the outset that there are three points on which I am resolved and will entertain no parley, but otherwise I will gladly meet any proposals you may make. The three exceptions are: (1) I will show no mercy to the Sicilians; (2) the walls of Rome shall not be left standing; (3) I will not give up the slaves who deserted to us from their Roman masters on the promise that we should never surrender them.'' Pelagius did not conceal his chagrin at these reservations and departed without making any proposals.

p238
It is difficult to suppose that this interview has been quite correctly recorded. Why should Totila have introduced the subject of Sicily, which had no apparent bearing on the surrender of Rome, unless it had been first introduced by Pelagius? If the report of Procopius is true so far as it goes, we must suppose that it is incomplete and that he has omitted to say that the ambassador opened the conversation by mentioning certain conditions for eventual surrender among which were the three points as to which Totila said he could make no concessions.

It is intelligible that Pelagius should have availed himself of this opportunity, whether with or without the authorisation of Bessas, to attempt to safeguard the Sicilians. For it is probable that Totila had made no concealment of his intention to punish them for what he regarded as their black ingratitude to the Goths. They had enjoyed a privileged position under Theoderic and his successors; for no Goths had been settled in the island. But when Belisarius landed they had welcomed him with unanimous enthusiasm, and smoothed the way for his conquest of Italy. Their conduct rankled in the minds of the Goths, and they might well shiver at the thought of the chastisement awaiting them when Totila should have his hands free after the capture of Rome.

The vindictiveness displayed by Totila towards Sicily seems to have been the reason which induced Pelagius to break off the negotiation without pressing for the truce which he had been sent to arrange. His failure drove the citizens to despair. Some of them appeared before Bessas and his officers and implored them either to give them food to keep them alive or to allow them to leave the city or to kill them. They received a cold, unsympathetic reply. ``We cannot agree to any of your suggestions. The first is impossible, the second would be dangerous, the third criminal. But Belisarius will soon be here to relieve the city.'' Throughout the siege Bessas and his subordinate commanders had been profiting by the dire necessity of the inhabitants to fill their own purses. At first they had plenty of corn º in their magazines, and they sold it to the richer people at an exorbitant price.\texthebrew{⁠}\footnote{About 58 shillings a bushel (7 solidi for a medimnus, \emph{i.e.} 6 modii or about 1½ bushels). Occasionally a stray ox was captured in the vicinity of the walls and was sold for over £30. Hodgkin (\emph{op. cit.} iv.532) misunderstands Procopius (III.17.12) and supposes this happened only once.

} Those who could not afford to buy had

p239
to content themselves with bran at a quarter of the price. The mass of the populace fed on cooked nettles, and when the supplies of corn and bran ran short, nettles became the food of all. On this fare, occasionally supplemented by the flesh of a dog or a rodent, they died, or, wasting away, moved about like ghosts. At last the heart of Bessas was moved by the offer of a sum of money to allow the civilians to leave the city. Nearly all took advantage of the permission. Many fell into the hands of the Goths and were cut to pieces, and of the rest it is said that the greater number dropped by the wayside exhausted and died where they lay. ``The fortunes of the Senate and Roman people had come to this.''\footnote{\emph{B. G.} III.17.25.

}

The next event was the landing of Belisarius at Portus. It was his intention on the arrival of John with reinforcements at Dyrrhachium to proceed immediately with all the forces he had to the relief of Rome. But John urged that it would be better to drive the Goths first out of Calabria and southern Italy, which they did not hold strongly, and then march on Rome. The result of these deliberations was a compromise.\texthebrew{⁠}\footnote{Hodgkin's interpretation of what happened is probably right (iv.535) .

} The generals divided their forces. The voyage of Belisarius, who, accompanied by Antonina, first set sail with part of the army, was interrupted by adverse winds which compelled him to put in at Otranto. This port was still being besieged by the Goths, who, on the approach of his fleet, fled to Brundusium, where John presently landed and put them to rout. This victory meant the definite recovery of Calabria. John then marched northwards into Apulia and took Canusium, then southwards into Bruttii, where he defeated the Gothic general who was in command at Rhegium. He appears to have been determined, for other than military reasons,\texthebrew{⁠}\footnote{In

\emph{H. A.} 5.13\texthebrew{‑}14 , Procopius explains John's conduct as due to the fear that Antonina, acting under the instructions of Theodora, might attempt to compass his death. His marriage with the daughter of Germanus had been opposed by Theodora, and Procopius asserts that she threatened to destroy him.

} not to join Belisarius, who was impatiently expecting him on the Tiber; for we cannot suppose that he was deterred from fulfilling his promise by a body of 300 cavalry which Totila had sent to Capua.

Having established himself at Portus, Belisarius decided that his forces were too weak to attack the Gothic camp with any

p240
chance of success, and that the only thing he could attempt was to provision the city. To prevent foodships from ascending the river, Totila had thrown across the stream, some miles above Portus, a wooden boom,\texthebrew{⁠}\footnote{\textgreek{γέφυρα.}} with a tower at either end in which guards were stationed, and below it he had stretched an iron chain from bank to bank. To overcome this obstacle, Belisarius bound together two broad boats on which he constructed a wooden tower higher than the towers of the boom, and on the top he placed a boat filled with pitch, sulphur, resin, and other combustibles. He loaded with provisions and manned with the best of his soldiers two hundred dromons or light warships, on the decks of which he had erected high wooden parapets pierced with holes through which his archers could shoot. When all was ready he stationed some troops near the mouth of the Tiber, in case the enemy should attack Portus. He left Isaac the Armenian in charge of Portus, entrusting Antonina to his care, with strict instructions not to leave the place on any plea, not even if he should hear that Belisarius had been slain. Other troops were ordered to advance along the right bank of the Tiber to co-operate with the ships, and a message had been sent to Bessas bidding him distract the enemy by a sortie. It was the one thing which Bessas was determined not to do.

Belisarius embarked in one of the dromons, and the double barge was slowly urged or hauled upstream.\texthebrew{⁠}\footnote{Procopius does not say how. Hodgkin supposes that it was tugged by some of the dromons.

} Unhindered by the enemy, who did not appear, they reached the iron chain. Here they had to deal with some Goths who had been set on either bank to guard it. Having killed or put them to flight, they hauled up the chain and advanced against the boom, where they were confronted by more serious resistance, for enemy soldiers rushed from their encampments to help the guards in the towers. The double barge was then guided close to the tower on the right bank,\texthebrew{⁠}\footnote{Evidently this side was chosen by Belisarius because his troops on shore could assist.

} the combustibles were set alight, and the boat was dropped on the tower. The tower was immediately wrapt in flame, and the two hundred Goths inside were consumed. Meanwhile the archers in the dromons rained arrows on the Gothic forces which had assembled on the bank till these, terrified at once by the conflagration and by the deadly shower,

p241
turned and fled. The men of Belisarius then set fire to the boom, and the way to Rome was clear.

But in the very moment in which he was rejoi­cing that his difficult enterprise, so skilfully planned and executed, was crowned with success, horsemen galloped up the road from Portus with the tidings that Isaac the Armenian was in the hands of the enemy. Belisarius lost his presence of mind. He did not wait to ask for details. He leaped to the conclusion that Portus had been captured, that his wife had fallen into the hands of the Goths, that he and his army had lost their base and refuge; and he decided that the only thing to be done was to return at once with all his forces and attempt to recover Portus before the enemy had time to organise its defence.\texthebrew{⁠}\footnote{\textgreek{\texthebrew{Ἐ}πιθησόμενος μ\texthebrew{ὲ}ν \texthebrew{ἀ}τάκτοις \texthebrew{ἔ}τι το\texthebrew{ῖ}ς πολεμίοις ο\texthebrew{ὖ}σι}, \emph{B. G.} III.19.31.

} The dromons sailed down the river, to find Portus unharmed and Antonina safe.

What had happened was this. The news of the breaking of the chain and the conflagration of the tower had come \texthebrew{—} perhaps it was signalled \texthebrew{—} to the ears of Isaac. He could not resist the temptation of doing something to win glory for himself, and, in flat disobedience to the express orders of his general, he left the fortress, crossed to the other bank of the channel, and, taking a hundred of the cavalry which Belisarius had posted on the Isola Sacra, attacked a Gothic encampment which was under the command of Roderick.\texthebrew{⁠}\footnote{Apparently in the Isola Sacra, probably opposite Ostia. Hodgkin assumes that it was at Ostia (which Procopius does not mention), but this would have involved crossing the river between Isola Sacra and Ostia. Nothing is said of this.

} The enemy were taken by surprise and retired; Roderick himself was wounded. But as Isaac and his men were plundering the camp, they were in turn surprised by the Goths who returned to attack them. Many were slain and Isaac was taken alive. Roderick died of his wound, and Totila, who valued him highly, avenged him by putting Isaac to death.

The misfortune might have been retrieved if Belisarius, on discovering his mistake, had promptly retraced his course upstream, before the enemy had time to replace the boom or construct new obstacles. But he had not the heart to make another attempt. The shock had been so great and the disappointment so grievous that his physical strength collapsed.

p242
Envious fortune seemed to have snatched the cup from his lips, and he must have felt that, if Isaac's unpardonable disobedience had originated the misfortune, it would have had no serious consequences but for his own precipitate action. He fell ill and a dangerous fever supervened.

It was not long after this that Rome was captured. Bessas as well as the soldiers grew negligent of the routine work on which the safety of a besieged city depends. Sentinels slept at their posts, and the patrols which used to go round the walls to see that watch was kept were discontinued. Four Isaurians, whose nightly post was close to the Asinarian Gate, took advantage of this laxity to betray the city. Letting themselves down from the battlements by a rope, they went in the darkness to the camp of Totila and offered to open the gate. He agreed to pay them well for their treachery and sent two of his followers back with them to report whether the scheme was practicable. But he did not altogether trust them, and it was not till they had twice returned to urge him to the enterprise that he finally decided to make use of their help. On the appointed night four strong Goths were hauled up by the Isaurians, and cleaving the wooden bolts of the gate with axes they admitted their king and the army (December 17, A.D. 546).\footnote{\emph{Contin. Marcell., sub} 547. John Mal. xviii p483 says \textgreek{μην\texthebrew{ὶ} Φεβρουαρί\texthebrew{ῳ}.}}

Bessas and the greater part of the garrison, with a few senators who still had horses , fled through another gate (perhaps the Flaminian). Bessas in his haste left behind him all the treasure which he had spent a year in wringing from the starving citizens. Of the civilian population there were only about 500 left. These took refuge in churches, and sixty of them were killed by the Gothic soldiery when Totila let his troops loose to slay and plunder. He went himself to pray in St. Peter's, where Pelagius, holding the Bible in his hands, accosted him with the words, ``Spare thy people, my lord.'' Totila, thinking of their last meeting, said, ``Now, O Pelagius, thou hast come to supplicate me.'' ``Yes,'' was the reply, ``as God has made me thy servant. But henceforth spare thy servants, my lord.'' Totila then issued an order to stay the slaughter, but he allowed the Goths to plunder at their will, reserving the most valuable treasures for himself. The fact that no acts of violence to women disgraced the capture of Rome redounded to his glory.

p243
Totila hoped that this success would end the war. He despatched Pelagius and another Roman to Constantinople bearing a letter to the Emperor, to the following effect:

``You have already heard what has happened to Rome, and you will learn from these envoys why I have sent them. We are asking you to accept yourself and accord to us the blessings of that peace which was enjoyed in the time of Anastasius and Theoderic. If you consent, I will call you my father and we Goths will be your allies against all your enemies.''

It is clear from this letter that Totila's idea was not to establish a completely independent power in Italy, like those of the Germanic kingdoms in Gaul and Spain, but to restore the constitutional system which had been in force under Theoderic and Athalaric. The capitulations of A.D. 497 were to be renewed, the Imperial authority was still to be nominally supreme. The ambassadors were instructed to intimate that, if the offer of peace were rejected, Totila would raze Rome to the ground and invade Illyricum. Justinian did not detain them long. He sent them back with a curt answer that as full powers for conducting the war and concluding peace had been committed to Belisarius, Totila might apply to him.\footnote{It may be conjectured that Procopius, directly or indirectly, gathered a good deal of his information about the siege of Rome from this embassy. We are not told whether Totila did make overtures to Belisarius. The envoys may have been back in Italy in the first half of February.

}

In the meantime the slow but steady progress of the Imperial cause in southern Italy, where, if John had not taken any risks or achieved any striking success, Lucania had been detached from Gothic rule, demanded Totila's presence in the south. He did not want to lock up a garrison in Rome or to leave it for his enemies to reoccupy, and he decided to demolish it. He began by pulling down various sections of the walls,\texthebrew{⁠}\footnote{Procopius estimates the razed portions as about one-third of the whole circuit \texthebrew{—} probably an excessive estimate (Hodgkin, iv.566) .

} and was about to burn the principal buildings and monuments when envoys arrived with a letter from Belisarius, who was recovering from his illness at Portus. The tenor of the letter is reported thus:

``As those to whom a city owes the construction of beautiful buildings are reputed wise and civilised, so those who cause their destruction are naturally regarded by posterity as persons devoid of intelligence, true to their own nature. Of all cities under

p244
the sun Rome is admitted universally to be the greatest and most important. She attained this hill-town not suddenly nor by the genius of one man, but in the course of a long history throughout which emperors and nobles by their vast resources and employing skilful artists from all parts of the world have gradually made her what you see her to\texthebrew{‑}day. Her monuments belong to posterity, and an outrage committed upon them will rightly be regarded as a great injustice to all future generations as well as to the memory of those who created them. Therefore consider well. Should you be victorious in this war, Rome destroyed will be your own loss, preserved it will be your fairest possession. Should it be your fortune to be defeated, the conqueror will owe you gratitude if you spare Rome, whereas if you demolish it, there will be no reason for clemency, while the act itself will have brought you no profit. And remember that your reputation in the eyes of the world is at stake.''

This is an interesting document, whether it reproduces closely or not the drift of the actual letter of Belisarius. Totila read that letter again and again; it gave him a new point of view; and the remonstrance of civilisation finally defeated in his breast the barbarous instincts of his race. He bade the work of vandalism cease.

The greater part of the Gothic army was left, at a place called Algedon, about eighteen miles west of Rome,\texthebrew{⁠}\footnote{The name \textgreek{\texthebrew{Ἀ}λγηδών} (? Alcedum) is otherwise unknown. Martroye (\emph{op. cit.} 466) thinks it may be identified with Castel Malnome, not far from Porto. Mount Algidus, which lies east of Rome, is out of the question.

} to watch Belisarius and prevent him from leaving Portus. With the rest Totila marched southward and soon recovered Lucania, Apulia, and Calabria, except Otranto and Taranto, in which John entrenched himself.\texthebrew{⁠}\footnote{Tarentum was unwalled, and he fortified only that part of it which lay on the isthmus. Totila probably departed for the south towards the end of February. We have to allow time for the return of the envoys who were despatched in December to Constantinople.

} Then leaving a detachment of 400 men in the hill-town of Acherontia,\texthebrew{⁠}\footnote{Celsae nidum Acherontiae (Horace); now Acerenza.

} on the borders of Apulia and Lucania, he marched northwards. Was his design to surprise Ravenna, as the historian

p245
intimates, or to re-establish his command of the Flaminian Way, which was threatened by a recent success of the Imperialists? They had recovered Spoletium.

But grave news from Rome compelled him to postpone his purpose. He had left Rome uninhabited, its walls partly destroyed, and all the gates removed. Belisarius, whose health was now returned, visited the desolate city and decided to occupy it and put it in a state of defence.\texthebrew{⁠}\footnote{In April. For Totila left Rome towards the end of February, and it remained uninhabited for forty days (\emph{Cont. Marcell.}, \emph{s.a.}).

} The plan seemed wild, but it was carried out. He transferred his army from Portus to Rome, where he was able to establish an abundant market, as there was no longer any obstacle to the importation of food from Sicily. The market attracted the people of the surrounding districts to come and settle in the deserted houses, and in less than four weeks the portions of the wall which the Goths had pulled down had been roughly reconstructed, though without mortar. New gates, however, could not be made so quickly, for lack of carpenters, and when Totila appeared in front of a gateless fortress he expected to capture it with ease. Belisarius placed in the gateways men of notable valour. Two days the Goths spent themselves in furious attacks, suffering great losses, but failed to carry any of the gates. After an interval of a few days, during which they cared for their wounded and mended their weapons, they renewed their assault. Totila's standard-bearer fell mortally wounded, and there was a fierce fight over the corpse. The Goths recovered the standard, but their whole army presently retreated in disorder. They were soon flying far afield pursued by the victors. Rome for the time was saved.\texthebrew{⁠}\footnote{Probably in May 547.

} Belisarius furnished it with new gates and sent the keys to the emperor.

This was the first check that Totila had experienced. While he won battles and captured cities, his followers regarded him as a god, but now in the hour of defeat they forgot all he had done and were immoderate in their criticism. The nobles reproached him bitterly with his blunder in leaving Rome in such a condition that the enemy could occupy it; he should either have utterly destroyed it or held it himself. But though there was open discontent, there was no thought of revolution. Having demolished the bridges across the Tiber, except the

p246
Milvian, Totila withdrew to Tibur, which he refounded and made his headquarters.

In the course of summer he went to press the siege of Perusia, which Gothic troops had been blockading for some time past and which was now distressed by shortage of food. But his attention was soon diverted to the south, which was to be the scene of the principal operations of the war during the winter. He had interred in Campania those senators and their families whom he had carried off after the capture of Rome. The general John, who had been engaged in besieging Acherontia, determined to rescue them while Totila was still occupied in the north. Moving rapidly, he defeated a squadron of Gothic cavalry at Capua, and successfully delivered many of the Roman captives,\texthebrew{⁠}\footnote{From \emph{Cont. Marcell. sub} 548, one would infer that only some of the women were rescued; nonnullas liberat senatrices.

} whom he immediately sent by sea to the safety of Sicily. It was a blow to Totila, for he regarded those prisoners as hostages who might be useful hereafter, and he marched hastily and stealthily from Perusia, with 10,000 men, into Lucania, where John was encamped. It was a complete surprise, and few of his enemies would have escaped if he had not committed the blunder of attacking the camp by night, for he outnumbered them by ten to one. But in the darkness about 900, including John, were able to escape, and 100 were slain. The prisoners taken were very few. Among them was an Armenian general, Gilak, who could speak no language but his own, and knew only one Greek word, stratêgos , ``general.'' The only intelligible answer which his captors could extract from him was Gilakios stratêgos . They put him to death a few days later.

The importance of holding Calabria had been realised by John, and Belisarius appears to have anticipated before the end of the year ( A.D. 547) that the main operations in spring would probably be in that region. Justinian, urged by his appeals, sent reinforcements of more than 2000 men at the beginning of the winter.\texthebrew{⁠}\footnote{Among these, more than 1000 guards (doryphoroi and hypaspistai) under Valerian, Master of Soldiers in Armenia, III.27.3.

} Early in the year, Belisarius committed the charge of Rome to Conon and sailed for Sicily with 900 men. Proceeding thence up the eastern coast of Bruttii he found his voyage impeded by strong north winds, and instead of making for Tarentum, as he had intended, he landed at Croton, an unwalled

p247
town. As the neighbourhood could not furnish provisions for his army, he sent his cavalry northward into the mountains to forage, expecting that if they met the enemy in the narrow defiles they would be able to repulse them.\texthebrew{⁠}\footnote{In \emph{Itin. Antonini} the distance of Rossano from Thurii is given as 12 Roman miles. Procopius gives its distance from the port of Thurii as 60 stades (nearly 9 miles). The name of the port was Ruscia (III.28.8 \textgreek{\texthebrew{Ῥ}ουσκια \texthebrew{ἐ}στ\texthebrew{ὶ} τ\texthebrew{ὸ} Θουρίων \texthebrew{ἐ}πίνειον}, 30.12 \textgreek{\texthebrew{ἐ}π\texthebrew{ὶ Ῥ}ουσκίαν \texthebrew{ἀ}νήγοντο}); the inland territory behind it was called (castrum) Ruscianum. Procopius designates the fort as \textgreek{τ\texthebrew{ὸ ἐ}π\texthebrew{ὶ Ῥ}ουσκιαν\texthebrew{ῆ}ς φρουρίον} (29.21, and 30.19), or \textgreek{τ\texthebrew{ῷ Ῥ}ουσκιαν\texthebrew{ῷ} φρουρί\texthebrew{ῳ}} (30.5, so one MS., but the other omits \textgreek{\texthebrew{Ῥ}}. and it may be a gloss). In 30.2 \textgreek{ε\texthebrew{ὐ}θ\texthebrew{ὺ Ῥ}ουσκιαν\texthebrew{ῆ}ς κατ\texthebrew{ὰ} τάχος \texthebrew{ἔ}πλει}, the text need not be altered, as ``sailed straight to the Ruscian territory'' is good sense. Haury, wrongly assuming that Rusciana was that name of the port, has perversely corrected the readings of the MSS. in three places (in 28.8 he reads \textgreek{\texthebrew{Ῥ}ουσκιανή}, in 30.12 \textgreek{\texthebrew{Ῥ}ουσκιαν\texthebrew{ὴ}ν}, and in 30.5 \textgreek{\texthebrew{ἐ}ν τ\texthebrew{ῷ ἐ}π\texthebrew{ὶ Ῥ}ουσκιαν\texthebrew{ῆ}ς φρουρί\texthebrew{ῳ}}). Martroye has made the opposite mistake and called the fort Ruscia. The port of Thurii can hardly have been so far from that town as the sea-board near Ruscia (cp.
Nissen, \emph{Ital. Landesk.} II.2, p923 ).

} Totila was there with his army bent on taking the hill-town of Ruscianum \texthebrew{—} the modern Rossano \texthebrew{—} in which John had placed a garrison.\texthebrew{⁠}\footnote{Procopius names two passes between Bruttii and Lucania, Labūla, and Petra Sanguinis (\textgreek{πέτρα Α\texthebrew{ἵ}ματος}, III.28.7). The former may mean the east coast road (cp. Nissen, \emph{op. cit.} p926). The latter must be on the inland road coming down from Neruli (=Rotonda, on the Laus) and Muranum (=Morano) to Interamnium, which corresponds to Castrovillari, and is situated on the two brooks which unite to form the Sybaris (now the Coscile). The road went on to Capraria (now Tarsia, east of Spezzano) and Consentia. The boundary between E. Lucania and E. Bruttii was the Crathis; on the west, I suspect that the division lay along the N-S road from Neruli to Interamnium.

} The disparity in numbers was immense, yet the small body of horse inflicted a severe defeat on the Gothic host, of whom more than 200 fell. But the Goths enjoyed a speedy revenge. The Romans, elated by their victory, neglected their night watches and did not pitch their tents in one place, so that Totila was able to surprise and nearly exterminate them. On hearing the news, Belisarius, his wife, and infantry º ``leapt'' into the ships\texthebrew{⁠}\footnote{\textgreek{\texthebrew{Ἐ}ς τ\texthebrew{ὰ}ς να\texthebrew{ῦ}ς \texthebrew{ἐ}σεπήδησεν}. Procopius gives the distance from Croton to Messina as 700 stades, about 92 English miles, which is virtually correct.

} and reached Messina in one day. Totila laid siege to Rossano (probably in May).

Soon afterwards a new contingent of about 2000 arrived in Sicily from the East. Much larger forces were needed against a leader of Totila's capacity; Belisarius was weary of conducting a war in which, though he might gain local successes, he was never strong enough to take full advantage of them; and he

p248
had decided that Antonia should return to Constantinople and implore the Empress to use all her influence to secure the sending of such an army as the situation required. They proceeded together to Otranto, and there she embarked on a journey which was to prove fruitless, for she arrived to find that Theodora was dead.\footnote{She died on June 28; see above,

p66 .

}

The garrison of Rossano, in dire need of food but expecting aid from Belisarius, promised Totila that they would surrender on a certain day,\texthebrew{⁠}\footnote{\textgreek{Μεσούσης μάλιστα τ\texthebrew{ῆ}ς το\texthebrew{ῦ} θέρους \texthebrew{ὤ}ρας} is generally taken with \textgreek{\texthebrew{ἐ}νδώσειν} in III.30.5. If so, it cannot mean, as it ought, the summer solstice. I think that it should be taken with \textgreek{\texthebrew{ὡ}μολόγησαν}, and that the bargain was made about the end of June.

} if no relief arrived, on condition that they should be allowed to depart in safety. The commanders of the garrison were Chalazar, a Hun, and the Thracian Gudilas. The attempts of Belisarius and John to bring help were frustrated, and they then hit on the plan of forcing Totila to raise the siege by diverting his attention elsewhere. Belisarius sailed to Rome, and John with Valerian \texthebrew{—} a general who had been sent to Italy six months before\texthebrew{⁠}\footnote{See above,

p246, n2 .

} \texthebrew{—} set out to relieve the fortresses in Picenum, which enemy forces were besieging. But Totila was bent upon the capture of the Bruttian fortress, and he contented himself with despatching 2000 cavalry in the rear of John.

The garrison of Rossano, confident that help was approaching, failed to keep their promise; the appointed day passed; and then, when they knew they could no longer hope, they threw themselves on Totila's mercy. He pardoned them all except Chalazar, whom he shamefully mutilated and put to death. Those soldiers who were willing to become Gothic subjects remained in the place; the rest were deprived of their property and went to Croton.\footnote{Eighty in number. The garrison consisted of 300 Illyrian cavalry and 100 foot soldiers. There were also in the place many Italians of good family; these Totila punished by confiscating their property.

}

Rome needed the presence of Belisarius. Some time before, the garrison had mutinied and slain Conon their commander. They had then sent some clergy to Constantinople to demand a free pardon for the murder and the payment of their arrears, with the threat that they would deliver the city to the Goths if these conditions were not accepted. The Emperor accepted them. Belisarius then arrived. He saw to it that the city was

p249
furnished with a good supply of provisions in case it should be again besieged, and he probably weeded out the garrison. When he left Italy for ever, early in A.D. 549, Rome was held by 3000 chosen troops, under the command of Diogenes, one of his own retainers, whose intelligence and military capacity he trusted.

Antonina had procured without difficulty her husband's recall. Theodora's death meant the ascendancy of the party which was attached to Germanus, and the enemies and critics of Belisarius could now make their influence felt. What had this great general accomplished in five years? He had simply navigated about the coasts of Italy, never venturing to land except when he had the refuge of a fortress. Totila desired nothing so much as to meet him in battle, but he had never taken the field. He had lost Rome, he had lost everything.\texthebrew{⁠}\footnote{Procopius, \emph{B. G.} III.35.1, 2, repeated with additions in \emph{H. A.} 5. 1\texthebrew{‑}6

and

17 .

} He might vanquish a general of mediocre capacity like Witigis, but it was a different story when he had to do with a foe of considerable talent and unflagging energy like Totila. Belisarius might have much to say in extenuation of his failure, but the broad fact was that he had failed. Knowing that there was no chance of his receiving such reinforcements as might enable him to retrieve his reputation, he was glad to bid farewell to Italy.

Soon after his departure, Perusia fell, after a siege of four years.\footnote{A.D. 545\texthebrew{‑}549.

}

In the summer after the departure of Belisarius, the king of the Goths appeared for the third time before the walls of Rome.\texthebrew{⁠}\footnote{At the beginning of the 15th year of the war (Proc. \emph{B. G.} III.36.1), \emph{i.e.} end of June or early in July, 549.

} He was determined to capture it, but he had abandoned all those thoughts of destroying it which had moved him when he first laid siege to it. He had laid to heart the letter of the Imperial general, which other opinions had perhaps reinforced;\texthebrew{⁠}\footnote{Cp. below,

§ 8, p258 .

} he had come to realise \texthebrew{—} as Theoderic and Alaric had realised \texthebrew{—} the meaning of Rome.

The garrison was valiant, and the commander Diogenes had made provident preparations for an eventual siege. As there was only a small population now, besides the garrison, there were large areas of waste land in the city, and these were sown

p250
with grain. When repeated attempts of the Goths to storm the walls were foiled by the valour of the soldiers, Totila resigned himself to the prospect of a long blockade. It was uncertain whether relief forces would arrive from the East under a new commander-in\texthebrew{‑}chief , but as he had captured Portus, he was in a much more favourable position for conducting a blockade than he had been three years before.

The blockade lasted a long time, but the city fell into his hands at last. The circumstances of the previous capture were repeated. Isaurian treachery again delivered Rome to the Goth. Some Isaurian soldiers, who were keeping watch in the south of the city at the Porta Ostiensis \texthebrew{—} which was already known by its modern designation from the Church of the Apostle Paul \texthebrew{—} discontented because they had received no pay for years, and remembering the large rewards which Totila had bestowed on their fellow-countrymen , offered to open the gate. On a pre-arranged night, two barques\texthebrew{⁠}\footnote{\textgreek{Μικρ\texthebrew{ὰ} πλο\texthebrew{ῖ}α}, ``little boats,'' is the reading of L (the Laurentian MS.), and is probably right, but K (the Vatican MS.) has \textgreek{π. μακρά} (long boats or war-vessels).

} were launched in the Tiber, probably to the north of the Porta Flaminia.\texthebrew{⁠}\footnote{This conjecture seems more likely than that of Hodgkin, who supposes that the boats were launched south of the city and ordered ``to creep up the river and blow a loud blast from their trumpets as near as possible to the centre of the City'' (p615) . In this case the garrison would have supposed that the attack was to be made, as Hodgkin says, near the Aventine. The vague words of Procopius are consistent with either interpretation; but he says nothing about ``the centre of the city''; on the contrary, his words are \textgreek{\texthebrew{ἐ}πειδ\texthebrew{ὰ}ν το\texthebrew{ῦ} περιβόλου \texthebrew{ἄ}γχιστα \texthebrew{ἤ}κωσι} (III.36.9) and \textgreek{\texthebrew{ἐ}πε\texthebrew{ὶ} τ\texthebrew{ῆ}ς πόλεως \texthebrew{ἄ}γχι \texthebrew{ἐ}γένοντο} (\emph{ib.} 12). As it was just as easy for Totila to raise the false alarm in the north, it seems highly improbable that he would have chosen to attract all the soldiers of the garrison to the Aventine, which is quite close to the Gate of St. Paul.

} They were rowed down as close to the city as possible, and then trumpeters who had been embarked in them sounded a loud blast. The alarm was given, and all sections of the garrison rushed to the defence of the walls in the threatened quarter, in the north-west. Meanwhile the Gothic army had been quietly assembled in front of the gate of St. Paul; the Isaurians unlocked it, and the army marched in (January 16, A.D. 550).\footnote{\emph{Cons. Ital.} p334. The year given here is 549, but we should doubtless read p. c. Basilii viiii. =550 (see Körbs, \emph{op. cit.} p54).

}

It was easy to anticipate that any of the garrison who succeeded in escaping would make for Centumcellae, the only fortress that remained to the Imperialists in the neighbourhood of Rome,

p251
and Totila had posted some troops along the western road to intercept fugitives. The precaution was effective; a few escaped the ambush, among whom was Diogenes. In Rome itself there was great slaughter, but a band of four hundred cavalry led by Paul, who had belonged to the household of Belisarius and was the right-hand man of Diogenes, occupied the Mausoleum of Hadrian and the adjacent Aelian Bridge. Here they held out for two days. Totila expected that the cravings of hunger would compel them to surrender, and kept troops posted on the eastern bank. They thought of eating their horses , but could not make up their minds to taste the unaccustomed food. On the evening of the second day they resolved to court a heroic death, to make a dash against the enemy and fall fighting. They embraced one another, said their last adieux, and prepared for the charge. Totila was watching them and divined their intention. He knew that desperate men, who had devoted themselves to death, would decimate his army. He therefore sent a messenger offering that if they would lay down their arms and take an oath never to fight against the Goths again, he would let them depart unharmed for Constantinople, or if they would fight for him, he would treat them on a perfect equality with the Goths. The offer was gladly accepted. At first all elected to go home, but on further reflexion they changed their minds. They could not bring themselves to undertake the long journey without horses or arms, they feared its perils, and if they had any hesitation about going over to the enemy, they remembered that the Imperial treasury had withheld their pay for years. Only Paul himself and one other resisted the lure of the barbarian and returned to Byzantium.

Totila had no longer any thought of destroying Rome or of leaving it undefended. His position was much stronger than it had been three years before, and he had come to realise the prestige which the possession of the Imperial city conferred in the eyes of the world.\texthebrew{⁠}\footnote{See below,

p258 .

} He was now bent on rebuilding and repopulating it. He sent for the senators and other Romans who were still kept under guard in the fortresses of Campania.\footnote{\emph{B. G.} III.36.29; 37.3. This shows that all the captives had not been rescued by John in 547.

}

He was planning to carry the war into Sicily, but he first made

p252
a new proposal of peace, just as he had done after his former capture of Rome. On this occasion his envoy was not even admitted to the presence of the Emperor, who had just appointed a new commander-in\texthebrew{‑}chief to succeed Belisarius. He had thought of entrusting the conduct of the war to his cousin, Germanus, but changed his intention and selected Liberius, the Roman senator, who fourteen years before had come to him as an envoy of Theodahad, and since then had remained in his service.\texthebrew{⁠}\footnote{See above,

c. xviii p164 . He was a Patrician, and had been appointed Prefect of Alexandria, c. 541; see below,

p380 .

} It was a curious appointment, for Liberius, who had served in civil capacities under Odovacar and Theoderic, had no military experience, and he was now an octogenarian; the ground of his nomination must have been that as an Italian he would inspire the Italians with confidence.

Totila meanwhile was making preparations for his next campaign. He collected a fleet of 400 ships of war and some large merchant vessels, which he had recently captured from the enemy, to convey his troops across the Sicilian Straits. It was perhaps about the end of March that, having presided at horse-races in the Circus Maximus, he left Rome. Before marching southwards he turned aside with the hope of redu­cing Centumcellae, which was now under the command of Diogenes. This valiant officer refused to surrender until he had communicated with Constantinople, but agreed that, if by a certain date no reinforcements should arrive, he would leave the city to the Goths. Totila consented, hostages were interchanged, and the Gothic army marched to Rhegium, which it may have reached early in May.\footnote{The distance of Cività Vecchia from Rome is about 75 kils.; from Rome to Rhegium less than 700.

}

On the southern coasts of Italy the most important places still held for the Empire were Hydruntum, Rhegium, Tarentum, and Croton. The Goths now laid siege to Rhegium and captured Tarentum. Without waiting for Rhegium to fall, Totila crossed to Messina, which he failed to take. But he was at last able to gratify one of the dearest desires of his heart and wreak vengeance upon the Sicilians for the welcome they had given

p253
Belisarius fifteen years before. His army ravaged the island without resistance. Meanwhile Rhegium, which was short of provisions, surrendered.

The news of these mena­cing successes seems to have made a greater impression at Constantinople than the recent capture of Rome. The Emperor reverted to his former plan of sending Germanus to the West as commander-in\texthebrew{‑}chief . But Germanus could not start until he had collected an army sufficiently strong to end the war, and in the meantime Liberius was despatched to defend Sicily. He had hardly set sail before it was recognised that he was too old and inexperienced, and Artabanes,\texthebrew{⁠}\footnote{Artabanes had recently been implicated in a conspiracy; see
p67

above.

} who was appointed Master of Soldiers in Thrace, was sent to supersede him.

Germanus was now regarded as the probable heir to the Imperial throne. The death of Theodora had removed the adverse influence which might have withheld the Emperor from favouring his claim.\texthebrew{⁠}\footnote{He was never formally designated, and Justinian never committed himself. But everything points to the fact that it was generally taken for granted that he would succeed, and that Justinian was contented that this opinion should prevail.

} He had already established his reputation by suppressing the Moorish rebellion of Stutzas, and he was ambitious of enhancing it by recovering Italy and succeeding where Belisarius had failed. As the prestige of the dynasty was involved, the Emperor was prepared to spend money in a less stinting spirit than he had shown hitherto in the conduct of the Italian war; and Germanus had considerable private resources which he did not hesitate to devote to the collection of troops. The raising of an army for services not connected with the defence of the frontiers had come to be the task of the commander who was to lead it. None of the standing troops in the East could be withdrawn, although some of the cavalry squadrons stationed in Thrace might be spared. Germanus, with his sons, Justin and Justinian, busily recruited volunteers in the highlands of Thrace and Illyricum, and bands of barbarians from the Danubian regions flocked to his standards. The king of the Lombards promised a thousand heavily armed warriors. The private retainers of many generals left their less illustrious masters to attach themselves to the service of Germanus.

But besides these preparations for a vigorous military offensive,

p254
the plan of Germanus included what might be called a moral offensive, on which he counted much and with good reason. He contracted a Gothic marriage. He took as his second wife the queen Matasuntha. As the reluctant consort of Witigis she had been once queen of the Goths, but it was as the granddaughter of king Theoderic and sister of king Athalaric that she had the strongest claims on their loyalty and affection. If her mother had brought her up in the ways of Roman civilisation, she was of the purest Amal lineage, and Germanus might confidently hope that the effect of his coming as her husband, and presumably in her company, would be to undermine the allegiance of many of the Goths to Totila, or at least to embarrass their minds in such a way as to impair their military vigour. They would feel that they were fighting no longer merely against Greeks, but against the granddaughter of their greatest king. And they would calculate that, as Germanus was marked out to succeed to the Empire on Justinian's death, Matasuntha would presently share the Imperial throne.\footnote{Wroth has an interesting conjecture on the coins of Matasuntha (bearing her name in monogram), which are generally assumed to have been issued in 536\texthebrew{‑}540, when she was queen of Witigis. Pointing out that there is nothing Italian about the coins, he conjectures that they may have been minted at Constantinople in 550. The issue of such coins would have been a very natural way of asserting the queen's claim to the Ostrogothic throne.

}

When the news of the marriage reached Italy it seems to have produced the effect which was anticipated. Many Goths began to ask themselves whether it would be well to continue their resistance. And the reports which arrived of the Imperial preparations for prosecuting the war affected the numerous soldiers who had deserted the Roman cause to serve under Totila. They managed to send messages to Germanus that as soon as he landed in Italy they would go over to him and fight again under the standards they had abandoned. Diogenes, who had agreed to surrender Centumcellae on a certain day, declared himself absolved from the covenant because Germanus was coming.

But Germanus was not to come. He was at Sardica, and his army was ready. It was the autumn of A.D. 550.\texthebrew{⁠}\footnote{Perhaps in September. Cp. Körbs, \emph{op. cit.} 48\texthebrew{‑}49.

} He had announced that he would start in two days, when he suddenly fell sick, and the disease proved fatal. His death meant much.

p255
It meant particularly the destruction of the hopes which were swaying opinion in Italy both among Italians and Goths.\footnote{Matasuntha bore a posthumous son who was called by the name of his father, and is mentioned in the \emph{Getica} of Jordanes (which seems to have been composed in A.D. 551, or at latest 552, before the issue of the war for Italy was decided). The words of Jordanes are (§ 314): post humatum patris Germani natus est filius idem Germanus, in quo coniuncta Aniciorum gens cum Amala stirpe spem adhuc utriusque generi domino praestante promittit. The exact connexion of Germanus with the Anician gens is unknown; Mommsen conjectured that his mother might have been a daughter of Juliana Anicia, daughter of the Emperor Olybrius and Placidia. From the fact that Jordanes looked upon this infant as the hope of the Ostrogothic race, Schirren (\emph{De ratione quae inter Iordanem et Cassiodorum intercedit}) drew the hardly justifiable conclusion that the motif of the book was political \texthebrew{—} to promote the idea of reconciling Goth and Italian under the rule of Germanus.

}

The plans for the prosecution of the war were disconcerted by the death of the commander-in\texthebrew{‑}chief , and Justinian appointed no one to replace him for some time. But in the meantime it was arranged that John, the nephew of Vitalian, who was now the Master of Soldiers in Illyricum, and was to have served under his father-in\texthebrew{‑}law , Germanus, should, with his brother-in\texthebrew{‑}law , Justinian, lead the army to Italy. John had proved himself an able soldier, and if he and Belisarius had been able to work cordially together, it is probable that the duration of the war would have been considerably curtailed. He was not appointed to the supreme power because it was felt that he did not possess the requisite prestige to command the obedience of the other generals.

When the troops were collected in Dalmatia it was late in the year, and it was thought better to spend the winter there than to march immediately to Venetia. There was no sufficient supply of ships at Salona to transport the army across the Hadriatic.

Meanwhile Totila had been wreaking his vengeance upon Sicily. When he was besieging Syracuse, Liberius arrived, and seeing that he was not strong enough to help the city he sailed on to Panormus. Artabanes, who, as we saw, had been appointed to replace Liberius, was already on his way, but his ships were caught off the coast of Calabria by a storm which drove them

p256
back to Greece. The Goths succeeded in capturing only four fortresses, probably places of secondary importance, in which they placed garrisons, and having lived in the island for many months,\texthebrew{⁠}\footnote{The dates of Totila's arrival in Sicily and his departure are not marked very clearly by Procopius, but we can deduce from the data that he must have reached the island in May and left it before the end of the year.

} they returned to Italy laden with booty and provisions.

During the summer and autumn of this year ( A.D. 550) the Imperial generals in Italy were inactive, though the absence of Totila in Sicily was an opportunity for an enterprising leader. Then the news arrived that the Emperor had appointed the Armenian eunuch Narses to the supreme command. The appointment was universally welcomed. Narses, the Grand Chamberlain, appears to have been one of the most popular ministers at Justinian's court. He was celebrated for his generosity, he did not make enemies, and such was his reputation for piety that it was believed that the Virgin Mother herself watched over his actions and suggested the right moment for engaging in battle.\texthebrew{⁠}\footnote{Evagrius, \emph{H. E.} IV.24.

} He was a friend of John, whom, as it will be remembered, he had followed Belisarius to rescue at Rimini, and of whose loyal co-operation he was assured. This fact, we may conjecture, had a good deal to do with his appointment. Narses had the qualities of a leader, but he had not much military experience; the advice of John would remedy this deficiency.

John had been ordered to await the arrival of the now commander-in\texthebrew{‑}chief at Salona, but Narses was delayed on his way, at Philippopolis, by an invasion of Kotrigur Huns,\texthebrew{⁠}\footnote{See below,

p303 . Narses had left Constantinople about June acc. to our text of John Malalas xviii p484; but Theophanes, who was copying Malalas, says April (A.M. 6043).

} and it was probably late in A.D. 551 that he arrived in Dalmatia.

Fortune had steadily favoured the Goths for the last four years. In A.D. 547 the Imperialists held in central Italy Ravenna, Ancona, and Ariminum, Spoletium and Perusia, Rome itself with Portus, Centumcellae; in the south Otranto, Taranto, the province of Bruttii, and Sicily. In A.D. 551 the only important places they held on the mainland were Ravenna, Ancona, Otranto, and Croton, while in Sicily they had lost four strongholds; and Totila, on returning from Sicily, had sent an army to besiege Ancona. This tide of success was now about to turn.

p257
Ever since Totila had crossed the Po after has accession, the war had been waged entirely to the south of that river, and the conditions which prevailed in northern Italy are obscure. Here the situation was complicated by the intervention of a third power. As all at forces of the Ostrogoths were demanded by the struggle in the south, the Franks had seized the opportunity to extend their power into Italy. Theodebert, who followed the progress of the war attentively, had occupied the province of the Cottian Alps, a part of Liguria, and the greater part of Venetia.\texthebrew{⁠}\footnote{Procopius, \emph{B. G.} III.33.7 ; IV.24. 4

and 6\texthebrew{‑}8 . The Romans still held coast places in Venetia.

} The only important cities which the Goths still retained seem to have been Verona and

Ticinum . Some time afterwards a treaty was concluded between the Franks and Goths, by which Totila acquiesced in the provisional occupation by the Franks of the territory which had been seized, and the two powers agreed, in case the war ended with a Gothic victory, to make a new permanent arrangement.\texthebrew{⁠}\footnote{Procopius, \emph{ib.} 9\texthebrew{‑}10

and 27 .

} Far-reaching plans are attributed to Theodebert. He was incensed at Justinian's assumption of the titles Francicus and Alamannicus , with the implication that the Franks and their subjects the Alamanni had been subjugated and were vassals of the Empire, and he expressed his formal independence by issuing gold coins with his own bust and his own name.\texthebrew{⁠}\footnote{Agathias, I.4 .

Procopius, \emph{ib.} III.33 . Cp. Keary, \emph{Coinages of Western Europe}, p22.

} He was the first German king to venture on this innovation, which from a commercial point of view was hazardous. It was said that he formed the project of leading the German nations, the Lombards, the Gepids, and others through the Illyrian countries and attacking Constantinople itself.

We possess one diplomatic document, belonging to this period, which records the Italian conquests of the Franks. The Emperor had written to Theodebert requesting information as to the extent of his dominions, and Theodebert's reply has been preserved,\texthebrew{⁠}\footnote{\emph{Epp. Mer. aeui}, III; \emph{Epp. Austras.} 20. Theodore and Solomon are mentioned as the Imperial envoys to the Frank court.

} in which he enumerates Pannonia and the northern parts of Italy among the countries which he has subjugated.

After his capture and abandonment of Rome in A.D. 547,

p258
Totila had proposed to espouse the daughter of one of the Merovingian kings who is not named, but we are entitled to presume that he was Theodebert. The offer was refused, on the ground that Totila would never succeed in the subjugation of Italy, seeing that he had shown himself so foolish as to let the great capital slip from his hands.\texthebrew{⁠}\footnote{Procopius, III.37.1\texthebrew{‑}2.

} This criticism helped to open the Ostrogoth's eyes to the importance of Rome. In the following year, Theodebert died and was succeeded by his son Theodebald.\texthebrew{⁠}\footnote{Gregory of Tours, \emph{Hist. Fr.} iii.36

(for date,

37 \emph{ad fin.} ), says he died of a long illness;

Agathias, I.4 , says he was killed, hunting, by a wild bull.

} To him Justinian sent an ambassador to complain of the encroachments of his father in northern Italy, to demand the evacuation of the cities, and to request him to fulfil the promises of Theodebert and co-operate in the Italian war.\texthebrew{⁠}\footnote{The date of this embassy (which Hodgkin assigns to 551) cannot be inferred from the fact that Procopius notices it \emph{after} the embassy of Totila in 551 (IV.24.11); for here he digresses, \emph{à propos} of the Franks, and goes back to Theodebert. I have little doubt that the embassy was sent in 549\texthebrew{‑}550, when the preparations were afoot for the expedition of Germanus, and that it bore to Theodebald Justinian's congratulatory letter on his accession, in which the Emperor took occasion to say hard words of his father's conduct, and to which Theodebald's reply is preserved (\emph{Epp. Austras.} 18).

} Theodebald promptly sent an embassy to Constantinople.\texthebrew{⁠}\footnote{It may be conjectured that it was on this occasion that Theodebald attached some Angles to his embassy, to show the Emperor that his authority extended to Britain. See Procopius, iv.20,10. These Angles doubtless supplied Procopius with the material for his curious account of that island.

\texthebrew{◂} previous

next \texthebrew{▸}Images with borders lead to more information.

The thicker the border, the more information.

(Details here.)

UP TO:

J. B. Bury
Homepage

Latin \& Greek

Texts

LacusCurtius

Home} The course of the negotiations is unknown, but the Franks remained in Italy.

\xchapter{Chapter XIX, Part C}

Short URL for this page:

bit.ly/BURLAT19C

\xsubsection{(Part 3 of 4)}

Totila realised that a supreme effort was now to be made to destroy the Ostrogothic power in Italy. The appointment of Narses was hardly less significant than the appointment of Germanus. He had always understood the importance of reconciling the Italians to his rule, and he now urgently pressed forward the rebuilding of Rome in order to ingratiate himself with the Romans. His immediate military objects were the capture of Ancona and Croton, two of the few valuable places that were still left to the Empire. In the autumn of A.D. 551, his forces, as we saw, were besieging Ancona, but it is probable that he had not yet sent an army against Croton. At the same

p259
time, he was employing his fleet. Three hundred vessels sailed to the shores of Greece. The rich island of Corcyra was ravaged, and on the mainland the districts around Nicopolis, Anchialus, and Dodona. Transports conveying supplies to the army of Narses at Salona were intercepted and captured.

The garrison of Ancona was hard pressed, for it was blockaded by sea as well as by land. Forty-seven Gothic warships hindered any provisions from reaching it by sea. The general, Valerian, who was stationed at Ravenna and was not strong enough to send relief, wrote to John at Salona an urgent letter on the gravity of the situation. John promptly manned thirty-eight warships with seasoned men,\texthebrew{⁠}\footnote{It is clear that Narses had not yet arrived at Salona, for Procopius says that John took upon himself to disobey the Imperial command that he should not stir till Narses came.

} and at Scardona, higher up the Dalmatian coast, they were joined by twelve more which came across from Ravenna with Valerian. The two generals and their fleet sailed to Sena Gallica, of which the distance by sea to Ancona is about seventeen miles. The two squadrons were practically equal in strength, and the Gothic commanders, Indulf and Gibal,\texthebrew{⁠}\footnote{Indulf had formerly been a retainer of Belisarius. The MSS. of Procopius (III.35.23 and 29; iv.12) vary between \textgreek{\texthebrew{Ἰ}λαούφ, \texthebrew{Ἰ}νδούλφ, \texthebrew{Ἰ}λδούφ} and \textgreek{Γουνδούλφ}. I have followed Haury.

} immediately determined to risk a naval battle, and sailed to Sena.

The action, as in a land battle, was begun by the archers; then some of the vessels closed with each other, and the crews fought with sword and lance. But the Goths were at a great disadvantage. They had not the natural aptitude of the Greeks for handling ships, and they can have had very little training in the operations of maritime warfare. They were unable, in the excitement of the action, to maintain a suitable distance between their ships. Some of these were too far from their neighbours and were easily sunk by the enemy, but most of them were too close together and had no room to manoeuvre. Their opponents, on the other hand, kept perfect order, and with cool readiness took advantage of all the blunders of the Goths, who at last, weary and helpless, gave up the contest and fled. Thirty-six Gothic ships were sent to the bottom and Gibal was captured; Indulf escaped with eleven ships, which he burned as soon as he landed, and reached the camp at Ancona. When

p260
the victorious fleet arrived, they found that enemy had abandoned the camp and taken refuge in Auximum. The crushing victory meant more than the safety of Ancona, it dealt a heavy blow to the power and prestige of the Goths.\footnote{\emph{B. G.} IV.24.42.

}

Soon after this Artabanes, who had arrived in Sicily, recovered the four fortresses which the Goths had captured. The tide seemed to have definitely turned, and the Goths were acutely conscious of the change in their prospects. They felt that if the enemy came in strength they would be unable to hold out. Once more Totila sent ambassadors to Constantinople to propose terms of peace, offering to resign the claim to Sicily and Dalmatia, and to pay the taxes to which the tenantless estates in those provinces were liable. But the Emperor refused to listen to the pleadings of the envoys. He was so bitter against the Ostrogoths that he had determined to expunge their name from the map of the Roman world.

One more success was achieved by Totila, though it was perhaps purchased too dearly. He had sent a fleet to Corsica and Sardinia with forces sufficient to overcome the Roman garrisons. As those islands belonged to the African Prefecture, it devolved upon John, the Master of Soldiers in Africa, to defend them, and he sent an army to Sardinia (autumn, A.D. 551). It was defeated near Cagliari, and sailed back to Carthage, to return in the spring in greater strength. Whatever prestige Totila gained by the occupation of the islands can hardly have counterbalanced the disadvantage of redu­cing the numbers of his fighting forces in Italy, when every man was needed for the approaching struggle with the armies of Narses.

During the spring Croton was hard pressed by the Goths who were blockading it. No one came to its relief until the Emperor, hearing that it would inevitably fall unless speedy help arrived, ordered the troops stationed at Thermopylae to embark immediately and sail thither. The mere appearance of the relief squadron in the harbour sufficed to terrify the besiegers, who hastily broke up their camp and fled. The effect of this bloodless victory was that the commanders of the Gothic garrisons in Tarentum and Acherontia offered to surrender those places on condition that their own safety was secured. Their proposals were referred to the Emperor.

In the spring of A.D. 552 Narses was at length ready to set out for Italy. He had collected large forces in addition to those which had been recruited two years before by Germanus, and which had remained at Salona under the command of John. We are not told what was the entire strength of the army, though we know the number of some of the particular contingents. The Lombard King Audoin sent more than 5500 fighting men;\texthebrew{⁠}\footnote{Two thousand five hundred warriors and a comitatus (\textgreek{θεραπεία}) of more than 3000.

} there were more than 3000 Heruls;\texthebrew{⁠}\footnote{More than 3000 under Philemuth and others; and many others under Aruth, a Romanised Herul.

} there were

p262
400 Gepids; there were Huns,\texthebrew{⁠}\footnote{Very numerous (\textgreek{παμπλήθεις}).

} of course, and there was a band of Persian deserters.\texthebrew{⁠}\footnote{Under Kavad, grandson of king Kavad, and nephew of Chosroes.

} All these foreign auxiliaries can hardly have amounted to less than 11,000. For the regular Imperial regiments which the Emperor placed at the disposal of Narses, for the Thracian and Illyrian troops which Germanus and Narses had specially recruited at their own expense we have Norman figures, but it will not be extravagant to suppose that they were more numerous than the foreign contingents, and to conjecture 25,000 as a probable figure for the strength of the whole army which marched with Narses from Salona along the Dalmatian coast road to the head of the Hadriatic.\footnote{Hartmann, \emph{Gesch. Italiens}, i. p346, conjectures 30,000.

}

The towns and forts which commanded the road from the east into Venetia were in possession of the Franks, and Narses, when he approached the Venetian borders, sent envoys to the commanders asking them to permit a friendly army to pass in peace. The request was refused on the pretext that Lombards who were bitter foes of the Franks accompanied the Imperial army. Then Narses learned that, even if the Franks did not oppose his passage, he would be held up when he reached the Adige, inasmuch as Teïas, one of the most capable of Totila's captains, had arrived at Verona with all the best Gothic troops, to hinder and embarrass his march. Every possible measure had been taken to make the road from the Adige to Ravenna impracticable. By the advice of John, who was acquainted with the country, it was decided that the troops should march along the sea coast from Istria, attended by a few ships and a large fleet of small boats to transport them across the mouths of the rivers. Time was lost, but Ravenna was safely reached. But it is curious that an expedition for which long preparations had been made should have been allowed to find itself in such a predicament. One would have thought that an adequate fleet of transports could have been collected at Salona to convey the whole army direct to Classis.\footnote{Totila was aware of the deficiency of transports, and had hoped that if the troops were conveyed in relays to Italian ports he would be easily able to oppose their landing. \emph{B. G.} IV.26.23. For the coast route,

marked in the \emph{Tabula Peut.},

from Ravenna to Altinum, via Ad Padum, Neronia, and Hadria, see Miller, \emph{Itin.} pp309\texthebrew{‑}311. The Goths had probably beset the road and destroyed the bridges of the Meduaco, the Adige, and the Po, and Narses made his way

(p263)
through the lagoons with the help of boats from Altinum as far as Ad Padum, which was \texthebrew{•} 24 Roman miles north of Ravenna. The date of his arrival in Ravenna was on a Thursday in June (see Agnellus, c62, where the month \emph{July} is a mistake); Körbs (\emph{op. cit.} 84) thinks it was on June 6, and this suits the probable date of the battle of Busta Gallorum, which cannot be placed later than the end of June or first days of July.

}

p263
At Ravenna the army rested for nine days and was reinforced by the troops of Justin and Valerian. Then, leaving Justin in charge of Ravenna, Narses pushed southward along the coast road. He was determined not to spend time or strength in lesser operations, but to come face to face with Totila and decide the issue of the war by a battle involving all the forces of both belligerents. Totila was in the neighbourhood of Rome, and therefore it was on the road to Rome that Narses hastened. When he reached Ariminum he found that the bridge across the river had been destroyed. His engineers bridged it and he might easily have taken the town, for the commander of the garrison, who had sallied out to see what the Romans were doing, was slain by a Herul. But Narses did not tarry; Ariminum could wait. In ordinary circumstances the quickest route for an army marching from Ariminum to Rome was along the coast as far as Fanum and thence by the Via Flaminia. But this way was not open to Narses, for the eastern end of the Via Flaminia was commanded by the enemy who were in possession of Petra Pertusa, a barrier which might be found insuperable. It was therefore necessary for him to strike the road a some point to the west of that fortress. We do not know whether he left the coast near Ariminum or further on, at Pisaurum. In either case he probably reached the Via Flaminia about five miles on the Romeward side of the gorge of Petra Pertusa, at a place which is now known as Acqualagna.

In the meantime Totila, learning that Narses had reached Ravenna, had recalled Teïas and his army from Venetia, and, as eager for battle as Narses, set out for the north. It is not clear where he expected to encounter the Imperialists,\texthebrew{⁠}\footnote{Waiting \textgreek{\texthebrew{ἐ}ν το\texthebrew{ῖ}ς \texthebrew{ἐ}π\texthebrew{ὶ Ῥ}ώμης χωρίοις} till the troops of Teïas arrived, he marched intending to choose his own ground for the battle, \textgreek{\texthebrew{ὡ}ς το\texthebrew{ῖ}ς πολεμίοις \texthebrew{ἐ}ν \texthebrew{ἐ}πιτηδεί\texthebrew{ῳ ὑ}παντιάσων \texthebrew{ἤ}λει}. \emph{B. G.} IV.29.2. Procopius does not record the number of his forces, but it is not probable that, even with the 2000 which arrived late, he had many more than 15,000.

} but when the news reached him that the enemy had left Ravenna and passed Ariminum he struck into the Apennines by the Via p264 Flaminia and encamped near (probably to the north of) Tadinum.\texthebrew{⁠}\footnote{Close to the modern village Gualdo Tadino. Procopius gives the name as \textgreek{Τάγιναι} but the identity is unquestionable. For the topographical questions see

Appendix to this chapter .

} Immediately afterwards the army of Narses reached the neighbourhood and encamped at a place of which the name, Busta Gallorum, preserved a tradition of the wars of the early Roman republic with the Celts of the north. The only other clue the historian gives us as to its position is the statement that it was about fourteen miles from the camp of Totila. We may conjecture that the place is to be sought to the east of the Flaminian Way, in the neighbourhood of Fabriano.

As soon as his army had encamped, Narses sent some trusted officers to Totila, to recommend him to make submission without attempting to oppose much superior forces, and, if he were determined to fight, to invite him to name a day for the battle. Totila would not hear of peace or submission. He said, ``Let us do battle in eight days.'' But Narses was too shrewd to trust the Goth's word. He guessed that Totila would attack him on the next day and made his preparations for battle. So it fell out. The Goths moved during the night, and at dawn the Romans saw their army drawn up within two bowshots of their own line.

Narses placed the Lombards, the Heruls, and the other barbarian auxiliaries in the centre. They were mounted troops but he made them dismount and used them as infantry. On the two wings he posted his regular troops, on the right, under himself and John, on the left under Valerian, Dagisthaeus, and John Phagas; and in front of each wing he stationed 4000 archers. Beyond the extremity of his left, he placed a reserve of 1500 cavalry. Of these one squadron of 500 was to bring help to any part of the line that might be hard pressed;\texthebrew{⁠}\footnote{The natural position for such a reserve would have been in the centre. Narses must have placed it on the left, because he anticipated that it would be most likely to be needed on that side.

} the other body of 1000 was to attempt, when the Gothic infantry were engaged, to ride round and take them in the rear.

Narses had chosen a strong defensive position. It was such that the only way by which the enemy could send a detachment to circumvent him and attack him from behind was a narrow path which ran by the slopes of a small hill close to his left wing.

p265
It was, therefore, important to hold this position, and before daybreak fifty men stationed themselves in the bed of a stream on the slope of the hill fa­cing the Goths. When Totila espied them he sent a squadron of horse to dislodge them, but the Romans held their ground against repeated attacks, performing prodigies of valour. Others were sent, but with the same result, and Totila abandoned the attempt. In the meantime the armies did not join battle. Narses, in his strong position, was determined not to attack first, and Totila had a reason for delaying the action. He was expecting every moment a force of 2000 cavalry under the command of Teïas, who had not arrived in time to march with the main army. Outmatched as he was in numbers by the enemy, this reinforcement was of supreme importance; it might decide the issue of the day. Accordingly he resorted to devices to gain time. Coccas, a horseman of great physical strength, who had deserted from the Imperial to the Gothic side, rode up to the Roman line within speaking distance and challenged the enemy to send out a champion to engage with him in single combat. Anzalas, an Armenian, one of the retainers of Narses, accepted the invitation. Coccas rode hard at him, aiming at his stomach, but Anzalas made his own horse swerve just in time to avoid the lance and at the same moment struck at his opponent's left side. Coccas fell mortally wounded, and cries of triumph rang out from the Roman ranks. After this interlude, Totila himself, caparisoned in shining armour, adorned with gold and purple trimmings, rode out into the space between the armies, on a huge steed, hurling his spear in the air and catching again as he galloped, and performing other feats of horseman ­ship. Finally he sent a message to Narses, proposing negotiations, but Narses knew that he was not in earnest.

By these devices Totila wore away the forenoon, and at length in the early afternoon the belated two thousand arrived. The Goths immediately dissolved their array of battle and retired within the precincts of their camp to dine. Apparently Totila was confident that Narses would not attack,\texthebrew{⁠}\footnote{It may be thought strange that Narses did not attack, but he was determined to avail himself of the strength of his defensive position. Cp. Delbrück, \emph{Kriegskunst}, ii.372.

} and that

p266
the Romans would likewise break their ranks for the purpose of a meal. He thought that he might possibly take them unawares. But their cautious commander did not allow them to move from their places or take off their armour or lie down to rest. They took food as they stood.

In the morning the array of the Goths had been much the same as that of the Imperial army, but when they returned to fight in the afternoon, Totila adopted an entirely different plan. He placed all his cavalry in front and all his infantry behind. His idea seems to have been that his best chance was to attempt to break the enemy ranks by a concentrated charge of all his horse, and then bring up his infantry (probably few in number) to take advantage of the confusion which his cavalry had wrought. And he issued the extraordinary command that all the troops alike should discard the use of all weapons except their spears.

To meet the tactic of the Goths Narses made a slight change in his dispositions. The two large bodies of archers on the wings,\texthebrew{⁠}\footnote{They must have been on elevated ground above the plain, or they would have been swept away by the Gothic cavalry. Delbrück, \emph{Kriegskunst}, ii.375.

} which had faced the enemy full front, were now turned half round so as to form crescents fa­cing each other; and when the Gothic cavalry charged they were assailed from both sides

p267
by showers of arrows and suffered considerable losses before they came to grips with the main line. The battle was fierce, but apparently short, and towards evening the Goths gave way and were gradually pressed back on the infantry who had hitherto taken no part in the fighting, and now, instead of opening a way for the cavalry to pass through their ranks and themselves fa­cing the enemy, turned and retreated with them. The retreat soon became a flight. About 6000 were slain; many were taken alive, to be put to death afterwards; all the rest fled as they could.

The description of the battle, which we owe to the historian Procopius, and which he doubtless derived from an eyewitness, is so deficient in details that it is difficult to form any definite opinion as to the merits of the combatants. Above all, we do not know the numbers of either army. We are not told how Totila and his ablest general Teïas behaved during the action, nor whether the wings or the centre of the Imperialists were the more heavily engaged. Praise is given to the bravery of the barbarian troops of Narses and of ``some of the Romans,'' but the military critics of the day seem to have ascribed the swift discomfiture of the Goths largely to the strange order of their king that the spear only was to be used. We can, however, divine that Totila's general ­ship was deficient and that, even if his forces were inferior in number, he might have made better use of them.

But in spite of the slightness of our information as to the course of the battle, it is clear that Narses displayed exceptional military talent and deserves full credit for his victory. His plan was original, differing entirely from the tactic employed by Belisarius in the Persian campaigns. He opposed unmounted troops to the mounted troops of the enemy, and used his bowmen to weaken and disconcert the charge of the cavalry. Thus aided, the barbarian auxiliaries did what the Roman infantry had failed to do on the field of Hadrianople, and resisted the shock of the Gothic horsemen. The battle has been described as ``the first experiment in the combination of pike and bow which modern history shows,'' and reminds us of the battle of Crecy º which was won by similar tactic .\footnote{Oman, \emph{Art of War}, p35. Tacticians of the time had contemplated battles fought by Roman cavalry against enemy cavalry; cp. the passage of Urbicius quoted in Pseudo-Maurice, \emph{Strat.} xii.24.

}

Totila himself had fled in the dark and there were various

p268
stories as to what befell him. According to one tale, accompanied by four or five of his followers he was pursued by Asbad the Gepid leader, and some others who were unaware of his identity. Overtaking him, Asbad was about to strike when a Gothic youth cried, ``Dog, will you smite your master?'' The Gepid drove his spear with all his might into Totila's body, but was himself wounded by one of the king's companions. The Goths dragged their wounded lord for about seven miles, not halting till they reached Caprae, a village not far from Tadinum. Here he died and was hastily buried. His fate and place of sepulture was revealed to the Romans by a Gothic woman; the body was exhumed and identified; the blood-stained garments and the cap adorned with gems which he had worn were taken to Narses, who sent them to Constantinople, where they were laid at the feet of the Emperor as a visible proof that the enemy who had so long defied his power was no more.\footnote{Here John Mal. xviii p486 supplements Procopius. The passage was transcribed by Theophanes, \emph{sub} A.M. 6044, who adds a sentence which evidently belonged to the original text of Malalas but is omitted in ours (\textgreek{κα\texthebrew{ὶ ἐ}ρρίφησαν ε\texthebrew{ἰ}ς το\texthebrew{ὺ}ς πόδας το\texthebrew{ῦ} βασιλέως \texthebrew{ἐ}π\texthebrew{ὶ} σεκρήτου}).

}

A leader who has fought a long fight in a not ignoble cause and failed in the end will always arouse some sympathy and pity, with whatever satisfaction we may view his failure. The sudden reversal of Totila's fortune after an almost unbroken career of success had just the elements of tragedy which appealed even to the imagination of his enemies. He had revived the cause of his nation when it seemed utterly lost and restored their hope, and in a struggle of nine years, in which he displayed untiring energy, unwavering confidence, and some political capacity, had reconquered the whole of Italy except three or four towns. But this long run of success does not argue that he possessed transcendent talents. He owed it to the fact that the Emperor starved his military forces in Italy, refused to send the necessary supplies of money and men, and at first did not even appoint a supreme commander. As soon as Justinian decided, after the return of Belisarius, to make a serious effort to end the war and adopted proper measures for the purpose, the situation began immediately to change, and all that Totila had achieved in nine years was undone in two. But though the weakness and mistakes of his enemies were chiefly responsible for Totila's fame, though he did not possess military genius of a

p269
high order, and was capable of such a political blunder as the abandonment of Rome when he had captured it, he will always be remembered as one of the great figures in the German heroic age.

Some modern writers have idealised him as a romantic hero, distinguished among all his barbarian fellows by chivalrous sentiments and noble behaviour towards his foes, gentle and humane in his instincts. It is difficult to find much in the record of his accounts to justify such a conception of the man. He was clear-sighted enough to realise that it was good policy to conciliate the Italians and to attract to his standards deserters from the Imperial army, and for these purposes he often showed a moderation which in time of war was unusual. Perhaps his considerate treatment of the inhabitants of Naples, which the historian Procopius ungrudgingly admired, has won for him a reputation which his conduct on other occasions can hardly be said to bear out. But his friendliness to the Neapolitans was plainly dictated by policy. It was to reward them for the obstinate resistance they had offered to Belisarius eight years before, and Totila intended it to be contrasted with the punishment which he hoped to inflict upon the Sicilians who had received Belisarius with open arms. In the practice of deliberate cruelties can it be said that there is much to choose between this Ostrogoth and other leaders of his race and age? What instinct of clemency can we attribute to the man who mutilated Demetrius at Naples, who cut off the hands of the bishop from Porto, who put Isaac the Armenian to death, who did not spare his unhappy captive Gilacius, who shamefully mutilated Chalazar? What are we to say of the assassination of Cyprian at Perusia? Can we call him humane who suffered the bishop and inhabitants of Tibur to be done to death in such atrocious fashion that the historian declines to describe the treatment? Did he treat the inhabitants of Rome as leniently as Alaric or Gaiseric? Narses had no illusions about his character, and it was well for him that, when Totila named a day for the great battle which was to be fought between them, he did not imagine him to be a pure chevalier, but knew him for an ordinary perfidious barbarian and took corresponding precautions.\footnote{Totila's reputation for cruelty is illustrated by the story told by Pope Gregory I (\emph{Dial.} 3, c11) that he condemned Cerbonius, bishop of Populonium, to be thrown to a bear because he had given shelter to some

(p270)
soldiers of the Imperial army. Totila expected to enjoy the spectacle of the execution, but the bear lay down and licked the bishop's feet. Though Gregory appeals to survivors at Rome who witnessed the amazing incident, we can hardly credit this version of Androcles and the lion. But it shows Totila's reputation.

}

The first act of Narses after his great success, for which he piously ascribed all the credit to the Deity, was to dismiss his savage allies, the Lombards, who, as soon as the victory was won, were devoting themselves to the congenial occupations of arson and rape. He rewarded them with large sums of gold, and committed to Valerian the task of conducting them to the Italian frontier. When Valerian had parted from these undesirable friends, he encamped outside Verona and parleyed with the Gothic garrison. The Goths were willing to capitulate, but the Franks who were firmly stationed in the Venetian province intervened and the negotiations were broken off. Valerian withdrew to the Po, and Narses ordered him to remain there to watch the movements of the Goths, who had not yet given up their cause as lost. The remnant of Totila's army had fled with Teïas northward to

Ticinum . There Teïas was elected king,\texthebrew{⁠}\footnote{His father's name was Fritigern. He is Theia (also Teia, Thela, Thila) on his coins, which have the head of Anastasius. Wroth, p95 \emph{sqq.}

} and he hoped with the help of the Franks to restore the fortunes of his people. He had at his disposal the treasures which Totila had prudently left in Ticinum.

In the meantime Narses himself had advanced on Rome. On his way he occupied Narnia and Spoleto, and sent a detachment to take Perusia. The Gothic garrison in Rome was much too small to attempt to defend the great circuit of the city, and Totila had constructed a little fortress round the Mausoleum of Hadrian by building a new wall attached to the external wall. When the army of Narses arrived, the Goths made some attempt to hold the fortifications wherever they were attacked, but the Imperialists soon succeeded in scaling the wall with ladders and opening the gates. The garrison then retreated into the inner fortress; some escaped to Portus. But seeing that further defence was useless they surrendered on condition that their lives were spared. This was the fifth time that Rome had been assaulted and captured during the war.

p271
Narses sent the keys to the emperors. Soon afterwards Portus surrendered.

The Goths now showed themselves, without any reserve, in their true colours. (1) In Campania they put to death the senators who had been sent there by Totila and now proposed to return to Rome. (2) Before Totila went forth to meet Narses he had selected three hundred boys from Roman families of repute and sent them to the north of Italy as hostages. Teïas seized them and slew them all. (3) It will be remembered that Ragnaris, the commander in Tarentum, had agreed to surrender on conditions which Pacurius, the commander of Hydruntum, had gone in person to submit to the Emperor; in the meantime he had given hostages. Learning that Teïas was resolved to renew the struggle and counted on the help of the Franks, Ragnaris changed his plans. When Pacurius returned from Constantinople,\texthebrew{⁠}\footnote{The journey to Constantinople and back need not have taken much more than a month, as the business need have caused no delay.

} he asked him to send a few Roman soldiers to conduct him safely to Hydruntum and thence by sea to Constantinople. Pacurius sent fifty men. Ragnaris imprisoned them and then informed Pacurius that they would not be released until the Gothic hostages had been restored. The Roman commander lost no time in marching to Tarentum with all his forces. At his approach Ragnaris put the fifty men to death and marched out to meet him. The Goths were defeated and Ragnaris fled to Acherontia. These circumstances of the recovery of Tarentum deserved to be recorded as an illustration of the character of the Ostrogoths.

Narses meanwhile had not been idle. He sent a force to reduce Centumcellae, and another into Campania to lay siege to Cumae. The importance of this fortress lay in the fact that Totila had deposited in it all the Gothic treasure that was not stored at Ticinum, and left it in the custody of his brother Aligern.\texthebrew{⁠}\footnote{Agathias, I.8 , records the name, but says he was the youngest brother of Teïas. Procopius says he was the brother of Totila, \emph{B. G.} IV.34.19. Cumae is about 200 kilometres from Rome, so that an army leaving Rome just after the middle of July would have arrived before the end of the month.

} When the news that this store was in immediate danger reached Teïas, who had been waiting in the vain hope that the Franks would provide an army to help him, he determined to make an attempt to rescue Cumae. It was a long

p272
march from Ticinum to Campania, and even a small army, moving more rapidly than usual, could not accomplish it in much less than a month. The shortest route was through Etruria, and Narses sent a force under John to watch the Etruscan roads. But Teïas did not choose the shortest route. His object was to avoid the enemy, and he went by devious and roundabout ways, finally following the coast road of the Hadriatic. He must have crossed the peninsula by Beneventum, where he could proceed either by Capua or by Salerno to the neighbourhood of Naples.\footnote{The possible route of Teïas is discussed by Körbs (\emph{op. cit.} p82), who calculates that he had about 814 kilometres to cover, that he probably did at least 30 kilometres a day, that he started soon after the middle of July and reached Campania after the middle of August.

}

Narses, when he found that the enemy had eluded both John, who was guarding the western roads, and Valerian, who had captured Petra Pertusa\texthebrew{⁠}\footnote{\emph{B. G.} IV.34.24; but the Petra Pertusa which is mentioned (\emph{ib.} 16) along with
Nepi and Porto seems to have been another place, also on the Via Flaminia, but quite close to Rome, north of Prima Porta. See Tomassetti, \emph{La Campagna Romana}, III pp138, 260\texthebrew{‑}261.

} and was thus master of the Via Flaminia, recalled both these generals and proceeded with all his forces to Campania. When Teïas at last reached the southern foothills of Vesuvius, near Nuceria, he found a Roman army drawn up on the bank of the Draco.

This river, now the Sarno, runs into the bay of Naples, north of the Sorrento peninsula. The remnant of the Gothic fleet was assembled in the bay of Naples. As Teïas might expect that the land approaches to Cumae, north of Naples, would be guarded, his plan probably was to embark his troops near Sorrento and reach Cumae by sea. There was no fleet at hand to oppose him, and the plan was only foiled by the vigilance and good intelligence service of the Roman general, who was just in time to prevent him from reaching the sea.

The armies remained for weeks\texthebrew{⁠}\footnote{Two months, acc. to Procopius, IV.35.11. This must be a considerable exaggeration if Agnellus is right in his dating of the battle to Oct. 1 (in Kal. Octobris, c79). Teïas could not at the earliest have reached Campania before the middle of August; it seems more likely that he arrived at the end of the month. If Procopius is accurate the date of Agnellus is wrong.

} fa­cing each other on either bank of the narrow stream, which neither infantry nor cavalry could ford on account of the steepness of the banks, the archers carrying on a desultory battle. There was indeed a bridge on

p273
which the Goths who held it erected towers and assailed their enemies with bolts from ballistae . Teïas succeeded in getting into touch with his fleet and it was able to supply him with provisions. The situation was changed when Imperial warships which Narses summoned began to come in great numbers from Sicily and other places. The Gothic naval commander, anticipating their arrival, surrendered his fleet. The food-supply of the army was thus cut off, and at the same time it began to suffer from the play of the engines which Narses installed in wooden towers along his bank of the stream. Teïas broke up his camp and retreated to the shelter of the mountain which over­looks the valley. This mountain, belonging to the St. Angelo range, was known as Mons Lactarius and still retains the name as Monte Lettere. On the slope of this hill the Goths were safe from attack, which the nature of the ground would have rendered too dangerous an enterprise, but they found themselves worse off for food, and they soon repented their change of ground. At length they resolved to make a surprise attack upon their foes. It was their only chance.

They appeared so unexpectedly in the valley that the Romans had no time to form themselves in the regular array prescribed by military handbooks.\texthebrew{⁠}\footnote{According to local tradition the scene of the battle was at Pozzo dei Goti, a kilometre west of Angri at the foot of Monte Lettere.

} The Goths had left their horses behind and advanced as a solid mass of infantry. The Romans received them in the same formation.\texthebrew{⁠}\footnote{Procopius does not explain why the Goths should have advanced on foot, or why Romans should have dismounted. Delbrück (\emph{op. cit.} 382) conjectures that the Romans had constructed fortifications (earthworks and ditch?) to blockade the Goths, and thus the Goths were obliged to attack on foot, and the Romans to defend the line on foot.

} In the battle there was no room for tactic , it was a sheer trial of personal strength, bravery, and skill. The Gothic king, a few warriors by his side, led the assault, and, the Romans recognising him and thinking that if he fell his followers who were formed in a very deep phalanx would not continue the contest, he became the mark for their most dexterous lancers and javelin-throwers . It was a Homeric combat, and the historian has described it vividly. Teïas stood covered by his shield, which received the spears that were hurled or thrust at him, and then suddenly attacking laid many of his assailants low. When he saw that the shield was full of spears he gave it to one of his squires, who handed him another. He is

p274
said to have fought thus for a third part of the day, then his strength failed. There were twelve spears sticking in his shield, and he found he could not move it as easily as he would. Without retreating a foot or moving to right or left, smiting his foes with his right hand, he called the name of a squire. A new shield was brought, but in the instant in which he was exchanging it for the old his chest was exposed, and a lucky javelin wounded him mortally.\footnote{The whole account of the first part of the battle seems fanci­ful and improbable. The deep phalanx of the Goths plays no part in the action, and Teïas alone occupies the stage. His heroism is assuredly a fact, but the narrative of Procopius cannot be accepted as a true account of the battle.

}

The head of the fallen hero was at once severed from his body and raised aloft on a pole that all his host might know that he had fallen. But the expectation of the Romans that their enemies would abandon the struggle was not fulfilled. The Goths did not flee like fawns, nor lay down their arms. They were animated by a spirit of desperation, and in a very different temper from that which they had displayed in the last battle of Totila. They fought on till nightfall, and on the next day the fray was resumed, and again lasted till the evening. Then, seeing that they could not win and recognising that God was against them, they sent some of their leaders to Narses to announce that they would yield, not, however, to live in subjection to the Emperor, but to retire somewhere outside the Roman frontiers where they could live independently. They asked to be allowed to retire in peace, and to take with them any money or belongings that they had individually deposited in Italian fortresses.

On the advice of John, who made a strong plea for moderation, these conditions were accepted, on the undertaking of the Goths that they would not again make war on the Empire.\footnote{The \emph{History} of Procopius ends with the victory of Mons Lactarius, and the story is taken up by Agathias.

}

The shields of Teïas had not availed to avert the doom of his people. He was their last king. The kingdom of the Ostrogoths went down on the hard-fought field under Mount Lactarius. But there was still fighting to be done. The great defeat did not lead to the immediate surrender of the strongholds which

p275
were still held by Gothic garrisons. There was Cumae, there was Centumcellae, there were a number of towns in Tuscany, and there was North Italy beyond the Po. Narses had still much strenuous military work before him. He might have hoped to complete the reduction of the land by the following summer, but his plans were disconcerted by the appearance of a new and more barbarous enemy upon the scene.

Teïas had invoked the assistance of the Franks. The answer of the young king Theodebald\texthebrew{⁠}\footnote{Theodebald married Vuldetrada, daughter of Wacho, king of the Lombards. Her mother was Austriguna, a Gepid princess.

Gregory of Tours, \emph{Hist. Fr.} iv.9 .

} to the pleadings of his envoys was unfavourable. The Franks had no mind to embark on a war for the sake of the Ostrogoths; they coveted Italy for themselves,\texthebrew{⁠}\footnote{Procopius, \emph{B. G.} IV.34.18.

} but at the moment they judged neutrality to be the best policy. But neutrality was only official. Two chieftains of their subjects the Alamanni, Leutharis and Buccelin, who were brothers, formed the plan of invading Italy. Ostensibly Theodebald did not approve of this act of aggression, but he took no steps to prevent it.\texthebrew{⁠}\footnote{Agathias says

(I.6)

that the king was opposed to the invasion. On the other hand, the invaders sent him a portion of the treasures which they collected in Italy ( Gregory of Tours, \emph{H. F.} III.32 ; Paulus Diac. \emph{H. Lang.} II.2). Agathias has much to tell us about the Franks. I conjecture that he gathered his information from the ambassadors of King Sigebert, who visited Constantinople in A.D. 566. It was their cue to represent the invasion of Italy as not countenanced by Theodebald. Rastia, in which the Alamanni had been settled by Theoderic, had been abandoned by the Goths (Totila?) and Theodebert had brought it under his rule (Agathias \emph{ib.}).

} The two adventurers raised a host of 75,000, in which Franks as well as Alamanni served, and descended into Italy in the spring of A.D. 553, confident that they could overwhelm Narses, for whose military talents, eunuch and chamberlain as he was, they professed supreme contempt.\footnote{Agathias is the main source for the invasion, but we have also brief accounts in western sources: Marius of Aventicum, \emph{Chron., sub} 555, 556; Gregory of Tours, \emph{Hist. Fr.}

iii.32 ,

iv.9 ; Paulus Diac. \emph{Hist. Lang.} ii.2.

}

Narses spent the winter months in besieging Cumae, but Aligern and the little fortress held out obstinately. When all his assaults and devices failed, he left a small investing force, and proceeded to Central Italy, where he found the Gothic garrisons ready to make terms. Centumcellae surrendered, and the Tuscan towns, Florence, Volaterrae, Pisa, and Luna did likewise. Lucca alone bargained for a delay; if no help came to them before thirty days expired, surrender was promised, and

p276
hostages were given. The help which the Goths of Lucca looked for was the arrival of the Franks, who had already crossed the Alps. It was the imminence of their invasion that had probably decided Narses to march northward, and he had sent the greater part of his army under John and Valerian to guard the passages of the Po.

The thirty days passed and the garrison of Lucca refused to abide by their agreement. Some of his officers, in their indignation at this breach of faith, suggested that the hostages should be put to death. Narses was not a Goth; he would not commit the injustice of executing innocent men. But he led them forth, with their hands bound across their bodies and their heads bowed, within sight of the walls, and proclaimed that they would be slain if the town were not surrendered. Thin pieces of wood, wrapped in pieces of cloth, had been fastened on the backs of the hostages from the neck to the waist, and, when the garrison gave no sign of yielding, guardsmen stepped forward and drawing their swords brought them down on the well-protected necks. The victims, who had been let into the secret, fell forward, as if they had been decapitated, and their bodies feigned the spams and contortions of death. The spectators on the wall set up howls and wails, for the hostages belonged to noblest families; mothers and affianced brides rushed along the battlements rending their garments. All cried shame on the bloody cruelty of Narses.

Narses sent a herald to address them.\texthebrew{⁠}\footnote{\textgreek{Τα\texthebrew{ῦ}τα δ\texthebrew{ὲ} α\texthebrew{ὐ}τ\texthebrew{ῶ}ν \texthebrew{ἐ}πιβοώντων . . . \texthebrew{ἔ}φη \texthebrew{ὁ} Ναρσ\texthebrew{ῆ}ς} (Agathias, I.43) ; a herald is not mentioned but is clearly to be presumed.

} ``You have yourselves to blame,'' he said, ``for the shameless violation of your oaths. But if you will come to your senses even now it will be well for you; these men will come to life again and you will suffer no harm.'' The Goths had no doubt that he was deceiving them, but they readily swore that if he showed them the hostages alive they would at once capitulate. Then at the general's command all the dead stood up together and shied themselves safe and sound to their friends, who were divided between incredulity and joy. But incredulity prevailed, and then Narses, with a magnanimity which was well calculated, set his prisoners free, and allowed them, without imposing any conditions, to return to their people in the town. They went back loud in his

p277
praises, but Lucca did not surrender. Oaths and solemn engagements were of no account in the eyes of the Goths, who were elated with new hopes by the success ­ful advance of the Franks.

For Buccelin and his Alamanni had won possession of Parma, and had cut to pieces a force of Heruls who, under a brave but rash leader, attempted to recover it. All the Goths in the Ligurian and Aemilian provinces had rallied to the invaders, and it is probable that these were in command at Ticinum itself. John and Valerian, upon whom Narses relied to keep them back from Etruria while he was engaged in redu­cing Lucca, had withdrawn to Faventia. Lucca, however, he was determined to take, and he prosecuted the siege with vigour. It would have surrendered soon if Frank officers had not succeeded in entering the town and stiffening the defence. But at length the will of the majority prevailed, and the Luccans opened their gates and received the army of Narses, who had agreed not to punish them for their ill-faith.

The siege had lasted three months, and it was now the end of autumn. Narses went to Ravenna to arrange the dispositions of the troops for the winter, and presently Aligern, the Gothic commander of Cumae, which had held out all this time, arrived at Classis and gave him the keys of the town, Aligern had come to the conclusion that the Franks had no intention of restoring to Ostrogothic power, and that whether they succeeded or not in conquering Italy, in neither event had he the least chance of inheriting the throne of Teïas. He therefore decided to resist no longer but to become a subject of the Empire.

Narses spent the winter in Rome, and in the spring ( A.D. 554) his army, which had been dispersed among the forts and towns in the Ravennate region for the winter, was collected and reunited at Rome. We do not know his reasons for this retreat, which meant the abandonment of Etruria and the Hadriatic provinces to the enemy. He could rely with some confidence on his garrisons in the great fortresses, but the open country and unwalled towns were at the mercy of the invader.

The host of Buccelin and Leutharis moved southward, without haste, plundering and destroying. When they approached Rome they divided into two separate armies, of which the larger under Buccelin, avoiding Rome itself, marched through Campania, Lucania, and Bruttii to the Straits of

p278
Messina,\texthebrew{⁠}\footnote{Gregory

(iii.32)

says that Buccelin defeated Narses in a battle, and then occupied Sicily. These statements may be due to exaggerated rumours derived from Buccelin's report of his successes. It is probable that when he reached Rhegium, he despatched a message to Theodebald.

} while Leutharis led the other through Apulia and Calabria as far as Hydruntum. The provinces were systematically plundered, and an enormous booty was collected. In this work of pillage and devastation there was a marked difference between the conduct of the Franks and their Alamannic comrades. The Franks, who were orthodox Christians, showed respect for churches, but the heathen Alamanni were restrained by no scruples from carrying off the ecclesiastical plate and pulling down the roofs of the sacred buildings.\footnote{Agathias

(I.7)

describes the nature worship of the Alamanni, their cults of trees, rivers, and hills, but thinks that it will soon disappear through the influence of the Franks.

}

When he had reached the limits of Calabria, Leutharis laden with spoils decided to return home to enjoy them. He had no political ambitions, and his one thought was to get safely away with his wealth and run no further risks. He marched along the coast as far as Fanum, but there his troops suffered considerable losses through an attack by the Roman garrison of Pisaurum, and the greater part of the booty was lost. Leaving the coast he struck into the Apennines and reached the Po safe but dispirited.\texthebrew{⁠}\footnote{\textgreek{Ο\texthebrew{ὕ}τω τε \texthebrew{ἰ}θ\texthebrew{ὺ} Α\texthebrew{ἰ}μιλείας κα\texthebrew{ὶ Ἀ}λπισκοτίας \texthebrew{ἐ}λθόντες}; here Agathias

(II.3)

betrays his ignorance of Italian geography. He supposed that the district of Alpes Cottiae was adjacent to Venetia.

} At the Venetian town of Ceneta,\texthebrew{⁠}\footnote{Paulus, \emph{loc. cit.}, says near lake Garda, between Verona and Trent. Ceneta, now Ceneda, lies between Oderzo and Feltre.

} where he took up his quarters to rest, a virulent plague broke out in the army and Leutharis himself was one of its victims.

His brother Buccelin was more enterprising and ambitious. He had professed to the Goths that his object was to restore their kingdom, and many of them doubtless attached themselves to his army in his southern march. He fell under the influence of their flatteries; they told him that they would proclaim him king if he drove Narses out of Italy; and he was finally persuaded to risk everything in a battle with the army which he had hitherto aimed at avoiding.

He returned to Campania and encamped on the banks of the Vulturnus\texthebrew{⁠}\footnote{Agathias

(II.4)

calls it the Casulinus. Paulus says the battle was fought at a place called Tannetum (\emph{al.} Cannetum).

} close to Casilinum and Capua, which are only a few

p279
miles apart. Casilinum is the modern Capua, and the ancient Capua is the modern village of S. Maria di Capua Vetere. On one side the river formed the wall of his camp, on the other side he fortified it securely.\texthebrew{⁠}\footnote{Agathias says that his army amounted to about 30,000, and that the numbers were considerably reduced by dysentery, attributed to the immoderate use of ripe grapes. The figure of 30,000 is probably too high.

} He had some hopes that he would soon be reinforced, for his brother had promised that when he had reached Venetia he would send back his troops. As soon as Narses learned that Buccelin had occupied this position at Capua he marched from Rome with his army, numbering about 18,000, and encamped not far from the enemy. The battle which ensued was probably fought across the Appian Way which passed through Capua and crossed the river at Casilinum.

The course of the battle was affected by an accident. One of the Herul captains killed his servant for some delinquency, and when Narses called him to account asserted that masters had the power of life and death over their slaves and that he would do the same thing again. He was put to death by the command of Narses, to the great indignation of the Heruls, who withdrew from the camp and said they would not fight. Narses drew up his line of battle without them. He placed his cavalry on the two wings and all the infantry in the centre. There was a wood on the left, and Valerian and Artabanes, who commanded on that side, were directed to keep a part of their forces concealed in the wood till the enemy attacked. Narses himself commanded on the right. The leader of the Heruls, Sindual, who was burning to fight, implored Narses to wait until he could persuade his followers to reserve a place for them, where they could fall in, if they arrived late. Accordingly he left an open space in the middle of the infantry.

Meanwhile two Heruls had deserted to the enemy, and persuaded Buccelin that his chance was to attack at once, as the Romans were in consternation at the defection of the Herul troops. Buccelin had drawn up his army, which consisted entirely of infantry, in the shape of a deep column, which should penetrate like a wedge through the hostile lines.\texthebrew{⁠}\footnote{Agathias

(II.8)

describes the formation as a triangular (\textgreek{δελτωτ\texthebrew{ῷ}}) wedge, with the point towards the enemy, and compares it to the ``head of a boar.'' It was simply the cuneus described by Vegetius (III.19): cuneus dicitur multitudo peditum,

(p280)
quae iuncta cum acie primo angustior deinde latior procedit et aduersariorum ordines rumpit, quia a pluribus in unum locum tela mittuntur. Quam rem milites nominant caput porcinum. There must have been the same number of men in each rank of the column, but in advancing the men of the first ranks drew closer together, and the columns became a trapezium instead of a rectangle, with the smallest side towards the foe.

} In this array

p280
the Franks arrived, armed with missile lances, swords, and axes,\texthebrew{⁠}\footnote{See

Agathias II.5

for Frank armour (cp. Sidonius, \emph{Epp.} IV.20 ). The axe was called francisca, the lance for hurling angon. The Franks generally fought naked to the waist, with leather trousers, without breastplate or greaves, and bareheaded, though a few had helmets.

} confident that they would sweep all before them at the first rush. They penetrated into the central space which was to have been occupied by the Heruls, dislodging the outer ranks of the Roman infantry on either side. Narses quietly issued orders to his wings to face about, and the enemy were caught between the cross fire of the cavalry, who were all armed with bows. The Franks were now fa­cing both ways. The archers on the right wing aimed at the backs of those who were fighting with the infantry on the left, the archers on the left wing at the backs of those who were engaged with the right. The barbarians did not understand what was happening. They saw the foemen just in front of them with whom they were fighting hand to hand, but they could not see the enemies who from far behind were raining arrows upon their backs. Their ranks were gradually mown down, and then Sindual and his Heruls appeared upon the scene. The defeat of the Franks was already certain; it was now to be annihilation. Buccelin was slain and only a handful escaped alive from the stricken field. The Roman losses were small.\texthebrew{⁠}\footnote{Agathias says only five Franks escaped, and that only eighty Romans were killed; Marius that Buccelin cum omni exercitu suo interiit.

} It will be noticed that Narses won this, his third victory, by a tactical plan similar to that which he had employed in the battle with Totila.

The Italians had been terror-stricken by the ruthless deeds of the northern barbarians, and they were wild with joy at the news of their utter destruction. Narses and thoughtful people had little hope that the brilliant victory of Capua head dispelled the danger. They reflected that the foes whose corpses were strewn on the banks or floated in the waters of the Vulturnus were such a small fraction of the Frank people and their dependents, that their fate would provoke rather than intimidate. They expected that a greater host would soon come down to

p281
avenge the fallen and restore German prestige.\texthebrew{⁠}\footnote{See the speech of Narses to his army in

Agathias, II.12 .

} These fears were not realised, as they might well have been if Theodebert had been still alive; his feeble son Theodebald, who suffered from a congenital disease, died in the following year. Narses was able to complete in peace the settlement of Italy.

The winter months which followed the battle of Capua were spent in besieging Capsa, a strong place in the Apennines, where seven thousand Goths had established themselves under the leader ­ship Ragnaris, the man who had behaved so treacherously at Tarentum.\texthebrew{⁠}\footnote{There can hardly be any doubt as to the identity.

} Campsa has been identified with Conza, \texthebrew{•} about fifty miles east of Naples. Its position defied assault and Narses sat down to blockade it, but a large stock of provisions had been laid in. At the beginning of spring ( A.D. 555), Ragnaris proposed to Narses that they should meet and discuss terms. Narses refused to agree to his proposals, and he retired in great wrath. When he was near the wall of the fort he turned round, drew his bow, and aimed an arrow at the general who was returning to his lines. It missed its mark, but one of the guardsmen who were with Narses had a surer aim, and transfixed the treacherous Goth. He fell dead, and the garrison surrendered immediately and were sent to Constantinople.

All Italy south of the Po was now restored to the Imperial authority. Of the subjugation of the Transpadane provinces, where Goths and Franks were still in possession, we have no record. It was a slow business, and Verona and Brixia were not recovered till A.D. 562. In November of that year Narses sent the keys of their gates to Justinian.\footnote{John Mal. xviii p492 =Theophanes, A.M. 6055. With the recovery of Verona and the end of the warfare in Venetia we may perhaps connect the defeat of the Frank Aming and the Goth Widin, of whom we hear in Paul Diac. \emph{(l.c.)} and Menander (\emph{De leg. Rom., fr.} 2, p171). Aming opposed a Roman army which was about to cross the Adige. Narses sent envoys warning him to depart, as a truce had been concluded

(p282)
between the Empire and the Franks. Aming replied that he would not retreat so long as his hand could wield a javelin. He had come to the assistance of a Goth named Widin (possibly the commander in Verona). A battle ensued; Aming was slain by the sword of Narses, Widin made prisoner and sent to Constantinople.



a
Thayer's Note: In the print version, this Appendix is given at the end of the chapter. In this Web transcription, I've moved it up to the end of the section it relates to; that also accounts for the discontinuity of the note numbering. Notes 120\texthebrew{‑}142 will be found in the last section of Chapter XIX,

on the next webpage .}

\xsubsection{On the Battle of Busta Gallorum}

The route taken by Narses after he crossed the river at Ariminum is not precisely indicated by Procopius. His words are (\emph{B. G.} IV .28.13): \textgreek{\texthebrew{ὁ}δο\texthebrew{ῦ} δ\texthebrew{ὲ} τ\texthebrew{ῆ}ς Φλαμηνίας \texthebrew{ἐ}νθένδε \texthebrew{ἀ}φέμενος \texthebrew{ἐ}ν \texthebrew{ἀ}ριστερ\texthebrew{ᾷ ἤ}λει} . The sentence seems to have been misunderstood by Hodgkin , who contended that Narses marched along the coast to a point south of Fanum and north of Sena Gallica, and then turned inland by a road ascending the valley of the Sena (Cesano), and reaching the Via Flaminia at Ad Calem (Cagli). Such a road is noticed in the \emph{Itinerarium Antonini.}\texthebrew{⁠}\footnote{As to this road Nissen observes

(\emph{Ital. Landeskunde}, II.1.392n) : ``Das It. Ant. 315 erwähnt eine Strasse von Helvillum über ad Calem und ad Pirum nach Sena und Ancona jedoch is dieselbe nicht nachgewiesen und die Entfernungen ganz entstellt.'' But see on the other hand Cuntz, in \emph{Jahresh. Österr. arch. Inst.} VII.61. Muratori

(\emph{Annali d'Italia}, III p433)

says: ``voltò Narsete a man destra per valicar l'Apennino,'' but does not specify at what point he turned.

} But if he had taken this route Narses would have had the Via Flaminia on his right, whereas Procopius plainly says the opposite: ``He marched having left the Flaminian Way on his left from this point.''\texthebrew{⁠}\footnote{The phrase is illustrated \emph{infra}, 34.23 \textgreek{Τείας \texthebrew{ὁ}δο\texthebrew{ὺ}ς μ\texthebrew{ὲ}ν \texthebrew{ἐ}ν δεξι\texthebrew{ᾷ} τ\texthebrew{ὰ}ς \texthebrew{ἐ}πιτομωτάτας \texthebrew{ἐ}π\texthebrew{ὶ} τ\texthebrew{ὸ} πλε\texthebrew{ῖ}στον \texthebrew{ἀ}φείς.}} In order to have the Flaminian Way on his left he must have turned inland between Ariminum and Fanum.

The word \textgreek{\texthebrew{ἐ}νθένδε} ``from this point,'' shows that Procopius supposed that Narses diverged from the coast road close to Ariminum. If this statement is correct, it might imply (1) that Narses passed San Marino, followed a road now well defined to Pieve di S. Stefano, crossed the watershed of the Apennines, reached the town now called Città di Castello:\texthebrew{⁠}\footnote{Situated on the upper waters of the Tiber, and identified with Tifernum Tiberinum.

} from which point he could proceed

p289
either (a) to Urbania, and thence to Acqualagna, or (b) to Iguvium (Gubbio), and thence to Aesis (Scheggia). But (2) it is not improbable that there was a direct road from Ariminum to Urbinum (by Coriano, Montefiore, Tavoleto, Schieti), and thence to Acqualagna by Fermignano. º The engineer, P . Montecchini, found traces of it north of Fermignano (\emph{La Strada Flaminia dall' Apennino all' Adriatico}, pp38 \emph{sqq.}, 1879, published at Pesaro).

The other alternative which was open to Narses was to proceed along the coast road as far as Pisaurum, and there to take the road to Urbinum, and it may be said that we cannot pass \textgreek{\texthebrew{ἐ}νθένδε} so strictly as to exclude this possibility from our consideration; if the informant of Procopius omitted to mention Pisaurum, the historian might easily have received a wrong impression.

It is safer to accept the statement of Procopius as it stands, but the question is not important for the subsequent course of events. By one of three routes the army could reach the Via Flaminia at Acqualagna, about five miles on the Roman side of Petra Pertusa. In any case Gibbon saw the truth as to the general direction taken by Narses: he ``traversed in a direct line the hills of Urbino and re-entered the Flaminian Way nine miles beyond the perforated rock.''

The situation of the camp of Narses is named by Procopius \texthebrew{—} Busta Gallorum; the difficulty is to identify it. The district here east of the Flaminian Way lay in the ager Sentinas , the limits of which are unknown;

Sentinum

itself was close to Sassoferrato. Somewhere in this district the consuls Fabius and Decius defeated the Gauls (in the ager Sentinas ,

Livy, X .27 ) in B.C. 295, and the name \texthebrew{—} Sepulchres of the Gauls \texthebrew{—} evidently commemorated that defeat, like the

Busta Gallica

at Rome which commemorated their repulse from the city nearly a century before

(Livy, V .42) .

Cluverius in \emph{Italia antiqua}, lib. II c. VI , identified Busta Gallorum with a small ``town'' in the Apennines called Bosta or Basta:

Extat hodie in Apennino inter Sentinum, Fabrianum, Matilicam et Sigillum oppida . . . oppidum vulgari vocabulo Bosta: quod plerique notiore vulgaris linguae vocabulo, quod Latine valet \emph{sufficit}, seu \emph{satis est}, adpellant Basta.

No town or village of the name seems to exist now: but Cluver doubtless meant the castle called Bastia, a few miles west of Fabriano, close to the railway connecting that town with Sassoferrato (it will be found marked on maps of the Italian Touring Club). If so, his description of its situation is incorrect, as it does not lie between Matelica and any of the other three towns.

Colucci in his \emph{Antichità Picene}, vol. VIII pp42\texthebrew{‑}106, has discussed at great length the two questions where the Romans defeated

p290
the Gauls and where Narses defeated Totila. He identifies Bastia with Cluver's Basta and with Busta Gallorum, and comes to the conclusion that the battle of Narses was fought in the plain south of Sassoferrato (Piano della Croce) and the battle of Fabius and Decius further south ``nella gran pianura in cui ora esiste Fabriano.''

A very different view from these was developed by Hodgkin,\texthebrew{⁠}\footnote{\emph{Italy and Her Invaders}, IV.710\texthebrew{‑}713 , and 726\texthebrew{‑}728 ; and a special memoir (1884) referred to p726.

} who, having brought Narses from Sena to Cagli, makes him advance along the Via Flaminia and encamp at Scheggia, ``or at some point south of that place where the valley is somewhat broader.''\texthebrew{⁠}\footnote{We have several lists of the stations on the Via Flaminia: \emph{Itin. Ant.} 125, 310;

\emph{Itin. Hier.} 613 ; \emph{Tab. Peut.} ; and

CIL XI.3281 \texthebrew{‑}3284 . They are set out together for the section we are concerned with in CIL XI p995. The distances are:

Ad Calem to Aesis (Scheggia)

Aesis to Helvillum (Sigillo)

6 miles

Helvillum to Tadinae

7 miles

Tadinae is 1½ mile from Gualdo Tadino, near the church of S. Maria Tadinae (Nissen, \emph{op. cit.} 392). Above Tadinae comes in from the west the road from Iguvium, from the east a road from the valley of the Aesis (Esino).

} As Totila encamped at Tadinae, the distance between the two camps would be about 15 miles. Procopius says that the distance was 100 stadia, which is about 14 Roman miles. The battle was fought, according to Hodgkin, somewhere south of Scheggia and north of Tadino. He thinks it ``safe to disregard the Busta Gallorum of Procopius altogether.''

The only other theory he considers is that the battlefield was near Sassoferrato, and he rules this out on the ground that Totila could not have marched thither ``consistently with the narrative of Procopius. The best of the roads between Tadino and Sassoferrato is a high mountain pass, somewhat resembling the Pass of Glencoe. The rest are little more than mountain paths carried through deep gorges in which no armies could manoeuvre.''

There are three objections to Hodgkin's reconstruction of the event. (1) The ground along the Via Flaminia between Tadinae and Aesis in unsuitable for such a battle as Procopius describes. (2) It is very difficult to believe that, if this had been the scene, Aesis or Helvillum, or the Via Flaminia itself, would not have been mentioned to Procopius by his informants, who knew the place where Totila encamped and the village to which his body was carried. (3) No account is taken of the name Busta Gallorum.

The data seem to me to point to the conclusion that the battle was fought somewhere between the Via Flaminia and the river Aesis (Esino), and probably in the neighbourhood of Fabriano. Narses, having reached Acqualagna, marched southward along the Via Flaminia as far as Cagli, and there diverging to the east proceeded by a road, passing (the village of) Frontone and Sassoferrato, to the valley of the Bono (a stream which joins the Esino some

p291
miles east of Fabriano). He expected that Totila would march along the Via Flaminia \texthebrew{—} probably news of the movements of the Goths reached him at Cagli \texthebrew{—} and he decided to choose his own battleground.

In order to reach the camping place of Narses, Totila (if his camp was actually at the station of Tadinae) would only have to march a few miles northward along the Via Flaminia to the place which is now Fossato and there diverge to his right and cross the Colle di Fossato, by the same route which leads to\texthebrew{‑}day from Fossato to Fabriano. Assuming that Narses was encamped west of Fabriano, Totila in descending the mountain pass could have turned to the left (a road is marked on modern maps) and reached Melano, which is about halfway between Fabriano and Bastia.

There seems, however, to be a considerable probability in the conjecture that the camp of Totila was not actually at the station of Tadinae, but some kilometres to the north of it, at Fossato, ``the Camp''; for if the Gothic camp was pitched here on the occasion of the memorable battle, the original of the place-name Fossatum is accounted for. As there was no Roman station here, the locality would naturally be associated by the informants of Procopius with the name of one of the nearest stations, either Tadinae or Helvillum.

Procopius gives two indications of distance. He says that (1) the distance between the camps, that is between ``Tadinae'' and Busta Gallorum, was 100 stades, \emph{i.e.} somewhat more than 14 Roman miles, and that (2) the distance from the place where Totila was wounded to the village of Caprae was 84 stades, \emph{i.e.} 12 Roman miles. Caprae has been identified, no doubt rightly, with the little village of Caprara which lies to the west of the Via Flaminia, about six kilometres to the south-west of Fossato.\texthebrew{⁠}\footnote{According to Nissen (\emph{op. cit.} p393) the Via Flaminia ran far to west of the present road: ``Die Strasse lief nicht wie jetzt an der östlichen Seite der Einsenkung über Fossato und Gualdo sondern an der entgegensgesetzten Seite an Caprara vorbei.''

\texthebrew{◂} previous

next \texthebrew{▸}Images with borders lead to more information.

The thicker the border, the more information.

(Details here.)

UP TO:

J. B. Bury
Homepage

Latin \& Greek

Texts

LacusCurtius

Home} But as we do not know how far Totila had fled from the battlefield before he was wounded, this second indication does not help us much. The first indication, however, is closely in accordance with my theory if the camp of Totila was at Fossato. For the distance from Fossato, by Melano, to the neighbourhood of Bastia is about 20 kilometres, which is equivalent to about 14 Roman miles, the distance given by Procopius.

I must acknowledge help which I have received from my friend Mr. E. H. Freshfield in investigating this subject. I have had the advantage of seeing in MS. a study (soon to be published) of the Via Flaminia by Mr. Ashby and Mr. Fell , and they introduced me to the book of Montecchini.

\xchapter{Chapter XIX, Part D}

Short URL for this page:

bit.ly/BURLAT19D

\xsubsection{(Part 4 of 4)}

In the meantime Narses had been engaged in establishing an ordered administration in Italy, and restoring the life of the

p282
provinces and their cities which had suffered so much through the long war. Though officially he held a military post, he acted as viceroy, and was evidently supreme over the civil functionaries as well as over the army. He had at his side a Prefect, Antiochus, at the head of the civil service, but it is significant that the title of Antiochus was not Praetorian Prefect, but simply Prefect of Italy.

The general lines for the reorganisation were laid down by the Emperor in a law which he addressed to Narses and Antiochus in August A.D. 554, and which he described as a Pragmatic Sanction.\texthebrew{⁠}\footnote{Considerable extracts are preserved, and will be found in editions of the Novellae (\emph{e.g.} App. vii in Kroll's ed.). Another law, relating to debts incurred before the Frank invasion, and the rights of creditors, is incompletely preserved (App. viii). For the title and powers of the Prefect of Italy see Diehl, \emph{Études sur l'adm. byz.} 157 \emph{seq.}

} It was supremely important for the Italians to know immediately how far the Imperial Government would recognize the acts of Gothic rulers, particularly in regard to property. This law provides that henceforward the enactments of the Imperial Code shall apply to Italy as well as to the other parts of the Empire. All grants that were made to individuals or corporations by Athalaric, Amalasuntha, and Theodahad shall be valid, but all grants made by the tyrant Totila are annulled. All contracts made between Romans in besieged towns during the war shall remain valid.\texthebrew{⁠}\footnote{§ 6 quod enim ritu perfectum est, per fortuitos belli casus subverti subtilitatis non patitur ratio.

} In many cases during the war and the Frank invasion people had been forced to flee from their homes and their property had been occupied by others; it is enacted that their property must be restored to them. The old regulations allocating funds for the repair of public buildings in Rome, for dredging the bed of the Tiber, for the repair of the aqueducts are all confirmed, and doles of food are to be supplied to the Roman populace as of old. A remarkable innovation is made in regard to provincial governors. They are no longer to be appointed from above, but to be elected for each province from among its residents by the bishops and magnates. This change may have had some arguments in its favour, but

p283
it was evidently conceived in the interests of the large landed proprietors and must have increased their local power. In other regulations we see the desire to relieve the burden of taxation so far as was deemed compatible with Imperial needs.

The boundaries of the provinces\texthebrew{⁠}\footnote{A new province, Alpes Cottiae, seems to have been cut off from Liguria; Diehl, \emph{op. cit.} p3.

} and the general system of the civil service\texthebrew{⁠}\footnote{As to the meagre evidence for the vicarius Italiae residing at Milan) and the vicarius urbis Romae see Diehl, \emph{op. cit.} p161.

} remained as they had been before the war. It is to be observed, however, that Sicily was not included in Italy. It remained under its own Praetor, who was independent of the Imperial authorities at Ravenna, and from whose courts the appeal was to the Quaestor of the Sacred Palace at Constantinople.\texthebrew{⁠}\footnote{Cp. above,

p216 ; Diehl, p169.

} Sardinia and Corsica were under the viceroy of Africa.

Narses administered Italy for thirteen years after the defeat of the Frank invaders, presiding over the work of reconstruction.\texthebrew{⁠}\footnote{His usual title was simply patricius (see next note; Pelagius, \emph{Ep.} 2, \emph{P. L.} lxix.393 patricius et dux in Italia.

} The walls and gates of Rome were restored, and one of the few memorials of the time records the rebuilding of a bridge across the Anio, which had been destroyed by the Goths, \texthebrew{•} about two miles from the city on the Via Salaria.\texthebrew{⁠}\footnote{It was destroyed in 1798. The inscription (CIL VI.1199 ) is dated in 565 and records how Narses expraeposito sacri palatii, excons. atque patricius after his Gothic victory, ipsis eorum regibus celeritate mirabili conflictu publico superatis atque prostratis, libertate urbis Romae ac totius Italiae restituta restored the bridge a nefandissimo Totila tyranno destructum. It concludes with eight verses of which the last two are

qui potuit rigidas Gothorum subdere mentes,

hic docuit durum flumina ferre iugum.

} Perhaps the most troublesome concern with which the Patrician was called upon to deal was the danger of ecclesiastical strife arising out of the Ecumenical Council of Constantinople in A.D. 553. The circumstances of that assembly will be described in another chapter. The Pope Vigilius who had been forced against his will to subscribe to its decisions died on his way back to Rome on January 7, A.D. 555, and his archdeacon Pelagius was, at the instance of the Emperor, consecrated as his successor on April 13. Pelagius was unpopular in Italy; he was suspected of having in some way caused the death of Vigilius, and only two Italian bishops could be found willing to consecrate him. Narses was present at the ceremony at St. Peter's, and Pelagius took the Gospels in his hand and swore that he was innocent.\texthebrew{⁠}\footnote{See his life in \emph{Lib. Pont.}

} His oath

p284
calmed the popular feeling, but, if he had had his way, he would soon have created a dangerous schism in the Italian Church. In northern Italy particularly, the opinion of the bishops was against the decisions of the recent Council, while the new Pope was determined to enforce them and expel from their sees those who refused to accept them. He wrote repeatedly to Narses requesting or rather requiring of him to use the secular arm against the contumacious bishops.\texthebrew{⁠}\footnote{Pelagius, \emph{Epp.} ii.4 (\emph{P. G.} lxix pp392, 397).

} Narses wisely declined to do anything, and the Imperial government, in the interests of peace, adopted throughout the Empire the policy of suspending the anathemas of the Council and allowing time to heal the discord which the controversy had caused. This unusual moderation, which we may probably attribute to the advice of Narses, was success ­ful. If the matter had rested with the Pope, the Church in Italy would have been rent in twain at a moment when concord and peace were imperatively needed.

The secluded city of the marshes continued to be the seat of government in Italy under Justinian and his successors until it was lost to the Empire in the eighth century. The Empress Placidia had lavished money in making it a treasure-house of art; the barbarian king Theoderic had lived up to her example; and after its recovery by their armies, Justinian and Theodora, who knew it only by reputation, were eager to associate their names with the artistic monuments of Ravenna.

The octagonal church of St. Vitalis, close to Placidia's mausoleum, had been designed and begun under the regency of Amalasuntha, and the building was continued during the war, perhaps by the Ostrogoths themselves. But it was completed and decorated under the auspices of Justinian and Theodora, who made it peculiarly their own, \texthebrew{—} a monument of the Imperial restoration. It was consecrated by the archbishop Maximian in A.D. 547, the year before the death of the Empress, and in the mosaic decoration of the apse the most striking pictures are those of the two sovrans fa­cing each other offering their gifts to the church. But it was not only by their portraits that they appropriated St. Vitalis. Justinian gave it his own impress in the scheme of the Scriptural scenes which are portrayed. They are not simply, as in the other Ravennate churches, intended to illustrate sacred history. The motive is theological, they are

p285
designed to inculcate doctrine, probably the orthodox view on the question which was agitating the world, the two natures of Christ.\texthebrew{⁠}\footnote{See Dalton, \emph{Byz. Art.} 357\texthebrew{‑}360.

} The effects are fine, but these mosaics are far from possessing the charm of those which adorn the sepulchral church of Placidia.

Another church which had been begun by the Goths during the war and was left to their conquerors to complete was dedicated two years later ( A.D. 549) to St. Apollinaris, not in the city itself but in the port of Classis. But many of the mosaics of this basilica, which still stands in the marshes, were executed at a later period; among them is the portrait of an Emperor who ascended the throne a hundred years after Justinian's death.\footnote{Constantine IV.

}

The decorations of Theoderic's basilica of St. Martin were completed under Justinian, and a mosaic representation of the Emperor's bust was put up on the façade,\texthebrew{⁠}\footnote{During the last years of the reign (553\texthebrew{‑}566).

} but was afterwards transferred to a chapel in the interior where it may still be seen. In his time the church was still St. Martin's; it was not till the ninth century that it received the remains of Apollinaris, the tutelary saint of Ravenna, and was re-dedicated to his name.\footnote{Another church of the Justinianean period with St. Michael's (in Affricisco), A.D. 545. Its mosaics were sold to the king of Prussia in 1847 and are preserved in the Berlin museum.

}

The island city, which was later to become the queen of the Hadriatic, had not yet been founded. But it is probable that long before the reign of Justinian inhabitants of the Venetian mainland had been settling in the islands of the lagoons, Malamocco and Rialto, as a secure retreat where they could escape such dangers as the invasions of Alaric and Attila. Under Gothic rule we find the people of this coast in possession of numerous ships, and they were employed to transport wine and oil from Istria to Ravenna. The minister Cassiodorus, in a picturesque despatch, calling upon them to perform this office, likens them to sea-birds.\texthebrew{⁠}\footnote{\emph{Variae}, xii.24.

} But though danger from Visigoth and Hun may have prepared the way for the rise of a city in the lagoons, it was not till three years after Justinian's death, when the Lombards descended into the land, that any such

p286
large and permanent settlements were made on the islands that they could properly be described as the foundation of Venice.\footnote{See Kretschmayr, \emph{Geschichte von Venedig}, p19. The foundation of Grado was older, but it was the Lombard invasion that transformed it into an important city and made it definitely the residence of the Patriarch of Aquileia.

}

It is impossible to say whether Justinian in the early years of his reign had formed any definite plan for reconquering Spain, but we may be sure that it was one of his ambitions, and that if the fall of Witigis had led immediately to the recovery of Italy, he would have sought a prize for carrying his victorious arms against the Visigoths. But before he had completed the subjugation of the Ostrogoths he was invited to intervene in Spain, and, although the issue of the Italian war was still far from certain, he did not hesitate to take advantage of the occasion.

Theoderic, who was regent of the Visigothic kingdom during the minority of his grandson Amalaric, had entrusted the conduct of affairs to Theudis, a capable general, and after the death of Theoderic and the end of the regency Theudis continued to be the virtual ruler. The young king, who had none of the qualities of either his father or his grandfather, married a Frank princess, and this mixed marriage proved unfortunate. Amalaric behaved so brutally to her because she refused to embrace his Arian faith that she invoked the aid of her brother king Childebert, and he advanced against Narbonne. Amalaric marched to defend his Gallic possessions, was defeated in battle, and was then slain in a mutiny of his own army ( A.D. 531).\texthebrew{⁠}\footnote{\emph{Chron. Caesarum} (in \emph{Chron. min.} ii p223).

} The throne was seized by Theudis, who reigned for seventeen years, and after a short intervening reign\texthebrew{⁠}\footnote{Thiudigisalus, 548\texthebrew{‑}59.

} was succeeded by Agila ( A.D. 549). But Agila was not universally acceptable to the people; civil war broke out, and after a struggle of five years he was overthrown by his opponent Athanagild, who ascended the throne ( A.D. 554).

In this struggle Athanagild sought the support of the Emperor, and the Emperor sent a fleet to the southern coasts of Spain. The commander of this expedition was the octogenarian patrician

p287
Liberius, who, it will be remembered, had set out to defend Sicily against Totila, and had hardly reached the island before a more experienced general was sent to take his place.\texthebrew{⁠}\footnote{See above,

p255 . He died in Italy (later than 554); and was buried at Ariminum
(CIL XI.382 ) .

} As he appears not to have returned to Constantinople till late in A.D. 551, it is probable that he received commands to sail directly to Spain with the troops who had accompanied him to Sicily, in A.D. 550, for the date of his expedition cannot have been later than in this year. As the armament must have been small, it achieved a remarkable success. Many maritime cities and forts were captured.\texthebrew{⁠}\footnote{Jordanes, \emph{Get.} c58 . As he was writing in 551, we cannot place the expedition later than in 550. Isidore, \emph{Chron.} 399 ;

\emph{Hist. Goth.} p286 . For the return of Liberius to Constantinople see Procopius, \emph{B. G.} IV.24.1.

} They were captured professedly in the interests of Athanagild, but when Athanagild's cause had triumphed, the Imperialists refused to hand them over and the Visigoths were unable to expel them. Athanagild recovered a few places,\texthebrew{⁠}\footnote{Isidore, \emph{ib.}; Greg. Tur. \emph{H. Fr.} iv.8. Athanagild reigned 554 to 567.

} but Liberius had established an Imperial province in Baetica which was to remain under the rule of Constantinople for about seventy years. There can be no doubt that this change of government was welcomed by the Spanish-Roman population.

We have very few details as to the extent of this Spanish province. It comprised districts and towns to the west as well as to the east of the Straits of Gades; it included the cities of New Carthage, Corduba, and Assionia;\texthebrew{⁠}\footnote{John Biclar. \emph{Chron., sub} 570. Dahn (\emph{Kön. der Germ.} v p178) defines two groups of towns, (1) eastern on the Mediterranean, from Colopona to Sucruna, and (2) western, including Lacobriga and Ossonoba. See also Altamira, in \emph{C. Med. H.} ii p164.

} we do not know whether at any time it included Hispalis. It was placed under a military governor who had the rank of Master of Soldiers, but we do not know whether he was independent or subordinate to the governor of Africa.\footnote{An inscription of New Carthage, of A.D. 589, records that Comentiolus, sent by the Emperor Maurice to defend the province, bore the title of magister militum Spaniae

(CIL II.3420) .

}

It is curious that the two well-informed historical writers who have narrated the fortunes of Justinian's armies in Italy in these years, Procopius and Agathias, should not have made even an incidental reference to this far-western extension of Roman rule. But Agathias was a poet as well as a historian,

p288
and in verses which describe how Justinian has girdled the world with his empire, he alludes to the conquest of which in his History he was silent. Let the Roman traveller, he says, follow the steps of Hercules over the blue western sea and rest on the sands of Spain, he will still be within the borders of the wise Emperor's sovranty.\footnote{In the introduction to the Anthology which he edited (\emph{Anth. Gr.} IV.3.82 \emph{sqq.}). The passage ends with:

\textgreek{ο\texthebrew{ὐ}δ\texthebrew{ὲ} γ\texthebrew{ὰ}ρ \texthebrew{ὀ}θνείης σε δεδέξεταιτ\texthebrew{’ ἤ}θεα γαίης, \texthebrew{ἀ}λλ\texthebrew{ὰ} σοφο\texthebrew{ῦ} κτεάνοισιν \texthebrew{ὁ}μιλήσεις βασιλ\texthebrew{ῆ}ος, \texthebrew{ἔ}νθα κεν \texthebrew{ἀ}ϊξειας, \texthebrew{ἐ}πε\texthebrew{ὶ} κυκλώσατο κόσμον κοιρανί\texthebrew{ῃ}.}In l. 82 \textgreek{κυανωπ\texthebrew{ὸ}ν \texthebrew{ὑ}π\texthebrew{ὲ}ρ δύσιν} means the waters of the west Mediterranean. There may be an allusion to the Spanish conquest in the poem of Paul the Silentiary on St. Sophia,

v. 228 :

\textgreek{\texthebrew{ἠ}ρεμέει κα\texthebrew{ὶ} Μ\texthebrew{ῆ}δος \texthebrew{ἄ}ναξ κα\texthebrew{ὶ} Κελτ\texthebrew{ὶ}ς \texthebrew{ὁ}μοκλή.}Cp. vv. 11\texthebrew{‑}13.

\texthebrew{◂} previous

next \texthebrew{▸}Images with borders lead to more information.

The thicker the border, the more information.

(Details here.)

UP TO:

J. B. Bury
Homepage

Latin \& Greek

Texts

LacusCurtius

Home}

\xchapter{Chapter XX}

Short URL for this page:

bit.ly/BURLAT20

Justinian was not less energetic in increasing the prestige and strengthening the power of the Empire by his diplomacy than by his arms. While his generals went forth to recover lost provinces, he and his agents were incessantly engaged in maintaining the Roman spheres of influence beyond the frontiers and drawing new peoples within the circle of Imperial client states. The methods were traditional and are familiar, but he pursued and developed them more systematically than any of his predecessors. Youths of the dynasties ruling in semi-dependent countries were educated at Constantinople, and sometimes married Roman wives. Barbarian kinglets constantly visited the capital, and Justinian spared no expense in impressing them with the majesty and splendour of the Imperial court. He gave them titles of Roman rank, often with salaries attached; above all, if they were heathen, he procured their conversion to Christianity. Baptism was virtually equivalent to an acknowledgment of Roman overlord ­ship. He used both merchants and missionaries for the purposes of peaceful penetration. And he understood and applied the art of stirring up one barbarian people against another.\texthebrew{⁠}\footnote{He is praised for his dexterity in this art in the contemporary anonymous treatise \emph{\textgreek{Περ\texthebrew{ὶ} στρατηγικ\texthebrew{ῆ}ς}}, II.4 p58.

} Perhaps no Emperor practised all these methods, which are conveniently comprehended under the name of diplomacy, on such a grand scale as Justinian, who was the last to aspire to the Imperial ideal expressed by the Augustan poet:

The objects aimed at varied in different quarters. On

p293
some frontiers they were mainly political, on others largely commercial. In the north, the European provinces against invasion by managing the rapacious barbarians who lived within striking distance. In the Caucasian regions, the chief concern was to contend against the influence of Persia. In the neighbourhood of the Red Sea commercial aims were predominant. In a general survey of these multifarious activities it will be convenient to notice the hostile invasions which afflicted the Balkan provinces during this reign and the system of fortifications which was constructed to protect them, and to describe the general conditions of commerce. We have already seen examples of the Emperor's diplomatic methods in his dealings with the Moors and with the Franks.\footnote{There is a comprehensive survey of ``the diplomatic work'' of Justinian in Diehl's monograph p367 \emph{sqq.} Commerce is treated separately (533 \emph{sqq.}).

}

The array of barbarous peoples against whom Justinian had to protect his European subjects by diplomacy or arms, from the Middle Danube to the Don, were of three different races. There were Germans and Huns as before, but a third group, the Slavs, were now coming upon the scene. The German group consisted of three East German peoples, the Gepids of Transylvania, and the Heruls and the Langobardi to the north-west of the Gepids. The Huns were represented by the Bulgarians of Bessarabia and Walachia, and the Kotrigurs further east. The Slavs lived in the neighbourhood of the Bulgarians on the banks of the lower Danube in Walachia.

This general disposition of peoples had resulted from the great battle of the Netad which dissolved the empire of Attila. One of the obscure but most important consequences of that event was the westward and southward expansion of the Slavs towards the Elbe and towards the Danube.

It has been made probable by recent research that the prehistoric home of the Slavs was in the marshlands of the river Pripet, which flows into the Dniepr north of Kiev.\texthebrew{⁠}\footnote{This theory is based on a combination of botanical and linguistic evidence. It was originated (1908) by Rostafinski and has been developed by Peisker. The Slavs have no

(p294)
native words for the beech, the larch, and the yew, but they have a word for the hornbeam; hence their original home must have lain in the hornbeam zone, but outside the zones of the other trees, and this consideration determines it as Polesia. The brief sketch I have given of the primitive Slavs is derived from the writings of Peisker (see

Bibliography ), especially from \emph{C. Med. H.} II chap. XIV. Rostafinski's article, \emph{Les demeures primitives des Slaves}, will be found in \emph{Bull. de l'Acad. des Sciences de Cracovie}, Cl. de phil. 1908.

} This unhealthy district, known as Polesia, hardly half as large as England,

p294
is now inhabited by White Russians. It could produce little corn º as it could only be cultivated in spots, and it was so entirely unsuitable for cattle that the Slavs had no native words for \emph{cattle} or \emph{milk.} They may have reared swine, but perhaps their food chiefly consisted of fish and the manna-grass which grows freely in the marshy soil. The nature of the territory, impeding free and constant intercourse, hindered the establishment of political unity. The Slavs of Polesia did not form a state; they had no king; they lived in small isolated village groups, under patriarchal government.

Their history, from the earliest times, was a tragedy. Their proximity to the steppes of Southern Russia exposed them as a prey to the Asiatic mounted nomads who successively invaded and occupied the lands between the Don and the Dniester. Living as they did, they could not combine against these enemies who plundered them and carried them off as slaves. They could only protect themselves by hiding in the forest or in the waters of their lakes and rivers. They built their huts with several doors to facilitate escape when danger threatened; they hid their belongings, which were as few as possible, in the earth. They could elude a foe by diving under the water and lying for hours on the bottom, breathing through a long reed, which only the most experienced pursuers could detect.\footnote{Pseudo-Maurice, \emph{Strateg.} XI.5. (It has been conjectured by Kulakovski, \emph{Viz. Vrem.} VII.108 \emph{sqq.}, that the word \textgreek{πλωταί} which occurs in this chapter for rafts or flosses is the Slavonic plot.) The accounts of the manners of the Slavs in this sixth-century treatise and in

Procopius, \emph{B. G.} III.14 , are in general agreement and supplement each other. For their religion (cult of fire, worship of nymphs and rivers) see

Peisker, \emph{op. cit.} p425 ; Jireček, \emph{Geschichte der Bulgaren}, 102 \emph{sqq.} The Slavs under this name are, I think, first mentioned in the fourth century by Caesarius (brother of Gregory of Nazianzus), \emph{Quaestiones}, \emph{P. G.} XXXVIII p985: \textgreek{ο\texthebrew{ἱ} Σκλαυηνο\texthebrew{ὶ} κα\texthebrew{ὶ} Φυσων\texthebrew{ῖ}ται, ο\texthebrew{ἱ} κα\texthebrew{ὶ} Δανούβιοι προσαγορευόμενοι, ο\texthebrew{ἱ} μ\texthebrew{ὲ}ν γυναικομαστοβορο\texthebrew{ῦ}σιν \texthebrew{ἡ}δε\texthebrew{ῶ}ς δι\texthebrew{ὰ} τ\texthebrew{ὸ} πεπληρ\texthebrew{ῶ}σθαι το\texthebrew{ῦ} γάλακτος, μυ\texthebrew{ῶ}ν δίκην το\texthebrew{ὺ}ς \texthebrew{ὑ}ποτίτθους τα\texthebrew{ῖ}ς πέτραις \texthebrew{ἐ}παράττοντες, ο\texthebrew{ἱ} δ\texthebrew{ὲ} κα\texthebrew{ὶ} τ\texthebrew{ῆ}ς νομίμης κα\texthebrew{ὶ ἀ}διαβλήτου κρεωβορίας \texthebrew{ἀ}πέχουσιν}. See Müllenhoff, \emph{Deutsche Altertumskunde}, II p367.

Thayer's Note: ``Long reed'' is suspect on technical grounds. A useful breathing tube or snorkel cannot be more than about 50 cm long, since beyond that human lungs, subjected to increasing water pressure, do not have the strength to expel the breathed air to the surface and exchange it for fresh. The physics \texthebrew{—} of the extreme case in which the lungs are assumed to have a 100\% oxygen absorption efficiency \texthebrew{—} are given in detail in

``Theoretical Limit to the Depth at Which a Snorkel May Be Used''

(paper by A. Toohie, J. McGuire and A. Pohl of the Department of Physics and Astronomy, University of Leicester).

}

At a time of which we have no record the Slavs began to spread silently beyond the borders of Polesia, northward, eastward, and southward. In the fourth century they were conquered

p295
by Hermanric, king of the Ostrogoths, and included in his extensive realm.\texthebrew{⁠}\footnote{Jordanes, \emph{Get.} 119 , where they are called Veneti (as in Pliny and Tacitus). They attempted to resist, numerositate pollentes \texthebrew{—} sed nihil valet multitudo imbellium. We can put no credence in what Jordanes (after Cassiodorus) tells us of Hermanric's immediate successors (which is at variance with statements of Ammian), and I cannot accept (as Peisker does,

\emph{op. cit.} p431 ) his statement that King Vinithar subdued the Antae soon after the Hunnic invasion (\emph{ib.} 247).

} They enjoyed a brief interlude of German tyranny instead of nomad raids; then the Huns appeared and they were exposed once more to the oppression which had been their secular lot. They had probably learned much from the Goths; but when they emerge at length into the full light of history in the sixth century, they still retained most of the characteristics which their life in Polesia had impressed upon them. They lived far apart from one another in wretched hovels;\texthebrew{⁠}\footnote{Procopius, \emph{ib.} 24.

} though they had learned to act together, they did not abandon their freedom to the authority of a king. Revolting against military discipline, they had no battle array and seldom met a foe in the open field.\texthebrew{⁠}\footnote{Pseudo-Maurice, who describes them as \textgreek{\texthebrew{ἄ}ναρχα κα\texthebrew{ὶ} μεισάλληλα} (pp275\texthebrew{‑}276).

} Their arms were a shield, darts, and poisoned arrows.\texthebrew{⁠}\footnote{\emph{Ib.}, and

Procopius, \emph{ib.} 25 , who says that some of them went into battle without tunic or cloak, and wearing only trousers. He describes them as tall and brave, and in complexion reddish.

} They were perfidious, for no compact could bind them all; but they are praised for their hospitality to strangers and for the fidelity of their women.

As we might expect, they had no common name. Slav , by which we designate all the various peoples who spread far and wide in Eastern Europe from the original Polesian home, comes from Slovene , which appears originally to have been a local name attached to a particular group dwelling at a place called Slovy; and the fortunes of the name are due to the fact that this group was among the first to come into contact with the Roman Empire. Before the reign of Justinian these Sclavenes, as the historian Procopius calls them,\texthebrew{⁠}\footnote{Jordanes also has Sclaveni (\emph{e.g.} \emph{Rom.} 388), distinct from Antae. In Pseudo-Maurice we get as the generic term \textgreek{Σκλάβοι}. Procopius says

(\emph{ib.} 29)

that Antae and Sclavenes had originally a common name \textgreek{Σπόροι}, which, according to Dobrovsky and Šafarik (\emph{Slav. Altertümer}, I.95), is a corruption of Srbi (Serbs). The thesis maintained by Šafarik and Drinov, and defended by Jireček, that Slavs had begun to settle into the Balkan Peninsula already in the third century A.D., and that the Carpi and Kostoboks were Slavonic peoples, must be rejected as resting on insufficient evidence. See Šafarik, \emph{op. cit.} I.213 \emph{sq.}, Jireček, \emph{op. cit.} ch. III.

} had along with another kindred people, the Antae, settled in the neighbourhood of the Bulgarians,

p296
along the banks of the Lower Danube. Antae is not a Slavonic name, and it is not unlikely that they were a Slavonic tribe which had been conquered and organised by a non-Slavonic people \texthebrew{—} somewhat as in later times the Slavs of Moesia were conquered by the Bulgarians and took their name. However this may be, these new neighbours of the Empire now began to exchange the rôle of victims for that of plunderers.

Like the Huns, the Antae and Sclavenes supplied auxiliaries for the Roman army.\texthebrew{⁠}\footnote{See

Procopius, \emph{B. G.} I.27.2 . They must have supplied recruits already in the fifth century, for in 468 we meet a man of Slavonic name (Anagast) who had risen to be Mag. mil. of Thrace. See above
vol. I p434 .

} And along with the Huns they were always watching for an opportunity to cross the Danube and plunder the Roman provinces. In the invasions which are recorded in the reign of Justinian, it is sometimes the Slavs, sometimes the Bulgarians who are mentioned, but it is probable that they often came together. In A.D. 529 the Bulgarians overran Lower Moesia and Scythia. They defeated Justin and Baduarius, the generals who opposed them, and crossing the Balkan passes, invaded Thrace.\texthebrew{⁠}\footnote{John Mal. XVIII.437. Theophanes, A.M. 6031. Justin was slain. Baduarius is not to be confused with his namesake, son-in\texthebrew{‑}law of the Emperor Justin II.

} There they captured another general, Constantiolus, and obtained from the Emperor ten thousand pieces of gold for his release. Another incursion in the following year was repulsed with numerous losses to the invaders by Mundus, the Master of Soldiers in Illyricum;\texthebrew{⁠}\footnote{\emph{Ib.} 451 \textgreek{Ο\texthebrew{ὗ}ννοι μετ\texthebrew{ὰ} πολλο\texthebrew{ῦ} πλ\texthebrew{ῆ}θους διαφόρων βαρβάρων}, Marcellinus, \emph{sub a.} (Bulgares).

} and Chilbudius, who was appointed Master of Soldiers in Thrace about the same time, not only prevented the barbarians from crossing the Danube for three years, but terrorised them by making raids into their own country. His success made him rash. Venturing to cross the river with too small a force, he was defeated and slain by the Sclavenes. No one of the same ability replaced him, and the provinces were once more at the mercy of the foe.\texthebrew{⁠}\footnote{Procopius, \emph{B. G.} III.14 . Chilbudius was appointed in the year of Justinian, A.D. 530\texthebrew{‑}531, and was slain three years later. Here the \textgreek{Ο\texthebrew{ὗ}ννοι, \texthebrew{Ἄ}νται} and \textgreek{Σκλαβηνοί} are associated as invaders.

} We hear, however, of no serious invasion till A.D. 540, when the Bulgarians, with a host exceptionally huge, devastated the peninsula from sea to sea.\texthebrew{⁠}\footnote{Procopius, \emph{B. V.} II.4 . John of Ephesus, who was then in Constantinople, speaks of Justinian barricading himself in his Palace, \emph{H. E.} Part II p485.

} They forced

p297
their way through the Long Wall and spread terror to the suburbs of the capital. They occupied the Chersonesus, and some of them even crossed the Hellespont and ravaged the opposite coast. They laid waste Thessaly and Northern Greece; the Peloponnesus was saved by the fortifications of the Isthmus. Many of the castles and walled towns fell into their hands,\texthebrew{⁠}\footnote{Thirty-two fortresses in Illyricum were taken, and the town of Cassandrea was captured by assault.

} and their captives were numbered by tens of thousands. This experience moved Justinian to undertake the construction of an extensive system of fortifications which will be described hereafter.

Soon after this invasion a quarrel broke out between the Sclavenes and the Antae, and Justinian seized the opportunity to inflame their rivalry by offering to the Antae a settlement at Turris, an old foundation of Trajan on the further side of the Lower Danube, where as federates of the Empire, in receipt of annual subsidies, they should act as a bulwark against the Bulgarians.\texthebrew{⁠}\footnote{Procopius, \emph{B. G.} III.14 . Turris had long been derelict; Justinian apparently proposed to have it restored at his expense.

} We are not told whether this plan was carried out, but we may infer that the proposal was accepted, from the fact that in the subsequent invasions the Antae appear to have taken no part.\texthebrew{⁠}\footnote{The Antae accepted, on condition that a captive, whom they believed to be Chilbudius (the general who was slain in A.D. 533\texthebrew{‑}534), should organise the settlement. The impostor was sent to Constantinople and captured by Narses in Thrace, and his pretensions were exposed. Procopius does not tell the sequel.

} In A.D. 545 the Sclavenes were thoroughly defeated in Thrace by Narses and a body of Heruls whom he had engaged for service in Italy.\texthebrew{⁠}\footnote{\emph{Ib.} III.13.

} Three years later the same marauders devastated Illyricum as far as Dyrrhachium,\texthebrew{⁠}\footnote{\emph{Ib.} III.29.1\texthebrew{‑}3.

} and in A.D. 549 a band of 3000 penetrated to the Hebrus, where they divided into two parties, of which one ravaged Illyricum and the other Thrace. The maritime city of Topirus was taken, and the cruelties committed by the barbarians exceed in atrocity all that is recorded of the invasions of the Huns of Attila.\texthebrew{⁠}\footnote{Procopius relates this invasion under the year 549\texthebrew{‑}550

(III.38) . I infer that it belongs to 549, from the fact that the next invasion was clearly in the summer of 550 ( III.40.1 ; cp. 39.29 ). It is often placed in 551 (as by Diehl, \emph{op. cit.} 220). The impalings which the Sclavenes practised may have been learned from the Huns.

} In the following summer the Sclavenes came again, intending to attack Thessalonica, but Germanus happened to be

p298
at Sardica, making preparations to take reinforcements to Italy. The terror of his name diverted the barbarians from their southward course and they invaded Dalmatia.\texthebrew{⁠}\footnote{Germanus had formerly inflicted a great defeat on the Antae, when he was Master of Soldiers in Thrace (\emph{ib.} 40.6); the date is unknown.

} Later in the year the Sclavenes, reinforced by newcomers, gained a bloody victory over an Imperial army at Hadrianople,\texthebrew{⁠}\footnote{The defeated army was under well-known leaders: Constantian (Count of the Stable), Aratius, Nazares (who was or had been mag. mil. Illyrici, \emph{B. G.} III.11.18), Justin, son of German us, and John Phagas, but the supreme command was entrusted to Scholasticus, a Palace eunuch, otherwise unknown. The soldiers forced their leaders to give battle against their wish.

} penetrated to the Long Wall, but were pursued and forced to give up much of their booty.

Two years later there was another inroad, and on this occasion the Gepids aided and abetted the Sclavenes, helping them, when they were hard pressed by Roman troops, to escape across the river, but exacting high fees from the booty-laden fugitives.\footnote{A piece of gold for every person they ferried into safety. \emph{Ib.} IV.25.5.

}

Permanent Slavonic settlements on Imperial soil were not to begin till about twenty years after Justinian's death, but the movements we have been following were the prelude to the territorial occupation which was to determine the future history of south-eastern Europe.

The most power ­ful of the barbarous peoples on the Danube frontier, against whom the Emperors had protect their European subjects, were the Gepids of Transylvania. The old policy of recognising them as federates and paying them yearly subsidies, seems to have been success ­ful until Sirmium was taken from the Ostrogoths by Justinian, and being weakly held was allowed to fall into their hands. Establishing themselves in this stronghold they occupied a portion of Dacia Ripensis and made raids into the southern provinces.\texthebrew{⁠}\footnote{Procopius, \emph{B. G.} II.14.

} Justinian immediately discontinued the payment of subsidies and sought a new method of checking their hostilities. He found it in the rivalry of another East-German nation, the Langobardi, who had recently appeared upon the scene of Danubian politics. Yet another people, the Heruls, who belonged to the same group, payed a

p299
minor part in the drama, in which the Gepids and Langobardi were the principal actors, and Justinian the director.

It was more than a century since the Langobardi, or Lombards, as we may call them in anticipation of the later and more familiar corruption of their name, had left their ancient homes on the Lower Elbe, where they were neighbours of the Saxons, whose customs resembled their own, but the details of their long migration are obscure.\texthebrew{⁠}\footnote{The original home of the Langobardi was in Scandinavia, but they had settled in the regions of the Lower Elbe before the time of Augustus. Their southward migration is dated by modern historians as not earlier than the beginning of the fourth century. It is probable that the old interpretation of their name (Long Beards) is the true one (see Blasel, \emph{Die Wanderzüge der Langobarden}, 129 \emph{sqq.}). The chief sources of their early history are the

\emph{Origo gentis Langobardorum}

(\emph{c.} A.D. 650); Fredegarius, \emph{Chron.} III.65 (embodying Lombard tradition); Paulus Diac. \emph{Hist. Lang.} Book I (based on the \emph{Origo}). See, on the difficult geographical and chronological questions connected with the movements of the Lombards,

Hodgkin, \emph{Italy and her Invaders}, vol. V ; Schmidt, \emph{Gesch. der deutschen Stämme}, I.427 \emph{sqq.}; Blasel, \emph{op. cit.} (where a bibliography will be found).

} Soon after the conquest of the Rugians by Odovacar, they took possession of the Rugian lands, to the north of the province of Noricum, but they remained here only for a few years and then settled in the plains between the Theiss and the Danube.\texthebrew{⁠}\footnote{Campi patentes = Feld (\emph{Origo} and Paul, \emph{Hist. Lang.} I.20); which in \emph{Chron. Gothorum}, ch. II, is called Tracia.

} At this time, it was in the reign of Anastasius, they lived as tributary subjects of another East-German people, more savage than themselves. We have already met the Heruls taking part in the overthrow of the Hunnic realm and contributing mercenary troops to the Imperial service. In the second half of the fifth century they seem to have fixed their abode somewhere in North-western Hungary, and when the Ostrogoths left Pannonia they became a considerable and aggressive power dominating the regions beyond the Upper Danube. They invaded the provinces of Noricum and Pannonia, and won overlord ­ship over the Lombards. Theoderic, following his general policy towards his German neighbours, allied himself with their king Rodulf, whom he adopted as a son.\texthebrew{⁠}\footnote{Cassiodorus, \emph{Var.} IV.2, a letter addressed to the king of the Heruls, whose name comes from the Lombard sources. Its date is between 507 and 511, so that the battle must be placed, not \emph{c.} 505 with Schmidt, but at earliest 507\texthebrew{‑}508, and at latest 511\texthebrew{‑}512 (see next note). If it is true that the Lombards moved from Rugia to the Campi patentes three years before the battle (Paul, \emph{ib.}), the earliest date for their change of abode is 504\texthebrew{‑}505. The name of the Lombard king at this time was Tato. Rodulf, the Herul, was slain in the battle (Paul, \emph{ib.}). The best source is

Procopius, \emph{B. G.} II.14 ; the fuller story of Lombard tradition is largely legendary.

} But soon

p300
afterwards ( A.D. 507\texthebrew{‑}512), they attacked the Lombards without provocation and were defeated in a sanguinary battle.

This defeat had important results. It led to the dissolution of the Herul nation into two portions, of which one migrated northward and returned to the old home of the people in Scandinavia. The rest moved first into the former territory of the Rugians, but finding the land a desert they begged the Gepids to allow them to settle in their country. The Gepids granted the request, but repaid themselves by carrying off the cattle and violating their women. Then the Heruls sought the protection of the Emperor, who readily granted them land in one of the Illyrian provinces.\texthebrew{⁠}\footnote{Probably Dacia Ripensis. Marcellinus, \emph{s. a.} 512; Procopius, \emph{ib.} XV.i.1.

} But their rapacious instincts soon drove them to plunder and maltreat the provincials, and Anastasius was compelled to send an army to chastise them. Many were killed off; the rest made complete submission, and were suffered to remain. No people quite so barbarous had ever yet been settled on Roman soil. It was their habit to put to death the old and the sick; and the women were expected to hang themselves when their husbands died. When Justinian came to the throne he effected their conversion to Christianity. Their king with his nobles was invited to Constantinople, where he was baptized with all his party, the Emperor standing sponsor, and was dismissed with handsome gifts. Larger subsidies were granted to them, and better lands in the neighbourhood of Singidunum, with the province of Second Pannonia ( A.D. 527\texthebrew{‑}528).\texthebrew{⁠}\footnote{Procopius, \emph{B. G.} II.14.33 ;

III.34.42 ; John Mal. XVIII.427; John Eph., \emph{H. E.} Part II p475 (\emph{sub a.} 844 = A.D. 533). Cp. Menander, \emph{fr.} 9 (\emph{F. H. G.} IV) .

} Henceforward, for some years, they fulfilled their duties as Federates, and supplied contingents to the Roman army. But though their savagery had been mitigated after they embraced the Christian faith, they were capricious and faithless; they had not even the merit of chaste manners, for which Tacitus and Salvian praise the Germanic peoples; they were the worst people in the whole world, in the opinion of a contemporary historian.\footnote{Procopius, \emph{ib.} II.14.36 \textgreek{κα\texthebrew{ὶ} μίξεις ο\texthebrew{ὐ}χ \texthebrew{ὁ}σίας τελο\texthebrew{ῦ}σιν, \texthebrew{ἄ}λλας τε κα\texthebrew{ὶ ἀ}νδρ\texthebrew{ῶ}ν κα\texthebrew{ὶ ὀ}ρνιθων.}}

Suddenly it occurred to them that they would prefer a republican form of government, though their kings enjoyed only a shadow of authority. Accordingly they slew their king, but

p301
very soon, for they were unstable as water, they repented, and decided to choose a ruler among the people of their own race who had settled in Scandinavia. Some of their leading men were sent on this distant errand and duly returned with a candidate for the throne.\texthebrew{⁠}\footnote{In his account of this episode Procopius (\emph{ib.} 15) designates Scandinavia as Thule and describes it as ten times larger than Britain. Among the peoples who inhabit it he knows of two, the Gauts and the Skrithifinoi. Of the Gauts in Sweden we otherwise know, and it is natural to identify the Skrithifinoi with the Finns.

} But in the meantime, during their long absence, the Heruls, with characteristic indecision, bethought themselves that they ought not to elect a king from Scandinavia without the consent of Justinian, and they invited him to choose a king for them. Justinian selected a certain Suartuas, a Herul who had long lived at Constantinople. He was welcomed and acclaimed by the Heruls, but not many days had passed when the news came that the envoys who had gone to Scandinavia would soon arrive. Suartuas ordered the Heruls to march forth and destroy them; they obeyed cheerfully; but one night they all left him and went over to the rival whom they had gone forth to slay. Suartuas returned alone to Constantinople.

The consequence of this escapade was that the Heruls split up again into two portions. The greater part attached themselves to the Gepids; the rest remained federates of the Empire.\texthebrew{⁠}\footnote{Procopius, \emph{B. G.} III.34.43, who says that the total fighting strength of the Heruls was 4500 men, of whom 3000 joined the Gepids. Cp. II.15.37.

} This was the position of affairs when about the middle of the sixth century war broke out between the Gepids and the Lombards.

The Lombards are represented as having been Christians while they were still under the yoke of the Heruls. After they had won their independence they lived north of the Danube in the neighbourhood of the Gepids.\texthebrew{⁠}\footnote{This was the period of the linguistic change, which is known as the second shifting of consonants and produced the High German language. It originated in southern Germany, and the Lombard language was affected by it.

} We hear nothing more of them until we find their king Wacho, in A.D. 539, refusing to send help to the Ostrogoths on the ground that he was a friend and ally of Justinian.\texthebrew{⁠}\footnote{Proc. \emph{B. G.} II.22. Above,
p205 .

} Some years later the Emperor assigned to them settlements in Noricum and Pannonia,\texthebrew{⁠}\footnote{\emph{Ib.} III.33.10 \textgreek{Ηωρικ\texthebrew{ῶ}ν τε πόλει} (Noreia = Neumarkt) \textgreek{κα\texthebrew{ὶ} το\texthebrew{ῖ}ς \texthebrew{ἐ}π\texthebrew{ὶ} Παννονίας \texthebrew{ὀ}χυρώμασί τε κα\texthebrew{ὶ ἄ}λλοις χωρίοις.}} and granted them the subsidies which it was usual to pay to federates. We

p302
may take it that he deliberately adopted this policy in order to use the Lombards as a counterpoise to the Gepids, with whom he had recently broken off relations.\footnote{Procopius relates two events together, under the fourteenth year of the Gothic War, \emph{i.e.} A.D. 548\texthebrew{‑}549, but in a digression which assigns only the loose date ``when Totila had gained the upper hand'' (\emph{ib.} 7). In the following chapter (III.34) he anticipates the chronology (\textgreek{χρόν\texthebrew{ῷ} δ\texthebrew{ὲ ὔ}στερον}) and narrates the war of the Gepids and Lombards, which was thus subsequent to A.D. 549.

}

It was not long before these two peoples quarrelled and prepared for war. Audoin at this time was king of the Lombards\texthebrew{⁠}\footnote{Audoin (half-brother of Wacho) married the daughter of Hermanfrid, king of the Thuringians. The marriage was arranged by Justinian. For after Hermanfrid's death, his wife Amalaberga (Theoderic's niece) had returned to Italy with her children, and they were after brought to Constantinople by Belisarius. See \emph{B. G.} I.31.2; IV.12.

} and Thorisin of the Gepids. They both sent ambassadors to Constantinople, the Lombards to beg for military aid,\texthebrew{⁠}\footnote{Procopius, who puts long speeches into the mouths of the envoys, makes the Lombards urge that they were Catholics, not Arians like the Gepids (III.34, 24). Yet when they subsequently conquered Italy, they were Arians. They seem to have been exceptionally indifferent to religion. Cp. Hodgkin, \emph{op. cit.} V.158 .

} the Gepids hardly hoping to do more than induce the Emperor to remain neutral. Justinian decided to assist the Lombards and sent a body of 10,000 horse, who were directed to proceed to Italy when they had dealt with the Gepids. These troops met an army of hostile Heruls and defeated them severely, but in the meantime the Lombards and Gepids had composed their differences, to the disappointment of Justinian. It was felt, however, by both sides that war was inevitable and was only postponed. The Gepids, fearing that their enemies, supported by Constantinople, would prove too strong for them, concluded an alliance with the Kotrigurs.

The Kotrigurs, who were a branch of the Hunnic race, occupied the steppes of South Russia, from the Don to the Dniester, and were probably closely allied to the Bulgarians\texthebrew{⁠}\footnote{The name Kotrigur is to be compared with Kotragos in the genealogy of the Bulgarians. Theophanes describes \textgreek{Κότραγοι} near L. Maeotis as \textgreek{\texthebrew{ὁ}μόφυλοι} of the Bulgarians (A.M. 6171).

} or Onogundurs \texthebrew{—} the descendants of Attila's Huns \texthebrew{—} who had their homes in Bessarabia and Walachia. They were a formidable people and Justinian had long ago taken precautions to keep them in check, in case they should threaten to attack the Empire, though it was probably for the Roman cities of the Crimea, Cherson and Bosporus, that he feared, rather than for the Danubian provinces. As his policy on the Danube was to

p303
use the Lombards as a check on the Gepids, so his policy in Scythia was to use another Hunnic people, the Utigurs, as a check on the Kotrigurs. The Utigurs lived beyond the Don, on the east of the Sea of Azov, and Justinian cultivated their friendship by yearly gifts.

When a host of 12,000 Kotrigurs, incited by the Gepids, crossed the Danube and ravaged the Illyrian lands, Justinian immediately despatched an envoy to Sandichl, king of the Utigurs, to bid him prove his friendship to the Empire by invading the territory of their neighbours. Sandichl, an experienced warrior, fulfilled the Emperor's expectations; he crossed the Don, routed the enemy, and carried their women and children into slavery. When the news reached Constantinople, Justinian sent one of his generals\texthebrew{⁠}\footnote{Aratius, the Armenian. The name of the Kotrigur leader was Chinialon.

} to the Kotrigurs who were still plundering the Balkan provinces, to inform them of what had happened in their own land, and to offer them a large sum of money to evacuate Roman territory. They accepted the proposal, and it was stipulated that if they found their own country occupied by the Utigurs, they should return and receive from the Emperor lands in Thrace. Soon afterwards another party of 2000 Kotrigurs, with their wives and children, arrived as fugitives on Roman soil. They were led by Sinnion, who had fought in Africa as a commander of Hunnic auxiliaries in the Vandal campaign of Belisarius. The Emperor accorded them a settlement in Thrace. This complacency shown to their foes excited the jealous indignation of the Utigurs, and king Sandichl sent envoys to remonstrate with Justinian on the injustice and impolicy of his action. They were appeased by large gifts, which it was obviously the purpose of their coming to obtain.\footnote{Procopius, \emph{B. G.} IV.18 and 19. The long speech which the author puts into the mouths of the envoys is, of course, his own criticism of Justinian's policy. The date of these events seems to be A.D. 551. Cp. \emph{ib.} 21, 4.

}

In the following year ( A.D. 552), the war so often threatened and so often postponed between the Lombards and Gepids broke out. The Gepids sought to renew their old alliance with the Empire, and Justinian consented,\texthebrew{⁠}\footnote{\emph{Ib.} 25.8\texthebrew{‑}9. Procopius obliquely criticises Justinian by emphasising the solemnity of the oaths with which the treaty was confirmed.

} but when the Lombards soon afterwards asked him to fulfil his engagements and send

p304
troops to help them he denounced his new treaty with the Gepids on the pretext that they had helped Sclavenes to cross the Danube. Among the leaders of the forces which marched to co-operate with the Lombards, were Justin and Justinian, the Emperor's cousins, but they were detained on their way to suppress a revolt at Ulpiana, and never arrived at their destination. Only those troops which were commanded by Amalfridas, the brother-in\texthebrew{‑}law of the Lombard king,\texthebrew{⁠}\footnote{He was son of Hermanfrid; see above,
p302, n2 .

} pursued their march and took part in the campaign. The Lombards won a complete victory over the Gepids, and Audoin, in announcing the good news to Justinian, reproached him for failing to furnish the help which they had a right to expect in consideration of the large force of Lombards which had recently gone forth to support Roman arms in Italy.

After this defeat the Gepids concluded treaties of perpetual peace with the Lombards and with the Empire,\texthebrew{⁠}\footnote{\emph{Ib.} 27.21.

} and peace seems to have been preserved so long as Justinian reigned. After his death the enmity between these two German peoples broke out again, and the Lombards, aided by other allies, eliminated the name of the Gepids from the political map of Europe.

In a few years the Kotrigurs recovered from the chastisement which had been inflicted upon them by their Utigur neighbours, and in the winter of A.D. 558\texthebrew{‑}559, under a chieftain whose name was Zabergan, a host of these barbarians crossed the frozen Danube, and passing unopposed through Scythia and Moesia, entered Thrace. These provinces would seem to have been entirely denuded of troops. In Thrace Zabergan divided his followers into three armies. One was sent to Greece, to ravage the unprotected country; the second invaded the Thracian Chersonese; the third army, consisting of seven thousand cavalry, rode under Zabergan himself to Constantinople.

The atrocities committed by the third body are thus described by a contemporary writer:\footnote{Agathias V.2 ; cp. John Mal. XVIII p490 ; Theophanes, A.M. 6051 . The Huns were almost a whole year in Roman territory. See Clinton, \emph{F. R.}, \emph{sub} A.D. 559.

}

p305
As no resistance was offered to their course, they overran the country and plundered without mercy, obtaining a great booty and large numbers of captives. Among the rest, well-born women of chaste life were most cruelly carried off to undergo the worst of all misfortunes, and minister to the unbridled lust of the barbarians; some who in early youth had renounced marriage and the cares and pleasures of this life, and had immured themselves in some religious retreat, deeming it of the highest importance to be free from cohabitation with men, were dragged from the chambers of their virginity and violated. Many married women who happened to be pregnant were dragged away, and when their hour was come brought forth children on the march, unable to conceal their throes, or to take up and swaddle the new-born babes; they were hauled along, in spite of all, hardly allowed even time to suffer, and the wretched infants were left where they fell, a prey for dogs and birds, as though this were the purpose of their appearance in the world.

To such a pass had the Roman Empire come that, even within the precincts of the districts surrounding the Imperial city, a \emph{very small} number of barbarians committed such enormities.\texthebrew{⁠}\footnote{Theophanes, \emph{ib.}, notices that two generals, Sergius and Edermas, were defeated by the Huns before they reached the Long Wall.

} Their audacity went so far as to pass the Long Walls and approach the inner fortifications. For time and neglect had in many places dilapidated the great wall, and other parts were easily thrown down by the barbarians, as there was no military garrison, no engines of defence. Not even the bark of a dog was to be heard; the wall was less efficiently protected than a pig-sty or a sheep-cot.

The Huns encamped at Melantias, a village on the small river Athyras, which flows into the Propontis. Their proximity created a panic in Constantinople, whose inhabitants saw in imagination the horrors of siege, conflagration, and famine. The terror was not confined to the lower classes; the nobles trembled in their palaces, the Emperor was alarmed on his throne. All the treasures of the churches, in the tract of country between the Euxine and the Golden Horn, were either carted into the city or shipped to the Asiatic side of the Bosphorus. The undisciplined corps of the Scholarian guards, ignorant of real warfare, did not inspire the citizens with much confidence.

On this critical occasion Justinian appealed to his veteran general Belisarius to save the seat of empire. In spite of his years and feebleness Belisarius put on his helmet and cuirass once more. He relied chiefly on a small body of three hundred men who had fought with him in Italy; the other troops that he mustered knew nothing of war, and they were more for appearance than for action. The peasants who had fled before

p306
the barbarians from their ruined homesteads in Thrace accompanied the little army. He encamped at the village of Chettus, and employed the peasants in digging a wide trench round the camp. Spies were sent out to discover the numbers of the enemy, and at night many beacons were kindled in the plain with the purpose of misleading the Huns as to the number of the forces sent out against them. For a while they were misled, but it was soon known that the Roman army was small, and two thousand cavalry selected by Zabergan rode forth to annihilate it. The spies informed Belisarius of the enemy's approach, and he made a skilful disposition of his troops. He concealed two hundred peltasts and javelin-men in the woods on either side of the plain, close to the place where he expected the attack of the barbarians; the ambuscaders, at a given signal, were to shower their missiles on the hostile ranks. The object of this was to compel the lines of the enemy to close in, in order to avoid the javelins on the flank, and thus to render their superior numbers useless through inability to deploy. Belisarius himself headed the rest of the army; in the rear followed the rustics, who were not to engage in the battle, but were to accompany it with loud shouts and cause a clatter with wooden beams, which they carried for that purpose.

All fell out as Belisarius had planned. The Huns, pressed by the peltasts, thronged together, and were hindered both from using their bows and arrows with effect, and from circumventing the Roman wings. The noise of the rustics in the rear, combined with the attack on the flanks, gave the foe the impression that the Roman army was immense, and that they were being surrounded; clouds of dust obscured the real situation, and the barbarians turned and fled. Four hundred perished before they reached their camp at Melantias, while not a single Roman was mortally wounded. The camp was immediately abandoned, and all the Kotrigurs hurried away, imagining that the victors were still on their track. But by the Emperor's orders Belisarius did not pursue them.

The fortunes of the Hunnic troops who were sent against the Chersonese were not happier. Germanus, a native of Prima Justiniana, had been appointed some time previously commandant in that peninsula, and he now proved himself a capable officer. As the Huns could make no breach in the great wall,

p307
which barred the approach to the peninsula and was skilfully defended by the dispositions of Germanus, they resorted to the expedient of manufacturing boats of reeds fastened together in sheaves; each boat was large enough to hold four men; one hundred and fifty were constructed, and six hundred men embarked secretly in the bay of Aenus (near the mouth of the Hebrus), in order to land on the south-western coast of the Chersonese. Germanus learned the news of their enterprise with delight, and immediately manned twenty galleys with armed men. The fleet of reed-built boats was easily annihilated, not a single barbarian escaping. This success was followed up by an excursion of the Romans from the wall against the army of the dispirited besiegers, who then abandoned their enterprise and joined Zabergan, now retreating after the defeat at Chettus.

The other division of the Huns, which had been sent in the direction of Greece, also returned without achieving any signal success. They had not penetrated farther than Thermopylae, where the garrison of the fortress prevented their advance.

Thus, although Thrace, Macedonia, and Thessaly suffered terribly from this invasion, Zabergan was frustrated in all three points of attack, by the ability of Belisarius, Germanus, and the garrison of Thermopylae. Justinian redeemed the captives for a considerable sum of money, and the Kotrigurs retreated beyond the Danube. But the wily Emperor laid a trap for their destruction. He despatched a characteristic letter to Sandichl, the king of the Utigurs, whose friendship he still cultivated by periodical presents of money. He informed Sandichl that the Kotrigurs had invaded Thrace and carried off all the gold that was destined to enrich the treasury of the Utigurs. ``It would have been easy for us,'' ran the Imperial letter, ``to have destroyed them utterly, or at least to have sent them empty away. But we did neither one thing nor the other, because we wished to test your sentiments. For if you are really valiant and wise, and not disposed to tolerate the appropriation by others of what belongs to you, you are not losers; for you have nothing to do but punish the enemy and receive from them your money at the sword's point, as though we had sent it to you by their hands.'' The Emperor further threatened that,

p308
if Sandichl proved himself craven enough to let the insult pass, he would transfer his amity to the Kotrigurs. The letter had the desired effect. The Utigurs were stirred up against their neighbours, and ceaseless hostilities wasted the strength of the two peoples.\footnote{Agathias, V.25 ;

John Ant., \emph{fr.} 217 (\emph{F. H. G.} IV) ; Menander, \emph{fr.} 1, \emph{De leg. Rom.} Another invasion of Huns is recorded in A.D. 562 (Theophanes, A.M. 6054); Anastasiopolis was captured.

}

The historian who recorded the expedition of Zabergan concludes his story by remarking that these two Hunnic peoples were soon so weakened by this continual warfare that though they were not wholly extinguished they were incorporated in larger empires and lost their individualities and even their names.\texthebrew{⁠}\footnote{Agathias, V.25 ; while the Kotrigurs were subjugated by the Avars, the Utigurs were conquered by the western Turks about 576 (cp. Menander, \emph{fr.} 14, \emph{De leg. Rom.} p208).

} The power which threatened them was already at the gate of Europe at the time of Zabergan's invasion.\footnote{Below,
§ 6 .

}

Unable to spare military forces adequate to protect the Balkan provinces against the inroads of the barbarians, Justinian endeavoured to mitigate the evil by an elaborate system of fortresses, which must have cost his treasury large sums. In Thrace, Macedonia, Dardania, Epirus, and Greece, new forts were built, old forts were restored and improved, about six hundred in all.\footnote{The source is Procopius, \emph{Aed.} IV \emph{(Thayer Note: in my Web transcription, in Parts B and C )}, where full lists of the forts (of which few can be identified) will be found. It seems probable that the fortifications were carried out on a general plan, after A.D. 540 (we know that Cassandrea and the Chersonese were fortified after that date, and Topirus after 549). The invasion of that year had displayed the deficiencies of the existing fortifications. Most of the old military forts had only one tower (they were called \textgreek{μονοπύργια}). Justinian's seem to have been larger and had several towers.

\emph{Ib.} 5.4 . On the general principles of the defensive fortifications of the provinces, as illustrated by the remains in Africa, see above,
p148 .

}

Thrace had always been defended by a line of fortresses on both sides of the Danube. They were now renovated and their number was increased. Behind them, in the provinces of Lower Moesia and Scythia, there were about fifty walled towns and castles. South of the Balkan range, the regions of Mount Rhodope and the Thracian plain were protected by 112 fortresses. The defences of Hadrianople and Philippopolis, Plotinopolis and

p309
Beroea (Stara Zagora) were restored, and Topirus, under Mount Rhodope, which the Sclavenes had taken by assault, was carefully fortified. Trajanopolis and Maximianopolis, in the same region, were secured by new walls, and the populous village of Ballurus was converted into a fortified town. On the Aegean coast, the walls of Aenus were raised in height, and Anastasiopolis strengthened by a new sea-wall. The wall, which hedged in the Thracian Chersonese but had proved too weak to keep out the Bulgarians, was demolished, and a new and stronger defence was built, which proved effective against the Kotrigurs. Sestos was made impregnable, and a high tower was erected at Elaeûs. On the Propontis, Justinian built a strong city at Rhaedestus and restored Heraclea. Finally, he repaired and strengthened the Long Wall of Anastasius.

The provinces belonging to the Prefecture of Illyricum were strewn with fortresses proportionate in number to the greater dimensions of the territory. The stations on the Danube from Singidunum to Novae were set in order. In Dardania, the Emperor's native province, eight new castles were built and sixty-one restored. Here he was concerned not only to provide for the defence of the province but to make it worthy of his own greatness by imposing and well-furnished cities. Scupi, near the village where he was born, began a new era in its history under the name of Justiniana Prima, though the old name refused to be displaced, and the town is now Üsküb. It was raised to high dignity as the ecclesiastical metropolis of Illyricum; the number of its churches, its municipal offices, the size of its porticoes, the beauty of its market-places impressed the visitor. Ulpiana (Lipljan), too, was embellished, and became Justiniana Secunda,\texthebrew{⁠}\footnote{For the identification of the two Justinianas see Evans, \emph{Antiquarian Researches in Illyricum}, Part III.62 \emph{sq.}; Part IV.134 \emph{sqq.} Scupi had been ruined by an earthquake in 518 (Marcellinus, \emph{s. a.}).

} and near it the Emperor founded a new town called Justinopolis in honour of his uncle. In the centre of the peninsula the walls of Sardica and Naissus were rebuilt.

The inhabitants of Macedonia were protected by forty-six forts and towns. Cassandrea, which had failed to withstand the Sclavenes, was made impregnable. In the two provinces of Epirus, forty-five new forts were built and fifty rehabilitated. In Thessaly, the decayed walls of Thebes, Pharsalus, Demetrias,

p310
Larissa, and Diocletianopolis on Lake Castoria, and other towns\texthebrew{⁠}\footnote{Metropolis (near the modern Karditsa), Gomphi, Tricca (now Trikkala), Caesarea, Centauropolis.

} were restored. The defences of Thermopylae were renewed and improved, and the historic barrier which had hitherto been guarded by the local farmers was entrusted to 2000 soldiers.\texthebrew{⁠}\footnote{Procopius, \emph{Aed.} IV.2.14 ;

\emph{H. A.} 26.33 , where it is said that, on the pretext of paying the garrison, the municipal rates of all the cities of Greece were appropriated to the treasury; this change is attributed to the logothete Alexander Psalidios.

} The Isthmus of Corinth was fortified anew,\texthebrew{⁠}\footnote{The fortress of Megara had been restored apparently under Anastasius, CIG IV.8622. Cp. Hertzberg, \emph{Gesch. Griechenlands}, III.469.

} and the walls of Athens and the Boeotian towns, which were dilapidated by age or earthquakes, were restored.

This immense work of defence did not avail to keep the barbarians out of the land. Writing in A.D. 550 Procopius sums up the situation: ``Illyricum and Thrace, from the Ionian Sea to the suburbs of Byzantium, were overrun almost every year since Justinian's accession to the throne by Huns, Sclavenes, and Antae, who dealt atrociously with the inhabitants. In every invasion I suppose that about 200,000 Roman subjects were killed or enslaved; the whole land became a sort of Scythian desert.''\texthebrew{⁠}\footnote{\emph{H. A.} 18.20 .

} The historian's supposition doubtless exaggerates the truth considerably, and he would have been more instructive if he had told us how far the improved fortifications mitigated the evils of the invasions. It is clear, however, that it was a great advantage for the inhabitants to have more numerous and safer refuges when the barbarians approached; and we may guess that if statistics had been kept they would have shown a decrease in the number of the victims.

No cities in the Roman Empire deserve greater credit for preserving Greek civilisation in barbarous surroundings than Cherson and Bosporus in the lonely Cimmerian peninsula. They were the great centres for the trade between the Mediterranean and the basins of the Volga and the Don. They were exposed to the attacks of the Huns both from the north and from the east, and the subsidies which Justinian paid to the Utigurs must have been chiefly designed to purchase immunity for these

p311
outposts of the Empire. They had always stood outside the provincial system, and the political position of Bosporus seems to have been more independent of the central power than that of Cherson, where the Emperors maintained a company of artillery ( ballistarii ).\texthebrew{⁠}\footnote{Cp. Kulakovski, \emph{Proshloe Tavridi}, c. viii. For the geography of the peninsula see E. H. Minns, \emph{Scythians and Greeks} (1913), where a sketch of the history of Cherson will be found.

} In the fifth century the bond between Bosporus and Constantinople was broken, a change which was doubtless a result of the Hunnic invasion, and during this period it was probably tributary to the neighbouring Huns. But in the reign of Justin the men of Bosporus sought the protection of the Empire and were restored to its fold.\texthebrew{⁠}\footnote{Procopius, \emph{B. P.} I.12.

It is a disputed question whether the inscription of the Caesar Tiberius Julius Diptunes, ``friend of Caesar, friend of the Romans,'' belongs to A.D. 522 (as Kulakovski maintains, \emph{ib.} 59). A count and an eparch are mentioned, raising the presumption that the stone was inscribed when Bosporus was subject to the Empire. The inscription is published in Latyshev, \emph{Sbornik}, II.39.

} They soon found that they would have to pay for the privilege. They were not indeed asked to pay the ordinary provincial taxes, but Cherson and Bosporus were required to contribute to the maintenance of a merchant fleet which we may suppose was intended exclusively for use in the Euxine waters. This ship-money was also imposed on Lazica, when that land was annexed to the Empire.\footnote{See the \emph{Novel} of Tiberius of A.D. 575 (= Justinian, \emph{Nov.} 163) \textgreek{\texthebrew{ἐ}π\texthebrew{ὶ} το\texthebrew{ῖ}ς λεγομένοις τ\texthebrew{ῶ}ν ε\texthebrew{ἰ}δ\texthebrew{ῶ}ν πλωίμοις γινομένοις \texthebrew{ἐ}πί τε τ\texthebrew{ῆ}ς Λάζων χώρας κα\texthebrew{ὶ} Βοσπόρου κα\texthebrew{ὶ} Χερσονήσου.}}

The Crimean Huns occupied the territory between the two cities. It is not clear whether they stood in the definite relation of federates to the Empire; but in A.D. 528 their king Grod was induced to come to Constantinople, where he was baptized, the Emperor acting as sponsor, and he undertook to defend Roman interests in the Crimea.\texthebrew{⁠}\footnote{John Mal. XVIII.431. John Eph. \emph{H. E.} Part II p475, where the king is called Gordian. Some time previously, Probus had been sent to Grod to induce him to send help to the Iberians against Persia. Procopius, \emph{B. P.} I.12.6.

} At the same time Justinian sent a garrison of soldiers to Bosporus under the command of a tribune. Grod, on returning home, took the images of his heathen gods \texthebrew{—} they were made of silver and electrum, \texthebrew{—} and melted them down. But the priests and the people were enraged by this impiety, and led by his brother, Mugel, they slew Grod, made Mugel king, and killed the garrison of Bosporus. The Emperor then sent considerable forces which intimidated the

p312
Huns and tranquillity was restored.\texthebrew{⁠}\footnote{John Mal. \emph{ib.}; cp. Theophanes, A.M. 6020. It would be interesting to know more about this expedition. According to John Mal., transports of soldiers were sent by sea, and a large force under Baduarius \emph{by land}, starting from Odessus. The march of a Roman army by the northern coast of the Euxine, through the territory of the Bulgarians and Kotrigurs, was a unique event. John of Ephesus (\emph{ib.}) says that Mugel and his followers fled to another country in fear of the Emperor.

} Bosporus was then strongly fortified, the walls of Cherson, which were old and weak, were rebuilt, and two new forts were erected in the south of the peninsula.\footnote{\textgreek{Τ\texthebrew{ὸἈ}λούστου κα\texthebrew{ὶ} τ\texthebrew{ὸ ἐ}ν Γορζουβίταις.}Procopius, \emph{Aed.} III.7.11 . In this passage Procopius clearly alludes to the events of 528. The walls of Cherson had been strengthened in the reign of Zeno, CIG IV.8621. Procopius (\emph{B. G.} IV.5.28) curiously describes Phanagoria as near Cherson and still subject to the Romans.

}

In the north of the Crimea there was a small Gothic settlement, apparently a remnant of the Ostrogothic kingdom which in the fourth century extended along the north coast of the Euxine. These Goths are described as few in number, but good soldiers, skilful in agriculture, and a people of hospitable habits. They were under the protection of the Empire and were ready, when the Emperor summoned them, to fight against his foes. Their chief place was Dory on the coast; they would have no walled towns or forts in their land, but Justinian built long walls at the points where it was most exposed to an invader.\footnote{They were Christians, and perhaps they had a bishop at as early a date as the Council of Nicaea (Mansi, II.696 provinciae Gothiae. Theophilus Gothiae metropolis). If this is so, they must have been distinct from the Ostrogoths of Hermanric. The chief source for this people is

Procopius, \emph{Aed.} III.7.13\texthebrew{‑}17 . He describes them as \textgreek{\texthebrew{Ῥ}ωμαίων \texthebrew{ἔ}νσπονδοι}. When he says that they numbered about 3000, he perhaps means the men of military age. For these Goths and the Tetraxites see Loewe, \emph{Die Reste der Germanen am Schwarzen Meere}, 22 \emph{sqq.} They are confused by Tomaschek, \emph{Die Goten in Taurien}, 12.

}

From these genuine Goths of the Crimea we must carefully distinguish another people, who were also described as Goths but perhaps erroneously. These were the Tetraxites (a name of mysterious origin) who lived in the peninsula of Taman over against Bosporus.\texthebrew{⁠}\footnote{A name for Taman, T'mutarakan, which occurs in old Russian sources and is evidently of Arabic or Turkish origin, supplied Vasil'evski with an ingenious interpretation of Tetraxie, which is approved by Loewe (\emph{op. cit.} 33\texthebrew{‑}34). He explains T'mutarakan as derived from \textgreek{τ\texthebrew{ὰ} Μάτραχα}*, which he identifies with \textgreek{τ\texthebrew{ὸ} Ταμάταρχα} (in Constantine Porph. \emph{De adm. imp.} c42), from which he gets \textgreek{Τμετραξ\texthebrew{ῖ}ται}* as a name of the inhabitants, and hence \textgreek{Τετραξ\texthebrew{ῖ}ται} (the corruption being influenced by \textgreek{τετραξός}). Loewe thinks that the Tetraxites were Heruls.

Thayer's Note: The sources are various manuscripts of the so\texthebrew{‑}called Hypatian Codex, a compendium of three early Kievan chronicles (not ``Russian'', as Hodgkin would have it, since Russia didn't exist at the time). For full details, see

Doroshenko, \emph{Survey of Ukrainian Historiography}, pp21\texthebrew{‑}31 ; also, a good summary article

at the Encyclopedia of Ukraine.

} They too were a small people, and their fate depended on the goodwill of the Utigurs,\texthebrew{⁠}\footnote{They supplied 2000 soldiers to the expedition of the Utigurs against the Kotrigurs in 551. Procopius, \emph{B. G.} IV.18.22.

} whose kingdom

p313
stretched from the Don as far south as the Hypanis. They engaged, however, in secret diplomacy with Justinian. Their bishop had died, and ( A.D. 548) they sent envoys to Constantinople to ask the Emperor to provide a successor. This was the ostensible object of the embassy, and nothing else was mentioned in the official audience, for they were afraid of the Utigurs; but they had a secret interview with the Emperor, at which they gave him useful information for the purpose of stirring up strife among the Huns.\footnote{Procopius, \emph{B. G.} IV.4.9\texthebrew{‑}13. Of their religion he says: ``I cannot say whether they were once Arians, like the other Gothic peoples, or held some other creed, for they do not know themselves, but now they adhere with simple sincerity to the (orthodox) religion'' (\textgreek{\texthebrew{ἀ}φελεί\texthebrew{ᾳ} κα\texthebrew{ὶ ἀ}πραγμοσύν\texthebrew{ῃ} πολλ\texthebrew{ῇ} τιμ\texthebrew{ῶ}σι τ\texthebrew{ὴ}ν δόξαν}). I cannot find in Procopius (\emph{ib.} IV.5)

the statement, ascribed to him by Loewe, that the Tetraxites lived in the Crimea before they settled in the Taman peninsula. It is to be noticed that the old Greek town of Phanagoria, opposite to Bosporus, was in their hands, and was probably the headquarters of their ecclesiastics. An inscription found at Taman, and doubtless brought there from Phanagoria, relates to the restoration of a church under the auspices of Justinian. It is dated to an eleventh indiction, which gives three possible dates, 533, 548, and 563. This stone is discussed by Latyshev, \emph{Viz. Vrem.} I.657 \emph{sqq.}, who decided for 533; by Kulakovski, \emph{ib.} 2.189 \emph{sqq.}, who argued for 548 (but he is now doubtful, \emph{Pamiatnik}, p10, n1); and by Semenov, \emph{B. Z.} VI.387 \emph{sqq.}, who denies that the year can be fixed.

}

To the south of the Utigurs, in the inland regions north of the Caucasian range, were the lands of the Alans, traditionally friends of the Romans, and further east the Sabirs, whose relations to the Empire have come before us in connexion with the Persian wars. On the coast south of the Hypanis, the Zichs, whose king used in old days to be nominated by the Emperor, were accounted of small importance.\texthebrew{⁠}\footnote{In his account of these regions in \emph{B. G.} IV.4, Procopius places the Saginae apparently to the north of the Zichs, though one might infer from another passage, \emph{ib.} 2.16, that they were nearer to Colchis. He also mentions the Bruchoi as dwelling between the Alans and Abasgians. The Sunitae were also neighbours of the Alans, \emph{B. P.} I.15.1. See also \emph{B. P.} II.29.15. It is difficult to identify all the names of the tribes enumerated as living north of Abasgia in the table of peoples in Zacharias Myt. XII.6, p328 (the Kotrigurs appear as Khorthrigor).

} But their southern neighbours, the Abasgians and the Apsilians, came, as we have already seen, within the sphere of political intrigue and military operations by which Rome and Persia fought for the control of Colchis. On the Abasgian coast the Romans had two fortresses, Sebastopolis (formerly called Dioscurias) and Pityus. On hearing that Chosroes intended to send an army to seize these places, Justinian ordered the garrisons to demolish the fortifications, burn the houses, and withdraw. But he afterwards rebuilt

p314
Sebastopolis on a scale worthy of his reputation as a great builder.\texthebrew{⁠}\footnote{Procopius, \emph{B. G.} IV.4\texthebrew{‑}6; and

\emph{Aed.} III.7.8\texthebrew{‑}9 .

} The fact that he thought it worth while to maintain this outpost shows how considerable were the political and commercial interests of the Empire in this region.

One of the disadvantages of the system of subsidising the barbarians on the frontiers or endowing them with territory was that fresh and formidable enemies were lured to the Roman borders from remote wilds and wastes by the hope of similar benefits. Towards the end of Justinian's reign, a new people of Hunnic race appeared on the frontier of Europe, north of the Caspian, and immediately fixed their covetous desires on the Empire, whose wealth and resources were probably exaggerated far beyond the truth among the barbarian tribes. They called themselves Avars, though it is alleged that they had usurped the name of another people better than themselves;\texthebrew{⁠}\footnote{Cp. Theophylactus Simocatta, \emph{Hist.} VII.7; Bury, App. 5 to Gibbon, vol. V.

} but they were destined to play a part on the European scene similar, if on a smaller scale, to that which had been played by the Huns.

Their westward migration was undoubtedly due to the revolution in Central Asia, which, about the middle of the sixth century, overthrew the power of the Zhu\texthebrew{‑}zhu\texthebrew{⁠}\footnote{See above,
chap. III § 3. The Zhu\texthebrew{‑}zhu are supposed to have been the true and original ``Avars.''

} and set in their place the Turks, who had been their despised vassals. Tu-men was the name of the leader who rose against his masters and founded the empire of the Turks. His successor, Mo-kan ( A.D. 553\texthebrew{‑}572), overthrew the kingdom of the Ephthalites and organised the vast Turkish empire which extended from China to the Caspian and southwards to the borders of Persia, dividing it into two khanates, of which the western was subordinate to the eastern.\footnote{See Bury, Appendix 17 to Gibbon, vol. IV. The scanty information supplied by Greek sources about the early Turkish Empire must be supplemented by Chinese records (cp. E. H. Parker's article in \emph{E. H. R.} XI.431 \emph{sqq.} , and \emph{A Thousand Years of the Tartars}, 1896; Marquart, \emph{Historische Glossen zu den alttürkischen Inschriften}, in \emph{Wiener Zeitschrift f. die Kunde des Morgenlandes}, VII.157 \emph{sqq.}, 1898). The later history of the Turks and their institutions (seventh and eighth centuries) have been illustrated by the Turkish inscriptions discovered in Eastern Mongolia (Thomsen, \emph{Inscriptions de l'Orkhondéchiffrées}, 1894; Radloff, \emph{Die alt-türkischen Inschriften der Mongolei}, 1895 (Neue Folge), 1897 (Zweite Folge),

(p315)
1899; Marquart, \emph{op. cit.}, and \emph{Die Chronologie der alt-türk. Inschriften}, 1898. Mokan may almost certainly be identified with Silzibulos in Menander.

}

p315
In A.D. 558 Justin, the son of Germanus, who was commanding the forces in Colchis, received a message from Sarus, king of the Alans, to the effect that Candich, king of the Avars, desired to enter into communications with the Emperor. Justin informed his country, who signified his readiness to receive an embassy. The envoys of Candich arrived at Constantinople. They vaunted the invincibility of the Avars and made large demands \texthebrew{—} land, gifts, annual subsidies. Justinian, having consulted the Imperial Council, gave them handsome gifts, couches, clothes, and gold chains, and sent an ambassador to Candich, who was informed that the Emperor might take his requests into consideration, if the Avars proved their worth by subduing his enemies. The Avars immediately made war upon the Sabirs and destroyed them, and fought with success against the Utigurs. Having cleared the way, they advanced through Kotrigur territory to the regions of the Bug and Seret, subjugated the Antae, and in A.D. 562 they made a great raid through Central Europe, appeared on the Elbe, and threatened the eastern marches of the Frank kingdom of Austrasia. But all these expeditions seem to have been carried out from their headquarters, somewhere between the Caspian and the Black Sea.\footnote{Menander, \emph{De leg. gent. frs.} 1\texthebrew{‑}13; John Eph. \emph{H. E.} Part III VI.24;

Gregory of Tours, \emph{Hist. Fr.} IV.23 . Cp. also Pseudo-Nestor, \emph{Chron.} c8 (p6, ed. Miklovich), on the subjugation of the Slavonic Dudlebians of Volkynia by the Avars (Obre), which Šafarik (\emph{Slaw. Altertümer}, II.60) and Marquart (\emph{Osteur. Streifzüge}, 147) refer to the invasion of 562. The comments of Marquart (\emph{Chronologie}, 78 \emph{sqq.}) on Menander, \emph{fr.} 3, based on the theory that the Antae at this time exercised overlord­ship over the Bulgarians, are very hazardous. John Eph. (\emph{ib.} cp. III.25) says the Avars were so called from wearing their hair long.

}

In the same year Baian, who had succeeded Candich and was afterwards to prove himself the Attila of the Avars, sent an embassy to Constantinople, demanding land in a Roman province. The ambassadors travelled by Colchis, and Justin, who arranged for their journey to the capital, gained the confidence of one of the party and was secretly informed by him that treachery was intended. He therefore advised Justinian to detain the barbarians as long as possible, since the Avars would not carry out their purpose of crossing the Danube till the envoys had departed. The Emperor acted on this advice, and Bonus, the Quaestor of Moesia and Scythia, was instructed

p316
to see to the defences of the river.\texthebrew{⁠}\footnote{Menander, \emph{ib.} \emph{fr.} 4. Bonus is described as \textgreek{πρωτοστάτης το\texthebrew{ῦ} θητικο\texthebrew{ῦ} κα\texthebrew{ὶ} ο\texthebrew{ἰ}κετικο\texthebrew{ῦ}}. He had been appointed quaestor exercitus of the Five Provinces (see below,
p340 ) in 536 (\emph{Nov.} 41).

} The policy succeeded, though we do not know exactly why; the Avars did not attempt to invade the Empire; and the envoys were at last dismissed. They received the usual gifts, which they employed in buying clothes and arms before they left Constantinople. The arms must have been furnished by the Imperial factories, and the Emperor apparently did not consider it politic to refuse to sell them. But he sent secret instructions to Justin to take the arms away from the barbarians when they arrived in Colchis. Justin obeyed, and this act is said to have been the beginning of enmity between the Romans and the Avars. Justinian did not live to see the sequel. But he had not been long in his grave before Baian led his people to the Danube, where they secured a permanent abode and were a scourge to the Balkan provinces for nearly sixty years.

In the efforts of the Imperial government to extend its influence in the Red Sea sphere, the interests of trade were the principal consideration. Before we examine the fragmentary and obscure record of Roman intervention in the affairs of Ethiopia and Southern Arabia, we may survey the commercial activities of the Empire abroad.

The trade of the Mediterranean was almost entirely in the hands of Syrians and Greeks. In Rome and Naples and Carthage, and not only in Marseilles and Bordeaux, but also in the chief inland cities of Gaul, we find settlements of oriental merchants.\texthebrew{⁠}\footnote{Cp. the article of Bréhier, \emph{Les Colonies d'orientaux en occident}, \emph{B. Z.} XII.1 \emph{sqq.} (1903). On commerce in general in the sixth century see Heyd, \emph{Hist. du commerce}, I. pp1\texthebrew{‑}24.

} Their ships conveyed to the west garments of silk and wrought linen from the factories of Tyre and Berytus, purple from Caesarea and Neapolis, pistachios from Damascus, the strong wines of Gaza and Ascalon, papyrus from Egypt, furs from Cappadocia.\footnote{See the \emph{Expositio totius mundi}, which was translated from the Greek soon after 345. Sidonius (\emph{Carm.} 17.15) speaks of the vina Gazetica, cp. Cassiodorus, \emph{Var.} 12.5; Greg. of Tours, \emph{Hist. Fr.} VII.29. Among the exports from Spain, the \emph{Expositio} enumerates oil, bacon, cloth, and mules; from Africa, oil, cattle, and clothing. On the multitude of Syrian merchants in Gaul cp. Salvian, \emph{De gub. Dei}, IV.14.

}

p317
There was a large demand for embroidered stuffs, especially for ecclesiastical use, cloths for altars, curtains for churches.\texthebrew{⁠}\footnote{Cp. the document of A.D. 471 given by Duchesne, \emph{Lib. pont.} I.cxlvii. Wealthy private persons also obtained from the East the artistic tapestries which they needed for the adornment of their houses, like that embroidered with hunting-scenes which is described by Sidonius, \emph{Epp.} IX.13. On figured textiles see Dalton, \emph{Byz. Art}, 577 \emph{sqq.}

} But the great centre to which the ships from all quarters converged was the Imperial capital, as the richest and most populous city of the world.\texthebrew{⁠}\footnote{Paulus Silentiarius (\emph{S. Sophia}, 232) makes Constantinople say:

\textgreek{ε\texthebrew{ἰ}ς \texthebrew{ἐ}μ\texthebrew{ὲ} φορτ\texthebrew{ὶ}ς \texthebrew{ἄ}πασα φερέσβιον \texthebrew{ἐ}λπιδα τείνει κύκλιον ε\texthebrew{ἰ}σορόωσα δρόμον διδυμάονος \texthebrew{ἄ}ρκτου.}} It seems probable that most of the imports which the Empire received from the countries bordering on the Euxine came directly across its waters to Constantinople and were distributed from there: the skins which the Huns exchanged at Cherson for stuffs and jewels,\texthebrew{⁠}\footnote{Jordanes, \emph{Get.} 37 ; cp. Procopius, \emph{B. G.} IV.20, 17.

} and the slaves, skins, corn, salt, wine, which were obtained from Lazica.

For the Empire trade with the East had always been mainly a trade in imports. The East supplied the Mediterranean peoples with many products which they could not do without, while they had themselves less produce to offer that was greatly desired by the orientals. There had, from of old, been a certain market in China for glass, enamelled work, and fine stuffs from Syria;\texthebrew{⁠}\footnote{See

Hirth, \emph{China and the Roman Orient} . M. Khvostov's Russian work, \emph{Istoriia vostochnoi torgovli Grekorimskago Egipta}, is indispensable for eastern trade down to the end of the third century.

} but whatever exports found their way thither or to India and Arabia were far from being a set-off to the supplies of silk, not to speak of spices, precious stones, and other things which the East sent to the West. The balance of trade was, therefore, decidedly against the Empire, and there was a constant drain of gold to the East.\footnote{We have no figures bearing on the amount of the trade except those furnished by Pliny in the first century. He says that India received annually from the Empire 55 million sesterces (\emph{c.} £600,000), and that China, India, and Arabia together took at least 100 million sesterces (\emph{c.} £1,000,000); \emph{Nat. Hist.}

VI.23, § 101 , and

XII.18, § 84 . If these sums represented the whole value of the imports, the volume of trade would have been small; but it probably means only the balance of trade \texthebrew{—} the amount of specie which was taken from the Empire; and to know the value of the imports, we should have also to know the amount of the exports which partly paid for them. So

Hirth, \emph{ib.} p227 , and Khvostov, \emph{op. cit.} p410, inclines to this view.

}

Under the early Roman Empire, the trade with India, the Persian Gulf, Arabia, and the eastern coast of Africa had been

p318
in the hands of Roman merchants, who sailed through the Red Sea and the Indian Ocean in their own vessels. Before the end of the third century this direct commerce seems to have ceased almost entirely. The trade between the Mediterranean and the East passed into the hands of intermediaries, the Persians, the Abyssinians, and the Himyarites of Yemen. This change may have been due to the anarchical conditions of the Empire, which followed on the death of Alexander Severus and were unfavourable to commercial enterprise. The energy of Persian merchants, under the orderly rule of the Sassanids, secured a monopoly of the silk trade, and the products of India were conveyed by Abyssinian traders to their own market at Adulis, or even to the Roman ports on the isthmus, Clysma (Suez)\texthebrew{⁠}\footnote{Clysma is Qulzum, \texthebrew{•} a quarter of a mile north of Suez. This is shown by its description in the account of her pilgrimage to Sinai by the Abbess Aetheria (of South Gaul), who travelled in the early years of Justinian's reign (between 533 and 540), according to K. Meister in \emph{Rhein. Mus.}, N. F., LXIV.337 \emph{sqq.} (1909), and Mommsen, \emph{Hist. Sch.} III.610 \emph{sqq.}, but according to others in the last years of the fourth century (see E. Weigand, \emph{B. Z.} XX.1 \emph{sqq.}, 1912). She saw many large ships there, and mentions that there was a resident agens in rebus, known as a logothete, who used to visit India every year by order of the Emperor. India, of course, must mean Ethiopia. The most important duty of this office was to see that the regulations as to exports were observed (cp.

\emph{C. Th.} VI.29.8 ).

} and Aila. The Red Sea trade itself seems to have been gradually abandoned, as time went on, to the Abyssinians and Himyarites, who grew more power ­ful and important as their commercial profits increased. The Abyssinians \texthebrew{—} as we may conveniently call the Ethiopians of the kingdom of Axum, from which modern Abyssinia descends \texthebrew{—} also profited by the disuse of the Nile as a trade route with East-Central Africa. The products of those regions (slaves, ivory, ebony, gold, gems, ochre, etc.) had come to Egypt by the Nile, as well as by the Red Sea, in the old days when the Ethiopic kingdom of Meroe flourished. Meroe declined in the second century, and in the third its organisation fell to pieces, and the Upper Nile, under the control of the barbarous Nubians and Blemmyes, became impracticable as a road for trade. With the shifting of power from Meroe to Axum, East African commerce passed entirely into the hands of the Abyssinians.\footnote{Khvostov, \emph{op. cit.} 29 \emph{sqq.} For the early trade of Roman Egypt with the East see also Mommsen, \emph{Röm. Gesch.} V chap. XII. It is to be observed that the cessation of direct trade with the East was reflected in the decline of geographical knowledge, illustrated by the misuse of India to designate Ethiopia, which is frequent in Greek and Latin writers from the fourth century.

}

p319
As to the traffic with India, we find much curious information in a remarkable book which was written about the middle of the sixth century, the \emph{Christian Cosmography} of Cosmas.

Cosmas, who is known as Indicopleustes, ``sailor of the Indian Sea,'' was an Egyptian merchant, but when he wrote his book he had probably abandoned his calling and become a monk. The \emph{Cosmography}, which was composed about A.D. 445\texthebrew{‑}450,\texthebrew{⁠}\footnote{Books I\texthebrew{‑}V appeared first, and VI\texthebrew{‑}X were separately added to answer objections and supply additional explanations. The chronological indications are as follows: (1) Timothy, Patriarch of Alexandria, who died in 535, is described as \textgreek{ν\texthebrew{ῦ}ν τετελευτηκώς}, X p315. (2) Theodosius of Alexandria is still alive at Constantinople, \emph{ib.} p314. Theodosius was confined, after his deposition in 536, in the Thracian fort of Derkos (John Eph. \emph{Comm. de B. Or.}, p14), but afterwards lived under Theodora's protection at Constantinople, where he died soon after Justinian's death (\emph{ib.} p159; cp. \emph{H. E.} Part II p248). These indications give the limits 536 and 565. (3) It is stated in VI p232 that a solar and a lunar eclipse, which occurred in the same year, on Mecheri 12 and Mesori 14, had been predicted. Two such eclipses did occur on February 6 and August 1 in 547. Hence Book VI was written after 547. (4) Book II was written (p72) twenty-five years after the Himyarite expedition of Elesboas, which occurred ``at the beginning of the reign of Justin.'' This implies about 544\texthebrew{‑}545 for Books I\texthebrew{‑}V.

} is unfortunately neither a treatise on geography like Strabo's or Ptolemy's, nor a plain account of his travels, but a theological work, designed to explain the true shape of the universe as proved by Scripture, and especially to refute the error of pagan science that the earth is spherical. His theory as to the shape of the world, which is based on the hypothesis that Moses, ``the great cosmographer,'' intended his tabernacle to be a miniature model of the universe, is not devoid of interest as an example of the fantastic speculations to which the interpretation of the Biblical documents as literally inspired inevitably leads.

The earth, according to Cosmas, is a flat rectangle, and its length is double its breadth. The heavens form a second story, welded to the extremities of the earth by four walls. The dry land which we inhabit is surrounded by the ocean, and beyond it is another land where men lived before the Deluge. The firmament is the ceiling between the two stories, and the earth, the lower story, lies at the bottom of our universe, to which it sank when it was created. There is nothing below it. Hence the pagan theory of the antipodes is a delusion. On its western side the earth rises into a great conical mountain, which hides

p320
the sun at night. The sun is not larger than the earth, as the pagans falsely imagine, but much smaller. The revolutions of all the celestial bodies are guided by angel pilots.\footnote{The extension of the work of creation over six days \texthebrew{—} whereas it could have been accomplished by a single fiat \texthebrew{—} is ingeniously explained as due to the Creator's wish to give a series of object-lessons to the angels, III p105 \emph{sqq.}

}

It would be a mistake to suppose that this strange reconstruction of the world, which contemptuously set aside all that Greek science had achieved, represented the current views of orthodox Christians or ever obtained any general credence. It was not indeed original. Cosmas derived his conceptions from hints which had been thrown out by theologians of the Syrian school, especially from Theodore of Mopsuestia.\texthebrew{⁠}\footnote{Cosmas was a Nestorian. Cp. M'Crindle, \emph{Introd.} to his translation, p. ix.

} But for us the value of the work lies in the scraps of information relating to his own travels which the author introduces incidentally, and in the contents of an appendix, which has no relation to his theme, and seems to have been part of another work of Cosmas, and to have been attached to the \emph{Cosmography} by some injudicious editor.\footnote{Books XI and XII, of which the latter is a series of fragments. He had written a general geography which is lost (\emph{Prologue, ad init.}), but it has been suggested that Book XI formed part of it (Winstedt, \emph{Introd.} p5).

}

Cosmas knew the Red Sea well. He visited Ethiopia in the reign of Justin,\texthebrew{⁠}\footnote{II p72.

} and he made at least one voyage to the Persian Gulf.\texthebrew{⁠}\footnote{He says that he had sailed in the Persian, Arabian, and Roman gulfs (the Roman means the Mediterranean), \emph{ib.} p62, where he relates that ``once having sailed to Inner India and crossed a little towards Barbaria \texthebrew{—} where farther on is situated Zingion, as they call the mouth of the Ocean \texthebrew{—} we saw a flight of albatrosses (suspha).'' Inner India is either South Arabia or Abyssinia, though in the same passage Ceylon is said to be in Inner India (if Inner is not an error for Outer). Barbaria is the African coast south of Abyssinia, and Zingion is Zanzibar.

} It is to this voyage that he probably referred when he wrote: ``I sailed along the coast of the island of Dioscorides (Socotra), but did not land, though in Ethiopia I met some of its Greek-speaking inhabitants.''\texthebrew{⁠}\footnote{III p119. There were Christian clergy here who received ordination from Persia. Cosmas also mentions that there were Christian churches in Ceylon, Mala (Malabar), and Calliana (near Bombay) under a bishop ordained in Persia.

} The Persian Gulf probably represents the limit of his eastern travel, for in all that he tells of Ceylon and India we are struck by the absence of any of those personal touches which could not fail to appear in the description of an eye-witness . It was only a rare Roman merchant

p321
who visited the markets of Ceylon.\texthebrew{⁠}\footnote{Cosmas knew one. See below,
p332 .

} The trade between the Red Sea and India was entirely in the hands of the Abyssinians, and the Roman merchants dealt with them.

Ceylon, which the ancients knew as Taprobane,\texthebrew{⁠}\footnote{Cosmas also calls it Sielediva (from which comes the modern Arabic Serendib). He tells us that there were two kings in the island, and there were many temples, on the top of one of which (perhaps the Buddhist temple of Anarajapura) shone a red stone, large as the cone of a pine, which Cosmas calls a hyacinth. This stone was known by report to Marco Polo (III.14), who calls it a ruby. It is supposed by some to be an amethyst. Cosmas, XI.321 \emph{sqq.}

} was the great centre of maritime commerce between the Far East and the West. In its ports congregated Persian, Ethiopian, and Indian merchants. Silk was brought from China to its markets, and continental India sent her products: Malabar, pepper; Calliana, copper; Sindu, musk and castor.\texthebrew{⁠}\footnote{Sindu has been identified with Diul-Sind, at the mouth of the Indus. China in Cosmas is \textgreek{Τζινίστα}. He says that from there and other emporia (probably Further India) Ceylon receives silk, aloes (\textgreek{\texthebrew{ἀ}λοήν}), cloves (\textgreek{καρυόφυλλα}), sandalwood (\textgreek{τζανδάναν}).

} The islanders exported their own products eastward and westward, and they had a merchant service themselves, but the significance of Ceylon was its position as an emporium for merchandise in transit. The Persians had an advantage over the Romans in that they traded directly with the island, and had a commercial colony there,\texthebrew{⁠}\footnote{XI p322.

} while the Roman trade, as we have seen, was carried on through the Ethiopians and intermediaries.

While it is probable that most of the Indian commodities which were consumed in the Empire travelled by this route, the Ethiopian traders did not carry silk. The large supplies of silk which reached the Romans were bought from Persian merchants, and most of it was probably conveyed overland from China to Persia, though part of it may also have come by sea, by way of Ceylon and the Persian Gulf. We do not know by what methods the Persians succeeded in establishing this monopoly and preventing the Abyssinians from trading in silk. It was highly inconvenient to the Empire to depend exclusively on a political rival for a produce of which the consumption was immense, and in time of war the inconvenience was grave. Justinian deemed it a matter of first importance to break the Persian monopoly, and for this purpose, during the first Persian War, he entered into negotiations with the king of Abyssinia.

p322

The kingdom of the Abyssinians or Ethiopians, who were also known as the Axumites, from the name of their capital city Axum, approached Suakim on the north, stretched westwards to the valley of the Nile, and southwards to the Somali coast. Their port of Adulis was reckoned as a journey of fifteen days from Axum where the king resided.\texthebrew{⁠}\footnote{So Nonnosus, who had made the journey

(\emph{F. H. G.} IV.179) . Procopius says 12. We first hear of this Ethiopian kingdom in the \emph{Periplus maris Erythraei}, § 2 \emph{sqq.} (\emph{Geogr. Gr. Minores}, vol. I), \emph{i.e.} in the first century A.D. Its history has been elucidated by Dillmann in his articles in \emph{Abh. Berliner Akad.}, 1878 and 1880.

} Roman merchants frequented Adulis, where there was a great market of the produces of Africa, slaves, spices, papyrus, ivory, and gold from Sasu.\footnote{The Ethiopians gave meat, salt, and iron in exchange for the gold they got from Sasu, Cosmas, II p70. The mention of iron is to be noticed in view of what Procopius says, \emph{B. P.} I.19.25: the Ethiopians have no iron; they cannot buy it from the Romans, for it is expressly forbidden by law.

}

The commercial relations of the Abyssinians with their neighbours across the straits, the Himyarites of Yemen, were naturally close, and from time to time they sought to obtain political control over South-western Arabia. Christian missionaries had been at work in both countries since the reign of Constantius II, when an Arian named Theophilus was appointed bishop of the new churches in Abyssinia, Yemen, and the island of Socotra.\texthebrew{⁠}\footnote{Frumentius, ordained by Athanasius, had been the first bishop in Abyssinia. Athanasius (\emph{Apologia ad Constantinum}, 31) quotes a letter from Constantius to the Ethiopian kings, Aizan and Sazan, asking them to send back Frumentius as a heretic. (An inscription of this Aizan is preserved, CIG III.5128, where he appears as sole king, but his brother Saiazan is mentioned. The mission of Theophilus is recorded by Philostorgius,

III.4\texthebrew{‑}6 .

} He is said to have founded churches at Safar and Aden.\texthebrew{⁠}\footnote{\textgreek{Τάφαρον, \texthebrew{Ἀ}δάνη.}} After this we lose sight of these countries for about a century and a half, during which Christianity probably made little way in either country, and Judaism established itself firmly in Yemen. Then we learn that in the reign of Anastasius a bishop was sent to the Himyarites.\texthebrew{⁠}\footnote{Theodore Lector, II.58 (his source was John Diakrinoumenos).

} We may conjecture that this step was the consequence of a war between the Himyarites and Abyssinians, which is misdated in our records, but apparently belongs to the reign of Zeno or of Anastasius.\footnote{The events connected with the names of Andas and Dimnos are related by John Mal. (XVIII.433 and 429) to the reign of Justinian (A.D. 529). But we know on unimpeachable

(p323)
authority that at that time the names of the kings were respectively Elesboas and Esimiphaios, and Elesboas had been on the throne since the beginning of Justin's reign. We must therefore suppose that John Malalas, misunderstanding his authority, made a chronological mistake, and refer the episode to an earlier period. Cp. Duchesne, \emph{Les Églises séparées}, 316\texthebrew{‑}317; Andöldeke, \emph{Tabari}, 175; Fell, \emph{Die Christenverfolgung in Südarabien}, in \emph{Z. D. M. G.} XXXV.19. The name \textgreek{\texthebrew{Ἄ}νδας} appears as \textgreek{\texthebrew{Ἀ}δάδ} in Theophanes (A.M. 6035), and as Aidug in John Eph. \emph{H. E.} Part II, extract in Assemani, \emph{Bibl. Or.} I p359. It is supposed to correspond to Ela-Amida in Ethiopic chronicles.

}

p323
Dimnos, king of the Himyarites, who was probably a convert to Judaism, massacred some Greek merchants, as a measure of reprisal for alleged ill-treatment of Jews in the Roman Empire. Thereupon, presumably at direct instigation from Constantinople, the Abyssinian king Andas invaded Yemen, put Dimnos to death, and doubtless left a viceroy in the country with an Ethiopian garrison. Andas had vowed that, if he were victorious, he would embrace Christianity. He fulfilled his vow, and the Emperor sent him a bishop from Alexandria. Andas was succeeded by Tazena, whose inscriptions describe him as ``King of Axum and Homer and Reidan and Saba and Salhen.''\texthebrew{⁠}\footnote{Homer is Himyar; Reidan has been explained as = Safar, and Salhen as the fortress of Ma'rib (Fell, \emph{ib.} 27); Saba (Sheba) is familiar. For the inscriptions of Tazena see Dillmann, \emph{Z. D. M. G.} VII.357 \emph{sqq.} Huart (\emph{Hist. des Arabes}, p53) suggests that the Ethiopians had no ships and that the Romans must have supplied them with transports for their expeditions to Yemen.

} He also was converted from paganism, and his son Elesboas, who was on the throne at the beginning of Justin's reign, was probably brought up a Christian.\footnote{Ela Atzbeha is the Ethiopic name (Nöldeke, \emph{ib.} 188). Of the Greeks, Cosmas gets nearest to it with his \textgreek{\texthebrew{Ἐ}λλατζβάς}, II p72; John Mal. has \textgreek{\texthebrew{Ἐ}λεσβόας}, \emph{Acta mart. Arethae}, p721 \textgreek{\texthebrew{Ἐ}λεσβάς}, Procopius \textgreek{\texthebrew{Ἐ}λλησθεα\texthebrew{ῖ}ος}. On his coins he was also called Chaleb (see Schlumberger in \emph{Rev. numism.} 1886, pl. XIX \textgreek{Χαλ\texthebrew{ὴ}β βασιλε\texthebrew{ὺ}ς υ\texthebrew{ἱὸ}ς Θεζενα} º ), and this, his ``throne-name,'' appears in the Ethiopic version of the \emph{Acta mart. Arethae.} Cp. Fell, \emph{ib.} 17 \emph{sqq.}

}

In the meantime a Himyarite leader, Dhu Novas, of Jewish faith, succeeded in over­powering the Ethiopian garrison, proclaimed himself king, and proceeded to persecute the Christians.\texthebrew{⁠}\footnote{See the consolatory letter written by Jacob of Sarug (who died November 29, 521) to the Himyarite Christians, edited by Guidi (see
Bibliogr. I.2, B ).

} It is not quite certain whether Elesboas immediately sent an army to re-establish his authority ( A.D. 519\texthebrew{‑}520),\texthebrew{⁠}\footnote{Guidi (\emph{op. cit.} pp476, 479) argues for two expeditions against Dhu Novas, the earlier in 519. There is a good deal to be said for this. Cosmas witnessed the preparations of Ela Atzbeha ``at the beginning of the reign of Justin'' (\emph{loc. cit.}), and, as Justin reigned only till 527, it would have been a strange misuse of words to speak of 524 as the beginning of his reign. See above,
p319, n1 . Cp. Mordtmann, \emph{Z. D. M. G.} XXXV.698.

} but if he did so, Dhu Novas recovered his power within the next two years

p324
and began systematically to exterminate the Christian communities of southern Arabia, if they refused to renounce their errors and embrace Judaism. Having killed all the Ethiopians in the land, he marched with a large army against the fortified town of Nejran, which was the headquarters of the Christians ( A.D. 523). The siege was long, but, when the king promised that he would spare all the inhabitants, the place capitulated. Dhu Novas, however, had no intention of keeping faith, and when the Christians refused to apostatise, he massacred them to the number of 280, among whom the most conspicuous was Harith, the emir of the tribe of Harith ibn\texthebrew{‑}Kaab . After having performed this service to the Jewish faith, Dhu Novas despatched envoys to Al\texthebrew{‑}Mundir of Hira, bearing a letter in which he described his exploits, boasted that he had not left a Christian in his land, and urged the Saracen emir to do likewise. When the envoys arrived at Al\texthebrew{‑}Mundhir's camp at Ramla (January 20, A.D. 524), Simeon Beth Arsham, the head of the Monophysites of the Persian empire, happened to be there, having come on the part of the Emperor Justin to negotiate peace with the Saracens. Horrified by the news, Simeon immediately transmitted it to Simeon, abbot of Gabula, asking him to arrange that the Monophysites of Antioch, Tarsus, and other cities should be informed of what had happened.\footnote{For these events the Syriac letter of Simeon, which has been edited by Guidi, and is generally recognised as genuine, is the most authentic source. Simeon, who was accompanied by Mar Abram and Sergius, bishop of Rosapha, had obtained further information as to the massacre when he returned to Hira from Al\texthebrew{‑}Mundhir's camp. It has been conjectured by Duchesne (\emph{op. cit.} p325) that Sergius was the author of the \emph{Martyrium Arethae et sociorum} which has come down in Greek. John Psaltes, in 524\texthebrew{‑}525, composed a Greek hymn on the martyrs, which was immediately translated into Syriac by Paul, bishop of Edessa (died October 30, 526), and his version is preserved (\emph{Z. D. M. G.} XXXI.400 \emph{sqq.}). He speaks not of 280, but of more than 200 martyrs. A verse in the Koran (Sura 85) is said to refer to the massacre: ``Cursed were the contrivers of the pit, of fire supplied with fuel; when they sat round the same, and were witnesses of what they did against the true believers.''

}

It is possible that Justin and the Patriarch of Alexandria\texthebrew{⁠}\footnote{So the Greek version of \emph{Mart. Arethae}, where a letter from Justin (doubtless an invention) is given. But as the Armenian version contains nothing of these negotiations, we have no guarantee that they were mentioned in the original Syriac work. Cp. Duchesne, \emph{loc. cit.}

} despatched messengers to Axum to incite the Abyssinians to avenge the slaughtered Christians and suppress the tyrant. In any case Ela Atzbeha invaded Yemen with a great army ( A.D. 524\texthebrew{‑}525), defeated and killed Dhu Novas, and set up in his

p325
stead a Himyarite Christian, whose name was Esimiphaios, as tributary king.\footnote{Procopius, \emph{B. P.} I.20.1. (John Mal. XVIII.457, says that the king of the Axumites made Anganes, a man of his own family, king of the Himyarites.) A Himyarite inscription found at Hisn-Gurab seems to record these events. It commemorates an Abyssinian invasion, and the defeat and death of the Himyarite king. The name of the man who set it up was read as Es\texthebrew{‑}Samaika, but Fell has plausibly suggested that the true reading may be Es\texthebrew{‑}Samaifa, which would correspond to \textgreek{\texthebrew{Ἐ}σιμιφα\texthebrew{ῖ}ος}. He does not, however, designate himself as king. The date is 640 of the Himyarite era, which (if the theory is correct) would be determined as 640\texthebrew{‑}525 = 115. See \emph{Z. D. M. G.} 7, 473, and 35, 36.

}

Such were the political relations of the two Red Sea kingdoms when, in A.D. 531, Justinian sent Julian, an agens in rebus , to the courts of Ela Atzbeha and Esimiphaios.\texthebrew{⁠}\footnote{There are two sources for this embassy, Procopius, \emph{ib.} 20.9 \emph{sqq.}, and John Mal. \emph{loc. cit.} (with the additions of Theophanes, who has placed it under a wrong year, A.M. 6064 = A.D. 571\texthebrew{‑}572). It can be inferred from the words of John Mal. (\textgreek{\texthebrew{ὠ}ς \texthebrew{ἐ}ξηγήσατο \texthebrew{ὁ} α\texthebrew{ὐ}τ\texthebrew{ὸ}ς πρεσβευτ\texthebrew{ὴ}ς}) that Julian published an account of this embassy, which was doubtless also known to Procopius.

} The purpose of the embassy was to win their co-operation against Persia in different ways. Julian travelled to Adulis by sea, and had an audience of Ela Atzbeha at Axum. The king stood on a four-wheeled car harnessed to four elephants. He was naked, except for a linen apron embroidered with gold and straps set with pearls over his stomach and shoulders.\texthebrew{⁠}\footnote{\textgreek{Σχιαστ\texthebrew{ὰ}ς δι\texthebrew{ὰ} μαργαριτ\texthebrew{ῶ}ν κα\texthebrew{ὶ} κλαβία \texthebrew{ἀ}ν\texthebrew{ὰ} πέντε} (John Mal. \emph{ib.}).

} He wore gold bracelets and held a gilt shield and two gilt lances. His councillors, who stood round him, were armed, and flute-players were performing. He kissed the seal of the Emperor's letter, and was amazed by the rich gifts which Julian brought him. He readily agreed to ally himself with the Empire against Persia. The chief service which the Abyssinians could render was to destroy the Persian monopoly in the silk trade by acting as carriers of silk between Ceylon and the Red Sea ports, a service which would also be highly profitable to themselves.

The consent of Ela Atzbeha, as over ­lord of Yemen, must also have been obtained to the proposals which Julian was instructed to lay before Esimiphaios. The Arabians of Maad (Nejd) were subject to the Himyarites, and their chieftain, Kais, who was a notable warrior, had slain a kinsman of the king and had been forced to flee into the desert. The plan of Justinian was to procure the pardon of Kais, in order that he, at the head of an army of Himyarites and Maadites, might invade the Persian empire.

p326
Although Julian was success ­ful in his negotiations and the kings promised to do what was required, they were unable to perform their promises. For men of Yemen to attack Persia meant long marches through the Arabian deserts, and the Himyarites shrank from such a difficult enterprise. In Ceylon the Abyssinian merchants were out-manoeuvred by the Persians, who bought up all the cargoes of silk as soon as they arrived in port.

It must have been soon after Julian's embassy that a revolt broke out in Yemen. Esimiphaios was dethroned and imprisoned, and a certain Abram, who was originally the slave of a Roman resident at Adulis, seized power.\texthebrew{⁠}\footnote{Procopius, \emph{ib.} 3\texthebrew{‑}8, who, as he says, anticipates events subsequent to Julian's embassy. According to \emph{Mart. Arethae} and the Gregentius documents (see below), which entirely ignore Esimiphaios, Abram was set up as king immediately after the overthrow of Dhu Novas. If Abram was commander of the resident Abyssinian troops, the error is explicable.\texthebrew{—} The name of Abram or Abraha was remembered in Arabic legend for his expedition against Mecca. His purpose was to destroy the Ka'ba. He was riding an elephant called Mahmud, and when he approached the city the animal knelt down and refused to advance. Then a flock of birds came flying from the sea with stones in their bills which they dropped on the heads of the troops. The legend is commemorated in the Koran (Sura 105): ``Hast thou not seen how the Lord dealt with the masters of the elephant?''

} It seems to have been a revolt of the Ethiopian garrison, not of the natives, and it is probable that Abram, who was a Christian, had been appointed commander of the garrison by Ela Atzbeha himself. Two expeditions were sent against Abram, but in both the Abyssinians were decisively defeated, and Ela Atzbeha then resigned himself to the recognition of Abram as viceroy.

Of the subsequent mission of Nonnosus, whom Justinian sent to Abyssinia, Yemen, and Maad, we only know that the ambassador on his journeys incurred many dangers from both men and beasts. The father of Nonnosus, Abram, was employed on similar business, and on two occasions conducted negotiations with Kais, the Arab chief of Nejd. Kais sent his son Muaviah as a hostage to Constantinople, and afterwards, having resigned the chieftaincy to his brother, visited the Imperial capital himself and was appointed phylarch of Palestine.\footnote{The fragments of the book of Nonnosus, preserved by Photius, are meagre and disappointing, and there are no chronological indications

(\emph{F. H. G.} IV.179) . We may conjecture that Abram was of Saracen race and that both he and Nonnosus could speak Arabic.

Thayer's Note: The link is to the Greek text and a Latin translation. The English translation by J. H. Freese is given

at Tertullian.Org .

}

p327
Historians and chroniclers tell us nothing of the revival of the Christian communities in the kingdom of the Himyarites after the fall of their persecutor Dhu Novas. There are other documents, however, which record the appointment of a bishop and describe his activities in Yemen. According to this tradition, Gregentius of Ulpiana was sent from Alexandria as bishop of Safar in the reign of Justin. He held a public disputation on the merits of Judaism and Christianity with a learned Jew and utterly discomfited him; and he drew up a Code of laws for Abram king of the Himyarites. As some of the historical statements in these documents are inconsistent with fact, the story of Gregentius has been regarded with scepticism and even his existence has been questioned.\texthebrew{⁠}\footnote{All the interesting parts of the \emph{Vita Gregentii}, preserved in a Sinaitic MS., have been published by Vasil'ev (see

Bibliography ). The disputation with the Jew Herbanus, which is included in the \emph{Vita}, is found by itself in many MSS., and had already been edited (see Migne, \emph{P. G.} LXXXVI). According to the Life, Gregentius was born at Ulpiana in Dardania. It is significant that this city, which is called \textgreek{Μπλιάρες} (cp. the Slavonic name Lipljan), is described as \textgreek{\texthebrew{ἐ}ν το\texthebrew{ῖ}ς μεθορίοις \texthebrew{Ἀ}βάρων} and \textgreek{τελο\texthebrew{ῦ}σα ε\texthebrew{ἰ}ς το α\texthebrew{ὐ}τ\texthebrew{ὸ} τ\texthebrew{ῶ}ν \texthebrew{Ἀ}βάρων γένος}, suggesting that the work was composed between 580 and 630. Gregentius travelled in Sicily, Italy, and Spain (\textgreek{Καρταγένα}), and finally went to Alexandria \textgreek{\texthebrew{ἐ}ν τα\texthebrew{ῖ}ς \texthebrew{ἡ}μέραις \texthebrew{Ἰ}ουστίνου βασιλέως \texthebrew{Ῥ}ωμαίων κα\texthebrew{ὶ Ἐ}λεσβο\texthebrew{ὰ}μ βασιλέως Α\texthebrew{ἰ}θιόπων κα\texthebrew{ὶ} Δουνα\texthebrew{ῦ}} (Dhu Novas) \textgreek{βασιλέως \texthebrew{Ὁ}μηριτ\texthebrew{ῶ}ν κα\texthebrew{ὶ} Προτερίου πάπα \texthebrew{Ἀ}λεξανδρείας}. When the Ethiopian king overthrew Dhu Novas, he wrote to Proterius asking him to ordain and send a bishop to the Himyarites. Proterius consecrated the deacon Gregentius, who travels apparently up the Nile and reaches the capital city of the Ethiopians, which is called \textgreek{\texthebrew{Ἀ}μλέμ}. Thence he proceeds to Safar, where he finds Ela Atzbeha, and under his auspices restores and founds churches at Nejran (\textgreek{Ηεγρά}), Safar, Akana, Legmia. Before Ela Atzbeha returned home (he had remained, \textgreek{\texthebrew{ὡ}ς \texthebrew{ἔ}φασάν τινες}, about three years in Yemen) he and Gregentius elevated Abram to the throne. In this narrative Abram appears as the successor of Dhu Novas; Esimiphaios is entirely ignored; and a Patriarch of Alexandria, Proterius, is introduced who is never mentioned in any records except in connexion with Gregentius. Another suspicious point should be noticed. Gregentius visits Agrigentum; the names of his parents are Agapius and Theodote. Now there exists the Life of a mysterious saint, Gregory of Agrigentum, and his parents were Chariton (which has much the same meaning as Agapius) and Theodote, while the name Gregentius itself suggests Agrigentum (cp. Vasil'ev, p67). In the Greek and Slavonic Menaea and Synaxaria Gregentius is noticed, sometimes under the name of Gregory. See \emph{Synax. eccl. Cplae.} (\emph{A. SS. Nov.}), p328.

} But there is no good reason to suppose that the story does not rest on a genuine tradition which was improved by legend and was written down when the historical details were forgotten. The Code of laws bears some internal marks of genuineness, though we may hope, for the sake of the Himyarites, that it was never enforced.\footnote{For the Code see below,

p413 .

}

The missionary zeal of Justinian and Theodora did not overlook the African peoples who lived on the Upper Nile between Egypt and Abyssinia. We have already seen how the hostility of the Blemyes, whose seats were above the First Cataract, and their southern neighbours the Nobadae, whose capital was at Dongola, constantly troubled the upper provinces of Egypt. The Nobadae and their king Silko were converted to Christianity about A.D. 540. The story of their conversion is curious. Theodora was determined that they should learn the Monophysitic doctrine; Justinian desired to make them Chalcedonians. In this competition for the souls of the Nobadae, Theodora was success ­ful. The episode is thus related by a Monophysitic historian:\footnote{John Eph. \emph{H. E.} Part III IV.6\texthebrew{‑}7. I have borrowed the version of Payne-Smith. A short account of the conversion of the Nobadae and Blemyes will be found in Duchesne, \emph{op. cit.} 287 \emph{sqq.}

}

Among the clergy in attendance on the Patriarch Theodosius was a proselyte named Julianus, an old man of great worth, who conceived an earnest spiritual desire to christianise the wandering people who dwell on the eastern borders of the Thebais beyond Egypt, and who are not only not subject to the authority of the Roman Empire, but even receive a subsidy on condition that they do not enter nor pillage Egypt. The blessed Julianus, therefore, being full of anxiety for this people, went and spoke about them to the late queen Theodora, in the hope of awakening in her a similar desire for their conversion; and as the queen was fervent in zeal for God, she received the proposal with joy, promised to do everything in her power for the conversion of these tribes from the errors of idolatry. In her joy, therefore, she informed the victorious King Justinian of the purposed undertaking, and promised and anxiously desired to send the blessed Julian thither. But when the king [Emperor] heard that the person he intended to send was opposed to the Council of Chalcedon, he was not pleased, and determined to write to the bishops of his own side in the Thebais, with orders for them to proceed thither and instruct the Nobadae, and plant among them the name of synod. And as he entered upon the matter with great zeal, he sent thither, without a moment's delay, ambassadors with gold and baptismal robes, and gifts of honour for the king of that people, and letters for the duke of the Thebais, enjoining him to take every care of the embassy and escort them to the territories of the Nobadae. When, however, the queen learnt these things, she quickly, with much cunning, wrote letters to the duke of the Thebais, which were as follows: ``Inasmuch as both his majesty and myself have purposed to send an embassy to the people of the Nobadae, and I am now

p329
despatching a blessed man named Julian; and further my will is that my ambassador should arrive at the aforesaid people before his majesty's; be warned, that if you permit his ambassador to arrive there before mine, and do not hinder him by various pretexts until mine shall have reached you and shall have passed through your province and arrived at his destination, your life shall answer for it; for I shall immediately send and take off your head.'' Soon after the receipt of this letter the king's ambassador also came, and the duke said to him, ``You must wait a little while we look out and procure beasts of burden and men who know the deserts, and then you will be able to proceed.'' And thus he delayed him until the arrival of the merci ­ful queen's embassy, who found horses and guides in waiting, and the same day, without loss of time, under a show of doing it by violence, they laid hands upon him, and were the first to proceed. As for the duke, he made his excuses to the king's ambassador, saying, ``Lo! when I had made my preparations and was desirous of sending you onward, ambassadors from the queen arrived and fell upon me with violence, and took away the beasts of burden I had got ready, and have passed onward; and I am too well acquainted with the fear in which the queen is held to venture to oppose them. But abide still with me until I can make fresh preparations for you, and then you also shall go in peace.'' And when he heard these things he rent his garments, and threatened him terribly and reviled him; and after some time he also was able to proceed, and followed the other's track without being aware of the fraud which had been practised upon him.

The blessed Julian meanwhile and the ambassadors who accompanied him had arrived at the confines of the Nobadae, whence they sent to the king and princes informing him of their coming; upon which an armed escort set out, who received them joyfully, and brought them into their land unto the king. And he too received them with pleasure, and her majesty's letter was presented and read to him, and the purport of it explained. They accepted also the magnificent honours sent them, and the numerous baptismal robes, and everything else richly provided for their use. And immediately with joy they yielded themselves up and utterly abjured the errors of their forefathers, and confessed the God of the Christians, saying, ``He is the one true God, and there is no other beside Him.'' And after Julian had given them much instruction, and taught them, he further told them about the council of Chalcedon, saying that ``inasmuch as certain disputes had sprung up among Christians touching the faith, and the blessed Theodosius being required to receive the council and having refused was ejected by the king [Emperor] from his throne, whereas the queen received him and rejoiced in him because he stood firm in the right faith and left his throne for its sake, on this account her majesty has sent us to you, that ye also may walk in the ways of Pope Theodosius, and stand in his faith and imitate his constancy. And moreover the king has sent unto you ambassadors, who are already on their way, in our footsteps.''

The Emperor's emissaries arrived soon afterwards, and were dismissed by Silko, who informed them that if his people embraced

p330
Christianity at all it would be the doctrine of the holy Theodosius of Alexandria, and not the ``wicked faith'' of the Emperor. The story, which is told by one who admired the Empress and lived under her protection, illustrates her unscrupulousness and her power.

The Nobadae, converted to Christianity, immediately co-operated with the Empire in chastising the Blemyes and forcing them to adopt the same faith. Roman troops under Narses made a demonstration on the frontier of the Thebaid, but the main work was done by Silko, who celebrated his victory by setting up an inscription in the temple of the Blemyes at Talmis (Dodekaschoinois, now Kelabsheh). The boast of this petty potentate might be appropriate in the mouth of Attila or Tamurlane: ``I do not allow my foes to rest in the shade but compel them to remain in the full sunlight, with no one to bring them water to their houses. I am a lion for the lands below, and a bear for the lands above.''\texthebrew{⁠}\footnote{Lefebvre, \emph{Recueil}, 628, p118; CIG III.5072. He describes himself as \textgreek{\texthebrew{ἐ}γ\texthebrew{ὼ} Σιλκ\texthebrew{ὼ} βασιλίσκος Ηουβάδων κα\texthebrew{ὶ ὅ}λων τ\texthebrew{ῶ}ν Α\texthebrew{ἰ}θιόπων}. The Greek who composed the inscription must have smiled to himself when he introduced the diminutive \textgreek{βασιλίσκος}, ``kinglet.'' Silko was succeeded by Eirpanomos (or Ergamenes). See Revillout, \emph{Mémoire sur les Blemmyes}, in \emph{Acad. des Inscr.}, sér. 1, VIII.2, pp371 \emph{sqq.} \texthebrew{—} References to the hostilities of the Blemyes in the sixth century will be found in \emph{Pap./Cairo} I.670004, 67007, and 67009. Here we find them plundering Omboi, and a pagan subject of the Empire reopening apparently a heathen temple for them (67004).

} The conversion of the Blemyes enabled Justinian to abolish the scandal of the pagan worship at Philae, which had been suffered to exist on account of an ancient convention with that people.\texthebrew{⁠}\footnote{See below,
p371 .

} A Greek agent was appointed to reside at Talmis and represent the Imperial authority.\footnote{Cp. CIG IV.8647\texthebrew{‑}8649 (posterior to Justinian's reign). Three interesting Greek inscriptions, found at Gebeleïn, are discussed by J. Krall, in \emph{Denkschr.} of the Vienna Academy, vol. XLVI (1900). The princes of the Blemyes, like those of the Nobadae, were styled \textgreek{βασιλίσκοι} by their Greek notaries. In the first of these texts (which date probably from the time of Anastasius or Justin) we find Charachen, \textgreek{βασιλείσκος τ\texthebrew{ῶ}ν Βλεμ\texthebrew{ῦ}ων}, giving orders as to an island in the Nile (perhaps near Gebeleïn), for which the Romans (\textgreek{ο\texthebrew{ἱ Ῥ}ώμεις}) paid a \textgreek{συνήθεια}. See Wessely, \emph{Gr. Papyrusurk.} No. 132.

}

The efforts of Justinian and his Abyssinian friends to break down the Persian monopoly of the silk trade had been frustrated by the superior organisation of Persian mercantile interests in

p331
the markets of Ceylon. There was one other route by which it might have been possible to import silk direct from China, namely overland through Central Asia and north of the Caspian Sea to Cherson. This possibility was no doubt considered. Justinian, however, does not seem to have made any attempt to realise it, but it was to be one of the political objects of his successor.

After the outbreak of the war with Persia in A.D. 540, the private silk factories of Berytus and Tyre suffered severely.\texthebrew{⁠}\footnote{Procopius, \emph{H. A.} 25.13 \emph{sqq.}

} It must be explained that, in order to prevent the Persian traders from taking advantage of competition to raise the price of silk, all the raw material was purchased from them by the commerciarii of the fisc, who then sold to private enterprises all that was not required by the public factories ( gynaecia ) which ministered to the needs of the court.\texthebrew{⁠}\footnote{\emph{C. J.} IV.40.2 . The most important study of the silk trade in the sixth century is that of Zachariä von Lingenthal, \emph{Eine Verordnung Justinians über den Seidenhandel} in \emph{Mém. de l'Acad. de St\texthebrew{‑}Pét.}, sér. VII vol. IX.6. (See also a paper of Herrmann, \emph{Die alern Seidenstrassen}, in Sieglins Forschungen, Heft 21.).

} Justinian instructed the commerciarii not to pay more than 15 gold pieces (£9:7:6) for a pound of silk, but he could not force the Persians to sell at this price, and they preferred not to sell at all or at least not to sell enough to serve the private as well as the public factories. It is not clear whether hostilities entirely suspended the trade, but at best they seriously embarrassed it, and as the supplies dwindled the industrial houses of Tyre and Berytus raised the prices of their manufactures. The Emperor intervened and fixed 8 gold pieces a pound as the maximum price of silk stuffs. The result was that many manufactures were ruined. Peter Barsymes, who was Count of the Sacred Largesses in A.D. 542, took advantage of the crisis to make the manufacture of silk a State monopoly, and some of the private industries which had failed were converted into government factories. This change created a new source of revenue for the treasury.\footnote{Procopius, \emph{ib.} The government sold silk stuffs with ordinary dye at 6 nomismata an ounce, but the Imperial dye, called \textgreek{\texthebrew{ὁ}λόβηρον}, at more than 24 nom. an ounce.

}

Chance came to the aid of Justinian ten years later and solved the problem more effectively than he could have hoped. Two monks, who had lived long in China or some adjacent

p332
country,\texthebrew{⁠}\footnote{Serinda, supposed by some to be Khotan. Procopius (\emph{B. G.} IV.17) describes it as \textgreek{\texthebrew{ὑ}π\texthebrew{ὲ}ρ \texthebrew{Ἰ}νδ\texthebrew{ῶ}ν \texthebrew{ἔ}θνη τ\texthebrew{ὰ} πολλά}. Possibly Cochin-China is meant. The sources are Procopius, \emph{ib.} (the order of whose narrative points to the year 552), and

Theophanes of Byzantium, \emph{F. H. G.} IV p270 , who ascribes the importation to a Persian.

} visited Constantinople ( A.D. 552) and explained to the Emperor the whole process of the cultivation of silkworms. Though the insect itself was too ephemeral to be carried a long distance, they suggested that it would be possible to transport eggs, and were convinced that they could be hatched in dung, and that the worms could thrive on mulberry leaves in Europe as successfully as in China. Justinian offered them large rewards if they procured eggs and smuggled them to Constantinople. They willingly undertook the adventure, and returned a second time from the East with the precious eggs concealed in a hollow cane. The worms were developed under their instructions, Syria was covered with mulberry trees, and a new industry was introduced into Europe. Years indeed must elapse before the home-grown silk sufficed for the needs of the Empire, and in the meantime importation through Persia continued,\texthebrew{⁠}\footnote{Compare the treaty of 562 (above,
p121 ).

} and Justinian's successor attempted to open a new way of supply with the help of the Turks.

If we regard commerce as a whole, there is no doubt that it prospered in the sixth century. Significant is the universal credit and currency which the Imperial gold nomisma enjoyed. Cosmas Indicopleustes, arguing that the ``Roman Empire participates in the dignity of Christ, transcending every other power, and will remain unconquered till the final consummation,'' mentions as a proof of its eminent position that all nations from one end of the earth to the other use the Imperial coinage in their mercantile transactions.\texthebrew{⁠}\footnote{Cosmas, II p81. Coins of Arcadius and Honorius, Theodosius II, Marcian, Leo I, Zeno, Anastasius and Justin I have been found in southern and western India; of Marcian and Leo in northern India. See

Sewell, \emph{Roman Coins found in India, Journal Asiat. Soc.} XXXVI.620\texthebrew{‑}635 (1904) .

} Illustrative anecdotes had been told of old by merchants who visited Ceylon. Pliny relates that a freedman who landed there exhibited Roman denarii to the king, who was deeply impressed by the fact that all were of equal weight though they bore the busts of different Emperors.\texthebrew{⁠}\footnote{\emph{Nat. Hist.} VI.22 .

} Sopatros, a Roman merchant who went to Ceylon in an Ethiopian vessel in the reign of Zeno or Anastasius, told Cosmas\texthebrew{⁠}\footnote{XI p323.

} that he

p333
had an audience of the king along with a Persian who had arrived at the same time. The king asked them, ``Which of your monarchs is the greater?'' The Persian promptly replied, ``Ours, he is the king of kings.'' When Sopatros was silent, the king said, ``And you, Roman, do you say nothing?'' Sopatros replied, ``If you would know the truth, both the kings are here.'' ``What do you mean?'' asked the king. ``Here you have their coins,'' said Sopatros, ``the nomisma of the one and the drachm of the other. Examine them.'' The Persian silver coin was good enough, but could not be compared to the bright and shapely gold piece. Though Sopatros was probably appropriating to himself an ancient traveller's tale,\texthebrew{⁠}\footnote{So Winstedt rightly (his ed. of Cosmas, p355).

} it illustrates the prestige of the Imperial mint.

The independent German kingdoms of the West still found it to their interest to preserve the images and superscriptions of the Emperors on their gold money. In the reign of Justinian the Gallic coins of the Merovingian Franks have the Emperor's bust and only the initials of the names of the kings.\texthebrew{⁠}\footnote{Only Theodebert, as a sign of defiance, substituted his own name for that of Justinian, but left the title PP Aug. This is what Procopius refers to, \emph{B. G.} III.33.5. Here he makes the curious statement that it is not lawful (\textgreek{θέμις}) for any barbarian potentates, including the Persian king, to stamp gold coins with their own images, because even their own merchants would not accept such money.

} The Suevians in Spain continued to reproduce the monetary types of Honorius and Avitus. The last two Ostrogothic kings struck Imperial coinage, only showing their hostility to Justinian by substituting for his image and inscriptions those of Anastasius.\footnote{Wroth, \emph{Catalogue of Coins of the Vandals}, etc., Plates X and XII, cp. Introd. p. xxxviii. The rule did not apply to silver and bronze coins.

\texthebrew{◂} previous

next \texthebrew{▸}Images with borders lead to more information.

The thicker the border, the more information.

(Details here.)

UP TO:

J. B. Bury
Homepage

Latin \& Greek

Texts

LacusCurtius

Home}

\xchapter{Chapter XXI}

The second Prefecture of John the Cappadocian ( A.D. 533\texthebrew{‑}540) was marked by a series of reforms in the administration of the Eastern provinces, and it would be interesting to know how far he was responsible for instigating them. Administrative laws affecting the provinces were probably, as a result, evoked by reports of the Praetorian Prefects calling attention to abuses or anomalies and suggesting changes.\texthebrew{⁠}\footnote{For instance, \emph{Nov.} 151 (A.D. 533\texthebrew{‑}534), intended to check the practice of senators and officials in the provinces coming to Constantinople on litigation business, was due to a relatio of John Capp.

} If half of what the writers of the time tell us of John's character is true, we should not expect to find him promoting legislation designed to relieve the lot of the provincial taxpayers. But we observe that, while the legislator is earnestly professing his sincere solicitude for the welfare of his subjects, he always has his eye on the interests of the revenue, and does not pretend to disguise it. The removal of abuses which diminished the power of the subjects to pay the taxes was in the interest of the treasury, and it was a capital blunder of the fiscal administration of the later Empire that this obvious truth was not kept steadily in view and made a governing principle of policy. It was fitfully recognised when the excessive burdens of the cultivators of the land led to an accumulation of arrears and the danger of bankruptcy, or when some glaring abuse came to light. John, clever as he was, could not extract money from an empty purse, and there is no reason to suppose that he may not have promoted some of the remedial laws which the Emperor directed him to administer.

p335
We need not doubt that the Emperor was thoroughly sincere when he asserts his own concern for the welfare of his subjects, nor suspect him of hypocrisy when he expresses indignation at the abuses which he strives to suppress. All the capable Roman Emperors honestly desired a pure administration and a contented people; but their good intentions were frustrated by defects of the fiscal system which they had inherited, and by the corruption of the vast army of officials who administered it.

We do not know how far Justinian's enactments may have been success ­ful, but they teach us the abuses which existed. There was none perhaps which he himself regarded as more important \texthebrew{—} if we may judge from his language \texthebrew{—} than the law which forbade the practice of buying the post of a provincial governor.\texthebrew{⁠}\footnote{\emph{Nov.} 8

(April 15, 535). Procopius refers to this law

(\emph{H. A.} 21.16)

and says that before a year had passed Justinian disregarded it and allowed the offices to be sold openly.

} Theodora, if she did not instigate the measure, had taken a deep interest in it, and the Emperor also expressly acknowledges that he had received some help from the Prefect. It had long been the custom to require the payment of considerable sums ( suffragia ) from those who received appointments as governors of provinces, and these sums went partly to the Emperor, partly to the Praetorian Prefect. Men who aspired to these posts were often obliged to borrow the money. The official salary was not sufficient to recompense them for the expense of obtaining the post, and they calculated on reimbursing themselves by irregular means at the cost of the provincials. The Emperor states that they used to extract from the taxpayers three or even ten times the amount they had paid for the office, and he shows how the system caused loss to the treasury, and led to the sale of justice and to general demoralisation in the provinces. The law abolishes the system of suffragia . Henceforward the governor must live on his salary, and when he is appointed he will only have to pay certain fixed fees for the ensigns and diploma of his office.\texthebrew{⁠}\footnote{The fees (\textgreek{συνήθειαι}) for the higher posts (like the comes Orientis, proconsul of Asia) amounted to £122:10s, for the posts of consular rank £47:10s, for those of praeses or corrector, about £39.

} Before he enters on his post he has to swear \texthebrew{—} the form of oath is prescribed \texthebrew{—} that he has paid no man any money as a suffragium and severe penalties

p336
are provided if the Prefect or any of his staff or any other person should be convicted of having received such bribes. The governor who has paid for his appointment or who receives bribes during his administration is liable to exile, confiscation of property, and corporal punishment. Justinian takes the opportunity of exhorting his subjects to pay their taxes loyally, ``inasmuch as the military preparations and the offensive measures against the enemy which are now engaging us are urgent and cannot be carried on without money; for we cannot allow Roman territory to be diminished, and having recovered Africa from the Vandals, we have greater acquisitions in view.''

Several other laws were passed in this period to protect the people from mal-administration .\texthebrew{⁠}\footnote{See \emph{Edicts} 2 and 12.

} The confirmation of the old rule that a governor should remain in his province for fifty days after vacating his office, in order to answer any charges against his actions, may specially be mentioned.\texthebrew{⁠}\footnote{\emph{Nov.} 95

(539); cp. \emph{Nov.} 8, § 9 .

} The office of defensor civitatis had become practically useless as a safeguard against injustice because it had come to be filled by persons of no standing or influence, who could not assume an independent attitude towards the governors. Justinian sought to restore its usefulness by a reform which can hardly have been welcomed by the municipalities. He ordained that the leading citizens in each town should fill the office for two years in rotation; and he imposed on the defensor , in addition to his former functions, the duty of deciding lawsuits not involving more than 300 nomismata and of judging in minor criminal cases. The work of the governor's court was thus lightened. We may suspect that the bishops who were authorised to intervene were more efficacious in defending the rights of the provincials because they were more independent of the governor's goodwill.\footnote{\emph{Nov.} 15

(535); \emph{Nov.} 86 (539).

}

Among the restrictions which the Roman autocrats placed upon the liberty of their subjects there is none perhaps that would appear more intolerable to a modern freeman than those which hindered freedom of movement. It was the desire of the Emperors to keep the provincials in their own native places and to discourage their changing their homes or visiting the

p337
capital. This policy was dictated by requirements of the system of taxation, and by the danger and inconvenience of increasing the proletariat of Constantinople. Impoverished provincials had played a great part in the Nika sedition, and the duties of the Prefect of the City were rendered more difficult and onerous by the arrival of multitudes of unemployed persons to seek a living by beggary or crime. Justinian created a new ministry of police for the special purpose of dealing with this problem.\texthebrew{⁠}\footnote{\emph{Nov.} 80 (A.D. 539), addressed to the Pr. Prefect of the East, because it concerned the provincials as well as the capital. The institution of the quaesitor is also mentioned by John Lyd. \emph{De mag.} II.29 (\textgreek{τ\texthebrew{ὸ}ν λεγόμενον κυαισίτορα \texthebrew{ἀ}ντι το\texthebrew{ῦ} τ\texthebrew{ῶ}ν βιωτικ\texthebrew{ῶ}ν \texthebrew{ἐ}γκλημάτων \texthebrew{ἐ}πρευνάδα σεμνότατον}), by John Mal. XVIII.479 (A.D. 539\texthebrew{‑}540), and by Procopius, \emph{H. A.} 20. 9

and

11 . The notice of Procopius ignores the principal functions of the quaesitor, and represents him as concerned with unnatural vice and offences against religion. Nothing is said of such duties in

\emph{Nov.} 80 , and Panchenko (\emph{O tain. ist.} 213) therefore thinks that the quaesitor of Procopius and John Lydus is a different official from that of

\emph{Nov.} 80 , whose title he supposes to have been quaestor. But he is wrong in supposing that \textgreek{κοιαισίτωρ} º has no manuscript authority in the Novel (see Kroll's ed. p391). The difficulty may easily be solved, I think, by supposing that the duties mentioned by Procopius were subsequently assigned to the quaesitor, who was already empowered (\emph{Nov.} 80, § 7) to try cases of forgery (\textgreek{πλαστογραφία}).

} The function of the quaesitor , as the minister was called, was to inquire into the circumstances and business of all persons who came from the provinces to take up their quarters in the capital, to assist those who came for legitimate reasons to get their business transacted quickly and speed them back to their homes,\texthebrew{⁠}\footnote{The business affairs, referred to in the law, are chiefly lawsuits (cp. \emph{Nov.} 69 ), or the affairs of agricultural tenants whose landlords resided in the capital. Persons who have once been dismissed from the city by the quaesitor and return are liable to punishment (§ 9).

} and to send back to the provinces those who had no valid excuse for having left their native soil. He was also empowered to deal with the unemployed class in the capital, and to force those who were physically fit into the service of some public industry (such as the bakeries), on pain of being expelled from the city if they refused to work. Judicial functions were also entrusted to him, and his court dealt with certain classes of crime, for instance forgery.

The Prefect of the City was further relieved of a part of his large responsibilities by the creation of another minister, who, like the quaesitor, was both a judge and a chief of police. The Praefectus Vigilum,\texthebrew{⁠}\footnote{The Greek terms was \textgreek{νυκτέπαρχος} ``night-prefect,'' which Justinian wastes many words in deriding.

} who was subordinate to the Prefect, was abolished, and his place was taken by an independent official

p338
who was named the Praetor of the Demes,\texthebrew{⁠}\footnote{\textgreek{Πραίτωρ δήμων,}\emph{Nov.} 13 , A.D. 535 (John Mal. \emph{ib.}, gives a wrong date, 539); John Lyd. \emph{ib.}; Procopius, \emph{ib.} We can infer from the law that at this time crime was particularly prevalent in the capital; and the Praefectus Vigilum employed agents who were in collusion with the criminals. Procopius complains that the Praetor and the Quaesitor were arbitrary in administering justice, condemning men without evidence, and, in accordance with his thesis, represents them as instituted for the purpose of oppression and extortion.

} and whose most important duty was to catch and punish thieves and robbers.

During the fifth century few changes had been made in the details of the provincial system as it was ordered by Diocletian and modified here and there by his successors. Such alterations as had been found advisable were in accordance with the principles which had inspired Diocletian's reform. Provinces were further subdivided, they were not enlarged. Theodosius II, for instance, broke up Epirus, Galatia, and Palestine, each into two provinces.\texthebrew{⁠}\footnote{John Mal. XIII.347\texthebrew{‑}348.

} Changes had also been made in Egypt. This diocese had at first consisted of five provinces, Aegyptus, Augustamnica, Thebais, and the two Libyas, but Theodosius I (after A.D. 486) cut off a part of Augustamnica (including the Oxyrhynchus district) to form the province of Arcadia.\texthebrew{⁠}\footnote{\emph{Not. dig., Or.} I.29 . See M. Gelzer, \emph{Stud. z. byz. Verw. Ägyptens}, 8\texthebrew{‑}9.

} At some later period Augustamnica was again divided into two provinces, Prima and Secunda.\texthebrew{⁠}\footnote{Before A.D. 535 (Hierocles, \emph{Synekd.} 726.3; 727.13; Justinian,

\emph{Nov.} 8 ).

} But the principal innovation was made by Theodosius II, who subdivided the Thebaid into the Upper and Lower provinces. The Upper or southern Thebaid was constituted under a duke, to whom the civil as well as the military administration was entrusted, along with a general authority over the Lower Thebaid, which had its own civil governor.\texthebrew{⁠}\footnote{We have a contemporary record of this change in an Imperial rescript preserved in a Leiden papyrus (\emph{Archiv für Papyrusforschung}, I.397), and discussed by Gelzer, \emph{ib.} 10 \emph{sqq.}, who shows that the innovation must be dated between 435 and 450. The Upper Thebaid extended from Ptolemais to Omboi.

} The motive of this arrangement was to strengthen the hands of the commander who was responsible for protecting the frontier against the Blemyes and Nobadae. Yet another alteration was made, perhaps early in the sixth century;

p339
the province of Aegyptus was divided into two, Prima and Secunda.\footnote{\emph{Nov.} 8

(April 15, 535); in Hierocles Egypt is still one province. The precise date of the notitia of Hierocles has not been fixed.

}

A charge º of a different kind, but based on the same principle of dividing responsibility, had been introduced by Anastasius in Thrace. When he constructed the Long Wall he established a new vicariate, at the expense of the vicariate or diocese of Thrace. We do not know its extent, or with powers the new official possessed, but as he was entitled ``Vicar of the Long Walls'' his diocese evidently stretched northwards from Constantinople.\footnote{The sources are the notitiae appended to

\emph{Nov.} 8 , and

\emph{Nov.} 26 . That the vicariate was created by Anastasius is not stated, but is a natural inference.

}

Justinian did not indeed attempt a complete revision of the existing system, but he made a great number of changes in which he departed from the principles of Diocletian. He combined in some cases small provinces to form larger circumscriptions; he did away with most of the diocesan governors, who formed the intermediate links in the hierarchical chain between the provincial governors and the Praetorian Prefect; and he united in many case the civil and military powers which had been so strictly divorced by Diocletian. The tendency of these changes anticipates to some extent the later system which was to come into being in the seventh century and was characterised by large provinces, the union of civil and military administration in the same hands, and the total disappearance of the dioceses. The reforms of Justinian, which belong to the years 535 and 536, were called forth by particular circumstances. Some of them were designed to avert conflicts between the civil and military authorities.

The Count of the East was deprived of his jurisdiction over the Orient diocese and, retaining his title, rank, and emoluments, became the civil governor of the province of Syria Prima. The Vicariate of Asiana was likewise abolished, and the vicar became the governor of Phrygia Pacatiana, exercising both civil and military powers, and adorned with the new title of comes Iustinianus .\footnote{\emph{Nov.} 8

(April 15, 535).

}

Similarly the Vicariate of Pontica was abolished, the vicar becoming the comes Iustinianus of Galatia Prima. But this arrangement was found to work badly, and at the end of thirteen

p340
years the vicariate was restored. We are told that the Pontic provinces were infested by robbers and assassins, who formed armed bands and escaped the justice which threatened them in one province by moving into another. No governor ventured to transgress the limits of his own province by pursuing them. It seemed that the difficulty could only be met by the appointment of a superior governor with jurisdiction over all the provinces, and the vicar of Pontica was reinstated, but with powers considerably larger than those which had belonged to him before. He was to have military and financial as well as civil functions. He was to be the vicar not only of the Praetorian Prefect, but also of the Master of Soldiers, and was to have authority over all the troops stationed in his diocese. He was also to represent the Master of Offices and the Counts of the Private Estate and the Sacred Patrimony; so that none of the officials who served these ministers could defy or evade his authority.\footnote{\emph{Edict} 8, A.D. 548.

}

In Thrace discord between the military and civil officials appears to have been incessant, and as the Thracian provinces constantly suffered from the incursions of the barbarians, want of harmony in the administration was more disastrous here than elsewhere. Justinian abolished the Vicar of Thrace and the Vicar of the Long Wall, and committed the civil and military power of the whole diocese to a single governor with the title of Praetor Justinianus of Thrace.\texthebrew{⁠}\footnote{\emph{Nov.} 26 , A.D. 535, May 18.

} Soon afterwards, however, the dominion of the new Praetor was curtailed by the withdrawal from his jurisdiction of the frontier provinces of Lower Moesia and Scythia. These, by a very curious arrangement, were associated with Caria, the Cyclades, and Cyprus, and placed under the control of a governor entitled \emph{Quaestor Iustinianus of the Army}, who enjoyed an authority independent of the Praetorian Prefect as well as of the Master of Soldiers. He was really a fourth Praetorian Prefect but with military functions, and his institution must have been deliberately intended to diminish the power of the Prefect of the East. The motive of this strange union of provinces so far apart and without any common interest to connect them is unknown;\texthebrew{⁠}\footnote{The body of the law (\emph{Nov.} 41) which created the office, May 18, 536, is lost, but information is supplied by

\emph{Nov.} 50

(August 537).

} but we may conjecture that the

p341
object was to place the financial expenses of administering the Danubian lands, exhausted by invasions, on provinces which were exceptionally rich.\footnote{This is rather suggested by the words of John Lyd. \emph{De mag.} II.29 \textgreek{προάγει \texthebrew{ἔ}παρχωον \texthebrew{ἐ}πόπτην τ\texthebrew{ῶ}ν Σκυθικ\texthebrew{ῶ}ν δυνάμεων, \texthebrew{ἀ}φορίσας α\texthebrew{ὐ}τ\texthebrew{ῷ ἐ}παρχίας τρε\texthebrew{ῖ}ς τ\texthebrew{ὰ}ς τασ\texthebrew{ῶ}ν \texthebrew{ἐ}γγ\texthebrew{ὺ}ς ε\texthebrew{ὐ}πορυτάτας.}}

These changes made a considerable breach in the hierarchical system which had been constructed by Diocletian and Constantine. The union of civil and military powers was also introduced in many of the Asiatic provinces, and in every case the new governor received the rank of spectabilis and a new title. Pisidia and Lycaonia were each placed under a Praetor Iustinianus. The Count of Isauria had already possessed the double authority under the old system; Justinian did not change his title, but gave him the rank of spectabilis.\texthebrew{⁠}\footnote{Pisidia,

\emph{Nov.} 24 ; Lycaonia,

\emph{Nov.} 25 ; Isauria,

\emph{Nov.} 27

(all in May 535).

} In three cases large provinces were created by the union of two smaller. Pontus Polemoniacus was joined to Helenopontus, and formed a new Helenopontus under a Moderator Iustinianus. Paphlagonia and Honorias were reunited as Paphlagonia under a Praetor Iustinianus. The Moderator and the Praetor possessed the double functions.\footnote{Helenopontus,

\emph{Nov.} 28 ; Paphlagonia,

\emph{Nov.} 29

(July 535). Other changes were the elevation of the praeses of Phoenicia Libanensis to the rank of Moderator (spectabilis), and that of the praeses of Palestine Salutaris to that of proconsul, with authority to supervise the government of Palestine Secunda. See \emph{Edict} 4 and \emph{Nov.} 103 (536). The civil governor of Arabia, whose authority had been reduced to a cipher by that of the duke, was made a moderator, \emph{Nov.} 102 (536).

}

The third case was the union of the two provinces of Cappadocia under a Proconsul Iustinianus. Cappadocia presented peculiar problems of its own. It had drifted into an almost anarchical condition which demanded special treatment. Here were the large Imperial domains, which were under the management of the Praepositus of the Sacred Bedchamber, and the rest of the land seems to have mainly consisted of large private estates. The wealthy landowners and their stewards kept bodies of armed retainers, and acted as if they were masters of the provinces. They even encroached upon the Imperial domains, and the Emperor complains that ``almost all the Imperial Estate has become private property.'' He declares that every day he and his ministers have to deal with the petitions of Cappadocians who have been deprived of their property, including clergy and especially women. The governors and officials were afraid to

p342
resist these power ­ful magnates, who stopped their mouths with gold. ``The crimes which are committed in that country,'' says Justinian, ``are so many that even the greatest man would find it difficult to check them.'' He therefore invested the new governor of united Cappadocia with exceptional powers and prestige. The Proconsul controlled the civil administration and the military forces, but he was also responsible for the revenue and controlled all the officials and agents of the Private Estate, and that not only in Cappadocia, but in other provinces of the Pontic diocese. He received a salary double that of the Moderator of Helenopontus or the Praetor of Paphlagonia.\footnote{\emph{Nov.} 30 . The salary of the Proconsul was 20 lbs. of gold (about £900); that of the Moderator and the Praetor was 725 nom. (about £450). The Praetors of Thrace, Lycaonia, and Pisidia received 300 nom. annually (£187).

}

Some changes were also made in the administration of Egypt.\texthebrew{⁠}\footnote{\emph{Edict} 13, A.D. 538\texthebrew{‑}539 (there can be little doubt about the date; it was addressed to John Capp. in a 2nd indiction, see § 15 and § 24. The reasons given by Zachariä von L. for ascribing it to 553\texthebrew{‑}554 are not convincing, and are confuted by papyrus evidence, see Gelzer, \emph{ib.} 23 \emph{sq.}). This long edict throws much light on the arrangements connected with the corn supplies.

} Here perhaps the chief preoccupation of the government was to secure the regular delivery of the grain with which the country of the Nile supplied Constantinople. Justinian found that the wheels of the administrative machinery were out of gear. For some time back, he says, things have been in such confusion in the Egyptian Diocese that the central authorities have not known what was going on there. ``The taxpayers asserted that all the legal dues were demanded in a lump, and that they had entirely fulfilled their liabilities, while we received nothing beyond the corn º supplies; and the curials, the pagarchs (mayors of the villages), the tax-collectors , and the governors arranged things in such a way as to obscure the true facts and to make profit for themselves.'' But there were other considerations, which, though not specially mentioned in the Imperial edict, must have influenced the legislator. In A.D. 536 and 537 Alexandria had been the scene of popular seditions, arising out of a contest between two heretical claimants to the Patriarchal throne. The military forces had been powerless to suppress the disorders.\footnote{Liberatus, \emph{Brev.} 19. Cp. Gelzer, \emph{ib.} 25 \emph{sqq.}

}

Justinian here adopted a policy opposite to that which he had pursued in Cappadocia. Instead of making one man responsible for the whole administration, he reduced the responsibilities

p343
of the Augustal Prefect, who had hitherto governed the Diocese. He made him governor of Alexandria and of the two provinces of Aegyptus Prima and Aegyptus Secunda, with civil and military powers.\texthebrew{⁠}\footnote{Mareotes and the city of Menelaites were separated from Aegyptus Prima and added to the province of Libya. The Prefect's staff, both civil (\textgreek{α\texthebrew{ὐ}γουσταλιανή}) and military (\textgreek{δουκική}), was to number 600, and he was to receive the large salary of 40 lbs. of gold (£1800).

} These provinces were not united; they still retained their civil governors, subordinate to the Prefect, who now bore the title of duke. The Emperor expressly justified this change by the consideration that the supervision of the whole Diocese was too much for one man. It is not quite clear whether the two provinces of Lower and Upper Libya were united under one civil prefect, or whether they continued to be distinct, but in either case the governors were placed under the control of the military duke of the Libyan frontier.\texthebrew{⁠}\footnote{The Edict speaks of \textgreek{\texthebrew{ἡ} Λιβύων \texthebrew{ἐ}παρχία} and \textgreek{το\texthebrew{ῦ} τ\texthebrew{ῆ}ς Λιβύων πολιτικο\texthebrew{ῦ}} (civil) \textgreek{\texthebrew{ἄ}ρχοντος}, as if there were only one province (cp. the note of Zachariä v. L. in his edition, p51). But in the \emph{Descriptio orbis Romani} of Georgius Cyprius (\emph{c.} A.D. 600) the two provinces appear (ed. Gelzer, p40). Paraeetonium in Lower Libya was the seat of the duke of the Libyan frontier. Upper Libya was the Cyrenaica.

} In Upper Egypt the duke of the Thebaid received the Augustal title and was endowed with both civil and military authority over the two Thebaid provinces whose governors were subordinate to him. The general result of these reforms was the completion of the policy of abolishing Diocesan governors in the Eastern Prefecture. In Egypt there were now eight (or nine) provinces grouped in five independent circumscriptions, Egypt, Augustamnica, Arcadia,\texthebrew{⁠}\footnote{Arcadia is not mentioned in the Edict, but there is some evidence from papyri which confirms the reasonable inference that it was independent (Gelzer, \emph{ib.} 29). From the same sources we learn that the governor was known as the count of Arcadia, and afterwards had the rank of Patrician (\emph{ib.} 33), a dignity which was also conferred on the Augustal dukes of Egypt and the Thebaid. For the title \textgreek{δο\texthebrew{ὺ}ξ κι\texthebrew{ὰ} α\texthebrew{ὐ}γουστάλιος} of the duke of Thebais see \emph{ib.} 23.

} Thebais, and Libya, of which the governors had each military as well as civil competence and were directly responsible to the Praetorian Prefect of the East.

The law which introduced these changes laid down minute regulations for the collection and transportation of the corn supplies both for Constantinople and for Alexandria, and for the gathering in of all other dues whether for the treasury of the Praetorian Prefect or for that of the Count of the Sacred Largesses. The several duties and responsibilities of all the authorities concerned were carefully distinguished.

p344
The treatment of the Armenian provinces, which embraced the most easterly districts of the Diocese of Pontus, stands apart. Here Justinian's policy was not to increase the size of the governments, but to rearrange.\texthebrew{⁠} He formed four provinces, partly by readjustments in the two old Armenian provinces, partly by taking districts from Helenopontus, and partly by converting new districts into provincial territory, suppressing the native satraps.

The new First Armenia, which had the privilege of being governed by a proconsul, included four towns of the old First Armenia, namely Theodosiopolis, Satala, Nicopolis, and Colonea, and two towns of the old Pontus Polemoniacus, Trapezus and Cerasus. The once important town of Bazanis or Leontopolis received the name of the Emperor, and was elevated to the rank of the metropolis.

The new Second Armenia, under a praeses , corresponded to the old First Armenia, and included its towns Sebastea and Sebastopolis. But in place of the towns which Justinian handed over to the new First Armenia, it received Comana, Zela, and Brisa from the new province of Helenopontus.

The Third Armenia, governed by a comes Iustinianus with military as well as civil authority, corresponded to the old Second Armenia, and included Melitene, Arca, Arabissus, Cucusus, Ariarathea, and Cappadocian Comana.

Fourth Armenia was a province new in fact as well as in name, consisting of the Roman districts beyond the Euphrates\texthebrew{⁠}\footnote{Sophanene, Anxitene, Sophene, Asthianene, and Belabitene. Cp. Procopius, \emph{Aed.} III.1 .

} (to the east of the Third Armenia), which had hitherto been governed by native satraps. It was placed under a consular, and the metropolis was Martyropolis.

The names appear to have been determined by the geographical order. The new trans-Euphratean province went naturally with the district of Melitene, and therefore the Second Armenia became the Third, because it was connected with what it was most natural to call the Fourth. For the consular of Fourth Armenia was to be in a certain way dependent on the count of Third Armenia, who was to hear appeals from the less important province. In the same way the new First and Second Armenias naturally went together, and therefore it was convenient

p345
that the numbers should be consecutive. The praeses of Second was dependent to a certain extent on the proconsul of First Armenia.\footnote{\emph{Nov.} 31

(March 536). The regulations are discussed at length by Adonts, \emph{Armeniia v Epokhu Iustiniana}, 157 \emph{sqq.} I cannot think that he is right in supposing that in selecting his First and Third Armenia as provinces of superior rank the Emperor was influenced by no better reason than ``to reward the Imperial favourites Acacius and Thomas,'' who at the time were governing in those districts (p176); for he could easily have transferred them. Both these persons were of Armenian origin, and Procopius gives Acacius a very bad character (\emph{B. P.} II.3).

}

In the case of these provinces, Justinian not only revised the administrative machinery, but also introduced changes of another kind. Hitherto the Armenians had lived according to their own laws and customs, and had not been called upon to regulate their private dealings according to the civil law of Rome. It was in the domain of real property that the divergence of Armenian from Roman law provoked the Emperor's special intervention. Armenian estates\texthebrew{⁠}\footnote{\textgreek{Γενεαρχικ\texthebrew{ὰ} χωρία.}} passed undivided from father to son, or in default of a son to the nearest male agnate. No proprietor could leave his property by will \texthebrew{—} wills, in fact, were unknown. No woman could inherit, nor did she receive a dowry when she married. Justinian determined to break down this system, which he professes to consider barbarous; and in two successive laws\texthebrew{⁠}\footnote{\emph{Edict} 3 (535): \emph{Nov.} 21 (March 536). These documents are discussed at length by Adonts, \emph{op. cit.} 184 \emph{sqq.}

} he ordained that henceforward the inheritance of property should be regulated by Roman law, that women should inherit their due shares, and should receive dowries. It is not probable that Justinian was moved to this reform solely by consideration of the female population of the Armenian provinces. Apart from the fact that it outraged his ideal of uniformity that Roman law should not prevail in any quarter of the Empire, we may suspect that it was his aim to break up the large estates of Armenia and thereby weaken the power of the princes and magnates, to force them to give up their national exclusiveness and draw them into the sphere of general Imperial interests.\texthebrew{⁠}\footnote{Adonts, \emph{ib.} 201.

} The policy was crowned with success. Constantinople and the Imperial service had already begun to attract many Armenians, and this movement towards the centre increased. In Justinian's reign men of this race began to come to the front in the Imperial service; Narses and Artabanes are the

p346
most eminent examples. Hereafter they would ascend the throne itself.\footnote{In the period after Justinian, and indirectly as a consequence of his policy, a westward expansion of the Armenian population began in two directions, from the Second Armenia westward towards Caesarea and north-westward towards the Black Sea, and from the Third Armenia south-westward towards Cilicia and the Mediterranean. See Adonts, \emph{ib.} 203 \emph{sq.}

}

The long list of administrative changes which we have surveyed shows that the Emperor addressed himself earnestly in A.D. 535 to the task of thoroughly overhauling the system of provincial government, and, in the appreciation of his work as a ruler, these reforms have hardly received due attention. He did not attempt, according to any general preconceived plan, to organise a new system, like Augustus or Diocletian, but sought to remedy, in each case according to its own circumstances, the defects of the existing scheme. It is characteristic of him that he likes to justify his innovations by appeals to history and antiquity. For example, when he bestows upon Lycaonia a governor of higher rank with the title of praetor, he pedantically recalls the legendary connexion of the country with Lycaon of Arcadia, who was also said to have colonised in Italy, thereby anticipating Aeneas the ancestor of Romulus. ``On this account, it would be just to decorate the province with the ancient symbols of Roman government, and therefore we give the governor the title of praetor, older even than that of consul.'' It was probably a consideration of public opinion as well as his own personal sentiments that made him seek to represent his innovations, whenever it was possible, as reversions to an older order. He wished it to be thought, and possibly thought himself, that he was ``reintrodu­cing antiquity with greater splendour.''\texthebrew{⁠}\footnote{\textgreek{Τ\texthebrew{ὴ}ν παλαιότητα πάλιν μετ\texthebrew{ὰ} μείζονος \texthebrew{ἄ}νθους ε\texthebrew{ἰ}ς τ\texthebrew{ὴ}ν πολιτείαν \texthebrew{ἐ}παναγαγόντες,}\emph{Nov.} 24, § 1 .

} He frequently speaks with pride of his own native language, Latin; yet it was in his reign that it definitely became the practice to issue the laws in Greek. The contrast between the innovator and the enthusiast for historical tradition stands out most conspicuously in the abolition of the consul ­ship.

It would be difficult to contend that Justinian in allowing the consul ­ship to lapse was not thoroughly justified by the

p347
circumstances. Before he finally took this step, he had made an effort to render possible the preservation of an institution ``which for nearly a thousand years had grown with the growth of the Roman state.'' For all political purposes the institution was obsolete. It was a distinction to a man to hold it, to give his name to the year and have it perpetuated in the Fasti Consulares. But the public spectacles, which the new consul exhibited in the first weeks of January, and the largesses which he was expected to distribute to the people, entailed a large outlay, which only the wealthiest could undertake. It became more and more difficult to find private persons ready to incur the expenditure, which amounted to at least 2000 lbs. of gold (£90,000), for the sake of the honour, and the Emperor was sometimes obliged to contribute from the treasury a large part of the money.\texthebrew{⁠}\footnote{Procopius, \emph{H. A.} 23.13 .

} Belisarius was consul in A.D. 535, and in the two following years no consul was elected, presumably because no one was willing to pay and the treasury could not afford the luxury. We can well imagine that there was much disappointment and discontent among the populace of the capital, and Justinian attempted to rescue the endangered institution by a legal curtailment of the expenses. The Praetorian Prefect, John of Cappadocia, had come forward to fill the consul ­ship for A.D. 538, perhaps on this condition, and a few days before the kalends of January the Emperor subscribed a law\texthebrew{⁠}\footnote{\emph{Nov.} 105, dated 537, December 28. It is addressed to Strategius, Count of the Sacred Largesses, presumably because the fisc had sometimes contributed to the expenses. Justinian says that the law is intended to secure that the consul­ship shall be perpetual and not beyond any suitable person's purse (\textgreek{\texthebrew{ὅ}πως \texthebrew{ἂ}ν διηνεκ\texthebrew{ὴ}ς μείν\texthebrew{ῃ Ῥ}ωμαίοις \texthebrew{ἅ}πασι δ\texthebrew{ὲ} το\texthebrew{ῖ}ς \texthebrew{ἀ}γαθο\texthebrew{ῖ}ς \texthebrew{ἀ}νδράσιν \texthebrew{ὑ}πάρχ\texthebrew{ῃ} βατή}). It provides that the consular festivities shall last for only seven days: January 1, inaugural procession and ceremony of investiture January 2, chariot races; January 3, theatrokynêgion (exhibition of wild beasts); January 4, combats of men with wild beasts; January 5, theatrical representations; January 6, chariot races; January 7, the consul lays down his office.

} which abbreviated the programme of consular spectacles, made it optional for the consul to distribute a largesse or not, but ordained that if there were a distribution it should be of silver not of gold.\texthebrew{⁠}\footnote{§ 2. The silver is to be in the form of miliaresia \textgreek{κα\texthebrew{ὶ} μήλοις κα\texthebrew{ὶ} καυκίοις} (cups) \textgreek{κα\texthebrew{ὶ} τετραγωνίοις κα\texthebrew{ὶ} το\texthebrew{ῖ}ς τοιούτοις.}} It is manifest that the permission to withhold the largesse was useless, as no consul could have ventured to face the unpopularity which such an economy would bring upon him. The people ought to be grateful to him, Justinian thinks, not grumble

p348
at this curtailment of the amusements and largesses to which they have been accustomed, for they are threatened with the alternative of enjoying neither one nor the other. He expressly exempted the Emperor from the provisions of the law.

The new regulations postponed the doom of the consul ­ship for just four years. Basilius was consul for A.D. 541,\texthebrew{⁠}\footnote{One leaf of the consular diptych of Anicius Faustus Albinus Basilius is preserved in Florence, and Meyer thinks that the second leaf may be identified with one in the Brera at Milan, with the inscription Et inlustris ex c. domesticorum patricius cons. ord. (\emph{Zwei ant. Elfenb.} pp74\texthebrew{‑}75).

} and he was the last private person to hold it. The practice of dating years officially by the consuls was not given up. During the rest of Justinian's reign the year was designated as ``such and such a year after the consul ­ship of Basilius.''\texthebrew{⁠}\footnote{A.D. 536 and 537, in which no consul was elected, had been designated as p. c. Belisarii and p. c. Bel. ann. II.

} Succeeding Emperors assumed the consular dignity in the first year of their reigns. But Justinian introduced a new system of dating state documents by three distinct indications, the consulate (or post-consulate ), the regnal year of the Emperor, and the indiction ( A.D. 537).\texthebrew{⁠}\footnote{\emph{Nov.} 47 (August 31). (\emph{E.g.} a law of March 545 is dated ``in the 18th year of Justinian, the 4th after the consul­ship of Basilius, the 8th indiction.'') Shortly afterwards in 538\texthebrew{‑}539, the practice came in of dating the issue of bronze coins by the number of the regnal year on the reverse.

} The innovation of using the regnal year as an official mark of time was perhaps suggested by the practice of the Vandal kings.\footnote{See Mommsen, \emph{Hist. Schr.} III.357.

}

The system of raising revenue in the later Roman Empire was so oppressive that there is perhaps no Emperor whom a hostile critic could not have made out a case for charging with a deliberate design to ruin his subjects. The lot of the provincials might have been tolerable if the ministers and governors and their hosts of subordinate officials had all been men of stainless integrity, but an incorruptible official seems to have been the exception. The laws show how the Emperors were always striving to secure a just and honest administration and imagining new devices to check corruption and oppression. In such endeavours Justinian was indefatigable, as his laws eloquently prove. But it was easy for an enemy to dwell on all the evils and abuses which

p349
existed, to represent them as due to his deliberate policy, and to ignore his remedial legislation or misinterpret its intention. This is the method of the author of the \emph{Secret History.} His statements as to the abuses and hardships and misery suffered by Justinian's subjects are borne out in general by Justinian's own statements in his laws, but the same laws disprove the historian's inferences as to the Emperor's intentions. Although, as has been observed, his policy of aggrandisement and the scale of his public expenditure placed a disastrous strain on the resources of his subjects, he was far from being indifferent to their welfare, and he fully understood that it was to the interest of the treasury that they should be protected from injustice and extortion. We have already seen some of his efforts in connexion with his reforms in the provincial administration. The fact remains, however, that he was inflexible in insisting on the regular exaction of legal dues and was less liberal and prudent than many of his predecessors in cancelling accumulated arrears, and remitting the taxation of provinces which had been devastated by hostile invasions.\footnote{Procopius, \emph{H. A.} 23

(where it is, however, acknowledged that a year's tribute was remitted to cities actually taken by an enemy). It is stated here that there was no remission of arrears throughout the thirty-two years of his rule from 518 to 550 (the author represents Justin's reign as virtually part of Justinian's). This is not accurate, as there was a remission in 522 (\emph{Nov.} 147, § 1). In 553, however, there was a general remission of all arrears to 554 inclusive (\emph{ib.}).

}

If we examine the principal charges of economic oppression which were preferred against him by his enemies, we shall find that the abuses which they stigmatise were for the most part not new inventions of Justinian but legacies from the past. There was nothing new, for instance, in the fact that the inhabitants of the provinces through which troops passed to the scene of war were bound to provide food for the soldiers and fodder for the horses , and to transport these supplies to the camps. Sometimes a province had not sufficient provisions and they had to be procured elsewhere. The system, which was known as coemption,\texthebrew{⁠}\footnote{\textgreek{Συνωνή}. There was another form of coemption of which Procopius

(\emph{H. A.} 22.17 \emph{sqq.})

complains and gives one example. One year when the harvest in Egypt was bad and the corn supply was insufficient for the needs of the people, the Praetorian Prefect bought up immense quantities in Thrace and Bithynia at low compulsory prices; the farmers were obliged to transport the corn to the capital, and were at a dead loss. This seems to have occurred in 545\texthebrew{‑}546, shortly before Peter Barsymes was deposed from the office of Prefect.

} lent itself to intolerable exactions,

p350
and Justinian in A.D. 545 issued a law to guard the interests of the inhabitants. It provided that they should be paid in full for all they furnished to the troops, and that no contributions in money should be demanded from them, and forbade them to give anything gratuitously or without a written receipt.\texthebrew{⁠}\footnote{\emph{Nov.} 130. Quartering soldiers in private houses was forbidden (§ 9).

Procopius, \emph{H. A.} 23.11 \emph{sqq.}

} Another burdensome institution was the epibole , which, it will be remembered, when lands fell out of cultivation, made, in certain cases, neighbouring landowners responsible for the taxes. Justinian maintained this principle, but he does not appear to have made it harsher than before, and he sought to guard against its abuse.\texthebrew{⁠}\footnote{\emph{Nov.} 128 (A.D. 545), §§ 7, 8. Here the persons responsible are described as \textgreek{ο\texthebrew{ἱ ὁ}μόδουλα \texthebrew{ἢ ὁ}μόκηνσα χωρία κεκτημένοι} (see above,

Chap. XIII p444 \emph{sq.} , where the general nature of the epibole is explained). It is expressly stated that they were not liable for arrears. Justin and Justinian relieved Church lands from liability to the epibole (see Cyril, \emph{Vita Sabae}, 294).

} It is probable, however, that in the oriental provinces during the Second Persian War the invasions of the enemy as well as the pestilence had caused the ruin of many proprietors, and that the application of the epibole was a frequent and serious grievance.\footnote{Procopius, \emph{ib.} 9 and 15\texthebrew{‑}16.

}

One tax is mentioned which seems to have been a novelty, and of which we can find no trace in the Imperial legislation. It was called the air-tax or sky-tax ( aërikon ), a name which suggests that it was a tax on high buildings, such as the insulae or apartment houses in cities. It was administered by the Praetorian Prefect and yield 3000 lbs. of gold (£135,000) a year to the treasury, while it is insinuated that the Prefects made much more out of it.\footnote{The only source is

Procopius, \emph{H. A.} 21.1 \emph{sqq.} , who states that it was an unusual tax and was so called \textgreek{\texthebrew{ὥ}σπεξτατορ \texthebrew{ἐ}ξ \texthebrew{ἀ}έρος \texthebrew{ἀ}ε\texthebrew{ὶ} α\texthebrew{ὐ}τ\texthebrew{ὴ}ν φερομένην}. In much later times we meet a tax of the same name (Leo VI \emph{Tactica}, XX.71; Alexius Comnenus,

\emph{Nov.} 27, § 4 ). Kalligas, improbably, explains the \textgreek{\texthebrew{ἀ}ερικόν} as a hearth-tax (\textgreek{καπνικόν}); such a tax would have produced a far larger sum. The discussion of Panchenko (\emph{O tainoi ist.} 149 \emph{sqq.}) throws little light on it, and he misinterprets Procopius, who says \textgreek{κα\texthebrew{ὶ} τα\texthebrew{ῦ}τα μ\texthebrew{ὲ}ν τ\texthebrew{ῷ} α\texthebrew{ὐ}τογράτροι \texthebrew{ἀ}ποφέρειν \texthebrew{ἠ}ξίουν, α\texthebrew{ὐ}τοι δ\texthebrew{ὲ} πλο\texthebrew{ῦ}τον βασιλικ\texthebrew{ὸ}ν περιεβάλλοντο ο\texthebrew{ὐ}δενί πόν\texthebrew{ῶ}}. He curiously refers \textgreek{α\texthebrew{ὐ}τοί} to the ministry of the comes larg., whereas it is the Praet. Prefects who are in question, and, explaining the words to mean that the proceeds of the tax were paid into the treasury and then paid out to the officials, he infers that the purpose of the tax was to supply the officials with sportulae (pp151, 153). But \textgreek{περιεβάλλοντο} refers to \textgreek{λ\texthebrew{ὴ}στείαις} in the preceding sentence, not to legal acquisition. Monnier supposed that the \textgreek{\texthebrew{ἀ}ερικόν} was a tax on houses (\emph{Étude de droit byz.} 508 \emph{sqq.}), and that it was so called from counting the openings, doors, and windows. But this is a far-fetched derivation and incongruous with the

(p351)
use of \textgreek{\texthebrew{ἀ}ήρ}; whereas ``aerial house'' would have been a natural Greek expression for ``sky-scraper.'' Moreover, it is clear from the amount of the yield that it cannot have been a tax on \emph{all} houses. Stein (\emph{Hermes}, LII.579) discusses it in connexion with \emph{Nov.} 43, and factories (\textgreek{\texthebrew{ἐ}ργαστήρια}) would probably have been liable to the tax.

}

p351
The decay of municipal life reached a further stage in the reign of Justinian, who describes its decline;\texthebrew{⁠}\footnote{\emph{Nov.} 38, \emph{Pref.}

} and increased interference on the part of the central government in the local finances seems to have been unavoidable. We saw how Anastasius took the supervision of the collection of taxes out of the hands of the decurions and appointed vindices , whose administration proved a failure. Justinian stigmatises them as pestilential and appears to have abolished them, though not entirely.\texthebrew{⁠}\footnote{\emph{Ib.} and \emph{Edict} 13.

} The rates, known as politika , which were imposed for municipal purposes and used to be altogether under the control of the local authorities, had already in the time of Anastasius been partly appropriated by the fisc. They were collected along with the other taxes, and were divided into two portions, of which one went to the treasury, the other to the cities. The same Emperor sometimes sent a special inspector to see that the necessary public works were carried out.\texthebrew{⁠}\footnote{\emph{C. J.} X.16.13 . Procopius

(\emph{H. A.} 26.6)

ascribes to Justinian the appropriation of \textgreek{πολιτικά} by the treasury, and although \emph{Nov.} 128 seems inconsistent with this charge, the evidence of \emph{Edict} 13 virtually bears it out. Cp. Maspéro in \emph{Archiv f. Papyrusforschung}, V.363 \emph{sqq.} (1913).

} In A.D. 530 Justinian placed the management of the public works, the local expenditure, and the control of the accounts in the hands of the bishops and the leading local dignitaries. But he reserved to himself the right of sending special accountants to exercise supervision.\texthebrew{⁠}\footnote{\emph{C. J.} X.39.4

discussores operum publicorum.

} These accountants ( discussores ) must be sent by his own personal mandate, and the local authorities are warned to recognise no one who comes with a mandate of the Praetorian Prefect. It would appear that the treatment of the politika as a due to the treasury had given the Praetorian Prefects and their officials additional opportunities of injustice and extortion, for the Emperor shows great concern to exclude any interference on the part of this ministry in the local administration. Some years later he committed to the provincial governors the general duty of seeing that the most necessary public works, such as repairs of bridges, roads, walls, harbours, were carried out, that the cities were properly provisioned with food, and that the

p352
accounts were duly audited. But they were to do this in person and not through subordinates.\footnote{See the Mandata principis to governors,

\emph{Nov.} 17, § 4 , A.D. 535; also \emph{Novv.} 24, § 3 ;

25, § 4 ;

26, § 4 ;

30, § 8 . In

\emph{Nov.} 128

(A.D. 545), however, the rights of the governor to interfere are carefully limited. He is to see that the portion of the \textgreek{πολιτικά} appropriated to the city is duly paid, that it is not diverted to improper purposes by the inhabitants, that the bishop and the curials duly elect a ``father of the city,'' a corn-commissioner (\textgreek{σιτώνης}), and other functionaries, and that the accounts are regularly audited (§ 16).

}

But the proceeds of the local taxes, diminished by the claims of the treasury, were frequently insufficient to defray the municipal upkeep, especially when exceptional expenses were incurred in consequence of earthquakes, for instance, or hostile invasions. In such cases, the matter was referred to the Emperor, who sometimes advanced large sums from the treasury to assist a city which had been visited by some grave disaster. But as a rule the method was to levy a special tax known as a \emph{description}, which was assessed in proportion to the amount of the land-tax. That this tax gave rise to abuses is shown by the fact that Justinian forbids governors to impose \emph{descriptions} on towns during their progresses through the provinces.\footnote{Imperial, senatorial, and church lands were exempt from these \textgreek{διαγραφαί} ( \emph{C. J.} XII.1.7 ; \emph{Nov.} 131, § 5), which Procopius

(\emph{H. A.} 23.17\texthebrew{‑}19)

describes as one of the principal burdens falling on the provincials. They had at one time been imposed on other classes as well as on the curials and landowners (cp. \emph{C. J.} XI.1.2 ). These \textgreek{διαγραφαί} extraordinariae are to be distinguished from \textgreek{διαγραφαί} lucrati\textgreek{υων}, º taxes on estates which changed hands. Church lands were exempted from this burden also (\emph{Nov.} 131, \emph{ib.}).

}

The decline of the municipal resources became more marked from A.D. 543 onwards in consequence of the ravages of the Plague, and it led to the decay of the liberal professions. The cities, forced to economise, withdrew the public salaries which they had hitherto paid to physicians and teachers. Advocates are said to have suffered because people were so impoverished that they could not afford the luxury of litigation. Some towns could not defray the cost of lighting the streets, and public amusements, theatres, and chariot races were curtailed.\footnote{Procopius, \emph{H. A.} 26.5\texthebrew{‑}11 .

}

On the whole, although he made alterations in detail, which were chiefly designed to check the abuse of their authority by officials and to diminish the power of the Praetorian Prefect, Justinian preserved the existing financial system in all its essential principles. He did not make it worse, and he endeavoured to arrest the progress of municipal decay. The ruin

p353
wrought by the inroads of the Persians and of the northern barbarians, and the effects of the Plague made, however, in many parts of the Empire the burdens more grievous than ever, and the Emperor may be blamed for not seeing that a fundamental and drastic reform of the whole system of taxation was demanded in the interest of the public welfare. The retrenchments which he might well have made in the early years of his reign, instead of embarking on large schemes of conquest and spending exorbitant sums on buildings, were almost impossible subsequently when he was involved simultaneously in the wars with Persia and the Ostrogoths. The measure to which he was forced in A.D. 552 of cancelling all arrears of taxation is an eloquent indication of the plight of the provinces, for his previous policy shows that he would not have forgone a fraction of the treasury's legal dues unless absolute necessity had compelled him.\footnote{Stein (\emph{Studien}, 143 \emph{sqq.}) has attempted to prove that the total revenue of the State in this reign cannot have much exceeded 7,000,000 solidi (£4,375,000), and that this was enough to cover the outgoings. He starts with a fallacious assumption, and leaves out of account many departments of expenditure. Even if Justinian had no more than 21,000,000 subjects, his conclusion would imply that the taxation only came to about 4s. 2d. per head of the population annually. The errors have been pointed out by Andreades in \emph{Revue des études grecques}, XXXIV. No. 156 (1921). The taxation accounts of Antaeopolis (\emph{Pap. Cairo}, I.67057) give information which may ultimately help us to estimate the contribution of Egypt to the revenue. Cp. the accounts of the village of Aphrodito (\emph{ib.} 67058) and those of the rich proprietor Ammonius (II.67138\texthebrew{‑}67139).

}

The conquest of Africa enabled him to make large additions to the Imperial estates,\texthebrew{⁠}\footnote{Procopius, \emph{B. V.} II.14. This policy led to sedition among the soldiers, who expected that the land would be distributed among themselves.

} but in the eastern provinces also the Private Estate and the crown domains appear to have been gradually and considerably extended, at the expense of adjacent private property. We have not much information as to the methods and pretexts by which this was effected, but about fifteen years after the death of Justinian complaints reached the Emperor Tiberius from almost all the provinces as to the unjust appropriation of private property by the officials of Imperial estates.\texthebrew{⁠}\footnote{Zachariä v. Lingenthal, \emph{Jus Graeco-Rom.} III, \emph{Nov.} 12.

} That this form of robbery was practised in Justinian's reign we have other evidence.\texthebrew{⁠}\footnote{Aetherius, curator domus divinae under Justinian, was brought to account under Justin II for his acts of plunder (\textgreek{τάς τε τ\texthebrew{ῶ}ν ζώντων τ\texthebrew{ῶ}ν τε τελευτώντων τ\texthebrew{ὰ}ς ο\texthebrew{ὐ}σίας ληιζόμενος,}Evagrius, V.3 ). The appropriations of Anatolius, another curator, are described by

Agathias, V.4 . Cp. Procopius, \emph{H. A.} 12.12 .

} In some cases on

p354
the death of a proprietor his will or the claims of his legal heirs were set aside and his possessions acquired by the fisc. It is only too likely that many unjust acts were deliberately committed by the help of legal quibbles, but we need not pay serious attention to the allegations that Justinian forged wills or acts of donation in order to acquire the possessions of rich subjects.\texthebrew{⁠}\footnote{Procopius, \emph{ib.} 1\texthebrew{‑}11, gives six instances with the names of the persons. We can accept them as cases in which the fisc inherited, but the charge of forgery is evidently asserted merely on hearsay, as indeed in one case the writer lets out (\textgreek{διαθήκην \texthebrew{ἥ}νπερ ο\texthebrew{ὐ} παρ\texthebrew{’ ἐ}κείνου ξυγκε\texthebrew{ῖ}σθαι διατεθρύλληται}). Against these cases and that of the inheritance of Anatolius

(\emph{H. A.} 29.17 \emph{sqq.})

we may set the generosity of Justinian in dealing with the daughters of Eulalius (John Mal. XVIII.439).

} Nor does the less improbable charge that he misused criminal justice for the purpose of confiscating property seem to be borne out by the facts. For instance, he restored, so far as he was able, to the disloyal senators their properties which had become forfeit to the State after the Nika rebellion. And in the later years of his reign, at a time when fiscal necessities were urgent, he abolished confiscation as a penalty for ordinary crimes.\footnote{\emph{Nov.} 134, § 13, A.D. 556.

}

The treatment of the private estates had varied, as we have seen, from time to time since the days of Septimius Severus. The last innovation had been that of Anastasius who, instead of incorporating recently confiscated lands in the res privata , had instituted a new minister, the Count of the Patrimony. This had simply meant a division of administration, for the Patrimony as well as the Private Estate was appropriated to public needs, not to the Emperor's private use. Justinian made yet another change. The Patrimony disappears,\texthebrew{⁠}\footnote{Stein (\emph{Studien}, 174 \emph{sqq.}) has established this, or at least made it highly probable. (1) There is no mention of the comes patr. in

\emph{C. J.} VII.37.3

(A.D. 531), where he ought to appear if he still existed, or in \emph{Nov.} 22 (A.D. 536). (2) In Justinian's laws sacr. patrimonium is sometimes used as an equivalent of s. larg.; this would have been confusing if the patrimony existed. (3) While we hear no more of the comes patr., we begin to hear a great deal of the divinae domus and nostri curatores per quos res divinarum domuum aguntur

(\emph{C. J.} \emph{ib.}) . \textgreek{Τ\texthebrew{ὸ} πατριμόνιον} in

Procopius, \emph{H. A.} 22.12 , Stein identifies with the estates in Sicily which were under the com. patr. who was instituted by Theoderic, and continued to function under Justinian. The passage in John Lydus, \emph{De mag.} II.27 , on \textgreek{\texthebrew{ὁ} λεγόμενος πατριμώνιος} of Anastasius, would certainly by itself suggest that the com. patr. still existed during Justinian's reign, but it cannot be pressed in view of the other evidence.

} and the domains which composed it are placed under the management of Curatores ( curatores divinae domus ). We do not know exactly what was involved in this change; more perhaps than a mere change of

p355
name. The domus divina was the patrimony,\texthebrew{⁠}\footnote{The equation will be found in Procopius, \emph{B. G.} I.4.1 and 6.26.

} and the Curator, subordinate to whom were the curators of the several domains,\texthebrew{⁠}\footnote{Stein asserts that the curators of the particular domains were independent and had no superior. This certainly was not the case under Justinian and Justin II. The title curator dominicae domus, without any limitation in \emph{C. J.} \emph{ib.}, and \textgreek{κουράτωρ τ\texthebrew{ῶ}ν ο\texthebrew{ἰ}κι\texthebrew{ῶ}ν} in \emph{Nov.} 148, § 1, point to a central controller. Anatolius, for instance, held this post

(Agathias, V.3) .

} discharged functions of the comes patrimonii . But the Curator seems to have been a court official rather than a State official, and Justinian's aim may have been to assert the principle that the administration of the patrimonial domains, consisting of confiscated properties, was the Emperor's own personal affair.

The policy of this reign in regard to trade is not very clear, and it is difficult to say how far it was responsible for the economic crises which arose and compelled the intervention of the government. Some changes were made in the custom-house arrangements at Constantinople. Hitherto the custom duties had been collected when ships reached the harbour of the capital. But there were posts of observation in the Hellespont and the Bosphorus to make sure that the public regulations were not evaded. An officer was stationed at Abydos to see that no vessel with a cargo of arms entered the straits without Imperial orders, and that no vessel passed through to the Aegean without papers duly signed by the Master of Offices. This officer was paid by fees levied on the owners of the ships.\texthebrew{⁠}\footnote{The source is

Procopius, \emph{H. A.} 25.2 \emph{sqq.}

Cp. the inscription of Abydos (probably belonging to the time of Anastasius) published in the \emph{Mittheilungen des deutschen arch. Inst.} (Athen), IV.307 \emph{sqq.} (1879).

} Another officer was posted at Hieron, at the northern issue of the Bosphorus, to examine the cargoes of craft sailing into the Euxine and prevent the export of certain wares which it was forbidden to furnish to the peoples of southern Russia and the Caucasian regions.\texthebrew{⁠}\footnote{For quae res exportari non debeant see

\emph{C. J.} IV.41

(wine, oil, lard, and arms are mentioned).

} He was paid a fixed salary and received no money from the shipowners. Justinian's innovation was to convert both these stations into custom-houses for imports, of which the officials were salaried but also received an additional bonus proportionate to the amount of the duties which they collected.\footnote{The title of the officer at Hieron was in later times at least \textgreek{κόμης το\texthebrew{ῦ Ἱ}ερο\texthebrew{ῦ} και το\texthebrew{ῦ} Πόντου}. John Mal. (XVIII.432) describes him as \textgreek{κόμης στεν\texthebrew{ῶ}ν τ\texthebrew{ῆ}ς Ποντικ\texthebrew{ῆ}ς θαλάσσης}, and dates his institution the A.D. 528\texthebrew{‑}529. For seals of commerciarii of Abydos see Schlumberger, \emph{Sig. byz.} 196 \emph{sqq.}

}

p356
But the tolls on exports were still collected at Constantinople, and these charges are said to have been so onerous that they forced the merchants to raise the prices of their wares enormously. But we have no information as to the tariff.\footnote{That the customs levied at Abydos and Hieron were only on imports to the capital, and those levied at Constantinople were on the cargoes of outgoing ships is my interpretation of the passage in Procopius, which is not very clear. The supervisor of the latter was probably entitled comes commerciorum (cp. Panchenko, \emph{op. cit.} 155), and the office was held by a Syrian named Addaeus (who was after, in 551, Pr. Prefect of the East, \emph{Nov.} 129).

}

Justinian is accused of having made necessaries as well as luxuries dearer not only by exorbitant duties on merchandise \texthebrew{—} a charge which we cannot control \texthebrew{—} but also by establishing ``monopolies'' for the benefit of the government.\texthebrew{⁠}\footnote{The motive of the monopoly of arms (\emph{Nov.} 85) was not financial. The sale of arms to private persons was forbidden.

} The restrictions which he imposed in the silk trade were considered when we surveyed the commercial relations of the Empire with foreign lands, and we saw that, though his policy in some respects was not happy, he deserves credit for his efforts to solve a difficult problem. It is far from clear how he made an income of 300 lbs. of gold from the sale of bread in the capital, as he is alleged to have done.\texthebrew{⁠}\footnote{Procopius, \emph{H. A.} 20.1 \emph{sqq.} , and 26.19 \emph{sqq.} He says that the bread was not only dearer but of worse quality. His \textgreek{τριπλάσιονα τιμήματα} agrees with \emph{Nov.} 122.

} Whatever new regulations were introduced cannot be described as a monopoly in the proper sense of the term. It is, however, certain that in the years after the Plague the price of labour rose considerably, and in A.D. 544 the Emperor issued an edict to re-establish the old prices. ``We have learned,'' he says, ``that since the visitation of God traders and artisans and husbandmen and sailors have yielded to a spirit of covetousness and are demanding prices and wages two or three times as great as they formerly received. We therefore forbid all such to demand higher wages or prices than before. We also forbid contractors for building and for agricultural and other works to pay the workmen more than was customary in old days.'' A fine of three times the additional profit was imposed on those who transgressed the edict. Justinian evidently assumes that there was no good reason for the higher rates. Unfortunately we have no information as to the effects of the edict, in which the interests of the customers are solely considered.\texthebrew{⁠}\footnote{\emph{Nov.} 122.

} That there was a fall of credit even before the Plague

p357
is indicated by measures which were taken to protect the interests of the power ­ful corporation of bankers against their debtors.\footnote{\emph{Nov.} 136 (A.D. 535); \emph{Edict} 9; \emph{Edict} 7 (A.D. 542).

}

It would probably be rash to infer from the tendency of interest on loans to rise since A.D. 472 that trade had been tending to decline.\texthebrew{⁠}\footnote{See Billeter, \emph{Gesch. des Zinsfusses} (pp219, 317).

} The ordinary commercial rate of interest in Justinian's reign was 8 per cent.\texthebrew{⁠}\footnote{This may be illustrated from papyri, \emph{e.g.} \emph{Pap. Cairo}, II. No. 67126. Cp. \emph{C. J.} IV.32.26 ; and \emph{Nov.} 110, repeating \emph{Nov.} 106.

} On good securities money could be borrowed at 5 or 6 per cent. Justinian paid attention to the question of interest and reduced the maximum 12 per cent, which had hitherto been legal, to 8, except in the case of maritime ventures, where 12 was allowed. But 8 was allowed only in the case of traders, and 6 was fixed as the maximum for loans between private persons. In the case of money advanced to peasants he enacted that only 4 per cent should be charged, and he forbade senators of illustrious or higher rank to exact more than 4 per cent.

The coinage of Justinian's reign, which is exceptionally abundant, may be taken as testifying to a flourishing condition of commerce. The curious statement in the \emph{Secret History} that he depreciated the gold coinage has no confirmation in the evidence of the extant nomismata .\texthebrew{⁠}\footnote{\emph{H. A.} 22.38 . In 25.12 , the Emperor is blamed for the practice of the money-changers to give only 10 folles for a nomisma instead of the normal 210.

} The number of Imperial mints was increased, not only in consequence of the conquest of Africa and Italy, but also by the establishment of a new centre in the East.\texthebrew{⁠}\footnote{Under Anastasius there were only three mints, Constantinople, Nicomedia, and Antioch; under Justin I Thessalonica and Cyzicus were added. Under Justinian money was coined also at Alexandria and Cherson; and in the west, at Carthage, in Sicily (Catana?), at Rome and Ravenna. See Wroth, \emph{Imp. Byz. Coins}, I. xv. \emph{sqq.}

} The minting of gold was confined to Constantinople, and silver was issued only there and in Carthage.

If Justinian was blamed for his expenditure on wars, for his extravagance in building, for the large sums with which he bought off the hostilities of the northern barbarians, he was blamed no less for his economies. Some of these may have been short-sighted and unwise, for instance the curtailments of the

p358
public Post, to which attention has already been called, and the reduction of the intelligence department.\texthebrew{⁠}\footnote{Procopius

(\emph{H. A.} 30.12 \emph{sqq.})

says that the secret service ceased to exist, but this is assuredly an exaggeration. These spies (\textgreek{κατάσκοποι}) used to penetrate into the palace of the Persian king as merchants or on other pretexts. Procopius ascribes the successes of Chosroes to the fact that he improved his secret service, while Justinian refused to spend money on his.

} But much greater dissatisfaction was caused by economies which to an impartial posterity seem unquestionably justified. Such, for instance, were the abolition of the consul ­ship, which had ceased to perform any useful function, the reduction of expenses on public amusements, the discontinuance of the large distribution of corn which, since the time of Diocletian, had pauperised the proletariate of Alexandria.\texthebrew{⁠}\footnote{\emph{Ib.} 26.40 \emph{sqq.} This was done when Hephaestus was Augustal prefect; he is said to have enriched himself and the treasury by monopolising in his own hands the sale of all provisions in the city. Hephaestus is probably the same person who was Pr. Pref. of the East \emph{c.} A.D. 550\texthebrew{‑}551 (John Lyd. \emph{De mag.} III.30) .

} Another economy was the diminution of the pensions of the officials serving in the central bureaux, which had hitherto cost the treasury about 10,000 lbs. of gold (£450,000), a measure which must have been extremely unpopular.\footnote{\emph{H. A.} 24.30 \emph{sqq.}

\textgreek{\texthebrew{ἐ}ν Βυζαντί\texthebrew{ῷ}} shows that the provincial bureaux are not included. We have no means of judging whether the pensions were excessive, and the reduction may not have been considerable; for we cannot trust Procopius when he says \emph{\textgreek{τούτων α\texthebrew{ὐ}το\texthebrew{ὺ}ς \texthebrew{ἀ}ποστερήσας σχεδόν τι \texthebrew{ἁ}πάντων}}, as hyperbole is a note of the \emph{Secret History}. But the complaints of John Lydus (\emph{De mag.} III.67) bear out the statement.

}

The parsimony of Justinian which seems most open to criticism was in the treatment of the army. He reduced its numbers and tried to reduce the expenses on its upkeep. The names of the dead remained on the lists, new soldiers were not recruited, and there was no promotion. The old practice of Imperial donatives every five years was discontinued. Pay was always in arrears, and was often refused altogether on various pretexts. No sooner had a soldier received his pay than the logothete appeared with a bill for taxes. We are told that Justinian appointed the worst sort of men as logothetes, and they received a commission of one-twelfth on all they managed to collect. After the peace of A.D. 545 there appears to have been a considerable reduction of the frontier forces in the East.\footnote{The decline of the army is the subject of

\emph{H. A.} 24 . Procopius speaks as if the limitanei were finally abolished after A.D. 545, but so far as his criticisms have a foundation they apply only to the East (see Hartmann, \emph{Byz. Verwaltung in Italien}, p151; Panchenko, \emph{op. cit.} 117 \emph{sqq.}). Maspéro has made it probable that the total of the forces stationed in Egypt was from 29,000 to 30,000

(p359)
men, of whom about 5000 were in Libya and Tripolitana. Only two or three thousand of these were \textgreek{καστρησινοί}, limitanei (\emph{Org. mil. de l'Égypte}, 117). Maspéro thinks that when Agathias

(V.13)

gives the total strength of the army as 150,000, he does not include the limitanei (\emph{ib.} 119). Justinian appears to have formed a new corps of Palace guards called Scribones; it is at least in his reign that we first hear of them

(Agathias, III.14) . They were often employed on special missions in the provinces. The Scholarian guards (3500 in number) had now ceased to have any military significance; they were employed purely for parade purposes. Young men who had a little money and desired to lead an idle life in splendid uniform invested it in purchasing a post in the guards, and the high pay was a satisfactory annuity for their capital (cp. Agathias, V.15 ). Procopius says that in Justin's reign 2000 ``supernumeraries'' were added in order to obtain the entrance fees, and that Justinian on his accession disbanded them without compensation (\emph{ib.} 20).

}

p359
That the efficiency of Justinian's administration degenerated in the latter part of his reign there is every sign. After the deaths of Theodora and Germanus he concentrated his attention more and more on theology \texthebrew{—} in caelum mens omnis erat \texthebrew{—} and was inclined to neglect public affairs and postpone decisions. When he died it was probably the general opinion that it was high time for a younger man to take the helm and restore, above all, the financial situation. For the fisc was exhausted.\footnote{See Corippus, \emph{In laud. Iust.} II.260 \emph{sqq.}:

plurima sunt vivo nimium neglecta parente

unde tot exhaustus contraxit debita fiscus

reddere quae miseris moti pietate paramus.

quod minus ob senium factumve actumve parentis

tempore Iustini correctum gaudeat orbis.

nulla fuit iam cura senis, etc.}

\xchapter{Chapter XXII}

Short URL for this page:

bit.ly/BURLAT22

Theoretically the Emperors were as completely competent to legislate in all religious as in all secular affairs. How far they made use of this right was a question of tact and policy. No Emperor attempted to order the whole province of sacred concerns. Questions of ritual, for instance, were left entirely to the clergy, and the rulers, however bent they might be on having their way in questions of doctrine, always recognised that doctrine must be decided by ecclesiastical councils. The theory, which was afterwards to prevail in western Europe, of a trenchant separation between the spiritual and temporal powers was still unborn, and ecclesiastical affairs were ordered as one department of the general civil legislation. In framing laws concerning the organisation of the Church, it was a matter of course that the Patriarch of Constantinople should be consulted, but it is significant that such contributions were often addressed not to the Patriarch or the bishops, but to the Praetorian Prefect of the East whose duty it was to make them publicly known throughout the Empire.

Justinian took his responsibilities as head of the Church more seriously than any ever had hitherto done, and asserted his authority in its internal affairs more constantly and systematically. It was his object to identify the Church and State more intimately, to blend them, as it were, into a single organism, of which he was himself the controlling brain. We must view in this light his important enactment that the Canons of the four great Ecumenical Councils should have the same validity as

p361
Imperial laws.\texthebrew{⁠}\footnote{\emph{C. J.} I.3.44 ; \emph{Nov.} 131.

} And we can see in his legislation against heretics and pagans that he set before himself the ideal of an Empire which should be populated only by orthodox Christians. He determined ``to close all the roads which lead to error and to place religion on the firm foundations of a single faith,''\texthebrew{⁠}\footnote{Procopius, \emph{Aed.} I.1 .

} and for this purpose he made orthodoxy a requisite condition of citizen ­ship. He declared that he considered himself responsible for the welfare of his subjects, and therefore, above all, for securing the salvation of their souls; from this he deduced the necessity of intolerance towards heterodox opinions.\texthebrew{⁠}\footnote{\emph{C. J.} I.5.18 .

} It was the principle of the Inquisition. None of his predecessors had taken such a deep personal interest in theology as Justinian, and he surpassed them all in religious bigotry and in the passion for uniformity.

The numerous ecclesiastical laws of Justinian, which do not concern doctrine or heresy, deal with such topics as the election of bishops, the ordination of priests and deacons, the appointment of the abbots of monasteries, the management of Church property, the administration of charitable institutions, such as orphanages, hostels, and poorhouses, the privileges and duties of the clergy.\texthebrew{⁠}\footnote{\emph{C. J.} VIII. titles 1\texthebrew{‑}13 , and 33 Novels are devoted to ecclesiastical laws.

} We learn from this legislation the existence of various abuses, simony,\texthebrew{⁠}\footnote{To remedy bribery at episcopal elections, a definite tariff of fees was fixed, to be paid by the new bishop to those who ordained him and their assistants, according to the income of the see. \emph{E.g.} if the income was not less than £1350, the total of the fees amounted to £250 (\emph{Nov.} 123, § 3).

} for instance, an illiterate priests and bishops. Little regard was shown for freedom in the restrictive enactments which were intended to prevent bishops from neglecting their sees;\texthebrew{⁠}\footnote{\emph{Ib.} § 9. Bishops were forbidden to leave their sees without permission of their metropolitan, and they could only visit Constantinople with leave of the Patriarch or by order of the Emperor.

} and the clergy were strictly forbidden to indulge in the pastimes of attending horse-races or visiting the theatres.\footnote{\emph{C. J.} I.4.34 .

}

But the most important feature in this section of Justinian's legislation is the increasing part which the bishops were called upon to play in civil and social administration. They were gradually taking the place of the defensores civitatis , and probably served as a more power ­ful check on unjust or rapacious provincial

p362
governors.\texthebrew{⁠}\footnote{\emph{Nov.} 86. The defensores (\textgreek{\texthebrew{ἔ}κδικοι}) still existed (\emph{ib.} and

\emph{Nov.} 15 ).

} In certain matters of business they could act instead of the governor himself.\texthebrew{⁠}\footnote{\emph{C. J.} I.4.21, and 31 . Other examples of the use which the government made of bishops will be found in this title.

} They were expected to take part in overseeing the execution of public works, to take charge of the rearing of exposed infants, to enforce the laws against gambling. When Justinian issued a law against the constraint of any woman, slave or free, to appear on the stage, it was to the bishops that he addressed it, and they were charged to see that it was enforced, even against a provincial governor.\texthebrew{⁠}\footnote{\emph{Ib.} 33.

} It was on their vigilance that the government chiefly replied for setting the law in motion against heretics.

On any theory of the relations of Church and State, it would have been reasonable that, as the State granted to the bishops judicial and administrative authority and to the clergy special privileges, it should insist on their fulfilling certain qualifications and should lay down rules binding on the clerical order. It was not so clear why the Emperor should consider it his business to regulate the conduct of monastic institutions,\texthebrew{⁠}\footnote{The chief regulations will be found in \emph{Nov.} 5, 123, 133.

} seeing that they discharged no function in the political organisation and were established only for those who desired to escape the temptations, the troubles, and the labours of social life. He justifies his action in one of his laws, where he expresses the superstitious belief that the prosperity of the State could be secured by the constant prayers of inmates of monasteries. ``If they, with their hands pure and their souls bare, offer to God prayers for the State, it is evident that it will be well with the army, and the cities will prosper and our land will bear fruits and the sea will yield us its products, for their prayers will propitiate God's favour towards the whole State.''\texthebrew{⁠}\footnote{\emph{Nov.} 133, § 5 (A.D. 539).

} The great pestilence and numerous earthquakes were a commentary on the Emperor's faith, which he was not likely to take to heart.

It has been observed that his legislation ``became in the Byzantine Empire the true foundation of monastic institutions.''\texthebrew{⁠}\footnote{Diehl, \emph{Justinien}, 503. For a general picture of monasticism in this age, see the whole chapter (499 \emph{sqq.}).

} During his reign the number of monasteries enormously increased,\texthebrew{⁠}\footnote{In A.D. 536 there were sixty-seven monasteries for men in Constantinople and its suburbs (Mansi, VIII.1007 \emph{sqq.}). For the plan of the monastery of Saint Simeon, between Antioch and

(p363)
Aleppo, see Vogüé, \emph{Syrie centrale}, Pl. 139\texthebrew{‑}150; for that of Tebessa in Byzacena, Diehl, \emph{L'Afrique byzantine}, Pl. XI.

} and in later times the growth of these parasitic institutions

p363
multiplied more and more. Rich men and women vied with each other in adding to their number.

In Syria and Palestine monastic houses were particularly numerous and power ­ful, and the oriental monks enjoyed and merited a higher reputation than any others for extreme asceticism. A certain number of cells were reserved in the Syrian convents for those who, not content with the ordinary rule and desiring a more rigorous mortification of the flesh, yet preferred the shelter of a monastery to the life of the recluses who lived isolated in deserts or mountains. The historian, John of Ephesus, has left us a gallery of contemporary eastern monks, who were distinguished by their piety or eccentricities, and his portraits are sufficiently repulsive. They exercised an extraordinary influence not only over the common people, but even at court, and could indulge with impunity in the most audacious language in the Imperial presence. For instance, when proceedings were taken against the Monophysites in Egypt in A.D. 536, Maras, a heretical anchoret of the most savage manners, arrived at Constantinople for the purpose of loading the Emperor and Empress with vituperation. Admitted to an audience he used language which would have been almost incredible if it had been flung at persons of low degree; his panegyrist declines to reproduce it. But Emperor and Empress, if astonished, did not resent the insults of the ragged hermit; they said that he was a truly spiritual philosopher.\footnote{John Eph. \emph{Comm. de b. or.} c37. Maras was a native of the region of Amida, but he set up a cell in the Thebaid. Another Monophysite monk, Zooras, about the same time used similarly audacious language to Justinian and made the Emperor very angry (\emph{ib.} c2).

}

One important change in diocesan administration was introduced by Justinian. He divided the ecclesiastical vicariate of Illyricum into two parts for the strategy of increasing the prestige and importance of Justiniana Prima, as he had renamed the town of Scupi, which was close to his own birthplace. Having first transferred the seat of the Praetorian Prefect of Illyricum from Thessalonica to Justiniana, he resolved to increase to prestige of his home\texthebrew{⁠}\footnote{Nostram patriam augere cupientes, \emph{Nov.} 11 (A.D. 535); 131, 3 (A.D. 545).

} by making it also a great ecclesiastical centre. The bishop of Justiniana was raised to the rank not only of a metropolitan but of an archbishop, and his diocese

p364
corresponded to the civil diocese of Dacia, with its seven provinces. He was independent of Thessalonica, but the see of Thessalonica retained its authority over the rest of Illyricum, the diocese of Macedonia. This arrangement, which was carried out with the consent of the Pope, did not change the position of ecclesiastical Illyricum as a vicariate of the Roman see. The only difference was that the Pope was now represented by two vicars.\footnote{For details see Zeiller, \emph{Origines chrét.} 385 \emph{sqq.}; Duchesne, \emph{Églises sép.} Scupi (Üsküb) remained a metropolitan see till 1914.

}

The measures which Justinian adopted to suppress heresy were marked by a consistency and uniformity which contrast with the somewhat hesitant and vacillating policy of previous Emperors. Laying down the principle that ``from those who are not orthodox in their worship of God, earthly goods should also be withheld,''\texthebrew{⁠}\footnote{\emph{C. J.} I.5.12, 5

(A.D. 527).

} he applied it ruthlessly. Right belief was made a condition for admission to the service of the State, and an attestation of orthodoxy from three witnesses was required.\texthebrew{⁠}\footnote{\emph{Ib.} I.4.20 .

} Heretics were debarred from practising the liberal professions of law and teaching.\texthebrew{⁠}\footnote{\emph{Ib.} I.5.12, and 18 .

} But Justinian went much further in the path of persecution. He deprived heretics of the common rights of citizen ­ship. They were not allowed to inherit property; their testamentary rights were strictly limited; they could not appear in court to bear witness against orthodox persons. On the other hand, they were liable to the burdens and obligations of the curiales.\texthebrew{⁠}\footnote{\emph{Ib.} I.5.12 ; \emph{Nov.} 45.

} The spirit of the Imperial bigot is shown by a law which deprived a woman, if she belonged to a heretical sect, of her legal rights in regard to her dowry and property. The local priests and officials were to decide whether she was orthodox, and attendance at Holy Communion was to be regarded as the test.\texthebrew{⁠}\footnote{\emph{Nov.} 109 (A.D. 541).

} Here we have a foretaste of the Inquisition.

It is noteworthy that the sect of the Montanists in Phrygia was singled out for particularly severe treatment. But the penalty of death was inflicted only on two classes, the Manichaeans, whom the government had always regarded as the worst enemies of humanity, and heretics who, having been converted to the

p365 true creed, relapsed into their errors.\texthebrew{⁠}\footnote{Cp. esp.

\emph{Cod. J.} I.5.20 , and Procopius, \emph{H. A.} 11 .

} Perhaps these severe laws were not executed thoroughly or consistently, but we have a contemporary account of a cruel persecution of Manichaeans, which occurred perhaps about A.D. 545.\footnote{John of Ephesus, \emph{H. E.} Part II (Nau), p481. The date is uncertain, but the notice precedes events dated to the nineteenth year of Justinian.

}

Many people adhered to the deadly error of the Manichaeans. They used to meet in houses and hear the mysteries of that impure doctrine. When they were arrested, they were taken into the presence of the Emperor who hoped to convert them. He disputed with them but could not convince them. With Satanic obstinacy they cried fearlessly that they were ready to face the stake for the religion of Manes and to suffer every torture. The Emperor commanded that their desire should be accomplished. They were burned on the sea that they might be buried in the waves, and their property was confiscated. There were among them illustrious women, nobles, and senators.

The most important of all the heretical sects, the Monophysites, were hardly affected by the general laws against heretics. Their numbers and influence in Egypt and in Syria would have rendered it impossible to inflict upon them the disabilities which the law imposed on heretics generally, and they were protected by the favour of the Empress. Moreover, the Emperor's policy vacillated; he was engaged throughout his reign with doctrinal questions arising from the Monophysitic controversy, and the position of the Monophysites will most conveniently be considered in that connexion.

The Jews and Samaritans were subject to the same disabilities as heretics.\texthebrew{⁠}\footnote{\emph{C. J.} I.5.12, 17, 18 .

} This severity was followed by the destruction of the Samaritan synagogues, and a dangerous revolt broke out in Samaria in the summer of A.D. 529.\texthebrew{⁠}\footnote{John Mal. XVIII p445; Procopius, \emph{ib.}

} Christians were massacred; a brigand named Julian was proclaimed Emperor; and the rising was bloodily suppressed.\texthebrew{⁠}\footnote{Malalas says that 20,000 were slain and 20,000 (of both sexes) given to the Saracens, who assisted the duke of Palestine in the war, to be sold into slavery. Procopius says that 100,000 were reported to have perished.

} The desperate remnant of the people then formed a plan to betray Palestine to the Persians,\texthebrew{⁠}\footnote{John Mal. XVIII p455, states that 50,000 fled to Persia and offered their aid to Kavad.

} but their treachery appears to have had no results. Twenty years later, at the intercession of Sergius, bishop of Caesarea, and his assurance that the Samaritans had been converted from their

p366
evil ways and would remain tranquil, the Emperor removed some of the civil disabilities which he had imposed.\texthebrew{⁠}\footnote{\emph{Nov.} 129 (A.D. 551).

} But the hopes of Sergius were not realised. Samaritans and Jews joined in a sanguinary revolt at Caesarea, and murdered Stephanus, the proconsul of Palestine.\texthebrew{⁠}\footnote{John Mal. \emph{ib.} p487.

} Their ringleaders were executed, but the Samaritans were refractory and abandoned the pretence of having been converted to Christianity. The civil disabilities which had been imposed on them by Justinian were renewed by his successor.\texthebrew{⁠}\footnote{\emph{Nov.} 144 .

} The Samaritan troubles are a black enough page in the history of persecution.

The Jews were treated less harshly. Though the lawgiver regarded them as ``abominable men who sit in darkness,'' and they were excluded from the State-service , they were not deprived of their civil rights. Justinian recognised their religion as legitimate and respectable so far as to dictate to them how they should conduct the services in their synagogues.\texthebrew{⁠}\footnote{\emph{Nov.} 146 (A.D. 553).

} He graciously permitted them to read aloud their Scriptures in Greek or Latin or other versions. If Greek was the language they were enjoined to use the Septuagint, ``which is more accurate than all others,'' but they were allowed to use also the translation of Aquila.\texthebrew{⁠}\footnote{Aquila, a native of Pontus, converted to Judaism (hence called by Justinian \textgreek{\texthebrew{ἀ}λλόφυλος}), published his translation in the second century A.D. His aim was to produce a version more literal and accurate than the Septuagint; it was so literal that it was often obscure.

} On the other hand, he strictly forbade the use of the ``Deuterosis,'' which he described as the invention of uninspired mortals.\texthebrew{⁠}\footnote{It is uncertain what precisely the Deuterosis means. The term occurs several times in Jerome. In his \emph{Comm. in Matth.} c22 (\emph{P. L.} XXVI p165) he explains \textgreek{δευτερώσεις} as the traditions and observations of the Pharisees. S. Krauss in

an art. on Justinian in the \emph{Jewish Encyclopedia} , and in his \emph{Studien zur byz-jüdischem Gesch.} (1914), contended that by deuterosis must be understood the whole of the oral teaching, and repeated this view in \emph{Jewish Guardian}, March 26, 1920, p4; but in the same periodical, April 2, 1920, p11, J. Abrahams maintains that what Justinian forbade under Deuterosis was ``the traditional Rabbinic translation of the Law \texthebrew{—} the Targum.'' The words of the legislator certainly seem to imply a book. Mr. F. Colson has suggested to me that the Mishna is meant: deuterosis has the same meaning, ``repetition.''

} This amazing law is thoroughly characteristic of the Imperial theologian.

We saw in a former chapter how throughout the fifth century the severe laws against paganism were not very strictly enforced.

p367
So long as there was no open scandal, men could still believe in the old religions and disseminate anti-Christian doctrine. This comparatively tolerant attitude of the State terminated with the accession of Justinian, who had firmly resolved to realise the conception of an empire in which there should be no differences of religious opinion. Paganism was already dying slowly, and it seemed no difficult task to extinguish it entirely. There were two distinct forms in which it survived. In a few outlying places, and in some wild districts where the work of conversion had been imperfectly done, the population still indulged with impunity in heathen practices. To suppress these was a matter of administration, reinforced by missionary zeal; no new laws were required. A more serious problem was presented by the Hellenism which prevailed widely enough among the educated classes, and consequently in the State-service itself. To cope with this Justinian saw that there was need not only of new administrative rigour, but of new legislation. He saw that Hellenism was kept alive by pagan instructors of youth, especially in teaching establishments which had preserved the Greek tradition of education. If the evil thing was to be eradicated, he must strike at these.

Not long after his accession, he reaffirmed the penalties which previous Emperors had enacted against the pagans, and forbade all donations or legacies for the purpose of maintaining ``Hellenic impiety,'' while in the same constitution he enjoined upon all the civil authorities and the bishops, in Constantinople and in the provinces, to inquire into cases of pagan superstition.\texthebrew{⁠}\footnote{\emph{C. J.} I.11.9 , evidently the law of Justinian which preceded the trials at Constantinople, which apparently began (John Mal. XVIII.449) in the last months of 528.

} This law was soon followed by another which made it illegal for any persons ``infected with the madness of the unholy Hellenes'' to teach any subject, and thereby under the pretext of education corrupt the souls of their pupils.\footnote{\emph{C. J.} I.11.10, probably A.D. 529. One of the provisions is the penalty of death for apostates. The constitution I.15.18, seems to be later (cp. § 4 \textgreek{μηδ\texthebrew{ὲ ἐ}ν σχήματι διδασκάλου παιδείας}).

}

The persecution began with an inquisition at Constantinople. Many persons of the highest position were accused and condemned.\texthebrew{⁠}\footnote{John Mal. \emph{ib.}, but Theophanes, A.M. 6022, gives the fuller text of Malalas. Asclepiodotus, ex-prefect, took poison. Malalas says that Thomas, Phocas, and the others

(p368)
\textgreek{\texthebrew{ἐ}τελεύτησαν}; but Theophanes has \textgreek{συνελήφθησαν}, evidently the right reading, leaving their fate open. The fact that Procopius does not mention executions in the brief passage where he refers to the persecution of pagans (\emph{H. A.} II \textgreek{α\texthebrew{ἰ}κιζόμενος τε τ\texthebrew{ὰ} σώματα κα\texthebrew{ὶ} τ\texthebrew{ὰ} χρήματα ληιζόμενος}) must make us hesitate to accept the text of Malalas. Those who were executed must have been condemned as apostates. It is to be observed that Thomas was still Quaestor in April 529, \emph{C. J.} \emph{De Just. cod. confirmando}; this shows that the investigations lasted a considerable time.

} Their property was confiscated, and some may have

p368
been put to death; one committed suicide. Among those who were involved were Thomas the Quaestor and Phocas, son of Craterus. But Phocas, a patrician of whose estimable character we have a portrait drawn by a contemporary,\texthebrew{⁠}\footnote{John Lyd. \emph{De mag.} III.72 ; cp. Procopius, \emph{H. A.} 21 . This Phocas is not to be confused with Phocas who was Praet. Prefect of Illyricum in 529 and took part in compiling the Code (\emph{C. J.} \emph{ib.}).

} was speedily pardoned, for, as we saw, he was appointed Praetorian Prefect of the East after the Nika riot.

Some of the accused escaped by pretending to embrace the Christian faith, but we are told that ``not long afterwards they were convicted of offering libations and sacrifices and other unholy practices.''\texthebrew{⁠}\footnote{Procopius, \emph{ib.}

} There was, in fact, a second inquisition in A.D. 546. On this occasion a heretic was set to catch the pagan. Through the zeal of John of Ephesus, a Monophysite, who was head of a Syrian monastery in the suburb of Sycae, a large number of senators, ``with a crowd of grammarians, sophists, lawyers, and physicians,'' were denounced, not without the use of torture, and suffered whippings and imprisonment. Then ``they were given to the churches to be instructed in the Christian faith.'' One name is mentioned: Phocas, a rich and power ­ful patrician, who, knowing that he had been denounced, took poison. The Emperor ordered that he should be buried like an ass without any rites. We may suspect that this was the same Phocas, son of Craterus, who had been involved in the earlier inquest and knew that death would be the penalty of his relapse.\texthebrew{⁠}\footnote{John Eph., \emph{H. E.} Part II (Nau), p481. The date is given as the nineteenth year of Justinian, which was 546; but as John places the association of Justinian with Justin in 531 (p474), it may be that 550 is the true date. Against this is the apparent reference of Procopius (see last note) in a work written in 550.

} There was yet another pagan scandal in the capital in A.D. 559; the condemned were exposed to popular derision in a mock procession and their books publicly burned.\footnote{John Mal. \emph{ib.} p491.

}

It may be considered certain that in all cases the condemned were found guilty of actual heathen practices, for instance of

p369
sacrifi­cing or pouring libations in their private houses, on the altars of pagan deities. Men could still cling to pagan beliefs, provided they did not express their faith in any overt act. There were many distinguished people of this kind in the highest circles at Constantinople, many lawyers and literary men, whose infidelity was well known and tolerated. The great jurist Tribonian, who was in high favour with the Emperor, was an eminent example. He seems to have made no pretence at disguising his opinions, but others feigned to conform to the State religion. We are told that John the Cappadocian used sometimes to go to church at night, but he went dressed in a rough cloak like an old pagan priest, and instead of behaving as a Christian worshipper he used to mumble impious words the whole night.\footnote{Procopius, \emph{B. P.} I.25.10.

}

It can hardly be doubted that by making the profession of orthodoxy a necessary condition for public teaching Justinian accelerated the extinction of ``Hellenism.'' Pagan traditions and a pagan atmosphere were still maintained, not only in the schools of philosophy, but in the schools of law, not only at Athens, but at Alexandria, Gaza, and elsewhere. The suppression of all law schools, except those of Constantinople and Berytus, though not intended for this purpose, must have affected the interests of paganism. But philosophical teaching was the great danger, and Athens was the most notorious home of uncompromising Hellenists. After the death of Proclus ( A.D. 485) the Athenian university declined, but there were teachers of considerable metaphysical ability, such as Simplicius and Damascius, the last scholarch,\texthebrew{⁠}\footnote{Proclus was succeeded by Marinus his biographer; then came Isidore, Hegias, Damascius.

} whose attainments can still be judged by their works.\footnote{The excellent commentaries of Simplicius on Aristotle are well known. Damascius wrote commentaries on Aristotle and a treatise, \textgreek{περ\texthebrew{ὶ} πρώτων \texthebrew{ἀ}ρχ\texthebrew{ῶ}ν}, which are extant. Their colleague, Priscian, also wrote on the Aristotelian philosophy, and we have in a Latin translation his ``Solutions of Questions proposed by King Chosroes.''

}

The edicts of Justinian sounded the doom of the Athenian schools, which had a continuous tradition since the days of Plato and Aristotle. We do not know exactly what happened in A.D. 529.\texthebrew{⁠}\footnote{John Mal. XVIII p451 \textgreek{\texthebrew{ὁ} α\texthebrew{ὐ}τ\texthebrew{ὸ}ς βασιλε\texthebrew{ὺ}ς θεσπίσας πρόσταξιν \texthebrew{ἔ}πεμψεν \texthebrew{ἐ}ν \texthebrew{Ἀ}θήναις, κλεύσας μηδένα διδάσκειν φιλοσοφίαν μήτε νόμιμα \texthebrew{ἐ}ξηγε\texthebrew{ῖ}σθαι. θεσπίσας} seems to refer to

\emph{C. J.} I.11.10 .

} We may suppose that the teachers were warned that

p370
unless they were baptized and publicly embraced Christianity, they would no longer be permitted to teach; and that when they refused, the property of the schools was confiscated and their means of livelihood withdrawn.\footnote{The provision in

\emph{C. J.} \emph{ib.} § 2

\textgreek{μηδ\texthebrew{ὲ ἐ}κ το\texthebrew{ῦ} δημοσίου σιτήσεως \texthebrew{ἀ}πολαύειν α\texthebrew{ὐ}τους} would hardly apply, as it refers only to grants from the fisc. But the seizure of the endowments might be covered by

I.11.9 , whereby it is provided that all property bequeathed or given for the maintenance of Hellenic impiety should be seized and handed over to the municipality. In case this law was applied, Athens at least had the benefit of the property which the Academy and the Lyceum had accumulated. Damascius, \emph{Vita Isidori}, § 158, says that in the time of Proclus the revenue of the Academy was 1000 solidi (£625), or somewhat more. The closing of these schools \texthebrew{—} which was the result of Justinian's general laws, not of any edict aimed specially at Athens \texthebrew{—} has been discussed by Hertzberg, \emph{Gesch. Griechenlands unter der röm. Herrschaft}, III.538 \emph{sqq.}; Gregorovius, \emph{Gesch. der Stadt Athen}, I.54 \emph{sqq.}; Paparrigopoulos, \textgreek{\texthebrew{Ἱ}στ. το\texthebrew{ῦ Ἑ}λλ. \texthebrew{ἔ}θνους}, III.174\texthebrew{‑}175; Diehl, \emph{Justinien}, 562 \emph{sqq.}

}

This event had a curious sequel. Some of the philosophers whose occupation was gone resolved to cast the dust of the Christian Empire from their feet and migrate to Persia. Of these the most illustrious were Damascius, the last scholarch of the Academy, Simplicius, and Priscian. The names of four others are mentioned, but we do not know whether they had taught at Athens or at some other seat of learning.\texthebrew{⁠}\footnote{It has been generally assumed that they were all Athenian professors, but Agathias, who is our authority, does not say so

(II.30) . The others were Eulamios of Phrygia, Hermeias and Diogenes of Phoenicia, and Isidore of Gaza, all otherwise unknown. Suidas, \emph{s.v.} \textgreek{πρέσβεις}, says that they accompanied Areobindus, who was sent as envoy to Persia; and Clinton (\emph{sub} A.D. 532) accepts this, not perceiving what Agathias had said about the impostor Uranias

(II.29) . As Chosroes came to the throne in September 531, Clinton perhaps rightly places their journey after that date. Gregorovius notes that the pythagorean number of the \emph{seven} philosophers is somewhat suspicious.

} These men had heard that king Chosroes was interested in philosophy, and they hoped, protected by his favour and supported by his generosity, to end their days in a more enlightened country than their own. But they were disappointed. Chosroes was flattered by their arrival and begged them to remain. But they soon found the strange conditions of life intolerable. They fell homesick, and felt that they would prefer death on Roman soil to the highest honours the Persian could confer. And so they returned. But the king did them a great service. In his treaty with Justinian in A.D. 532 he stipulated that they should not be molested or forced to embrace the Christian faith. We are told that they lived comfortably for the rest of their lives, and we know that

p371
Simplicius was still writing philosophical works in the later years of Justinian.\footnote{Agathias, II.31

\textgreek{\texthebrew{ὁ ἐ}φεξ\texthebrew{ῆ}ς βίος ε\texthebrew{ἰ}ς τ\texthebrew{ὸ} θυμ\texthebrew{ῆ}ρές τε κα\texthebrew{ὶ ἥ}διστον \texthebrew{ἀ}πετελεύτησεν}. For the date of the commentaries of Simplicius on Aristotle's \emph{De caelo} and \emph{Physica} cp. Clinton, \emph{F. R.} II. pp328\texthebrew{‑}329.

}

In western Asia Minor, in the provinces of Asia, Phrygia, Lydia, and Caria, there was still a considerable survival of pagan cults, not only in the country regions, but in some of the towns, for instance in Tralles. In A.D. 542 John of Ephesus, the Monophysite whose activity in hunting down the Hellenes at Constantinople has already been noticed, was sent as a missionary to these provinces to convert the heathen and to put an end to idolatrous practices. He tells us in his \emph{Ecclesiastical History} that he converted 70,000 souls. The temples were destroyed; 96 churches and 12 monasteries were founded. Justinian paid for the baptismal vestments of the converts and gave each a small sum of money (about 4s.).\footnote{Of the new churches 55 were paid for by the treasury, 41 by the proselytes. John Eph., \emph{H. E.} Part II (Nau), p482; Part III III. cc36\texthebrew{‑}37; \emph{Comm. de b. or.} cc40, 43, 51, etc.

}

In Egypt, in the oasis of Augila, the temple dedicated to Zeus Ammon and Alexander the Great still stood, and sacrifices were still offered. Justinian put an end to this worship and built a church to the Mother of God.\texthebrew{⁠}\footnote{Procopius, \emph{Aed.} VI.2 .

} At Philae the cult of Osiris and Isis had been permitted to continue undisturbed. This toleration was chiefly due to the fact that the Blemyes and Nobadae, the southern neighbours of Egypt, had a vested interest in the temples by virtue of a treaty which they had made with Diocletian. Every year they came down the river to worship Isis in the island of Elephantine; and at fixed times the image of the goddess was brought back to the temple.\texthebrew{⁠}\footnote{Priscus, \emph{De leg. gent.}, \emph{fr.} 11, p583.

} Justinian would tolerate this indulgence no longer. Early in his reign he sent Narses the Persarmenian to destroy the sanctuaries. The priests were arrested and the divine images sent to Constantinople.\texthebrew{⁠}\footnote{Procopius, \emph{B. P.} I.19. The death of Narses in 543 gives a posterior limit of date.

} Much about the same time the Christian conversion of the Nobadae and Blemyes began.

Justinian was undoubtedly success ­ful in hastening the disappearance of open heathen practices and in suppressing anti-Christian

p372
philosophy. Although in some places, like Heliopolis,\texthebrew{⁠}\footnote{The great temple of Baal had been converted into a church; but sacrifices were still performed there in the sixth century. In A.D. 555 it was ruined by lightning. John Eph., \emph{H. E.} Part II p490. See Diehl, \emph{Justinien}, pp550\texthebrew{‑}551.

} paganism may have survived for another generation, and although there were inquisitions under his immediate successors, it may be said that by the close of the sixth century the old faiths were virtually extinct throughout the Empire.

Throughout his reign one of Justinian's chief preoccupations was to find an issue from the dilemma in which the controversy over the natures of Christ had placed the Imperial government. Concord with Rome and the western churches meant discord in the East; toleration in the East meant separation from Rome. The solution of the problem was not rendered easier by the fact that the Emperor was a theologian and took a deep interest in the questions at issue on their own account apart from the political consequences which were involved.

In the abandonment of the ecclesiastical policy of Zeno and Anastasius, in order to heal the schism with Rome, Justinian, co-operating with Vitalian and the patriarch John, had been a moving spirit. The greater part of the correspondence between Pope Hormisdas and the personages at Constantinople who took part in the negotiations has been preserved.\texthebrew{⁠}\footnote{It will be found in the \emph{Collectio Avellana, Epp.} 142\texthebrew{‑}181.

} The main question was settled by a synod which met in the capital in 518 and decided that the Monophysite bishops should be expelled from their sees. The only difficulty which occurred in the negotiations with the Pope regarded the removal of the name of Acacius from the diptychs of the Church. There was a desire at Constantinople to spare the memory of the Patriarch, but Hormisdas was firm,\texthebrew{⁠}\footnote{See \emph{ib.} \emph{Ep.} 147. Justinian's letter of September 7, received December 20, and the Pope's letters, \emph{Epp.} 145, 148, 149.

} and in April A.D. 519 the Patriarch despatched to the Pope a memorandum, in which he anathematised Acacius and all those who had participated with him, and confessed that ``the Catholic faith is always kept inviolable in the Apostolic see.''\footnote{\emph{Ep.} 159. The Pope had already sent a deputation of bishops to Constantinople (January), and the deacon, Dioscorus, who attended them describes their journey by the Egnatian Way, their reception at the tenth milestone from the capital by

Vitalian, Pompeius, Justinian, and many other illustrious persons, and their presentation to the Emperor in the presence of the Senate (\emph{Ep.} 167).

}

p373
The names of five Patriarchs, Acacius, Fravitta, Euphemius, Macedonius, Timotheus, and of two Emperors, Zeno and Anastasius, were solemnly erased from the diptychs of the Church of Constantinople, and it only remained for the Pope to remind the Emperor that he had still to take measures to ``correct'' the Churches of Antioch and Alexandria.\footnote{\emph{Ep.} 168 (July 9, 519).

}

``Correction'' meant persecution, and the Emperor did not hesitate. The great Monophysite leader Severus had already been expelled from Antioch, and more than fifty other bishops driven into exile, including Julian of Halicarnassus, Peter of Apamea, and Thomas of Daras. The heretical monastic communities in Syria were dispersed and the convents closed. Resistance led to imprisonment and massacres. Such measures did not extirpate the heresy. In Egypt, Palestine, Syria, and the Mesopotamian deserts the Monophysites persisted in their errors, hoping for better days. Severus himself was able to live quietly in Alexandria.\footnote{For the persecutions see Zacharias Myt. IX.4, 5; John Eph. \emph{H. E.} Part II. pp467\texthebrew{‑}468, and \emph{Comm.} c5, p35 \emph{sqq.}, c8, p46 \emph{sqq.} (cp. pp 217\texthebrew{‑}220); \emph{Chron. Edess.} p124\texthebrew{‑}128.

}

The persecution continued throughout the reign of Justin. But Justinian determined to essay a different policy. He did not despair of finding a theological formula which would reconcile the views of moderate Monophysites with the adherents of the dogma of Chalcedon. For there was after all a common basis in the doctrine of Cyril, which the Monophysites acknowledged and the Dyophysites could not repudiate. For the Council of Chalcedon had approved of the views of Cyril, and Severus would hardly have admitted that his own doctrine diverged from Cyril's, if rightly interpreted.

The whole question was being studied anew by a theologian whom modern authorities regard as the ablest interpreter of the Chalcedonian Creed, Leontius of Byzantium.\texthebrew{⁠}\footnote{The study of Loofs, \emph{Leontius von Byzanz} (1887), rescued this theologian from neglect, and was followed in 1894 by a monograph with the same title by Rügamer, written from the Catholic point of view and traversing successfully some of the conclusions of Loofs. The earliest work of Leontius was probably the Three Books against the Nestorians and Eutychians (\emph{P. G.} LXXXVI.1267 \emph{sqq.}), which may be dated to A.D. 529\texthebrew{‑}530 (cp. Rügamer, 9 \emph{sqq.}). The \emph{Epilysis of the Syllogisms of Severus} (\emph{ib.} 1915 \emph{sqq.}) and the \emph{Thirty Chapters against Severus} (\emph{P. G.} cxxx.1068 \emph{sqq.}) may have been composed in the years immediately following. Other works against the Monophysites (\emph{P. G.} LXXXVI.1769 \emph{sqq.}) and the Nestorians (\emph{ib.} 1399 \emph{sqq.}), and the \emph{Scholia}, a treatise on Sects, which bears the marks of later editing (\emph{ib.} 1193 \emph{sqq.}), may be later than A.D. 544.

} In his youth he

p374
had been ensnared in the errors of Nestorianism, but, happily, guided into the ways of orthodoxy, he lived to write with equal zeal against Nestorians and Monophysites. He has the distinction of introdu­cing a new technical term into Greek theology, enhypostasis , which magically solved the difficulty that had led Nestorians and Monophysites into their opposite heresies. Admitting the axiom that there is no nature without a hypostasis, Leontius said: it does not follow that the subsistence of two natures in Christ involves two hypostaseis (as the Nestorians say), nor yet that to avoid the assumption of two hypostaseis we must assume only one nature with the Monophysites. The truth is that both natures, the human like the divine, subsist in the same hypostasis of the Logos; and to this relation he gave the name of enhypostasis .\footnote{Cp. \emph{ib.} 1277 \emph{sqq.}; 1944.

}

Of much greater interest is the fact that in his theological discussions he resorts to a new instrument, the categories and distinctions of the Aristotelian philosophy.\texthebrew{⁠}\footnote{See Loofs, \emph{op. cit.} 60 \emph{sqq.}

} Substance, genus, species, qualities, play their parts as in the western scholasticism of a later age. It is not probable that Leontius himself was a student of Aristotle, but at this period there was a revival of Aristotelian thought which influenced Christian as well as secular learning. The ablest exponent of this movement was indeed in the camp of the heretics, John Philoponus of Alexandria, a philosopher, and a Monophysite. His writings are said to have been partly responsible for the development of a theory about the Trinity, known as Tritheism, which had some vogue at this period and was ardently supported by Athanasius, a grandson of the Empress. The Tritheites held the persons of the Trinity to be of the same substance and One God; but they explained the identity of substance as purely generic, in the Aristotelian sense. Numerically, they said, there are three substances and three natures, though these are one and equal by virtue of the unchangeable identity of the Godhead.\footnote{John Ascosnaghes is said to have been the founder of the Tritheite sect, to which bishops Conon of Tarsus and Eugenius of Seleucia belonged. See John Eph. \emph{H. E.} Part III V.1\texthebrew{‑}12; Timotheus, \emph{De recept. haereticorum}; Migne, \emph{P. G.} LXXXVI. p64. John Philoponus in his \emph{\textgreek{Διαιτητ\texthebrew{ῆ}ς}} (``Arbiter'') discussed the bases of Tritheism and Monophysitism. We possess a philosophical work by him \emph{On the Eternity of the World} (against Proclus), written in A.D. 529

(cp. XVI.4), and an \emph{Exegesis of the Cosmogony of Moses}, dedicated to Sergius, Patriarch of Antioch (\emph{c.} 546\texthebrew{‑}549).

}

p375
To return to Leontius, it is a curious fact that notwithstanding the importance and considerable number of his theological works contemporary writers never mention him. Modern writers have indeed proposed to identify him with other persons of the same name who played minor parts in the ecclesiastical history of his time, but these conjectures are extremely doubtful.\texthebrew{⁠}\footnote{We may decidedly reject the identification, maintained by Loofs (\emph{op. cit.}) with Leontius, the Origenist of Palestine, who visited Constantinople with St. Sabas in A.D. 531 and was repudiated by him. Cp. the criticisms of Rügamer, \emph{op. cit.} 58 \emph{sqq.} The other proposition of Loofs that Leontius of Byzantium is the same as Leontius, a relative of Vitalian, and one of the Scythian monks who raised the Theopaschite question at the beginning of Justin's reign, can neither be proved nor disproved.

} His works were composed during a period of fifteen or twenty years (about A.D. 530\texthebrew{‑}550), and it is probable that they helped Justinian in his efforts to interpret the creed of Chalcedon in such a way as to win Monophysites of the school of Severus.

The Monophysites were far from being a united body. The ground common to all was the repudiation of the Council of Chalcedon and the reception of the Patriarch Dioscorus. There were ultimately twelve different sections,\texthebrew{⁠}\footnote{Timotheus, \emph{op. cit.} 52 \emph{sqq.}

} but the only division of much importance was that between the followers of Severus of Antioch and those of Julian of Halicarnassus. Julian, identifying the substance and qualities of the divinity and humanity of Christ, deduced that his body was indestructible from the moment at which it was assumed by the Logos. This doctrine, which was known as aphthartodocetism, called forth the polemic of Leontius; but no Chalcedonian could have attacked it with more energy than Severus.\footnote{Leontius, \emph{Contra Nest. et Eut.} book II; Zach. Myt. IX.9\texthebrew{‑}13.

}

Justinian began his policy of conciliation by allowing the heretical bishops and monks to return from exile, about A.D. 529.\texthebrew{⁠}\footnote{John Eph. \emph{H. E.} Part II p469; Zacharias Myt. IX.15.

} His plan was to hold a conference, not a formal synod, at Constantinople, and to have the whole question discussed. Severus himself resisted all the Emperor's efforts to induce him to attend it, but some of his followers came and the conference was held

p376
in A.D. 531.\texthebrew{⁠}\footnote{Mansi, VIII.817 \emph{sq.}; John Eph. \emph{Comm.} 203. 245. For the attempt to win Severus cp. Evagrius, IV.11 ; Zacharias reproduces the letters of Severus, \emph{ib.} 16. The date of the conference, at which Hypatius, bishop of Ephesus, presided, is 531 (not 533), cp. Loofs, \emph{op. cit.} 283.

} Leontius, representative of the orthodox monks of Jerusalem, took part in it, and we may possibly identify him with Leontius of Byzantium, the theologian.\texthebrew{⁠}\footnote{Mansi, \emph{ib.} Leontius vir venerabilis monachus et apocrisiarius patrum in sancta civitate constitutorum. In the MSS. of some of his works Leontius Byz. is described as \textgreek{μοναχός} and \textgreek{\texthebrew{Ἱ}εροσολυμίτης}. The same Leontius was present at the Synod of 536.

} The conference led to no results.

The failure of his first attempt did not deter Justinian from making a second, and he sought a formula of conciliation in what is known as the Theopaschite doctrine. The thesis that it was orthodox to hold that ``one of the Holy Trinity suffered in the flesh'' had been defended in A.D. 519 by four Scythian monks, in the presence of John the Patriarch and the Papal legates who had come to restore peace to the Church.\texthebrew{⁠}\footnote{The sources for the affair of the Scythian monks are their joint letter to some African bishops exiled in Sardinia (\emph{P. L.} LXV.442 \emph{sqq.}), the writings of their patron John Maxentius (\emph{P. G.} LXXXVI.73 \emph{sqq.}), and the correspondence of Justin, Justinian, and others with Hormisdas concerning them (\emph{Coll. Avellana, Epp.} 187, 188, 196, 216, 232\texthebrew{‑}239). The Papal legates were unfavourably impressed by them, and their secretary Dioscorus reported to the Pope that the monks, whom he describes as de domo Vitaliani, were ``adversaries of the prayers of all Christians.'' The monks went to Rome (summer 519) to submit their views to the Pope and remained there till August 520. It has been already mentioned that Loofs identifies Leontius, one of these monks, with Leontius of Byzantium, but in this theologian's voluminous works he can only find one or two allusions to the Theopaschite dogma (\emph{Contra Nest. et Eut.} 1289, 1377); see Loofs, \emph{op. cit.} 228.

} The formula was denounced as heretical by the Sleepless monks, who had been so active in opposing the Trisagion , to which it had suspicious resemblance. Justinian was interested in the question, and he wrote to Pope Hormisdas repeatedly, urging him to pronounce a decision.\texthebrew{⁠}\footnote{\emph{Coll. Avell.} cxcvi.235.

} But the Pope evaded a definite reply. Justinian recurred to the subject in A.D. 533, with a political object. He issued an edict which implicitly asserted that one of the Trinity suffered in the flesh, and he procured a confirmation of the edict by Pope John II.\texthebrew{⁠}\footnote{\emph{C. J.} I.1.6 ; Mansi, VIII.798. Cp. Loofs, \emph{ib.} 260.

} The Sleepless monks, who refused to accept the doctrine, were excommunicated.

The recognition of a formula which did not touch the main issue\texthebrew{⁠}\footnote{Loofs says (255): ``The Theopaschite controversy is an event in the history of the doctrines of the Trinity far more than in that of Christology: it is one of the first signs of the victory of Aristotle over Plato.''

} could not deceive Severus and the Monophysites, and

p377
having suffered two defeats the Emperor seems to have been persuaded by Theodora to allow her to deal with the situation on other lines. At least it is difficult otherwise to explain what happened. When Epiphanius, the Patriarch of Constantinople, died (June, A.D. 535), she procured the election of Anthimus, bishop of Trapezus, who was secretly a Monophysite. He addressed to Severus a letter containing a Monophysitic confession of faith; he communicated with Theodosius, the Monophysite Patriarch of Alexandria,\texthebrew{⁠}\footnote{Theodosius succeeded Timothy IV on February 9 or 11, 535, but at the same time a rival Patriarch, Gaian (whose views agreed with those of Julian of Halicarnassus), was elected and was Patriarch in possession for 103 days (till May 23 or 25). See Brooks, \emph{B. Z.} XII.494 \emph{sqq.} For the letter of Anthimus see Zach. Myt. IX.21.

} and induced the Patriarch of Jerusalem to follow his example. Severus was invited to the capital and Theodora lodged him in the Palace.\texthebrew{⁠}\footnote{\emph{Ib.} IX.15, 19.

} The Patriarch of Antioch, Ephraim, was a firm adherent of Chalcedon, and he sent a message to Pope Agapetus warning him that heresy was again in the ascendant. Agapetus, arriving at Constantinople early in A.D. 536 and received with great honour by Justinian,\texthebrew{⁠}\footnote{\emph{Ib.} 19, where it is said that Justinian was particularly pleased to see Agapetus because he spoke the same language (Latin).

} refused to communicate with Anthimus, procured his deposition (March 12), and consecrated Menas as his successor.\texthebrew{⁠}\footnote{Anthimus had held office for ten months. Menas succeeded on March 13. Cp. Andreev, \emph{Kpl. Patr.} 170\texthebrew{‑}173.

} The Pope died suddenly a few weeks later, but in May Menas summoned a synod which anathematised Anthimus, Severus, and others, and condemned their writings.\texthebrew{⁠}\footnote{Mansi, VIII.877 \emph{sqq.} Peter of Apamea and Zooras were the others who were condemned.

} The Emperor then issued a law confirming the acts of the synod, and forbidding Anthimus, Severus, and the others to reside in any large city.\texthebrew{⁠}\footnote{\emph{Nov.} 42.

} Severus spent the last years of his life in the Egyptian desert. Anthimus lived in concealment in Theodora's palace, along with other Monophysites like Theodosius of Alexandria.\footnote{See John Eph. \emph{Comm.} pp247\texthebrew{‑}248. At the time of the Council Anthimus could not be found anywhere (Mansi, VIII.941). The date of the death of Severus is probably February 8, 538 (Michael Mel. \emph{Chron.} IX.29). Cp. Brooks, \emph{B. Z.} XII.497; Krüger places it in 543 (art. \emph{Monophysiten} in \emph{Realenc. f. protest. Theologie}).

}

A new persecution was now let loose in the East. It was organised by Ephraim of Antioch, who acted as grand inquisitor, and the Monophysite historians have their tale to tell of

p378
imprisonments, tortures, and burnings.\texthebrew{⁠}\footnote{Zacharias Myt. (X.1) and John Eph. (cp. \emph{Comm.} pp111, 134, etc.; 221 \emph{sqq.}) are the chief sources on the persecution. John, bishop of Tella, died in consequence of the tortures which he underwent (John Eph. \emph{Comm.} c24).

} The Emperor, abandoning his policy of conciliation, was perhaps principally moved by the consideration of his designs on Italy. It was important at this juncture to make it quite clear that his own zeal for orthodoxy was above cavil and to dispel in the minds of the Italians any suspicion that he was inclined to coquette with the Monophysites.\footnote{Dante

(\emph{Paradiso}, VI.13 \emph{sqq.})

represents Justinian as holding Monophysitic opinions and converted by Agapetus:

E prima ch' io all' opra fossi attento,

Una natura in Cristo esser, non piue

Credeva, e di tal fede era contento.

Ma il benedetto Agabito, che fue

Sommo pastore, alla fede sincera

Mi dirizzò con le parole sue.

}

The fall of Anthimus, the ensuing synod, and the Imperial edict which confirmed it, were deeply displeasing to Theodora. But she did not lose heart. She not only protected the heretical leaders, but she formed the bold design of counteracting her husband's policy from Rome itself. The deacon Vigilius was at this time the apocrisiarius or nuncio of the Roman see at Constantinople. He was a man of old senatorial family, the son of a consul, and he had been a favourite of Boniface II, who had desired to secure his succession to the pontifical throne. On the death of Agapetus he saw his chance, and Theodora, who though º she knew what manner of man he was, saw her opportunity. An arrangement was made between them. Theodora promised to place at his disposal 200 lbs. of gold (£9000) and provided him with letters to Belisarius and Antonina, and on his part, if he did not definitely promise, he led her to believe that he would repudiate the Council of Chalcedon and re-establish Anthimus in the see of Constantinople.\texthebrew{⁠}\footnote{Liberatus, \emph{Brev.} 22; \emph{Lib. Pont., Vita Silv.} p292; \emph{Vita Vigil.} p297. Victor Tonn. \emph{Chron., sub} 542, says that Vigilius undertook to condemn the Three Chapters occulto chirographo which he gave to Theodora; that he was made Pope by Antonina; and he quotes a letter of Vigilius (probably spurious) in which Antonina is mentioned. The reference to the Three Chapters here is an anachronism.

} He hastened to Italy, but he arrived too late. King Theodahad had received early notice of the sudden death of Agapetus, and under his auspices Silverius had been elected Pope (in June).\footnote{June 8. In \emph{Lib. Pont.} p290, it is erroneously said that Silverius held the pontifical chair one year, five months, eleven days, which would give November 18, 537, for his deposition.

}

The Empress then wrote to Silverius asking him to procure the restoration of Anthimus, and on his refusal she determined to avail herself of the military occupation of Rome by Belisarius

p379
to intimidate or, if necessary, to remove him. She sent secret instructions to Antonina, probably leaving it to her ingenuity to concoct a plot against the Pope. Silverius resided at the Lateran beside the Asinarian Gate, and a letter was fabricated as evidence that he was in treacherous communication with the Goths and proposed to admit them into the city. Belisarius summoned him to the Pincian palace, showed him the danger of his position, and intimated that he could save himself by obeying the wishes of the Empress. Silverius refused to yield, and was suffered to depart, but he took the precaution of withdrawing from the Lateran to º St. Sabina on the Aventine at a safe distance from the walls, to prove that he had no desire to communicate with the enemy. He was called a second time to the general's presence and went attended by a numerous retinue, including the deacon Vigilius, who had come to Rome with Belisarius and was eagerly awaiting the development of events. The chief hall in the Pincian palace was divided by curtains into three apartments. The Roman clergy remained in the two outer rooms; only Silverius and Vigilius were admitted into the presence of Belisarius. When the Pope again proved inflexible, two subdeacons entered, removed his pallium, and clothed him in the garb of a monk. He was banished to Patara in Lycia. This perfidious act occurred about the middle of March, and was followed by the election of Vigilius, who was undoubtedly accessory to it. He was ordained bishop of Rome on March 29, A.D. 537.\footnote{The story is told in \emph{Lib. Pont., Vita Silv.} 292\texthebrew{‑}293, and is noticed briefly by Procopius, \emph{B. G.} I.25.13, who only gives the publicly alleged ground for the action against Silverius, namely \textgreek{\texthebrew{ὡ}ς δ\texthebrew{ὴ} προδοσίαν \texthebrew{ἐ}ς Γότθους πράσσοι}, and says that he was immediately sent to Greece. But in \emph{H. A.} (see next note) the hand of Theodora is recognised. \emph{Cont. Marcellini, sub} 537, gives the official story: cui (Vitigi) tunc faventem papam Silverium Belisarius ab episcopatu summovit.

}

There is a certain mystery about the subsequent fate of the unhappy Silverius. The government of Constantinople deemed it expedient that he should leave Patara and return to Italy. It is not clear whether Theodora approved or not, but Pelagius, the Papal nuncio, protested. It would be difficult to believe that Pelagius was not perfectly aware of the scandalous intrigue to which Vigilius owed his elevation, and it was certainly in the interest of Vigilius that he desired to keep Silverius far from Italy. When Silverius returned, Vigilius appealed to Belisarius and Antonina. With their permission, he caused his victim to

p380
be conveyed to the island of Palmaria, where according to one account he died of hunger and exhaustion, while there is another record that he was done to death by a creature of Antonia.\footnote{\emph{H. A.} 1.14 . \textgreek{Σιλβέριον διαχρησαμένη} (Antonina, in the interest of Theodora), and 27 \textgreek{τ\texthebrew{ῶ}ν τινος ο\texthebrew{ἰ}κετ\texthebrew{ῶ}ν Ε\texthebrew{ὐ}γενίου \texthebrew{ὄ}νομα \texthebrew{ὑ}πουργήσαντός ο\texthebrew{ἱ ἐ}ς \texthebrew{ἅ}παν τ\texthebrew{ὸ ἄ}γος, \texthebrew{ᾧ} δ\texthebrew{ὴ} κα\texthebrew{ὶ} τ\texthebrew{ὸ ἐ}ς Σιλβέριον ε\texthebrew{ἴ}ργασται μίασμα}. Nothing is said of the motive. In \emph{Lib. Pont.} 293, the other account will be found.

}

This intrigue of the Empress did not profit her much. The theological convictions of Pope Vigilius were stronger than his respect for his plighted word, and, when he had attained the goal of his ambition by her help, his robust conscience had no scruples in evading the fulfilment of his promises. By evasions and postponements, and by the assistance of his loyal and tactful nuncio, Pelagius, who had succeeded in ingratiating himself with Theodora as well as Justinian, he managed to avoid a breach with the Empress, while he addressed to the Emperor and to the Patriarch letters in which he maintained the condemnation of the opponents of the Council of Chalcedon.\footnote{\emph{Coll. Avell., Epp.} 92, 93 (also in Migne, \emph{P. L.} LXIX p21 and p25). Pelagius, who was a man of large fortune and high social position, had gone to Constantinople with Agapetus, where he remained for the Council of 536, and was appointed nuncio, in succession to Vigilius, in the same year.

}

Theodosius, the Monophysite Patriarch of Alexandria, who had been deprived of his see in A.D. 536, was succeeded by Paul, a monk of Tabenna, who was ordained by Menas and went to Egypt with full powers to cleanse the sees of the Patriarchate from heretical bishops. Rhodon, the Augustal Prefect, received instructions to support him in all the measures he thought fit to take.\texthebrew{⁠}\footnote{Procopius, \emph{H. A.} 27.3 \emph{sqq.} ; and Liberatus, \emph{Brev.} 22, 23, are the sources.

} The submission seems to have been general; the treatment of Theodosius, who had not been popular, excited little resentment. But a certain deacon, named Psoes, headed an opposition to the new Patriarch, at whose instance he was arrested by Rhodon and died under torture. Theodora was furious and insisted on an investigation; Liberius, the Roman senator who had held high offices in Italy under the Ostrogothic kings and came to Constantinople as ambassador of Theodahad,\texthebrew{⁠}\footnote{See above,

p164 .

} was appointed to succeed Rhodon, and a clerical commission, including

p381
the nuncio Pelagius, was sent with him to Alexandria to pronounce on the conduct of Paul. The clergy proceeded to Gaza, where they held a synod (about Easter, A.D. 542),\texthebrew{⁠}\footnote{See Mansi, VIII.1164. Ephraim, Patriarch of Antioch, attended. For the date see Diekamp, \emph{Die origenistichen Streitigkeiten}, 41 \emph{sqq.}

} at which Pelagius presided, and Paul was found guilty for the death of Psoes and deposed. Rhodon, who fled to Constantinople, was beheaded, though it is said that he produced thirteen letters of the Emperor authorising all that he had done.\footnote{Procopius, \emph{ib.} 15 and 18. Probably Justinian's letters only instructed the Prefect to obey Paul in all things, though Procopius seems to wish to suggest that they authorised a particular act. Arsenius, a converted Samaritan, who had co-operated with Paul, was hanged by Liberius at the instance of Theodora, \emph{ib.} 19.

}

Pelagius returned from Gaza through Palestine, where he fell in with some monks of Jerusalem who were on the point of starting for Constaninople for the purpose of indu­cing Justinian to condemn the opinions of Origen, which were infecting the monasteries of Palestine.

The revival of Origenistic doctrine in the sixth century was closely connected with a mystical movement which seems to have originated in eastern Syria and threatened to taint Christian theology with speculations of a pronounced pantheistic tendency. The teacher who was principally responsible for propagating a Christian pantheism, seductive to many minds, was Stephen bar-Sudaili , of Edessa, who in consequence of his advanced opinions was compelled to leave Edessa and betake himself to Palestine.\texthebrew{⁠}\footnote{A letter of Jacob of Sarug (died 521) to Stephen himself, and another from Philoxenus (Xenaias) the Monophysite bishop of Mabug (485\texthebrew{‑}518) to two priests of Edessa, are the chief sources for the little we know of Stephen. They will be found with English translations

in A. L. Frothingham, \emph{Stephen bar Sudaili} .

} He seems to have been the author of a book which pretended to have been composed by Hierotheus, an Athenian who was alleged to have been a follower of St. Paul and to have taught Dionysius the Areopagite.\texthebrew{⁠}\footnote{A summary of the work (extant in a Syriac MS. of the British Museum) is given by Frothingham, \emph{op. cit.} 91 \emph{sqq.}

} If this is so, Stephen was the spiritual father of the famous mystical treatises which, professing to be the works of Dionysius, were given to the world early in the sixth century. The author of these fabrications emphasises his debt to ``Hierotheus,'' but he was also profoundly influenced by the writings of Proclus, the Neoplatonic philosopher, though this was an influence which naturally he

p382
could not acknowledge.\texthebrew{⁠}\footnote{See Hugo Koch, \emph{Pseudo-Dionysius Areopagita in seinen Beziehungen zum Neuplatonismus und Mysterienwesen}, 1900. The Pseudo-Dionysius works (of which the chief are entitled \emph{On the Heavenly Hierarchy; On the Ecclesiastical Hierarchy; On Divine Names; On Mystic Theology}) will be found in \emph{P. G.} 3 and 4. These works were referred to at the conference with the Severian Monophysites in 531, when Hypatius, bishop of Ephesus, pointed out that they could not be genuine (Mansi, VIII.821). But they were soon generally accepted in the Eastern Church. In the ninth century Michael II sent a copy to the Emperor Lewis the Pious, and soon afterwards they were translated into Latin by Joannes Erigena. The strong influence which they exerted on Thomas Aquinas is well known.

} The learned physician Sergius of Resaina translated these mystical treatises into Syriac, and it is noteworthy that Sergius is described as versed in the teaching of Origen.\footnote{Zach. Myt. IX.19.

}

Stephen bar-Sudaili , spending the later years of his life in a convent near Jerusalem, seems to have provoked by his teaching the return to Origen's speculations, which was to be for half a century the burning interest in the monasteries of Palestine. The ablest of the Origenist party and their leading spirit was a monk named Nonnus. It is not probable that they went so far in their speculations as Stephen himself, whose views are briefly summed up in the treatise of ``Hierotheus'' in the following words:

``All nature will be confused with the Father; nothing will perish, but all will return, be sanctified, united and confused. Thus God will be all in all. Even hell will pass away and the damned return. All orders and distinctions will cease. God will pass away, and Christ will cease to be, and the Spirit will no longer be called spirit. Essence alone will remain.''\footnote{Frothingham, \emph{op. cit.} p99.

}

Origen could not have endorsed such doctrine, but it is easy to understand that any one who entertained these ideas would find his writings more congenial than those of any other Christian theologian. There was common ground especially in the rejection of eternal damnation.\texthebrew{⁠}\footnote{This view is dealt with mercilessly in the letter of Philoxenus referred to above,

p381, n3 .

} Among the other heterodox opinions which the Palestinian heretics derived from Origen were the persistence of the soul, the creation of the world not by the Trinity but by creative Nous, the similarity of Christ to men in strength and substance, the doctrines that in the resurrection our bodies will be of circular form, that ultimately matter will entirely disappear and that the kingdom of Christ will have an end.\footnote{See the denunciations of Justinian and of the assembly of bishops in 553 (below,

p389, n2 ).

}

p383
After the death of St. Sabas (December 5, A.D. 532), the number and influence of the Origenists grew in the monasteries of Palestine. Two of the most prominent, Theodore Ascidas and Domitian,\texthebrew{⁠}\footnote{Theodore was a deacon of the New Laura, Domitian was abbot of the convent of Martyrius. Mansi, VIII.910, 911; Cyril, \emph{Vita Sabae} (the principal source for the Origenistic movement in Palestine), p518.

} visited the capital in A.D. 536 to attend the synod which condemned the Monophysites, and gaining the favour of the Emperor they were appointed to fill the sees, Domitian of Ancyra and Theodore of Caesarea in Cappadocia. Both Pelagius and the Patriarch Menas were anxious to break the influence which Theodore Ascidas, a man of considerable astuteness and not over-scrupulous , exerted over Justinian; and they eagerly took up the cause of the monks who desired to purge Palestine of the heresy. Ephraim, the Patriarch of Antioch, held a synod in summer A.D. 542 to condemn the doctrines of Origen, but the heretics were so power ­ful that they induced the Patriarch of Jerusalem to strike out Ephraim's name from the diptychs.

Pelagius and Menas convinced Justinian that it was imperative to take action, and in A.D. 543 the Emperor issued an edict condemning ten opinions of Origen.\texthebrew{⁠}\footnote{The text will be found in Mansi, IX.488 \emph{sqq.}; and \emph{P. G.} LXXXVI.945 \emph{sqq.}

} It was subscribed by Menas, and the Pope and the other Patriarchs, including Peter of Jerusalem, signed it also.\texthebrew{⁠}\footnote{Liberatus, 22; Cassiodorus, \emph{De inst. div. litt.} c1 (\emph{P. L.} LXX p1111).

} Theodore Ascidas was in a difficult position. To refuse to accept the edict would have cost him his bishopric and his influence at court. He sacrificed his opinions and affixed his signature,\texthebrew{⁠}\footnote{Domitian of Ancyra also signed but afterwards retracted.

} but he had his revenge by raising a new theological question which was to occupy the stage of ecclesiastical politics for more than ten years.\footnote{After the death of Nonnus in 547, a schism arose among the Origenists. The Isochristoi of the New Laura, who held that in the \textgreek{\texthebrew{ἀ}προκατάστασις} or restitution after the general Resurrection men will be united with God as Christ is, were opposed to the Proto­ktistai or Tetradites (of the Laura of Firminus), whose names and views are obscure. Se Diekamp, \emph{op. cit.} 60 \emph{sqq.}

}

There was no theologian whose writings were more offensive to the Monophysites than Theodore of Mopsuestia, who was esteemed the spiritual father of Nestorianism. He had also

p384
written against Origen and was detested by the Origenists. To Theodore Ascidas, who was apparently a secret Monophysite as well as an Origenist, there could hardly be a greater triumph than to procure his condemnation by the Church.

Ascidas, warmly seconded by Theodora, persuaded the Emperor that he might solve the problem which had hitherto baffled him of restoring unity to the Church, by anathematising Theodore of Mopsuestia and his writings. This, he urged, would remove the chief stumbling-block that the Monophysites found in the Council of Chalcedon. For their objection to that Council was based far less on its dogmatic formula than on the countenance which it gave to a Nestorian like Theodore. For if the formula were consistent with Theodore's opinions, it would not be consistent with the doctrine of Cyril, and therefore could not admit of an interpretation that could ever be acceptable to the Monophysites. What Ascidas proposed was a rectification of the acts of Chalcedon, so as to make it clear that Chalcedonian orthodoxy had no leanings to Nestorianism. There were some other documents which it would be necessary to condemn at the same time: certain writings of Theodoret against Cyril, and a letter of Ibas, bishop of Edessa, in which Cyril was censured. Justinian was impressed by the idea, and acted promptly. In A.D. 546 he promulgated an Edict of Three Chapters, condemning (1) Theodore of Mopsuestia and his works; (2) specified works of Theodoret; and (3) the letter of Ibas.\texthebrew{⁠}\footnote{The Edict has not been preserved. Diekamp (\emph{op. cit.} 54) would date it to 543.

} In the subsequent controversy the expression ``Three Chapters'' was perverted to mean the condemned opinions, so that those who opposed the edict were said to defend the Chapters.

The eastern Patriarchs were at first unwilling to subscribe to this edict. It seemed a dangerous precedent to condemn the dead who could not speak for themselves. And was there any prospect that anything short of the repudiation of the Council of Chalcedon would satisfy the Monophysites?\texthebrew{⁠}\footnote{Cp. Duchesne, \emph{op. cit.} 395\texthebrew{‑}396.

} But the pressure of the Emperor induced the four Patriarchs to sign on the express condition that the Pope should be consulted.

On November 22, A.D. 545, during Totila's siege of Rome, Pope Vigilius was in the church of St. Cecilia in Trastevere, celebrating the anniversary of its dedication. In the middle of

p385
the ceremony a body of soldiers arrived, and an officer entered the church and presented Vigilius with a mandate to start immediately for Constantinople. He did not stay to finish the service, but accompanied the soldiers to the Tiber, where a ship was waiting. The congregation followed him and he pronounced the blessing which concluded the liturgy, but when the ship started, the crowd hurled missiles and maledictions. It looked as if the Pope were being carried off against his will, and general rumour ascribed his departure from Rome to the machinations of Theodora. But the sequel does not bear out this explanation. Vigilius was not taken to Constantinople under constraint. He went to Sicily, where he remained for ten months and made arrangements for sending provisions to Rome from the lands belonging to the pontifical patrimony. The truth seemed to be that the Emperor wanted Vigilius, that Vigilius was not reluctant to leave the besieged city, and that the scene in St. Cecilia was concerted in order to protect him from the reproach that he was voluntarily abandoning Rome.\footnote{This is the view of Duchesne.

}

In Sicily, the Pope was able to learn the opinion of western ecclesiastics on the Three Chapters of Justinian. They were unanimously opposed to the edict. Dacius, the archbishop of Milan, arrived from Constantinople, where he had lived for some years, and informed him that he had broken off communion with Menas. Supported by western opinion the Pope resolved to oppose the edict, and in autumn A.D. 546\texthebrew{⁠}\footnote{For the date see Procopius, \emph{B. G.} III.16.1; \emph{Lib. Pont., Vita Vig.} 297; \emph{Cont. Marcellini}, \emph{s.a.} 547; John Mal. XVIII.483.

} he set sail for Patrae, accompanied by Dacius. He travelled slowly, and when he reached Thessalonica he wrote a letter to Menas explaining his views and threatening to break off communion with him if he continued to support the Three Chapters.\texthebrew{⁠}\footnote{Facundus, \emph{Contra Mocianum}, \emph{P. L.} p862; \emph{Pro def. trium capit.} IV.3; \emph{ib.} p623.

} On January 25 ( A.D. 547) he arrived at the capital, where he was honourably and cordially received by the Emperor. He took up his quarters in the palace of Placidia, the residence of the Roman nuncios.

It was unfortunate for him that Pelagius was no longer at Constantinople. He sorely needed the guidance of a man of ability and tact. He had a learned adviser in Facundus, bishop of Hermiane in Africa, who was well acquainted with Greek, but the disposition and manners of Facundus were far from

p386
conciliating.\texthebrew{⁠}\footnote{Duchesne, p402.

} Vigilius himself was not much of a theologian, and he seems never to have been quite sure as to the merits of the controversy. He was pressed on one side by the Emperor and the Patriarch, on the other by western opinion. His vacillations, due both to intellectual and to moral weakness, presented a pitiable spectacle. In view of his past record, he cannot excite much compassion, but it is not uninteresting to read the story of a Pope trailing in the dust the dignity of the Roman see.

When the Patriarch Menas, who, notwithstanding his first hesitations, had become a warm supporter of the Imperial policy, refused to withdraw his subscription to the Three Chapters, Vigilius excommunicated him and his followers, but a reconciliation was soon effected by the intervention of Theodora,\texthebrew{⁠}\footnote{June 9, 547, Theophanes, A.M. 6039 (source, John Mal. cp. XVIII p483).

} and presently the Pope was assailed with doubts whether the Three Chapters were not justifiable. He read extracts from the works of Theodore of Mopsuestia, which the Greeks translated for him, and came to the conclusion that his doctrines were extremely dangerous. He would not indeed sign the edict; to do so, would concede to the Emperor the right to dogmatise on matters of faith. But he promised to declare an independent judgment, and in the meantime gave the Emperor and Empress written assurances that he intended to pronounce in the sense of the edict.\texthebrew{⁠}\footnote{\emph{Acta V. Conc.}, Mansi, IX.350.

} On Easter-eve , A.D. 548, he issued a Iudicatum \texthebrew{⁠}\footnote{Only some fragments are preserved, Mansi, IX.104\texthebrew{‑}105, 181.

} or pronouncement, addressed to Menas, condemning Theodore and the writings condemned in the edict, but carefully protecting the authority of Chalcedon.

The Papal decision created consternation in western Christendom. Facundus, bishop of Hermiane, published a learned treatise against the Three Chapters, on which he had been engaged.\texthebrew{⁠}\footnote{\emph{Pro defensione trium capitulorum.} In its original form it consisted of two books, which were presented to the Emperor, but was afterwards expanded into twelve. Facundus returned to Africa, took part in the African Council of 550 which excommunicated Vigilius, and avoided imprisonment by flight. In 571 he wrote the treatise \emph{Contra Mocianum} on the same subject. These two works are important sources for the story of the Three Chapters.

} The African Church dissolved communion with the Pope, and even Zoilus, Patriarch of Alexandria, who had provisionally subscribed to the edict, withdrew his signature and refused to accept

p387
the Iudicatum. The good opinion of the west was of more importance to Vigilius than the Emperor's favour, and, alarmed by the general outcry which his decision had provoked, he sought refuge in the expedient of a General Council. He told the Emperor that this was the only way of averting a schism, and persuaded him to consent to the withdrawal of the Iudicatum. But Justinian, before he agreed, made him swear on the Gospels and the nails of the Cross that he would use all his influence to procure the confirmation of the edict.\footnote{\emph{P. L.} LXIX.121\texthebrew{‑}122.

}

Justinian, however, took further measures before the meeting of the Council. He deposed from their sees the Patriarchs of Alexandria and Jerusalem, who refused to approve the Three Chapters, and he issued another edict ( A.D. 551) to the same purport as the former one. On the morning of its publication Theodore Ascidas and other Greek clergy visited the Pope in the Placidian palace. He urged them not to commit themselves to any judgment on the Imperial decree, but to await the decision of the Council. When they refused, he declined to receive them or enter their churches, and he excommunicated Menas and Ascidas. The rumour reached him that it was proposed to remove him by force from his residence, and he took refuge, along with the archbishop of Milan, in the sanctuary of SS. Peter and Paul near the palace of Hormisdas. Soldiers were sent to drag them away, and they clung to the altar. Vigilius was seized by his feet and beard, but he was a man of power ­ful build and in the struggle the altar gave way and fell to the ground crushing him under its weight. There was a cry of horror from the crowd which had gathered in the church, and the soldiers and their commander\texthebrew{⁠}\footnote{Vigilius calls him praetor, and if this is not a loose term, the Praetor of the Demes must be meant.

} retreated, abandoning their purpose (August).

The Emperor comprehended that he had gone too far. He sent assurances to the Pope and his clergy that they would be safe if they returned to the Placidian palace. They went back, and though no further violence was offered, the house was guarded like a prison. This became so intolerable that, two days before Christmas, Vigilius resolved to escape and fled under cover of darkness to the church of St. Euphemia in Chalcedon, the scene of the council which had been the origin of so many

p388
troubles. The Emperor then sent Belisarius, chosen doubtless on account of his old relations with the Pope, at the head of a distinguished deputation, to offer him sworn guarantees that he would be honourably treated. The Pope replied that the time for oaths was past; let the Emperor abstain from holding relations with Menas and Ascidas. His tone enraged Justinian, who wrote him a long unsigned letter full of menaces. Vigilius enjoyed the days of his sojourn at St. Euphemia in composing an Encyclical Epistle, addressed ``to the whole people of God,'' describing the violent treatment he had received, and declaring a profession of faith in which no mention was made of the Three Chapters. At length a new message arrived from the Emperor, again offering guarantees (Feb. 4, A.D. 552), but nothing came of it.\texthebrew{⁠}\footnote{Vigilius added an account of the interview to his Encyclical, which he signed on the following day. The text of the Encyclical will be found in \emph{P. L.} LXIX.53 \emph{sqq.}; the condemnation of Ascidas (dated in August 551), \emph{ib.} 59 \emph{sqq.}

} Some time afterwards the Pope published his sentence of excommunication against Menas and Ascidas and their followers.

This obstinate attitude wore out Justinian, and, not seeing how he could find any one to put in the place of Vigilius, he agreed with the Patriarch and his clergy that they should make submission to the Pope. They presented him with a declaration, couched in sufficiently humble terms, of their reverence for the Council of Chalcedon and the dogmatic Epistle of Leo, and he then returned to the Placidian palace.

The Emperor had hoped to avoid the convocation of a Council, but he resigned himself to the necessity before the end of the year. Menas died in August ( A.D. 552),\texthebrew{⁠}\footnote{John Mal. XVIII.486. This agrees with the Lists of the Patriarchs which assign sixteen years six months to Menas. Eutychius was consecrated immediately after his death. See Andreev, \emph{op. cit.} 175.

} and his successor Eutychius addressed a letter on the subject to the Pope, who replied favourably.\texthebrew{⁠}\footnote{\emph{P. L.} LXIX.63 \emph{sq.}, and 65 \emph{sqq.} (January 6, 553).

} Then the Emperor proceeded to issue notes convoking the bishops. From Gaul and Spain, from Illyricum and Dalmatia none came; and from Africa only those were allowed to attend on whom the Emperor thought he could count. It was clear that the Council would consist almost entirely of bishops of the Eastern Patriarchates.

The bishops duly arrived, but they were kept waiting at Constantinople for months before the Council met. The delay

p389
was due to the Pope, who, though he had originated the proposal of a Council, now declared that he would not take part in it. Afraid, at the last moment, of injuring irrevocably his authority in the eyes of the western churches, he had bethought himself of a \emph{via media.}\texthebrew{⁠}\footnote{Duchesne thinks that this compromise was probably suggested by Pelagius.

} He would condemn certain doctrines of Theodore of Mopsuestia without anathematising his person; but he would refuse to pass any judgment on the writings of Theodoret and Ibas, on the ground that their condemnation would bring discredit on the council of Chalcedon which had defended them. But he did not imagine that he would be able to induce the Council to adopt this compromise, and he therefore decided not to attend it but to issue his own judgment independently.

The meeting of the Council could not be indefinitely postponed, and at last the first session was held in the Secretariat of St. Sophia, on May 5, A.D. 553.\texthebrew{⁠}\footnote{The Acts are in Mansi, IX.173 \emph{sqq.} There were eight sessions, the last on June 2. The Origenistic heresy seems to have been discussed at meetings of the bishops previous to May 5, which did not form part of the proceedings of the Council proper, and at one of these conferences a letter of the Emperor was read which is preserved in the chronicle of George Monachus (ed. de Boor, 630 \emph{sqq.}). This is the result of Diekamp's investigations (\emph{op. cit.}). Origen is mentioned in the eleventh Anathema of the Council, but the fifteen canons against Origenistic doctrines (Mansi, IX.395 \emph{sqq.}) were drawn up at the previous meetings, and apparently were not specially confirmed by the Council.

} The proceedings opened by the reading of a letter of the Emperor reviewing the question of the Three Chapters. The assembly sent many deputations to the Pope requesting him to appear; he replied that he would send a written judgment on the question at issue.\texthebrew{⁠}\footnote{It is known as the \emph{Constitutum Vigilii}, and will be found in \emph{P. L.} LXVII \emph{sqq.}, and \emph{Coll. Avel., Ep.} 83.

} On May 14 it was ready, and Belisarius proceeded to the Placidian palace, but only to decline to transmit the document. A messenger of Vigilius then carried it to the Great Palace, but the Emperor refused to receive it, on the ground that if it confirmed the Three Chapters, it merely repeated what Vigilius had already declared and was therefore superfluous; and if it was unfavourable, it was inconsistent with his previous utterances and could carry no weight.

At a subsequent session, Justinian presented to the Council documents in which the Pope had approved of the Chapters of his edict, and then laid before the assembly an Imperial decree directing that the name of Vigilius should be struck out of the

p390
diptychs on account of his tergiversation and because he refused to attend the council. This was done.

The decrees of the Fifth Ecumenical council, which condemned a Pope, as well as Theodore of Mopsuestia and works of Theodoret and Ibas, were accepted without opposition. In the west they led to the banishment of some bishops,\texthebrew{⁠}\footnote{Like Victor Tonnennensis and Liberatus. The \emph{Breviarium} of Liberatus, which has been often quoted in the foregoing pages, was written with polemical intention against the Three Chapters. Primasius, who had succeeded the exiled Reparatus in the see of Carthage, yielded.

} and Pelagius, who had signed the document of the Pope, was imprisoned.

Vigilius found himself alone, and once more he revoked his latest decision. In yielding to the Emperor's wishes, he may have been moved by the fact that Narses had just completed the subjugation of the Ostrogoths and that his own place was at Rome. ``He chose among the different opinions which he had successively defended that which appeared most favourable to his personal interests and undoubtedly to those of his flock long deprived of its shepherd.''\texthebrew{⁠}\footnote{Duchesne, \emph{ib.} 422.

} At the end of six months he addressed to the Patriarch Eutychius a letter signifying his acceptance of the decrees of the Council (December 8, A.D. 553),\texthebrew{⁠}\footnote{Mansi, IX.413 \emph{sqq.}; \emph{P. L.} LXIX.121 \emph{sqq.}

} and then prepared a formal judgment in which he refuted the arguments alleged against the condemnations of the Three Chapters (February 26, A.D. 554).\texthebrew{⁠}\footnote{Mansi, IX.457 \emph{sqq.}

} The Emperor showed his satisfaction by conferring benefits on the Roman see,\texthebrew{⁠}\footnote{By a so\texthebrew{‑}called \emph{Pragmatica}, dated August 13, 554; \emph{Novellae}, App. VII (= \emph{Nov.} 164, ed. Zachariä).

} and in the following year the Pope set out for Italy. But he never saw Rome again. He died at Syracuse (June 7, A.D. 555), and his body was conveyed to Rome and buried in the Church of St. Silvester on the Via Salaria.

Pelagius had refused to follow Vigilius in his last recantation, and had written pamphlets against the Council. But he was soon to do even as Vigilius, when the Emperor, who valued his qualities, told him that he might succeed to the pontifical throne if he would accept the Council. He revised his opinions with little delay and was chosen and consecrated bishop of Rome.\texthebrew{⁠}\footnote{Duchesne's dates for the pontificate of Pelagius are April 16, 556 (not 555) to March 4, 561 (not 560). Cp. \emph{op. cit.} 428. At his consecration Pelagius declared a profession of faith in which he entirely ignored the Fifth Council.

} On this occasion, Justinian assumed the right, which had been

p391
exercised by the Ostrogothic kings, of confirming elections to the Roman see.\footnote{Cp. \emph{Liber Diurnus Rom. Pont.}, ed. Rozière (1869), LVIII p103 \emph{sqq.}

}

The fifth Ecumenical Council failed utterly in its main object of bringing about unity in the east, and it caused a schism in the west. Milan and Aquileia would know nothing of its decrees, and though political events, when the Lombards invaded Italy, forced Milan to resume communion with Rome, the see of Aquileia maintained its secession for more than a hundred and forty years.

It is possible that under the stress of persecution the Monophysitic faith might have expired, had it not been for the indefatigable labours of one devoted zealot, who not only kept the heresy alive, but founded a permanent Monophysitic Church. This was Jacob Baradaeus, who was ordained bishop of Edessa (about 541) by the Monophysitic bishops who were hiding at Constantinople under the protection of Theodora. Endowed with an exceptionally strong physical constitution, he spent the rest of his life in wandering through the provinces of the East, Syria, Mesopotamia, and Asia Minor, disguised as a beggar, and he derived the name Baradaeus\texthebrew{⁠}\footnote{From the Arabic form al\texthebrew{‑}Barādi'ā, which corresponds to the Syriac Burde'ayā. The chief sources for Jacob are the \emph{Life} by John of Ephesus, \emph{Comm.} c49; another \emph{Life} wrongly ascribed to the same writer, \emph{ib.} p203 \emph{sqq.}; and John Eph. \emph{Hist. ecc.} Part III. B. IV.13 \emph{sqq.}

} from his dress, which was made of the saddle-cloths of asses stitched together. His disguise was so effective, and his fellow-heretics were so faithful, that all the efforts of the Imperial authorities to arrest him were vain, and he lived until A.D. 578. His work was not only to confirm the Monophysites in their faith and maintain their drooping spirit, but also to ordain bishops and clergy and provide them with a secret organisation. His name has been perpetuated in that of the Jacobite Church which he founded.

The Fifth Ecumenical Council differed from the four which preceded it in that, while they pronounced on issues which divided Christendom, and which called for an authoritative decision of the Church, the Fifth dealt with a question that had been artificially created. Constantine, Theodosius the Great, his grandson, and Marcian had convoked ecclesiastical assemblies

p392
to settle successive controversies which had arisen in the natural course of theological speculation and which threatened to break up the Church into sects; the purpose of the Council which Justinian summoned was to confirm a theological decision of his own which was incidental indeed to a vital controversy, but only incidental. His object was to repair the failure of Chalcedon and to smooth the way to reunion with the Monophysites; and it may be said that the Three Chapters were entirely in the spirit of the orthodox theological school of his time.\texthebrew{⁠}\footnote{Cp. Loofs, \emph{op. cit.} 316. In one of his laws (\emph{Nov.} 132, A.D. 544) Justinian appeals to his books and edicts to prove his zeal for orthodox doctrine. His extant works (see \emph{P. G.} 86) include, besides the edicts against Origen and on the Three Chapters and some minor writings, a dogmatic refutation of the Monophysite doctrine, addressed to monks of Alexandria, \emph{c}. 542\texthebrew{‑}543; and an open letter against the defenders of Theodore of Mopsuestia. See Loofs, \emph{op. cit.} 310 \emph{sqq.}; W. H. Hutton, \emph{The Church of the Sixth Century}, 189 \emph{sqq.}

} But the question was provoked by himself; it was not one on which the decree of a General Council was imperatively required.

The importance of this episode of ecclesiastical history lies in the claim which Justinian successfully made to the theological guidance of the Church, a claim which went far beyond the rights of control exercised by previous emperors. Zeno had indeed taken a step in this direction by his Henotikon, but the purpose of the Henotikon was to suppress controversy, not to dictate doctrine. Justinian asserted the principle that doctrinal decisions could be made by Imperial edicts. An edict imposed upon the Church the orthodoxy of the Theopaschite formula; an edict condemned opinions of Origen; and, though the behaviour of Pope Vigilius forced the Emperor to summon a Council, the Council did no more than confirm the two edicts which he had issued on the Three Chapters. Justinian seems to have regarded it as merely a matter of policy and expediency whether theological questions should be settled by ecclesiastical synods or by Imperial legislation. Eastern ecclesiastics acquiesced in the claims of the Emperor when they adhered to the first edict on the Three Chapters, even though they made their adhesion conditional on the attitude of Rome; and at the synod of A.D. 536, while the assembled bishops said ``We both follow and obey the apostolic throne,'' it was also laid down by the Patriarch that nothing should be done in the Church contrary to the will of the Emperor.\footnote{Mansi, IX.970 \textgreek{προσήκει μηδ\texthebrew{ὲ}ν τ\texthebrew{ῶ}ν \texthebrew{ἐ}ν τ\texthebrew{ῇ ἁ}γιωτάτη \texthebrew{ἐ}κκλησί\texthebrew{ᾲ} κινουμένων παρ\texthebrew{ὰ} γνώμην α\texthebrew{ὐ}του κα\texthebrew{ὶ} κέλευσιν γενέσθαι.}}

p393
This Caesaro-papism , as it has been called, or Erastianism, to use the word by which the same principle has been known in modern history, was the logical result of the position of the Church as a State institution.

The Three Chapters was not the last theological enterprise of Justinian. In the last years of his life he adopted the dogma of aphthartodocetism, which had been propagated, as we have seen,\texthebrew{⁠}\footnote{See above,

p375 .

} by Julian of Halicarnassus, and had sown strife among the Monophysites of Egypt. This change of opinion is generally considered an aberration due to senility; but when we find a learned modern theologian asserting that the aphthartodocetic dogma is a logical development of the Greek doctrine of salvation,\texthebrew{⁠}\footnote{Harnack, \emph{History of Dogma}, IV.238, where it is suggested that Justinian was inclined to this heresy in the early years of his reign.

} we may hesitate to take Justinian's conversion to it as a sign that his intellectual power had been enfeebled by old age. The Imperial edict in which he dictated the dogma has not been preserved. The Patriarch Eutychius firmly refused to accept it, and the Emperor, not forgetting his success in breaking the will of Vigilius, caused him to be arrested (January 22, A.D. 565). He was first sent to the Island of the Prince and then banished to a monastery at Amasea.\texthebrew{⁠}\footnote{Eustathius, \emph{Vita Eutychii}, c5. Eutychius was succeeded in the Patriarchate by John of Sirmium, who held the see for twelve and a half years, and then Eutychius was restored (A.D. 577).

} The other Patriarchs were unanimous in rejecting the Imperial dogma. Anastasius of Antioch and his bishops addressed to the Emperor a reasoned protest against the edict. Their bold remonstrances enraged Justinian, and he was preparing to deal with them, as he had dealt with Eutychius, when his death relieved the Church from the prospect of a new persecution.\footnote{The sources for Justinian's heresy are:

Evagrius, \emph{H. E.} IV.39, 40 ; Eustathius, \emph{ib.} cc4, 5; Theophanes, A.M. 6057 (= Cramer, \emph{Anecd.} II.111), whose notice is probably derived from John Malalas; Michael Mel. \emph{Chron.} IX.34, who gives the texts of the Antiochene document (his ultimate source was probably John of Ephesus, \emph{H. E.} Part II);

John of Nikiu, \emph{Chron.} c94 , p399; Nicephorus Patriarcha, \emph{Chronogr.} p117, ed. de Boor. Victor Tonn., \emph{s.a.} 566, notices the deposition of Eutychius, but does not assign the reason. There is also a letter of Nicetius, bishop of Trier, to Justinian, reproaching him with his lapse into heresy in his old age (in ultima senectute tua), though the bishop appears to have had no clear idea as to the nature of the heresy (\emph{P. L.} LXVIII.378). An attempt to rescue the reputation of Justinian for unblemished orthodoxy was made by Richard Crakanthorp, who in 1616 published a pamphlet, \emph{Justinian the Emperor defended against Cardinal Baronius}, in which he sought to invalidate the evidence (which at that time mainly consisted of Evagrius

(p394)
and Eustratius). In 1693 Humphrey Hody refuted his arguments in a treatise entitled \emph{The Case of Sees Vacant by an Unjust or Uncanonical Deprivation Stated.} The thesis of Crakanthorp has been revived in modern times, with much greater ingenuity and learning, by W. H. Hutton (\emph{op. cit.} 205 \emph{sqq.}, and articles in \emph{The Guardian}, August 12, 1891, April 22, 1897; cp. my articles, \emph{ib.}, March 4, 1896, January 13, 1897). But the testimonies are too strong and circumstantial to be set aside or evaded. It may be noted that, according to Michael Mel. (\emph{loc. cit.}), Justinian was perverted by a monk of Joppa; and according to some he returned to the path of orthodoxy just before he died.

\texthebrew{◂} previous

next \texthebrew{▸}Images with borders lead to more information.

The thicker the border, the more information.

(Details here.)

UP TO:

J. B. Bury
Homepage

Latin \& Greek

Texts

LacusCurtius

Home}

\xchapter{Chapter XXIII}

Short URL for this page:

bit.ly/BURLAT23

Justinian is the only Emperor after Constantine, or at least after Julian the apostate, whose name is familiar to many who have never read a line about the history of the later Empire. He owes this fame to the great legal works which are associated with his name; and it may be suspected that some of those who have heard of the Digest and Institutions of Justinian think of him as a jurist and are hardly aware that he was an Emperor.

Justinian's legal achievements were twofold. By new legislation he brought to completion, in several important domains of civil law, the tendencies towards simplicity and equity which had been steadily developing for many centuries. This alone would have made his name remembered in the history of European law. But his chief work did not consist in legislative improvements. It consisted in redu­cing to order and arranging in manageable form the enormous and unwieldy body of Roman law as it existed.\footnote{Good general accounts of the legal works of Justinian will be found in Roby's \emph{Introduction to the Digest}, and in Bryce's article on Justinian in the \emph{Dict. of Christian Biography.}

}

Roman law, at this time, was of two kinds, which we may distinguish as statute law and jurisprudence; the statute law consisting of the Imperial constitution, and the jurisprudence of the works of the authoritative jurists who had written in the second and third centuries. Codification of the statue law was not a novelty. There had already been three Codes, the last of which, as we saw, was issued under the auspices of Theodosius II

p396
and his western colleague in A.D. 438. But a new collection, more compendious and up-to\texthebrew{‑}date, was a pressing need. The book of Theodosius was bulky, and was not always at hand to consult in the courts. Many of the enactments contained in it were wholly obsolete or had suffered modification, and in the seventy years which had elapsed since its appearance, a large number of new laws had been made.

It seems almost certain that Justinian had conceived the idea of compiling a new Code before he ascended the throne, not many months after his accession to power he issued a constitution addressed to the Senate, in which he announced the plan of a new collection of laws, edited up to date, with contradictions carefully eliminated, obsolete constitutions expunged, superfluous preambles or explanations omitted, words altered, eliminated, or added for the strategy of clearness; and appointed a commission of ten expert jurists to execute the work.\texthebrew{⁠}\footnote{February 13, A.D. 528. This constitution (Haec quae \texthebrew{—}) is prefixed to the Code.

} Of these ten, the pagan Tribonian, afterwards Quaestor, the Emperor's right hand in his great legal enterprises, and perhaps partly their inspirer, and Theophilus, professor of law at Constantine, were the most distinguished.

The commission must have worked hard, for the Code was completed and published in little more than a year. The Imperial constitution which introduced it to the world and made it authoritative was dated on April 7, A.D. 529.\texthebrew{⁠}\footnote{Prefixed to Code (Summa \texthebrew{—}); addressed to the Prefect of the City. [A fragment of the \emph{Index titulorum} of the 1st ed., found in Egypt, has been published this year by A. S. Hunt (\emph{Oxyrh. Papyri}, XV p217 \emph{sqq.}). It contains the Index to Bk. I, titles 11\texthebrew{‑}18, but 12 and 13 are omitted. Among other interesting points it supplies the name of the Emperor who enacted I.11.9: it was Anastasius.]

} But the Code which then appeared, and was arranged in ten Books, has not come down to us. Five and a half years later, an amended edition was issued,\texthebrew{⁠}\footnote{Prefixed to Code (Cordi nobis \texthebrew{—}); addressed to the Senate.

} arranged in twelve Books and including the new constitutions of the intervening period. It is this edition which we possess, and it contains 4652 laws.

In the meantime a more original and far more difficult work had been planned and completed. This was to reduce to order and consistency, and to present in a convenient form, the admirable body of jurisprudence which had been built up in the second and third centuries, the classical period of Roman law. The great

p397
lawyers of that age, who were licensed to give opinions and whose ``answers'' carried the weight of Imperial authority, explained and developed the rules of law which had been finally embodied in the Perpetual Edict of Hadrian. Their opinions ( responsa prudentium ) were scattered in many treatises, and they often differed. On many points antagonists might produce two opposite opinions, and on almost any the judge might be perplexed by inconsistent citations. The writings of five jurists soon came to obtain a predominant influence. These were Gaius, Papinian, Ulpian, Paulus, and Modestinus. The Emperor Constantine sought to diminish the practical inconvenience caused through the disagreements of these lawyers by exalting the authority of Papinian above Paulus and Ulpian. Valentinian III and Theodosius II passed an important measure, known as the Law of Citations, ordaining that the majority of opinions should determine the decision of the judge, and that, if they were equally divided, the ruling of Papinian should prevail.\footnote{A.D. 426, \emph{C. Th.} 1.4.3.

}

But the treatises of the recognised experts were so voluminous that in practice it was very difficult to administer good law. At most courts there was probably neither the necessary library nor the necessary learning available. It was a crying need, in the interests of justice, to make the opinions of the jurists easily accessible, and the idea was conceived and carried out by Justinian of meeting this want by ``enucleating the old jurisprudence.''

But one thing had to be done first. The Law of Citations imposed upon each judge the task of examining and correcting the opinions of the authorities when they disagreed. Plainly it would be much more convenient and satisfactory to have all important cases of disagreement settled once for all, so that the judge should have one clear ruling to guide him. Accordingly Justinian's lawyers drew up Fifty Decisions, which settled principal points of dispute.

This cleared the way for compiling an authoritative and consistent body of Jurisprudence. In the last month of A.D. 530, Justinian authorised a commission of sixteen lawyers, under the presidency of Tribonian, to set about the work.\texthebrew{⁠}\footnote{\emph{C. J.} 1.17.1 .

} They were to eliminate all contradictions and omit all repetitions, and when they had thus prepared the vast material they were to arrange

p398
it in one fair work, as it were a holy temple of Justice, containing in fifty Books the law of 1300 years. Tribonian seems to have adopted the practical expedient of dividing the commission into three committees, each of which digested and prepared a portion of the material.\texthebrew{⁠}\footnote{The composition of the Digest was elucidated by F. Bluhme, \emph{Die Ordnung der Fragmente in den Pandecten-titeln.} Bluhme showed that the arrangement of the Books was determined by Pythagorean theories of numbers. This was quite in the spirit of the time. Cp. Roby, preface to \emph{Introduction to the Digest.}

} Immense as the task was, it was completed in less than three years, and was published in December, A.D. 533.\texthebrew{⁠}\footnote{See the constitutions prefixed to the work (Omnem reipublicae \texthebrew{—}, and Dedit nobis \texthebrew{—}); and

\emph{C. J.} I.17.2 .

} The work was known as the Digest or the Pandects.\footnote{Digesta; \textgreek{πανδέκται.}\emph{C. J.} I.17.1 § 12 .

}

The Code and the Digest were each promulgated as an Imperial statute. They were to compound the whole body of valid law, except such Imperial constitutions as might subsequently be issued. All the books of the jurists were herewith rendered obsolete, as well as the Twelve Tables and the older Codes.

During the compilation of the Digest, Tribonian and his two most learned coadjutors, Theophilus, professor at Constantinople, and Dorotheus, professor at Berytus, prepared and published an official handbook of civil law for the use of students, the famous \emph{Institutions.} This manual reproduces the Commentaries of Gaius, the great jurist of the second century, but brings that work up to date by numerous changes, omissions, and additions. Like the Code and Digest, it is published with all the authority of a statute.\footnote{See the prefixed constitution (dated A.D. 533, November 21) addressed cupidae legum iuventuti. Gaius and the Institutions can be conveniently compared in Gneist's ed., in parallel columns (2nd ed., 1880).

}

With the publication of the Institutions and the Digest, the Emperor announced a reform of legal studies. The education of a student in the legal schools extended over five years. Justinian prescribed a rearrangement of the course,\texthebrew{⁠}\footnote{See the constitution Omnem (addressed to the professors of law, antecessoribus). The students of each year were distinguished by special names: 1st year, Dupondii; 2nd, Edictales; 3rd, Papinianistae; 4th, \textgreek{λύται}; 5th, \textgreek{προλύται}. Justinian stigmatises dupondii as a frivolous and ridiculous name, and orders that the freshmen shall henceforth be designated Iustiniani novi. In Zacharias, \emph{Life of Severus} (ed. Kugener), there is some interesting information about the life of law students at Berytus. It was a regular practice for the Edictales to ``rag'' the Dupondii.

} which was now to be confined to his own law books, and he abolished all

p399
the laws schools of the Empire except those of Constantinople and Berytus. This was intended to secure that the teaching should be in the hands of entirely competent persons.\footnote{Const. \emph{Omnem}, § 7 audivimus etiam in Alexandrina splendidissima civitate et in Caesariensium et in aliis quosdam imperitos homines devagare et doctrinam discipulis adulterinam tradere.

}

The Code, Digest, and Institutions form the principal parts of the Corpus Juris Civilis, on which the law of most European countries is based, and which has influenced English law, although it was never accepted in England. The fourth part of the Corpus consists of the later laws of Justinian, published after the second edition of the Code, and known as the Novels. It is perhaps surprising that the Emperor did not, in the course of the last thirty years of his reign, issue another edition of the Code, including the new constitutions. He promised to publish a collection of his Novels, but he never did so;\texthebrew{⁠}\footnote{See Const. \emph{Cordi nobis} (prefixed to Code): aliam congregationem \texthebrew{—} quae novellarum nomine constitutionum significetur. It has been conjectured that the death of Tribonian may have had something to do with the non-fulfilment of this purpose.

} and it was left to private jurists to collect them after his death.\texthebrew{⁠}\footnote{The basis of our text is a collection of 168 Novels, of which four are duplicates, seven belong to Justin II and Tiberius II, and four are edicts of Praetorian Prefects; so that it contains 153 laws of Justinian. Another collection formed the basis of the Latin epitome of Julian, which we possess (125 Novels). We have also the Authenticum, a Latin version used in the Middle Ages (134 Novels). But we know of other collections too (cp. Heimbach, \emph{Gr.-röm. Recht}, p199). At the end of the first text are eleven edicts (not counted as Novels); and we have also some other constitutions of Justinian, derived from various sources. These additional texts will be found as Appendices I and II to the edition of Schoell and Kroll.

} Thus the fourth part of the Corpus has not the same official character as the other three. The Novels testify to the growing disuse of Latin as the official language. Private Emperors, even Theodosius II, had occasionally issued constitutions in Greek; but in the reign of Justinian, Greek became the rule and Latin the exception. Nearly all the Novels (except those intended for publication in Africa and Italy) were drawn up in Greek.\footnote{Zachariä von Lingenthal attributes this change to the influence of John the Cappadocian (appendix to his ed. of \emph{Novellae}, pp6, 8). John was ignorant of Latin, and most of the laws were addressed to the Praet. Pref. of the East.

}

Many of the laws of Justinian are concerned with administrative reforms, which have claimed our attention in other pieces. Here we may consider how civil jurisprudence and criminal law developed under his predecessors and were completed or modified by him.

The civil legislation of Justinian forms in many respects the logical term of the development of Roman law. The old law of the Twelve Tables had undergone profound modifications, first by the judgments of the Praetors under the influence of the Ius gentium , and then by Imperial statutes. We may say that this development was marked by two general features. The law was simplified in form, and it was humanised in substance. Both these processes were mainly a consequence of the Imperial expansion of Rome. The acquisition of strange territories, the subjection of foreign peoples, had led to the formation of a second system of jurisprudence, the praetorian law; and this, which had the merit of greater elasticity, reacted upon the native civil law of Rome and eventually wrought considerable changes in it, both by mitigating some of its harsher features and by superseding some of its cumbrous forms. At later stages the process of simplification progressed, first by Caracalla's grant of Roman citizen ­ship to all the free subjects of the Empire (in A.D. 212), and secondly, at a later time, by the disappearance of the distinction between Roman and provincial soil, whereby it became possible to simplify the law of real property. The gradual changes in the spirit of Roman law responded generally to changes in public opinion, and the chief agency in educating Roman opinion and humanising the Roman attitude to life was undoubtedly Greek thought. The spirit of the \emph{De officiis} of Cicero illustrates how far Roman educated opinion had travelled during the last two centuries of the Republic.

The extension of Roman citizen ­ship to all freemen in the Empire did away with the ius Latii and the legal distinctions appertaining to it. But between the slaves and the citizens there still remained some intermediate classes, who were less than citizens and more than slaves. There were the Latini Iuniani, slaves who had been manumitted, but through some flaw in the process had not become citizens in the full sense, having neither the right to hold public office nor to marry a free person, and being unable to make a will or to inherit under a will. There were the dediticii , slaves who had undergone punishment for crime and were afterwards manumitted, but

p401
who, in consequence of their old offences, did not enjoy the full rights of a citizen and could not live within a hundred miles of Rome. a And there were persons in mancipii causa ; children whom their fathers had surrendered into slavery, in consequence of some misdemeanour which they had committed, and whose status differed from that of true slaves in that, if they were manumitted, they became not freedmen but freemen ( ingenui ). These three classes had little importance in the time of Justinian, but he finally did away with them, and thus consummated the simplification of personal status. There were now, in the eyes of the law, only two classes, citizens and slaves. Among citizens indeed the class of freedmen was still distinguished, but only by the obligations which a freedman owed to his patron, not by any civil disabilities. Formerly he could not be a senator or a magistrate, unless the restrictions were removed by imperial favour,\texthebrew{⁠}\footnote{Namely, by a decree granting restitutio natalium (which cancelled the rights of the patron), or ius anulorum aureorum. Justinian granted these to all freedmen, but reserved the patron's rights, \emph{Nov.} 78.

} nor could he marry a lady of senatorial family. Justinian abolished these disabilities.

In regard to slavery itself, the legislation of Justinian was also progressive. He repealed the Lex Fufia Caninia ( A.D. 8), which limited the number of slaves a master might manumit, and he abolished the restrictions which the Lex Aelia Sentia had imposed on the liberation of slaves under thirty years of age. The solemn forms of manumission\texthebrew{⁠}\footnote{Under the Empire, till the fourth century manumission could in general only be effected by
vindicta
or by testament. The recognition of Christianity introduced a new form, declaration in ecclesiis.

} ceased to be necessary; any signification of the intention to manumit was legally valid.

The

patria potestas

was one of the fundamental principles which underlay the fabric of Roman law, and nothing better illustrates the influence which the gradual humanising of public opinion exercised on legislation than the limitations which were successively placed upon the authority of the paterfamilias of the persons and property of those who were under his potestas . One of the last severities to disappear was the right of a father to surrender his children as slaves to any one whom they had wronged, a right of which he might be tempted to avail himself if he were unable or unwilling to pay compensation. This practice ( noxae deditio ) had practically disappeared before

p402
the sixth century, but was still legally recognised. Justinian abolished it formally, and his observations on the subject illustrate the tendency of Roman legislation.\texthebrew{⁠}\footnote{\emph{Inst.} IV.8 § 7. The influence of ``oriental'' (Greek) law on many of Justinian's reforms is traced and emphasised in P. Collinet's valuable \emph{Études historiques sur le droit de Justinien}, º vol. I, 1912. The Syro-Roman lawbook, edited by Bruns and Sachau (1880), is of much service in elucidating this subject.

} ``According to the just opinion of modern society, harshness ( asperitas ) of this kind must be rejected, and this practice has fallen utterly into disuse. Who will consent to give his son, far less his daughter, into noxal servitude? For a father will suffer, through his son, far more than the son himself, and in the case of a daughter such a thing is barred still more by consideration for her chastity.''

By the harshness of early law, all property acquired by persons in potestas belonged to the father. This was modified by successive provisions under the Empire, and, before Justinian, the father was entitled to the usufruct of the property which his son had independently acquired. If he emancipated his son, he retained one third as absolute owner. Justinian changed this law to the advantage of the children. He gave the father a life interest in half the son's property; but when the father died, it reverted to the son.\footnote{\emph{C. J.} VI.61.6 .

}

Justinian also simplified the process of emancipation. The ancient elaborate method of emancipating persons in potestas by fictitious sales was still in use. The Emperor Anastasius introduced, as an alternative method, emancipation by Imperial rescript, but this did not make the process easier, though it was highly convenient when the person to be emancipated was not residing in the same place as the paterfamilias. Justinian ``exploded'' the old fictitious process\texthebrew{⁠}\footnote{Fictione pristina explosa, \emph{Inst.} I.12 § 6. For the influence of Greek law on the simplification of form in this case and in that of adoption see Collinet, \emph{op. cit.} 52 \emph{sqq.}

} and enacted that a simple declaration of both parties in the presence of a magistrate or judge should be legally valid.

The history of marriage shows the same tendency to simplification. In early times a legal marriage between Roman citizens could be contracted in one of three ways: by a religious ceremony, which was confined to patricians ( confarreatio ); by a process of fictitious sale ( coemptio ); and by cohabitation for a year ( usus ). In each of these ways, the wife came under the power ( manus )

p403
of her husband; this power, in fact, was the fundamental feature in the legal conception of marriage. Towards the end of the second century A.D. , these old forms of contracting civil marriage had fallen into disuse.\texthebrew{⁠}\footnote{Confarreatio survived, but was only used by certain families who held certain religious offices.

} In other words, manus was obsolete. The Romans had adopted the matrimonium iure gentium , which had formerly been used by those who did not possess the right of marriage with Roman citizens ( ius connubii ).\texthebrew{⁠}\footnote{Caracalla's law abolished, ipso facto, all restrictions on the ius connubii.

} This union did not produce manus , nor did it originally give the father potestas over his children. It was quite informal; consent only was required. But as it came into use among Roman citizens, it was allowed to carry with it the patria potestas . Divorce by consent was the logical result of marriage by consent and the disuse of manus . So long as manus had constituted the legal relation, the husband had to emancipate his wife in order to effect a divorce.

But the disuse of manus , which had placed the wife in the position of a daughter, did not make her legally independent ( sui iuris ). She remained either under patria potestas or under guardian ­ship

(tutela) . The old theory was that a woman was not a person capable of legal action, and that if she were under neither potestas nor manus she must be legally represented by a guardian. Exceptions were made to this rule even in the time of Augustus;\texthebrew{⁠}\footnote{Lex Iulia et Papia Poppaea, A.D. 9.

} and the result of the growing belief that women were capable of acting for themselves was that by the fourth century perpetual guardian ­ship of females had disappeared.

If we turn from the law of persons to the law of property, we notice similar tendencies. When the distinction between Italian and provincial soil disappeared,\texthebrew{⁠}\footnote{\emph{Inst.} II.8.40; \emph{C. J.} VII.25 ;

VII.31 .

} the distinction also fell away between the full \emph{quiritary} owner ­ship, which applied only to Italian land, and the \emph{bonitary} owner ­ship granted to the actual proprietors of provincial land, of which the supreme owner was the Roman people. The curious classification of property, which had played a great part in the old law, as

res mancipi

(real property in Italy, slaves, the chief domestic animals)\texthebrew{⁠}\footnote{Also rural servitudes. The word res has of course a much wider connotation than property.

} and

res nec mancipi , was abandoned and abolished; and the conveyance

p404
of property was simplified by the disappearance of the ancient and cumbrous civil methods ( mancipatio

and

in iure cessio ) which were superseded by the natural process of simple delivery ( traditio ). Full owner ­ship

(dominium)

could now be acquired by delivery. It could also be acquired by long possession or

usucapio . This method of acquisition had formerly been inapplicable to provincial land (because the dominium belonged to the Roman people), and the praetors had introduced an equivalent institution ( longi temporis praescriptio ), which was extended to all kinds of property. Justinian simplified the law by applying the second method to land which could be acquired by prescription after ten or (in some cases) twenty years, and the first method to moveables, possession of which for three years produced full owner ­ship.\footnote{\emph{Inst.} II.6; \emph{C. J.} VII.31.1 ;

VII.33.12 .

}

The governing conception in the Roman jurisprudence which concerned the family was the relation ­ship known as

agnatio . This untranslatable term is defined by Roman lawyers as kinship

(cognatio)

through males, but perhaps its scope is more clearly explained by saying that agnates were those who were under the patria potestas of the same person, or would have been so, if he were alive.\footnote{For the benefit of readers who are not familiar with the conception, it may be explained that if we start with a paterfamilias (A), a first group of agnates is formed by his sons and daughters, the children of his sons, the children of the sons of his sons, etc. This group is determined by a common relation to A. A second group of agnati proximi is constituted by those who were in the potestas of A's father; namely, A's brothers and sisters, the sons and daughters of the brothers, and so on. A third and more remote group is formed by those who were in the potestas of A's grandfather. And so on in ascending scale. Adoption carried with it inclusion among the agnates.

}

The most important sphere in which agnatio operated was the law of inheritance. When a man died without making a will, his heirs at law were in the first instance those persons (children, grandchildren, etc.) who were in his potestas and whom his death automatically rendered independent ( sui juris ). These were called sui heredes and did not include sons whom he had emancipated before his death or married daughters. If there were no sui heredes the inheritance passed to the nearest agnates; and if these failed to the gens . The two most serious defects in this system were the exclusion of sons and daughters who had passed out of the potestas of the deceased before his death, and the disqualification of cognates who were not also agnates. The

p405
Praetors devised expedients to remedy these hardships and to introduce new rules of succession which favoured cognation at the expense of agnation; but it was reserved for Justinian finally to lay down a scheme of intestate succession, which prevails in most European countries to\texthebrew{‑}day.\footnote{\emph{Nov.} 118; \emph{Nov.} 127.

}

By this reform the first heirs to an estate are the cognate descendants of the deceased, that is, his sons and daughters, their children, grandchildren, etc. The children inherit in equal shares; grandchildren only come in if their parent is dead, and divide his or her share. One trace, indeed, of the agnate system remains; adopted children count as natural.

If there are no descendants, the full brothers and sisters of the deceased, or the next nearest cognates, inherit, and dead brothers and sisters are represented by their issue, in the same way as in the former case. Failing heirs of this group, half-brothers and half-sisters have the next claim; after this, other collaterals.\footnote{In the case of an intestate freedman, the patron's rights had to be considered and the law was different. His heirs were (1) natural descendants; (2) his patron; (3) the patron's children; (4) the patron's collaterals.

}

In this legislation, there is no recognition of the claim of a wife or of a husband. The theory was that the wife was adequately provided for by her dowry; but Justinian enacted that a poor widow should inherit a quarter of her husband's estate.

Nor was the law of inheritance under wills left unaltered. Hitherto if a testator failed to make any provision for his near kin, the aggrieved relatives had to seek a remedy by a process known as ``complaint against an undutiful will.''\texthebrew{⁠}\footnote{Querela inofficiosi testamenti.

Thayer's Note: For details and references, see under the article

Testamentum

in Smith's \emph{Dictionary of Greek and Roman Antiquities}.

} Justinian obliged a testator to leave his children, if they were four or fewer, at least one third; if they were more than four, at least one half of his estate;\texthebrew{⁠}\footnote{\emph{Nov.} 18.

} and bound him further to institute as his heirs those descendants who would be his heirs in case of an intestacy, unless he could specify some cause for disinheriting them which would appear reasonable in the eyes of the law.\footnote{\emph{Nov.} 115. It may be added that Justinian also simplified the law of Legacies, and the law of Trusts, by assimilating all legacies to one another (there used to be four different kinds), and by pla­cing legacies and trusts on the much the same footing. Cp. \emph{Inst.} II.20 § 2, and 23 § 12.

}

The jurists of Justinian also introduced a simple and final remedy for the hardship of the ancient law, by which the heir

p406
was made responsible for the liabilities of the deceased even if they exceeded the value of the estate which he inherited. He might of course refuse to accept the inheritance, but if he accepted it he assumed, as it were, the person of the deceased, and any property he otherwise possessed was liable for debts which the estate could not meet. Of this law, which may well be considered one of the asperities of Roman antiquity, various modifications had been devised to meet particular cases, but they were inadequate. Justinian's ``benefice of inventory'' solved the difficulties. The heir was required to make an inventory of the estate and complete it within two months of the decease; if he did this, the estate alone was liable.\footnote{\emph{Ib.} 19.6.

}

The history of the law of divorce may be considered separately, for the legislation on this subject under the autocracy forms a remarkable and unpleasing exception to the general course of the logical and reasonable development of Roman jurisprudence. Here ecclesiastical influence was active, and the Emperors from Constantine to Justinian fluctuated between the wishes of the Church on one side, and on the other common sense and Roman tradition. The result was a confusion, no less absurd to a lawyer's sense of fitness than offensive to the reason of ordinary men. The uncertainty and vacillation which marked the Imperial attempts at compromise was aggravated by the fact that the ecclesiastics themselves had not yet arrived at a clear and definite doctrine, and were guided now, as later, not by any considerations of the earthly welfare of mankind, but by inconsistent texts in the New Testament\texthebrew{⁠}\footnote{Matthew v.32, against Mark x.2\texthebrew{‑}12, Luke xvi.18. Gregory of Nazianzus expressed the general ecclesiastical disapprobation of divorce when he wrote (\emph{Ep.} 176), ``Divorce is altogether displeasing to our law, though the laws of the Romans ordain otherwise.''

} which they were at some loss to reconcile.

Roman law recognised two ways in which a marriage could be dissolved \texthebrew{—} divorce by mutual consent, and the repudiation of one spouse by the other.\texthebrew{⁠}\footnote{Divorce bona gratia and repudium.

} Divorce by mutual consent was always regarded as a purely private matter and was never submitted to a legal form, and even the Christian Emperors before Justinian did not attempt to violate the spirit of the Roman law of contract by imposing any limitations. It was reserved for Justinian to prohibit it, unless the motive was to allow one of the spouses

p407
to embrace a life of asceticism.\texthebrew{⁠}\footnote{\emph{Nov.} 117, § 10, A.D. 542. Six years before he had confirmed the old practice on the ground that all ties can be dissolved, \textgreek{τ\texthebrew{ὸ} δεθ\texthebrew{ὲ}ν \texthebrew{ἅ}παν λυτόν}, \emph{Nov.} 22, § 3.\texthebrew{—} It is interesting to read records of actual divorces. Thus in \emph{P. Cairo}, II. No. 67154, we have a contract of divorce by mutual consent between Fl. Callinicus, a notary, and Aurelia Cyra. They state that they have quarrelled, they ascribe their quarrel to a malignant demon (\textgreek{σκαιο\texthebrew{ῦ} δαίμονος}), and agree to separate permanently by the ``written contract of dissolution which has the force of a repudium.'' It is agreed that they are to have the care of their son Anastasius in common (\textgreek{\texthebrew{ἀ}ν\texthebrew{ὰ} μέσον \texthebrew{ἀ}μφοτέρων}). This amicable document is signed for the wife (who cannot write) and by several witnesses. It belongs to the reign of Justinian, evidently before the law of A.D. 542. \texthebrew{—} No. 67153 and No. 67155 are examples of one-sided divorce, in the form of a letter of repudiation from the husband to the wife, the former in the year 568, the latter undated.

} This arbitrary and rigorous innovation was intolerable to his subjects, and after his death his successor was assailed by numerous petitions for its repeal. The domestic misery resulting from incompatibility of temper was forcibly represented to him, and he restored the ancient freedom as a concession owing to the frailty of human nature.\footnote{This law of Justin is included in the collection of Justinian's Novels: \emph{Nov.} 140, A.D. 566.

}

One-sided divorce had been equally unfettered; Augustus only required that the partner who decided to dissolve the marriage should make a formal declaration to this effect in the presence of seven citizens. Constantine introduced a new and despotic policy. He forbade one-sided divorce entirely except for a very few specified reasons. A woman was only permitted to divorce her husband, if he was found guilty of murder, poisoning, or the violation of tombs. If she separated herself from him for any other reason, she forfeited her dowry and all her property to the very bodkin of her hair, and was condemned to be deported to an island. A man might divorce his wife for adultery, or if she were guilty of preparing poisons, or of acting as a procuress. If he repudiated her for any other reason he was declared incapable of contracting a second marriage.\texthebrew{⁠}\footnote{\emph{C. Th.} III.16.1

(A.D. 331).

} This cruel law was but slightly softened by Honorius,\texthebrew{⁠}\footnote{\emph{Ib.} 2 (A.D. 421). If a woman divorced her husband for immorality or faults that were not legally heinous, she forfeited the dos and donatio, and could not remarry; a man who divorced his wife for culpa morum non criminum, might remarry after two years. In other cases, Constantine's penalties remained in force.

} but in the reign of Theodosius II reason and Roman legality prevailed for a moment. The legal advisers of that Emperor persuaded him that in the matter of divorce ``it is harsh to depart from the governing principle of the ancient laws,'' and he abolished all the restrictions and penalties which his Christian predecessors had imposed.\texthebrew{⁠}\footnote{Theodosius, \emph{Nov.} 12, A.D. 439.

} But

p408
this triumph of reason and tradition was precarious and brief. Ten years later, the same Emperor, under contrary influence, did not indeed venture to revive the stringent laws which he had abolished, but attempted a compromise between the old Roman practice and the wishes of the Church.\texthebrew{⁠}\footnote{\emph{C. J.} V.17.8 , A.D. 449.

} He multiplied the legitimate grounds for divorce. If a man was condemned for any one of nine or ten serious crimes, if he introduced immodest women into his home, if he attempted to take the life of his consort or chastise her like a slave, she was justified in repudiating him. If she dissolved marriage on any other ground, she was forbidden to remarry for five years.\texthebrew{⁠}\footnote{Anastasius reduced this to one year, A.D. 497, \emph{ib.} 9. The woman of course forfeited both dowry and donation. It may be mentioned that Justinian enacted (\emph{Nov.} 97) that the amount of the donation shall be exactly equal to that of the dowry; this was due to the influence of Greek law, Collinet, \emph{op. cit.} 145 \emph{sqq.}

} A woman, guilty of similar crimes, might be divorced, or if she sought her husband's life, or spent a night abroad without good cause, or attended public spectacles against his command. He might divorce her for adultery, but she could not divorce him. The husband who dissolved the marriage for any other than the specified reasons, was obliged to restore the dowry and the donation.

In his early legislation Justinian made no serious change in the law of Theodosius, but he added some new grounds for divorce, permitting a marriage to be dissolved if the husband proved to be impotent, or if either partner desired to embrace an ascetic life.\texthebrew{⁠}\footnote{\emph{C. J.} V.17.10 and 11 ; \emph{Nov.} 22 (536).

} But the Emperor soon repented of the comparative liberality of these enactments, and his final law, which deals comprehensively with the whole subject, exhibits a new spirit of rigour, though it does not altogether revive the tyrannical policy of Constantine.\texthebrew{⁠}\footnote{\emph{Nov.} 117 (542).

} The causes for which a husband may dissolve the union and retain the dowry, and for which a wife may dissolve it and receive the dowry with the donation, are reduced in number;\texthebrew{⁠}\footnote{But some new causes are added in the wife's favour; she may dissolve the marriage if her husband is persistently and flagrantly unfaithful, if he falsely accused her of adultery, or if he is detained in captivity in a foreign country.

} no release is allowed for a partner guilty of a public crime, except in the case of treason. A woman who repudiates her husband on other than the legal grounds is to be delivered to the bishop and consigned to a monastery. A man, in the same case, suffers only in his pocket. He forfeits the

p409
dowry, the donation, and a further sum equal to one-third of the donation. But this disparity of treatment was afterwards altered, and the husband was also liable to incarceration in a monastery.\footnote{\emph{Nov.} 134, § 11 (556).

}

The general tenor of these enactments of Justinian, though they were temporarily set aside in the eighth and ninth centuries, remained in force throughout the later period of the Empire, and the ecclesiastics never succeeded in bringing the civil into harmony with the canonical law which pronounced marriage indissoluble, and penalised a divorced person who married again as guilty of adultery.

This was perhaps the only department in which the Church exercised an influence on the civil law. It did not aim at nor desire any change in the laws concerning slaves, for slavery was an institution which it accepted and approved.\texthebrew{⁠}\footnote{For the views of the Fathers and the Church on slavery see Carlyle (R. W. and A. J.), \emph{A History of the Mediaeval Political Theory in the West}, I.116 \emph{sqq.} A full history of the \emph{Roman Law of Slavery} will be found in W. W. Buckland's treatise with that title (1908).

} In practice, of course, it encouraged mitigation of the slave's lot, but there it was merely in accord with general public opinion. Enlightened pagans had been just as emphatic in their pleas for humanity to slaves as enlightened Christians,\texthebrew{⁠}\footnote{Take, for instance, the views expressed by the pagan Praetextatus in the \emph{Saturnalia} of Macrobius

Macrob. \emph{Sat.}
(I.11) . Dill observes (\emph{Roman Society}, p136): ``The contempt for slaves expressed by S. Jerome and Salvianus is not shared by the characters of Macrobius.''

} and for the growing improvement in the conditions of slavery since the days of Cicero, the Stoics are perhaps more responsible than any other teachers. In this connexion it may be added, though it does not concern the civil law, that the Church happily failed to force upon the State its unpractical policy of prohibiting the lending of money at interest.\texthebrew{⁠}\footnote{On laws on interest see above,

Vol. I p55, n1 ;

Vol. II p357 .

} In the sphere of criminal law, as we shall now see, it intervened effectively.

The criminal law of the Empire, which was chiefly based on the legislation of Sulla, Pompey, and Augustus, had been little altered or developed under the Principate;\texthebrew{⁠}\footnote{Roman jurisprudence did not draw a capital distinction between civil and criminal law. What corresponds to our criminal law came partly under private (privata delicta)

(p410)
partly under public law (publica iudicia). See
\emph{Digest}, Bks. 47, 48 ; \emph{Instit.} IV.18;

\emph{C. J.} Bk. 9 ; \emph{C. Th.} Bk. 9. The subject is treated exhaustively in Mommsen's \emph{Römisches Stafrecht.}

} and the Cornelian

p410
laws on murder and forgery, the Pompeian law on parricide, the Julian laws on treason, adultery, violence, and peculation, were still the foundation of the law which was in force in the reign of Justinian. Such minor changes as had been made before the reign of Constantine were generally in the direction of increased severity. This tendency became more pronounced under the Christian Emperors. Two fundamental changes were introduced by these rulers by the addition of two new items to the list of \emph{public} crimes, seduction and heresy; but in those domains of crime which we should consider the gravest there were no important alterations.\footnote{Three legislative acts may be noticed. (1) For
plagium

(which may be roughly translated ``kidnapping''), the penalty before Constantine had been exile with confiscation of half the property of the condemned. Constantine made it death,

\emph{C. Th.} IX.18.1 , but Justinian allowed this punishment only in aggravated cases (\emph{Inst.} IV.18). (2) Anonymous libellous publications: Constantine made the penalty death instead of deportation,

\emph{C. Th.} IX.34 , but Justinian rejected this law. (3) Arcadius

(\emph{C. Th.} IX.14.3)

extended the law of maiestas (treason) to include attacks on the persons of senators, members of the Imperial Council, and all Imperial officials. Previously only magistrates (consuls, praetors, etc.) were protected from personal violence by the law of
perduellio .

}

Ordinary murder,\texthebrew{⁠}\footnote{Except in the case of parricide (murder of near relatives), for which Constantine

(\emph{C. Th.} IX.15.1)

revived the ancient punishment of the culleus, or sack, in which the criminal was sewn up in the society of snakes (serpentum contuberniis misceatur) and drowned.

} for instance, was punished by banishment\texthebrew{⁠}\footnote{Originally banishment from Italy with ``interdiction of fire and water.'' Under the Principate deportation to an island replaced this penalty. But the statement in the text applies only to freemen of the better classes; death was inflicted not only on slaves, but also on men of the lower classes \texthebrew{—} a distinction which will be explained below.

} under Justinian as under Augustus, and in the penalties for treason, arson, sorcery, forgery and kindred offences, theft and robbery in their various forms, violence, false witness, there was little change. In contrast with this conservatism, a new spirit animated Constantine and his successors in their legislation on sexual offences, and the inhuman rigour of the laws by which they attempted to suppress sexual immorality amazes a modern reader of the Codes of Theodosius and Justinian. Adultery, which in civilised countries to\texthebrew{‑}day is regarded as a private wrong for which satisfaction must be obtained in the civil courts, had been elevated by Augustus to the rank of a public offence, and the injured husband who let the adulterer go free or compounded

p411
with him for the injury, was liable to the same penalty as if he had himself committed the crime. The penalty consisted in the deportation of the guilty partners to separate islands. Augustus assuredly did not err on the side of leniency, but his severity did not satisfy Constantine, who made death the penalty of adultery.\texthebrew{⁠}\footnote{\emph{C. Th.} IX.7.1, 2 ;

II.36.4 . In \emph{Inst.} IV.18 the death penalty is erroneously ascribed to the Lex Iulia. The texts which Mommsen cites to show that this penalty had been introduced in the third century (\emph{op. cit.} 699, n3) do not appear to be decisive.

} Perhaps this law was seldom enforced;\texthebrew{⁠}\footnote{Ammianus

(XXVIII.1.16)

mentions one case, but as an instance of the inhuman cruelty of Valentinian I.

} and Justinian relaxed it by condemning the guilty female to be immured in a nunnery.\texthebrew{⁠}\footnote{\emph{Nov.} 77 ;

141 .

} The crime of incest, or marriage of persons within forbidden degrees, was usually punished by deportation; the Christian Emperors sought both to aggravate the penalty and to extend the prohibitions. Constantine imposed the penalty of death on marriage with a niece,\texthebrew{⁠}\footnote{\emph{C. Th.} III.12.1 . Claudius had permitted, and set the example of, the union of a niece with her \emph{father's} brother.

} and forbade unions with a deceased wife's sister or a deceased husband's brother.\texthebrew{⁠}\footnote{\emph{Ib.} 2.

} The savage legislator Theodosius I prohibited the marriage of first cousins, and decreed for those who were guilty of this or any of the other forbidden alliances, the penalty of being burned alive and the confiscation of their property.\texthebrew{⁠}\footnote{Ambrose, \emph{Ep.} 60;

\emph{C. Th.} III.12.3

(cp. Augustine, \emph{De civ. Dei}, XV.16 ).

} There were limits to the patience of the Roman public under the autocracy. Theodosius was not long in his grave before his son Arcadius cancelled these atrocious penalties,\texthebrew{⁠}\footnote{\emph{C. Th.} \emph{ib.} (A.D. 396).

} and some years later the same Emperor rescinded the prohibition of the marriage of cousins.\footnote{\emph{C. Th.} V.4.19

(A.D. 405); cp. \emph{Instit.} I.10.4.

}

The abduction of a female for immoral purposes, if not accompanied by violence, was, under the Principate, regarded as a private injury which entitled the father or husband to bring an action. Constantine made the abduction of women a public crime of the most heinous kind,\texthebrew{⁠}\footnote{See Mommsen, \emph{op. cit.} 701\texthebrew{‑}702;

\emph{C. Th.} IX.24.1 .

} to be punished by death in a painful form. The woman, if she consented, was liable to the same penalty as her seducer; if she attempted to resist, the lenient lawgiver only disqualified her from inheriting. If the nurse who was in charge of a girl were proved to have encouraged her to yield to a seducer, molten lead was to be poured into

p412
her mouth and throat, to close the aperture through which the wicked suggestions had emanated. Parents who connived at abduction were punished by deportation. This astonishing law, with slight mitigation,\texthebrew{⁠}\footnote{\emph{C. Th.} IX.24.2

(Julian punished the offence by banishment,

Ammianus, XVI.5.12 .)

} remained in force, and was extended to the seduction of women who had taken vows of chastity. Justinian made a new law on the subject, but the essential provisions were the same.\footnote{\emph{C. J.} IX.13.1 .

}

Unnatural vice was pursued by the Christian monarchs with the utmost severity. Constantius imposed the death penalty on both culprits, and Theodosius the Great condemned persons guilty of this enormity to death by fire.\texthebrew{⁠}\footnote{\emph{C. Th.} IX.7.3 and 6 . In the third century the seduction of a puer praetextatus was punished by death (\emph{Digest}, XLVII.11.1, 2). Measures, but we do not know what, were taken by the Emperor Philip to suppress this vice

(\emph{Hist. Aug.} XVIII.24.4) .

} Justinian, inspired by the example of the chastisement which befell ``those who formerly lived in Sodom,'' and firmly believing that such crimes were the immediate causes of famines, plagues, and earthquakes, was particularly active and cruel in dealing with this vice. In his laws,\texthebrew{⁠}\footnote{\emph{Nov.} 77 ;

141 .

} he contented himself with imposing the penalty of death, but in practice he did not scruple to resort to extraordinary punishments. It is recorded that senators and bishops who were found guilty were shamefully mutilated, or exquisitely tortured, and paraded through the streets of the capital before their execution.\footnote{John Malalas, XVIII p436 ;

Procopius, \emph{H. A.} 11, \emph{ad fin.} ; Theodosius of Melitene, \emph{Chron.} p90 \textgreek{το\texthebrew{ὺ}ς μ\texthebrew{ὲ}ν \texthebrew{ἐ}καυλοτόμησε το\texthebrew{ῖ}ς δ\texthebrew{ὲ} καλάμους \texthebrew{ὀ}ξε\texthebrew{ῖ}ς \texthebrew{ἐ}μβάλλεσθαι ε\texthebrew{ἰ}ς το\texthebrew{ὺ}ς πόρους τ\texthebrew{ῶ}ν α\texthebrew{ἰ}δοίων \texthebrew{ἐ}κέλευσεν κα\texthebrew{ὶ} γυμνο\texthebrew{ὺ}ς κατ\texthebrew{ὰ} τ\texthebrew{ὴ}ν \texthebrew{ἀ}γορ\texthebrew{ὰ}ν θριαμβεύεσθαι}. Cp. Michael Syrus, IX.26. The enemies of Justinian alleged that the trials were a farce, that men were condemned on the single testimony of one man or boy, who was often a slave; and the victims were selected because they belonged to the Green Faction, or were very rich, or had given some offence to the Emperor (Procop. \emph{loc. cit.}). No account was taken by the law of female homosexuality. Abbesses were obliged to take precautions against it in nunneries. See

the \emph{Epistle} of Paul Helladicus, abbot of Elusa (p21) , one of the most unsavoury documents of Christian monasticism. He lived in the sixth century.

}

The disproportion and cruelty of the punishments, which mark the legislation of the autocracy in regard to sexual crimes, and are eminently unworthy of the legal reason of Rome, were due to ecclesiastical influence and the prevalence of extravagant ascetic ideals. That these bloodthirsty laws were in accord with ecclesiastical opinion is shown by the code which a Christian

p413
missionary, untrammelled by Roman law, is reported to have imposed on the unfortunate inhabitants of Southern Arabia.

We saw how in the reign of Justin, Christianity was established in the kingdom of the Himyarites by the efforts of the Christian king of Ethiopia. When Abram was set upon the throne, Gregentius was sent from Alexandria to be the bishop of Safar, the chief city of the Himyarites.\texthebrew{⁠}\footnote{See above,

p327 .

} The laws which Gregentius drew up in the name of Abram are preserved. Doubts of their authenticity have been entertained; but even if they were never issued or enforced, they illustrate the kind of legislation at which the ecclesiastical spirit, unchecked, would have aimed. It is characteristic that sexual offences occupy a wholly disproportionate part of the code. Fornication was punished by a hundred stripes, the amputation of the left ear, and confiscation of property. If the crime was committed with a woman who was in the potestas of a man, her left breast was cut off and the male sinner was emasculated. Similar but rather severer penalties were inflicted on adulterers. Procurers were liable to amputation of the tongue. Public singers, harp-players , actors, dancers, were suppressed, and any one found practising these acts was punished by whipping and a year's hard labour. To be burned alive was the fate of a sorcerer. Severe penalties were imposed for failing to inform the public authorities of a neighbour's misconduct. On the ground of St. Paul's dictum that the man is the head of the woman, cruel punishments were meted out to women who ventured to deride men.\footnote{\emph{Homeritarum Leges}, \emph{P. G.} LXXXVI.1.581 \emph{sqq.} One curious provision regards early marriages; parents who do not arrange for the future marriage of their children when they are between the ages of ten and twelve are liable to a fine.

}

Perhaps the greatest blot in Roman criminal law under the Empire, judged by modern ideas, was the distinction which it drew, in the apportionment of penalties, between different classes of freemen. There was one law for the rich, and another for the poor. A distinction between the honourable or respectable, and the humble or plebeian classes was legalised,\texthebrew{⁠}\footnote{Honesti, and humiles (tenuiores, plebeii; \textgreek{ε\texthebrew{ὐ}τελε\texthebrew{ῖ}ς}). Mommsen, \emph{op. cit.} 1032 \emph{sqq.} A member of the higher classes could be degraded to the lower (reiectus in plebem),

\emph{C. Th.} VI.22.1 . The distinction was based on status, not on income, but it virtually discriminated the richer from the poorer.

} and different treatment was meted out in punishing criminals according to the class to which they belonged. The privileged group

p414
included persons of senatorial and equestrian birth, soldiers, veterans, decurions and the children of decurions; and on such persons milder penalties were inflicted than on their fellow-citizens of inferior status. They were, in general, exempt from the degrading and painful punishments which were originally reserved for slaves. If a man of the higher status, for instance, issued a forged document, he was deported, while the same crime committed by a poor man was punished by servitude in the mines.\texthebrew{⁠}\footnote{In many cases death was inflicted on the humiles where the honesti were only deported. See the table showing the disparity of punishments in the first half of the third century (based on the \emph{Sententiae} of Paulus), in Mommsen, \emph{op. cit.} 1045 \emph{sqq.}

} The general principle, indeed, of this disparity of treatment was the extension of servile punishments to the free proletariate, and it appears also in the use of torture for the extraction of evidence. Under the Republic freemen could not legally be tortured, but under the Empire the question was applied to men of the lower classes as well as to slaves.\footnote{Mommsen, 406 \emph{sq.} In cases of treason, and magic, and forgery, indeed, any citizen might be tortured, and many instances are recorded.

}

The normal mode of inflicting death on freemen was decapitation by the sword. But more painful modes of execution were also prescribed for certain offences.\texthebrew{⁠}\footnote{Such modes are often designated in the laws as summum supplicium. Four are recognised under the autocracy: (1) the gallows, furca; (2) fire; (3) the sack, for parricides; this had fallen out of use but was revived by Constantine; (4) exposure to wild beasts; this was an alternative to (1) or (2) and could be inflicted only when there happened to be a popular spectacle immediately after the condemnation. Crucifixion was discontinued under the Christian emperors.

} Sorcerers, for instance, were buried alive, and deserters to the enemy incurred the same penalty, or the gallows. In some cases, as for treason, the painful death was inflicted only on people of the lower class;\texthebrew{⁠}\footnote{Cp. \emph{Digest}, XLVIII.19.28. The milder form of banishment (relegatio) allowed to the condemned a choice of the place of his domicile, and did not involve disfranchisement.

} and in some, persons of this status were put to death while persons of higher rank got off with a sentence of deportation. The privileged classes were also exempted from the punishment of being destroyed by wild beasts in the arena. Next to death, the severest penalty was servitude in the mines for life, or for a limited period. This horrible fate was never inflicted on the better classes. They were punished by deportation to an island, or an oasis in the desert.\footnote{Penal servitude and deportation, when they involved disfranchisement, were classified as \emph{capital} punishments. All capital punishments involved confiscation of property.

}

p415
Mutilation does not appear to have been recognised as a legal penalty under the Principate, but it may sometimes have been resorted to as an extraordinary measure by the express sentence of an Emperor. It first appears in an enactment of Constantine ordaining that the tongue of an informer be torn out by the root.\texthebrew{⁠}\footnote{\emph{C. Th.} X.10.2 ; it is not quite clear whether this was to be done before or after death (strangulation on the gallows). A later law reduced the penalty to death by the sword, \emph{ib.} 10.

} Leo I condemned persons who were implicated in the murder of Proterius, patriarch of Alexandria, to excision of the tongue and deportation.\texthebrew{⁠}\footnote{Theophanes, A.M. 5951.

} In the sixth century, mutilation became more common, and Justinian recognises amputation of the hands as a legal punishment in some of his later enactments. Tax-collectors who falsify their accounts and persons who copy the writings of Monophysites are threatened with this pain.\texthebrew{⁠}\footnote{\emph{Nov.} 17, § 8 ; 42, § 1. Cp. \emph{Nov.} 134, § 13, where it is forbidden to punish theft by cutting off \textgreek{ο\texthebrew{ἱ}ονδήποτε μέλος.}} And we have records of the infliction of a like punishment on other criminals.\texthebrew{⁠}\footnote{John Mal. XVIII pp451 (gambling, A.D. 429); but no physical penalty is enacted in the law of this year against gambling in

\emph{C. J.} III.43.2 ), 483, 488.

} This practice seems to have been prompted by the rather childish idea that, if the member which sinned suffered, the punishment was fitly adjusted to the crime.\texthebrew{⁠}\footnote{Cp. \emph{C. Th.} X.10.2 ; and Zonaras, XIV.7.2, where Justinian is reported to have justified his punishment of paederasty on this principle. The same idea may have dictated the punishment of incendiaries by fire.

} Amputation of the nose or tongue was frequently practised, and such penalties afterwards became a leading feature in Byzantine criminal law, and were often inflicted as a mitigation of the death penalty. When these punishments and that of blinding are pointed to as one of the barbarous and repulsive characters of Byzantine civilisation, it should not be forgotten that in the seventeenth century it was still the practice in England to lop off hands and ears.

It must be remembered that a considerable latitude was allowed to the judges (praetors, prefects, provincial governors) in passing sentences on culprits.\texthebrew{⁠}\footnote{On this subject see Mommsen, \emph{op. cit.} 1037 \emph{sqq.}} The penalties prescribed in the laws were rather directions for their guidance than hard and fast sanctions. They were expected to take into account circumstances which aggravated the guilt, and still more circumstances which extenuated it. For instance, youth, intoxication, an º ethical motive were considered good reasons for mitigating

p416
penalties, and women were generally treated more leniently than men.

On the whole, the Roman system, from Augustus to Justinian, of protecting society against evil-doers and correcting the delinquencies of frail humanity, can hardly arouse much admiration. It was, indeed, more reasonable and humane than the criminal law of England before its reform in the nineteenth century. Its barbaric features were due either indirectly to the institution of slavery, or to the influence of the Church in those domains which especially engaged the interest of ecclesiastics. Augustus and his successors definitely stemmed the current of tendency which in the last period of the Republic promised entirely to do away with capital punishment, but they did not introduce any new reasonable principle into the theory or practice of criminal law. Wider extension of the field of public crimes, increasing severity in the penalties, and differential treatment of citizens of the lower classes, are the most conspicuous features of the development of criminal justice under the Empire.

\xchapter{Chapter XXIV}

Short URL for this page:

bit.ly/BURLAT24

Throughout the fifth century there were Greek historians writing the history of their own times, and, if their writings had survived, we should possess a fairly full record of events, particularly in the East, from the accession of Arcadius to the reign of Zeno. And it would have been a consecutive record, or at least there would have been only one or two short gaps. The bitter pagan sophist Eunapius of Sardis carried down his history, composed in his old age, to A.D. 404.\texthebrew{⁠}\footnote{Eunapius designed his history as a continuation of that of Dexippus which ended at 270. The evidence does not point to any continuity between him and Olympiodorus, or between Olympiodorus and Priscus.

} Olympiodorus, a native of Egyptian Thebes, began his book at A.D. 407 and went down to A.D. 425.\texthebrew{⁠}\footnote{Olympiodorus was a traveller and a poet. He was sent on an embassy in 412 to Donatus, a Hunnic prince of whom otherwise we know nothing.

} The work of Priscus of Panion (near Heraclea on the Propontis) probably began about A.D. 434 and ended with the death of Leo I. Malchus, of the Syrian Philadelphia, continued Priscus, and embraced in his work either the whole or a part of the reign of Zeno.\texthebrew{⁠}\footnote{From the evidence of the excerpts and the notice of Photius,

\emph{Bibl.} 78 , we should conclude that he stopped at 480, but Suidas \emph{sub nomine} says that he went down to the reign of Anastasius. He certainly wrote after the death of Zeno, to whom he was hostile.

} But all these histories have perished. Some of the information they contained passed into later writers; for instance,

Zosimus , who wrote towards the end of the fifth century, derived much of his material for the later portion of his work from Eunapius and Olympiodorus.\footnote{After a brief survey of the earlier Empire, Zosimus began his fuller narrative about A.D. 270. He is the only important historian who wrote non-contemporary history in the fifth and sixth centuries, except Peter the Patrician \texthebrew{—} the diplomatist and Master of Offices \texthebrew{—} who composed a History of the Roman Empire from

(p418)
Augustus to Julian. Of this we have fragments, some of which \texthebrew{—} known as the anonymous continuation of Dio Cassius \texthebrew{—} have only recently been connected with Peter's work, by C. de Boor, \emph{Römische Kaisersgeschichte}, \emph{B. Z.} I.13 \emph{sqq.}, 1892.

}

p418
But of the original text we possess only excerpts which in many case are mere summaries. The ecclesiastical histories written by the orthodox laymen, Socrates and Sozomen, about the middle of the century, have fared better; we have them intact and fortunately they include notices of secular events rather capriciously selected.\footnote{For secular history Theodoret's \emph{Ecclesiastical History} is almost negligible; it has hardly anything that is not in Socrates or Sozomen. The work of Philostorgius, the Eunomian, is a real loss, as

the fragments

show; the works of heretics had little chance of surviving.

}

The fragments of these lost historians enable us to judge that Priscus is a greater loss than any of the rest. The long fragment on Attila and his court, of which a translation was given

in an earlier chapter , shows that he was a master of narrative, and the general impression we get is that he was the ablest Roman historian between Ammian and Procopius.

Why did all these works disappear? Some of them survived till the ninth and tenth centuries, but were doubtless extremely rare then, and no more copies were made from the one or two copies that existed. The probability is that they never had a wide circulation, and it is fair to ascribe this partly to the fact that their authors were pagans.\texthebrew{⁠}\footnote{This is not so clear in the case of Priscus, but his friendship with the pagan Maximin establishes a presumption. The sympathies of Malchus were plainly Neoplatonic, as is shown by his treatment of Pamprepius, and his designation of Proclus as ``the great Proclus.'' \texthebrew{—} It is curious that the aggressively pagan work of Zosimus survived.

} But there is another reason which may partly account for the loss of some of these historians, and may also explain the character of the excerpts which have come down. In the ninth and following centuries the Greeks were interested in the past history of the Illyrian peninsula and in the oriental wars with the Persians, which were fought on the same ground as the contemporary wars with the Moslems; they were not interested in the history of Italy and the West. Now in the fifth century, with the exception of one or two short and unimportant episodes of hostility, there was hardly anything to tell of the oriental frontier, so that the portions of historians like Priscus and Malchus that had a living interest for readers were those which dealt with the invaders and devastators of the Balkan provinces. And so we find that the most considerable

p419
fragments of Priscus, preserved in the summaries and selections that were made in the ninth and tenth centuries, relate to the doings of the Huns, and the most considerable fragments of Malchus to the doings of the Ostrogoths. It is, in fact, probable that these extracts represent pretty fully the information on these topics given by both writers. On the other hand, it is significant that the Gallic and Italian campaigns, which Priscus must certainly have described, were passed over by those who made the selections.\footnote{Perhaps the clearest illustration of the point is that, if the works of Procopius had been lost, we should know a great deal about his Persian Wars, but very little about his Vandalic and Gothic Wars; for Photius,

in his \emph{Bibliotheca} , gives a long account of the former, but none of the latter.

}

That there is almost as much to tell about thirty years of the sixth century, as there is about the whole of the fifth, is due partly to Justinian's activity as a legislator, but chiefly to the pen of Procopius. It was one of the glories of Justinian's age to have produced a writer who must be accounted the most excellent Greek historian since Polybius. Procopius was a native of Caesarea, the metropolis of the First Palestine. He was trained to be a jurist, and we have seen how he was appointed in A.D. 527 councillor to Belisarius,\texthebrew{⁠}\footnote{Haury argues that he was not a jurist (\emph{Zur Beurteilung}, etc. p20), but he is not convincing. For the post of consiliarius (\textgreek{ξύμβουλος}) cp. \emph{C. J.} I.51.11 . Cp. Dahn, \emph{Prokopius}, pp19\texthebrew{‑}20.

} how he accompanied him in his Persian, African, and Italian campaigns, and how he was in Constantinople when the city was ravaged by the plague. He was not with Belisarius in his later campaigns in the East, and it is improbable that he revisited Italy.\footnote{It is indeed possible that Procopius was with Belisarius in 541, though it seems more probable that his place had been taken by George (\emph{B. P.} II.19.22; cp. Haury, \emph{ib.}). Dahn leaves it open (\emph{ib.} p30) whether he was in Italy after 542. Haury has argued that he went to Italy in 546, because he relates the events of 546\texthebrew{‑}547 at greater length than those of the other years of the Second Italian War (\emph{Procopiana}, I.8\texthebrew{‑}9). But this is not very convincing. It appears probable to me that he lived at Constantinople continuously after his return from Italy, and there collected from officers, ambassadors, etc., the material which he required for his \emph{History}, both in the East and in the West from 541 to 553. Haury thinks that he wrote his \emph{History} in Caesarea (\emph{ib.} 26), but it would have been difficult for him there to have obtained regularly first-hand and detailed information about the war in Italy.

}

His writings attest that Procopius had received an excellent literary education. There is nothing which would lead us to suppose that he had studied either at Athens or at Alexandria, and it seems most probable that he owed his attainments to the

p420
professors of the university of Constantinople. It has indeed been held that he was educated at Gaza,\texthebrew{⁠}\footnote{This theory of Haury (\emph{op. cit.}), is closely connected with a theory as to his parentage. Haury argues that he was the son of Stephanus, a leading citizen of Caesarea, who, before A.D. 526, held the post of astynomos or commissioner of public works, won distinction by restoring the aqueduct which supplied the city, and in 536 was appointed proconsul of the First Palestine (Choricius, \emph{Epithal.} p22; \emph{In Arat. et Stephanum}, § 10; Justinian,

\emph{Nov.} 103, § 1 ; Aeneas Gaz. \emph{Ep.} 11). Stephanus was a friend of Procopius of Gaza, and sent his son to be educated there (Procopius, \emph{Ep.} 18). This unnamed son Haury identifies with a Procopius who married a young woman of Ascalon, wealthy and of good family; two of his fellow-students married at the same time; and Choricius wrote an (extant) epithalamion for the occasion. In the Samaritan revolt of A.D. 556, Stephanus was murdered by the rebels in his court house. His wife went to Constantinople and besought Justinian to punish the murderers. The Emperor did her justice promptly, ordering Amantius, Master of Soldiers in the East, to search them out, and they were executed (John Mal. \emph{fr.} 48, \emph{De ins.}). In this incident Haury finds the explanation of the revolution in the attitude of Procopius towards Justinian between 550 and 560. The theory is ingenious, but much more evidence would be needed to make it probable. It is difficult to believe that, if Procopius had been the son of such a notable civil servant as Stephanus, the fact would not have been recorded. The known facts do not point to any connexion with Gaza. And it is to be observed that the theory involves \emph{two} conjectural identifications: that of the historian with the Procopius who married the girl of Ascalon, and that of this Procopius with the son of Stephanus.

} but this theory rests on no convincing external evidence, and the internal evidence of his style does not bear out the hypothesis that he ever sat at the feet of his namesake Procopius and the other sophists of Gaza. We know a good deal about that euphuistic literary school.

It may be conjectured that Procopius formed the design of writing a history of the wars, of which Belisarius was the hero, at the time of the expedition against the Vandals, and that he commenced then to keep a written record of events.\texthebrew{⁠}\footnote{The account of the First Persian War (\emph{B. P.} I.12\texthebrew{‑}22) is so (comparatively) brief and incomplete that we may perhaps infer that he had not formed the plan of writing a history of it during its progress. The chronology of the composition of his works has been cleared up by the researches of Haury (\emph{Procopiana}, I). In the latter part of 545 he was writing \emph{B. P.} I.25.43; \emph{B. G.} I.24.32 (cp. \emph{Procopiana}, II p5); and II.5.26\texthebrew{‑}27. These passages indicate that he contemplated this year as the termination of his history, and if he had then published it, it would have consisted of \emph{B. P.} I\texthebrew{‑}II.28, 11; \emph{B. V.}, except the last two pages; \emph{B. G.} I\texthebrew{‑}III.15. He may have circulated what he had written among his acquaintances, but he continued to add to it from year to year, without changing what he had already written, and finally published the seven Books of the Wars, as they stand, in 550. In the same year he wrote the \emph{Secret History.} In 553 he published the eighth and last Book of the Wars. In 560 he wrote the \emph{De aedificiis}, as is proved by the mention of the building of the bridge over the Sangarius

(V.3.10)

which was built in 559\texthebrew{‑}560 (Theoph. A.M. 6052). On this bridge cp. Anderson, \emph{J. H S.} XIX.66 \emph{sq.}

} He had

p421
certainly begun to compose his history immediately after his return from Italy with Belisarius, for he states that, as that general's councillor, he had personal knowledge of almost all the events which he is about to relate, a statement which would not be true of the later campaigns. While he is studiously careful to suppress feeling, he gives the impression, in his narrative of the early wars, that he sympathised with Justinian's military enterprises, and viewed with satisfaction the exploits of Belisarius. As to Justinian himself he is reticent; he may sometimes imply blame; he never awards praise. The sequel disillusioned him. He was disappointed by the inglorious struggles with Chosroes and Totila, and by the tedious troubles in Africa; and his attitude of critical approbation changed into one of bitter hostility towards the government. His vision of his own age as a period of unexpected glory for the Empire faded away, and was replaced by a nightmare, in which Justinian's reign appeared to him as an era of universal ruin. Seeking to explain the defeat of the early prospects of the reign, he found the causes in the general system of government and the personalities of the rulers. The defects of the Imperial administration, especially in the domain of finance, were indeed so grave, that it would have been easy to frame a formidable indictment without transcending the truth, or setting down aught in malice; but with Procopius the abuses and injustices which came to his notice worked like madness on his brain, and regarding the Emperor as the common enemy of mankind he was ready to impute the worst of motives to all his acts.

We may divine that the historian went through a mental process of this nature between A.D. 540 and 550, but we cannot believe that pure concern for the public interests is sufficient to explain the singular and almost grotesque malignity of the impeachment of Justinian and all his works, which he drew up at the end of the decade. Any writer who indulges in such an orgy of hatred as that which amazes us in the \emph{Secret History}, exposes himself to the fair suspicion that he has personal reasons for spite. We hardly run much risk of doing an injustice to Procopius if we assume that he was a disappointed man. One who had occupied a position of intimate trust by the side of the conqueror of Africa and Italy could not fail to entertain hopes of preferment to some administrative post. But he was

p422
passed over. The influence of Belisarius, if it was exerted in his favour, did not avail, and from being a friendly admirer of his old patron he became a merciless critic.

In a book which was intended to be published and to establish a literary reputation, Procopius could not venture to say openly what was in his mind. But to an attentive reader of his narrative of the later wars there are many indications that he disapproved of the Imperial policy and the general conduct of affairs. His \emph{History of the Wars} was divided into seven books, and the material, on the model of Appian, was arranged geographically. Two books on the Persian War brought the story down to A.D. 548, two on the Vandalic War embraced events in Africa subsequent to the conquest and reached the same date. The three books of the Gothic War terminated in A.D. 551. It was probably the final defeat of Totila in A.D. 552 that moved the author afterwards to complete his work by adding an eighth book, in which, abandoning the geographical arrangement, he not only concludes the story of the Italian War, but deals with military operations on every front from A.D. 548 to 553.

We shall presently see that in the later parts of this work the historian went as far as prudence permitted in condemning the policy of Justinian. But in A.D. 550 he secretly committed to writing a sweeping indictment of the Emperor and the late Empress, of their private lives and their public actions. It was a document which he must have preserved in his most secret hiding-place , and which he could read only to the most faithful and discreet of his friends. It could never see the light till Justinian was safely dead, and if he were succeeded by a nephew or cousin, its publication even then might be impossible. As a matter of fact we may suspect that his heir withheld it from circulation, and that it was not published till a considerable time had elapsed. For it was unknown to the writers of the next generation, unless we suppose that they deliberately ignored it.\footnote{Some have supposed that Evagrius was acquainted with it, but the evidence is quite unconvincing. The earliest reference to the work is in Suidas (tenth century), \emph{s.v.} \textgreek{Προκόπιος}. He calls it the Anecdota.

}

The introduction to the \emph{Secret History}\texthebrew{⁠}\footnote{It used to be supposed that the date of composition was 559, because it was written in the thirty-second year of Justinian's régime (\emph{H. A.} 18.33 ;

23.1 ;

24.29 ). But it is a fundamental part of the thesis of Procopius that Justinian was the real governor (\textgreek{δι\texthebrew{ῳ}κήσατο τ\texthebrew{ὴ}ν πολιτείαν})

(p423)
throughout the reign of Justin, and the thirty-second year is to be reckoned from 518, not from 527. This was first established by Haury (Procop. I) . On the old view it would be impossible to explain why Procopius should have ignored all the material which the events between 550 and 559 supplied for his purpose. Cp. also Haury, \emph{Zur Beurteilung}, p37. The revised date of the \emph{Secret History} shows that the first lines of the last Book of the Wars (\emph{B. G.} IV) were modelled on the opening lines of \emph{H. A.} and not \emph{vice versa.} Haury has suggested that the author originally began \emph{B. G.} IV. with \textgreek{\texthebrew{ἤ}δη μ\texthebrew{ὲ}ν ο\texthebrew{ὗ}ν} (c. i.3), and afterwards added the preceding sentences for the purpose of misleading posterity, and suggesting that \emph{H. A.} was the work of an imitator. Conversely it might be argued that the motive was to indicate identity of author­ship.

} states that its object

p423
is to supplement the \emph{History of the Wars} by an account of things that happened in all parts of the Empire, and to explain certain occurrences which in that work had been barely recorded, as it was impossible to reveal the intrigues which lay behind them. It is hinted that so long as Theodora was alive,\texthebrew{⁠}\footnote{1.2 \textgreek{περιόντων \texthebrew{ἔ}τι τ\texthebrew{ῶ}ν α\texthebrew{ὐ}τ\texthebrew{ὰ} ε\texthebrew{ἰ}ργασμένων} must mean simply Theodora (cp. 16.3).

} it would have been dangerous even to commit the truth to writing, for her spies were ubiquitous, and discovery would have meant a miserable death. This reinforces other evidence which goes to prove that Theodora was held in much greater fear than Justinian.

The thesis of the \emph{Secret History}\texthebrew{⁠}\footnote{The book is badly arranged. A preliminary section (1\texthebrew{‑}5) is devoted to Belisarius and the scandals connected with his private life; his piti­ful uxoriousness, his weakness, which leads him into breach of faith, and his military failures. The next subject is the family and character of Justinian (6\texthebrew{‑}8), and then the author goes on to tell the scandalous story of Theodora's early life (9, 10). He proceeds to characterise the revolutionary policy of the Emperor, and to give a summary account of his persecutions, his avarice, and his unjust judgments (11\texthebrew{‑}14). Then he reverts again to Theodora, and illustrates her power, her crimes, and her cruelties (15\texthebrew{‑}17). He goes on to review the calamities and loss of life brought not only upon his own subjects, but also on the barbarians by Justinian's wars, as well as by the pestilence, earthquakes, and inundations, for which he holds him responsible (18); and then enters upon a merciless criticism of his financial administration (19\texthebrew{‑}23). Various classes of society \texthebrew{—} the army, the merchants, the professions, the proletariate \texthebrew{—} are then passed in review, and it is shown that they are all grievously oppressed (24\texthebrew{‑}26). The last chapters are occupied with miscellaneous instances of cruelty and injustice (27\texthebrew{‑}29), the decline of the cursus publicus, and new servile customs in court etiquette (30).

} is that in all the acts of his public policy Justinian was actuated by two motives, rapacity and an inhuman delight in evil-doing and destruction. In this policy he was aided by Theodora, and if they appeared in certain matters, such as religion, to pursue different ends, this was merely a plot designed to hoodwink the public.\texthebrew{⁠}\footnote{10.14.

} Procopius gravely asserts that he himself and ``most of us'' had come to the conclusion that the Emperor and Empress were demons

p424
in human form, and he did not mean this as a figure of speech.\texthebrew{⁠}\footnote{12.14 \emph{sqq.} ;

18.1 . Theodora's influence over her husband is ascribed to magic practices, 22.27 .

} He tells a number of anecdotes to substantiate the idea. Justinian's mother had once said that she conceived of a demon. He had been seen in the palace at night walking about without a head, and a clairvoyant monk had once refused to enter the presence chamber because he saw the chief of the demons sitting on the throne. Before her marriage, Theodora had dreamt that she would cohabit with the prince of the devils. Even Justinian's abstemious diet is adduced as a proof of his non-human nature. It was a theory which did not sound so ludicrous in the age of Procopius as in ours, and it enabled him to enlarge the field of the Emperor's mischievous work, by imputing to his direct agency the natural calamities like earthquakes and plagues which afflicted mankind during his reign.

In elaborating his indictment Procopius adopted two sophistic tricks. One of these was to represent Justinian as responsible for institutions and administrative methods which he had inherited from his predecessors. The other was to seize upon incidental hardships and abuses arising out of Imperial measures, and to suggest that these were the objects at which the Emperor had deliberately aimed. The unfairness of the particular criticisms can in many cases be proved, and in others reasonably suspected. But it may be asked whether the book deserves any serious consideration as an historical document, except so far as it illustrates the intense dissatisfaction prevailing in some circles against the government. The daemonic theory, the pornographic story of Theodora's early career, the self-defeating maliciousness of the whole performance discredit the work, and have even suggested doubts whether it could have been written at all by the sober and responsible historian of the wars.\texthebrew{⁠}\footnote{L. von Ranke (\emph{Weltgeschichte}, IV.2.300 \emph{sqq.}) argued that it was not written by Procopius, but was partly based on a Procopian diary. The Procopian style of the \emph{Secret History} is unmistakable and has been well illustrated by Dahn, 257 \emph{sqq.}, 416 \emph{sqq.}

} The author ­ship, however, is indisputable. No imitator could have achieved the Procopian style of the \emph{Secret History}, and a comparison with the \emph{History of the Wars} shows that in that work after A.D. 541 the author makes or suggests criticisms which are found, in a more explicit and lurid form, in the libel.

For in the public \emph{History} he sometimes used the device of

p425
putting criticism into the mouths of foreigners. One of the prominent points in the \emph{Secret History} is Justinian's love of innovations; he upset established order, and broke with the traditions of the past. The same character is given him in the public \emph{History} by the Gothic ambassadors who went to the Persian Court.\texthebrew{⁠}\footnote{\emph{B. P.} II.2.6 \textgreek{νεωτεροποιός τε \texthebrew{ὢ}ν φύσει . . . μένειν τε ο\texthebrew{ὐ} δυνάμενος \texthebrew{ἐ}ν το\texthebrew{ῖ}ς καθεστ\texthebrew{ῶ}σι}, cp. \emph{H. A.} 11.1\texthebrew{‑}2 . See also the remark in a speech of Armenian envoys to Chosroes, \emph{B. P.} II.3.38 \textgreek{τί ο\texthebrew{ὐ}κ \texthebrew{ἐ}κίνησε τ\texthebrew{ῶ}ν καθεστώτων;}} The motive of the speech which is attributed to the Utigur envoys in A.D. 552 is to censure the policy of giving large grants of money to the trans-Danubian barbarians, which is bitterly assailed in the \emph{Secret History}.\texthebrew{⁠}\footnote{\emph{B. G.} IV.19.9 \emph{sqq.};

\emph{H. A.} 11.5-7 .

} Procopius indeed criticises it directly by an irony which is hardly veiled. The Kotrigurs, he says, ``receive many gifts every year from the Emperor, and, even so, crossing the Danube they overrun the Emperor's territory continually.''\texthebrew{⁠}\footnote{\emph{B. G.} IV.5.16.

} Although Justinian was here only pursuing, though perhaps on a larger scale, the inveterate practice of Roman policy, his critic speaks as if it were a new method which he had discovered for exhausting the resources of the Empire. ``It is a subject of discussion,'' he says, ``what has happened to the wealth of the Romans. Some assert that it has all passed into the hands of the barbarians, others think that the Emperor retains it locked up in many treasure chambers. When Justinian dies, supposing him to be human, or when he renounces his incarnate existence if he is the lord of the demons, survivors will learn the truth.''\footnote{\emph{H. A.} 30.33\texthebrew{‑}34 , the last words of the book.

}

In the \emph{Secret History} the Emperor is arraigned as the guilty party in causing the outbreak of the second Persian War. In the published \emph{History} the author could not say so, but goes as far as he dares by refusing to say a word in his favour. Having stated the charges made by Chosroes that Justinian had violated the treaty of A.D. 532, he adds, ``Whether he was telling the truth, I cannot say.''\texthebrew{⁠}\footnote{\emph{B. P.} II.1.15;

\emph{H. A.} 11.12 .

} On the peace of A.D. 551, which evidently excited his indignation, he resorts to the same formula. ``Most of the Romans were annoyed at this treaty, not unnaturally. But whether their criticism was just or unreasonable I cannot say.''\texthebrew{⁠}\footnote{\emph{B. G.} IV.15.13 (the author's opinion is clearly suggested \emph{ib.} 7). His dissatisfaction with the particular favour shown by Justinian to the Persian envoy Isdigunas is not disguised, \emph{ib.} 19, 20; \emph{B. P.} II.28.40\texthebrew{‑}44.

} Nor did the historian of the Vandalic War fail to

p426
suggest the same conclusion which is drawn in the \emph{Secret History} as to the consequences of the Imperial conquest. Having recorded the victory of John, the brother of Pappus, over the Moors in A.D. 548, he terminates the story with the remark, ``Thus, at last and hardly, to the survivors of the Libyans, few and very destitute, there came a period of peace.''\footnote{\emph{B. V.} II.28.52. Other points of comparison between the public and the \emph{Secret History} may be noted. (1) The view of Justin's reign as virtually part of Justinian's, \emph{B. V.} I.9.5 (Justin is \textgreek{\texthebrew{ὑ}περγήρως}):

\emph{H. A.} 6.11

(Justin is \textgreek{τυμβογέρων}). (2) Pessimistic utterances as to the general situation, A.D. 541\texthebrew{‑}549; \emph{B. P.} II.21.34; \emph{B. G.} III.33.1. (3) Justinian's slackness in prosecuting his wars is ascribed, in

\emph{H. A.} 18.29 , partly to his avarice and partly to his occupation with theological studies. The latter reason is plainly assigned in

\emph{B. G.} III.35.11

\textgreek{βασιλε\texthebrew{ὺ}ς δ\texthebrew{ὲ Ἰ}ταλίας μ\texthebrew{ὲ}ν \texthebrew{ἐ}πηγγέλλετο προνοήσειν α\texthebrew{ὐ}τός, \texthebrew{ἀ}μφ\texthebrew{ὶ} δ\texthebrew{ὲ} τ\texthebrew{ὰ} Χριστιαν\texthebrew{ῶ}ν δόγματα \texthebrew{ἐ}κ το\texthebrew{ῦ ἐ}π\texthebrew{ὶ} πλε\texthebrew{ὶ}στον διατριβ\texthebrew{ὴ}ν ε\texthebrew{ἶ}χεν}, and is perhaps hinted at

\emph{ib.} 36.6

\textgreek{\texthebrew{ἀ}σχολίας ο\texthebrew{ἱ ἴ}σως \texthebrew{ἐ}πιγενομένης \texthebrew{ἑ}τέρας τιν\texthebrew{ὸ}ς τ\texthebrew{ὴ}ν προθυμίαν κατέπαυσε}. (The same ironical formula is employed when the Emperor failed to send money due to Gubazes at the right time, \textgreek{\texthebrew{ἐ}πιγενομένης ο\texthebrew{ἱ ἀ}σχολίας τινός,}\emph{B. P.} II.29.32 ). Cp. also

\emph{B. G.} III.32.9 . (4) The financial oppression of the government is exposed in the case of the logothete Alexander

(\emph{ib.} III.1.29) . More pointed is the remark that John Tzibos was created a general because he was the worst of men, and understood the art of raising money,

\emph{B. P.} II.15.10 . The criticism of Justinian that, though well aware of the unpopularity of Sergius in Africa, \textgreek{ο\texthebrew{ὐ}δ\texthebrew{’ ὥ}ς} would he recall him, may be noted. (5) In holding up to reprobation the conduct of the second Italian War by Belisarius, the author of the libel has only to repeat

(\emph{H. A.} 5.1)

what was said in

\emph{B. G.} III.35.1 . Compare also the remarks,

\emph{ib.} xiii.15\texthebrew{‑}19 .

}

In fact, the attitude of Procopius towards the government, as it is guardedly displayed in the \emph{History of the Wars}, is not inconsistent with the general drift of the \emph{Secret History}, and the only reason for doubting the genuineness of the libel was the presumption that the political views in the two works were irreconcilable. It is another question whether the statements of the \emph{Secret History} are credible. Here we must carefully distinguish between the facts which the author records, and the interpretation which he places upon them. Malice need not resort to invention. It can serve its purpose far more successfully by adhering to facts, misrepresenting motives, and suppressing circumstances which point to a different interpretation. That this was the method followed by Procopius is certain. For we find that in a large number of cases his facts are borne out by other contemporary sources,\texthebrew{⁠}\footnote{\emph{E.g.} (1) The unscrupulousness of Theodora is illustrated by the episode of the Nobadae, above,

p328 \emph{sqq.} ; (2) the statements as to the intrigue against Amalasuntha fit into the other evidence, above,

p165 \emph{sqq.} ; (3) the connexion of Antonina with the episode of Silverius, cp. p378 ; (4) John Eph. (\emph{Hist. ecc.} I c32) confirms the information that Antonina's son Photius was a monk; (5) the

(p427)
story of Callinicus,

\emph{H. A.} 17.2 , is confirmed by

Evagrius, IV.32 ; (6) that of Priscus,

\emph{H. A.} 16.7 , by John Mal. XVIII.449, and \emph{fr.} 46, \emph{De ins.}; (7) that of Theodotus by John Mal. XVIII.416; while (8) the story of Psoes,

\emph{H. A.} 27.14 , is consistent with ecclesiastical records. (9) The laws referred to in the \emph{H. A.} can be verified in the legal monuments of the reign; (10) the statements about religious persecutions accord with facts; and (11) in what is said,

\emph{H. A.} 17.5 \emph{sqq.} , of the attempts to repress prostitution, other evidence shows that the author is only representing facts in a light unfavourable to the policy.

} while in no instance can we

p427
convict him of a statement which has no basis in fact.\texthebrew{⁠}\footnote{There is indeed an exception in the statement,

\emph{H. A.} 22.38 , that the gold nomisma was reduced in value, which is not in accordance with the numismatic evidence. But even here it may be doubted whether Procopius had not some actual temporary or local fact in mind.

} We have seen that even in the case of Theodora's career, where his charges have been thought particularly open to suspicion, there is other evidence which suggests that she was not a model of virtue in her youth. The \emph{Secret History} therefore is a document of which the historian is entitled to avail himself, but he must remember that here the author has probably used, to a greater extent than elsewhere, material derived from gossip which he could not verify himself.

Procopius entertained the design of writing another book dealing especially with the ecclesiastical policy of the reign.\texthebrew{⁠}\footnote{This intention can be inferred from three passages in \emph{H. A.} ( 1.14 ;

11.33 ;

26.18 ), and from \emph{B. G.} IV.25.13. The fulfilment of the promise in

\emph{H. A.} 17.14 , to tell the sequel of the story of two young women who had been unhappily married, may possibly have been also reserved for this book, though there is no hint that it had anything to do with ecclesiastical affairs.

} If the work was ever executed it was lost, but as there is no reference to it in subsequent literature, it seems most probable that it was never written. Among other things which the historian promised to relate in it was the fate of Pope Silverius, concerning which our extant records leave us in doubt as to the respective responsibilities of Vigilius and Antonina. Apart from the facts which it would have preserved to posterity, the book would have been of singular interest on account of the Laodicean attitude of the author, who, whatever may have been his general opinion of Christian revelation, was a Gallio in regard to the theological questions which agitated the Church. ``I am acquainted with these controversial questions,'' he says somewhere, referring to the Monophysite disputes, ``but I will not go into them. For I consider it a sort of insane folly to investigate the nature of God. Man cannot accurately apprehend the constitution of man, how much less that of the Deity.''\texthebrew{⁠}\footnote{\emph{B. G.} I.3.6.

} The words imply an

p428
oblique hit at the Emperor who in the \emph{Secret History} is described as gratuitously busy about the nature of God.\texthebrew{⁠}\footnote{\emph{H. A.} 18.29 .

} That the book would also have been a document of some significance in the literature of toleration we may infer from a general remark which Procopius makes on Justinian's ecclesiastical place. ``Anxious to unite all men in the same opinion about Christ, he destroyed dissidents indiscriminately, and that under the pretext of piety; for he did not think that the slaying of men was murder unless they happened to share his own religious opinions.''\footnote{\emph{Ib.} 13.7 .

}

An amazing change came to pass in the attitude of Procopius between the year in which he composed the \emph{Secret History} and ten years later when he wrote his work on the \emph{Buildings}, in which he bestows on the policy and acts of the Emperor superlative praise which would astonish us as coming from the author of the \emph{History of the Wars}, even if the \emph{Secret History} had been lost or never written. The victories of Narses had probably mitigated the pessimism into which he had fallen through the failure of Belisarius and the long series of Totila's successes; but it is difficult to avoid the conjecture that he had received some preferment or recognition from the Emperor.\texthebrew{⁠}\footnote{There is some evidence (see Suidas \emph{sub} \textgreek{Προκόπιος}) that he attained the rank of Illustrious, and there is a possibility that he was the Procopius who was Prefect of the City in 562 (Theophanes, A.M. 6055). But the name was a common one.

} In the opening paragraph of the \emph{Buildings} there is a hint at private motives of gratitude. ``Subjects who have been well treated feel goodwill towards their benefactors, and may express thanks by immortalising their virtues.'' The author goes on to review and appreciate briefly Justinian's achievements in augmenting the size and prestige of the Empire, in imposing theological unity on its inhabitants, in ordering and classifying its laws, in strengthening its defences, and, noting particularly his indulgent treatment of conspirators, praises his general beneficence. This was a wonder ­ful recantation of the unpublished libel, and we may doubt whether it was entirely sincere. Procopius did not take the \emph{Secret History} out of its hiding-place and burn it, but he abstained from writing the book on ecclesiastical history which he had planned.

Wherever he was educated, Procopius had been saturated

p429
with Herodotus and Thucydides. His works are full of phrases which come from their works, and his descriptions of military operations sometimes appear to be modelled on passages in Thucydides. This fact has in modern days suggested the suspicion that some of his accounts of battles or sieges are the literary exercise of an imitator bearing little relation to what actually occurred.\texthebrew{⁠}\footnote{See the tracts of Braun and Brückner, mentioned in

Bibliogr. II.2, B , and the refutation of their suggestions by Haury, \emph{Zur Beurt.} 1 \emph{sqq.} \texthebrew{—} The fatalistic remarks which Procopius introduces from time to time are Herodotean.

} But when we find that in some cases, which we can control, other sources bear out his accounts of operations at which he was not present (for instance, of the siege of Amida in the reign of Anastasius), we see that he did not misconceive the duty of a historian to record facts, and was able through his familiarity with Thucydides and Herodotus to choose phrases from their writings suitable to a particular case. It is remarkable that he does not seem to have read the \emph{History} of Priscus, for, where he relates events of the fifth century, he seems to have derived his information not directly from that historian, but from intermediate writers who had used Priscus and perhaps distorted his statements.\texthebrew{⁠}\footnote{One of the intermediaries may have been Eustathius of Epiphania, who (see

Evagrius, V.24 ) wrote a universal history, for the latter part of which his sources were Zosimus and Priscus, whom he abbreviated. It ended with the year 503. He was one of the principal sources of Evagrius. See further Haury's preface to his ed. of Procopius.

} He appears to have known the Syriac tongue, and it has been suggested that this knowledge recommended him to Belisarius when he selected him as his assessor in his first Persian campaigns.\footnote{Haury, \emph{Zur Beurt.} p20. Haury thinks that he made use of Zacharias of Mytilene, and perhaps of the

History of Armenia by Faustus of Byzantium

in a Syriac version, \emph{ib.} pp4 and 21.

}

For his own time he derived information as to events and transactions, with which he was not in contact himself by virtue of his office on the staff of Belisarius, from people who had personal knowledge of them. It is probable that Peter, the Master of Offices, and possible that John, the nephew of Vitalian, were among his informants on Italian affairs.\texthebrew{⁠}\footnote{See above,

Chap. XVIII § 2 .

} And he seems to have lost no opportunity of making the acquaintance of ambassadors who came from foreign courts of Constantinople, and questioning them about the history of their countries.\footnote{See above,

Chap. XIX § 8 .

}

He wrote in the literary Greek which had developed in a

p430
direct line from the classical writers of antiquity, and had hardly been affected by the ordinary spoken language, from which it was far removed. His prose is straightforward and unadorned; his only affectation is that he liked to imitate Thucydides. For it would be unfair to describe as an affectation the avoidance of current terms of his own day,\texthebrew{⁠}\footnote{For instance he regularly describes a bishop as \textgreek{\texthebrew{ἱ}ερεύς}, not \textgreek{\texthebrew{ἐ}πίσκοπος.}} especially when they were of Latin origin, or the introduction of them with an explanation which is almost an apology.\texthebrew{⁠}\footnote{\textgreek{Φοιδεράτοι} is one of the few words of this class that he employs \emph{sans phrase.}

} For that was common form with all authors who aimed at writing dignified prose. His ``so\texthebrew{‑}called'' is simply equivalent to our inverted commas. But he did not conform to the technical rules which governed the prose of the more pretentious stylists of his time. He did not contort his sentences in order to avoid hiatus, and he ignored the rule which had recently been coming into fashion as to the fall of the accents in the last words of a clause. This rule was that the last accented syllable in a clause must be preceded by at least two unaccented syllables.\texthebrew{⁠}\footnote{Attention was drawn to it by W. Meyer (\emph{Der accentuirte Satzschluss}), but his conclusions have been considerably modified by C. Litzica (\emph{Das Meyersche Satzschlussgesetz}), who pointed out that, in all Greek prose writers, the majority of sentences conform to the rule. This is due to the nature of the language. His conclusion is that unless the exceptions do not exceed 10 or 11 per cent the writer was unconscious of the rule.

} Thus a sentence ending with the words pánton anthrópōn would be right, but one ending with anthrópōn pánton would be wrong. This rule is observed by Zosimus, and was strictly adopted by the two chief sophists of the school of Gaza, Procopius and Choricius.\texthebrew{⁠}\footnote{The fact that the historian Procopius did not observe it would not prove that he was not educated at Gaza, for the contemporary sophist Aeneas of Gaza does not seem to have adopted it. But it is an argument against Haury's attempt to associate him, on grounds of style, with the Gaza school of prose.

} Some writers observed it in a modified form, allowing occasional exceptions.

The history of Procopius breaks off in A.D. 552, and Agathias of Myrina takes up the story.\texthebrew{⁠}\footnote{Born \emph{c.} 536, he died in 582. He probably intended to continue his history till 565. It was continued by his admirer, Menander, who was trained as a lawyer, but spent his early life as a man about town. The fragments of Menander, whose work terminated at 582, suggest that he was a better historian than Agathias. It is to be noticed that Agathias made some use of Persian chronicles, from which his friend Sergius, the official interpreter, made translations for him.

} Agathias is a much less interesting person. By profession he was a lawyer, and his ambition was to be a poet. He was inferior to Procopius as a historian, and

p431
modern readers will judge him inferior as a writer, though this would not have been the opinion of his contemporaries, to whose taste his affected style, with its abundance of metaphors and its preciosity, strongly appealed. His clauses carefully observed that accentual law which Procopius had wisely neglected.

Agathias occupies a place in the history of Greek poetry both for his own compositions and for the anthology which he compiled of short poems by contemporary writers, including some of his own and some of his friend Paul the Silentiary.\texthebrew{⁠}\footnote{It was arranged in seven books according to subject. See his Prooemium (\emph{Anth. Gr.} IV.3).

} This, like the earlier collections of Meleager and Philip, passed, perhaps almost entire, into the Anthology of Constantine Cephalas which has been preserved.\texthebrew{⁠}\footnote{\emph{Anthologia Palatina} (so called because the sole MS. is preserved in the Palatine library of Heidelberg). Agathias also wrote a volume of epyllia, love stories from Greek mythology, which he called \emph{Daphniaca.}

} His talent was considerable, and he was a master of metrical technique in the style which was then fashionable, and of which the best example from the age of Justinian is the poem of Paul on the church of St. Sophia. This technique had been elaborated in the previous century by Nonnus of Panopolis.\footnote{The Egyptians, it was said (see next note), are mad about poetry. And so in the fifth and sixth centuries we have a long procession of Egyptian poets: Palladas, Claudian,

Synesius , Cyrus, and Nonnus; then Christodorus of Coptus,

Colluthus , Tryphiodorus, Julian (writer of many epigrams in the \emph{Anth. Gr.}); and finally we have the horrible scribblings with which Dioscorus of Aphrodito plagued the life of a duke of the Thebaid (\emph{Pap. Cairo}, I.67055, 67097, etc.; II.67177\texthebrew{‑}67188).

}

The \emph{Dionysiaca} of Nonnus\texthebrew{⁠} a is the most interesting Greek poem that was written since the days of the great Alexandrines, Theocritus, Callimachus, and Apollonius. Published perhaps after the middle of the fifth century,\texthebrew{⁠}\footnote{The date of Nonnus is disputed. We know that he was posterior to Gregory of Nazianzus whom he echoes in \emph{Dionys.} I.310, and wrote before the reign of Anastasius, during which Christodorus and Colluthus, who were influenced by him, lived (Suidas, \emph{sub nn.}). Ludwich dates the \emph{Dionysiaca} about 390\texthebrew{‑}405 on the ground of a passage in Eunapius, \emph{Vit. Proaeresii}, p92, who says that the Egyptians \textgreek{\texthebrew{ἐ}π\texthebrew{ὶ} ποιητικ\texthebrew{ῇ} σφόδρα μαίνονται}. But why should this necessarily allude to the poem of Nonnus? When Eunapius was writing before 405, there were at least two distinguished contemporary Egyptian poets, Palladas and Claudian, as well as Synesius. In two places (XVI.321, XX.372) Nonnus has \textgreek{α\texthebrew{ἴ}θε πατήρ με δίδαξε}, identical with the first words of a short poem of his fellow-townsman Cyrus, who rose to be Praetorian Prefect (see above,

Vol. I p228 ). Ludwich holds that Cyrus took them from Nonnus; Friedländer, and those who place Nonnus in the second half of the fifth century, hold that Nonnus was the borrower.\texthebrew{—} Christodorus was a prolific poet. He wrote an epic on the Isaurian war of Anastasius and much else, but the only extant work is the \emph{Ecphrasis} of the statues in the baths of Zeuxippus, \emph{Anth. Gr.} Book II. The \emph{Capture (p432)
of Ilion} by Tryphiodorus and the \emph{Rape of Helen} by Colluthus are preserved. They exhibit Nonnian influence, but do not strictly observe his rules. John of Gaza, who wrote a description of a picture representing the cosmos in the winter-bath of Gaza (opened about A.D. 536), is a servile imitator of the master.\texthebrew{—} On the technique of this school of poetry see Ludwich, \emph{Beiträge zur Kritik des Nonnus}, 1873; Friedländer, in \emph{Hermes}, XLVII.43 \emph{sqq.}; Tiedke, \emph{ib.} XLIX.214 \emph{sqq.} and XII.445 \emph{sqq.}

} it arrested the attention

p432
of all young men who were addicted to writing verse, and for the next three or four generations poets imitated his manner, and observed, some more and some less, the technical rules which had made his heroic metre seem a new revelation. Of these rules the most important were that a spondee is never admitted in the fifth foot; that of the first four feet two at least must be dactyls, and when there are two spondees they should not be successive;\texthebrew{⁠}\footnote{This is a rule to which he allows very occasional exceptions.

} that hiatus is forbidden; and that elision is allowed only in the case of some particles and prepositions. If we add to these restrictions the fact that the caesura after the second syllable of the third foot predominates far more than in earlier poets, it is evident that the hexameters of Nonnus produce an entirely different poetical effect from those of Homer or of Apollonius. But Nonnus introduced another rule of a different kind which points to the direction in which Greek versification was to develop in later times. He strictly excludes proparoxytone words from the ends of his verses. This consideration of accent, which was a complete departure from classical tradition, was due doubtless to the influence of popular poetry; and may be set side by side with the consideration of accent which, as we saw, was affecting Greek prose. The truth is, that in this age the Greeks had ceased to feel instinctively the difference between long and short syllables, and only those whose ear was educated by classical studies could appreciate poems written in the old metres. All vowels had the same value, and the new Christian hymnography, which was at its best in the sixth and seventh centuries,\texthebrew{⁠}\footnote{For the controversy on the date of Romanus, whom the admirers of hymnography consider the greatest of the Greek hymn writers, see Krumbacher, \emph{G.B.L.}2 663 \emph{sqq.}; \emph{Studien zu Romanos}, 1898; \emph{Umarbeitungen bei Romanos}, 1879; \emph{Romanos und Kyriakos}, 1901; C. de Boor, \emph{Die Lebenszeit des Dichters R.}, in \emph{B. Z.} IX.633 \emph{sqq.}; P. van den Ven, \emph{Encore Romanos le mélode}, \emph{ib.} XII.153 \emph{sqq.} The last two studies seem to establish that he lived in the sixth century, not in the eighth.

} took no account of quantity, but was governed by the simple rules that corresponding verses should have the same number of syllables and should have the final accent on the same syllable.

p433
By these metrical innovations the character of the epic metre was changed and made a suitable instrument for a Dionysiac theme. In order to achieve a whirling breathless speed Nonnus bound it in fetters which excluded the variety of metrical effects that the unrestricted use of spondees had enabled the Homeric hexameter to compass.\texthebrew{⁠}\footnote{Of every five lines of Homer, probably four would have been rejected by Nonnus.

} His harmonious dactyls, with the procession of long compound words which is almost a necessary consequence of the predominance of this foot,\texthebrew{⁠}\footnote{Thus the average number of words in his verses is small. It has been calculated that it is 6 to 7·2 in Homer. The great majority of the uncommon compounds in Nonnus are not of his own coinage. They are to be found here and there in earlier poets, whom he carefully searched.

} however pleasing and effective in a short poem, become, in a long epic like the \emph{Dionysiaca} which has forty-eight cantos, monotonous and wearisome.

The poem begins with the rape of Europa. The fiery birth of the hero is not reached till the eighth book, and the proper subject of the poem, the expedition of Dionysus to India, begins only in the thirteenth. Such is the scale of the work. We are carried along throughout the whole range of mythology in a sort of corybantic dance, \texthebrew{—} a dance of words. The interest for us lies in the unclassical, one is tempted to say romantic, treatment of classical themes. Astraeus takes the horoscope of Persephone for her mother.\texthebrew{⁠}\footnote{VI.58 \emph{sqq.}

} We are taken aback by the surprising modesty of Zeus when he is gazing at Semele bathing.\texthebrew{⁠}\footnote{VII.265\texthebrew{‑}268.

} As an example of the poet's dexterity take the verses in which he describes the invention of the alphabet by Cadmus.\footnote{IV.259 \emph{sqq.}

}

The world of this poet's imagination has not the clear-cut lines of classical art. He produces his effects by reflexions,

p434
correspondences, indirections, and has a whole vocabulary of words for this purpose.\texthebrew{⁠}\footnote{Such as \textgreek{νοθός, \texthebrew{ἀ}ντίτυπος, \texthebrew{ἰ}σότυπος, μιμηλός, \texthebrew{ἀ}ντικέλευθος, \texthebrew{ἀ}λλοπρόσαλλος}. One of his favourite words is \textgreek{φειδόμενος} in the sense of \emph{forbearing, discreet, gentle}; and adjectives in -\textgreek{αλέος} (like \textgreek{φρικαλέος, σιγαλέος, \texthebrew{ὑ}πναλέος}), of which he has a great number, are a feature of his poetry. Another characteristic of his technique is the repetition of words, \emph{e.g.} (XXXV.42)

\textgreek{\texthebrew{ἐ}πε\texthebrew{ὶ} σέο μ\texthebrew{ᾶ}λλον \texthebrew{ὀ}ιστ\texthebrew{ῶ}ν μαζο\texthebrew{ὶ ὀ}ιστεύουσιν \texthebrew{ὀ}ιστευτ\texthebrew{ῆ}ρες \texthebrew{Ἐ}ρώτων.}} But he could have achieved distinction in simple pastoral poetry, as some idyllic passages show; for instance the song with the refrain

And occasionally he strikes off a verse which stays in the memory, like

That Nonnus was a pagan or quite indifferent to religion when he wrote the \emph{Dionysiaca} is always taken for granted, on the ground that a believing Christian of that age would not have revelled in such a theme. But he was converted, and he composed a free paraphrase of the Gospel of St. John, which he strangely thought suitable for his dactyls. That he should have spent his extraordinary skill on such an experiment illustrates the curious defect in literary taste common to most of the poets of the age.

Of the poets of his school, all of whom are vastly inferior to the master, Paul, the poet of St. Sophia, was the most talented, and there was something to be said for employing the new hexameter in a description of the aerial creation of Justinian. But the poet whose name, though never mentioned by contemporaries, is best known to posterity is Musaeus.\texthebrew{⁠}\footnote{Of his date we have no direct information, and the place of his birth is unknown. Agathias in the epigram \emph{Anth. Gr.} V.263 seems to be thinking of his poem. It has been guessed that he may be the Musaeus who was a correspondent of Procopius of Gaza (\emph{Epp.} 48). In any case he probably lived about that time. The contention of Rohde (\emph{Der gr. Roman}, 502) that he was imitated by Achilles Tatius has been disproved by the discovery that Achilles cannot have composed his romance long after A.D. 300, as a fragment of it, written in the fourth century, was found in Egypt (\emph{Pap. Oxyrh.} X.1250).\texthebrew{—} The text of Musaeus is very corrupt. It is to be hoped that in the last line but one he wrote

\textgreek{καδ δ\texthebrew{’ Ἡ}ρ\texthebrew{ῶ} τέθνηκε σ\texthebrew{ὺ}ν \texthebrew{ὀ}λλυμέν\texthebrew{ῳ} παρακοίτ\texthebrew{ῃ }}and not \textgreek{κα\texthebrew{ὶ} διερ\texthebrew{ὴ}}, as his latest editor amends. º The double spondee is the one good point in the verse.

} His \emph{Hero and Leander} caught the fancy of modern poets, more for the romance of the subject than for his treatment. The lamp of

p435
Hero gives it a certain charm, but it shows no more distinguished poetical talent than the little epics of Tryphiodorus and Colluthus, and the Nonnian metre is as little suitable to the subject.

To return from this digression, the historians like Procopius, Agathias, and Menander, who kept up the unbroken line of literary tradition and believed they wrote Attic Greek, could not be read except by highly educated people. So far had the spoken language drifted away from literary prose. For a larger public there was need of a popular history, written simply in the vulgar tongue. For this purpose John Malalas\texthebrew{⁠}\footnote{Malalas is the Syriac for rhetor, and the author is called John Rhetor by Evagrius, who used the work in its first form. The text we possess is an abridgment, and mutilated at the end, but it can be supplemented by many excerpts and fragments, by its use in later chronicles, especially Theophanes and the \emph{Pascal Chronicle}, and there are also fragments of the Slavonic translation. There are several difficult problems connected with Malalas which cannot be discussed here. See, for a general account, Krumbacher, \emph{G.B.L.} 325 \emph{sqq.}, and Bury, App. I to Gibbon, vol. IV, where the special studies on the subject are mentioned. Cp. works of Patzig and Gleye cited below,

Bibliography, II.2, B .} of Antioch compiled, perhaps about A.D. 550, a chronicle of the history of the world, coming down to the first year of Justinian. In a new edition there was added, whether by the author or by another hand, a continuation treating Justinian's reign on a much larger scale than the reigns of his predecessors. In the earlier part of the work there is no sense of proportion, and there are many blunders. It was written down to the level of the masses, and was nicely calculated to give them what would interest them. Pages and pages are occupied with descriptions of the personal appearances of the heroes of the Trojan War. It hit the popular taste, was largely used by subsequent writers, was in a later age translated into Slavonic, and was the first of a long series of popular Byzantine chronicles.

It is an unfortunate gap in our knowledge that we have no information as to the activities of the book trade. It would be interesting to know whether the booksellers of Constantinople received regular announcements of the works produced at Alexandria, Athens, and other places, and how many copies were circulated of a book like the \emph{Dionysiaca} of Nonnus, or the \emph{Wars} of Procopius, or the \emph{Chronicle} of Malalas, during the lives of their authors. We should then have some idea what these works meant for their own times.

In literature, as in law, the age of Justinian witnessed the

p436
culmination of the old Graeco-Roman tradition, and at the same time the signs were quite clear that the world was turning in a new direction. While his talented lawyers were shaping the greatest creation of Rome, its jurisprudence, into a final form, Latin was being definitely abandoned for Greek as the language of the legislator and the jurist; and from the same age which produced the best Greek historian since the time of Scipio Africanus comes the first of the popular chronicles which reflected the ignorance and superstition of the Middle Ages. It must not, however, be supposed that the old Greek tradition in literature disappeared. It was attenuated and modified in many tasteless ways, but the literary language was always learned as a second tongue, and never fell into disuse. The educated laity never ceased to read the ancient classics, and while in western Europe the writing of books was almost confined to ecclesiastics, in Greek lands the best books were generally written by laymen.
